{
\tolerance=4000
\hbadness=10000
\vbadness=10000
\subsection{BCU Instructions}
\label{sec:BCU-instructions}

\subsubsection[{\small AWAIT: Wait for any event or interrupt.}]{AWAIT\hfill Wait for any event or interrupt.}
\label{instr:AWAIT}\index{await}
 
\paragraph{Encoding} Format \textsf{BCU\_IPC} (Section~\ref{sec:format-BCU-IPC}), \verb|bcucode3 = 01000|\\

\bitfields{ 1/P, 2/00, 2/01, 4/1111, 5/01000, 6/-- -- -- -- -- --, 6/-- -- -- -- -- --, 6/-- -- -- -- -- -- }


\paragraph{Syntax} \texttt{await}

\paragraph{Description} Wait until \textbf{WS.WU0} equals one (cf Idle Modes chapter). The main core clock is gated during the waiting period to reduce power consumption.


\paragraph{Execution} {Reservation table \textsf{ALL} (Section~\ref{sec:reservation-ALL})}
~
{\begin{lstlisting}
stage RR:
  idle(0);
\end{lstlisting}}
\newpage
\subsubsection[{\small BARRIER: Instruction Pipeline Barrier}]{BARRIER\hfill Instruction Pipeline Barrier}
\label{instr:BARRIER}\index{barrier}
 
\paragraph{Encoding} Format \textsf{BCU\_IPC} (Section~\ref{sec:format-BCU-IPC}), \verb|bcucode3 = 01011|\\

\bitfields{ 1/P, 2/00, 2/01, 4/1111, 5/01011, 6/-- -- -- -- -- --, 6/-- -- -- -- -- --, 6/-- -- -- -- -- -- }


\paragraph{Syntax} \texttt{barrier}

\paragraph{Description} Flush the core instruction pipeline and memory system.


\paragraph{Execution} {Reservation table \textsf{ALL} (Section~\ref{sec:reservation-ALL})}
~
{\begin{lstlisting}
stage RR:
  barrier();
\end{lstlisting}}
\newpage
\subsubsection[{\small BLEND: Blended execution prefix}]{BLEND\hfill Blended execution prefix}
\label{instr:BLEND}\index{blend}
 
\paragraph{Encoding} Format \textsf{BCU\_BLEND} (Section~\ref{sec:format-BCU-BLEND}), \verb|bcucode3 = 00001|\\

\bitfields{ 1/P, 2/00, 2/01, 4/1111, 5/00001, 2/lanetodo, 2/lanesize, 8/activate, 6/registerZ }


\paragraph{Syntax} \texttt{blend}\emph{lanetodo}\emph{lanesize} \emph{registerZ}? \emph{activate}

\paragraph{Description} The results of the execution units specified by the bit mask \emph{registerZ} are blended under the mask \emph{lanetodo} applied to the \emph{lanesize}.


\paragraph{Modifier(s)}
\noindent\begin{itemize}
\item lanetodo (cf. lanetodo in section \ref{par:modifier-lanetodo} on page \pageref{par:modifier-lanetodo})
\item lanesize (cf. lanesize in section \ref{par:modifier-lanesize} on page \pageref{par:modifier-lanesize})
\end{itemize}

\paragraph{Execution} {Reservation table \textsf{BCU\_BRRP} (Section~\ref{sec:reservation-BCU-BRRP})}
~
{\begin{lstlisting}
stage ID:
  new lanetodo = _ZX_2(lanetodo);
  new lanesize = _ZX_2(lanesize);
  new activate = _ZX_6(activate);
stage RR:
  new argument = GPR[registerZ];
  blend(lanetodo, lanesize, argument, activate);
\end{lstlisting}}
\newpage
\subsubsection[{\small BREAK: Debugger Breakpoint}]{BREAK\hfill Debugger Breakpoint}
\label{instr:BREAK}\index{break}
 
\paragraph{Encoding} Format \textsf{BCU\_BREAK} (Section~\ref{sec:format-BCU-BREAK}), \verb|bcucode2 = 1111|\\

\bitfields{ 1/P, 2/00, 2/01, 4/1111, 5/11010, 6/-- -- -- -- -- --, 6/-- -- -- -- -- --, 4/-- -- -- --, 2/brknumber }


\paragraph{Syntax} \texttt{break} \emph{brknumber}

\paragraph{Description} Breakpoint for the debugger.


\paragraph{Execution} {Reservation table \textsf{ALL} (Section~\ref{sec:reservation-ALL})}
~
{\begin{lstlisting}
stage ID:
  new argument1 = _ZX_2(brknumber);
stage RR:
  break(argument1);
\end{lstlisting}}
\newpage
\subsubsection[{\small CALL: Call Subroutine PC-Relative}]{CALL\hfill Call Subroutine PC-Relative}
\label{instr:CALL}\index{call}
 
\paragraph{Encoding} Format \textsf{BCU\_UB} (Section~\ref{sec:format-BCU-UB}), \verb|bcucode1 = 11|\\

\bitfields{ 1/P, 2/00, 2/11, 27/pcrel27s2 }


\paragraph{Syntax} \texttt{call} \emph{pcrel27s2}

\paragraph{Description} The sequential next bundle \textbf{PC} is copied into the \textbf{RA} register. The next bundle \textbf{PC} is computed by adding to this instruction \textbf{PC} the \emph{pcrel27s2} after 27-bit sign extension and scaling by 4.


\paragraph{Execution} {Reservation table \textsf{BCU\_XFER} (Section~\ref{sec:reservation-BCU-XFER})}
~
{\begin{lstlisting}
stage ID:
  new argument1 = (_SX_27(pcrel27s2)<<2);
  new next = NPC;
  new address = PC + argument1;
  NPC = address;
  branch_info(1, address);
stage RR:
  RA = next;
\end{lstlisting}}
\newpage
\subsubsection[{\small CALLX: Call Subroutine Extended PC-Relative}]{CALLX\hfill Call Subroutine Extended PC-Relative}
\label{instr:CALLX}\index{callx}
 
\paragraph{Encoding} Format \textsf{BCU\_UB.X} (Section~\ref{sec:format-BCU-UB.X}), \verb|bcucode1 = 11|\\

\bitfields{ 1/P, 2/00, 2/00, 27/upper27, 1/1, 2/00, 2/11, 27/lower27 }


\paragraph{Syntax} \texttt{callx} \emph{upper27\_\discretionary{}{}{}lower27}

\paragraph{Description} The sequential next bundle \textbf{PC} is copied into the \textbf{RA} register. The next bundle \textbf{PC} is computed by adding to this instruction \textbf{PC} the \emph{upper27\_\discretionary{}{}{}lower27} after 54-bit sign extension and scaling by 4.


\paragraph{Execution} {Reservation table \textsf{BCU2.X} (Section~\ref{sec:reservation-BCU2.X})}
~
{\begin{lstlisting}
stage ID:
  new argument1 = (_SX_54(upper27_lower27)<<2);
  new next = NPC;
  new address = PC + argument1;
  NPC = address;
  branch_info(1, address);
stage RR:
  RA = next;
\end{lstlisting}}
\newpage
\subsubsection[{\small CB: Conditional Branch PC-Relative}]{CB\hfill Conditional Branch PC-Relative}
\label{instr:CB}\index{cb}
 
\paragraph{Encoding} Format \textsf{BCU\_CB} (Section~\ref{sec:format-BCU-CB}), \verb|bcucode1 = 01|\\

\bitfields{ 1/P, 2/00, 2/01, 4/bcucond, 17/pcrel17s2, 6/registerZ }


\paragraph{Syntax} \texttt{cb}\emph{bcucond} \emph{registerZ}? \emph{pcrel17s2}

\paragraph{Description} If \emph{registerZ} satisfies condition \emph{bcucond}, the next bundle \textbf{PC} is computed by adding to this instruction \textbf{PC} the \emph{pcrel17s2} after 17-bit sign extension and scaling by 4. \textbf{SRHPC} register may be updated according to current privilege level and owners configuration.


\paragraph{Modifier(s)}
\noindent\begin{itemize}
\item bcucond (cf. bcucond in section \ref{par:modifier-bcucond} on page \pageref{par:modifier-bcucond})
\end{itemize}

\paragraph{Execution} {Reservation table \textsf{BCU\_BRRP} (Section~\ref{sec:reservation-BCU-BRRP})}
~
{\begin{lstlisting}
stage RR:
  new argument1 = _ZX_4(bcucond);
  new argument2 = GPR[registerZ];
  new argument3 = (_SX_17(pcrel17s2)<<2);
  if (bcucond(argument1, argument2)) {
    new address = PC + argument3;
    NPC = address;
    branch_info(1, address);
    srhpc_update();
  } else {
    branch_info(0, 0);
  }
\end{lstlisting}}
\newpage
\subsubsection[{\small CBX: Conditional Branch Extended PC-Relative}]{CBX\hfill Conditional Branch Extended PC-Relative}
\label{instr:CBX}\index{cbx}
 
\paragraph{Encoding} Format \textsf{BCU\_CB.X} (Section~\ref{sec:format-BCU-CB.X}), \verb|bcucode1 = 01|\\

\bitfields{ 1/P, 2/00, 2/00, 27/upper27, 1/1, 2/00, 2/01, 4/bcucond, 17/lower17, 6/registerZ }


\paragraph{Syntax} \texttt{cbx}\emph{bcucond} \emph{registerZ}? \emph{upper27\_\discretionary{}{}{}lower17}

\paragraph{Description} If \emph{registerZ} satisfies condition \emph{bcucond}, the next bundle \textbf{PC} is computed by adding to this instruction \textbf{PC} the \emph{upper27\_\discretionary{}{}{}lower17} after 44-bit sign extension and scaling by 4. \textbf{SRHPC} register may be updated according to current privilege level and owners configuration.


\paragraph{Modifier(s)}
\noindent\begin{itemize}
\item bcucond (cf. bcucond in section \ref{par:modifier-bcucond} on page \pageref{par:modifier-bcucond})
\end{itemize}

\paragraph{Execution} {Reservation table \textsf{BCU2.X} (Section~\ref{sec:reservation-BCU2.X})}
~
{\begin{lstlisting}
stage RR:
  new argument1 = _ZX_4(bcucond);
  new argument2 = GPR[registerZ];
  new argument3 = (_SX_44(upper27_lower17)<<2);
  if (bcucond(argument1, argument2)) {
    new address = PC + argument3;
    NPC = address;
    branch_info(1, address);
    srhpc_update();
  } else {
    branch_info(0, 0);
  }
\end{lstlisting}}
\newpage
\subsubsection[{\small CCB: Compare Conditional Branch PC-Relative}]{CCB\hfill Compare Conditional Branch PC-Relative}
\label{instr:CCB}\index{ccb}
 
\paragraph{Encoding} Format \textsf{BCU\_CCB} (Section~\ref{sec:format-BCU-CCB}), \verb|bcucode1 = 00|\\

\bitfields{ 1/P, 2/00, 2/00, 4/ccbcomp, 11/pcrel11s2, 6/registerY, 6/registerZ }


\paragraph{Syntax} \texttt{ccb}\emph{ccbcomp} \emph{registerZ}, \emph{registerY} ? \emph{pcrel11s2}

\paragraph{Description} If \emph{registerZ} and \emph{registerY} satisfy condition \emph{ccbcomp}, the next bundle \textbf{PC} is computed by adding to this instruction \textbf{PC} the \emph{pcrel11s2} after 11-bit sign extension and scaling by 4. \textbf{SRHPC} register may be updated according to current privilege level and owners configuration.


\paragraph{Modifier(s)}
\noindent\begin{itemize}
\item ccbcomp (cf. ccbcomp in section \ref{par:modifier-ccbcomp} on page \pageref{par:modifier-ccbcomp})
\end{itemize}

\paragraph{Execution} {Reservation table \textsf{BCU\_BRRP2} (Section~\ref{sec:reservation-BCU-BRRP2})}
~
{\begin{lstlisting}
stage ID:
  if (ccbcomp(argument1, argument2, argument3)) {
    new address = PC + argument4;
    NPC = address;
    branch_info(1, address);
    srhpc_update();
  } else {
    branch_info(0, 0);
  }
stage RR:
  new argument1 = _ZX_4(ccbcomp);
  new argument2 = GPR[registerZ];
  new argument3 = GPR[registerY];
  new argument4 = (_SX_11(pcrel11s2)<<2);
\end{lstlisting}}
\newpage
\subsubsection[{\small CCBX: Compare Conditional Extended Branch PC-Relative}]{CCBX\hfill Compare Conditional Extended Branch PC-Relative}
\label{instr:CCBX}\index{ccbx}
 
\paragraph{Encoding} Format \textsf{BCU\_CCB.X} (Section~\ref{sec:format-BCU-CCB.X}), \verb|bcucode1 = 00|\\

\bitfields{ 1/P, 2/00, 2/00, 27/upper27, 1/1, 2/00, 2/00, 4/ccbcomp, 11/lower11, 6/registerY, 6/registerZ }


\paragraph{Syntax} \texttt{ccbx}\emph{ccbcomp} \emph{registerZ}, \emph{registerY} ? \emph{upper27\_\discretionary{}{}{}lower11}

\paragraph{Description} If \emph{registerZ} and \emph{registerY} satisfy condition \emph{ccbcomp}, the next bundle \textbf{PC} is computed by adding to this instruction \textbf{PC} the \emph{upper27\_\discretionary{}{}{}lower11} after 38-bit sign extension and scaling by 4. \textbf{SRHPC} register may be updated according to current privilege level and owners configuration.


\paragraph{Modifier(s)}
\noindent\begin{itemize}
\item ccbcomp (cf. ccbcomp in section \ref{par:modifier-ccbcomp} on page \pageref{par:modifier-ccbcomp})
\end{itemize}

\paragraph{Execution} {Reservation table \textsf{BCU2.X} (Section~\ref{sec:reservation-BCU2.X})}
~
{\begin{lstlisting}
stage ID:
  if (ccbcomp(argument1, argument2, argument3)) {
    new address = PC + argument4;
    NPC = address;
    branch_info(1, address);
    srhpc_update();
  } else {
    branch_info(0, 0);
  }
stage RR:
  new argument1 = _ZX_4(ccbcomp);
  new argument2 = GPR[registerZ];
  new argument3 = GPR[registerY];
  new argument4 = (_SX_38(upper27_lower11)<<2);
\end{lstlisting}}
\newpage
\subsubsection[{\small GET: Get System Register}]{GET\hfill Get System Register}
\label{instr:GET}\index{get}
 
\paragraph{Encoding} Format \textsf{BCU\_GSR} (Section~\ref{sec:format-BCU-GSR}), \verb|bcucode3 = 10001|\\

\bitfields{ 1/P, 2/00, 2/01, 4/1111, 5/10001, 3/-- -- --, 9/systemS2, 6/registerZ }


\paragraph{Syntax} \texttt{get} \emph{registerZ} = \emph{systemS2}

\paragraph{Description} The \emph{registerZ} is copied from the \emph{systemS2} register.


\paragraph{Execution} {Reservation table \textsf{BCU2\_TINY\_LSU} (Section~\ref{sec:reservation-BCU2-TINY-LSU})}
~
{\begin{lstlisting}
stage RR:
  new argument2 = SFR[systemS2];
  if (get_check_access(systemS2, registerZ)) {
    new result1 = get(systemS2, _ZX_64(argument2));
  }
stage E1:
  GPR[registerZ] = result1;
\end{lstlisting}}
\newpage
\subsubsection[{\small GOTO: Branch Unconditional PC-Relative}]{GOTO\hfill Branch Unconditional PC-Relative}
\label{instr:GOTO}\index{goto}
 
\paragraph{Encoding} Format \textsf{BCU\_UB} (Section~\ref{sec:format-BCU-UB}), \verb|bcucode1 = 10|\\

\bitfields{ 1/P, 2/00, 2/10, 27/pcrel27s2 }


\paragraph{Syntax} \texttt{goto} \emph{pcrel27s2}

\paragraph{Description} The next bundle \textbf{PC} is computed by adding to this instruction \textbf{PC} the \emph{pcrel27s2} after 27-bit sign extension and scaling by 4. \textbf{SRHPC} register may be updated according to current privilege level and owners configuration.


\paragraph{Execution} {Reservation table \textsf{BCU\_XFER} (Section~\ref{sec:reservation-BCU-XFER})}
~
{\begin{lstlisting}
stage ID:
  new argument1 = (_SX_27(pcrel27s2)<<2);
  new address = PC + argument1;
  NPC = address;
  branch_info(1, address);
stage RR:
  srhpc_update();
\end{lstlisting}}
\newpage
\subsubsection[{\small GOTOX: Branch Unconditional Extended PC-Relative}]{GOTOX\hfill Branch Unconditional Extended PC-Relative}
\label{instr:GOTOX}\index{gotox}
 
\paragraph{Encoding} Format \textsf{BCU\_UB.X} (Section~\ref{sec:format-BCU-UB.X}), \verb|bcucode1 = 10|\\

\bitfields{ 1/P, 2/00, 2/00, 27/upper27, 1/1, 2/00, 2/10, 27/lower27 }


\paragraph{Syntax} \texttt{gotox} \emph{upper27\_\discretionary{}{}{}lower27}

\paragraph{Description} The next bundle \textbf{PC} is computed by adding to this instruction \textbf{PC} the \emph{upper27\_\discretionary{}{}{}lower27} after 54-bit sign extension and scaling by 4. \textbf{SRHPC} register may be updated according to current privilege level and owners configuration.


\paragraph{Execution} {Reservation table \textsf{BCU2.X} (Section~\ref{sec:reservation-BCU2.X})}
~
{\begin{lstlisting}
stage ID:
  new argument1 = (_SX_54(upper27_lower27)<<2);
  new address = PC + argument1;
  NPC = address;
  branch_info(1, address);
stage RR:
  srhpc_update();
\end{lstlisting}}
\newpage
\subsubsection[{\small GUARD: Guarded execution prefix.}]{GUARD\hfill Guarded execution prefix.}
\label{instr:GUARD}\index{guard}
 
\paragraph{Encoding} Format \textsf{BCU\_GUARD} (Section~\ref{sec:format-BCU-GUARD}), \verb|bcucode3 = 00000|\\

\bitfields{ 1/P, 2/00, 2/01, 4/1111, 5/00000, 4/execpred, 8/activate, 6/registerZ }


\paragraph{Syntax} \texttt{guard}\emph{execpred} \emph{registerZ}? \emph{activate}

\paragraph{Description} The execution units specified by the bit mask \emph{activate} are guarded by the condition \emph{execpred} applied to the \emph{registerZ}.


\paragraph{Modifier(s)}
\noindent\begin{itemize}
\item execpred (cf. bcucond in section \ref{par:modifier-bcucond} on page \pageref{par:modifier-bcucond})
\end{itemize}

\paragraph{Execution} {Reservation table \textsf{BCU\_BRRP} (Section~\ref{sec:reservation-BCU-BRRP})}
~
{\begin{lstlisting}
stage ID:
  new bcucond = _ZX_4(execpred);
  new activate = _ZX_6(activate);
stage RR:
  new argument = GPR[registerZ];
  guard(bcucond, argument, activate);
\end{lstlisting}}
\newpage
\subsubsection[{\small ICALL: Indirect Call from Register}]{ICALL\hfill Indirect Call from Register}
\label{instr:ICALL}\index{icall}
 
\paragraph{Encoding} Format \textsf{BCU\_IBC} (Section~\ref{sec:format-BCU-IBC}), \verb|bcucode3 = 10111|\\

\bitfields{ 1/P, 2/00, 2/01, 4/1111, 5/10111, 6/-- -- -- -- -- --, 6/-- -- -- -- -- --, 6/registerZ }


\paragraph{Syntax} \texttt{icall} \emph{registerZ}

\paragraph{Description} The sequential next bundle \textbf{PC} is copied into the \textbf{RA} register. The next bundle \textbf{PC} is computed by reading the \emph{registerZ}.


\paragraph{Execution} {Reservation table \textsf{BCU\_XFER\_BRRP} (Section~\ref{sec:reservation-BCU-XFER-BRRP})}
~
{\begin{lstlisting}
stage ID:
  new argument1 = GPR[registerZ];
  new next = NPC;
stage RR:
  NPC = argument1;
  RA = next;
  branch_info(1, argument1);
\end{lstlisting}}
\newpage
\subsubsection[{\small IGET: Indirect Get System Register}]{IGET\hfill Indirect Get System Register}
\label{instr:IGET}\index{iget}
 
\paragraph{Encoding} Format \textsf{BCU\_IGSR} (Section~\ref{sec:format-BCU-IGSR}), \verb|bcucode3 = 10011|\\

\bitfields{ 1/P, 2/00, 2/01, 4/1111, 5/10011, 6/-- -- -- -- -- --, 6/-- -- -- -- -- --, 6/registerZ }


\paragraph{Syntax} \texttt{iget} \emph{registerZ}

\paragraph{Description} The \emph{registerZ} is copied from the system register whose index is in the \emph{registerZ}.


\paragraph{Execution} {Reservation table \textsf{BCU2\_TINY\_LSU} (Section~\ref{sec:reservation-BCU2-TINY-LSU})}
~
{\begin{lstlisting}
stage ID:
  new index = _ZX_9(GPR[registerZ]);
stage RR:
  if (get_check_access(index, registerZ)) {
    new result1 = get(index, SFR[index]);
  }
stage E1:
  GPR[registerZ] = result1;
\end{lstlisting}}
\newpage
\subsubsection[{\small IGOTO: Indirect Branch from Register}]{IGOTO\hfill Indirect Branch from Register}
\label{instr:IGOTO}\index{igoto}
 
\paragraph{Encoding} Format \textsf{BCU\_IBC} (Section~\ref{sec:format-BCU-IBC}), \verb|bcucode3 = 10110|\\

\bitfields{ 1/P, 2/00, 2/01, 4/1111, 5/10110, 6/-- -- -- -- -- --, 6/-- -- -- -- -- --, 6/registerZ }


\paragraph{Syntax} \texttt{igoto} \emph{registerZ}

\paragraph{Description} The next bundle \textbf{PC} is computed by reading the \emph{registerZ}. \textbf{SRHPC} register may be updated according to current privilege level and owners configuration.


\paragraph{Execution} {Reservation table \textsf{BCU\_XFER\_BRRP} (Section~\ref{sec:reservation-BCU-XFER-BRRP})}
~
{\begin{lstlisting}
stage ID:
  new argument1 = GPR[registerZ];
stage RR:
  NPC = argument1;
  branch_info(1, argument1);
  srhpc_update();
\end{lstlisting}}
\newpage
\subsubsection[{\small LOOPDO: Hardware Loop setup and Do execute}]{LOOPDO\hfill Hardware Loop setup and Do execute}
\label{instr:LOOPDO}\index{loopdo}
 
\paragraph{Encoding} Format \textsf{BCU\_HLS} (Section~\ref{sec:format-BCU-HLS}), \verb|bcucode2 = 1110|\\

\bitfields{ 1/P, 2/00, 2/01, 4/1110, 17/pcrel17s2, 6/registerZ }


\paragraph{Syntax} \texttt{loopdo} \emph{registerZ}, \emph{pcrel17s2}

\paragraph{Description} If \textbf{PS.HLE} is cleared, this instruction throws an OPCODE trap. The \textbf{LS} register is set to the sequential next bundle \textbf{PC}. The \textbf{LE} register is set to the current bundle \textbf{PC} plus the \emph{pcrel17s2} after sign extension and scaling by 4. The \textbf{LC} register is set with the \emph{registerZ}. If the \textbf{PS.HLE} bit is set in the \textbf{PS} register, hardware loop execution proceeds by comparing the sequential next bundle \textbf{PC} to the value of the \textbf{LE} register: \begin{itemize} \item If different, the next bundle \textbf{PC} is computed as usual. \item If the sequential next bundle \textbf{PC} equals the value of the \textbf{LE} register: \begin{itemize}
  \item If the value of the \textbf{LC} register ${}-1$ is non-zero, the value of the next bundle \textbf{PC}
  is read from the \textbf{LS} register, unless a branch is taken in the current bundle (in which case the next bundle \textbf{PC} is the branch target bundle \textbf{PC}).
  \item If the value of the \textbf{LC} register ${}-1$ is zero, the value of the next bundle \textbf{PC} is
  the value of the \textbf{LE} register, unless a branch is taken in the current bundle (in which case the next bundle \textbf{PC} is the branch target bundle \textbf{PC}).
  \item The \textbf{LC} register is decremented if its value is not already zero.
  \end{itemize}
\end{itemize} According to these rules, the LOOPDO instruction loops forever whenever the value of the \emph{registerZ} is zero, as the the value of the \textbf{LC} register ${}-1$ is always non-zero, while the \textbf{LC} register is never decremented. This instruction is provided for counted loops which iterate at least once, and while loops by setting the counter to zero.


\paragraph{Execution} {Reservation table \textsf{ALL} (Section~\ref{sec:reservation-ALL})}
~
{\begin{lstlisting}
stage ID:
  new argument1 = GPR[registerZ];
  new argument2 = (_SX_17(pcrel17s2)<<2);
  if (!PS.HLE) {
    _THROW(OPCODE);
  }
  invalpfb();
  LC = argument1;
stage E2:
  LS = NPC;
  LE = PC + _SX_32(argument2);
\end{lstlisting}}
\newpage
\subsubsection[{\small RET: Return from Call}]{RET\hfill Return from Call}
\label{instr:RET}\index{ret}
 
\paragraph{Encoding} Format \textsf{BCU\_RTS} (Section~\ref{sec:format-BCU-RTS}), \verb|bcucode3 = 10100|\\

\bitfields{ 1/P, 2/00, 2/01, 4/1111, 5/10100, 6/-- -- -- -- -- --, 6/-- -- -- -- -- --, 6/-- -- -- -- -- -- }


\paragraph{Syntax} \texttt{ret}

\paragraph{Description} The next bundle \textbf{PC} is obtained by reading the contents of the \textbf{RA} register. \textbf{SRHPC} register may be updated according to current privilege level and owners configuration.


\paragraph{Execution} {Reservation table \textsf{BCU\_XFER} (Section~\ref{sec:reservation-BCU-XFER})}
~
{\begin{lstlisting}
stage ID:
  NPC = RA;
stage RR:
  branch_info(1, RA);
  srhpc_update();
\end{lstlisting}}
\newpage
\subsubsection[{\small RFE: Return from Exception}]{RFE\hfill Return from Exception}
\label{instr:RFE}\index{rfe}
 
\paragraph{Encoding} Format \textsf{BCU\_RTS} (Section~\ref{sec:format-BCU-RTS}), \verb|bcucode3 = 10101|\\

\bitfields{ 1/P, 2/00, 2/01, 4/1111, 5/10101, 6/-- -- -- -- -- --, 6/-- -- -- -- -- --, 6/-- -- -- -- -- -- }


\paragraph{Syntax} \texttt{rfe}

\paragraph{Description} The \textbf{PS} is restored from the \textbf{SPS}. The \textbf{SPS} is restored from the \textbf{SSPS}. The next bundle \textbf{PC} is obtained by reading the contents of the \textbf{SPC} register. The \textbf{SPC} is restored from the \textbf{SSPC}. \textbf{SRHPC} register may be updated according to current privilege level and owners configuration.


\paragraph{Execution} {Reservation table \textsf{ALL} (Section~\ref{sec:reservation-ALL})}
~
{\begin{lstlisting}
stage RR:
  if (rfe_owner()) {
    rfe();
    branch_info(1, SPC);
    srhpc_update();
  }
\end{lstlisting}}
\newpage
\subsubsection[{\small RSWAP: General and System Register Swap}]{RSWAP\hfill General and System Register Swap}
\label{instr:RSWAP}\index{rswap}
 
\paragraph{Encoding} Format \textsf{BCU\_RSWAP} (Section~\ref{sec:format-BCU-RSWAP}), \verb|bcucode3 = 10010|\\

\bitfields{ 1/P, 2/00, 2/01, 4/1111, 5/10010, 3/-- -- --, 9/systemS4, 6/registerZ }


\paragraph{Syntax} \texttt{rswap} \emph{registerZ} = \emph{systemS4}

\paragraph{Description} The contents of the \emph{registerZ} are swapped set with the contents of the \emph{systemS4}.


\paragraph{Execution} {Reservation table \textsf{BCU2\_TINY\_LSU} (Section~\ref{sec:reservation-BCU2-TINY-LSU})}
~
{\begin{lstlisting}
stage RR:
  new argument2 = SFR[systemS4];
  new argument1 = GPR[registerZ];
stage E1:
  if (get_check_access(systemS4, registerZ) && set_check_access(systemS4, _ZX_64(argument1), registerZ)) {
    new result2 = _ZX_64(argument1);
    new result1 = _ZX_64(argument2);
    GPR[registerZ] = result1;
    SFR[systemS4] = result2;
  }
\end{lstlisting}}
\newpage

\paragraph{Encoding} Format \textsf{BCU\_RSWAPA} (Section~\ref{sec:format-BCU-RSWAPA}), \verb|bcucode3 = 10010|\\

\bitfields{ 1/P, 2/00, 2/01, 4/1111, 5/10010, 3/-- -- --, 9/systemAlone, 6/registerZ }


\paragraph{Syntax} \texttt{rswap} \emph{registerZ} = \emph{systemAlone}

\paragraph{Description} The contents of the \emph{registerZ} are swapped set with the contents of the \emph{systemAlone}.


\paragraph{Execution} {Reservation table \textsf{ALL} (Section~\ref{sec:reservation-ALL})}
~
{\begin{lstlisting}
stage RR:
  new argument2 = SFR[systemAlone];
  new argument1 = GPR[registerZ];
stage E1:
  if (get_check_access(systemAlone, registerZ) && set_check_access(systemAlone, _ZX_64(argument1), registerZ)) {
    new result2 = _ZX_64(argument1);
    new result1 = _ZX_64(argument2);
    GPR[registerZ] = result1;
    SFR[systemAlone] = result2;
  }
\end{lstlisting}}
\newpage
\subsubsection[{\small SCALL: System Call}]{SCALL\hfill System Call}
\label{instr:SCALL}\index{scall}
 
\paragraph{Encoding} Format \textsf{BCU\_SCR} (Section~\ref{sec:format-BCU-SCR}), \verb|bcucode2 = 1111|\\

\bitfields{ 1/P, 2/00, 2/01, 4/1111, 5/11001, 6/-- -- -- -- -- --, 6/-- -- -- -- -- --, 6/registerZ }


\paragraph{Syntax} \texttt{scall} \emph{registerZ}

\paragraph{Description} System call with number \emph{}. Interrupts are disabled.


\paragraph{Execution} {Reservation table \textsf{ALL} (Section~\ref{sec:reservation-ALL})}
~
{\begin{lstlisting}
stage ID:
  new argument1 = GPR[registerZ];
  syscall(argument1);
  branch_info(1, argument1);
\end{lstlisting}}
\newpage

\paragraph{Encoding} Format \textsf{BCU\_SCI} (Section~\ref{sec:format-BCU-SCI}), \verb|bcucode2 = 1111|\\

\bitfields{ 1/P, 2/00, 2/01, 4/1111, 5/11000, 6/-- -- -- -- -- --, 12/sysnumber }


\paragraph{Syntax} \texttt{scall} \emph{sysnumber}

\paragraph{Description} System call with number \emph{}. Interrupts are disabled.


\paragraph{Execution} {Reservation table \textsf{ALL} (Section~\ref{sec:reservation-ALL})}
~
{\begin{lstlisting}
stage ID:
  new argument1 = _ZX_12(sysnumber);
  syscall(argument1);
  branch_info(1, argument1);
\end{lstlisting}}
\newpage
\subsubsection[{\small SET: Set System Register}]{SET\hfill Set System Register}
\label{instr:SET}\index{set}
 
\paragraph{Encoding} Format \textsf{BCU\_SET} (Section~\ref{sec:format-BCU-SET}), \verb|bcucode3 = 10000|\\

\bitfields{ 1/P, 2/00, 2/01, 4/1111, 5/10000, 3/-- -- --, 9/systemT3, 6/registerZ }


\paragraph{Syntax} \texttt{set} \emph{systemT3} = \emph{registerZ}

\paragraph{Description} The \emph{systemT3} register is set with the \emph{registerZ}.


\paragraph{Execution} {Reservation table \textsf{BCU2} (Section~\ref{sec:reservation-BCU2})}
~
{\begin{lstlisting}
stage ID:
  new argument2 = GPR[registerZ];
stage E3:
  if (set_check_access(systemT3, _ZX_64(argument2), registerZ)) {
    new result1 = _ZX_64(argument2);
    SFR[systemT3] = result1;
  }
\end{lstlisting}}
\newpage

\paragraph{Encoding} Format \textsf{BCU\_SETA} (Section~\ref{sec:format-BCU-SETA}), \verb|bcucode3 = 10000|\\

\bitfields{ 1/P, 2/00, 2/01, 4/1111, 5/10000, 3/-- -- --, 9/systemAlone, 6/registerZ }


\paragraph{Syntax} \texttt{set} \emph{systemAlone} = \emph{registerZ}

\paragraph{Description} The \emph{systemAlone} register is set with the \emph{registerZ}.


\paragraph{Execution} {Reservation table \textsf{ALL} (Section~\ref{sec:reservation-ALL})}
~
{\begin{lstlisting}
stage ID:
  new argument2 = GPR[registerZ];
stage E3:
  if (set_check_access(systemAlone, _ZX_64(argument2), registerZ)) {
    new result1 = _ZX_64(argument2);
    SFR[systemAlone] = result1;
  }
\end{lstlisting}}
\newpage

\paragraph{Encoding} Format \textsf{BCU\_SETRA} (Section~\ref{sec:format-BCU-SETRA}), \verb|bcucode3 = 10000|\\

\bitfields{ 1/P, 2/00, 2/01, 4/1111, 5/10000, 3/-- -- --, 9/systemRA, 6/registerZ }


\paragraph{Syntax} \texttt{set} \emph{systemRA} = \emph{registerZ}

\paragraph{Description} The \emph{systemRA} register is set with the \emph{registerZ}.


\paragraph{Execution} {Reservation table \textsf{BCU2} (Section~\ref{sec:reservation-BCU2})}
~
{\begin{lstlisting}
stage ID:
  new argument2 = GPR[registerZ];
stage E3:
  if (set_check_access(systemRA, _ZX_64(argument2), registerZ)) {
    new result1 = _ZX_64(argument2);
    SFR[systemRA] = result1;
  }
\end{lstlisting}}
\newpage
\subsubsection[{\small SLEEP: Wait for an eligible interrupt.}]{SLEEP\hfill Wait for an eligible interrupt.}
\label{instr:SLEEP}\index{sleep}
 
\paragraph{Encoding} Format \textsf{BCU\_IPC} (Section~\ref{sec:format-BCU-IPC}), \verb|bcucode3 = 01001|\\

\bitfields{ 1/P, 2/00, 2/01, 4/1111, 5/01001, 6/-- -- -- -- -- --, 6/-- -- -- -- -- --, 6/-- -- -- -- -- -- }


\paragraph{Syntax} \texttt{sleep}

\paragraph{Description} Wait until \textbf{WS.WU1} equals one (cf Idle Modes chapter). The main core clock is gated during the waiting period to reduce power consumption.


\paragraph{Execution} {Reservation table \textsf{ALL} (Section~\ref{sec:reservation-ALL})}
~
{\begin{lstlisting}
stage RR:
  idle(1);
\end{lstlisting}}
\newpage
\subsubsection[{\small STOP: Wait for power controller wake-up.}]{STOP\hfill Wait for power controller wake-up.}
\label{instr:STOP}\index{stop}
 
\paragraph{Encoding} Format \textsf{BCU\_IPC} (Section~\ref{sec:format-BCU-IPC}), \verb|bcucode3 = 01010|\\

\bitfields{ 1/P, 2/00, 2/01, 4/1111, 5/01010, 6/-- -- -- -- -- --, 6/-- -- -- -- -- --, 6/-- -- -- -- -- -- }


\paragraph{Syntax} \texttt{stop}

\paragraph{Description} Wait until \textbf{WS.WU2} equals one (cf Idle Modes chapter). The main core clock is gated during the waiting period to reduce power consumption.


\paragraph{Execution} {Reservation table \textsf{ALL} (Section~\ref{sec:reservation-ALL})}
~
{\begin{lstlisting}
stage RR:
  if (stop_owner()) {
    idle(2);
  } else {
    _THROW(PRIVILEGE);
  }
\end{lstlisting}}
\newpage
\subsubsection[{\small SYNCGROUP: Synchronize Group of Cores}]{SYNCGROUP\hfill Synchronize Group of Cores}
\label{instr:SYNCGROUP}\index{syncgroup}
 
\paragraph{Encoding} Format \textsf{BCU\_PGI} (Section~\ref{sec:format-BCU-PGI}), \verb|bcucode3 = 01101|\\

\bitfields{ 1/P, 2/00, 2/01, 4/1111, 5/01101, 6/-- -- -- -- -- --, 6/-- -- -- -- -- --, 6/registerZ }


\paragraph{Syntax} \texttt{syncgroup} \emph{registerZ}

\paragraph{Description} The \emph{registerZ} is interpreted as four 16-bit fields: \textbf{waitclrF}, \textbf{waitclrB}, \textbf{notifyF}, \textbf{notifyB}. Upon executing the instruction, the core waits until all the events designated by bits in \textbf{waitclrF} and \textbf{waitclrB} in the 32 least-significant bits of \textbf{IPE} are set. When all the designated events are set, the core clears them, then notifies the event lines corresponding to the bit sets in fields \textbf{notifyF} and and \textbf{notifyB}. This notification of events is not registered in the core, it is a signal sent to set bits into the \textbf{IPE.FE} and \textbf{IPE.BE} fields of the destination cores. The successor core always receives the \textbf{notifyF} events in its \textbf{IPE.FE} field; likewise, the predecessor core always receives the \textbf{notifyB} events in its \textbf{IPE.BE} field. Furthermore, other PEs may receive these events, depending on the setting of the \textbf{IPE.FM} and \textbf{IPE.BM} fields in these neighbor PE cores.


\paragraph{Execution} {Reservation table \textsf{BCU2} (Section~\ref{sec:reservation-BCU2})}
~
{\begin{lstlisting}
stage ID:
  new argument1 = GPR[registerZ];
stage RR:
  if (!syncgroup_owner()) {
    _THROW(PRIVILEGE);
  } else {
    new waitclrF = argument1 & 0xFFFF;
    new waitclrB = (argument1 >> 16) & 0xFFFF;
    new notifyF = (argument1 >> 32) & 0xFFFF;
    new notifyB = (argument1 >> 48) & 0xFFFF;
    syncgroup(waitclrF, waitclrB, notifyF, notifyB);
  }
\end{lstlisting}}
\newpage
\subsubsection[{\small TLBDINVAL: Invalidate Data Micro-TLB.}]{TLBDINVAL\hfill Invalidate Data Micro-TLB.}
\label{instr:TLBDINVAL}\index{tlbdinval}
 
\paragraph{Encoding} Format \textsf{BCU\_TLB} (Section~\ref{sec:format-BCU-TLB}), \verb|bcucode3 = 00110|\\

\bitfields{ 1/P, 2/00, 2/01, 4/1111, 5/00110, 6/-- -- -- -- -- --, 6/-- -- -- -- -- --, 6/-- -- -- -- -- -- }


\paragraph{Syntax} \texttt{tlbdinval}

\paragraph{Description} Invalidate Data Micro-TLB.


\paragraph{Execution} {Reservation table \textsf{ALL} (Section~\ref{sec:reservation-ALL})}
~
{\begin{lstlisting}
stage RR:
  invaldtlb();
\end{lstlisting}}
\newpage
\subsubsection[{\small TLBIINVAL: Invalidate Instruction Micro-TLB.}]{TLBIINVAL\hfill Invalidate Instruction Micro-TLB.}
\label{instr:TLBIINVAL}\index{tlbiinval}
 
\paragraph{Encoding} Format \textsf{BCU\_TLB} (Section~\ref{sec:format-BCU-TLB}), \verb|bcucode3 = 00111|\\

\bitfields{ 1/P, 2/00, 2/01, 4/1111, 5/00111, 6/-- -- -- -- -- --, 6/-- -- -- -- -- --, 6/-- -- -- -- -- -- }


\paragraph{Syntax} \texttt{tlbiinval}

\paragraph{Description} Invalidate Instruction Micro-TLB.


\paragraph{Execution} {Reservation table \textsf{ALL} (Section~\ref{sec:reservation-ALL})}
~
{\begin{lstlisting}
stage RR:
  invalitlb();
\end{lstlisting}}
\newpage
\subsubsection[{\small TLBPROBE: Probe TLB}]{TLBPROBE\hfill Probe TLB}
\label{instr:TLBPROBE}\index{tlbprobe}
 
\paragraph{Encoding} Format \textsf{BCU\_TLB} (Section~\ref{sec:format-BCU-TLB}), \verb|bcucode3 = 00011|\\

\bitfields{ 1/P, 2/00, 2/01, 4/1111, 5/00011, 6/-- -- -- -- -- --, 6/-- -- -- -- -- --, 6/-- -- -- -- -- -- }


\paragraph{Syntax} \texttt{tlbprobe}

\paragraph{Description} Probe the TLB for virtual address match in the TEL:TEH register pair and write to MMC.IDX.


\paragraph{Execution} {Reservation table \textsf{ALL} (Section~\ref{sec:reservation-ALL})}
~
{\begin{lstlisting}
stage ID:
  if (!mmi_owner()) {
    _THROW(PRIVILEGE);
  } else {
    probetlb();
  }
\end{lstlisting}}
\newpage
\subsubsection[{\small TLBREAD: Read TLB Entry}]{TLBREAD\hfill Read TLB Entry}
\label{instr:TLBREAD}\index{tlbread}
 
\paragraph{Encoding} Format \textsf{BCU\_TLB} (Section~\ref{sec:format-BCU-TLB}), \verb|bcucode3 = 00010|\\

\bitfields{ 1/P, 2/00, 2/01, 4/1111, 5/00010, 6/-- -- -- -- -- --, 6/-- -- -- -- -- --, 6/-- -- -- -- -- -- }


\paragraph{Syntax} \texttt{tlbread}

\paragraph{Description} Read the TLB entry into the TEL:TEH register pair using MMC.IDX.


\paragraph{Execution} {Reservation table \textsf{ALL} (Section~\ref{sec:reservation-ALL})}
~
{\begin{lstlisting}
stage ID:
  if (!mmi_owner()) {
    _THROW(PRIVILEGE);
  } else {
    readtlb();
  }
\end{lstlisting}}
\newpage
\subsubsection[{\small TLBWRITE: TLB Entry Write}]{TLBWRITE\hfill TLB Entry Write}
\label{instr:TLBWRITE}\index{tlbwrite}
 
\paragraph{Encoding} Format \textsf{BCU\_TLB} (Section~\ref{sec:format-BCU-TLB}), \verb|bcucode3 = 00100|\\

\bitfields{ 1/P, 2/00, 2/01, 4/1111, 5/00100, 6/-- -- -- -- -- --, 6/-- -- -- -- -- --, 6/-- -- -- -- -- -- }


\paragraph{Syntax} \texttt{tlbwrite}

\paragraph{Description} Write the TLB entry from the TEL:TEH register pair using MMC.IDX.


\paragraph{Execution} {Reservation table \textsf{ALL} (Section~\ref{sec:reservation-ALL})}
~
{\begin{lstlisting}
stage ID:
  if (!mmi_owner()) {
    _THROW(PRIVILEGE);
  } else {
    writetlb();
  }
\end{lstlisting}}
\newpage
\subsubsection[{\small WAITIT: Wait for Interrupt}]{WAITIT\hfill Wait for Interrupt}
\label{instr:WAITIT}\index{waitit}
 
\paragraph{Encoding} Format \textsf{BCU\_PGI} (Section~\ref{sec:format-BCU-PGI}), \verb|bcucode3 = 01100|\\

\bitfields{ 1/P, 2/00, 2/01, 4/1111, 5/01100, 6/-- -- -- -- -- --, 6/-- -- -- -- -- --, 6/registerZ }


\paragraph{Syntax} \texttt{waitit} \emph{registerZ}

\paragraph{Description} Execution waits for selected interrupts in the \textbf{ILR} register (masking non-owned interrupts, i.e. interrupts that belong to a more privileged PL. The \emph{registerZ} is interpreted as an OR-mask in the least significant word and a AND-mask in the most significant word. The wait ends if all the bits in \textbf{ILR} match the non-zero AND-mask, or if any bit in \textbf{ILR} matches the non-zero OR-mask. When the wait is over, the value of \textbf{ILR} (masked in the same way) is written to the \emph{registerZ}. When both the AND-mask and OR-mask are zero, no wait happens.


\paragraph{Execution} {Reservation table \textsf{BCU2\_TINY\_LSU} (Section~\ref{sec:reservation-BCU2-TINY-LSU})}
~
{\begin{lstlisting}
stage ID:
  new argument1 = GPR[registerZ];
stage RR:
  new ormask = argument1 & 0xFFFFFFFF;
  new andmask = argument1.32[1];
  new result1 = waitit(ormask, andmask);
  GPR[registerZ] = result1;
\end{lstlisting}}
\newpage
\subsubsection[{\small WFXL: Effects on Least Significant Word System Register}]{WFXL\hfill Effects on Least Significant Word System Register}
\label{instr:WFXL}\index{wfxl}
 
\paragraph{Encoding} Format \textsf{BCU\_WFX} (Section~\ref{sec:format-BCU-WFX}), \verb|bcucode3 = 01110|\\

\bitfields{ 1/P, 2/00, 2/01, 4/1111, 5/01110, 3/-- -- --, 9/systemT2, 6/registerZ }


\paragraph{Syntax} \texttt{wfxl} \emph{systemT2}, \emph{registerZ}

\paragraph{Description} The least significant word of the \emph{systemT2} is operated, with the most significant word of the \emph{registerZ} as a set mask and the least significant word of the \emph{registerZ} as a clear mask.


\paragraph{Execution} {Reservation table \textsf{BCU2} (Section~\ref{sec:reservation-BCU2})}
~
{\begin{lstlisting}
stage ID:
  new argument2 = GPR[registerZ];
stage E4:
  new argument1 = SFR[systemT2];
  if (wfxl_check_access(systemT2, registerZ)) {
    new result1 = wfxl(systemT2, _ZX_64(argument2));
  }
  SFR[systemT2] = result1;
\end{lstlisting}}
\newpage

\paragraph{Encoding} Format \textsf{BCU\_WFXA} (Section~\ref{sec:format-BCU-WFXA}), \verb|bcucode3 = 01110|\\

\bitfields{ 1/P, 2/00, 2/01, 4/1111, 5/01110, 3/-- -- --, 9/systemAlone, 6/registerZ }


\paragraph{Syntax} \texttt{wfxl} \emph{systemAlone}, \emph{registerZ}

\paragraph{Description} The least significant word of the \emph{systemAlone} is operated, with the most significant word of the \emph{registerZ} as a set mask and the least significant word of the \emph{registerZ} as a clear mask.


\paragraph{Execution} {Reservation table \textsf{ALL} (Section~\ref{sec:reservation-ALL})}
~
{\begin{lstlisting}
stage ID:
  new argument1 = SFR[systemAlone];
  new argument2 = GPR[registerZ];
stage RR:
  if (wfxl_check_access(systemAlone, registerZ)) {
    new result1 = wfxl(systemAlone, _ZX_64(argument2));
  }
stage E4:
  SFR[systemAlone] = result1;
\end{lstlisting}}
\newpage
\subsubsection[{\small WFXM: Word Effects on Most Significant Word System Register}]{WFXM\hfill Word Effects on Most Significant Word System Register}
\label{instr:WFXM}\index{wfxm}
 
\paragraph{Encoding} Format \textsf{BCU\_WFX} (Section~\ref{sec:format-BCU-WFX}), \verb|bcucode3 = 01111|\\

\bitfields{ 1/P, 2/00, 2/01, 4/1111, 5/01111, 3/-- -- --, 9/systemT2, 6/registerZ }


\paragraph{Syntax} \texttt{wfxm} \emph{systemT2}, \emph{registerZ}

\paragraph{Description} The most significant word of the \emph{systemT2} is operated, with the most significant word of the \emph{registerZ} as a set mask and the least significant word of the \emph{registerZ} as a clear mask.


\paragraph{Execution} {Reservation table \textsf{BCU2} (Section~\ref{sec:reservation-BCU2})}
~
{\begin{lstlisting}
stage ID:
  new argument2 = GPR[registerZ];
stage E4:
  new argument1 = SFR[systemT2];
  if (wfxm_check_access(systemT2, registerZ)) {
    new result1 = wfxm(systemT2, _ZX_64(argument2));
  }
  SFR[systemT2] = result1;
\end{lstlisting}}
\newpage

\paragraph{Encoding} Format \textsf{BCU\_WFXA} (Section~\ref{sec:format-BCU-WFXA}), \verb|bcucode3 = 01111|\\

\bitfields{ 1/P, 2/00, 2/01, 4/1111, 5/01111, 3/-- -- --, 9/systemAlone, 6/registerZ }


\paragraph{Syntax} \texttt{wfxm} \emph{systemAlone}, \emph{registerZ}

\paragraph{Description} The most significant word of the \emph{systemAlone} is operated, with the most significant word of the \emph{registerZ} as a set mask and the least significant word of the \emph{registerZ} as a clear mask.


\paragraph{Execution} {Reservation table \textsf{ALL} (Section~\ref{sec:reservation-ALL})}
~
{\begin{lstlisting}
stage ID:
  new argument1 = SFR[systemAlone];
  new argument2 = GPR[registerZ];
stage RR:
  if (wfxm_check_access(systemAlone, registerZ)) {
    new result1 = wfxm(systemAlone, _ZX_64(argument2));
  }
stage E4:
  SFR[systemAlone] = result1;
\end{lstlisting}}
\newpage
\subsection{LSU Instructions}
\label{sec:LSU-instructions}

\subsubsection[{\small ACSWAPB: Atomic Compare and Swap Byte}]{ACSWAPB\hfill Atomic Compare and Swap Byte}
\label{instr:ACSWAPB}\index{acswapb}
 
\paragraph{Encoding} Format \textsf{LSU\_ACSWAPSB} (Section~\ref{sec:format-LSU-ACSWAPSB}), \verb|lsucode5 = 000|\\

\bitfields{ 1/P, 2/01, 3/111, 2/coherency, 6/registerW, 2/11, 3/000, 1/boolcas, 5/registerO, 1/1, 6/registerZ }


\paragraph{Syntax} \texttt{acswapb}\emph{boolcas}\emph{coherency} \emph{registerW}, [\emph{registerZ}] = \emph{registerO}

\paragraph{Description} The effective address is the value of the \emph{registerZ}.
This instruction implements Compare-And-Swap (CAS) on memory byte (8 bits). The \emph{registerO} is interpreted as the concatenation of \texttt{update} (low 64 bits) and \texttt{expected} (high 64 bits). The memory byte at the effective address is read into \texttt{current}. If \texttt{current} equals \texttt{expected}, the swap is successful and \texttt{update} is stored into the memory byte at the effective address. Otherwise, the swap has failed and the memory is not updated. The instruction return value is stored into the \emph{registerW}. It is true/false in case \texttt{boolcas} is true, or the \texttt{current} value in case \texttt{boolcas} is false.


\paragraph{Modifier(s)}
\noindent\begin{itemize}
\item coherency (cf. coherency in section \ref{par:modifier-coherency} on page \pageref{par:modifier-coherency})
\item boolcas (cf. boolcas in section \ref{par:modifier-boolcas} on page \pageref{par:modifier-boolcas})
\end{itemize}

\paragraph{Execution} {Reservation table \textsf{LSU2\_MEMW\_AUXR\_AUXW} (Section~\ref{sec:reservation-LSU2-MEMW-AUXR-AUXW})}
~
{\begin{lstlisting}
stage ID:
  new boolcas = _ZX_1(boolcas);
  new coherency = _ZX_2(coherency);
stage RR:
  new address = GPR[registerZ];
stage E1:
  new arguments = PGR[registerO];
  new result1 = MEM.atomic_cas(address, 0x1, boolcas | (coherency << 32), arguments.8[0], arguments.8[8], registerW);
stage SR:
  GPR[registerW] = result1;
\end{lstlisting}}
\newpage

\paragraph{Encoding} Format \textsf{LSU\_ACSWAPSB.O} (Section~\ref{sec:format-LSU-ACSWAPSB.O}), \verb|lsucode5 = 000|\\

\bitfields{ 1/P, 2/00, 2/-- --, 27/offset27, 1/1, 2/01, 3/111, 2/coherency, 6/registerW, 2/11, 3/000, 1/boolcas, 5/registerO, 1/1, 6/registerZ }


\paragraph{Syntax} \texttt{acswapb}\emph{boolcas}\emph{coherency} \emph{registerW}, \emph{offset27}[\emph{registerZ}] = \emph{registerO}

\paragraph{Description} The effective address is computed by adding the \emph{offset27} to the \emph{registerZ}.
This instruction implements Compare-And-Swap (CAS) on memory byte (8 bits). The \emph{registerO} is interpreted as the concatenation of \texttt{update} (low 64 bits) and \texttt{expected} (high 64 bits). The memory byte at the effective address is read into \texttt{current}. If \texttt{current} equals \texttt{expected}, the swap is successful and \texttt{update} is stored into the memory byte at the effective address. Otherwise, the swap has failed and the memory is not updated. The instruction return value is stored into the \emph{registerW}. It is true/false in case \texttt{boolcas} is true, or the \texttt{current} value in case \texttt{boolcas} is false.


\paragraph{Modifier(s)}
\noindent\begin{itemize}
\item coherency (cf. coherency in section \ref{par:modifier-coherency} on page \pageref{par:modifier-coherency})
\item boolcas (cf. boolcas in section \ref{par:modifier-boolcas} on page \pageref{par:modifier-boolcas})
\end{itemize}

\paragraph{Execution} {Reservation table \textsf{LSU2\_MEMW\_AUXR\_AUXW.X} (Section~\ref{sec:reservation-LSU2-MEMW-AUXR-AUXW.X})}
~
{\begin{lstlisting}
stage ID:
  new offset = _SX_27(offset27);
  new boolcas = _ZX_1(boolcas);
  new coherency = _ZX_2(coherency);
stage RR:
  new base = GPR[registerZ];
  new address = base + offset;
stage E1:
  new arguments = PGR[registerO];
  new result1 = MEM.atomic_cas(address, 0x1, boolcas | (coherency << 32), arguments.8[0], arguments.8[8], registerW);
stage SR:
  GPR[registerW] = result1;
\end{lstlisting}}
\newpage

\paragraph{Encoding} Format \textsf{LSU\_ACSWAPSB.D} (Section~\ref{sec:format-LSU-ACSWAPSB.D}), \verb|lsucode5 = 000|\\

\bitfields{ 1/P, 2/00, 2/-- --, 27/extend27, 1/1, 2/00, 2/-- --, 27/offset27, 1/1, 2/01, 3/111, 2/coherency, 6/registerW, 2/11, 3/000, 1/boolcas, 5/registerO, 1/1, 6/registerZ }


\paragraph{Syntax} \texttt{acswapb}\emph{boolcas}\emph{coherency} \emph{registerW}, \emph{extend27\_\discretionary{}{}{}offset27}[\emph{registerZ}] = \emph{registerO}

\paragraph{Description} The effective address is computed by adding the \emph{extend27\_\discretionary{}{}{}offset27} to the \emph{registerZ}.
This instruction implements Compare-And-Swap (CAS) on memory byte (8 bits). The \emph{registerO} is interpreted as the concatenation of \texttt{update} (low 64 bits) and \texttt{expected} (high 64 bits). The memory byte at the effective address is read into \texttt{current}. If \texttt{current} equals \texttt{expected}, the swap is successful and \texttt{update} is stored into the memory byte at the effective address. Otherwise, the swap has failed and the memory is not updated. The instruction return value is stored into the \emph{registerW}. It is true/false in case \texttt{boolcas} is true, or the \texttt{current} value in case \texttt{boolcas} is false.


\paragraph{Modifier(s)}
\noindent\begin{itemize}
\item coherency (cf. coherency in section \ref{par:modifier-coherency} on page \pageref{par:modifier-coherency})
\item boolcas (cf. boolcas in section \ref{par:modifier-boolcas} on page \pageref{par:modifier-boolcas})
\end{itemize}

\paragraph{Execution} {Reservation table \textsf{LSU2\_MEMW\_AUXR\_AUXW.Y} (Section~\ref{sec:reservation-LSU2-MEMW-AUXR-AUXW.Y})}
~
{\begin{lstlisting}
stage ID:
  new offset = _SX_54(extend27_offset27);
  new boolcas = _ZX_1(boolcas);
  new coherency = _ZX_2(coherency);
stage RR:
  new base = GPR[registerZ];
  new address = base + offset;
stage E1:
  new arguments = PGR[registerO];
  new result1 = MEM.atomic_cas(address, 0x1, boolcas | (coherency << 32), arguments.8[0], arguments.8[8], registerW);
stage SR:
  GPR[registerW] = result1;
\end{lstlisting}}
\newpage
\subsubsection[{\small ACSWAPD: Atomic Compare and Swap Double Word}]{ACSWAPD\hfill Atomic Compare and Swap Double Word}
\label{instr:ACSWAPD}\index{acswapd}
 
\paragraph{Encoding} Format \textsf{LSU\_ACSWAPSB} (Section~\ref{sec:format-LSU-ACSWAPSB}), \verb|lsucode5 = 011|\\

\bitfields{ 1/P, 2/01, 3/111, 2/coherency, 6/registerW, 2/11, 3/011, 1/boolcas, 5/registerO, 1/1, 6/registerZ }


\paragraph{Syntax} \texttt{acswapd}\emph{boolcas}\emph{coherency} \emph{registerW}, [\emph{registerZ}] = \emph{registerO}

\paragraph{Description} The effective address is the value of the \emph{registerZ}.
This instruction implements Compare-And-Swap (CAS) on memory double word (64 bits). The \emph{registerO} is interpreted as the concatenation of \texttt{update} (low 64 bits) and \texttt{expected} (high 64 bits). The memory double word at the effective address is read into \texttt{current}. If \texttt{current} equals \texttt{expected}, the swap is successful and \texttt{update} is stored into the memory double word at the effective address. Otherwise, the swap has failed and the memory is not updated. The instruction return value is stored into the \emph{registerW}. It is true/false in case \texttt{boolcas} is true, or the \texttt{current} value in case \texttt{boolcas} is false. The effective address must be double word aligned (multiple of 8).


\paragraph{Modifier(s)}
\noindent\begin{itemize}
\item coherency (cf. coherency in section \ref{par:modifier-coherency} on page \pageref{par:modifier-coherency})
\item boolcas (cf. boolcas in section \ref{par:modifier-boolcas} on page \pageref{par:modifier-boolcas})
\end{itemize}

\paragraph{Execution} {Reservation table \textsf{LSU2\_MEMW\_AUXR\_AUXW} (Section~\ref{sec:reservation-LSU2-MEMW-AUXR-AUXW})}
~
{\begin{lstlisting}
stage ID:
  new boolcas = _ZX_1(boolcas);
  new coherency = _ZX_2(coherency);
stage RR:
  new address = GPR[registerZ];
stage E1:
  new arguments = PGR[registerO];
  new result1 = MEM.atomic_cas(address, 0xFF, boolcas | (coherency << 32), arguments.64[0], arguments.64[1], registerW);
stage SR:
  GPR[registerW] = result1;
\end{lstlisting}}
\newpage

\paragraph{Encoding} Format \textsf{LSU\_ACSWAPSB.O} (Section~\ref{sec:format-LSU-ACSWAPSB.O}), \verb|lsucode5 = 011|\\

\bitfields{ 1/P, 2/00, 2/-- --, 27/offset27, 1/1, 2/01, 3/111, 2/coherency, 6/registerW, 2/11, 3/011, 1/boolcas, 5/registerO, 1/1, 6/registerZ }


\paragraph{Syntax} \texttt{acswapd}\emph{boolcas}\emph{coherency} \emph{registerW}, \emph{offset27}[\emph{registerZ}] = \emph{registerO}

\paragraph{Description} The effective address is computed by adding the \emph{offset27} to the \emph{registerZ}.
This instruction implements Compare-And-Swap (CAS) on memory double word (64 bits). The \emph{registerO} is interpreted as the concatenation of \texttt{update} (low 64 bits) and \texttt{expected} (high 64 bits). The memory double word at the effective address is read into \texttt{current}. If \texttt{current} equals \texttt{expected}, the swap is successful and \texttt{update} is stored into the memory double word at the effective address. Otherwise, the swap has failed and the memory is not updated. The instruction return value is stored into the \emph{registerW}. It is true/false in case \texttt{boolcas} is true, or the \texttt{current} value in case \texttt{boolcas} is false. The effective address must be double word aligned (multiple of 8).


\paragraph{Modifier(s)}
\noindent\begin{itemize}
\item coherency (cf. coherency in section \ref{par:modifier-coherency} on page \pageref{par:modifier-coherency})
\item boolcas (cf. boolcas in section \ref{par:modifier-boolcas} on page \pageref{par:modifier-boolcas})
\end{itemize}

\paragraph{Execution} {Reservation table \textsf{LSU2\_MEMW\_AUXR\_AUXW.X} (Section~\ref{sec:reservation-LSU2-MEMW-AUXR-AUXW.X})}
~
{\begin{lstlisting}
stage ID:
  new offset = _SX_27(offset27);
  new boolcas = _ZX_1(boolcas);
  new coherency = _ZX_2(coherency);
stage RR:
  new base = GPR[registerZ];
  new address = base + offset;
stage E1:
  new arguments = PGR[registerO];
  new result1 = MEM.atomic_cas(address, 0xFF, boolcas | (coherency << 32), arguments.64[0], arguments.64[1], registerW);
stage SR:
  GPR[registerW] = result1;
\end{lstlisting}}
\newpage

\paragraph{Encoding} Format \textsf{LSU\_ACSWAPSB.D} (Section~\ref{sec:format-LSU-ACSWAPSB.D}), \verb|lsucode5 = 011|\\

\bitfields{ 1/P, 2/00, 2/-- --, 27/extend27, 1/1, 2/00, 2/-- --, 27/offset27, 1/1, 2/01, 3/111, 2/coherency, 6/registerW, 2/11, 3/011, 1/boolcas, 5/registerO, 1/1, 6/registerZ }


\paragraph{Syntax} \texttt{acswapd}\emph{boolcas}\emph{coherency} \emph{registerW}, \emph{extend27\_\discretionary{}{}{}offset27}[\emph{registerZ}] = \emph{registerO}

\paragraph{Description} The effective address is computed by adding the \emph{extend27\_\discretionary{}{}{}offset27} to the \emph{registerZ}.
This instruction implements Compare-And-Swap (CAS) on memory double word (64 bits). The \emph{registerO} is interpreted as the concatenation of \texttt{update} (low 64 bits) and \texttt{expected} (high 64 bits). The memory double word at the effective address is read into \texttt{current}. If \texttt{current} equals \texttt{expected}, the swap is successful and \texttt{update} is stored into the memory double word at the effective address. Otherwise, the swap has failed and the memory is not updated. The instruction return value is stored into the \emph{registerW}. It is true/false in case \texttt{boolcas} is true, or the \texttt{current} value in case \texttt{boolcas} is false. The effective address must be double word aligned (multiple of 8).


\paragraph{Modifier(s)}
\noindent\begin{itemize}
\item coherency (cf. coherency in section \ref{par:modifier-coherency} on page \pageref{par:modifier-coherency})
\item boolcas (cf. boolcas in section \ref{par:modifier-boolcas} on page \pageref{par:modifier-boolcas})
\end{itemize}

\paragraph{Execution} {Reservation table \textsf{LSU2\_MEMW\_AUXR\_AUXW.Y} (Section~\ref{sec:reservation-LSU2-MEMW-AUXR-AUXW.Y})}
~
{\begin{lstlisting}
stage ID:
  new offset = _SX_54(extend27_offset27);
  new boolcas = _ZX_1(boolcas);
  new coherency = _ZX_2(coherency);
stage RR:
  new base = GPR[registerZ];
  new address = base + offset;
stage E1:
  new arguments = PGR[registerO];
  new result1 = MEM.atomic_cas(address, 0xFF, boolcas | (coherency << 32), arguments.64[0], arguments.64[1], registerW);
stage SR:
  GPR[registerW] = result1;
\end{lstlisting}}
\newpage
\subsubsection[{\small ACSWAPH: Atomic Compare and Swap Half Word}]{ACSWAPH\hfill Atomic Compare and Swap Half Word}
\label{instr:ACSWAPH}\index{acswaph}
 
\paragraph{Encoding} Format \textsf{LSU\_ACSWAPSB} (Section~\ref{sec:format-LSU-ACSWAPSB}), \verb|lsucode5 = 001|\\

\bitfields{ 1/P, 2/01, 3/111, 2/coherency, 6/registerW, 2/11, 3/001, 1/boolcas, 5/registerO, 1/1, 6/registerZ }


\paragraph{Syntax} \texttt{acswaph}\emph{boolcas}\emph{coherency} \emph{registerW}, [\emph{registerZ}] = \emph{registerO}

\paragraph{Description} The effective address is the value of the \emph{registerZ}.
This instruction implements Compare-And-Swap (CAS) on memory half word (16 bits). The \emph{registerO} is interpreted as the concatenation of \texttt{update} (low 64 bits) and \texttt{expected} (high 64 bits). The memory half word at the effective address is read into \texttt{current}. If \texttt{current} equals \texttt{expected}, the swap is successful and \texttt{update} is stored into the memory half word at the effective address. Otherwise, the swap has failed and the memory is not updated. The instruction return value is stored into the \emph{registerW}. It is true/false in case \texttt{boolcas} is true, or the \texttt{current} value in case \texttt{boolcas} is false. The effective address must be half word aligned (multiple of 2).


\paragraph{Modifier(s)}
\noindent\begin{itemize}
\item coherency (cf. coherency in section \ref{par:modifier-coherency} on page \pageref{par:modifier-coherency})
\item boolcas (cf. boolcas in section \ref{par:modifier-boolcas} on page \pageref{par:modifier-boolcas})
\end{itemize}

\paragraph{Execution} {Reservation table \textsf{LSU2\_MEMW\_AUXR\_AUXW} (Section~\ref{sec:reservation-LSU2-MEMW-AUXR-AUXW})}
~
{\begin{lstlisting}
stage ID:
  new boolcas = _ZX_1(boolcas);
  new coherency = _ZX_2(coherency);
stage RR:
  new address = GPR[registerZ];
stage E1:
  new arguments = PGR[registerO];
  new result1 = MEM.atomic_cas(address, 0x3, boolcas | (coherency << 32), arguments.16[0], arguments.16[4], registerW);
stage SR:
  GPR[registerW] = result1;
\end{lstlisting}}
\newpage

\paragraph{Encoding} Format \textsf{LSU\_ACSWAPSB.O} (Section~\ref{sec:format-LSU-ACSWAPSB.O}), \verb|lsucode5 = 001|\\

\bitfields{ 1/P, 2/00, 2/-- --, 27/offset27, 1/1, 2/01, 3/111, 2/coherency, 6/registerW, 2/11, 3/001, 1/boolcas, 5/registerO, 1/1, 6/registerZ }


\paragraph{Syntax} \texttt{acswaph}\emph{boolcas}\emph{coherency} \emph{registerW}, \emph{offset27}[\emph{registerZ}] = \emph{registerO}

\paragraph{Description} The effective address is computed by adding the \emph{offset27} to the \emph{registerZ}.
This instruction implements Compare-And-Swap (CAS) on memory half word (16 bits). The \emph{registerO} is interpreted as the concatenation of \texttt{update} (low 64 bits) and \texttt{expected} (high 64 bits). The memory half word at the effective address is read into \texttt{current}. If \texttt{current} equals \texttt{expected}, the swap is successful and \texttt{update} is stored into the memory half word at the effective address. Otherwise, the swap has failed and the memory is not updated. The instruction return value is stored into the \emph{registerW}. It is true/false in case \texttt{boolcas} is true, or the \texttt{current} value in case \texttt{boolcas} is false. The effective address must be half word aligned (multiple of 2).


\paragraph{Modifier(s)}
\noindent\begin{itemize}
\item coherency (cf. coherency in section \ref{par:modifier-coherency} on page \pageref{par:modifier-coherency})
\item boolcas (cf. boolcas in section \ref{par:modifier-boolcas} on page \pageref{par:modifier-boolcas})
\end{itemize}

\paragraph{Execution} {Reservation table \textsf{LSU2\_MEMW\_AUXR\_AUXW.X} (Section~\ref{sec:reservation-LSU2-MEMW-AUXR-AUXW.X})}
~
{\begin{lstlisting}
stage ID:
  new offset = _SX_27(offset27);
  new boolcas = _ZX_1(boolcas);
  new coherency = _ZX_2(coherency);
stage RR:
  new base = GPR[registerZ];
  new address = base + offset;
stage E1:
  new arguments = PGR[registerO];
  new result1 = MEM.atomic_cas(address, 0x3, boolcas | (coherency << 32), arguments.16[0], arguments.16[4], registerW);
stage SR:
  GPR[registerW] = result1;
\end{lstlisting}}
\newpage

\paragraph{Encoding} Format \textsf{LSU\_ACSWAPSB.D} (Section~\ref{sec:format-LSU-ACSWAPSB.D}), \verb|lsucode5 = 001|\\

\bitfields{ 1/P, 2/00, 2/-- --, 27/extend27, 1/1, 2/00, 2/-- --, 27/offset27, 1/1, 2/01, 3/111, 2/coherency, 6/registerW, 2/11, 3/001, 1/boolcas, 5/registerO, 1/1, 6/registerZ }


\paragraph{Syntax} \texttt{acswaph}\emph{boolcas}\emph{coherency} \emph{registerW}, \emph{extend27\_\discretionary{}{}{}offset27}[\emph{registerZ}] = \emph{registerO}

\paragraph{Description} The effective address is computed by adding the \emph{extend27\_\discretionary{}{}{}offset27} to the \emph{registerZ}.
This instruction implements Compare-And-Swap (CAS) on memory half word (16 bits). The \emph{registerO} is interpreted as the concatenation of \texttt{update} (low 64 bits) and \texttt{expected} (high 64 bits). The memory half word at the effective address is read into \texttt{current}. If \texttt{current} equals \texttt{expected}, the swap is successful and \texttt{update} is stored into the memory half word at the effective address. Otherwise, the swap has failed and the memory is not updated. The instruction return value is stored into the \emph{registerW}. It is true/false in case \texttt{boolcas} is true, or the \texttt{current} value in case \texttt{boolcas} is false. The effective address must be half word aligned (multiple of 2).


\paragraph{Modifier(s)}
\noindent\begin{itemize}
\item coherency (cf. coherency in section \ref{par:modifier-coherency} on page \pageref{par:modifier-coherency})
\item boolcas (cf. boolcas in section \ref{par:modifier-boolcas} on page \pageref{par:modifier-boolcas})
\end{itemize}

\paragraph{Execution} {Reservation table \textsf{LSU2\_MEMW\_AUXR\_AUXW.Y} (Section~\ref{sec:reservation-LSU2-MEMW-AUXR-AUXW.Y})}
~
{\begin{lstlisting}
stage ID:
  new offset = _SX_54(extend27_offset27);
  new boolcas = _ZX_1(boolcas);
  new coherency = _ZX_2(coherency);
stage RR:
  new base = GPR[registerZ];
  new address = base + offset;
stage E1:
  new arguments = PGR[registerO];
  new result1 = MEM.atomic_cas(address, 0x3, boolcas | (coherency << 32), arguments.16[0], arguments.16[4], registerW);
stage SR:
  GPR[registerW] = result1;
\end{lstlisting}}
\newpage
\subsubsection[{\small ACSWAPQ: Atomic Compare and Swap Quadruple Word}]{ACSWAPQ\hfill Atomic Compare and Swap Quadruple Word}
\label{instr:ACSWAPQ}\index{acswapq}
 
\paragraph{Encoding} Format \textsf{LSU\_ACSWAPQB} (Section~\ref{sec:format-LSU-ACSWAPQB}), \verb|lsucode5 = 100|\\

\bitfields{ 1/P, 2/01, 3/111, 2/coherency, 5/registerM, 1/0, 2/11, 3/100, 1/boolcas, 4/registerQ, 1/--, 1/1, 6/registerZ }


\paragraph{Syntax} \texttt{acswapq}\emph{boolcas}\emph{coherency} \emph{registerM}, [\emph{registerZ}] = \emph{registerQ}

\paragraph{Description} The effective address is the value of the \emph{registerZ}.
This instruction implements Compare-And-Swap (CAS) on memory quadruple word (128 bits). The \emph{registerQ} is interpreted as the concatenation of \texttt{update} (low 128 bits) and \texttt{expected} (high 128 bits). The memory double word at the effective address is read into \texttt{current}. If \texttt{current} equals \texttt{expected}, the swap is successful and \texttt{update} is stored into the memory double word at the effective address. Otherwise, the swap has failed and the memory is not updated. The instruction return value is stored into the \emph{registerM}. It is true/false in case \texttt{boolcas} is true, or the \texttt{current} value in case \texttt{boolcas} is false. The effective address must be quadruple word aligned (multiple of 16).


\paragraph{Modifier(s)}
\noindent\begin{itemize}
\item coherency (cf. coherency in section \ref{par:modifier-coherency} on page \pageref{par:modifier-coherency})
\item boolcas (cf. boolcas in section \ref{par:modifier-boolcas} on page \pageref{par:modifier-boolcas})
\end{itemize}

\paragraph{Execution} {Reservation table \textsf{LSU2\_MEMW\_AUXR\_AUXW} (Section~\ref{sec:reservation-LSU2-MEMW-AUXR-AUXW})}
~
{\begin{lstlisting}
stage ID:
  new boolcas = _ZX_1(boolcas);
  new coherency = _ZX_2(coherency);
stage RR:
  new address = GPR[registerZ];
stage E1:
  new arguments = QGR[registerQ];
  new result1 = MEM.atomic_cas(address, 0xFFFF, boolcas | (coherency << 32), arguments.128[0], arguments.128[1], registerM << 1);
stage SR:
  PGR[registerM] = result1;
\end{lstlisting}}
\newpage

\paragraph{Encoding} Format \textsf{LSU\_ACSWAPQB.O} (Section~\ref{sec:format-LSU-ACSWAPQB.O}), \verb|lsucode5 = 100|\\

\bitfields{ 1/P, 2/00, 2/-- --, 27/offset27, 1/1, 2/01, 3/111, 2/coherency, 5/registerM, 1/0, 2/11, 3/100, 1/boolcas, 4/registerQ, 1/--, 1/1, 6/registerZ }


\paragraph{Syntax} \texttt{acswapq}\emph{boolcas}\emph{coherency} \emph{registerM}, \emph{offset27}[\emph{registerZ}] = \emph{registerQ}

\paragraph{Description} The effective address is computed by adding the \emph{offset27} to the \emph{registerZ}.
This instruction implements Compare-And-Swap (CAS) on memory quadruple word (128 bits). The \emph{registerQ} is interpreted as the concatenation of \texttt{update} (low 128 bits) and \texttt{expected} (high 128 bits). The memory double word at the effective address is read into \texttt{current}. If \texttt{current} equals \texttt{expected}, the swap is successful and \texttt{update} is stored into the memory double word at the effective address. Otherwise, the swap has failed and the memory is not updated. The instruction return value is stored into the \emph{registerM}. It is true/false in case \texttt{boolcas} is true, or the \texttt{current} value in case \texttt{boolcas} is false. The effective address must be quadruple word aligned (multiple of 16).


\paragraph{Modifier(s)}
\noindent\begin{itemize}
\item coherency (cf. coherency in section \ref{par:modifier-coherency} on page \pageref{par:modifier-coherency})
\item boolcas (cf. boolcas in section \ref{par:modifier-boolcas} on page \pageref{par:modifier-boolcas})
\end{itemize}

\paragraph{Execution} {Reservation table \textsf{LSU2\_MEMW\_AUXR\_AUXW.X} (Section~\ref{sec:reservation-LSU2-MEMW-AUXR-AUXW.X})}
~
{\begin{lstlisting}
stage ID:
  new offset = _SX_27(offset27);
  new boolcas = _ZX_1(boolcas);
  new coherency = _ZX_2(coherency);
stage RR:
  new base = GPR[registerZ];
  new address = base + offset;
stage E1:
  new arguments = QGR[registerQ];
  new result1 = MEM.atomic_cas(address, 0xFFFF, boolcas | (coherency << 32), arguments.128[0], arguments.128[1], registerM << 1);
stage SR:
  PGR[registerM] = result1;
\end{lstlisting}}
\newpage

\paragraph{Encoding} Format \textsf{LSU\_ACSWAPQB.D} (Section~\ref{sec:format-LSU-ACSWAPQB.D}), \verb|lsucode5 = 100|\\

\bitfields{ 1/P, 2/00, 2/-- --, 27/extend27, 1/1, 2/00, 2/-- --, 27/offset27, 1/1, 2/01, 3/111, 2/coherency, 5/registerM, 1/0, 2/11, 3/100, 1/boolcas, 4/registerQ, 1/--, 1/1, 6/registerZ }


\paragraph{Syntax} \texttt{acswapq}\emph{boolcas}\emph{coherency} \emph{registerM}, \emph{extend27\_\discretionary{}{}{}offset27}[\emph{registerZ}] = \emph{registerQ}

\paragraph{Description} The effective address is computed by adding the \emph{extend27\_\discretionary{}{}{}offset27} to the \emph{registerZ}.
This instruction implements Compare-And-Swap (CAS) on memory quadruple word (128 bits). The \emph{registerQ} is interpreted as the concatenation of \texttt{update} (low 128 bits) and \texttt{expected} (high 128 bits). The memory double word at the effective address is read into \texttt{current}. If \texttt{current} equals \texttt{expected}, the swap is successful and \texttt{update} is stored into the memory double word at the effective address. Otherwise, the swap has failed and the memory is not updated. The instruction return value is stored into the \emph{registerM}. It is true/false in case \texttt{boolcas} is true, or the \texttt{current} value in case \texttt{boolcas} is false. The effective address must be quadruple word aligned (multiple of 16).


\paragraph{Modifier(s)}
\noindent\begin{itemize}
\item coherency (cf. coherency in section \ref{par:modifier-coherency} on page \pageref{par:modifier-coherency})
\item boolcas (cf. boolcas in section \ref{par:modifier-boolcas} on page \pageref{par:modifier-boolcas})
\end{itemize}

\paragraph{Execution} {Reservation table \textsf{LSU2\_MEMW\_AUXR\_AUXW.Y} (Section~\ref{sec:reservation-LSU2-MEMW-AUXR-AUXW.Y})}
~
{\begin{lstlisting}
stage ID:
  new offset = _SX_54(extend27_offset27);
  new boolcas = _ZX_1(boolcas);
  new coherency = _ZX_2(coherency);
stage RR:
  new base = GPR[registerZ];
  new address = base + offset;
stage E1:
  new arguments = QGR[registerQ];
  new result1 = MEM.atomic_cas(address, 0xFFFF, boolcas | (coherency << 32), arguments.128[0], arguments.128[1], registerM << 1);
stage SR:
  PGR[registerM] = result1;
\end{lstlisting}}
\newpage
\subsubsection[{\small ACSWAPW: Atomic Compare and Swap Word}]{ACSWAPW\hfill Atomic Compare and Swap Word}
\label{instr:ACSWAPW}\index{acswapw}
 
\paragraph{Encoding} Format \textsf{LSU\_ACSWAPSB} (Section~\ref{sec:format-LSU-ACSWAPSB}), \verb|lsucode5 = 010|\\

\bitfields{ 1/P, 2/01, 3/111, 2/coherency, 6/registerW, 2/11, 3/010, 1/boolcas, 5/registerO, 1/1, 6/registerZ }


\paragraph{Syntax} \texttt{acswapw}\emph{boolcas}\emph{coherency} \emph{registerW}, [\emph{registerZ}] = \emph{registerO}

\paragraph{Description} The effective address is the value of the \emph{registerZ}.
This instruction implements Compare-And-Swap (CAS) on memory word (32 bits). The \emph{registerO} is interpreted as the concatenation of \texttt{update} (low 64 bits) and \texttt{expected} (high 64 bits). The memory word at the effective address is read into \texttt{current}. If \texttt{current} equals \texttt{expected}, the swap is successful and \texttt{update} is stored into the memory word at the effective address. Otherwise, the swap has failed and the memory is not updated. The instruction return value is stored into the \emph{registerW}. It is true/false in case \texttt{boolcas} is true, or the \texttt{current} value in case \texttt{boolcas} is false. The effective address must be word aligned (multiple of 4).


\paragraph{Modifier(s)}
\noindent\begin{itemize}
\item coherency (cf. coherency in section \ref{par:modifier-coherency} on page \pageref{par:modifier-coherency})
\item boolcas (cf. boolcas in section \ref{par:modifier-boolcas} on page \pageref{par:modifier-boolcas})
\end{itemize}

\paragraph{Execution} {Reservation table \textsf{LSU2\_MEMW\_AUXR\_AUXW} (Section~\ref{sec:reservation-LSU2-MEMW-AUXR-AUXW})}
~
{\begin{lstlisting}
stage ID:
  new boolcas = _ZX_1(boolcas);
  new coherency = _ZX_2(coherency);
stage RR:
  new address = GPR[registerZ];
stage E1:
  new arguments = PGR[registerO];
  new result1 = MEM.atomic_cas(address, 0xF, boolcas | (coherency << 32), arguments.32[0], arguments.32[2], registerW);
stage SR:
  GPR[registerW] = result1;
\end{lstlisting}}
\newpage

\paragraph{Encoding} Format \textsf{LSU\_ACSWAPSB.O} (Section~\ref{sec:format-LSU-ACSWAPSB.O}), \verb|lsucode5 = 010|\\

\bitfields{ 1/P, 2/00, 2/-- --, 27/offset27, 1/1, 2/01, 3/111, 2/coherency, 6/registerW, 2/11, 3/010, 1/boolcas, 5/registerO, 1/1, 6/registerZ }


\paragraph{Syntax} \texttt{acswapw}\emph{boolcas}\emph{coherency} \emph{registerW}, \emph{offset27}[\emph{registerZ}] = \emph{registerO}

\paragraph{Description} The effective address is computed by adding the \emph{offset27} to the \emph{registerZ}.
This instruction implements Compare-And-Swap (CAS) on memory word (32 bits). The \emph{registerO} is interpreted as the concatenation of \texttt{update} (low 64 bits) and \texttt{expected} (high 64 bits). The memory word at the effective address is read into \texttt{current}. If \texttt{current} equals \texttt{expected}, the swap is successful and \texttt{update} is stored into the memory word at the effective address. Otherwise, the swap has failed and the memory is not updated. The instruction return value is stored into the \emph{registerW}. It is true/false in case \texttt{boolcas} is true, or the \texttt{current} value in case \texttt{boolcas} is false. The effective address must be word aligned (multiple of 4).


\paragraph{Modifier(s)}
\noindent\begin{itemize}
\item coherency (cf. coherency in section \ref{par:modifier-coherency} on page \pageref{par:modifier-coherency})
\item boolcas (cf. boolcas in section \ref{par:modifier-boolcas} on page \pageref{par:modifier-boolcas})
\end{itemize}

\paragraph{Execution} {Reservation table \textsf{LSU2\_MEMW\_AUXR\_AUXW.X} (Section~\ref{sec:reservation-LSU2-MEMW-AUXR-AUXW.X})}
~
{\begin{lstlisting}
stage ID:
  new offset = _SX_27(offset27);
  new boolcas = _ZX_1(boolcas);
  new coherency = _ZX_2(coherency);
stage RR:
  new base = GPR[registerZ];
  new address = base + offset;
stage E1:
  new arguments = PGR[registerO];
  new result1 = MEM.atomic_cas(address, 0xF, boolcas | (coherency << 32), arguments.32[0], arguments.32[2], registerW);
stage SR:
  GPR[registerW] = result1;
\end{lstlisting}}
\newpage

\paragraph{Encoding} Format \textsf{LSU\_ACSWAPSB.D} (Section~\ref{sec:format-LSU-ACSWAPSB.D}), \verb|lsucode5 = 010|\\

\bitfields{ 1/P, 2/00, 2/-- --, 27/extend27, 1/1, 2/00, 2/-- --, 27/offset27, 1/1, 2/01, 3/111, 2/coherency, 6/registerW, 2/11, 3/010, 1/boolcas, 5/registerO, 1/1, 6/registerZ }


\paragraph{Syntax} \texttt{acswapw}\emph{boolcas}\emph{coherency} \emph{registerW}, \emph{extend27\_\discretionary{}{}{}offset27}[\emph{registerZ}] = \emph{registerO}

\paragraph{Description} The effective address is computed by adding the \emph{extend27\_\discretionary{}{}{}offset27} to the \emph{registerZ}.
This instruction implements Compare-And-Swap (CAS) on memory word (32 bits). The \emph{registerO} is interpreted as the concatenation of \texttt{update} (low 64 bits) and \texttt{expected} (high 64 bits). The memory word at the effective address is read into \texttt{current}. If \texttt{current} equals \texttt{expected}, the swap is successful and \texttt{update} is stored into the memory word at the effective address. Otherwise, the swap has failed and the memory is not updated. The instruction return value is stored into the \emph{registerW}. It is true/false in case \texttt{boolcas} is true, or the \texttt{current} value in case \texttt{boolcas} is false. The effective address must be word aligned (multiple of 4).


\paragraph{Modifier(s)}
\noindent\begin{itemize}
\item coherency (cf. coherency in section \ref{par:modifier-coherency} on page \pageref{par:modifier-coherency})
\item boolcas (cf. boolcas in section \ref{par:modifier-boolcas} on page \pageref{par:modifier-boolcas})
\end{itemize}

\paragraph{Execution} {Reservation table \textsf{LSU2\_MEMW\_AUXR\_AUXW.Y} (Section~\ref{sec:reservation-LSU2-MEMW-AUXR-AUXW.Y})}
~
{\begin{lstlisting}
stage ID:
  new offset = _SX_54(extend27_offset27);
  new boolcas = _ZX_1(boolcas);
  new coherency = _ZX_2(coherency);
stage RR:
  new base = GPR[registerZ];
  new address = base + offset;
stage E1:
  new arguments = PGR[registerO];
  new result1 = MEM.atomic_cas(address, 0xF, boolcas | (coherency << 32), arguments.32[0], arguments.32[2], registerW);
stage SR:
  GPR[registerW] = result1;
\end{lstlisting}}
\newpage
\subsubsection[{\small ALADDB: Atomic Load and Add Byte}]{ALADDB\hfill Atomic Load and Add Byte}
\label{instr:ALADDB}\index{aladdb}
 
\paragraph{Encoding} Format \textsf{LSU\_ALADDSB} (Section~\ref{sec:format-LSU-ALADDSB}), \verb|lsucode5 = 000|\\

\bitfields{ 1/P, 2/01, 3/111, 2/coherency, 6/registerT, 2/11, 3/000, 1/--, 5/00100, 1/0, 6/registerZ }


\paragraph{Syntax} \texttt{aladdb}\emph{coherency} [\emph{registerZ}] = \emph{registerT}

\paragraph{Description} The effective address is the value of the \emph{registerZ}.
This instruction implements atomic load-add-store on byte. The byte at effective address is added with the \emph{registerT}, and its previous value is returned into the \emph{registerT}.


\paragraph{Modifier(s)}
\noindent\begin{itemize}
\item coherency (cf. coherency in section \ref{par:modifier-coherency} on page \pageref{par:modifier-coherency})
\end{itemize}

\paragraph{Execution} {Reservation table \textsf{LSU2\_MEMW\_AUXR\_AUXW} (Section~\ref{sec:reservation-LSU2-MEMW-AUXR-AUXW})}
~
{\begin{lstlisting}
stage ID:
  new coherency = _ZX_2(coherency);
stage RR:
  new address = GPR[registerZ];
stage E1:
  new argument2 = GPR[registerT];
  new result2 = _ZX_8(MEM.atomic_add(address, 0x1, coherency << 32, argument2.8[0], registerT));
stage SR:
  GPR[registerT] = result2;
\end{lstlisting}}
\newpage

\paragraph{Encoding} Format \textsf{LSU\_ALADDSB.O} (Section~\ref{sec:format-LSU-ALADDSB.O}), \verb|lsucode5 = 000|\\

\bitfields{ 1/P, 2/00, 2/-- --, 27/offset27, 1/1, 2/01, 3/111, 2/coherency, 6/registerT, 2/11, 3/000, 1/--, 5/00100, 1/0, 6/registerZ }


\paragraph{Syntax} \texttt{aladdb}\emph{coherency} \emph{offset27}[\emph{registerZ}] = \emph{registerT}

\paragraph{Description} The effective address is computed by adding the \emph{offset27} to the \emph{registerZ}.
This instruction implements atomic load-add-store on byte. The byte at effective address is added with the \emph{registerT}, and its previous value is returned into the \emph{registerT}.


\paragraph{Modifier(s)}
\noindent\begin{itemize}
\item coherency (cf. coherency in section \ref{par:modifier-coherency} on page \pageref{par:modifier-coherency})
\end{itemize}

\paragraph{Execution} {Reservation table \textsf{LSU2\_MEMW\_AUXR\_AUXW.X} (Section~\ref{sec:reservation-LSU2-MEMW-AUXR-AUXW.X})}
~
{\begin{lstlisting}
stage ID:
  new offset = _SX_27(offset27);
  new coherency = _ZX_2(coherency);
stage RR:
  new base = GPR[registerZ];
  new address = base + offset;
stage E1:
  new argument2 = GPR[registerT];
  new result2 = _ZX_8(MEM.atomic_add(address, 0x1, coherency << 32, argument2.8[0], registerT));
stage SR:
  GPR[registerT] = result2;
\end{lstlisting}}
\newpage

\paragraph{Encoding} Format \textsf{LSU\_ALADDSB.D} (Section~\ref{sec:format-LSU-ALADDSB.D}), \verb|lsucode5 = 000|\\

\bitfields{ 1/P, 2/00, 2/-- --, 27/extend27, 1/1, 2/00, 2/-- --, 27/offset27, 1/1, 2/01, 3/111, 2/coherency, 6/registerT, 2/11, 3/000, 1/--, 5/00100, 1/0, 6/registerZ }


\paragraph{Syntax} \texttt{aladdb}\emph{coherency} \emph{extend27\_\discretionary{}{}{}offset27}[\emph{registerZ}] = \emph{registerT}

\paragraph{Description} The effective address is computed by adding the \emph{extend27\_\discretionary{}{}{}offset27} to the \emph{registerZ}.
This instruction implements atomic load-add-store on byte. The byte at effective address is added with the \emph{registerT}, and its previous value is returned into the \emph{registerT}.


\paragraph{Modifier(s)}
\noindent\begin{itemize}
\item coherency (cf. coherency in section \ref{par:modifier-coherency} on page \pageref{par:modifier-coherency})
\end{itemize}

\paragraph{Execution} {Reservation table \textsf{LSU2\_MEMW\_AUXR\_AUXW.Y} (Section~\ref{sec:reservation-LSU2-MEMW-AUXR-AUXW.Y})}
~
{\begin{lstlisting}
stage ID:
  new offset = _SX_54(extend27_offset27);
  new coherency = _ZX_2(coherency);
stage RR:
  new base = GPR[registerZ];
  new address = base + offset;
stage E1:
  new argument2 = GPR[registerT];
  new result2 = _ZX_8(MEM.atomic_add(address, 0x1, coherency << 32, argument2.8[0], registerT));
stage SR:
  GPR[registerT] = result2;
\end{lstlisting}}
\newpage
\subsubsection[{\small ALADDD: Atomic Load and Add Double Word}]{ALADDD\hfill Atomic Load and Add Double Word}
\label{instr:ALADDD}\index{aladdd}
 
\paragraph{Encoding} Format \textsf{LSU\_ALADDSB} (Section~\ref{sec:format-LSU-ALADDSB}), \verb|lsucode5 = 011|\\

\bitfields{ 1/P, 2/01, 3/111, 2/coherency, 6/registerT, 2/11, 3/011, 1/--, 5/00100, 1/0, 6/registerZ }


\paragraph{Syntax} \texttt{aladdd}\emph{coherency} [\emph{registerZ}] = \emph{registerT}

\paragraph{Description} The effective address is the value of the \emph{registerZ}.
This instruction implements atomic load-add-store on double word. The double word at effective address is added with the \emph{registerT}, and its previous value is returned into the \emph{registerT}. The effective address must be double word aligned (multiple of 8).


\paragraph{Modifier(s)}
\noindent\begin{itemize}
\item coherency (cf. coherency in section \ref{par:modifier-coherency} on page \pageref{par:modifier-coherency})
\end{itemize}

\paragraph{Execution} {Reservation table \textsf{LSU2\_MEMW\_AUXR\_AUXW} (Section~\ref{sec:reservation-LSU2-MEMW-AUXR-AUXW})}
~
{\begin{lstlisting}
stage ID:
  new coherency = _ZX_2(coherency);
stage RR:
  new address = GPR[registerZ];
stage E1:
  new argument2 = GPR[registerT];
  new result2 = MEM.atomic_add(address, 0xFF, coherency << 32, argument2.64[0], registerT);
stage SR:
  GPR[registerT] = result2;
\end{lstlisting}}
\newpage

\paragraph{Encoding} Format \textsf{LSU\_ALADDSB.O} (Section~\ref{sec:format-LSU-ALADDSB.O}), \verb|lsucode5 = 011|\\

\bitfields{ 1/P, 2/00, 2/-- --, 27/offset27, 1/1, 2/01, 3/111, 2/coherency, 6/registerT, 2/11, 3/011, 1/--, 5/00100, 1/0, 6/registerZ }


\paragraph{Syntax} \texttt{aladdd}\emph{coherency} \emph{offset27}[\emph{registerZ}] = \emph{registerT}

\paragraph{Description} The effective address is computed by adding the \emph{offset27} to the \emph{registerZ}.
This instruction implements atomic load-add-store on double word. The double word at effective address is added with the \emph{registerT}, and its previous value is returned into the \emph{registerT}. The effective address must be double word aligned (multiple of 8).


\paragraph{Modifier(s)}
\noindent\begin{itemize}
\item coherency (cf. coherency in section \ref{par:modifier-coherency} on page \pageref{par:modifier-coherency})
\end{itemize}

\paragraph{Execution} {Reservation table \textsf{LSU2\_MEMW\_AUXR\_AUXW.X} (Section~\ref{sec:reservation-LSU2-MEMW-AUXR-AUXW.X})}
~
{\begin{lstlisting}
stage ID:
  new offset = _SX_27(offset27);
  new coherency = _ZX_2(coherency);
stage RR:
  new base = GPR[registerZ];
  new address = base + offset;
stage E1:
  new argument2 = GPR[registerT];
  new result2 = MEM.atomic_add(address, 0xFF, coherency << 32, argument2.64[0], registerT);
stage SR:
  GPR[registerT] = result2;
\end{lstlisting}}
\newpage

\paragraph{Encoding} Format \textsf{LSU\_ALADDSB.D} (Section~\ref{sec:format-LSU-ALADDSB.D}), \verb|lsucode5 = 011|\\

\bitfields{ 1/P, 2/00, 2/-- --, 27/extend27, 1/1, 2/00, 2/-- --, 27/offset27, 1/1, 2/01, 3/111, 2/coherency, 6/registerT, 2/11, 3/011, 1/--, 5/00100, 1/0, 6/registerZ }


\paragraph{Syntax} \texttt{aladdd}\emph{coherency} \emph{extend27\_\discretionary{}{}{}offset27}[\emph{registerZ}] = \emph{registerT}

\paragraph{Description} The effective address is computed by adding the \emph{extend27\_\discretionary{}{}{}offset27} to the \emph{registerZ}.
This instruction implements atomic load-add-store on double word. The double word at effective address is added with the \emph{registerT}, and its previous value is returned into the \emph{registerT}. The effective address must be double word aligned (multiple of 8).


\paragraph{Modifier(s)}
\noindent\begin{itemize}
\item coherency (cf. coherency in section \ref{par:modifier-coherency} on page \pageref{par:modifier-coherency})
\end{itemize}

\paragraph{Execution} {Reservation table \textsf{LSU2\_MEMW\_AUXR\_AUXW.Y} (Section~\ref{sec:reservation-LSU2-MEMW-AUXR-AUXW.Y})}
~
{\begin{lstlisting}
stage ID:
  new offset = _SX_54(extend27_offset27);
  new coherency = _ZX_2(coherency);
stage RR:
  new base = GPR[registerZ];
  new address = base + offset;
stage E1:
  new argument2 = GPR[registerT];
  new result2 = MEM.atomic_add(address, 0xFF, coherency << 32, argument2.64[0], registerT);
stage SR:
  GPR[registerT] = result2;
\end{lstlisting}}
\newpage
\subsubsection[{\small ALADDH: Atomic Load and Add Half Word}]{ALADDH\hfill Atomic Load and Add Half Word}
\label{instr:ALADDH}\index{aladdh}
 
\paragraph{Encoding} Format \textsf{LSU\_ALADDSB} (Section~\ref{sec:format-LSU-ALADDSB}), \verb|lsucode5 = 001|\\

\bitfields{ 1/P, 2/01, 3/111, 2/coherency, 6/registerT, 2/11, 3/001, 1/--, 5/00100, 1/0, 6/registerZ }


\paragraph{Syntax} \texttt{aladdh}\emph{coherency} [\emph{registerZ}] = \emph{registerT}

\paragraph{Description} The effective address is the value of the \emph{registerZ}.
This instruction implements atomic load-add-store on half word. The half word at effective address is added with the \emph{registerT}, and its previous value is returned into the \emph{registerT}. The effective address must be half word aligned (multiple of 2).


\paragraph{Modifier(s)}
\noindent\begin{itemize}
\item coherency (cf. coherency in section \ref{par:modifier-coherency} on page \pageref{par:modifier-coherency})
\end{itemize}

\paragraph{Execution} {Reservation table \textsf{LSU2\_MEMW\_AUXR\_AUXW} (Section~\ref{sec:reservation-LSU2-MEMW-AUXR-AUXW})}
~
{\begin{lstlisting}
stage ID:
  new coherency = _ZX_2(coherency);
stage RR:
  new address = GPR[registerZ];
stage E1:
  new argument2 = GPR[registerT];
  new result2 = _ZX_16(MEM.atomic_add(address, 0x3, coherency << 32, argument2.16[0], registerT));
stage SR:
  GPR[registerT] = result2;
\end{lstlisting}}
\newpage

\paragraph{Encoding} Format \textsf{LSU\_ALADDSB.O} (Section~\ref{sec:format-LSU-ALADDSB.O}), \verb|lsucode5 = 001|\\

\bitfields{ 1/P, 2/00, 2/-- --, 27/offset27, 1/1, 2/01, 3/111, 2/coherency, 6/registerT, 2/11, 3/001, 1/--, 5/00100, 1/0, 6/registerZ }


\paragraph{Syntax} \texttt{aladdh}\emph{coherency} \emph{offset27}[\emph{registerZ}] = \emph{registerT}

\paragraph{Description} The effective address is computed by adding the \emph{offset27} to the \emph{registerZ}.
This instruction implements atomic load-add-store on half word. The half word at effective address is added with the \emph{registerT}, and its previous value is returned into the \emph{registerT}. The effective address must be half word aligned (multiple of 2).


\paragraph{Modifier(s)}
\noindent\begin{itemize}
\item coherency (cf. coherency in section \ref{par:modifier-coherency} on page \pageref{par:modifier-coherency})
\end{itemize}

\paragraph{Execution} {Reservation table \textsf{LSU2\_MEMW\_AUXR\_AUXW.X} (Section~\ref{sec:reservation-LSU2-MEMW-AUXR-AUXW.X})}
~
{\begin{lstlisting}
stage ID:
  new offset = _SX_27(offset27);
  new coherency = _ZX_2(coherency);
stage RR:
  new base = GPR[registerZ];
  new address = base + offset;
stage E1:
  new argument2 = GPR[registerT];
  new result2 = _ZX_16(MEM.atomic_add(address, 0x3, coherency << 32, argument2.16[0], registerT));
stage SR:
  GPR[registerT] = result2;
\end{lstlisting}}
\newpage

\paragraph{Encoding} Format \textsf{LSU\_ALADDSB.D} (Section~\ref{sec:format-LSU-ALADDSB.D}), \verb|lsucode5 = 001|\\

\bitfields{ 1/P, 2/00, 2/-- --, 27/extend27, 1/1, 2/00, 2/-- --, 27/offset27, 1/1, 2/01, 3/111, 2/coherency, 6/registerT, 2/11, 3/001, 1/--, 5/00100, 1/0, 6/registerZ }


\paragraph{Syntax} \texttt{aladdh}\emph{coherency} \emph{extend27\_\discretionary{}{}{}offset27}[\emph{registerZ}] = \emph{registerT}

\paragraph{Description} The effective address is computed by adding the \emph{extend27\_\discretionary{}{}{}offset27} to the \emph{registerZ}.
This instruction implements atomic load-add-store on half word. The half word at effective address is added with the \emph{registerT}, and its previous value is returned into the \emph{registerT}. The effective address must be half word aligned (multiple of 2).


\paragraph{Modifier(s)}
\noindent\begin{itemize}
\item coherency (cf. coherency in section \ref{par:modifier-coherency} on page \pageref{par:modifier-coherency})
\end{itemize}

\paragraph{Execution} {Reservation table \textsf{LSU2\_MEMW\_AUXR\_AUXW.Y} (Section~\ref{sec:reservation-LSU2-MEMW-AUXR-AUXW.Y})}
~
{\begin{lstlisting}
stage ID:
  new offset = _SX_54(extend27_offset27);
  new coherency = _ZX_2(coherency);
stage RR:
  new base = GPR[registerZ];
  new address = base + offset;
stage E1:
  new argument2 = GPR[registerT];
  new result2 = _ZX_16(MEM.atomic_add(address, 0x3, coherency << 32, argument2.16[0], registerT));
stage SR:
  GPR[registerT] = result2;
\end{lstlisting}}
\newpage
\subsubsection[{\small ALADDW: Atomic Load and Add Word}]{ALADDW\hfill Atomic Load and Add Word}
\label{instr:ALADDW}\index{aladdw}
 
\paragraph{Encoding} Format \textsf{LSU\_ALADDSB} (Section~\ref{sec:format-LSU-ALADDSB}), \verb|lsucode5 = 010|\\

\bitfields{ 1/P, 2/01, 3/111, 2/coherency, 6/registerT, 2/11, 3/010, 1/--, 5/00100, 1/0, 6/registerZ }


\paragraph{Syntax} \texttt{aladdw}\emph{coherency} [\emph{registerZ}] = \emph{registerT}

\paragraph{Description} The effective address is the value of the \emph{registerZ}.
This instruction implements atomic load-add-store on word. The word at effective address is added with the \emph{registerT}, and its previous value is returned into the \emph{registerT}. The effective address must be word aligned (multiple of 4).


\paragraph{Modifier(s)}
\noindent\begin{itemize}
\item coherency (cf. coherency in section \ref{par:modifier-coherency} on page \pageref{par:modifier-coherency})
\end{itemize}

\paragraph{Execution} {Reservation table \textsf{LSU2\_MEMW\_AUXR\_AUXW} (Section~\ref{sec:reservation-LSU2-MEMW-AUXR-AUXW})}
~
{\begin{lstlisting}
stage ID:
  new coherency = _ZX_2(coherency);
stage RR:
  new address = GPR[registerZ];
stage E1:
  new argument2 = GPR[registerT];
  new result2 = _ZX_32(MEM.atomic_add(address, 0xF, coherency << 32, argument2.32[0], registerT));
stage SR:
  GPR[registerT] = result2;
\end{lstlisting}}
\newpage

\paragraph{Encoding} Format \textsf{LSU\_ALADDSB.O} (Section~\ref{sec:format-LSU-ALADDSB.O}), \verb|lsucode5 = 010|\\

\bitfields{ 1/P, 2/00, 2/-- --, 27/offset27, 1/1, 2/01, 3/111, 2/coherency, 6/registerT, 2/11, 3/010, 1/--, 5/00100, 1/0, 6/registerZ }


\paragraph{Syntax} \texttt{aladdw}\emph{coherency} \emph{offset27}[\emph{registerZ}] = \emph{registerT}

\paragraph{Description} The effective address is computed by adding the \emph{offset27} to the \emph{registerZ}.
This instruction implements atomic load-add-store on word. The word at effective address is added with the \emph{registerT}, and its previous value is returned into the \emph{registerT}. The effective address must be word aligned (multiple of 4).


\paragraph{Modifier(s)}
\noindent\begin{itemize}
\item coherency (cf. coherency in section \ref{par:modifier-coherency} on page \pageref{par:modifier-coherency})
\end{itemize}

\paragraph{Execution} {Reservation table \textsf{LSU2\_MEMW\_AUXR\_AUXW.X} (Section~\ref{sec:reservation-LSU2-MEMW-AUXR-AUXW.X})}
~
{\begin{lstlisting}
stage ID:
  new offset = _SX_27(offset27);
  new coherency = _ZX_2(coherency);
stage RR:
  new base = GPR[registerZ];
  new address = base + offset;
stage E1:
  new argument2 = GPR[registerT];
  new result2 = _ZX_32(MEM.atomic_add(address, 0xF, coherency << 32, argument2.32[0], registerT));
stage SR:
  GPR[registerT] = result2;
\end{lstlisting}}
\newpage

\paragraph{Encoding} Format \textsf{LSU\_ALADDSB.D} (Section~\ref{sec:format-LSU-ALADDSB.D}), \verb|lsucode5 = 010|\\

\bitfields{ 1/P, 2/00, 2/-- --, 27/extend27, 1/1, 2/00, 2/-- --, 27/offset27, 1/1, 2/01, 3/111, 2/coherency, 6/registerT, 2/11, 3/010, 1/--, 5/00100, 1/0, 6/registerZ }


\paragraph{Syntax} \texttt{aladdw}\emph{coherency} \emph{extend27\_\discretionary{}{}{}offset27}[\emph{registerZ}] = \emph{registerT}

\paragraph{Description} The effective address is computed by adding the \emph{extend27\_\discretionary{}{}{}offset27} to the \emph{registerZ}.
This instruction implements atomic load-add-store on word. The word at effective address is added with the \emph{registerT}, and its previous value is returned into the \emph{registerT}. The effective address must be word aligned (multiple of 4).


\paragraph{Modifier(s)}
\noindent\begin{itemize}
\item coherency (cf. coherency in section \ref{par:modifier-coherency} on page \pageref{par:modifier-coherency})
\end{itemize}

\paragraph{Execution} {Reservation table \textsf{LSU2\_MEMW\_AUXR\_AUXW.Y} (Section~\ref{sec:reservation-LSU2-MEMW-AUXR-AUXW.Y})}
~
{\begin{lstlisting}
stage ID:
  new offset = _SX_54(extend27_offset27);
  new coherency = _ZX_2(coherency);
stage RR:
  new base = GPR[registerZ];
  new address = base + offset;
stage E1:
  new argument2 = GPR[registerT];
  new result2 = _ZX_32(MEM.atomic_add(address, 0xF, coherency << 32, argument2.32[0], registerT));
stage SR:
  GPR[registerT] = result2;
\end{lstlisting}}
\newpage
\subsubsection[{\small ALANDB: Atomic Load and AND Byte}]{ALANDB\hfill Atomic Load and AND Byte}
\label{instr:ALANDB}\index{alandb}
 
\paragraph{Encoding} Format \textsf{LSU\_ALANDSB} (Section~\ref{sec:format-LSU-ALANDSB}), \verb|lsucode5 = 000|\\

\bitfields{ 1/P, 2/01, 3/111, 2/coherency, 6/registerT, 2/11, 3/000, 1/--, 5/00101, 1/0, 6/registerZ }


\paragraph{Syntax} \texttt{alandb}\emph{coherency} [\emph{registerZ}] = \emph{registerT}

\paragraph{Description} The effective address is the value of the \emph{registerZ}.
This instruction implements atomic load-AND-store on byte. The byte at effective address is ANDed with the \emph{registerT}, and its previous value is returned into the \emph{registerT}.


\paragraph{Modifier(s)}
\noindent\begin{itemize}
\item coherency (cf. coherency in section \ref{par:modifier-coherency} on page \pageref{par:modifier-coherency})
\end{itemize}

\paragraph{Execution} {Reservation table \textsf{LSU2\_MEMW\_AUXR\_AUXW} (Section~\ref{sec:reservation-LSU2-MEMW-AUXR-AUXW})}
~
{\begin{lstlisting}
stage ID:
  new coherency = _ZX_2(coherency);
stage RR:
  new address = GPR[registerZ];
stage E1:
  new argument2 = GPR[registerT];
  new result2 = _ZX_8(MEM.atomic_and(address, 0x1, coherency << 32, argument2.8[0], registerT));
stage SR:
  GPR[registerT] = result2;
\end{lstlisting}}
\newpage

\paragraph{Encoding} Format \textsf{LSU\_ALANDSB.O} (Section~\ref{sec:format-LSU-ALANDSB.O}), \verb|lsucode5 = 000|\\

\bitfields{ 1/P, 2/00, 2/-- --, 27/offset27, 1/1, 2/01, 3/111, 2/coherency, 6/registerT, 2/11, 3/000, 1/--, 5/00101, 1/0, 6/registerZ }


\paragraph{Syntax} \texttt{alandb}\emph{coherency} \emph{offset27}[\emph{registerZ}] = \emph{registerT}

\paragraph{Description} The effective address is computed by adding the \emph{offset27} to the \emph{registerZ}.
This instruction implements atomic load-AND-store on byte. The byte at effective address is ANDed with the \emph{registerT}, and its previous value is returned into the \emph{registerT}.


\paragraph{Modifier(s)}
\noindent\begin{itemize}
\item coherency (cf. coherency in section \ref{par:modifier-coherency} on page \pageref{par:modifier-coherency})
\end{itemize}

\paragraph{Execution} {Reservation table \textsf{LSU2\_MEMW\_AUXR\_AUXW.X} (Section~\ref{sec:reservation-LSU2-MEMW-AUXR-AUXW.X})}
~
{\begin{lstlisting}
stage ID:
  new offset = _SX_27(offset27);
  new coherency = _ZX_2(coherency);
stage RR:
  new base = GPR[registerZ];
  new address = base + offset;
stage E1:
  new argument2 = GPR[registerT];
  new result2 = _ZX_8(MEM.atomic_and(address, 0x1, coherency << 32, argument2.8[0], registerT));
stage SR:
  GPR[registerT] = result2;
\end{lstlisting}}
\newpage

\paragraph{Encoding} Format \textsf{LSU\_ALANDSB.D} (Section~\ref{sec:format-LSU-ALANDSB.D}), \verb|lsucode5 = 000|\\

\bitfields{ 1/P, 2/00, 2/-- --, 27/extend27, 1/1, 2/00, 2/-- --, 27/offset27, 1/1, 2/01, 3/111, 2/coherency, 6/registerT, 2/11, 3/000, 1/--, 5/00101, 1/0, 6/registerZ }


\paragraph{Syntax} \texttt{alandb}\emph{coherency} \emph{extend27\_\discretionary{}{}{}offset27}[\emph{registerZ}] = \emph{registerT}

\paragraph{Description} The effective address is computed by adding the \emph{extend27\_\discretionary{}{}{}offset27} to the \emph{registerZ}.
This instruction implements atomic load-AND-store on byte. The byte at effective address is ANDed with the \emph{registerT}, and its previous value is returned into the \emph{registerT}.


\paragraph{Modifier(s)}
\noindent\begin{itemize}
\item coherency (cf. coherency in section \ref{par:modifier-coherency} on page \pageref{par:modifier-coherency})
\end{itemize}

\paragraph{Execution} {Reservation table \textsf{LSU2\_MEMW\_AUXR\_AUXW.Y} (Section~\ref{sec:reservation-LSU2-MEMW-AUXR-AUXW.Y})}
~
{\begin{lstlisting}
stage ID:
  new offset = _SX_54(extend27_offset27);
  new coherency = _ZX_2(coherency);
stage RR:
  new base = GPR[registerZ];
  new address = base + offset;
stage E1:
  new argument2 = GPR[registerT];
  new result2 = _ZX_8(MEM.atomic_and(address, 0x1, coherency << 32, argument2.8[0], registerT));
stage SR:
  GPR[registerT] = result2;
\end{lstlisting}}
\newpage
\subsubsection[{\small ALANDD: Atomic Load and AND Double Word}]{ALANDD\hfill Atomic Load and AND Double Word}
\label{instr:ALANDD}\index{alandd}
 
\paragraph{Encoding} Format \textsf{LSU\_ALANDSB} (Section~\ref{sec:format-LSU-ALANDSB}), \verb|lsucode5 = 011|\\

\bitfields{ 1/P, 2/01, 3/111, 2/coherency, 6/registerT, 2/11, 3/011, 1/--, 5/00101, 1/0, 6/registerZ }


\paragraph{Syntax} \texttt{alandd}\emph{coherency} [\emph{registerZ}] = \emph{registerT}

\paragraph{Description} The effective address is the value of the \emph{registerZ}.
This instruction implements atomic load-AND-store on double word. The double word at effective address is ANDed with the \emph{registerT}, and its previous value is returned into the \emph{registerT}. The effective address must be double word aligned (multiple of 8).


\paragraph{Modifier(s)}
\noindent\begin{itemize}
\item coherency (cf. coherency in section \ref{par:modifier-coherency} on page \pageref{par:modifier-coherency})
\end{itemize}

\paragraph{Execution} {Reservation table \textsf{LSU2\_MEMW\_AUXR\_AUXW} (Section~\ref{sec:reservation-LSU2-MEMW-AUXR-AUXW})}
~
{\begin{lstlisting}
stage ID:
  new coherency = _ZX_2(coherency);
stage RR:
  new address = GPR[registerZ];
stage E1:
  new argument2 = GPR[registerT];
  new result2 = MEM.atomic_and(address, 0xFF, coherency << 32, argument2.64[0], registerT);
stage SR:
  GPR[registerT] = result2;
\end{lstlisting}}
\newpage

\paragraph{Encoding} Format \textsf{LSU\_ALANDSB.O} (Section~\ref{sec:format-LSU-ALANDSB.O}), \verb|lsucode5 = 011|\\

\bitfields{ 1/P, 2/00, 2/-- --, 27/offset27, 1/1, 2/01, 3/111, 2/coherency, 6/registerT, 2/11, 3/011, 1/--, 5/00101, 1/0, 6/registerZ }


\paragraph{Syntax} \texttt{alandd}\emph{coherency} \emph{offset27}[\emph{registerZ}] = \emph{registerT}

\paragraph{Description} The effective address is computed by adding the \emph{offset27} to the \emph{registerZ}.
This instruction implements atomic load-AND-store on double word. The double word at effective address is ANDed with the \emph{registerT}, and its previous value is returned into the \emph{registerT}. The effective address must be double word aligned (multiple of 8).


\paragraph{Modifier(s)}
\noindent\begin{itemize}
\item coherency (cf. coherency in section \ref{par:modifier-coherency} on page \pageref{par:modifier-coherency})
\end{itemize}

\paragraph{Execution} {Reservation table \textsf{LSU2\_MEMW\_AUXR\_AUXW.X} (Section~\ref{sec:reservation-LSU2-MEMW-AUXR-AUXW.X})}
~
{\begin{lstlisting}
stage ID:
  new offset = _SX_27(offset27);
  new coherency = _ZX_2(coherency);
stage RR:
  new base = GPR[registerZ];
  new address = base + offset;
stage E1:
  new argument2 = GPR[registerT];
  new result2 = MEM.atomic_and(address, 0xFF, coherency << 32, argument2.64[0], registerT);
stage SR:
  GPR[registerT] = result2;
\end{lstlisting}}
\newpage

\paragraph{Encoding} Format \textsf{LSU\_ALANDSB.D} (Section~\ref{sec:format-LSU-ALANDSB.D}), \verb|lsucode5 = 011|\\

\bitfields{ 1/P, 2/00, 2/-- --, 27/extend27, 1/1, 2/00, 2/-- --, 27/offset27, 1/1, 2/01, 3/111, 2/coherency, 6/registerT, 2/11, 3/011, 1/--, 5/00101, 1/0, 6/registerZ }


\paragraph{Syntax} \texttt{alandd}\emph{coherency} \emph{extend27\_\discretionary{}{}{}offset27}[\emph{registerZ}] = \emph{registerT}

\paragraph{Description} The effective address is computed by adding the \emph{extend27\_\discretionary{}{}{}offset27} to the \emph{registerZ}.
This instruction implements atomic load-AND-store on double word. The double word at effective address is ANDed with the \emph{registerT}, and its previous value is returned into the \emph{registerT}. The effective address must be double word aligned (multiple of 8).


\paragraph{Modifier(s)}
\noindent\begin{itemize}
\item coherency (cf. coherency in section \ref{par:modifier-coherency} on page \pageref{par:modifier-coherency})
\end{itemize}

\paragraph{Execution} {Reservation table \textsf{LSU2\_MEMW\_AUXR\_AUXW.Y} (Section~\ref{sec:reservation-LSU2-MEMW-AUXR-AUXW.Y})}
~
{\begin{lstlisting}
stage ID:
  new offset = _SX_54(extend27_offset27);
  new coherency = _ZX_2(coherency);
stage RR:
  new base = GPR[registerZ];
  new address = base + offset;
stage E1:
  new argument2 = GPR[registerT];
  new result2 = MEM.atomic_and(address, 0xFF, coherency << 32, argument2.64[0], registerT);
stage SR:
  GPR[registerT] = result2;
\end{lstlisting}}
\newpage
\subsubsection[{\small ALANDH: Atomic Load and AND Half Word}]{ALANDH\hfill Atomic Load and AND Half Word}
\label{instr:ALANDH}\index{alandh}
 
\paragraph{Encoding} Format \textsf{LSU\_ALANDSB} (Section~\ref{sec:format-LSU-ALANDSB}), \verb|lsucode5 = 001|\\

\bitfields{ 1/P, 2/01, 3/111, 2/coherency, 6/registerT, 2/11, 3/001, 1/--, 5/00101, 1/0, 6/registerZ }


\paragraph{Syntax} \texttt{alandh}\emph{coherency} [\emph{registerZ}] = \emph{registerT}

\paragraph{Description} The effective address is the value of the \emph{registerZ}.
This instruction implements atomic load-AND-store on half word. The half word at effective address is ANDed with the \emph{registerT}, and its previous value is returned into the \emph{registerT}. The effective address must be half word aligned (multiple of 2).


\paragraph{Modifier(s)}
\noindent\begin{itemize}
\item coherency (cf. coherency in section \ref{par:modifier-coherency} on page \pageref{par:modifier-coherency})
\end{itemize}

\paragraph{Execution} {Reservation table \textsf{LSU2\_MEMW\_AUXR\_AUXW} (Section~\ref{sec:reservation-LSU2-MEMW-AUXR-AUXW})}
~
{\begin{lstlisting}
stage ID:
  new coherency = _ZX_2(coherency);
stage RR:
  new address = GPR[registerZ];
stage E1:
  new argument2 = GPR[registerT];
  new result2 = _ZX_16(MEM.atomic_and(address, 0x3, coherency << 32, argument2.16[0], registerT));
stage SR:
  GPR[registerT] = result2;
\end{lstlisting}}
\newpage

\paragraph{Encoding} Format \textsf{LSU\_ALANDSB.O} (Section~\ref{sec:format-LSU-ALANDSB.O}), \verb|lsucode5 = 001|\\

\bitfields{ 1/P, 2/00, 2/-- --, 27/offset27, 1/1, 2/01, 3/111, 2/coherency, 6/registerT, 2/11, 3/001, 1/--, 5/00101, 1/0, 6/registerZ }


\paragraph{Syntax} \texttt{alandh}\emph{coherency} \emph{offset27}[\emph{registerZ}] = \emph{registerT}

\paragraph{Description} The effective address is computed by adding the \emph{offset27} to the \emph{registerZ}.
This instruction implements atomic load-AND-store on half word. The half word at effective address is ANDed with the \emph{registerT}, and its previous value is returned into the \emph{registerT}. The effective address must be half word aligned (multiple of 2).


\paragraph{Modifier(s)}
\noindent\begin{itemize}
\item coherency (cf. coherency in section \ref{par:modifier-coherency} on page \pageref{par:modifier-coherency})
\end{itemize}

\paragraph{Execution} {Reservation table \textsf{LSU2\_MEMW\_AUXR\_AUXW.X} (Section~\ref{sec:reservation-LSU2-MEMW-AUXR-AUXW.X})}
~
{\begin{lstlisting}
stage ID:
  new offset = _SX_27(offset27);
  new coherency = _ZX_2(coherency);
stage RR:
  new base = GPR[registerZ];
  new address = base + offset;
stage E1:
  new argument2 = GPR[registerT];
  new result2 = _ZX_16(MEM.atomic_and(address, 0x3, coherency << 32, argument2.16[0], registerT));
stage SR:
  GPR[registerT] = result2;
\end{lstlisting}}
\newpage

\paragraph{Encoding} Format \textsf{LSU\_ALANDSB.D} (Section~\ref{sec:format-LSU-ALANDSB.D}), \verb|lsucode5 = 001|\\

\bitfields{ 1/P, 2/00, 2/-- --, 27/extend27, 1/1, 2/00, 2/-- --, 27/offset27, 1/1, 2/01, 3/111, 2/coherency, 6/registerT, 2/11, 3/001, 1/--, 5/00101, 1/0, 6/registerZ }


\paragraph{Syntax} \texttt{alandh}\emph{coherency} \emph{extend27\_\discretionary{}{}{}offset27}[\emph{registerZ}] = \emph{registerT}

\paragraph{Description} The effective address is computed by adding the \emph{extend27\_\discretionary{}{}{}offset27} to the \emph{registerZ}.
This instruction implements atomic load-AND-store on half word. The half word at effective address is ANDed with the \emph{registerT}, and its previous value is returned into the \emph{registerT}. The effective address must be half word aligned (multiple of 2).


\paragraph{Modifier(s)}
\noindent\begin{itemize}
\item coherency (cf. coherency in section \ref{par:modifier-coherency} on page \pageref{par:modifier-coherency})
\end{itemize}

\paragraph{Execution} {Reservation table \textsf{LSU2\_MEMW\_AUXR\_AUXW.Y} (Section~\ref{sec:reservation-LSU2-MEMW-AUXR-AUXW.Y})}
~
{\begin{lstlisting}
stage ID:
  new offset = _SX_54(extend27_offset27);
  new coherency = _ZX_2(coherency);
stage RR:
  new base = GPR[registerZ];
  new address = base + offset;
stage E1:
  new argument2 = GPR[registerT];
  new result2 = _ZX_16(MEM.atomic_and(address, 0x3, coherency << 32, argument2.16[0], registerT));
stage SR:
  GPR[registerT] = result2;
\end{lstlisting}}
\newpage
\subsubsection[{\small ALANDW: Atomic Load and AND Word}]{ALANDW\hfill Atomic Load and AND Word}
\label{instr:ALANDW}\index{alandw}
 
\paragraph{Encoding} Format \textsf{LSU\_ALANDSB} (Section~\ref{sec:format-LSU-ALANDSB}), \verb|lsucode5 = 010|\\

\bitfields{ 1/P, 2/01, 3/111, 2/coherency, 6/registerT, 2/11, 3/010, 1/--, 5/00101, 1/0, 6/registerZ }


\paragraph{Syntax} \texttt{alandw}\emph{coherency} [\emph{registerZ}] = \emph{registerT}

\paragraph{Description} The effective address is the value of the \emph{registerZ}.
This instruction implements atomic load-AND-store on word. The word at effective address is ANDed with the \emph{registerT}, and its previous value is returned into the \emph{registerT}. The effective address must be word aligned (multiple of 4).


\paragraph{Modifier(s)}
\noindent\begin{itemize}
\item coherency (cf. coherency in section \ref{par:modifier-coherency} on page \pageref{par:modifier-coherency})
\end{itemize}

\paragraph{Execution} {Reservation table \textsf{LSU2\_MEMW\_AUXR\_AUXW} (Section~\ref{sec:reservation-LSU2-MEMW-AUXR-AUXW})}
~
{\begin{lstlisting}
stage ID:
  new coherency = _ZX_2(coherency);
stage RR:
  new address = GPR[registerZ];
stage E1:
  new argument2 = GPR[registerT];
  new result2 = _ZX_32(MEM.atomic_and(address, 0xF, coherency << 32, argument2.32[0], registerT));
stage SR:
  GPR[registerT] = result2;
\end{lstlisting}}
\newpage

\paragraph{Encoding} Format \textsf{LSU\_ALANDSB.O} (Section~\ref{sec:format-LSU-ALANDSB.O}), \verb|lsucode5 = 010|\\

\bitfields{ 1/P, 2/00, 2/-- --, 27/offset27, 1/1, 2/01, 3/111, 2/coherency, 6/registerT, 2/11, 3/010, 1/--, 5/00101, 1/0, 6/registerZ }


\paragraph{Syntax} \texttt{alandw}\emph{coherency} \emph{offset27}[\emph{registerZ}] = \emph{registerT}

\paragraph{Description} The effective address is computed by adding the \emph{offset27} to the \emph{registerZ}.
This instruction implements atomic load-AND-store on word. The word at effective address is ANDed with the \emph{registerT}, and its previous value is returned into the \emph{registerT}. The effective address must be word aligned (multiple of 4).


\paragraph{Modifier(s)}
\noindent\begin{itemize}
\item coherency (cf. coherency in section \ref{par:modifier-coherency} on page \pageref{par:modifier-coherency})
\end{itemize}

\paragraph{Execution} {Reservation table \textsf{LSU2\_MEMW\_AUXR\_AUXW.X} (Section~\ref{sec:reservation-LSU2-MEMW-AUXR-AUXW.X})}
~
{\begin{lstlisting}
stage ID:
  new offset = _SX_27(offset27);
  new coherency = _ZX_2(coherency);
stage RR:
  new base = GPR[registerZ];
  new address = base + offset;
stage E1:
  new argument2 = GPR[registerT];
  new result2 = _ZX_32(MEM.atomic_and(address, 0xF, coherency << 32, argument2.32[0], registerT));
stage SR:
  GPR[registerT] = result2;
\end{lstlisting}}
\newpage

\paragraph{Encoding} Format \textsf{LSU\_ALANDSB.D} (Section~\ref{sec:format-LSU-ALANDSB.D}), \verb|lsucode5 = 010|\\

\bitfields{ 1/P, 2/00, 2/-- --, 27/extend27, 1/1, 2/00, 2/-- --, 27/offset27, 1/1, 2/01, 3/111, 2/coherency, 6/registerT, 2/11, 3/010, 1/--, 5/00101, 1/0, 6/registerZ }


\paragraph{Syntax} \texttt{alandw}\emph{coherency} \emph{extend27\_\discretionary{}{}{}offset27}[\emph{registerZ}] = \emph{registerT}

\paragraph{Description} The effective address is computed by adding the \emph{extend27\_\discretionary{}{}{}offset27} to the \emph{registerZ}.
This instruction implements atomic load-AND-store on word. The word at effective address is ANDed with the \emph{registerT}, and its previous value is returned into the \emph{registerT}. The effective address must be word aligned (multiple of 4).


\paragraph{Modifier(s)}
\noindent\begin{itemize}
\item coherency (cf. coherency in section \ref{par:modifier-coherency} on page \pageref{par:modifier-coherency})
\end{itemize}

\paragraph{Execution} {Reservation table \textsf{LSU2\_MEMW\_AUXR\_AUXW.Y} (Section~\ref{sec:reservation-LSU2-MEMW-AUXR-AUXW.Y})}
~
{\begin{lstlisting}
stage ID:
  new offset = _SX_54(extend27_offset27);
  new coherency = _ZX_2(coherency);
stage RR:
  new base = GPR[registerZ];
  new address = base + offset;
stage E1:
  new argument2 = GPR[registerT];
  new result2 = _ZX_32(MEM.atomic_and(address, 0xF, coherency << 32, argument2.32[0], registerT));
stage SR:
  GPR[registerT] = result2;
\end{lstlisting}}
\newpage
\subsubsection[{\small ALB: Atomic Load Byte}]{ALB\hfill Atomic Load Byte}
\label{instr:ALB}\index{alb}
 
\paragraph{Encoding} Format \textsf{LSU\_ALSB} (Section~\ref{sec:format-LSU-ALSB}), \verb|lsucode5 = 000|\\

\bitfields{ 1/P, 2/01, 3/111, 2/coherency, 6/registerW, 2/11, 3/000, 1/--, 5/00000, 1/0, 6/registerZ }


\paragraph{Syntax} \texttt{alb}\emph{coherency} \emph{registerW} = [\emph{registerZ}]

\paragraph{Description} The effective address is the value of the \emph{registerZ}.
The \emph{registerW} is atomically loaded with zero extension from the memory byte starting at the effective address. This instruction is provided to directly operate on the last level of the memory hierarchy, which is typically a processor DDR or a cluster TCM, irrespective of the MMU DCP (Data Cache Policy) settings.


\paragraph{Modifier(s)}
\noindent\begin{itemize}
\item coherency (cf. coherency in section \ref{par:modifier-coherency} on page \pageref{par:modifier-coherency})
\end{itemize}

\paragraph{Execution} {Reservation table \textsf{LSU\_MEMW\_AUXW} (Section~\ref{sec:reservation-LSU-MEMW-AUXW})}
~
{\begin{lstlisting}
stage ID:
  new coherency = _ZX_2(coherency);
stage RR:
  new address = GPR[registerZ];
  new result1 = _ZX_8(MEM.atomic_load(address, 0x1, coherency << 32, registerW));
stage CR:
  GPR[registerW] = result1;
\end{lstlisting}}
\newpage

\paragraph{Encoding} Format \textsf{LSU\_ALSB.O} (Section~\ref{sec:format-LSU-ALSB.O}), \verb|lsucode5 = 000|\\

\bitfields{ 1/P, 2/00, 2/-- --, 27/offset27, 1/1, 2/01, 3/111, 2/coherency, 6/registerW, 2/11, 3/000, 1/--, 5/00000, 1/0, 6/registerZ }


\paragraph{Syntax} \texttt{alb}\emph{coherency} \emph{registerW} = \emph{offset27}[\emph{registerZ}]

\paragraph{Description} The effective address is computed by adding the \emph{offset27} to the \emph{registerZ}.
The \emph{registerW} is atomically loaded with zero extension from the memory byte starting at the effective address. This instruction is provided to directly operate on the last level of the memory hierarchy, which is typically a processor DDR or a cluster TCM, irrespective of the MMU DCP (Data Cache Policy) settings.


\paragraph{Modifier(s)}
\noindent\begin{itemize}
\item coherency (cf. coherency in section \ref{par:modifier-coherency} on page \pageref{par:modifier-coherency})
\end{itemize}

\paragraph{Execution} {Reservation table \textsf{LSU\_MEMW\_AUXW.X} (Section~\ref{sec:reservation-LSU-MEMW-AUXW.X})}
~
{\begin{lstlisting}
stage ID:
  new offset = _SX_27(offset27);
  new coherency = _ZX_2(coherency);
stage RR:
  new base = GPR[registerZ];
  new address = base + offset;
  new result1 = _ZX_8(MEM.atomic_load(address, 0x1, coherency << 32, registerW));
stage CR:
  GPR[registerW] = result1;
\end{lstlisting}}
\newpage

\paragraph{Encoding} Format \textsf{LSU\_ALSB.D} (Section~\ref{sec:format-LSU-ALSB.D}), \verb|lsucode5 = 000|\\

\bitfields{ 1/P, 2/00, 2/-- --, 27/extend27, 1/1, 2/00, 2/-- --, 27/offset27, 1/1, 2/01, 3/111, 2/coherency, 6/registerW, 2/11, 3/000, 1/--, 5/00000, 1/0, 6/registerZ }


\paragraph{Syntax} \texttt{alb}\emph{coherency} \emph{registerW} = \emph{extend27\_\discretionary{}{}{}offset27}[\emph{registerZ}]

\paragraph{Description} The effective address is computed by adding the \emph{extend27\_\discretionary{}{}{}offset27} to the \emph{registerZ}.
The \emph{registerW} is atomically loaded with zero extension from the memory byte starting at the effective address. This instruction is provided to directly operate on the last level of the memory hierarchy, which is typically a processor DDR or a cluster TCM, irrespective of the MMU DCP (Data Cache Policy) settings.


\paragraph{Modifier(s)}
\noindent\begin{itemize}
\item coherency (cf. coherency in section \ref{par:modifier-coherency} on page \pageref{par:modifier-coherency})
\end{itemize}

\paragraph{Execution} {Reservation table \textsf{LSU\_MEMW\_AUXW.Y} (Section~\ref{sec:reservation-LSU-MEMW-AUXW.Y})}
~
{\begin{lstlisting}
stage ID:
  new offset = _SX_54(extend27_offset27);
  new coherency = _ZX_2(coherency);
stage RR:
  new base = GPR[registerZ];
  new address = base + offset;
  new result1 = _ZX_8(MEM.atomic_load(address, 0x1, coherency << 32, registerW));
stage CR:
  GPR[registerW] = result1;
\end{lstlisting}}
\newpage
\subsubsection[{\small ALCLRB: Atomic Load and Clear Byte}]{ALCLRB\hfill Atomic Load and Clear Byte}
\label{instr:ALCLRB}\index{alclrb}
 
\paragraph{Encoding} Format \textsf{LSU\_ALCLRSB} (Section~\ref{sec:format-LSU-ALCLRSB}), \verb|lsucode5 = 000|\\

\bitfields{ 1/P, 2/01, 3/111, 2/coherency, 6/registerW, 2/11, 3/000, 1/--, 5/00001, 1/0, 6/registerZ }


\paragraph{Syntax} \texttt{alclrb}\emph{coherency} \emph{registerW} = [\emph{registerZ}]

\paragraph{Description} The effective address is the value of the \emph{registerZ}.
The \emph{registerW} is loaded from the memory byte starting at the effective address. The value 0 is stored into the memory byte starting at the effective address. This atomic instruction implements Test-and-Clear in memory. A zero byte in memory means that the lock at the effective address is taken. Unlocking is achieved by storing a non-zero byte at the effective address.


\paragraph{Modifier(s)}
\noindent\begin{itemize}
\item coherency (cf. coherency in section \ref{par:modifier-coherency} on page \pageref{par:modifier-coherency})
\end{itemize}

\paragraph{Execution} {Reservation table \textsf{LSU2\_MEMW\_AUXW} (Section~\ref{sec:reservation-LSU2-MEMW-AUXW})}
~
{\begin{lstlisting}
stage ID:
  new coherency = _ZX_2(coherency);
stage RR:
  new address = GPR[registerZ];
  new result1 = _ZX_8(MEM.atomic_swap(address, 0x1, coherency << 32, 0, registerW));
stage CR:
  GPR[registerW] = result1;
\end{lstlisting}}
\newpage

\paragraph{Encoding} Format \textsf{LSU\_ALCLRSB.O} (Section~\ref{sec:format-LSU-ALCLRSB.O}), \verb|lsucode5 = 000|\\

\bitfields{ 1/P, 2/00, 2/-- --, 27/offset27, 1/1, 2/01, 3/111, 2/coherency, 6/registerW, 2/11, 3/000, 1/--, 5/00001, 1/0, 6/registerZ }


\paragraph{Syntax} \texttt{alclrb}\emph{coherency} \emph{registerW} = \emph{offset27}[\emph{registerZ}]

\paragraph{Description} The effective address is computed by adding the \emph{offset27} to the \emph{registerZ}.
The \emph{registerW} is loaded from the memory byte starting at the effective address. The value 0 is stored into the memory byte starting at the effective address. This atomic instruction implements Test-and-Clear in memory. A zero byte in memory means that the lock at the effective address is taken. Unlocking is achieved by storing a non-zero byte at the effective address.


\paragraph{Modifier(s)}
\noindent\begin{itemize}
\item coherency (cf. coherency in section \ref{par:modifier-coherency} on page \pageref{par:modifier-coherency})
\end{itemize}

\paragraph{Execution} {Reservation table \textsf{LSU2\_MEMW\_AUXW.X} (Section~\ref{sec:reservation-LSU2-MEMW-AUXW.X})}
~
{\begin{lstlisting}
stage ID:
  new offset = _SX_27(offset27);
  new coherency = _ZX_2(coherency);
stage RR:
  new base = GPR[registerZ];
  new address = base + offset;
  new result1 = _ZX_8(MEM.atomic_swap(address, 0x1, coherency << 32, 0, registerW));
stage CR:
  GPR[registerW] = result1;
\end{lstlisting}}
\newpage

\paragraph{Encoding} Format \textsf{LSU\_ALCLRSB.D} (Section~\ref{sec:format-LSU-ALCLRSB.D}), \verb|lsucode5 = 000|\\

\bitfields{ 1/P, 2/00, 2/-- --, 27/extend27, 1/1, 2/00, 2/-- --, 27/offset27, 1/1, 2/01, 3/111, 2/coherency, 6/registerW, 2/11, 3/000, 1/--, 5/00001, 1/0, 6/registerZ }


\paragraph{Syntax} \texttt{alclrb}\emph{coherency} \emph{registerW} = \emph{extend27\_\discretionary{}{}{}offset27}[\emph{registerZ}]

\paragraph{Description} The effective address is computed by adding the \emph{extend27\_\discretionary{}{}{}offset27} to the \emph{registerZ}.
The \emph{registerW} is loaded from the memory byte starting at the effective address. The value 0 is stored into the memory byte starting at the effective address. This atomic instruction implements Test-and-Clear in memory. A zero byte in memory means that the lock at the effective address is taken. Unlocking is achieved by storing a non-zero byte at the effective address.


\paragraph{Modifier(s)}
\noindent\begin{itemize}
\item coherency (cf. coherency in section \ref{par:modifier-coherency} on page \pageref{par:modifier-coherency})
\end{itemize}

\paragraph{Execution} {Reservation table \textsf{LSU2\_MEMW\_AUXW.Y} (Section~\ref{sec:reservation-LSU2-MEMW-AUXW.Y})}
~
{\begin{lstlisting}
stage ID:
  new offset = _SX_54(extend27_offset27);
  new coherency = _ZX_2(coherency);
stage RR:
  new base = GPR[registerZ];
  new address = base + offset;
  new result1 = _ZX_8(MEM.atomic_swap(address, 0x1, coherency << 32, 0, registerW));
stage CR:
  GPR[registerW] = result1;
\end{lstlisting}}
\newpage
\subsubsection[{\small ALCLRD: Atomic Load and Clear Double Word}]{ALCLRD\hfill Atomic Load and Clear Double Word}
\label{instr:ALCLRD}\index{alclrd}
 
\paragraph{Encoding} Format \textsf{LSU\_ALCLRSB} (Section~\ref{sec:format-LSU-ALCLRSB}), \verb|lsucode5 = 011|\\

\bitfields{ 1/P, 2/01, 3/111, 2/coherency, 6/registerW, 2/11, 3/011, 1/--, 5/00001, 1/0, 6/registerZ }


\paragraph{Syntax} \texttt{alclrd}\emph{coherency} \emph{registerW} = [\emph{registerZ}]

\paragraph{Description} The effective address is the value of the \emph{registerZ}.
The \emph{registerW} is loaded from the memory double word starting at the effective address. The value 0 is stored into the memory double word starting at the effective address. This atomic instruction implements Test-and-Clear in memory. A zero double word in memory means that the lock at the effective address is taken. Unlocking is achieved by storing a non-zero double word at the effective address.


\paragraph{Modifier(s)}
\noindent\begin{itemize}
\item coherency (cf. coherency in section \ref{par:modifier-coherency} on page \pageref{par:modifier-coherency})
\end{itemize}

\paragraph{Execution} {Reservation table \textsf{LSU2\_MEMW\_AUXW} (Section~\ref{sec:reservation-LSU2-MEMW-AUXW})}
~
{\begin{lstlisting}
stage ID:
  new coherency = _ZX_2(coherency);
stage RR:
  new address = GPR[registerZ];
  new result1 = MEM.atomic_swap(address, 0xFF, coherency << 32, 0, registerW);
stage CR:
  GPR[registerW] = result1;
\end{lstlisting}}
\newpage

\paragraph{Encoding} Format \textsf{LSU\_ALCLRSB.O} (Section~\ref{sec:format-LSU-ALCLRSB.O}), \verb|lsucode5 = 011|\\

\bitfields{ 1/P, 2/00, 2/-- --, 27/offset27, 1/1, 2/01, 3/111, 2/coherency, 6/registerW, 2/11, 3/011, 1/--, 5/00001, 1/0, 6/registerZ }


\paragraph{Syntax} \texttt{alclrd}\emph{coherency} \emph{registerW} = \emph{offset27}[\emph{registerZ}]

\paragraph{Description} The effective address is computed by adding the \emph{offset27} to the \emph{registerZ}.
The \emph{registerW} is loaded from the memory double word starting at the effective address. The value 0 is stored into the memory double word starting at the effective address. This atomic instruction implements Test-and-Clear in memory. A zero double word in memory means that the lock at the effective address is taken. Unlocking is achieved by storing a non-zero double word at the effective address.


\paragraph{Modifier(s)}
\noindent\begin{itemize}
\item coherency (cf. coherency in section \ref{par:modifier-coherency} on page \pageref{par:modifier-coherency})
\end{itemize}

\paragraph{Execution} {Reservation table \textsf{LSU2\_MEMW\_AUXW.X} (Section~\ref{sec:reservation-LSU2-MEMW-AUXW.X})}
~
{\begin{lstlisting}
stage ID:
  new offset = _SX_27(offset27);
  new coherency = _ZX_2(coherency);
stage RR:
  new base = GPR[registerZ];
  new address = base + offset;
  new result1 = MEM.atomic_swap(address, 0xFF, coherency << 32, 0, registerW);
stage CR:
  GPR[registerW] = result1;
\end{lstlisting}}
\newpage

\paragraph{Encoding} Format \textsf{LSU\_ALCLRSB.D} (Section~\ref{sec:format-LSU-ALCLRSB.D}), \verb|lsucode5 = 011|\\

\bitfields{ 1/P, 2/00, 2/-- --, 27/extend27, 1/1, 2/00, 2/-- --, 27/offset27, 1/1, 2/01, 3/111, 2/coherency, 6/registerW, 2/11, 3/011, 1/--, 5/00001, 1/0, 6/registerZ }


\paragraph{Syntax} \texttt{alclrd}\emph{coherency} \emph{registerW} = \emph{extend27\_\discretionary{}{}{}offset27}[\emph{registerZ}]

\paragraph{Description} The effective address is computed by adding the \emph{extend27\_\discretionary{}{}{}offset27} to the \emph{registerZ}.
The \emph{registerW} is loaded from the memory double word starting at the effective address. The value 0 is stored into the memory double word starting at the effective address. This atomic instruction implements Test-and-Clear in memory. A zero double word in memory means that the lock at the effective address is taken. Unlocking is achieved by storing a non-zero double word at the effective address.


\paragraph{Modifier(s)}
\noindent\begin{itemize}
\item coherency (cf. coherency in section \ref{par:modifier-coherency} on page \pageref{par:modifier-coherency})
\end{itemize}

\paragraph{Execution} {Reservation table \textsf{LSU2\_MEMW\_AUXW.Y} (Section~\ref{sec:reservation-LSU2-MEMW-AUXW.Y})}
~
{\begin{lstlisting}
stage ID:
  new offset = _SX_54(extend27_offset27);
  new coherency = _ZX_2(coherency);
stage RR:
  new base = GPR[registerZ];
  new address = base + offset;
  new result1 = MEM.atomic_swap(address, 0xFF, coherency << 32, 0, registerW);
stage CR:
  GPR[registerW] = result1;
\end{lstlisting}}
\newpage
\subsubsection[{\small ALCLRH: Atomic Load and Clear Half Word}]{ALCLRH\hfill Atomic Load and Clear Half Word}
\label{instr:ALCLRH}\index{alclrh}
 
\paragraph{Encoding} Format \textsf{LSU\_ALCLRSB} (Section~\ref{sec:format-LSU-ALCLRSB}), \verb|lsucode5 = 001|\\

\bitfields{ 1/P, 2/01, 3/111, 2/coherency, 6/registerW, 2/11, 3/001, 1/--, 5/00001, 1/0, 6/registerZ }


\paragraph{Syntax} \texttt{alclrh}\emph{coherency} \emph{registerW} = [\emph{registerZ}]

\paragraph{Description} The effective address is the value of the \emph{registerZ}.
The \emph{registerW} is loaded from the memory half word starting at the effective address. The value 0 is stored into the memory half word starting at the effective address. This atomic instruction implements Test-and-Clear in memory. A zero half word in memory means that the lock at the effective address is taken. Unlocking is achieved by storing a non-zero half word at the effective address.


\paragraph{Modifier(s)}
\noindent\begin{itemize}
\item coherency (cf. coherency in section \ref{par:modifier-coherency} on page \pageref{par:modifier-coherency})
\end{itemize}

\paragraph{Execution} {Reservation table \textsf{LSU2\_MEMW\_AUXW} (Section~\ref{sec:reservation-LSU2-MEMW-AUXW})}
~
{\begin{lstlisting}
stage ID:
  new coherency = _ZX_2(coherency);
stage RR:
  new address = GPR[registerZ];
  new result1 = _ZX_16(MEM.atomic_swap(address, 0x3, coherency << 32, 0, registerW));
stage CR:
  GPR[registerW] = result1;
\end{lstlisting}}
\newpage

\paragraph{Encoding} Format \textsf{LSU\_ALCLRSB.O} (Section~\ref{sec:format-LSU-ALCLRSB.O}), \verb|lsucode5 = 001|\\

\bitfields{ 1/P, 2/00, 2/-- --, 27/offset27, 1/1, 2/01, 3/111, 2/coherency, 6/registerW, 2/11, 3/001, 1/--, 5/00001, 1/0, 6/registerZ }


\paragraph{Syntax} \texttt{alclrh}\emph{coherency} \emph{registerW} = \emph{offset27}[\emph{registerZ}]

\paragraph{Description} The effective address is computed by adding the \emph{offset27} to the \emph{registerZ}.
The \emph{registerW} is loaded from the memory half word starting at the effective address. The value 0 is stored into the memory half word starting at the effective address. This atomic instruction implements Test-and-Clear in memory. A zero half word in memory means that the lock at the effective address is taken. Unlocking is achieved by storing a non-zero half word at the effective address.


\paragraph{Modifier(s)}
\noindent\begin{itemize}
\item coherency (cf. coherency in section \ref{par:modifier-coherency} on page \pageref{par:modifier-coherency})
\end{itemize}

\paragraph{Execution} {Reservation table \textsf{LSU2\_MEMW\_AUXW.X} (Section~\ref{sec:reservation-LSU2-MEMW-AUXW.X})}
~
{\begin{lstlisting}
stage ID:
  new offset = _SX_27(offset27);
  new coherency = _ZX_2(coherency);
stage RR:
  new base = GPR[registerZ];
  new address = base + offset;
  new result1 = _ZX_16(MEM.atomic_swap(address, 0x3, coherency << 32, 0, registerW));
stage CR:
  GPR[registerW] = result1;
\end{lstlisting}}
\newpage

\paragraph{Encoding} Format \textsf{LSU\_ALCLRSB.D} (Section~\ref{sec:format-LSU-ALCLRSB.D}), \verb|lsucode5 = 001|\\

\bitfields{ 1/P, 2/00, 2/-- --, 27/extend27, 1/1, 2/00, 2/-- --, 27/offset27, 1/1, 2/01, 3/111, 2/coherency, 6/registerW, 2/11, 3/001, 1/--, 5/00001, 1/0, 6/registerZ }


\paragraph{Syntax} \texttt{alclrh}\emph{coherency} \emph{registerW} = \emph{extend27\_\discretionary{}{}{}offset27}[\emph{registerZ}]

\paragraph{Description} The effective address is computed by adding the \emph{extend27\_\discretionary{}{}{}offset27} to the \emph{registerZ}.
The \emph{registerW} is loaded from the memory half word starting at the effective address. The value 0 is stored into the memory half word starting at the effective address. This atomic instruction implements Test-and-Clear in memory. A zero half word in memory means that the lock at the effective address is taken. Unlocking is achieved by storing a non-zero half word at the effective address.


\paragraph{Modifier(s)}
\noindent\begin{itemize}
\item coherency (cf. coherency in section \ref{par:modifier-coherency} on page \pageref{par:modifier-coherency})
\end{itemize}

\paragraph{Execution} {Reservation table \textsf{LSU2\_MEMW\_AUXW.Y} (Section~\ref{sec:reservation-LSU2-MEMW-AUXW.Y})}
~
{\begin{lstlisting}
stage ID:
  new offset = _SX_54(extend27_offset27);
  new coherency = _ZX_2(coherency);
stage RR:
  new base = GPR[registerZ];
  new address = base + offset;
  new result1 = _ZX_16(MEM.atomic_swap(address, 0x3, coherency << 32, 0, registerW));
stage CR:
  GPR[registerW] = result1;
\end{lstlisting}}
\newpage
\subsubsection[{\small ALCLRW: Atomic Load and Clear Word}]{ALCLRW\hfill Atomic Load and Clear Word}
\label{instr:ALCLRW}\index{alclrw}
 
\paragraph{Encoding} Format \textsf{LSU\_ALCLRSB} (Section~\ref{sec:format-LSU-ALCLRSB}), \verb|lsucode5 = 010|\\

\bitfields{ 1/P, 2/01, 3/111, 2/coherency, 6/registerW, 2/11, 3/010, 1/--, 5/00001, 1/0, 6/registerZ }


\paragraph{Syntax} \texttt{alclrw}\emph{coherency} \emph{registerW} = [\emph{registerZ}]

\paragraph{Description} The effective address is the value of the \emph{registerZ}.
The \emph{registerW} is loaded from the memory word starting at the effective address. The value 0 is stored into the memory word starting at the effective address. This atomic instruction implements Test-and-Clear in memory. A zero word in memory means that the lock at the effective address is taken. Unlocking is achieved by storing a non-zero word at the effective address.


\paragraph{Modifier(s)}
\noindent\begin{itemize}
\item coherency (cf. coherency in section \ref{par:modifier-coherency} on page \pageref{par:modifier-coherency})
\end{itemize}

\paragraph{Execution} {Reservation table \textsf{LSU2\_MEMW\_AUXW} (Section~\ref{sec:reservation-LSU2-MEMW-AUXW})}
~
{\begin{lstlisting}
stage ID:
  new coherency = _ZX_2(coherency);
stage RR:
  new address = GPR[registerZ];
  new result1 = _ZX_32(MEM.atomic_swap(address, 0xF, coherency << 32, 0, registerW));
stage CR:
  GPR[registerW] = result1;
\end{lstlisting}}
\newpage

\paragraph{Encoding} Format \textsf{LSU\_ALCLRSB.O} (Section~\ref{sec:format-LSU-ALCLRSB.O}), \verb|lsucode5 = 010|\\

\bitfields{ 1/P, 2/00, 2/-- --, 27/offset27, 1/1, 2/01, 3/111, 2/coherency, 6/registerW, 2/11, 3/010, 1/--, 5/00001, 1/0, 6/registerZ }


\paragraph{Syntax} \texttt{alclrw}\emph{coherency} \emph{registerW} = \emph{offset27}[\emph{registerZ}]

\paragraph{Description} The effective address is computed by adding the \emph{offset27} to the \emph{registerZ}.
The \emph{registerW} is loaded from the memory word starting at the effective address. The value 0 is stored into the memory word starting at the effective address. This atomic instruction implements Test-and-Clear in memory. A zero word in memory means that the lock at the effective address is taken. Unlocking is achieved by storing a non-zero word at the effective address.


\paragraph{Modifier(s)}
\noindent\begin{itemize}
\item coherency (cf. coherency in section \ref{par:modifier-coherency} on page \pageref{par:modifier-coherency})
\end{itemize}

\paragraph{Execution} {Reservation table \textsf{LSU2\_MEMW\_AUXW.X} (Section~\ref{sec:reservation-LSU2-MEMW-AUXW.X})}
~
{\begin{lstlisting}
stage ID:
  new offset = _SX_27(offset27);
  new coherency = _ZX_2(coherency);
stage RR:
  new base = GPR[registerZ];
  new address = base + offset;
  new result1 = _ZX_32(MEM.atomic_swap(address, 0xF, coherency << 32, 0, registerW));
stage CR:
  GPR[registerW] = result1;
\end{lstlisting}}
\newpage

\paragraph{Encoding} Format \textsf{LSU\_ALCLRSB.D} (Section~\ref{sec:format-LSU-ALCLRSB.D}), \verb|lsucode5 = 010|\\

\bitfields{ 1/P, 2/00, 2/-- --, 27/extend27, 1/1, 2/00, 2/-- --, 27/offset27, 1/1, 2/01, 3/111, 2/coherency, 6/registerW, 2/11, 3/010, 1/--, 5/00001, 1/0, 6/registerZ }


\paragraph{Syntax} \texttt{alclrw}\emph{coherency} \emph{registerW} = \emph{extend27\_\discretionary{}{}{}offset27}[\emph{registerZ}]

\paragraph{Description} The effective address is computed by adding the \emph{extend27\_\discretionary{}{}{}offset27} to the \emph{registerZ}.
The \emph{registerW} is loaded from the memory word starting at the effective address. The value 0 is stored into the memory word starting at the effective address. This atomic instruction implements Test-and-Clear in memory. A zero word in memory means that the lock at the effective address is taken. Unlocking is achieved by storing a non-zero word at the effective address.


\paragraph{Modifier(s)}
\noindent\begin{itemize}
\item coherency (cf. coherency in section \ref{par:modifier-coherency} on page \pageref{par:modifier-coherency})
\end{itemize}

\paragraph{Execution} {Reservation table \textsf{LSU2\_MEMW\_AUXW.Y} (Section~\ref{sec:reservation-LSU2-MEMW-AUXW.Y})}
~
{\begin{lstlisting}
stage ID:
  new offset = _SX_54(extend27_offset27);
  new coherency = _ZX_2(coherency);
stage RR:
  new base = GPR[registerZ];
  new address = base + offset;
  new result1 = _ZX_32(MEM.atomic_swap(address, 0xF, coherency << 32, 0, registerW));
stage CR:
  GPR[registerW] = result1;
\end{lstlisting}}
\newpage
\subsubsection[{\small ALD: Atomic Load Double Word}]{ALD\hfill Atomic Load Double Word}
\label{instr:ALD}\index{ald}
 
\paragraph{Encoding} Format \textsf{LSU\_ALSB} (Section~\ref{sec:format-LSU-ALSB}), \verb|lsucode5 = 011|\\

\bitfields{ 1/P, 2/01, 3/111, 2/coherency, 6/registerW, 2/11, 3/011, 1/--, 5/00000, 1/0, 6/registerZ }


\paragraph{Syntax} \texttt{ald}\emph{coherency} \emph{registerW} = [\emph{registerZ}]

\paragraph{Description} The effective address is the value of the \emph{registerZ}.
The \emph{registerW} is atomically loaded from the memory double word starting at the effective address. This instruction is provided to directly operate on the last level of the memory hierarchy, which is typically a processor DDR or a cluster TCM, irrespective of the MMU DCP (Data Cache Policy) settings. The effective address must be double word aligned (multiple of 8).


\paragraph{Modifier(s)}
\noindent\begin{itemize}
\item coherency (cf. coherency in section \ref{par:modifier-coherency} on page \pageref{par:modifier-coherency})
\end{itemize}

\paragraph{Execution} {Reservation table \textsf{LSU\_MEMW\_AUXW} (Section~\ref{sec:reservation-LSU-MEMW-AUXW})}
~
{\begin{lstlisting}
stage ID:
  new coherency = _ZX_2(coherency);
stage RR:
  new address = GPR[registerZ];
  new result1 = MEM.atomic_load(address, 0xFF, coherency << 32, registerW);
stage CR:
  GPR[registerW] = result1;
\end{lstlisting}}
\newpage

\paragraph{Encoding} Format \textsf{LSU\_ALSB.O} (Section~\ref{sec:format-LSU-ALSB.O}), \verb|lsucode5 = 011|\\

\bitfields{ 1/P, 2/00, 2/-- --, 27/offset27, 1/1, 2/01, 3/111, 2/coherency, 6/registerW, 2/11, 3/011, 1/--, 5/00000, 1/0, 6/registerZ }


\paragraph{Syntax} \texttt{ald}\emph{coherency} \emph{registerW} = \emph{offset27}[\emph{registerZ}]

\paragraph{Description} The effective address is computed by adding the \emph{offset27} to the \emph{registerZ}.
The \emph{registerW} is atomically loaded from the memory double word starting at the effective address. This instruction is provided to directly operate on the last level of the memory hierarchy, which is typically a processor DDR or a cluster TCM, irrespective of the MMU DCP (Data Cache Policy) settings. The effective address must be double word aligned (multiple of 8).


\paragraph{Modifier(s)}
\noindent\begin{itemize}
\item coherency (cf. coherency in section \ref{par:modifier-coherency} on page \pageref{par:modifier-coherency})
\end{itemize}

\paragraph{Execution} {Reservation table \textsf{LSU\_MEMW\_AUXW.X} (Section~\ref{sec:reservation-LSU-MEMW-AUXW.X})}
~
{\begin{lstlisting}
stage ID:
  new offset = _SX_27(offset27);
  new coherency = _ZX_2(coherency);
stage RR:
  new base = GPR[registerZ];
  new address = base + offset;
  new result1 = MEM.atomic_load(address, 0xFF, coherency << 32, registerW);
stage CR:
  GPR[registerW] = result1;
\end{lstlisting}}
\newpage

\paragraph{Encoding} Format \textsf{LSU\_ALSB.D} (Section~\ref{sec:format-LSU-ALSB.D}), \verb|lsucode5 = 011|\\

\bitfields{ 1/P, 2/00, 2/-- --, 27/extend27, 1/1, 2/00, 2/-- --, 27/offset27, 1/1, 2/01, 3/111, 2/coherency, 6/registerW, 2/11, 3/011, 1/--, 5/00000, 1/0, 6/registerZ }


\paragraph{Syntax} \texttt{ald}\emph{coherency} \emph{registerW} = \emph{extend27\_\discretionary{}{}{}offset27}[\emph{registerZ}]

\paragraph{Description} The effective address is computed by adding the \emph{extend27\_\discretionary{}{}{}offset27} to the \emph{registerZ}.
The \emph{registerW} is atomically loaded from the memory double word starting at the effective address. This instruction is provided to directly operate on the last level of the memory hierarchy, which is typically a processor DDR or a cluster TCM, irrespective of the MMU DCP (Data Cache Policy) settings. The effective address must be double word aligned (multiple of 8).


\paragraph{Modifier(s)}
\noindent\begin{itemize}
\item coherency (cf. coherency in section \ref{par:modifier-coherency} on page \pageref{par:modifier-coherency})
\end{itemize}

\paragraph{Execution} {Reservation table \textsf{LSU\_MEMW\_AUXW.Y} (Section~\ref{sec:reservation-LSU-MEMW-AUXW.Y})}
~
{\begin{lstlisting}
stage ID:
  new offset = _SX_54(extend27_offset27);
  new coherency = _ZX_2(coherency);
stage RR:
  new base = GPR[registerZ];
  new address = base + offset;
  new result1 = MEM.atomic_load(address, 0xFF, coherency << 32, registerW);
stage CR:
  GPR[registerW] = result1;
\end{lstlisting}}
\newpage
\subsubsection[{\small ALDUSB: Atomic Load and Decrement Unsigned Saturating Byte}]{ALDUSB\hfill Atomic Load and Decrement Unsigned Saturating Byte}
\label{instr:ALDUSB}\index{aldusb}
 
\paragraph{Encoding} Format \textsf{LSU\_ALDUSSB} (Section~\ref{sec:format-LSU-ALDUSSB}), \verb|lsucode5 = 000|\\

\bitfields{ 1/P, 2/01, 3/111, 2/coherency, 6/registerT, 2/11, 3/000, 1/--, 5/01110, 1/0, 6/registerZ }


\paragraph{Syntax} \texttt{aldusb}\emph{coherency} [\emph{registerZ}] = \emph{registerT}

\paragraph{Description} The effective address is the value of the \emph{registerZ}.
This instruction implements atomic load-DUS-store on signed byte (Decrement Unsigned Saturating). The byte at effective address is DUSed with the \emph{registerT}, and its previous value is returned into the \emph{registerT}.


\paragraph{Modifier(s)}
\noindent\begin{itemize}
\item coherency (cf. coherency in section \ref{par:modifier-coherency} on page \pageref{par:modifier-coherency})
\end{itemize}

\paragraph{Execution} {Reservation table \textsf{LSU2\_MEMW\_AUXR\_AUXW} (Section~\ref{sec:reservation-LSU2-MEMW-AUXR-AUXW})}
~
{\begin{lstlisting}
stage ID:
  new coherency = _ZX_2(coherency);
stage RR:
  new address = GPR[registerZ];
stage E1:
  new argument2 = GPR[registerT];
  new result2 = _ZX_8(MEM.atomic_dus(address, 0x1, coherency << 32, argument2.8[0], registerT));
stage SR:
  GPR[registerT] = result2;
\end{lstlisting}}
\newpage

\paragraph{Encoding} Format \textsf{LSU\_ALDUSSB.O} (Section~\ref{sec:format-LSU-ALDUSSB.O}), \verb|lsucode5 = 000|\\

\bitfields{ 1/P, 2/00, 2/-- --, 27/offset27, 1/1, 2/01, 3/111, 2/coherency, 6/registerT, 2/11, 3/000, 1/--, 5/01110, 1/0, 6/registerZ }


\paragraph{Syntax} \texttt{aldusb}\emph{coherency} \emph{offset27}[\emph{registerZ}] = \emph{registerT}

\paragraph{Description} The effective address is computed by adding the \emph{offset27} to the \emph{registerZ}.
This instruction implements atomic load-DUS-store on signed byte (Decrement Unsigned Saturating). The byte at effective address is DUSed with the \emph{registerT}, and its previous value is returned into the \emph{registerT}.


\paragraph{Modifier(s)}
\noindent\begin{itemize}
\item coherency (cf. coherency in section \ref{par:modifier-coherency} on page \pageref{par:modifier-coherency})
\end{itemize}

\paragraph{Execution} {Reservation table \textsf{LSU2\_MEMW\_AUXR\_AUXW.X} (Section~\ref{sec:reservation-LSU2-MEMW-AUXR-AUXW.X})}
~
{\begin{lstlisting}
stage ID:
  new offset = _SX_27(offset27);
  new coherency = _ZX_2(coherency);
stage RR:
  new base = GPR[registerZ];
  new address = base + offset;
stage E1:
  new argument2 = GPR[registerT];
  new result2 = _ZX_8(MEM.atomic_dus(address, 0x1, coherency << 32, argument2.8[0], registerT));
stage SR:
  GPR[registerT] = result2;
\end{lstlisting}}
\newpage

\paragraph{Encoding} Format \textsf{LSU\_ALDUSSB.D} (Section~\ref{sec:format-LSU-ALDUSSB.D}), \verb|lsucode5 = 000|\\

\bitfields{ 1/P, 2/00, 2/-- --, 27/extend27, 1/1, 2/00, 2/-- --, 27/offset27, 1/1, 2/01, 3/111, 2/coherency, 6/registerT, 2/11, 3/000, 1/--, 5/01110, 1/0, 6/registerZ }


\paragraph{Syntax} \texttt{aldusb}\emph{coherency} \emph{extend27\_\discretionary{}{}{}offset27}[\emph{registerZ}] = \emph{registerT}

\paragraph{Description} The effective address is computed by adding the \emph{extend27\_\discretionary{}{}{}offset27} to the \emph{registerZ}.
This instruction implements atomic load-DUS-store on signed byte (Decrement Unsigned Saturating). The byte at effective address is DUSed with the \emph{registerT}, and its previous value is returned into the \emph{registerT}.


\paragraph{Modifier(s)}
\noindent\begin{itemize}
\item coherency (cf. coherency in section \ref{par:modifier-coherency} on page \pageref{par:modifier-coherency})
\end{itemize}

\paragraph{Execution} {Reservation table \textsf{LSU2\_MEMW\_AUXR\_AUXW.Y} (Section~\ref{sec:reservation-LSU2-MEMW-AUXR-AUXW.Y})}
~
{\begin{lstlisting}
stage ID:
  new offset = _SX_54(extend27_offset27);
  new coherency = _ZX_2(coherency);
stage RR:
  new base = GPR[registerZ];
  new address = base + offset;
stage E1:
  new argument2 = GPR[registerT];
  new result2 = _ZX_8(MEM.atomic_dus(address, 0x1, coherency << 32, argument2.8[0], registerT));
stage SR:
  GPR[registerT] = result2;
\end{lstlisting}}
\newpage
\subsubsection[{\small ALDUSD: Atomic Load and DUS Double Word}]{ALDUSD\hfill Atomic Load and DUS Double Word}
\label{instr:ALDUSD}\index{aldusd}
 
\paragraph{Encoding} Format \textsf{LSU\_ALDUSSB} (Section~\ref{sec:format-LSU-ALDUSSB}), \verb|lsucode5 = 011|\\

\bitfields{ 1/P, 2/01, 3/111, 2/coherency, 6/registerT, 2/11, 3/011, 1/--, 5/01110, 1/0, 6/registerZ }


\paragraph{Syntax} \texttt{aldusd}\emph{coherency} [\emph{registerZ}] = \emph{registerT}

\paragraph{Description} The effective address is the value of the \emph{registerZ}.
This instruction implements atomic load-DUS-store on signed double word (Decrement Unsigned Saturating). The double word at effective address is DUSed with the \emph{registerT}, and its previous value is returned into the \emph{registerT}. The effective address must be double word aligned (multiple of 8).


\paragraph{Modifier(s)}
\noindent\begin{itemize}
\item coherency (cf. coherency in section \ref{par:modifier-coherency} on page \pageref{par:modifier-coherency})
\end{itemize}

\paragraph{Execution} {Reservation table \textsf{LSU2\_MEMW\_AUXR\_AUXW} (Section~\ref{sec:reservation-LSU2-MEMW-AUXR-AUXW})}
~
{\begin{lstlisting}
stage ID:
  new coherency = _ZX_2(coherency);
stage RR:
  new address = GPR[registerZ];
stage E1:
  new argument2 = GPR[registerT];
  new result2 = MEM.atomic_dus(address, 0xFF, coherency << 32, argument2.64[0], registerT);
stage SR:
  GPR[registerT] = result2;
\end{lstlisting}}
\newpage

\paragraph{Encoding} Format \textsf{LSU\_ALDUSSB.O} (Section~\ref{sec:format-LSU-ALDUSSB.O}), \verb|lsucode5 = 011|\\

\bitfields{ 1/P, 2/00, 2/-- --, 27/offset27, 1/1, 2/01, 3/111, 2/coherency, 6/registerT, 2/11, 3/011, 1/--, 5/01110, 1/0, 6/registerZ }


\paragraph{Syntax} \texttt{aldusd}\emph{coherency} \emph{offset27}[\emph{registerZ}] = \emph{registerT}

\paragraph{Description} The effective address is computed by adding the \emph{offset27} to the \emph{registerZ}.
This instruction implements atomic load-DUS-store on signed double word (Decrement Unsigned Saturating). The double word at effective address is DUSed with the \emph{registerT}, and its previous value is returned into the \emph{registerT}. The effective address must be double word aligned (multiple of 8).


\paragraph{Modifier(s)}
\noindent\begin{itemize}
\item coherency (cf. coherency in section \ref{par:modifier-coherency} on page \pageref{par:modifier-coherency})
\end{itemize}

\paragraph{Execution} {Reservation table \textsf{LSU2\_MEMW\_AUXR\_AUXW.X} (Section~\ref{sec:reservation-LSU2-MEMW-AUXR-AUXW.X})}
~
{\begin{lstlisting}
stage ID:
  new offset = _SX_27(offset27);
  new coherency = _ZX_2(coherency);
stage RR:
  new base = GPR[registerZ];
  new address = base + offset;
stage E1:
  new argument2 = GPR[registerT];
  new result2 = MEM.atomic_dus(address, 0xFF, coherency << 32, argument2.64[0], registerT);
stage SR:
  GPR[registerT] = result2;
\end{lstlisting}}
\newpage

\paragraph{Encoding} Format \textsf{LSU\_ALDUSSB.D} (Section~\ref{sec:format-LSU-ALDUSSB.D}), \verb|lsucode5 = 011|\\

\bitfields{ 1/P, 2/00, 2/-- --, 27/extend27, 1/1, 2/00, 2/-- --, 27/offset27, 1/1, 2/01, 3/111, 2/coherency, 6/registerT, 2/11, 3/011, 1/--, 5/01110, 1/0, 6/registerZ }


\paragraph{Syntax} \texttt{aldusd}\emph{coherency} \emph{extend27\_\discretionary{}{}{}offset27}[\emph{registerZ}] = \emph{registerT}

\paragraph{Description} The effective address is computed by adding the \emph{extend27\_\discretionary{}{}{}offset27} to the \emph{registerZ}.
This instruction implements atomic load-DUS-store on signed double word (Decrement Unsigned Saturating). The double word at effective address is DUSed with the \emph{registerT}, and its previous value is returned into the \emph{registerT}. The effective address must be double word aligned (multiple of 8).


\paragraph{Modifier(s)}
\noindent\begin{itemize}
\item coherency (cf. coherency in section \ref{par:modifier-coherency} on page \pageref{par:modifier-coherency})
\end{itemize}

\paragraph{Execution} {Reservation table \textsf{LSU2\_MEMW\_AUXR\_AUXW.Y} (Section~\ref{sec:reservation-LSU2-MEMW-AUXR-AUXW.Y})}
~
{\begin{lstlisting}
stage ID:
  new offset = _SX_54(extend27_offset27);
  new coherency = _ZX_2(coherency);
stage RR:
  new base = GPR[registerZ];
  new address = base + offset;
stage E1:
  new argument2 = GPR[registerT];
  new result2 = MEM.atomic_dus(address, 0xFF, coherency << 32, argument2.64[0], registerT);
stage SR:
  GPR[registerT] = result2;
\end{lstlisting}}
\newpage
\subsubsection[{\small ALDUSH: Atomic Load and DUS Half Word}]{ALDUSH\hfill Atomic Load and DUS Half Word}
\label{instr:ALDUSH}\index{aldush}
 
\paragraph{Encoding} Format \textsf{LSU\_ALDUSSB} (Section~\ref{sec:format-LSU-ALDUSSB}), \verb|lsucode5 = 001|\\

\bitfields{ 1/P, 2/01, 3/111, 2/coherency, 6/registerT, 2/11, 3/001, 1/--, 5/01110, 1/0, 6/registerZ }


\paragraph{Syntax} \texttt{aldush}\emph{coherency} [\emph{registerZ}] = \emph{registerT}

\paragraph{Description} The effective address is the value of the \emph{registerZ}.
This instruction implements atomic load-DUS-store on signed half word (Decrement Unsigned Saturating). The half word at effective address is DUSed with the \emph{registerT}, and its previous value is returned into the \emph{registerT}. The effective address must be half word aligned (multiple of 2).


\paragraph{Modifier(s)}
\noindent\begin{itemize}
\item coherency (cf. coherency in section \ref{par:modifier-coherency} on page \pageref{par:modifier-coherency})
\end{itemize}

\paragraph{Execution} {Reservation table \textsf{LSU2\_MEMW\_AUXR\_AUXW} (Section~\ref{sec:reservation-LSU2-MEMW-AUXR-AUXW})}
~
{\begin{lstlisting}
stage ID:
  new coherency = _ZX_2(coherency);
stage RR:
  new address = GPR[registerZ];
stage E1:
  new argument2 = GPR[registerT];
  new result2 = _ZX_16(MEM.atomic_dus(address, 0x3, coherency << 32, argument2.16[0], registerT));
stage SR:
  GPR[registerT] = result2;
\end{lstlisting}}
\newpage

\paragraph{Encoding} Format \textsf{LSU\_ALDUSSB.O} (Section~\ref{sec:format-LSU-ALDUSSB.O}), \verb|lsucode5 = 001|\\

\bitfields{ 1/P, 2/00, 2/-- --, 27/offset27, 1/1, 2/01, 3/111, 2/coherency, 6/registerT, 2/11, 3/001, 1/--, 5/01110, 1/0, 6/registerZ }


\paragraph{Syntax} \texttt{aldush}\emph{coherency} \emph{offset27}[\emph{registerZ}] = \emph{registerT}

\paragraph{Description} The effective address is computed by adding the \emph{offset27} to the \emph{registerZ}.
This instruction implements atomic load-DUS-store on signed half word (Decrement Unsigned Saturating). The half word at effective address is DUSed with the \emph{registerT}, and its previous value is returned into the \emph{registerT}. The effective address must be half word aligned (multiple of 2).


\paragraph{Modifier(s)}
\noindent\begin{itemize}
\item coherency (cf. coherency in section \ref{par:modifier-coherency} on page \pageref{par:modifier-coherency})
\end{itemize}

\paragraph{Execution} {Reservation table \textsf{LSU2\_MEMW\_AUXR\_AUXW.X} (Section~\ref{sec:reservation-LSU2-MEMW-AUXR-AUXW.X})}
~
{\begin{lstlisting}
stage ID:
  new offset = _SX_27(offset27);
  new coherency = _ZX_2(coherency);
stage RR:
  new base = GPR[registerZ];
  new address = base + offset;
stage E1:
  new argument2 = GPR[registerT];
  new result2 = _ZX_16(MEM.atomic_dus(address, 0x3, coherency << 32, argument2.16[0], registerT));
stage SR:
  GPR[registerT] = result2;
\end{lstlisting}}
\newpage

\paragraph{Encoding} Format \textsf{LSU\_ALDUSSB.D} (Section~\ref{sec:format-LSU-ALDUSSB.D}), \verb|lsucode5 = 001|\\

\bitfields{ 1/P, 2/00, 2/-- --, 27/extend27, 1/1, 2/00, 2/-- --, 27/offset27, 1/1, 2/01, 3/111, 2/coherency, 6/registerT, 2/11, 3/001, 1/--, 5/01110, 1/0, 6/registerZ }


\paragraph{Syntax} \texttt{aldush}\emph{coherency} \emph{extend27\_\discretionary{}{}{}offset27}[\emph{registerZ}] = \emph{registerT}

\paragraph{Description} The effective address is computed by adding the \emph{extend27\_\discretionary{}{}{}offset27} to the \emph{registerZ}.
This instruction implements atomic load-DUS-store on signed half word (Decrement Unsigned Saturating). The half word at effective address is DUSed with the \emph{registerT}, and its previous value is returned into the \emph{registerT}. The effective address must be half word aligned (multiple of 2).


\paragraph{Modifier(s)}
\noindent\begin{itemize}
\item coherency (cf. coherency in section \ref{par:modifier-coherency} on page \pageref{par:modifier-coherency})
\end{itemize}

\paragraph{Execution} {Reservation table \textsf{LSU2\_MEMW\_AUXR\_AUXW.Y} (Section~\ref{sec:reservation-LSU2-MEMW-AUXR-AUXW.Y})}
~
{\begin{lstlisting}
stage ID:
  new offset = _SX_54(extend27_offset27);
  new coherency = _ZX_2(coherency);
stage RR:
  new base = GPR[registerZ];
  new address = base + offset;
stage E1:
  new argument2 = GPR[registerT];
  new result2 = _ZX_16(MEM.atomic_dus(address, 0x3, coherency << 32, argument2.16[0], registerT));
stage SR:
  GPR[registerT] = result2;
\end{lstlisting}}
\newpage
\subsubsection[{\small ALDUSW: Atomic Load and DUS Word}]{ALDUSW\hfill Atomic Load and DUS Word}
\label{instr:ALDUSW}\index{aldusw}
 
\paragraph{Encoding} Format \textsf{LSU\_ALDUSSB} (Section~\ref{sec:format-LSU-ALDUSSB}), \verb|lsucode5 = 010|\\

\bitfields{ 1/P, 2/01, 3/111, 2/coherency, 6/registerT, 2/11, 3/010, 1/--, 5/01110, 1/0, 6/registerZ }


\paragraph{Syntax} \texttt{aldusw}\emph{coherency} [\emph{registerZ}] = \emph{registerT}

\paragraph{Description} The effective address is the value of the \emph{registerZ}.
This instruction implements atomic load-DUS-store on signed word (Decrement Unsigned Saturating). The word at effective address is DUSed with the \emph{registerT}, and its previous value is returned into the \emph{registerT}. The effective address must be word aligned (multiple of 4).


\paragraph{Modifier(s)}
\noindent\begin{itemize}
\item coherency (cf. coherency in section \ref{par:modifier-coherency} on page \pageref{par:modifier-coherency})
\end{itemize}

\paragraph{Execution} {Reservation table \textsf{LSU2\_MEMW\_AUXR\_AUXW} (Section~\ref{sec:reservation-LSU2-MEMW-AUXR-AUXW})}
~
{\begin{lstlisting}
stage ID:
  new coherency = _ZX_2(coherency);
stage RR:
  new address = GPR[registerZ];
stage E1:
  new argument2 = GPR[registerT];
  new result2 = _ZX_32(MEM.atomic_dus(address, 0xF, coherency << 32, argument2.32[0], registerT));
stage SR:
  GPR[registerT] = result2;
\end{lstlisting}}
\newpage

\paragraph{Encoding} Format \textsf{LSU\_ALDUSSB.O} (Section~\ref{sec:format-LSU-ALDUSSB.O}), \verb|lsucode5 = 010|\\

\bitfields{ 1/P, 2/00, 2/-- --, 27/offset27, 1/1, 2/01, 3/111, 2/coherency, 6/registerT, 2/11, 3/010, 1/--, 5/01110, 1/0, 6/registerZ }


\paragraph{Syntax} \texttt{aldusw}\emph{coherency} \emph{offset27}[\emph{registerZ}] = \emph{registerT}

\paragraph{Description} The effective address is computed by adding the \emph{offset27} to the \emph{registerZ}.
This instruction implements atomic load-DUS-store on signed word (Decrement Unsigned Saturating). The word at effective address is DUSed with the \emph{registerT}, and its previous value is returned into the \emph{registerT}. The effective address must be word aligned (multiple of 4).


\paragraph{Modifier(s)}
\noindent\begin{itemize}
\item coherency (cf. coherency in section \ref{par:modifier-coherency} on page \pageref{par:modifier-coherency})
\end{itemize}

\paragraph{Execution} {Reservation table \textsf{LSU2\_MEMW\_AUXR\_AUXW.X} (Section~\ref{sec:reservation-LSU2-MEMW-AUXR-AUXW.X})}
~
{\begin{lstlisting}
stage ID:
  new offset = _SX_27(offset27);
  new coherency = _ZX_2(coherency);
stage RR:
  new base = GPR[registerZ];
  new address = base + offset;
stage E1:
  new argument2 = GPR[registerT];
  new result2 = _ZX_32(MEM.atomic_dus(address, 0xF, coherency << 32, argument2.32[0], registerT));
stage SR:
  GPR[registerT] = result2;
\end{lstlisting}}
\newpage

\paragraph{Encoding} Format \textsf{LSU\_ALDUSSB.D} (Section~\ref{sec:format-LSU-ALDUSSB.D}), \verb|lsucode5 = 010|\\

\bitfields{ 1/P, 2/00, 2/-- --, 27/extend27, 1/1, 2/00, 2/-- --, 27/offset27, 1/1, 2/01, 3/111, 2/coherency, 6/registerT, 2/11, 3/010, 1/--, 5/01110, 1/0, 6/registerZ }


\paragraph{Syntax} \texttt{aldusw}\emph{coherency} \emph{extend27\_\discretionary{}{}{}offset27}[\emph{registerZ}] = \emph{registerT}

\paragraph{Description} The effective address is computed by adding the \emph{extend27\_\discretionary{}{}{}offset27} to the \emph{registerZ}.
This instruction implements atomic load-DUS-store on signed word (Decrement Unsigned Saturating). The word at effective address is DUSed with the \emph{registerT}, and its previous value is returned into the \emph{registerT}. The effective address must be word aligned (multiple of 4).


\paragraph{Modifier(s)}
\noindent\begin{itemize}
\item coherency (cf. coherency in section \ref{par:modifier-coherency} on page \pageref{par:modifier-coherency})
\end{itemize}

\paragraph{Execution} {Reservation table \textsf{LSU2\_MEMW\_AUXR\_AUXW.Y} (Section~\ref{sec:reservation-LSU2-MEMW-AUXR-AUXW.Y})}
~
{\begin{lstlisting}
stage ID:
  new offset = _SX_54(extend27_offset27);
  new coherency = _ZX_2(coherency);
stage RR:
  new base = GPR[registerZ];
  new address = base + offset;
stage E1:
  new argument2 = GPR[registerT];
  new result2 = _ZX_32(MEM.atomic_dus(address, 0xF, coherency << 32, argument2.32[0], registerT));
stage SR:
  GPR[registerT] = result2;
\end{lstlisting}}
\newpage
\subsubsection[{\small ALEORB: Atomic Load and EOR Byte}]{ALEORB\hfill Atomic Load and EOR Byte}
\label{instr:ALEORB}\index{aleorb}
 
\paragraph{Encoding} Format \textsf{LSU\_ALEORSB} (Section~\ref{sec:format-LSU-ALEORSB}), \verb|lsucode5 = 000|\\

\bitfields{ 1/P, 2/01, 3/111, 2/coherency, 6/registerT, 2/11, 3/000, 1/--, 5/00111, 1/0, 6/registerZ }


\paragraph{Syntax} \texttt{aleorb}\emph{coherency} [\emph{registerZ}] = \emph{registerT}

\paragraph{Description} The effective address is the value of the \emph{registerZ}.
This instruction implements atomic load-EOR-store on byte. The byte at effective address is EORed with the \emph{registerT}, and its previous value is returned into the \emph{registerT}.


\paragraph{Modifier(s)}
\noindent\begin{itemize}
\item coherency (cf. coherency in section \ref{par:modifier-coherency} on page \pageref{par:modifier-coherency})
\end{itemize}

\paragraph{Execution} {Reservation table \textsf{LSU2\_MEMW\_AUXR\_AUXW} (Section~\ref{sec:reservation-LSU2-MEMW-AUXR-AUXW})}
~
{\begin{lstlisting}
stage ID:
  new coherency = _ZX_2(coherency);
stage RR:
  new address = GPR[registerZ];
stage E1:
  new argument2 = GPR[registerT];
  new result2 = _ZX_8(MEM.atomic_eor(address, 0x1, coherency << 32, argument2.8[0], registerT));
stage SR:
  GPR[registerT] = result2;
\end{lstlisting}}
\newpage

\paragraph{Encoding} Format \textsf{LSU\_ALEORSB.O} (Section~\ref{sec:format-LSU-ALEORSB.O}), \verb|lsucode5 = 000|\\

\bitfields{ 1/P, 2/00, 2/-- --, 27/offset27, 1/1, 2/01, 3/111, 2/coherency, 6/registerT, 2/11, 3/000, 1/--, 5/00111, 1/0, 6/registerZ }


\paragraph{Syntax} \texttt{aleorb}\emph{coherency} \emph{offset27}[\emph{registerZ}] = \emph{registerT}

\paragraph{Description} The effective address is computed by adding the \emph{offset27} to the \emph{registerZ}.
This instruction implements atomic load-EOR-store on byte. The byte at effective address is EORed with the \emph{registerT}, and its previous value is returned into the \emph{registerT}.


\paragraph{Modifier(s)}
\noindent\begin{itemize}
\item coherency (cf. coherency in section \ref{par:modifier-coherency} on page \pageref{par:modifier-coherency})
\end{itemize}

\paragraph{Execution} {Reservation table \textsf{LSU2\_MEMW\_AUXR\_AUXW.X} (Section~\ref{sec:reservation-LSU2-MEMW-AUXR-AUXW.X})}
~
{\begin{lstlisting}
stage ID:
  new offset = _SX_27(offset27);
  new coherency = _ZX_2(coherency);
stage RR:
  new base = GPR[registerZ];
  new address = base + offset;
stage E1:
  new argument2 = GPR[registerT];
  new result2 = _ZX_8(MEM.atomic_eor(address, 0x1, coherency << 32, argument2.8[0], registerT));
stage SR:
  GPR[registerT] = result2;
\end{lstlisting}}
\newpage

\paragraph{Encoding} Format \textsf{LSU\_ALEORSB.D} (Section~\ref{sec:format-LSU-ALEORSB.D}), \verb|lsucode5 = 000|\\

\bitfields{ 1/P, 2/00, 2/-- --, 27/extend27, 1/1, 2/00, 2/-- --, 27/offset27, 1/1, 2/01, 3/111, 2/coherency, 6/registerT, 2/11, 3/000, 1/--, 5/00111, 1/0, 6/registerZ }


\paragraph{Syntax} \texttt{aleorb}\emph{coherency} \emph{extend27\_\discretionary{}{}{}offset27}[\emph{registerZ}] = \emph{registerT}

\paragraph{Description} The effective address is computed by adding the \emph{extend27\_\discretionary{}{}{}offset27} to the \emph{registerZ}.
This instruction implements atomic load-EOR-store on byte. The byte at effective address is EORed with the \emph{registerT}, and its previous value is returned into the \emph{registerT}.


\paragraph{Modifier(s)}
\noindent\begin{itemize}
\item coherency (cf. coherency in section \ref{par:modifier-coherency} on page \pageref{par:modifier-coherency})
\end{itemize}

\paragraph{Execution} {Reservation table \textsf{LSU2\_MEMW\_AUXR\_AUXW.Y} (Section~\ref{sec:reservation-LSU2-MEMW-AUXR-AUXW.Y})}
~
{\begin{lstlisting}
stage ID:
  new offset = _SX_54(extend27_offset27);
  new coherency = _ZX_2(coherency);
stage RR:
  new base = GPR[registerZ];
  new address = base + offset;
stage E1:
  new argument2 = GPR[registerT];
  new result2 = _ZX_8(MEM.atomic_eor(address, 0x1, coherency << 32, argument2.8[0], registerT));
stage SR:
  GPR[registerT] = result2;
\end{lstlisting}}
\newpage
\subsubsection[{\small ALEORD: Atomic Load and EOR Double Word}]{ALEORD\hfill Atomic Load and EOR Double Word}
\label{instr:ALEORD}\index{aleord}
 
\paragraph{Encoding} Format \textsf{LSU\_ALEORSB} (Section~\ref{sec:format-LSU-ALEORSB}), \verb|lsucode5 = 011|\\

\bitfields{ 1/P, 2/01, 3/111, 2/coherency, 6/registerT, 2/11, 3/011, 1/--, 5/00111, 1/0, 6/registerZ }


\paragraph{Syntax} \texttt{aleord}\emph{coherency} [\emph{registerZ}] = \emph{registerT}

\paragraph{Description} The effective address is the value of the \emph{registerZ}.
This instruction implements atomic load-EOR-store on double word. The double word at effective address is EORed with the \emph{registerT}, and its previous value is returned into the \emph{registerT}. The effective address must be double word aligned (multiple of 8).


\paragraph{Modifier(s)}
\noindent\begin{itemize}
\item coherency (cf. coherency in section \ref{par:modifier-coherency} on page \pageref{par:modifier-coherency})
\end{itemize}

\paragraph{Execution} {Reservation table \textsf{LSU2\_MEMW\_AUXR\_AUXW} (Section~\ref{sec:reservation-LSU2-MEMW-AUXR-AUXW})}
~
{\begin{lstlisting}
stage ID:
  new coherency = _ZX_2(coherency);
stage RR:
  new address = GPR[registerZ];
stage E1:
  new argument2 = GPR[registerT];
  new result2 = MEM.atomic_eor(address, 0xFF, coherency << 32, argument2.64[0], registerT);
stage SR:
  GPR[registerT] = result2;
\end{lstlisting}}
\newpage

\paragraph{Encoding} Format \textsf{LSU\_ALEORSB.O} (Section~\ref{sec:format-LSU-ALEORSB.O}), \verb|lsucode5 = 011|\\

\bitfields{ 1/P, 2/00, 2/-- --, 27/offset27, 1/1, 2/01, 3/111, 2/coherency, 6/registerT, 2/11, 3/011, 1/--, 5/00111, 1/0, 6/registerZ }


\paragraph{Syntax} \texttt{aleord}\emph{coherency} \emph{offset27}[\emph{registerZ}] = \emph{registerT}

\paragraph{Description} The effective address is computed by adding the \emph{offset27} to the \emph{registerZ}.
This instruction implements atomic load-EOR-store on double word. The double word at effective address is EORed with the \emph{registerT}, and its previous value is returned into the \emph{registerT}. The effective address must be double word aligned (multiple of 8).


\paragraph{Modifier(s)}
\noindent\begin{itemize}
\item coherency (cf. coherency in section \ref{par:modifier-coherency} on page \pageref{par:modifier-coherency})
\end{itemize}

\paragraph{Execution} {Reservation table \textsf{LSU2\_MEMW\_AUXR\_AUXW.X} (Section~\ref{sec:reservation-LSU2-MEMW-AUXR-AUXW.X})}
~
{\begin{lstlisting}
stage ID:
  new offset = _SX_27(offset27);
  new coherency = _ZX_2(coherency);
stage RR:
  new base = GPR[registerZ];
  new address = base + offset;
stage E1:
  new argument2 = GPR[registerT];
  new result2 = MEM.atomic_eor(address, 0xFF, coherency << 32, argument2.64[0], registerT);
stage SR:
  GPR[registerT] = result2;
\end{lstlisting}}
\newpage

\paragraph{Encoding} Format \textsf{LSU\_ALEORSB.D} (Section~\ref{sec:format-LSU-ALEORSB.D}), \verb|lsucode5 = 011|\\

\bitfields{ 1/P, 2/00, 2/-- --, 27/extend27, 1/1, 2/00, 2/-- --, 27/offset27, 1/1, 2/01, 3/111, 2/coherency, 6/registerT, 2/11, 3/011, 1/--, 5/00111, 1/0, 6/registerZ }


\paragraph{Syntax} \texttt{aleord}\emph{coherency} \emph{extend27\_\discretionary{}{}{}offset27}[\emph{registerZ}] = \emph{registerT}

\paragraph{Description} The effective address is computed by adding the \emph{extend27\_\discretionary{}{}{}offset27} to the \emph{registerZ}.
This instruction implements atomic load-EOR-store on double word. The double word at effective address is EORed with the \emph{registerT}, and its previous value is returned into the \emph{registerT}. The effective address must be double word aligned (multiple of 8).


\paragraph{Modifier(s)}
\noindent\begin{itemize}
\item coherency (cf. coherency in section \ref{par:modifier-coherency} on page \pageref{par:modifier-coherency})
\end{itemize}

\paragraph{Execution} {Reservation table \textsf{LSU2\_MEMW\_AUXR\_AUXW.Y} (Section~\ref{sec:reservation-LSU2-MEMW-AUXR-AUXW.Y})}
~
{\begin{lstlisting}
stage ID:
  new offset = _SX_54(extend27_offset27);
  new coherency = _ZX_2(coherency);
stage RR:
  new base = GPR[registerZ];
  new address = base + offset;
stage E1:
  new argument2 = GPR[registerT];
  new result2 = MEM.atomic_eor(address, 0xFF, coherency << 32, argument2.64[0], registerT);
stage SR:
  GPR[registerT] = result2;
\end{lstlisting}}
\newpage
\subsubsection[{\small ALEORH: Atomic Load and EOR Half Word}]{ALEORH\hfill Atomic Load and EOR Half Word}
\label{instr:ALEORH}\index{aleorh}
 
\paragraph{Encoding} Format \textsf{LSU\_ALEORSB} (Section~\ref{sec:format-LSU-ALEORSB}), \verb|lsucode5 = 001|\\

\bitfields{ 1/P, 2/01, 3/111, 2/coherency, 6/registerT, 2/11, 3/001, 1/--, 5/00111, 1/0, 6/registerZ }


\paragraph{Syntax} \texttt{aleorh}\emph{coherency} [\emph{registerZ}] = \emph{registerT}

\paragraph{Description} The effective address is the value of the \emph{registerZ}.
This instruction implements atomic load-EOR-store on half word. The half word at effective address is EORed with the \emph{registerT}, and its previous value is returned into the \emph{registerT}. The effective address must be half word aligned (multiple of 2).


\paragraph{Modifier(s)}
\noindent\begin{itemize}
\item coherency (cf. coherency in section \ref{par:modifier-coherency} on page \pageref{par:modifier-coherency})
\end{itemize}

\paragraph{Execution} {Reservation table \textsf{LSU2\_MEMW\_AUXR\_AUXW} (Section~\ref{sec:reservation-LSU2-MEMW-AUXR-AUXW})}
~
{\begin{lstlisting}
stage ID:
  new coherency = _ZX_2(coherency);
stage RR:
  new address = GPR[registerZ];
stage E1:
  new argument2 = GPR[registerT];
  new result2 = _ZX_16(MEM.atomic_eor(address, 0x3, coherency << 32, argument2.16[0], registerT));
stage SR:
  GPR[registerT] = result2;
\end{lstlisting}}
\newpage

\paragraph{Encoding} Format \textsf{LSU\_ALEORSB.O} (Section~\ref{sec:format-LSU-ALEORSB.O}), \verb|lsucode5 = 001|\\

\bitfields{ 1/P, 2/00, 2/-- --, 27/offset27, 1/1, 2/01, 3/111, 2/coherency, 6/registerT, 2/11, 3/001, 1/--, 5/00111, 1/0, 6/registerZ }


\paragraph{Syntax} \texttt{aleorh}\emph{coherency} \emph{offset27}[\emph{registerZ}] = \emph{registerT}

\paragraph{Description} The effective address is computed by adding the \emph{offset27} to the \emph{registerZ}.
This instruction implements atomic load-EOR-store on half word. The half word at effective address is EORed with the \emph{registerT}, and its previous value is returned into the \emph{registerT}. The effective address must be half word aligned (multiple of 2).


\paragraph{Modifier(s)}
\noindent\begin{itemize}
\item coherency (cf. coherency in section \ref{par:modifier-coherency} on page \pageref{par:modifier-coherency})
\end{itemize}

\paragraph{Execution} {Reservation table \textsf{LSU2\_MEMW\_AUXR\_AUXW.X} (Section~\ref{sec:reservation-LSU2-MEMW-AUXR-AUXW.X})}
~
{\begin{lstlisting}
stage ID:
  new offset = _SX_27(offset27);
  new coherency = _ZX_2(coherency);
stage RR:
  new base = GPR[registerZ];
  new address = base + offset;
stage E1:
  new argument2 = GPR[registerT];
  new result2 = _ZX_16(MEM.atomic_eor(address, 0x3, coherency << 32, argument2.16[0], registerT));
stage SR:
  GPR[registerT] = result2;
\end{lstlisting}}
\newpage

\paragraph{Encoding} Format \textsf{LSU\_ALEORSB.D} (Section~\ref{sec:format-LSU-ALEORSB.D}), \verb|lsucode5 = 001|\\

\bitfields{ 1/P, 2/00, 2/-- --, 27/extend27, 1/1, 2/00, 2/-- --, 27/offset27, 1/1, 2/01, 3/111, 2/coherency, 6/registerT, 2/11, 3/001, 1/--, 5/00111, 1/0, 6/registerZ }


\paragraph{Syntax} \texttt{aleorh}\emph{coherency} \emph{extend27\_\discretionary{}{}{}offset27}[\emph{registerZ}] = \emph{registerT}

\paragraph{Description} The effective address is computed by adding the \emph{extend27\_\discretionary{}{}{}offset27} to the \emph{registerZ}.
This instruction implements atomic load-EOR-store on half word. The half word at effective address is EORed with the \emph{registerT}, and its previous value is returned into the \emph{registerT}. The effective address must be half word aligned (multiple of 2).


\paragraph{Modifier(s)}
\noindent\begin{itemize}
\item coherency (cf. coherency in section \ref{par:modifier-coherency} on page \pageref{par:modifier-coherency})
\end{itemize}

\paragraph{Execution} {Reservation table \textsf{LSU2\_MEMW\_AUXR\_AUXW.Y} (Section~\ref{sec:reservation-LSU2-MEMW-AUXR-AUXW.Y})}
~
{\begin{lstlisting}
stage ID:
  new offset = _SX_54(extend27_offset27);
  new coherency = _ZX_2(coherency);
stage RR:
  new base = GPR[registerZ];
  new address = base + offset;
stage E1:
  new argument2 = GPR[registerT];
  new result2 = _ZX_16(MEM.atomic_eor(address, 0x3, coherency << 32, argument2.16[0], registerT));
stage SR:
  GPR[registerT] = result2;
\end{lstlisting}}
\newpage
\subsubsection[{\small ALEORW: Atomic Load and EOR Word}]{ALEORW\hfill Atomic Load and EOR Word}
\label{instr:ALEORW}\index{aleorw}
 
\paragraph{Encoding} Format \textsf{LSU\_ALEORSB} (Section~\ref{sec:format-LSU-ALEORSB}), \verb|lsucode5 = 010|\\

\bitfields{ 1/P, 2/01, 3/111, 2/coherency, 6/registerT, 2/11, 3/010, 1/--, 5/00111, 1/0, 6/registerZ }


\paragraph{Syntax} \texttt{aleorw}\emph{coherency} [\emph{registerZ}] = \emph{registerT}

\paragraph{Description} The effective address is the value of the \emph{registerZ}.
This instruction implements atomic load-EOR-store on word. The word at effective address is EORed with the \emph{registerT}, and its previous value is returned into the \emph{registerT}. The effective address must be word aligned (multiple of 4).


\paragraph{Modifier(s)}
\noindent\begin{itemize}
\item coherency (cf. coherency in section \ref{par:modifier-coherency} on page \pageref{par:modifier-coherency})
\end{itemize}

\paragraph{Execution} {Reservation table \textsf{LSU2\_MEMW\_AUXR\_AUXW} (Section~\ref{sec:reservation-LSU2-MEMW-AUXR-AUXW})}
~
{\begin{lstlisting}
stage ID:
  new coherency = _ZX_2(coherency);
stage RR:
  new address = GPR[registerZ];
stage E1:
  new argument2 = GPR[registerT];
  new result2 = _ZX_32(MEM.atomic_eor(address, 0xF, coherency << 32, argument2.32[0], registerT));
stage SR:
  GPR[registerT] = result2;
\end{lstlisting}}
\newpage

\paragraph{Encoding} Format \textsf{LSU\_ALEORSB.O} (Section~\ref{sec:format-LSU-ALEORSB.O}), \verb|lsucode5 = 010|\\

\bitfields{ 1/P, 2/00, 2/-- --, 27/offset27, 1/1, 2/01, 3/111, 2/coherency, 6/registerT, 2/11, 3/010, 1/--, 5/00111, 1/0, 6/registerZ }


\paragraph{Syntax} \texttt{aleorw}\emph{coherency} \emph{offset27}[\emph{registerZ}] = \emph{registerT}

\paragraph{Description} The effective address is computed by adding the \emph{offset27} to the \emph{registerZ}.
This instruction implements atomic load-EOR-store on word. The word at effective address is EORed with the \emph{registerT}, and its previous value is returned into the \emph{registerT}. The effective address must be word aligned (multiple of 4).


\paragraph{Modifier(s)}
\noindent\begin{itemize}
\item coherency (cf. coherency in section \ref{par:modifier-coherency} on page \pageref{par:modifier-coherency})
\end{itemize}

\paragraph{Execution} {Reservation table \textsf{LSU2\_MEMW\_AUXR\_AUXW.X} (Section~\ref{sec:reservation-LSU2-MEMW-AUXR-AUXW.X})}
~
{\begin{lstlisting}
stage ID:
  new offset = _SX_27(offset27);
  new coherency = _ZX_2(coherency);
stage RR:
  new base = GPR[registerZ];
  new address = base + offset;
stage E1:
  new argument2 = GPR[registerT];
  new result2 = _ZX_32(MEM.atomic_eor(address, 0xF, coherency << 32, argument2.32[0], registerT));
stage SR:
  GPR[registerT] = result2;
\end{lstlisting}}
\newpage

\paragraph{Encoding} Format \textsf{LSU\_ALEORSB.D} (Section~\ref{sec:format-LSU-ALEORSB.D}), \verb|lsucode5 = 010|\\

\bitfields{ 1/P, 2/00, 2/-- --, 27/extend27, 1/1, 2/00, 2/-- --, 27/offset27, 1/1, 2/01, 3/111, 2/coherency, 6/registerT, 2/11, 3/010, 1/--, 5/00111, 1/0, 6/registerZ }


\paragraph{Syntax} \texttt{aleorw}\emph{coherency} \emph{extend27\_\discretionary{}{}{}offset27}[\emph{registerZ}] = \emph{registerT}

\paragraph{Description} The effective address is computed by adding the \emph{extend27\_\discretionary{}{}{}offset27} to the \emph{registerZ}.
This instruction implements atomic load-EOR-store on word. The word at effective address is EORed with the \emph{registerT}, and its previous value is returned into the \emph{registerT}. The effective address must be word aligned (multiple of 4).


\paragraph{Modifier(s)}
\noindent\begin{itemize}
\item coherency (cf. coherency in section \ref{par:modifier-coherency} on page \pageref{par:modifier-coherency})
\end{itemize}

\paragraph{Execution} {Reservation table \textsf{LSU2\_MEMW\_AUXR\_AUXW.Y} (Section~\ref{sec:reservation-LSU2-MEMW-AUXR-AUXW.Y})}
~
{\begin{lstlisting}
stage ID:
  new offset = _SX_54(extend27_offset27);
  new coherency = _ZX_2(coherency);
stage RR:
  new base = GPR[registerZ];
  new address = base + offset;
stage E1:
  new argument2 = GPR[registerT];
  new result2 = _ZX_32(MEM.atomic_eor(address, 0xF, coherency << 32, argument2.32[0], registerT));
stage SR:
  GPR[registerT] = result2;
\end{lstlisting}}
\newpage
\subsubsection[{\small ALH: Atomic Load Half Word}]{ALH\hfill Atomic Load Half Word}
\label{instr:ALH}\index{alh}
 
\paragraph{Encoding} Format \textsf{LSU\_ALSB} (Section~\ref{sec:format-LSU-ALSB}), \verb|lsucode5 = 001|\\

\bitfields{ 1/P, 2/01, 3/111, 2/coherency, 6/registerW, 2/11, 3/001, 1/--, 5/00000, 1/0, 6/registerZ }


\paragraph{Syntax} \texttt{alh}\emph{coherency} \emph{registerW} = [\emph{registerZ}]

\paragraph{Description} The effective address is the value of the \emph{registerZ}.
The \emph{registerW} is atomically loaded with zero extension from the memory half word starting at the effective address. This instruction is provided to directly operate on the last level of the memory hierarchy, which is typically a processor DDR or a cluster TCM, irrespective of the MMU DCP (Data Cache Policy) settings. The effective address must be half word aligned (multiple of 2).


\paragraph{Modifier(s)}
\noindent\begin{itemize}
\item coherency (cf. coherency in section \ref{par:modifier-coherency} on page \pageref{par:modifier-coherency})
\end{itemize}

\paragraph{Execution} {Reservation table \textsf{LSU\_MEMW\_AUXW} (Section~\ref{sec:reservation-LSU-MEMW-AUXW})}
~
{\begin{lstlisting}
stage ID:
  new coherency = _ZX_2(coherency);
stage RR:
  new address = GPR[registerZ];
  new result1 = _ZX_16(MEM.atomic_load(address, 0x3, coherency << 32, registerW));
stage CR:
  GPR[registerW] = result1;
\end{lstlisting}}
\newpage

\paragraph{Encoding} Format \textsf{LSU\_ALSB.O} (Section~\ref{sec:format-LSU-ALSB.O}), \verb|lsucode5 = 001|\\

\bitfields{ 1/P, 2/00, 2/-- --, 27/offset27, 1/1, 2/01, 3/111, 2/coherency, 6/registerW, 2/11, 3/001, 1/--, 5/00000, 1/0, 6/registerZ }


\paragraph{Syntax} \texttt{alh}\emph{coherency} \emph{registerW} = \emph{offset27}[\emph{registerZ}]

\paragraph{Description} The effective address is computed by adding the \emph{offset27} to the \emph{registerZ}.
The \emph{registerW} is atomically loaded with zero extension from the memory half word starting at the effective address. This instruction is provided to directly operate on the last level of the memory hierarchy, which is typically a processor DDR or a cluster TCM, irrespective of the MMU DCP (Data Cache Policy) settings. The effective address must be half word aligned (multiple of 2).


\paragraph{Modifier(s)}
\noindent\begin{itemize}
\item coherency (cf. coherency in section \ref{par:modifier-coherency} on page \pageref{par:modifier-coherency})
\end{itemize}

\paragraph{Execution} {Reservation table \textsf{LSU\_MEMW\_AUXW.X} (Section~\ref{sec:reservation-LSU-MEMW-AUXW.X})}
~
{\begin{lstlisting}
stage ID:
  new offset = _SX_27(offset27);
  new coherency = _ZX_2(coherency);
stage RR:
  new base = GPR[registerZ];
  new address = base + offset;
  new result1 = _ZX_16(MEM.atomic_load(address, 0x3, coherency << 32, registerW));
stage CR:
  GPR[registerW] = result1;
\end{lstlisting}}
\newpage

\paragraph{Encoding} Format \textsf{LSU\_ALSB.D} (Section~\ref{sec:format-LSU-ALSB.D}), \verb|lsucode5 = 001|\\

\bitfields{ 1/P, 2/00, 2/-- --, 27/extend27, 1/1, 2/00, 2/-- --, 27/offset27, 1/1, 2/01, 3/111, 2/coherency, 6/registerW, 2/11, 3/001, 1/--, 5/00000, 1/0, 6/registerZ }


\paragraph{Syntax} \texttt{alh}\emph{coherency} \emph{registerW} = \emph{extend27\_\discretionary{}{}{}offset27}[\emph{registerZ}]

\paragraph{Description} The effective address is computed by adding the \emph{extend27\_\discretionary{}{}{}offset27} to the \emph{registerZ}.
The \emph{registerW} is atomically loaded with zero extension from the memory half word starting at the effective address. This instruction is provided to directly operate on the last level of the memory hierarchy, which is typically a processor DDR or a cluster TCM, irrespective of the MMU DCP (Data Cache Policy) settings. The effective address must be half word aligned (multiple of 2).


\paragraph{Modifier(s)}
\noindent\begin{itemize}
\item coherency (cf. coherency in section \ref{par:modifier-coherency} on page \pageref{par:modifier-coherency})
\end{itemize}

\paragraph{Execution} {Reservation table \textsf{LSU\_MEMW\_AUXW.Y} (Section~\ref{sec:reservation-LSU-MEMW-AUXW.Y})}
~
{\begin{lstlisting}
stage ID:
  new offset = _SX_54(extend27_offset27);
  new coherency = _ZX_2(coherency);
stage RR:
  new base = GPR[registerZ];
  new address = base + offset;
  new result1 = _ZX_16(MEM.atomic_load(address, 0x3, coherency << 32, registerW));
stage CR:
  GPR[registerW] = result1;
\end{lstlisting}}
\newpage
\subsubsection[{\small ALIORB: Atomic Load and IOR Byte}]{ALIORB\hfill Atomic Load and IOR Byte}
\label{instr:ALIORB}\index{aliorb}
 
\paragraph{Encoding} Format \textsf{LSU\_ALIORSB} (Section~\ref{sec:format-LSU-ALIORSB}), \verb|lsucode5 = 000|\\

\bitfields{ 1/P, 2/01, 3/111, 2/coherency, 6/registerT, 2/11, 3/000, 1/--, 5/00110, 1/0, 6/registerZ }


\paragraph{Syntax} \texttt{aliorb}\emph{coherency} [\emph{registerZ}] = \emph{registerT}

\paragraph{Description} The effective address is the value of the \emph{registerZ}.
This instruction implements atomic load-IOR-store on byte. The byte at effective address is IORed with the \emph{registerT}, and its previous value is returned into the \emph{registerT}.


\paragraph{Modifier(s)}
\noindent\begin{itemize}
\item coherency (cf. coherency in section \ref{par:modifier-coherency} on page \pageref{par:modifier-coherency})
\end{itemize}

\paragraph{Execution} {Reservation table \textsf{LSU2\_MEMW\_AUXR\_AUXW} (Section~\ref{sec:reservation-LSU2-MEMW-AUXR-AUXW})}
~
{\begin{lstlisting}
stage ID:
  new coherency = _ZX_2(coherency);
stage RR:
  new address = GPR[registerZ];
stage E1:
  new argument2 = GPR[registerT];
  new result2 = _ZX_8(MEM.atomic_ior(address, 0x1, coherency << 32, argument2.8[0], registerT));
stage SR:
  GPR[registerT] = result2;
\end{lstlisting}}
\newpage

\paragraph{Encoding} Format \textsf{LSU\_ALIORSB.O} (Section~\ref{sec:format-LSU-ALIORSB.O}), \verb|lsucode5 = 000|\\

\bitfields{ 1/P, 2/00, 2/-- --, 27/offset27, 1/1, 2/01, 3/111, 2/coherency, 6/registerT, 2/11, 3/000, 1/--, 5/00110, 1/0, 6/registerZ }


\paragraph{Syntax} \texttt{aliorb}\emph{coherency} \emph{offset27}[\emph{registerZ}] = \emph{registerT}

\paragraph{Description} The effective address is computed by adding the \emph{offset27} to the \emph{registerZ}.
This instruction implements atomic load-IOR-store on byte. The byte at effective address is IORed with the \emph{registerT}, and its previous value is returned into the \emph{registerT}.


\paragraph{Modifier(s)}
\noindent\begin{itemize}
\item coherency (cf. coherency in section \ref{par:modifier-coherency} on page \pageref{par:modifier-coherency})
\end{itemize}

\paragraph{Execution} {Reservation table \textsf{LSU2\_MEMW\_AUXR\_AUXW.X} (Section~\ref{sec:reservation-LSU2-MEMW-AUXR-AUXW.X})}
~
{\begin{lstlisting}
stage ID:
  new offset = _SX_27(offset27);
  new coherency = _ZX_2(coherency);
stage RR:
  new base = GPR[registerZ];
  new address = base + offset;
stage E1:
  new argument2 = GPR[registerT];
  new result2 = _ZX_8(MEM.atomic_ior(address, 0x1, coherency << 32, argument2.8[0], registerT));
stage SR:
  GPR[registerT] = result2;
\end{lstlisting}}
\newpage

\paragraph{Encoding} Format \textsf{LSU\_ALIORSB.D} (Section~\ref{sec:format-LSU-ALIORSB.D}), \verb|lsucode5 = 000|\\

\bitfields{ 1/P, 2/00, 2/-- --, 27/extend27, 1/1, 2/00, 2/-- --, 27/offset27, 1/1, 2/01, 3/111, 2/coherency, 6/registerT, 2/11, 3/000, 1/--, 5/00110, 1/0, 6/registerZ }


\paragraph{Syntax} \texttt{aliorb}\emph{coherency} \emph{extend27\_\discretionary{}{}{}offset27}[\emph{registerZ}] = \emph{registerT}

\paragraph{Description} The effective address is computed by adding the \emph{extend27\_\discretionary{}{}{}offset27} to the \emph{registerZ}.
This instruction implements atomic load-IOR-store on byte. The byte at effective address is IORed with the \emph{registerT}, and its previous value is returned into the \emph{registerT}.


\paragraph{Modifier(s)}
\noindent\begin{itemize}
\item coherency (cf. coherency in section \ref{par:modifier-coherency} on page \pageref{par:modifier-coherency})
\end{itemize}

\paragraph{Execution} {Reservation table \textsf{LSU2\_MEMW\_AUXR\_AUXW.Y} (Section~\ref{sec:reservation-LSU2-MEMW-AUXR-AUXW.Y})}
~
{\begin{lstlisting}
stage ID:
  new offset = _SX_54(extend27_offset27);
  new coherency = _ZX_2(coherency);
stage RR:
  new base = GPR[registerZ];
  new address = base + offset;
stage E1:
  new argument2 = GPR[registerT];
  new result2 = _ZX_8(MEM.atomic_ior(address, 0x1, coherency << 32, argument2.8[0], registerT));
stage SR:
  GPR[registerT] = result2;
\end{lstlisting}}
\newpage
\subsubsection[{\small ALIORD: Atomic Load and IOR Double Word}]{ALIORD\hfill Atomic Load and IOR Double Word}
\label{instr:ALIORD}\index{aliord}
 
\paragraph{Encoding} Format \textsf{LSU\_ALIORSB} (Section~\ref{sec:format-LSU-ALIORSB}), \verb|lsucode5 = 011|\\

\bitfields{ 1/P, 2/01, 3/111, 2/coherency, 6/registerT, 2/11, 3/011, 1/--, 5/00110, 1/0, 6/registerZ }


\paragraph{Syntax} \texttt{aliord}\emph{coherency} [\emph{registerZ}] = \emph{registerT}

\paragraph{Description} The effective address is the value of the \emph{registerZ}.
This instruction implements atomic load-IOR-store on double word. The double word at effective address is IORed with the \emph{registerT}, and its previous value is returned into the \emph{registerT}. The effective address must be double word aligned (multiple of 8).


\paragraph{Modifier(s)}
\noindent\begin{itemize}
\item coherency (cf. coherency in section \ref{par:modifier-coherency} on page \pageref{par:modifier-coherency})
\end{itemize}

\paragraph{Execution} {Reservation table \textsf{LSU2\_MEMW\_AUXR\_AUXW} (Section~\ref{sec:reservation-LSU2-MEMW-AUXR-AUXW})}
~
{\begin{lstlisting}
stage ID:
  new coherency = _ZX_2(coherency);
stage RR:
  new address = GPR[registerZ];
stage E1:
  new argument2 = GPR[registerT];
  new result2 = MEM.atomic_ior(address, 0xFF, coherency << 32, argument2.64[0], registerT);
stage SR:
  GPR[registerT] = result2;
\end{lstlisting}}
\newpage

\paragraph{Encoding} Format \textsf{LSU\_ALIORSB.O} (Section~\ref{sec:format-LSU-ALIORSB.O}), \verb|lsucode5 = 011|\\

\bitfields{ 1/P, 2/00, 2/-- --, 27/offset27, 1/1, 2/01, 3/111, 2/coherency, 6/registerT, 2/11, 3/011, 1/--, 5/00110, 1/0, 6/registerZ }


\paragraph{Syntax} \texttt{aliord}\emph{coherency} \emph{offset27}[\emph{registerZ}] = \emph{registerT}

\paragraph{Description} The effective address is computed by adding the \emph{offset27} to the \emph{registerZ}.
This instruction implements atomic load-IOR-store on double word. The double word at effective address is IORed with the \emph{registerT}, and its previous value is returned into the \emph{registerT}. The effective address must be double word aligned (multiple of 8).


\paragraph{Modifier(s)}
\noindent\begin{itemize}
\item coherency (cf. coherency in section \ref{par:modifier-coherency} on page \pageref{par:modifier-coherency})
\end{itemize}

\paragraph{Execution} {Reservation table \textsf{LSU2\_MEMW\_AUXR\_AUXW.X} (Section~\ref{sec:reservation-LSU2-MEMW-AUXR-AUXW.X})}
~
{\begin{lstlisting}
stage ID:
  new offset = _SX_27(offset27);
  new coherency = _ZX_2(coherency);
stage RR:
  new base = GPR[registerZ];
  new address = base + offset;
stage E1:
  new argument2 = GPR[registerT];
  new result2 = MEM.atomic_ior(address, 0xFF, coherency << 32, argument2.64[0], registerT);
stage SR:
  GPR[registerT] = result2;
\end{lstlisting}}
\newpage

\paragraph{Encoding} Format \textsf{LSU\_ALIORSB.D} (Section~\ref{sec:format-LSU-ALIORSB.D}), \verb|lsucode5 = 011|\\

\bitfields{ 1/P, 2/00, 2/-- --, 27/extend27, 1/1, 2/00, 2/-- --, 27/offset27, 1/1, 2/01, 3/111, 2/coherency, 6/registerT, 2/11, 3/011, 1/--, 5/00110, 1/0, 6/registerZ }


\paragraph{Syntax} \texttt{aliord}\emph{coherency} \emph{extend27\_\discretionary{}{}{}offset27}[\emph{registerZ}] = \emph{registerT}

\paragraph{Description} The effective address is computed by adding the \emph{extend27\_\discretionary{}{}{}offset27} to the \emph{registerZ}.
This instruction implements atomic load-IOR-store on double word. The double word at effective address is IORed with the \emph{registerT}, and its previous value is returned into the \emph{registerT}. The effective address must be double word aligned (multiple of 8).


\paragraph{Modifier(s)}
\noindent\begin{itemize}
\item coherency (cf. coherency in section \ref{par:modifier-coherency} on page \pageref{par:modifier-coherency})
\end{itemize}

\paragraph{Execution} {Reservation table \textsf{LSU2\_MEMW\_AUXR\_AUXW.Y} (Section~\ref{sec:reservation-LSU2-MEMW-AUXR-AUXW.Y})}
~
{\begin{lstlisting}
stage ID:
  new offset = _SX_54(extend27_offset27);
  new coherency = _ZX_2(coherency);
stage RR:
  new base = GPR[registerZ];
  new address = base + offset;
stage E1:
  new argument2 = GPR[registerT];
  new result2 = MEM.atomic_ior(address, 0xFF, coherency << 32, argument2.64[0], registerT);
stage SR:
  GPR[registerT] = result2;
\end{lstlisting}}
\newpage
\subsubsection[{\small ALIORH: Atomic Load and IOR Half Word}]{ALIORH\hfill Atomic Load and IOR Half Word}
\label{instr:ALIORH}\index{aliorh}
 
\paragraph{Encoding} Format \textsf{LSU\_ALIORSB} (Section~\ref{sec:format-LSU-ALIORSB}), \verb|lsucode5 = 001|\\

\bitfields{ 1/P, 2/01, 3/111, 2/coherency, 6/registerT, 2/11, 3/001, 1/--, 5/00110, 1/0, 6/registerZ }


\paragraph{Syntax} \texttt{aliorh}\emph{coherency} [\emph{registerZ}] = \emph{registerT}

\paragraph{Description} The effective address is the value of the \emph{registerZ}.
This instruction implements atomic load-IOR-store on half word. The half word at effective address is IORed with the \emph{registerT}, and its previous value is returned into the \emph{registerT}. The effective address must be half word aligned (multiple of 2).


\paragraph{Modifier(s)}
\noindent\begin{itemize}
\item coherency (cf. coherency in section \ref{par:modifier-coherency} on page \pageref{par:modifier-coherency})
\end{itemize}

\paragraph{Execution} {Reservation table \textsf{LSU2\_MEMW\_AUXR\_AUXW} (Section~\ref{sec:reservation-LSU2-MEMW-AUXR-AUXW})}
~
{\begin{lstlisting}
stage ID:
  new coherency = _ZX_2(coherency);
stage RR:
  new address = GPR[registerZ];
stage E1:
  new argument2 = GPR[registerT];
  new result2 = _ZX_16(MEM.atomic_ior(address, 0x3, coherency << 32, argument2.16[0], registerT));
stage SR:
  GPR[registerT] = result2;
\end{lstlisting}}
\newpage

\paragraph{Encoding} Format \textsf{LSU\_ALIORSB.O} (Section~\ref{sec:format-LSU-ALIORSB.O}), \verb|lsucode5 = 001|\\

\bitfields{ 1/P, 2/00, 2/-- --, 27/offset27, 1/1, 2/01, 3/111, 2/coherency, 6/registerT, 2/11, 3/001, 1/--, 5/00110, 1/0, 6/registerZ }


\paragraph{Syntax} \texttt{aliorh}\emph{coherency} \emph{offset27}[\emph{registerZ}] = \emph{registerT}

\paragraph{Description} The effective address is computed by adding the \emph{offset27} to the \emph{registerZ}.
This instruction implements atomic load-IOR-store on half word. The half word at effective address is IORed with the \emph{registerT}, and its previous value is returned into the \emph{registerT}. The effective address must be half word aligned (multiple of 2).


\paragraph{Modifier(s)}
\noindent\begin{itemize}
\item coherency (cf. coherency in section \ref{par:modifier-coherency} on page \pageref{par:modifier-coherency})
\end{itemize}

\paragraph{Execution} {Reservation table \textsf{LSU2\_MEMW\_AUXR\_AUXW.X} (Section~\ref{sec:reservation-LSU2-MEMW-AUXR-AUXW.X})}
~
{\begin{lstlisting}
stage ID:
  new offset = _SX_27(offset27);
  new coherency = _ZX_2(coherency);
stage RR:
  new base = GPR[registerZ];
  new address = base + offset;
stage E1:
  new argument2 = GPR[registerT];
  new result2 = _ZX_16(MEM.atomic_ior(address, 0x3, coherency << 32, argument2.16[0], registerT));
stage SR:
  GPR[registerT] = result2;
\end{lstlisting}}
\newpage

\paragraph{Encoding} Format \textsf{LSU\_ALIORSB.D} (Section~\ref{sec:format-LSU-ALIORSB.D}), \verb|lsucode5 = 001|\\

\bitfields{ 1/P, 2/00, 2/-- --, 27/extend27, 1/1, 2/00, 2/-- --, 27/offset27, 1/1, 2/01, 3/111, 2/coherency, 6/registerT, 2/11, 3/001, 1/--, 5/00110, 1/0, 6/registerZ }


\paragraph{Syntax} \texttt{aliorh}\emph{coherency} \emph{extend27\_\discretionary{}{}{}offset27}[\emph{registerZ}] = \emph{registerT}

\paragraph{Description} The effective address is computed by adding the \emph{extend27\_\discretionary{}{}{}offset27} to the \emph{registerZ}.
This instruction implements atomic load-IOR-store on half word. The half word at effective address is IORed with the \emph{registerT}, and its previous value is returned into the \emph{registerT}. The effective address must be half word aligned (multiple of 2).


\paragraph{Modifier(s)}
\noindent\begin{itemize}
\item coherency (cf. coherency in section \ref{par:modifier-coherency} on page \pageref{par:modifier-coherency})
\end{itemize}

\paragraph{Execution} {Reservation table \textsf{LSU2\_MEMW\_AUXR\_AUXW.Y} (Section~\ref{sec:reservation-LSU2-MEMW-AUXR-AUXW.Y})}
~
{\begin{lstlisting}
stage ID:
  new offset = _SX_54(extend27_offset27);
  new coherency = _ZX_2(coherency);
stage RR:
  new base = GPR[registerZ];
  new address = base + offset;
stage E1:
  new argument2 = GPR[registerT];
  new result2 = _ZX_16(MEM.atomic_ior(address, 0x3, coherency << 32, argument2.16[0], registerT));
stage SR:
  GPR[registerT] = result2;
\end{lstlisting}}
\newpage
\subsubsection[{\small ALIORW: Atomic Load and IOR Word}]{ALIORW\hfill Atomic Load and IOR Word}
\label{instr:ALIORW}\index{aliorw}
 
\paragraph{Encoding} Format \textsf{LSU\_ALIORSB} (Section~\ref{sec:format-LSU-ALIORSB}), \verb|lsucode5 = 010|\\

\bitfields{ 1/P, 2/01, 3/111, 2/coherency, 6/registerT, 2/11, 3/010, 1/--, 5/00110, 1/0, 6/registerZ }


\paragraph{Syntax} \texttt{aliorw}\emph{coherency} [\emph{registerZ}] = \emph{registerT}

\paragraph{Description} The effective address is the value of the \emph{registerZ}.
This instruction implements atomic load-IOR-store on word. The word at effective address is IORed with the \emph{registerT}, and its previous value is returned into the \emph{registerT}. The effective address must be word aligned (multiple of 4).


\paragraph{Modifier(s)}
\noindent\begin{itemize}
\item coherency (cf. coherency in section \ref{par:modifier-coherency} on page \pageref{par:modifier-coherency})
\end{itemize}

\paragraph{Execution} {Reservation table \textsf{LSU2\_MEMW\_AUXR\_AUXW} (Section~\ref{sec:reservation-LSU2-MEMW-AUXR-AUXW})}
~
{\begin{lstlisting}
stage ID:
  new coherency = _ZX_2(coherency);
stage RR:
  new address = GPR[registerZ];
stage E1:
  new argument2 = GPR[registerT];
  new result2 = _ZX_32(MEM.atomic_ior(address, 0xF, coherency << 32, argument2.32[0], registerT));
stage SR:
  GPR[registerT] = result2;
\end{lstlisting}}
\newpage

\paragraph{Encoding} Format \textsf{LSU\_ALIORSB.O} (Section~\ref{sec:format-LSU-ALIORSB.O}), \verb|lsucode5 = 010|\\

\bitfields{ 1/P, 2/00, 2/-- --, 27/offset27, 1/1, 2/01, 3/111, 2/coherency, 6/registerT, 2/11, 3/010, 1/--, 5/00110, 1/0, 6/registerZ }


\paragraph{Syntax} \texttt{aliorw}\emph{coherency} \emph{offset27}[\emph{registerZ}] = \emph{registerT}

\paragraph{Description} The effective address is computed by adding the \emph{offset27} to the \emph{registerZ}.
This instruction implements atomic load-IOR-store on word. The word at effective address is IORed with the \emph{registerT}, and its previous value is returned into the \emph{registerT}. The effective address must be word aligned (multiple of 4).


\paragraph{Modifier(s)}
\noindent\begin{itemize}
\item coherency (cf. coherency in section \ref{par:modifier-coherency} on page \pageref{par:modifier-coherency})
\end{itemize}

\paragraph{Execution} {Reservation table \textsf{LSU2\_MEMW\_AUXR\_AUXW.X} (Section~\ref{sec:reservation-LSU2-MEMW-AUXR-AUXW.X})}
~
{\begin{lstlisting}
stage ID:
  new offset = _SX_27(offset27);
  new coherency = _ZX_2(coherency);
stage RR:
  new base = GPR[registerZ];
  new address = base + offset;
stage E1:
  new argument2 = GPR[registerT];
  new result2 = _ZX_32(MEM.atomic_ior(address, 0xF, coherency << 32, argument2.32[0], registerT));
stage SR:
  GPR[registerT] = result2;
\end{lstlisting}}
\newpage

\paragraph{Encoding} Format \textsf{LSU\_ALIORSB.D} (Section~\ref{sec:format-LSU-ALIORSB.D}), \verb|lsucode5 = 010|\\

\bitfields{ 1/P, 2/00, 2/-- --, 27/extend27, 1/1, 2/00, 2/-- --, 27/offset27, 1/1, 2/01, 3/111, 2/coherency, 6/registerT, 2/11, 3/010, 1/--, 5/00110, 1/0, 6/registerZ }


\paragraph{Syntax} \texttt{aliorw}\emph{coherency} \emph{extend27\_\discretionary{}{}{}offset27}[\emph{registerZ}] = \emph{registerT}

\paragraph{Description} The effective address is computed by adding the \emph{extend27\_\discretionary{}{}{}offset27} to the \emph{registerZ}.
This instruction implements atomic load-IOR-store on word. The word at effective address is IORed with the \emph{registerT}, and its previous value is returned into the \emph{registerT}. The effective address must be word aligned (multiple of 4).


\paragraph{Modifier(s)}
\noindent\begin{itemize}
\item coherency (cf. coherency in section \ref{par:modifier-coherency} on page \pageref{par:modifier-coherency})
\end{itemize}

\paragraph{Execution} {Reservation table \textsf{LSU2\_MEMW\_AUXR\_AUXW.Y} (Section~\ref{sec:reservation-LSU2-MEMW-AUXR-AUXW.Y})}
~
{\begin{lstlisting}
stage ID:
  new offset = _SX_54(extend27_offset27);
  new coherency = _ZX_2(coherency);
stage RR:
  new base = GPR[registerZ];
  new address = base + offset;
stage E1:
  new argument2 = GPR[registerT];
  new result2 = _ZX_32(MEM.atomic_ior(address, 0xF, coherency << 32, argument2.32[0], registerT));
stage SR:
  GPR[registerT] = result2;
\end{lstlisting}}
\newpage
\subsubsection[{\small ALMAXB: Atomic Load and MAX Byte}]{ALMAXB\hfill Atomic Load and MAX Byte}
\label{instr:ALMAXB}\index{almaxb}
 
\paragraph{Encoding} Format \textsf{LSU\_ALMAXSB} (Section~\ref{sec:format-LSU-ALMAXSB}), \verb|lsucode5 = 000|\\

\bitfields{ 1/P, 2/01, 3/111, 2/coherency, 6/registerT, 2/11, 3/000, 1/--, 5/01001, 1/0, 6/registerZ }


\paragraph{Syntax} \texttt{almaxb}\emph{coherency} [\emph{registerZ}] = \emph{registerT}

\paragraph{Description} The effective address is the value of the \emph{registerZ}.
This instruction implements atomic load-MAX-store on signed byte. The byte at effective address is MAXed with the \emph{registerT}, and its previous value is returned into the \emph{registerT}.


\paragraph{Modifier(s)}
\noindent\begin{itemize}
\item coherency (cf. coherency in section \ref{par:modifier-coherency} on page \pageref{par:modifier-coherency})
\end{itemize}

\paragraph{Execution} {Reservation table \textsf{LSU2\_MEMW\_AUXR\_AUXW} (Section~\ref{sec:reservation-LSU2-MEMW-AUXR-AUXW})}
~
{\begin{lstlisting}
stage ID:
  new coherency = _ZX_2(coherency);
stage RR:
  new address = GPR[registerZ];
stage E1:
  new argument2 = GPR[registerT];
  new result2 = _ZX_8(MEM.atomic_max(address, 0x1, coherency << 32, argument2.8[0], registerT));
stage SR:
  GPR[registerT] = result2;
\end{lstlisting}}
\newpage

\paragraph{Encoding} Format \textsf{LSU\_ALMAXSB.O} (Section~\ref{sec:format-LSU-ALMAXSB.O}), \verb|lsucode5 = 000|\\

\bitfields{ 1/P, 2/00, 2/-- --, 27/offset27, 1/1, 2/01, 3/111, 2/coherency, 6/registerT, 2/11, 3/000, 1/--, 5/01001, 1/0, 6/registerZ }


\paragraph{Syntax} \texttt{almaxb}\emph{coherency} \emph{offset27}[\emph{registerZ}] = \emph{registerT}

\paragraph{Description} The effective address is computed by adding the \emph{offset27} to the \emph{registerZ}.
This instruction implements atomic load-MAX-store on signed byte. The byte at effective address is MAXed with the \emph{registerT}, and its previous value is returned into the \emph{registerT}.


\paragraph{Modifier(s)}
\noindent\begin{itemize}
\item coherency (cf. coherency in section \ref{par:modifier-coherency} on page \pageref{par:modifier-coherency})
\end{itemize}

\paragraph{Execution} {Reservation table \textsf{LSU2\_MEMW\_AUXR\_AUXW.X} (Section~\ref{sec:reservation-LSU2-MEMW-AUXR-AUXW.X})}
~
{\begin{lstlisting}
stage ID:
  new offset = _SX_27(offset27);
  new coherency = _ZX_2(coherency);
stage RR:
  new base = GPR[registerZ];
  new address = base + offset;
stage E1:
  new argument2 = GPR[registerT];
  new result2 = _ZX_8(MEM.atomic_max(address, 0x1, coherency << 32, argument2.8[0], registerT));
stage SR:
  GPR[registerT] = result2;
\end{lstlisting}}
\newpage

\paragraph{Encoding} Format \textsf{LSU\_ALMAXSB.D} (Section~\ref{sec:format-LSU-ALMAXSB.D}), \verb|lsucode5 = 000|\\

\bitfields{ 1/P, 2/00, 2/-- --, 27/extend27, 1/1, 2/00, 2/-- --, 27/offset27, 1/1, 2/01, 3/111, 2/coherency, 6/registerT, 2/11, 3/000, 1/--, 5/01001, 1/0, 6/registerZ }


\paragraph{Syntax} \texttt{almaxb}\emph{coherency} \emph{extend27\_\discretionary{}{}{}offset27}[\emph{registerZ}] = \emph{registerT}

\paragraph{Description} The effective address is computed by adding the \emph{extend27\_\discretionary{}{}{}offset27} to the \emph{registerZ}.
This instruction implements atomic load-MAX-store on signed byte. The byte at effective address is MAXed with the \emph{registerT}, and its previous value is returned into the \emph{registerT}.


\paragraph{Modifier(s)}
\noindent\begin{itemize}
\item coherency (cf. coherency in section \ref{par:modifier-coherency} on page \pageref{par:modifier-coherency})
\end{itemize}

\paragraph{Execution} {Reservation table \textsf{LSU2\_MEMW\_AUXR\_AUXW.Y} (Section~\ref{sec:reservation-LSU2-MEMW-AUXR-AUXW.Y})}
~
{\begin{lstlisting}
stage ID:
  new offset = _SX_54(extend27_offset27);
  new coherency = _ZX_2(coherency);
stage RR:
  new base = GPR[registerZ];
  new address = base + offset;
stage E1:
  new argument2 = GPR[registerT];
  new result2 = _ZX_8(MEM.atomic_max(address, 0x1, coherency << 32, argument2.8[0], registerT));
stage SR:
  GPR[registerT] = result2;
\end{lstlisting}}
\newpage
\subsubsection[{\small ALMAXD: Atomic Load and MAX Double Word}]{ALMAXD\hfill Atomic Load and MAX Double Word}
\label{instr:ALMAXD}\index{almaxd}
 
\paragraph{Encoding} Format \textsf{LSU\_ALMAXSB} (Section~\ref{sec:format-LSU-ALMAXSB}), \verb|lsucode5 = 011|\\

\bitfields{ 1/P, 2/01, 3/111, 2/coherency, 6/registerT, 2/11, 3/011, 1/--, 5/01001, 1/0, 6/registerZ }


\paragraph{Syntax} \texttt{almaxd}\emph{coherency} [\emph{registerZ}] = \emph{registerT}

\paragraph{Description} The effective address is the value of the \emph{registerZ}.
This instruction implements atomic load-MAX-store on signed double word. The double word at effective address is MAXed with the \emph{registerT}, and its previous value is returned into the \emph{registerT}. The effective address must be double word aligned (multiple of 8).


\paragraph{Modifier(s)}
\noindent\begin{itemize}
\item coherency (cf. coherency in section \ref{par:modifier-coherency} on page \pageref{par:modifier-coherency})
\end{itemize}

\paragraph{Execution} {Reservation table \textsf{LSU2\_MEMW\_AUXR\_AUXW} (Section~\ref{sec:reservation-LSU2-MEMW-AUXR-AUXW})}
~
{\begin{lstlisting}
stage ID:
  new coherency = _ZX_2(coherency);
stage RR:
  new address = GPR[registerZ];
stage E1:
  new argument2 = GPR[registerT];
  new result2 = MEM.atomic_max(address, 0xFF, coherency << 32, argument2.64[0], registerT);
stage SR:
  GPR[registerT] = result2;
\end{lstlisting}}
\newpage

\paragraph{Encoding} Format \textsf{LSU\_ALMAXSB.O} (Section~\ref{sec:format-LSU-ALMAXSB.O}), \verb|lsucode5 = 011|\\

\bitfields{ 1/P, 2/00, 2/-- --, 27/offset27, 1/1, 2/01, 3/111, 2/coherency, 6/registerT, 2/11, 3/011, 1/--, 5/01001, 1/0, 6/registerZ }


\paragraph{Syntax} \texttt{almaxd}\emph{coherency} \emph{offset27}[\emph{registerZ}] = \emph{registerT}

\paragraph{Description} The effective address is computed by adding the \emph{offset27} to the \emph{registerZ}.
This instruction implements atomic load-MAX-store on signed double word. The double word at effective address is MAXed with the \emph{registerT}, and its previous value is returned into the \emph{registerT}. The effective address must be double word aligned (multiple of 8).


\paragraph{Modifier(s)}
\noindent\begin{itemize}
\item coherency (cf. coherency in section \ref{par:modifier-coherency} on page \pageref{par:modifier-coherency})
\end{itemize}

\paragraph{Execution} {Reservation table \textsf{LSU2\_MEMW\_AUXR\_AUXW.X} (Section~\ref{sec:reservation-LSU2-MEMW-AUXR-AUXW.X})}
~
{\begin{lstlisting}
stage ID:
  new offset = _SX_27(offset27);
  new coherency = _ZX_2(coherency);
stage RR:
  new base = GPR[registerZ];
  new address = base + offset;
stage E1:
  new argument2 = GPR[registerT];
  new result2 = MEM.atomic_max(address, 0xFF, coherency << 32, argument2.64[0], registerT);
stage SR:
  GPR[registerT] = result2;
\end{lstlisting}}
\newpage

\paragraph{Encoding} Format \textsf{LSU\_ALMAXSB.D} (Section~\ref{sec:format-LSU-ALMAXSB.D}), \verb|lsucode5 = 011|\\

\bitfields{ 1/P, 2/00, 2/-- --, 27/extend27, 1/1, 2/00, 2/-- --, 27/offset27, 1/1, 2/01, 3/111, 2/coherency, 6/registerT, 2/11, 3/011, 1/--, 5/01001, 1/0, 6/registerZ }


\paragraph{Syntax} \texttt{almaxd}\emph{coherency} \emph{extend27\_\discretionary{}{}{}offset27}[\emph{registerZ}] = \emph{registerT}

\paragraph{Description} The effective address is computed by adding the \emph{extend27\_\discretionary{}{}{}offset27} to the \emph{registerZ}.
This instruction implements atomic load-MAX-store on signed double word. The double word at effective address is MAXed with the \emph{registerT}, and its previous value is returned into the \emph{registerT}. The effective address must be double word aligned (multiple of 8).


\paragraph{Modifier(s)}
\noindent\begin{itemize}
\item coherency (cf. coherency in section \ref{par:modifier-coherency} on page \pageref{par:modifier-coherency})
\end{itemize}

\paragraph{Execution} {Reservation table \textsf{LSU2\_MEMW\_AUXR\_AUXW.Y} (Section~\ref{sec:reservation-LSU2-MEMW-AUXR-AUXW.Y})}
~
{\begin{lstlisting}
stage ID:
  new offset = _SX_54(extend27_offset27);
  new coherency = _ZX_2(coherency);
stage RR:
  new base = GPR[registerZ];
  new address = base + offset;
stage E1:
  new argument2 = GPR[registerT];
  new result2 = MEM.atomic_max(address, 0xFF, coherency << 32, argument2.64[0], registerT);
stage SR:
  GPR[registerT] = result2;
\end{lstlisting}}
\newpage
\subsubsection[{\small ALMAXH: Atomic Load and MAX Half Word}]{ALMAXH\hfill Atomic Load and MAX Half Word}
\label{instr:ALMAXH}\index{almaxh}
 
\paragraph{Encoding} Format \textsf{LSU\_ALMAXSB} (Section~\ref{sec:format-LSU-ALMAXSB}), \verb|lsucode5 = 001|\\

\bitfields{ 1/P, 2/01, 3/111, 2/coherency, 6/registerT, 2/11, 3/001, 1/--, 5/01001, 1/0, 6/registerZ }


\paragraph{Syntax} \texttt{almaxh}\emph{coherency} [\emph{registerZ}] = \emph{registerT}

\paragraph{Description} The effective address is the value of the \emph{registerZ}.
This instruction implements atomic load-MAX-store on signed half word. The half word at effective address is MAXed with the \emph{registerT}, and its previous value is returned into the \emph{registerT}. The effective address must be half word aligned (multiple of 2).


\paragraph{Modifier(s)}
\noindent\begin{itemize}
\item coherency (cf. coherency in section \ref{par:modifier-coherency} on page \pageref{par:modifier-coherency})
\end{itemize}

\paragraph{Execution} {Reservation table \textsf{LSU2\_MEMW\_AUXR\_AUXW} (Section~\ref{sec:reservation-LSU2-MEMW-AUXR-AUXW})}
~
{\begin{lstlisting}
stage ID:
  new coherency = _ZX_2(coherency);
stage RR:
  new address = GPR[registerZ];
stage E1:
  new argument2 = GPR[registerT];
  new result2 = _ZX_16(MEM.atomic_max(address, 0x3, coherency << 32, argument2.16[0], registerT));
stage SR:
  GPR[registerT] = result2;
\end{lstlisting}}
\newpage

\paragraph{Encoding} Format \textsf{LSU\_ALMAXSB.O} (Section~\ref{sec:format-LSU-ALMAXSB.O}), \verb|lsucode5 = 001|\\

\bitfields{ 1/P, 2/00, 2/-- --, 27/offset27, 1/1, 2/01, 3/111, 2/coherency, 6/registerT, 2/11, 3/001, 1/--, 5/01001, 1/0, 6/registerZ }


\paragraph{Syntax} \texttt{almaxh}\emph{coherency} \emph{offset27}[\emph{registerZ}] = \emph{registerT}

\paragraph{Description} The effective address is computed by adding the \emph{offset27} to the \emph{registerZ}.
This instruction implements atomic load-MAX-store on signed half word. The half word at effective address is MAXed with the \emph{registerT}, and its previous value is returned into the \emph{registerT}. The effective address must be half word aligned (multiple of 2).


\paragraph{Modifier(s)}
\noindent\begin{itemize}
\item coherency (cf. coherency in section \ref{par:modifier-coherency} on page \pageref{par:modifier-coherency})
\end{itemize}

\paragraph{Execution} {Reservation table \textsf{LSU2\_MEMW\_AUXR\_AUXW.X} (Section~\ref{sec:reservation-LSU2-MEMW-AUXR-AUXW.X})}
~
{\begin{lstlisting}
stage ID:
  new offset = _SX_27(offset27);
  new coherency = _ZX_2(coherency);
stage RR:
  new base = GPR[registerZ];
  new address = base + offset;
stage E1:
  new argument2 = GPR[registerT];
  new result2 = _ZX_16(MEM.atomic_max(address, 0x3, coherency << 32, argument2.16[0], registerT));
stage SR:
  GPR[registerT] = result2;
\end{lstlisting}}
\newpage

\paragraph{Encoding} Format \textsf{LSU\_ALMAXSB.D} (Section~\ref{sec:format-LSU-ALMAXSB.D}), \verb|lsucode5 = 001|\\

\bitfields{ 1/P, 2/00, 2/-- --, 27/extend27, 1/1, 2/00, 2/-- --, 27/offset27, 1/1, 2/01, 3/111, 2/coherency, 6/registerT, 2/11, 3/001, 1/--, 5/01001, 1/0, 6/registerZ }


\paragraph{Syntax} \texttt{almaxh}\emph{coherency} \emph{extend27\_\discretionary{}{}{}offset27}[\emph{registerZ}] = \emph{registerT}

\paragraph{Description} The effective address is computed by adding the \emph{extend27\_\discretionary{}{}{}offset27} to the \emph{registerZ}.
This instruction implements atomic load-MAX-store on signed half word. The half word at effective address is MAXed with the \emph{registerT}, and its previous value is returned into the \emph{registerT}. The effective address must be half word aligned (multiple of 2).


\paragraph{Modifier(s)}
\noindent\begin{itemize}
\item coherency (cf. coherency in section \ref{par:modifier-coherency} on page \pageref{par:modifier-coherency})
\end{itemize}

\paragraph{Execution} {Reservation table \textsf{LSU2\_MEMW\_AUXR\_AUXW.Y} (Section~\ref{sec:reservation-LSU2-MEMW-AUXR-AUXW.Y})}
~
{\begin{lstlisting}
stage ID:
  new offset = _SX_54(extend27_offset27);
  new coherency = _ZX_2(coherency);
stage RR:
  new base = GPR[registerZ];
  new address = base + offset;
stage E1:
  new argument2 = GPR[registerT];
  new result2 = _ZX_16(MEM.atomic_max(address, 0x3, coherency << 32, argument2.16[0], registerT));
stage SR:
  GPR[registerT] = result2;
\end{lstlisting}}
\newpage
\subsubsection[{\small ALMAXUB: Atomic Load and MAX Unsigned Byte}]{ALMAXUB\hfill Atomic Load and MAX Unsigned Byte}
\label{instr:ALMAXUB}\index{almaxub}
 
\paragraph{Encoding} Format \textsf{LSU\_ALMAXUSB} (Section~\ref{sec:format-LSU-ALMAXUSB}), \verb|lsucode5 = 000|\\

\bitfields{ 1/P, 2/01, 3/111, 2/coherency, 6/registerT, 2/11, 3/000, 1/--, 5/01011, 1/0, 6/registerZ }


\paragraph{Syntax} \texttt{almaxub}\emph{coherency} [\emph{registerZ}] = \emph{registerT}

\paragraph{Description} The effective address is the value of the \emph{registerZ}.
This instruction implements atomic load-MAXU-store on signed byte. The byte at effective address is MAXUed with the \emph{registerT}, and its previous value is returned into the \emph{registerT}.


\paragraph{Modifier(s)}
\noindent\begin{itemize}
\item coherency (cf. coherency in section \ref{par:modifier-coherency} on page \pageref{par:modifier-coherency})
\end{itemize}

\paragraph{Execution} {Reservation table \textsf{LSU2\_MEMW\_AUXR\_AUXW} (Section~\ref{sec:reservation-LSU2-MEMW-AUXR-AUXW})}
~
{\begin{lstlisting}
stage ID:
  new coherency = _ZX_2(coherency);
stage RR:
  new address = GPR[registerZ];
stage E1:
  new argument2 = GPR[registerT];
  new result2 = _ZX_8(MEM.atomic_maxu(address, 0x1, coherency << 32, argument2.8[0], registerT));
stage SR:
  GPR[registerT] = result2;
\end{lstlisting}}
\newpage

\paragraph{Encoding} Format \textsf{LSU\_ALMAXUSB.O} (Section~\ref{sec:format-LSU-ALMAXUSB.O}), \verb|lsucode5 = 000|\\

\bitfields{ 1/P, 2/00, 2/-- --, 27/offset27, 1/1, 2/01, 3/111, 2/coherency, 6/registerT, 2/11, 3/000, 1/--, 5/01011, 1/0, 6/registerZ }


\paragraph{Syntax} \texttt{almaxub}\emph{coherency} \emph{offset27}[\emph{registerZ}] = \emph{registerT}

\paragraph{Description} The effective address is computed by adding the \emph{offset27} to the \emph{registerZ}.
This instruction implements atomic load-MAXU-store on signed byte. The byte at effective address is MAXUed with the \emph{registerT}, and its previous value is returned into the \emph{registerT}.


\paragraph{Modifier(s)}
\noindent\begin{itemize}
\item coherency (cf. coherency in section \ref{par:modifier-coherency} on page \pageref{par:modifier-coherency})
\end{itemize}

\paragraph{Execution} {Reservation table \textsf{LSU2\_MEMW\_AUXR\_AUXW.X} (Section~\ref{sec:reservation-LSU2-MEMW-AUXR-AUXW.X})}
~
{\begin{lstlisting}
stage ID:
  new offset = _SX_27(offset27);
  new coherency = _ZX_2(coherency);
stage RR:
  new base = GPR[registerZ];
  new address = base + offset;
stage E1:
  new argument2 = GPR[registerT];
  new result2 = _ZX_8(MEM.atomic_maxu(address, 0x1, coherency << 32, argument2.8[0], registerT));
stage SR:
  GPR[registerT] = result2;
\end{lstlisting}}
\newpage

\paragraph{Encoding} Format \textsf{LSU\_ALMAXUSB.D} (Section~\ref{sec:format-LSU-ALMAXUSB.D}), \verb|lsucode5 = 000|\\

\bitfields{ 1/P, 2/00, 2/-- --, 27/extend27, 1/1, 2/00, 2/-- --, 27/offset27, 1/1, 2/01, 3/111, 2/coherency, 6/registerT, 2/11, 3/000, 1/--, 5/01011, 1/0, 6/registerZ }


\paragraph{Syntax} \texttt{almaxub}\emph{coherency} \emph{extend27\_\discretionary{}{}{}offset27}[\emph{registerZ}] = \emph{registerT}

\paragraph{Description} The effective address is computed by adding the \emph{extend27\_\discretionary{}{}{}offset27} to the \emph{registerZ}.
This instruction implements atomic load-MAXU-store on signed byte. The byte at effective address is MAXUed with the \emph{registerT}, and its previous value is returned into the \emph{registerT}.


\paragraph{Modifier(s)}
\noindent\begin{itemize}
\item coherency (cf. coherency in section \ref{par:modifier-coherency} on page \pageref{par:modifier-coherency})
\end{itemize}

\paragraph{Execution} {Reservation table \textsf{LSU2\_MEMW\_AUXR\_AUXW.Y} (Section~\ref{sec:reservation-LSU2-MEMW-AUXR-AUXW.Y})}
~
{\begin{lstlisting}
stage ID:
  new offset = _SX_54(extend27_offset27);
  new coherency = _ZX_2(coherency);
stage RR:
  new base = GPR[registerZ];
  new address = base + offset;
stage E1:
  new argument2 = GPR[registerT];
  new result2 = _ZX_8(MEM.atomic_maxu(address, 0x1, coherency << 32, argument2.8[0], registerT));
stage SR:
  GPR[registerT] = result2;
\end{lstlisting}}
\newpage
\subsubsection[{\small ALMAXUD: Atomic Load and MAX Unsigned Double Word}]{ALMAXUD\hfill Atomic Load and MAX Unsigned Double Word}
\label{instr:ALMAXUD}\index{almaxud}
 
\paragraph{Encoding} Format \textsf{LSU\_ALMAXUSB} (Section~\ref{sec:format-LSU-ALMAXUSB}), \verb|lsucode5 = 011|\\

\bitfields{ 1/P, 2/01, 3/111, 2/coherency, 6/registerT, 2/11, 3/011, 1/--, 5/01011, 1/0, 6/registerZ }


\paragraph{Syntax} \texttt{almaxud}\emph{coherency} [\emph{registerZ}] = \emph{registerT}

\paragraph{Description} The effective address is the value of the \emph{registerZ}.
This instruction implements atomic load-MAXU-store on signed double word. The double word at effective address is MAXUed with the \emph{registerT}, and its previous value is returned into the \emph{registerT}. The effective address must be double word aligned (multiple of 8).


\paragraph{Modifier(s)}
\noindent\begin{itemize}
\item coherency (cf. coherency in section \ref{par:modifier-coherency} on page \pageref{par:modifier-coherency})
\end{itemize}

\paragraph{Execution} {Reservation table \textsf{LSU2\_MEMW\_AUXR\_AUXW} (Section~\ref{sec:reservation-LSU2-MEMW-AUXR-AUXW})}
~
{\begin{lstlisting}
stage ID:
  new coherency = _ZX_2(coherency);
stage RR:
  new address = GPR[registerZ];
stage E1:
  new argument2 = GPR[registerT];
  new result2 = MEM.atomic_maxu(address, 0xFF, coherency << 32, argument2.64[0], registerT);
stage SR:
  GPR[registerT] = result2;
\end{lstlisting}}
\newpage

\paragraph{Encoding} Format \textsf{LSU\_ALMAXUSB.O} (Section~\ref{sec:format-LSU-ALMAXUSB.O}), \verb|lsucode5 = 011|\\

\bitfields{ 1/P, 2/00, 2/-- --, 27/offset27, 1/1, 2/01, 3/111, 2/coherency, 6/registerT, 2/11, 3/011, 1/--, 5/01011, 1/0, 6/registerZ }


\paragraph{Syntax} \texttt{almaxud}\emph{coherency} \emph{offset27}[\emph{registerZ}] = \emph{registerT}

\paragraph{Description} The effective address is computed by adding the \emph{offset27} to the \emph{registerZ}.
This instruction implements atomic load-MAXU-store on signed double word. The double word at effective address is MAXUed with the \emph{registerT}, and its previous value is returned into the \emph{registerT}. The effective address must be double word aligned (multiple of 8).


\paragraph{Modifier(s)}
\noindent\begin{itemize}
\item coherency (cf. coherency in section \ref{par:modifier-coherency} on page \pageref{par:modifier-coherency})
\end{itemize}

\paragraph{Execution} {Reservation table \textsf{LSU2\_MEMW\_AUXR\_AUXW.X} (Section~\ref{sec:reservation-LSU2-MEMW-AUXR-AUXW.X})}
~
{\begin{lstlisting}
stage ID:
  new offset = _SX_27(offset27);
  new coherency = _ZX_2(coherency);
stage RR:
  new base = GPR[registerZ];
  new address = base + offset;
stage E1:
  new argument2 = GPR[registerT];
  new result2 = MEM.atomic_maxu(address, 0xFF, coherency << 32, argument2.64[0], registerT);
stage SR:
  GPR[registerT] = result2;
\end{lstlisting}}
\newpage

\paragraph{Encoding} Format \textsf{LSU\_ALMAXUSB.D} (Section~\ref{sec:format-LSU-ALMAXUSB.D}), \verb|lsucode5 = 011|\\

\bitfields{ 1/P, 2/00, 2/-- --, 27/extend27, 1/1, 2/00, 2/-- --, 27/offset27, 1/1, 2/01, 3/111, 2/coherency, 6/registerT, 2/11, 3/011, 1/--, 5/01011, 1/0, 6/registerZ }


\paragraph{Syntax} \texttt{almaxud}\emph{coherency} \emph{extend27\_\discretionary{}{}{}offset27}[\emph{registerZ}] = \emph{registerT}

\paragraph{Description} The effective address is computed by adding the \emph{extend27\_\discretionary{}{}{}offset27} to the \emph{registerZ}.
This instruction implements atomic load-MAXU-store on signed double word. The double word at effective address is MAXUed with the \emph{registerT}, and its previous value is returned into the \emph{registerT}. The effective address must be double word aligned (multiple of 8).


\paragraph{Modifier(s)}
\noindent\begin{itemize}
\item coherency (cf. coherency in section \ref{par:modifier-coherency} on page \pageref{par:modifier-coherency})
\end{itemize}

\paragraph{Execution} {Reservation table \textsf{LSU2\_MEMW\_AUXR\_AUXW.Y} (Section~\ref{sec:reservation-LSU2-MEMW-AUXR-AUXW.Y})}
~
{\begin{lstlisting}
stage ID:
  new offset = _SX_54(extend27_offset27);
  new coherency = _ZX_2(coherency);
stage RR:
  new base = GPR[registerZ];
  new address = base + offset;
stage E1:
  new argument2 = GPR[registerT];
  new result2 = MEM.atomic_maxu(address, 0xFF, coherency << 32, argument2.64[0], registerT);
stage SR:
  GPR[registerT] = result2;
\end{lstlisting}}
\newpage
\subsubsection[{\small ALMAXUH: Atomic Load and MAX Unsigned Half Word}]{ALMAXUH\hfill Atomic Load and MAX Unsigned Half Word}
\label{instr:ALMAXUH}\index{almaxuh}
 
\paragraph{Encoding} Format \textsf{LSU\_ALMAXUSB} (Section~\ref{sec:format-LSU-ALMAXUSB}), \verb|lsucode5 = 001|\\

\bitfields{ 1/P, 2/01, 3/111, 2/coherency, 6/registerT, 2/11, 3/001, 1/--, 5/01011, 1/0, 6/registerZ }


\paragraph{Syntax} \texttt{almaxuh}\emph{coherency} [\emph{registerZ}] = \emph{registerT}

\paragraph{Description} The effective address is the value of the \emph{registerZ}.
This instruction implements atomic load-MAXU-store on signed half word. The half word at effective address is MAXUed with the \emph{registerT}, and its previous value is returned into the \emph{registerT}. The effective address must be half word aligned (multiple of 2).


\paragraph{Modifier(s)}
\noindent\begin{itemize}
\item coherency (cf. coherency in section \ref{par:modifier-coherency} on page \pageref{par:modifier-coherency})
\end{itemize}

\paragraph{Execution} {Reservation table \textsf{LSU2\_MEMW\_AUXR\_AUXW} (Section~\ref{sec:reservation-LSU2-MEMW-AUXR-AUXW})}
~
{\begin{lstlisting}
stage ID:
  new coherency = _ZX_2(coherency);
stage RR:
  new address = GPR[registerZ];
stage E1:
  new argument2 = GPR[registerT];
  new result2 = _ZX_16(MEM.atomic_maxu(address, 0x3, coherency << 32, argument2.16[0], registerT));
stage SR:
  GPR[registerT] = result2;
\end{lstlisting}}
\newpage

\paragraph{Encoding} Format \textsf{LSU\_ALMAXUSB.O} (Section~\ref{sec:format-LSU-ALMAXUSB.O}), \verb|lsucode5 = 001|\\

\bitfields{ 1/P, 2/00, 2/-- --, 27/offset27, 1/1, 2/01, 3/111, 2/coherency, 6/registerT, 2/11, 3/001, 1/--, 5/01011, 1/0, 6/registerZ }


\paragraph{Syntax} \texttt{almaxuh}\emph{coherency} \emph{offset27}[\emph{registerZ}] = \emph{registerT}

\paragraph{Description} The effective address is computed by adding the \emph{offset27} to the \emph{registerZ}.
This instruction implements atomic load-MAXU-store on signed half word. The half word at effective address is MAXUed with the \emph{registerT}, and its previous value is returned into the \emph{registerT}. The effective address must be half word aligned (multiple of 2).


\paragraph{Modifier(s)}
\noindent\begin{itemize}
\item coherency (cf. coherency in section \ref{par:modifier-coherency} on page \pageref{par:modifier-coherency})
\end{itemize}

\paragraph{Execution} {Reservation table \textsf{LSU2\_MEMW\_AUXR\_AUXW.X} (Section~\ref{sec:reservation-LSU2-MEMW-AUXR-AUXW.X})}
~
{\begin{lstlisting}
stage ID:
  new offset = _SX_27(offset27);
  new coherency = _ZX_2(coherency);
stage RR:
  new base = GPR[registerZ];
  new address = base + offset;
stage E1:
  new argument2 = GPR[registerT];
  new result2 = _ZX_16(MEM.atomic_maxu(address, 0x3, coherency << 32, argument2.16[0], registerT));
stage SR:
  GPR[registerT] = result2;
\end{lstlisting}}
\newpage

\paragraph{Encoding} Format \textsf{LSU\_ALMAXUSB.D} (Section~\ref{sec:format-LSU-ALMAXUSB.D}), \verb|lsucode5 = 001|\\

\bitfields{ 1/P, 2/00, 2/-- --, 27/extend27, 1/1, 2/00, 2/-- --, 27/offset27, 1/1, 2/01, 3/111, 2/coherency, 6/registerT, 2/11, 3/001, 1/--, 5/01011, 1/0, 6/registerZ }


\paragraph{Syntax} \texttt{almaxuh}\emph{coherency} \emph{extend27\_\discretionary{}{}{}offset27}[\emph{registerZ}] = \emph{registerT}

\paragraph{Description} The effective address is computed by adding the \emph{extend27\_\discretionary{}{}{}offset27} to the \emph{registerZ}.
This instruction implements atomic load-MAXU-store on signed half word. The half word at effective address is MAXUed with the \emph{registerT}, and its previous value is returned into the \emph{registerT}. The effective address must be half word aligned (multiple of 2).


\paragraph{Modifier(s)}
\noindent\begin{itemize}
\item coherency (cf. coherency in section \ref{par:modifier-coherency} on page \pageref{par:modifier-coherency})
\end{itemize}

\paragraph{Execution} {Reservation table \textsf{LSU2\_MEMW\_AUXR\_AUXW.Y} (Section~\ref{sec:reservation-LSU2-MEMW-AUXR-AUXW.Y})}
~
{\begin{lstlisting}
stage ID:
  new offset = _SX_54(extend27_offset27);
  new coherency = _ZX_2(coherency);
stage RR:
  new base = GPR[registerZ];
  new address = base + offset;
stage E1:
  new argument2 = GPR[registerT];
  new result2 = _ZX_16(MEM.atomic_maxu(address, 0x3, coherency << 32, argument2.16[0], registerT));
stage SR:
  GPR[registerT] = result2;
\end{lstlisting}}
\newpage
\subsubsection[{\small ALMAXUW: Atomic Load and MAX Unsigned Word}]{ALMAXUW\hfill Atomic Load and MAX Unsigned Word}
\label{instr:ALMAXUW}\index{almaxuw}
 
\paragraph{Encoding} Format \textsf{LSU\_ALMAXUSB} (Section~\ref{sec:format-LSU-ALMAXUSB}), \verb|lsucode5 = 010|\\

\bitfields{ 1/P, 2/01, 3/111, 2/coherency, 6/registerT, 2/11, 3/010, 1/--, 5/01011, 1/0, 6/registerZ }


\paragraph{Syntax} \texttt{almaxuw}\emph{coherency} [\emph{registerZ}] = \emph{registerT}

\paragraph{Description} The effective address is the value of the \emph{registerZ}.
This instruction implements atomic load-MAXU-store on signed word. The word at effective address is MAXUed with the \emph{registerT}, and its previous value is returned into the \emph{registerT}. The effective address must be word aligned (multiple of 4).


\paragraph{Modifier(s)}
\noindent\begin{itemize}
\item coherency (cf. coherency in section \ref{par:modifier-coherency} on page \pageref{par:modifier-coherency})
\end{itemize}

\paragraph{Execution} {Reservation table \textsf{LSU2\_MEMW\_AUXR\_AUXW} (Section~\ref{sec:reservation-LSU2-MEMW-AUXR-AUXW})}
~
{\begin{lstlisting}
stage ID:
  new coherency = _ZX_2(coherency);
stage RR:
  new address = GPR[registerZ];
stage E1:
  new argument2 = GPR[registerT];
  new result2 = _ZX_32(MEM.atomic_maxu(address, 0xF, coherency << 32, argument2.32[0], registerT));
stage SR:
  GPR[registerT] = result2;
\end{lstlisting}}
\newpage

\paragraph{Encoding} Format \textsf{LSU\_ALMAXUSB.O} (Section~\ref{sec:format-LSU-ALMAXUSB.O}), \verb|lsucode5 = 010|\\

\bitfields{ 1/P, 2/00, 2/-- --, 27/offset27, 1/1, 2/01, 3/111, 2/coherency, 6/registerT, 2/11, 3/010, 1/--, 5/01011, 1/0, 6/registerZ }


\paragraph{Syntax} \texttt{almaxuw}\emph{coherency} \emph{offset27}[\emph{registerZ}] = \emph{registerT}

\paragraph{Description} The effective address is computed by adding the \emph{offset27} to the \emph{registerZ}.
This instruction implements atomic load-MAXU-store on signed word. The word at effective address is MAXUed with the \emph{registerT}, and its previous value is returned into the \emph{registerT}. The effective address must be word aligned (multiple of 4).


\paragraph{Modifier(s)}
\noindent\begin{itemize}
\item coherency (cf. coherency in section \ref{par:modifier-coherency} on page \pageref{par:modifier-coherency})
\end{itemize}

\paragraph{Execution} {Reservation table \textsf{LSU2\_MEMW\_AUXR\_AUXW.X} (Section~\ref{sec:reservation-LSU2-MEMW-AUXR-AUXW.X})}
~
{\begin{lstlisting}
stage ID:
  new offset = _SX_27(offset27);
  new coherency = _ZX_2(coherency);
stage RR:
  new base = GPR[registerZ];
  new address = base + offset;
stage E1:
  new argument2 = GPR[registerT];
  new result2 = _ZX_32(MEM.atomic_maxu(address, 0xF, coherency << 32, argument2.32[0], registerT));
stage SR:
  GPR[registerT] = result2;
\end{lstlisting}}
\newpage

\paragraph{Encoding} Format \textsf{LSU\_ALMAXUSB.D} (Section~\ref{sec:format-LSU-ALMAXUSB.D}), \verb|lsucode5 = 010|\\

\bitfields{ 1/P, 2/00, 2/-- --, 27/extend27, 1/1, 2/00, 2/-- --, 27/offset27, 1/1, 2/01, 3/111, 2/coherency, 6/registerT, 2/11, 3/010, 1/--, 5/01011, 1/0, 6/registerZ }


\paragraph{Syntax} \texttt{almaxuw}\emph{coherency} \emph{extend27\_\discretionary{}{}{}offset27}[\emph{registerZ}] = \emph{registerT}

\paragraph{Description} The effective address is computed by adding the \emph{extend27\_\discretionary{}{}{}offset27} to the \emph{registerZ}.
This instruction implements atomic load-MAXU-store on signed word. The word at effective address is MAXUed with the \emph{registerT}, and its previous value is returned into the \emph{registerT}. The effective address must be word aligned (multiple of 4).


\paragraph{Modifier(s)}
\noindent\begin{itemize}
\item coherency (cf. coherency in section \ref{par:modifier-coherency} on page \pageref{par:modifier-coherency})
\end{itemize}

\paragraph{Execution} {Reservation table \textsf{LSU2\_MEMW\_AUXR\_AUXW.Y} (Section~\ref{sec:reservation-LSU2-MEMW-AUXR-AUXW.Y})}
~
{\begin{lstlisting}
stage ID:
  new offset = _SX_54(extend27_offset27);
  new coherency = _ZX_2(coherency);
stage RR:
  new base = GPR[registerZ];
  new address = base + offset;
stage E1:
  new argument2 = GPR[registerT];
  new result2 = _ZX_32(MEM.atomic_maxu(address, 0xF, coherency << 32, argument2.32[0], registerT));
stage SR:
  GPR[registerT] = result2;
\end{lstlisting}}
\newpage
\subsubsection[{\small ALMAXW: Atomic Load and MAX Word}]{ALMAXW\hfill Atomic Load and MAX Word}
\label{instr:ALMAXW}\index{almaxw}
 
\paragraph{Encoding} Format \textsf{LSU\_ALMAXSB} (Section~\ref{sec:format-LSU-ALMAXSB}), \verb|lsucode5 = 010|\\

\bitfields{ 1/P, 2/01, 3/111, 2/coherency, 6/registerT, 2/11, 3/010, 1/--, 5/01001, 1/0, 6/registerZ }


\paragraph{Syntax} \texttt{almaxw}\emph{coherency} [\emph{registerZ}] = \emph{registerT}

\paragraph{Description} The effective address is the value of the \emph{registerZ}.
This instruction implements atomic load-MAX-store on signed word. The word at effective address is MAXed with the \emph{registerT}, and its previous value is returned into the \emph{registerT}. The effective address must be word aligned (multiple of 4).


\paragraph{Modifier(s)}
\noindent\begin{itemize}
\item coherency (cf. coherency in section \ref{par:modifier-coherency} on page \pageref{par:modifier-coherency})
\end{itemize}

\paragraph{Execution} {Reservation table \textsf{LSU2\_MEMW\_AUXR\_AUXW} (Section~\ref{sec:reservation-LSU2-MEMW-AUXR-AUXW})}
~
{\begin{lstlisting}
stage ID:
  new coherency = _ZX_2(coherency);
stage RR:
  new address = GPR[registerZ];
stage E1:
  new argument2 = GPR[registerT];
  new result2 = _ZX_32(MEM.atomic_max(address, 0xF, coherency << 32, argument2.32[0], registerT));
stage SR:
  GPR[registerT] = result2;
\end{lstlisting}}
\newpage

\paragraph{Encoding} Format \textsf{LSU\_ALMAXSB.O} (Section~\ref{sec:format-LSU-ALMAXSB.O}), \verb|lsucode5 = 010|\\

\bitfields{ 1/P, 2/00, 2/-- --, 27/offset27, 1/1, 2/01, 3/111, 2/coherency, 6/registerT, 2/11, 3/010, 1/--, 5/01001, 1/0, 6/registerZ }


\paragraph{Syntax} \texttt{almaxw}\emph{coherency} \emph{offset27}[\emph{registerZ}] = \emph{registerT}

\paragraph{Description} The effective address is computed by adding the \emph{offset27} to the \emph{registerZ}.
This instruction implements atomic load-MAX-store on signed word. The word at effective address is MAXed with the \emph{registerT}, and its previous value is returned into the \emph{registerT}. The effective address must be word aligned (multiple of 4).


\paragraph{Modifier(s)}
\noindent\begin{itemize}
\item coherency (cf. coherency in section \ref{par:modifier-coherency} on page \pageref{par:modifier-coherency})
\end{itemize}

\paragraph{Execution} {Reservation table \textsf{LSU2\_MEMW\_AUXR\_AUXW.X} (Section~\ref{sec:reservation-LSU2-MEMW-AUXR-AUXW.X})}
~
{\begin{lstlisting}
stage ID:
  new offset = _SX_27(offset27);
  new coherency = _ZX_2(coherency);
stage RR:
  new base = GPR[registerZ];
  new address = base + offset;
stage E1:
  new argument2 = GPR[registerT];
  new result2 = _ZX_32(MEM.atomic_max(address, 0xF, coherency << 32, argument2.32[0], registerT));
stage SR:
  GPR[registerT] = result2;
\end{lstlisting}}
\newpage

\paragraph{Encoding} Format \textsf{LSU\_ALMAXSB.D} (Section~\ref{sec:format-LSU-ALMAXSB.D}), \verb|lsucode5 = 010|\\

\bitfields{ 1/P, 2/00, 2/-- --, 27/extend27, 1/1, 2/00, 2/-- --, 27/offset27, 1/1, 2/01, 3/111, 2/coherency, 6/registerT, 2/11, 3/010, 1/--, 5/01001, 1/0, 6/registerZ }


\paragraph{Syntax} \texttt{almaxw}\emph{coherency} \emph{extend27\_\discretionary{}{}{}offset27}[\emph{registerZ}] = \emph{registerT}

\paragraph{Description} The effective address is computed by adding the \emph{extend27\_\discretionary{}{}{}offset27} to the \emph{registerZ}.
This instruction implements atomic load-MAX-store on signed word. The word at effective address is MAXed with the \emph{registerT}, and its previous value is returned into the \emph{registerT}. The effective address must be word aligned (multiple of 4).


\paragraph{Modifier(s)}
\noindent\begin{itemize}
\item coherency (cf. coherency in section \ref{par:modifier-coherency} on page \pageref{par:modifier-coherency})
\end{itemize}

\paragraph{Execution} {Reservation table \textsf{LSU2\_MEMW\_AUXR\_AUXW.Y} (Section~\ref{sec:reservation-LSU2-MEMW-AUXR-AUXW.Y})}
~
{\begin{lstlisting}
stage ID:
  new offset = _SX_54(extend27_offset27);
  new coherency = _ZX_2(coherency);
stage RR:
  new base = GPR[registerZ];
  new address = base + offset;
stage E1:
  new argument2 = GPR[registerT];
  new result2 = _ZX_32(MEM.atomic_max(address, 0xF, coherency << 32, argument2.32[0], registerT));
stage SR:
  GPR[registerT] = result2;
\end{lstlisting}}
\newpage
\subsubsection[{\small ALMINB: Atomic Load and MIN Byte}]{ALMINB\hfill Atomic Load and MIN Byte}
\label{instr:ALMINB}\index{alminb}
 
\paragraph{Encoding} Format \textsf{LSU\_ALMINSB} (Section~\ref{sec:format-LSU-ALMINSB}), \verb|lsucode5 = 000|\\

\bitfields{ 1/P, 2/01, 3/111, 2/coherency, 6/registerT, 2/11, 3/000, 1/--, 5/01000, 1/0, 6/registerZ }


\paragraph{Syntax} \texttt{alminb}\emph{coherency} [\emph{registerZ}] = \emph{registerT}

\paragraph{Description} The effective address is the value of the \emph{registerZ}.
This instruction implements atomic load-MIN-store on signed byte. The byte at effective address is MINed with the \emph{registerT}, and its previous value is returned into the \emph{registerT}.


\paragraph{Modifier(s)}
\noindent\begin{itemize}
\item coherency (cf. coherency in section \ref{par:modifier-coherency} on page \pageref{par:modifier-coherency})
\end{itemize}

\paragraph{Execution} {Reservation table \textsf{LSU2\_MEMW\_AUXR\_AUXW} (Section~\ref{sec:reservation-LSU2-MEMW-AUXR-AUXW})}
~
{\begin{lstlisting}
stage ID:
  new coherency = _ZX_2(coherency);
stage RR:
  new address = GPR[registerZ];
stage E1:
  new argument2 = GPR[registerT];
  new result2 = _ZX_8(MEM.atomic_min(address, 0x1, coherency << 32, argument2.8[0], registerT));
stage SR:
  GPR[registerT] = result2;
\end{lstlisting}}
\newpage

\paragraph{Encoding} Format \textsf{LSU\_ALMINSB.O} (Section~\ref{sec:format-LSU-ALMINSB.O}), \verb|lsucode5 = 000|\\

\bitfields{ 1/P, 2/00, 2/-- --, 27/offset27, 1/1, 2/01, 3/111, 2/coherency, 6/registerT, 2/11, 3/000, 1/--, 5/01000, 1/0, 6/registerZ }


\paragraph{Syntax} \texttt{alminb}\emph{coherency} \emph{offset27}[\emph{registerZ}] = \emph{registerT}

\paragraph{Description} The effective address is computed by adding the \emph{offset27} to the \emph{registerZ}.
This instruction implements atomic load-MIN-store on signed byte. The byte at effective address is MINed with the \emph{registerT}, and its previous value is returned into the \emph{registerT}.


\paragraph{Modifier(s)}
\noindent\begin{itemize}
\item coherency (cf. coherency in section \ref{par:modifier-coherency} on page \pageref{par:modifier-coherency})
\end{itemize}

\paragraph{Execution} {Reservation table \textsf{LSU2\_MEMW\_AUXR\_AUXW.X} (Section~\ref{sec:reservation-LSU2-MEMW-AUXR-AUXW.X})}
~
{\begin{lstlisting}
stage ID:
  new offset = _SX_27(offset27);
  new coherency = _ZX_2(coherency);
stage RR:
  new base = GPR[registerZ];
  new address = base + offset;
stage E1:
  new argument2 = GPR[registerT];
  new result2 = _ZX_8(MEM.atomic_min(address, 0x1, coherency << 32, argument2.8[0], registerT));
stage SR:
  GPR[registerT] = result2;
\end{lstlisting}}
\newpage

\paragraph{Encoding} Format \textsf{LSU\_ALMINSB.D} (Section~\ref{sec:format-LSU-ALMINSB.D}), \verb|lsucode5 = 000|\\

\bitfields{ 1/P, 2/00, 2/-- --, 27/extend27, 1/1, 2/00, 2/-- --, 27/offset27, 1/1, 2/01, 3/111, 2/coherency, 6/registerT, 2/11, 3/000, 1/--, 5/01000, 1/0, 6/registerZ }


\paragraph{Syntax} \texttt{alminb}\emph{coherency} \emph{extend27\_\discretionary{}{}{}offset27}[\emph{registerZ}] = \emph{registerT}

\paragraph{Description} The effective address is computed by adding the \emph{extend27\_\discretionary{}{}{}offset27} to the \emph{registerZ}.
This instruction implements atomic load-MIN-store on signed byte. The byte at effective address is MINed with the \emph{registerT}, and its previous value is returned into the \emph{registerT}.


\paragraph{Modifier(s)}
\noindent\begin{itemize}
\item coherency (cf. coherency in section \ref{par:modifier-coherency} on page \pageref{par:modifier-coherency})
\end{itemize}

\paragraph{Execution} {Reservation table \textsf{LSU2\_MEMW\_AUXR\_AUXW.Y} (Section~\ref{sec:reservation-LSU2-MEMW-AUXR-AUXW.Y})}
~
{\begin{lstlisting}
stage ID:
  new offset = _SX_54(extend27_offset27);
  new coherency = _ZX_2(coherency);
stage RR:
  new base = GPR[registerZ];
  new address = base + offset;
stage E1:
  new argument2 = GPR[registerT];
  new result2 = _ZX_8(MEM.atomic_min(address, 0x1, coherency << 32, argument2.8[0], registerT));
stage SR:
  GPR[registerT] = result2;
\end{lstlisting}}
\newpage
\subsubsection[{\small ALMIND: Atomic Load and MIN Double Word}]{ALMIND\hfill Atomic Load and MIN Double Word}
\label{instr:ALMIND}\index{almind}
 
\paragraph{Encoding} Format \textsf{LSU\_ALMINSB} (Section~\ref{sec:format-LSU-ALMINSB}), \verb|lsucode5 = 011|\\

\bitfields{ 1/P, 2/01, 3/111, 2/coherency, 6/registerT, 2/11, 3/011, 1/--, 5/01000, 1/0, 6/registerZ }


\paragraph{Syntax} \texttt{almind}\emph{coherency} [\emph{registerZ}] = \emph{registerT}

\paragraph{Description} The effective address is the value of the \emph{registerZ}.
This instruction implements atomic load-MIN-store on signed double word. The double word at effective address is MINed with the \emph{registerT}, and its previous value is returned into the \emph{registerT}. The effective address must be double word aligned (multiple of 8).


\paragraph{Modifier(s)}
\noindent\begin{itemize}
\item coherency (cf. coherency in section \ref{par:modifier-coherency} on page \pageref{par:modifier-coherency})
\end{itemize}

\paragraph{Execution} {Reservation table \textsf{LSU2\_MEMW\_AUXR\_AUXW} (Section~\ref{sec:reservation-LSU2-MEMW-AUXR-AUXW})}
~
{\begin{lstlisting}
stage ID:
  new coherency = _ZX_2(coherency);
stage RR:
  new address = GPR[registerZ];
stage E1:
  new argument2 = GPR[registerT];
  new result2 = MEM.atomic_min(address, 0xFF, coherency << 32, argument2.64[0], registerT);
stage SR:
  GPR[registerT] = result2;
\end{lstlisting}}
\newpage

\paragraph{Encoding} Format \textsf{LSU\_ALMINSB.O} (Section~\ref{sec:format-LSU-ALMINSB.O}), \verb|lsucode5 = 011|\\

\bitfields{ 1/P, 2/00, 2/-- --, 27/offset27, 1/1, 2/01, 3/111, 2/coherency, 6/registerT, 2/11, 3/011, 1/--, 5/01000, 1/0, 6/registerZ }


\paragraph{Syntax} \texttt{almind}\emph{coherency} \emph{offset27}[\emph{registerZ}] = \emph{registerT}

\paragraph{Description} The effective address is computed by adding the \emph{offset27} to the \emph{registerZ}.
This instruction implements atomic load-MIN-store on signed double word. The double word at effective address is MINed with the \emph{registerT}, and its previous value is returned into the \emph{registerT}. The effective address must be double word aligned (multiple of 8).


\paragraph{Modifier(s)}
\noindent\begin{itemize}
\item coherency (cf. coherency in section \ref{par:modifier-coherency} on page \pageref{par:modifier-coherency})
\end{itemize}

\paragraph{Execution} {Reservation table \textsf{LSU2\_MEMW\_AUXR\_AUXW.X} (Section~\ref{sec:reservation-LSU2-MEMW-AUXR-AUXW.X})}
~
{\begin{lstlisting}
stage ID:
  new offset = _SX_27(offset27);
  new coherency = _ZX_2(coherency);
stage RR:
  new base = GPR[registerZ];
  new address = base + offset;
stage E1:
  new argument2 = GPR[registerT];
  new result2 = MEM.atomic_min(address, 0xFF, coherency << 32, argument2.64[0], registerT);
stage SR:
  GPR[registerT] = result2;
\end{lstlisting}}
\newpage

\paragraph{Encoding} Format \textsf{LSU\_ALMINSB.D} (Section~\ref{sec:format-LSU-ALMINSB.D}), \verb|lsucode5 = 011|\\

\bitfields{ 1/P, 2/00, 2/-- --, 27/extend27, 1/1, 2/00, 2/-- --, 27/offset27, 1/1, 2/01, 3/111, 2/coherency, 6/registerT, 2/11, 3/011, 1/--, 5/01000, 1/0, 6/registerZ }


\paragraph{Syntax} \texttt{almind}\emph{coherency} \emph{extend27\_\discretionary{}{}{}offset27}[\emph{registerZ}] = \emph{registerT}

\paragraph{Description} The effective address is computed by adding the \emph{extend27\_\discretionary{}{}{}offset27} to the \emph{registerZ}.
This instruction implements atomic load-MIN-store on signed double word. The double word at effective address is MINed with the \emph{registerT}, and its previous value is returned into the \emph{registerT}. The effective address must be double word aligned (multiple of 8).


\paragraph{Modifier(s)}
\noindent\begin{itemize}
\item coherency (cf. coherency in section \ref{par:modifier-coherency} on page \pageref{par:modifier-coherency})
\end{itemize}

\paragraph{Execution} {Reservation table \textsf{LSU2\_MEMW\_AUXR\_AUXW.Y} (Section~\ref{sec:reservation-LSU2-MEMW-AUXR-AUXW.Y})}
~
{\begin{lstlisting}
stage ID:
  new offset = _SX_54(extend27_offset27);
  new coherency = _ZX_2(coherency);
stage RR:
  new base = GPR[registerZ];
  new address = base + offset;
stage E1:
  new argument2 = GPR[registerT];
  new result2 = MEM.atomic_min(address, 0xFF, coherency << 32, argument2.64[0], registerT);
stage SR:
  GPR[registerT] = result2;
\end{lstlisting}}
\newpage
\subsubsection[{\small ALMINH: Atomic Load and MIN Half Word}]{ALMINH\hfill Atomic Load and MIN Half Word}
\label{instr:ALMINH}\index{alminh}
 
\paragraph{Encoding} Format \textsf{LSU\_ALMINSB} (Section~\ref{sec:format-LSU-ALMINSB}), \verb|lsucode5 = 001|\\

\bitfields{ 1/P, 2/01, 3/111, 2/coherency, 6/registerT, 2/11, 3/001, 1/--, 5/01000, 1/0, 6/registerZ }


\paragraph{Syntax} \texttt{alminh}\emph{coherency} [\emph{registerZ}] = \emph{registerT}

\paragraph{Description} The effective address is the value of the \emph{registerZ}.
This instruction implements atomic load-MIN-store on signed half word. The half word at effective address is MINed with the \emph{registerT}, and its previous value is returned into the \emph{registerT}. The effective address must be half word aligned (multiple of 2).


\paragraph{Modifier(s)}
\noindent\begin{itemize}
\item coherency (cf. coherency in section \ref{par:modifier-coherency} on page \pageref{par:modifier-coherency})
\end{itemize}

\paragraph{Execution} {Reservation table \textsf{LSU2\_MEMW\_AUXR\_AUXW} (Section~\ref{sec:reservation-LSU2-MEMW-AUXR-AUXW})}
~
{\begin{lstlisting}
stage ID:
  new coherency = _ZX_2(coherency);
stage RR:
  new address = GPR[registerZ];
stage E1:
  new argument2 = GPR[registerT];
  new result2 = _ZX_16(MEM.atomic_min(address, 0x3, coherency << 32, argument2.16[0], registerT));
stage SR:
  GPR[registerT] = result2;
\end{lstlisting}}
\newpage

\paragraph{Encoding} Format \textsf{LSU\_ALMINSB.O} (Section~\ref{sec:format-LSU-ALMINSB.O}), \verb|lsucode5 = 001|\\

\bitfields{ 1/P, 2/00, 2/-- --, 27/offset27, 1/1, 2/01, 3/111, 2/coherency, 6/registerT, 2/11, 3/001, 1/--, 5/01000, 1/0, 6/registerZ }


\paragraph{Syntax} \texttt{alminh}\emph{coherency} \emph{offset27}[\emph{registerZ}] = \emph{registerT}

\paragraph{Description} The effective address is computed by adding the \emph{offset27} to the \emph{registerZ}.
This instruction implements atomic load-MIN-store on signed half word. The half word at effective address is MINed with the \emph{registerT}, and its previous value is returned into the \emph{registerT}. The effective address must be half word aligned (multiple of 2).


\paragraph{Modifier(s)}
\noindent\begin{itemize}
\item coherency (cf. coherency in section \ref{par:modifier-coherency} on page \pageref{par:modifier-coherency})
\end{itemize}

\paragraph{Execution} {Reservation table \textsf{LSU2\_MEMW\_AUXR\_AUXW.X} (Section~\ref{sec:reservation-LSU2-MEMW-AUXR-AUXW.X})}
~
{\begin{lstlisting}
stage ID:
  new offset = _SX_27(offset27);
  new coherency = _ZX_2(coherency);
stage RR:
  new base = GPR[registerZ];
  new address = base + offset;
stage E1:
  new argument2 = GPR[registerT];
  new result2 = _ZX_16(MEM.atomic_min(address, 0x3, coherency << 32, argument2.16[0], registerT));
stage SR:
  GPR[registerT] = result2;
\end{lstlisting}}
\newpage

\paragraph{Encoding} Format \textsf{LSU\_ALMINSB.D} (Section~\ref{sec:format-LSU-ALMINSB.D}), \verb|lsucode5 = 001|\\

\bitfields{ 1/P, 2/00, 2/-- --, 27/extend27, 1/1, 2/00, 2/-- --, 27/offset27, 1/1, 2/01, 3/111, 2/coherency, 6/registerT, 2/11, 3/001, 1/--, 5/01000, 1/0, 6/registerZ }


\paragraph{Syntax} \texttt{alminh}\emph{coherency} \emph{extend27\_\discretionary{}{}{}offset27}[\emph{registerZ}] = \emph{registerT}

\paragraph{Description} The effective address is computed by adding the \emph{extend27\_\discretionary{}{}{}offset27} to the \emph{registerZ}.
This instruction implements atomic load-MIN-store on signed half word. The half word at effective address is MINed with the \emph{registerT}, and its previous value is returned into the \emph{registerT}. The effective address must be half word aligned (multiple of 2).


\paragraph{Modifier(s)}
\noindent\begin{itemize}
\item coherency (cf. coherency in section \ref{par:modifier-coherency} on page \pageref{par:modifier-coherency})
\end{itemize}

\paragraph{Execution} {Reservation table \textsf{LSU2\_MEMW\_AUXR\_AUXW.Y} (Section~\ref{sec:reservation-LSU2-MEMW-AUXR-AUXW.Y})}
~
{\begin{lstlisting}
stage ID:
  new offset = _SX_54(extend27_offset27);
  new coherency = _ZX_2(coherency);
stage RR:
  new base = GPR[registerZ];
  new address = base + offset;
stage E1:
  new argument2 = GPR[registerT];
  new result2 = _ZX_16(MEM.atomic_min(address, 0x3, coherency << 32, argument2.16[0], registerT));
stage SR:
  GPR[registerT] = result2;
\end{lstlisting}}
\newpage
\subsubsection[{\small ALMINUB: Atomic Load and MIN Unsigned Byte}]{ALMINUB\hfill Atomic Load and MIN Unsigned Byte}
\label{instr:ALMINUB}\index{alminub}
 
\paragraph{Encoding} Format \textsf{LSU\_ALMINUSB} (Section~\ref{sec:format-LSU-ALMINUSB}), \verb|lsucode5 = 000|\\

\bitfields{ 1/P, 2/01, 3/111, 2/coherency, 6/registerT, 2/11, 3/000, 1/--, 5/01010, 1/0, 6/registerZ }


\paragraph{Syntax} \texttt{alminub}\emph{coherency} [\emph{registerZ}] = \emph{registerT}

\paragraph{Description} The effective address is the value of the \emph{registerZ}.
This instruction implements atomic load-MINU-store on signed byte. The byte at effective address is MINUed with the \emph{registerT}, and its previous value is returned into the \emph{registerT}.


\paragraph{Modifier(s)}
\noindent\begin{itemize}
\item coherency (cf. coherency in section \ref{par:modifier-coherency} on page \pageref{par:modifier-coherency})
\end{itemize}

\paragraph{Execution} {Reservation table \textsf{LSU2\_MEMW\_AUXR\_AUXW} (Section~\ref{sec:reservation-LSU2-MEMW-AUXR-AUXW})}
~
{\begin{lstlisting}
stage ID:
  new coherency = _ZX_2(coherency);
stage RR:
  new address = GPR[registerZ];
stage E1:
  new argument2 = GPR[registerT];
  new result2 = _ZX_8(MEM.atomic_minu(address, 0x1, coherency << 32, argument2.8[0], registerT));
stage SR:
  GPR[registerT] = result2;
\end{lstlisting}}
\newpage

\paragraph{Encoding} Format \textsf{LSU\_ALMINUSB.O} (Section~\ref{sec:format-LSU-ALMINUSB.O}), \verb|lsucode5 = 000|\\

\bitfields{ 1/P, 2/00, 2/-- --, 27/offset27, 1/1, 2/01, 3/111, 2/coherency, 6/registerT, 2/11, 3/000, 1/--, 5/01010, 1/0, 6/registerZ }


\paragraph{Syntax} \texttt{alminub}\emph{coherency} \emph{offset27}[\emph{registerZ}] = \emph{registerT}

\paragraph{Description} The effective address is computed by adding the \emph{offset27} to the \emph{registerZ}.
This instruction implements atomic load-MINU-store on signed byte. The byte at effective address is MINUed with the \emph{registerT}, and its previous value is returned into the \emph{registerT}.


\paragraph{Modifier(s)}
\noindent\begin{itemize}
\item coherency (cf. coherency in section \ref{par:modifier-coherency} on page \pageref{par:modifier-coherency})
\end{itemize}

\paragraph{Execution} {Reservation table \textsf{LSU2\_MEMW\_AUXR\_AUXW.X} (Section~\ref{sec:reservation-LSU2-MEMW-AUXR-AUXW.X})}
~
{\begin{lstlisting}
stage ID:
  new offset = _SX_27(offset27);
  new coherency = _ZX_2(coherency);
stage RR:
  new base = GPR[registerZ];
  new address = base + offset;
stage E1:
  new argument2 = GPR[registerT];
  new result2 = _ZX_8(MEM.atomic_minu(address, 0x1, coherency << 32, argument2.8[0], registerT));
stage SR:
  GPR[registerT] = result2;
\end{lstlisting}}
\newpage

\paragraph{Encoding} Format \textsf{LSU\_ALMINUSB.D} (Section~\ref{sec:format-LSU-ALMINUSB.D}), \verb|lsucode5 = 000|\\

\bitfields{ 1/P, 2/00, 2/-- --, 27/extend27, 1/1, 2/00, 2/-- --, 27/offset27, 1/1, 2/01, 3/111, 2/coherency, 6/registerT, 2/11, 3/000, 1/--, 5/01010, 1/0, 6/registerZ }


\paragraph{Syntax} \texttt{alminub}\emph{coherency} \emph{extend27\_\discretionary{}{}{}offset27}[\emph{registerZ}] = \emph{registerT}

\paragraph{Description} The effective address is computed by adding the \emph{extend27\_\discretionary{}{}{}offset27} to the \emph{registerZ}.
This instruction implements atomic load-MINU-store on signed byte. The byte at effective address is MINUed with the \emph{registerT}, and its previous value is returned into the \emph{registerT}.


\paragraph{Modifier(s)}
\noindent\begin{itemize}
\item coherency (cf. coherency in section \ref{par:modifier-coherency} on page \pageref{par:modifier-coherency})
\end{itemize}

\paragraph{Execution} {Reservation table \textsf{LSU2\_MEMW\_AUXR\_AUXW.Y} (Section~\ref{sec:reservation-LSU2-MEMW-AUXR-AUXW.Y})}
~
{\begin{lstlisting}
stage ID:
  new offset = _SX_54(extend27_offset27);
  new coherency = _ZX_2(coherency);
stage RR:
  new base = GPR[registerZ];
  new address = base + offset;
stage E1:
  new argument2 = GPR[registerT];
  new result2 = _ZX_8(MEM.atomic_minu(address, 0x1, coherency << 32, argument2.8[0], registerT));
stage SR:
  GPR[registerT] = result2;
\end{lstlisting}}
\newpage
\subsubsection[{\small ALMINUD: Atomic Load and MIN Unsigned Double Word}]{ALMINUD\hfill Atomic Load and MIN Unsigned Double Word}
\label{instr:ALMINUD}\index{alminud}
 
\paragraph{Encoding} Format \textsf{LSU\_ALMINUSB} (Section~\ref{sec:format-LSU-ALMINUSB}), \verb|lsucode5 = 011|\\

\bitfields{ 1/P, 2/01, 3/111, 2/coherency, 6/registerT, 2/11, 3/011, 1/--, 5/01010, 1/0, 6/registerZ }


\paragraph{Syntax} \texttt{alminud}\emph{coherency} [\emph{registerZ}] = \emph{registerT}

\paragraph{Description} The effective address is the value of the \emph{registerZ}.
This instruction implements atomic load-MINU-store on signed double word. The double word at effective address is MINUed with the \emph{registerT}, and its previous value is returned into the \emph{registerT}. The effective address must be double word aligned (multiple of 8).


\paragraph{Modifier(s)}
\noindent\begin{itemize}
\item coherency (cf. coherency in section \ref{par:modifier-coherency} on page \pageref{par:modifier-coherency})
\end{itemize}

\paragraph{Execution} {Reservation table \textsf{LSU2\_MEMW\_AUXR\_AUXW} (Section~\ref{sec:reservation-LSU2-MEMW-AUXR-AUXW})}
~
{\begin{lstlisting}
stage ID:
  new coherency = _ZX_2(coherency);
stage RR:
  new address = GPR[registerZ];
stage E1:
  new argument2 = GPR[registerT];
  new result2 = MEM.atomic_minu(address, 0xFF, coherency << 32, argument2.64[0], registerT);
stage SR:
  GPR[registerT] = result2;
\end{lstlisting}}
\newpage

\paragraph{Encoding} Format \textsf{LSU\_ALMINUSB.O} (Section~\ref{sec:format-LSU-ALMINUSB.O}), \verb|lsucode5 = 011|\\

\bitfields{ 1/P, 2/00, 2/-- --, 27/offset27, 1/1, 2/01, 3/111, 2/coherency, 6/registerT, 2/11, 3/011, 1/--, 5/01010, 1/0, 6/registerZ }


\paragraph{Syntax} \texttt{alminud}\emph{coherency} \emph{offset27}[\emph{registerZ}] = \emph{registerT}

\paragraph{Description} The effective address is computed by adding the \emph{offset27} to the \emph{registerZ}.
This instruction implements atomic load-MINU-store on signed double word. The double word at effective address is MINUed with the \emph{registerT}, and its previous value is returned into the \emph{registerT}. The effective address must be double word aligned (multiple of 8).


\paragraph{Modifier(s)}
\noindent\begin{itemize}
\item coherency (cf. coherency in section \ref{par:modifier-coherency} on page \pageref{par:modifier-coherency})
\end{itemize}

\paragraph{Execution} {Reservation table \textsf{LSU2\_MEMW\_AUXR\_AUXW.X} (Section~\ref{sec:reservation-LSU2-MEMW-AUXR-AUXW.X})}
~
{\begin{lstlisting}
stage ID:
  new offset = _SX_27(offset27);
  new coherency = _ZX_2(coherency);
stage RR:
  new base = GPR[registerZ];
  new address = base + offset;
stage E1:
  new argument2 = GPR[registerT];
  new result2 = MEM.atomic_minu(address, 0xFF, coherency << 32, argument2.64[0], registerT);
stage SR:
  GPR[registerT] = result2;
\end{lstlisting}}
\newpage

\paragraph{Encoding} Format \textsf{LSU\_ALMINUSB.D} (Section~\ref{sec:format-LSU-ALMINUSB.D}), \verb|lsucode5 = 011|\\

\bitfields{ 1/P, 2/00, 2/-- --, 27/extend27, 1/1, 2/00, 2/-- --, 27/offset27, 1/1, 2/01, 3/111, 2/coherency, 6/registerT, 2/11, 3/011, 1/--, 5/01010, 1/0, 6/registerZ }


\paragraph{Syntax} \texttt{alminud}\emph{coherency} \emph{extend27\_\discretionary{}{}{}offset27}[\emph{registerZ}] = \emph{registerT}

\paragraph{Description} The effective address is computed by adding the \emph{extend27\_\discretionary{}{}{}offset27} to the \emph{registerZ}.
This instruction implements atomic load-MINU-store on signed double word. The double word at effective address is MINUed with the \emph{registerT}, and its previous value is returned into the \emph{registerT}. The effective address must be double word aligned (multiple of 8).


\paragraph{Modifier(s)}
\noindent\begin{itemize}
\item coherency (cf. coherency in section \ref{par:modifier-coherency} on page \pageref{par:modifier-coherency})
\end{itemize}

\paragraph{Execution} {Reservation table \textsf{LSU2\_MEMW\_AUXR\_AUXW.Y} (Section~\ref{sec:reservation-LSU2-MEMW-AUXR-AUXW.Y})}
~
{\begin{lstlisting}
stage ID:
  new offset = _SX_54(extend27_offset27);
  new coherency = _ZX_2(coherency);
stage RR:
  new base = GPR[registerZ];
  new address = base + offset;
stage E1:
  new argument2 = GPR[registerT];
  new result2 = MEM.atomic_minu(address, 0xFF, coherency << 32, argument2.64[0], registerT);
stage SR:
  GPR[registerT] = result2;
\end{lstlisting}}
\newpage
\subsubsection[{\small ALMINUH: Atomic Load and MIN Unsigned Half Word}]{ALMINUH\hfill Atomic Load and MIN Unsigned Half Word}
\label{instr:ALMINUH}\index{alminuh}
 
\paragraph{Encoding} Format \textsf{LSU\_ALMINUSB} (Section~\ref{sec:format-LSU-ALMINUSB}), \verb|lsucode5 = 001|\\

\bitfields{ 1/P, 2/01, 3/111, 2/coherency, 6/registerT, 2/11, 3/001, 1/--, 5/01010, 1/0, 6/registerZ }


\paragraph{Syntax} \texttt{alminuh}\emph{coherency} [\emph{registerZ}] = \emph{registerT}

\paragraph{Description} The effective address is the value of the \emph{registerZ}.
This instruction implements atomic load-MINU-store on signed half word. The half word at effective address is MINUed with the \emph{registerT}, and its previous value is returned into the \emph{registerT}. The effective address must be half word aligned (multiple of 2).


\paragraph{Modifier(s)}
\noindent\begin{itemize}
\item coherency (cf. coherency in section \ref{par:modifier-coherency} on page \pageref{par:modifier-coherency})
\end{itemize}

\paragraph{Execution} {Reservation table \textsf{LSU2\_MEMW\_AUXR\_AUXW} (Section~\ref{sec:reservation-LSU2-MEMW-AUXR-AUXW})}
~
{\begin{lstlisting}
stage ID:
  new coherency = _ZX_2(coherency);
stage RR:
  new address = GPR[registerZ];
stage E1:
  new argument2 = GPR[registerT];
  new result2 = _ZX_16(MEM.atomic_minu(address, 0x3, coherency << 32, argument2.16[0], registerT));
stage SR:
  GPR[registerT] = result2;
\end{lstlisting}}
\newpage

\paragraph{Encoding} Format \textsf{LSU\_ALMINUSB.O} (Section~\ref{sec:format-LSU-ALMINUSB.O}), \verb|lsucode5 = 001|\\

\bitfields{ 1/P, 2/00, 2/-- --, 27/offset27, 1/1, 2/01, 3/111, 2/coherency, 6/registerT, 2/11, 3/001, 1/--, 5/01010, 1/0, 6/registerZ }


\paragraph{Syntax} \texttt{alminuh}\emph{coherency} \emph{offset27}[\emph{registerZ}] = \emph{registerT}

\paragraph{Description} The effective address is computed by adding the \emph{offset27} to the \emph{registerZ}.
This instruction implements atomic load-MINU-store on signed half word. The half word at effective address is MINUed with the \emph{registerT}, and its previous value is returned into the \emph{registerT}. The effective address must be half word aligned (multiple of 2).


\paragraph{Modifier(s)}
\noindent\begin{itemize}
\item coherency (cf. coherency in section \ref{par:modifier-coherency} on page \pageref{par:modifier-coherency})
\end{itemize}

\paragraph{Execution} {Reservation table \textsf{LSU2\_MEMW\_AUXR\_AUXW.X} (Section~\ref{sec:reservation-LSU2-MEMW-AUXR-AUXW.X})}
~
{\begin{lstlisting}
stage ID:
  new offset = _SX_27(offset27);
  new coherency = _ZX_2(coherency);
stage RR:
  new base = GPR[registerZ];
  new address = base + offset;
stage E1:
  new argument2 = GPR[registerT];
  new result2 = _ZX_16(MEM.atomic_minu(address, 0x3, coherency << 32, argument2.16[0], registerT));
stage SR:
  GPR[registerT] = result2;
\end{lstlisting}}
\newpage

\paragraph{Encoding} Format \textsf{LSU\_ALMINUSB.D} (Section~\ref{sec:format-LSU-ALMINUSB.D}), \verb|lsucode5 = 001|\\

\bitfields{ 1/P, 2/00, 2/-- --, 27/extend27, 1/1, 2/00, 2/-- --, 27/offset27, 1/1, 2/01, 3/111, 2/coherency, 6/registerT, 2/11, 3/001, 1/--, 5/01010, 1/0, 6/registerZ }


\paragraph{Syntax} \texttt{alminuh}\emph{coherency} \emph{extend27\_\discretionary{}{}{}offset27}[\emph{registerZ}] = \emph{registerT}

\paragraph{Description} The effective address is computed by adding the \emph{extend27\_\discretionary{}{}{}offset27} to the \emph{registerZ}.
This instruction implements atomic load-MINU-store on signed half word. The half word at effective address is MINUed with the \emph{registerT}, and its previous value is returned into the \emph{registerT}. The effective address must be half word aligned (multiple of 2).


\paragraph{Modifier(s)}
\noindent\begin{itemize}
\item coherency (cf. coherency in section \ref{par:modifier-coherency} on page \pageref{par:modifier-coherency})
\end{itemize}

\paragraph{Execution} {Reservation table \textsf{LSU2\_MEMW\_AUXR\_AUXW.Y} (Section~\ref{sec:reservation-LSU2-MEMW-AUXR-AUXW.Y})}
~
{\begin{lstlisting}
stage ID:
  new offset = _SX_54(extend27_offset27);
  new coherency = _ZX_2(coherency);
stage RR:
  new base = GPR[registerZ];
  new address = base + offset;
stage E1:
  new argument2 = GPR[registerT];
  new result2 = _ZX_16(MEM.atomic_minu(address, 0x3, coherency << 32, argument2.16[0], registerT));
stage SR:
  GPR[registerT] = result2;
\end{lstlisting}}
\newpage
\subsubsection[{\small ALMINUW: Atomic Load and MIN Unsigned Word}]{ALMINUW\hfill Atomic Load and MIN Unsigned Word}
\label{instr:ALMINUW}\index{alminuw}
 
\paragraph{Encoding} Format \textsf{LSU\_ALMINUSB} (Section~\ref{sec:format-LSU-ALMINUSB}), \verb|lsucode5 = 010|\\

\bitfields{ 1/P, 2/01, 3/111, 2/coherency, 6/registerT, 2/11, 3/010, 1/--, 5/01010, 1/0, 6/registerZ }


\paragraph{Syntax} \texttt{alminuw}\emph{coherency} [\emph{registerZ}] = \emph{registerT}

\paragraph{Description} The effective address is the value of the \emph{registerZ}.
This instruction implements atomic load-MINU-store on signed word. The word at effective address is MINUed with the \emph{registerT}, and its previous value is returned into the \emph{registerT}. The effective address must be word aligned (multiple of 4).


\paragraph{Modifier(s)}
\noindent\begin{itemize}
\item coherency (cf. coherency in section \ref{par:modifier-coherency} on page \pageref{par:modifier-coherency})
\end{itemize}

\paragraph{Execution} {Reservation table \textsf{LSU2\_MEMW\_AUXR\_AUXW} (Section~\ref{sec:reservation-LSU2-MEMW-AUXR-AUXW})}
~
{\begin{lstlisting}
stage ID:
  new coherency = _ZX_2(coherency);
stage RR:
  new address = GPR[registerZ];
stage E1:
  new argument2 = GPR[registerT];
  new result2 = _ZX_32(MEM.atomic_minu(address, 0xF, coherency << 32, argument2.32[0], registerT));
stage SR:
  GPR[registerT] = result2;
\end{lstlisting}}
\newpage

\paragraph{Encoding} Format \textsf{LSU\_ALMINUSB.O} (Section~\ref{sec:format-LSU-ALMINUSB.O}), \verb|lsucode5 = 010|\\

\bitfields{ 1/P, 2/00, 2/-- --, 27/offset27, 1/1, 2/01, 3/111, 2/coherency, 6/registerT, 2/11, 3/010, 1/--, 5/01010, 1/0, 6/registerZ }


\paragraph{Syntax} \texttt{alminuw}\emph{coherency} \emph{offset27}[\emph{registerZ}] = \emph{registerT}

\paragraph{Description} The effective address is computed by adding the \emph{offset27} to the \emph{registerZ}.
This instruction implements atomic load-MINU-store on signed word. The word at effective address is MINUed with the \emph{registerT}, and its previous value is returned into the \emph{registerT}. The effective address must be word aligned (multiple of 4).


\paragraph{Modifier(s)}
\noindent\begin{itemize}
\item coherency (cf. coherency in section \ref{par:modifier-coherency} on page \pageref{par:modifier-coherency})
\end{itemize}

\paragraph{Execution} {Reservation table \textsf{LSU2\_MEMW\_AUXR\_AUXW.X} (Section~\ref{sec:reservation-LSU2-MEMW-AUXR-AUXW.X})}
~
{\begin{lstlisting}
stage ID:
  new offset = _SX_27(offset27);
  new coherency = _ZX_2(coherency);
stage RR:
  new base = GPR[registerZ];
  new address = base + offset;
stage E1:
  new argument2 = GPR[registerT];
  new result2 = _ZX_32(MEM.atomic_minu(address, 0xF, coherency << 32, argument2.32[0], registerT));
stage SR:
  GPR[registerT] = result2;
\end{lstlisting}}
\newpage

\paragraph{Encoding} Format \textsf{LSU\_ALMINUSB.D} (Section~\ref{sec:format-LSU-ALMINUSB.D}), \verb|lsucode5 = 010|\\

\bitfields{ 1/P, 2/00, 2/-- --, 27/extend27, 1/1, 2/00, 2/-- --, 27/offset27, 1/1, 2/01, 3/111, 2/coherency, 6/registerT, 2/11, 3/010, 1/--, 5/01010, 1/0, 6/registerZ }


\paragraph{Syntax} \texttt{alminuw}\emph{coherency} \emph{extend27\_\discretionary{}{}{}offset27}[\emph{registerZ}] = \emph{registerT}

\paragraph{Description} The effective address is computed by adding the \emph{extend27\_\discretionary{}{}{}offset27} to the \emph{registerZ}.
This instruction implements atomic load-MINU-store on signed word. The word at effective address is MINUed with the \emph{registerT}, and its previous value is returned into the \emph{registerT}. The effective address must be word aligned (multiple of 4).


\paragraph{Modifier(s)}
\noindent\begin{itemize}
\item coherency (cf. coherency in section \ref{par:modifier-coherency} on page \pageref{par:modifier-coherency})
\end{itemize}

\paragraph{Execution} {Reservation table \textsf{LSU2\_MEMW\_AUXR\_AUXW.Y} (Section~\ref{sec:reservation-LSU2-MEMW-AUXR-AUXW.Y})}
~
{\begin{lstlisting}
stage ID:
  new offset = _SX_54(extend27_offset27);
  new coherency = _ZX_2(coherency);
stage RR:
  new base = GPR[registerZ];
  new address = base + offset;
stage E1:
  new argument2 = GPR[registerT];
  new result2 = _ZX_32(MEM.atomic_minu(address, 0xF, coherency << 32, argument2.32[0], registerT));
stage SR:
  GPR[registerT] = result2;
\end{lstlisting}}
\newpage
\subsubsection[{\small ALMINW: Atomic Load and MIN Word}]{ALMINW\hfill Atomic Load and MIN Word}
\label{instr:ALMINW}\index{alminw}
 
\paragraph{Encoding} Format \textsf{LSU\_ALMINSB} (Section~\ref{sec:format-LSU-ALMINSB}), \verb|lsucode5 = 010|\\

\bitfields{ 1/P, 2/01, 3/111, 2/coherency, 6/registerT, 2/11, 3/010, 1/--, 5/01000, 1/0, 6/registerZ }


\paragraph{Syntax} \texttt{alminw}\emph{coherency} [\emph{registerZ}] = \emph{registerT}

\paragraph{Description} The effective address is the value of the \emph{registerZ}.
This instruction implements atomic load-MIN-store on signed word. The word at effective address is MINed with the \emph{registerT}, and its previous value is returned into the \emph{registerT}. The effective address must be word aligned (multiple of 4).


\paragraph{Modifier(s)}
\noindent\begin{itemize}
\item coherency (cf. coherency in section \ref{par:modifier-coherency} on page \pageref{par:modifier-coherency})
\end{itemize}

\paragraph{Execution} {Reservation table \textsf{LSU2\_MEMW\_AUXR\_AUXW} (Section~\ref{sec:reservation-LSU2-MEMW-AUXR-AUXW})}
~
{\begin{lstlisting}
stage ID:
  new coherency = _ZX_2(coherency);
stage RR:
  new address = GPR[registerZ];
stage E1:
  new argument2 = GPR[registerT];
  new result2 = _ZX_32(MEM.atomic_min(address, 0xF, coherency << 32, argument2.32[0], registerT));
stage SR:
  GPR[registerT] = result2;
\end{lstlisting}}
\newpage

\paragraph{Encoding} Format \textsf{LSU\_ALMINSB.O} (Section~\ref{sec:format-LSU-ALMINSB.O}), \verb|lsucode5 = 010|\\

\bitfields{ 1/P, 2/00, 2/-- --, 27/offset27, 1/1, 2/01, 3/111, 2/coherency, 6/registerT, 2/11, 3/010, 1/--, 5/01000, 1/0, 6/registerZ }


\paragraph{Syntax} \texttt{alminw}\emph{coherency} \emph{offset27}[\emph{registerZ}] = \emph{registerT}

\paragraph{Description} The effective address is computed by adding the \emph{offset27} to the \emph{registerZ}.
This instruction implements atomic load-MIN-store on signed word. The word at effective address is MINed with the \emph{registerT}, and its previous value is returned into the \emph{registerT}. The effective address must be word aligned (multiple of 4).


\paragraph{Modifier(s)}
\noindent\begin{itemize}
\item coherency (cf. coherency in section \ref{par:modifier-coherency} on page \pageref{par:modifier-coherency})
\end{itemize}

\paragraph{Execution} {Reservation table \textsf{LSU2\_MEMW\_AUXR\_AUXW.X} (Section~\ref{sec:reservation-LSU2-MEMW-AUXR-AUXW.X})}
~
{\begin{lstlisting}
stage ID:
  new offset = _SX_27(offset27);
  new coherency = _ZX_2(coherency);
stage RR:
  new base = GPR[registerZ];
  new address = base + offset;
stage E1:
  new argument2 = GPR[registerT];
  new result2 = _ZX_32(MEM.atomic_min(address, 0xF, coherency << 32, argument2.32[0], registerT));
stage SR:
  GPR[registerT] = result2;
\end{lstlisting}}
\newpage

\paragraph{Encoding} Format \textsf{LSU\_ALMINSB.D} (Section~\ref{sec:format-LSU-ALMINSB.D}), \verb|lsucode5 = 010|\\

\bitfields{ 1/P, 2/00, 2/-- --, 27/extend27, 1/1, 2/00, 2/-- --, 27/offset27, 1/1, 2/01, 3/111, 2/coherency, 6/registerT, 2/11, 3/010, 1/--, 5/01000, 1/0, 6/registerZ }


\paragraph{Syntax} \texttt{alminw}\emph{coherency} \emph{extend27\_\discretionary{}{}{}offset27}[\emph{registerZ}] = \emph{registerT}

\paragraph{Description} The effective address is computed by adding the \emph{extend27\_\discretionary{}{}{}offset27} to the \emph{registerZ}.
This instruction implements atomic load-MIN-store on signed word. The word at effective address is MINed with the \emph{registerT}, and its previous value is returned into the \emph{registerT}. The effective address must be word aligned (multiple of 4).


\paragraph{Modifier(s)}
\noindent\begin{itemize}
\item coherency (cf. coherency in section \ref{par:modifier-coherency} on page \pageref{par:modifier-coherency})
\end{itemize}

\paragraph{Execution} {Reservation table \textsf{LSU2\_MEMW\_AUXR\_AUXW.Y} (Section~\ref{sec:reservation-LSU2-MEMW-AUXR-AUXW.Y})}
~
{\begin{lstlisting}
stage ID:
  new offset = _SX_54(extend27_offset27);
  new coherency = _ZX_2(coherency);
stage RR:
  new base = GPR[registerZ];
  new address = base + offset;
stage E1:
  new argument2 = GPR[registerT];
  new result2 = _ZX_32(MEM.atomic_min(address, 0xF, coherency << 32, argument2.32[0], registerT));
stage SR:
  GPR[registerT] = result2;
\end{lstlisting}}
\newpage
\subsubsection[{\small ALW: Atomic Load Word}]{ALW\hfill Atomic Load Word}
\label{instr:ALW}\index{alw}
 
\paragraph{Encoding} Format \textsf{LSU\_ALSB} (Section~\ref{sec:format-LSU-ALSB}), \verb|lsucode5 = 010|\\

\bitfields{ 1/P, 2/01, 3/111, 2/coherency, 6/registerW, 2/11, 3/010, 1/--, 5/00000, 1/0, 6/registerZ }


\paragraph{Syntax} \texttt{alw}\emph{coherency} \emph{registerW} = [\emph{registerZ}]

\paragraph{Description} The effective address is the value of the \emph{registerZ}.
The \emph{registerW} is atomically loaded with zero extension from the memory word starting at the effective address. This instruction is provided to directly operate on the last level of the memory hierarchy, which is typically a processor DDR or a cluster TCM, irrespective of the MMU DCP (Data Cache Policy) settings. The effective address must be word aligned (multiple of 4).


\paragraph{Modifier(s)}
\noindent\begin{itemize}
\item coherency (cf. coherency in section \ref{par:modifier-coherency} on page \pageref{par:modifier-coherency})
\end{itemize}

\paragraph{Execution} {Reservation table \textsf{LSU\_MEMW\_AUXW} (Section~\ref{sec:reservation-LSU-MEMW-AUXW})}
~
{\begin{lstlisting}
stage ID:
  new coherency = _ZX_2(coherency);
stage RR:
  new address = GPR[registerZ];
  new result1 = _ZX_32(MEM.atomic_load(address, 0xF, coherency << 32, registerW));
stage CR:
  GPR[registerW] = result1;
\end{lstlisting}}
\newpage

\paragraph{Encoding} Format \textsf{LSU\_ALSB.O} (Section~\ref{sec:format-LSU-ALSB.O}), \verb|lsucode5 = 010|\\

\bitfields{ 1/P, 2/00, 2/-- --, 27/offset27, 1/1, 2/01, 3/111, 2/coherency, 6/registerW, 2/11, 3/010, 1/--, 5/00000, 1/0, 6/registerZ }


\paragraph{Syntax} \texttt{alw}\emph{coherency} \emph{registerW} = \emph{offset27}[\emph{registerZ}]

\paragraph{Description} The effective address is computed by adding the \emph{offset27} to the \emph{registerZ}.
The \emph{registerW} is atomically loaded with zero extension from the memory word starting at the effective address. This instruction is provided to directly operate on the last level of the memory hierarchy, which is typically a processor DDR or a cluster TCM, irrespective of the MMU DCP (Data Cache Policy) settings. The effective address must be word aligned (multiple of 4).


\paragraph{Modifier(s)}
\noindent\begin{itemize}
\item coherency (cf. coherency in section \ref{par:modifier-coherency} on page \pageref{par:modifier-coherency})
\end{itemize}

\paragraph{Execution} {Reservation table \textsf{LSU\_MEMW\_AUXW.X} (Section~\ref{sec:reservation-LSU-MEMW-AUXW.X})}
~
{\begin{lstlisting}
stage ID:
  new offset = _SX_27(offset27);
  new coherency = _ZX_2(coherency);
stage RR:
  new base = GPR[registerZ];
  new address = base + offset;
  new result1 = _ZX_32(MEM.atomic_load(address, 0xF, coherency << 32, registerW));
stage CR:
  GPR[registerW] = result1;
\end{lstlisting}}
\newpage

\paragraph{Encoding} Format \textsf{LSU\_ALSB.D} (Section~\ref{sec:format-LSU-ALSB.D}), \verb|lsucode5 = 010|\\

\bitfields{ 1/P, 2/00, 2/-- --, 27/extend27, 1/1, 2/00, 2/-- --, 27/offset27, 1/1, 2/01, 3/111, 2/coherency, 6/registerW, 2/11, 3/010, 1/--, 5/00000, 1/0, 6/registerZ }


\paragraph{Syntax} \texttt{alw}\emph{coherency} \emph{registerW} = \emph{extend27\_\discretionary{}{}{}offset27}[\emph{registerZ}]

\paragraph{Description} The effective address is computed by adding the \emph{extend27\_\discretionary{}{}{}offset27} to the \emph{registerZ}.
The \emph{registerW} is atomically loaded with zero extension from the memory word starting at the effective address. This instruction is provided to directly operate on the last level of the memory hierarchy, which is typically a processor DDR or a cluster TCM, irrespective of the MMU DCP (Data Cache Policy) settings. The effective address must be word aligned (multiple of 4).


\paragraph{Modifier(s)}
\noindent\begin{itemize}
\item coherency (cf. coherency in section \ref{par:modifier-coherency} on page \pageref{par:modifier-coherency})
\end{itemize}

\paragraph{Execution} {Reservation table \textsf{LSU\_MEMW\_AUXW.Y} (Section~\ref{sec:reservation-LSU-MEMW-AUXW.Y})}
~
{\begin{lstlisting}
stage ID:
  new offset = _SX_54(extend27_offset27);
  new coherency = _ZX_2(coherency);
stage RR:
  new base = GPR[registerZ];
  new address = base + offset;
  new result1 = _ZX_32(MEM.atomic_load(address, 0xF, coherency << 32, registerW));
stage CR:
  GPR[registerW] = result1;
\end{lstlisting}}
\newpage
\subsubsection[{\small ASADDB: Atomic Store Add Byte}]{ASADDB\hfill Atomic Store Add Byte}
\label{instr:ASADDB}\index{asaddb}
 
\paragraph{Encoding} Format \textsf{LSU\_ASADDSB} (Section~\ref{sec:format-LSU-ASADDSB}), \verb|lsucode5 = 000|\\

\bitfields{ 1/P, 2/01, 3/111, 2/coherency, 6/registerT, 2/11, 3/000, 1/--, 5/10100, 1/0, 6/registerZ }


\paragraph{Syntax} \texttt{asaddb}\emph{coherency} [\emph{registerZ}] = \emph{registerT}

\paragraph{Description} The effective address is the value of the \emph{registerZ}.
This instruction implements atomic add-store on byte. The byte at effective address is added with the \emph{registerT}.


\paragraph{Modifier(s)}
\noindent\begin{itemize}
\item coherency (cf. coherency in section \ref{par:modifier-coherency} on page \pageref{par:modifier-coherency})
\end{itemize}

\paragraph{Execution} {Reservation table \textsf{LSU2\_MEMW\_AUXR} (Section~\ref{sec:reservation-LSU2-MEMW-AUXR})}
~
{\begin{lstlisting}
stage ID:
  new coherency = _ZX_2(coherency);
stage RR:
  new address = GPR[registerZ];
stage E1:
  new argument2 = GPR[registerT];
stage E3:
  MEM.atomic_add(address, 0x1, coherency << 32, argument2.8[0], registerT);
\end{lstlisting}}
\newpage

\paragraph{Encoding} Format \textsf{LSU\_ASADDSB.O} (Section~\ref{sec:format-LSU-ASADDSB.O}), \verb|lsucode5 = 000|\\

\bitfields{ 1/P, 2/00, 2/-- --, 27/offset27, 1/1, 2/01, 3/111, 2/coherency, 6/registerT, 2/11, 3/000, 1/--, 5/10100, 1/0, 6/registerZ }


\paragraph{Syntax} \texttt{asaddb}\emph{coherency} \emph{offset27}[\emph{registerZ}] = \emph{registerT}

\paragraph{Description} The effective address is computed by adding the \emph{offset27} to the \emph{registerZ}.
This instruction implements atomic add-store on byte. The byte at effective address is added with the \emph{registerT}.


\paragraph{Modifier(s)}
\noindent\begin{itemize}
\item coherency (cf. coherency in section \ref{par:modifier-coherency} on page \pageref{par:modifier-coherency})
\end{itemize}

\paragraph{Execution} {Reservation table \textsf{LSU2\_MEMW\_AUXR.X} (Section~\ref{sec:reservation-LSU2-MEMW-AUXR.X})}
~
{\begin{lstlisting}
stage ID:
  new offset = _SX_27(offset27);
  new coherency = _ZX_2(coherency);
stage RR:
  new base = GPR[registerZ];
  new address = base + offset;
stage E1:
  new argument2 = GPR[registerT];
stage E3:
  MEM.atomic_add(address, 0x1, coherency << 32, argument2.8[0], registerT);
\end{lstlisting}}
\newpage

\paragraph{Encoding} Format \textsf{LSU\_ASADDSB.D} (Section~\ref{sec:format-LSU-ASADDSB.D}), \verb|lsucode5 = 000|\\

\bitfields{ 1/P, 2/00, 2/-- --, 27/extend27, 1/1, 2/00, 2/-- --, 27/offset27, 1/1, 2/01, 3/111, 2/coherency, 6/registerT, 2/11, 3/000, 1/--, 5/10100, 1/0, 6/registerZ }


\paragraph{Syntax} \texttt{asaddb}\emph{coherency} \emph{extend27\_\discretionary{}{}{}offset27}[\emph{registerZ}] = \emph{registerT}

\paragraph{Description} The effective address is computed by adding the \emph{extend27\_\discretionary{}{}{}offset27} to the \emph{registerZ}.
This instruction implements atomic add-store on byte. The byte at effective address is added with the \emph{registerT}.


\paragraph{Modifier(s)}
\noindent\begin{itemize}
\item coherency (cf. coherency in section \ref{par:modifier-coherency} on page \pageref{par:modifier-coherency})
\end{itemize}

\paragraph{Execution} {Reservation table \textsf{LSU2\_MEMW\_AUXR.Y} (Section~\ref{sec:reservation-LSU2-MEMW-AUXR.Y})}
~
{\begin{lstlisting}
stage ID:
  new offset = _SX_54(extend27_offset27);
  new coherency = _ZX_2(coherency);
stage RR:
  new base = GPR[registerZ];
  new address = base + offset;
stage E1:
  new argument2 = GPR[registerT];
stage E3:
  MEM.atomic_add(address, 0x1, coherency << 32, argument2.8[0], registerT);
\end{lstlisting}}
\newpage
\subsubsection[{\small ASADDD: Atomic Store Add Double Word}]{ASADDD\hfill Atomic Store Add Double Word}
\label{instr:ASADDD}\index{asaddd}
 
\paragraph{Encoding} Format \textsf{LSU\_ASADDSB} (Section~\ref{sec:format-LSU-ASADDSB}), \verb|lsucode5 = 011|\\

\bitfields{ 1/P, 2/01, 3/111, 2/coherency, 6/registerT, 2/11, 3/011, 1/--, 5/10100, 1/0, 6/registerZ }


\paragraph{Syntax} \texttt{asaddd}\emph{coherency} [\emph{registerZ}] = \emph{registerT}

\paragraph{Description} The effective address is the value of the \emph{registerZ}.
This instruction implements atomic add-store on double word. The double word at effective address is added with the \emph{registerT}. The effective address must be double word aligned (multiple of 8).


\paragraph{Modifier(s)}
\noindent\begin{itemize}
\item coherency (cf. coherency in section \ref{par:modifier-coherency} on page \pageref{par:modifier-coherency})
\end{itemize}

\paragraph{Execution} {Reservation table \textsf{LSU2\_MEMW\_AUXR} (Section~\ref{sec:reservation-LSU2-MEMW-AUXR})}
~
{\begin{lstlisting}
stage ID:
  new coherency = _ZX_2(coherency);
stage RR:
  new address = GPR[registerZ];
stage E1:
  new argument2 = GPR[registerT];
stage E3:
  MEM.atomic_add(address, 0xFF, coherency << 32, argument2.64[0], registerT);
\end{lstlisting}}
\newpage

\paragraph{Encoding} Format \textsf{LSU\_ASADDSB.O} (Section~\ref{sec:format-LSU-ASADDSB.O}), \verb|lsucode5 = 011|\\

\bitfields{ 1/P, 2/00, 2/-- --, 27/offset27, 1/1, 2/01, 3/111, 2/coherency, 6/registerT, 2/11, 3/011, 1/--, 5/10100, 1/0, 6/registerZ }


\paragraph{Syntax} \texttt{asaddd}\emph{coherency} \emph{offset27}[\emph{registerZ}] = \emph{registerT}

\paragraph{Description} The effective address is computed by adding the \emph{offset27} to the \emph{registerZ}.
This instruction implements atomic add-store on double word. The double word at effective address is added with the \emph{registerT}. The effective address must be double word aligned (multiple of 8).


\paragraph{Modifier(s)}
\noindent\begin{itemize}
\item coherency (cf. coherency in section \ref{par:modifier-coherency} on page \pageref{par:modifier-coherency})
\end{itemize}

\paragraph{Execution} {Reservation table \textsf{LSU2\_MEMW\_AUXR.X} (Section~\ref{sec:reservation-LSU2-MEMW-AUXR.X})}
~
{\begin{lstlisting}
stage ID:
  new offset = _SX_27(offset27);
  new coherency = _ZX_2(coherency);
stage RR:
  new base = GPR[registerZ];
  new address = base + offset;
stage E1:
  new argument2 = GPR[registerT];
stage E3:
  MEM.atomic_add(address, 0xFF, coherency << 32, argument2.64[0], registerT);
\end{lstlisting}}
\newpage

\paragraph{Encoding} Format \textsf{LSU\_ASADDSB.D} (Section~\ref{sec:format-LSU-ASADDSB.D}), \verb|lsucode5 = 011|\\

\bitfields{ 1/P, 2/00, 2/-- --, 27/extend27, 1/1, 2/00, 2/-- --, 27/offset27, 1/1, 2/01, 3/111, 2/coherency, 6/registerT, 2/11, 3/011, 1/--, 5/10100, 1/0, 6/registerZ }


\paragraph{Syntax} \texttt{asaddd}\emph{coherency} \emph{extend27\_\discretionary{}{}{}offset27}[\emph{registerZ}] = \emph{registerT}

\paragraph{Description} The effective address is computed by adding the \emph{extend27\_\discretionary{}{}{}offset27} to the \emph{registerZ}.
This instruction implements atomic add-store on double word. The double word at effective address is added with the \emph{registerT}. The effective address must be double word aligned (multiple of 8).


\paragraph{Modifier(s)}
\noindent\begin{itemize}
\item coherency (cf. coherency in section \ref{par:modifier-coherency} on page \pageref{par:modifier-coherency})
\end{itemize}

\paragraph{Execution} {Reservation table \textsf{LSU2\_MEMW\_AUXR.Y} (Section~\ref{sec:reservation-LSU2-MEMW-AUXR.Y})}
~
{\begin{lstlisting}
stage ID:
  new offset = _SX_54(extend27_offset27);
  new coherency = _ZX_2(coherency);
stage RR:
  new base = GPR[registerZ];
  new address = base + offset;
stage E1:
  new argument2 = GPR[registerT];
stage E3:
  MEM.atomic_add(address, 0xFF, coherency << 32, argument2.64[0], registerT);
\end{lstlisting}}
\newpage
\subsubsection[{\small ASADDH: Atomic Store Add Half Word}]{ASADDH\hfill Atomic Store Add Half Word}
\label{instr:ASADDH}\index{asaddh}
 
\paragraph{Encoding} Format \textsf{LSU\_ASADDSB} (Section~\ref{sec:format-LSU-ASADDSB}), \verb|lsucode5 = 001|\\

\bitfields{ 1/P, 2/01, 3/111, 2/coherency, 6/registerT, 2/11, 3/001, 1/--, 5/10100, 1/0, 6/registerZ }


\paragraph{Syntax} \texttt{asaddh}\emph{coherency} [\emph{registerZ}] = \emph{registerT}

\paragraph{Description} The effective address is the value of the \emph{registerZ}.
This instruction implements atomic add-store on half word. The half word at effective address is added with the \emph{registerT}. The effective address must be half word aligned (multiple of 2).


\paragraph{Modifier(s)}
\noindent\begin{itemize}
\item coherency (cf. coherency in section \ref{par:modifier-coherency} on page \pageref{par:modifier-coherency})
\end{itemize}

\paragraph{Execution} {Reservation table \textsf{LSU2\_MEMW\_AUXR} (Section~\ref{sec:reservation-LSU2-MEMW-AUXR})}
~
{\begin{lstlisting}
stage ID:
  new coherency = _ZX_2(coherency);
stage RR:
  new address = GPR[registerZ];
stage E1:
  new argument2 = GPR[registerT];
stage E3:
  MEM.atomic_add(address, 0x3, coherency << 32, argument2.16[0], registerT);
\end{lstlisting}}
\newpage

\paragraph{Encoding} Format \textsf{LSU\_ASADDSB.O} (Section~\ref{sec:format-LSU-ASADDSB.O}), \verb|lsucode5 = 001|\\

\bitfields{ 1/P, 2/00, 2/-- --, 27/offset27, 1/1, 2/01, 3/111, 2/coherency, 6/registerT, 2/11, 3/001, 1/--, 5/10100, 1/0, 6/registerZ }


\paragraph{Syntax} \texttt{asaddh}\emph{coherency} \emph{offset27}[\emph{registerZ}] = \emph{registerT}

\paragraph{Description} The effective address is computed by adding the \emph{offset27} to the \emph{registerZ}.
This instruction implements atomic add-store on half word. The half word at effective address is added with the \emph{registerT}. The effective address must be half word aligned (multiple of 2).


\paragraph{Modifier(s)}
\noindent\begin{itemize}
\item coherency (cf. coherency in section \ref{par:modifier-coherency} on page \pageref{par:modifier-coherency})
\end{itemize}

\paragraph{Execution} {Reservation table \textsf{LSU2\_MEMW\_AUXR.X} (Section~\ref{sec:reservation-LSU2-MEMW-AUXR.X})}
~
{\begin{lstlisting}
stage ID:
  new offset = _SX_27(offset27);
  new coherency = _ZX_2(coherency);
stage RR:
  new base = GPR[registerZ];
  new address = base + offset;
stage E1:
  new argument2 = GPR[registerT];
stage E3:
  MEM.atomic_add(address, 0x3, coherency << 32, argument2.16[0], registerT);
\end{lstlisting}}
\newpage

\paragraph{Encoding} Format \textsf{LSU\_ASADDSB.D} (Section~\ref{sec:format-LSU-ASADDSB.D}), \verb|lsucode5 = 001|\\

\bitfields{ 1/P, 2/00, 2/-- --, 27/extend27, 1/1, 2/00, 2/-- --, 27/offset27, 1/1, 2/01, 3/111, 2/coherency, 6/registerT, 2/11, 3/001, 1/--, 5/10100, 1/0, 6/registerZ }


\paragraph{Syntax} \texttt{asaddh}\emph{coherency} \emph{extend27\_\discretionary{}{}{}offset27}[\emph{registerZ}] = \emph{registerT}

\paragraph{Description} The effective address is computed by adding the \emph{extend27\_\discretionary{}{}{}offset27} to the \emph{registerZ}.
This instruction implements atomic add-store on half word. The half word at effective address is added with the \emph{registerT}. The effective address must be half word aligned (multiple of 2).


\paragraph{Modifier(s)}
\noindent\begin{itemize}
\item coherency (cf. coherency in section \ref{par:modifier-coherency} on page \pageref{par:modifier-coherency})
\end{itemize}

\paragraph{Execution} {Reservation table \textsf{LSU2\_MEMW\_AUXR.Y} (Section~\ref{sec:reservation-LSU2-MEMW-AUXR.Y})}
~
{\begin{lstlisting}
stage ID:
  new offset = _SX_54(extend27_offset27);
  new coherency = _ZX_2(coherency);
stage RR:
  new base = GPR[registerZ];
  new address = base + offset;
stage E1:
  new argument2 = GPR[registerT];
stage E3:
  MEM.atomic_add(address, 0x3, coherency << 32, argument2.16[0], registerT);
\end{lstlisting}}
\newpage
\subsubsection[{\small ASADDW: Atomic Store Add Word}]{ASADDW\hfill Atomic Store Add Word}
\label{instr:ASADDW}\index{asaddw}
 
\paragraph{Encoding} Format \textsf{LSU\_ASADDSB} (Section~\ref{sec:format-LSU-ASADDSB}), \verb|lsucode5 = 010|\\

\bitfields{ 1/P, 2/01, 3/111, 2/coherency, 6/registerT, 2/11, 3/010, 1/--, 5/10100, 1/0, 6/registerZ }


\paragraph{Syntax} \texttt{asaddw}\emph{coherency} [\emph{registerZ}] = \emph{registerT}

\paragraph{Description} The effective address is the value of the \emph{registerZ}.
This instruction implements atomic add-store on word. The word at effective address is added with the \emph{registerT}. The effective address must be word aligned (multiple of 4).


\paragraph{Modifier(s)}
\noindent\begin{itemize}
\item coherency (cf. coherency in section \ref{par:modifier-coherency} on page \pageref{par:modifier-coherency})
\end{itemize}

\paragraph{Execution} {Reservation table \textsf{LSU2\_MEMW\_AUXR} (Section~\ref{sec:reservation-LSU2-MEMW-AUXR})}
~
{\begin{lstlisting}
stage ID:
  new coherency = _ZX_2(coherency);
stage RR:
  new address = GPR[registerZ];
stage E1:
  new argument2 = GPR[registerT];
stage E3:
  MEM.atomic_add(address, 0xF, coherency << 32, argument2.32[0], registerT);
\end{lstlisting}}
\newpage

\paragraph{Encoding} Format \textsf{LSU\_ASADDSB.O} (Section~\ref{sec:format-LSU-ASADDSB.O}), \verb|lsucode5 = 010|\\

\bitfields{ 1/P, 2/00, 2/-- --, 27/offset27, 1/1, 2/01, 3/111, 2/coherency, 6/registerT, 2/11, 3/010, 1/--, 5/10100, 1/0, 6/registerZ }


\paragraph{Syntax} \texttt{asaddw}\emph{coherency} \emph{offset27}[\emph{registerZ}] = \emph{registerT}

\paragraph{Description} The effective address is computed by adding the \emph{offset27} to the \emph{registerZ}.
This instruction implements atomic add-store on word. The word at effective address is added with the \emph{registerT}. The effective address must be word aligned (multiple of 4).


\paragraph{Modifier(s)}
\noindent\begin{itemize}
\item coherency (cf. coherency in section \ref{par:modifier-coherency} on page \pageref{par:modifier-coherency})
\end{itemize}

\paragraph{Execution} {Reservation table \textsf{LSU2\_MEMW\_AUXR.X} (Section~\ref{sec:reservation-LSU2-MEMW-AUXR.X})}
~
{\begin{lstlisting}
stage ID:
  new offset = _SX_27(offset27);
  new coherency = _ZX_2(coherency);
stage RR:
  new base = GPR[registerZ];
  new address = base + offset;
stage E1:
  new argument2 = GPR[registerT];
stage E3:
  MEM.atomic_add(address, 0xF, coherency << 32, argument2.32[0], registerT);
\end{lstlisting}}
\newpage

\paragraph{Encoding} Format \textsf{LSU\_ASADDSB.D} (Section~\ref{sec:format-LSU-ASADDSB.D}), \verb|lsucode5 = 010|\\

\bitfields{ 1/P, 2/00, 2/-- --, 27/extend27, 1/1, 2/00, 2/-- --, 27/offset27, 1/1, 2/01, 3/111, 2/coherency, 6/registerT, 2/11, 3/010, 1/--, 5/10100, 1/0, 6/registerZ }


\paragraph{Syntax} \texttt{asaddw}\emph{coherency} \emph{extend27\_\discretionary{}{}{}offset27}[\emph{registerZ}] = \emph{registerT}

\paragraph{Description} The effective address is computed by adding the \emph{extend27\_\discretionary{}{}{}offset27} to the \emph{registerZ}.
This instruction implements atomic add-store on word. The word at effective address is added with the \emph{registerT}. The effective address must be word aligned (multiple of 4).


\paragraph{Modifier(s)}
\noindent\begin{itemize}
\item coherency (cf. coherency in section \ref{par:modifier-coherency} on page \pageref{par:modifier-coherency})
\end{itemize}

\paragraph{Execution} {Reservation table \textsf{LSU2\_MEMW\_AUXR.Y} (Section~\ref{sec:reservation-LSU2-MEMW-AUXR.Y})}
~
{\begin{lstlisting}
stage ID:
  new offset = _SX_54(extend27_offset27);
  new coherency = _ZX_2(coherency);
stage RR:
  new base = GPR[registerZ];
  new address = base + offset;
stage E1:
  new argument2 = GPR[registerT];
stage E3:
  MEM.atomic_add(address, 0xF, coherency << 32, argument2.32[0], registerT);
\end{lstlisting}}
\newpage
\subsubsection[{\small ASANDB: Atomic Store AND Byte}]{ASANDB\hfill Atomic Store AND Byte}
\label{instr:ASANDB}\index{asandb}
 
\paragraph{Encoding} Format \textsf{LSU\_ASANDSB} (Section~\ref{sec:format-LSU-ASANDSB}), \verb|lsucode5 = 000|\\

\bitfields{ 1/P, 2/01, 3/111, 2/coherency, 6/registerT, 2/11, 3/000, 1/--, 5/10101, 1/0, 6/registerZ }


\paragraph{Syntax} \texttt{asandb}\emph{coherency} [\emph{registerZ}] = \emph{registerT}

\paragraph{Description} The effective address is the value of the \emph{registerZ}.
This instruction implements atomic AND-store on byte. The byte at effective address is ANDed with the \emph{registerT}.


\paragraph{Modifier(s)}
\noindent\begin{itemize}
\item coherency (cf. coherency in section \ref{par:modifier-coherency} on page \pageref{par:modifier-coherency})
\end{itemize}

\paragraph{Execution} {Reservation table \textsf{LSU2\_MEMW\_AUXR} (Section~\ref{sec:reservation-LSU2-MEMW-AUXR})}
~
{\begin{lstlisting}
stage ID:
  new coherency = _ZX_2(coherency);
stage RR:
  new address = GPR[registerZ];
stage E1:
  new argument2 = GPR[registerT];
stage E3:
  MEM.atomic_and(address, 0x1, coherency << 32, argument2.8[0], registerT);
\end{lstlisting}}
\newpage

\paragraph{Encoding} Format \textsf{LSU\_ASANDSB.O} (Section~\ref{sec:format-LSU-ASANDSB.O}), \verb|lsucode5 = 000|\\

\bitfields{ 1/P, 2/00, 2/-- --, 27/offset27, 1/1, 2/01, 3/111, 2/coherency, 6/registerT, 2/11, 3/000, 1/--, 5/10101, 1/0, 6/registerZ }


\paragraph{Syntax} \texttt{asandb}\emph{coherency} \emph{offset27}[\emph{registerZ}] = \emph{registerT}

\paragraph{Description} The effective address is computed by adding the \emph{offset27} to the \emph{registerZ}.
This instruction implements atomic AND-store on byte. The byte at effective address is ANDed with the \emph{registerT}.


\paragraph{Modifier(s)}
\noindent\begin{itemize}
\item coherency (cf. coherency in section \ref{par:modifier-coherency} on page \pageref{par:modifier-coherency})
\end{itemize}

\paragraph{Execution} {Reservation table \textsf{LSU2\_MEMW\_AUXR.X} (Section~\ref{sec:reservation-LSU2-MEMW-AUXR.X})}
~
{\begin{lstlisting}
stage ID:
  new offset = _SX_27(offset27);
  new coherency = _ZX_2(coherency);
stage RR:
  new base = GPR[registerZ];
  new address = base + offset;
stage E1:
  new argument2 = GPR[registerT];
stage E3:
  MEM.atomic_and(address, 0x1, coherency << 32, argument2.8[0], registerT);
\end{lstlisting}}
\newpage

\paragraph{Encoding} Format \textsf{LSU\_ASANDSB.D} (Section~\ref{sec:format-LSU-ASANDSB.D}), \verb|lsucode5 = 000|\\

\bitfields{ 1/P, 2/00, 2/-- --, 27/extend27, 1/1, 2/00, 2/-- --, 27/offset27, 1/1, 2/01, 3/111, 2/coherency, 6/registerT, 2/11, 3/000, 1/--, 5/10101, 1/0, 6/registerZ }


\paragraph{Syntax} \texttt{asandb}\emph{coherency} \emph{extend27\_\discretionary{}{}{}offset27}[\emph{registerZ}] = \emph{registerT}

\paragraph{Description} The effective address is computed by adding the \emph{extend27\_\discretionary{}{}{}offset27} to the \emph{registerZ}.
This instruction implements atomic AND-store on byte. The byte at effective address is ANDed with the \emph{registerT}.


\paragraph{Modifier(s)}
\noindent\begin{itemize}
\item coherency (cf. coherency in section \ref{par:modifier-coherency} on page \pageref{par:modifier-coherency})
\end{itemize}

\paragraph{Execution} {Reservation table \textsf{LSU2\_MEMW\_AUXR.Y} (Section~\ref{sec:reservation-LSU2-MEMW-AUXR.Y})}
~
{\begin{lstlisting}
stage ID:
  new offset = _SX_54(extend27_offset27);
  new coherency = _ZX_2(coherency);
stage RR:
  new base = GPR[registerZ];
  new address = base + offset;
stage E1:
  new argument2 = GPR[registerT];
stage E3:
  MEM.atomic_and(address, 0x1, coherency << 32, argument2.8[0], registerT);
\end{lstlisting}}
\newpage
\subsubsection[{\small ASANDD: Atomic Store AND Double Word}]{ASANDD\hfill Atomic Store AND Double Word}
\label{instr:ASANDD}\index{asandd}
 
\paragraph{Encoding} Format \textsf{LSU\_ASANDSB} (Section~\ref{sec:format-LSU-ASANDSB}), \verb|lsucode5 = 011|\\

\bitfields{ 1/P, 2/01, 3/111, 2/coherency, 6/registerT, 2/11, 3/011, 1/--, 5/10101, 1/0, 6/registerZ }


\paragraph{Syntax} \texttt{asandd}\emph{coherency} [\emph{registerZ}] = \emph{registerT}

\paragraph{Description} The effective address is the value of the \emph{registerZ}.
This instruction implements atomic AND-store on double word. The double word at effective address is ANDed with the \emph{registerT}. The effective address must be double word aligned (multiple of 8).


\paragraph{Modifier(s)}
\noindent\begin{itemize}
\item coherency (cf. coherency in section \ref{par:modifier-coherency} on page \pageref{par:modifier-coherency})
\end{itemize}

\paragraph{Execution} {Reservation table \textsf{LSU2\_MEMW\_AUXR} (Section~\ref{sec:reservation-LSU2-MEMW-AUXR})}
~
{\begin{lstlisting}
stage ID:
  new coherency = _ZX_2(coherency);
stage RR:
  new address = GPR[registerZ];
stage E1:
  new argument2 = GPR[registerT];
stage E3:
  MEM.atomic_and(address, 0xFF, coherency << 32, argument2.64[0], registerT);
\end{lstlisting}}
\newpage

\paragraph{Encoding} Format \textsf{LSU\_ASANDSB.O} (Section~\ref{sec:format-LSU-ASANDSB.O}), \verb|lsucode5 = 011|\\

\bitfields{ 1/P, 2/00, 2/-- --, 27/offset27, 1/1, 2/01, 3/111, 2/coherency, 6/registerT, 2/11, 3/011, 1/--, 5/10101, 1/0, 6/registerZ }


\paragraph{Syntax} \texttt{asandd}\emph{coherency} \emph{offset27}[\emph{registerZ}] = \emph{registerT}

\paragraph{Description} The effective address is computed by adding the \emph{offset27} to the \emph{registerZ}.
This instruction implements atomic AND-store on double word. The double word at effective address is ANDed with the \emph{registerT}. The effective address must be double word aligned (multiple of 8).


\paragraph{Modifier(s)}
\noindent\begin{itemize}
\item coherency (cf. coherency in section \ref{par:modifier-coherency} on page \pageref{par:modifier-coherency})
\end{itemize}

\paragraph{Execution} {Reservation table \textsf{LSU2\_MEMW\_AUXR.X} (Section~\ref{sec:reservation-LSU2-MEMW-AUXR.X})}
~
{\begin{lstlisting}
stage ID:
  new offset = _SX_27(offset27);
  new coherency = _ZX_2(coherency);
stage RR:
  new base = GPR[registerZ];
  new address = base + offset;
stage E1:
  new argument2 = GPR[registerT];
stage E3:
  MEM.atomic_and(address, 0xFF, coherency << 32, argument2.64[0], registerT);
\end{lstlisting}}
\newpage

\paragraph{Encoding} Format \textsf{LSU\_ASANDSB.D} (Section~\ref{sec:format-LSU-ASANDSB.D}), \verb|lsucode5 = 011|\\

\bitfields{ 1/P, 2/00, 2/-- --, 27/extend27, 1/1, 2/00, 2/-- --, 27/offset27, 1/1, 2/01, 3/111, 2/coherency, 6/registerT, 2/11, 3/011, 1/--, 5/10101, 1/0, 6/registerZ }


\paragraph{Syntax} \texttt{asandd}\emph{coherency} \emph{extend27\_\discretionary{}{}{}offset27}[\emph{registerZ}] = \emph{registerT}

\paragraph{Description} The effective address is computed by adding the \emph{extend27\_\discretionary{}{}{}offset27} to the \emph{registerZ}.
This instruction implements atomic AND-store on double word. The double word at effective address is ANDed with the \emph{registerT}. The effective address must be double word aligned (multiple of 8).


\paragraph{Modifier(s)}
\noindent\begin{itemize}
\item coherency (cf. coherency in section \ref{par:modifier-coherency} on page \pageref{par:modifier-coherency})
\end{itemize}

\paragraph{Execution} {Reservation table \textsf{LSU2\_MEMW\_AUXR.Y} (Section~\ref{sec:reservation-LSU2-MEMW-AUXR.Y})}
~
{\begin{lstlisting}
stage ID:
  new offset = _SX_54(extend27_offset27);
  new coherency = _ZX_2(coherency);
stage RR:
  new base = GPR[registerZ];
  new address = base + offset;
stage E1:
  new argument2 = GPR[registerT];
stage E3:
  MEM.atomic_and(address, 0xFF, coherency << 32, argument2.64[0], registerT);
\end{lstlisting}}
\newpage
\subsubsection[{\small ASANDH: Atomic Store AND Half Word}]{ASANDH\hfill Atomic Store AND Half Word}
\label{instr:ASANDH}\index{asandh}
 
\paragraph{Encoding} Format \textsf{LSU\_ASANDSB} (Section~\ref{sec:format-LSU-ASANDSB}), \verb|lsucode5 = 001|\\

\bitfields{ 1/P, 2/01, 3/111, 2/coherency, 6/registerT, 2/11, 3/001, 1/--, 5/10101, 1/0, 6/registerZ }


\paragraph{Syntax} \texttt{asandh}\emph{coherency} [\emph{registerZ}] = \emph{registerT}

\paragraph{Description} The effective address is the value of the \emph{registerZ}.
This instruction implements atomic AND-store on half word. The half word at effective address is ANDed with the \emph{registerT}. The effective address must be half word aligned (multiple of 2).


\paragraph{Modifier(s)}
\noindent\begin{itemize}
\item coherency (cf. coherency in section \ref{par:modifier-coherency} on page \pageref{par:modifier-coherency})
\end{itemize}

\paragraph{Execution} {Reservation table \textsf{LSU2\_MEMW\_AUXR} (Section~\ref{sec:reservation-LSU2-MEMW-AUXR})}
~
{\begin{lstlisting}
stage ID:
  new coherency = _ZX_2(coherency);
stage RR:
  new address = GPR[registerZ];
stage E1:
  new argument2 = GPR[registerT];
stage E3:
  MEM.atomic_and(address, 0x3, coherency << 32, argument2.16[0], registerT);
\end{lstlisting}}
\newpage

\paragraph{Encoding} Format \textsf{LSU\_ASANDSB.O} (Section~\ref{sec:format-LSU-ASANDSB.O}), \verb|lsucode5 = 001|\\

\bitfields{ 1/P, 2/00, 2/-- --, 27/offset27, 1/1, 2/01, 3/111, 2/coherency, 6/registerT, 2/11, 3/001, 1/--, 5/10101, 1/0, 6/registerZ }


\paragraph{Syntax} \texttt{asandh}\emph{coherency} \emph{offset27}[\emph{registerZ}] = \emph{registerT}

\paragraph{Description} The effective address is computed by adding the \emph{offset27} to the \emph{registerZ}.
This instruction implements atomic AND-store on half word. The half word at effective address is ANDed with the \emph{registerT}. The effective address must be half word aligned (multiple of 2).


\paragraph{Modifier(s)}
\noindent\begin{itemize}
\item coherency (cf. coherency in section \ref{par:modifier-coherency} on page \pageref{par:modifier-coherency})
\end{itemize}

\paragraph{Execution} {Reservation table \textsf{LSU2\_MEMW\_AUXR.X} (Section~\ref{sec:reservation-LSU2-MEMW-AUXR.X})}
~
{\begin{lstlisting}
stage ID:
  new offset = _SX_27(offset27);
  new coherency = _ZX_2(coherency);
stage RR:
  new base = GPR[registerZ];
  new address = base + offset;
stage E1:
  new argument2 = GPR[registerT];
stage E3:
  MEM.atomic_and(address, 0x3, coherency << 32, argument2.16[0], registerT);
\end{lstlisting}}
\newpage

\paragraph{Encoding} Format \textsf{LSU\_ASANDSB.D} (Section~\ref{sec:format-LSU-ASANDSB.D}), \verb|lsucode5 = 001|\\

\bitfields{ 1/P, 2/00, 2/-- --, 27/extend27, 1/1, 2/00, 2/-- --, 27/offset27, 1/1, 2/01, 3/111, 2/coherency, 6/registerT, 2/11, 3/001, 1/--, 5/10101, 1/0, 6/registerZ }


\paragraph{Syntax} \texttt{asandh}\emph{coherency} \emph{extend27\_\discretionary{}{}{}offset27}[\emph{registerZ}] = \emph{registerT}

\paragraph{Description} The effective address is computed by adding the \emph{extend27\_\discretionary{}{}{}offset27} to the \emph{registerZ}.
This instruction implements atomic AND-store on half word. The half word at effective address is ANDed with the \emph{registerT}. The effective address must be half word aligned (multiple of 2).


\paragraph{Modifier(s)}
\noindent\begin{itemize}
\item coherency (cf. coherency in section \ref{par:modifier-coherency} on page \pageref{par:modifier-coherency})
\end{itemize}

\paragraph{Execution} {Reservation table \textsf{LSU2\_MEMW\_AUXR.Y} (Section~\ref{sec:reservation-LSU2-MEMW-AUXR.Y})}
~
{\begin{lstlisting}
stage ID:
  new offset = _SX_54(extend27_offset27);
  new coherency = _ZX_2(coherency);
stage RR:
  new base = GPR[registerZ];
  new address = base + offset;
stage E1:
  new argument2 = GPR[registerT];
stage E3:
  MEM.atomic_and(address, 0x3, coherency << 32, argument2.16[0], registerT);
\end{lstlisting}}
\newpage
\subsubsection[{\small ASANDW: Atomic Store AND Word}]{ASANDW\hfill Atomic Store AND Word}
\label{instr:ASANDW}\index{asandw}
 
\paragraph{Encoding} Format \textsf{LSU\_ASANDSB} (Section~\ref{sec:format-LSU-ASANDSB}), \verb|lsucode5 = 010|\\

\bitfields{ 1/P, 2/01, 3/111, 2/coherency, 6/registerT, 2/11, 3/010, 1/--, 5/10101, 1/0, 6/registerZ }


\paragraph{Syntax} \texttt{asandw}\emph{coherency} [\emph{registerZ}] = \emph{registerT}

\paragraph{Description} The effective address is the value of the \emph{registerZ}.
This instruction implements atomic AND-store on word. The word at effective address is ANDed with the \emph{registerT}. The effective address must be word aligned (multiple of 4).


\paragraph{Modifier(s)}
\noindent\begin{itemize}
\item coherency (cf. coherency in section \ref{par:modifier-coherency} on page \pageref{par:modifier-coherency})
\end{itemize}

\paragraph{Execution} {Reservation table \textsf{LSU2\_MEMW\_AUXR} (Section~\ref{sec:reservation-LSU2-MEMW-AUXR})}
~
{\begin{lstlisting}
stage ID:
  new coherency = _ZX_2(coherency);
stage RR:
  new address = GPR[registerZ];
stage E1:
  new argument2 = GPR[registerT];
stage E3:
  MEM.atomic_and(address, 0xF, coherency << 32, argument2.32[0], registerT);
\end{lstlisting}}
\newpage

\paragraph{Encoding} Format \textsf{LSU\_ASANDSB.O} (Section~\ref{sec:format-LSU-ASANDSB.O}), \verb|lsucode5 = 010|\\

\bitfields{ 1/P, 2/00, 2/-- --, 27/offset27, 1/1, 2/01, 3/111, 2/coherency, 6/registerT, 2/11, 3/010, 1/--, 5/10101, 1/0, 6/registerZ }


\paragraph{Syntax} \texttt{asandw}\emph{coherency} \emph{offset27}[\emph{registerZ}] = \emph{registerT}

\paragraph{Description} The effective address is computed by adding the \emph{offset27} to the \emph{registerZ}.
This instruction implements atomic AND-store on word. The word at effective address is ANDed with the \emph{registerT}. The effective address must be word aligned (multiple of 4).


\paragraph{Modifier(s)}
\noindent\begin{itemize}
\item coherency (cf. coherency in section \ref{par:modifier-coherency} on page \pageref{par:modifier-coherency})
\end{itemize}

\paragraph{Execution} {Reservation table \textsf{LSU2\_MEMW\_AUXR.X} (Section~\ref{sec:reservation-LSU2-MEMW-AUXR.X})}
~
{\begin{lstlisting}
stage ID:
  new offset = _SX_27(offset27);
  new coherency = _ZX_2(coherency);
stage RR:
  new base = GPR[registerZ];
  new address = base + offset;
stage E1:
  new argument2 = GPR[registerT];
stage E3:
  MEM.atomic_and(address, 0xF, coherency << 32, argument2.32[0], registerT);
\end{lstlisting}}
\newpage

\paragraph{Encoding} Format \textsf{LSU\_ASANDSB.D} (Section~\ref{sec:format-LSU-ASANDSB.D}), \verb|lsucode5 = 010|\\

\bitfields{ 1/P, 2/00, 2/-- --, 27/extend27, 1/1, 2/00, 2/-- --, 27/offset27, 1/1, 2/01, 3/111, 2/coherency, 6/registerT, 2/11, 3/010, 1/--, 5/10101, 1/0, 6/registerZ }


\paragraph{Syntax} \texttt{asandw}\emph{coherency} \emph{extend27\_\discretionary{}{}{}offset27}[\emph{registerZ}] = \emph{registerT}

\paragraph{Description} The effective address is computed by adding the \emph{extend27\_\discretionary{}{}{}offset27} to the \emph{registerZ}.
This instruction implements atomic AND-store on word. The word at effective address is ANDed with the \emph{registerT}. The effective address must be word aligned (multiple of 4).


\paragraph{Modifier(s)}
\noindent\begin{itemize}
\item coherency (cf. coherency in section \ref{par:modifier-coherency} on page \pageref{par:modifier-coherency})
\end{itemize}

\paragraph{Execution} {Reservation table \textsf{LSU2\_MEMW\_AUXR.Y} (Section~\ref{sec:reservation-LSU2-MEMW-AUXR.Y})}
~
{\begin{lstlisting}
stage ID:
  new offset = _SX_54(extend27_offset27);
  new coherency = _ZX_2(coherency);
stage RR:
  new base = GPR[registerZ];
  new address = base + offset;
stage E1:
  new argument2 = GPR[registerT];
stage E3:
  MEM.atomic_and(address, 0xF, coherency << 32, argument2.32[0], registerT);
\end{lstlisting}}
\newpage
\subsubsection[{\small ASB: Atomic Store Byte}]{ASB\hfill Atomic Store Byte}
\label{instr:ASB}\index{asb}
 
\paragraph{Encoding} Format \textsf{LSU\_ASSB} (Section~\ref{sec:format-LSU-ASSB}), \verb|lsucode5 = 000|\\

\bitfields{ 1/P, 2/01, 3/111, 2/coherency, 6/registerT, 2/11, 3/000, 1/--, 5/10000, 1/0, 6/registerZ }


\paragraph{Syntax} \texttt{asb}\emph{coherency} [\emph{registerZ}] = \emph{registerT}

\paragraph{Description} The effective address is the value of the \emph{registerZ}.
The \emph{registerT} is atomically stored to the memory byte starting at the effective address. This instruction is provided to directly operate on the last level of the memory hierarchy, which is typically a processor DDR or a cluster TCM, irrespective of the MMU DCP (Data Cache Policy) settings.


\paragraph{Modifier(s)}
\noindent\begin{itemize}
\item coherency (cf. coherency in section \ref{par:modifier-coherency} on page \pageref{par:modifier-coherency})
\end{itemize}

\paragraph{Execution} {Reservation table \textsf{LSU\_MEMW\_AUXR} (Section~\ref{sec:reservation-LSU-MEMW-AUXR})}
~
{\begin{lstlisting}
stage ID:
  new coherency = _ZX_2(coherency);
stage RR:
  new address = GPR[registerZ];
stage E1:
  new argument2 = GPR[registerT];
stage E3:
  MEM.atomic_store(address, 0x1, coherency << 32, argument2, registerT);
\end{lstlisting}}
\newpage

\paragraph{Encoding} Format \textsf{LSU\_ASSB.O} (Section~\ref{sec:format-LSU-ASSB.O}), \verb|lsucode5 = 000|\\

\bitfields{ 1/P, 2/00, 2/-- --, 27/offset27, 1/1, 2/01, 3/111, 2/coherency, 6/registerT, 2/11, 3/000, 1/--, 5/10000, 1/0, 6/registerZ }


\paragraph{Syntax} \texttt{asb}\emph{coherency} \emph{offset27}[\emph{registerZ}] = \emph{registerT}

\paragraph{Description} The effective address is computed by adding the \emph{offset27} to the \emph{registerZ}.
The \emph{registerT} is atomically stored to the memory byte starting at the effective address. This instruction is provided to directly operate on the last level of the memory hierarchy, which is typically a processor DDR or a cluster TCM, irrespective of the MMU DCP (Data Cache Policy) settings.


\paragraph{Modifier(s)}
\noindent\begin{itemize}
\item coherency (cf. coherency in section \ref{par:modifier-coherency} on page \pageref{par:modifier-coherency})
\end{itemize}

\paragraph{Execution} {Reservation table \textsf{LSU\_MEMW\_AUXR.X} (Section~\ref{sec:reservation-LSU-MEMW-AUXR.X})}
~
{\begin{lstlisting}
stage ID:
  new offset = _SX_27(offset27);
  new coherency = _ZX_2(coherency);
stage RR:
  new base = GPR[registerZ];
  new address = base + offset;
stage E1:
  new argument2 = GPR[registerT];
stage E3:
  MEM.atomic_store(address, 0x1, coherency << 32, argument2, registerT);
\end{lstlisting}}
\newpage

\paragraph{Encoding} Format \textsf{LSU\_ASSB.D} (Section~\ref{sec:format-LSU-ASSB.D}), \verb|lsucode5 = 000|\\

\bitfields{ 1/P, 2/00, 2/-- --, 27/extend27, 1/1, 2/00, 2/-- --, 27/offset27, 1/1, 2/01, 3/111, 2/coherency, 6/registerT, 2/11, 3/000, 1/--, 5/10000, 1/0, 6/registerZ }


\paragraph{Syntax} \texttt{asb}\emph{coherency} \emph{extend27\_\discretionary{}{}{}offset27}[\emph{registerZ}] = \emph{registerT}

\paragraph{Description} The effective address is computed by adding the \emph{extend27\_\discretionary{}{}{}offset27} to the \emph{registerZ}.
The \emph{registerT} is atomically stored to the memory byte starting at the effective address. This instruction is provided to directly operate on the last level of the memory hierarchy, which is typically a processor DDR or a cluster TCM, irrespective of the MMU DCP (Data Cache Policy) settings.


\paragraph{Modifier(s)}
\noindent\begin{itemize}
\item coherency (cf. coherency in section \ref{par:modifier-coherency} on page \pageref{par:modifier-coherency})
\end{itemize}

\paragraph{Execution} {Reservation table \textsf{LSU\_MEMW\_AUXR.Y} (Section~\ref{sec:reservation-LSU-MEMW-AUXR.Y})}
~
{\begin{lstlisting}
stage ID:
  new offset = _SX_54(extend27_offset27);
  new coherency = _ZX_2(coherency);
stage RR:
  new base = GPR[registerZ];
  new address = base + offset;
stage E1:
  new argument2 = GPR[registerT];
stage E3:
  MEM.atomic_store(address, 0x1, coherency << 32, argument2, registerT);
\end{lstlisting}}
\newpage
\subsubsection[{\small ASD: Atomic Store Double Word}]{ASD\hfill Atomic Store Double Word}
\label{instr:ASD}\index{asd}
 
\paragraph{Encoding} Format \textsf{LSU\_ASSB} (Section~\ref{sec:format-LSU-ASSB}), \verb|lsucode5 = 011|\\

\bitfields{ 1/P, 2/01, 3/111, 2/coherency, 6/registerT, 2/11, 3/011, 1/--, 5/10000, 1/0, 6/registerZ }


\paragraph{Syntax} \texttt{asd}\emph{coherency} [\emph{registerZ}] = \emph{registerT}

\paragraph{Description} The effective address is the value of the \emph{registerZ}.
The \emph{registerT} is atomically stored to the memory double word starting at the effective address. This instruction is provided to directly operate on the last level of the memory hierarchy, which is typically a processor DDR or a cluster TCM, irrespective of the MMU DCP (Data Cache Policy) settings. The effective address must be double word aligned (multiple of 8).


\paragraph{Modifier(s)}
\noindent\begin{itemize}
\item coherency (cf. coherency in section \ref{par:modifier-coherency} on page \pageref{par:modifier-coherency})
\end{itemize}

\paragraph{Execution} {Reservation table \textsf{LSU\_MEMW\_AUXR} (Section~\ref{sec:reservation-LSU-MEMW-AUXR})}
~
{\begin{lstlisting}
stage ID:
  new coherency = _ZX_2(coherency);
stage RR:
  new address = GPR[registerZ];
stage E1:
  new argument2 = GPR[registerT];
stage E3:
  MEM.atomic_store(address, 0xFF, coherency << 32, argument2, registerT);
\end{lstlisting}}
\newpage

\paragraph{Encoding} Format \textsf{LSU\_ASSB.O} (Section~\ref{sec:format-LSU-ASSB.O}), \verb|lsucode5 = 011|\\

\bitfields{ 1/P, 2/00, 2/-- --, 27/offset27, 1/1, 2/01, 3/111, 2/coherency, 6/registerT, 2/11, 3/011, 1/--, 5/10000, 1/0, 6/registerZ }


\paragraph{Syntax} \texttt{asd}\emph{coherency} \emph{offset27}[\emph{registerZ}] = \emph{registerT}

\paragraph{Description} The effective address is computed by adding the \emph{offset27} to the \emph{registerZ}.
The \emph{registerT} is atomically stored to the memory double word starting at the effective address. This instruction is provided to directly operate on the last level of the memory hierarchy, which is typically a processor DDR or a cluster TCM, irrespective of the MMU DCP (Data Cache Policy) settings. The effective address must be double word aligned (multiple of 8).


\paragraph{Modifier(s)}
\noindent\begin{itemize}
\item coherency (cf. coherency in section \ref{par:modifier-coherency} on page \pageref{par:modifier-coherency})
\end{itemize}

\paragraph{Execution} {Reservation table \textsf{LSU\_MEMW\_AUXR.X} (Section~\ref{sec:reservation-LSU-MEMW-AUXR.X})}
~
{\begin{lstlisting}
stage ID:
  new offset = _SX_27(offset27);
  new coherency = _ZX_2(coherency);
stage RR:
  new base = GPR[registerZ];
  new address = base + offset;
stage E1:
  new argument2 = GPR[registerT];
stage E3:
  MEM.atomic_store(address, 0xFF, coherency << 32, argument2, registerT);
\end{lstlisting}}
\newpage

\paragraph{Encoding} Format \textsf{LSU\_ASSB.D} (Section~\ref{sec:format-LSU-ASSB.D}), \verb|lsucode5 = 011|\\

\bitfields{ 1/P, 2/00, 2/-- --, 27/extend27, 1/1, 2/00, 2/-- --, 27/offset27, 1/1, 2/01, 3/111, 2/coherency, 6/registerT, 2/11, 3/011, 1/--, 5/10000, 1/0, 6/registerZ }


\paragraph{Syntax} \texttt{asd}\emph{coherency} \emph{extend27\_\discretionary{}{}{}offset27}[\emph{registerZ}] = \emph{registerT}

\paragraph{Description} The effective address is computed by adding the \emph{extend27\_\discretionary{}{}{}offset27} to the \emph{registerZ}.
The \emph{registerT} is atomically stored to the memory double word starting at the effective address. This instruction is provided to directly operate on the last level of the memory hierarchy, which is typically a processor DDR or a cluster TCM, irrespective of the MMU DCP (Data Cache Policy) settings. The effective address must be double word aligned (multiple of 8).


\paragraph{Modifier(s)}
\noindent\begin{itemize}
\item coherency (cf. coherency in section \ref{par:modifier-coherency} on page \pageref{par:modifier-coherency})
\end{itemize}

\paragraph{Execution} {Reservation table \textsf{LSU\_MEMW\_AUXR.Y} (Section~\ref{sec:reservation-LSU-MEMW-AUXR.Y})}
~
{\begin{lstlisting}
stage ID:
  new offset = _SX_54(extend27_offset27);
  new coherency = _ZX_2(coherency);
stage RR:
  new base = GPR[registerZ];
  new address = base + offset;
stage E1:
  new argument2 = GPR[registerT];
stage E3:
  MEM.atomic_store(address, 0xFF, coherency << 32, argument2, registerT);
\end{lstlisting}}
\newpage
\subsubsection[{\small ASDUSB: Atomic Store Decrement Unsigned Saturating Byte}]{ASDUSB\hfill Atomic Store Decrement Unsigned Saturating Byte}
\label{instr:ASDUSB}\index{asdusb}
 
\paragraph{Encoding} Format \textsf{LSU\_ASDUSSB} (Section~\ref{sec:format-LSU-ASDUSSB}), \verb|lsucode5 = 000|\\

\bitfields{ 1/P, 2/01, 3/111, 2/coherency, 6/registerT, 2/11, 3/000, 1/--, 5/11110, 1/0, 6/registerZ }


\paragraph{Syntax} \texttt{asdusb}\emph{coherency} [\emph{registerZ}] = \emph{registerT}

\paragraph{Description} The effective address is the value of the \emph{registerZ}.
This instruction implements atomic DUS-store on signed byte (Decrement Unsigned Saturating). The byte at effective address is DUSed with the \emph{registerT}.


\paragraph{Modifier(s)}
\noindent\begin{itemize}
\item coherency (cf. coherency in section \ref{par:modifier-coherency} on page \pageref{par:modifier-coherency})
\end{itemize}

\paragraph{Execution} {Reservation table \textsf{LSU2\_MEMW\_AUXR} (Section~\ref{sec:reservation-LSU2-MEMW-AUXR})}
~
{\begin{lstlisting}
stage ID:
  new coherency = _ZX_2(coherency);
stage RR:
  new address = GPR[registerZ];
stage E1:
  new argument2 = GPR[registerT];
stage E3:
  MEM.atomic_dus(address, 0x1, coherency << 32, argument2.8[0], registerT);
\end{lstlisting}}
\newpage

\paragraph{Encoding} Format \textsf{LSU\_ASDUSSB.O} (Section~\ref{sec:format-LSU-ASDUSSB.O}), \verb|lsucode5 = 000|\\

\bitfields{ 1/P, 2/00, 2/-- --, 27/offset27, 1/1, 2/01, 3/111, 2/coherency, 6/registerT, 2/11, 3/000, 1/--, 5/11110, 1/0, 6/registerZ }


\paragraph{Syntax} \texttt{asdusb}\emph{coherency} \emph{offset27}[\emph{registerZ}] = \emph{registerT}

\paragraph{Description} The effective address is computed by adding the \emph{offset27} to the \emph{registerZ}.
This instruction implements atomic DUS-store on signed byte (Decrement Unsigned Saturating). The byte at effective address is DUSed with the \emph{registerT}.


\paragraph{Modifier(s)}
\noindent\begin{itemize}
\item coherency (cf. coherency in section \ref{par:modifier-coherency} on page \pageref{par:modifier-coherency})
\end{itemize}

\paragraph{Execution} {Reservation table \textsf{LSU2\_MEMW\_AUXR.X} (Section~\ref{sec:reservation-LSU2-MEMW-AUXR.X})}
~
{\begin{lstlisting}
stage ID:
  new offset = _SX_27(offset27);
  new coherency = _ZX_2(coherency);
stage RR:
  new base = GPR[registerZ];
  new address = base + offset;
stage E1:
  new argument2 = GPR[registerT];
stage E3:
  MEM.atomic_dus(address, 0x1, coherency << 32, argument2.8[0], registerT);
\end{lstlisting}}
\newpage

\paragraph{Encoding} Format \textsf{LSU\_ASDUSSB.D} (Section~\ref{sec:format-LSU-ASDUSSB.D}), \verb|lsucode5 = 000|\\

\bitfields{ 1/P, 2/00, 2/-- --, 27/extend27, 1/1, 2/00, 2/-- --, 27/offset27, 1/1, 2/01, 3/111, 2/coherency, 6/registerT, 2/11, 3/000, 1/--, 5/11110, 1/0, 6/registerZ }


\paragraph{Syntax} \texttt{asdusb}\emph{coherency} \emph{extend27\_\discretionary{}{}{}offset27}[\emph{registerZ}] = \emph{registerT}

\paragraph{Description} The effective address is computed by adding the \emph{extend27\_\discretionary{}{}{}offset27} to the \emph{registerZ}.
This instruction implements atomic DUS-store on signed byte (Decrement Unsigned Saturating). The byte at effective address is DUSed with the \emph{registerT}.


\paragraph{Modifier(s)}
\noindent\begin{itemize}
\item coherency (cf. coherency in section \ref{par:modifier-coherency} on page \pageref{par:modifier-coherency})
\end{itemize}

\paragraph{Execution} {Reservation table \textsf{LSU2\_MEMW\_AUXR.Y} (Section~\ref{sec:reservation-LSU2-MEMW-AUXR.Y})}
~
{\begin{lstlisting}
stage ID:
  new offset = _SX_54(extend27_offset27);
  new coherency = _ZX_2(coherency);
stage RR:
  new base = GPR[registerZ];
  new address = base + offset;
stage E1:
  new argument2 = GPR[registerT];
stage E3:
  MEM.atomic_dus(address, 0x1, coherency << 32, argument2.8[0], registerT);
\end{lstlisting}}
\newpage
\subsubsection[{\small ASDUSD: Atomic Store DUS Double Word}]{ASDUSD\hfill Atomic Store DUS Double Word}
\label{instr:ASDUSD}\index{asdusd}
 
\paragraph{Encoding} Format \textsf{LSU\_ASDUSSB} (Section~\ref{sec:format-LSU-ASDUSSB}), \verb|lsucode5 = 011|\\

\bitfields{ 1/P, 2/01, 3/111, 2/coherency, 6/registerT, 2/11, 3/011, 1/--, 5/11110, 1/0, 6/registerZ }


\paragraph{Syntax} \texttt{asdusd}\emph{coherency} [\emph{registerZ}] = \emph{registerT}

\paragraph{Description} The effective address is the value of the \emph{registerZ}.
This instruction implements atomic DUS-store on signed double word (Decrement Unsigned Saturating). The double word at effective address is DUSed with the \emph{registerT}. The effective address must be double word aligned (multiple of 8).


\paragraph{Modifier(s)}
\noindent\begin{itemize}
\item coherency (cf. coherency in section \ref{par:modifier-coherency} on page \pageref{par:modifier-coherency})
\end{itemize}

\paragraph{Execution} {Reservation table \textsf{LSU2\_MEMW\_AUXR} (Section~\ref{sec:reservation-LSU2-MEMW-AUXR})}
~
{\begin{lstlisting}
stage ID:
  new coherency = _ZX_2(coherency);
stage RR:
  new address = GPR[registerZ];
stage E1:
  new argument2 = GPR[registerT];
stage E3:
  MEM.atomic_dus(address, 0xFF, coherency << 32, argument2.64[0], registerT);
\end{lstlisting}}
\newpage

\paragraph{Encoding} Format \textsf{LSU\_ASDUSSB.O} (Section~\ref{sec:format-LSU-ASDUSSB.O}), \verb|lsucode5 = 011|\\

\bitfields{ 1/P, 2/00, 2/-- --, 27/offset27, 1/1, 2/01, 3/111, 2/coherency, 6/registerT, 2/11, 3/011, 1/--, 5/11110, 1/0, 6/registerZ }


\paragraph{Syntax} \texttt{asdusd}\emph{coherency} \emph{offset27}[\emph{registerZ}] = \emph{registerT}

\paragraph{Description} The effective address is computed by adding the \emph{offset27} to the \emph{registerZ}.
This instruction implements atomic DUS-store on signed double word (Decrement Unsigned Saturating). The double word at effective address is DUSed with the \emph{registerT}. The effective address must be double word aligned (multiple of 8).


\paragraph{Modifier(s)}
\noindent\begin{itemize}
\item coherency (cf. coherency in section \ref{par:modifier-coherency} on page \pageref{par:modifier-coherency})
\end{itemize}

\paragraph{Execution} {Reservation table \textsf{LSU2\_MEMW\_AUXR.X} (Section~\ref{sec:reservation-LSU2-MEMW-AUXR.X})}
~
{\begin{lstlisting}
stage ID:
  new offset = _SX_27(offset27);
  new coherency = _ZX_2(coherency);
stage RR:
  new base = GPR[registerZ];
  new address = base + offset;
stage E1:
  new argument2 = GPR[registerT];
stage E3:
  MEM.atomic_dus(address, 0xFF, coherency << 32, argument2.64[0], registerT);
\end{lstlisting}}
\newpage

\paragraph{Encoding} Format \textsf{LSU\_ASDUSSB.D} (Section~\ref{sec:format-LSU-ASDUSSB.D}), \verb|lsucode5 = 011|\\

\bitfields{ 1/P, 2/00, 2/-- --, 27/extend27, 1/1, 2/00, 2/-- --, 27/offset27, 1/1, 2/01, 3/111, 2/coherency, 6/registerT, 2/11, 3/011, 1/--, 5/11110, 1/0, 6/registerZ }


\paragraph{Syntax} \texttt{asdusd}\emph{coherency} \emph{extend27\_\discretionary{}{}{}offset27}[\emph{registerZ}] = \emph{registerT}

\paragraph{Description} The effective address is computed by adding the \emph{extend27\_\discretionary{}{}{}offset27} to the \emph{registerZ}.
This instruction implements atomic DUS-store on signed double word (Decrement Unsigned Saturating). The double word at effective address is DUSed with the \emph{registerT}. The effective address must be double word aligned (multiple of 8).


\paragraph{Modifier(s)}
\noindent\begin{itemize}
\item coherency (cf. coherency in section \ref{par:modifier-coherency} on page \pageref{par:modifier-coherency})
\end{itemize}

\paragraph{Execution} {Reservation table \textsf{LSU2\_MEMW\_AUXR.Y} (Section~\ref{sec:reservation-LSU2-MEMW-AUXR.Y})}
~
{\begin{lstlisting}
stage ID:
  new offset = _SX_54(extend27_offset27);
  new coherency = _ZX_2(coherency);
stage RR:
  new base = GPR[registerZ];
  new address = base + offset;
stage E1:
  new argument2 = GPR[registerT];
stage E3:
  MEM.atomic_dus(address, 0xFF, coherency << 32, argument2.64[0], registerT);
\end{lstlisting}}
\newpage
\subsubsection[{\small ASDUSH: Atomic Store DUS Half Word}]{ASDUSH\hfill Atomic Store DUS Half Word}
\label{instr:ASDUSH}\index{asdush}
 
\paragraph{Encoding} Format \textsf{LSU\_ASDUSSB} (Section~\ref{sec:format-LSU-ASDUSSB}), \verb|lsucode5 = 001|\\

\bitfields{ 1/P, 2/01, 3/111, 2/coherency, 6/registerT, 2/11, 3/001, 1/--, 5/11110, 1/0, 6/registerZ }


\paragraph{Syntax} \texttt{asdush}\emph{coherency} [\emph{registerZ}] = \emph{registerT}

\paragraph{Description} The effective address is the value of the \emph{registerZ}.
This instruction implements atomic DUS-store on signed half word (Decrement Unsigned Saturating). The half word at effective address is DUSed with the \emph{registerT}. The effective address must be half word aligned (multiple of 2).


\paragraph{Modifier(s)}
\noindent\begin{itemize}
\item coherency (cf. coherency in section \ref{par:modifier-coherency} on page \pageref{par:modifier-coherency})
\end{itemize}

\paragraph{Execution} {Reservation table \textsf{LSU2\_MEMW\_AUXR} (Section~\ref{sec:reservation-LSU2-MEMW-AUXR})}
~
{\begin{lstlisting}
stage ID:
  new coherency = _ZX_2(coherency);
stage RR:
  new address = GPR[registerZ];
stage E1:
  new argument2 = GPR[registerT];
stage E3:
  MEM.atomic_dus(address, 0x3, coherency << 32, argument2.16[0], registerT);
\end{lstlisting}}
\newpage

\paragraph{Encoding} Format \textsf{LSU\_ASDUSSB.O} (Section~\ref{sec:format-LSU-ASDUSSB.O}), \verb|lsucode5 = 001|\\

\bitfields{ 1/P, 2/00, 2/-- --, 27/offset27, 1/1, 2/01, 3/111, 2/coherency, 6/registerT, 2/11, 3/001, 1/--, 5/11110, 1/0, 6/registerZ }


\paragraph{Syntax} \texttt{asdush}\emph{coherency} \emph{offset27}[\emph{registerZ}] = \emph{registerT}

\paragraph{Description} The effective address is computed by adding the \emph{offset27} to the \emph{registerZ}.
This instruction implements atomic DUS-store on signed half word (Decrement Unsigned Saturating). The half word at effective address is DUSed with the \emph{registerT}. The effective address must be half word aligned (multiple of 2).


\paragraph{Modifier(s)}
\noindent\begin{itemize}
\item coherency (cf. coherency in section \ref{par:modifier-coherency} on page \pageref{par:modifier-coherency})
\end{itemize}

\paragraph{Execution} {Reservation table \textsf{LSU2\_MEMW\_AUXR.X} (Section~\ref{sec:reservation-LSU2-MEMW-AUXR.X})}
~
{\begin{lstlisting}
stage ID:
  new offset = _SX_27(offset27);
  new coherency = _ZX_2(coherency);
stage RR:
  new base = GPR[registerZ];
  new address = base + offset;
stage E1:
  new argument2 = GPR[registerT];
stage E3:
  MEM.atomic_dus(address, 0x3, coherency << 32, argument2.16[0], registerT);
\end{lstlisting}}
\newpage

\paragraph{Encoding} Format \textsf{LSU\_ASDUSSB.D} (Section~\ref{sec:format-LSU-ASDUSSB.D}), \verb|lsucode5 = 001|\\

\bitfields{ 1/P, 2/00, 2/-- --, 27/extend27, 1/1, 2/00, 2/-- --, 27/offset27, 1/1, 2/01, 3/111, 2/coherency, 6/registerT, 2/11, 3/001, 1/--, 5/11110, 1/0, 6/registerZ }


\paragraph{Syntax} \texttt{asdush}\emph{coherency} \emph{extend27\_\discretionary{}{}{}offset27}[\emph{registerZ}] = \emph{registerT}

\paragraph{Description} The effective address is computed by adding the \emph{extend27\_\discretionary{}{}{}offset27} to the \emph{registerZ}.
This instruction implements atomic DUS-store on signed half word (Decrement Unsigned Saturating). The half word at effective address is DUSed with the \emph{registerT}. The effective address must be half word aligned (multiple of 2).


\paragraph{Modifier(s)}
\noindent\begin{itemize}
\item coherency (cf. coherency in section \ref{par:modifier-coherency} on page \pageref{par:modifier-coherency})
\end{itemize}

\paragraph{Execution} {Reservation table \textsf{LSU2\_MEMW\_AUXR.Y} (Section~\ref{sec:reservation-LSU2-MEMW-AUXR.Y})}
~
{\begin{lstlisting}
stage ID:
  new offset = _SX_54(extend27_offset27);
  new coherency = _ZX_2(coherency);
stage RR:
  new base = GPR[registerZ];
  new address = base + offset;
stage E1:
  new argument2 = GPR[registerT];
stage E3:
  MEM.atomic_dus(address, 0x3, coherency << 32, argument2.16[0], registerT);
\end{lstlisting}}
\newpage
\subsubsection[{\small ASDUSW: Atomic Store DUS Word}]{ASDUSW\hfill Atomic Store DUS Word}
\label{instr:ASDUSW}\index{asdusw}
 
\paragraph{Encoding} Format \textsf{LSU\_ASDUSSB} (Section~\ref{sec:format-LSU-ASDUSSB}), \verb|lsucode5 = 010|\\

\bitfields{ 1/P, 2/01, 3/111, 2/coherency, 6/registerT, 2/11, 3/010, 1/--, 5/11110, 1/0, 6/registerZ }


\paragraph{Syntax} \texttt{asdusw}\emph{coherency} [\emph{registerZ}] = \emph{registerT}

\paragraph{Description} The effective address is the value of the \emph{registerZ}.
This instruction implements atomic DUS-store on signed word (Decrement Unsigned Saturating). The word at effective address is DUSed with the \emph{registerT}. The effective address must be word aligned (multiple of 4).


\paragraph{Modifier(s)}
\noindent\begin{itemize}
\item coherency (cf. coherency in section \ref{par:modifier-coherency} on page \pageref{par:modifier-coherency})
\end{itemize}

\paragraph{Execution} {Reservation table \textsf{LSU2\_MEMW\_AUXR} (Section~\ref{sec:reservation-LSU2-MEMW-AUXR})}
~
{\begin{lstlisting}
stage ID:
  new coherency = _ZX_2(coherency);
stage RR:
  new address = GPR[registerZ];
stage E1:
  new argument2 = GPR[registerT];
stage E3:
  MEM.atomic_dus(address, 0xF, coherency << 32, argument2.32[0], registerT);
\end{lstlisting}}
\newpage

\paragraph{Encoding} Format \textsf{LSU\_ASDUSSB.O} (Section~\ref{sec:format-LSU-ASDUSSB.O}), \verb|lsucode5 = 010|\\

\bitfields{ 1/P, 2/00, 2/-- --, 27/offset27, 1/1, 2/01, 3/111, 2/coherency, 6/registerT, 2/11, 3/010, 1/--, 5/11110, 1/0, 6/registerZ }


\paragraph{Syntax} \texttt{asdusw}\emph{coherency} \emph{offset27}[\emph{registerZ}] = \emph{registerT}

\paragraph{Description} The effective address is computed by adding the \emph{offset27} to the \emph{registerZ}.
This instruction implements atomic DUS-store on signed word (Decrement Unsigned Saturating). The word at effective address is DUSed with the \emph{registerT}. The effective address must be word aligned (multiple of 4).


\paragraph{Modifier(s)}
\noindent\begin{itemize}
\item coherency (cf. coherency in section \ref{par:modifier-coherency} on page \pageref{par:modifier-coherency})
\end{itemize}

\paragraph{Execution} {Reservation table \textsf{LSU2\_MEMW\_AUXR.X} (Section~\ref{sec:reservation-LSU2-MEMW-AUXR.X})}
~
{\begin{lstlisting}
stage ID:
  new offset = _SX_27(offset27);
  new coherency = _ZX_2(coherency);
stage RR:
  new base = GPR[registerZ];
  new address = base + offset;
stage E1:
  new argument2 = GPR[registerT];
stage E3:
  MEM.atomic_dus(address, 0xF, coherency << 32, argument2.32[0], registerT);
\end{lstlisting}}
\newpage

\paragraph{Encoding} Format \textsf{LSU\_ASDUSSB.D} (Section~\ref{sec:format-LSU-ASDUSSB.D}), \verb|lsucode5 = 010|\\

\bitfields{ 1/P, 2/00, 2/-- --, 27/extend27, 1/1, 2/00, 2/-- --, 27/offset27, 1/1, 2/01, 3/111, 2/coherency, 6/registerT, 2/11, 3/010, 1/--, 5/11110, 1/0, 6/registerZ }


\paragraph{Syntax} \texttt{asdusw}\emph{coherency} \emph{extend27\_\discretionary{}{}{}offset27}[\emph{registerZ}] = \emph{registerT}

\paragraph{Description} The effective address is computed by adding the \emph{extend27\_\discretionary{}{}{}offset27} to the \emph{registerZ}.
This instruction implements atomic DUS-store on signed word (Decrement Unsigned Saturating). The word at effective address is DUSed with the \emph{registerT}. The effective address must be word aligned (multiple of 4).


\paragraph{Modifier(s)}
\noindent\begin{itemize}
\item coherency (cf. coherency in section \ref{par:modifier-coherency} on page \pageref{par:modifier-coherency})
\end{itemize}

\paragraph{Execution} {Reservation table \textsf{LSU2\_MEMW\_AUXR.Y} (Section~\ref{sec:reservation-LSU2-MEMW-AUXR.Y})}
~
{\begin{lstlisting}
stage ID:
  new offset = _SX_54(extend27_offset27);
  new coherency = _ZX_2(coherency);
stage RR:
  new base = GPR[registerZ];
  new address = base + offset;
stage E1:
  new argument2 = GPR[registerT];
stage E3:
  MEM.atomic_dus(address, 0xF, coherency << 32, argument2.32[0], registerT);
\end{lstlisting}}
\newpage
\subsubsection[{\small ASEORB: Atomic Store EOR Byte}]{ASEORB\hfill Atomic Store EOR Byte}
\label{instr:ASEORB}\index{aseorb}
 
\paragraph{Encoding} Format \textsf{LSU\_ASEORSB} (Section~\ref{sec:format-LSU-ASEORSB}), \verb|lsucode5 = 000|\\

\bitfields{ 1/P, 2/01, 3/111, 2/coherency, 6/registerT, 2/11, 3/000, 1/--, 5/10111, 1/0, 6/registerZ }


\paragraph{Syntax} \texttt{aseorb}\emph{coherency} [\emph{registerZ}] = \emph{registerT}

\paragraph{Description} The effective address is the value of the \emph{registerZ}.
This instruction implements atomic EOR-store on byte. The byte at effective address is EORed with the \emph{registerT}.


\paragraph{Modifier(s)}
\noindent\begin{itemize}
\item coherency (cf. coherency in section \ref{par:modifier-coherency} on page \pageref{par:modifier-coherency})
\end{itemize}

\paragraph{Execution} {Reservation table \textsf{LSU2\_MEMW\_AUXR} (Section~\ref{sec:reservation-LSU2-MEMW-AUXR})}
~
{\begin{lstlisting}
stage ID:
  new coherency = _ZX_2(coherency);
stage RR:
  new address = GPR[registerZ];
stage E1:
  new argument2 = GPR[registerT];
stage E3:
  MEM.atomic_eor(address, 0x1, coherency << 32, argument2.8[0], registerT);
\end{lstlisting}}
\newpage

\paragraph{Encoding} Format \textsf{LSU\_ASEORSB.O} (Section~\ref{sec:format-LSU-ASEORSB.O}), \verb|lsucode5 = 000|\\

\bitfields{ 1/P, 2/00, 2/-- --, 27/offset27, 1/1, 2/01, 3/111, 2/coherency, 6/registerT, 2/11, 3/000, 1/--, 5/10111, 1/0, 6/registerZ }


\paragraph{Syntax} \texttt{aseorb}\emph{coherency} \emph{offset27}[\emph{registerZ}] = \emph{registerT}

\paragraph{Description} The effective address is computed by adding the \emph{offset27} to the \emph{registerZ}.
This instruction implements atomic EOR-store on byte. The byte at effective address is EORed with the \emph{registerT}.


\paragraph{Modifier(s)}
\noindent\begin{itemize}
\item coherency (cf. coherency in section \ref{par:modifier-coherency} on page \pageref{par:modifier-coherency})
\end{itemize}

\paragraph{Execution} {Reservation table \textsf{LSU2\_MEMW\_AUXR.X} (Section~\ref{sec:reservation-LSU2-MEMW-AUXR.X})}
~
{\begin{lstlisting}
stage ID:
  new offset = _SX_27(offset27);
  new coherency = _ZX_2(coherency);
stage RR:
  new base = GPR[registerZ];
  new address = base + offset;
stage E1:
  new argument2 = GPR[registerT];
stage E3:
  MEM.atomic_eor(address, 0x1, coherency << 32, argument2.8[0], registerT);
\end{lstlisting}}
\newpage

\paragraph{Encoding} Format \textsf{LSU\_ASEORSB.D} (Section~\ref{sec:format-LSU-ASEORSB.D}), \verb|lsucode5 = 000|\\

\bitfields{ 1/P, 2/00, 2/-- --, 27/extend27, 1/1, 2/00, 2/-- --, 27/offset27, 1/1, 2/01, 3/111, 2/coherency, 6/registerT, 2/11, 3/000, 1/--, 5/10111, 1/0, 6/registerZ }


\paragraph{Syntax} \texttt{aseorb}\emph{coherency} \emph{extend27\_\discretionary{}{}{}offset27}[\emph{registerZ}] = \emph{registerT}

\paragraph{Description} The effective address is computed by adding the \emph{extend27\_\discretionary{}{}{}offset27} to the \emph{registerZ}.
This instruction implements atomic EOR-store on byte. The byte at effective address is EORed with the \emph{registerT}.


\paragraph{Modifier(s)}
\noindent\begin{itemize}
\item coherency (cf. coherency in section \ref{par:modifier-coherency} on page \pageref{par:modifier-coherency})
\end{itemize}

\paragraph{Execution} {Reservation table \textsf{LSU2\_MEMW\_AUXR.Y} (Section~\ref{sec:reservation-LSU2-MEMW-AUXR.Y})}
~
{\begin{lstlisting}
stage ID:
  new offset = _SX_54(extend27_offset27);
  new coherency = _ZX_2(coherency);
stage RR:
  new base = GPR[registerZ];
  new address = base + offset;
stage E1:
  new argument2 = GPR[registerT];
stage E3:
  MEM.atomic_eor(address, 0x1, coherency << 32, argument2.8[0], registerT);
\end{lstlisting}}
\newpage
\subsubsection[{\small ASEORD: Atomic Store EOR Double Word}]{ASEORD\hfill Atomic Store EOR Double Word}
\label{instr:ASEORD}\index{aseord}
 
\paragraph{Encoding} Format \textsf{LSU\_ASEORSB} (Section~\ref{sec:format-LSU-ASEORSB}), \verb|lsucode5 = 011|\\

\bitfields{ 1/P, 2/01, 3/111, 2/coherency, 6/registerT, 2/11, 3/011, 1/--, 5/10111, 1/0, 6/registerZ }


\paragraph{Syntax} \texttt{aseord}\emph{coherency} [\emph{registerZ}] = \emph{registerT}

\paragraph{Description} The effective address is the value of the \emph{registerZ}.
This instruction implements atomic EOR-store on double word. The double word at effective address is EORed with the \emph{registerT}. The effective address must be double word aligned (multiple of 8).


\paragraph{Modifier(s)}
\noindent\begin{itemize}
\item coherency (cf. coherency in section \ref{par:modifier-coherency} on page \pageref{par:modifier-coherency})
\end{itemize}

\paragraph{Execution} {Reservation table \textsf{LSU2\_MEMW\_AUXR} (Section~\ref{sec:reservation-LSU2-MEMW-AUXR})}
~
{\begin{lstlisting}
stage ID:
  new coherency = _ZX_2(coherency);
stage RR:
  new address = GPR[registerZ];
stage E1:
  new argument2 = GPR[registerT];
stage E3:
  MEM.atomic_eor(address, 0xFF, coherency << 32, argument2.64[0], registerT);
\end{lstlisting}}
\newpage

\paragraph{Encoding} Format \textsf{LSU\_ASEORSB.O} (Section~\ref{sec:format-LSU-ASEORSB.O}), \verb|lsucode5 = 011|\\

\bitfields{ 1/P, 2/00, 2/-- --, 27/offset27, 1/1, 2/01, 3/111, 2/coherency, 6/registerT, 2/11, 3/011, 1/--, 5/10111, 1/0, 6/registerZ }


\paragraph{Syntax} \texttt{aseord}\emph{coherency} \emph{offset27}[\emph{registerZ}] = \emph{registerT}

\paragraph{Description} The effective address is computed by adding the \emph{offset27} to the \emph{registerZ}.
This instruction implements atomic EOR-store on double word. The double word at effective address is EORed with the \emph{registerT}. The effective address must be double word aligned (multiple of 8).


\paragraph{Modifier(s)}
\noindent\begin{itemize}
\item coherency (cf. coherency in section \ref{par:modifier-coherency} on page \pageref{par:modifier-coherency})
\end{itemize}

\paragraph{Execution} {Reservation table \textsf{LSU2\_MEMW\_AUXR.X} (Section~\ref{sec:reservation-LSU2-MEMW-AUXR.X})}
~
{\begin{lstlisting}
stage ID:
  new offset = _SX_27(offset27);
  new coherency = _ZX_2(coherency);
stage RR:
  new base = GPR[registerZ];
  new address = base + offset;
stage E1:
  new argument2 = GPR[registerT];
stage E3:
  MEM.atomic_eor(address, 0xFF, coherency << 32, argument2.64[0], registerT);
\end{lstlisting}}
\newpage

\paragraph{Encoding} Format \textsf{LSU\_ASEORSB.D} (Section~\ref{sec:format-LSU-ASEORSB.D}), \verb|lsucode5 = 011|\\

\bitfields{ 1/P, 2/00, 2/-- --, 27/extend27, 1/1, 2/00, 2/-- --, 27/offset27, 1/1, 2/01, 3/111, 2/coherency, 6/registerT, 2/11, 3/011, 1/--, 5/10111, 1/0, 6/registerZ }


\paragraph{Syntax} \texttt{aseord}\emph{coherency} \emph{extend27\_\discretionary{}{}{}offset27}[\emph{registerZ}] = \emph{registerT}

\paragraph{Description} The effective address is computed by adding the \emph{extend27\_\discretionary{}{}{}offset27} to the \emph{registerZ}.
This instruction implements atomic EOR-store on double word. The double word at effective address is EORed with the \emph{registerT}. The effective address must be double word aligned (multiple of 8).


\paragraph{Modifier(s)}
\noindent\begin{itemize}
\item coherency (cf. coherency in section \ref{par:modifier-coherency} on page \pageref{par:modifier-coherency})
\end{itemize}

\paragraph{Execution} {Reservation table \textsf{LSU2\_MEMW\_AUXR.Y} (Section~\ref{sec:reservation-LSU2-MEMW-AUXR.Y})}
~
{\begin{lstlisting}
stage ID:
  new offset = _SX_54(extend27_offset27);
  new coherency = _ZX_2(coherency);
stage RR:
  new base = GPR[registerZ];
  new address = base + offset;
stage E1:
  new argument2 = GPR[registerT];
stage E3:
  MEM.atomic_eor(address, 0xFF, coherency << 32, argument2.64[0], registerT);
\end{lstlisting}}
\newpage
\subsubsection[{\small ASEORH: Atomic Store EOR Half Word}]{ASEORH\hfill Atomic Store EOR Half Word}
\label{instr:ASEORH}\index{aseorh}
 
\paragraph{Encoding} Format \textsf{LSU\_ASEORSB} (Section~\ref{sec:format-LSU-ASEORSB}), \verb|lsucode5 = 001|\\

\bitfields{ 1/P, 2/01, 3/111, 2/coherency, 6/registerT, 2/11, 3/001, 1/--, 5/10111, 1/0, 6/registerZ }


\paragraph{Syntax} \texttt{aseorh}\emph{coherency} [\emph{registerZ}] = \emph{registerT}

\paragraph{Description} The effective address is the value of the \emph{registerZ}.
This instruction implements atomic EOR-store on half word. The half word at effective address is EORed with the \emph{registerT}. The effective address must be half word aligned (multiple of 2).


\paragraph{Modifier(s)}
\noindent\begin{itemize}
\item coherency (cf. coherency in section \ref{par:modifier-coherency} on page \pageref{par:modifier-coherency})
\end{itemize}

\paragraph{Execution} {Reservation table \textsf{LSU2\_MEMW\_AUXR} (Section~\ref{sec:reservation-LSU2-MEMW-AUXR})}
~
{\begin{lstlisting}
stage ID:
  new coherency = _ZX_2(coherency);
stage RR:
  new address = GPR[registerZ];
stage E1:
  new argument2 = GPR[registerT];
stage E3:
  MEM.atomic_eor(address, 0x3, coherency << 32, argument2.16[0], registerT);
\end{lstlisting}}
\newpage

\paragraph{Encoding} Format \textsf{LSU\_ASEORSB.O} (Section~\ref{sec:format-LSU-ASEORSB.O}), \verb|lsucode5 = 001|\\

\bitfields{ 1/P, 2/00, 2/-- --, 27/offset27, 1/1, 2/01, 3/111, 2/coherency, 6/registerT, 2/11, 3/001, 1/--, 5/10111, 1/0, 6/registerZ }


\paragraph{Syntax} \texttt{aseorh}\emph{coherency} \emph{offset27}[\emph{registerZ}] = \emph{registerT}

\paragraph{Description} The effective address is computed by adding the \emph{offset27} to the \emph{registerZ}.
This instruction implements atomic EOR-store on half word. The half word at effective address is EORed with the \emph{registerT}. The effective address must be half word aligned (multiple of 2).


\paragraph{Modifier(s)}
\noindent\begin{itemize}
\item coherency (cf. coherency in section \ref{par:modifier-coherency} on page \pageref{par:modifier-coherency})
\end{itemize}

\paragraph{Execution} {Reservation table \textsf{LSU2\_MEMW\_AUXR.X} (Section~\ref{sec:reservation-LSU2-MEMW-AUXR.X})}
~
{\begin{lstlisting}
stage ID:
  new offset = _SX_27(offset27);
  new coherency = _ZX_2(coherency);
stage RR:
  new base = GPR[registerZ];
  new address = base + offset;
stage E1:
  new argument2 = GPR[registerT];
stage E3:
  MEM.atomic_eor(address, 0x3, coherency << 32, argument2.16[0], registerT);
\end{lstlisting}}
\newpage

\paragraph{Encoding} Format \textsf{LSU\_ASEORSB.D} (Section~\ref{sec:format-LSU-ASEORSB.D}), \verb|lsucode5 = 001|\\

\bitfields{ 1/P, 2/00, 2/-- --, 27/extend27, 1/1, 2/00, 2/-- --, 27/offset27, 1/1, 2/01, 3/111, 2/coherency, 6/registerT, 2/11, 3/001, 1/--, 5/10111, 1/0, 6/registerZ }


\paragraph{Syntax} \texttt{aseorh}\emph{coherency} \emph{extend27\_\discretionary{}{}{}offset27}[\emph{registerZ}] = \emph{registerT}

\paragraph{Description} The effective address is computed by adding the \emph{extend27\_\discretionary{}{}{}offset27} to the \emph{registerZ}.
This instruction implements atomic EOR-store on half word. The half word at effective address is EORed with the \emph{registerT}. The effective address must be half word aligned (multiple of 2).


\paragraph{Modifier(s)}
\noindent\begin{itemize}
\item coherency (cf. coherency in section \ref{par:modifier-coherency} on page \pageref{par:modifier-coherency})
\end{itemize}

\paragraph{Execution} {Reservation table \textsf{LSU2\_MEMW\_AUXR.Y} (Section~\ref{sec:reservation-LSU2-MEMW-AUXR.Y})}
~
{\begin{lstlisting}
stage ID:
  new offset = _SX_54(extend27_offset27);
  new coherency = _ZX_2(coherency);
stage RR:
  new base = GPR[registerZ];
  new address = base + offset;
stage E1:
  new argument2 = GPR[registerT];
stage E3:
  MEM.atomic_eor(address, 0x3, coherency << 32, argument2.16[0], registerT);
\end{lstlisting}}
\newpage
\subsubsection[{\small ASEORW: Atomic Store EOR Word}]{ASEORW\hfill Atomic Store EOR Word}
\label{instr:ASEORW}\index{aseorw}
 
\paragraph{Encoding} Format \textsf{LSU\_ASEORSB} (Section~\ref{sec:format-LSU-ASEORSB}), \verb|lsucode5 = 010|\\

\bitfields{ 1/P, 2/01, 3/111, 2/coherency, 6/registerT, 2/11, 3/010, 1/--, 5/10111, 1/0, 6/registerZ }


\paragraph{Syntax} \texttt{aseorw}\emph{coherency} [\emph{registerZ}] = \emph{registerT}

\paragraph{Description} The effective address is the value of the \emph{registerZ}.
This instruction implements atomic EOR-store on word. The word at effective address is EORed with the \emph{registerT}. The effective address must be word aligned (multiple of 4).


\paragraph{Modifier(s)}
\noindent\begin{itemize}
\item coherency (cf. coherency in section \ref{par:modifier-coherency} on page \pageref{par:modifier-coherency})
\end{itemize}

\paragraph{Execution} {Reservation table \textsf{LSU2\_MEMW\_AUXR} (Section~\ref{sec:reservation-LSU2-MEMW-AUXR})}
~
{\begin{lstlisting}
stage ID:
  new coherency = _ZX_2(coherency);
stage RR:
  new address = GPR[registerZ];
stage E1:
  new argument2 = GPR[registerT];
stage E3:
  MEM.atomic_eor(address, 0xF, coherency << 32, argument2.32[0], registerT);
\end{lstlisting}}
\newpage

\paragraph{Encoding} Format \textsf{LSU\_ASEORSB.O} (Section~\ref{sec:format-LSU-ASEORSB.O}), \verb|lsucode5 = 010|\\

\bitfields{ 1/P, 2/00, 2/-- --, 27/offset27, 1/1, 2/01, 3/111, 2/coherency, 6/registerT, 2/11, 3/010, 1/--, 5/10111, 1/0, 6/registerZ }


\paragraph{Syntax} \texttt{aseorw}\emph{coherency} \emph{offset27}[\emph{registerZ}] = \emph{registerT}

\paragraph{Description} The effective address is computed by adding the \emph{offset27} to the \emph{registerZ}.
This instruction implements atomic EOR-store on word. The word at effective address is EORed with the \emph{registerT}. The effective address must be word aligned (multiple of 4).


\paragraph{Modifier(s)}
\noindent\begin{itemize}
\item coherency (cf. coherency in section \ref{par:modifier-coherency} on page \pageref{par:modifier-coherency})
\end{itemize}

\paragraph{Execution} {Reservation table \textsf{LSU2\_MEMW\_AUXR.X} (Section~\ref{sec:reservation-LSU2-MEMW-AUXR.X})}
~
{\begin{lstlisting}
stage ID:
  new offset = _SX_27(offset27);
  new coherency = _ZX_2(coherency);
stage RR:
  new base = GPR[registerZ];
  new address = base + offset;
stage E1:
  new argument2 = GPR[registerT];
stage E3:
  MEM.atomic_eor(address, 0xF, coherency << 32, argument2.32[0], registerT);
\end{lstlisting}}
\newpage

\paragraph{Encoding} Format \textsf{LSU\_ASEORSB.D} (Section~\ref{sec:format-LSU-ASEORSB.D}), \verb|lsucode5 = 010|\\

\bitfields{ 1/P, 2/00, 2/-- --, 27/extend27, 1/1, 2/00, 2/-- --, 27/offset27, 1/1, 2/01, 3/111, 2/coherency, 6/registerT, 2/11, 3/010, 1/--, 5/10111, 1/0, 6/registerZ }


\paragraph{Syntax} \texttt{aseorw}\emph{coherency} \emph{extend27\_\discretionary{}{}{}offset27}[\emph{registerZ}] = \emph{registerT}

\paragraph{Description} The effective address is computed by adding the \emph{extend27\_\discretionary{}{}{}offset27} to the \emph{registerZ}.
This instruction implements atomic EOR-store on word. The word at effective address is EORed with the \emph{registerT}. The effective address must be word aligned (multiple of 4).


\paragraph{Modifier(s)}
\noindent\begin{itemize}
\item coherency (cf. coherency in section \ref{par:modifier-coherency} on page \pageref{par:modifier-coherency})
\end{itemize}

\paragraph{Execution} {Reservation table \textsf{LSU2\_MEMW\_AUXR.Y} (Section~\ref{sec:reservation-LSU2-MEMW-AUXR.Y})}
~
{\begin{lstlisting}
stage ID:
  new offset = _SX_54(extend27_offset27);
  new coherency = _ZX_2(coherency);
stage RR:
  new base = GPR[registerZ];
  new address = base + offset;
stage E1:
  new argument2 = GPR[registerT];
stage E3:
  MEM.atomic_eor(address, 0xF, coherency << 32, argument2.32[0], registerT);
\end{lstlisting}}
\newpage
\subsubsection[{\small ASH: Atomic Store Half Word}]{ASH\hfill Atomic Store Half Word}
\label{instr:ASH}\index{ash}
 
\paragraph{Encoding} Format \textsf{LSU\_ASSB} (Section~\ref{sec:format-LSU-ASSB}), \verb|lsucode5 = 001|\\

\bitfields{ 1/P, 2/01, 3/111, 2/coherency, 6/registerT, 2/11, 3/001, 1/--, 5/10000, 1/0, 6/registerZ }


\paragraph{Syntax} \texttt{ash}\emph{coherency} [\emph{registerZ}] = \emph{registerT}

\paragraph{Description} The effective address is the value of the \emph{registerZ}.
The \emph{registerT} is atomically stored to the memory half word starting at the effective address. This instruction is provided to directly operate on the last level of the memory hierarchy, which is typically a processor DDR or a cluster TCM, irrespective of the MMU DCP (Data Cache Policy) settings. The effective address must be half word aligned (multiple of 2).


\paragraph{Modifier(s)}
\noindent\begin{itemize}
\item coherency (cf. coherency in section \ref{par:modifier-coherency} on page \pageref{par:modifier-coherency})
\end{itemize}

\paragraph{Execution} {Reservation table \textsf{LSU\_MEMW\_AUXR} (Section~\ref{sec:reservation-LSU-MEMW-AUXR})}
~
{\begin{lstlisting}
stage ID:
  new coherency = _ZX_2(coherency);
stage RR:
  new address = GPR[registerZ];
stage E1:
  new argument2 = GPR[registerT];
stage E3:
  MEM.atomic_store(address, 0x3, coherency << 32, argument2, registerT);
\end{lstlisting}}
\newpage

\paragraph{Encoding} Format \textsf{LSU\_ASSB.O} (Section~\ref{sec:format-LSU-ASSB.O}), \verb|lsucode5 = 001|\\

\bitfields{ 1/P, 2/00, 2/-- --, 27/offset27, 1/1, 2/01, 3/111, 2/coherency, 6/registerT, 2/11, 3/001, 1/--, 5/10000, 1/0, 6/registerZ }


\paragraph{Syntax} \texttt{ash}\emph{coherency} \emph{offset27}[\emph{registerZ}] = \emph{registerT}

\paragraph{Description} The effective address is computed by adding the \emph{offset27} to the \emph{registerZ}.
The \emph{registerT} is atomically stored to the memory half word starting at the effective address. This instruction is provided to directly operate on the last level of the memory hierarchy, which is typically a processor DDR or a cluster TCM, irrespective of the MMU DCP (Data Cache Policy) settings. The effective address must be half word aligned (multiple of 2).


\paragraph{Modifier(s)}
\noindent\begin{itemize}
\item coherency (cf. coherency in section \ref{par:modifier-coherency} on page \pageref{par:modifier-coherency})
\end{itemize}

\paragraph{Execution} {Reservation table \textsf{LSU\_MEMW\_AUXR.X} (Section~\ref{sec:reservation-LSU-MEMW-AUXR.X})}
~
{\begin{lstlisting}
stage ID:
  new offset = _SX_27(offset27);
  new coherency = _ZX_2(coherency);
stage RR:
  new base = GPR[registerZ];
  new address = base + offset;
stage E1:
  new argument2 = GPR[registerT];
stage E3:
  MEM.atomic_store(address, 0x3, coherency << 32, argument2, registerT);
\end{lstlisting}}
\newpage

\paragraph{Encoding} Format \textsf{LSU\_ASSB.D} (Section~\ref{sec:format-LSU-ASSB.D}), \verb|lsucode5 = 001|\\

\bitfields{ 1/P, 2/00, 2/-- --, 27/extend27, 1/1, 2/00, 2/-- --, 27/offset27, 1/1, 2/01, 3/111, 2/coherency, 6/registerT, 2/11, 3/001, 1/--, 5/10000, 1/0, 6/registerZ }


\paragraph{Syntax} \texttt{ash}\emph{coherency} \emph{extend27\_\discretionary{}{}{}offset27}[\emph{registerZ}] = \emph{registerT}

\paragraph{Description} The effective address is computed by adding the \emph{extend27\_\discretionary{}{}{}offset27} to the \emph{registerZ}.
The \emph{registerT} is atomically stored to the memory half word starting at the effective address. This instruction is provided to directly operate on the last level of the memory hierarchy, which is typically a processor DDR or a cluster TCM, irrespective of the MMU DCP (Data Cache Policy) settings. The effective address must be half word aligned (multiple of 2).


\paragraph{Modifier(s)}
\noindent\begin{itemize}
\item coherency (cf. coherency in section \ref{par:modifier-coherency} on page \pageref{par:modifier-coherency})
\end{itemize}

\paragraph{Execution} {Reservation table \textsf{LSU\_MEMW\_AUXR.Y} (Section~\ref{sec:reservation-LSU-MEMW-AUXR.Y})}
~
{\begin{lstlisting}
stage ID:
  new offset = _SX_54(extend27_offset27);
  new coherency = _ZX_2(coherency);
stage RR:
  new base = GPR[registerZ];
  new address = base + offset;
stage E1:
  new argument2 = GPR[registerT];
stage E3:
  MEM.atomic_store(address, 0x3, coherency << 32, argument2, registerT);
\end{lstlisting}}
\newpage
\subsubsection[{\small ASIORB: Atomic Store IOR Byte}]{ASIORB\hfill Atomic Store IOR Byte}
\label{instr:ASIORB}\index{asiorb}
 
\paragraph{Encoding} Format \textsf{LSU\_ASIORSB} (Section~\ref{sec:format-LSU-ASIORSB}), \verb|lsucode5 = 000|\\

\bitfields{ 1/P, 2/01, 3/111, 2/coherency, 6/registerT, 2/11, 3/000, 1/--, 5/10110, 1/0, 6/registerZ }


\paragraph{Syntax} \texttt{asiorb}\emph{coherency} [\emph{registerZ}] = \emph{registerT}

\paragraph{Description} The effective address is the value of the \emph{registerZ}.
This instruction implements atomic IOR-store on byte. The byte at effective address is IORed with the \emph{registerT}.


\paragraph{Modifier(s)}
\noindent\begin{itemize}
\item coherency (cf. coherency in section \ref{par:modifier-coherency} on page \pageref{par:modifier-coherency})
\end{itemize}

\paragraph{Execution} {Reservation table \textsf{LSU2\_MEMW\_AUXR} (Section~\ref{sec:reservation-LSU2-MEMW-AUXR})}
~
{\begin{lstlisting}
stage ID:
  new coherency = _ZX_2(coherency);
stage RR:
  new address = GPR[registerZ];
stage E1:
  new argument2 = GPR[registerT];
stage E3:
  MEM.atomic_ior(address, 0x1, coherency << 32, argument2.8[0], registerT);
\end{lstlisting}}
\newpage

\paragraph{Encoding} Format \textsf{LSU\_ASIORSB.O} (Section~\ref{sec:format-LSU-ASIORSB.O}), \verb|lsucode5 = 000|\\

\bitfields{ 1/P, 2/00, 2/-- --, 27/offset27, 1/1, 2/01, 3/111, 2/coherency, 6/registerT, 2/11, 3/000, 1/--, 5/10110, 1/0, 6/registerZ }


\paragraph{Syntax} \texttt{asiorb}\emph{coherency} \emph{offset27}[\emph{registerZ}] = \emph{registerT}

\paragraph{Description} The effective address is computed by adding the \emph{offset27} to the \emph{registerZ}.
This instruction implements atomic IOR-store on byte. The byte at effective address is IORed with the \emph{registerT}.


\paragraph{Modifier(s)}
\noindent\begin{itemize}
\item coherency (cf. coherency in section \ref{par:modifier-coherency} on page \pageref{par:modifier-coherency})
\end{itemize}

\paragraph{Execution} {Reservation table \textsf{LSU2\_MEMW\_AUXR.X} (Section~\ref{sec:reservation-LSU2-MEMW-AUXR.X})}
~
{\begin{lstlisting}
stage ID:
  new offset = _SX_27(offset27);
  new coherency = _ZX_2(coherency);
stage RR:
  new base = GPR[registerZ];
  new address = base + offset;
stage E1:
  new argument2 = GPR[registerT];
stage E3:
  MEM.atomic_ior(address, 0x1, coherency << 32, argument2.8[0], registerT);
\end{lstlisting}}
\newpage

\paragraph{Encoding} Format \textsf{LSU\_ASIORSB.D} (Section~\ref{sec:format-LSU-ASIORSB.D}), \verb|lsucode5 = 000|\\

\bitfields{ 1/P, 2/00, 2/-- --, 27/extend27, 1/1, 2/00, 2/-- --, 27/offset27, 1/1, 2/01, 3/111, 2/coherency, 6/registerT, 2/11, 3/000, 1/--, 5/10110, 1/0, 6/registerZ }


\paragraph{Syntax} \texttt{asiorb}\emph{coherency} \emph{extend27\_\discretionary{}{}{}offset27}[\emph{registerZ}] = \emph{registerT}

\paragraph{Description} The effective address is computed by adding the \emph{extend27\_\discretionary{}{}{}offset27} to the \emph{registerZ}.
This instruction implements atomic IOR-store on byte. The byte at effective address is IORed with the \emph{registerT}.


\paragraph{Modifier(s)}
\noindent\begin{itemize}
\item coherency (cf. coherency in section \ref{par:modifier-coherency} on page \pageref{par:modifier-coherency})
\end{itemize}

\paragraph{Execution} {Reservation table \textsf{LSU2\_MEMW\_AUXR.Y} (Section~\ref{sec:reservation-LSU2-MEMW-AUXR.Y})}
~
{\begin{lstlisting}
stage ID:
  new offset = _SX_54(extend27_offset27);
  new coherency = _ZX_2(coherency);
stage RR:
  new base = GPR[registerZ];
  new address = base + offset;
stage E1:
  new argument2 = GPR[registerT];
stage E3:
  MEM.atomic_ior(address, 0x1, coherency << 32, argument2.8[0], registerT);
\end{lstlisting}}
\newpage
\subsubsection[{\small ASIORD: Atomic Store IOR Double Word}]{ASIORD\hfill Atomic Store IOR Double Word}
\label{instr:ASIORD}\index{asiord}
 
\paragraph{Encoding} Format \textsf{LSU\_ASIORSB} (Section~\ref{sec:format-LSU-ASIORSB}), \verb|lsucode5 = 011|\\

\bitfields{ 1/P, 2/01, 3/111, 2/coherency, 6/registerT, 2/11, 3/011, 1/--, 5/10110, 1/0, 6/registerZ }


\paragraph{Syntax} \texttt{asiord}\emph{coherency} [\emph{registerZ}] = \emph{registerT}

\paragraph{Description} The effective address is the value of the \emph{registerZ}.
This instruction implements atomic IOR-store on double word. The double word at effective address is IORed with the \emph{registerT}. The effective address must be double word aligned (multiple of 8).


\paragraph{Modifier(s)}
\noindent\begin{itemize}
\item coherency (cf. coherency in section \ref{par:modifier-coherency} on page \pageref{par:modifier-coherency})
\end{itemize}

\paragraph{Execution} {Reservation table \textsf{LSU2\_MEMW\_AUXR} (Section~\ref{sec:reservation-LSU2-MEMW-AUXR})}
~
{\begin{lstlisting}
stage ID:
  new coherency = _ZX_2(coherency);
stage RR:
  new address = GPR[registerZ];
stage E1:
  new argument2 = GPR[registerT];
stage E3:
  MEM.atomic_ior(address, 0xFF, coherency << 32, argument2.64[0], registerT);
\end{lstlisting}}
\newpage

\paragraph{Encoding} Format \textsf{LSU\_ASIORSB.O} (Section~\ref{sec:format-LSU-ASIORSB.O}), \verb|lsucode5 = 011|\\

\bitfields{ 1/P, 2/00, 2/-- --, 27/offset27, 1/1, 2/01, 3/111, 2/coherency, 6/registerT, 2/11, 3/011, 1/--, 5/10110, 1/0, 6/registerZ }


\paragraph{Syntax} \texttt{asiord}\emph{coherency} \emph{offset27}[\emph{registerZ}] = \emph{registerT}

\paragraph{Description} The effective address is computed by adding the \emph{offset27} to the \emph{registerZ}.
This instruction implements atomic IOR-store on double word. The double word at effective address is IORed with the \emph{registerT}. The effective address must be double word aligned (multiple of 8).


\paragraph{Modifier(s)}
\noindent\begin{itemize}
\item coherency (cf. coherency in section \ref{par:modifier-coherency} on page \pageref{par:modifier-coherency})
\end{itemize}

\paragraph{Execution} {Reservation table \textsf{LSU2\_MEMW\_AUXR.X} (Section~\ref{sec:reservation-LSU2-MEMW-AUXR.X})}
~
{\begin{lstlisting}
stage ID:
  new offset = _SX_27(offset27);
  new coherency = _ZX_2(coherency);
stage RR:
  new base = GPR[registerZ];
  new address = base + offset;
stage E1:
  new argument2 = GPR[registerT];
stage E3:
  MEM.atomic_ior(address, 0xFF, coherency << 32, argument2.64[0], registerT);
\end{lstlisting}}
\newpage

\paragraph{Encoding} Format \textsf{LSU\_ASIORSB.D} (Section~\ref{sec:format-LSU-ASIORSB.D}), \verb|lsucode5 = 011|\\

\bitfields{ 1/P, 2/00, 2/-- --, 27/extend27, 1/1, 2/00, 2/-- --, 27/offset27, 1/1, 2/01, 3/111, 2/coherency, 6/registerT, 2/11, 3/011, 1/--, 5/10110, 1/0, 6/registerZ }


\paragraph{Syntax} \texttt{asiord}\emph{coherency} \emph{extend27\_\discretionary{}{}{}offset27}[\emph{registerZ}] = \emph{registerT}

\paragraph{Description} The effective address is computed by adding the \emph{extend27\_\discretionary{}{}{}offset27} to the \emph{registerZ}.
This instruction implements atomic IOR-store on double word. The double word at effective address is IORed with the \emph{registerT}. The effective address must be double word aligned (multiple of 8).


\paragraph{Modifier(s)}
\noindent\begin{itemize}
\item coherency (cf. coherency in section \ref{par:modifier-coherency} on page \pageref{par:modifier-coherency})
\end{itemize}

\paragraph{Execution} {Reservation table \textsf{LSU2\_MEMW\_AUXR.Y} (Section~\ref{sec:reservation-LSU2-MEMW-AUXR.Y})}
~
{\begin{lstlisting}
stage ID:
  new offset = _SX_54(extend27_offset27);
  new coherency = _ZX_2(coherency);
stage RR:
  new base = GPR[registerZ];
  new address = base + offset;
stage E1:
  new argument2 = GPR[registerT];
stage E3:
  MEM.atomic_ior(address, 0xFF, coherency << 32, argument2.64[0], registerT);
\end{lstlisting}}
\newpage
\subsubsection[{\small ASIORH: Atomic Store IOR Half Word}]{ASIORH\hfill Atomic Store IOR Half Word}
\label{instr:ASIORH}\index{asiorh}
 
\paragraph{Encoding} Format \textsf{LSU\_ASIORSB} (Section~\ref{sec:format-LSU-ASIORSB}), \verb|lsucode5 = 001|\\

\bitfields{ 1/P, 2/01, 3/111, 2/coherency, 6/registerT, 2/11, 3/001, 1/--, 5/10110, 1/0, 6/registerZ }


\paragraph{Syntax} \texttt{asiorh}\emph{coherency} [\emph{registerZ}] = \emph{registerT}

\paragraph{Description} The effective address is the value of the \emph{registerZ}.
This instruction implements atomic IOR-store on half word. The half word at effective address is IORed with the \emph{registerT}. The effective address must be half word aligned (multiple of 2).


\paragraph{Modifier(s)}
\noindent\begin{itemize}
\item coherency (cf. coherency in section \ref{par:modifier-coherency} on page \pageref{par:modifier-coherency})
\end{itemize}

\paragraph{Execution} {Reservation table \textsf{LSU2\_MEMW\_AUXR} (Section~\ref{sec:reservation-LSU2-MEMW-AUXR})}
~
{\begin{lstlisting}
stage ID:
  new coherency = _ZX_2(coherency);
stage RR:
  new address = GPR[registerZ];
stage E1:
  new argument2 = GPR[registerT];
stage E3:
  MEM.atomic_ior(address, 0x3, coherency << 32, argument2.16[0], registerT);
\end{lstlisting}}
\newpage

\paragraph{Encoding} Format \textsf{LSU\_ASIORSB.O} (Section~\ref{sec:format-LSU-ASIORSB.O}), \verb|lsucode5 = 001|\\

\bitfields{ 1/P, 2/00, 2/-- --, 27/offset27, 1/1, 2/01, 3/111, 2/coherency, 6/registerT, 2/11, 3/001, 1/--, 5/10110, 1/0, 6/registerZ }


\paragraph{Syntax} \texttt{asiorh}\emph{coherency} \emph{offset27}[\emph{registerZ}] = \emph{registerT}

\paragraph{Description} The effective address is computed by adding the \emph{offset27} to the \emph{registerZ}.
This instruction implements atomic IOR-store on half word. The half word at effective address is IORed with the \emph{registerT}. The effective address must be half word aligned (multiple of 2).


\paragraph{Modifier(s)}
\noindent\begin{itemize}
\item coherency (cf. coherency in section \ref{par:modifier-coherency} on page \pageref{par:modifier-coherency})
\end{itemize}

\paragraph{Execution} {Reservation table \textsf{LSU2\_MEMW\_AUXR.X} (Section~\ref{sec:reservation-LSU2-MEMW-AUXR.X})}
~
{\begin{lstlisting}
stage ID:
  new offset = _SX_27(offset27);
  new coherency = _ZX_2(coherency);
stage RR:
  new base = GPR[registerZ];
  new address = base + offset;
stage E1:
  new argument2 = GPR[registerT];
stage E3:
  MEM.atomic_ior(address, 0x3, coherency << 32, argument2.16[0], registerT);
\end{lstlisting}}
\newpage

\paragraph{Encoding} Format \textsf{LSU\_ASIORSB.D} (Section~\ref{sec:format-LSU-ASIORSB.D}), \verb|lsucode5 = 001|\\

\bitfields{ 1/P, 2/00, 2/-- --, 27/extend27, 1/1, 2/00, 2/-- --, 27/offset27, 1/1, 2/01, 3/111, 2/coherency, 6/registerT, 2/11, 3/001, 1/--, 5/10110, 1/0, 6/registerZ }


\paragraph{Syntax} \texttt{asiorh}\emph{coherency} \emph{extend27\_\discretionary{}{}{}offset27}[\emph{registerZ}] = \emph{registerT}

\paragraph{Description} The effective address is computed by adding the \emph{extend27\_\discretionary{}{}{}offset27} to the \emph{registerZ}.
This instruction implements atomic IOR-store on half word. The half word at effective address is IORed with the \emph{registerT}. The effective address must be half word aligned (multiple of 2).


\paragraph{Modifier(s)}
\noindent\begin{itemize}
\item coherency (cf. coherency in section \ref{par:modifier-coherency} on page \pageref{par:modifier-coherency})
\end{itemize}

\paragraph{Execution} {Reservation table \textsf{LSU2\_MEMW\_AUXR.Y} (Section~\ref{sec:reservation-LSU2-MEMW-AUXR.Y})}
~
{\begin{lstlisting}
stage ID:
  new offset = _SX_54(extend27_offset27);
  new coherency = _ZX_2(coherency);
stage RR:
  new base = GPR[registerZ];
  new address = base + offset;
stage E1:
  new argument2 = GPR[registerT];
stage E3:
  MEM.atomic_ior(address, 0x3, coherency << 32, argument2.16[0], registerT);
\end{lstlisting}}
\newpage
\subsubsection[{\small ASIORW: Atomic Store IOR Word}]{ASIORW\hfill Atomic Store IOR Word}
\label{instr:ASIORW}\index{asiorw}
 
\paragraph{Encoding} Format \textsf{LSU\_ASIORSB} (Section~\ref{sec:format-LSU-ASIORSB}), \verb|lsucode5 = 010|\\

\bitfields{ 1/P, 2/01, 3/111, 2/coherency, 6/registerT, 2/11, 3/010, 1/--, 5/10110, 1/0, 6/registerZ }


\paragraph{Syntax} \texttt{asiorw}\emph{coherency} [\emph{registerZ}] = \emph{registerT}

\paragraph{Description} The effective address is the value of the \emph{registerZ}.
This instruction implements atomic IOR-store on word. The word at effective address is IORed with the \emph{registerT}. The effective address must be word aligned (multiple of 4).


\paragraph{Modifier(s)}
\noindent\begin{itemize}
\item coherency (cf. coherency in section \ref{par:modifier-coherency} on page \pageref{par:modifier-coherency})
\end{itemize}

\paragraph{Execution} {Reservation table \textsf{LSU2\_MEMW\_AUXR} (Section~\ref{sec:reservation-LSU2-MEMW-AUXR})}
~
{\begin{lstlisting}
stage ID:
  new coherency = _ZX_2(coherency);
stage RR:
  new address = GPR[registerZ];
stage E1:
  new argument2 = GPR[registerT];
stage E3:
  MEM.atomic_ior(address, 0xF, coherency << 32, argument2.32[0], registerT);
\end{lstlisting}}
\newpage

\paragraph{Encoding} Format \textsf{LSU\_ASIORSB.O} (Section~\ref{sec:format-LSU-ASIORSB.O}), \verb|lsucode5 = 010|\\

\bitfields{ 1/P, 2/00, 2/-- --, 27/offset27, 1/1, 2/01, 3/111, 2/coherency, 6/registerT, 2/11, 3/010, 1/--, 5/10110, 1/0, 6/registerZ }


\paragraph{Syntax} \texttt{asiorw}\emph{coherency} \emph{offset27}[\emph{registerZ}] = \emph{registerT}

\paragraph{Description} The effective address is computed by adding the \emph{offset27} to the \emph{registerZ}.
This instruction implements atomic IOR-store on word. The word at effective address is IORed with the \emph{registerT}. The effective address must be word aligned (multiple of 4).


\paragraph{Modifier(s)}
\noindent\begin{itemize}
\item coherency (cf. coherency in section \ref{par:modifier-coherency} on page \pageref{par:modifier-coherency})
\end{itemize}

\paragraph{Execution} {Reservation table \textsf{LSU2\_MEMW\_AUXR.X} (Section~\ref{sec:reservation-LSU2-MEMW-AUXR.X})}
~
{\begin{lstlisting}
stage ID:
  new offset = _SX_27(offset27);
  new coherency = _ZX_2(coherency);
stage RR:
  new base = GPR[registerZ];
  new address = base + offset;
stage E1:
  new argument2 = GPR[registerT];
stage E3:
  MEM.atomic_ior(address, 0xF, coherency << 32, argument2.32[0], registerT);
\end{lstlisting}}
\newpage

\paragraph{Encoding} Format \textsf{LSU\_ASIORSB.D} (Section~\ref{sec:format-LSU-ASIORSB.D}), \verb|lsucode5 = 010|\\

\bitfields{ 1/P, 2/00, 2/-- --, 27/extend27, 1/1, 2/00, 2/-- --, 27/offset27, 1/1, 2/01, 3/111, 2/coherency, 6/registerT, 2/11, 3/010, 1/--, 5/10110, 1/0, 6/registerZ }


\paragraph{Syntax} \texttt{asiorw}\emph{coherency} \emph{extend27\_\discretionary{}{}{}offset27}[\emph{registerZ}] = \emph{registerT}

\paragraph{Description} The effective address is computed by adding the \emph{extend27\_\discretionary{}{}{}offset27} to the \emph{registerZ}.
This instruction implements atomic IOR-store on word. The word at effective address is IORed with the \emph{registerT}. The effective address must be word aligned (multiple of 4).


\paragraph{Modifier(s)}
\noindent\begin{itemize}
\item coherency (cf. coherency in section \ref{par:modifier-coherency} on page \pageref{par:modifier-coherency})
\end{itemize}

\paragraph{Execution} {Reservation table \textsf{LSU2\_MEMW\_AUXR.Y} (Section~\ref{sec:reservation-LSU2-MEMW-AUXR.Y})}
~
{\begin{lstlisting}
stage ID:
  new offset = _SX_54(extend27_offset27);
  new coherency = _ZX_2(coherency);
stage RR:
  new base = GPR[registerZ];
  new address = base + offset;
stage E1:
  new argument2 = GPR[registerT];
stage E3:
  MEM.atomic_ior(address, 0xF, coherency << 32, argument2.32[0], registerT);
\end{lstlisting}}
\newpage
\subsubsection[{\small ASMAXB: Atomic Store MAX Byte}]{ASMAXB\hfill Atomic Store MAX Byte}
\label{instr:ASMAXB}\index{asmaxb}
 
\paragraph{Encoding} Format \textsf{LSU\_ASMAXSB} (Section~\ref{sec:format-LSU-ASMAXSB}), \verb|lsucode5 = 000|\\

\bitfields{ 1/P, 2/01, 3/111, 2/coherency, 6/registerT, 2/11, 3/000, 1/--, 5/11001, 1/0, 6/registerZ }


\paragraph{Syntax} \texttt{asmaxb}\emph{coherency} [\emph{registerZ}] = \emph{registerT}

\paragraph{Description} The effective address is the value of the \emph{registerZ}.
This instruction implements atomic MAX-store on signed byte. The byte at effective address is MAXed with the \emph{registerT}.


\paragraph{Modifier(s)}
\noindent\begin{itemize}
\item coherency (cf. coherency in section \ref{par:modifier-coherency} on page \pageref{par:modifier-coherency})
\end{itemize}

\paragraph{Execution} {Reservation table \textsf{LSU2\_MEMW\_AUXR} (Section~\ref{sec:reservation-LSU2-MEMW-AUXR})}
~
{\begin{lstlisting}
stage ID:
  new coherency = _ZX_2(coherency);
stage RR:
  new address = GPR[registerZ];
stage E1:
  new argument2 = GPR[registerT];
stage E3:
  MEM.atomic_max(address, 0x1, coherency << 32, argument2.8[0], registerT);
\end{lstlisting}}
\newpage

\paragraph{Encoding} Format \textsf{LSU\_ASMAXSB.O} (Section~\ref{sec:format-LSU-ASMAXSB.O}), \verb|lsucode5 = 000|\\

\bitfields{ 1/P, 2/00, 2/-- --, 27/offset27, 1/1, 2/01, 3/111, 2/coherency, 6/registerT, 2/11, 3/000, 1/--, 5/11001, 1/0, 6/registerZ }


\paragraph{Syntax} \texttt{asmaxb}\emph{coherency} \emph{offset27}[\emph{registerZ}] = \emph{registerT}

\paragraph{Description} The effective address is computed by adding the \emph{offset27} to the \emph{registerZ}.
This instruction implements atomic MAX-store on signed byte. The byte at effective address is MAXed with the \emph{registerT}.


\paragraph{Modifier(s)}
\noindent\begin{itemize}
\item coherency (cf. coherency in section \ref{par:modifier-coherency} on page \pageref{par:modifier-coherency})
\end{itemize}

\paragraph{Execution} {Reservation table \textsf{LSU2\_MEMW\_AUXR.X} (Section~\ref{sec:reservation-LSU2-MEMW-AUXR.X})}
~
{\begin{lstlisting}
stage ID:
  new offset = _SX_27(offset27);
  new coherency = _ZX_2(coherency);
stage RR:
  new base = GPR[registerZ];
  new address = base + offset;
stage E1:
  new argument2 = GPR[registerT];
stage E3:
  MEM.atomic_max(address, 0x1, coherency << 32, argument2.8[0], registerT);
\end{lstlisting}}
\newpage

\paragraph{Encoding} Format \textsf{LSU\_ASMAXSB.D} (Section~\ref{sec:format-LSU-ASMAXSB.D}), \verb|lsucode5 = 000|\\

\bitfields{ 1/P, 2/00, 2/-- --, 27/extend27, 1/1, 2/00, 2/-- --, 27/offset27, 1/1, 2/01, 3/111, 2/coherency, 6/registerT, 2/11, 3/000, 1/--, 5/11001, 1/0, 6/registerZ }


\paragraph{Syntax} \texttt{asmaxb}\emph{coherency} \emph{extend27\_\discretionary{}{}{}offset27}[\emph{registerZ}] = \emph{registerT}

\paragraph{Description} The effective address is computed by adding the \emph{extend27\_\discretionary{}{}{}offset27} to the \emph{registerZ}.
This instruction implements atomic MAX-store on signed byte. The byte at effective address is MAXed with the \emph{registerT}.


\paragraph{Modifier(s)}
\noindent\begin{itemize}
\item coherency (cf. coherency in section \ref{par:modifier-coherency} on page \pageref{par:modifier-coherency})
\end{itemize}

\paragraph{Execution} {Reservation table \textsf{LSU2\_MEMW\_AUXR.Y} (Section~\ref{sec:reservation-LSU2-MEMW-AUXR.Y})}
~
{\begin{lstlisting}
stage ID:
  new offset = _SX_54(extend27_offset27);
  new coherency = _ZX_2(coherency);
stage RR:
  new base = GPR[registerZ];
  new address = base + offset;
stage E1:
  new argument2 = GPR[registerT];
stage E3:
  MEM.atomic_max(address, 0x1, coherency << 32, argument2.8[0], registerT);
\end{lstlisting}}
\newpage
\subsubsection[{\small ASMAXD: Atomic Store MAX Double Word}]{ASMAXD\hfill Atomic Store MAX Double Word}
\label{instr:ASMAXD}\index{asmaxd}
 
\paragraph{Encoding} Format \textsf{LSU\_ASMAXSB} (Section~\ref{sec:format-LSU-ASMAXSB}), \verb|lsucode5 = 011|\\

\bitfields{ 1/P, 2/01, 3/111, 2/coherency, 6/registerT, 2/11, 3/011, 1/--, 5/11001, 1/0, 6/registerZ }


\paragraph{Syntax} \texttt{asmaxd}\emph{coherency} [\emph{registerZ}] = \emph{registerT}

\paragraph{Description} The effective address is the value of the \emph{registerZ}.
This instruction implements atomic MAX-store on signed double word. The double word at effective address is MAXed with the \emph{registerT}. The effective address must be double word aligned (multiple of 8).


\paragraph{Modifier(s)}
\noindent\begin{itemize}
\item coherency (cf. coherency in section \ref{par:modifier-coherency} on page \pageref{par:modifier-coherency})
\end{itemize}

\paragraph{Execution} {Reservation table \textsf{LSU2\_MEMW\_AUXR} (Section~\ref{sec:reservation-LSU2-MEMW-AUXR})}
~
{\begin{lstlisting}
stage ID:
  new coherency = _ZX_2(coherency);
stage RR:
  new address = GPR[registerZ];
stage E1:
  new argument2 = GPR[registerT];
stage E3:
  MEM.atomic_max(address, 0xFF, coherency << 32, argument2.64[0], registerT);
\end{lstlisting}}
\newpage

\paragraph{Encoding} Format \textsf{LSU\_ASMAXSB.O} (Section~\ref{sec:format-LSU-ASMAXSB.O}), \verb|lsucode5 = 011|\\

\bitfields{ 1/P, 2/00, 2/-- --, 27/offset27, 1/1, 2/01, 3/111, 2/coherency, 6/registerT, 2/11, 3/011, 1/--, 5/11001, 1/0, 6/registerZ }


\paragraph{Syntax} \texttt{asmaxd}\emph{coherency} \emph{offset27}[\emph{registerZ}] = \emph{registerT}

\paragraph{Description} The effective address is computed by adding the \emph{offset27} to the \emph{registerZ}.
This instruction implements atomic MAX-store on signed double word. The double word at effective address is MAXed with the \emph{registerT}. The effective address must be double word aligned (multiple of 8).


\paragraph{Modifier(s)}
\noindent\begin{itemize}
\item coherency (cf. coherency in section \ref{par:modifier-coherency} on page \pageref{par:modifier-coherency})
\end{itemize}

\paragraph{Execution} {Reservation table \textsf{LSU2\_MEMW\_AUXR.X} (Section~\ref{sec:reservation-LSU2-MEMW-AUXR.X})}
~
{\begin{lstlisting}
stage ID:
  new offset = _SX_27(offset27);
  new coherency = _ZX_2(coherency);
stage RR:
  new base = GPR[registerZ];
  new address = base + offset;
stage E1:
  new argument2 = GPR[registerT];
stage E3:
  MEM.atomic_max(address, 0xFF, coherency << 32, argument2.64[0], registerT);
\end{lstlisting}}
\newpage

\paragraph{Encoding} Format \textsf{LSU\_ASMAXSB.D} (Section~\ref{sec:format-LSU-ASMAXSB.D}), \verb|lsucode5 = 011|\\

\bitfields{ 1/P, 2/00, 2/-- --, 27/extend27, 1/1, 2/00, 2/-- --, 27/offset27, 1/1, 2/01, 3/111, 2/coherency, 6/registerT, 2/11, 3/011, 1/--, 5/11001, 1/0, 6/registerZ }


\paragraph{Syntax} \texttt{asmaxd}\emph{coherency} \emph{extend27\_\discretionary{}{}{}offset27}[\emph{registerZ}] = \emph{registerT}

\paragraph{Description} The effective address is computed by adding the \emph{extend27\_\discretionary{}{}{}offset27} to the \emph{registerZ}.
This instruction implements atomic MAX-store on signed double word. The double word at effective address is MAXed with the \emph{registerT}. The effective address must be double word aligned (multiple of 8).


\paragraph{Modifier(s)}
\noindent\begin{itemize}
\item coherency (cf. coherency in section \ref{par:modifier-coherency} on page \pageref{par:modifier-coherency})
\end{itemize}

\paragraph{Execution} {Reservation table \textsf{LSU2\_MEMW\_AUXR.Y} (Section~\ref{sec:reservation-LSU2-MEMW-AUXR.Y})}
~
{\begin{lstlisting}
stage ID:
  new offset = _SX_54(extend27_offset27);
  new coherency = _ZX_2(coherency);
stage RR:
  new base = GPR[registerZ];
  new address = base + offset;
stage E1:
  new argument2 = GPR[registerT];
stage E3:
  MEM.atomic_max(address, 0xFF, coherency << 32, argument2.64[0], registerT);
\end{lstlisting}}
\newpage
\subsubsection[{\small ASMAXH: Atomic Store MAX Half Word}]{ASMAXH\hfill Atomic Store MAX Half Word}
\label{instr:ASMAXH}\index{asmaxh}
 
\paragraph{Encoding} Format \textsf{LSU\_ASMAXSB} (Section~\ref{sec:format-LSU-ASMAXSB}), \verb|lsucode5 = 001|\\

\bitfields{ 1/P, 2/01, 3/111, 2/coherency, 6/registerT, 2/11, 3/001, 1/--, 5/11001, 1/0, 6/registerZ }


\paragraph{Syntax} \texttt{asmaxh}\emph{coherency} [\emph{registerZ}] = \emph{registerT}

\paragraph{Description} The effective address is the value of the \emph{registerZ}.
This instruction implements atomic MAX-store on signed half word. The half word at effective address is MAXed with the \emph{registerT}. The effective address must be half word aligned (multiple of 2).


\paragraph{Modifier(s)}
\noindent\begin{itemize}
\item coherency (cf. coherency in section \ref{par:modifier-coherency} on page \pageref{par:modifier-coherency})
\end{itemize}

\paragraph{Execution} {Reservation table \textsf{LSU2\_MEMW\_AUXR} (Section~\ref{sec:reservation-LSU2-MEMW-AUXR})}
~
{\begin{lstlisting}
stage ID:
  new coherency = _ZX_2(coherency);
stage RR:
  new address = GPR[registerZ];
stage E1:
  new argument2 = GPR[registerT];
stage E3:
  MEM.atomic_max(address, 0x3, coherency << 32, argument2.16[0], registerT);
\end{lstlisting}}
\newpage

\paragraph{Encoding} Format \textsf{LSU\_ASMAXSB.O} (Section~\ref{sec:format-LSU-ASMAXSB.O}), \verb|lsucode5 = 001|\\

\bitfields{ 1/P, 2/00, 2/-- --, 27/offset27, 1/1, 2/01, 3/111, 2/coherency, 6/registerT, 2/11, 3/001, 1/--, 5/11001, 1/0, 6/registerZ }


\paragraph{Syntax} \texttt{asmaxh}\emph{coherency} \emph{offset27}[\emph{registerZ}] = \emph{registerT}

\paragraph{Description} The effective address is computed by adding the \emph{offset27} to the \emph{registerZ}.
This instruction implements atomic MAX-store on signed half word. The half word at effective address is MAXed with the \emph{registerT}. The effective address must be half word aligned (multiple of 2).


\paragraph{Modifier(s)}
\noindent\begin{itemize}
\item coherency (cf. coherency in section \ref{par:modifier-coherency} on page \pageref{par:modifier-coherency})
\end{itemize}

\paragraph{Execution} {Reservation table \textsf{LSU2\_MEMW\_AUXR.X} (Section~\ref{sec:reservation-LSU2-MEMW-AUXR.X})}
~
{\begin{lstlisting}
stage ID:
  new offset = _SX_27(offset27);
  new coherency = _ZX_2(coherency);
stage RR:
  new base = GPR[registerZ];
  new address = base + offset;
stage E1:
  new argument2 = GPR[registerT];
stage E3:
  MEM.atomic_max(address, 0x3, coherency << 32, argument2.16[0], registerT);
\end{lstlisting}}
\newpage

\paragraph{Encoding} Format \textsf{LSU\_ASMAXSB.D} (Section~\ref{sec:format-LSU-ASMAXSB.D}), \verb|lsucode5 = 001|\\

\bitfields{ 1/P, 2/00, 2/-- --, 27/extend27, 1/1, 2/00, 2/-- --, 27/offset27, 1/1, 2/01, 3/111, 2/coherency, 6/registerT, 2/11, 3/001, 1/--, 5/11001, 1/0, 6/registerZ }


\paragraph{Syntax} \texttt{asmaxh}\emph{coherency} \emph{extend27\_\discretionary{}{}{}offset27}[\emph{registerZ}] = \emph{registerT}

\paragraph{Description} The effective address is computed by adding the \emph{extend27\_\discretionary{}{}{}offset27} to the \emph{registerZ}.
This instruction implements atomic MAX-store on signed half word. The half word at effective address is MAXed with the \emph{registerT}. The effective address must be half word aligned (multiple of 2).


\paragraph{Modifier(s)}
\noindent\begin{itemize}
\item coherency (cf. coherency in section \ref{par:modifier-coherency} on page \pageref{par:modifier-coherency})
\end{itemize}

\paragraph{Execution} {Reservation table \textsf{LSU2\_MEMW\_AUXR.Y} (Section~\ref{sec:reservation-LSU2-MEMW-AUXR.Y})}
~
{\begin{lstlisting}
stage ID:
  new offset = _SX_54(extend27_offset27);
  new coherency = _ZX_2(coherency);
stage RR:
  new base = GPR[registerZ];
  new address = base + offset;
stage E1:
  new argument2 = GPR[registerT];
stage E3:
  MEM.atomic_max(address, 0x3, coherency << 32, argument2.16[0], registerT);
\end{lstlisting}}
\newpage
\subsubsection[{\small ASMAXUB: Atomic Store MAX Unsigned Byte}]{ASMAXUB\hfill Atomic Store MAX Unsigned Byte}
\label{instr:ASMAXUB}\index{asmaxub}
 
\paragraph{Encoding} Format \textsf{LSU\_ASMAXUSB} (Section~\ref{sec:format-LSU-ASMAXUSB}), \verb|lsucode5 = 000|\\

\bitfields{ 1/P, 2/01, 3/111, 2/coherency, 6/registerT, 2/11, 3/000, 1/--, 5/11011, 1/0, 6/registerZ }


\paragraph{Syntax} \texttt{asmaxub}\emph{coherency} [\emph{registerZ}] = \emph{registerT}

\paragraph{Description} The effective address is the value of the \emph{registerZ}.
This instruction implements atomic MAXU-store on signed byte. The byte at effective address is MAXUed with the \emph{registerT}.


\paragraph{Modifier(s)}
\noindent\begin{itemize}
\item coherency (cf. coherency in section \ref{par:modifier-coherency} on page \pageref{par:modifier-coherency})
\end{itemize}

\paragraph{Execution} {Reservation table \textsf{LSU2\_MEMW\_AUXR} (Section~\ref{sec:reservation-LSU2-MEMW-AUXR})}
~
{\begin{lstlisting}
stage ID:
  new coherency = _ZX_2(coherency);
stage RR:
  new address = GPR[registerZ];
stage E1:
  new argument2 = GPR[registerT];
stage E3:
  MEM.atomic_maxu(address, 0x1, coherency << 32, argument2.8[0], registerT);
\end{lstlisting}}
\newpage

\paragraph{Encoding} Format \textsf{LSU\_ASMAXUSB.O} (Section~\ref{sec:format-LSU-ASMAXUSB.O}), \verb|lsucode5 = 000|\\

\bitfields{ 1/P, 2/00, 2/-- --, 27/offset27, 1/1, 2/01, 3/111, 2/coherency, 6/registerT, 2/11, 3/000, 1/--, 5/11011, 1/0, 6/registerZ }


\paragraph{Syntax} \texttt{asmaxub}\emph{coherency} \emph{offset27}[\emph{registerZ}] = \emph{registerT}

\paragraph{Description} The effective address is computed by adding the \emph{offset27} to the \emph{registerZ}.
This instruction implements atomic MAXU-store on signed byte. The byte at effective address is MAXUed with the \emph{registerT}.


\paragraph{Modifier(s)}
\noindent\begin{itemize}
\item coherency (cf. coherency in section \ref{par:modifier-coherency} on page \pageref{par:modifier-coherency})
\end{itemize}

\paragraph{Execution} {Reservation table \textsf{LSU2\_MEMW\_AUXR.X} (Section~\ref{sec:reservation-LSU2-MEMW-AUXR.X})}
~
{\begin{lstlisting}
stage ID:
  new offset = _SX_27(offset27);
  new coherency = _ZX_2(coherency);
stage RR:
  new base = GPR[registerZ];
  new address = base + offset;
stage E1:
  new argument2 = GPR[registerT];
stage E3:
  MEM.atomic_maxu(address, 0x1, coherency << 32, argument2.8[0], registerT);
\end{lstlisting}}
\newpage

\paragraph{Encoding} Format \textsf{LSU\_ASMAXUSB.D} (Section~\ref{sec:format-LSU-ASMAXUSB.D}), \verb|lsucode5 = 000|\\

\bitfields{ 1/P, 2/00, 2/-- --, 27/extend27, 1/1, 2/00, 2/-- --, 27/offset27, 1/1, 2/01, 3/111, 2/coherency, 6/registerT, 2/11, 3/000, 1/--, 5/11011, 1/0, 6/registerZ }


\paragraph{Syntax} \texttt{asmaxub}\emph{coherency} \emph{extend27\_\discretionary{}{}{}offset27}[\emph{registerZ}] = \emph{registerT}

\paragraph{Description} The effective address is computed by adding the \emph{extend27\_\discretionary{}{}{}offset27} to the \emph{registerZ}.
This instruction implements atomic MAXU-store on signed byte. The byte at effective address is MAXUed with the \emph{registerT}.


\paragraph{Modifier(s)}
\noindent\begin{itemize}
\item coherency (cf. coherency in section \ref{par:modifier-coherency} on page \pageref{par:modifier-coherency})
\end{itemize}

\paragraph{Execution} {Reservation table \textsf{LSU2\_MEMW\_AUXR.Y} (Section~\ref{sec:reservation-LSU2-MEMW-AUXR.Y})}
~
{\begin{lstlisting}
stage ID:
  new offset = _SX_54(extend27_offset27);
  new coherency = _ZX_2(coherency);
stage RR:
  new base = GPR[registerZ];
  new address = base + offset;
stage E1:
  new argument2 = GPR[registerT];
stage E3:
  MEM.atomic_maxu(address, 0x1, coherency << 32, argument2.8[0], registerT);
\end{lstlisting}}
\newpage
\subsubsection[{\small ASMAXUD: Atomic Store MAX Unsigned Double Word}]{ASMAXUD\hfill Atomic Store MAX Unsigned Double Word}
\label{instr:ASMAXUD}\index{asmaxud}
 
\paragraph{Encoding} Format \textsf{LSU\_ASMAXUSB} (Section~\ref{sec:format-LSU-ASMAXUSB}), \verb|lsucode5 = 011|\\

\bitfields{ 1/P, 2/01, 3/111, 2/coherency, 6/registerT, 2/11, 3/011, 1/--, 5/11011, 1/0, 6/registerZ }


\paragraph{Syntax} \texttt{asmaxud}\emph{coherency} [\emph{registerZ}] = \emph{registerT}

\paragraph{Description} The effective address is the value of the \emph{registerZ}.
This instruction implements atomic MAXU-store on signed double word. The double word at effective address is MAXUed with the \emph{registerT}. The effective address must be double word aligned (multiple of 8).


\paragraph{Modifier(s)}
\noindent\begin{itemize}
\item coherency (cf. coherency in section \ref{par:modifier-coherency} on page \pageref{par:modifier-coherency})
\end{itemize}

\paragraph{Execution} {Reservation table \textsf{LSU2\_MEMW\_AUXR} (Section~\ref{sec:reservation-LSU2-MEMW-AUXR})}
~
{\begin{lstlisting}
stage ID:
  new coherency = _ZX_2(coherency);
stage RR:
  new address = GPR[registerZ];
stage E1:
  new argument2 = GPR[registerT];
stage E3:
  MEM.atomic_maxu(address, 0xFF, coherency << 32, argument2.64[0], registerT);
\end{lstlisting}}
\newpage

\paragraph{Encoding} Format \textsf{LSU\_ASMAXUSB.O} (Section~\ref{sec:format-LSU-ASMAXUSB.O}), \verb|lsucode5 = 011|\\

\bitfields{ 1/P, 2/00, 2/-- --, 27/offset27, 1/1, 2/01, 3/111, 2/coherency, 6/registerT, 2/11, 3/011, 1/--, 5/11011, 1/0, 6/registerZ }


\paragraph{Syntax} \texttt{asmaxud}\emph{coherency} \emph{offset27}[\emph{registerZ}] = \emph{registerT}

\paragraph{Description} The effective address is computed by adding the \emph{offset27} to the \emph{registerZ}.
This instruction implements atomic MAXU-store on signed double word. The double word at effective address is MAXUed with the \emph{registerT}. The effective address must be double word aligned (multiple of 8).


\paragraph{Modifier(s)}
\noindent\begin{itemize}
\item coherency (cf. coherency in section \ref{par:modifier-coherency} on page \pageref{par:modifier-coherency})
\end{itemize}

\paragraph{Execution} {Reservation table \textsf{LSU2\_MEMW\_AUXR.X} (Section~\ref{sec:reservation-LSU2-MEMW-AUXR.X})}
~
{\begin{lstlisting}
stage ID:
  new offset = _SX_27(offset27);
  new coherency = _ZX_2(coherency);
stage RR:
  new base = GPR[registerZ];
  new address = base + offset;
stage E1:
  new argument2 = GPR[registerT];
stage E3:
  MEM.atomic_maxu(address, 0xFF, coherency << 32, argument2.64[0], registerT);
\end{lstlisting}}
\newpage

\paragraph{Encoding} Format \textsf{LSU\_ASMAXUSB.D} (Section~\ref{sec:format-LSU-ASMAXUSB.D}), \verb|lsucode5 = 011|\\

\bitfields{ 1/P, 2/00, 2/-- --, 27/extend27, 1/1, 2/00, 2/-- --, 27/offset27, 1/1, 2/01, 3/111, 2/coherency, 6/registerT, 2/11, 3/011, 1/--, 5/11011, 1/0, 6/registerZ }


\paragraph{Syntax} \texttt{asmaxud}\emph{coherency} \emph{extend27\_\discretionary{}{}{}offset27}[\emph{registerZ}] = \emph{registerT}

\paragraph{Description} The effective address is computed by adding the \emph{extend27\_\discretionary{}{}{}offset27} to the \emph{registerZ}.
This instruction implements atomic MAXU-store on signed double word. The double word at effective address is MAXUed with the \emph{registerT}. The effective address must be double word aligned (multiple of 8).


\paragraph{Modifier(s)}
\noindent\begin{itemize}
\item coherency (cf. coherency in section \ref{par:modifier-coherency} on page \pageref{par:modifier-coherency})
\end{itemize}

\paragraph{Execution} {Reservation table \textsf{LSU2\_MEMW\_AUXR.Y} (Section~\ref{sec:reservation-LSU2-MEMW-AUXR.Y})}
~
{\begin{lstlisting}
stage ID:
  new offset = _SX_54(extend27_offset27);
  new coherency = _ZX_2(coherency);
stage RR:
  new base = GPR[registerZ];
  new address = base + offset;
stage E1:
  new argument2 = GPR[registerT];
stage E3:
  MEM.atomic_maxu(address, 0xFF, coherency << 32, argument2.64[0], registerT);
\end{lstlisting}}
\newpage
\subsubsection[{\small ASMAXUH: Atomic Store MAX Unsigned Half Word}]{ASMAXUH\hfill Atomic Store MAX Unsigned Half Word}
\label{instr:ASMAXUH}\index{asmaxuh}
 
\paragraph{Encoding} Format \textsf{LSU\_ASMAXUSB} (Section~\ref{sec:format-LSU-ASMAXUSB}), \verb|lsucode5 = 001|\\

\bitfields{ 1/P, 2/01, 3/111, 2/coherency, 6/registerT, 2/11, 3/001, 1/--, 5/11011, 1/0, 6/registerZ }


\paragraph{Syntax} \texttt{asmaxuh}\emph{coherency} [\emph{registerZ}] = \emph{registerT}

\paragraph{Description} The effective address is the value of the \emph{registerZ}.
This instruction implements atomic MAXU-store on signed half word. The half word at effective address is MAXUed with the \emph{registerT}. The effective address must be half word aligned (multiple of 2).


\paragraph{Modifier(s)}
\noindent\begin{itemize}
\item coherency (cf. coherency in section \ref{par:modifier-coherency} on page \pageref{par:modifier-coherency})
\end{itemize}

\paragraph{Execution} {Reservation table \textsf{LSU2\_MEMW\_AUXR} (Section~\ref{sec:reservation-LSU2-MEMW-AUXR})}
~
{\begin{lstlisting}
stage ID:
  new coherency = _ZX_2(coherency);
stage RR:
  new address = GPR[registerZ];
stage E1:
  new argument2 = GPR[registerT];
stage E3:
  MEM.atomic_maxu(address, 0x3, coherency << 32, argument2.16[0], registerT);
\end{lstlisting}}
\newpage

\paragraph{Encoding} Format \textsf{LSU\_ASMAXUSB.O} (Section~\ref{sec:format-LSU-ASMAXUSB.O}), \verb|lsucode5 = 001|\\

\bitfields{ 1/P, 2/00, 2/-- --, 27/offset27, 1/1, 2/01, 3/111, 2/coherency, 6/registerT, 2/11, 3/001, 1/--, 5/11011, 1/0, 6/registerZ }


\paragraph{Syntax} \texttt{asmaxuh}\emph{coherency} \emph{offset27}[\emph{registerZ}] = \emph{registerT}

\paragraph{Description} The effective address is computed by adding the \emph{offset27} to the \emph{registerZ}.
This instruction implements atomic MAXU-store on signed half word. The half word at effective address is MAXUed with the \emph{registerT}. The effective address must be half word aligned (multiple of 2).


\paragraph{Modifier(s)}
\noindent\begin{itemize}
\item coherency (cf. coherency in section \ref{par:modifier-coherency} on page \pageref{par:modifier-coherency})
\end{itemize}

\paragraph{Execution} {Reservation table \textsf{LSU2\_MEMW\_AUXR.X} (Section~\ref{sec:reservation-LSU2-MEMW-AUXR.X})}
~
{\begin{lstlisting}
stage ID:
  new offset = _SX_27(offset27);
  new coherency = _ZX_2(coherency);
stage RR:
  new base = GPR[registerZ];
  new address = base + offset;
stage E1:
  new argument2 = GPR[registerT];
stage E3:
  MEM.atomic_maxu(address, 0x3, coherency << 32, argument2.16[0], registerT);
\end{lstlisting}}
\newpage

\paragraph{Encoding} Format \textsf{LSU\_ASMAXUSB.D} (Section~\ref{sec:format-LSU-ASMAXUSB.D}), \verb|lsucode5 = 001|\\

\bitfields{ 1/P, 2/00, 2/-- --, 27/extend27, 1/1, 2/00, 2/-- --, 27/offset27, 1/1, 2/01, 3/111, 2/coherency, 6/registerT, 2/11, 3/001, 1/--, 5/11011, 1/0, 6/registerZ }


\paragraph{Syntax} \texttt{asmaxuh}\emph{coherency} \emph{extend27\_\discretionary{}{}{}offset27}[\emph{registerZ}] = \emph{registerT}

\paragraph{Description} The effective address is computed by adding the \emph{extend27\_\discretionary{}{}{}offset27} to the \emph{registerZ}.
This instruction implements atomic MAXU-store on signed half word. The half word at effective address is MAXUed with the \emph{registerT}. The effective address must be half word aligned (multiple of 2).


\paragraph{Modifier(s)}
\noindent\begin{itemize}
\item coherency (cf. coherency in section \ref{par:modifier-coherency} on page \pageref{par:modifier-coherency})
\end{itemize}

\paragraph{Execution} {Reservation table \textsf{LSU2\_MEMW\_AUXR.Y} (Section~\ref{sec:reservation-LSU2-MEMW-AUXR.Y})}
~
{\begin{lstlisting}
stage ID:
  new offset = _SX_54(extend27_offset27);
  new coherency = _ZX_2(coherency);
stage RR:
  new base = GPR[registerZ];
  new address = base + offset;
stage E1:
  new argument2 = GPR[registerT];
stage E3:
  MEM.atomic_maxu(address, 0x3, coherency << 32, argument2.16[0], registerT);
\end{lstlisting}}
\newpage
\subsubsection[{\small ASMAXUW: Atomic Store MAX Unsigned Word}]{ASMAXUW\hfill Atomic Store MAX Unsigned Word}
\label{instr:ASMAXUW}\index{asmaxuw}
 
\paragraph{Encoding} Format \textsf{LSU\_ASMAXUSB} (Section~\ref{sec:format-LSU-ASMAXUSB}), \verb|lsucode5 = 010|\\

\bitfields{ 1/P, 2/01, 3/111, 2/coherency, 6/registerT, 2/11, 3/010, 1/--, 5/11011, 1/0, 6/registerZ }


\paragraph{Syntax} \texttt{asmaxuw}\emph{coherency} [\emph{registerZ}] = \emph{registerT}

\paragraph{Description} The effective address is the value of the \emph{registerZ}.
This instruction implements atomic MAXU-store on signed word. The word at effective address is MAXUed with the \emph{registerT}. The effective address must be word aligned (multiple of 4).


\paragraph{Modifier(s)}
\noindent\begin{itemize}
\item coherency (cf. coherency in section \ref{par:modifier-coherency} on page \pageref{par:modifier-coherency})
\end{itemize}

\paragraph{Execution} {Reservation table \textsf{LSU2\_MEMW\_AUXR} (Section~\ref{sec:reservation-LSU2-MEMW-AUXR})}
~
{\begin{lstlisting}
stage ID:
  new coherency = _ZX_2(coherency);
stage RR:
  new address = GPR[registerZ];
stage E1:
  new argument2 = GPR[registerT];
stage E3:
  MEM.atomic_maxu(address, 0xF, coherency << 32, argument2.32[0], registerT);
\end{lstlisting}}
\newpage

\paragraph{Encoding} Format \textsf{LSU\_ASMAXUSB.O} (Section~\ref{sec:format-LSU-ASMAXUSB.O}), \verb|lsucode5 = 010|\\

\bitfields{ 1/P, 2/00, 2/-- --, 27/offset27, 1/1, 2/01, 3/111, 2/coherency, 6/registerT, 2/11, 3/010, 1/--, 5/11011, 1/0, 6/registerZ }


\paragraph{Syntax} \texttt{asmaxuw}\emph{coherency} \emph{offset27}[\emph{registerZ}] = \emph{registerT}

\paragraph{Description} The effective address is computed by adding the \emph{offset27} to the \emph{registerZ}.
This instruction implements atomic MAXU-store on signed word. The word at effective address is MAXUed with the \emph{registerT}. The effective address must be word aligned (multiple of 4).


\paragraph{Modifier(s)}
\noindent\begin{itemize}
\item coherency (cf. coherency in section \ref{par:modifier-coherency} on page \pageref{par:modifier-coherency})
\end{itemize}

\paragraph{Execution} {Reservation table \textsf{LSU2\_MEMW\_AUXR.X} (Section~\ref{sec:reservation-LSU2-MEMW-AUXR.X})}
~
{\begin{lstlisting}
stage ID:
  new offset = _SX_27(offset27);
  new coherency = _ZX_2(coherency);
stage RR:
  new base = GPR[registerZ];
  new address = base + offset;
stage E1:
  new argument2 = GPR[registerT];
stage E3:
  MEM.atomic_maxu(address, 0xF, coherency << 32, argument2.32[0], registerT);
\end{lstlisting}}
\newpage

\paragraph{Encoding} Format \textsf{LSU\_ASMAXUSB.D} (Section~\ref{sec:format-LSU-ASMAXUSB.D}), \verb|lsucode5 = 010|\\

\bitfields{ 1/P, 2/00, 2/-- --, 27/extend27, 1/1, 2/00, 2/-- --, 27/offset27, 1/1, 2/01, 3/111, 2/coherency, 6/registerT, 2/11, 3/010, 1/--, 5/11011, 1/0, 6/registerZ }


\paragraph{Syntax} \texttt{asmaxuw}\emph{coherency} \emph{extend27\_\discretionary{}{}{}offset27}[\emph{registerZ}] = \emph{registerT}

\paragraph{Description} The effective address is computed by adding the \emph{extend27\_\discretionary{}{}{}offset27} to the \emph{registerZ}.
This instruction implements atomic MAXU-store on signed word. The word at effective address is MAXUed with the \emph{registerT}. The effective address must be word aligned (multiple of 4).


\paragraph{Modifier(s)}
\noindent\begin{itemize}
\item coherency (cf. coherency in section \ref{par:modifier-coherency} on page \pageref{par:modifier-coherency})
\end{itemize}

\paragraph{Execution} {Reservation table \textsf{LSU2\_MEMW\_AUXR.Y} (Section~\ref{sec:reservation-LSU2-MEMW-AUXR.Y})}
~
{\begin{lstlisting}
stage ID:
  new offset = _SX_54(extend27_offset27);
  new coherency = _ZX_2(coherency);
stage RR:
  new base = GPR[registerZ];
  new address = base + offset;
stage E1:
  new argument2 = GPR[registerT];
stage E3:
  MEM.atomic_maxu(address, 0xF, coherency << 32, argument2.32[0], registerT);
\end{lstlisting}}
\newpage
\subsubsection[{\small ASMAXW: Atomic Store MAX Word}]{ASMAXW\hfill Atomic Store MAX Word}
\label{instr:ASMAXW}\index{asmaxw}
 
\paragraph{Encoding} Format \textsf{LSU\_ASMAXSB} (Section~\ref{sec:format-LSU-ASMAXSB}), \verb|lsucode5 = 010|\\

\bitfields{ 1/P, 2/01, 3/111, 2/coherency, 6/registerT, 2/11, 3/010, 1/--, 5/11001, 1/0, 6/registerZ }


\paragraph{Syntax} \texttt{asmaxw}\emph{coherency} [\emph{registerZ}] = \emph{registerT}

\paragraph{Description} The effective address is the value of the \emph{registerZ}.
This instruction implements atomic MAX-store on signed word. The word at effective address is MAXed with the \emph{registerT}. The effective address must be word aligned (multiple of 4).


\paragraph{Modifier(s)}
\noindent\begin{itemize}
\item coherency (cf. coherency in section \ref{par:modifier-coherency} on page \pageref{par:modifier-coherency})
\end{itemize}

\paragraph{Execution} {Reservation table \textsf{LSU2\_MEMW\_AUXR} (Section~\ref{sec:reservation-LSU2-MEMW-AUXR})}
~
{\begin{lstlisting}
stage ID:
  new coherency = _ZX_2(coherency);
stage RR:
  new address = GPR[registerZ];
stage E1:
  new argument2 = GPR[registerT];
stage E3:
  MEM.atomic_max(address, 0xF, coherency << 32, argument2.32[0], registerT);
\end{lstlisting}}
\newpage

\paragraph{Encoding} Format \textsf{LSU\_ASMAXSB.O} (Section~\ref{sec:format-LSU-ASMAXSB.O}), \verb|lsucode5 = 010|\\

\bitfields{ 1/P, 2/00, 2/-- --, 27/offset27, 1/1, 2/01, 3/111, 2/coherency, 6/registerT, 2/11, 3/010, 1/--, 5/11001, 1/0, 6/registerZ }


\paragraph{Syntax} \texttt{asmaxw}\emph{coherency} \emph{offset27}[\emph{registerZ}] = \emph{registerT}

\paragraph{Description} The effective address is computed by adding the \emph{offset27} to the \emph{registerZ}.
This instruction implements atomic MAX-store on signed word. The word at effective address is MAXed with the \emph{registerT}. The effective address must be word aligned (multiple of 4).


\paragraph{Modifier(s)}
\noindent\begin{itemize}
\item coherency (cf. coherency in section \ref{par:modifier-coherency} on page \pageref{par:modifier-coherency})
\end{itemize}

\paragraph{Execution} {Reservation table \textsf{LSU2\_MEMW\_AUXR.X} (Section~\ref{sec:reservation-LSU2-MEMW-AUXR.X})}
~
{\begin{lstlisting}
stage ID:
  new offset = _SX_27(offset27);
  new coherency = _ZX_2(coherency);
stage RR:
  new base = GPR[registerZ];
  new address = base + offset;
stage E1:
  new argument2 = GPR[registerT];
stage E3:
  MEM.atomic_max(address, 0xF, coherency << 32, argument2.32[0], registerT);
\end{lstlisting}}
\newpage

\paragraph{Encoding} Format \textsf{LSU\_ASMAXSB.D} (Section~\ref{sec:format-LSU-ASMAXSB.D}), \verb|lsucode5 = 010|\\

\bitfields{ 1/P, 2/00, 2/-- --, 27/extend27, 1/1, 2/00, 2/-- --, 27/offset27, 1/1, 2/01, 3/111, 2/coherency, 6/registerT, 2/11, 3/010, 1/--, 5/11001, 1/0, 6/registerZ }


\paragraph{Syntax} \texttt{asmaxw}\emph{coherency} \emph{extend27\_\discretionary{}{}{}offset27}[\emph{registerZ}] = \emph{registerT}

\paragraph{Description} The effective address is computed by adding the \emph{extend27\_\discretionary{}{}{}offset27} to the \emph{registerZ}.
This instruction implements atomic MAX-store on signed word. The word at effective address is MAXed with the \emph{registerT}. The effective address must be word aligned (multiple of 4).


\paragraph{Modifier(s)}
\noindent\begin{itemize}
\item coherency (cf. coherency in section \ref{par:modifier-coherency} on page \pageref{par:modifier-coherency})
\end{itemize}

\paragraph{Execution} {Reservation table \textsf{LSU2\_MEMW\_AUXR.Y} (Section~\ref{sec:reservation-LSU2-MEMW-AUXR.Y})}
~
{\begin{lstlisting}
stage ID:
  new offset = _SX_54(extend27_offset27);
  new coherency = _ZX_2(coherency);
stage RR:
  new base = GPR[registerZ];
  new address = base + offset;
stage E1:
  new argument2 = GPR[registerT];
stage E3:
  MEM.atomic_max(address, 0xF, coherency << 32, argument2.32[0], registerT);
\end{lstlisting}}
\newpage
\subsubsection[{\small ASMINB: Atomic Store MIN Byte}]{ASMINB\hfill Atomic Store MIN Byte}
\label{instr:ASMINB}\index{asminb}
 
\paragraph{Encoding} Format \textsf{LSU\_ASMINSB} (Section~\ref{sec:format-LSU-ASMINSB}), \verb|lsucode5 = 000|\\

\bitfields{ 1/P, 2/01, 3/111, 2/coherency, 6/registerT, 2/11, 3/000, 1/--, 5/11000, 1/0, 6/registerZ }


\paragraph{Syntax} \texttt{asminb}\emph{coherency} [\emph{registerZ}] = \emph{registerT}

\paragraph{Description} The effective address is the value of the \emph{registerZ}.
This instruction implements atomic MIN-store on signed byte. The byte at effective address is MINed with the \emph{registerT}.


\paragraph{Modifier(s)}
\noindent\begin{itemize}
\item coherency (cf. coherency in section \ref{par:modifier-coherency} on page \pageref{par:modifier-coherency})
\end{itemize}

\paragraph{Execution} {Reservation table \textsf{LSU2\_MEMW\_AUXR} (Section~\ref{sec:reservation-LSU2-MEMW-AUXR})}
~
{\begin{lstlisting}
stage ID:
  new coherency = _ZX_2(coherency);
stage RR:
  new address = GPR[registerZ];
stage E1:
  new argument2 = GPR[registerT];
stage E3:
  MEM.atomic_min(address, 0x1, coherency << 32, argument2.8[0], registerT);
\end{lstlisting}}
\newpage

\paragraph{Encoding} Format \textsf{LSU\_ASMINSB.O} (Section~\ref{sec:format-LSU-ASMINSB.O}), \verb|lsucode5 = 000|\\

\bitfields{ 1/P, 2/00, 2/-- --, 27/offset27, 1/1, 2/01, 3/111, 2/coherency, 6/registerT, 2/11, 3/000, 1/--, 5/11000, 1/0, 6/registerZ }


\paragraph{Syntax} \texttt{asminb}\emph{coherency} \emph{offset27}[\emph{registerZ}] = \emph{registerT}

\paragraph{Description} The effective address is computed by adding the \emph{offset27} to the \emph{registerZ}.
This instruction implements atomic MIN-store on signed byte. The byte at effective address is MINed with the \emph{registerT}.


\paragraph{Modifier(s)}
\noindent\begin{itemize}
\item coherency (cf. coherency in section \ref{par:modifier-coherency} on page \pageref{par:modifier-coherency})
\end{itemize}

\paragraph{Execution} {Reservation table \textsf{LSU2\_MEMW\_AUXR.X} (Section~\ref{sec:reservation-LSU2-MEMW-AUXR.X})}
~
{\begin{lstlisting}
stage ID:
  new offset = _SX_27(offset27);
  new coherency = _ZX_2(coherency);
stage RR:
  new base = GPR[registerZ];
  new address = base + offset;
stage E1:
  new argument2 = GPR[registerT];
stage E3:
  MEM.atomic_min(address, 0x1, coherency << 32, argument2.8[0], registerT);
\end{lstlisting}}
\newpage

\paragraph{Encoding} Format \textsf{LSU\_ASMINSB.D} (Section~\ref{sec:format-LSU-ASMINSB.D}), \verb|lsucode5 = 000|\\

\bitfields{ 1/P, 2/00, 2/-- --, 27/extend27, 1/1, 2/00, 2/-- --, 27/offset27, 1/1, 2/01, 3/111, 2/coherency, 6/registerT, 2/11, 3/000, 1/--, 5/11000, 1/0, 6/registerZ }


\paragraph{Syntax} \texttt{asminb}\emph{coherency} \emph{extend27\_\discretionary{}{}{}offset27}[\emph{registerZ}] = \emph{registerT}

\paragraph{Description} The effective address is computed by adding the \emph{extend27\_\discretionary{}{}{}offset27} to the \emph{registerZ}.
This instruction implements atomic MIN-store on signed byte. The byte at effective address is MINed with the \emph{registerT}.


\paragraph{Modifier(s)}
\noindent\begin{itemize}
\item coherency (cf. coherency in section \ref{par:modifier-coherency} on page \pageref{par:modifier-coherency})
\end{itemize}

\paragraph{Execution} {Reservation table \textsf{LSU2\_MEMW\_AUXR.Y} (Section~\ref{sec:reservation-LSU2-MEMW-AUXR.Y})}
~
{\begin{lstlisting}
stage ID:
  new offset = _SX_54(extend27_offset27);
  new coherency = _ZX_2(coherency);
stage RR:
  new base = GPR[registerZ];
  new address = base + offset;
stage E1:
  new argument2 = GPR[registerT];
stage E3:
  MEM.atomic_min(address, 0x1, coherency << 32, argument2.8[0], registerT);
\end{lstlisting}}
\newpage
\subsubsection[{\small ASMIND: Atomic Store MIN Double Word}]{ASMIND\hfill Atomic Store MIN Double Word}
\label{instr:ASMIND}\index{asmind}
 
\paragraph{Encoding} Format \textsf{LSU\_ASMINSB} (Section~\ref{sec:format-LSU-ASMINSB}), \verb|lsucode5 = 011|\\

\bitfields{ 1/P, 2/01, 3/111, 2/coherency, 6/registerT, 2/11, 3/011, 1/--, 5/11000, 1/0, 6/registerZ }


\paragraph{Syntax} \texttt{asmind}\emph{coherency} [\emph{registerZ}] = \emph{registerT}

\paragraph{Description} The effective address is the value of the \emph{registerZ}.
This instruction implements atomic MIN-store on signed double word. The double word at effective address is MINed with the \emph{registerT}. The effective address must be double word aligned (multiple of 8).


\paragraph{Modifier(s)}
\noindent\begin{itemize}
\item coherency (cf. coherency in section \ref{par:modifier-coherency} on page \pageref{par:modifier-coherency})
\end{itemize}

\paragraph{Execution} {Reservation table \textsf{LSU2\_MEMW\_AUXR} (Section~\ref{sec:reservation-LSU2-MEMW-AUXR})}
~
{\begin{lstlisting}
stage ID:
  new coherency = _ZX_2(coherency);
stage RR:
  new address = GPR[registerZ];
stage E1:
  new argument2 = GPR[registerT];
stage E3:
  MEM.atomic_min(address, 0xFF, coherency << 32, argument2.64[0], registerT);
\end{lstlisting}}
\newpage

\paragraph{Encoding} Format \textsf{LSU\_ASMINSB.O} (Section~\ref{sec:format-LSU-ASMINSB.O}), \verb|lsucode5 = 011|\\

\bitfields{ 1/P, 2/00, 2/-- --, 27/offset27, 1/1, 2/01, 3/111, 2/coherency, 6/registerT, 2/11, 3/011, 1/--, 5/11000, 1/0, 6/registerZ }


\paragraph{Syntax} \texttt{asmind}\emph{coherency} \emph{offset27}[\emph{registerZ}] = \emph{registerT}

\paragraph{Description} The effective address is computed by adding the \emph{offset27} to the \emph{registerZ}.
This instruction implements atomic MIN-store on signed double word. The double word at effective address is MINed with the \emph{registerT}. The effective address must be double word aligned (multiple of 8).


\paragraph{Modifier(s)}
\noindent\begin{itemize}
\item coherency (cf. coherency in section \ref{par:modifier-coherency} on page \pageref{par:modifier-coherency})
\end{itemize}

\paragraph{Execution} {Reservation table \textsf{LSU2\_MEMW\_AUXR.X} (Section~\ref{sec:reservation-LSU2-MEMW-AUXR.X})}
~
{\begin{lstlisting}
stage ID:
  new offset = _SX_27(offset27);
  new coherency = _ZX_2(coherency);
stage RR:
  new base = GPR[registerZ];
  new address = base + offset;
stage E1:
  new argument2 = GPR[registerT];
stage E3:
  MEM.atomic_min(address, 0xFF, coherency << 32, argument2.64[0], registerT);
\end{lstlisting}}
\newpage

\paragraph{Encoding} Format \textsf{LSU\_ASMINSB.D} (Section~\ref{sec:format-LSU-ASMINSB.D}), \verb|lsucode5 = 011|\\

\bitfields{ 1/P, 2/00, 2/-- --, 27/extend27, 1/1, 2/00, 2/-- --, 27/offset27, 1/1, 2/01, 3/111, 2/coherency, 6/registerT, 2/11, 3/011, 1/--, 5/11000, 1/0, 6/registerZ }


\paragraph{Syntax} \texttt{asmind}\emph{coherency} \emph{extend27\_\discretionary{}{}{}offset27}[\emph{registerZ}] = \emph{registerT}

\paragraph{Description} The effective address is computed by adding the \emph{extend27\_\discretionary{}{}{}offset27} to the \emph{registerZ}.
This instruction implements atomic MIN-store on signed double word. The double word at effective address is MINed with the \emph{registerT}. The effective address must be double word aligned (multiple of 8).


\paragraph{Modifier(s)}
\noindent\begin{itemize}
\item coherency (cf. coherency in section \ref{par:modifier-coherency} on page \pageref{par:modifier-coherency})
\end{itemize}

\paragraph{Execution} {Reservation table \textsf{LSU2\_MEMW\_AUXR.Y} (Section~\ref{sec:reservation-LSU2-MEMW-AUXR.Y})}
~
{\begin{lstlisting}
stage ID:
  new offset = _SX_54(extend27_offset27);
  new coherency = _ZX_2(coherency);
stage RR:
  new base = GPR[registerZ];
  new address = base + offset;
stage E1:
  new argument2 = GPR[registerT];
stage E3:
  MEM.atomic_min(address, 0xFF, coherency << 32, argument2.64[0], registerT);
\end{lstlisting}}
\newpage
\subsubsection[{\small ASMINH: Atomic Store MIN Half Word}]{ASMINH\hfill Atomic Store MIN Half Word}
\label{instr:ASMINH}\index{asminh}
 
\paragraph{Encoding} Format \textsf{LSU\_ASMINSB} (Section~\ref{sec:format-LSU-ASMINSB}), \verb|lsucode5 = 001|\\

\bitfields{ 1/P, 2/01, 3/111, 2/coherency, 6/registerT, 2/11, 3/001, 1/--, 5/11000, 1/0, 6/registerZ }


\paragraph{Syntax} \texttt{asminh}\emph{coherency} [\emph{registerZ}] = \emph{registerT}

\paragraph{Description} The effective address is the value of the \emph{registerZ}.
This instruction implements atomic MIN-store on signed half word. The half word at effective address is MINed with the \emph{registerT}. The effective address must be half word aligned (multiple of 2).


\paragraph{Modifier(s)}
\noindent\begin{itemize}
\item coherency (cf. coherency in section \ref{par:modifier-coherency} on page \pageref{par:modifier-coherency})
\end{itemize}

\paragraph{Execution} {Reservation table \textsf{LSU2\_MEMW\_AUXR} (Section~\ref{sec:reservation-LSU2-MEMW-AUXR})}
~
{\begin{lstlisting}
stage ID:
  new coherency = _ZX_2(coherency);
stage RR:
  new address = GPR[registerZ];
stage E1:
  new argument2 = GPR[registerT];
stage E3:
  MEM.atomic_min(address, 0x3, coherency << 32, argument2.16[0], registerT);
\end{lstlisting}}
\newpage

\paragraph{Encoding} Format \textsf{LSU\_ASMINSB.O} (Section~\ref{sec:format-LSU-ASMINSB.O}), \verb|lsucode5 = 001|\\

\bitfields{ 1/P, 2/00, 2/-- --, 27/offset27, 1/1, 2/01, 3/111, 2/coherency, 6/registerT, 2/11, 3/001, 1/--, 5/11000, 1/0, 6/registerZ }


\paragraph{Syntax} \texttt{asminh}\emph{coherency} \emph{offset27}[\emph{registerZ}] = \emph{registerT}

\paragraph{Description} The effective address is computed by adding the \emph{offset27} to the \emph{registerZ}.
This instruction implements atomic MIN-store on signed half word. The half word at effective address is MINed with the \emph{registerT}. The effective address must be half word aligned (multiple of 2).


\paragraph{Modifier(s)}
\noindent\begin{itemize}
\item coherency (cf. coherency in section \ref{par:modifier-coherency} on page \pageref{par:modifier-coherency})
\end{itemize}

\paragraph{Execution} {Reservation table \textsf{LSU2\_MEMW\_AUXR.X} (Section~\ref{sec:reservation-LSU2-MEMW-AUXR.X})}
~
{\begin{lstlisting}
stage ID:
  new offset = _SX_27(offset27);
  new coherency = _ZX_2(coherency);
stage RR:
  new base = GPR[registerZ];
  new address = base + offset;
stage E1:
  new argument2 = GPR[registerT];
stage E3:
  MEM.atomic_min(address, 0x3, coherency << 32, argument2.16[0], registerT);
\end{lstlisting}}
\newpage

\paragraph{Encoding} Format \textsf{LSU\_ASMINSB.D} (Section~\ref{sec:format-LSU-ASMINSB.D}), \verb|lsucode5 = 001|\\

\bitfields{ 1/P, 2/00, 2/-- --, 27/extend27, 1/1, 2/00, 2/-- --, 27/offset27, 1/1, 2/01, 3/111, 2/coherency, 6/registerT, 2/11, 3/001, 1/--, 5/11000, 1/0, 6/registerZ }


\paragraph{Syntax} \texttt{asminh}\emph{coherency} \emph{extend27\_\discretionary{}{}{}offset27}[\emph{registerZ}] = \emph{registerT}

\paragraph{Description} The effective address is computed by adding the \emph{extend27\_\discretionary{}{}{}offset27} to the \emph{registerZ}.
This instruction implements atomic MIN-store on signed half word. The half word at effective address is MINed with the \emph{registerT}. The effective address must be half word aligned (multiple of 2).


\paragraph{Modifier(s)}
\noindent\begin{itemize}
\item coherency (cf. coherency in section \ref{par:modifier-coherency} on page \pageref{par:modifier-coherency})
\end{itemize}

\paragraph{Execution} {Reservation table \textsf{LSU2\_MEMW\_AUXR.Y} (Section~\ref{sec:reservation-LSU2-MEMW-AUXR.Y})}
~
{\begin{lstlisting}
stage ID:
  new offset = _SX_54(extend27_offset27);
  new coherency = _ZX_2(coherency);
stage RR:
  new base = GPR[registerZ];
  new address = base + offset;
stage E1:
  new argument2 = GPR[registerT];
stage E3:
  MEM.atomic_min(address, 0x3, coherency << 32, argument2.16[0], registerT);
\end{lstlisting}}
\newpage
\subsubsection[{\small ASMINUB: Atomic Store MIN Unsigned Byte}]{ASMINUB\hfill Atomic Store MIN Unsigned Byte}
\label{instr:ASMINUB}\index{asminub}
 
\paragraph{Encoding} Format \textsf{LSU\_ASMINUSB} (Section~\ref{sec:format-LSU-ASMINUSB}), \verb|lsucode5 = 000|\\

\bitfields{ 1/P, 2/01, 3/111, 2/coherency, 6/registerT, 2/11, 3/000, 1/--, 5/11010, 1/0, 6/registerZ }


\paragraph{Syntax} \texttt{asminub}\emph{coherency} [\emph{registerZ}] = \emph{registerT}

\paragraph{Description} The effective address is the value of the \emph{registerZ}.
This instruction implements atomic MINU-store on signed byte. The byte at effective address is MINUed with the \emph{registerT}.


\paragraph{Modifier(s)}
\noindent\begin{itemize}
\item coherency (cf. coherency in section \ref{par:modifier-coherency} on page \pageref{par:modifier-coherency})
\end{itemize}

\paragraph{Execution} {Reservation table \textsf{LSU2\_MEMW\_AUXR} (Section~\ref{sec:reservation-LSU2-MEMW-AUXR})}
~
{\begin{lstlisting}
stage ID:
  new coherency = _ZX_2(coherency);
stage RR:
  new address = GPR[registerZ];
stage E1:
  new argument2 = GPR[registerT];
stage E3:
  MEM.atomic_minu(address, 0x1, coherency << 32, argument2.8[0], registerT);
\end{lstlisting}}
\newpage

\paragraph{Encoding} Format \textsf{LSU\_ASMINUSB.O} (Section~\ref{sec:format-LSU-ASMINUSB.O}), \verb|lsucode5 = 000|\\

\bitfields{ 1/P, 2/00, 2/-- --, 27/offset27, 1/1, 2/01, 3/111, 2/coherency, 6/registerT, 2/11, 3/000, 1/--, 5/11010, 1/0, 6/registerZ }


\paragraph{Syntax} \texttt{asminub}\emph{coherency} \emph{offset27}[\emph{registerZ}] = \emph{registerT}

\paragraph{Description} The effective address is computed by adding the \emph{offset27} to the \emph{registerZ}.
This instruction implements atomic MINU-store on signed byte. The byte at effective address is MINUed with the \emph{registerT}.


\paragraph{Modifier(s)}
\noindent\begin{itemize}
\item coherency (cf. coherency in section \ref{par:modifier-coherency} on page \pageref{par:modifier-coherency})
\end{itemize}

\paragraph{Execution} {Reservation table \textsf{LSU2\_MEMW\_AUXR.X} (Section~\ref{sec:reservation-LSU2-MEMW-AUXR.X})}
~
{\begin{lstlisting}
stage ID:
  new offset = _SX_27(offset27);
  new coherency = _ZX_2(coherency);
stage RR:
  new base = GPR[registerZ];
  new address = base + offset;
stage E1:
  new argument2 = GPR[registerT];
stage E3:
  MEM.atomic_minu(address, 0x1, coherency << 32, argument2.8[0], registerT);
\end{lstlisting}}
\newpage

\paragraph{Encoding} Format \textsf{LSU\_ASMINUSB.D} (Section~\ref{sec:format-LSU-ASMINUSB.D}), \verb|lsucode5 = 000|\\

\bitfields{ 1/P, 2/00, 2/-- --, 27/extend27, 1/1, 2/00, 2/-- --, 27/offset27, 1/1, 2/01, 3/111, 2/coherency, 6/registerT, 2/11, 3/000, 1/--, 5/11010, 1/0, 6/registerZ }


\paragraph{Syntax} \texttt{asminub}\emph{coherency} \emph{extend27\_\discretionary{}{}{}offset27}[\emph{registerZ}] = \emph{registerT}

\paragraph{Description} The effective address is computed by adding the \emph{extend27\_\discretionary{}{}{}offset27} to the \emph{registerZ}.
This instruction implements atomic MINU-store on signed byte. The byte at effective address is MINUed with the \emph{registerT}.


\paragraph{Modifier(s)}
\noindent\begin{itemize}
\item coherency (cf. coherency in section \ref{par:modifier-coherency} on page \pageref{par:modifier-coherency})
\end{itemize}

\paragraph{Execution} {Reservation table \textsf{LSU2\_MEMW\_AUXR.Y} (Section~\ref{sec:reservation-LSU2-MEMW-AUXR.Y})}
~
{\begin{lstlisting}
stage ID:
  new offset = _SX_54(extend27_offset27);
  new coherency = _ZX_2(coherency);
stage RR:
  new base = GPR[registerZ];
  new address = base + offset;
stage E1:
  new argument2 = GPR[registerT];
stage E3:
  MEM.atomic_minu(address, 0x1, coherency << 32, argument2.8[0], registerT);
\end{lstlisting}}
\newpage
\subsubsection[{\small ASMINUD: Atomic Store MIN Unsigned Double Word}]{ASMINUD\hfill Atomic Store MIN Unsigned Double Word}
\label{instr:ASMINUD}\index{asminud}
 
\paragraph{Encoding} Format \textsf{LSU\_ASMINUSB} (Section~\ref{sec:format-LSU-ASMINUSB}), \verb|lsucode5 = 011|\\

\bitfields{ 1/P, 2/01, 3/111, 2/coherency, 6/registerT, 2/11, 3/011, 1/--, 5/11010, 1/0, 6/registerZ }


\paragraph{Syntax} \texttt{asminud}\emph{coherency} [\emph{registerZ}] = \emph{registerT}

\paragraph{Description} The effective address is the value of the \emph{registerZ}.
This instruction implements atomic MINU-store on signed double word. The double word at effective address is MINUed with the \emph{registerT}. The effective address must be double word aligned (multiple of 8).


\paragraph{Modifier(s)}
\noindent\begin{itemize}
\item coherency (cf. coherency in section \ref{par:modifier-coherency} on page \pageref{par:modifier-coherency})
\end{itemize}

\paragraph{Execution} {Reservation table \textsf{LSU2\_MEMW\_AUXR} (Section~\ref{sec:reservation-LSU2-MEMW-AUXR})}
~
{\begin{lstlisting}
stage ID:
  new coherency = _ZX_2(coherency);
stage RR:
  new address = GPR[registerZ];
stage E1:
  new argument2 = GPR[registerT];
stage E3:
  MEM.atomic_minu(address, 0xFF, coherency << 32, argument2.64[0], registerT);
\end{lstlisting}}
\newpage

\paragraph{Encoding} Format \textsf{LSU\_ASMINUSB.O} (Section~\ref{sec:format-LSU-ASMINUSB.O}), \verb|lsucode5 = 011|\\

\bitfields{ 1/P, 2/00, 2/-- --, 27/offset27, 1/1, 2/01, 3/111, 2/coherency, 6/registerT, 2/11, 3/011, 1/--, 5/11010, 1/0, 6/registerZ }


\paragraph{Syntax} \texttt{asminud}\emph{coherency} \emph{offset27}[\emph{registerZ}] = \emph{registerT}

\paragraph{Description} The effective address is computed by adding the \emph{offset27} to the \emph{registerZ}.
This instruction implements atomic MINU-store on signed double word. The double word at effective address is MINUed with the \emph{registerT}. The effective address must be double word aligned (multiple of 8).


\paragraph{Modifier(s)}
\noindent\begin{itemize}
\item coherency (cf. coherency in section \ref{par:modifier-coherency} on page \pageref{par:modifier-coherency})
\end{itemize}

\paragraph{Execution} {Reservation table \textsf{LSU2\_MEMW\_AUXR.X} (Section~\ref{sec:reservation-LSU2-MEMW-AUXR.X})}
~
{\begin{lstlisting}
stage ID:
  new offset = _SX_27(offset27);
  new coherency = _ZX_2(coherency);
stage RR:
  new base = GPR[registerZ];
  new address = base + offset;
stage E1:
  new argument2 = GPR[registerT];
stage E3:
  MEM.atomic_minu(address, 0xFF, coherency << 32, argument2.64[0], registerT);
\end{lstlisting}}
\newpage

\paragraph{Encoding} Format \textsf{LSU\_ASMINUSB.D} (Section~\ref{sec:format-LSU-ASMINUSB.D}), \verb|lsucode5 = 011|\\

\bitfields{ 1/P, 2/00, 2/-- --, 27/extend27, 1/1, 2/00, 2/-- --, 27/offset27, 1/1, 2/01, 3/111, 2/coherency, 6/registerT, 2/11, 3/011, 1/--, 5/11010, 1/0, 6/registerZ }


\paragraph{Syntax} \texttt{asminud}\emph{coherency} \emph{extend27\_\discretionary{}{}{}offset27}[\emph{registerZ}] = \emph{registerT}

\paragraph{Description} The effective address is computed by adding the \emph{extend27\_\discretionary{}{}{}offset27} to the \emph{registerZ}.
This instruction implements atomic MINU-store on signed double word. The double word at effective address is MINUed with the \emph{registerT}. The effective address must be double word aligned (multiple of 8).


\paragraph{Modifier(s)}
\noindent\begin{itemize}
\item coherency (cf. coherency in section \ref{par:modifier-coherency} on page \pageref{par:modifier-coherency})
\end{itemize}

\paragraph{Execution} {Reservation table \textsf{LSU2\_MEMW\_AUXR.Y} (Section~\ref{sec:reservation-LSU2-MEMW-AUXR.Y})}
~
{\begin{lstlisting}
stage ID:
  new offset = _SX_54(extend27_offset27);
  new coherency = _ZX_2(coherency);
stage RR:
  new base = GPR[registerZ];
  new address = base + offset;
stage E1:
  new argument2 = GPR[registerT];
stage E3:
  MEM.atomic_minu(address, 0xFF, coherency << 32, argument2.64[0], registerT);
\end{lstlisting}}
\newpage
\subsubsection[{\small ASMINUH: Atomic Store MIN Unsigned Half Word}]{ASMINUH\hfill Atomic Store MIN Unsigned Half Word}
\label{instr:ASMINUH}\index{asminuh}
 
\paragraph{Encoding} Format \textsf{LSU\_ASMINUSB} (Section~\ref{sec:format-LSU-ASMINUSB}), \verb|lsucode5 = 001|\\

\bitfields{ 1/P, 2/01, 3/111, 2/coherency, 6/registerT, 2/11, 3/001, 1/--, 5/11010, 1/0, 6/registerZ }


\paragraph{Syntax} \texttt{asminuh}\emph{coherency} [\emph{registerZ}] = \emph{registerT}

\paragraph{Description} The effective address is the value of the \emph{registerZ}.
This instruction implements atomic MINU-store on signed half word. The half word at effective address is MINUed with the \emph{registerT}. The effective address must be half word aligned (multiple of 2).


\paragraph{Modifier(s)}
\noindent\begin{itemize}
\item coherency (cf. coherency in section \ref{par:modifier-coherency} on page \pageref{par:modifier-coherency})
\end{itemize}

\paragraph{Execution} {Reservation table \textsf{LSU2\_MEMW\_AUXR} (Section~\ref{sec:reservation-LSU2-MEMW-AUXR})}
~
{\begin{lstlisting}
stage ID:
  new coherency = _ZX_2(coherency);
stage RR:
  new address = GPR[registerZ];
stage E1:
  new argument2 = GPR[registerT];
stage E3:
  MEM.atomic_minu(address, 0x3, coherency << 32, argument2.16[0], registerT);
\end{lstlisting}}
\newpage

\paragraph{Encoding} Format \textsf{LSU\_ASMINUSB.O} (Section~\ref{sec:format-LSU-ASMINUSB.O}), \verb|lsucode5 = 001|\\

\bitfields{ 1/P, 2/00, 2/-- --, 27/offset27, 1/1, 2/01, 3/111, 2/coherency, 6/registerT, 2/11, 3/001, 1/--, 5/11010, 1/0, 6/registerZ }


\paragraph{Syntax} \texttt{asminuh}\emph{coherency} \emph{offset27}[\emph{registerZ}] = \emph{registerT}

\paragraph{Description} The effective address is computed by adding the \emph{offset27} to the \emph{registerZ}.
This instruction implements atomic MINU-store on signed half word. The half word at effective address is MINUed with the \emph{registerT}. The effective address must be half word aligned (multiple of 2).


\paragraph{Modifier(s)}
\noindent\begin{itemize}
\item coherency (cf. coherency in section \ref{par:modifier-coherency} on page \pageref{par:modifier-coherency})
\end{itemize}

\paragraph{Execution} {Reservation table \textsf{LSU2\_MEMW\_AUXR.X} (Section~\ref{sec:reservation-LSU2-MEMW-AUXR.X})}
~
{\begin{lstlisting}
stage ID:
  new offset = _SX_27(offset27);
  new coherency = _ZX_2(coherency);
stage RR:
  new base = GPR[registerZ];
  new address = base + offset;
stage E1:
  new argument2 = GPR[registerT];
stage E3:
  MEM.atomic_minu(address, 0x3, coherency << 32, argument2.16[0], registerT);
\end{lstlisting}}
\newpage

\paragraph{Encoding} Format \textsf{LSU\_ASMINUSB.D} (Section~\ref{sec:format-LSU-ASMINUSB.D}), \verb|lsucode5 = 001|\\

\bitfields{ 1/P, 2/00, 2/-- --, 27/extend27, 1/1, 2/00, 2/-- --, 27/offset27, 1/1, 2/01, 3/111, 2/coherency, 6/registerT, 2/11, 3/001, 1/--, 5/11010, 1/0, 6/registerZ }


\paragraph{Syntax} \texttt{asminuh}\emph{coherency} \emph{extend27\_\discretionary{}{}{}offset27}[\emph{registerZ}] = \emph{registerT}

\paragraph{Description} The effective address is computed by adding the \emph{extend27\_\discretionary{}{}{}offset27} to the \emph{registerZ}.
This instruction implements atomic MINU-store on signed half word. The half word at effective address is MINUed with the \emph{registerT}. The effective address must be half word aligned (multiple of 2).


\paragraph{Modifier(s)}
\noindent\begin{itemize}
\item coherency (cf. coherency in section \ref{par:modifier-coherency} on page \pageref{par:modifier-coherency})
\end{itemize}

\paragraph{Execution} {Reservation table \textsf{LSU2\_MEMW\_AUXR.Y} (Section~\ref{sec:reservation-LSU2-MEMW-AUXR.Y})}
~
{\begin{lstlisting}
stage ID:
  new offset = _SX_54(extend27_offset27);
  new coherency = _ZX_2(coherency);
stage RR:
  new base = GPR[registerZ];
  new address = base + offset;
stage E1:
  new argument2 = GPR[registerT];
stage E3:
  MEM.atomic_minu(address, 0x3, coherency << 32, argument2.16[0], registerT);
\end{lstlisting}}
\newpage
\subsubsection[{\small ASMINUW: Atomic Store MIN Unsigned Word}]{ASMINUW\hfill Atomic Store MIN Unsigned Word}
\label{instr:ASMINUW}\index{asminuw}
 
\paragraph{Encoding} Format \textsf{LSU\_ASMINUSB} (Section~\ref{sec:format-LSU-ASMINUSB}), \verb|lsucode5 = 010|\\

\bitfields{ 1/P, 2/01, 3/111, 2/coherency, 6/registerT, 2/11, 3/010, 1/--, 5/11010, 1/0, 6/registerZ }


\paragraph{Syntax} \texttt{asminuw}\emph{coherency} [\emph{registerZ}] = \emph{registerT}

\paragraph{Description} The effective address is the value of the \emph{registerZ}.
This instruction implements atomic MINU-store on signed word. The word at effective address is MINUed with the \emph{registerT}. The effective address must be word aligned (multiple of 4).


\paragraph{Modifier(s)}
\noindent\begin{itemize}
\item coherency (cf. coherency in section \ref{par:modifier-coherency} on page \pageref{par:modifier-coherency})
\end{itemize}

\paragraph{Execution} {Reservation table \textsf{LSU2\_MEMW\_AUXR} (Section~\ref{sec:reservation-LSU2-MEMW-AUXR})}
~
{\begin{lstlisting}
stage ID:
  new coherency = _ZX_2(coherency);
stage RR:
  new address = GPR[registerZ];
stage E1:
  new argument2 = GPR[registerT];
stage E3:
  MEM.atomic_minu(address, 0xF, coherency << 32, argument2.32[0], registerT);
\end{lstlisting}}
\newpage

\paragraph{Encoding} Format \textsf{LSU\_ASMINUSB.O} (Section~\ref{sec:format-LSU-ASMINUSB.O}), \verb|lsucode5 = 010|\\

\bitfields{ 1/P, 2/00, 2/-- --, 27/offset27, 1/1, 2/01, 3/111, 2/coherency, 6/registerT, 2/11, 3/010, 1/--, 5/11010, 1/0, 6/registerZ }


\paragraph{Syntax} \texttt{asminuw}\emph{coherency} \emph{offset27}[\emph{registerZ}] = \emph{registerT}

\paragraph{Description} The effective address is computed by adding the \emph{offset27} to the \emph{registerZ}.
This instruction implements atomic MINU-store on signed word. The word at effective address is MINUed with the \emph{registerT}. The effective address must be word aligned (multiple of 4).


\paragraph{Modifier(s)}
\noindent\begin{itemize}
\item coherency (cf. coherency in section \ref{par:modifier-coherency} on page \pageref{par:modifier-coherency})
\end{itemize}

\paragraph{Execution} {Reservation table \textsf{LSU2\_MEMW\_AUXR.X} (Section~\ref{sec:reservation-LSU2-MEMW-AUXR.X})}
~
{\begin{lstlisting}
stage ID:
  new offset = _SX_27(offset27);
  new coherency = _ZX_2(coherency);
stage RR:
  new base = GPR[registerZ];
  new address = base + offset;
stage E1:
  new argument2 = GPR[registerT];
stage E3:
  MEM.atomic_minu(address, 0xF, coherency << 32, argument2.32[0], registerT);
\end{lstlisting}}
\newpage

\paragraph{Encoding} Format \textsf{LSU\_ASMINUSB.D} (Section~\ref{sec:format-LSU-ASMINUSB.D}), \verb|lsucode5 = 010|\\

\bitfields{ 1/P, 2/00, 2/-- --, 27/extend27, 1/1, 2/00, 2/-- --, 27/offset27, 1/1, 2/01, 3/111, 2/coherency, 6/registerT, 2/11, 3/010, 1/--, 5/11010, 1/0, 6/registerZ }


\paragraph{Syntax} \texttt{asminuw}\emph{coherency} \emph{extend27\_\discretionary{}{}{}offset27}[\emph{registerZ}] = \emph{registerT}

\paragraph{Description} The effective address is computed by adding the \emph{extend27\_\discretionary{}{}{}offset27} to the \emph{registerZ}.
This instruction implements atomic MINU-store on signed word. The word at effective address is MINUed with the \emph{registerT}. The effective address must be word aligned (multiple of 4).


\paragraph{Modifier(s)}
\noindent\begin{itemize}
\item coherency (cf. coherency in section \ref{par:modifier-coherency} on page \pageref{par:modifier-coherency})
\end{itemize}

\paragraph{Execution} {Reservation table \textsf{LSU2\_MEMW\_AUXR.Y} (Section~\ref{sec:reservation-LSU2-MEMW-AUXR.Y})}
~
{\begin{lstlisting}
stage ID:
  new offset = _SX_54(extend27_offset27);
  new coherency = _ZX_2(coherency);
stage RR:
  new base = GPR[registerZ];
  new address = base + offset;
stage E1:
  new argument2 = GPR[registerT];
stage E3:
  MEM.atomic_minu(address, 0xF, coherency << 32, argument2.32[0], registerT);
\end{lstlisting}}
\newpage
\subsubsection[{\small ASMINW: Atomic Store MIN Word}]{ASMINW\hfill Atomic Store MIN Word}
\label{instr:ASMINW}\index{asminw}
 
\paragraph{Encoding} Format \textsf{LSU\_ASMINSB} (Section~\ref{sec:format-LSU-ASMINSB}), \verb|lsucode5 = 010|\\

\bitfields{ 1/P, 2/01, 3/111, 2/coherency, 6/registerT, 2/11, 3/010, 1/--, 5/11000, 1/0, 6/registerZ }


\paragraph{Syntax} \texttt{asminw}\emph{coherency} [\emph{registerZ}] = \emph{registerT}

\paragraph{Description} The effective address is the value of the \emph{registerZ}.
This instruction implements atomic MIN-store on signed word. The word at effective address is MINed with the \emph{registerT}. The effective address must be word aligned (multiple of 4).


\paragraph{Modifier(s)}
\noindent\begin{itemize}
\item coherency (cf. coherency in section \ref{par:modifier-coherency} on page \pageref{par:modifier-coherency})
\end{itemize}

\paragraph{Execution} {Reservation table \textsf{LSU2\_MEMW\_AUXR} (Section~\ref{sec:reservation-LSU2-MEMW-AUXR})}
~
{\begin{lstlisting}
stage ID:
  new coherency = _ZX_2(coherency);
stage RR:
  new address = GPR[registerZ];
stage E1:
  new argument2 = GPR[registerT];
stage E3:
  MEM.atomic_min(address, 0xF, coherency << 32, argument2.32[0], registerT);
\end{lstlisting}}
\newpage

\paragraph{Encoding} Format \textsf{LSU\_ASMINSB.O} (Section~\ref{sec:format-LSU-ASMINSB.O}), \verb|lsucode5 = 010|\\

\bitfields{ 1/P, 2/00, 2/-- --, 27/offset27, 1/1, 2/01, 3/111, 2/coherency, 6/registerT, 2/11, 3/010, 1/--, 5/11000, 1/0, 6/registerZ }


\paragraph{Syntax} \texttt{asminw}\emph{coherency} \emph{offset27}[\emph{registerZ}] = \emph{registerT}

\paragraph{Description} The effective address is computed by adding the \emph{offset27} to the \emph{registerZ}.
This instruction implements atomic MIN-store on signed word. The word at effective address is MINed with the \emph{registerT}. The effective address must be word aligned (multiple of 4).


\paragraph{Modifier(s)}
\noindent\begin{itemize}
\item coherency (cf. coherency in section \ref{par:modifier-coherency} on page \pageref{par:modifier-coherency})
\end{itemize}

\paragraph{Execution} {Reservation table \textsf{LSU2\_MEMW\_AUXR.X} (Section~\ref{sec:reservation-LSU2-MEMW-AUXR.X})}
~
{\begin{lstlisting}
stage ID:
  new offset = _SX_27(offset27);
  new coherency = _ZX_2(coherency);
stage RR:
  new base = GPR[registerZ];
  new address = base + offset;
stage E1:
  new argument2 = GPR[registerT];
stage E3:
  MEM.atomic_min(address, 0xF, coherency << 32, argument2.32[0], registerT);
\end{lstlisting}}
\newpage

\paragraph{Encoding} Format \textsf{LSU\_ASMINSB.D} (Section~\ref{sec:format-LSU-ASMINSB.D}), \verb|lsucode5 = 010|\\

\bitfields{ 1/P, 2/00, 2/-- --, 27/extend27, 1/1, 2/00, 2/-- --, 27/offset27, 1/1, 2/01, 3/111, 2/coherency, 6/registerT, 2/11, 3/010, 1/--, 5/11000, 1/0, 6/registerZ }


\paragraph{Syntax} \texttt{asminw}\emph{coherency} \emph{extend27\_\discretionary{}{}{}offset27}[\emph{registerZ}] = \emph{registerT}

\paragraph{Description} The effective address is computed by adding the \emph{extend27\_\discretionary{}{}{}offset27} to the \emph{registerZ}.
This instruction implements atomic MIN-store on signed word. The word at effective address is MINed with the \emph{registerT}. The effective address must be word aligned (multiple of 4).


\paragraph{Modifier(s)}
\noindent\begin{itemize}
\item coherency (cf. coherency in section \ref{par:modifier-coherency} on page \pageref{par:modifier-coherency})
\end{itemize}

\paragraph{Execution} {Reservation table \textsf{LSU2\_MEMW\_AUXR.Y} (Section~\ref{sec:reservation-LSU2-MEMW-AUXR.Y})}
~
{\begin{lstlisting}
stage ID:
  new offset = _SX_54(extend27_offset27);
  new coherency = _ZX_2(coherency);
stage RR:
  new base = GPR[registerZ];
  new address = base + offset;
stage E1:
  new argument2 = GPR[registerT];
stage E3:
  MEM.atomic_min(address, 0xF, coherency << 32, argument2.32[0], registerT);
\end{lstlisting}}
\newpage
\subsubsection[{\small ASW: Atomic Store Word}]{ASW\hfill Atomic Store Word}
\label{instr:ASW}\index{asw}
 
\paragraph{Encoding} Format \textsf{LSU\_ASSB} (Section~\ref{sec:format-LSU-ASSB}), \verb|lsucode5 = 010|\\

\bitfields{ 1/P, 2/01, 3/111, 2/coherency, 6/registerT, 2/11, 3/010, 1/--, 5/10000, 1/0, 6/registerZ }


\paragraph{Syntax} \texttt{asw}\emph{coherency} [\emph{registerZ}] = \emph{registerT}

\paragraph{Description} The effective address is the value of the \emph{registerZ}.
The \emph{registerT} is atomically stored to the memory word starting at the effective address. This instruction is provided to directly operate on the last level of the memory hierarchy, which is typically a processor DDR or a cluster TCM, irrespective of the MMU DCP (Data Cache Policy) settings. The effective address must be word aligned (multiple of 4).


\paragraph{Modifier(s)}
\noindent\begin{itemize}
\item coherency (cf. coherency in section \ref{par:modifier-coherency} on page \pageref{par:modifier-coherency})
\end{itemize}

\paragraph{Execution} {Reservation table \textsf{LSU\_MEMW\_AUXR} (Section~\ref{sec:reservation-LSU-MEMW-AUXR})}
~
{\begin{lstlisting}
stage ID:
  new coherency = _ZX_2(coherency);
stage RR:
  new address = GPR[registerZ];
stage E1:
  new argument2 = GPR[registerT];
stage E3:
  MEM.atomic_store(address, 0xF, coherency << 32, argument2, registerT);
\end{lstlisting}}
\newpage

\paragraph{Encoding} Format \textsf{LSU\_ASSB.O} (Section~\ref{sec:format-LSU-ASSB.O}), \verb|lsucode5 = 010|\\

\bitfields{ 1/P, 2/00, 2/-- --, 27/offset27, 1/1, 2/01, 3/111, 2/coherency, 6/registerT, 2/11, 3/010, 1/--, 5/10000, 1/0, 6/registerZ }


\paragraph{Syntax} \texttt{asw}\emph{coherency} \emph{offset27}[\emph{registerZ}] = \emph{registerT}

\paragraph{Description} The effective address is computed by adding the \emph{offset27} to the \emph{registerZ}.
The \emph{registerT} is atomically stored to the memory word starting at the effective address. This instruction is provided to directly operate on the last level of the memory hierarchy, which is typically a processor DDR or a cluster TCM, irrespective of the MMU DCP (Data Cache Policy) settings. The effective address must be word aligned (multiple of 4).


\paragraph{Modifier(s)}
\noindent\begin{itemize}
\item coherency (cf. coherency in section \ref{par:modifier-coherency} on page \pageref{par:modifier-coherency})
\end{itemize}

\paragraph{Execution} {Reservation table \textsf{LSU\_MEMW\_AUXR.X} (Section~\ref{sec:reservation-LSU-MEMW-AUXR.X})}
~
{\begin{lstlisting}
stage ID:
  new offset = _SX_27(offset27);
  new coherency = _ZX_2(coherency);
stage RR:
  new base = GPR[registerZ];
  new address = base + offset;
stage E1:
  new argument2 = GPR[registerT];
stage E3:
  MEM.atomic_store(address, 0xF, coherency << 32, argument2, registerT);
\end{lstlisting}}
\newpage

\paragraph{Encoding} Format \textsf{LSU\_ASSB.D} (Section~\ref{sec:format-LSU-ASSB.D}), \verb|lsucode5 = 010|\\

\bitfields{ 1/P, 2/00, 2/-- --, 27/extend27, 1/1, 2/00, 2/-- --, 27/offset27, 1/1, 2/01, 3/111, 2/coherency, 6/registerT, 2/11, 3/010, 1/--, 5/10000, 1/0, 6/registerZ }


\paragraph{Syntax} \texttt{asw}\emph{coherency} \emph{extend27\_\discretionary{}{}{}offset27}[\emph{registerZ}] = \emph{registerT}

\paragraph{Description} The effective address is computed by adding the \emph{extend27\_\discretionary{}{}{}offset27} to the \emph{registerZ}.
The \emph{registerT} is atomically stored to the memory word starting at the effective address. This instruction is provided to directly operate on the last level of the memory hierarchy, which is typically a processor DDR or a cluster TCM, irrespective of the MMU DCP (Data Cache Policy) settings. The effective address must be word aligned (multiple of 4).


\paragraph{Modifier(s)}
\noindent\begin{itemize}
\item coherency (cf. coherency in section \ref{par:modifier-coherency} on page \pageref{par:modifier-coherency})
\end{itemize}

\paragraph{Execution} {Reservation table \textsf{LSU\_MEMW\_AUXR.Y} (Section~\ref{sec:reservation-LSU-MEMW-AUXR.Y})}
~
{\begin{lstlisting}
stage ID:
  new offset = _SX_54(extend27_offset27);
  new coherency = _ZX_2(coherency);
stage RR:
  new base = GPR[registerZ];
  new address = base + offset;
stage E1:
  new argument2 = GPR[registerT];
stage E3:
  MEM.atomic_store(address, 0xF, coherency << 32, argument2, registerT);
\end{lstlisting}}
\newpage
\subsubsection[{\small ASWAPB: Atomic Swap Byte}]{ASWAPB\hfill Atomic Swap Byte}
\label{instr:ASWAPB}\index{aswapb}
 
\paragraph{Encoding} Format \textsf{LSU\_ASWAPSB} (Section~\ref{sec:format-LSU-ASWAPSB}), \verb|lsucode5 = 000|\\

\bitfields{ 1/P, 2/01, 3/111, 2/coherency, 6/registerT, 2/11, 3/000, 1/--, 5/00010, 1/0, 6/registerZ }


\paragraph{Syntax} \texttt{aswapb}\emph{coherency} [\emph{registerZ}] = \emph{registerT}

\paragraph{Description} The effective address is the value of the \emph{registerZ}.
This instruction implements an atomic memory swap of byte. The byte at effective address is written with the \emph{registerT} value, and its previous value is returned into the \emph{registerT}.


\paragraph{Modifier(s)}
\noindent\begin{itemize}
\item coherency (cf. coherency in section \ref{par:modifier-coherency} on page \pageref{par:modifier-coherency})
\end{itemize}

\paragraph{Execution} {Reservation table \textsf{LSU2\_MEMW\_AUXR\_AUXW} (Section~\ref{sec:reservation-LSU2-MEMW-AUXR-AUXW})}
~
{\begin{lstlisting}
stage ID:
  new coherency = _ZX_2(coherency);
stage RR:
  new address = GPR[registerZ];
stage E1:
  new argument2 = GPR[registerT];
  new result2 = _ZX_8(MEM.atomic_swap(address, 0x1, coherency << 32, argument2, registerT));
stage SR:
  GPR[registerT] = result2;
\end{lstlisting}}
\newpage

\paragraph{Encoding} Format \textsf{LSU\_ASWAPSB.O} (Section~\ref{sec:format-LSU-ASWAPSB.O}), \verb|lsucode5 = 000|\\

\bitfields{ 1/P, 2/00, 2/-- --, 27/offset27, 1/1, 2/01, 3/111, 2/coherency, 6/registerT, 2/11, 3/000, 1/--, 5/00010, 1/0, 6/registerZ }


\paragraph{Syntax} \texttt{aswapb}\emph{coherency} \emph{offset27}[\emph{registerZ}] = \emph{registerT}

\paragraph{Description} The effective address is computed by adding the \emph{offset27} to the \emph{registerZ}.
This instruction implements an atomic memory swap of byte. The byte at effective address is written with the \emph{registerT} value, and its previous value is returned into the \emph{registerT}.


\paragraph{Modifier(s)}
\noindent\begin{itemize}
\item coherency (cf. coherency in section \ref{par:modifier-coherency} on page \pageref{par:modifier-coherency})
\end{itemize}

\paragraph{Execution} {Reservation table \textsf{LSU2\_MEMW\_AUXR\_AUXW.X} (Section~\ref{sec:reservation-LSU2-MEMW-AUXR-AUXW.X})}
~
{\begin{lstlisting}
stage ID:
  new offset = _SX_27(offset27);
  new coherency = _ZX_2(coherency);
stage RR:
  new base = GPR[registerZ];
  new address = base + offset;
stage E1:
  new argument2 = GPR[registerT];
  new result2 = _ZX_8(MEM.atomic_swap(address, 0x1, coherency << 32, argument2, registerT));
stage SR:
  GPR[registerT] = result2;
\end{lstlisting}}
\newpage

\paragraph{Encoding} Format \textsf{LSU\_ASWAPSB.D} (Section~\ref{sec:format-LSU-ASWAPSB.D}), \verb|lsucode5 = 000|\\

\bitfields{ 1/P, 2/00, 2/-- --, 27/extend27, 1/1, 2/00, 2/-- --, 27/offset27, 1/1, 2/01, 3/111, 2/coherency, 6/registerT, 2/11, 3/000, 1/--, 5/00010, 1/0, 6/registerZ }


\paragraph{Syntax} \texttt{aswapb}\emph{coherency} \emph{extend27\_\discretionary{}{}{}offset27}[\emph{registerZ}] = \emph{registerT}

\paragraph{Description} The effective address is computed by adding the \emph{extend27\_\discretionary{}{}{}offset27} to the \emph{registerZ}.
This instruction implements an atomic memory swap of byte. The byte at effective address is written with the \emph{registerT} value, and its previous value is returned into the \emph{registerT}.


\paragraph{Modifier(s)}
\noindent\begin{itemize}
\item coherency (cf. coherency in section \ref{par:modifier-coherency} on page \pageref{par:modifier-coherency})
\end{itemize}

\paragraph{Execution} {Reservation table \textsf{LSU2\_MEMW\_AUXR\_AUXW.Y} (Section~\ref{sec:reservation-LSU2-MEMW-AUXR-AUXW.Y})}
~
{\begin{lstlisting}
stage ID:
  new offset = _SX_54(extend27_offset27);
  new coherency = _ZX_2(coherency);
stage RR:
  new base = GPR[registerZ];
  new address = base + offset;
stage E1:
  new argument2 = GPR[registerT];
  new result2 = _ZX_8(MEM.atomic_swap(address, 0x1, coherency << 32, argument2, registerT));
stage SR:
  GPR[registerT] = result2;
\end{lstlisting}}
\newpage
\subsubsection[{\small ASWAPD: Atomic Swap Double Word}]{ASWAPD\hfill Atomic Swap Double Word}
\label{instr:ASWAPD}\index{aswapd}
 
\paragraph{Encoding} Format \textsf{LSU\_ASWAPSB} (Section~\ref{sec:format-LSU-ASWAPSB}), \verb|lsucode5 = 011|\\

\bitfields{ 1/P, 2/01, 3/111, 2/coherency, 6/registerT, 2/11, 3/011, 1/--, 5/00010, 1/0, 6/registerZ }


\paragraph{Syntax} \texttt{aswapd}\emph{coherency} [\emph{registerZ}] = \emph{registerT}

\paragraph{Description} The effective address is the value of the \emph{registerZ}.
This instruction implements an atomic memory swap of double word. The double word at effective address is written with the \emph{registerT} value, and its previous value is returned into the \emph{registerT}.


\paragraph{Modifier(s)}
\noindent\begin{itemize}
\item coherency (cf. coherency in section \ref{par:modifier-coherency} on page \pageref{par:modifier-coherency})
\end{itemize}

\paragraph{Execution} {Reservation table \textsf{LSU2\_MEMW\_AUXR\_AUXW} (Section~\ref{sec:reservation-LSU2-MEMW-AUXR-AUXW})}
~
{\begin{lstlisting}
stage ID:
  new coherency = _ZX_2(coherency);
stage RR:
  new address = GPR[registerZ];
stage E1:
  new argument2 = GPR[registerT];
  new result2 = MEM.atomic_swap(address, 0xFF, coherency << 32, argument2, registerT);
stage SR:
  GPR[registerT] = result2;
\end{lstlisting}}
\newpage

\paragraph{Encoding} Format \textsf{LSU\_ASWAPSB.O} (Section~\ref{sec:format-LSU-ASWAPSB.O}), \verb|lsucode5 = 011|\\

\bitfields{ 1/P, 2/00, 2/-- --, 27/offset27, 1/1, 2/01, 3/111, 2/coherency, 6/registerT, 2/11, 3/011, 1/--, 5/00010, 1/0, 6/registerZ }


\paragraph{Syntax} \texttt{aswapd}\emph{coherency} \emph{offset27}[\emph{registerZ}] = \emph{registerT}

\paragraph{Description} The effective address is computed by adding the \emph{offset27} to the \emph{registerZ}.
This instruction implements an atomic memory swap of double word. The double word at effective address is written with the \emph{registerT} value, and its previous value is returned into the \emph{registerT}.


\paragraph{Modifier(s)}
\noindent\begin{itemize}
\item coherency (cf. coherency in section \ref{par:modifier-coherency} on page \pageref{par:modifier-coherency})
\end{itemize}

\paragraph{Execution} {Reservation table \textsf{LSU2\_MEMW\_AUXR\_AUXW.X} (Section~\ref{sec:reservation-LSU2-MEMW-AUXR-AUXW.X})}
~
{\begin{lstlisting}
stage ID:
  new offset = _SX_27(offset27);
  new coherency = _ZX_2(coherency);
stage RR:
  new base = GPR[registerZ];
  new address = base + offset;
stage E1:
  new argument2 = GPR[registerT];
  new result2 = MEM.atomic_swap(address, 0xFF, coherency << 32, argument2, registerT);
stage SR:
  GPR[registerT] = result2;
\end{lstlisting}}
\newpage

\paragraph{Encoding} Format \textsf{LSU\_ASWAPSB.D} (Section~\ref{sec:format-LSU-ASWAPSB.D}), \verb|lsucode5 = 011|\\

\bitfields{ 1/P, 2/00, 2/-- --, 27/extend27, 1/1, 2/00, 2/-- --, 27/offset27, 1/1, 2/01, 3/111, 2/coherency, 6/registerT, 2/11, 3/011, 1/--, 5/00010, 1/0, 6/registerZ }


\paragraph{Syntax} \texttt{aswapd}\emph{coherency} \emph{extend27\_\discretionary{}{}{}offset27}[\emph{registerZ}] = \emph{registerT}

\paragraph{Description} The effective address is computed by adding the \emph{extend27\_\discretionary{}{}{}offset27} to the \emph{registerZ}.
This instruction implements an atomic memory swap of double word. The double word at effective address is written with the \emph{registerT} value, and its previous value is returned into the \emph{registerT}.


\paragraph{Modifier(s)}
\noindent\begin{itemize}
\item coherency (cf. coherency in section \ref{par:modifier-coherency} on page \pageref{par:modifier-coherency})
\end{itemize}

\paragraph{Execution} {Reservation table \textsf{LSU2\_MEMW\_AUXR\_AUXW.Y} (Section~\ref{sec:reservation-LSU2-MEMW-AUXR-AUXW.Y})}
~
{\begin{lstlisting}
stage ID:
  new offset = _SX_54(extend27_offset27);
  new coherency = _ZX_2(coherency);
stage RR:
  new base = GPR[registerZ];
  new address = base + offset;
stage E1:
  new argument2 = GPR[registerT];
  new result2 = MEM.atomic_swap(address, 0xFF, coherency << 32, argument2, registerT);
stage SR:
  GPR[registerT] = result2;
\end{lstlisting}}
\newpage
\subsubsection[{\small ASWAPH: Atomic Swap Half Word}]{ASWAPH\hfill Atomic Swap Half Word}
\label{instr:ASWAPH}\index{aswaph}
 
\paragraph{Encoding} Format \textsf{LSU\_ASWAPSB} (Section~\ref{sec:format-LSU-ASWAPSB}), \verb|lsucode5 = 001|\\

\bitfields{ 1/P, 2/01, 3/111, 2/coherency, 6/registerT, 2/11, 3/001, 1/--, 5/00010, 1/0, 6/registerZ }


\paragraph{Syntax} \texttt{aswaph}\emph{coherency} [\emph{registerZ}] = \emph{registerT}

\paragraph{Description} The effective address is the value of the \emph{registerZ}.
This instruction implements an atomic memory swap of half word. The half word at effective address is written with the \emph{registerT} value, and its previous value is returned into the \emph{registerT}.


\paragraph{Modifier(s)}
\noindent\begin{itemize}
\item coherency (cf. coherency in section \ref{par:modifier-coherency} on page \pageref{par:modifier-coherency})
\end{itemize}

\paragraph{Execution} {Reservation table \textsf{LSU2\_MEMW\_AUXR\_AUXW} (Section~\ref{sec:reservation-LSU2-MEMW-AUXR-AUXW})}
~
{\begin{lstlisting}
stage ID:
  new coherency = _ZX_2(coherency);
stage RR:
  new address = GPR[registerZ];
stage E1:
  new argument2 = GPR[registerT];
  new result2 = _ZX_16(MEM.atomic_swap(address, 0x3, coherency << 32, argument2, registerT));
stage SR:
  GPR[registerT] = result2;
\end{lstlisting}}
\newpage

\paragraph{Encoding} Format \textsf{LSU\_ASWAPSB.O} (Section~\ref{sec:format-LSU-ASWAPSB.O}), \verb|lsucode5 = 001|\\

\bitfields{ 1/P, 2/00, 2/-- --, 27/offset27, 1/1, 2/01, 3/111, 2/coherency, 6/registerT, 2/11, 3/001, 1/--, 5/00010, 1/0, 6/registerZ }


\paragraph{Syntax} \texttt{aswaph}\emph{coherency} \emph{offset27}[\emph{registerZ}] = \emph{registerT}

\paragraph{Description} The effective address is computed by adding the \emph{offset27} to the \emph{registerZ}.
This instruction implements an atomic memory swap of half word. The half word at effective address is written with the \emph{registerT} value, and its previous value is returned into the \emph{registerT}.


\paragraph{Modifier(s)}
\noindent\begin{itemize}
\item coherency (cf. coherency in section \ref{par:modifier-coherency} on page \pageref{par:modifier-coherency})
\end{itemize}

\paragraph{Execution} {Reservation table \textsf{LSU2\_MEMW\_AUXR\_AUXW.X} (Section~\ref{sec:reservation-LSU2-MEMW-AUXR-AUXW.X})}
~
{\begin{lstlisting}
stage ID:
  new offset = _SX_27(offset27);
  new coherency = _ZX_2(coherency);
stage RR:
  new base = GPR[registerZ];
  new address = base + offset;
stage E1:
  new argument2 = GPR[registerT];
  new result2 = _ZX_16(MEM.atomic_swap(address, 0x3, coherency << 32, argument2, registerT));
stage SR:
  GPR[registerT] = result2;
\end{lstlisting}}
\newpage

\paragraph{Encoding} Format \textsf{LSU\_ASWAPSB.D} (Section~\ref{sec:format-LSU-ASWAPSB.D}), \verb|lsucode5 = 001|\\

\bitfields{ 1/P, 2/00, 2/-- --, 27/extend27, 1/1, 2/00, 2/-- --, 27/offset27, 1/1, 2/01, 3/111, 2/coherency, 6/registerT, 2/11, 3/001, 1/--, 5/00010, 1/0, 6/registerZ }


\paragraph{Syntax} \texttt{aswaph}\emph{coherency} \emph{extend27\_\discretionary{}{}{}offset27}[\emph{registerZ}] = \emph{registerT}

\paragraph{Description} The effective address is computed by adding the \emph{extend27\_\discretionary{}{}{}offset27} to the \emph{registerZ}.
This instruction implements an atomic memory swap of half word. The half word at effective address is written with the \emph{registerT} value, and its previous value is returned into the \emph{registerT}.


\paragraph{Modifier(s)}
\noindent\begin{itemize}
\item coherency (cf. coherency in section \ref{par:modifier-coherency} on page \pageref{par:modifier-coherency})
\end{itemize}

\paragraph{Execution} {Reservation table \textsf{LSU2\_MEMW\_AUXR\_AUXW.Y} (Section~\ref{sec:reservation-LSU2-MEMW-AUXR-AUXW.Y})}
~
{\begin{lstlisting}
stage ID:
  new offset = _SX_54(extend27_offset27);
  new coherency = _ZX_2(coherency);
stage RR:
  new base = GPR[registerZ];
  new address = base + offset;
stage E1:
  new argument2 = GPR[registerT];
  new result2 = _ZX_16(MEM.atomic_swap(address, 0x3, coherency << 32, argument2, registerT));
stage SR:
  GPR[registerT] = result2;
\end{lstlisting}}
\newpage
\subsubsection[{\small ASWAPW: Atomic Swap Word}]{ASWAPW\hfill Atomic Swap Word}
\label{instr:ASWAPW}\index{aswapw}
 
\paragraph{Encoding} Format \textsf{LSU\_ASWAPSB} (Section~\ref{sec:format-LSU-ASWAPSB}), \verb|lsucode5 = 010|\\

\bitfields{ 1/P, 2/01, 3/111, 2/coherency, 6/registerT, 2/11, 3/010, 1/--, 5/00010, 1/0, 6/registerZ }


\paragraph{Syntax} \texttt{aswapw}\emph{coherency} [\emph{registerZ}] = \emph{registerT}

\paragraph{Description} The effective address is the value of the \emph{registerZ}.
This instruction implements an atomic memory swap of word. The word at effective address is written with the \emph{registerT} value, and its previous value is returned into the \emph{registerT}.


\paragraph{Modifier(s)}
\noindent\begin{itemize}
\item coherency (cf. coherency in section \ref{par:modifier-coherency} on page \pageref{par:modifier-coherency})
\end{itemize}

\paragraph{Execution} {Reservation table \textsf{LSU2\_MEMW\_AUXR\_AUXW} (Section~\ref{sec:reservation-LSU2-MEMW-AUXR-AUXW})}
~
{\begin{lstlisting}
stage ID:
  new coherency = _ZX_2(coherency);
stage RR:
  new address = GPR[registerZ];
stage E1:
  new argument2 = GPR[registerT];
  new result2 = _ZX_32(MEM.atomic_swap(address, 0xF, coherency << 32, argument2, registerT));
stage SR:
  GPR[registerT] = result2;
\end{lstlisting}}
\newpage

\paragraph{Encoding} Format \textsf{LSU\_ASWAPSB.O} (Section~\ref{sec:format-LSU-ASWAPSB.O}), \verb|lsucode5 = 010|\\

\bitfields{ 1/P, 2/00, 2/-- --, 27/offset27, 1/1, 2/01, 3/111, 2/coherency, 6/registerT, 2/11, 3/010, 1/--, 5/00010, 1/0, 6/registerZ }


\paragraph{Syntax} \texttt{aswapw}\emph{coherency} \emph{offset27}[\emph{registerZ}] = \emph{registerT}

\paragraph{Description} The effective address is computed by adding the \emph{offset27} to the \emph{registerZ}.
This instruction implements an atomic memory swap of word. The word at effective address is written with the \emph{registerT} value, and its previous value is returned into the \emph{registerT}.


\paragraph{Modifier(s)}
\noindent\begin{itemize}
\item coherency (cf. coherency in section \ref{par:modifier-coherency} on page \pageref{par:modifier-coherency})
\end{itemize}

\paragraph{Execution} {Reservation table \textsf{LSU2\_MEMW\_AUXR\_AUXW.X} (Section~\ref{sec:reservation-LSU2-MEMW-AUXR-AUXW.X})}
~
{\begin{lstlisting}
stage ID:
  new offset = _SX_27(offset27);
  new coherency = _ZX_2(coherency);
stage RR:
  new base = GPR[registerZ];
  new address = base + offset;
stage E1:
  new argument2 = GPR[registerT];
  new result2 = _ZX_32(MEM.atomic_swap(address, 0xF, coherency << 32, argument2, registerT));
stage SR:
  GPR[registerT] = result2;
\end{lstlisting}}
\newpage

\paragraph{Encoding} Format \textsf{LSU\_ASWAPSB.D} (Section~\ref{sec:format-LSU-ASWAPSB.D}), \verb|lsucode5 = 010|\\

\bitfields{ 1/P, 2/00, 2/-- --, 27/extend27, 1/1, 2/00, 2/-- --, 27/offset27, 1/1, 2/01, 3/111, 2/coherency, 6/registerT, 2/11, 3/010, 1/--, 5/00010, 1/0, 6/registerZ }


\paragraph{Syntax} \texttt{aswapw}\emph{coherency} \emph{extend27\_\discretionary{}{}{}offset27}[\emph{registerZ}] = \emph{registerT}

\paragraph{Description} The effective address is computed by adding the \emph{extend27\_\discretionary{}{}{}offset27} to the \emph{registerZ}.
This instruction implements an atomic memory swap of word. The word at effective address is written with the \emph{registerT} value, and its previous value is returned into the \emph{registerT}.


\paragraph{Modifier(s)}
\noindent\begin{itemize}
\item coherency (cf. coherency in section \ref{par:modifier-coherency} on page \pageref{par:modifier-coherency})
\end{itemize}

\paragraph{Execution} {Reservation table \textsf{LSU2\_MEMW\_AUXR\_AUXW.Y} (Section~\ref{sec:reservation-LSU2-MEMW-AUXR-AUXW.Y})}
~
{\begin{lstlisting}
stage ID:
  new offset = _SX_54(extend27_offset27);
  new coherency = _ZX_2(coherency);
stage RR:
  new base = GPR[registerZ];
  new address = base + offset;
stage E1:
  new argument2 = GPR[registerT];
  new result2 = _ZX_32(MEM.atomic_swap(address, 0xF, coherency << 32, argument2, registerT));
stage SR:
  GPR[registerT] = result2;
\end{lstlisting}}
\newpage
\subsubsection[{\small COPYO: Copy Octuple Word}]{COPYO\hfill Copy Octuple Word}
\label{instr:COPYO}\index{copyo}
 
\paragraph{Encoding} Format \textsf{LSU\_COPYO} (Section~\ref{sec:format-LSU-COPYO}), \verb|lsucode4 = 01|\\

\bitfields{ 1/P, 2/01, 5/10100, 4/registerN, 1/1, 1/1, 2/01, 10/-- -- -- -- -- -- -- -- -- --, 4/registerR, 1/--, 1/-- }


\paragraph{Syntax} \texttt{copyo} \emph{registerN} = \emph{registerR}

\paragraph{Description} The \emph{registerN} operand receives the \emph{registerR} value.


\paragraph{Execution} {Reservation table \textsf{LSU\_AUXR\_AUXW} (Section~\ref{sec:reservation-LSU-AUXR-AUXW})}
~
{\begin{lstlisting}
stage RR:
  new argument2 = QGR[registerR];
  new result1 = argument2;
stage E3:
  QGR[registerN] = result1;
\end{lstlisting}}
\newpage
\subsubsection[{\small D1INVAL: Data Cache 1 Invalidate}]{D1INVAL\hfill Data Cache 1 Invalidate}
\label{instr:D1INVAL}\index{d1inval}
 
\paragraph{Encoding} Format \textsf{LSU\_MCC} (Section~\ref{sec:format-LSU-MCC}), \verb|lsucode2 = 1000|\\

\bitfields{ 1/P, 2/01, 3/111, 2/00, 4/1000, 1/1, 1/1, 2/00, 16/-- -- -- -- -- -- -- -- -- -- -- -- -- -- -- -- }


\paragraph{Syntax} \texttt{d1inval}

\paragraph{Description} Request to invalidate the data cache. Cache contents is invalidated even if it is modified.


\paragraph{Execution} {Reservation table \textsf{LSU2\_MEMW} (Section~\ref{sec:reservation-LSU2-MEMW})}
~
{\begin{lstlisting}
stage E3:
  MEM.d1inval();
\end{lstlisting}}
\newpage
\subsubsection[{\small DFLUSHL: Data Cache Flush Line}]{DFLUSHL\hfill Data Cache Flush Line}
\label{instr:DFLUSHL}\index{dflushl}
 
\paragraph{Encoding} Format \textsf{LSU\_FXBO} (Section~\ref{sec:format-LSU-FXBO}), \verb|lsucode2 = 0011|\\

\bitfields{ 1/P, 2/01, 3/111, 2/00, 4/0011, 1/1, 1/1, 2/00, 10/signed10, 6/registerZ }


\paragraph{Syntax} \texttt{dflushl} \emph{signed10}[\emph{registerZ}]

\paragraph{Description} The effective address is computed by adding the \emph{signed10} to the \emph{registerZ}.
The effective address is sent to the data cache, with request to evict the corresponding cache line. Has no effects if the line is not cached.


\paragraph{Execution} {Reservation table \textsf{LSU} (Section~\ref{sec:reservation-LSU})}
~
{\begin{lstlisting}
stage ID:
  new offset = _SX_10(signed10);
stage RR:
  new base = GPR[registerZ];
  new address = base + offset;
stage E3:
  MEM.dflushl(address);
\end{lstlisting}}
\newpage

\paragraph{Encoding} Format \textsf{LSU\_FXBO.X} (Section~\ref{sec:format-LSU-FXBO.X}), \verb|lsucode2 = 0011|\\

\bitfields{ 1/P, 2/00, 2/-- --, 27/upper27, 1/1, 2/01, 3/111, 2/00, 4/0011, 1/1, 1/1, 2/00, 10/lower10, 6/registerZ }


\paragraph{Syntax} \texttt{dflushl} \emph{upper27\_\discretionary{}{}{}lower10}[\emph{registerZ}]

\paragraph{Description} The effective address is computed by adding the \emph{upper27\_\discretionary{}{}{}lower10} to the \emph{registerZ}.
The effective address is sent to the data cache, with request to evict the corresponding cache line. Has no effects if the line is not cached.


\paragraph{Execution} {Reservation table \textsf{LSU.X} (Section~\ref{sec:reservation-LSU.X})}
~
{\begin{lstlisting}
stage ID:
  new offset = _SX_37(upper27_lower10);
stage RR:
  new base = GPR[registerZ];
  new address = base + offset;
stage E3:
  MEM.dflushl(address);
\end{lstlisting}}
\newpage

\paragraph{Encoding} Format \textsf{LSU\_FXBO.Y} (Section~\ref{sec:format-LSU-FXBO.Y}), \verb|lsucode2 = 0011|\\

\bitfields{ 1/P, 2/00, 2/-- --, 27/extend27, 1/1, 2/00, 2/-- --, 27/upper27, 1/1, 2/01, 3/111, 2/00, 4/0011, 1/1, 1/1, 2/00, 10/lower10, 6/registerZ }


\paragraph{Syntax} \texttt{dflushl} \emph{extend27\_\discretionary{}{}{}upper27\_\discretionary{}{}{}lower10}[\emph{registerZ}]

\paragraph{Description} The effective address is computed by adding the \emph{extend27\_\discretionary{}{}{}upper27\_\discretionary{}{}{}lower10} to the \emph{registerZ}.
The effective address is sent to the data cache, with request to evict the corresponding cache line. Has no effects if the line is not cached.


\paragraph{Execution} {Reservation table \textsf{LSU.Y} (Section~\ref{sec:reservation-LSU.Y})}
~
{\begin{lstlisting}
stage ID:
  new offset = _WX_64(extend27_upper27_lower10);
stage RR:
  new base = GPR[registerZ];
  new address = base + offset;
stage E3:
  MEM.dflushl(address);
\end{lstlisting}}
\newpage

\paragraph{Encoding} Format \textsf{LSU\_FXBI} (Section~\ref{sec:format-LSU-FXBI}), \verb|lsucode2 = 0011|\\

\bitfields{ 1/P, 2/01, 3/111, 2/00, 4/0011, 1/1, 1/1, 2/10, 3/111, 1/--, 6/registerY, 6/registerZ }


\paragraph{Syntax} \texttt{dflushl} \emph{registerY}[\emph{registerZ}]

\paragraph{Description} The effective address is computed by adding the \emph{registerZ} to the \emph{registerY}.
The effective address is sent to the data cache, with request to evict the corresponding cache line. Has no effects if the line is not cached.


\paragraph{Execution} {Reservation table \textsf{LSU} (Section~\ref{sec:reservation-LSU})}
~
{\begin{lstlisting}
stage RR:
  new doscale = 0;
  new base = GPR[registerZ];
  new index = GPR[registerY];
  new scaling = doscale ? 64 : 1;
  new address = base + (index * scaling);
stage E3:
  MEM.dflushl(address);
\end{lstlisting}}
\newpage
\subsubsection[{\small DFLUSHSW: Data Cache Flush Line by Set and Way}]{DFLUSHSW\hfill Data Cache Flush Line by Set and Way}
\label{instr:DFLUSHSW}\index{dflushsw}
 
\paragraph{Encoding} Format \textsf{LSU\_FXSW} (Section~\ref{sec:format-LSU-FXSW}), \verb|lsucode2 = 1011|\\

\bitfields{ 1/P, 2/01, 3/111, 2/cachelev, 4/1011, 1/1, 1/1, 2/10, 3/111, 1/--, 6/registerY, 6/registerZ }


\paragraph{Syntax} \texttt{dflushsw}\emph{cachelev} \emph{registerY}, \emph{registerZ}

\paragraph{Description} The set and way are sent to the data cache, with request to evict the corresponding cache line. Has no effects if the line is not cached.


\paragraph{Modifier(s)}
\noindent\begin{itemize}
\item cachelev (cf. cachelev in section \ref{par:modifier-cachelev} on page \pageref{par:modifier-cachelev})
\end{itemize}

\paragraph{Execution} {Reservation table \textsf{LSU2\_MEMW} (Section~\ref{sec:reservation-LSU2-MEMW})}
~
{\begin{lstlisting}
stage ID:
  new cachelev = _ZX_2(cachelev);
stage RR:
  new way = GPR[registerZ];
  new set = GPR[registerY];
stage E3:
  MEM.dflushlsw(set, way, cachelev);
\end{lstlisting}}
\newpage
\subsubsection[{\small DINVALL: Data Cache Invalidate Line}]{DINVALL\hfill Data Cache Invalidate Line}
\label{instr:DINVALL}\index{dinvall}
 
\paragraph{Encoding} Format \textsf{LSU\_FXBO} (Section~\ref{sec:format-LSU-FXBO}), \verb|lsucode2 = 0001|\\

\bitfields{ 1/P, 2/01, 3/111, 2/00, 4/0001, 1/1, 1/1, 2/00, 10/signed10, 6/registerZ }


\paragraph{Syntax} \texttt{dinvall} \emph{signed10}[\emph{registerZ}]

\paragraph{Description} The effective address is computed by adding the \emph{signed10} to the \emph{registerZ}.
The effective address is sent to the data cache, with request to invalidate the corresponding cache line. Has no effects if the line is not cached.


\paragraph{Execution} {Reservation table \textsf{LSU} (Section~\ref{sec:reservation-LSU})}
~
{\begin{lstlisting}
stage ID:
  new offset = _SX_10(signed10);
stage RR:
  new base = GPR[registerZ];
  new address = base + offset;
stage E3:
  MEM.dinvall(address);
\end{lstlisting}}
\newpage

\paragraph{Encoding} Format \textsf{LSU\_FXBO.X} (Section~\ref{sec:format-LSU-FXBO.X}), \verb|lsucode2 = 0001|\\

\bitfields{ 1/P, 2/00, 2/-- --, 27/upper27, 1/1, 2/01, 3/111, 2/00, 4/0001, 1/1, 1/1, 2/00, 10/lower10, 6/registerZ }


\paragraph{Syntax} \texttt{dinvall} \emph{upper27\_\discretionary{}{}{}lower10}[\emph{registerZ}]

\paragraph{Description} The effective address is computed by adding the \emph{upper27\_\discretionary{}{}{}lower10} to the \emph{registerZ}.
The effective address is sent to the data cache, with request to invalidate the corresponding cache line. Has no effects if the line is not cached.


\paragraph{Execution} {Reservation table \textsf{LSU.X} (Section~\ref{sec:reservation-LSU.X})}
~
{\begin{lstlisting}
stage ID:
  new offset = _SX_37(upper27_lower10);
stage RR:
  new base = GPR[registerZ];
  new address = base + offset;
stage E3:
  MEM.dinvall(address);
\end{lstlisting}}
\newpage

\paragraph{Encoding} Format \textsf{LSU\_FXBO.Y} (Section~\ref{sec:format-LSU-FXBO.Y}), \verb|lsucode2 = 0001|\\

\bitfields{ 1/P, 2/00, 2/-- --, 27/extend27, 1/1, 2/00, 2/-- --, 27/upper27, 1/1, 2/01, 3/111, 2/00, 4/0001, 1/1, 1/1, 2/00, 10/lower10, 6/registerZ }


\paragraph{Syntax} \texttt{dinvall} \emph{extend27\_\discretionary{}{}{}upper27\_\discretionary{}{}{}lower10}[\emph{registerZ}]

\paragraph{Description} The effective address is computed by adding the \emph{extend27\_\discretionary{}{}{}upper27\_\discretionary{}{}{}lower10} to the \emph{registerZ}.
The effective address is sent to the data cache, with request to invalidate the corresponding cache line. Has no effects if the line is not cached.


\paragraph{Execution} {Reservation table \textsf{LSU.Y} (Section~\ref{sec:reservation-LSU.Y})}
~
{\begin{lstlisting}
stage ID:
  new offset = _WX_64(extend27_upper27_lower10);
stage RR:
  new base = GPR[registerZ];
  new address = base + offset;
stage E3:
  MEM.dinvall(address);
\end{lstlisting}}
\newpage

\paragraph{Encoding} Format \textsf{LSU\_FXBI} (Section~\ref{sec:format-LSU-FXBI}), \verb|lsucode2 = 0001|\\

\bitfields{ 1/P, 2/01, 3/111, 2/00, 4/0001, 1/1, 1/1, 2/10, 3/111, 1/--, 6/registerY, 6/registerZ }


\paragraph{Syntax} \texttt{dinvall} \emph{registerY}[\emph{registerZ}]

\paragraph{Description} The effective address is computed by adding the \emph{registerZ} to the \emph{registerY}.
The effective address is sent to the data cache, with request to invalidate the corresponding cache line. Has no effects if the line is not cached.


\paragraph{Execution} {Reservation table \textsf{LSU} (Section~\ref{sec:reservation-LSU})}
~
{\begin{lstlisting}
stage RR:
  new doscale = 0;
  new base = GPR[registerZ];
  new index = GPR[registerY];
  new scaling = doscale ? 64 : 1;
  new address = base + (index * scaling);
stage E3:
  MEM.dinvall(address);
\end{lstlisting}}
\newpage
\subsubsection[{\small DINVALSW: Data Cache Invalidate Line by Set and Way}]{DINVALSW\hfill Data Cache Invalidate Line by Set and Way}
\label{instr:DINVALSW}\index{dinvalsw}
 
\paragraph{Encoding} Format \textsf{LSU\_FXSW} (Section~\ref{sec:format-LSU-FXSW}), \verb|lsucode2 = 1001|\\

\bitfields{ 1/P, 2/01, 3/111, 2/cachelev, 4/1001, 1/1, 1/1, 2/10, 3/111, 1/--, 6/registerY, 6/registerZ }


\paragraph{Syntax} \texttt{dinvalsw}\emph{cachelev} \emph{registerY}, \emph{registerZ}

\paragraph{Description} The set and way are sent to the data cache, with request to invalidate the corresponding cache line. Has no effects if the line is not cached.


\paragraph{Modifier(s)}
\noindent\begin{itemize}
\item cachelev (cf. cachelev in section \ref{par:modifier-cachelev} on page \pageref{par:modifier-cachelev})
\end{itemize}

\paragraph{Execution} {Reservation table \textsf{LSU2\_MEMW} (Section~\ref{sec:reservation-LSU2-MEMW})}
~
{\begin{lstlisting}
stage ID:
  new cachelev = _ZX_2(cachelev);
stage RR:
  new way = GPR[registerZ];
  new set = GPR[registerY];
stage E3:
  if (!dinvallsw_owner()) {
    _THROW(PRIVILEGE);
  } else {
    MEM.dinvallsw(set, way, cachelev);
  }
\end{lstlisting}}
\newpage
\subsubsection[{\small DPURGEL: Data Cache Purge Line}]{DPURGEL\hfill Data Cache Purge Line}
\label{instr:DPURGEL}\index{dpurgel}
 
\paragraph{Encoding} Format \textsf{LSU\_FXBO} (Section~\ref{sec:format-LSU-FXBO}), \verb|lsucode2 = 0010|\\

\bitfields{ 1/P, 2/01, 3/111, 2/00, 4/0010, 1/1, 1/1, 2/00, 10/signed10, 6/registerZ }


\paragraph{Syntax} \texttt{dpurgel} \emph{signed10}[\emph{registerZ}]

\paragraph{Description} The effective address is computed by adding the \emph{signed10} to the \emph{registerZ}.
The effective address is sent to the data cache, with request to purge the corresponding cache line. Has no effects if the line is not cached.


\paragraph{Execution} {Reservation table \textsf{LSU} (Section~\ref{sec:reservation-LSU})}
~
{\begin{lstlisting}
stage ID:
  new offset = _SX_10(signed10);
stage RR:
  new base = GPR[registerZ];
  new address = base + offset;
stage E3:
  MEM.dpurgel(address);
\end{lstlisting}}
\newpage

\paragraph{Encoding} Format \textsf{LSU\_FXBO.X} (Section~\ref{sec:format-LSU-FXBO.X}), \verb|lsucode2 = 0010|\\

\bitfields{ 1/P, 2/00, 2/-- --, 27/upper27, 1/1, 2/01, 3/111, 2/00, 4/0010, 1/1, 1/1, 2/00, 10/lower10, 6/registerZ }


\paragraph{Syntax} \texttt{dpurgel} \emph{upper27\_\discretionary{}{}{}lower10}[\emph{registerZ}]

\paragraph{Description} The effective address is computed by adding the \emph{upper27\_\discretionary{}{}{}lower10} to the \emph{registerZ}.
The effective address is sent to the data cache, with request to purge the corresponding cache line. Has no effects if the line is not cached.


\paragraph{Execution} {Reservation table \textsf{LSU.X} (Section~\ref{sec:reservation-LSU.X})}
~
{\begin{lstlisting}
stage ID:
  new offset = _SX_37(upper27_lower10);
stage RR:
  new base = GPR[registerZ];
  new address = base + offset;
stage E3:
  MEM.dpurgel(address);
\end{lstlisting}}
\newpage

\paragraph{Encoding} Format \textsf{LSU\_FXBO.Y} (Section~\ref{sec:format-LSU-FXBO.Y}), \verb|lsucode2 = 0010|\\

\bitfields{ 1/P, 2/00, 2/-- --, 27/extend27, 1/1, 2/00, 2/-- --, 27/upper27, 1/1, 2/01, 3/111, 2/00, 4/0010, 1/1, 1/1, 2/00, 10/lower10, 6/registerZ }


\paragraph{Syntax} \texttt{dpurgel} \emph{extend27\_\discretionary{}{}{}upper27\_\discretionary{}{}{}lower10}[\emph{registerZ}]

\paragraph{Description} The effective address is computed by adding the \emph{extend27\_\discretionary{}{}{}upper27\_\discretionary{}{}{}lower10} to the \emph{registerZ}.
The effective address is sent to the data cache, with request to purge the corresponding cache line. Has no effects if the line is not cached.


\paragraph{Execution} {Reservation table \textsf{LSU.Y} (Section~\ref{sec:reservation-LSU.Y})}
~
{\begin{lstlisting}
stage ID:
  new offset = _WX_64(extend27_upper27_lower10);
stage RR:
  new base = GPR[registerZ];
  new address = base + offset;
stage E3:
  MEM.dpurgel(address);
\end{lstlisting}}
\newpage

\paragraph{Encoding} Format \textsf{LSU\_FXBI} (Section~\ref{sec:format-LSU-FXBI}), \verb|lsucode2 = 0010|\\

\bitfields{ 1/P, 2/01, 3/111, 2/00, 4/0010, 1/1, 1/1, 2/10, 3/111, 1/--, 6/registerY, 6/registerZ }


\paragraph{Syntax} \texttt{dpurgel} \emph{registerY}[\emph{registerZ}]

\paragraph{Description} The effective address is computed by adding the \emph{registerZ} to the \emph{registerY}.
The effective address is sent to the data cache, with request to purge the corresponding cache line. Has no effects if the line is not cached.


\paragraph{Execution} {Reservation table \textsf{LSU} (Section~\ref{sec:reservation-LSU})}
~
{\begin{lstlisting}
stage RR:
  new doscale = 0;
  new base = GPR[registerZ];
  new index = GPR[registerY];
  new scaling = doscale ? 64 : 1;
  new address = base + (index * scaling);
stage E3:
  MEM.dpurgel(address);
\end{lstlisting}}
\newpage
\subsubsection[{\small DPURGESW: Data Cache Purge Line by Set and Way}]{DPURGESW\hfill Data Cache Purge Line by Set and Way}
\label{instr:DPURGESW}\index{dpurgesw}
 
\paragraph{Encoding} Format \textsf{LSU\_FXSW} (Section~\ref{sec:format-LSU-FXSW}), \verb|lsucode2 = 1010|\\

\bitfields{ 1/P, 2/01, 3/111, 2/cachelev, 4/1010, 1/1, 1/1, 2/10, 3/111, 1/--, 6/registerY, 6/registerZ }


\paragraph{Syntax} \texttt{dpurgesw}\emph{cachelev} \emph{registerY}, \emph{registerZ}

\paragraph{Description} The set and way are sent to the data cache, with request to purge the corresponding cache line. Has no effects if the line is not cached.


\paragraph{Modifier(s)}
\noindent\begin{itemize}
\item cachelev (cf. cachelev in section \ref{par:modifier-cachelev} on page \pageref{par:modifier-cachelev})
\end{itemize}

\paragraph{Execution} {Reservation table \textsf{LSU2\_MEMW} (Section~\ref{sec:reservation-LSU2-MEMW})}
~
{\begin{lstlisting}
stage ID:
  new cachelev = _ZX_2(cachelev);
stage RR:
  new way = GPR[registerZ];
  new set = GPR[registerY];
stage E3:
  MEM.dpurgelsw(set, way, cachelev);
\end{lstlisting}}
\newpage
\subsubsection[{\small DTOUCHL: Data Cache Touch Line}]{DTOUCHL\hfill Data Cache Touch Line}
\label{instr:DTOUCHL}\index{dtouchl}
 
\paragraph{Encoding} Format \textsf{LSU\_FXBO} (Section~\ref{sec:format-LSU-FXBO}), \verb|lsucode2 = 0000|\\

\bitfields{ 1/P, 2/01, 3/111, 2/00, 4/0000, 1/1, 1/1, 2/00, 10/signed10, 6/registerZ }


\paragraph{Syntax} \texttt{dtouchl} \emph{signed10}[\emph{registerZ}]

\paragraph{Description} The effective address is computed by adding the \emph{signed10} to the \emph{registerZ}.
The effective address is sent to the data cache, with request to prefetch the corresponding cache line.


\paragraph{Execution} {Reservation table \textsf{LSU} (Section~\ref{sec:reservation-LSU})}
~
{\begin{lstlisting}
stage ID:
  new offset = _SX_10(signed10);
stage RR:
  new base = GPR[registerZ];
  new address = base + offset;
stage E3:
  MEM.dtouchl(address);
\end{lstlisting}}
\newpage

\paragraph{Encoding} Format \textsf{LSU\_FXBO.X} (Section~\ref{sec:format-LSU-FXBO.X}), \verb|lsucode2 = 0000|\\

\bitfields{ 1/P, 2/00, 2/-- --, 27/upper27, 1/1, 2/01, 3/111, 2/00, 4/0000, 1/1, 1/1, 2/00, 10/lower10, 6/registerZ }


\paragraph{Syntax} \texttt{dtouchl} \emph{upper27\_\discretionary{}{}{}lower10}[\emph{registerZ}]

\paragraph{Description} The effective address is computed by adding the \emph{upper27\_\discretionary{}{}{}lower10} to the \emph{registerZ}.
The effective address is sent to the data cache, with request to prefetch the corresponding cache line.


\paragraph{Execution} {Reservation table \textsf{LSU.X} (Section~\ref{sec:reservation-LSU.X})}
~
{\begin{lstlisting}
stage ID:
  new offset = _SX_37(upper27_lower10);
stage RR:
  new base = GPR[registerZ];
  new address = base + offset;
stage E3:
  MEM.dtouchl(address);
\end{lstlisting}}
\newpage

\paragraph{Encoding} Format \textsf{LSU\_FXBO.Y} (Section~\ref{sec:format-LSU-FXBO.Y}), \verb|lsucode2 = 0000|\\

\bitfields{ 1/P, 2/00, 2/-- --, 27/extend27, 1/1, 2/00, 2/-- --, 27/upper27, 1/1, 2/01, 3/111, 2/00, 4/0000, 1/1, 1/1, 2/00, 10/lower10, 6/registerZ }


\paragraph{Syntax} \texttt{dtouchl} \emph{extend27\_\discretionary{}{}{}upper27\_\discretionary{}{}{}lower10}[\emph{registerZ}]

\paragraph{Description} The effective address is computed by adding the \emph{extend27\_\discretionary{}{}{}upper27\_\discretionary{}{}{}lower10} to the \emph{registerZ}.
The effective address is sent to the data cache, with request to prefetch the corresponding cache line.


\paragraph{Execution} {Reservation table \textsf{LSU.Y} (Section~\ref{sec:reservation-LSU.Y})}
~
{\begin{lstlisting}
stage ID:
  new offset = _WX_64(extend27_upper27_lower10);
stage RR:
  new base = GPR[registerZ];
  new address = base + offset;
stage E3:
  MEM.dtouchl(address);
\end{lstlisting}}
\newpage

\paragraph{Encoding} Format \textsf{LSU\_FXBI} (Section~\ref{sec:format-LSU-FXBI}), \verb|lsucode2 = 0000|\\

\bitfields{ 1/P, 2/01, 3/111, 2/00, 4/0000, 1/1, 1/1, 2/10, 3/111, 1/--, 6/registerY, 6/registerZ }


\paragraph{Syntax} \texttt{dtouchl} \emph{registerY}[\emph{registerZ}]

\paragraph{Description} The effective address is computed by adding the \emph{registerZ} to the \emph{registerY}.
The effective address is sent to the data cache, with request to prefetch the corresponding cache line.


\paragraph{Execution} {Reservation table \textsf{LSU} (Section~\ref{sec:reservation-LSU})}
~
{\begin{lstlisting}
stage RR:
  new doscale = 0;
  new base = GPR[registerZ];
  new index = GPR[registerY];
  new scaling = doscale ? 64 : 1;
  new address = base + (index * scaling);
stage E3:
  MEM.dtouchl(address);
\end{lstlisting}}
\newpage
\subsubsection[{\small FENCE: Memory Fence}]{FENCE\hfill Memory Fence}
\label{instr:FENCE}\index{fence}
 
\paragraph{Encoding} Format \textsf{LSU\_FENCE} (Section~\ref{sec:format-LSU-FENCE}), \verb|lsucode2 = 1111|\\

\bitfields{ 1/P, 2/01, 3/111, 2/accesses, 4/1111, 1/1, 1/1, 2/00, 16/-- -- -- -- -- -- -- -- -- -- -- -- -- -- -- -- }


\paragraph{Syntax} \texttt{fence}\emph{accesses}

\paragraph{Description} Ensures that all issued memory accesses are committed to memory. The \emph{accesses} specifies for what types of access. FENCE.R instructions cannot pass earlier load instructions. FENCE.W instructions cannot pass earlier store, atomic and cache instructions. FENCE instructions produces the effects of FENCE.R, FENCE.W and in addition wait for the completion of outstanding prefetch instructions.


\paragraph{Modifier(s)}
\noindent\begin{itemize}
\item accesses (cf. accesses in section \ref{par:modifier-accesses} on page \pageref{par:modifier-accesses})
\end{itemize}

\paragraph{Execution} {Reservation table \textsf{LSU2\_MEMW} (Section~\ref{sec:reservation-LSU2-MEMW})}
~
{\begin{lstlisting}
stage ID:
  new accesses = _ZX_2(accesses);
stage E3:
  MEM.fence(accesses);
\end{lstlisting}}
\newpage
\subsubsection[{\small I1INVAL: Instruction Cache 1 Invalidate}]{I1INVAL\hfill Instruction Cache 1 Invalidate}
\label{instr:I1INVAL}\index{i1inval}
 
\paragraph{Encoding} Format \textsf{LSU\_MCC} (Section~\ref{sec:format-LSU-MCC}), \verb|lsucode2 = 1100|\\

\bitfields{ 1/P, 2/01, 3/111, 2/00, 4/1100, 1/1, 1/1, 2/00, 16/-- -- -- -- -- -- -- -- -- -- -- -- -- -- -- -- }


\paragraph{Syntax} \texttt{i1inval}

\paragraph{Description} Request to invalidate the instruction cache.


\paragraph{Execution} {Reservation table \textsf{LSU2\_MEMW} (Section~\ref{sec:reservation-LSU2-MEMW})}
~
{\begin{lstlisting}
stage E3:
  MEM.i1inval();
\end{lstlisting}}
\newpage
\subsubsection[{\small I1INVALS: Instruction Cache 1 Invalidate Set}]{I1INVALS\hfill Instruction Cache 1 Invalidate Set}
\label{instr:I1INVALS}\index{i1invals}
 
\paragraph{Encoding} Format \textsf{LSU\_FXBO} (Section~\ref{sec:format-LSU-FXBO}), \verb|lsucode2 = 0101|\\

\bitfields{ 1/P, 2/01, 3/111, 2/00, 4/0101, 1/1, 1/1, 2/00, 10/signed10, 6/registerZ }


\paragraph{Syntax} \texttt{i1invals} \emph{signed10}[\emph{registerZ}]

\paragraph{Description} The effective address is computed by adding the \emph{signed10} to the \emph{registerZ}.
The effective address is sent to the instruction cache, with request to invalidate the corresponding L1 cache set. (All the cache lines with the same index.)


\paragraph{Execution} {Reservation table \textsf{LSU2\_MEMW} (Section~\ref{sec:reservation-LSU2-MEMW})}
~
{\begin{lstlisting}
stage ID:
  new offset = _SX_10(signed10);
stage RR:
  new base = GPR[registerZ];
  new address = base + offset;
stage E3:
  MEM.i1invals(address);
\end{lstlisting}}
\newpage

\paragraph{Encoding} Format \textsf{LSU\_FXBO.X} (Section~\ref{sec:format-LSU-FXBO.X}), \verb|lsucode2 = 0101|\\

\bitfields{ 1/P, 2/00, 2/-- --, 27/upper27, 1/1, 2/01, 3/111, 2/00, 4/0101, 1/1, 1/1, 2/00, 10/lower10, 6/registerZ }


\paragraph{Syntax} \texttt{i1invals} \emph{upper27\_\discretionary{}{}{}lower10}[\emph{registerZ}]

\paragraph{Description} The effective address is computed by adding the \emph{upper27\_\discretionary{}{}{}lower10} to the \emph{registerZ}.
The effective address is sent to the instruction cache, with request to invalidate the corresponding L1 cache set. (All the cache lines with the same index.)


\paragraph{Execution} {Reservation table \textsf{LSU2\_MEMW.X} (Section~\ref{sec:reservation-LSU2-MEMW.X})}
~
{\begin{lstlisting}
stage ID:
  new offset = _SX_37(upper27_lower10);
stage RR:
  new base = GPR[registerZ];
  new address = base + offset;
stage E3:
  MEM.i1invals(address);
\end{lstlisting}}
\newpage

\paragraph{Encoding} Format \textsf{LSU\_FXBO.Y} (Section~\ref{sec:format-LSU-FXBO.Y}), \verb|lsucode2 = 0101|\\

\bitfields{ 1/P, 2/00, 2/-- --, 27/extend27, 1/1, 2/00, 2/-- --, 27/upper27, 1/1, 2/01, 3/111, 2/00, 4/0101, 1/1, 1/1, 2/00, 10/lower10, 6/registerZ }


\paragraph{Syntax} \texttt{i1invals} \emph{extend27\_\discretionary{}{}{}upper27\_\discretionary{}{}{}lower10}[\emph{registerZ}]

\paragraph{Description} The effective address is computed by adding the \emph{extend27\_\discretionary{}{}{}upper27\_\discretionary{}{}{}lower10} to the \emph{registerZ}.
The effective address is sent to the instruction cache, with request to invalidate the corresponding L1 cache set. (All the cache lines with the same index.)


\paragraph{Execution} {Reservation table \textsf{LSU2\_MEMW.Y} (Section~\ref{sec:reservation-LSU2-MEMW.Y})}
~
{\begin{lstlisting}
stage ID:
  new offset = _WX_64(extend27_upper27_lower10);
stage RR:
  new base = GPR[registerZ];
  new address = base + offset;
stage E3:
  MEM.i1invals(address);
\end{lstlisting}}
\newpage

\paragraph{Encoding} Format \textsf{LSU\_FXBI} (Section~\ref{sec:format-LSU-FXBI}), \verb|lsucode2 = 0101|\\

\bitfields{ 1/P, 2/01, 3/111, 2/00, 4/0101, 1/1, 1/1, 2/10, 3/111, 1/--, 6/registerY, 6/registerZ }


\paragraph{Syntax} \texttt{i1invals} \emph{registerY}[\emph{registerZ}]

\paragraph{Description} The effective address is computed by adding the \emph{registerZ} to the \emph{registerY}.
The effective address is sent to the instruction cache, with request to invalidate the corresponding L1 cache set. (All the cache lines with the same index.)


\paragraph{Execution} {Reservation table \textsf{LSU2\_MEMW} (Section~\ref{sec:reservation-LSU2-MEMW})}
~
{\begin{lstlisting}
stage RR:
  new doscale = 0;
  new base = GPR[registerZ];
  new index = GPR[registerY];
  new scaling = doscale ? 64 : 1;
  new address = base + (index * scaling);
stage E3:
  MEM.i1invals(address);
\end{lstlisting}}
\newpage
\subsubsection[{\small LBS: Load Byte Sign Extended}]{LBS\hfill Load Byte Sign Extended}
\label{instr:LBS}\index{lbs}
 
\paragraph{Encoding} Format \textsf{LSU\_LSBO} (Section~\ref{sec:format-LSU-LSBO}), \verb|loadcode = 001|\\

\bitfields{ 1/P, 2/01, 3/001, 2/variant, 6/registerW, 2/00, 10/signed10, 6/registerZ }


\paragraph{Syntax} \texttt{lbs}\emph{variant} \emph{registerW} = \emph{signed10}[\emph{registerZ}]

\paragraph{Description} The effective address is computed by adding the \emph{signed10} to the \emph{registerZ}.
The \emph{registerW} is loaded from the memory byte at the effective address after sign extension.


\paragraph{Modifier(s)}
\noindent\begin{itemize}
\item variant (cf. variant in section \ref{par:modifier-variant} on page \pageref{par:modifier-variant})
\end{itemize}

\paragraph{Execution} {Reservation table \textsf{LSU\_AUXW} (Section~\ref{sec:reservation-LSU-AUXW})}
~
{\begin{lstlisting}
stage ID:
  new variant = _ZX_2(variant);
  new offset = _SX_10(signed10);
stage RR:
  new base = GPR[registerZ];
  new address = base + offset;
  new result1 = _SX_8(MEM.load(address, 0x1, variant, registerW));
stage CR:
  GPR[registerW] = result1;
\end{lstlisting}}
\newpage

\paragraph{Encoding} Format \textsf{LSU\_LSBO.X} (Section~\ref{sec:format-LSU-LSBO.X}), \verb|loadcode = 001|\\

\bitfields{ 1/P, 2/00, 2/-- --, 27/upper27, 1/1, 2/01, 3/001, 2/variant, 6/registerW, 2/00, 10/lower10, 6/registerZ }


\paragraph{Syntax} \texttt{lbs}\emph{variant} \emph{registerW} = \emph{upper27\_\discretionary{}{}{}lower10}[\emph{registerZ}]

\paragraph{Description} The effective address is computed by adding the \emph{upper27\_\discretionary{}{}{}lower10} to the \emph{registerZ}.
The \emph{registerW} is loaded from the memory byte at the effective address after sign extension.


\paragraph{Modifier(s)}
\noindent\begin{itemize}
\item variant (cf. variant in section \ref{par:modifier-variant} on page \pageref{par:modifier-variant})
\end{itemize}

\paragraph{Execution} {Reservation table \textsf{LSU\_AUXW.X} (Section~\ref{sec:reservation-LSU-AUXW.X})}
~
{\begin{lstlisting}
stage ID:
  new variant = _ZX_2(variant);
  new offset = _SX_37(upper27_lower10);
stage RR:
  new base = GPR[registerZ];
  new address = base + offset;
  new result1 = _SX_8(MEM.load(address, 0x1, variant, registerW));
stage CR:
  GPR[registerW] = result1;
\end{lstlisting}}
\newpage

\paragraph{Encoding} Format \textsf{LSU\_LSBO.Y} (Section~\ref{sec:format-LSU-LSBO.Y}), \verb|loadcode = 001|\\

\bitfields{ 1/P, 2/00, 2/-- --, 27/extend27, 1/1, 2/00, 2/-- --, 27/upper27, 1/1, 2/01, 3/001, 2/variant, 6/registerW, 2/00, 10/lower10, 6/registerZ }


\paragraph{Syntax} \texttt{lbs}\emph{variant} \emph{registerW} = \emph{extend27\_\discretionary{}{}{}upper27\_\discretionary{}{}{}lower10}[\emph{registerZ}]

\paragraph{Description} The effective address is computed by adding the \emph{extend27\_\discretionary{}{}{}upper27\_\discretionary{}{}{}lower10} to the \emph{registerZ}.
The \emph{registerW} is loaded from the memory byte at the effective address after sign extension.


\paragraph{Modifier(s)}
\noindent\begin{itemize}
\item variant (cf. variant in section \ref{par:modifier-variant} on page \pageref{par:modifier-variant})
\end{itemize}

\paragraph{Execution} {Reservation table \textsf{LSU\_AUXW.Y} (Section~\ref{sec:reservation-LSU-AUXW.Y})}
~
{\begin{lstlisting}
stage ID:
  new variant = _ZX_2(variant);
  new offset = _WX_64(extend27_upper27_lower10);
stage RR:
  new base = GPR[registerZ];
  new address = base + offset;
  new result1 = _SX_8(MEM.load(address, 0x1, variant, registerW));
stage CR:
  GPR[registerW] = result1;
\end{lstlisting}}
\newpage

\paragraph{Encoding} Format \textsf{LSU\_LSBI} (Section~\ref{sec:format-LSU-LSBI}), \verb|loadcode = 001|\\

\bitfields{ 1/P, 2/01, 3/001, 2/variant, 6/registerW, 2/10, 3/111, 1/--, 6/registerY, 6/registerZ }


\paragraph{Syntax} \texttt{lbs}\emph{variant} \emph{registerW} = \emph{registerY}[\emph{registerZ}]

\paragraph{Description} The effective address is computed by adding the \emph{registerZ} to the \emph{registerY} and scaled by the \emph{variant}.
The \emph{registerW} is loaded from the memory byte at the effective address after sign extension.


\paragraph{Modifier(s)}
\noindent\begin{itemize}
\item variant (cf. variant in section \ref{par:modifier-variant} on page \pageref{par:modifier-variant})
\end{itemize}

\paragraph{Execution} {Reservation table \textsf{LSU\_AUXW} (Section~\ref{sec:reservation-LSU-AUXW})}
~
{\begin{lstlisting}
stage ID:
  new variant = _ZX_2(variant);
stage RR:
  new doscale = 0;
  new base = GPR[registerZ];
  new index = GPR[registerY];
  new scaling = doscale ? 1 : 1;
  new address = base + (index * scaling);
  new result1 = _SX_8(MEM.load(address, 0x1, variant, registerW));
stage CR:
  GPR[registerW] = result1;
\end{lstlisting}}
\newpage
\subsubsection[{\small LBZ: Load Byte Zero Extended}]{LBZ\hfill Load Byte Zero Extended}
\label{instr:LBZ}\index{lbz}
 
\paragraph{Encoding} Format \textsf{LSU\_LSBO} (Section~\ref{sec:format-LSU-LSBO}), \verb|loadcode = 000|\\

\bitfields{ 1/P, 2/01, 3/000, 2/variant, 6/registerW, 2/00, 10/signed10, 6/registerZ }


\paragraph{Syntax} \texttt{lbz}\emph{variant} \emph{registerW} = \emph{signed10}[\emph{registerZ}]

\paragraph{Description} The effective address is computed by adding the \emph{signed10} to the \emph{registerZ}.
The \emph{registerW} is loaded from the memory byte at the effective address after zero extension.


\paragraph{Modifier(s)}
\noindent\begin{itemize}
\item variant (cf. variant in section \ref{par:modifier-variant} on page \pageref{par:modifier-variant})
\end{itemize}

\paragraph{Execution} {Reservation table \textsf{LSU\_AUXW} (Section~\ref{sec:reservation-LSU-AUXW})}
~
{\begin{lstlisting}
stage ID:
  new variant = _ZX_2(variant);
  new offset = _SX_10(signed10);
stage RR:
  new base = GPR[registerZ];
  new address = base + offset;
  new result1 = _ZX_8(MEM.load(address, 0x1, variant, registerW));
stage CR:
  GPR[registerW] = result1;
\end{lstlisting}}
\newpage

\paragraph{Encoding} Format \textsf{LSU\_LSBO.X} (Section~\ref{sec:format-LSU-LSBO.X}), \verb|loadcode = 000|\\

\bitfields{ 1/P, 2/00, 2/-- --, 27/upper27, 1/1, 2/01, 3/000, 2/variant, 6/registerW, 2/00, 10/lower10, 6/registerZ }


\paragraph{Syntax} \texttt{lbz}\emph{variant} \emph{registerW} = \emph{upper27\_\discretionary{}{}{}lower10}[\emph{registerZ}]

\paragraph{Description} The effective address is computed by adding the \emph{upper27\_\discretionary{}{}{}lower10} to the \emph{registerZ}.
The \emph{registerW} is loaded from the memory byte at the effective address after zero extension.


\paragraph{Modifier(s)}
\noindent\begin{itemize}
\item variant (cf. variant in section \ref{par:modifier-variant} on page \pageref{par:modifier-variant})
\end{itemize}

\paragraph{Execution} {Reservation table \textsf{LSU\_AUXW.X} (Section~\ref{sec:reservation-LSU-AUXW.X})}
~
{\begin{lstlisting}
stage ID:
  new variant = _ZX_2(variant);
  new offset = _SX_37(upper27_lower10);
stage RR:
  new base = GPR[registerZ];
  new address = base + offset;
  new result1 = _ZX_8(MEM.load(address, 0x1, variant, registerW));
stage CR:
  GPR[registerW] = result1;
\end{lstlisting}}
\newpage

\paragraph{Encoding} Format \textsf{LSU\_LSBO.Y} (Section~\ref{sec:format-LSU-LSBO.Y}), \verb|loadcode = 000|\\

\bitfields{ 1/P, 2/00, 2/-- --, 27/extend27, 1/1, 2/00, 2/-- --, 27/upper27, 1/1, 2/01, 3/000, 2/variant, 6/registerW, 2/00, 10/lower10, 6/registerZ }


\paragraph{Syntax} \texttt{lbz}\emph{variant} \emph{registerW} = \emph{extend27\_\discretionary{}{}{}upper27\_\discretionary{}{}{}lower10}[\emph{registerZ}]

\paragraph{Description} The effective address is computed by adding the \emph{extend27\_\discretionary{}{}{}upper27\_\discretionary{}{}{}lower10} to the \emph{registerZ}.
The \emph{registerW} is loaded from the memory byte at the effective address after zero extension.


\paragraph{Modifier(s)}
\noindent\begin{itemize}
\item variant (cf. variant in section \ref{par:modifier-variant} on page \pageref{par:modifier-variant})
\end{itemize}

\paragraph{Execution} {Reservation table \textsf{LSU\_AUXW.Y} (Section~\ref{sec:reservation-LSU-AUXW.Y})}
~
{\begin{lstlisting}
stage ID:
  new variant = _ZX_2(variant);
  new offset = _WX_64(extend27_upper27_lower10);
stage RR:
  new base = GPR[registerZ];
  new address = base + offset;
  new result1 = _ZX_8(MEM.load(address, 0x1, variant, registerW));
stage CR:
  GPR[registerW] = result1;
\end{lstlisting}}
\newpage

\paragraph{Encoding} Format \textsf{LSU\_LSBI} (Section~\ref{sec:format-LSU-LSBI}), \verb|loadcode = 000|\\

\bitfields{ 1/P, 2/01, 3/000, 2/variant, 6/registerW, 2/10, 3/111, 1/--, 6/registerY, 6/registerZ }


\paragraph{Syntax} \texttt{lbz}\emph{variant} \emph{registerW} = \emph{registerY}[\emph{registerZ}]

\paragraph{Description} The effective address is computed by adding the \emph{registerZ} to the \emph{registerY} and scaled by the \emph{variant}.
The \emph{registerW} is loaded from the memory byte at the effective address after zero extension.


\paragraph{Modifier(s)}
\noindent\begin{itemize}
\item variant (cf. variant in section \ref{par:modifier-variant} on page \pageref{par:modifier-variant})
\end{itemize}

\paragraph{Execution} {Reservation table \textsf{LSU\_AUXW} (Section~\ref{sec:reservation-LSU-AUXW})}
~
{\begin{lstlisting}
stage ID:
  new variant = _ZX_2(variant);
stage RR:
  new doscale = 0;
  new base = GPR[registerZ];
  new index = GPR[registerY];
  new scaling = doscale ? 1 : 1;
  new address = base + (index * scaling);
  new result1 = _ZX_8(MEM.load(address, 0x1, variant, registerW));
stage CR:
  GPR[registerW] = result1;
\end{lstlisting}}
\newpage
\subsubsection[{\small LD: Load Double Word}]{LD\hfill Load Double Word}
\label{instr:LD}\index{ld}
 
\paragraph{Encoding} Format \textsf{LSU\_LSBO} (Section~\ref{sec:format-LSU-LSBO}), \verb|loadcode = 110|\\

\bitfields{ 1/P, 2/01, 3/110, 2/variant, 6/registerW, 2/00, 10/signed10, 6/registerZ }


\paragraph{Syntax} \texttt{ld}\emph{variant} \emph{registerW} = \emph{signed10}[\emph{registerZ}]

\paragraph{Description} The effective address is computed by adding the \emph{signed10} to the \emph{registerZ}.
The \emph{registerW} is loaded from the memory double word at the effective address.


\paragraph{Modifier(s)}
\noindent\begin{itemize}
\item variant (cf. variant in section \ref{par:modifier-variant} on page \pageref{par:modifier-variant})
\end{itemize}

\paragraph{Execution} {Reservation table \textsf{LSU\_AUXW} (Section~\ref{sec:reservation-LSU-AUXW})}
~
{\begin{lstlisting}
stage ID:
  new variant = _ZX_2(variant);
  new offset = _SX_10(signed10);
stage RR:
  new base = GPR[registerZ];
  new address = base + offset;
  new result1 = MEM.load(address, 0xFF, variant, registerW);
stage CR:
  GPR[registerW] = result1;
\end{lstlisting}}
\newpage

\paragraph{Encoding} Format \textsf{LSU\_LSBO.X} (Section~\ref{sec:format-LSU-LSBO.X}), \verb|loadcode = 110|\\

\bitfields{ 1/P, 2/00, 2/-- --, 27/upper27, 1/1, 2/01, 3/110, 2/variant, 6/registerW, 2/00, 10/lower10, 6/registerZ }


\paragraph{Syntax} \texttt{ld}\emph{variant} \emph{registerW} = \emph{upper27\_\discretionary{}{}{}lower10}[\emph{registerZ}]

\paragraph{Description} The effective address is computed by adding the \emph{upper27\_\discretionary{}{}{}lower10} to the \emph{registerZ}.
The \emph{registerW} is loaded from the memory double word at the effective address.


\paragraph{Modifier(s)}
\noindent\begin{itemize}
\item variant (cf. variant in section \ref{par:modifier-variant} on page \pageref{par:modifier-variant})
\end{itemize}

\paragraph{Execution} {Reservation table \textsf{LSU\_AUXW.X} (Section~\ref{sec:reservation-LSU-AUXW.X})}
~
{\begin{lstlisting}
stage ID:
  new variant = _ZX_2(variant);
  new offset = _SX_37(upper27_lower10);
stage RR:
  new base = GPR[registerZ];
  new address = base + offset;
  new result1 = MEM.load(address, 0xFF, variant, registerW);
stage CR:
  GPR[registerW] = result1;
\end{lstlisting}}
\newpage

\paragraph{Encoding} Format \textsf{LSU\_LSBO.Y} (Section~\ref{sec:format-LSU-LSBO.Y}), \verb|loadcode = 110|\\

\bitfields{ 1/P, 2/00, 2/-- --, 27/extend27, 1/1, 2/00, 2/-- --, 27/upper27, 1/1, 2/01, 3/110, 2/variant, 6/registerW, 2/00, 10/lower10, 6/registerZ }


\paragraph{Syntax} \texttt{ld}\emph{variant} \emph{registerW} = \emph{extend27\_\discretionary{}{}{}upper27\_\discretionary{}{}{}lower10}[\emph{registerZ}]

\paragraph{Description} The effective address is computed by adding the \emph{extend27\_\discretionary{}{}{}upper27\_\discretionary{}{}{}lower10} to the \emph{registerZ}.
The \emph{registerW} is loaded from the memory double word at the effective address.


\paragraph{Modifier(s)}
\noindent\begin{itemize}
\item variant (cf. variant in section \ref{par:modifier-variant} on page \pageref{par:modifier-variant})
\end{itemize}

\paragraph{Execution} {Reservation table \textsf{LSU\_AUXW.Y} (Section~\ref{sec:reservation-LSU-AUXW.Y})}
~
{\begin{lstlisting}
stage ID:
  new variant = _ZX_2(variant);
  new offset = _WX_64(extend27_upper27_lower10);
stage RR:
  new base = GPR[registerZ];
  new address = base + offset;
  new result1 = MEM.load(address, 0xFF, variant, registerW);
stage CR:
  GPR[registerW] = result1;
\end{lstlisting}}
\newpage

\paragraph{Encoding} Format \textsf{LSU\_LSBI} (Section~\ref{sec:format-LSU-LSBI}), \verb|loadcode = 110|\\

\bitfields{ 1/P, 2/01, 3/110, 2/variant, 6/registerW, 2/10, 3/111, 1/--, 6/registerY, 6/registerZ }


\paragraph{Syntax} \texttt{ld}\emph{variant} \emph{registerW} = \emph{registerY}[\emph{registerZ}]

\paragraph{Description} The effective address is computed by adding the \emph{registerZ} to the \emph{registerY} and scaled by the \emph{variant}.
The \emph{registerW} is loaded from the memory double word at the effective address.


\paragraph{Modifier(s)}
\noindent\begin{itemize}
\item variant (cf. variant in section \ref{par:modifier-variant} on page \pageref{par:modifier-variant})
\end{itemize}

\paragraph{Execution} {Reservation table \textsf{LSU\_AUXW} (Section~\ref{sec:reservation-LSU-AUXW})}
~
{\begin{lstlisting}
stage ID:
  new variant = _ZX_2(variant);
stage RR:
  new doscale = 0;
  new base = GPR[registerZ];
  new index = GPR[registerY];
  new scaling = doscale ? 8 : 1;
  new address = base + (index * scaling);
  new result1 = MEM.load(address, 0xFF, variant, registerW);
stage CR:
  GPR[registerW] = result1;
\end{lstlisting}}
\newpage
\subsubsection[{\small LHS: Load Half Word Sign Extended}]{LHS\hfill Load Half Word Sign Extended}
\label{instr:LHS}\index{lhs}
 
\paragraph{Encoding} Format \textsf{LSU\_LSBO} (Section~\ref{sec:format-LSU-LSBO}), \verb|loadcode = 011|\\

\bitfields{ 1/P, 2/01, 3/011, 2/variant, 6/registerW, 2/00, 10/signed10, 6/registerZ }


\paragraph{Syntax} \texttt{lhs}\emph{variant} \emph{registerW} = \emph{signed10}[\emph{registerZ}]

\paragraph{Description} The effective address is computed by adding the \emph{signed10} to the \emph{registerZ}.
The \emph{registerW} is loaded from the memory half word at the effective address after sign extension.


\paragraph{Modifier(s)}
\noindent\begin{itemize}
\item variant (cf. variant in section \ref{par:modifier-variant} on page \pageref{par:modifier-variant})
\end{itemize}

\paragraph{Execution} {Reservation table \textsf{LSU\_AUXW} (Section~\ref{sec:reservation-LSU-AUXW})}
~
{\begin{lstlisting}
stage ID:
  new variant = _ZX_2(variant);
  new offset = _SX_10(signed10);
stage RR:
  new base = GPR[registerZ];
  new address = base + offset;
  new result1 = _SX_16(MEM.load(address, 0x3, variant, registerW));
stage CR:
  GPR[registerW] = result1;
\end{lstlisting}}
\newpage

\paragraph{Encoding} Format \textsf{LSU\_LSBO.X} (Section~\ref{sec:format-LSU-LSBO.X}), \verb|loadcode = 011|\\

\bitfields{ 1/P, 2/00, 2/-- --, 27/upper27, 1/1, 2/01, 3/011, 2/variant, 6/registerW, 2/00, 10/lower10, 6/registerZ }


\paragraph{Syntax} \texttt{lhs}\emph{variant} \emph{registerW} = \emph{upper27\_\discretionary{}{}{}lower10}[\emph{registerZ}]

\paragraph{Description} The effective address is computed by adding the \emph{upper27\_\discretionary{}{}{}lower10} to the \emph{registerZ}.
The \emph{registerW} is loaded from the memory half word at the effective address after sign extension.


\paragraph{Modifier(s)}
\noindent\begin{itemize}
\item variant (cf. variant in section \ref{par:modifier-variant} on page \pageref{par:modifier-variant})
\end{itemize}

\paragraph{Execution} {Reservation table \textsf{LSU\_AUXW.X} (Section~\ref{sec:reservation-LSU-AUXW.X})}
~
{\begin{lstlisting}
stage ID:
  new variant = _ZX_2(variant);
  new offset = _SX_37(upper27_lower10);
stage RR:
  new base = GPR[registerZ];
  new address = base + offset;
  new result1 = _SX_16(MEM.load(address, 0x3, variant, registerW));
stage CR:
  GPR[registerW] = result1;
\end{lstlisting}}
\newpage

\paragraph{Encoding} Format \textsf{LSU\_LSBO.Y} (Section~\ref{sec:format-LSU-LSBO.Y}), \verb|loadcode = 011|\\

\bitfields{ 1/P, 2/00, 2/-- --, 27/extend27, 1/1, 2/00, 2/-- --, 27/upper27, 1/1, 2/01, 3/011, 2/variant, 6/registerW, 2/00, 10/lower10, 6/registerZ }


\paragraph{Syntax} \texttt{lhs}\emph{variant} \emph{registerW} = \emph{extend27\_\discretionary{}{}{}upper27\_\discretionary{}{}{}lower10}[\emph{registerZ}]

\paragraph{Description} The effective address is computed by adding the \emph{extend27\_\discretionary{}{}{}upper27\_\discretionary{}{}{}lower10} to the \emph{registerZ}.
The \emph{registerW} is loaded from the memory half word at the effective address after sign extension.


\paragraph{Modifier(s)}
\noindent\begin{itemize}
\item variant (cf. variant in section \ref{par:modifier-variant} on page \pageref{par:modifier-variant})
\end{itemize}

\paragraph{Execution} {Reservation table \textsf{LSU\_AUXW.Y} (Section~\ref{sec:reservation-LSU-AUXW.Y})}
~
{\begin{lstlisting}
stage ID:
  new variant = _ZX_2(variant);
  new offset = _WX_64(extend27_upper27_lower10);
stage RR:
  new base = GPR[registerZ];
  new address = base + offset;
  new result1 = _SX_16(MEM.load(address, 0x3, variant, registerW));
stage CR:
  GPR[registerW] = result1;
\end{lstlisting}}
\newpage

\paragraph{Encoding} Format \textsf{LSU\_LSBI} (Section~\ref{sec:format-LSU-LSBI}), \verb|loadcode = 011|\\

\bitfields{ 1/P, 2/01, 3/011, 2/variant, 6/registerW, 2/10, 3/111, 1/--, 6/registerY, 6/registerZ }


\paragraph{Syntax} \texttt{lhs}\emph{variant} \emph{registerW} = \emph{registerY}[\emph{registerZ}]

\paragraph{Description} The effective address is computed by adding the \emph{registerZ} to the \emph{registerY} and scaled by the \emph{variant}.
The \emph{registerW} is loaded from the memory half word at the effective address after sign extension.


\paragraph{Modifier(s)}
\noindent\begin{itemize}
\item variant (cf. variant in section \ref{par:modifier-variant} on page \pageref{par:modifier-variant})
\end{itemize}

\paragraph{Execution} {Reservation table \textsf{LSU\_AUXW} (Section~\ref{sec:reservation-LSU-AUXW})}
~
{\begin{lstlisting}
stage ID:
  new variant = _ZX_2(variant);
stage RR:
  new doscale = 0;
  new base = GPR[registerZ];
  new index = GPR[registerY];
  new scaling = doscale ? 2 : 1;
  new address = base + (index * scaling);
  new result1 = _SX_16(MEM.load(address, 0x3, variant, registerW));
stage CR:
  GPR[registerW] = result1;
\end{lstlisting}}
\newpage
\subsubsection[{\small LHZ: Load Half Word Zero Extended}]{LHZ\hfill Load Half Word Zero Extended}
\label{instr:LHZ}\index{lhz}
 
\paragraph{Encoding} Format \textsf{LSU\_LSBO} (Section~\ref{sec:format-LSU-LSBO}), \verb|loadcode = 010|\\

\bitfields{ 1/P, 2/01, 3/010, 2/variant, 6/registerW, 2/00, 10/signed10, 6/registerZ }


\paragraph{Syntax} \texttt{lhz}\emph{variant} \emph{registerW} = \emph{signed10}[\emph{registerZ}]

\paragraph{Description} The effective address is computed by adding the \emph{signed10} to the \emph{registerZ}.
The \emph{registerW} is loaded from the memory half word at the effective address after zero extension.


\paragraph{Modifier(s)}
\noindent\begin{itemize}
\item variant (cf. variant in section \ref{par:modifier-variant} on page \pageref{par:modifier-variant})
\end{itemize}

\paragraph{Execution} {Reservation table \textsf{LSU\_AUXW} (Section~\ref{sec:reservation-LSU-AUXW})}
~
{\begin{lstlisting}
stage ID:
  new variant = _ZX_2(variant);
  new offset = _SX_10(signed10);
stage RR:
  new base = GPR[registerZ];
  new address = base + offset;
  new result1 = _ZX_16(MEM.load(address, 0x3, variant, registerW));
stage CR:
  GPR[registerW] = result1;
\end{lstlisting}}
\newpage

\paragraph{Encoding} Format \textsf{LSU\_LSBO.X} (Section~\ref{sec:format-LSU-LSBO.X}), \verb|loadcode = 010|\\

\bitfields{ 1/P, 2/00, 2/-- --, 27/upper27, 1/1, 2/01, 3/010, 2/variant, 6/registerW, 2/00, 10/lower10, 6/registerZ }


\paragraph{Syntax} \texttt{lhz}\emph{variant} \emph{registerW} = \emph{upper27\_\discretionary{}{}{}lower10}[\emph{registerZ}]

\paragraph{Description} The effective address is computed by adding the \emph{upper27\_\discretionary{}{}{}lower10} to the \emph{registerZ}.
The \emph{registerW} is loaded from the memory half word at the effective address after zero extension.


\paragraph{Modifier(s)}
\noindent\begin{itemize}
\item variant (cf. variant in section \ref{par:modifier-variant} on page \pageref{par:modifier-variant})
\end{itemize}

\paragraph{Execution} {Reservation table \textsf{LSU\_AUXW.X} (Section~\ref{sec:reservation-LSU-AUXW.X})}
~
{\begin{lstlisting}
stage ID:
  new variant = _ZX_2(variant);
  new offset = _SX_37(upper27_lower10);
stage RR:
  new base = GPR[registerZ];
  new address = base + offset;
  new result1 = _ZX_16(MEM.load(address, 0x3, variant, registerW));
stage CR:
  GPR[registerW] = result1;
\end{lstlisting}}
\newpage

\paragraph{Encoding} Format \textsf{LSU\_LSBO.Y} (Section~\ref{sec:format-LSU-LSBO.Y}), \verb|loadcode = 010|\\

\bitfields{ 1/P, 2/00, 2/-- --, 27/extend27, 1/1, 2/00, 2/-- --, 27/upper27, 1/1, 2/01, 3/010, 2/variant, 6/registerW, 2/00, 10/lower10, 6/registerZ }


\paragraph{Syntax} \texttt{lhz}\emph{variant} \emph{registerW} = \emph{extend27\_\discretionary{}{}{}upper27\_\discretionary{}{}{}lower10}[\emph{registerZ}]

\paragraph{Description} The effective address is computed by adding the \emph{extend27\_\discretionary{}{}{}upper27\_\discretionary{}{}{}lower10} to the \emph{registerZ}.
The \emph{registerW} is loaded from the memory half word at the effective address after zero extension.


\paragraph{Modifier(s)}
\noindent\begin{itemize}
\item variant (cf. variant in section \ref{par:modifier-variant} on page \pageref{par:modifier-variant})
\end{itemize}

\paragraph{Execution} {Reservation table \textsf{LSU\_AUXW.Y} (Section~\ref{sec:reservation-LSU-AUXW.Y})}
~
{\begin{lstlisting}
stage ID:
  new variant = _ZX_2(variant);
  new offset = _WX_64(extend27_upper27_lower10);
stage RR:
  new base = GPR[registerZ];
  new address = base + offset;
  new result1 = _ZX_16(MEM.load(address, 0x3, variant, registerW));
stage CR:
  GPR[registerW] = result1;
\end{lstlisting}}
\newpage

\paragraph{Encoding} Format \textsf{LSU\_LSBI} (Section~\ref{sec:format-LSU-LSBI}), \verb|loadcode = 010|\\

\bitfields{ 1/P, 2/01, 3/010, 2/variant, 6/registerW, 2/10, 3/111, 1/--, 6/registerY, 6/registerZ }


\paragraph{Syntax} \texttt{lhz}\emph{variant} \emph{registerW} = \emph{registerY}[\emph{registerZ}]

\paragraph{Description} The effective address is computed by adding the \emph{registerZ} to the \emph{registerY} and scaled by the \emph{variant}.
The \emph{registerW} is loaded from the memory half word at the effective address after zero extension.


\paragraph{Modifier(s)}
\noindent\begin{itemize}
\item variant (cf. variant in section \ref{par:modifier-variant} on page \pageref{par:modifier-variant})
\end{itemize}

\paragraph{Execution} {Reservation table \textsf{LSU\_AUXW} (Section~\ref{sec:reservation-LSU-AUXW})}
~
{\begin{lstlisting}
stage ID:
  new variant = _ZX_2(variant);
stage RR:
  new doscale = 0;
  new base = GPR[registerZ];
  new index = GPR[registerY];
  new scaling = doscale ? 2 : 1;
  new address = base + (index * scaling);
  new result1 = _ZX_16(MEM.load(address, 0x3, variant, registerW));
stage CR:
  GPR[registerW] = result1;
\end{lstlisting}}
\newpage
\subsubsection[{\small LO: Load Octuple Word}]{LO\hfill Load Octuple Word}
\label{instr:LO}\index{lo}
 
\paragraph{Encoding} Format \textsf{LSU\_LOBO} (Section~\ref{sec:format-LSU-LOBO}), \verb|loadcode = 111|\\

\bitfields{ 1/P, 2/01, 3/111, 2/variant, 4/registerN, 1/0, 1/1, 2/00, 10/signed10, 6/registerZ }


\paragraph{Syntax} \texttt{lo}\emph{variant} \emph{registerN} = \emph{signed10}[\emph{registerZ}]

\paragraph{Description} The effective address is computed by adding the \emph{signed10} to the \emph{registerZ}.
The \emph{registerN} is loaded from the memory octuple word at the effective address.


\paragraph{Modifier(s)}
\noindent\begin{itemize}
\item variant (cf. variant in section \ref{par:modifier-variant} on page \pageref{par:modifier-variant})
\end{itemize}

\paragraph{Execution} {Reservation table \textsf{LSU\_AUXW} (Section~\ref{sec:reservation-LSU-AUXW})}
~
{\begin{lstlisting}
stage ID:
  new variant = _ZX_2(variant);
  new offset = _SX_10(signed10);
stage RR:
  new bytemask = 0xFFFFFFFF;
  new base = GPR[registerZ];
  new address = base + offset;
  new result1 = MEM.load(address, 0xFFFFFFFF & bytemask, variant, registerN << 2);
stage CR:
  QGR[registerN] = result1;
\end{lstlisting}}
\newpage

\paragraph{Encoding} Format \textsf{LSU\_LOBO.X} (Section~\ref{sec:format-LSU-LOBO.X}), \verb|loadcode = 111|\\

\bitfields{ 1/P, 2/00, 2/-- --, 27/upper27, 1/1, 2/01, 3/111, 2/variant, 4/registerN, 1/0, 1/1, 2/00, 10/lower10, 6/registerZ }


\paragraph{Syntax} \texttt{lo}\emph{variant} \emph{registerN} = \emph{upper27\_\discretionary{}{}{}lower10}[\emph{registerZ}]

\paragraph{Description} The effective address is computed by adding the \emph{upper27\_\discretionary{}{}{}lower10} to the \emph{registerZ}.
The \emph{registerN} is loaded from the memory octuple word at the effective address.


\paragraph{Modifier(s)}
\noindent\begin{itemize}
\item variant (cf. variant in section \ref{par:modifier-variant} on page \pageref{par:modifier-variant})
\end{itemize}

\paragraph{Execution} {Reservation table \textsf{LSU\_AUXW.X} (Section~\ref{sec:reservation-LSU-AUXW.X})}
~
{\begin{lstlisting}
stage ID:
  new variant = _ZX_2(variant);
  new offset = _SX_37(upper27_lower10);
stage RR:
  new bytemask = 0xFFFFFFFF;
  new base = GPR[registerZ];
  new address = base + offset;
  new result1 = MEM.load(address, 0xFFFFFFFF & bytemask, variant, registerN << 2);
stage CR:
  QGR[registerN] = result1;
\end{lstlisting}}
\newpage

\paragraph{Encoding} Format \textsf{LSU\_LOBO.Y} (Section~\ref{sec:format-LSU-LOBO.Y}), \verb|loadcode = 111|\\

\bitfields{ 1/P, 2/00, 2/-- --, 27/extend27, 1/1, 2/00, 2/-- --, 27/upper27, 1/1, 2/01, 3/111, 2/variant, 4/registerN, 1/0, 1/1, 2/00, 10/lower10, 6/registerZ }


\paragraph{Syntax} \texttt{lo}\emph{variant} \emph{registerN} = \emph{extend27\_\discretionary{}{}{}upper27\_\discretionary{}{}{}lower10}[\emph{registerZ}]

\paragraph{Description} The effective address is computed by adding the \emph{extend27\_\discretionary{}{}{}upper27\_\discretionary{}{}{}lower10} to the \emph{registerZ}.
The \emph{registerN} is loaded from the memory octuple word at the effective address.


\paragraph{Modifier(s)}
\noindent\begin{itemize}
\item variant (cf. variant in section \ref{par:modifier-variant} on page \pageref{par:modifier-variant})
\end{itemize}

\paragraph{Execution} {Reservation table \textsf{LSU\_AUXW.Y} (Section~\ref{sec:reservation-LSU-AUXW.Y})}
~
{\begin{lstlisting}
stage ID:
  new variant = _ZX_2(variant);
  new offset = _WX_64(extend27_upper27_lower10);
stage RR:
  new bytemask = 0xFFFFFFFF;
  new base = GPR[registerZ];
  new address = base + offset;
  new result1 = MEM.load(address, 0xFFFFFFFF & bytemask, variant, registerN << 2);
stage CR:
  QGR[registerN] = result1;
\end{lstlisting}}
\newpage

\paragraph{Encoding} Format \textsf{LSU\_LOBI} (Section~\ref{sec:format-LSU-LOBI}), \verb|loadcode = 111|\\

\bitfields{ 1/P, 2/01, 3/111, 2/variant, 4/registerN, 1/0, 1/1, 2/10, 3/111, 1/--, 6/registerY, 6/registerZ }


\paragraph{Syntax} \texttt{lo}\emph{variant} \emph{registerN} = \emph{registerY}[\emph{registerZ}]

\paragraph{Description} The effective address is computed by adding the \emph{registerZ} to the \emph{registerY} and scaled by the \emph{variant}.
The \emph{registerN} is loaded from the memory octuple word at the effective address.


\paragraph{Modifier(s)}
\noindent\begin{itemize}
\item variant (cf. variant in section \ref{par:modifier-variant} on page \pageref{par:modifier-variant})
\end{itemize}

\paragraph{Execution} {Reservation table \textsf{LSU\_AUXW} (Section~\ref{sec:reservation-LSU-AUXW})}
~
{\begin{lstlisting}
stage ID:
  new variant = _ZX_2(variant);
stage RR:
  new bytemask = 0xFFFFFFFF;
  new doscale = 0;
  new base = GPR[registerZ];
  new index = GPR[registerY];
  new scaling = doscale ? 32 : 1;
  new address = base + (index * scaling);
  new result1 = MEM.load(address, 0xFFFFFFFF & bytemask, variant, registerN << 2);
stage CR:
  QGR[registerN] = result1;
\end{lstlisting}}
\newpage
\subsubsection[{\small LQ: Load Quadruple Word}]{LQ\hfill Load Quadruple Word}
\label{instr:LQ}\index{lq}
 
\paragraph{Encoding} Format \textsf{LSU\_LQBO} (Section~\ref{sec:format-LSU-LQBO}), \verb|loadcode = 111|\\

\bitfields{ 1/P, 2/01, 3/111, 2/variant, 5/registerM, 1/0, 2/00, 10/signed10, 6/registerZ }


\paragraph{Syntax} \texttt{lq}\emph{variant} \emph{registerM} = \emph{signed10}[\emph{registerZ}]

\paragraph{Description} The effective address is computed by adding the \emph{signed10} to the \emph{registerZ}.
The \emph{registerM} is loaded from the memory quadruple word at the effective address.


\paragraph{Modifier(s)}
\noindent\begin{itemize}
\item variant (cf. variant in section \ref{par:modifier-variant} on page \pageref{par:modifier-variant})
\end{itemize}

\paragraph{Execution} {Reservation table \textsf{LSU\_AUXW} (Section~\ref{sec:reservation-LSU-AUXW})}
~
{\begin{lstlisting}
stage ID:
  new variant = _ZX_2(variant);
  new offset = _SX_10(signed10);
stage RR:
  new base = GPR[registerZ];
  new address = base + offset;
  new result1 = MEM.load(address, 0xFFFF, variant, registerM << 1);
stage CR:
  PGR[registerM] = result1;
\end{lstlisting}}
\newpage

\paragraph{Encoding} Format \textsf{LSU\_LQBO.X} (Section~\ref{sec:format-LSU-LQBO.X}), \verb|loadcode = 111|\\

\bitfields{ 1/P, 2/00, 2/-- --, 27/upper27, 1/1, 2/01, 3/111, 2/variant, 5/registerM, 1/0, 2/00, 10/lower10, 6/registerZ }


\paragraph{Syntax} \texttt{lq}\emph{variant} \emph{registerM} = \emph{upper27\_\discretionary{}{}{}lower10}[\emph{registerZ}]

\paragraph{Description} The effective address is computed by adding the \emph{upper27\_\discretionary{}{}{}lower10} to the \emph{registerZ}.
The \emph{registerM} is loaded from the memory quadruple word at the effective address.


\paragraph{Modifier(s)}
\noindent\begin{itemize}
\item variant (cf. variant in section \ref{par:modifier-variant} on page \pageref{par:modifier-variant})
\end{itemize}

\paragraph{Execution} {Reservation table \textsf{LSU\_AUXW.X} (Section~\ref{sec:reservation-LSU-AUXW.X})}
~
{\begin{lstlisting}
stage ID:
  new variant = _ZX_2(variant);
  new offset = _SX_37(upper27_lower10);
stage RR:
  new base = GPR[registerZ];
  new address = base + offset;
  new result1 = MEM.load(address, 0xFFFF, variant, registerM << 1);
stage CR:
  PGR[registerM] = result1;
\end{lstlisting}}
\newpage

\paragraph{Encoding} Format \textsf{LSU\_LQBO.Y} (Section~\ref{sec:format-LSU-LQBO.Y}), \verb|loadcode = 111|\\

\bitfields{ 1/P, 2/00, 2/-- --, 27/extend27, 1/1, 2/00, 2/-- --, 27/upper27, 1/1, 2/01, 3/111, 2/variant, 5/registerM, 1/0, 2/00, 10/lower10, 6/registerZ }


\paragraph{Syntax} \texttt{lq}\emph{variant} \emph{registerM} = \emph{extend27\_\discretionary{}{}{}upper27\_\discretionary{}{}{}lower10}[\emph{registerZ}]

\paragraph{Description} The effective address is computed by adding the \emph{extend27\_\discretionary{}{}{}upper27\_\discretionary{}{}{}lower10} to the \emph{registerZ}.
The \emph{registerM} is loaded from the memory quadruple word at the effective address.


\paragraph{Modifier(s)}
\noindent\begin{itemize}
\item variant (cf. variant in section \ref{par:modifier-variant} on page \pageref{par:modifier-variant})
\end{itemize}

\paragraph{Execution} {Reservation table \textsf{LSU\_AUXW.Y} (Section~\ref{sec:reservation-LSU-AUXW.Y})}
~
{\begin{lstlisting}
stage ID:
  new variant = _ZX_2(variant);
  new offset = _WX_64(extend27_upper27_lower10);
stage RR:
  new base = GPR[registerZ];
  new address = base + offset;
  new result1 = MEM.load(address, 0xFFFF, variant, registerM << 1);
stage CR:
  PGR[registerM] = result1;
\end{lstlisting}}
\newpage

\paragraph{Encoding} Format \textsf{LSU\_LQBI} (Section~\ref{sec:format-LSU-LQBI}), \verb|loadcode = 111|\\

\bitfields{ 1/P, 2/01, 3/111, 2/variant, 5/registerM, 1/0, 2/10, 3/111, 1/--, 6/registerY, 6/registerZ }


\paragraph{Syntax} \texttt{lq}\emph{variant} \emph{registerM} = \emph{registerY}[\emph{registerZ}]

\paragraph{Description} The effective address is computed by adding the \emph{registerZ} to the \emph{registerY} and scaled by the \emph{variant}.
The \emph{registerM} is loaded from the memory quadruple word at the effective address.


\paragraph{Modifier(s)}
\noindent\begin{itemize}
\item variant (cf. variant in section \ref{par:modifier-variant} on page \pageref{par:modifier-variant})
\end{itemize}

\paragraph{Execution} {Reservation table \textsf{LSU\_AUXW} (Section~\ref{sec:reservation-LSU-AUXW})}
~
{\begin{lstlisting}
stage ID:
  new variant = _ZX_2(variant);
stage RR:
  new doscale = 0;
  new base = GPR[registerZ];
  new index = GPR[registerY];
  new scaling = doscale ? 16 : 1;
  new address = base + (index * scaling);
  new result1 = MEM.load(address, 0xFFFF, variant, registerM << 1);
stage CR:
  PGR[registerM] = result1;
\end{lstlisting}}
\newpage
\subsubsection[{\small LWS: Load Word Sign Extended}]{LWS\hfill Load Word Sign Extended}
\label{instr:LWS}\index{lws}
 
\paragraph{Encoding} Format \textsf{LSU\_LSBO} (Section~\ref{sec:format-LSU-LSBO}), \verb|loadcode = 101|\\

\bitfields{ 1/P, 2/01, 3/101, 2/variant, 6/registerW, 2/00, 10/signed10, 6/registerZ }


\paragraph{Syntax} \texttt{lws}\emph{variant} \emph{registerW} = \emph{signed10}[\emph{registerZ}]

\paragraph{Description} The effective address is computed by adding the \emph{signed10} to the \emph{registerZ}.
The \emph{registerW} is loaded from the memory word at the effective address after sign extension.


\paragraph{Modifier(s)}
\noindent\begin{itemize}
\item variant (cf. variant in section \ref{par:modifier-variant} on page \pageref{par:modifier-variant})
\end{itemize}

\paragraph{Execution} {Reservation table \textsf{LSU\_AUXW} (Section~\ref{sec:reservation-LSU-AUXW})}
~
{\begin{lstlisting}
stage ID:
  new variant = _ZX_2(variant);
  new offset = _SX_10(signed10);
stage RR:
  new base = GPR[registerZ];
  new address = base + offset;
  new result1 = _SX_32(MEM.load(address, 0xF, variant, registerW));
stage CR:
  GPR[registerW] = result1;
\end{lstlisting}}
\newpage

\paragraph{Encoding} Format \textsf{LSU\_LSBO.X} (Section~\ref{sec:format-LSU-LSBO.X}), \verb|loadcode = 101|\\

\bitfields{ 1/P, 2/00, 2/-- --, 27/upper27, 1/1, 2/01, 3/101, 2/variant, 6/registerW, 2/00, 10/lower10, 6/registerZ }


\paragraph{Syntax} \texttt{lws}\emph{variant} \emph{registerW} = \emph{upper27\_\discretionary{}{}{}lower10}[\emph{registerZ}]

\paragraph{Description} The effective address is computed by adding the \emph{upper27\_\discretionary{}{}{}lower10} to the \emph{registerZ}.
The \emph{registerW} is loaded from the memory word at the effective address after sign extension.


\paragraph{Modifier(s)}
\noindent\begin{itemize}
\item variant (cf. variant in section \ref{par:modifier-variant} on page \pageref{par:modifier-variant})
\end{itemize}

\paragraph{Execution} {Reservation table \textsf{LSU\_AUXW.X} (Section~\ref{sec:reservation-LSU-AUXW.X})}
~
{\begin{lstlisting}
stage ID:
  new variant = _ZX_2(variant);
  new offset = _SX_37(upper27_lower10);
stage RR:
  new base = GPR[registerZ];
  new address = base + offset;
  new result1 = _SX_32(MEM.load(address, 0xF, variant, registerW));
stage CR:
  GPR[registerW] = result1;
\end{lstlisting}}
\newpage

\paragraph{Encoding} Format \textsf{LSU\_LSBO.Y} (Section~\ref{sec:format-LSU-LSBO.Y}), \verb|loadcode = 101|\\

\bitfields{ 1/P, 2/00, 2/-- --, 27/extend27, 1/1, 2/00, 2/-- --, 27/upper27, 1/1, 2/01, 3/101, 2/variant, 6/registerW, 2/00, 10/lower10, 6/registerZ }


\paragraph{Syntax} \texttt{lws}\emph{variant} \emph{registerW} = \emph{extend27\_\discretionary{}{}{}upper27\_\discretionary{}{}{}lower10}[\emph{registerZ}]

\paragraph{Description} The effective address is computed by adding the \emph{extend27\_\discretionary{}{}{}upper27\_\discretionary{}{}{}lower10} to the \emph{registerZ}.
The \emph{registerW} is loaded from the memory word at the effective address after sign extension.


\paragraph{Modifier(s)}
\noindent\begin{itemize}
\item variant (cf. variant in section \ref{par:modifier-variant} on page \pageref{par:modifier-variant})
\end{itemize}

\paragraph{Execution} {Reservation table \textsf{LSU\_AUXW.Y} (Section~\ref{sec:reservation-LSU-AUXW.Y})}
~
{\begin{lstlisting}
stage ID:
  new variant = _ZX_2(variant);
  new offset = _WX_64(extend27_upper27_lower10);
stage RR:
  new base = GPR[registerZ];
  new address = base + offset;
  new result1 = _SX_32(MEM.load(address, 0xF, variant, registerW));
stage CR:
  GPR[registerW] = result1;
\end{lstlisting}}
\newpage

\paragraph{Encoding} Format \textsf{LSU\_LSBI} (Section~\ref{sec:format-LSU-LSBI}), \verb|loadcode = 101|\\

\bitfields{ 1/P, 2/01, 3/101, 2/variant, 6/registerW, 2/10, 3/111, 1/--, 6/registerY, 6/registerZ }


\paragraph{Syntax} \texttt{lws}\emph{variant} \emph{registerW} = \emph{registerY}[\emph{registerZ}]

\paragraph{Description} The effective address is computed by adding the \emph{registerZ} to the \emph{registerY} and scaled by the \emph{variant}.
The \emph{registerW} is loaded from the memory word at the effective address after sign extension.


\paragraph{Modifier(s)}
\noindent\begin{itemize}
\item variant (cf. variant in section \ref{par:modifier-variant} on page \pageref{par:modifier-variant})
\end{itemize}

\paragraph{Execution} {Reservation table \textsf{LSU\_AUXW} (Section~\ref{sec:reservation-LSU-AUXW})}
~
{\begin{lstlisting}
stage ID:
  new variant = _ZX_2(variant);
stage RR:
  new doscale = 0;
  new base = GPR[registerZ];
  new index = GPR[registerY];
  new scaling = doscale ? 4 : 1;
  new address = base + (index * scaling);
  new result1 = _SX_32(MEM.load(address, 0xF, variant, registerW));
stage CR:
  GPR[registerW] = result1;
\end{lstlisting}}
\newpage
\subsubsection[{\small LWZ: Load Word Zero Extended}]{LWZ\hfill Load Word Zero Extended}
\label{instr:LWZ}\index{lwz}
 
\paragraph{Encoding} Format \textsf{LSU\_LSBO} (Section~\ref{sec:format-LSU-LSBO}), \verb|loadcode = 100|\\

\bitfields{ 1/P, 2/01, 3/100, 2/variant, 6/registerW, 2/00, 10/signed10, 6/registerZ }


\paragraph{Syntax} \texttt{lwz}\emph{variant} \emph{registerW} = \emph{signed10}[\emph{registerZ}]

\paragraph{Description} The effective address is computed by adding the \emph{signed10} to the \emph{registerZ}.
The \emph{registerW} is loaded from the memory word at the effective address after zero extension.


\paragraph{Modifier(s)}
\noindent\begin{itemize}
\item variant (cf. variant in section \ref{par:modifier-variant} on page \pageref{par:modifier-variant})
\end{itemize}

\paragraph{Execution} {Reservation table \textsf{LSU\_AUXW} (Section~\ref{sec:reservation-LSU-AUXW})}
~
{\begin{lstlisting}
stage ID:
  new variant = _ZX_2(variant);
  new offset = _SX_10(signed10);
stage RR:
  new base = GPR[registerZ];
  new address = base + offset;
  new result1 = _ZX_32(MEM.load(address, 0xF, variant, registerW));
stage CR:
  GPR[registerW] = result1;
\end{lstlisting}}
\newpage

\paragraph{Encoding} Format \textsf{LSU\_LSBO.X} (Section~\ref{sec:format-LSU-LSBO.X}), \verb|loadcode = 100|\\

\bitfields{ 1/P, 2/00, 2/-- --, 27/upper27, 1/1, 2/01, 3/100, 2/variant, 6/registerW, 2/00, 10/lower10, 6/registerZ }


\paragraph{Syntax} \texttt{lwz}\emph{variant} \emph{registerW} = \emph{upper27\_\discretionary{}{}{}lower10}[\emph{registerZ}]

\paragraph{Description} The effective address is computed by adding the \emph{upper27\_\discretionary{}{}{}lower10} to the \emph{registerZ}.
The \emph{registerW} is loaded from the memory word at the effective address after zero extension.


\paragraph{Modifier(s)}
\noindent\begin{itemize}
\item variant (cf. variant in section \ref{par:modifier-variant} on page \pageref{par:modifier-variant})
\end{itemize}

\paragraph{Execution} {Reservation table \textsf{LSU\_AUXW.X} (Section~\ref{sec:reservation-LSU-AUXW.X})}
~
{\begin{lstlisting}
stage ID:
  new variant = _ZX_2(variant);
  new offset = _SX_37(upper27_lower10);
stage RR:
  new base = GPR[registerZ];
  new address = base + offset;
  new result1 = _ZX_32(MEM.load(address, 0xF, variant, registerW));
stage CR:
  GPR[registerW] = result1;
\end{lstlisting}}
\newpage

\paragraph{Encoding} Format \textsf{LSU\_LSBO.Y} (Section~\ref{sec:format-LSU-LSBO.Y}), \verb|loadcode = 100|\\

\bitfields{ 1/P, 2/00, 2/-- --, 27/extend27, 1/1, 2/00, 2/-- --, 27/upper27, 1/1, 2/01, 3/100, 2/variant, 6/registerW, 2/00, 10/lower10, 6/registerZ }


\paragraph{Syntax} \texttt{lwz}\emph{variant} \emph{registerW} = \emph{extend27\_\discretionary{}{}{}upper27\_\discretionary{}{}{}lower10}[\emph{registerZ}]

\paragraph{Description} The effective address is computed by adding the \emph{extend27\_\discretionary{}{}{}upper27\_\discretionary{}{}{}lower10} to the \emph{registerZ}.
The \emph{registerW} is loaded from the memory word at the effective address after zero extension.


\paragraph{Modifier(s)}
\noindent\begin{itemize}
\item variant (cf. variant in section \ref{par:modifier-variant} on page \pageref{par:modifier-variant})
\end{itemize}

\paragraph{Execution} {Reservation table \textsf{LSU\_AUXW.Y} (Section~\ref{sec:reservation-LSU-AUXW.Y})}
~
{\begin{lstlisting}
stage ID:
  new variant = _ZX_2(variant);
  new offset = _WX_64(extend27_upper27_lower10);
stage RR:
  new base = GPR[registerZ];
  new address = base + offset;
  new result1 = _ZX_32(MEM.load(address, 0xF, variant, registerW));
stage CR:
  GPR[registerW] = result1;
\end{lstlisting}}
\newpage

\paragraph{Encoding} Format \textsf{LSU\_LSBI} (Section~\ref{sec:format-LSU-LSBI}), \verb|loadcode = 100|\\

\bitfields{ 1/P, 2/01, 3/100, 2/variant, 6/registerW, 2/10, 3/111, 1/--, 6/registerY, 6/registerZ }


\paragraph{Syntax} \texttt{lwz}\emph{variant} \emph{registerW} = \emph{registerY}[\emph{registerZ}]

\paragraph{Description} The effective address is computed by adding the \emph{registerZ} to the \emph{registerY} and scaled by the \emph{variant}.
The \emph{registerW} is loaded from the memory word at the effective address after zero extension.


\paragraph{Modifier(s)}
\noindent\begin{itemize}
\item variant (cf. variant in section \ref{par:modifier-variant} on page \pageref{par:modifier-variant})
\end{itemize}

\paragraph{Execution} {Reservation table \textsf{LSU\_AUXW} (Section~\ref{sec:reservation-LSU-AUXW})}
~
{\begin{lstlisting}
stage ID:
  new variant = _ZX_2(variant);
stage RR:
  new doscale = 0;
  new base = GPR[registerZ];
  new index = GPR[registerY];
  new scaling = doscale ? 4 : 1;
  new address = base + (index * scaling);
  new result1 = _ZX_32(MEM.load(address, 0xF, variant, registerW));
stage CR:
  GPR[registerW] = result1;
\end{lstlisting}}
\newpage
\subsubsection[{\small SB: Store Byte}]{SB\hfill Store Byte}
\label{instr:SB}\index{sb}
 
\paragraph{Encoding} Format \textsf{LSU\_SSBO} (Section~\ref{sec:format-LSU-SSBO}), \verb|lsucode0 = 10000|\\

\bitfields{ 1/P, 2/01, 5/10000, 6/registerT, 2/01, 10/signed10, 6/registerZ }


\paragraph{Syntax} \texttt{sb} \emph{signed10}[\emph{registerZ}] = \emph{registerT}

\paragraph{Description} The effective address is computed by adding the \emph{signed10} to the \emph{registerZ}.
The \emph{registerT} is stored into the memory byte at the effective address.


\paragraph{Execution} {Reservation table \textsf{LSU\_MEMW\_AUXR} (Section~\ref{sec:reservation-LSU-MEMW-AUXR})}
~
{\begin{lstlisting}
stage ID:
  new offset = _SX_10(signed10);
stage RR:
  new base = GPR[registerZ];
  new address = base + offset;
stage E1:
  new argument3 = GPR[registerT];
stage E3:
  MEM.store(address, 0x1, 0, argument3, registerT);
\end{lstlisting}}
\newpage

\paragraph{Encoding} Format \textsf{LSU\_SSBO.X} (Section~\ref{sec:format-LSU-SSBO.X}), \verb|lsucode0 = 10000|\\

\bitfields{ 1/P, 2/00, 2/-- --, 27/upper27, 1/1, 2/01, 5/10000, 6/registerT, 2/01, 10/lower10, 6/registerZ }


\paragraph{Syntax} \texttt{sb} \emph{upper27\_\discretionary{}{}{}lower10}[\emph{registerZ}] = \emph{registerT}

\paragraph{Description} The effective address is computed by adding the \emph{upper27\_\discretionary{}{}{}lower10} to the \emph{registerZ}.
The \emph{registerT} is stored into the memory byte at the effective address.


\paragraph{Execution} {Reservation table \textsf{LSU\_MEMW\_AUXR.X} (Section~\ref{sec:reservation-LSU-MEMW-AUXR.X})}
~
{\begin{lstlisting}
stage ID:
  new offset = _SX_37(upper27_lower10);
stage RR:
  new base = GPR[registerZ];
  new address = base + offset;
stage E1:
  new argument3 = GPR[registerT];
stage E3:
  MEM.store(address, 0x1, 0, argument3, registerT);
\end{lstlisting}}
\newpage

\paragraph{Encoding} Format \textsf{LSU\_SSBO.Y} (Section~\ref{sec:format-LSU-SSBO.Y}), \verb|lsucode0 = 10000|\\

\bitfields{ 1/P, 2/00, 2/-- --, 27/extend27, 1/1, 2/00, 2/-- --, 27/upper27, 1/1, 2/01, 5/10000, 6/registerT, 2/01, 10/lower10, 6/registerZ }


\paragraph{Syntax} \texttt{sb} \emph{extend27\_\discretionary{}{}{}upper27\_\discretionary{}{}{}lower10}[\emph{registerZ}] = \emph{registerT}

\paragraph{Description} The effective address is computed by adding the \emph{extend27\_\discretionary{}{}{}upper27\_\discretionary{}{}{}lower10} to the \emph{registerZ}.
The \emph{registerT} is stored into the memory byte at the effective address.


\paragraph{Execution} {Reservation table \textsf{LSU\_MEMW\_AUXR.Y} (Section~\ref{sec:reservation-LSU-MEMW-AUXR.Y})}
~
{\begin{lstlisting}
stage ID:
  new offset = _WX_64(extend27_upper27_lower10);
stage RR:
  new base = GPR[registerZ];
  new address = base + offset;
stage E1:
  new argument3 = GPR[registerT];
stage E3:
  MEM.store(address, 0x1, 0, argument3, registerT);
\end{lstlisting}}
\newpage

\paragraph{Encoding} Format \textsf{LSU\_SSBI} (Section~\ref{sec:format-LSU-SSBI}), \verb|lsucode0 = 10000|\\

\bitfields{ 1/P, 2/01, 5/10000, 6/registerT, 2/11, 3/111, 1/--, 6/registerY, 6/registerZ }


\paragraph{Syntax} \texttt{sb} \emph{registerY}[\emph{registerZ}] = \emph{registerT}

\paragraph{Description} The effective address is computed by adding the \emph{registerZ} to the \emph{registerY} and scaled by the \emph{}.
The \emph{registerT} is stored into the memory byte at the effective address.


\paragraph{Execution} {Reservation table \textsf{LSU\_MEMW\_AUXR} (Section~\ref{sec:reservation-LSU-MEMW-AUXR})}
~
{\begin{lstlisting}
stage RR:
  new doscale = 0;
  new index = GPR[registerY];
  new base = GPR[registerZ];
  new scaling = doscale ? 1 : 1;
  new address = base + (index * scaling);
stage E1:
  new argument3 = GPR[registerT];
stage E3:
  MEM.store(address, 0x1, 0, argument3, registerT);
\end{lstlisting}}
\newpage
\subsubsection[{\small SD: Store Double Word}]{SD\hfill Store Double Word}
\label{instr:SD}\index{sd}
 
\paragraph{Encoding} Format \textsf{LSU\_SSBO} (Section~\ref{sec:format-LSU-SSBO}), \verb|lsucode0 = 10011|\\

\bitfields{ 1/P, 2/01, 5/10011, 6/registerT, 2/01, 10/signed10, 6/registerZ }


\paragraph{Syntax} \texttt{sd} \emph{signed10}[\emph{registerZ}] = \emph{registerT}

\paragraph{Description} The effective address is computed by adding the \emph{signed10} to the \emph{registerZ}.
The \emph{registerT} is stored into the memory double word at the effective address.


\paragraph{Execution} {Reservation table \textsf{LSU\_MEMW\_AUXR} (Section~\ref{sec:reservation-LSU-MEMW-AUXR})}
~
{\begin{lstlisting}
stage ID:
  new offset = _SX_10(signed10);
stage RR:
  new base = GPR[registerZ];
  new address = base + offset;
stage E1:
  new argument3 = GPR[registerT];
stage E3:
  MEM.store(address, 0xFF, 0, argument3, registerT);
\end{lstlisting}}
\newpage

\paragraph{Encoding} Format \textsf{LSU\_SSBO.X} (Section~\ref{sec:format-LSU-SSBO.X}), \verb|lsucode0 = 10011|\\

\bitfields{ 1/P, 2/00, 2/-- --, 27/upper27, 1/1, 2/01, 5/10011, 6/registerT, 2/01, 10/lower10, 6/registerZ }


\paragraph{Syntax} \texttt{sd} \emph{upper27\_\discretionary{}{}{}lower10}[\emph{registerZ}] = \emph{registerT}

\paragraph{Description} The effective address is computed by adding the \emph{upper27\_\discretionary{}{}{}lower10} to the \emph{registerZ}.
The \emph{registerT} is stored into the memory double word at the effective address.


\paragraph{Execution} {Reservation table \textsf{LSU\_MEMW\_AUXR.X} (Section~\ref{sec:reservation-LSU-MEMW-AUXR.X})}
~
{\begin{lstlisting}
stage ID:
  new offset = _SX_37(upper27_lower10);
stage RR:
  new base = GPR[registerZ];
  new address = base + offset;
stage E1:
  new argument3 = GPR[registerT];
stage E3:
  MEM.store(address, 0xFF, 0, argument3, registerT);
\end{lstlisting}}
\newpage

\paragraph{Encoding} Format \textsf{LSU\_SSBO.Y} (Section~\ref{sec:format-LSU-SSBO.Y}), \verb|lsucode0 = 10011|\\

\bitfields{ 1/P, 2/00, 2/-- --, 27/extend27, 1/1, 2/00, 2/-- --, 27/upper27, 1/1, 2/01, 5/10011, 6/registerT, 2/01, 10/lower10, 6/registerZ }


\paragraph{Syntax} \texttt{sd} \emph{extend27\_\discretionary{}{}{}upper27\_\discretionary{}{}{}lower10}[\emph{registerZ}] = \emph{registerT}

\paragraph{Description} The effective address is computed by adding the \emph{extend27\_\discretionary{}{}{}upper27\_\discretionary{}{}{}lower10} to the \emph{registerZ}.
The \emph{registerT} is stored into the memory double word at the effective address.


\paragraph{Execution} {Reservation table \textsf{LSU\_MEMW\_AUXR.Y} (Section~\ref{sec:reservation-LSU-MEMW-AUXR.Y})}
~
{\begin{lstlisting}
stage ID:
  new offset = _WX_64(extend27_upper27_lower10);
stage RR:
  new base = GPR[registerZ];
  new address = base + offset;
stage E1:
  new argument3 = GPR[registerT];
stage E3:
  MEM.store(address, 0xFF, 0, argument3, registerT);
\end{lstlisting}}
\newpage

\paragraph{Encoding} Format \textsf{LSU\_SSBI} (Section~\ref{sec:format-LSU-SSBI}), \verb|lsucode0 = 10011|\\

\bitfields{ 1/P, 2/01, 5/10011, 6/registerT, 2/11, 3/111, 1/--, 6/registerY, 6/registerZ }


\paragraph{Syntax} \texttt{sd} \emph{registerY}[\emph{registerZ}] = \emph{registerT}

\paragraph{Description} The effective address is computed by adding the \emph{registerZ} to the \emph{registerY} and scaled by the \emph{}.
The \emph{registerT} is stored into the memory double word at the effective address.


\paragraph{Execution} {Reservation table \textsf{LSU\_MEMW\_AUXR} (Section~\ref{sec:reservation-LSU-MEMW-AUXR})}
~
{\begin{lstlisting}
stage RR:
  new doscale = 0;
  new index = GPR[registerY];
  new base = GPR[registerZ];
  new scaling = doscale ? 8 : 1;
  new address = base + (index * scaling);
stage E1:
  new argument3 = GPR[registerT];
stage E3:
  MEM.store(address, 0xFF, 0, argument3, registerT);
\end{lstlisting}}
\newpage
\subsubsection[{\small SH: Store Half Word}]{SH\hfill Store Half Word}
\label{instr:SH}\index{sh}
 
\paragraph{Encoding} Format \textsf{LSU\_SSBO} (Section~\ref{sec:format-LSU-SSBO}), \verb|lsucode0 = 10001|\\

\bitfields{ 1/P, 2/01, 5/10001, 6/registerT, 2/01, 10/signed10, 6/registerZ }


\paragraph{Syntax} \texttt{sh} \emph{signed10}[\emph{registerZ}] = \emph{registerT}

\paragraph{Description} The effective address is computed by adding the \emph{signed10} to the \emph{registerZ}.
The \emph{registerT} is stored into the memory half word at the effective address.


\paragraph{Execution} {Reservation table \textsf{LSU\_MEMW\_AUXR} (Section~\ref{sec:reservation-LSU-MEMW-AUXR})}
~
{\begin{lstlisting}
stage ID:
  new offset = _SX_10(signed10);
stage RR:
  new base = GPR[registerZ];
  new address = base + offset;
stage E1:
  new argument3 = GPR[registerT];
stage E3:
  MEM.store(address, 0x3, 0, argument3, registerT);
\end{lstlisting}}
\newpage

\paragraph{Encoding} Format \textsf{LSU\_SSBO.X} (Section~\ref{sec:format-LSU-SSBO.X}), \verb|lsucode0 = 10001|\\

\bitfields{ 1/P, 2/00, 2/-- --, 27/upper27, 1/1, 2/01, 5/10001, 6/registerT, 2/01, 10/lower10, 6/registerZ }


\paragraph{Syntax} \texttt{sh} \emph{upper27\_\discretionary{}{}{}lower10}[\emph{registerZ}] = \emph{registerT}

\paragraph{Description} The effective address is computed by adding the \emph{upper27\_\discretionary{}{}{}lower10} to the \emph{registerZ}.
The \emph{registerT} is stored into the memory half word at the effective address.


\paragraph{Execution} {Reservation table \textsf{LSU\_MEMW\_AUXR.X} (Section~\ref{sec:reservation-LSU-MEMW-AUXR.X})}
~
{\begin{lstlisting}
stage ID:
  new offset = _SX_37(upper27_lower10);
stage RR:
  new base = GPR[registerZ];
  new address = base + offset;
stage E1:
  new argument3 = GPR[registerT];
stage E3:
  MEM.store(address, 0x3, 0, argument3, registerT);
\end{lstlisting}}
\newpage

\paragraph{Encoding} Format \textsf{LSU\_SSBO.Y} (Section~\ref{sec:format-LSU-SSBO.Y}), \verb|lsucode0 = 10001|\\

\bitfields{ 1/P, 2/00, 2/-- --, 27/extend27, 1/1, 2/00, 2/-- --, 27/upper27, 1/1, 2/01, 5/10001, 6/registerT, 2/01, 10/lower10, 6/registerZ }


\paragraph{Syntax} \texttt{sh} \emph{extend27\_\discretionary{}{}{}upper27\_\discretionary{}{}{}lower10}[\emph{registerZ}] = \emph{registerT}

\paragraph{Description} The effective address is computed by adding the \emph{extend27\_\discretionary{}{}{}upper27\_\discretionary{}{}{}lower10} to the \emph{registerZ}.
The \emph{registerT} is stored into the memory half word at the effective address.


\paragraph{Execution} {Reservation table \textsf{LSU\_MEMW\_AUXR.Y} (Section~\ref{sec:reservation-LSU-MEMW-AUXR.Y})}
~
{\begin{lstlisting}
stage ID:
  new offset = _WX_64(extend27_upper27_lower10);
stage RR:
  new base = GPR[registerZ];
  new address = base + offset;
stage E1:
  new argument3 = GPR[registerT];
stage E3:
  MEM.store(address, 0x3, 0, argument3, registerT);
\end{lstlisting}}
\newpage

\paragraph{Encoding} Format \textsf{LSU\_SSBI} (Section~\ref{sec:format-LSU-SSBI}), \verb|lsucode0 = 10001|\\

\bitfields{ 1/P, 2/01, 5/10001, 6/registerT, 2/11, 3/111, 1/--, 6/registerY, 6/registerZ }


\paragraph{Syntax} \texttt{sh} \emph{registerY}[\emph{registerZ}] = \emph{registerT}

\paragraph{Description} The effective address is computed by adding the \emph{registerZ} to the \emph{registerY} and scaled by the \emph{}.
The \emph{registerT} is stored into the memory half word at the effective address.


\paragraph{Execution} {Reservation table \textsf{LSU\_MEMW\_AUXR} (Section~\ref{sec:reservation-LSU-MEMW-AUXR})}
~
{\begin{lstlisting}
stage RR:
  new doscale = 0;
  new index = GPR[registerY];
  new base = GPR[registerZ];
  new scaling = doscale ? 2 : 1;
  new address = base + (index * scaling);
stage E1:
  new argument3 = GPR[registerT];
stage E3:
  MEM.store(address, 0x3, 0, argument3, registerT);
\end{lstlisting}}
\newpage
\subsubsection[{\small SO: Store Octuple Word}]{SO\hfill Store Octuple Word}
\label{instr:SO}\index{so}
 
\paragraph{Encoding} Format \textsf{LSU\_SOBO} (Section~\ref{sec:format-LSU-SOBO}), \verb|lsucode0 = 10100|\\

\bitfields{ 1/P, 2/01, 5/10100, 4/registerV, 1/0, 1/1, 2/01, 10/signed10, 6/registerZ }


\paragraph{Syntax} \texttt{so} \emph{signed10}[\emph{registerZ}] = \emph{registerV}

\paragraph{Description} The effective address is computed by adding the \emph{signed10} to the \emph{registerZ}.
The \emph{registerV} is stored into the memory octuple word at the effective address.


\paragraph{Execution} {Reservation table \textsf{LSU\_MEMW\_AUXR} (Section~\ref{sec:reservation-LSU-MEMW-AUXR})}
~
{\begin{lstlisting}
stage ID:
  new offset = _SX_10(signed10);
stage RR:
  new bytemask = 0xFFFFFFFF;
  new base = GPR[registerZ];
  new address = base + offset;
stage E1:
  new argument3 = QGR[registerV];
stage E3:
  MEM.store(address, 0xFFFFFFFF & bytemask, 0, argument3, registerV << 2);
\end{lstlisting}}
\newpage

\paragraph{Encoding} Format \textsf{LSU\_SOBO.X} (Section~\ref{sec:format-LSU-SOBO.X}), \verb|lsucode0 = 10100|\\

\bitfields{ 1/P, 2/00, 2/-- --, 27/upper27, 1/1, 2/01, 5/10100, 4/registerV, 1/0, 1/1, 2/01, 10/lower10, 6/registerZ }


\paragraph{Syntax} \texttt{so} \emph{upper27\_\discretionary{}{}{}lower10}[\emph{registerZ}] = \emph{registerV}

\paragraph{Description} The effective address is computed by adding the \emph{upper27\_\discretionary{}{}{}lower10} to the \emph{registerZ}.
The \emph{registerV} is stored into the memory octuple word at the effective address.


\paragraph{Execution} {Reservation table \textsf{LSU\_MEMW\_AUXR.X} (Section~\ref{sec:reservation-LSU-MEMW-AUXR.X})}
~
{\begin{lstlisting}
stage ID:
  new offset = _SX_37(upper27_lower10);
stage RR:
  new bytemask = 0xFFFFFFFF;
  new base = GPR[registerZ];
  new address = base + offset;
stage E1:
  new argument3 = QGR[registerV];
stage E3:
  MEM.store(address, 0xFFFFFFFF & bytemask, 0, argument3, registerV << 2);
\end{lstlisting}}
\newpage

\paragraph{Encoding} Format \textsf{LSU\_SOBO.Y} (Section~\ref{sec:format-LSU-SOBO.Y}), \verb|lsucode0 = 10100|\\

\bitfields{ 1/P, 2/00, 2/-- --, 27/extend27, 1/1, 2/00, 2/-- --, 27/upper27, 1/1, 2/01, 5/10100, 4/registerV, 1/0, 1/1, 2/01, 10/lower10, 6/registerZ }


\paragraph{Syntax} \texttt{so} \emph{extend27\_\discretionary{}{}{}upper27\_\discretionary{}{}{}lower10}[\emph{registerZ}] = \emph{registerV}

\paragraph{Description} The effective address is computed by adding the \emph{extend27\_\discretionary{}{}{}upper27\_\discretionary{}{}{}lower10} to the \emph{registerZ}.
The \emph{registerV} is stored into the memory octuple word at the effective address.


\paragraph{Execution} {Reservation table \textsf{LSU\_MEMW\_AUXR.Y} (Section~\ref{sec:reservation-LSU-MEMW-AUXR.Y})}
~
{\begin{lstlisting}
stage ID:
  new offset = _WX_64(extend27_upper27_lower10);
stage RR:
  new bytemask = 0xFFFFFFFF;
  new base = GPR[registerZ];
  new address = base + offset;
stage E1:
  new argument3 = QGR[registerV];
stage E3:
  MEM.store(address, 0xFFFFFFFF & bytemask, 0, argument3, registerV << 2);
\end{lstlisting}}
\newpage

\paragraph{Encoding} Format \textsf{LSU\_SOBI} (Section~\ref{sec:format-LSU-SOBI}), \verb|lsucode0 = 10100|\\

\bitfields{ 1/P, 2/01, 5/10100, 4/registerV, 1/0, 1/1, 2/11, 3/111, 1/--, 6/registerY, 6/registerZ }


\paragraph{Syntax} \texttt{so} \emph{registerY}[\emph{registerZ}] = \emph{registerV}

\paragraph{Description} The effective address is computed by adding the \emph{registerZ} to the \emph{registerY} and scaled by the \emph{}.
The \emph{registerV} is stored into the memory octuple word at the effective address.


\paragraph{Execution} {Reservation table \textsf{LSU\_MEMW\_AUXR} (Section~\ref{sec:reservation-LSU-MEMW-AUXR})}
~
{\begin{lstlisting}
stage RR:
  new bytemask = 0xFFFFFFFF;
  new doscale = 0;
  new index = GPR[registerY];
  new base = GPR[registerZ];
  new scaling = doscale ? 32 : 1;
  new address = base + (index * scaling);
stage E1:
  new argument3 = QGR[registerV];
stage E3:
  MEM.store(address, 0xFFFFFFFF & bytemask, 0, argument3, registerV << 2);
\end{lstlisting}}
\newpage
\subsubsection[{\small SQ: Store Quadruple Word}]{SQ\hfill Store Quadruple Word}
\label{instr:SQ}\index{sq}
 
\paragraph{Encoding} Format \textsf{LSU\_SQBO} (Section~\ref{sec:format-LSU-SQBO}), \verb|lsucode0 = 10100|\\

\bitfields{ 1/P, 2/01, 5/10100, 5/registerU, 1/0, 2/01, 10/signed10, 6/registerZ }


\paragraph{Syntax} \texttt{sq} \emph{signed10}[\emph{registerZ}] = \emph{registerU}

\paragraph{Description} The effective address is computed by adding the \emph{signed10} to the \emph{registerZ}.
The \emph{registerU} is stored into the memory quadruple word at the effective address.


\paragraph{Execution} {Reservation table \textsf{LSU\_MEMW\_AUXR} (Section~\ref{sec:reservation-LSU-MEMW-AUXR})}
~
{\begin{lstlisting}
stage ID:
  new offset = _SX_10(signed10);
stage RR:
  new base = GPR[registerZ];
  new address = base + offset;
stage E1:
  new argument3 = PGR[registerU];
stage E3:
  MEM.store(address, 0xFFFF, 0, argument3, registerU << 1);
\end{lstlisting}}
\newpage

\paragraph{Encoding} Format \textsf{LSU\_SQBO.X} (Section~\ref{sec:format-LSU-SQBO.X}), \verb|lsucode0 = 10100|\\

\bitfields{ 1/P, 2/00, 2/-- --, 27/upper27, 1/1, 2/01, 5/10100, 5/registerU, 1/0, 2/01, 10/lower10, 6/registerZ }


\paragraph{Syntax} \texttt{sq} \emph{upper27\_\discretionary{}{}{}lower10}[\emph{registerZ}] = \emph{registerU}

\paragraph{Description} The effective address is computed by adding the \emph{upper27\_\discretionary{}{}{}lower10} to the \emph{registerZ}.
The \emph{registerU} is stored into the memory quadruple word at the effective address.


\paragraph{Execution} {Reservation table \textsf{LSU\_MEMW\_AUXR.X} (Section~\ref{sec:reservation-LSU-MEMW-AUXR.X})}
~
{\begin{lstlisting}
stage ID:
  new offset = _SX_37(upper27_lower10);
stage RR:
  new base = GPR[registerZ];
  new address = base + offset;
stage E1:
  new argument3 = PGR[registerU];
stage E3:
  MEM.store(address, 0xFFFF, 0, argument3, registerU << 1);
\end{lstlisting}}
\newpage

\paragraph{Encoding} Format \textsf{LSU\_SQBO.Y} (Section~\ref{sec:format-LSU-SQBO.Y}), \verb|lsucode0 = 10100|\\

\bitfields{ 1/P, 2/00, 2/-- --, 27/extend27, 1/1, 2/00, 2/-- --, 27/upper27, 1/1, 2/01, 5/10100, 5/registerU, 1/0, 2/01, 10/lower10, 6/registerZ }


\paragraph{Syntax} \texttt{sq} \emph{extend27\_\discretionary{}{}{}upper27\_\discretionary{}{}{}lower10}[\emph{registerZ}] = \emph{registerU}

\paragraph{Description} The effective address is computed by adding the \emph{extend27\_\discretionary{}{}{}upper27\_\discretionary{}{}{}lower10} to the \emph{registerZ}.
The \emph{registerU} is stored into the memory quadruple word at the effective address.


\paragraph{Execution} {Reservation table \textsf{LSU\_MEMW\_AUXR.Y} (Section~\ref{sec:reservation-LSU-MEMW-AUXR.Y})}
~
{\begin{lstlisting}
stage ID:
  new offset = _WX_64(extend27_upper27_lower10);
stage RR:
  new base = GPR[registerZ];
  new address = base + offset;
stage E1:
  new argument3 = PGR[registerU];
stage E3:
  MEM.store(address, 0xFFFF, 0, argument3, registerU << 1);
\end{lstlisting}}
\newpage

\paragraph{Encoding} Format \textsf{LSU\_SQBI} (Section~\ref{sec:format-LSU-SQBI}), \verb|lsucode0 = 10100|\\

\bitfields{ 1/P, 2/01, 5/10100, 5/registerU, 1/0, 2/11, 3/111, 1/--, 6/registerY, 6/registerZ }


\paragraph{Syntax} \texttt{sq} \emph{registerY}[\emph{registerZ}] = \emph{registerU}

\paragraph{Description} The effective address is computed by adding the \emph{registerZ} to the \emph{registerY} and scaled by the \emph{}.
The \emph{registerU} is stored into the memory quadruple word at the effective address.


\paragraph{Execution} {Reservation table \textsf{LSU\_MEMW\_AUXR} (Section~\ref{sec:reservation-LSU-MEMW-AUXR})}
~
{\begin{lstlisting}
stage RR:
  new doscale = 0;
  new index = GPR[registerY];
  new base = GPR[registerZ];
  new scaling = doscale ? 16 : 1;
  new address = base + (index * scaling);
stage E1:
  new argument3 = PGR[registerU];
stage E3:
  MEM.store(address, 0xFFFF, 0, argument3, registerU << 1);
\end{lstlisting}}
\newpage
\subsubsection[{\small SW: Store Word}]{SW\hfill Store Word}
\label{instr:SW}\index{sw}
 
\paragraph{Encoding} Format \textsf{LSU\_SSBO} (Section~\ref{sec:format-LSU-SSBO}), \verb|lsucode0 = 10010|\\

\bitfields{ 1/P, 2/01, 5/10010, 6/registerT, 2/01, 10/signed10, 6/registerZ }


\paragraph{Syntax} \texttt{sw} \emph{signed10}[\emph{registerZ}] = \emph{registerT}

\paragraph{Description} The effective address is computed by adding the \emph{signed10} to the \emph{registerZ}.
The \emph{registerT} is stored into the memory word at the effective address.


\paragraph{Execution} {Reservation table \textsf{LSU\_MEMW\_AUXR} (Section~\ref{sec:reservation-LSU-MEMW-AUXR})}
~
{\begin{lstlisting}
stage ID:
  new offset = _SX_10(signed10);
stage RR:
  new base = GPR[registerZ];
  new address = base + offset;
stage E1:
  new argument3 = GPR[registerT];
stage E3:
  MEM.store(address, 0xF, 0, argument3, registerT);
\end{lstlisting}}
\newpage

\paragraph{Encoding} Format \textsf{LSU\_SSBO.X} (Section~\ref{sec:format-LSU-SSBO.X}), \verb|lsucode0 = 10010|\\

\bitfields{ 1/P, 2/00, 2/-- --, 27/upper27, 1/1, 2/01, 5/10010, 6/registerT, 2/01, 10/lower10, 6/registerZ }


\paragraph{Syntax} \texttt{sw} \emph{upper27\_\discretionary{}{}{}lower10}[\emph{registerZ}] = \emph{registerT}

\paragraph{Description} The effective address is computed by adding the \emph{upper27\_\discretionary{}{}{}lower10} to the \emph{registerZ}.
The \emph{registerT} is stored into the memory word at the effective address.


\paragraph{Execution} {Reservation table \textsf{LSU\_MEMW\_AUXR.X} (Section~\ref{sec:reservation-LSU-MEMW-AUXR.X})}
~
{\begin{lstlisting}
stage ID:
  new offset = _SX_37(upper27_lower10);
stage RR:
  new base = GPR[registerZ];
  new address = base + offset;
stage E1:
  new argument3 = GPR[registerT];
stage E3:
  MEM.store(address, 0xF, 0, argument3, registerT);
\end{lstlisting}}
\newpage

\paragraph{Encoding} Format \textsf{LSU\_SSBO.Y} (Section~\ref{sec:format-LSU-SSBO.Y}), \verb|lsucode0 = 10010|\\

\bitfields{ 1/P, 2/00, 2/-- --, 27/extend27, 1/1, 2/00, 2/-- --, 27/upper27, 1/1, 2/01, 5/10010, 6/registerT, 2/01, 10/lower10, 6/registerZ }


\paragraph{Syntax} \texttt{sw} \emph{extend27\_\discretionary{}{}{}upper27\_\discretionary{}{}{}lower10}[\emph{registerZ}] = \emph{registerT}

\paragraph{Description} The effective address is computed by adding the \emph{extend27\_\discretionary{}{}{}upper27\_\discretionary{}{}{}lower10} to the \emph{registerZ}.
The \emph{registerT} is stored into the memory word at the effective address.


\paragraph{Execution} {Reservation table \textsf{LSU\_MEMW\_AUXR.Y} (Section~\ref{sec:reservation-LSU-MEMW-AUXR.Y})}
~
{\begin{lstlisting}
stage ID:
  new offset = _WX_64(extend27_upper27_lower10);
stage RR:
  new base = GPR[registerZ];
  new address = base + offset;
stage E1:
  new argument3 = GPR[registerT];
stage E3:
  MEM.store(address, 0xF, 0, argument3, registerT);
\end{lstlisting}}
\newpage

\paragraph{Encoding} Format \textsf{LSU\_SSBI} (Section~\ref{sec:format-LSU-SSBI}), \verb|lsucode0 = 10010|\\

\bitfields{ 1/P, 2/01, 5/10010, 6/registerT, 2/11, 3/111, 1/--, 6/registerY, 6/registerZ }


\paragraph{Syntax} \texttt{sw} \emph{registerY}[\emph{registerZ}] = \emph{registerT}

\paragraph{Description} The effective address is computed by adding the \emph{registerZ} to the \emph{registerY} and scaled by the \emph{}.
The \emph{registerT} is stored into the memory word at the effective address.


\paragraph{Execution} {Reservation table \textsf{LSU\_MEMW\_AUXR} (Section~\ref{sec:reservation-LSU-MEMW-AUXR})}
~
{\begin{lstlisting}
stage RR:
  new doscale = 0;
  new index = GPR[registerY];
  new base = GPR[registerZ];
  new scaling = doscale ? 4 : 1;
  new address = base + (index * scaling);
stage E1:
  new argument3 = GPR[registerT];
stage E3:
  MEM.store(address, 0xF, 0, argument3, registerT);
\end{lstlisting}}
\newpage
\subsubsection[{\small XLO: Extension Load Octuple Word}]{XLO\hfill Extension Load Octuple Word}
\label{instr:XLO}\index{xlo}
 
\paragraph{Encoding} Format \textsf{LSU\_XLOBO} (Section~\ref{sec:format-LSU-XLOBO}), \verb|steering = 01|\\

\bitfields{ 1/P, 2/01, 3/000, 2/variant, 6/registerG, 2/01, 10/signed10, 6/registerZ }


\paragraph{Syntax} \texttt{xlo}\emph{variant} \emph{registerG} = \emph{signed10}[\emph{registerZ}]

\paragraph{Description} The effective address is computed by adding the \emph{signed10} to the \emph{registerZ}.
The \emph{registerG} is loaded from the memory octuple word at the effective address.


\paragraph{Modifier(s)}
\noindent\begin{itemize}
\item variant (cf. variant in section \ref{par:modifier-variant} on page \pageref{par:modifier-variant})
\end{itemize}

\paragraph{Execution} {Reservation table \textsf{LSU} (Section~\ref{sec:reservation-LSU})}
~
{\begin{lstlisting}
stage ID:
  new variant = _ZX_2(variant);
  new offset = _SX_10(signed10);
stage RR:
  new bytemask = 0xFFFFFFFF;
  new dri = registerG;
  new base = GPR[registerZ];
  new address = base + offset;
stage E1:
  CS.XMF = 1;
  new result1 = MEM.load(address, 0xFFFFFFFF & bytemask, variant, dri);
stage CR:
  XVR[registerG] = result1;
\end{lstlisting}}
\newpage

\paragraph{Encoding} Format \textsf{LSU\_XLOBO.X} (Section~\ref{sec:format-LSU-XLOBO.X}), \verb|steering = 01|\\

\bitfields{ 1/P, 2/00, 2/-- --, 27/upper27, 1/1, 2/01, 3/000, 2/variant, 6/registerG, 2/01, 10/lower10, 6/registerZ }


\paragraph{Syntax} \texttt{xlo}\emph{variant} \emph{registerG} = \emph{upper27\_\discretionary{}{}{}lower10}[\emph{registerZ}]

\paragraph{Description} The effective address is computed by adding the \emph{upper27\_\discretionary{}{}{}lower10} to the \emph{registerZ}.
The \emph{registerG} is loaded from the memory octuple word at the effective address.


\paragraph{Modifier(s)}
\noindent\begin{itemize}
\item variant (cf. variant in section \ref{par:modifier-variant} on page \pageref{par:modifier-variant})
\end{itemize}

\paragraph{Execution} {Reservation table \textsf{LSU.X} (Section~\ref{sec:reservation-LSU.X})}
~
{\begin{lstlisting}
stage ID:
  new variant = _ZX_2(variant);
  new offset = _SX_37(upper27_lower10);
stage RR:
  new bytemask = 0xFFFFFFFF;
  new dri = registerG;
  new base = GPR[registerZ];
  new address = base + offset;
stage E1:
  CS.XMF = 1;
  new result1 = MEM.load(address, 0xFFFFFFFF & bytemask, variant, dri);
stage CR:
  XVR[registerG] = result1;
\end{lstlisting}}
\newpage

\paragraph{Encoding} Format \textsf{LSU\_XLOBO.Y} (Section~\ref{sec:format-LSU-XLOBO.Y}), \verb|steering = 01|\\

\bitfields{ 1/P, 2/00, 2/-- --, 27/extend27, 1/1, 2/00, 2/-- --, 27/upper27, 1/1, 2/01, 3/000, 2/variant, 6/registerG, 2/01, 10/lower10, 6/registerZ }


\paragraph{Syntax} \texttt{xlo}\emph{variant} \emph{registerG} = \emph{extend27\_\discretionary{}{}{}upper27\_\discretionary{}{}{}lower10}[\emph{registerZ}]

\paragraph{Description} The effective address is computed by adding the \emph{extend27\_\discretionary{}{}{}upper27\_\discretionary{}{}{}lower10} to the \emph{registerZ}.
The \emph{registerG} is loaded from the memory octuple word at the effective address.


\paragraph{Modifier(s)}
\noindent\begin{itemize}
\item variant (cf. variant in section \ref{par:modifier-variant} on page \pageref{par:modifier-variant})
\end{itemize}

\paragraph{Execution} {Reservation table \textsf{LSU.Y} (Section~\ref{sec:reservation-LSU.Y})}
~
{\begin{lstlisting}
stage ID:
  new variant = _ZX_2(variant);
  new offset = _WX_64(extend27_upper27_lower10);
stage RR:
  new bytemask = 0xFFFFFFFF;
  new dri = registerG;
  new base = GPR[registerZ];
  new address = base + offset;
stage E1:
  CS.XMF = 1;
  new result1 = MEM.load(address, 0xFFFFFFFF & bytemask, variant, dri);
stage CR:
  XVR[registerG] = result1;
\end{lstlisting}}
\newpage

\paragraph{Encoding} Format \textsf{LSU\_XLOC4BO} (Section~\ref{sec:format-LSU-XLOC4BO}), \verb|steering = 01|\\

\bitfields{ 1/P, 2/01, 3/010, 2/variant, 4/registerGq, 2/qindex, 2/01, 10/signed10, 6/registerZ }


\paragraph{Syntax} \texttt{xlo}\emph{variant}\emph{qindex} \emph{registerGq} = \emph{signed10}[\emph{registerZ}]

\paragraph{Description} The effective address is computed by adding the \emph{signed10} to the \emph{registerZ}.
The \emph{registerGq} is loaded from the memory octuple word at the effective address.


\paragraph{Modifier(s)}
\noindent\begin{itemize}
\item qindex (cf. qindex in section \ref{par:modifier-qindex} on page \pageref{par:modifier-qindex})
\item variant (cf. variant in section \ref{par:modifier-variant} on page \pageref{par:modifier-variant})
\end{itemize}

\paragraph{Execution} {Reservation table \textsf{LSU} (Section~\ref{sec:reservation-LSU})}
~
{\begin{lstlisting}
stage ID:
  new variant = _ZX_2(variant);
  new qindex = _ZX_2(qindex);
  new offset = _SX_10(signed10);
stage RR:
  new bytemask = 0xFFFFFFFF;
  new dri = registerGq << 2;
  new base = GPR[registerZ];
  new address = base + offset;
stage E1:
  new argument1_0 = XMR[registerGq]:0;
  new argument1_1 = XMR[registerGq]:1;
  new argument1_2 = XMR[registerGq]:2;
  new argument1_3 = XMR[registerGq]:3;
  CS.XMF = 1;
  new result1 = MEM.load(address, 0xFFFFFFFF & bytemask, variant, dri);
  new result1_0 = _ZX_64(result1);
  new result1_1 = result1 >> 64;
  new result1_2 = result1 >> 128;
  new result1_3 = result1 >> 192;
stage E3:
  XMR[registerGq]:0 = insert_64(argument1_0, result1_0, qindex);
  XMR[registerGq]:1 = insert_64(argument1_1, result1_1, qindex);
  XMR[registerGq]:2 = insert_64(argument1_2, result1_2, qindex);
  XMR[registerGq]:3 = insert_64(argument1_3, result1_3, qindex);
\end{lstlisting}}
\newpage

\paragraph{Encoding} Format \textsf{LSU\_XLOC4BO.X} (Section~\ref{sec:format-LSU-XLOC4BO.X}), \verb|steering = 01|\\

\bitfields{ 1/P, 2/00, 2/-- --, 27/upper27, 1/1, 2/01, 3/010, 2/variant, 4/registerGq, 2/qindex, 2/01, 10/lower10, 6/registerZ }


\paragraph{Syntax} \texttt{xlo}\emph{variant}\emph{qindex} \emph{registerGq} = \emph{upper27\_\discretionary{}{}{}lower10}[\emph{registerZ}]

\paragraph{Description} The effective address is computed by adding the \emph{upper27\_\discretionary{}{}{}lower10} to the \emph{registerZ}.
The \emph{registerGq} is loaded from the memory octuple word at the effective address.


\paragraph{Modifier(s)}
\noindent\begin{itemize}
\item qindex (cf. qindex in section \ref{par:modifier-qindex} on page \pageref{par:modifier-qindex})
\item variant (cf. variant in section \ref{par:modifier-variant} on page \pageref{par:modifier-variant})
\end{itemize}

\paragraph{Execution} {Reservation table \textsf{LSU.X} (Section~\ref{sec:reservation-LSU.X})}
~
{\begin{lstlisting}
stage ID:
  new variant = _ZX_2(variant);
  new qindex = _ZX_2(qindex);
  new offset = _SX_37(upper27_lower10);
stage RR:
  new bytemask = 0xFFFFFFFF;
  new dri = registerGq << 2;
  new base = GPR[registerZ];
  new address = base + offset;
stage E1:
  new argument1_0 = XMR[registerGq]:0;
  new argument1_1 = XMR[registerGq]:1;
  new argument1_2 = XMR[registerGq]:2;
  new argument1_3 = XMR[registerGq]:3;
  CS.XMF = 1;
  new result1 = MEM.load(address, 0xFFFFFFFF & bytemask, variant, dri);
  new result1_0 = _ZX_64(result1);
  new result1_1 = result1 >> 64;
  new result1_2 = result1 >> 128;
  new result1_3 = result1 >> 192;
stage E3:
  XMR[registerGq]:0 = insert_64(argument1_0, result1_0, qindex);
  XMR[registerGq]:1 = insert_64(argument1_1, result1_1, qindex);
  XMR[registerGq]:2 = insert_64(argument1_2, result1_2, qindex);
  XMR[registerGq]:3 = insert_64(argument1_3, result1_3, qindex);
\end{lstlisting}}
\newpage

\paragraph{Encoding} Format \textsf{LSU\_XLOC4BO.Y} (Section~\ref{sec:format-LSU-XLOC4BO.Y}), \verb|steering = 01|\\

\bitfields{ 1/P, 2/00, 2/-- --, 27/extend27, 1/1, 2/00, 2/-- --, 27/upper27, 1/1, 2/01, 3/010, 2/variant, 4/registerGq, 2/qindex, 2/01, 10/lower10, 6/registerZ }


\paragraph{Syntax} \texttt{xlo}\emph{variant}\emph{qindex} \emph{registerGq} = \emph{extend27\_\discretionary{}{}{}upper27\_\discretionary{}{}{}lower10}[\emph{registerZ}]

\paragraph{Description} The effective address is computed by adding the \emph{extend27\_\discretionary{}{}{}upper27\_\discretionary{}{}{}lower10} to the \emph{registerZ}.
The \emph{registerGq} is loaded from the memory octuple word at the effective address.


\paragraph{Modifier(s)}
\noindent\begin{itemize}
\item qindex (cf. qindex in section \ref{par:modifier-qindex} on page \pageref{par:modifier-qindex})
\item variant (cf. variant in section \ref{par:modifier-variant} on page \pageref{par:modifier-variant})
\end{itemize}

\paragraph{Execution} {Reservation table \textsf{LSU.Y} (Section~\ref{sec:reservation-LSU.Y})}
~
{\begin{lstlisting}
stage ID:
  new variant = _ZX_2(variant);
  new qindex = _ZX_2(qindex);
  new offset = _WX_64(extend27_upper27_lower10);
stage RR:
  new bytemask = 0xFFFFFFFF;
  new dri = registerGq << 2;
  new base = GPR[registerZ];
  new address = base + offset;
stage E1:
  new argument1_0 = XMR[registerGq]:0;
  new argument1_1 = XMR[registerGq]:1;
  new argument1_2 = XMR[registerGq]:2;
  new argument1_3 = XMR[registerGq]:3;
  CS.XMF = 1;
  new result1 = MEM.load(address, 0xFFFFFFFF & bytemask, variant, dri);
  new result1_0 = _ZX_64(result1);
  new result1_1 = result1 >> 64;
  new result1_2 = result1 >> 128;
  new result1_3 = result1 >> 192;
stage E3:
  XMR[registerGq]:0 = insert_64(argument1_0, result1_0, qindex);
  XMR[registerGq]:1 = insert_64(argument1_1, result1_1, qindex);
  XMR[registerGq]:2 = insert_64(argument1_2, result1_2, qindex);
  XMR[registerGq]:3 = insert_64(argument1_3, result1_3, qindex);
\end{lstlisting}}
\newpage

\paragraph{Encoding} Format \textsf{LSU\_XLBI} (Section~\ref{sec:format-LSU-XLBI}), \verb|steering = 01|\\

\bitfields{ 1/P, 2/01, 3/000, 2/variant, 6/registerG, 2/11, 3/111, 1/--, 6/registerY, 6/registerZ }


\paragraph{Syntax} \texttt{xlo}\emph{variant} \emph{registerG} = \emph{registerY}[\emph{registerZ}]

\paragraph{Description} The effective address is computed by adding the \emph{registerZ} to the \emph{registerY} and scaled by the \emph{variant}.
The \emph{registerG} is loaded from the memory octuple word at the effective address.


\paragraph{Modifier(s)}
\noindent\begin{itemize}
\item variant (cf. variant in section \ref{par:modifier-variant} on page \pageref{par:modifier-variant})
\end{itemize}

\paragraph{Execution} {Reservation table \textsf{LSU} (Section~\ref{sec:reservation-LSU})}
~
{\begin{lstlisting}
stage ID:
  new variant = _ZX_2(variant);
stage RR:
  new bytemask = 0xFFFFFFFF;
  new doscale = 0;
  new dri = registerG;
  new base = GPR[registerZ];
  new index = GPR[registerY];
  new scaling = doscale ? 32 : 1;
  new address = base + (index * scaling);
stage E1:
  CS.XMF = 1;
  new result1 = MEM.load(address, 0xFFFFFFFF & bytemask, variant, dri);
stage CR:
  XVR[registerG] = result1;
\end{lstlisting}}
\newpage

\paragraph{Encoding} Format \textsf{LSU\_XLOC4BI} (Section~\ref{sec:format-LSU-XLOC4BI}), \verb|steering = 01|\\

\bitfields{ 1/P, 2/01, 3/010, 2/variant, 4/registerGq, 2/qindex, 2/11, 3/111, 1/--, 6/registerY, 6/registerZ }


\paragraph{Syntax} \texttt{xlo}\emph{variant}\emph{qindex} \emph{registerGq} = \emph{registerY}[\emph{registerZ}]

\paragraph{Description} The effective address is computed by adding the \emph{registerZ} to the \emph{registerY} and scaled by the \emph{qindex}.
The \emph{registerGq} is loaded from the memory octuple word at the effective address.


\paragraph{Modifier(s)}
\noindent\begin{itemize}
\item qindex (cf. qindex in section \ref{par:modifier-qindex} on page \pageref{par:modifier-qindex})
\item variant (cf. variant in section \ref{par:modifier-variant} on page \pageref{par:modifier-variant})
\end{itemize}

\paragraph{Execution} {Reservation table \textsf{LSU} (Section~\ref{sec:reservation-LSU})}
~
{\begin{lstlisting}
stage ID:
  new variant = _ZX_2(variant);
  new qindex = _ZX_2(qindex);
stage RR:
  new bytemask = 0xFFFFFFFF;
  new doscale = 0;
  new dri = registerGq << 2;
  new base = GPR[registerZ];
  new index = GPR[registerY];
  new scaling = doscale ? 32 : 1;
  new address = base + (index * scaling);
stage E1:
  new argument1_0 = XMR[registerGq]:0;
  new argument1_1 = XMR[registerGq]:1;
  new argument1_2 = XMR[registerGq]:2;
  new argument1_3 = XMR[registerGq]:3;
  CS.XMF = 1;
  new result1 = MEM.load(address, 0xFFFFFFFF & bytemask, variant, dri);
  new result1_0 = _ZX_64(result1);
  new result1_1 = result1 >> 64;
  new result1_2 = result1 >> 128;
  new result1_3 = result1 >> 192;
stage E3:
  XMR[registerGq]:0 = insert_64(argument1_0, result1_0, qindex);
  XMR[registerGq]:1 = insert_64(argument1_1, result1_1, qindex);
  XMR[registerGq]:2 = insert_64(argument1_2, result1_2, qindex);
  XMR[registerGq]:3 = insert_64(argument1_3, result1_3, qindex);
\end{lstlisting}}
\newpage
\subsubsection[{\small XPLB: Extension Preload Byte}]{XPLB\hfill Extension Preload Byte}
\label{instr:XPLB}\index{xplb}
 
\paragraph{Encoding} Format \textsf{LSU\_XPL2RBB} (Section~\ref{sec:format-LSU-XPL2RBB}), \verb|lsucode5 = 000|\\

\bitfields{ 1/P, 2/01, 3/011, 2/variant, 5/registerGg, 1/0, 2/11, 3/000, 1/--, 6/registerY, 6/registerZ }


\paragraph{Syntax} \texttt{xplb}\emph{variant} \emph{registerGg}, \emph{registerY} = [\emph{registerZ}]

\paragraph{Description} The \emph{registerGg} is preloaded from the memory byte at the effective address.


\paragraph{Modifier(s)}
\noindent\begin{itemize}
\item variant (cf. variant in section \ref{par:modifier-variant} on page \pageref{par:modifier-variant})
\end{itemize}

\paragraph{Execution} {Reservation table \textsf{LSU} (Section~\ref{sec:reservation-LSU})}
~
{\begin{lstlisting}
stage ID:
  new variant = _ZX_2(variant);
stage RR:
  new address = GPR[registerZ];
  new target = GPR[registerY];
  new shift = (target & 31) * 8;
  new index = (target >> 5) & 1;
  new where = (registerGg << 1) + index;
  new dri = where;
stage E1:
  new argument1 = XVR[where];
  new bytemask = 1;
  new targetbytes = _ZX_256(bits2bytes(bytemask) << shift);
  CS.XMF = 1;
  new result1 = MEM.load(address, 0x1 & bytemask, variant, dri);
  result1 = (_ZX_256(result1 << shift) & targetbytes) | (argument1 & _ZX_256(~targetbytes));
stage CR:
  XVR[where] = result1;
\end{lstlisting}}
\newpage

\paragraph{Encoding} Format \textsf{LSU\_XPL2RBB.O} (Section~\ref{sec:format-LSU-XPL2RBB.O}), \verb|lsucode5 = 000|\\

\bitfields{ 1/P, 2/00, 2/-- --, 27/offset27, 1/1, 2/01, 3/011, 2/variant, 5/registerGg, 1/0, 2/11, 3/000, 1/--, 6/registerY, 6/registerZ }


\paragraph{Syntax} \texttt{xplb}\emph{variant} \emph{registerGg}, \emph{registerY} = \emph{offset27}[\emph{registerZ}]

\paragraph{Description} The \emph{registerGg} is preloaded from the memory byte at the effective address.


\paragraph{Modifier(s)}
\noindent\begin{itemize}
\item variant (cf. variant in section \ref{par:modifier-variant} on page \pageref{par:modifier-variant})
\end{itemize}

\paragraph{Execution} {Reservation table \textsf{LSU.X} (Section~\ref{sec:reservation-LSU.X})}
~
{\begin{lstlisting}
stage ID:
  new offset = _SX_27(offset27);
  new variant = _ZX_2(variant);
stage RR:
  new base = GPR[registerZ];
  new target = GPR[registerY];
  new address = base + offset;
  new shift = (target & 31) * 8;
  new index = (target >> 5) & 1;
  new where = (registerGg << 1) + index;
  new dri = where;
stage E1:
  new argument1 = XVR[where];
  new bytemask = 1;
  new targetbytes = _ZX_256(bits2bytes(bytemask) << shift);
  CS.XMF = 1;
  new result1 = MEM.load(address, 0x1 & bytemask, variant, dri);
  result1 = (_ZX_256(result1 << shift) & targetbytes) | (argument1 & _ZX_256(~targetbytes));
stage CR:
  XVR[where] = result1;
\end{lstlisting}}
\newpage

\paragraph{Encoding} Format \textsf{LSU\_XPL2RBB.D} (Section~\ref{sec:format-LSU-XPL2RBB.D}), \verb|lsucode5 = 000|\\

\bitfields{ 1/P, 2/00, 2/-- --, 27/extend27, 1/1, 2/00, 2/-- --, 27/offset27, 1/1, 2/01, 3/011, 2/variant, 5/registerGg, 1/0, 2/11, 3/000, 1/--, 6/registerY, 6/registerZ }


\paragraph{Syntax} \texttt{xplb}\emph{variant} \emph{registerGg}, \emph{registerY} = \emph{extend27\_\discretionary{}{}{}offset27}[\emph{registerZ}]

\paragraph{Description} The \emph{registerGg} is preloaded from the memory byte at the effective address.


\paragraph{Modifier(s)}
\noindent\begin{itemize}
\item variant (cf. variant in section \ref{par:modifier-variant} on page \pageref{par:modifier-variant})
\end{itemize}

\paragraph{Execution} {Reservation table \textsf{LSU.Y} (Section~\ref{sec:reservation-LSU.Y})}
~
{\begin{lstlisting}
stage ID:
  new offset = _SX_54(extend27_offset27);
  new variant = _ZX_2(variant);
stage RR:
  new base = GPR[registerZ];
  new target = GPR[registerY];
  new address = base + offset;
  new shift = (target & 31) * 8;
  new index = (target >> 5) & 1;
  new where = (registerGg << 1) + index;
  new dri = where;
stage E1:
  new argument1 = XVR[where];
  new bytemask = 1;
  new targetbytes = _ZX_256(bits2bytes(bytemask) << shift);
  CS.XMF = 1;
  new result1 = MEM.load(address, 0x1 & bytemask, variant, dri);
  result1 = (_ZX_256(result1 << shift) & targetbytes) | (argument1 & _ZX_256(~targetbytes));
stage CR:
  XVR[where] = result1;
\end{lstlisting}}
\newpage

\paragraph{Encoding} Format \textsf{LSU\_XPL4RBB} (Section~\ref{sec:format-LSU-XPL4RBB}), \verb|lsucode5 = 000|\\

\bitfields{ 1/P, 2/01, 3/011, 2/variant, 4/registerGh, 2/01, 2/11, 3/000, 1/--, 6/registerY, 6/registerZ }


\paragraph{Syntax} \texttt{xplb}\emph{variant} \emph{registerGh}, \emph{registerY} = [\emph{registerZ}]

\paragraph{Description} The \emph{registerGh} is preloaded from the memory byte at the effective address.


\paragraph{Modifier(s)}
\noindent\begin{itemize}
\item variant (cf. variant in section \ref{par:modifier-variant} on page \pageref{par:modifier-variant})
\end{itemize}

\paragraph{Execution} {Reservation table \textsf{LSU} (Section~\ref{sec:reservation-LSU})}
~
{\begin{lstlisting}
stage ID:
  new variant = _ZX_2(variant);
stage RR:
  new address = GPR[registerZ];
  new target = GPR[registerY];
  new shift = (target & 31) * 8;
  new index = (target >> 5) & 3;
  new where = (registerGh << 2) + index;
  new dri = where;
stage E1:
  new argument1 = XVR[where];
  new bytemask = 1;
  new targetbytes = _ZX_256(bits2bytes(bytemask) << shift);
  CS.XMF = 1;
  new result1 = MEM.load(address, 0x1 & bytemask, variant, dri);
  result1 = (_ZX_256(result1 << shift) & targetbytes) | (argument1 & _ZX_256(~targetbytes));
stage CR:
  XVR[where] = result1;
\end{lstlisting}}
\newpage

\paragraph{Encoding} Format \textsf{LSU\_XPL4RBB.O} (Section~\ref{sec:format-LSU-XPL4RBB.O}), \verb|lsucode5 = 000|\\

\bitfields{ 1/P, 2/00, 2/-- --, 27/offset27, 1/1, 2/01, 3/011, 2/variant, 4/registerGh, 2/01, 2/11, 3/000, 1/--, 6/registerY, 6/registerZ }


\paragraph{Syntax} \texttt{xplb}\emph{variant} \emph{registerGh}, \emph{registerY} = \emph{offset27}[\emph{registerZ}]

\paragraph{Description} The \emph{registerGh} is preloaded from the memory byte at the effective address.


\paragraph{Modifier(s)}
\noindent\begin{itemize}
\item variant (cf. variant in section \ref{par:modifier-variant} on page \pageref{par:modifier-variant})
\end{itemize}

\paragraph{Execution} {Reservation table \textsf{LSU.X} (Section~\ref{sec:reservation-LSU.X})}
~
{\begin{lstlisting}
stage ID:
  new offset = _SX_27(offset27);
  new variant = _ZX_2(variant);
stage RR:
  new base = GPR[registerZ];
  new target = GPR[registerY];
  new address = base + offset;
  new shift = (target & 31) * 8;
  new index = (target >> 5) & 3;
  new where = (registerGh << 2) + index;
  new dri = where;
stage E1:
  new argument1 = XVR[where];
  new bytemask = 1;
  new targetbytes = _ZX_256(bits2bytes(bytemask) << shift);
  CS.XMF = 1;
  new result1 = MEM.load(address, 0x1 & bytemask, variant, dri);
  result1 = (_ZX_256(result1 << shift) & targetbytes) | (argument1 & _ZX_256(~targetbytes));
stage CR:
  XVR[where] = result1;
\end{lstlisting}}
\newpage

\paragraph{Encoding} Format \textsf{LSU\_XPL4RBB.D} (Section~\ref{sec:format-LSU-XPL4RBB.D}), \verb|lsucode5 = 000|\\

\bitfields{ 1/P, 2/00, 2/-- --, 27/extend27, 1/1, 2/00, 2/-- --, 27/offset27, 1/1, 2/01, 3/011, 2/variant, 4/registerGh, 2/01, 2/11, 3/000, 1/--, 6/registerY, 6/registerZ }


\paragraph{Syntax} \texttt{xplb}\emph{variant} \emph{registerGh}, \emph{registerY} = \emph{extend27\_\discretionary{}{}{}offset27}[\emph{registerZ}]

\paragraph{Description} The \emph{registerGh} is preloaded from the memory byte at the effective address.


\paragraph{Modifier(s)}
\noindent\begin{itemize}
\item variant (cf. variant in section \ref{par:modifier-variant} on page \pageref{par:modifier-variant})
\end{itemize}

\paragraph{Execution} {Reservation table \textsf{LSU.Y} (Section~\ref{sec:reservation-LSU.Y})}
~
{\begin{lstlisting}
stage ID:
  new offset = _SX_54(extend27_offset27);
  new variant = _ZX_2(variant);
stage RR:
  new base = GPR[registerZ];
  new target = GPR[registerY];
  new address = base + offset;
  new shift = (target & 31) * 8;
  new index = (target >> 5) & 3;
  new where = (registerGh << 2) + index;
  new dri = where;
stage E1:
  new argument1 = XVR[where];
  new bytemask = 1;
  new targetbytes = _ZX_256(bits2bytes(bytemask) << shift);
  CS.XMF = 1;
  new result1 = MEM.load(address, 0x1 & bytemask, variant, dri);
  result1 = (_ZX_256(result1 << shift) & targetbytes) | (argument1 & _ZX_256(~targetbytes));
stage CR:
  XVR[where] = result1;
\end{lstlisting}}
\newpage

\paragraph{Encoding} Format \textsf{LSU\_XPL8RBB} (Section~\ref{sec:format-LSU-XPL8RBB}), \verb|lsucode5 = 000|\\

\bitfields{ 1/P, 2/01, 3/011, 2/variant, 3/registerGi, 3/011, 2/11, 3/000, 1/--, 6/registerY, 6/registerZ }


\paragraph{Syntax} \texttt{xplb}\emph{variant} \emph{registerGi}, \emph{registerY} = [\emph{registerZ}]

\paragraph{Description} The \emph{registerGi} is preloaded from the memory byte at the effective address.


\paragraph{Modifier(s)}
\noindent\begin{itemize}
\item variant (cf. variant in section \ref{par:modifier-variant} on page \pageref{par:modifier-variant})
\end{itemize}

\paragraph{Execution} {Reservation table \textsf{LSU} (Section~\ref{sec:reservation-LSU})}
~
{\begin{lstlisting}
stage ID:
  new variant = _ZX_2(variant);
stage RR:
  new address = GPR[registerZ];
  new target = GPR[registerY];
  new shift = (target & 31) * 8;
  new index = (target >> 5) & 7;
  new where = (registerGi << 3) + index;
  new dri = where;
stage E1:
  new argument1 = XVR[where];
  new bytemask = 1;
  new targetbytes = _ZX_256(bits2bytes(bytemask) << shift);
  CS.XMF = 1;
  new result1 = MEM.load(address, 0x1 & bytemask, variant, dri);
  result1 = (_ZX_256(result1 << shift) & targetbytes) | (argument1 & _ZX_256(~targetbytes));
stage CR:
  XVR[where] = result1;
\end{lstlisting}}
\newpage

\paragraph{Encoding} Format \textsf{LSU\_XPL8RBB.O} (Section~\ref{sec:format-LSU-XPL8RBB.O}), \verb|lsucode5 = 000|\\

\bitfields{ 1/P, 2/00, 2/-- --, 27/offset27, 1/1, 2/01, 3/011, 2/variant, 3/registerGi, 3/011, 2/11, 3/000, 1/--, 6/registerY, 6/registerZ }


\paragraph{Syntax} \texttt{xplb}\emph{variant} \emph{registerGi}, \emph{registerY} = \emph{offset27}[\emph{registerZ}]

\paragraph{Description} The \emph{registerGi} is preloaded from the memory byte at the effective address.


\paragraph{Modifier(s)}
\noindent\begin{itemize}
\item variant (cf. variant in section \ref{par:modifier-variant} on page \pageref{par:modifier-variant})
\end{itemize}

\paragraph{Execution} {Reservation table \textsf{LSU.X} (Section~\ref{sec:reservation-LSU.X})}
~
{\begin{lstlisting}
stage ID:
  new offset = _SX_27(offset27);
  new variant = _ZX_2(variant);
stage RR:
  new base = GPR[registerZ];
  new target = GPR[registerY];
  new address = base + offset;
  new shift = (target & 31) * 8;
  new index = (target >> 5) & 7;
  new where = (registerGi << 3) + index;
  new dri = where;
stage E1:
  new argument1 = XVR[where];
  new bytemask = 1;
  new targetbytes = _ZX_256(bits2bytes(bytemask) << shift);
  CS.XMF = 1;
  new result1 = MEM.load(address, 0x1 & bytemask, variant, dri);
  result1 = (_ZX_256(result1 << shift) & targetbytes) | (argument1 & _ZX_256(~targetbytes));
stage CR:
  XVR[where] = result1;
\end{lstlisting}}
\newpage

\paragraph{Encoding} Format \textsf{LSU\_XPL8RBB.D} (Section~\ref{sec:format-LSU-XPL8RBB.D}), \verb|lsucode5 = 000|\\

\bitfields{ 1/P, 2/00, 2/-- --, 27/extend27, 1/1, 2/00, 2/-- --, 27/offset27, 1/1, 2/01, 3/011, 2/variant, 3/registerGi, 3/011, 2/11, 3/000, 1/--, 6/registerY, 6/registerZ }


\paragraph{Syntax} \texttt{xplb}\emph{variant} \emph{registerGi}, \emph{registerY} = \emph{extend27\_\discretionary{}{}{}offset27}[\emph{registerZ}]

\paragraph{Description} The \emph{registerGi} is preloaded from the memory byte at the effective address.


\paragraph{Modifier(s)}
\noindent\begin{itemize}
\item variant (cf. variant in section \ref{par:modifier-variant} on page \pageref{par:modifier-variant})
\end{itemize}

\paragraph{Execution} {Reservation table \textsf{LSU.Y} (Section~\ref{sec:reservation-LSU.Y})}
~
{\begin{lstlisting}
stage ID:
  new offset = _SX_54(extend27_offset27);
  new variant = _ZX_2(variant);
stage RR:
  new base = GPR[registerZ];
  new target = GPR[registerY];
  new address = base + offset;
  new shift = (target & 31) * 8;
  new index = (target >> 5) & 7;
  new where = (registerGi << 3) + index;
  new dri = where;
stage E1:
  new argument1 = XVR[where];
  new bytemask = 1;
  new targetbytes = _ZX_256(bits2bytes(bytemask) << shift);
  CS.XMF = 1;
  new result1 = MEM.load(address, 0x1 & bytemask, variant, dri);
  result1 = (_ZX_256(result1 << shift) & targetbytes) | (argument1 & _ZX_256(~targetbytes));
stage CR:
  XVR[where] = result1;
\end{lstlisting}}
\newpage

\paragraph{Encoding} Format \textsf{LSU\_XPL16RBB} (Section~\ref{sec:format-LSU-XPL16RBB}), \verb|lsucode5 = 000|\\

\bitfields{ 1/P, 2/01, 3/011, 2/variant, 2/registerGj, 4/0111, 2/11, 3/000, 1/--, 6/registerY, 6/registerZ }


\paragraph{Syntax} \texttt{xplb}\emph{variant} \emph{registerGj}, \emph{registerY} = [\emph{registerZ}]

\paragraph{Description} The \emph{registerGj} is preloaded from the memory byte at the effective address.


\paragraph{Modifier(s)}
\noindent\begin{itemize}
\item variant (cf. variant in section \ref{par:modifier-variant} on page \pageref{par:modifier-variant})
\end{itemize}

\paragraph{Execution} {Reservation table \textsf{LSU} (Section~\ref{sec:reservation-LSU})}
~
{\begin{lstlisting}
stage ID:
  new variant = _ZX_2(variant);
stage RR:
  new address = GPR[registerZ];
  new target = GPR[registerY];
  new shift = (target & 31) * 8;
  new index = (target >> 5) & 15;
  new where = (registerGj << 4) + index;
  new dri = where;
stage E1:
  new argument1 = XVR[where];
  new bytemask = 1;
  new targetbytes = _ZX_256(bits2bytes(bytemask) << shift);
  CS.XMF = 1;
  new result1 = MEM.load(address, 0x1 & bytemask, variant, dri);
  result1 = (_ZX_256(result1 << shift) & targetbytes) | (argument1 & _ZX_256(~targetbytes));
stage CR:
  XVR[where] = result1;
\end{lstlisting}}
\newpage

\paragraph{Encoding} Format \textsf{LSU\_XPL16RBB.O} (Section~\ref{sec:format-LSU-XPL16RBB.O}), \verb|lsucode5 = 000|\\

\bitfields{ 1/P, 2/00, 2/-- --, 27/offset27, 1/1, 2/01, 3/011, 2/variant, 2/registerGj, 4/0111, 2/11, 3/000, 1/--, 6/registerY, 6/registerZ }


\paragraph{Syntax} \texttt{xplb}\emph{variant} \emph{registerGj}, \emph{registerY} = \emph{offset27}[\emph{registerZ}]

\paragraph{Description} The \emph{registerGj} is preloaded from the memory byte at the effective address.


\paragraph{Modifier(s)}
\noindent\begin{itemize}
\item variant (cf. variant in section \ref{par:modifier-variant} on page \pageref{par:modifier-variant})
\end{itemize}

\paragraph{Execution} {Reservation table \textsf{LSU.X} (Section~\ref{sec:reservation-LSU.X})}
~
{\begin{lstlisting}
stage ID:
  new offset = _SX_27(offset27);
  new variant = _ZX_2(variant);
stage RR:
  new base = GPR[registerZ];
  new target = GPR[registerY];
  new address = base + offset;
  new shift = (target & 31) * 8;
  new index = (target >> 5) & 15;
  new where = (registerGj << 4) + index;
  new dri = where;
stage E1:
  new argument1 = XVR[where];
  new bytemask = 1;
  new targetbytes = _ZX_256(bits2bytes(bytemask) << shift);
  CS.XMF = 1;
  new result1 = MEM.load(address, 0x1 & bytemask, variant, dri);
  result1 = (_ZX_256(result1 << shift) & targetbytes) | (argument1 & _ZX_256(~targetbytes));
stage CR:
  XVR[where] = result1;
\end{lstlisting}}
\newpage

\paragraph{Encoding} Format \textsf{LSU\_XPL16RBB.D} (Section~\ref{sec:format-LSU-XPL16RBB.D}), \verb|lsucode5 = 000|\\

\bitfields{ 1/P, 2/00, 2/-- --, 27/extend27, 1/1, 2/00, 2/-- --, 27/offset27, 1/1, 2/01, 3/011, 2/variant, 2/registerGj, 4/0111, 2/11, 3/000, 1/--, 6/registerY, 6/registerZ }


\paragraph{Syntax} \texttt{xplb}\emph{variant} \emph{registerGj}, \emph{registerY} = \emph{extend27\_\discretionary{}{}{}offset27}[\emph{registerZ}]

\paragraph{Description} The \emph{registerGj} is preloaded from the memory byte at the effective address.


\paragraph{Modifier(s)}
\noindent\begin{itemize}
\item variant (cf. variant in section \ref{par:modifier-variant} on page \pageref{par:modifier-variant})
\end{itemize}

\paragraph{Execution} {Reservation table \textsf{LSU.Y} (Section~\ref{sec:reservation-LSU.Y})}
~
{\begin{lstlisting}
stage ID:
  new offset = _SX_54(extend27_offset27);
  new variant = _ZX_2(variant);
stage RR:
  new base = GPR[registerZ];
  new target = GPR[registerY];
  new address = base + offset;
  new shift = (target & 31) * 8;
  new index = (target >> 5) & 15;
  new where = (registerGj << 4) + index;
  new dri = where;
stage E1:
  new argument1 = XVR[where];
  new bytemask = 1;
  new targetbytes = _ZX_256(bits2bytes(bytemask) << shift);
  CS.XMF = 1;
  new result1 = MEM.load(address, 0x1 & bytemask, variant, dri);
  result1 = (_ZX_256(result1 << shift) & targetbytes) | (argument1 & _ZX_256(~targetbytes));
stage CR:
  XVR[where] = result1;
\end{lstlisting}}
\newpage

\paragraph{Encoding} Format \textsf{LSU\_XPL32RBB} (Section~\ref{sec:format-LSU-XPL32RBB}), \verb|lsucode5 = 000|\\

\bitfields{ 1/P, 2/01, 3/011, 2/variant, 1/registerGk, 5/01111, 2/11, 3/000, 1/--, 6/registerY, 6/registerZ }


\paragraph{Syntax} \texttt{xplb}\emph{variant} \emph{registerGk}, \emph{registerY} = [\emph{registerZ}]

\paragraph{Description} The \emph{registerGk} is preloaded from the memory byte at the effective address.


\paragraph{Modifier(s)}
\noindent\begin{itemize}
\item variant (cf. variant in section \ref{par:modifier-variant} on page \pageref{par:modifier-variant})
\end{itemize}

\paragraph{Execution} {Reservation table \textsf{LSU} (Section~\ref{sec:reservation-LSU})}
~
{\begin{lstlisting}
stage ID:
  new variant = _ZX_2(variant);
stage RR:
  new address = GPR[registerZ];
  new target = GPR[registerY];
  new shift = (target & 31) * 8;
  new index = (target >> 5) & 31;
  new where = (registerGk << 5) + index;
  new dri = where;
stage E1:
  new argument1 = XVR[where];
  new bytemask = 1;
  new targetbytes = _ZX_256(bits2bytes(bytemask) << shift);
  CS.XMF = 1;
  new result1 = MEM.load(address, 0x1 & bytemask, variant, dri);
  result1 = (_ZX_256(result1 << shift) & targetbytes) | (argument1 & _ZX_256(~targetbytes));
stage CR:
  XVR[where] = result1;
\end{lstlisting}}
\newpage

\paragraph{Encoding} Format \textsf{LSU\_XPL32RBB.O} (Section~\ref{sec:format-LSU-XPL32RBB.O}), \verb|lsucode5 = 000|\\

\bitfields{ 1/P, 2/00, 2/-- --, 27/offset27, 1/1, 2/01, 3/011, 2/variant, 1/registerGk, 5/01111, 2/11, 3/000, 1/--, 6/registerY, 6/registerZ }


\paragraph{Syntax} \texttt{xplb}\emph{variant} \emph{registerGk}, \emph{registerY} = \emph{offset27}[\emph{registerZ}]

\paragraph{Description} The \emph{registerGk} is preloaded from the memory byte at the effective address.


\paragraph{Modifier(s)}
\noindent\begin{itemize}
\item variant (cf. variant in section \ref{par:modifier-variant} on page \pageref{par:modifier-variant})
\end{itemize}

\paragraph{Execution} {Reservation table \textsf{LSU.X} (Section~\ref{sec:reservation-LSU.X})}
~
{\begin{lstlisting}
stage ID:
  new offset = _SX_27(offset27);
  new variant = _ZX_2(variant);
stage RR:
  new base = GPR[registerZ];
  new target = GPR[registerY];
  new address = base + offset;
  new shift = (target & 31) * 8;
  new index = (target >> 5) & 31;
  new where = (registerGk << 5) + index;
  new dri = where;
stage E1:
  new argument1 = XVR[where];
  new bytemask = 1;
  new targetbytes = _ZX_256(bits2bytes(bytemask) << shift);
  CS.XMF = 1;
  new result1 = MEM.load(address, 0x1 & bytemask, variant, dri);
  result1 = (_ZX_256(result1 << shift) & targetbytes) | (argument1 & _ZX_256(~targetbytes));
stage CR:
  XVR[where] = result1;
\end{lstlisting}}
\newpage

\paragraph{Encoding} Format \textsf{LSU\_XPL32RBB.D} (Section~\ref{sec:format-LSU-XPL32RBB.D}), \verb|lsucode5 = 000|\\

\bitfields{ 1/P, 2/00, 2/-- --, 27/extend27, 1/1, 2/00, 2/-- --, 27/offset27, 1/1, 2/01, 3/011, 2/variant, 1/registerGk, 5/01111, 2/11, 3/000, 1/--, 6/registerY, 6/registerZ }


\paragraph{Syntax} \texttt{xplb}\emph{variant} \emph{registerGk}, \emph{registerY} = \emph{extend27\_\discretionary{}{}{}offset27}[\emph{registerZ}]

\paragraph{Description} The \emph{registerGk} is preloaded from the memory byte at the effective address.


\paragraph{Modifier(s)}
\noindent\begin{itemize}
\item variant (cf. variant in section \ref{par:modifier-variant} on page \pageref{par:modifier-variant})
\end{itemize}

\paragraph{Execution} {Reservation table \textsf{LSU.Y} (Section~\ref{sec:reservation-LSU.Y})}
~
{\begin{lstlisting}
stage ID:
  new offset = _SX_54(extend27_offset27);
  new variant = _ZX_2(variant);
stage RR:
  new base = GPR[registerZ];
  new target = GPR[registerY];
  new address = base + offset;
  new shift = (target & 31) * 8;
  new index = (target >> 5) & 31;
  new where = (registerGk << 5) + index;
  new dri = where;
stage E1:
  new argument1 = XVR[where];
  new bytemask = 1;
  new targetbytes = _ZX_256(bits2bytes(bytemask) << shift);
  CS.XMF = 1;
  new result1 = MEM.load(address, 0x1 & bytemask, variant, dri);
  result1 = (_ZX_256(result1 << shift) & targetbytes) | (argument1 & _ZX_256(~targetbytes));
stage CR:
  XVR[where] = result1;
\end{lstlisting}}
\newpage

\paragraph{Encoding} Format \textsf{LSU\_XPL64RBB} (Section~\ref{sec:format-LSU-XPL64RBB}), \verb|lsucode5 = 000|\\

\bitfields{ 1/P, 2/01, 3/011, 2/variant, 1/registerGl, 5/11111, 2/11, 3/000, 1/--, 6/registerY, 6/registerZ }


\paragraph{Syntax} \texttt{xplb}\emph{variant} \emph{registerGl}, \emph{registerY} = [\emph{registerZ}]

\paragraph{Description} The \emph{registerGl} is preloaded from the memory byte at the effective address.


\paragraph{Modifier(s)}
\noindent\begin{itemize}
\item variant (cf. variant in section \ref{par:modifier-variant} on page \pageref{par:modifier-variant})
\end{itemize}

\paragraph{Execution} {Reservation table \textsf{LSU} (Section~\ref{sec:reservation-LSU})}
~
{\begin{lstlisting}
stage ID:
  new variant = _ZX_2(variant);
stage RR:
  new address = GPR[registerZ];
  new target = GPR[registerY];
  new shift = (target & 31) * 8;
  new index = (target >> 5) & 63;
  new where = (registerGl << 6) + index;
  new dri = where;
stage E1:
  new argument1 = XVR[where];
  new bytemask = 1;
  new targetbytes = _ZX_256(bits2bytes(bytemask) << shift);
  CS.XMF = 1;
  new result1 = MEM.load(address, 0x1 & bytemask, variant, dri);
  result1 = (_ZX_256(result1 << shift) & targetbytes) | (argument1 & _ZX_256(~targetbytes));
stage CR:
  XVR[where] = result1;
\end{lstlisting}}
\newpage

\paragraph{Encoding} Format \textsf{LSU\_XPL64RBB.O} (Section~\ref{sec:format-LSU-XPL64RBB.O}), \verb|lsucode5 = 000|\\

\bitfields{ 1/P, 2/00, 2/-- --, 27/offset27, 1/1, 2/01, 3/011, 2/variant, 1/registerGl, 5/11111, 2/11, 3/000, 1/--, 6/registerY, 6/registerZ }


\paragraph{Syntax} \texttt{xplb}\emph{variant} \emph{registerGl}, \emph{registerY} = \emph{offset27}[\emph{registerZ}]

\paragraph{Description} The \emph{registerGl} is preloaded from the memory byte at the effective address.


\paragraph{Modifier(s)}
\noindent\begin{itemize}
\item variant (cf. variant in section \ref{par:modifier-variant} on page \pageref{par:modifier-variant})
\end{itemize}

\paragraph{Execution} {Reservation table \textsf{LSU.X} (Section~\ref{sec:reservation-LSU.X})}
~
{\begin{lstlisting}
stage ID:
  new offset = _SX_27(offset27);
  new variant = _ZX_2(variant);
stage RR:
  new base = GPR[registerZ];
  new target = GPR[registerY];
  new address = base + offset;
  new shift = (target & 31) * 8;
  new index = (target >> 5) & 63;
  new where = (registerGl << 6) + index;
  new dri = where;
stage E1:
  new argument1 = XVR[where];
  new bytemask = 1;
  new targetbytes = _ZX_256(bits2bytes(bytemask) << shift);
  CS.XMF = 1;
  new result1 = MEM.load(address, 0x1 & bytemask, variant, dri);
  result1 = (_ZX_256(result1 << shift) & targetbytes) | (argument1 & _ZX_256(~targetbytes));
stage CR:
  XVR[where] = result1;
\end{lstlisting}}
\newpage

\paragraph{Encoding} Format \textsf{LSU\_XPL64RBB.D} (Section~\ref{sec:format-LSU-XPL64RBB.D}), \verb|lsucode5 = 000|\\

\bitfields{ 1/P, 2/00, 2/-- --, 27/extend27, 1/1, 2/00, 2/-- --, 27/offset27, 1/1, 2/01, 3/011, 2/variant, 1/registerGl, 5/11111, 2/11, 3/000, 1/--, 6/registerY, 6/registerZ }


\paragraph{Syntax} \texttt{xplb}\emph{variant} \emph{registerGl}, \emph{registerY} = \emph{extend27\_\discretionary{}{}{}offset27}[\emph{registerZ}]

\paragraph{Description} The \emph{registerGl} is preloaded from the memory byte at the effective address.


\paragraph{Modifier(s)}
\noindent\begin{itemize}
\item variant (cf. variant in section \ref{par:modifier-variant} on page \pageref{par:modifier-variant})
\end{itemize}

\paragraph{Execution} {Reservation table \textsf{LSU.Y} (Section~\ref{sec:reservation-LSU.Y})}
~
{\begin{lstlisting}
stage ID:
  new offset = _SX_54(extend27_offset27);
  new variant = _ZX_2(variant);
stage RR:
  new base = GPR[registerZ];
  new target = GPR[registerY];
  new address = base + offset;
  new shift = (target & 31) * 8;
  new index = (target >> 5) & 63;
  new where = (registerGl << 6) + index;
  new dri = where;
stage E1:
  new argument1 = XVR[where];
  new bytemask = 1;
  new targetbytes = _ZX_256(bits2bytes(bytemask) << shift);
  CS.XMF = 1;
  new result1 = MEM.load(address, 0x1 & bytemask, variant, dri);
  result1 = (_ZX_256(result1 << shift) & targetbytes) | (argument1 & _ZX_256(~targetbytes));
stage CR:
  XVR[where] = result1;
\end{lstlisting}}
\newpage
\subsubsection[{\small XPLD: Extension Preload Double Word}]{XPLD\hfill Extension Preload Double Word}
\label{instr:XPLD}\index{xpld}
 
\paragraph{Encoding} Format \textsf{LSU\_XPL2RBB} (Section~\ref{sec:format-LSU-XPL2RBB}), \verb|lsucode5 = 011|\\

\bitfields{ 1/P, 2/01, 3/011, 2/variant, 5/registerGg, 1/0, 2/11, 3/011, 1/--, 6/registerY, 6/registerZ }


\paragraph{Syntax} \texttt{xpld}\emph{variant} \emph{registerGg}, \emph{registerY} = [\emph{registerZ}]

\paragraph{Description} The \emph{registerGg} is preloaded from the memory double word at the effective address.


\paragraph{Modifier(s)}
\noindent\begin{itemize}
\item variant (cf. variant in section \ref{par:modifier-variant} on page \pageref{par:modifier-variant})
\end{itemize}

\paragraph{Execution} {Reservation table \textsf{LSU} (Section~\ref{sec:reservation-LSU})}
~
{\begin{lstlisting}
stage ID:
  new variant = _ZX_2(variant);
stage RR:
  new address = GPR[registerZ];
  new target = GPR[registerY];
  new shift = (target & 31) * 8;
  new index = (target >> 5) & 1;
  new where = (registerGg << 1) + index;
  new dri = where;
stage E1:
  new argument1 = XVR[where];
  new bytemask = 255;
  new targetbytes = _ZX_256(bits2bytes(bytemask) << shift);
  CS.XMF = 1;
  new result1 = MEM.load(address, 0xFF & bytemask, variant, dri);
  result1 = (_ZX_256(result1 << shift) & targetbytes) | (argument1 & _ZX_256(~targetbytes));
stage CR:
  XVR[where] = result1;
\end{lstlisting}}
\newpage

\paragraph{Encoding} Format \textsf{LSU\_XPL2RBB.O} (Section~\ref{sec:format-LSU-XPL2RBB.O}), \verb|lsucode5 = 011|\\

\bitfields{ 1/P, 2/00, 2/-- --, 27/offset27, 1/1, 2/01, 3/011, 2/variant, 5/registerGg, 1/0, 2/11, 3/011, 1/--, 6/registerY, 6/registerZ }


\paragraph{Syntax} \texttt{xpld}\emph{variant} \emph{registerGg}, \emph{registerY} = \emph{offset27}[\emph{registerZ}]

\paragraph{Description} The \emph{registerGg} is preloaded from the memory double word at the effective address.


\paragraph{Modifier(s)}
\noindent\begin{itemize}
\item variant (cf. variant in section \ref{par:modifier-variant} on page \pageref{par:modifier-variant})
\end{itemize}

\paragraph{Execution} {Reservation table \textsf{LSU.X} (Section~\ref{sec:reservation-LSU.X})}
~
{\begin{lstlisting}
stage ID:
  new offset = _SX_27(offset27);
  new variant = _ZX_2(variant);
stage RR:
  new base = GPR[registerZ];
  new target = GPR[registerY];
  new address = base + offset;
  new shift = (target & 31) * 8;
  new index = (target >> 5) & 1;
  new where = (registerGg << 1) + index;
  new dri = where;
stage E1:
  new argument1 = XVR[where];
  new bytemask = 255;
  new targetbytes = _ZX_256(bits2bytes(bytemask) << shift);
  CS.XMF = 1;
  new result1 = MEM.load(address, 0xFF & bytemask, variant, dri);
  result1 = (_ZX_256(result1 << shift) & targetbytes) | (argument1 & _ZX_256(~targetbytes));
stage CR:
  XVR[where] = result1;
\end{lstlisting}}
\newpage

\paragraph{Encoding} Format \textsf{LSU\_XPL2RBB.D} (Section~\ref{sec:format-LSU-XPL2RBB.D}), \verb|lsucode5 = 011|\\

\bitfields{ 1/P, 2/00, 2/-- --, 27/extend27, 1/1, 2/00, 2/-- --, 27/offset27, 1/1, 2/01, 3/011, 2/variant, 5/registerGg, 1/0, 2/11, 3/011, 1/--, 6/registerY, 6/registerZ }


\paragraph{Syntax} \texttt{xpld}\emph{variant} \emph{registerGg}, \emph{registerY} = \emph{extend27\_\discretionary{}{}{}offset27}[\emph{registerZ}]

\paragraph{Description} The \emph{registerGg} is preloaded from the memory double word at the effective address.


\paragraph{Modifier(s)}
\noindent\begin{itemize}
\item variant (cf. variant in section \ref{par:modifier-variant} on page \pageref{par:modifier-variant})
\end{itemize}

\paragraph{Execution} {Reservation table \textsf{LSU.Y} (Section~\ref{sec:reservation-LSU.Y})}
~
{\begin{lstlisting}
stage ID:
  new offset = _SX_54(extend27_offset27);
  new variant = _ZX_2(variant);
stage RR:
  new base = GPR[registerZ];
  new target = GPR[registerY];
  new address = base + offset;
  new shift = (target & 31) * 8;
  new index = (target >> 5) & 1;
  new where = (registerGg << 1) + index;
  new dri = where;
stage E1:
  new argument1 = XVR[where];
  new bytemask = 255;
  new targetbytes = _ZX_256(bits2bytes(bytemask) << shift);
  CS.XMF = 1;
  new result1 = MEM.load(address, 0xFF & bytemask, variant, dri);
  result1 = (_ZX_256(result1 << shift) & targetbytes) | (argument1 & _ZX_256(~targetbytes));
stage CR:
  XVR[where] = result1;
\end{lstlisting}}
\newpage

\paragraph{Encoding} Format \textsf{LSU\_XPL4RBB} (Section~\ref{sec:format-LSU-XPL4RBB}), \verb|lsucode5 = 011|\\

\bitfields{ 1/P, 2/01, 3/011, 2/variant, 4/registerGh, 2/01, 2/11, 3/011, 1/--, 6/registerY, 6/registerZ }


\paragraph{Syntax} \texttt{xpld}\emph{variant} \emph{registerGh}, \emph{registerY} = [\emph{registerZ}]

\paragraph{Description} The \emph{registerGh} is preloaded from the memory double word at the effective address.


\paragraph{Modifier(s)}
\noindent\begin{itemize}
\item variant (cf. variant in section \ref{par:modifier-variant} on page \pageref{par:modifier-variant})
\end{itemize}

\paragraph{Execution} {Reservation table \textsf{LSU} (Section~\ref{sec:reservation-LSU})}
~
{\begin{lstlisting}
stage ID:
  new variant = _ZX_2(variant);
stage RR:
  new address = GPR[registerZ];
  new target = GPR[registerY];
  new shift = (target & 31) * 8;
  new index = (target >> 5) & 3;
  new where = (registerGh << 2) + index;
  new dri = where;
stage E1:
  new argument1 = XVR[where];
  new bytemask = 255;
  new targetbytes = _ZX_256(bits2bytes(bytemask) << shift);
  CS.XMF = 1;
  new result1 = MEM.load(address, 0xFF & bytemask, variant, dri);
  result1 = (_ZX_256(result1 << shift) & targetbytes) | (argument1 & _ZX_256(~targetbytes));
stage CR:
  XVR[where] = result1;
\end{lstlisting}}
\newpage

\paragraph{Encoding} Format \textsf{LSU\_XPL4RBB.O} (Section~\ref{sec:format-LSU-XPL4RBB.O}), \verb|lsucode5 = 011|\\

\bitfields{ 1/P, 2/00, 2/-- --, 27/offset27, 1/1, 2/01, 3/011, 2/variant, 4/registerGh, 2/01, 2/11, 3/011, 1/--, 6/registerY, 6/registerZ }


\paragraph{Syntax} \texttt{xpld}\emph{variant} \emph{registerGh}, \emph{registerY} = \emph{offset27}[\emph{registerZ}]

\paragraph{Description} The \emph{registerGh} is preloaded from the memory double word at the effective address.


\paragraph{Modifier(s)}
\noindent\begin{itemize}
\item variant (cf. variant in section \ref{par:modifier-variant} on page \pageref{par:modifier-variant})
\end{itemize}

\paragraph{Execution} {Reservation table \textsf{LSU.X} (Section~\ref{sec:reservation-LSU.X})}
~
{\begin{lstlisting}
stage ID:
  new offset = _SX_27(offset27);
  new variant = _ZX_2(variant);
stage RR:
  new base = GPR[registerZ];
  new target = GPR[registerY];
  new address = base + offset;
  new shift = (target & 31) * 8;
  new index = (target >> 5) & 3;
  new where = (registerGh << 2) + index;
  new dri = where;
stage E1:
  new argument1 = XVR[where];
  new bytemask = 255;
  new targetbytes = _ZX_256(bits2bytes(bytemask) << shift);
  CS.XMF = 1;
  new result1 = MEM.load(address, 0xFF & bytemask, variant, dri);
  result1 = (_ZX_256(result1 << shift) & targetbytes) | (argument1 & _ZX_256(~targetbytes));
stage CR:
  XVR[where] = result1;
\end{lstlisting}}
\newpage

\paragraph{Encoding} Format \textsf{LSU\_XPL4RBB.D} (Section~\ref{sec:format-LSU-XPL4RBB.D}), \verb|lsucode5 = 011|\\

\bitfields{ 1/P, 2/00, 2/-- --, 27/extend27, 1/1, 2/00, 2/-- --, 27/offset27, 1/1, 2/01, 3/011, 2/variant, 4/registerGh, 2/01, 2/11, 3/011, 1/--, 6/registerY, 6/registerZ }


\paragraph{Syntax} \texttt{xpld}\emph{variant} \emph{registerGh}, \emph{registerY} = \emph{extend27\_\discretionary{}{}{}offset27}[\emph{registerZ}]

\paragraph{Description} The \emph{registerGh} is preloaded from the memory double word at the effective address.


\paragraph{Modifier(s)}
\noindent\begin{itemize}
\item variant (cf. variant in section \ref{par:modifier-variant} on page \pageref{par:modifier-variant})
\end{itemize}

\paragraph{Execution} {Reservation table \textsf{LSU.Y} (Section~\ref{sec:reservation-LSU.Y})}
~
{\begin{lstlisting}
stage ID:
  new offset = _SX_54(extend27_offset27);
  new variant = _ZX_2(variant);
stage RR:
  new base = GPR[registerZ];
  new target = GPR[registerY];
  new address = base + offset;
  new shift = (target & 31) * 8;
  new index = (target >> 5) & 3;
  new where = (registerGh << 2) + index;
  new dri = where;
stage E1:
  new argument1 = XVR[where];
  new bytemask = 255;
  new targetbytes = _ZX_256(bits2bytes(bytemask) << shift);
  CS.XMF = 1;
  new result1 = MEM.load(address, 0xFF & bytemask, variant, dri);
  result1 = (_ZX_256(result1 << shift) & targetbytes) | (argument1 & _ZX_256(~targetbytes));
stage CR:
  XVR[where] = result1;
\end{lstlisting}}
\newpage

\paragraph{Encoding} Format \textsf{LSU\_XPL8RBB} (Section~\ref{sec:format-LSU-XPL8RBB}), \verb|lsucode5 = 011|\\

\bitfields{ 1/P, 2/01, 3/011, 2/variant, 3/registerGi, 3/011, 2/11, 3/011, 1/--, 6/registerY, 6/registerZ }


\paragraph{Syntax} \texttt{xpld}\emph{variant} \emph{registerGi}, \emph{registerY} = [\emph{registerZ}]

\paragraph{Description} The \emph{registerGi} is preloaded from the memory double word at the effective address.


\paragraph{Modifier(s)}
\noindent\begin{itemize}
\item variant (cf. variant in section \ref{par:modifier-variant} on page \pageref{par:modifier-variant})
\end{itemize}

\paragraph{Execution} {Reservation table \textsf{LSU} (Section~\ref{sec:reservation-LSU})}
~
{\begin{lstlisting}
stage ID:
  new variant = _ZX_2(variant);
stage RR:
  new address = GPR[registerZ];
  new target = GPR[registerY];
  new shift = (target & 31) * 8;
  new index = (target >> 5) & 7;
  new where = (registerGi << 3) + index;
  new dri = where;
stage E1:
  new argument1 = XVR[where];
  new bytemask = 255;
  new targetbytes = _ZX_256(bits2bytes(bytemask) << shift);
  CS.XMF = 1;
  new result1 = MEM.load(address, 0xFF & bytemask, variant, dri);
  result1 = (_ZX_256(result1 << shift) & targetbytes) | (argument1 & _ZX_256(~targetbytes));
stage CR:
  XVR[where] = result1;
\end{lstlisting}}
\newpage

\paragraph{Encoding} Format \textsf{LSU\_XPL8RBB.O} (Section~\ref{sec:format-LSU-XPL8RBB.O}), \verb|lsucode5 = 011|\\

\bitfields{ 1/P, 2/00, 2/-- --, 27/offset27, 1/1, 2/01, 3/011, 2/variant, 3/registerGi, 3/011, 2/11, 3/011, 1/--, 6/registerY, 6/registerZ }


\paragraph{Syntax} \texttt{xpld}\emph{variant} \emph{registerGi}, \emph{registerY} = \emph{offset27}[\emph{registerZ}]

\paragraph{Description} The \emph{registerGi} is preloaded from the memory double word at the effective address.


\paragraph{Modifier(s)}
\noindent\begin{itemize}
\item variant (cf. variant in section \ref{par:modifier-variant} on page \pageref{par:modifier-variant})
\end{itemize}

\paragraph{Execution} {Reservation table \textsf{LSU.X} (Section~\ref{sec:reservation-LSU.X})}
~
{\begin{lstlisting}
stage ID:
  new offset = _SX_27(offset27);
  new variant = _ZX_2(variant);
stage RR:
  new base = GPR[registerZ];
  new target = GPR[registerY];
  new address = base + offset;
  new shift = (target & 31) * 8;
  new index = (target >> 5) & 7;
  new where = (registerGi << 3) + index;
  new dri = where;
stage E1:
  new argument1 = XVR[where];
  new bytemask = 255;
  new targetbytes = _ZX_256(bits2bytes(bytemask) << shift);
  CS.XMF = 1;
  new result1 = MEM.load(address, 0xFF & bytemask, variant, dri);
  result1 = (_ZX_256(result1 << shift) & targetbytes) | (argument1 & _ZX_256(~targetbytes));
stage CR:
  XVR[where] = result1;
\end{lstlisting}}
\newpage

\paragraph{Encoding} Format \textsf{LSU\_XPL8RBB.D} (Section~\ref{sec:format-LSU-XPL8RBB.D}), \verb|lsucode5 = 011|\\

\bitfields{ 1/P, 2/00, 2/-- --, 27/extend27, 1/1, 2/00, 2/-- --, 27/offset27, 1/1, 2/01, 3/011, 2/variant, 3/registerGi, 3/011, 2/11, 3/011, 1/--, 6/registerY, 6/registerZ }


\paragraph{Syntax} \texttt{xpld}\emph{variant} \emph{registerGi}, \emph{registerY} = \emph{extend27\_\discretionary{}{}{}offset27}[\emph{registerZ}]

\paragraph{Description} The \emph{registerGi} is preloaded from the memory double word at the effective address.


\paragraph{Modifier(s)}
\noindent\begin{itemize}
\item variant (cf. variant in section \ref{par:modifier-variant} on page \pageref{par:modifier-variant})
\end{itemize}

\paragraph{Execution} {Reservation table \textsf{LSU.Y} (Section~\ref{sec:reservation-LSU.Y})}
~
{\begin{lstlisting}
stage ID:
  new offset = _SX_54(extend27_offset27);
  new variant = _ZX_2(variant);
stage RR:
  new base = GPR[registerZ];
  new target = GPR[registerY];
  new address = base + offset;
  new shift = (target & 31) * 8;
  new index = (target >> 5) & 7;
  new where = (registerGi << 3) + index;
  new dri = where;
stage E1:
  new argument1 = XVR[where];
  new bytemask = 255;
  new targetbytes = _ZX_256(bits2bytes(bytemask) << shift);
  CS.XMF = 1;
  new result1 = MEM.load(address, 0xFF & bytemask, variant, dri);
  result1 = (_ZX_256(result1 << shift) & targetbytes) | (argument1 & _ZX_256(~targetbytes));
stage CR:
  XVR[where] = result1;
\end{lstlisting}}
\newpage

\paragraph{Encoding} Format \textsf{LSU\_XPL16RBB} (Section~\ref{sec:format-LSU-XPL16RBB}), \verb|lsucode5 = 011|\\

\bitfields{ 1/P, 2/01, 3/011, 2/variant, 2/registerGj, 4/0111, 2/11, 3/011, 1/--, 6/registerY, 6/registerZ }


\paragraph{Syntax} \texttt{xpld}\emph{variant} \emph{registerGj}, \emph{registerY} = [\emph{registerZ}]

\paragraph{Description} The \emph{registerGj} is preloaded from the memory double word at the effective address.


\paragraph{Modifier(s)}
\noindent\begin{itemize}
\item variant (cf. variant in section \ref{par:modifier-variant} on page \pageref{par:modifier-variant})
\end{itemize}

\paragraph{Execution} {Reservation table \textsf{LSU} (Section~\ref{sec:reservation-LSU})}
~
{\begin{lstlisting}
stage ID:
  new variant = _ZX_2(variant);
stage RR:
  new address = GPR[registerZ];
  new target = GPR[registerY];
  new shift = (target & 31) * 8;
  new index = (target >> 5) & 15;
  new where = (registerGj << 4) + index;
  new dri = where;
stage E1:
  new argument1 = XVR[where];
  new bytemask = 255;
  new targetbytes = _ZX_256(bits2bytes(bytemask) << shift);
  CS.XMF = 1;
  new result1 = MEM.load(address, 0xFF & bytemask, variant, dri);
  result1 = (_ZX_256(result1 << shift) & targetbytes) | (argument1 & _ZX_256(~targetbytes));
stage CR:
  XVR[where] = result1;
\end{lstlisting}}
\newpage

\paragraph{Encoding} Format \textsf{LSU\_XPL16RBB.O} (Section~\ref{sec:format-LSU-XPL16RBB.O}), \verb|lsucode5 = 011|\\

\bitfields{ 1/P, 2/00, 2/-- --, 27/offset27, 1/1, 2/01, 3/011, 2/variant, 2/registerGj, 4/0111, 2/11, 3/011, 1/--, 6/registerY, 6/registerZ }


\paragraph{Syntax} \texttt{xpld}\emph{variant} \emph{registerGj}, \emph{registerY} = \emph{offset27}[\emph{registerZ}]

\paragraph{Description} The \emph{registerGj} is preloaded from the memory double word at the effective address.


\paragraph{Modifier(s)}
\noindent\begin{itemize}
\item variant (cf. variant in section \ref{par:modifier-variant} on page \pageref{par:modifier-variant})
\end{itemize}

\paragraph{Execution} {Reservation table \textsf{LSU.X} (Section~\ref{sec:reservation-LSU.X})}
~
{\begin{lstlisting}
stage ID:
  new offset = _SX_27(offset27);
  new variant = _ZX_2(variant);
stage RR:
  new base = GPR[registerZ];
  new target = GPR[registerY];
  new address = base + offset;
  new shift = (target & 31) * 8;
  new index = (target >> 5) & 15;
  new where = (registerGj << 4) + index;
  new dri = where;
stage E1:
  new argument1 = XVR[where];
  new bytemask = 255;
  new targetbytes = _ZX_256(bits2bytes(bytemask) << shift);
  CS.XMF = 1;
  new result1 = MEM.load(address, 0xFF & bytemask, variant, dri);
  result1 = (_ZX_256(result1 << shift) & targetbytes) | (argument1 & _ZX_256(~targetbytes));
stage CR:
  XVR[where] = result1;
\end{lstlisting}}
\newpage

\paragraph{Encoding} Format \textsf{LSU\_XPL16RBB.D} (Section~\ref{sec:format-LSU-XPL16RBB.D}), \verb|lsucode5 = 011|\\

\bitfields{ 1/P, 2/00, 2/-- --, 27/extend27, 1/1, 2/00, 2/-- --, 27/offset27, 1/1, 2/01, 3/011, 2/variant, 2/registerGj, 4/0111, 2/11, 3/011, 1/--, 6/registerY, 6/registerZ }


\paragraph{Syntax} \texttt{xpld}\emph{variant} \emph{registerGj}, \emph{registerY} = \emph{extend27\_\discretionary{}{}{}offset27}[\emph{registerZ}]

\paragraph{Description} The \emph{registerGj} is preloaded from the memory double word at the effective address.


\paragraph{Modifier(s)}
\noindent\begin{itemize}
\item variant (cf. variant in section \ref{par:modifier-variant} on page \pageref{par:modifier-variant})
\end{itemize}

\paragraph{Execution} {Reservation table \textsf{LSU.Y} (Section~\ref{sec:reservation-LSU.Y})}
~
{\begin{lstlisting}
stage ID:
  new offset = _SX_54(extend27_offset27);
  new variant = _ZX_2(variant);
stage RR:
  new base = GPR[registerZ];
  new target = GPR[registerY];
  new address = base + offset;
  new shift = (target & 31) * 8;
  new index = (target >> 5) & 15;
  new where = (registerGj << 4) + index;
  new dri = where;
stage E1:
  new argument1 = XVR[where];
  new bytemask = 255;
  new targetbytes = _ZX_256(bits2bytes(bytemask) << shift);
  CS.XMF = 1;
  new result1 = MEM.load(address, 0xFF & bytemask, variant, dri);
  result1 = (_ZX_256(result1 << shift) & targetbytes) | (argument1 & _ZX_256(~targetbytes));
stage CR:
  XVR[where] = result1;
\end{lstlisting}}
\newpage

\paragraph{Encoding} Format \textsf{LSU\_XPL32RBB} (Section~\ref{sec:format-LSU-XPL32RBB}), \verb|lsucode5 = 011|\\

\bitfields{ 1/P, 2/01, 3/011, 2/variant, 1/registerGk, 5/01111, 2/11, 3/011, 1/--, 6/registerY, 6/registerZ }


\paragraph{Syntax} \texttt{xpld}\emph{variant} \emph{registerGk}, \emph{registerY} = [\emph{registerZ}]

\paragraph{Description} The \emph{registerGk} is preloaded from the memory double word at the effective address.


\paragraph{Modifier(s)}
\noindent\begin{itemize}
\item variant (cf. variant in section \ref{par:modifier-variant} on page \pageref{par:modifier-variant})
\end{itemize}

\paragraph{Execution} {Reservation table \textsf{LSU} (Section~\ref{sec:reservation-LSU})}
~
{\begin{lstlisting}
stage ID:
  new variant = _ZX_2(variant);
stage RR:
  new address = GPR[registerZ];
  new target = GPR[registerY];
  new shift = (target & 31) * 8;
  new index = (target >> 5) & 31;
  new where = (registerGk << 5) + index;
  new dri = where;
stage E1:
  new argument1 = XVR[where];
  new bytemask = 255;
  new targetbytes = _ZX_256(bits2bytes(bytemask) << shift);
  CS.XMF = 1;
  new result1 = MEM.load(address, 0xFF & bytemask, variant, dri);
  result1 = (_ZX_256(result1 << shift) & targetbytes) | (argument1 & _ZX_256(~targetbytes));
stage CR:
  XVR[where] = result1;
\end{lstlisting}}
\newpage

\paragraph{Encoding} Format \textsf{LSU\_XPL32RBB.O} (Section~\ref{sec:format-LSU-XPL32RBB.O}), \verb|lsucode5 = 011|\\

\bitfields{ 1/P, 2/00, 2/-- --, 27/offset27, 1/1, 2/01, 3/011, 2/variant, 1/registerGk, 5/01111, 2/11, 3/011, 1/--, 6/registerY, 6/registerZ }


\paragraph{Syntax} \texttt{xpld}\emph{variant} \emph{registerGk}, \emph{registerY} = \emph{offset27}[\emph{registerZ}]

\paragraph{Description} The \emph{registerGk} is preloaded from the memory double word at the effective address.


\paragraph{Modifier(s)}
\noindent\begin{itemize}
\item variant (cf. variant in section \ref{par:modifier-variant} on page \pageref{par:modifier-variant})
\end{itemize}

\paragraph{Execution} {Reservation table \textsf{LSU.X} (Section~\ref{sec:reservation-LSU.X})}
~
{\begin{lstlisting}
stage ID:
  new offset = _SX_27(offset27);
  new variant = _ZX_2(variant);
stage RR:
  new base = GPR[registerZ];
  new target = GPR[registerY];
  new address = base + offset;
  new shift = (target & 31) * 8;
  new index = (target >> 5) & 31;
  new where = (registerGk << 5) + index;
  new dri = where;
stage E1:
  new argument1 = XVR[where];
  new bytemask = 255;
  new targetbytes = _ZX_256(bits2bytes(bytemask) << shift);
  CS.XMF = 1;
  new result1 = MEM.load(address, 0xFF & bytemask, variant, dri);
  result1 = (_ZX_256(result1 << shift) & targetbytes) | (argument1 & _ZX_256(~targetbytes));
stage CR:
  XVR[where] = result1;
\end{lstlisting}}
\newpage

\paragraph{Encoding} Format \textsf{LSU\_XPL32RBB.D} (Section~\ref{sec:format-LSU-XPL32RBB.D}), \verb|lsucode5 = 011|\\

\bitfields{ 1/P, 2/00, 2/-- --, 27/extend27, 1/1, 2/00, 2/-- --, 27/offset27, 1/1, 2/01, 3/011, 2/variant, 1/registerGk, 5/01111, 2/11, 3/011, 1/--, 6/registerY, 6/registerZ }


\paragraph{Syntax} \texttt{xpld}\emph{variant} \emph{registerGk}, \emph{registerY} = \emph{extend27\_\discretionary{}{}{}offset27}[\emph{registerZ}]

\paragraph{Description} The \emph{registerGk} is preloaded from the memory double word at the effective address.


\paragraph{Modifier(s)}
\noindent\begin{itemize}
\item variant (cf. variant in section \ref{par:modifier-variant} on page \pageref{par:modifier-variant})
\end{itemize}

\paragraph{Execution} {Reservation table \textsf{LSU.Y} (Section~\ref{sec:reservation-LSU.Y})}
~
{\begin{lstlisting}
stage ID:
  new offset = _SX_54(extend27_offset27);
  new variant = _ZX_2(variant);
stage RR:
  new base = GPR[registerZ];
  new target = GPR[registerY];
  new address = base + offset;
  new shift = (target & 31) * 8;
  new index = (target >> 5) & 31;
  new where = (registerGk << 5) + index;
  new dri = where;
stage E1:
  new argument1 = XVR[where];
  new bytemask = 255;
  new targetbytes = _ZX_256(bits2bytes(bytemask) << shift);
  CS.XMF = 1;
  new result1 = MEM.load(address, 0xFF & bytemask, variant, dri);
  result1 = (_ZX_256(result1 << shift) & targetbytes) | (argument1 & _ZX_256(~targetbytes));
stage CR:
  XVR[where] = result1;
\end{lstlisting}}
\newpage

\paragraph{Encoding} Format \textsf{LSU\_XPL64RBB} (Section~\ref{sec:format-LSU-XPL64RBB}), \verb|lsucode5 = 011|\\

\bitfields{ 1/P, 2/01, 3/011, 2/variant, 1/registerGl, 5/11111, 2/11, 3/011, 1/--, 6/registerY, 6/registerZ }


\paragraph{Syntax} \texttt{xpld}\emph{variant} \emph{registerGl}, \emph{registerY} = [\emph{registerZ}]

\paragraph{Description} The \emph{registerGl} is preloaded from the memory double word at the effective address.


\paragraph{Modifier(s)}
\noindent\begin{itemize}
\item variant (cf. variant in section \ref{par:modifier-variant} on page \pageref{par:modifier-variant})
\end{itemize}

\paragraph{Execution} {Reservation table \textsf{LSU} (Section~\ref{sec:reservation-LSU})}
~
{\begin{lstlisting}
stage ID:
  new variant = _ZX_2(variant);
stage RR:
  new address = GPR[registerZ];
  new target = GPR[registerY];
  new shift = (target & 31) * 8;
  new index = (target >> 5) & 63;
  new where = (registerGl << 6) + index;
  new dri = where;
stage E1:
  new argument1 = XVR[where];
  new bytemask = 255;
  new targetbytes = _ZX_256(bits2bytes(bytemask) << shift);
  CS.XMF = 1;
  new result1 = MEM.load(address, 0xFF & bytemask, variant, dri);
  result1 = (_ZX_256(result1 << shift) & targetbytes) | (argument1 & _ZX_256(~targetbytes));
stage CR:
  XVR[where] = result1;
\end{lstlisting}}
\newpage

\paragraph{Encoding} Format \textsf{LSU\_XPL64RBB.O} (Section~\ref{sec:format-LSU-XPL64RBB.O}), \verb|lsucode5 = 011|\\

\bitfields{ 1/P, 2/00, 2/-- --, 27/offset27, 1/1, 2/01, 3/011, 2/variant, 1/registerGl, 5/11111, 2/11, 3/011, 1/--, 6/registerY, 6/registerZ }


\paragraph{Syntax} \texttt{xpld}\emph{variant} \emph{registerGl}, \emph{registerY} = \emph{offset27}[\emph{registerZ}]

\paragraph{Description} The \emph{registerGl} is preloaded from the memory double word at the effective address.


\paragraph{Modifier(s)}
\noindent\begin{itemize}
\item variant (cf. variant in section \ref{par:modifier-variant} on page \pageref{par:modifier-variant})
\end{itemize}

\paragraph{Execution} {Reservation table \textsf{LSU.X} (Section~\ref{sec:reservation-LSU.X})}
~
{\begin{lstlisting}
stage ID:
  new offset = _SX_27(offset27);
  new variant = _ZX_2(variant);
stage RR:
  new base = GPR[registerZ];
  new target = GPR[registerY];
  new address = base + offset;
  new shift = (target & 31) * 8;
  new index = (target >> 5) & 63;
  new where = (registerGl << 6) + index;
  new dri = where;
stage E1:
  new argument1 = XVR[where];
  new bytemask = 255;
  new targetbytes = _ZX_256(bits2bytes(bytemask) << shift);
  CS.XMF = 1;
  new result1 = MEM.load(address, 0xFF & bytemask, variant, dri);
  result1 = (_ZX_256(result1 << shift) & targetbytes) | (argument1 & _ZX_256(~targetbytes));
stage CR:
  XVR[where] = result1;
\end{lstlisting}}
\newpage

\paragraph{Encoding} Format \textsf{LSU\_XPL64RBB.D} (Section~\ref{sec:format-LSU-XPL64RBB.D}), \verb|lsucode5 = 011|\\

\bitfields{ 1/P, 2/00, 2/-- --, 27/extend27, 1/1, 2/00, 2/-- --, 27/offset27, 1/1, 2/01, 3/011, 2/variant, 1/registerGl, 5/11111, 2/11, 3/011, 1/--, 6/registerY, 6/registerZ }


\paragraph{Syntax} \texttt{xpld}\emph{variant} \emph{registerGl}, \emph{registerY} = \emph{extend27\_\discretionary{}{}{}offset27}[\emph{registerZ}]

\paragraph{Description} The \emph{registerGl} is preloaded from the memory double word at the effective address.


\paragraph{Modifier(s)}
\noindent\begin{itemize}
\item variant (cf. variant in section \ref{par:modifier-variant} on page \pageref{par:modifier-variant})
\end{itemize}

\paragraph{Execution} {Reservation table \textsf{LSU.Y} (Section~\ref{sec:reservation-LSU.Y})}
~
{\begin{lstlisting}
stage ID:
  new offset = _SX_54(extend27_offset27);
  new variant = _ZX_2(variant);
stage RR:
  new base = GPR[registerZ];
  new target = GPR[registerY];
  new address = base + offset;
  new shift = (target & 31) * 8;
  new index = (target >> 5) & 63;
  new where = (registerGl << 6) + index;
  new dri = where;
stage E1:
  new argument1 = XVR[where];
  new bytemask = 255;
  new targetbytes = _ZX_256(bits2bytes(bytemask) << shift);
  CS.XMF = 1;
  new result1 = MEM.load(address, 0xFF & bytemask, variant, dri);
  result1 = (_ZX_256(result1 << shift) & targetbytes) | (argument1 & _ZX_256(~targetbytes));
stage CR:
  XVR[where] = result1;
\end{lstlisting}}
\newpage
\subsubsection[{\small XPLH: Extension Preload Half}]{XPLH\hfill Extension Preload Half}
\label{instr:XPLH}\index{xplh}
 
\paragraph{Encoding} Format \textsf{LSU\_XPL2RBB} (Section~\ref{sec:format-LSU-XPL2RBB}), \verb|lsucode5 = 001|\\

\bitfields{ 1/P, 2/01, 3/011, 2/variant, 5/registerGg, 1/0, 2/11, 3/001, 1/--, 6/registerY, 6/registerZ }


\paragraph{Syntax} \texttt{xplh}\emph{variant} \emph{registerGg}, \emph{registerY} = [\emph{registerZ}]

\paragraph{Description} The \emph{registerGg} is preloaded from the memory half word at the effective address.


\paragraph{Modifier(s)}
\noindent\begin{itemize}
\item variant (cf. variant in section \ref{par:modifier-variant} on page \pageref{par:modifier-variant})
\end{itemize}

\paragraph{Execution} {Reservation table \textsf{LSU} (Section~\ref{sec:reservation-LSU})}
~
{\begin{lstlisting}
stage ID:
  new variant = _ZX_2(variant);
stage RR:
  new address = GPR[registerZ];
  new target = GPR[registerY];
  new shift = (target & 31) * 8;
  new index = (target >> 5) & 1;
  new where = (registerGg << 1) + index;
  new dri = where;
stage E1:
  new argument1 = XVR[where];
  new bytemask = 3;
  new targetbytes = _ZX_256(bits2bytes(bytemask) << shift);
  CS.XMF = 1;
  new result1 = MEM.load(address, 0x3 & bytemask, variant, dri);
  result1 = (_ZX_256(result1 << shift) & targetbytes) | (argument1 & _ZX_256(~targetbytes));
stage CR:
  XVR[where] = result1;
\end{lstlisting}}
\newpage

\paragraph{Encoding} Format \textsf{LSU\_XPL2RBB.O} (Section~\ref{sec:format-LSU-XPL2RBB.O}), \verb|lsucode5 = 001|\\

\bitfields{ 1/P, 2/00, 2/-- --, 27/offset27, 1/1, 2/01, 3/011, 2/variant, 5/registerGg, 1/0, 2/11, 3/001, 1/--, 6/registerY, 6/registerZ }


\paragraph{Syntax} \texttt{xplh}\emph{variant} \emph{registerGg}, \emph{registerY} = \emph{offset27}[\emph{registerZ}]

\paragraph{Description} The \emph{registerGg} is preloaded from the memory half word at the effective address.


\paragraph{Modifier(s)}
\noindent\begin{itemize}
\item variant (cf. variant in section \ref{par:modifier-variant} on page \pageref{par:modifier-variant})
\end{itemize}

\paragraph{Execution} {Reservation table \textsf{LSU.X} (Section~\ref{sec:reservation-LSU.X})}
~
{\begin{lstlisting}
stage ID:
  new offset = _SX_27(offset27);
  new variant = _ZX_2(variant);
stage RR:
  new base = GPR[registerZ];
  new target = GPR[registerY];
  new address = base + offset;
  new shift = (target & 31) * 8;
  new index = (target >> 5) & 1;
  new where = (registerGg << 1) + index;
  new dri = where;
stage E1:
  new argument1 = XVR[where];
  new bytemask = 3;
  new targetbytes = _ZX_256(bits2bytes(bytemask) << shift);
  CS.XMF = 1;
  new result1 = MEM.load(address, 0x3 & bytemask, variant, dri);
  result1 = (_ZX_256(result1 << shift) & targetbytes) | (argument1 & _ZX_256(~targetbytes));
stage CR:
  XVR[where] = result1;
\end{lstlisting}}
\newpage

\paragraph{Encoding} Format \textsf{LSU\_XPL2RBB.D} (Section~\ref{sec:format-LSU-XPL2RBB.D}), \verb|lsucode5 = 001|\\

\bitfields{ 1/P, 2/00, 2/-- --, 27/extend27, 1/1, 2/00, 2/-- --, 27/offset27, 1/1, 2/01, 3/011, 2/variant, 5/registerGg, 1/0, 2/11, 3/001, 1/--, 6/registerY, 6/registerZ }


\paragraph{Syntax} \texttt{xplh}\emph{variant} \emph{registerGg}, \emph{registerY} = \emph{extend27\_\discretionary{}{}{}offset27}[\emph{registerZ}]

\paragraph{Description} The \emph{registerGg} is preloaded from the memory half word at the effective address.


\paragraph{Modifier(s)}
\noindent\begin{itemize}
\item variant (cf. variant in section \ref{par:modifier-variant} on page \pageref{par:modifier-variant})
\end{itemize}

\paragraph{Execution} {Reservation table \textsf{LSU.Y} (Section~\ref{sec:reservation-LSU.Y})}
~
{\begin{lstlisting}
stage ID:
  new offset = _SX_54(extend27_offset27);
  new variant = _ZX_2(variant);
stage RR:
  new base = GPR[registerZ];
  new target = GPR[registerY];
  new address = base + offset;
  new shift = (target & 31) * 8;
  new index = (target >> 5) & 1;
  new where = (registerGg << 1) + index;
  new dri = where;
stage E1:
  new argument1 = XVR[where];
  new bytemask = 3;
  new targetbytes = _ZX_256(bits2bytes(bytemask) << shift);
  CS.XMF = 1;
  new result1 = MEM.load(address, 0x3 & bytemask, variant, dri);
  result1 = (_ZX_256(result1 << shift) & targetbytes) | (argument1 & _ZX_256(~targetbytes));
stage CR:
  XVR[where] = result1;
\end{lstlisting}}
\newpage

\paragraph{Encoding} Format \textsf{LSU\_XPL4RBB} (Section~\ref{sec:format-LSU-XPL4RBB}), \verb|lsucode5 = 001|\\

\bitfields{ 1/P, 2/01, 3/011, 2/variant, 4/registerGh, 2/01, 2/11, 3/001, 1/--, 6/registerY, 6/registerZ }


\paragraph{Syntax} \texttt{xplh}\emph{variant} \emph{registerGh}, \emph{registerY} = [\emph{registerZ}]

\paragraph{Description} The \emph{registerGh} is preloaded from the memory half word at the effective address.


\paragraph{Modifier(s)}
\noindent\begin{itemize}
\item variant (cf. variant in section \ref{par:modifier-variant} on page \pageref{par:modifier-variant})
\end{itemize}

\paragraph{Execution} {Reservation table \textsf{LSU} (Section~\ref{sec:reservation-LSU})}
~
{\begin{lstlisting}
stage ID:
  new variant = _ZX_2(variant);
stage RR:
  new address = GPR[registerZ];
  new target = GPR[registerY];
  new shift = (target & 31) * 8;
  new index = (target >> 5) & 3;
  new where = (registerGh << 2) + index;
  new dri = where;
stage E1:
  new argument1 = XVR[where];
  new bytemask = 3;
  new targetbytes = _ZX_256(bits2bytes(bytemask) << shift);
  CS.XMF = 1;
  new result1 = MEM.load(address, 0x3 & bytemask, variant, dri);
  result1 = (_ZX_256(result1 << shift) & targetbytes) | (argument1 & _ZX_256(~targetbytes));
stage CR:
  XVR[where] = result1;
\end{lstlisting}}
\newpage

\paragraph{Encoding} Format \textsf{LSU\_XPL4RBB.O} (Section~\ref{sec:format-LSU-XPL4RBB.O}), \verb|lsucode5 = 001|\\

\bitfields{ 1/P, 2/00, 2/-- --, 27/offset27, 1/1, 2/01, 3/011, 2/variant, 4/registerGh, 2/01, 2/11, 3/001, 1/--, 6/registerY, 6/registerZ }


\paragraph{Syntax} \texttt{xplh}\emph{variant} \emph{registerGh}, \emph{registerY} = \emph{offset27}[\emph{registerZ}]

\paragraph{Description} The \emph{registerGh} is preloaded from the memory half word at the effective address.


\paragraph{Modifier(s)}
\noindent\begin{itemize}
\item variant (cf. variant in section \ref{par:modifier-variant} on page \pageref{par:modifier-variant})
\end{itemize}

\paragraph{Execution} {Reservation table \textsf{LSU.X} (Section~\ref{sec:reservation-LSU.X})}
~
{\begin{lstlisting}
stage ID:
  new offset = _SX_27(offset27);
  new variant = _ZX_2(variant);
stage RR:
  new base = GPR[registerZ];
  new target = GPR[registerY];
  new address = base + offset;
  new shift = (target & 31) * 8;
  new index = (target >> 5) & 3;
  new where = (registerGh << 2) + index;
  new dri = where;
stage E1:
  new argument1 = XVR[where];
  new bytemask = 3;
  new targetbytes = _ZX_256(bits2bytes(bytemask) << shift);
  CS.XMF = 1;
  new result1 = MEM.load(address, 0x3 & bytemask, variant, dri);
  result1 = (_ZX_256(result1 << shift) & targetbytes) | (argument1 & _ZX_256(~targetbytes));
stage CR:
  XVR[where] = result1;
\end{lstlisting}}
\newpage

\paragraph{Encoding} Format \textsf{LSU\_XPL4RBB.D} (Section~\ref{sec:format-LSU-XPL4RBB.D}), \verb|lsucode5 = 001|\\

\bitfields{ 1/P, 2/00, 2/-- --, 27/extend27, 1/1, 2/00, 2/-- --, 27/offset27, 1/1, 2/01, 3/011, 2/variant, 4/registerGh, 2/01, 2/11, 3/001, 1/--, 6/registerY, 6/registerZ }


\paragraph{Syntax} \texttt{xplh}\emph{variant} \emph{registerGh}, \emph{registerY} = \emph{extend27\_\discretionary{}{}{}offset27}[\emph{registerZ}]

\paragraph{Description} The \emph{registerGh} is preloaded from the memory half word at the effective address.


\paragraph{Modifier(s)}
\noindent\begin{itemize}
\item variant (cf. variant in section \ref{par:modifier-variant} on page \pageref{par:modifier-variant})
\end{itemize}

\paragraph{Execution} {Reservation table \textsf{LSU.Y} (Section~\ref{sec:reservation-LSU.Y})}
~
{\begin{lstlisting}
stage ID:
  new offset = _SX_54(extend27_offset27);
  new variant = _ZX_2(variant);
stage RR:
  new base = GPR[registerZ];
  new target = GPR[registerY];
  new address = base + offset;
  new shift = (target & 31) * 8;
  new index = (target >> 5) & 3;
  new where = (registerGh << 2) + index;
  new dri = where;
stage E1:
  new argument1 = XVR[where];
  new bytemask = 3;
  new targetbytes = _ZX_256(bits2bytes(bytemask) << shift);
  CS.XMF = 1;
  new result1 = MEM.load(address, 0x3 & bytemask, variant, dri);
  result1 = (_ZX_256(result1 << shift) & targetbytes) | (argument1 & _ZX_256(~targetbytes));
stage CR:
  XVR[where] = result1;
\end{lstlisting}}
\newpage

\paragraph{Encoding} Format \textsf{LSU\_XPL8RBB} (Section~\ref{sec:format-LSU-XPL8RBB}), \verb|lsucode5 = 001|\\

\bitfields{ 1/P, 2/01, 3/011, 2/variant, 3/registerGi, 3/011, 2/11, 3/001, 1/--, 6/registerY, 6/registerZ }


\paragraph{Syntax} \texttt{xplh}\emph{variant} \emph{registerGi}, \emph{registerY} = [\emph{registerZ}]

\paragraph{Description} The \emph{registerGi} is preloaded from the memory half word at the effective address.


\paragraph{Modifier(s)}
\noindent\begin{itemize}
\item variant (cf. variant in section \ref{par:modifier-variant} on page \pageref{par:modifier-variant})
\end{itemize}

\paragraph{Execution} {Reservation table \textsf{LSU} (Section~\ref{sec:reservation-LSU})}
~
{\begin{lstlisting}
stage ID:
  new variant = _ZX_2(variant);
stage RR:
  new address = GPR[registerZ];
  new target = GPR[registerY];
  new shift = (target & 31) * 8;
  new index = (target >> 5) & 7;
  new where = (registerGi << 3) + index;
  new dri = where;
stage E1:
  new argument1 = XVR[where];
  new bytemask = 3;
  new targetbytes = _ZX_256(bits2bytes(bytemask) << shift);
  CS.XMF = 1;
  new result1 = MEM.load(address, 0x3 & bytemask, variant, dri);
  result1 = (_ZX_256(result1 << shift) & targetbytes) | (argument1 & _ZX_256(~targetbytes));
stage CR:
  XVR[where] = result1;
\end{lstlisting}}
\newpage

\paragraph{Encoding} Format \textsf{LSU\_XPL8RBB.O} (Section~\ref{sec:format-LSU-XPL8RBB.O}), \verb|lsucode5 = 001|\\

\bitfields{ 1/P, 2/00, 2/-- --, 27/offset27, 1/1, 2/01, 3/011, 2/variant, 3/registerGi, 3/011, 2/11, 3/001, 1/--, 6/registerY, 6/registerZ }


\paragraph{Syntax} \texttt{xplh}\emph{variant} \emph{registerGi}, \emph{registerY} = \emph{offset27}[\emph{registerZ}]

\paragraph{Description} The \emph{registerGi} is preloaded from the memory half word at the effective address.


\paragraph{Modifier(s)}
\noindent\begin{itemize}
\item variant (cf. variant in section \ref{par:modifier-variant} on page \pageref{par:modifier-variant})
\end{itemize}

\paragraph{Execution} {Reservation table \textsf{LSU.X} (Section~\ref{sec:reservation-LSU.X})}
~
{\begin{lstlisting}
stage ID:
  new offset = _SX_27(offset27);
  new variant = _ZX_2(variant);
stage RR:
  new base = GPR[registerZ];
  new target = GPR[registerY];
  new address = base + offset;
  new shift = (target & 31) * 8;
  new index = (target >> 5) & 7;
  new where = (registerGi << 3) + index;
  new dri = where;
stage E1:
  new argument1 = XVR[where];
  new bytemask = 3;
  new targetbytes = _ZX_256(bits2bytes(bytemask) << shift);
  CS.XMF = 1;
  new result1 = MEM.load(address, 0x3 & bytemask, variant, dri);
  result1 = (_ZX_256(result1 << shift) & targetbytes) | (argument1 & _ZX_256(~targetbytes));
stage CR:
  XVR[where] = result1;
\end{lstlisting}}
\newpage

\paragraph{Encoding} Format \textsf{LSU\_XPL8RBB.D} (Section~\ref{sec:format-LSU-XPL8RBB.D}), \verb|lsucode5 = 001|\\

\bitfields{ 1/P, 2/00, 2/-- --, 27/extend27, 1/1, 2/00, 2/-- --, 27/offset27, 1/1, 2/01, 3/011, 2/variant, 3/registerGi, 3/011, 2/11, 3/001, 1/--, 6/registerY, 6/registerZ }


\paragraph{Syntax} \texttt{xplh}\emph{variant} \emph{registerGi}, \emph{registerY} = \emph{extend27\_\discretionary{}{}{}offset27}[\emph{registerZ}]

\paragraph{Description} The \emph{registerGi} is preloaded from the memory half word at the effective address.


\paragraph{Modifier(s)}
\noindent\begin{itemize}
\item variant (cf. variant in section \ref{par:modifier-variant} on page \pageref{par:modifier-variant})
\end{itemize}

\paragraph{Execution} {Reservation table \textsf{LSU.Y} (Section~\ref{sec:reservation-LSU.Y})}
~
{\begin{lstlisting}
stage ID:
  new offset = _SX_54(extend27_offset27);
  new variant = _ZX_2(variant);
stage RR:
  new base = GPR[registerZ];
  new target = GPR[registerY];
  new address = base + offset;
  new shift = (target & 31) * 8;
  new index = (target >> 5) & 7;
  new where = (registerGi << 3) + index;
  new dri = where;
stage E1:
  new argument1 = XVR[where];
  new bytemask = 3;
  new targetbytes = _ZX_256(bits2bytes(bytemask) << shift);
  CS.XMF = 1;
  new result1 = MEM.load(address, 0x3 & bytemask, variant, dri);
  result1 = (_ZX_256(result1 << shift) & targetbytes) | (argument1 & _ZX_256(~targetbytes));
stage CR:
  XVR[where] = result1;
\end{lstlisting}}
\newpage

\paragraph{Encoding} Format \textsf{LSU\_XPL16RBB} (Section~\ref{sec:format-LSU-XPL16RBB}), \verb|lsucode5 = 001|\\

\bitfields{ 1/P, 2/01, 3/011, 2/variant, 2/registerGj, 4/0111, 2/11, 3/001, 1/--, 6/registerY, 6/registerZ }


\paragraph{Syntax} \texttt{xplh}\emph{variant} \emph{registerGj}, \emph{registerY} = [\emph{registerZ}]

\paragraph{Description} The \emph{registerGj} is preloaded from the memory half word at the effective address.


\paragraph{Modifier(s)}
\noindent\begin{itemize}
\item variant (cf. variant in section \ref{par:modifier-variant} on page \pageref{par:modifier-variant})
\end{itemize}

\paragraph{Execution} {Reservation table \textsf{LSU} (Section~\ref{sec:reservation-LSU})}
~
{\begin{lstlisting}
stage ID:
  new variant = _ZX_2(variant);
stage RR:
  new address = GPR[registerZ];
  new target = GPR[registerY];
  new shift = (target & 31) * 8;
  new index = (target >> 5) & 15;
  new where = (registerGj << 4) + index;
  new dri = where;
stage E1:
  new argument1 = XVR[where];
  new bytemask = 3;
  new targetbytes = _ZX_256(bits2bytes(bytemask) << shift);
  CS.XMF = 1;
  new result1 = MEM.load(address, 0x3 & bytemask, variant, dri);
  result1 = (_ZX_256(result1 << shift) & targetbytes) | (argument1 & _ZX_256(~targetbytes));
stage CR:
  XVR[where] = result1;
\end{lstlisting}}
\newpage

\paragraph{Encoding} Format \textsf{LSU\_XPL16RBB.O} (Section~\ref{sec:format-LSU-XPL16RBB.O}), \verb|lsucode5 = 001|\\

\bitfields{ 1/P, 2/00, 2/-- --, 27/offset27, 1/1, 2/01, 3/011, 2/variant, 2/registerGj, 4/0111, 2/11, 3/001, 1/--, 6/registerY, 6/registerZ }


\paragraph{Syntax} \texttt{xplh}\emph{variant} \emph{registerGj}, \emph{registerY} = \emph{offset27}[\emph{registerZ}]

\paragraph{Description} The \emph{registerGj} is preloaded from the memory half word at the effective address.


\paragraph{Modifier(s)}
\noindent\begin{itemize}
\item variant (cf. variant in section \ref{par:modifier-variant} on page \pageref{par:modifier-variant})
\end{itemize}

\paragraph{Execution} {Reservation table \textsf{LSU.X} (Section~\ref{sec:reservation-LSU.X})}
~
{\begin{lstlisting}
stage ID:
  new offset = _SX_27(offset27);
  new variant = _ZX_2(variant);
stage RR:
  new base = GPR[registerZ];
  new target = GPR[registerY];
  new address = base + offset;
  new shift = (target & 31) * 8;
  new index = (target >> 5) & 15;
  new where = (registerGj << 4) + index;
  new dri = where;
stage E1:
  new argument1 = XVR[where];
  new bytemask = 3;
  new targetbytes = _ZX_256(bits2bytes(bytemask) << shift);
  CS.XMF = 1;
  new result1 = MEM.load(address, 0x3 & bytemask, variant, dri);
  result1 = (_ZX_256(result1 << shift) & targetbytes) | (argument1 & _ZX_256(~targetbytes));
stage CR:
  XVR[where] = result1;
\end{lstlisting}}
\newpage

\paragraph{Encoding} Format \textsf{LSU\_XPL16RBB.D} (Section~\ref{sec:format-LSU-XPL16RBB.D}), \verb|lsucode5 = 001|\\

\bitfields{ 1/P, 2/00, 2/-- --, 27/extend27, 1/1, 2/00, 2/-- --, 27/offset27, 1/1, 2/01, 3/011, 2/variant, 2/registerGj, 4/0111, 2/11, 3/001, 1/--, 6/registerY, 6/registerZ }


\paragraph{Syntax} \texttt{xplh}\emph{variant} \emph{registerGj}, \emph{registerY} = \emph{extend27\_\discretionary{}{}{}offset27}[\emph{registerZ}]

\paragraph{Description} The \emph{registerGj} is preloaded from the memory half word at the effective address.


\paragraph{Modifier(s)}
\noindent\begin{itemize}
\item variant (cf. variant in section \ref{par:modifier-variant} on page \pageref{par:modifier-variant})
\end{itemize}

\paragraph{Execution} {Reservation table \textsf{LSU.Y} (Section~\ref{sec:reservation-LSU.Y})}
~
{\begin{lstlisting}
stage ID:
  new offset = _SX_54(extend27_offset27);
  new variant = _ZX_2(variant);
stage RR:
  new base = GPR[registerZ];
  new target = GPR[registerY];
  new address = base + offset;
  new shift = (target & 31) * 8;
  new index = (target >> 5) & 15;
  new where = (registerGj << 4) + index;
  new dri = where;
stage E1:
  new argument1 = XVR[where];
  new bytemask = 3;
  new targetbytes = _ZX_256(bits2bytes(bytemask) << shift);
  CS.XMF = 1;
  new result1 = MEM.load(address, 0x3 & bytemask, variant, dri);
  result1 = (_ZX_256(result1 << shift) & targetbytes) | (argument1 & _ZX_256(~targetbytes));
stage CR:
  XVR[where] = result1;
\end{lstlisting}}
\newpage

\paragraph{Encoding} Format \textsf{LSU\_XPL32RBB} (Section~\ref{sec:format-LSU-XPL32RBB}), \verb|lsucode5 = 001|\\

\bitfields{ 1/P, 2/01, 3/011, 2/variant, 1/registerGk, 5/01111, 2/11, 3/001, 1/--, 6/registerY, 6/registerZ }


\paragraph{Syntax} \texttt{xplh}\emph{variant} \emph{registerGk}, \emph{registerY} = [\emph{registerZ}]

\paragraph{Description} The \emph{registerGk} is preloaded from the memory half word at the effective address.


\paragraph{Modifier(s)}
\noindent\begin{itemize}
\item variant (cf. variant in section \ref{par:modifier-variant} on page \pageref{par:modifier-variant})
\end{itemize}

\paragraph{Execution} {Reservation table \textsf{LSU} (Section~\ref{sec:reservation-LSU})}
~
{\begin{lstlisting}
stage ID:
  new variant = _ZX_2(variant);
stage RR:
  new address = GPR[registerZ];
  new target = GPR[registerY];
  new shift = (target & 31) * 8;
  new index = (target >> 5) & 31;
  new where = (registerGk << 5) + index;
  new dri = where;
stage E1:
  new argument1 = XVR[where];
  new bytemask = 3;
  new targetbytes = _ZX_256(bits2bytes(bytemask) << shift);
  CS.XMF = 1;
  new result1 = MEM.load(address, 0x3 & bytemask, variant, dri);
  result1 = (_ZX_256(result1 << shift) & targetbytes) | (argument1 & _ZX_256(~targetbytes));
stage CR:
  XVR[where] = result1;
\end{lstlisting}}
\newpage

\paragraph{Encoding} Format \textsf{LSU\_XPL32RBB.O} (Section~\ref{sec:format-LSU-XPL32RBB.O}), \verb|lsucode5 = 001|\\

\bitfields{ 1/P, 2/00, 2/-- --, 27/offset27, 1/1, 2/01, 3/011, 2/variant, 1/registerGk, 5/01111, 2/11, 3/001, 1/--, 6/registerY, 6/registerZ }


\paragraph{Syntax} \texttt{xplh}\emph{variant} \emph{registerGk}, \emph{registerY} = \emph{offset27}[\emph{registerZ}]

\paragraph{Description} The \emph{registerGk} is preloaded from the memory half word at the effective address.


\paragraph{Modifier(s)}
\noindent\begin{itemize}
\item variant (cf. variant in section \ref{par:modifier-variant} on page \pageref{par:modifier-variant})
\end{itemize}

\paragraph{Execution} {Reservation table \textsf{LSU.X} (Section~\ref{sec:reservation-LSU.X})}
~
{\begin{lstlisting}
stage ID:
  new offset = _SX_27(offset27);
  new variant = _ZX_2(variant);
stage RR:
  new base = GPR[registerZ];
  new target = GPR[registerY];
  new address = base + offset;
  new shift = (target & 31) * 8;
  new index = (target >> 5) & 31;
  new where = (registerGk << 5) + index;
  new dri = where;
stage E1:
  new argument1 = XVR[where];
  new bytemask = 3;
  new targetbytes = _ZX_256(bits2bytes(bytemask) << shift);
  CS.XMF = 1;
  new result1 = MEM.load(address, 0x3 & bytemask, variant, dri);
  result1 = (_ZX_256(result1 << shift) & targetbytes) | (argument1 & _ZX_256(~targetbytes));
stage CR:
  XVR[where] = result1;
\end{lstlisting}}
\newpage

\paragraph{Encoding} Format \textsf{LSU\_XPL32RBB.D} (Section~\ref{sec:format-LSU-XPL32RBB.D}), \verb|lsucode5 = 001|\\

\bitfields{ 1/P, 2/00, 2/-- --, 27/extend27, 1/1, 2/00, 2/-- --, 27/offset27, 1/1, 2/01, 3/011, 2/variant, 1/registerGk, 5/01111, 2/11, 3/001, 1/--, 6/registerY, 6/registerZ }


\paragraph{Syntax} \texttt{xplh}\emph{variant} \emph{registerGk}, \emph{registerY} = \emph{extend27\_\discretionary{}{}{}offset27}[\emph{registerZ}]

\paragraph{Description} The \emph{registerGk} is preloaded from the memory half word at the effective address.


\paragraph{Modifier(s)}
\noindent\begin{itemize}
\item variant (cf. variant in section \ref{par:modifier-variant} on page \pageref{par:modifier-variant})
\end{itemize}

\paragraph{Execution} {Reservation table \textsf{LSU.Y} (Section~\ref{sec:reservation-LSU.Y})}
~
{\begin{lstlisting}
stage ID:
  new offset = _SX_54(extend27_offset27);
  new variant = _ZX_2(variant);
stage RR:
  new base = GPR[registerZ];
  new target = GPR[registerY];
  new address = base + offset;
  new shift = (target & 31) * 8;
  new index = (target >> 5) & 31;
  new where = (registerGk << 5) + index;
  new dri = where;
stage E1:
  new argument1 = XVR[where];
  new bytemask = 3;
  new targetbytes = _ZX_256(bits2bytes(bytemask) << shift);
  CS.XMF = 1;
  new result1 = MEM.load(address, 0x3 & bytemask, variant, dri);
  result1 = (_ZX_256(result1 << shift) & targetbytes) | (argument1 & _ZX_256(~targetbytes));
stage CR:
  XVR[where] = result1;
\end{lstlisting}}
\newpage

\paragraph{Encoding} Format \textsf{LSU\_XPL64RBB} (Section~\ref{sec:format-LSU-XPL64RBB}), \verb|lsucode5 = 001|\\

\bitfields{ 1/P, 2/01, 3/011, 2/variant, 1/registerGl, 5/11111, 2/11, 3/001, 1/--, 6/registerY, 6/registerZ }


\paragraph{Syntax} \texttt{xplh}\emph{variant} \emph{registerGl}, \emph{registerY} = [\emph{registerZ}]

\paragraph{Description} The \emph{registerGl} is preloaded from the memory half word at the effective address.


\paragraph{Modifier(s)}
\noindent\begin{itemize}
\item variant (cf. variant in section \ref{par:modifier-variant} on page \pageref{par:modifier-variant})
\end{itemize}

\paragraph{Execution} {Reservation table \textsf{LSU} (Section~\ref{sec:reservation-LSU})}
~
{\begin{lstlisting}
stage ID:
  new variant = _ZX_2(variant);
stage RR:
  new address = GPR[registerZ];
  new target = GPR[registerY];
  new shift = (target & 31) * 8;
  new index = (target >> 5) & 63;
  new where = (registerGl << 6) + index;
  new dri = where;
stage E1:
  new argument1 = XVR[where];
  new bytemask = 3;
  new targetbytes = _ZX_256(bits2bytes(bytemask) << shift);
  CS.XMF = 1;
  new result1 = MEM.load(address, 0x3 & bytemask, variant, dri);
  result1 = (_ZX_256(result1 << shift) & targetbytes) | (argument1 & _ZX_256(~targetbytes));
stage CR:
  XVR[where] = result1;
\end{lstlisting}}
\newpage

\paragraph{Encoding} Format \textsf{LSU\_XPL64RBB.O} (Section~\ref{sec:format-LSU-XPL64RBB.O}), \verb|lsucode5 = 001|\\

\bitfields{ 1/P, 2/00, 2/-- --, 27/offset27, 1/1, 2/01, 3/011, 2/variant, 1/registerGl, 5/11111, 2/11, 3/001, 1/--, 6/registerY, 6/registerZ }


\paragraph{Syntax} \texttt{xplh}\emph{variant} \emph{registerGl}, \emph{registerY} = \emph{offset27}[\emph{registerZ}]

\paragraph{Description} The \emph{registerGl} is preloaded from the memory half word at the effective address.


\paragraph{Modifier(s)}
\noindent\begin{itemize}
\item variant (cf. variant in section \ref{par:modifier-variant} on page \pageref{par:modifier-variant})
\end{itemize}

\paragraph{Execution} {Reservation table \textsf{LSU.X} (Section~\ref{sec:reservation-LSU.X})}
~
{\begin{lstlisting}
stage ID:
  new offset = _SX_27(offset27);
  new variant = _ZX_2(variant);
stage RR:
  new base = GPR[registerZ];
  new target = GPR[registerY];
  new address = base + offset;
  new shift = (target & 31) * 8;
  new index = (target >> 5) & 63;
  new where = (registerGl << 6) + index;
  new dri = where;
stage E1:
  new argument1 = XVR[where];
  new bytemask = 3;
  new targetbytes = _ZX_256(bits2bytes(bytemask) << shift);
  CS.XMF = 1;
  new result1 = MEM.load(address, 0x3 & bytemask, variant, dri);
  result1 = (_ZX_256(result1 << shift) & targetbytes) | (argument1 & _ZX_256(~targetbytes));
stage CR:
  XVR[where] = result1;
\end{lstlisting}}
\newpage

\paragraph{Encoding} Format \textsf{LSU\_XPL64RBB.D} (Section~\ref{sec:format-LSU-XPL64RBB.D}), \verb|lsucode5 = 001|\\

\bitfields{ 1/P, 2/00, 2/-- --, 27/extend27, 1/1, 2/00, 2/-- --, 27/offset27, 1/1, 2/01, 3/011, 2/variant, 1/registerGl, 5/11111, 2/11, 3/001, 1/--, 6/registerY, 6/registerZ }


\paragraph{Syntax} \texttt{xplh}\emph{variant} \emph{registerGl}, \emph{registerY} = \emph{extend27\_\discretionary{}{}{}offset27}[\emph{registerZ}]

\paragraph{Description} The \emph{registerGl} is preloaded from the memory half word at the effective address.


\paragraph{Modifier(s)}
\noindent\begin{itemize}
\item variant (cf. variant in section \ref{par:modifier-variant} on page \pageref{par:modifier-variant})
\end{itemize}

\paragraph{Execution} {Reservation table \textsf{LSU.Y} (Section~\ref{sec:reservation-LSU.Y})}
~
{\begin{lstlisting}
stage ID:
  new offset = _SX_54(extend27_offset27);
  new variant = _ZX_2(variant);
stage RR:
  new base = GPR[registerZ];
  new target = GPR[registerY];
  new address = base + offset;
  new shift = (target & 31) * 8;
  new index = (target >> 5) & 63;
  new where = (registerGl << 6) + index;
  new dri = where;
stage E1:
  new argument1 = XVR[where];
  new bytemask = 3;
  new targetbytes = _ZX_256(bits2bytes(bytemask) << shift);
  CS.XMF = 1;
  new result1 = MEM.load(address, 0x3 & bytemask, variant, dri);
  result1 = (_ZX_256(result1 << shift) & targetbytes) | (argument1 & _ZX_256(~targetbytes));
stage CR:
  XVR[where] = result1;
\end{lstlisting}}
\newpage
\subsubsection[{\small XPLO: Extension Preload Octuple Word}]{XPLO\hfill Extension Preload Octuple Word}
\label{instr:XPLO}\index{xplo}
 
\paragraph{Encoding} Format \textsf{LSU\_XPL2RBB} (Section~\ref{sec:format-LSU-XPL2RBB}), \verb|lsucode5 = 101|\\

\bitfields{ 1/P, 2/01, 3/011, 2/variant, 5/registerGg, 1/0, 2/11, 3/101, 1/--, 6/registerY, 6/registerZ }


\paragraph{Syntax} \texttt{xplo}\emph{variant} \emph{registerGg}, \emph{registerY} = [\emph{registerZ}]

\paragraph{Description} The \emph{registerGg} is preloaded from the memory octuple word at the effective address.


\paragraph{Modifier(s)}
\noindent\begin{itemize}
\item variant (cf. variant in section \ref{par:modifier-variant} on page \pageref{par:modifier-variant})
\end{itemize}

\paragraph{Execution} {Reservation table \textsf{LSU} (Section~\ref{sec:reservation-LSU})}
~
{\begin{lstlisting}
stage ID:
  new variant = _ZX_2(variant);
stage RR:
  new address = GPR[registerZ];
  new target = GPR[registerY];
  new shift = (target & 31) * 8;
  new index = (target >> 5) & 1;
  new where = (registerGg << 1) + index;
  new dri = where;
stage E1:
  new argument1 = XVR[where];
  new bytemask = 4294967295;
  new targetbytes = _ZX_256(bits2bytes(bytemask) << shift);
  CS.XMF = 1;
  new result1 = MEM.load(address, 0xFFFFFFFF & bytemask, variant, dri);
  result1 = (_ZX_256(result1 << shift) & targetbytes) | (argument1 & _ZX_256(~targetbytes));
stage CR:
  XVR[where] = result1;
\end{lstlisting}}
\newpage

\paragraph{Encoding} Format \textsf{LSU\_XPL2RBB.O} (Section~\ref{sec:format-LSU-XPL2RBB.O}), \verb|lsucode5 = 101|\\

\bitfields{ 1/P, 2/00, 2/-- --, 27/offset27, 1/1, 2/01, 3/011, 2/variant, 5/registerGg, 1/0, 2/11, 3/101, 1/--, 6/registerY, 6/registerZ }


\paragraph{Syntax} \texttt{xplo}\emph{variant} \emph{registerGg}, \emph{registerY} = \emph{offset27}[\emph{registerZ}]

\paragraph{Description} The \emph{registerGg} is preloaded from the memory octuple word at the effective address.


\paragraph{Modifier(s)}
\noindent\begin{itemize}
\item variant (cf. variant in section \ref{par:modifier-variant} on page \pageref{par:modifier-variant})
\end{itemize}

\paragraph{Execution} {Reservation table \textsf{LSU.X} (Section~\ref{sec:reservation-LSU.X})}
~
{\begin{lstlisting}
stage ID:
  new offset = _SX_27(offset27);
  new variant = _ZX_2(variant);
stage RR:
  new base = GPR[registerZ];
  new target = GPR[registerY];
  new address = base + offset;
  new shift = (target & 31) * 8;
  new index = (target >> 5) & 1;
  new where = (registerGg << 1) + index;
  new dri = where;
stage E1:
  new argument1 = XVR[where];
  new bytemask = 4294967295;
  new targetbytes = _ZX_256(bits2bytes(bytemask) << shift);
  CS.XMF = 1;
  new result1 = MEM.load(address, 0xFFFFFFFF & bytemask, variant, dri);
  result1 = (_ZX_256(result1 << shift) & targetbytes) | (argument1 & _ZX_256(~targetbytes));
stage CR:
  XVR[where] = result1;
\end{lstlisting}}
\newpage

\paragraph{Encoding} Format \textsf{LSU\_XPL2RBB.D} (Section~\ref{sec:format-LSU-XPL2RBB.D}), \verb|lsucode5 = 101|\\

\bitfields{ 1/P, 2/00, 2/-- --, 27/extend27, 1/1, 2/00, 2/-- --, 27/offset27, 1/1, 2/01, 3/011, 2/variant, 5/registerGg, 1/0, 2/11, 3/101, 1/--, 6/registerY, 6/registerZ }


\paragraph{Syntax} \texttt{xplo}\emph{variant} \emph{registerGg}, \emph{registerY} = \emph{extend27\_\discretionary{}{}{}offset27}[\emph{registerZ}]

\paragraph{Description} The \emph{registerGg} is preloaded from the memory octuple word at the effective address.


\paragraph{Modifier(s)}
\noindent\begin{itemize}
\item variant (cf. variant in section \ref{par:modifier-variant} on page \pageref{par:modifier-variant})
\end{itemize}

\paragraph{Execution} {Reservation table \textsf{LSU.Y} (Section~\ref{sec:reservation-LSU.Y})}
~
{\begin{lstlisting}
stage ID:
  new offset = _SX_54(extend27_offset27);
  new variant = _ZX_2(variant);
stage RR:
  new base = GPR[registerZ];
  new target = GPR[registerY];
  new address = base + offset;
  new shift = (target & 31) * 8;
  new index = (target >> 5) & 1;
  new where = (registerGg << 1) + index;
  new dri = where;
stage E1:
  new argument1 = XVR[where];
  new bytemask = 4294967295;
  new targetbytes = _ZX_256(bits2bytes(bytemask) << shift);
  CS.XMF = 1;
  new result1 = MEM.load(address, 0xFFFFFFFF & bytemask, variant, dri);
  result1 = (_ZX_256(result1 << shift) & targetbytes) | (argument1 & _ZX_256(~targetbytes));
stage CR:
  XVR[where] = result1;
\end{lstlisting}}
\newpage

\paragraph{Encoding} Format \textsf{LSU\_XPL4RBB} (Section~\ref{sec:format-LSU-XPL4RBB}), \verb|lsucode5 = 101|\\

\bitfields{ 1/P, 2/01, 3/011, 2/variant, 4/registerGh, 2/01, 2/11, 3/101, 1/--, 6/registerY, 6/registerZ }


\paragraph{Syntax} \texttt{xplo}\emph{variant} \emph{registerGh}, \emph{registerY} = [\emph{registerZ}]

\paragraph{Description} The \emph{registerGh} is preloaded from the memory octuple word at the effective address.


\paragraph{Modifier(s)}
\noindent\begin{itemize}
\item variant (cf. variant in section \ref{par:modifier-variant} on page \pageref{par:modifier-variant})
\end{itemize}

\paragraph{Execution} {Reservation table \textsf{LSU} (Section~\ref{sec:reservation-LSU})}
~
{\begin{lstlisting}
stage ID:
  new variant = _ZX_2(variant);
stage RR:
  new address = GPR[registerZ];
  new target = GPR[registerY];
  new shift = (target & 31) * 8;
  new index = (target >> 5) & 3;
  new where = (registerGh << 2) + index;
  new dri = where;
stage E1:
  new argument1 = XVR[where];
  new bytemask = 4294967295;
  new targetbytes = _ZX_256(bits2bytes(bytemask) << shift);
  CS.XMF = 1;
  new result1 = MEM.load(address, 0xFFFFFFFF & bytemask, variant, dri);
  result1 = (_ZX_256(result1 << shift) & targetbytes) | (argument1 & _ZX_256(~targetbytes));
stage CR:
  XVR[where] = result1;
\end{lstlisting}}
\newpage

\paragraph{Encoding} Format \textsf{LSU\_XPL4RBB.O} (Section~\ref{sec:format-LSU-XPL4RBB.O}), \verb|lsucode5 = 101|\\

\bitfields{ 1/P, 2/00, 2/-- --, 27/offset27, 1/1, 2/01, 3/011, 2/variant, 4/registerGh, 2/01, 2/11, 3/101, 1/--, 6/registerY, 6/registerZ }


\paragraph{Syntax} \texttt{xplo}\emph{variant} \emph{registerGh}, \emph{registerY} = \emph{offset27}[\emph{registerZ}]

\paragraph{Description} The \emph{registerGh} is preloaded from the memory octuple word at the effective address.


\paragraph{Modifier(s)}
\noindent\begin{itemize}
\item variant (cf. variant in section \ref{par:modifier-variant} on page \pageref{par:modifier-variant})
\end{itemize}

\paragraph{Execution} {Reservation table \textsf{LSU.X} (Section~\ref{sec:reservation-LSU.X})}
~
{\begin{lstlisting}
stage ID:
  new offset = _SX_27(offset27);
  new variant = _ZX_2(variant);
stage RR:
  new base = GPR[registerZ];
  new target = GPR[registerY];
  new address = base + offset;
  new shift = (target & 31) * 8;
  new index = (target >> 5) & 3;
  new where = (registerGh << 2) + index;
  new dri = where;
stage E1:
  new argument1 = XVR[where];
  new bytemask = 4294967295;
  new targetbytes = _ZX_256(bits2bytes(bytemask) << shift);
  CS.XMF = 1;
  new result1 = MEM.load(address, 0xFFFFFFFF & bytemask, variant, dri);
  result1 = (_ZX_256(result1 << shift) & targetbytes) | (argument1 & _ZX_256(~targetbytes));
stage CR:
  XVR[where] = result1;
\end{lstlisting}}
\newpage

\paragraph{Encoding} Format \textsf{LSU\_XPL4RBB.D} (Section~\ref{sec:format-LSU-XPL4RBB.D}), \verb|lsucode5 = 101|\\

\bitfields{ 1/P, 2/00, 2/-- --, 27/extend27, 1/1, 2/00, 2/-- --, 27/offset27, 1/1, 2/01, 3/011, 2/variant, 4/registerGh, 2/01, 2/11, 3/101, 1/--, 6/registerY, 6/registerZ }


\paragraph{Syntax} \texttt{xplo}\emph{variant} \emph{registerGh}, \emph{registerY} = \emph{extend27\_\discretionary{}{}{}offset27}[\emph{registerZ}]

\paragraph{Description} The \emph{registerGh} is preloaded from the memory octuple word at the effective address.


\paragraph{Modifier(s)}
\noindent\begin{itemize}
\item variant (cf. variant in section \ref{par:modifier-variant} on page \pageref{par:modifier-variant})
\end{itemize}

\paragraph{Execution} {Reservation table \textsf{LSU.Y} (Section~\ref{sec:reservation-LSU.Y})}
~
{\begin{lstlisting}
stage ID:
  new offset = _SX_54(extend27_offset27);
  new variant = _ZX_2(variant);
stage RR:
  new base = GPR[registerZ];
  new target = GPR[registerY];
  new address = base + offset;
  new shift = (target & 31) * 8;
  new index = (target >> 5) & 3;
  new where = (registerGh << 2) + index;
  new dri = where;
stage E1:
  new argument1 = XVR[where];
  new bytemask = 4294967295;
  new targetbytes = _ZX_256(bits2bytes(bytemask) << shift);
  CS.XMF = 1;
  new result1 = MEM.load(address, 0xFFFFFFFF & bytemask, variant, dri);
  result1 = (_ZX_256(result1 << shift) & targetbytes) | (argument1 & _ZX_256(~targetbytes));
stage CR:
  XVR[where] = result1;
\end{lstlisting}}
\newpage

\paragraph{Encoding} Format \textsf{LSU\_XPL8RBB} (Section~\ref{sec:format-LSU-XPL8RBB}), \verb|lsucode5 = 101|\\

\bitfields{ 1/P, 2/01, 3/011, 2/variant, 3/registerGi, 3/011, 2/11, 3/101, 1/--, 6/registerY, 6/registerZ }


\paragraph{Syntax} \texttt{xplo}\emph{variant} \emph{registerGi}, \emph{registerY} = [\emph{registerZ}]

\paragraph{Description} The \emph{registerGi} is preloaded from the memory octuple word at the effective address.


\paragraph{Modifier(s)}
\noindent\begin{itemize}
\item variant (cf. variant in section \ref{par:modifier-variant} on page \pageref{par:modifier-variant})
\end{itemize}

\paragraph{Execution} {Reservation table \textsf{LSU} (Section~\ref{sec:reservation-LSU})}
~
{\begin{lstlisting}
stage ID:
  new variant = _ZX_2(variant);
stage RR:
  new address = GPR[registerZ];
  new target = GPR[registerY];
  new shift = (target & 31) * 8;
  new index = (target >> 5) & 7;
  new where = (registerGi << 3) + index;
  new dri = where;
stage E1:
  new argument1 = XVR[where];
  new bytemask = 4294967295;
  new targetbytes = _ZX_256(bits2bytes(bytemask) << shift);
  CS.XMF = 1;
  new result1 = MEM.load(address, 0xFFFFFFFF & bytemask, variant, dri);
  result1 = (_ZX_256(result1 << shift) & targetbytes) | (argument1 & _ZX_256(~targetbytes));
stage CR:
  XVR[where] = result1;
\end{lstlisting}}
\newpage

\paragraph{Encoding} Format \textsf{LSU\_XPL8RBB.O} (Section~\ref{sec:format-LSU-XPL8RBB.O}), \verb|lsucode5 = 101|\\

\bitfields{ 1/P, 2/00, 2/-- --, 27/offset27, 1/1, 2/01, 3/011, 2/variant, 3/registerGi, 3/011, 2/11, 3/101, 1/--, 6/registerY, 6/registerZ }


\paragraph{Syntax} \texttt{xplo}\emph{variant} \emph{registerGi}, \emph{registerY} = \emph{offset27}[\emph{registerZ}]

\paragraph{Description} The \emph{registerGi} is preloaded from the memory octuple word at the effective address.


\paragraph{Modifier(s)}
\noindent\begin{itemize}
\item variant (cf. variant in section \ref{par:modifier-variant} on page \pageref{par:modifier-variant})
\end{itemize}

\paragraph{Execution} {Reservation table \textsf{LSU.X} (Section~\ref{sec:reservation-LSU.X})}
~
{\begin{lstlisting}
stage ID:
  new offset = _SX_27(offset27);
  new variant = _ZX_2(variant);
stage RR:
  new base = GPR[registerZ];
  new target = GPR[registerY];
  new address = base + offset;
  new shift = (target & 31) * 8;
  new index = (target >> 5) & 7;
  new where = (registerGi << 3) + index;
  new dri = where;
stage E1:
  new argument1 = XVR[where];
  new bytemask = 4294967295;
  new targetbytes = _ZX_256(bits2bytes(bytemask) << shift);
  CS.XMF = 1;
  new result1 = MEM.load(address, 0xFFFFFFFF & bytemask, variant, dri);
  result1 = (_ZX_256(result1 << shift) & targetbytes) | (argument1 & _ZX_256(~targetbytes));
stage CR:
  XVR[where] = result1;
\end{lstlisting}}
\newpage

\paragraph{Encoding} Format \textsf{LSU\_XPL8RBB.D} (Section~\ref{sec:format-LSU-XPL8RBB.D}), \verb|lsucode5 = 101|\\

\bitfields{ 1/P, 2/00, 2/-- --, 27/extend27, 1/1, 2/00, 2/-- --, 27/offset27, 1/1, 2/01, 3/011, 2/variant, 3/registerGi, 3/011, 2/11, 3/101, 1/--, 6/registerY, 6/registerZ }


\paragraph{Syntax} \texttt{xplo}\emph{variant} \emph{registerGi}, \emph{registerY} = \emph{extend27\_\discretionary{}{}{}offset27}[\emph{registerZ}]

\paragraph{Description} The \emph{registerGi} is preloaded from the memory octuple word at the effective address.


\paragraph{Modifier(s)}
\noindent\begin{itemize}
\item variant (cf. variant in section \ref{par:modifier-variant} on page \pageref{par:modifier-variant})
\end{itemize}

\paragraph{Execution} {Reservation table \textsf{LSU.Y} (Section~\ref{sec:reservation-LSU.Y})}
~
{\begin{lstlisting}
stage ID:
  new offset = _SX_54(extend27_offset27);
  new variant = _ZX_2(variant);
stage RR:
  new base = GPR[registerZ];
  new target = GPR[registerY];
  new address = base + offset;
  new shift = (target & 31) * 8;
  new index = (target >> 5) & 7;
  new where = (registerGi << 3) + index;
  new dri = where;
stage E1:
  new argument1 = XVR[where];
  new bytemask = 4294967295;
  new targetbytes = _ZX_256(bits2bytes(bytemask) << shift);
  CS.XMF = 1;
  new result1 = MEM.load(address, 0xFFFFFFFF & bytemask, variant, dri);
  result1 = (_ZX_256(result1 << shift) & targetbytes) | (argument1 & _ZX_256(~targetbytes));
stage CR:
  XVR[where] = result1;
\end{lstlisting}}
\newpage

\paragraph{Encoding} Format \textsf{LSU\_XPL16RBB} (Section~\ref{sec:format-LSU-XPL16RBB}), \verb|lsucode5 = 101|\\

\bitfields{ 1/P, 2/01, 3/011, 2/variant, 2/registerGj, 4/0111, 2/11, 3/101, 1/--, 6/registerY, 6/registerZ }


\paragraph{Syntax} \texttt{xplo}\emph{variant} \emph{registerGj}, \emph{registerY} = [\emph{registerZ}]

\paragraph{Description} The \emph{registerGj} is preloaded from the memory octuple word at the effective address.


\paragraph{Modifier(s)}
\noindent\begin{itemize}
\item variant (cf. variant in section \ref{par:modifier-variant} on page \pageref{par:modifier-variant})
\end{itemize}

\paragraph{Execution} {Reservation table \textsf{LSU} (Section~\ref{sec:reservation-LSU})}
~
{\begin{lstlisting}
stage ID:
  new variant = _ZX_2(variant);
stage RR:
  new address = GPR[registerZ];
  new target = GPR[registerY];
  new shift = (target & 31) * 8;
  new index = (target >> 5) & 15;
  new where = (registerGj << 4) + index;
  new dri = where;
stage E1:
  new argument1 = XVR[where];
  new bytemask = 4294967295;
  new targetbytes = _ZX_256(bits2bytes(bytemask) << shift);
  CS.XMF = 1;
  new result1 = MEM.load(address, 0xFFFFFFFF & bytemask, variant, dri);
  result1 = (_ZX_256(result1 << shift) & targetbytes) | (argument1 & _ZX_256(~targetbytes));
stage CR:
  XVR[where] = result1;
\end{lstlisting}}
\newpage

\paragraph{Encoding} Format \textsf{LSU\_XPL16RBB.O} (Section~\ref{sec:format-LSU-XPL16RBB.O}), \verb|lsucode5 = 101|\\

\bitfields{ 1/P, 2/00, 2/-- --, 27/offset27, 1/1, 2/01, 3/011, 2/variant, 2/registerGj, 4/0111, 2/11, 3/101, 1/--, 6/registerY, 6/registerZ }


\paragraph{Syntax} \texttt{xplo}\emph{variant} \emph{registerGj}, \emph{registerY} = \emph{offset27}[\emph{registerZ}]

\paragraph{Description} The \emph{registerGj} is preloaded from the memory octuple word at the effective address.


\paragraph{Modifier(s)}
\noindent\begin{itemize}
\item variant (cf. variant in section \ref{par:modifier-variant} on page \pageref{par:modifier-variant})
\end{itemize}

\paragraph{Execution} {Reservation table \textsf{LSU.X} (Section~\ref{sec:reservation-LSU.X})}
~
{\begin{lstlisting}
stage ID:
  new offset = _SX_27(offset27);
  new variant = _ZX_2(variant);
stage RR:
  new base = GPR[registerZ];
  new target = GPR[registerY];
  new address = base + offset;
  new shift = (target & 31) * 8;
  new index = (target >> 5) & 15;
  new where = (registerGj << 4) + index;
  new dri = where;
stage E1:
  new argument1 = XVR[where];
  new bytemask = 4294967295;
  new targetbytes = _ZX_256(bits2bytes(bytemask) << shift);
  CS.XMF = 1;
  new result1 = MEM.load(address, 0xFFFFFFFF & bytemask, variant, dri);
  result1 = (_ZX_256(result1 << shift) & targetbytes) | (argument1 & _ZX_256(~targetbytes));
stage CR:
  XVR[where] = result1;
\end{lstlisting}}
\newpage

\paragraph{Encoding} Format \textsf{LSU\_XPL16RBB.D} (Section~\ref{sec:format-LSU-XPL16RBB.D}), \verb|lsucode5 = 101|\\

\bitfields{ 1/P, 2/00, 2/-- --, 27/extend27, 1/1, 2/00, 2/-- --, 27/offset27, 1/1, 2/01, 3/011, 2/variant, 2/registerGj, 4/0111, 2/11, 3/101, 1/--, 6/registerY, 6/registerZ }


\paragraph{Syntax} \texttt{xplo}\emph{variant} \emph{registerGj}, \emph{registerY} = \emph{extend27\_\discretionary{}{}{}offset27}[\emph{registerZ}]

\paragraph{Description} The \emph{registerGj} is preloaded from the memory octuple word at the effective address.


\paragraph{Modifier(s)}
\noindent\begin{itemize}
\item variant (cf. variant in section \ref{par:modifier-variant} on page \pageref{par:modifier-variant})
\end{itemize}

\paragraph{Execution} {Reservation table \textsf{LSU.Y} (Section~\ref{sec:reservation-LSU.Y})}
~
{\begin{lstlisting}
stage ID:
  new offset = _SX_54(extend27_offset27);
  new variant = _ZX_2(variant);
stage RR:
  new base = GPR[registerZ];
  new target = GPR[registerY];
  new address = base + offset;
  new shift = (target & 31) * 8;
  new index = (target >> 5) & 15;
  new where = (registerGj << 4) + index;
  new dri = where;
stage E1:
  new argument1 = XVR[where];
  new bytemask = 4294967295;
  new targetbytes = _ZX_256(bits2bytes(bytemask) << shift);
  CS.XMF = 1;
  new result1 = MEM.load(address, 0xFFFFFFFF & bytemask, variant, dri);
  result1 = (_ZX_256(result1 << shift) & targetbytes) | (argument1 & _ZX_256(~targetbytes));
stage CR:
  XVR[where] = result1;
\end{lstlisting}}
\newpage

\paragraph{Encoding} Format \textsf{LSU\_XPL32RBB} (Section~\ref{sec:format-LSU-XPL32RBB}), \verb|lsucode5 = 101|\\

\bitfields{ 1/P, 2/01, 3/011, 2/variant, 1/registerGk, 5/01111, 2/11, 3/101, 1/--, 6/registerY, 6/registerZ }


\paragraph{Syntax} \texttt{xplo}\emph{variant} \emph{registerGk}, \emph{registerY} = [\emph{registerZ}]

\paragraph{Description} The \emph{registerGk} is preloaded from the memory octuple word at the effective address.


\paragraph{Modifier(s)}
\noindent\begin{itemize}
\item variant (cf. variant in section \ref{par:modifier-variant} on page \pageref{par:modifier-variant})
\end{itemize}

\paragraph{Execution} {Reservation table \textsf{LSU} (Section~\ref{sec:reservation-LSU})}
~
{\begin{lstlisting}
stage ID:
  new variant = _ZX_2(variant);
stage RR:
  new address = GPR[registerZ];
  new target = GPR[registerY];
  new shift = (target & 31) * 8;
  new index = (target >> 5) & 31;
  new where = (registerGk << 5) + index;
  new dri = where;
stage E1:
  new argument1 = XVR[where];
  new bytemask = 4294967295;
  new targetbytes = _ZX_256(bits2bytes(bytemask) << shift);
  CS.XMF = 1;
  new result1 = MEM.load(address, 0xFFFFFFFF & bytemask, variant, dri);
  result1 = (_ZX_256(result1 << shift) & targetbytes) | (argument1 & _ZX_256(~targetbytes));
stage CR:
  XVR[where] = result1;
\end{lstlisting}}
\newpage

\paragraph{Encoding} Format \textsf{LSU\_XPL32RBB.O} (Section~\ref{sec:format-LSU-XPL32RBB.O}), \verb|lsucode5 = 101|\\

\bitfields{ 1/P, 2/00, 2/-- --, 27/offset27, 1/1, 2/01, 3/011, 2/variant, 1/registerGk, 5/01111, 2/11, 3/101, 1/--, 6/registerY, 6/registerZ }


\paragraph{Syntax} \texttt{xplo}\emph{variant} \emph{registerGk}, \emph{registerY} = \emph{offset27}[\emph{registerZ}]

\paragraph{Description} The \emph{registerGk} is preloaded from the memory octuple word at the effective address.


\paragraph{Modifier(s)}
\noindent\begin{itemize}
\item variant (cf. variant in section \ref{par:modifier-variant} on page \pageref{par:modifier-variant})
\end{itemize}

\paragraph{Execution} {Reservation table \textsf{LSU.X} (Section~\ref{sec:reservation-LSU.X})}
~
{\begin{lstlisting}
stage ID:
  new offset = _SX_27(offset27);
  new variant = _ZX_2(variant);
stage RR:
  new base = GPR[registerZ];
  new target = GPR[registerY];
  new address = base + offset;
  new shift = (target & 31) * 8;
  new index = (target >> 5) & 31;
  new where = (registerGk << 5) + index;
  new dri = where;
stage E1:
  new argument1 = XVR[where];
  new bytemask = 4294967295;
  new targetbytes = _ZX_256(bits2bytes(bytemask) << shift);
  CS.XMF = 1;
  new result1 = MEM.load(address, 0xFFFFFFFF & bytemask, variant, dri);
  result1 = (_ZX_256(result1 << shift) & targetbytes) | (argument1 & _ZX_256(~targetbytes));
stage CR:
  XVR[where] = result1;
\end{lstlisting}}
\newpage

\paragraph{Encoding} Format \textsf{LSU\_XPL32RBB.D} (Section~\ref{sec:format-LSU-XPL32RBB.D}), \verb|lsucode5 = 101|\\

\bitfields{ 1/P, 2/00, 2/-- --, 27/extend27, 1/1, 2/00, 2/-- --, 27/offset27, 1/1, 2/01, 3/011, 2/variant, 1/registerGk, 5/01111, 2/11, 3/101, 1/--, 6/registerY, 6/registerZ }


\paragraph{Syntax} \texttt{xplo}\emph{variant} \emph{registerGk}, \emph{registerY} = \emph{extend27\_\discretionary{}{}{}offset27}[\emph{registerZ}]

\paragraph{Description} The \emph{registerGk} is preloaded from the memory octuple word at the effective address.


\paragraph{Modifier(s)}
\noindent\begin{itemize}
\item variant (cf. variant in section \ref{par:modifier-variant} on page \pageref{par:modifier-variant})
\end{itemize}

\paragraph{Execution} {Reservation table \textsf{LSU.Y} (Section~\ref{sec:reservation-LSU.Y})}
~
{\begin{lstlisting}
stage ID:
  new offset = _SX_54(extend27_offset27);
  new variant = _ZX_2(variant);
stage RR:
  new base = GPR[registerZ];
  new target = GPR[registerY];
  new address = base + offset;
  new shift = (target & 31) * 8;
  new index = (target >> 5) & 31;
  new where = (registerGk << 5) + index;
  new dri = where;
stage E1:
  new argument1 = XVR[where];
  new bytemask = 4294967295;
  new targetbytes = _ZX_256(bits2bytes(bytemask) << shift);
  CS.XMF = 1;
  new result1 = MEM.load(address, 0xFFFFFFFF & bytemask, variant, dri);
  result1 = (_ZX_256(result1 << shift) & targetbytes) | (argument1 & _ZX_256(~targetbytes));
stage CR:
  XVR[where] = result1;
\end{lstlisting}}
\newpage

\paragraph{Encoding} Format \textsf{LSU\_XPL64RBB} (Section~\ref{sec:format-LSU-XPL64RBB}), \verb|lsucode5 = 101|\\

\bitfields{ 1/P, 2/01, 3/011, 2/variant, 1/registerGl, 5/11111, 2/11, 3/101, 1/--, 6/registerY, 6/registerZ }


\paragraph{Syntax} \texttt{xplo}\emph{variant} \emph{registerGl}, \emph{registerY} = [\emph{registerZ}]

\paragraph{Description} The \emph{registerGl} is preloaded from the memory octuple word at the effective address.


\paragraph{Modifier(s)}
\noindent\begin{itemize}
\item variant (cf. variant in section \ref{par:modifier-variant} on page \pageref{par:modifier-variant})
\end{itemize}

\paragraph{Execution} {Reservation table \textsf{LSU} (Section~\ref{sec:reservation-LSU})}
~
{\begin{lstlisting}
stage ID:
  new variant = _ZX_2(variant);
stage RR:
  new address = GPR[registerZ];
  new target = GPR[registerY];
  new shift = (target & 31) * 8;
  new index = (target >> 5) & 63;
  new where = (registerGl << 6) + index;
  new dri = where;
stage E1:
  new argument1 = XVR[where];
  new bytemask = 4294967295;
  new targetbytes = _ZX_256(bits2bytes(bytemask) << shift);
  CS.XMF = 1;
  new result1 = MEM.load(address, 0xFFFFFFFF & bytemask, variant, dri);
  result1 = (_ZX_256(result1 << shift) & targetbytes) | (argument1 & _ZX_256(~targetbytes));
stage CR:
  XVR[where] = result1;
\end{lstlisting}}
\newpage

\paragraph{Encoding} Format \textsf{LSU\_XPL64RBB.O} (Section~\ref{sec:format-LSU-XPL64RBB.O}), \verb|lsucode5 = 101|\\

\bitfields{ 1/P, 2/00, 2/-- --, 27/offset27, 1/1, 2/01, 3/011, 2/variant, 1/registerGl, 5/11111, 2/11, 3/101, 1/--, 6/registerY, 6/registerZ }


\paragraph{Syntax} \texttt{xplo}\emph{variant} \emph{registerGl}, \emph{registerY} = \emph{offset27}[\emph{registerZ}]

\paragraph{Description} The \emph{registerGl} is preloaded from the memory octuple word at the effective address.


\paragraph{Modifier(s)}
\noindent\begin{itemize}
\item variant (cf. variant in section \ref{par:modifier-variant} on page \pageref{par:modifier-variant})
\end{itemize}

\paragraph{Execution} {Reservation table \textsf{LSU.X} (Section~\ref{sec:reservation-LSU.X})}
~
{\begin{lstlisting}
stage ID:
  new offset = _SX_27(offset27);
  new variant = _ZX_2(variant);
stage RR:
  new base = GPR[registerZ];
  new target = GPR[registerY];
  new address = base + offset;
  new shift = (target & 31) * 8;
  new index = (target >> 5) & 63;
  new where = (registerGl << 6) + index;
  new dri = where;
stage E1:
  new argument1 = XVR[where];
  new bytemask = 4294967295;
  new targetbytes = _ZX_256(bits2bytes(bytemask) << shift);
  CS.XMF = 1;
  new result1 = MEM.load(address, 0xFFFFFFFF & bytemask, variant, dri);
  result1 = (_ZX_256(result1 << shift) & targetbytes) | (argument1 & _ZX_256(~targetbytes));
stage CR:
  XVR[where] = result1;
\end{lstlisting}}
\newpage

\paragraph{Encoding} Format \textsf{LSU\_XPL64RBB.D} (Section~\ref{sec:format-LSU-XPL64RBB.D}), \verb|lsucode5 = 101|\\

\bitfields{ 1/P, 2/00, 2/-- --, 27/extend27, 1/1, 2/00, 2/-- --, 27/offset27, 1/1, 2/01, 3/011, 2/variant, 1/registerGl, 5/11111, 2/11, 3/101, 1/--, 6/registerY, 6/registerZ }


\paragraph{Syntax} \texttt{xplo}\emph{variant} \emph{registerGl}, \emph{registerY} = \emph{extend27\_\discretionary{}{}{}offset27}[\emph{registerZ}]

\paragraph{Description} The \emph{registerGl} is preloaded from the memory octuple word at the effective address.


\paragraph{Modifier(s)}
\noindent\begin{itemize}
\item variant (cf. variant in section \ref{par:modifier-variant} on page \pageref{par:modifier-variant})
\end{itemize}

\paragraph{Execution} {Reservation table \textsf{LSU.Y} (Section~\ref{sec:reservation-LSU.Y})}
~
{\begin{lstlisting}
stage ID:
  new offset = _SX_54(extend27_offset27);
  new variant = _ZX_2(variant);
stage RR:
  new base = GPR[registerZ];
  new target = GPR[registerY];
  new address = base + offset;
  new shift = (target & 31) * 8;
  new index = (target >> 5) & 63;
  new where = (registerGl << 6) + index;
  new dri = where;
stage E1:
  new argument1 = XVR[where];
  new bytemask = 4294967295;
  new targetbytes = _ZX_256(bits2bytes(bytemask) << shift);
  CS.XMF = 1;
  new result1 = MEM.load(address, 0xFFFFFFFF & bytemask, variant, dri);
  result1 = (_ZX_256(result1 << shift) & targetbytes) | (argument1 & _ZX_256(~targetbytes));
stage CR:
  XVR[where] = result1;
\end{lstlisting}}
\newpage
\subsubsection[{\small XPLQ: Extension Preload Quadruple Word}]{XPLQ\hfill Extension Preload Quadruple Word}
\label{instr:XPLQ}\index{xplq}
 
\paragraph{Encoding} Format \textsf{LSU\_XPL2RBB} (Section~\ref{sec:format-LSU-XPL2RBB}), \verb|lsucode5 = 100|\\

\bitfields{ 1/P, 2/01, 3/011, 2/variant, 5/registerGg, 1/0, 2/11, 3/100, 1/--, 6/registerY, 6/registerZ }


\paragraph{Syntax} \texttt{xplq}\emph{variant} \emph{registerGg}, \emph{registerY} = [\emph{registerZ}]

\paragraph{Description} The \emph{registerGg} is preloaded from the memory quadruple word at the effective address.


\paragraph{Modifier(s)}
\noindent\begin{itemize}
\item variant (cf. variant in section \ref{par:modifier-variant} on page \pageref{par:modifier-variant})
\end{itemize}

\paragraph{Execution} {Reservation table \textsf{LSU} (Section~\ref{sec:reservation-LSU})}
~
{\begin{lstlisting}
stage ID:
  new variant = _ZX_2(variant);
stage RR:
  new address = GPR[registerZ];
  new target = GPR[registerY];
  new shift = (target & 31) * 8;
  new index = (target >> 5) & 1;
  new where = (registerGg << 1) + index;
  new dri = where;
stage E1:
  new argument1 = XVR[where];
  new bytemask = 65535;
  new targetbytes = _ZX_256(bits2bytes(bytemask) << shift);
  CS.XMF = 1;
  new result1 = MEM.load(address, 0xFFFF & bytemask, variant, dri);
  result1 = (_ZX_256(result1 << shift) & targetbytes) | (argument1 & _ZX_256(~targetbytes));
stage CR:
  XVR[where] = result1;
\end{lstlisting}}
\newpage

\paragraph{Encoding} Format \textsf{LSU\_XPL2RBB.O} (Section~\ref{sec:format-LSU-XPL2RBB.O}), \verb|lsucode5 = 100|\\

\bitfields{ 1/P, 2/00, 2/-- --, 27/offset27, 1/1, 2/01, 3/011, 2/variant, 5/registerGg, 1/0, 2/11, 3/100, 1/--, 6/registerY, 6/registerZ }


\paragraph{Syntax} \texttt{xplq}\emph{variant} \emph{registerGg}, \emph{registerY} = \emph{offset27}[\emph{registerZ}]

\paragraph{Description} The \emph{registerGg} is preloaded from the memory quadruple word at the effective address.


\paragraph{Modifier(s)}
\noindent\begin{itemize}
\item variant (cf. variant in section \ref{par:modifier-variant} on page \pageref{par:modifier-variant})
\end{itemize}

\paragraph{Execution} {Reservation table \textsf{LSU.X} (Section~\ref{sec:reservation-LSU.X})}
~
{\begin{lstlisting}
stage ID:
  new offset = _SX_27(offset27);
  new variant = _ZX_2(variant);
stage RR:
  new base = GPR[registerZ];
  new target = GPR[registerY];
  new address = base + offset;
  new shift = (target & 31) * 8;
  new index = (target >> 5) & 1;
  new where = (registerGg << 1) + index;
  new dri = where;
stage E1:
  new argument1 = XVR[where];
  new bytemask = 65535;
  new targetbytes = _ZX_256(bits2bytes(bytemask) << shift);
  CS.XMF = 1;
  new result1 = MEM.load(address, 0xFFFF & bytemask, variant, dri);
  result1 = (_ZX_256(result1 << shift) & targetbytes) | (argument1 & _ZX_256(~targetbytes));
stage CR:
  XVR[where] = result1;
\end{lstlisting}}
\newpage

\paragraph{Encoding} Format \textsf{LSU\_XPL2RBB.D} (Section~\ref{sec:format-LSU-XPL2RBB.D}), \verb|lsucode5 = 100|\\

\bitfields{ 1/P, 2/00, 2/-- --, 27/extend27, 1/1, 2/00, 2/-- --, 27/offset27, 1/1, 2/01, 3/011, 2/variant, 5/registerGg, 1/0, 2/11, 3/100, 1/--, 6/registerY, 6/registerZ }


\paragraph{Syntax} \texttt{xplq}\emph{variant} \emph{registerGg}, \emph{registerY} = \emph{extend27\_\discretionary{}{}{}offset27}[\emph{registerZ}]

\paragraph{Description} The \emph{registerGg} is preloaded from the memory quadruple word at the effective address.


\paragraph{Modifier(s)}
\noindent\begin{itemize}
\item variant (cf. variant in section \ref{par:modifier-variant} on page \pageref{par:modifier-variant})
\end{itemize}

\paragraph{Execution} {Reservation table \textsf{LSU.Y} (Section~\ref{sec:reservation-LSU.Y})}
~
{\begin{lstlisting}
stage ID:
  new offset = _SX_54(extend27_offset27);
  new variant = _ZX_2(variant);
stage RR:
  new base = GPR[registerZ];
  new target = GPR[registerY];
  new address = base + offset;
  new shift = (target & 31) * 8;
  new index = (target >> 5) & 1;
  new where = (registerGg << 1) + index;
  new dri = where;
stage E1:
  new argument1 = XVR[where];
  new bytemask = 65535;
  new targetbytes = _ZX_256(bits2bytes(bytemask) << shift);
  CS.XMF = 1;
  new result1 = MEM.load(address, 0xFFFF & bytemask, variant, dri);
  result1 = (_ZX_256(result1 << shift) & targetbytes) | (argument1 & _ZX_256(~targetbytes));
stage CR:
  XVR[where] = result1;
\end{lstlisting}}
\newpage

\paragraph{Encoding} Format \textsf{LSU\_XPL4RBB} (Section~\ref{sec:format-LSU-XPL4RBB}), \verb|lsucode5 = 100|\\

\bitfields{ 1/P, 2/01, 3/011, 2/variant, 4/registerGh, 2/01, 2/11, 3/100, 1/--, 6/registerY, 6/registerZ }


\paragraph{Syntax} \texttt{xplq}\emph{variant} \emph{registerGh}, \emph{registerY} = [\emph{registerZ}]

\paragraph{Description} The \emph{registerGh} is preloaded from the memory quadruple word at the effective address.


\paragraph{Modifier(s)}
\noindent\begin{itemize}
\item variant (cf. variant in section \ref{par:modifier-variant} on page \pageref{par:modifier-variant})
\end{itemize}

\paragraph{Execution} {Reservation table \textsf{LSU} (Section~\ref{sec:reservation-LSU})}
~
{\begin{lstlisting}
stage ID:
  new variant = _ZX_2(variant);
stage RR:
  new address = GPR[registerZ];
  new target = GPR[registerY];
  new shift = (target & 31) * 8;
  new index = (target >> 5) & 3;
  new where = (registerGh << 2) + index;
  new dri = where;
stage E1:
  new argument1 = XVR[where];
  new bytemask = 65535;
  new targetbytes = _ZX_256(bits2bytes(bytemask) << shift);
  CS.XMF = 1;
  new result1 = MEM.load(address, 0xFFFF & bytemask, variant, dri);
  result1 = (_ZX_256(result1 << shift) & targetbytes) | (argument1 & _ZX_256(~targetbytes));
stage CR:
  XVR[where] = result1;
\end{lstlisting}}
\newpage

\paragraph{Encoding} Format \textsf{LSU\_XPL4RBB.O} (Section~\ref{sec:format-LSU-XPL4RBB.O}), \verb|lsucode5 = 100|\\

\bitfields{ 1/P, 2/00, 2/-- --, 27/offset27, 1/1, 2/01, 3/011, 2/variant, 4/registerGh, 2/01, 2/11, 3/100, 1/--, 6/registerY, 6/registerZ }


\paragraph{Syntax} \texttt{xplq}\emph{variant} \emph{registerGh}, \emph{registerY} = \emph{offset27}[\emph{registerZ}]

\paragraph{Description} The \emph{registerGh} is preloaded from the memory quadruple word at the effective address.


\paragraph{Modifier(s)}
\noindent\begin{itemize}
\item variant (cf. variant in section \ref{par:modifier-variant} on page \pageref{par:modifier-variant})
\end{itemize}

\paragraph{Execution} {Reservation table \textsf{LSU.X} (Section~\ref{sec:reservation-LSU.X})}
~
{\begin{lstlisting}
stage ID:
  new offset = _SX_27(offset27);
  new variant = _ZX_2(variant);
stage RR:
  new base = GPR[registerZ];
  new target = GPR[registerY];
  new address = base + offset;
  new shift = (target & 31) * 8;
  new index = (target >> 5) & 3;
  new where = (registerGh << 2) + index;
  new dri = where;
stage E1:
  new argument1 = XVR[where];
  new bytemask = 65535;
  new targetbytes = _ZX_256(bits2bytes(bytemask) << shift);
  CS.XMF = 1;
  new result1 = MEM.load(address, 0xFFFF & bytemask, variant, dri);
  result1 = (_ZX_256(result1 << shift) & targetbytes) | (argument1 & _ZX_256(~targetbytes));
stage CR:
  XVR[where] = result1;
\end{lstlisting}}
\newpage

\paragraph{Encoding} Format \textsf{LSU\_XPL4RBB.D} (Section~\ref{sec:format-LSU-XPL4RBB.D}), \verb|lsucode5 = 100|\\

\bitfields{ 1/P, 2/00, 2/-- --, 27/extend27, 1/1, 2/00, 2/-- --, 27/offset27, 1/1, 2/01, 3/011, 2/variant, 4/registerGh, 2/01, 2/11, 3/100, 1/--, 6/registerY, 6/registerZ }


\paragraph{Syntax} \texttt{xplq}\emph{variant} \emph{registerGh}, \emph{registerY} = \emph{extend27\_\discretionary{}{}{}offset27}[\emph{registerZ}]

\paragraph{Description} The \emph{registerGh} is preloaded from the memory quadruple word at the effective address.


\paragraph{Modifier(s)}
\noindent\begin{itemize}
\item variant (cf. variant in section \ref{par:modifier-variant} on page \pageref{par:modifier-variant})
\end{itemize}

\paragraph{Execution} {Reservation table \textsf{LSU.Y} (Section~\ref{sec:reservation-LSU.Y})}
~
{\begin{lstlisting}
stage ID:
  new offset = _SX_54(extend27_offset27);
  new variant = _ZX_2(variant);
stage RR:
  new base = GPR[registerZ];
  new target = GPR[registerY];
  new address = base + offset;
  new shift = (target & 31) * 8;
  new index = (target >> 5) & 3;
  new where = (registerGh << 2) + index;
  new dri = where;
stage E1:
  new argument1 = XVR[where];
  new bytemask = 65535;
  new targetbytes = _ZX_256(bits2bytes(bytemask) << shift);
  CS.XMF = 1;
  new result1 = MEM.load(address, 0xFFFF & bytemask, variant, dri);
  result1 = (_ZX_256(result1 << shift) & targetbytes) | (argument1 & _ZX_256(~targetbytes));
stage CR:
  XVR[where] = result1;
\end{lstlisting}}
\newpage

\paragraph{Encoding} Format \textsf{LSU\_XPL8RBB} (Section~\ref{sec:format-LSU-XPL8RBB}), \verb|lsucode5 = 100|\\

\bitfields{ 1/P, 2/01, 3/011, 2/variant, 3/registerGi, 3/011, 2/11, 3/100, 1/--, 6/registerY, 6/registerZ }


\paragraph{Syntax} \texttt{xplq}\emph{variant} \emph{registerGi}, \emph{registerY} = [\emph{registerZ}]

\paragraph{Description} The \emph{registerGi} is preloaded from the memory quadruple word at the effective address.


\paragraph{Modifier(s)}
\noindent\begin{itemize}
\item variant (cf. variant in section \ref{par:modifier-variant} on page \pageref{par:modifier-variant})
\end{itemize}

\paragraph{Execution} {Reservation table \textsf{LSU} (Section~\ref{sec:reservation-LSU})}
~
{\begin{lstlisting}
stage ID:
  new variant = _ZX_2(variant);
stage RR:
  new address = GPR[registerZ];
  new target = GPR[registerY];
  new shift = (target & 31) * 8;
  new index = (target >> 5) & 7;
  new where = (registerGi << 3) + index;
  new dri = where;
stage E1:
  new argument1 = XVR[where];
  new bytemask = 65535;
  new targetbytes = _ZX_256(bits2bytes(bytemask) << shift);
  CS.XMF = 1;
  new result1 = MEM.load(address, 0xFFFF & bytemask, variant, dri);
  result1 = (_ZX_256(result1 << shift) & targetbytes) | (argument1 & _ZX_256(~targetbytes));
stage CR:
  XVR[where] = result1;
\end{lstlisting}}
\newpage

\paragraph{Encoding} Format \textsf{LSU\_XPL8RBB.O} (Section~\ref{sec:format-LSU-XPL8RBB.O}), \verb|lsucode5 = 100|\\

\bitfields{ 1/P, 2/00, 2/-- --, 27/offset27, 1/1, 2/01, 3/011, 2/variant, 3/registerGi, 3/011, 2/11, 3/100, 1/--, 6/registerY, 6/registerZ }


\paragraph{Syntax} \texttt{xplq}\emph{variant} \emph{registerGi}, \emph{registerY} = \emph{offset27}[\emph{registerZ}]

\paragraph{Description} The \emph{registerGi} is preloaded from the memory quadruple word at the effective address.


\paragraph{Modifier(s)}
\noindent\begin{itemize}
\item variant (cf. variant in section \ref{par:modifier-variant} on page \pageref{par:modifier-variant})
\end{itemize}

\paragraph{Execution} {Reservation table \textsf{LSU.X} (Section~\ref{sec:reservation-LSU.X})}
~
{\begin{lstlisting}
stage ID:
  new offset = _SX_27(offset27);
  new variant = _ZX_2(variant);
stage RR:
  new base = GPR[registerZ];
  new target = GPR[registerY];
  new address = base + offset;
  new shift = (target & 31) * 8;
  new index = (target >> 5) & 7;
  new where = (registerGi << 3) + index;
  new dri = where;
stage E1:
  new argument1 = XVR[where];
  new bytemask = 65535;
  new targetbytes = _ZX_256(bits2bytes(bytemask) << shift);
  CS.XMF = 1;
  new result1 = MEM.load(address, 0xFFFF & bytemask, variant, dri);
  result1 = (_ZX_256(result1 << shift) & targetbytes) | (argument1 & _ZX_256(~targetbytes));
stage CR:
  XVR[where] = result1;
\end{lstlisting}}
\newpage

\paragraph{Encoding} Format \textsf{LSU\_XPL8RBB.D} (Section~\ref{sec:format-LSU-XPL8RBB.D}), \verb|lsucode5 = 100|\\

\bitfields{ 1/P, 2/00, 2/-- --, 27/extend27, 1/1, 2/00, 2/-- --, 27/offset27, 1/1, 2/01, 3/011, 2/variant, 3/registerGi, 3/011, 2/11, 3/100, 1/--, 6/registerY, 6/registerZ }


\paragraph{Syntax} \texttt{xplq}\emph{variant} \emph{registerGi}, \emph{registerY} = \emph{extend27\_\discretionary{}{}{}offset27}[\emph{registerZ}]

\paragraph{Description} The \emph{registerGi} is preloaded from the memory quadruple word at the effective address.


\paragraph{Modifier(s)}
\noindent\begin{itemize}
\item variant (cf. variant in section \ref{par:modifier-variant} on page \pageref{par:modifier-variant})
\end{itemize}

\paragraph{Execution} {Reservation table \textsf{LSU.Y} (Section~\ref{sec:reservation-LSU.Y})}
~
{\begin{lstlisting}
stage ID:
  new offset = _SX_54(extend27_offset27);
  new variant = _ZX_2(variant);
stage RR:
  new base = GPR[registerZ];
  new target = GPR[registerY];
  new address = base + offset;
  new shift = (target & 31) * 8;
  new index = (target >> 5) & 7;
  new where = (registerGi << 3) + index;
  new dri = where;
stage E1:
  new argument1 = XVR[where];
  new bytemask = 65535;
  new targetbytes = _ZX_256(bits2bytes(bytemask) << shift);
  CS.XMF = 1;
  new result1 = MEM.load(address, 0xFFFF & bytemask, variant, dri);
  result1 = (_ZX_256(result1 << shift) & targetbytes) | (argument1 & _ZX_256(~targetbytes));
stage CR:
  XVR[where] = result1;
\end{lstlisting}}
\newpage

\paragraph{Encoding} Format \textsf{LSU\_XPL16RBB} (Section~\ref{sec:format-LSU-XPL16RBB}), \verb|lsucode5 = 100|\\

\bitfields{ 1/P, 2/01, 3/011, 2/variant, 2/registerGj, 4/0111, 2/11, 3/100, 1/--, 6/registerY, 6/registerZ }


\paragraph{Syntax} \texttt{xplq}\emph{variant} \emph{registerGj}, \emph{registerY} = [\emph{registerZ}]

\paragraph{Description} The \emph{registerGj} is preloaded from the memory quadruple word at the effective address.


\paragraph{Modifier(s)}
\noindent\begin{itemize}
\item variant (cf. variant in section \ref{par:modifier-variant} on page \pageref{par:modifier-variant})
\end{itemize}

\paragraph{Execution} {Reservation table \textsf{LSU} (Section~\ref{sec:reservation-LSU})}
~
{\begin{lstlisting}
stage ID:
  new variant = _ZX_2(variant);
stage RR:
  new address = GPR[registerZ];
  new target = GPR[registerY];
  new shift = (target & 31) * 8;
  new index = (target >> 5) & 15;
  new where = (registerGj << 4) + index;
  new dri = where;
stage E1:
  new argument1 = XVR[where];
  new bytemask = 65535;
  new targetbytes = _ZX_256(bits2bytes(bytemask) << shift);
  CS.XMF = 1;
  new result1 = MEM.load(address, 0xFFFF & bytemask, variant, dri);
  result1 = (_ZX_256(result1 << shift) & targetbytes) | (argument1 & _ZX_256(~targetbytes));
stage CR:
  XVR[where] = result1;
\end{lstlisting}}
\newpage

\paragraph{Encoding} Format \textsf{LSU\_XPL16RBB.O} (Section~\ref{sec:format-LSU-XPL16RBB.O}), \verb|lsucode5 = 100|\\

\bitfields{ 1/P, 2/00, 2/-- --, 27/offset27, 1/1, 2/01, 3/011, 2/variant, 2/registerGj, 4/0111, 2/11, 3/100, 1/--, 6/registerY, 6/registerZ }


\paragraph{Syntax} \texttt{xplq}\emph{variant} \emph{registerGj}, \emph{registerY} = \emph{offset27}[\emph{registerZ}]

\paragraph{Description} The \emph{registerGj} is preloaded from the memory quadruple word at the effective address.


\paragraph{Modifier(s)}
\noindent\begin{itemize}
\item variant (cf. variant in section \ref{par:modifier-variant} on page \pageref{par:modifier-variant})
\end{itemize}

\paragraph{Execution} {Reservation table \textsf{LSU.X} (Section~\ref{sec:reservation-LSU.X})}
~
{\begin{lstlisting}
stage ID:
  new offset = _SX_27(offset27);
  new variant = _ZX_2(variant);
stage RR:
  new base = GPR[registerZ];
  new target = GPR[registerY];
  new address = base + offset;
  new shift = (target & 31) * 8;
  new index = (target >> 5) & 15;
  new where = (registerGj << 4) + index;
  new dri = where;
stage E1:
  new argument1 = XVR[where];
  new bytemask = 65535;
  new targetbytes = _ZX_256(bits2bytes(bytemask) << shift);
  CS.XMF = 1;
  new result1 = MEM.load(address, 0xFFFF & bytemask, variant, dri);
  result1 = (_ZX_256(result1 << shift) & targetbytes) | (argument1 & _ZX_256(~targetbytes));
stage CR:
  XVR[where] = result1;
\end{lstlisting}}
\newpage

\paragraph{Encoding} Format \textsf{LSU\_XPL16RBB.D} (Section~\ref{sec:format-LSU-XPL16RBB.D}), \verb|lsucode5 = 100|\\

\bitfields{ 1/P, 2/00, 2/-- --, 27/extend27, 1/1, 2/00, 2/-- --, 27/offset27, 1/1, 2/01, 3/011, 2/variant, 2/registerGj, 4/0111, 2/11, 3/100, 1/--, 6/registerY, 6/registerZ }


\paragraph{Syntax} \texttt{xplq}\emph{variant} \emph{registerGj}, \emph{registerY} = \emph{extend27\_\discretionary{}{}{}offset27}[\emph{registerZ}]

\paragraph{Description} The \emph{registerGj} is preloaded from the memory quadruple word at the effective address.


\paragraph{Modifier(s)}
\noindent\begin{itemize}
\item variant (cf. variant in section \ref{par:modifier-variant} on page \pageref{par:modifier-variant})
\end{itemize}

\paragraph{Execution} {Reservation table \textsf{LSU.Y} (Section~\ref{sec:reservation-LSU.Y})}
~
{\begin{lstlisting}
stage ID:
  new offset = _SX_54(extend27_offset27);
  new variant = _ZX_2(variant);
stage RR:
  new base = GPR[registerZ];
  new target = GPR[registerY];
  new address = base + offset;
  new shift = (target & 31) * 8;
  new index = (target >> 5) & 15;
  new where = (registerGj << 4) + index;
  new dri = where;
stage E1:
  new argument1 = XVR[where];
  new bytemask = 65535;
  new targetbytes = _ZX_256(bits2bytes(bytemask) << shift);
  CS.XMF = 1;
  new result1 = MEM.load(address, 0xFFFF & bytemask, variant, dri);
  result1 = (_ZX_256(result1 << shift) & targetbytes) | (argument1 & _ZX_256(~targetbytes));
stage CR:
  XVR[where] = result1;
\end{lstlisting}}
\newpage

\paragraph{Encoding} Format \textsf{LSU\_XPL32RBB} (Section~\ref{sec:format-LSU-XPL32RBB}), \verb|lsucode5 = 100|\\

\bitfields{ 1/P, 2/01, 3/011, 2/variant, 1/registerGk, 5/01111, 2/11, 3/100, 1/--, 6/registerY, 6/registerZ }


\paragraph{Syntax} \texttt{xplq}\emph{variant} \emph{registerGk}, \emph{registerY} = [\emph{registerZ}]

\paragraph{Description} The \emph{registerGk} is preloaded from the memory quadruple word at the effective address.


\paragraph{Modifier(s)}
\noindent\begin{itemize}
\item variant (cf. variant in section \ref{par:modifier-variant} on page \pageref{par:modifier-variant})
\end{itemize}

\paragraph{Execution} {Reservation table \textsf{LSU} (Section~\ref{sec:reservation-LSU})}
~
{\begin{lstlisting}
stage ID:
  new variant = _ZX_2(variant);
stage RR:
  new address = GPR[registerZ];
  new target = GPR[registerY];
  new shift = (target & 31) * 8;
  new index = (target >> 5) & 31;
  new where = (registerGk << 5) + index;
  new dri = where;
stage E1:
  new argument1 = XVR[where];
  new bytemask = 65535;
  new targetbytes = _ZX_256(bits2bytes(bytemask) << shift);
  CS.XMF = 1;
  new result1 = MEM.load(address, 0xFFFF & bytemask, variant, dri);
  result1 = (_ZX_256(result1 << shift) & targetbytes) | (argument1 & _ZX_256(~targetbytes));
stage CR:
  XVR[where] = result1;
\end{lstlisting}}
\newpage

\paragraph{Encoding} Format \textsf{LSU\_XPL32RBB.O} (Section~\ref{sec:format-LSU-XPL32RBB.O}), \verb|lsucode5 = 100|\\

\bitfields{ 1/P, 2/00, 2/-- --, 27/offset27, 1/1, 2/01, 3/011, 2/variant, 1/registerGk, 5/01111, 2/11, 3/100, 1/--, 6/registerY, 6/registerZ }


\paragraph{Syntax} \texttt{xplq}\emph{variant} \emph{registerGk}, \emph{registerY} = \emph{offset27}[\emph{registerZ}]

\paragraph{Description} The \emph{registerGk} is preloaded from the memory quadruple word at the effective address.


\paragraph{Modifier(s)}
\noindent\begin{itemize}
\item variant (cf. variant in section \ref{par:modifier-variant} on page \pageref{par:modifier-variant})
\end{itemize}

\paragraph{Execution} {Reservation table \textsf{LSU.X} (Section~\ref{sec:reservation-LSU.X})}
~
{\begin{lstlisting}
stage ID:
  new offset = _SX_27(offset27);
  new variant = _ZX_2(variant);
stage RR:
  new base = GPR[registerZ];
  new target = GPR[registerY];
  new address = base + offset;
  new shift = (target & 31) * 8;
  new index = (target >> 5) & 31;
  new where = (registerGk << 5) + index;
  new dri = where;
stage E1:
  new argument1 = XVR[where];
  new bytemask = 65535;
  new targetbytes = _ZX_256(bits2bytes(bytemask) << shift);
  CS.XMF = 1;
  new result1 = MEM.load(address, 0xFFFF & bytemask, variant, dri);
  result1 = (_ZX_256(result1 << shift) & targetbytes) | (argument1 & _ZX_256(~targetbytes));
stage CR:
  XVR[where] = result1;
\end{lstlisting}}
\newpage

\paragraph{Encoding} Format \textsf{LSU\_XPL32RBB.D} (Section~\ref{sec:format-LSU-XPL32RBB.D}), \verb|lsucode5 = 100|\\

\bitfields{ 1/P, 2/00, 2/-- --, 27/extend27, 1/1, 2/00, 2/-- --, 27/offset27, 1/1, 2/01, 3/011, 2/variant, 1/registerGk, 5/01111, 2/11, 3/100, 1/--, 6/registerY, 6/registerZ }


\paragraph{Syntax} \texttt{xplq}\emph{variant} \emph{registerGk}, \emph{registerY} = \emph{extend27\_\discretionary{}{}{}offset27}[\emph{registerZ}]

\paragraph{Description} The \emph{registerGk} is preloaded from the memory quadruple word at the effective address.


\paragraph{Modifier(s)}
\noindent\begin{itemize}
\item variant (cf. variant in section \ref{par:modifier-variant} on page \pageref{par:modifier-variant})
\end{itemize}

\paragraph{Execution} {Reservation table \textsf{LSU.Y} (Section~\ref{sec:reservation-LSU.Y})}
~
{\begin{lstlisting}
stage ID:
  new offset = _SX_54(extend27_offset27);
  new variant = _ZX_2(variant);
stage RR:
  new base = GPR[registerZ];
  new target = GPR[registerY];
  new address = base + offset;
  new shift = (target & 31) * 8;
  new index = (target >> 5) & 31;
  new where = (registerGk << 5) + index;
  new dri = where;
stage E1:
  new argument1 = XVR[where];
  new bytemask = 65535;
  new targetbytes = _ZX_256(bits2bytes(bytemask) << shift);
  CS.XMF = 1;
  new result1 = MEM.load(address, 0xFFFF & bytemask, variant, dri);
  result1 = (_ZX_256(result1 << shift) & targetbytes) | (argument1 & _ZX_256(~targetbytes));
stage CR:
  XVR[where] = result1;
\end{lstlisting}}
\newpage

\paragraph{Encoding} Format \textsf{LSU\_XPL64RBB} (Section~\ref{sec:format-LSU-XPL64RBB}), \verb|lsucode5 = 100|\\

\bitfields{ 1/P, 2/01, 3/011, 2/variant, 1/registerGl, 5/11111, 2/11, 3/100, 1/--, 6/registerY, 6/registerZ }


\paragraph{Syntax} \texttt{xplq}\emph{variant} \emph{registerGl}, \emph{registerY} = [\emph{registerZ}]

\paragraph{Description} The \emph{registerGl} is preloaded from the memory quadruple word at the effective address.


\paragraph{Modifier(s)}
\noindent\begin{itemize}
\item variant (cf. variant in section \ref{par:modifier-variant} on page \pageref{par:modifier-variant})
\end{itemize}

\paragraph{Execution} {Reservation table \textsf{LSU} (Section~\ref{sec:reservation-LSU})}
~
{\begin{lstlisting}
stage ID:
  new variant = _ZX_2(variant);
stage RR:
  new address = GPR[registerZ];
  new target = GPR[registerY];
  new shift = (target & 31) * 8;
  new index = (target >> 5) & 63;
  new where = (registerGl << 6) + index;
  new dri = where;
stage E1:
  new argument1 = XVR[where];
  new bytemask = 65535;
  new targetbytes = _ZX_256(bits2bytes(bytemask) << shift);
  CS.XMF = 1;
  new result1 = MEM.load(address, 0xFFFF & bytemask, variant, dri);
  result1 = (_ZX_256(result1 << shift) & targetbytes) | (argument1 & _ZX_256(~targetbytes));
stage CR:
  XVR[where] = result1;
\end{lstlisting}}
\newpage

\paragraph{Encoding} Format \textsf{LSU\_XPL64RBB.O} (Section~\ref{sec:format-LSU-XPL64RBB.O}), \verb|lsucode5 = 100|\\

\bitfields{ 1/P, 2/00, 2/-- --, 27/offset27, 1/1, 2/01, 3/011, 2/variant, 1/registerGl, 5/11111, 2/11, 3/100, 1/--, 6/registerY, 6/registerZ }


\paragraph{Syntax} \texttt{xplq}\emph{variant} \emph{registerGl}, \emph{registerY} = \emph{offset27}[\emph{registerZ}]

\paragraph{Description} The \emph{registerGl} is preloaded from the memory quadruple word at the effective address.


\paragraph{Modifier(s)}
\noindent\begin{itemize}
\item variant (cf. variant in section \ref{par:modifier-variant} on page \pageref{par:modifier-variant})
\end{itemize}

\paragraph{Execution} {Reservation table \textsf{LSU.X} (Section~\ref{sec:reservation-LSU.X})}
~
{\begin{lstlisting}
stage ID:
  new offset = _SX_27(offset27);
  new variant = _ZX_2(variant);
stage RR:
  new base = GPR[registerZ];
  new target = GPR[registerY];
  new address = base + offset;
  new shift = (target & 31) * 8;
  new index = (target >> 5) & 63;
  new where = (registerGl << 6) + index;
  new dri = where;
stage E1:
  new argument1 = XVR[where];
  new bytemask = 65535;
  new targetbytes = _ZX_256(bits2bytes(bytemask) << shift);
  CS.XMF = 1;
  new result1 = MEM.load(address, 0xFFFF & bytemask, variant, dri);
  result1 = (_ZX_256(result1 << shift) & targetbytes) | (argument1 & _ZX_256(~targetbytes));
stage CR:
  XVR[where] = result1;
\end{lstlisting}}
\newpage

\paragraph{Encoding} Format \textsf{LSU\_XPL64RBB.D} (Section~\ref{sec:format-LSU-XPL64RBB.D}), \verb|lsucode5 = 100|\\

\bitfields{ 1/P, 2/00, 2/-- --, 27/extend27, 1/1, 2/00, 2/-- --, 27/offset27, 1/1, 2/01, 3/011, 2/variant, 1/registerGl, 5/11111, 2/11, 3/100, 1/--, 6/registerY, 6/registerZ }


\paragraph{Syntax} \texttt{xplq}\emph{variant} \emph{registerGl}, \emph{registerY} = \emph{extend27\_\discretionary{}{}{}offset27}[\emph{registerZ}]

\paragraph{Description} The \emph{registerGl} is preloaded from the memory quadruple word at the effective address.


\paragraph{Modifier(s)}
\noindent\begin{itemize}
\item variant (cf. variant in section \ref{par:modifier-variant} on page \pageref{par:modifier-variant})
\end{itemize}

\paragraph{Execution} {Reservation table \textsf{LSU.Y} (Section~\ref{sec:reservation-LSU.Y})}
~
{\begin{lstlisting}
stage ID:
  new offset = _SX_54(extend27_offset27);
  new variant = _ZX_2(variant);
stage RR:
  new base = GPR[registerZ];
  new target = GPR[registerY];
  new address = base + offset;
  new shift = (target & 31) * 8;
  new index = (target >> 5) & 63;
  new where = (registerGl << 6) + index;
  new dri = where;
stage E1:
  new argument1 = XVR[where];
  new bytemask = 65535;
  new targetbytes = _ZX_256(bits2bytes(bytemask) << shift);
  CS.XMF = 1;
  new result1 = MEM.load(address, 0xFFFF & bytemask, variant, dri);
  result1 = (_ZX_256(result1 << shift) & targetbytes) | (argument1 & _ZX_256(~targetbytes));
stage CR:
  XVR[where] = result1;
\end{lstlisting}}
\newpage
\subsubsection[{\small XPLW: Extension Preload Word}]{XPLW\hfill Extension Preload Word}
\label{instr:XPLW}\index{xplw}
 
\paragraph{Encoding} Format \textsf{LSU\_XPL2RBB} (Section~\ref{sec:format-LSU-XPL2RBB}), \verb|lsucode5 = 010|\\

\bitfields{ 1/P, 2/01, 3/011, 2/variant, 5/registerGg, 1/0, 2/11, 3/010, 1/--, 6/registerY, 6/registerZ }


\paragraph{Syntax} \texttt{xplw}\emph{variant} \emph{registerGg}, \emph{registerY} = [\emph{registerZ}]

\paragraph{Description} The \emph{registerGg} is preloaded from the memory word at the effective address.


\paragraph{Modifier(s)}
\noindent\begin{itemize}
\item variant (cf. variant in section \ref{par:modifier-variant} on page \pageref{par:modifier-variant})
\end{itemize}

\paragraph{Execution} {Reservation table \textsf{LSU} (Section~\ref{sec:reservation-LSU})}
~
{\begin{lstlisting}
stage ID:
  new variant = _ZX_2(variant);
stage RR:
  new address = GPR[registerZ];
  new target = GPR[registerY];
  new shift = (target & 31) * 8;
  new index = (target >> 5) & 1;
  new where = (registerGg << 1) + index;
  new dri = where;
stage E1:
  new argument1 = XVR[where];
  new bytemask = 15;
  new targetbytes = _ZX_256(bits2bytes(bytemask) << shift);
  CS.XMF = 1;
  new result1 = MEM.load(address, 0xF & bytemask, variant, dri);
  result1 = (_ZX_256(result1 << shift) & targetbytes) | (argument1 & _ZX_256(~targetbytes));
stage CR:
  XVR[where] = result1;
\end{lstlisting}}
\newpage

\paragraph{Encoding} Format \textsf{LSU\_XPL2RBB.O} (Section~\ref{sec:format-LSU-XPL2RBB.O}), \verb|lsucode5 = 010|\\

\bitfields{ 1/P, 2/00, 2/-- --, 27/offset27, 1/1, 2/01, 3/011, 2/variant, 5/registerGg, 1/0, 2/11, 3/010, 1/--, 6/registerY, 6/registerZ }


\paragraph{Syntax} \texttt{xplw}\emph{variant} \emph{registerGg}, \emph{registerY} = \emph{offset27}[\emph{registerZ}]

\paragraph{Description} The \emph{registerGg} is preloaded from the memory word at the effective address.


\paragraph{Modifier(s)}
\noindent\begin{itemize}
\item variant (cf. variant in section \ref{par:modifier-variant} on page \pageref{par:modifier-variant})
\end{itemize}

\paragraph{Execution} {Reservation table \textsf{LSU.X} (Section~\ref{sec:reservation-LSU.X})}
~
{\begin{lstlisting}
stage ID:
  new offset = _SX_27(offset27);
  new variant = _ZX_2(variant);
stage RR:
  new base = GPR[registerZ];
  new target = GPR[registerY];
  new address = base + offset;
  new shift = (target & 31) * 8;
  new index = (target >> 5) & 1;
  new where = (registerGg << 1) + index;
  new dri = where;
stage E1:
  new argument1 = XVR[where];
  new bytemask = 15;
  new targetbytes = _ZX_256(bits2bytes(bytemask) << shift);
  CS.XMF = 1;
  new result1 = MEM.load(address, 0xF & bytemask, variant, dri);
  result1 = (_ZX_256(result1 << shift) & targetbytes) | (argument1 & _ZX_256(~targetbytes));
stage CR:
  XVR[where] = result1;
\end{lstlisting}}
\newpage

\paragraph{Encoding} Format \textsf{LSU\_XPL2RBB.D} (Section~\ref{sec:format-LSU-XPL2RBB.D}), \verb|lsucode5 = 010|\\

\bitfields{ 1/P, 2/00, 2/-- --, 27/extend27, 1/1, 2/00, 2/-- --, 27/offset27, 1/1, 2/01, 3/011, 2/variant, 5/registerGg, 1/0, 2/11, 3/010, 1/--, 6/registerY, 6/registerZ }


\paragraph{Syntax} \texttt{xplw}\emph{variant} \emph{registerGg}, \emph{registerY} = \emph{extend27\_\discretionary{}{}{}offset27}[\emph{registerZ}]

\paragraph{Description} The \emph{registerGg} is preloaded from the memory word at the effective address.


\paragraph{Modifier(s)}
\noindent\begin{itemize}
\item variant (cf. variant in section \ref{par:modifier-variant} on page \pageref{par:modifier-variant})
\end{itemize}

\paragraph{Execution} {Reservation table \textsf{LSU.Y} (Section~\ref{sec:reservation-LSU.Y})}
~
{\begin{lstlisting}
stage ID:
  new offset = _SX_54(extend27_offset27);
  new variant = _ZX_2(variant);
stage RR:
  new base = GPR[registerZ];
  new target = GPR[registerY];
  new address = base + offset;
  new shift = (target & 31) * 8;
  new index = (target >> 5) & 1;
  new where = (registerGg << 1) + index;
  new dri = where;
stage E1:
  new argument1 = XVR[where];
  new bytemask = 15;
  new targetbytes = _ZX_256(bits2bytes(bytemask) << shift);
  CS.XMF = 1;
  new result1 = MEM.load(address, 0xF & bytemask, variant, dri);
  result1 = (_ZX_256(result1 << shift) & targetbytes) | (argument1 & _ZX_256(~targetbytes));
stage CR:
  XVR[where] = result1;
\end{lstlisting}}
\newpage

\paragraph{Encoding} Format \textsf{LSU\_XPL4RBB} (Section~\ref{sec:format-LSU-XPL4RBB}), \verb|lsucode5 = 010|\\

\bitfields{ 1/P, 2/01, 3/011, 2/variant, 4/registerGh, 2/01, 2/11, 3/010, 1/--, 6/registerY, 6/registerZ }


\paragraph{Syntax} \texttt{xplw}\emph{variant} \emph{registerGh}, \emph{registerY} = [\emph{registerZ}]

\paragraph{Description} The \emph{registerGh} is preloaded from the memory word at the effective address.


\paragraph{Modifier(s)}
\noindent\begin{itemize}
\item variant (cf. variant in section \ref{par:modifier-variant} on page \pageref{par:modifier-variant})
\end{itemize}

\paragraph{Execution} {Reservation table \textsf{LSU} (Section~\ref{sec:reservation-LSU})}
~
{\begin{lstlisting}
stage ID:
  new variant = _ZX_2(variant);
stage RR:
  new address = GPR[registerZ];
  new target = GPR[registerY];
  new shift = (target & 31) * 8;
  new index = (target >> 5) & 3;
  new where = (registerGh << 2) + index;
  new dri = where;
stage E1:
  new argument1 = XVR[where];
  new bytemask = 15;
  new targetbytes = _ZX_256(bits2bytes(bytemask) << shift);
  CS.XMF = 1;
  new result1 = MEM.load(address, 0xF & bytemask, variant, dri);
  result1 = (_ZX_256(result1 << shift) & targetbytes) | (argument1 & _ZX_256(~targetbytes));
stage CR:
  XVR[where] = result1;
\end{lstlisting}}
\newpage

\paragraph{Encoding} Format \textsf{LSU\_XPL4RBB.O} (Section~\ref{sec:format-LSU-XPL4RBB.O}), \verb|lsucode5 = 010|\\

\bitfields{ 1/P, 2/00, 2/-- --, 27/offset27, 1/1, 2/01, 3/011, 2/variant, 4/registerGh, 2/01, 2/11, 3/010, 1/--, 6/registerY, 6/registerZ }


\paragraph{Syntax} \texttt{xplw}\emph{variant} \emph{registerGh}, \emph{registerY} = \emph{offset27}[\emph{registerZ}]

\paragraph{Description} The \emph{registerGh} is preloaded from the memory word at the effective address.


\paragraph{Modifier(s)}
\noindent\begin{itemize}
\item variant (cf. variant in section \ref{par:modifier-variant} on page \pageref{par:modifier-variant})
\end{itemize}

\paragraph{Execution} {Reservation table \textsf{LSU.X} (Section~\ref{sec:reservation-LSU.X})}
~
{\begin{lstlisting}
stage ID:
  new offset = _SX_27(offset27);
  new variant = _ZX_2(variant);
stage RR:
  new base = GPR[registerZ];
  new target = GPR[registerY];
  new address = base + offset;
  new shift = (target & 31) * 8;
  new index = (target >> 5) & 3;
  new where = (registerGh << 2) + index;
  new dri = where;
stage E1:
  new argument1 = XVR[where];
  new bytemask = 15;
  new targetbytes = _ZX_256(bits2bytes(bytemask) << shift);
  CS.XMF = 1;
  new result1 = MEM.load(address, 0xF & bytemask, variant, dri);
  result1 = (_ZX_256(result1 << shift) & targetbytes) | (argument1 & _ZX_256(~targetbytes));
stage CR:
  XVR[where] = result1;
\end{lstlisting}}
\newpage

\paragraph{Encoding} Format \textsf{LSU\_XPL4RBB.D} (Section~\ref{sec:format-LSU-XPL4RBB.D}), \verb|lsucode5 = 010|\\

\bitfields{ 1/P, 2/00, 2/-- --, 27/extend27, 1/1, 2/00, 2/-- --, 27/offset27, 1/1, 2/01, 3/011, 2/variant, 4/registerGh, 2/01, 2/11, 3/010, 1/--, 6/registerY, 6/registerZ }


\paragraph{Syntax} \texttt{xplw}\emph{variant} \emph{registerGh}, \emph{registerY} = \emph{extend27\_\discretionary{}{}{}offset27}[\emph{registerZ}]

\paragraph{Description} The \emph{registerGh} is preloaded from the memory word at the effective address.


\paragraph{Modifier(s)}
\noindent\begin{itemize}
\item variant (cf. variant in section \ref{par:modifier-variant} on page \pageref{par:modifier-variant})
\end{itemize}

\paragraph{Execution} {Reservation table \textsf{LSU.Y} (Section~\ref{sec:reservation-LSU.Y})}
~
{\begin{lstlisting}
stage ID:
  new offset = _SX_54(extend27_offset27);
  new variant = _ZX_2(variant);
stage RR:
  new base = GPR[registerZ];
  new target = GPR[registerY];
  new address = base + offset;
  new shift = (target & 31) * 8;
  new index = (target >> 5) & 3;
  new where = (registerGh << 2) + index;
  new dri = where;
stage E1:
  new argument1 = XVR[where];
  new bytemask = 15;
  new targetbytes = _ZX_256(bits2bytes(bytemask) << shift);
  CS.XMF = 1;
  new result1 = MEM.load(address, 0xF & bytemask, variant, dri);
  result1 = (_ZX_256(result1 << shift) & targetbytes) | (argument1 & _ZX_256(~targetbytes));
stage CR:
  XVR[where] = result1;
\end{lstlisting}}
\newpage

\paragraph{Encoding} Format \textsf{LSU\_XPL8RBB} (Section~\ref{sec:format-LSU-XPL8RBB}), \verb|lsucode5 = 010|\\

\bitfields{ 1/P, 2/01, 3/011, 2/variant, 3/registerGi, 3/011, 2/11, 3/010, 1/--, 6/registerY, 6/registerZ }


\paragraph{Syntax} \texttt{xplw}\emph{variant} \emph{registerGi}, \emph{registerY} = [\emph{registerZ}]

\paragraph{Description} The \emph{registerGi} is preloaded from the memory word at the effective address.


\paragraph{Modifier(s)}
\noindent\begin{itemize}
\item variant (cf. variant in section \ref{par:modifier-variant} on page \pageref{par:modifier-variant})
\end{itemize}

\paragraph{Execution} {Reservation table \textsf{LSU} (Section~\ref{sec:reservation-LSU})}
~
{\begin{lstlisting}
stage ID:
  new variant = _ZX_2(variant);
stage RR:
  new address = GPR[registerZ];
  new target = GPR[registerY];
  new shift = (target & 31) * 8;
  new index = (target >> 5) & 7;
  new where = (registerGi << 3) + index;
  new dri = where;
stage E1:
  new argument1 = XVR[where];
  new bytemask = 15;
  new targetbytes = _ZX_256(bits2bytes(bytemask) << shift);
  CS.XMF = 1;
  new result1 = MEM.load(address, 0xF & bytemask, variant, dri);
  result1 = (_ZX_256(result1 << shift) & targetbytes) | (argument1 & _ZX_256(~targetbytes));
stage CR:
  XVR[where] = result1;
\end{lstlisting}}
\newpage

\paragraph{Encoding} Format \textsf{LSU\_XPL8RBB.O} (Section~\ref{sec:format-LSU-XPL8RBB.O}), \verb|lsucode5 = 010|\\

\bitfields{ 1/P, 2/00, 2/-- --, 27/offset27, 1/1, 2/01, 3/011, 2/variant, 3/registerGi, 3/011, 2/11, 3/010, 1/--, 6/registerY, 6/registerZ }


\paragraph{Syntax} \texttt{xplw}\emph{variant} \emph{registerGi}, \emph{registerY} = \emph{offset27}[\emph{registerZ}]

\paragraph{Description} The \emph{registerGi} is preloaded from the memory word at the effective address.


\paragraph{Modifier(s)}
\noindent\begin{itemize}
\item variant (cf. variant in section \ref{par:modifier-variant} on page \pageref{par:modifier-variant})
\end{itemize}

\paragraph{Execution} {Reservation table \textsf{LSU.X} (Section~\ref{sec:reservation-LSU.X})}
~
{\begin{lstlisting}
stage ID:
  new offset = _SX_27(offset27);
  new variant = _ZX_2(variant);
stage RR:
  new base = GPR[registerZ];
  new target = GPR[registerY];
  new address = base + offset;
  new shift = (target & 31) * 8;
  new index = (target >> 5) & 7;
  new where = (registerGi << 3) + index;
  new dri = where;
stage E1:
  new argument1 = XVR[where];
  new bytemask = 15;
  new targetbytes = _ZX_256(bits2bytes(bytemask) << shift);
  CS.XMF = 1;
  new result1 = MEM.load(address, 0xF & bytemask, variant, dri);
  result1 = (_ZX_256(result1 << shift) & targetbytes) | (argument1 & _ZX_256(~targetbytes));
stage CR:
  XVR[where] = result1;
\end{lstlisting}}
\newpage

\paragraph{Encoding} Format \textsf{LSU\_XPL8RBB.D} (Section~\ref{sec:format-LSU-XPL8RBB.D}), \verb|lsucode5 = 010|\\

\bitfields{ 1/P, 2/00, 2/-- --, 27/extend27, 1/1, 2/00, 2/-- --, 27/offset27, 1/1, 2/01, 3/011, 2/variant, 3/registerGi, 3/011, 2/11, 3/010, 1/--, 6/registerY, 6/registerZ }


\paragraph{Syntax} \texttt{xplw}\emph{variant} \emph{registerGi}, \emph{registerY} = \emph{extend27\_\discretionary{}{}{}offset27}[\emph{registerZ}]

\paragraph{Description} The \emph{registerGi} is preloaded from the memory word at the effective address.


\paragraph{Modifier(s)}
\noindent\begin{itemize}
\item variant (cf. variant in section \ref{par:modifier-variant} on page \pageref{par:modifier-variant})
\end{itemize}

\paragraph{Execution} {Reservation table \textsf{LSU.Y} (Section~\ref{sec:reservation-LSU.Y})}
~
{\begin{lstlisting}
stage ID:
  new offset = _SX_54(extend27_offset27);
  new variant = _ZX_2(variant);
stage RR:
  new base = GPR[registerZ];
  new target = GPR[registerY];
  new address = base + offset;
  new shift = (target & 31) * 8;
  new index = (target >> 5) & 7;
  new where = (registerGi << 3) + index;
  new dri = where;
stage E1:
  new argument1 = XVR[where];
  new bytemask = 15;
  new targetbytes = _ZX_256(bits2bytes(bytemask) << shift);
  CS.XMF = 1;
  new result1 = MEM.load(address, 0xF & bytemask, variant, dri);
  result1 = (_ZX_256(result1 << shift) & targetbytes) | (argument1 & _ZX_256(~targetbytes));
stage CR:
  XVR[where] = result1;
\end{lstlisting}}
\newpage

\paragraph{Encoding} Format \textsf{LSU\_XPL16RBB} (Section~\ref{sec:format-LSU-XPL16RBB}), \verb|lsucode5 = 010|\\

\bitfields{ 1/P, 2/01, 3/011, 2/variant, 2/registerGj, 4/0111, 2/11, 3/010, 1/--, 6/registerY, 6/registerZ }


\paragraph{Syntax} \texttt{xplw}\emph{variant} \emph{registerGj}, \emph{registerY} = [\emph{registerZ}]

\paragraph{Description} The \emph{registerGj} is preloaded from the memory word at the effective address.


\paragraph{Modifier(s)}
\noindent\begin{itemize}
\item variant (cf. variant in section \ref{par:modifier-variant} on page \pageref{par:modifier-variant})
\end{itemize}

\paragraph{Execution} {Reservation table \textsf{LSU} (Section~\ref{sec:reservation-LSU})}
~
{\begin{lstlisting}
stage ID:
  new variant = _ZX_2(variant);
stage RR:
  new address = GPR[registerZ];
  new target = GPR[registerY];
  new shift = (target & 31) * 8;
  new index = (target >> 5) & 15;
  new where = (registerGj << 4) + index;
  new dri = where;
stage E1:
  new argument1 = XVR[where];
  new bytemask = 15;
  new targetbytes = _ZX_256(bits2bytes(bytemask) << shift);
  CS.XMF = 1;
  new result1 = MEM.load(address, 0xF & bytemask, variant, dri);
  result1 = (_ZX_256(result1 << shift) & targetbytes) | (argument1 & _ZX_256(~targetbytes));
stage CR:
  XVR[where] = result1;
\end{lstlisting}}
\newpage

\paragraph{Encoding} Format \textsf{LSU\_XPL16RBB.O} (Section~\ref{sec:format-LSU-XPL16RBB.O}), \verb|lsucode5 = 010|\\

\bitfields{ 1/P, 2/00, 2/-- --, 27/offset27, 1/1, 2/01, 3/011, 2/variant, 2/registerGj, 4/0111, 2/11, 3/010, 1/--, 6/registerY, 6/registerZ }


\paragraph{Syntax} \texttt{xplw}\emph{variant} \emph{registerGj}, \emph{registerY} = \emph{offset27}[\emph{registerZ}]

\paragraph{Description} The \emph{registerGj} is preloaded from the memory word at the effective address.


\paragraph{Modifier(s)}
\noindent\begin{itemize}
\item variant (cf. variant in section \ref{par:modifier-variant} on page \pageref{par:modifier-variant})
\end{itemize}

\paragraph{Execution} {Reservation table \textsf{LSU.X} (Section~\ref{sec:reservation-LSU.X})}
~
{\begin{lstlisting}
stage ID:
  new offset = _SX_27(offset27);
  new variant = _ZX_2(variant);
stage RR:
  new base = GPR[registerZ];
  new target = GPR[registerY];
  new address = base + offset;
  new shift = (target & 31) * 8;
  new index = (target >> 5) & 15;
  new where = (registerGj << 4) + index;
  new dri = where;
stage E1:
  new argument1 = XVR[where];
  new bytemask = 15;
  new targetbytes = _ZX_256(bits2bytes(bytemask) << shift);
  CS.XMF = 1;
  new result1 = MEM.load(address, 0xF & bytemask, variant, dri);
  result1 = (_ZX_256(result1 << shift) & targetbytes) | (argument1 & _ZX_256(~targetbytes));
stage CR:
  XVR[where] = result1;
\end{lstlisting}}
\newpage

\paragraph{Encoding} Format \textsf{LSU\_XPL16RBB.D} (Section~\ref{sec:format-LSU-XPL16RBB.D}), \verb|lsucode5 = 010|\\

\bitfields{ 1/P, 2/00, 2/-- --, 27/extend27, 1/1, 2/00, 2/-- --, 27/offset27, 1/1, 2/01, 3/011, 2/variant, 2/registerGj, 4/0111, 2/11, 3/010, 1/--, 6/registerY, 6/registerZ }


\paragraph{Syntax} \texttt{xplw}\emph{variant} \emph{registerGj}, \emph{registerY} = \emph{extend27\_\discretionary{}{}{}offset27}[\emph{registerZ}]

\paragraph{Description} The \emph{registerGj} is preloaded from the memory word at the effective address.


\paragraph{Modifier(s)}
\noindent\begin{itemize}
\item variant (cf. variant in section \ref{par:modifier-variant} on page \pageref{par:modifier-variant})
\end{itemize}

\paragraph{Execution} {Reservation table \textsf{LSU.Y} (Section~\ref{sec:reservation-LSU.Y})}
~
{\begin{lstlisting}
stage ID:
  new offset = _SX_54(extend27_offset27);
  new variant = _ZX_2(variant);
stage RR:
  new base = GPR[registerZ];
  new target = GPR[registerY];
  new address = base + offset;
  new shift = (target & 31) * 8;
  new index = (target >> 5) & 15;
  new where = (registerGj << 4) + index;
  new dri = where;
stage E1:
  new argument1 = XVR[where];
  new bytemask = 15;
  new targetbytes = _ZX_256(bits2bytes(bytemask) << shift);
  CS.XMF = 1;
  new result1 = MEM.load(address, 0xF & bytemask, variant, dri);
  result1 = (_ZX_256(result1 << shift) & targetbytes) | (argument1 & _ZX_256(~targetbytes));
stage CR:
  XVR[where] = result1;
\end{lstlisting}}
\newpage

\paragraph{Encoding} Format \textsf{LSU\_XPL32RBB} (Section~\ref{sec:format-LSU-XPL32RBB}), \verb|lsucode5 = 010|\\

\bitfields{ 1/P, 2/01, 3/011, 2/variant, 1/registerGk, 5/01111, 2/11, 3/010, 1/--, 6/registerY, 6/registerZ }


\paragraph{Syntax} \texttt{xplw}\emph{variant} \emph{registerGk}, \emph{registerY} = [\emph{registerZ}]

\paragraph{Description} The \emph{registerGk} is preloaded from the memory word at the effective address.


\paragraph{Modifier(s)}
\noindent\begin{itemize}
\item variant (cf. variant in section \ref{par:modifier-variant} on page \pageref{par:modifier-variant})
\end{itemize}

\paragraph{Execution} {Reservation table \textsf{LSU} (Section~\ref{sec:reservation-LSU})}
~
{\begin{lstlisting}
stage ID:
  new variant = _ZX_2(variant);
stage RR:
  new address = GPR[registerZ];
  new target = GPR[registerY];
  new shift = (target & 31) * 8;
  new index = (target >> 5) & 31;
  new where = (registerGk << 5) + index;
  new dri = where;
stage E1:
  new argument1 = XVR[where];
  new bytemask = 15;
  new targetbytes = _ZX_256(bits2bytes(bytemask) << shift);
  CS.XMF = 1;
  new result1 = MEM.load(address, 0xF & bytemask, variant, dri);
  result1 = (_ZX_256(result1 << shift) & targetbytes) | (argument1 & _ZX_256(~targetbytes));
stage CR:
  XVR[where] = result1;
\end{lstlisting}}
\newpage

\paragraph{Encoding} Format \textsf{LSU\_XPL32RBB.O} (Section~\ref{sec:format-LSU-XPL32RBB.O}), \verb|lsucode5 = 010|\\

\bitfields{ 1/P, 2/00, 2/-- --, 27/offset27, 1/1, 2/01, 3/011, 2/variant, 1/registerGk, 5/01111, 2/11, 3/010, 1/--, 6/registerY, 6/registerZ }


\paragraph{Syntax} \texttt{xplw}\emph{variant} \emph{registerGk}, \emph{registerY} = \emph{offset27}[\emph{registerZ}]

\paragraph{Description} The \emph{registerGk} is preloaded from the memory word at the effective address.


\paragraph{Modifier(s)}
\noindent\begin{itemize}
\item variant (cf. variant in section \ref{par:modifier-variant} on page \pageref{par:modifier-variant})
\end{itemize}

\paragraph{Execution} {Reservation table \textsf{LSU.X} (Section~\ref{sec:reservation-LSU.X})}
~
{\begin{lstlisting}
stage ID:
  new offset = _SX_27(offset27);
  new variant = _ZX_2(variant);
stage RR:
  new base = GPR[registerZ];
  new target = GPR[registerY];
  new address = base + offset;
  new shift = (target & 31) * 8;
  new index = (target >> 5) & 31;
  new where = (registerGk << 5) + index;
  new dri = where;
stage E1:
  new argument1 = XVR[where];
  new bytemask = 15;
  new targetbytes = _ZX_256(bits2bytes(bytemask) << shift);
  CS.XMF = 1;
  new result1 = MEM.load(address, 0xF & bytemask, variant, dri);
  result1 = (_ZX_256(result1 << shift) & targetbytes) | (argument1 & _ZX_256(~targetbytes));
stage CR:
  XVR[where] = result1;
\end{lstlisting}}
\newpage

\paragraph{Encoding} Format \textsf{LSU\_XPL32RBB.D} (Section~\ref{sec:format-LSU-XPL32RBB.D}), \verb|lsucode5 = 010|\\

\bitfields{ 1/P, 2/00, 2/-- --, 27/extend27, 1/1, 2/00, 2/-- --, 27/offset27, 1/1, 2/01, 3/011, 2/variant, 1/registerGk, 5/01111, 2/11, 3/010, 1/--, 6/registerY, 6/registerZ }


\paragraph{Syntax} \texttt{xplw}\emph{variant} \emph{registerGk}, \emph{registerY} = \emph{extend27\_\discretionary{}{}{}offset27}[\emph{registerZ}]

\paragraph{Description} The \emph{registerGk} is preloaded from the memory word at the effective address.


\paragraph{Modifier(s)}
\noindent\begin{itemize}
\item variant (cf. variant in section \ref{par:modifier-variant} on page \pageref{par:modifier-variant})
\end{itemize}

\paragraph{Execution} {Reservation table \textsf{LSU.Y} (Section~\ref{sec:reservation-LSU.Y})}
~
{\begin{lstlisting}
stage ID:
  new offset = _SX_54(extend27_offset27);
  new variant = _ZX_2(variant);
stage RR:
  new base = GPR[registerZ];
  new target = GPR[registerY];
  new address = base + offset;
  new shift = (target & 31) * 8;
  new index = (target >> 5) & 31;
  new where = (registerGk << 5) + index;
  new dri = where;
stage E1:
  new argument1 = XVR[where];
  new bytemask = 15;
  new targetbytes = _ZX_256(bits2bytes(bytemask) << shift);
  CS.XMF = 1;
  new result1 = MEM.load(address, 0xF & bytemask, variant, dri);
  result1 = (_ZX_256(result1 << shift) & targetbytes) | (argument1 & _ZX_256(~targetbytes));
stage CR:
  XVR[where] = result1;
\end{lstlisting}}
\newpage

\paragraph{Encoding} Format \textsf{LSU\_XPL64RBB} (Section~\ref{sec:format-LSU-XPL64RBB}), \verb|lsucode5 = 010|\\

\bitfields{ 1/P, 2/01, 3/011, 2/variant, 1/registerGl, 5/11111, 2/11, 3/010, 1/--, 6/registerY, 6/registerZ }


\paragraph{Syntax} \texttt{xplw}\emph{variant} \emph{registerGl}, \emph{registerY} = [\emph{registerZ}]

\paragraph{Description} The \emph{registerGl} is preloaded from the memory word at the effective address.


\paragraph{Modifier(s)}
\noindent\begin{itemize}
\item variant (cf. variant in section \ref{par:modifier-variant} on page \pageref{par:modifier-variant})
\end{itemize}

\paragraph{Execution} {Reservation table \textsf{LSU} (Section~\ref{sec:reservation-LSU})}
~
{\begin{lstlisting}
stage ID:
  new variant = _ZX_2(variant);
stage RR:
  new address = GPR[registerZ];
  new target = GPR[registerY];
  new shift = (target & 31) * 8;
  new index = (target >> 5) & 63;
  new where = (registerGl << 6) + index;
  new dri = where;
stage E1:
  new argument1 = XVR[where];
  new bytemask = 15;
  new targetbytes = _ZX_256(bits2bytes(bytemask) << shift);
  CS.XMF = 1;
  new result1 = MEM.load(address, 0xF & bytemask, variant, dri);
  result1 = (_ZX_256(result1 << shift) & targetbytes) | (argument1 & _ZX_256(~targetbytes));
stage CR:
  XVR[where] = result1;
\end{lstlisting}}
\newpage

\paragraph{Encoding} Format \textsf{LSU\_XPL64RBB.O} (Section~\ref{sec:format-LSU-XPL64RBB.O}), \verb|lsucode5 = 010|\\

\bitfields{ 1/P, 2/00, 2/-- --, 27/offset27, 1/1, 2/01, 3/011, 2/variant, 1/registerGl, 5/11111, 2/11, 3/010, 1/--, 6/registerY, 6/registerZ }


\paragraph{Syntax} \texttt{xplw}\emph{variant} \emph{registerGl}, \emph{registerY} = \emph{offset27}[\emph{registerZ}]

\paragraph{Description} The \emph{registerGl} is preloaded from the memory word at the effective address.


\paragraph{Modifier(s)}
\noindent\begin{itemize}
\item variant (cf. variant in section \ref{par:modifier-variant} on page \pageref{par:modifier-variant})
\end{itemize}

\paragraph{Execution} {Reservation table \textsf{LSU.X} (Section~\ref{sec:reservation-LSU.X})}
~
{\begin{lstlisting}
stage ID:
  new offset = _SX_27(offset27);
  new variant = _ZX_2(variant);
stage RR:
  new base = GPR[registerZ];
  new target = GPR[registerY];
  new address = base + offset;
  new shift = (target & 31) * 8;
  new index = (target >> 5) & 63;
  new where = (registerGl << 6) + index;
  new dri = where;
stage E1:
  new argument1 = XVR[where];
  new bytemask = 15;
  new targetbytes = _ZX_256(bits2bytes(bytemask) << shift);
  CS.XMF = 1;
  new result1 = MEM.load(address, 0xF & bytemask, variant, dri);
  result1 = (_ZX_256(result1 << shift) & targetbytes) | (argument1 & _ZX_256(~targetbytes));
stage CR:
  XVR[where] = result1;
\end{lstlisting}}
\newpage

\paragraph{Encoding} Format \textsf{LSU\_XPL64RBB.D} (Section~\ref{sec:format-LSU-XPL64RBB.D}), \verb|lsucode5 = 010|\\

\bitfields{ 1/P, 2/00, 2/-- --, 27/extend27, 1/1, 2/00, 2/-- --, 27/offset27, 1/1, 2/01, 3/011, 2/variant, 1/registerGl, 5/11111, 2/11, 3/010, 1/--, 6/registerY, 6/registerZ }


\paragraph{Syntax} \texttt{xplw}\emph{variant} \emph{registerGl}, \emph{registerY} = \emph{extend27\_\discretionary{}{}{}offset27}[\emph{registerZ}]

\paragraph{Description} The \emph{registerGl} is preloaded from the memory word at the effective address.


\paragraph{Modifier(s)}
\noindent\begin{itemize}
\item variant (cf. variant in section \ref{par:modifier-variant} on page \pageref{par:modifier-variant})
\end{itemize}

\paragraph{Execution} {Reservation table \textsf{LSU.Y} (Section~\ref{sec:reservation-LSU.Y})}
~
{\begin{lstlisting}
stage ID:
  new offset = _SX_54(extend27_offset27);
  new variant = _ZX_2(variant);
stage RR:
  new base = GPR[registerZ];
  new target = GPR[registerY];
  new address = base + offset;
  new shift = (target & 31) * 8;
  new index = (target >> 5) & 63;
  new where = (registerGl << 6) + index;
  new dri = where;
stage E1:
  new argument1 = XVR[where];
  new bytemask = 15;
  new targetbytes = _ZX_256(bits2bytes(bytemask) << shift);
  CS.XMF = 1;
  new result1 = MEM.load(address, 0xF & bytemask, variant, dri);
  result1 = (_ZX_256(result1 << shift) & targetbytes) | (argument1 & _ZX_256(~targetbytes));
stage CR:
  XVR[where] = result1;
\end{lstlisting}}
\newpage
\subsubsection[{\small XSO: Extension Store Octuple Word}]{XSO\hfill Extension Store Octuple Word}
\label{instr:XSO}\index{xso}
 
\paragraph{Encoding} Format \textsf{LSU\_XSOBO} (Section~\ref{sec:format-LSU-XSOBO}), \verb|steering = 01|\\

\bitfields{ 1/P, 2/01, 5/10101, 6/registerE, 2/01, 10/signed10, 6/registerZ }


\paragraph{Syntax} \texttt{xso} \emph{signed10}[\emph{registerZ}] = \emph{registerE}

\paragraph{Description} The effective address is computed by adding the \emph{signed10} to the \emph{registerZ}.
The \emph{registerE} is stored into the memory octuple word at the effective address.


\paragraph{Execution} {Reservation table \textsf{LSU\_MEMW\_ACCR} (Section~\ref{sec:reservation-LSU-MEMW-ACCR})}
~
{\begin{lstlisting}
stage ID:
  new offset = _SX_10(signed10);
stage RR:
  new bytemask = 0xFFFFFFFF;
  new dri = registerE;
  new base = GPR[registerZ];
  new address = base + offset;
stage E1:
  new argument3 = XVR[registerE];
  CS.XMF = 1;
stage E3:
  MEM.store(address, 0xFFFFFFFF & bytemask, 0, argument3, dri);
\end{lstlisting}}
\newpage

\paragraph{Encoding} Format \textsf{LSU\_XSOBO.X} (Section~\ref{sec:format-LSU-XSOBO.X}), \verb|steering = 01|\\

\bitfields{ 1/P, 2/00, 2/-- --, 27/upper27, 1/1, 2/01, 5/10101, 6/registerE, 2/01, 10/lower10, 6/registerZ }


\paragraph{Syntax} \texttt{xso} \emph{upper27\_\discretionary{}{}{}lower10}[\emph{registerZ}] = \emph{registerE}

\paragraph{Description} The effective address is computed by adding the \emph{upper27\_\discretionary{}{}{}lower10} to the \emph{registerZ}.
The \emph{registerE} is stored into the memory octuple word at the effective address.


\paragraph{Execution} {Reservation table \textsf{LSU\_MEMW\_ACCR.X} (Section~\ref{sec:reservation-LSU-MEMW-ACCR.X})}
~
{\begin{lstlisting}
stage ID:
  new offset = _SX_37(upper27_lower10);
stage RR:
  new bytemask = 0xFFFFFFFF;
  new dri = registerE;
  new base = GPR[registerZ];
  new address = base + offset;
stage E1:
  new argument3 = XVR[registerE];
  CS.XMF = 1;
stage E3:
  MEM.store(address, 0xFFFFFFFF & bytemask, 0, argument3, dri);
\end{lstlisting}}
\newpage

\paragraph{Encoding} Format \textsf{LSU\_XSOBO.Y} (Section~\ref{sec:format-LSU-XSOBO.Y}), \verb|steering = 01|\\

\bitfields{ 1/P, 2/00, 2/-- --, 27/extend27, 1/1, 2/00, 2/-- --, 27/upper27, 1/1, 2/01, 5/10101, 6/registerE, 2/01, 10/lower10, 6/registerZ }


\paragraph{Syntax} \texttt{xso} \emph{extend27\_\discretionary{}{}{}upper27\_\discretionary{}{}{}lower10}[\emph{registerZ}] = \emph{registerE}

\paragraph{Description} The effective address is computed by adding the \emph{extend27\_\discretionary{}{}{}upper27\_\discretionary{}{}{}lower10} to the \emph{registerZ}.
The \emph{registerE} is stored into the memory octuple word at the effective address.


\paragraph{Execution} {Reservation table \textsf{LSU\_MEMW\_ACCR.Y} (Section~\ref{sec:reservation-LSU-MEMW-ACCR.Y})}
~
{\begin{lstlisting}
stage ID:
  new offset = _WX_64(extend27_upper27_lower10);
stage RR:
  new bytemask = 0xFFFFFFFF;
  new dri = registerE;
  new base = GPR[registerZ];
  new address = base + offset;
stage E1:
  new argument3 = XVR[registerE];
  CS.XMF = 1;
stage E3:
  MEM.store(address, 0xFFFFFFFF & bytemask, 0, argument3, dri);
\end{lstlisting}}
\newpage

\paragraph{Encoding} Format \textsf{LSU\_XSOBI} (Section~\ref{sec:format-LSU-XSOBI}), \verb|steering = 01|\\

\bitfields{ 1/P, 2/01, 5/10101, 6/registerE, 2/11, 3/111, 1/--, 6/registerY, 6/registerZ }


\paragraph{Syntax} \texttt{xso} \emph{registerY}[\emph{registerZ}] = \emph{registerE}

\paragraph{Description} The effective address is computed by adding the \emph{registerZ} to the \emph{registerY} and scaled by the \emph{}.
The \emph{registerE} is stored into the memory octuple word at the effective address.


\paragraph{Execution} {Reservation table \textsf{LSU\_MEMW\_ACCR} (Section~\ref{sec:reservation-LSU-MEMW-ACCR})}
~
{\begin{lstlisting}
stage RR:
  new bytemask = 0xFFFFFFFF;
  new doscale = 0;
  new dri = registerE;
  new index = GPR[registerY];
  new base = GPR[registerZ];
  new scaling = doscale ? 32 : 1;
  new address = base + (index * scaling);
stage E1:
  new argument3 = XVR[registerE];
  CS.XMF = 1;
stage E3:
  MEM.store(address, 0xFFFFFFFF & bytemask, 0, argument3, dri);
\end{lstlisting}}
\newpage

\paragraph{Encoding} Format \textsf{LSU\_XSOC4BO} (Section~\ref{sec:format-LSU-XSOC4BO}), \verb|steering = 01|\\

\bitfields{ 1/P, 2/01, 5/10111, 4/registerEq, 2/qindex, 2/01, 10/signed10, 6/registerZ }


\paragraph{Syntax} \texttt{xso}\emph{qindex} \emph{signed10}[\emph{registerZ}] = \emph{registerEq}

\paragraph{Description} The effective address is computed by adding the \emph{signed10} to the \emph{registerZ}.
The \emph{registerEq} is stored into the memory octuple word at the effective address.


\paragraph{Modifier(s)}
\noindent\begin{itemize}
\item qindex (cf. qindex in section \ref{par:modifier-qindex} on page \pageref{par:modifier-qindex})
\end{itemize}

\paragraph{Execution} {Reservation table \textsf{LSU\_MEMW\_ACCR} (Section~\ref{sec:reservation-LSU-MEMW-ACCR})}
~
{\begin{lstlisting}
stage ID:
  new qindex = _ZX_2(qindex);
  new offset = _SX_10(signed10);
stage RR:
  new bytemask = 0xFFFFFFFF;
  new dri = registerEq << 2;
  new base = GPR[registerZ];
  new address = base + offset;
stage E1:
  new argument3_0 = XMR[registerEq]:0;
  new argument3_0_0 = argument3_0.64[0];
  new argument3_0_1 = argument3_0.64[1];
  new argument3_0_2 = argument3_0.64[2];
  new argument3_0_3 = argument3_0.64[3];
  new argument3_1 = XMR[registerEq]:1;
  new argument3_1_0 = argument3_1.64[0];
  new argument3_1_1 = argument3_1.64[1];
  new argument3_1_2 = argument3_1.64[2];
  new argument3_1_3 = argument3_1.64[3];
  new argument3_2 = XMR[registerEq]:2;
  new argument3_2_0 = argument3_2.64[0];
  new argument3_2_1 = argument3_2.64[1];
  new argument3_2_2 = argument3_2.64[2];
  new argument3_2_3 = argument3_2.64[3];
  new argument3_3 = XMR[registerEq]:3;
  new argument3_3_0 = argument3_3.64[0];
  new argument3_3_1 = argument3_3.64[1];
  new argument3_3_2 = argument3_3.64[2];
  new argument3_3_3 = argument3_3.64[3];
  new argument3 = 0;
  if (qindex == 0) {
    argument3 = join_64_x4(argument3_0_0, argument3_1_0, argument3_2_0, argument3_3_0);
  }
  if (qindex == 1) {
    argument3 = join_64_x4(argument3_0_1, argument3_1_1, argument3_2_1, argument3_3_1);
  }
  if (qindex == 2) {
    argument3 = join_64_x4(argument3_0_2, argument3_1_2, argument3_2_2, argument3_3_2);
  }
  if (qindex == 3) {
    argument3 = join_64_x4(argument3_0_3, argument3_1_3, argument3_2_3, argument3_3_3);
  }
  CS.XMF = 1;
stage E3:
  MEM.store(address, 0xFFFFFFFF & bytemask, 0, argument3, dri);
\end{lstlisting}}
\newpage

\paragraph{Encoding} Format \textsf{LSU\_XSOC4BO.X} (Section~\ref{sec:format-LSU-XSOC4BO.X}), \verb|steering = 01|\\

\bitfields{ 1/P, 2/00, 2/-- --, 27/upper27, 1/1, 2/01, 5/10111, 4/registerEq, 2/qindex, 2/01, 10/lower10, 6/registerZ }


\paragraph{Syntax} \texttt{xso}\emph{qindex} \emph{upper27\_\discretionary{}{}{}lower10}[\emph{registerZ}] = \emph{registerEq}

\paragraph{Description} The effective address is computed by adding the \emph{upper27\_\discretionary{}{}{}lower10} to the \emph{registerZ}.
The \emph{registerEq} is stored into the memory octuple word at the effective address.


\paragraph{Modifier(s)}
\noindent\begin{itemize}
\item qindex (cf. qindex in section \ref{par:modifier-qindex} on page \pageref{par:modifier-qindex})
\end{itemize}

\paragraph{Execution} {Reservation table \textsf{LSU\_MEMW\_ACCR.X} (Section~\ref{sec:reservation-LSU-MEMW-ACCR.X})}
~
{\begin{lstlisting}
stage ID:
  new qindex = _ZX_2(qindex);
  new offset = _SX_37(upper27_lower10);
stage RR:
  new bytemask = 0xFFFFFFFF;
  new dri = registerEq << 2;
  new base = GPR[registerZ];
  new address = base + offset;
stage E1:
  new argument3_0 = XMR[registerEq]:0;
  new argument3_0_0 = argument3_0.64[0];
  new argument3_0_1 = argument3_0.64[1];
  new argument3_0_2 = argument3_0.64[2];
  new argument3_0_3 = argument3_0.64[3];
  new argument3_1 = XMR[registerEq]:1;
  new argument3_1_0 = argument3_1.64[0];
  new argument3_1_1 = argument3_1.64[1];
  new argument3_1_2 = argument3_1.64[2];
  new argument3_1_3 = argument3_1.64[3];
  new argument3_2 = XMR[registerEq]:2;
  new argument3_2_0 = argument3_2.64[0];
  new argument3_2_1 = argument3_2.64[1];
  new argument3_2_2 = argument3_2.64[2];
  new argument3_2_3 = argument3_2.64[3];
  new argument3_3 = XMR[registerEq]:3;
  new argument3_3_0 = argument3_3.64[0];
  new argument3_3_1 = argument3_3.64[1];
  new argument3_3_2 = argument3_3.64[2];
  new argument3_3_3 = argument3_3.64[3];
  new argument3 = 0;
  if (qindex == 0) {
    argument3 = join_64_x4(argument3_0_0, argument3_1_0, argument3_2_0, argument3_3_0);
  }
  if (qindex == 1) {
    argument3 = join_64_x4(argument3_0_1, argument3_1_1, argument3_2_1, argument3_3_1);
  }
  if (qindex == 2) {
    argument3 = join_64_x4(argument3_0_2, argument3_1_2, argument3_2_2, argument3_3_2);
  }
  if (qindex == 3) {
    argument3 = join_64_x4(argument3_0_3, argument3_1_3, argument3_2_3, argument3_3_3);
  }
  CS.XMF = 1;
stage E3:
  MEM.store(address, 0xFFFFFFFF & bytemask, 0, argument3, dri);
\end{lstlisting}}
\newpage

\paragraph{Encoding} Format \textsf{LSU\_XSOC4BO.Y} (Section~\ref{sec:format-LSU-XSOC4BO.Y}), \verb|steering = 01|\\

\bitfields{ 1/P, 2/00, 2/-- --, 27/extend27, 1/1, 2/00, 2/-- --, 27/upper27, 1/1, 2/01, 5/10111, 4/registerEq, 2/qindex, 2/01, 10/lower10, 6/registerZ }


\paragraph{Syntax} \texttt{xso}\emph{qindex} \emph{extend27\_\discretionary{}{}{}upper27\_\discretionary{}{}{}lower10}[\emph{registerZ}] = \emph{registerEq}

\paragraph{Description} The effective address is computed by adding the \emph{extend27\_\discretionary{}{}{}upper27\_\discretionary{}{}{}lower10} to the \emph{registerZ}.
The \emph{registerEq} is stored into the memory octuple word at the effective address.


\paragraph{Modifier(s)}
\noindent\begin{itemize}
\item qindex (cf. qindex in section \ref{par:modifier-qindex} on page \pageref{par:modifier-qindex})
\end{itemize}

\paragraph{Execution} {Reservation table \textsf{LSU\_MEMW\_ACCR.Y} (Section~\ref{sec:reservation-LSU-MEMW-ACCR.Y})}
~
{\begin{lstlisting}
stage ID:
  new qindex = _ZX_2(qindex);
  new offset = _WX_64(extend27_upper27_lower10);
stage RR:
  new bytemask = 0xFFFFFFFF;
  new dri = registerEq << 2;
  new base = GPR[registerZ];
  new address = base + offset;
stage E1:
  new argument3_0 = XMR[registerEq]:0;
  new argument3_0_0 = argument3_0.64[0];
  new argument3_0_1 = argument3_0.64[1];
  new argument3_0_2 = argument3_0.64[2];
  new argument3_0_3 = argument3_0.64[3];
  new argument3_1 = XMR[registerEq]:1;
  new argument3_1_0 = argument3_1.64[0];
  new argument3_1_1 = argument3_1.64[1];
  new argument3_1_2 = argument3_1.64[2];
  new argument3_1_3 = argument3_1.64[3];
  new argument3_2 = XMR[registerEq]:2;
  new argument3_2_0 = argument3_2.64[0];
  new argument3_2_1 = argument3_2.64[1];
  new argument3_2_2 = argument3_2.64[2];
  new argument3_2_3 = argument3_2.64[3];
  new argument3_3 = XMR[registerEq]:3;
  new argument3_3_0 = argument3_3.64[0];
  new argument3_3_1 = argument3_3.64[1];
  new argument3_3_2 = argument3_3.64[2];
  new argument3_3_3 = argument3_3.64[3];
  new argument3 = 0;
  if (qindex == 0) {
    argument3 = join_64_x4(argument3_0_0, argument3_1_0, argument3_2_0, argument3_3_0);
  }
  if (qindex == 1) {
    argument3 = join_64_x4(argument3_0_1, argument3_1_1, argument3_2_1, argument3_3_1);
  }
  if (qindex == 2) {
    argument3 = join_64_x4(argument3_0_2, argument3_1_2, argument3_2_2, argument3_3_2);
  }
  if (qindex == 3) {
    argument3 = join_64_x4(argument3_0_3, argument3_1_3, argument3_2_3, argument3_3_3);
  }
  CS.XMF = 1;
stage E3:
  MEM.store(address, 0xFFFFFFFF & bytemask, 0, argument3, dri);
\end{lstlisting}}
\newpage

\paragraph{Encoding} Format \textsf{LSU\_XSOC4BI} (Section~\ref{sec:format-LSU-XSOC4BI}), \verb|steering = 01|\\

\bitfields{ 1/P, 2/01, 5/10111, 4/registerEq, 2/qindex, 2/11, 3/111, 1/--, 6/registerY, 6/registerZ }


\paragraph{Syntax} \texttt{xso}\emph{qindex} \emph{registerY}[\emph{registerZ}] = \emph{registerEq}

\paragraph{Description} The effective address is computed by adding the \emph{registerZ} to the \emph{registerY} and scaled by the \emph{qindex}.
The \emph{registerEq} is stored into the memory octuple word at the effective address.


\paragraph{Modifier(s)}
\noindent\begin{itemize}
\item qindex (cf. qindex in section \ref{par:modifier-qindex} on page \pageref{par:modifier-qindex})
\end{itemize}

\paragraph{Execution} {Reservation table \textsf{LSU\_MEMW\_ACCR} (Section~\ref{sec:reservation-LSU-MEMW-ACCR})}
~
{\begin{lstlisting}
stage ID:
  new qindex = _ZX_2(qindex);
stage RR:
  new bytemask = 0xFFFFFFFF;
  new doscale = 0;
  new dri = registerEq >> 2;
  new index = GPR[registerY];
  new base = GPR[registerZ];
  new scaling = doscale ? 32 : 1;
  new address = base + (index * scaling);
stage E1:
  new argument3_0 = XMR[registerEq]:0;
  new argument3_0_0 = argument3_0.64[0];
  new argument3_0_1 = argument3_0.64[1];
  new argument3_0_2 = argument3_0.64[2];
  new argument3_0_3 = argument3_0.64[3];
  new argument3_1 = XMR[registerEq]:1;
  new argument3_1_0 = argument3_1.64[0];
  new argument3_1_1 = argument3_1.64[1];
  new argument3_1_2 = argument3_1.64[2];
  new argument3_1_3 = argument3_1.64[3];
  new argument3_2 = XMR[registerEq]:2;
  new argument3_2_0 = argument3_2.64[0];
  new argument3_2_1 = argument3_2.64[1];
  new argument3_2_2 = argument3_2.64[2];
  new argument3_2_3 = argument3_2.64[3];
  new argument3_3 = XMR[registerEq]:3;
  new argument3_3_0 = argument3_3.64[0];
  new argument3_3_1 = argument3_3.64[1];
  new argument3_3_2 = argument3_3.64[2];
  new argument3_3_3 = argument3_3.64[3];
  new argument3 = 0;
  if (qindex == 0) {
    argument3 = join_64_x4(argument3_0_0, argument3_1_0, argument3_2_0, argument3_3_0);
  }
  if (qindex == 1) {
    argument3 = join_64_x4(argument3_0_1, argument3_1_1, argument3_2_1, argument3_3_1);
  }
  if (qindex == 2) {
    argument3 = join_64_x4(argument3_0_2, argument3_1_2, argument3_2_2, argument3_3_2);
  }
  if (qindex == 3) {
    argument3 = join_64_x4(argument3_0_3, argument3_1_3, argument3_2_3, argument3_3_3);
  }
  CS.XMF = 1;
stage E3:
  MEM.store(address, 0xFFFFFFFF & bytemask, 0, argument3, dri);
\end{lstlisting}}
\newpage
\subsection{ALU Instructions}
\label{sec:ALU-instructions}

\subsubsection[{\small ABDBX: Absolute Difference Byte Hexadecuple}]{ABDBX\hfill Absolute Difference Byte Hexadecuple}
\label{instr:ABDBX}\index{abdbx}
 
\paragraph{Encoding} Format \textsf{ALU\_BXWRRS} (Section~\ref{sec:format-ALU-BXWRRS}), \verb|exucode2 = 1100|\\

\bitfields{ 1/P, 2/11, 1/1, 4/1100, 5/registerM, 1/--, 2/10, 3/010, 1/1, 5/registerO, 1/--, 5/registerP, 1/1 }


\paragraph{Syntax} \texttt{abdbx} \emph{registerM} = \emph{registerP}, \emph{registerO}

\paragraph{Description} The \emph{registerP} and the \emph{registerO} are interpreted as sixteen bytes packed into 128 bits. The \emph{registerP} bytes are subtracted from the \emph{registerO} bytes. The absolute value of the sixteen resulting bytes are packed and stored into the \emph{registerM}.


\paragraph{Execution} {Reservation table \textsf{ALU\_LITE} (Section~\ref{sec:reservation-ALU-LITE})}
~
{\begin{lstlisting}
stage RR:
  new argument3 = PGR[registerO];
  new argument2 = PGR[registerP];
  for (I = 0; I < 16; I += 1) {
    result1.8[I] = _ABS(_SX_8(argument3.8[I]) - _SX_8(argument2.8[I]))
  };
stage E1:
  PGR[registerM] = result1;
\end{lstlisting}}
\newpage

\paragraph{Encoding} Format \textsf{ALU\_BXWRRS.M} (Section~\ref{sec:format-ALU-BXWRRS.M}), \verb|exucode2 = 1100|\\

\bitfields{ 1/P, 2/00, 2/-- --, 27/upper27, 1/1, 2/11, 1/1, 4/1100, 5/registerM, 1/--, 2/10, 3/010, 1/1, 1/splat32, 5/lower5, 5/registerP, 1/1 }


\paragraph{Syntax} \texttt{abdbx} \emph{registerM} = \emph{registerP}, \emph{upper27\_\discretionary{}{}{}lower5}\emph{splat32}

\paragraph{Description} The \emph{registerP} and the \emph{upper27\_\discretionary{}{}{}lower5} are interpreted as sixteen bytes packed into 128 bits. The \emph{registerP} bytes are subtracted from the \emph{upper27\_\discretionary{}{}{}lower5} bytes. The absolute value of the sixteen resulting bytes are packed and stored into the \emph{registerM}.


\paragraph{Modifier(s)}
\noindent\begin{itemize}
\item splat32 (cf. splat32 in section \ref{par:modifier-splat32} on page \pageref{par:modifier-splat32})
\end{itemize}

\paragraph{Execution} {Reservation table \textsf{ALU\_LITE.X} (Section~\ref{sec:reservation-ALU-LITE.X})}
~
{\begin{lstlisting}
stage ID:
  new splat32 = _ZX_1(splat32);
stage RR:
  new argument3 = splat32 ? (_WX_32(upper27_lower5) << 32) | _ZX_32(_WX_32(upper27_lower5)) : _SX_32(_WX_32(upper27_lower5));
  new argument2 = PGR[registerP];
  for (I = 0; I < 16; I += 1) {
    result1.8[I] = _ABS(_SX_8(argument3.8[I]) - _SX_8(argument2.8[I]))
  };
stage E1:
  PGR[registerM] = result1;
\end{lstlisting}}
\newpage
\subsubsection[{\small ABDD: Absolute Difference Double Word}]{ABDD\hfill Absolute Difference Double Word}
\label{instr:ABDD}\index{abdd}
 
\paragraph{Encoding} Format \textsf{ALU\_DWRRS} (Section~\ref{sec:format-ALU-DWRRS}), \verb|exucode2 = 1100|\\

\bitfields{ 1/P, 2/11, 1/0, 4/1100, 6/registerW, 2/10, 3/010, 1/0, 6/registerY, 6/registerZ }


\paragraph{Syntax} \texttt{abdd} \emph{registerW} = \emph{registerZ}, \emph{registerY}

\paragraph{Description} The \emph{registerZ} is subtracted from the \emph{registerY}. The absolute value of the result is stored into the \emph{registerW}.


\paragraph{Execution} {Reservation table \textsf{ALU\_TINY} (Section~\ref{sec:reservation-ALU-TINY})}
~
{\begin{lstlisting}
stage RR:
  new argument3 = GPR[registerY];
  new argument2 = GPR[registerZ];
  new result1 = _ABS(_SX_64(argument3) - _SX_64(argument2));
stage E1:
  GPR[registerW] = result1;
\end{lstlisting}}
\newpage

\paragraph{Encoding} Format \textsf{ALU\_DWRRS.M} (Section~\ref{sec:format-ALU-DWRRS.M}), \verb|exucode2 = 1100|\\

\bitfields{ 1/P, 2/00, 2/-- --, 27/upper27, 1/1, 2/11, 1/0, 4/1100, 6/registerW, 2/10, 3/010, 1/0, 1/splat32, 5/lower5, 6/registerZ }


\paragraph{Syntax} \texttt{abdd} \emph{registerW} = \emph{registerZ}, \emph{upper27\_\discretionary{}{}{}lower5}\emph{splat32}

\paragraph{Description} The \emph{registerZ} is subtracted from the \emph{upper27\_\discretionary{}{}{}lower5}. The absolute value of the result is stored into the \emph{registerW}.


\paragraph{Modifier(s)}
\noindent\begin{itemize}
\item splat32 (cf. splat32 in section \ref{par:modifier-splat32} on page \pageref{par:modifier-splat32})
\end{itemize}

\paragraph{Execution} {Reservation table \textsf{ALU\_TINY.X} (Section~\ref{sec:reservation-ALU-TINY.X})}
~
{\begin{lstlisting}
stage ID:
  new splat32 = _ZX_1(splat32);
stage RR:
  new argument3 = splat32 ? (_WX_32(upper27_lower5) << 32) | _ZX_32(_WX_32(upper27_lower5)) : _SX_32(_WX_32(upper27_lower5));
  new argument2 = GPR[registerZ];
  new result1 = _ABS(_SX_64(argument3) - _SX_64(argument2));
stage E1:
  GPR[registerW] = result1;
\end{lstlisting}}
\newpage
\subsubsection[{\small ABDDP: Absolute Difference Double Word Pair}]{ABDDP\hfill Absolute Difference Double Word Pair}
\label{instr:ABDDP}\index{abddp}
 
\paragraph{Encoding} Format \textsf{ALU\_DPWRRS} (Section~\ref{sec:format-ALU-DPWRRS}), \verb|exucode2 = 1100|\\

\bitfields{ 1/P, 2/11, 1/1, 4/1100, 5/registerM, 1/--, 2/10, 3/010, 1/0, 5/registerO, 1/--, 5/registerP, 1/0 }


\paragraph{Syntax} \texttt{abddp} \emph{registerM} = \emph{registerP}, \emph{registerO}

\paragraph{Description} The \emph{registerP} and the \emph{registerO} are interpreted as two double word pairs packed into 128 bits. The \emph{registerP} double word pairs are subtracted from the \emph{registerO} double word pairs. The absolute value of the two resulting double words are packed and stored into the \emph{registerM}.


\paragraph{Execution} {Reservation table \textsf{ALU\_LITE} (Section~\ref{sec:reservation-ALU-LITE})}
~
{\begin{lstlisting}
stage RR:
  new argument3 = PGR[registerO];
  new argument2 = PGR[registerP];
  for (I = 0; I < 2; I += 1) {
    result1.64[I] = _ABS(_SX_64(argument3.64[I]) - _SX_64(argument2.64[I]))
  };
stage E1:
  PGR[registerM] = result1;
\end{lstlisting}}
\newpage

\paragraph{Encoding} Format \textsf{ALU\_DPWRRS.M} (Section~\ref{sec:format-ALU-DPWRRS.M}), \verb|exucode2 = 1100|\\

\bitfields{ 1/P, 2/00, 2/-- --, 27/upper27, 1/1, 2/11, 1/1, 4/1100, 5/registerM, 1/--, 2/10, 3/010, 1/0, 1/splat32, 5/lower5, 5/registerP, 1/0 }


\paragraph{Syntax} \texttt{abddp} \emph{registerM} = \emph{registerP}, \emph{upper27\_\discretionary{}{}{}lower5}\emph{splat32}

\paragraph{Description} The \emph{registerP} and the \emph{upper27\_\discretionary{}{}{}lower5} are interpreted as two double word pairs packed into 128 bits. The \emph{registerP} double word pairs are subtracted from the \emph{upper27\_\discretionary{}{}{}lower5} double word pairs. The absolute value of the two resulting double words are packed and stored into the \emph{registerM}.


\paragraph{Modifier(s)}
\noindent\begin{itemize}
\item splat32 (cf. splat32 in section \ref{par:modifier-splat32} on page \pageref{par:modifier-splat32})
\end{itemize}

\paragraph{Execution} {Reservation table \textsf{ALU\_LITE.X} (Section~\ref{sec:reservation-ALU-LITE.X})}
~
{\begin{lstlisting}
stage ID:
  new splat32 = _ZX_1(splat32);
stage RR:
  new argument3 = splat32 ? (_WX_32(upper27_lower5) << 32) | _ZX_32(_WX_32(upper27_lower5)) : _SX_32(_WX_32(upper27_lower5));
  new argument2 = PGR[registerP];
  for (I = 0; I < 2; I += 1) {
    result1.64[I] = _ABS(_SX_64(argument3.64[I]) - _SX_64(argument2.64[I]))
  };
stage E1:
  PGR[registerM] = result1;
\end{lstlisting}}
\newpage
\subsubsection[{\small ABDHO: Absolute Difference Half Word Octuple}]{ABDHO\hfill Absolute Difference Half Word Octuple}
\label{instr:ABDHO}\index{abdho}
 
\paragraph{Encoding} Format \textsf{ALU\_HOWRRS} (Section~\ref{sec:format-ALU-HOWRRS}), \verb|exucode2 = 1100|\\

\bitfields{ 1/P, 2/11, 1/1, 4/1100, 5/registerM, 1/--, 2/10, 3/010, 1/1, 5/registerO, 1/--, 5/registerP, 1/0 }


\paragraph{Syntax} \texttt{abdho} \emph{registerM} = \emph{registerP}, \emph{registerO}

\paragraph{Description} The \emph{registerP} and the \emph{registerO} are interpreted as eight half words packed into 128 bits. The \emph{registerP} half words are subtracted from the \emph{registerO} half words. The absolute value of the eight resulting half words are packed and stored into the \emph{registerM}.


\paragraph{Execution} {Reservation table \textsf{ALU\_LITE} (Section~\ref{sec:reservation-ALU-LITE})}
~
{\begin{lstlisting}
stage RR:
  new argument3 = PGR[registerO];
  new argument2 = PGR[registerP];
  for (I = 0; I < 8; I += 1) {
    result1.16[I] = _ABS(_SX_16(argument3.16[I]) - _SX_16(argument2.16[I]))
  };
stage E1:
  PGR[registerM] = result1;
\end{lstlisting}}
\newpage

\paragraph{Encoding} Format \textsf{ALU\_HOWRRS.M} (Section~\ref{sec:format-ALU-HOWRRS.M}), \verb|exucode2 = 1100|\\

\bitfields{ 1/P, 2/00, 2/-- --, 27/upper27, 1/1, 2/11, 1/1, 4/1100, 5/registerM, 1/--, 2/10, 3/010, 1/1, 1/splat32, 5/lower5, 5/registerP, 1/0 }


\paragraph{Syntax} \texttt{abdho} \emph{registerM} = \emph{registerP}, \emph{upper27\_\discretionary{}{}{}lower5}\emph{splat32}

\paragraph{Description} The \emph{registerP} and the \emph{upper27\_\discretionary{}{}{}lower5} are interpreted as eight half words packed into 128 bits. The \emph{registerP} half words are subtracted from the \emph{upper27\_\discretionary{}{}{}lower5} half words. The absolute value of the eight resulting half words are packed and stored into the \emph{registerM}.


\paragraph{Modifier(s)}
\noindent\begin{itemize}
\item splat32 (cf. splat32 in section \ref{par:modifier-splat32} on page \pageref{par:modifier-splat32})
\end{itemize}

\paragraph{Execution} {Reservation table \textsf{ALU\_LITE.X} (Section~\ref{sec:reservation-ALU-LITE.X})}
~
{\begin{lstlisting}
stage ID:
  new splat32 = _ZX_1(splat32);
stage RR:
  new argument3 = splat32 ? (_WX_32(upper27_lower5) << 32) | _ZX_32(_WX_32(upper27_lower5)) : _SX_32(_WX_32(upper27_lower5));
  new argument2 = PGR[registerP];
  for (I = 0; I < 8; I += 1) {
    result1.16[I] = _ABS(_SX_16(argument3.16[I]) - _SX_16(argument2.16[I]))
  };
stage E1:
  PGR[registerM] = result1;
\end{lstlisting}}
\newpage
\subsubsection[{\small ABDSBX: Absolute Difference Saturated Byte Hexadecuple}]{ABDSBX\hfill Absolute Difference Saturated Byte Hexadecuple}
\label{instr:ABDSBX}\index{abdsbx}
 
\paragraph{Encoding} Format \textsf{ALU\_BXWRRS} (Section~\ref{sec:format-ALU-BXWRRS}), \verb|exucode2 = 1110|\\

\bitfields{ 1/P, 2/11, 1/1, 4/1110, 5/registerM, 1/--, 2/10, 3/010, 1/1, 5/registerO, 1/--, 5/registerP, 1/1 }


\paragraph{Syntax} \texttt{abdsbx} \emph{registerM} = \emph{registerP}, \emph{registerO}

\paragraph{Description} The \emph{registerP} and the \emph{registerO} are interpreted as sixteen bytes packed into 128 bits. The word quadruple \emph{registerP} is subtracted from the word quadruple \emph{registerO}. The absolute values of the sixteen resulting bytes are saturated and stored into the \emph{registerM}.


\paragraph{Execution} {Reservation table \textsf{ALU\_LITE} (Section~\ref{sec:reservation-ALU-LITE})}
~
{\begin{lstlisting}
stage RR:
  new argument3 = PGR[registerO];
  new argument2 = PGR[registerP];
  for (I = 0; I < 16; I += 1) {
    result1.8[I] = _SAT_8(_ABS(_SX_8(argument3.8[I]) - _SX_8(argument2.8[I])))
  };
stage E1:
  PGR[registerM] = result1;
\end{lstlisting}}
\newpage

\paragraph{Encoding} Format \textsf{ALU\_BXWRRS.M} (Section~\ref{sec:format-ALU-BXWRRS.M}), \verb|exucode2 = 1110|\\

\bitfields{ 1/P, 2/00, 2/-- --, 27/upper27, 1/1, 2/11, 1/1, 4/1110, 5/registerM, 1/--, 2/10, 3/010, 1/1, 1/splat32, 5/lower5, 5/registerP, 1/1 }


\paragraph{Syntax} \texttt{abdsbx} \emph{registerM} = \emph{registerP}, \emph{upper27\_\discretionary{}{}{}lower5}\emph{splat32}

\paragraph{Description} The \emph{registerP} and the \emph{upper27\_\discretionary{}{}{}lower5} are interpreted as sixteen bytes packed into 128 bits. The word quadruple \emph{registerP} is subtracted from the word quadruple \emph{upper27\_\discretionary{}{}{}lower5}. The absolute values of the sixteen resulting bytes are saturated and stored into the \emph{registerM}.


\paragraph{Modifier(s)}
\noindent\begin{itemize}
\item splat32 (cf. splat32 in section \ref{par:modifier-splat32} on page \pageref{par:modifier-splat32})
\end{itemize}

\paragraph{Execution} {Reservation table \textsf{ALU\_LITE.X} (Section~\ref{sec:reservation-ALU-LITE.X})}
~
{\begin{lstlisting}
stage ID:
  new splat32 = _ZX_1(splat32);
stage RR:
  new argument3 = splat32 ? (_WX_32(upper27_lower5) << 32) | _ZX_32(_WX_32(upper27_lower5)) : _SX_32(_WX_32(upper27_lower5));
  new argument2 = PGR[registerP];
  for (I = 0; I < 16; I += 1) {
    result1.8[I] = _SAT_8(_ABS(_SX_8(argument3.8[I]) - _SX_8(argument2.8[I])))
  };
stage E1:
  PGR[registerM] = result1;
\end{lstlisting}}
\newpage
\subsubsection[{\small ABDSD: Absolute Difference Saturated Double Word}]{ABDSD\hfill Absolute Difference Saturated Double Word}
\label{instr:ABDSD}\index{abdsd}
 
\paragraph{Encoding} Format \textsf{ALU\_DWRRS} (Section~\ref{sec:format-ALU-DWRRS}), \verb|exucode2 = 1110|\\

\bitfields{ 1/P, 2/11, 1/0, 4/1110, 6/registerW, 2/10, 3/010, 1/0, 6/registerY, 6/registerZ }


\paragraph{Syntax} \texttt{abdsd} \emph{registerW} = \emph{registerZ}, \emph{registerY}

\paragraph{Description} The double word \emph{registerZ} is subtracted from the double word \emph{registerY}. The absolute value of the resulting double word is saturated and stored into the \emph{registerW}.


\paragraph{Execution} {Reservation table \textsf{ALU\_TINY} (Section~\ref{sec:reservation-ALU-TINY})}
~
{\begin{lstlisting}
stage RR:
  new argument3 = GPR[registerY];
  new argument2 = GPR[registerZ];
  new result1 = _SAT_64(_ABS(_SX_64(argument3) - _SX_64(argument2)));
stage E1:
  GPR[registerW] = result1;
\end{lstlisting}}
\newpage

\paragraph{Encoding} Format \textsf{ALU\_DWRRS.M} (Section~\ref{sec:format-ALU-DWRRS.M}), \verb|exucode2 = 1110|\\

\bitfields{ 1/P, 2/00, 2/-- --, 27/upper27, 1/1, 2/11, 1/0, 4/1110, 6/registerW, 2/10, 3/010, 1/0, 1/splat32, 5/lower5, 6/registerZ }


\paragraph{Syntax} \texttt{abdsd} \emph{registerW} = \emph{registerZ}, \emph{upper27\_\discretionary{}{}{}lower5}\emph{splat32}

\paragraph{Description} The double word \emph{registerZ} is subtracted from the double word \emph{upper27\_\discretionary{}{}{}lower5}. The absolute value of the resulting double word is saturated and stored into the \emph{registerW}.


\paragraph{Modifier(s)}
\noindent\begin{itemize}
\item splat32 (cf. splat32 in section \ref{par:modifier-splat32} on page \pageref{par:modifier-splat32})
\end{itemize}

\paragraph{Execution} {Reservation table \textsf{ALU\_TINY.X} (Section~\ref{sec:reservation-ALU-TINY.X})}
~
{\begin{lstlisting}
stage ID:
  new splat32 = _ZX_1(splat32);
stage RR:
  new argument3 = splat32 ? (_WX_32(upper27_lower5) << 32) | _ZX_32(_WX_32(upper27_lower5)) : _SX_32(_WX_32(upper27_lower5));
  new argument2 = GPR[registerZ];
  new result1 = _SAT_64(_ABS(_SX_64(argument3) - _SX_64(argument2)));
stage E1:
  GPR[registerW] = result1;
\end{lstlisting}}
\newpage
\subsubsection[{\small ABDSDP: Absolute Difference Saturated Double Word Pair}]{ABDSDP\hfill Absolute Difference Saturated Double Word Pair}
\label{instr:ABDSDP}\index{abdsdp}
 
\paragraph{Encoding} Format \textsf{ALU\_DPWRRS} (Section~\ref{sec:format-ALU-DPWRRS}), \verb|exucode2 = 1110|\\

\bitfields{ 1/P, 2/11, 1/1, 4/1110, 5/registerM, 1/--, 2/10, 3/010, 1/0, 5/registerO, 1/--, 5/registerP, 1/0 }


\paragraph{Syntax} \texttt{abdsdp} \emph{registerM} = \emph{registerP}, \emph{registerO}

\paragraph{Description} The double word pair \emph{registerP} is subtracted from the double word pair \emph{registerO}. The absolute value of the resultingdouble word pair is saturated and stored into the \emph{registerM}.


\paragraph{Execution} {Reservation table \textsf{ALU\_LITE} (Section~\ref{sec:reservation-ALU-LITE})}
~
{\begin{lstlisting}
stage RR:
  new argument3 = PGR[registerO];
  new argument2 = PGR[registerP];
  for (I = 0; I < 2; I += 1) {
    result1.64[I] = _SAT_64(_ABS(_SX_64(argument3.64[I]) - _SX_64(argument2.64[I])))
  };
stage E1:
  PGR[registerM] = result1;
\end{lstlisting}}
\newpage

\paragraph{Encoding} Format \textsf{ALU\_DPWRRS.M} (Section~\ref{sec:format-ALU-DPWRRS.M}), \verb|exucode2 = 1110|\\

\bitfields{ 1/P, 2/00, 2/-- --, 27/upper27, 1/1, 2/11, 1/1, 4/1110, 5/registerM, 1/--, 2/10, 3/010, 1/0, 1/splat32, 5/lower5, 5/registerP, 1/0 }


\paragraph{Syntax} \texttt{abdsdp} \emph{registerM} = \emph{registerP}, \emph{upper27\_\discretionary{}{}{}lower5}\emph{splat32}

\paragraph{Description} The double word pair \emph{registerP} is subtracted from the double word pair \emph{upper27\_\discretionary{}{}{}lower5}. The absolute value of the resultingdouble word pair is saturated and stored into the \emph{registerM}.


\paragraph{Modifier(s)}
\noindent\begin{itemize}
\item splat32 (cf. splat32 in section \ref{par:modifier-splat32} on page \pageref{par:modifier-splat32})
\end{itemize}

\paragraph{Execution} {Reservation table \textsf{ALU\_LITE.X} (Section~\ref{sec:reservation-ALU-LITE.X})}
~
{\begin{lstlisting}
stage ID:
  new splat32 = _ZX_1(splat32);
stage RR:
  new argument3 = splat32 ? (_WX_32(upper27_lower5) << 32) | _ZX_32(_WX_32(upper27_lower5)) : _SX_32(_WX_32(upper27_lower5));
  new argument2 = PGR[registerP];
  for (I = 0; I < 2; I += 1) {
    result1.64[I] = _SAT_64(_ABS(_SX_64(argument3.64[I]) - _SX_64(argument2.64[I])))
  };
stage E1:
  PGR[registerM] = result1;
\end{lstlisting}}
\newpage
\subsubsection[{\small ABDSHO: Absolute Difference Saturated Half Word Octuple}]{ABDSHO\hfill Absolute Difference Saturated Half Word Octuple}
\label{instr:ABDSHO}\index{abdsho}
 
\paragraph{Encoding} Format \textsf{ALU\_HOWRRS} (Section~\ref{sec:format-ALU-HOWRRS}), \verb|exucode2 = 1110|\\

\bitfields{ 1/P, 2/11, 1/1, 4/1110, 5/registerM, 1/--, 2/10, 3/010, 1/1, 5/registerO, 1/--, 5/registerP, 1/0 }


\paragraph{Syntax} \texttt{abdsho} \emph{registerM} = \emph{registerP}, \emph{registerO}

\paragraph{Description} The \emph{registerP} and the \emph{registerO} are interpreted as eight half words packed into 128 bits. The %e half words are subtracted from the \emph{registerO} half words. The absolute values of the eight resulting half words are saturated, packed and stored into the \emph{registerM}.


\paragraph{Execution} {Reservation table \textsf{ALU\_LITE} (Section~\ref{sec:reservation-ALU-LITE})}
~
{\begin{lstlisting}
stage RR:
  new argument3 = PGR[registerO];
  new argument2 = PGR[registerP];
  for (I = 0; I < 8; I += 1) {
    result1.16[I] = _SAT_16(_ABS(_SX_16(argument3.16[I]) - _SX_16(argument2.16[I])))
  };
stage E1:
  PGR[registerM] = result1;
\end{lstlisting}}
\newpage

\paragraph{Encoding} Format \textsf{ALU\_HOWRRS.M} (Section~\ref{sec:format-ALU-HOWRRS.M}), \verb|exucode2 = 1110|\\

\bitfields{ 1/P, 2/00, 2/-- --, 27/upper27, 1/1, 2/11, 1/1, 4/1110, 5/registerM, 1/--, 2/10, 3/010, 1/1, 1/splat32, 5/lower5, 5/registerP, 1/0 }


\paragraph{Syntax} \texttt{abdsho} \emph{registerM} = \emph{registerP}, \emph{upper27\_\discretionary{}{}{}lower5}\emph{splat32}

\paragraph{Description} The \emph{registerP} and the \emph{upper27\_\discretionary{}{}{}lower5} are interpreted as eight half words packed into 128 bits. The %e half words are subtracted from the \emph{upper27\_\discretionary{}{}{}lower5} half words. The absolute values of the eight resulting half words are saturated, packed and stored into the \emph{registerM}.


\paragraph{Modifier(s)}
\noindent\begin{itemize}
\item splat32 (cf. splat32 in section \ref{par:modifier-splat32} on page \pageref{par:modifier-splat32})
\end{itemize}

\paragraph{Execution} {Reservation table \textsf{ALU\_LITE.X} (Section~\ref{sec:reservation-ALU-LITE.X})}
~
{\begin{lstlisting}
stage ID:
  new splat32 = _ZX_1(splat32);
stage RR:
  new argument3 = splat32 ? (_WX_32(upper27_lower5) << 32) | _ZX_32(_WX_32(upper27_lower5)) : _SX_32(_WX_32(upper27_lower5));
  new argument2 = PGR[registerP];
  for (I = 0; I < 8; I += 1) {
    result1.16[I] = _SAT_16(_ABS(_SX_16(argument3.16[I]) - _SX_16(argument2.16[I])))
  };
stage E1:
  PGR[registerM] = result1;
\end{lstlisting}}
\newpage
\subsubsection[{\small ABDSW: Absolute Difference Saturated Word}]{ABDSW\hfill Absolute Difference Saturated Word}
\label{instr:ABDSW}\index{abdsw}
 
\paragraph{Encoding} Format \textsf{ALU\_WWRRS} (Section~\ref{sec:format-ALU-WWRRS}), \verb|exucode2 = 1110|\\

\bitfields{ 1/P, 2/11, 1/0, 4/1110, 6/registerW, 2/11, 3/010, 1/signextw, 6/registerY, 6/registerZ }


\paragraph{Syntax} \texttt{abdsw}\emph{signextw} \emph{registerW} = \emph{registerZ}, \emph{registerY}

\paragraph{Description} The word \emph{registerZ} is subtracted from the word \emph{registerY}. The absolute value of the resulting word is saturated and stored into the \emph{registerW}.


\paragraph{Modifier(s)}
\noindent\begin{itemize}
\item signextw (cf. signextw in section \ref{par:modifier-signextw} on page \pageref{par:modifier-signextw})
\end{itemize}

\paragraph{Execution} {Reservation table \textsf{ALU\_TINY} (Section~\ref{sec:reservation-ALU-TINY})}
~
{\begin{lstlisting}
stage ID:
  new signextw = _ZX_1(signextw);
stage RR:
  new argument3 = GPR[registerY];
  new argument2 = GPR[registerZ];
  new result1 = _SAT_32(_ABS(_SX_32(argument3) - _SX_32(argument2)));
stage E1:
  GPR[registerW] = signextw ? _SX_32(result1) : _ZX_32(result1);
\end{lstlisting}}
\newpage

\paragraph{Encoding} Format \textsf{ALU\_WWRRS.M} (Section~\ref{sec:format-ALU-WWRRS.M}), \verb|exucode2 = 1110|\\

\bitfields{ 1/P, 2/00, 2/-- --, 27/upper27, 1/1, 2/11, 1/0, 4/1110, 6/registerW, 2/11, 3/010, 1/signextw, 1/splat32, 5/lower5, 6/registerZ }


\paragraph{Syntax} \texttt{abdsw}\emph{signextw} \emph{registerW} = \emph{registerZ}, \emph{upper27\_\discretionary{}{}{}lower5}

\paragraph{Description} The word \emph{registerZ} is subtracted from the word \emph{upper27\_\discretionary{}{}{}lower5}. The absolute value of the resulting word is saturated and stored into the \emph{registerW}.


\paragraph{Modifier(s)}
\noindent\begin{itemize}
\item signextw (cf. signextw in section \ref{par:modifier-signextw} on page \pageref{par:modifier-signextw})
\end{itemize}

\paragraph{Execution} {Reservation table \textsf{ALU\_TINY.X} (Section~\ref{sec:reservation-ALU-TINY.X})}
~
{\begin{lstlisting}
stage ID:
  new signextw = _ZX_1(signextw);
stage RR:
  new argument3 = _WX_32(upper27_lower5);
  new argument2 = GPR[registerZ];
  new result1 = _SAT_32(_ABS(_SX_32(argument3) - _SX_32(argument2)));
stage E1:
  GPR[registerW] = signextw ? _SX_32(result1) : _ZX_32(result1);
\end{lstlisting}}
\newpage
\subsubsection[{\small ABDSWQ: Absolute Difference Saturated Word Quadruple}]{ABDSWQ\hfill Absolute Difference Saturated Word Quadruple}
\label{instr:ABDSWQ}\index{abdswq}
 
\paragraph{Encoding} Format \textsf{ALU\_WQWRRS} (Section~\ref{sec:format-ALU-WQWRRS}), \verb|exucode2 = 1110|\\

\bitfields{ 1/P, 2/11, 1/1, 4/1110, 5/registerM, 1/--, 2/10, 3/010, 1/0, 5/registerO, 1/--, 5/registerP, 1/1 }


\paragraph{Syntax} \texttt{abdswq} \emph{registerM} = \emph{registerP}, \emph{registerO}

\paragraph{Description} The \emph{registerP} and the \emph{registerO} are interpreted as four words packed into 128 bits. The \emph{registerP} words are subtracted from the \emph{registerO} words. The absolute values of the four resulting words are saturated, packed and stored into the \emph{registerM}.


\paragraph{Execution} {Reservation table \textsf{ALU\_LITE} (Section~\ref{sec:reservation-ALU-LITE})}
~
{\begin{lstlisting}
stage RR:
  new argument3 = PGR[registerO];
  new argument2 = PGR[registerP];
  for (I = 0; I < 4; I += 1) {
    result1.32[I] = _SAT_32(_ABS(_SX_32(argument3.32[I]) - _SX_32(argument2.32[I])))
  };
stage E1:
  PGR[registerM] = result1;
\end{lstlisting}}
\newpage

\paragraph{Encoding} Format \textsf{ALU\_WQWRRS.M} (Section~\ref{sec:format-ALU-WQWRRS.M}), \verb|exucode2 = 1110|\\

\bitfields{ 1/P, 2/00, 2/-- --, 27/upper27, 1/1, 2/11, 1/1, 4/1110, 5/registerM, 1/--, 2/10, 3/010, 1/0, 1/splat32, 5/lower5, 5/registerP, 1/1 }


\paragraph{Syntax} \texttt{abdswq} \emph{registerM} = \emph{registerP}, \emph{upper27\_\discretionary{}{}{}lower5}\emph{splat32}

\paragraph{Description} The \emph{registerP} and the \emph{upper27\_\discretionary{}{}{}lower5} are interpreted as four words packed into 128 bits. The \emph{registerP} words are subtracted from the \emph{upper27\_\discretionary{}{}{}lower5} words. The absolute values of the four resulting words are saturated, packed and stored into the \emph{registerM}.


\paragraph{Modifier(s)}
\noindent\begin{itemize}
\item splat32 (cf. splat32 in section \ref{par:modifier-splat32} on page \pageref{par:modifier-splat32})
\end{itemize}

\paragraph{Execution} {Reservation table \textsf{ALU\_LITE.X} (Section~\ref{sec:reservation-ALU-LITE.X})}
~
{\begin{lstlisting}
stage ID:
  new splat32 = _ZX_1(splat32);
stage RR:
  new argument3 = splat32 ? (_WX_32(upper27_lower5) << 32) | _ZX_32(_WX_32(upper27_lower5)) : _SX_32(_WX_32(upper27_lower5));
  new argument2 = PGR[registerP];
  for (I = 0; I < 4; I += 1) {
    result1.32[I] = _SAT_32(_ABS(_SX_32(argument3.32[I]) - _SX_32(argument2.32[I])))
  };
stage E1:
  PGR[registerM] = result1;
\end{lstlisting}}
\newpage
\subsubsection[{\small ABDUBX: Absolute Difference Unsigned Byte Hexadecuple}]{ABDUBX\hfill Absolute Difference Unsigned Byte Hexadecuple}
\label{instr:ABDUBX}\index{abdubx}
 
\paragraph{Encoding} Format \textsf{ALU\_BXWRRS} (Section~\ref{sec:format-ALU-BXWRRS}), \verb|exucode2 = 1101|\\

\bitfields{ 1/P, 2/11, 1/1, 4/1101, 5/registerM, 1/--, 2/10, 3/010, 1/1, 5/registerO, 1/--, 5/registerP, 1/1 }


\paragraph{Syntax} \texttt{abdubx} \emph{registerM} = \emph{registerP}, \emph{registerO}

\paragraph{Description} The \emph{registerP} and the \emph{registerO} are interpreted as sixteen unsigned bytes packed into 128 bits. The \emph{registerP} bytes are subtracted from the \emph{registerO} bytes. The absolute value of the sixteen resulting bytes are packed and stored into the \emph{registerM}.


\paragraph{Execution} {Reservation table \textsf{ALU\_LITE} (Section~\ref{sec:reservation-ALU-LITE})}
~
{\begin{lstlisting}
stage RR:
  new argument3 = PGR[registerO];
  new argument2 = PGR[registerP];
  for (I = 0; I < 16; I += 1) {
    result1.8[I] = _ABS(_ZX_8(argument3.8[I]) - _ZX_8(argument2.8[I]))
  };
stage E1:
  PGR[registerM] = result1;
\end{lstlisting}}
\newpage

\paragraph{Encoding} Format \textsf{ALU\_BXWRRS.M} (Section~\ref{sec:format-ALU-BXWRRS.M}), \verb|exucode2 = 1101|\\

\bitfields{ 1/P, 2/00, 2/-- --, 27/upper27, 1/1, 2/11, 1/1, 4/1101, 5/registerM, 1/--, 2/10, 3/010, 1/1, 1/splat32, 5/lower5, 5/registerP, 1/1 }


\paragraph{Syntax} \texttt{abdubx} \emph{registerM} = \emph{registerP}, \emph{upper27\_\discretionary{}{}{}lower5}\emph{splat32}

\paragraph{Description} The \emph{registerP} and the \emph{upper27\_\discretionary{}{}{}lower5} are interpreted as sixteen unsigned bytes packed into 128 bits. The \emph{registerP} bytes are subtracted from the \emph{upper27\_\discretionary{}{}{}lower5} bytes. The absolute value of the sixteen resulting bytes are packed and stored into the \emph{registerM}.


\paragraph{Modifier(s)}
\noindent\begin{itemize}
\item splat32 (cf. splat32 in section \ref{par:modifier-splat32} on page \pageref{par:modifier-splat32})
\end{itemize}

\paragraph{Execution} {Reservation table \textsf{ALU\_LITE.X} (Section~\ref{sec:reservation-ALU-LITE.X})}
~
{\begin{lstlisting}
stage ID:
  new splat32 = _ZX_1(splat32);
stage RR:
  new argument3 = splat32 ? (_WX_32(upper27_lower5) << 32) | _ZX_32(_WX_32(upper27_lower5)) : _SX_32(_WX_32(upper27_lower5));
  new argument2 = PGR[registerP];
  for (I = 0; I < 16; I += 1) {
    result1.8[I] = _ABS(_ZX_8(argument3.8[I]) - _ZX_8(argument2.8[I]))
  };
stage E1:
  PGR[registerM] = result1;
\end{lstlisting}}
\newpage
\subsubsection[{\small ABDUD: Absolute Difference Unsigned Double Word}]{ABDUD\hfill Absolute Difference Unsigned Double Word}
\label{instr:ABDUD}\index{abdud}
 
\paragraph{Encoding} Format \textsf{ALU\_DWRRS} (Section~\ref{sec:format-ALU-DWRRS}), \verb|exucode2 = 1101|\\

\bitfields{ 1/P, 2/11, 1/0, 4/1101, 6/registerW, 2/10, 3/010, 1/0, 6/registerY, 6/registerZ }


\paragraph{Syntax} \texttt{abdud} \emph{registerW} = \emph{registerZ}, \emph{registerY}

\paragraph{Description} The \emph{registerZ} is subtracted from the \emph{registerY}, assuming unsigned integers. The absolute value of the result is stored into the \emph{registerW}.


\paragraph{Execution} {Reservation table \textsf{ALU\_TINY} (Section~\ref{sec:reservation-ALU-TINY})}
~
{\begin{lstlisting}
stage RR:
  new argument3 = GPR[registerY];
  new argument2 = GPR[registerZ];
  new result1 = _ABS(_ZX_64(argument3) - _ZX_64(argument2));
stage E1:
  GPR[registerW] = result1;
\end{lstlisting}}
\newpage

\paragraph{Encoding} Format \textsf{ALU\_DWRRS.M} (Section~\ref{sec:format-ALU-DWRRS.M}), \verb|exucode2 = 1101|\\

\bitfields{ 1/P, 2/00, 2/-- --, 27/upper27, 1/1, 2/11, 1/0, 4/1101, 6/registerW, 2/10, 3/010, 1/0, 1/splat32, 5/lower5, 6/registerZ }


\paragraph{Syntax} \texttt{abdud} \emph{registerW} = \emph{registerZ}, \emph{upper27\_\discretionary{}{}{}lower5}\emph{splat32}

\paragraph{Description} The \emph{registerZ} is subtracted from the \emph{upper27\_\discretionary{}{}{}lower5}, assuming unsigned integers. The absolute value of the result is stored into the \emph{registerW}.


\paragraph{Modifier(s)}
\noindent\begin{itemize}
\item splat32 (cf. splat32 in section \ref{par:modifier-splat32} on page \pageref{par:modifier-splat32})
\end{itemize}

\paragraph{Execution} {Reservation table \textsf{ALU\_TINY.X} (Section~\ref{sec:reservation-ALU-TINY.X})}
~
{\begin{lstlisting}
stage ID:
  new splat32 = _ZX_1(splat32);
stage RR:
  new argument3 = splat32 ? (_WX_32(upper27_lower5) << 32) | _ZX_32(_WX_32(upper27_lower5)) : _SX_32(_WX_32(upper27_lower5));
  new argument2 = GPR[registerZ];
  new result1 = _ABS(_ZX_64(argument3) - _ZX_64(argument2));
stage E1:
  GPR[registerW] = result1;
\end{lstlisting}}
\newpage
\subsubsection[{\small ABDUDP: Absolute Difference Unsigned Double Word Pair}]{ABDUDP\hfill Absolute Difference Unsigned Double Word Pair}
\label{instr:ABDUDP}\index{abdudp}
 
\paragraph{Encoding} Format \textsf{ALU\_DPWRRS} (Section~\ref{sec:format-ALU-DPWRRS}), \verb|exucode2 = 1101|\\

\bitfields{ 1/P, 2/11, 1/1, 4/1101, 5/registerM, 1/--, 2/10, 3/010, 1/0, 5/registerO, 1/--, 5/registerP, 1/0 }


\paragraph{Syntax} \texttt{abdudp} \emph{registerM} = \emph{registerP}, \emph{registerO}

\paragraph{Description} The \emph{registerP} and the \emph{registerO} are interpreted as two unsigned double words packed into 64 bits. The \emph{registerP} double words are subtracted from the \emph{registerO} double words. The absolute value of the two resulting double words are packed and stored into the \emph{registerM}.


\paragraph{Execution} {Reservation table \textsf{ALU\_LITE} (Section~\ref{sec:reservation-ALU-LITE})}
~
{\begin{lstlisting}
stage RR:
  new argument3 = PGR[registerO];
  new argument2 = PGR[registerP];
  for (I = 0; I < 2; I += 1) {
    result1.64[I] = _ABS(_ZX_64(argument3.64[I]) - _ZX_64(argument2.64[I]))
  };
stage E1:
  PGR[registerM] = result1;
\end{lstlisting}}
\newpage

\paragraph{Encoding} Format \textsf{ALU\_DPWRRS.M} (Section~\ref{sec:format-ALU-DPWRRS.M}), \verb|exucode2 = 1101|\\

\bitfields{ 1/P, 2/00, 2/-- --, 27/upper27, 1/1, 2/11, 1/1, 4/1101, 5/registerM, 1/--, 2/10, 3/010, 1/0, 1/splat32, 5/lower5, 5/registerP, 1/0 }


\paragraph{Syntax} \texttt{abdudp} \emph{registerM} = \emph{registerP}, \emph{upper27\_\discretionary{}{}{}lower5}\emph{splat32}

\paragraph{Description} The \emph{registerP} and the \emph{upper27\_\discretionary{}{}{}lower5} are interpreted as two unsigned double words packed into 64 bits. The \emph{registerP} double words are subtracted from the \emph{upper27\_\discretionary{}{}{}lower5} double words. The absolute value of the two resulting double words are packed and stored into the \emph{registerM}.


\paragraph{Modifier(s)}
\noindent\begin{itemize}
\item splat32 (cf. splat32 in section \ref{par:modifier-splat32} on page \pageref{par:modifier-splat32})
\end{itemize}

\paragraph{Execution} {Reservation table \textsf{ALU\_LITE.X} (Section~\ref{sec:reservation-ALU-LITE.X})}
~
{\begin{lstlisting}
stage ID:
  new splat32 = _ZX_1(splat32);
stage RR:
  new argument3 = splat32 ? (_WX_32(upper27_lower5) << 32) | _ZX_32(_WX_32(upper27_lower5)) : _SX_32(_WX_32(upper27_lower5));
  new argument2 = PGR[registerP];
  for (I = 0; I < 2; I += 1) {
    result1.64[I] = _ABS(_ZX_64(argument3.64[I]) - _ZX_64(argument2.64[I]))
  };
stage E1:
  PGR[registerM] = result1;
\end{lstlisting}}
\newpage
\subsubsection[{\small ABDUHO: Absolute Difference Unsigned Half Word Octuple}]{ABDUHO\hfill Absolute Difference Unsigned Half Word Octuple}
\label{instr:ABDUHO}\index{abduho}
 
\paragraph{Encoding} Format \textsf{ALU\_HOWRRS} (Section~\ref{sec:format-ALU-HOWRRS}), \verb|exucode2 = 1101|\\

\bitfields{ 1/P, 2/11, 1/1, 4/1101, 5/registerM, 1/--, 2/10, 3/010, 1/1, 5/registerO, 1/--, 5/registerP, 1/0 }


\paragraph{Syntax} \texttt{abduho} \emph{registerM} = \emph{registerP}, \emph{registerO}

\paragraph{Description} The \emph{registerP} and the \emph{registerO} are interpreted as eight unsigned half words packed into 128 bits. The \emph{registerP} half words are subtracted from the \emph{registerO} half words. The absolute value of the eight resulting half words are packed and stored into the \emph{registerM}.


\paragraph{Execution} {Reservation table \textsf{ALU\_LITE} (Section~\ref{sec:reservation-ALU-LITE})}
~
{\begin{lstlisting}
stage RR:
  new argument3 = PGR[registerO];
  new argument2 = PGR[registerP];
  for (I = 0; I < 8; I += 1) {
    result1.16[I] = _ABS(_ZX_16(argument3.16[I]) - _ZX_16(argument2.16[I]))
  };
stage E1:
  PGR[registerM] = result1;
\end{lstlisting}}
\newpage

\paragraph{Encoding} Format \textsf{ALU\_HOWRRS.M} (Section~\ref{sec:format-ALU-HOWRRS.M}), \verb|exucode2 = 1101|\\

\bitfields{ 1/P, 2/00, 2/-- --, 27/upper27, 1/1, 2/11, 1/1, 4/1101, 5/registerM, 1/--, 2/10, 3/010, 1/1, 1/splat32, 5/lower5, 5/registerP, 1/0 }


\paragraph{Syntax} \texttt{abduho} \emph{registerM} = \emph{registerP}, \emph{upper27\_\discretionary{}{}{}lower5}\emph{splat32}

\paragraph{Description} The \emph{registerP} and the \emph{upper27\_\discretionary{}{}{}lower5} are interpreted as eight unsigned half words packed into 128 bits. The \emph{registerP} half words are subtracted from the \emph{upper27\_\discretionary{}{}{}lower5} half words. The absolute value of the eight resulting half words are packed and stored into the \emph{registerM}.


\paragraph{Modifier(s)}
\noindent\begin{itemize}
\item splat32 (cf. splat32 in section \ref{par:modifier-splat32} on page \pageref{par:modifier-splat32})
\end{itemize}

\paragraph{Execution} {Reservation table \textsf{ALU\_LITE.X} (Section~\ref{sec:reservation-ALU-LITE.X})}
~
{\begin{lstlisting}
stage ID:
  new splat32 = _ZX_1(splat32);
stage RR:
  new argument3 = splat32 ? (_WX_32(upper27_lower5) << 32) | _ZX_32(_WX_32(upper27_lower5)) : _SX_32(_WX_32(upper27_lower5));
  new argument2 = PGR[registerP];
  for (I = 0; I < 8; I += 1) {
    result1.16[I] = _ABS(_ZX_16(argument3.16[I]) - _ZX_16(argument2.16[I]))
  };
stage E1:
  PGR[registerM] = result1;
\end{lstlisting}}
\newpage
\subsubsection[{\small ABDUW: Absolute Difference Unsigned Word}]{ABDUW\hfill Absolute Difference Unsigned Word}
\label{instr:ABDUW}\index{abduw}
 
\paragraph{Encoding} Format \textsf{ALU\_WWRRS} (Section~\ref{sec:format-ALU-WWRRS}), \verb|exucode2 = 1101|\\

\bitfields{ 1/P, 2/11, 1/0, 4/1101, 6/registerW, 2/11, 3/010, 1/signextw, 6/registerY, 6/registerZ }


\paragraph{Syntax} \texttt{abduw}\emph{signextw} \emph{registerW} = \emph{registerZ}, \emph{registerY}

\paragraph{Description} The \emph{registerZ} is subtracted from the \emph{registerY}, assuming unsigned integers. The absolute value of the result with the 32 upper bits cleared is stored into the \emph{registerW}.


\paragraph{Modifier(s)}
\noindent\begin{itemize}
\item signextw (cf. signextw in section \ref{par:modifier-signextw} on page \pageref{par:modifier-signextw})
\end{itemize}

\paragraph{Execution} {Reservation table \textsf{ALU\_TINY} (Section~\ref{sec:reservation-ALU-TINY})}
~
{\begin{lstlisting}
stage ID:
  new signextw = _ZX_1(signextw);
stage RR:
  new argument3 = GPR[registerY];
  new argument2 = GPR[registerZ];
  new result1 = _ABS(_ZX_32(argument3) - _ZX_32(argument2));
stage E1:
  GPR[registerW] = signextw ? _SX_32(result1) : _ZX_32(result1);
\end{lstlisting}}
\newpage

\paragraph{Encoding} Format \textsf{ALU\_WWRRS.M} (Section~\ref{sec:format-ALU-WWRRS.M}), \verb|exucode2 = 1101|\\

\bitfields{ 1/P, 2/00, 2/-- --, 27/upper27, 1/1, 2/11, 1/0, 4/1101, 6/registerW, 2/11, 3/010, 1/signextw, 1/splat32, 5/lower5, 6/registerZ }


\paragraph{Syntax} \texttt{abduw}\emph{signextw} \emph{registerW} = \emph{registerZ}, \emph{upper27\_\discretionary{}{}{}lower5}

\paragraph{Description} The \emph{registerZ} is subtracted from the \emph{upper27\_\discretionary{}{}{}lower5}, assuming unsigned integers. The absolute value of the result with the 32 upper bits cleared is stored into the \emph{registerW}.


\paragraph{Modifier(s)}
\noindent\begin{itemize}
\item signextw (cf. signextw in section \ref{par:modifier-signextw} on page \pageref{par:modifier-signextw})
\end{itemize}

\paragraph{Execution} {Reservation table \textsf{ALU\_TINY.X} (Section~\ref{sec:reservation-ALU-TINY.X})}
~
{\begin{lstlisting}
stage ID:
  new signextw = _ZX_1(signextw);
stage RR:
  new argument3 = _WX_32(upper27_lower5);
  new argument2 = GPR[registerZ];
  new result1 = _ABS(_ZX_32(argument3) - _ZX_32(argument2));
stage E1:
  GPR[registerW] = signextw ? _SX_32(result1) : _ZX_32(result1);
\end{lstlisting}}
\newpage
\subsubsection[{\small ABDUWQ: Absolute Difference Unsigned Word Quadruple}]{ABDUWQ\hfill Absolute Difference Unsigned Word Quadruple}
\label{instr:ABDUWQ}\index{abduwq}
 
\paragraph{Encoding} Format \textsf{ALU\_WQWRRS} (Section~\ref{sec:format-ALU-WQWRRS}), \verb|exucode2 = 1101|\\

\bitfields{ 1/P, 2/11, 1/1, 4/1101, 5/registerM, 1/--, 2/10, 3/010, 1/0, 5/registerO, 1/--, 5/registerP, 1/1 }


\paragraph{Syntax} \texttt{abduwq} \emph{registerM} = \emph{registerP}, \emph{registerO}

\paragraph{Description} The \emph{registerP} and the \emph{registerO} are interpreted as four unsigned words packed into 128 bits. The \emph{registerP} words are subtracted from the \emph{registerO} words. The absolute value of the four resulting words are packed and stored into the \emph{registerM}.


\paragraph{Execution} {Reservation table \textsf{ALU\_LITE} (Section~\ref{sec:reservation-ALU-LITE})}
~
{\begin{lstlisting}
stage RR:
  new argument3 = PGR[registerO];
  new argument2 = PGR[registerP];
  for (I = 0; I < 4; I += 1) {
    result1.32[I] = _ABS(_ZX_32(argument3.32[I]) - _ZX_32(argument2.32[I]))
  };
stage E1:
  PGR[registerM] = result1;
\end{lstlisting}}
\newpage

\paragraph{Encoding} Format \textsf{ALU\_WQWRRS.M} (Section~\ref{sec:format-ALU-WQWRRS.M}), \verb|exucode2 = 1101|\\

\bitfields{ 1/P, 2/00, 2/-- --, 27/upper27, 1/1, 2/11, 1/1, 4/1101, 5/registerM, 1/--, 2/10, 3/010, 1/0, 1/splat32, 5/lower5, 5/registerP, 1/1 }


\paragraph{Syntax} \texttt{abduwq} \emph{registerM} = \emph{registerP}, \emph{upper27\_\discretionary{}{}{}lower5}\emph{splat32}

\paragraph{Description} The \emph{registerP} and the \emph{upper27\_\discretionary{}{}{}lower5} are interpreted as four unsigned words packed into 128 bits. The \emph{registerP} words are subtracted from the \emph{upper27\_\discretionary{}{}{}lower5} words. The absolute value of the four resulting words are packed and stored into the \emph{registerM}.


\paragraph{Modifier(s)}
\noindent\begin{itemize}
\item splat32 (cf. splat32 in section \ref{par:modifier-splat32} on page \pageref{par:modifier-splat32})
\end{itemize}

\paragraph{Execution} {Reservation table \textsf{ALU\_LITE.X} (Section~\ref{sec:reservation-ALU-LITE.X})}
~
{\begin{lstlisting}
stage ID:
  new splat32 = _ZX_1(splat32);
stage RR:
  new argument3 = splat32 ? (_WX_32(upper27_lower5) << 32) | _ZX_32(_WX_32(upper27_lower5)) : _SX_32(_WX_32(upper27_lower5));
  new argument2 = PGR[registerP];
  for (I = 0; I < 4; I += 1) {
    result1.32[I] = _ABS(_ZX_32(argument3.32[I]) - _ZX_32(argument2.32[I]))
  };
stage E1:
  PGR[registerM] = result1;
\end{lstlisting}}
\newpage
\subsubsection[{\small ABDW: Absolute Difference Word}]{ABDW\hfill Absolute Difference Word}
\label{instr:ABDW}\index{abdw}
 
\paragraph{Encoding} Format \textsf{ALU\_WWRRS} (Section~\ref{sec:format-ALU-WWRRS}), \verb|exucode2 = 1100|\\

\bitfields{ 1/P, 2/11, 1/0, 4/1100, 6/registerW, 2/11, 3/010, 1/signextw, 6/registerY, 6/registerZ }


\paragraph{Syntax} \texttt{abdw}\emph{signextw} \emph{registerW} = \emph{registerZ}, \emph{registerY}

\paragraph{Description} The \emph{registerZ} is subtracted from the \emph{registerY}. The absolute value of the result with the 32 upper bits cleared is stored into the \emph{registerW}.


\paragraph{Modifier(s)}
\noindent\begin{itemize}
\item signextw (cf. signextw in section \ref{par:modifier-signextw} on page \pageref{par:modifier-signextw})
\end{itemize}

\paragraph{Execution} {Reservation table \textsf{ALU\_TINY} (Section~\ref{sec:reservation-ALU-TINY})}
~
{\begin{lstlisting}
stage ID:
  new signextw = _ZX_1(signextw);
stage RR:
  new argument3 = GPR[registerY];
  new argument2 = GPR[registerZ];
  new result1 = _ABS(_SX_32(argument3) - _SX_32(argument2));
stage E1:
  GPR[registerW] = signextw ? _SX_32(result1) : _ZX_32(result1);
\end{lstlisting}}
\newpage

\paragraph{Encoding} Format \textsf{ALU\_WWRRS.M} (Section~\ref{sec:format-ALU-WWRRS.M}), \verb|exucode2 = 1100|\\

\bitfields{ 1/P, 2/00, 2/-- --, 27/upper27, 1/1, 2/11, 1/0, 4/1100, 6/registerW, 2/11, 3/010, 1/signextw, 1/splat32, 5/lower5, 6/registerZ }


\paragraph{Syntax} \texttt{abdw}\emph{signextw} \emph{registerW} = \emph{registerZ}, \emph{upper27\_\discretionary{}{}{}lower5}

\paragraph{Description} The \emph{registerZ} is subtracted from the \emph{upper27\_\discretionary{}{}{}lower5}. The absolute value of the result with the 32 upper bits cleared is stored into the \emph{registerW}.


\paragraph{Modifier(s)}
\noindent\begin{itemize}
\item signextw (cf. signextw in section \ref{par:modifier-signextw} on page \pageref{par:modifier-signextw})
\end{itemize}

\paragraph{Execution} {Reservation table \textsf{ALU\_TINY.X} (Section~\ref{sec:reservation-ALU-TINY.X})}
~
{\begin{lstlisting}
stage ID:
  new signextw = _ZX_1(signextw);
stage RR:
  new argument3 = _WX_32(upper27_lower5);
  new argument2 = GPR[registerZ];
  new result1 = _ABS(_SX_32(argument3) - _SX_32(argument2));
stage E1:
  GPR[registerW] = signextw ? _SX_32(result1) : _ZX_32(result1);
\end{lstlisting}}
\newpage
\subsubsection[{\small ABDWQ: Absolute Difference Word Quadruple}]{ABDWQ\hfill Absolute Difference Word Quadruple}
\label{instr:ABDWQ}\index{abdwq}
 
\paragraph{Encoding} Format \textsf{ALU\_WQWRRS} (Section~\ref{sec:format-ALU-WQWRRS}), \verb|exucode2 = 1100|\\

\bitfields{ 1/P, 2/11, 1/1, 4/1100, 5/registerM, 1/--, 2/10, 3/010, 1/0, 5/registerO, 1/--, 5/registerP, 1/1 }


\paragraph{Syntax} \texttt{abdwq} \emph{registerM} = \emph{registerP}, \emph{registerO}

\paragraph{Description} The \emph{registerP} and the \emph{registerO} are interpreted as four words packed into 128 bits. The \emph{registerP} words are subtracted from the \emph{registerO} words. The absolute value of the four resulting words are packed and stored into the \emph{registerM}.


\paragraph{Execution} {Reservation table \textsf{ALU\_LITE} (Section~\ref{sec:reservation-ALU-LITE})}
~
{\begin{lstlisting}
stage RR:
  new argument3 = PGR[registerO];
  new argument2 = PGR[registerP];
  for (I = 0; I < 4; I += 1) {
    result1.32[I] = _ABS(_SX_32(argument3.32[I]) - _SX_32(argument2.32[I]))
  };
stage E1:
  PGR[registerM] = result1;
\end{lstlisting}}
\newpage

\paragraph{Encoding} Format \textsf{ALU\_WQWRRS.M} (Section~\ref{sec:format-ALU-WQWRRS.M}), \verb|exucode2 = 1100|\\

\bitfields{ 1/P, 2/00, 2/-- --, 27/upper27, 1/1, 2/11, 1/1, 4/1100, 5/registerM, 1/--, 2/10, 3/010, 1/0, 1/splat32, 5/lower5, 5/registerP, 1/1 }


\paragraph{Syntax} \texttt{abdwq} \emph{registerM} = \emph{registerP}, \emph{upper27\_\discretionary{}{}{}lower5}\emph{splat32}

\paragraph{Description} The \emph{registerP} and the \emph{upper27\_\discretionary{}{}{}lower5} are interpreted as four words packed into 128 bits. The \emph{registerP} words are subtracted from the \emph{upper27\_\discretionary{}{}{}lower5} words. The absolute value of the four resulting words are packed and stored into the \emph{registerM}.


\paragraph{Modifier(s)}
\noindent\begin{itemize}
\item splat32 (cf. splat32 in section \ref{par:modifier-splat32} on page \pageref{par:modifier-splat32})
\end{itemize}

\paragraph{Execution} {Reservation table \textsf{ALU\_LITE.X} (Section~\ref{sec:reservation-ALU-LITE.X})}
~
{\begin{lstlisting}
stage ID:
  new splat32 = _ZX_1(splat32);
stage RR:
  new argument3 = splat32 ? (_WX_32(upper27_lower5) << 32) | _ZX_32(_WX_32(upper27_lower5)) : _SX_32(_WX_32(upper27_lower5));
  new argument2 = PGR[registerP];
  for (I = 0; I < 4; I += 1) {
    result1.32[I] = _ABS(_SX_32(argument3.32[I]) - _SX_32(argument2.32[I]))
  };
stage E1:
  PGR[registerM] = result1;
\end{lstlisting}}
\newpage
\subsubsection[{\small ABSBX: Absolute Value Byte Hexadecuple}]{ABSBX\hfill Absolute Value Byte Hexadecuple}
\label{instr:ABSBX}\index{absbx}
 
\paragraph{Encoding} Format \textsf{ALU\_BXBWR} (Section~\ref{sec:format-ALU-BXBWR}), \verb|alucode8 = 0001|\\

\bitfields{ 1/P, 2/11, 1/1, 4/1111, 5/registerM, 1/--, 2/10, 3/101, 1/1, 4/0001, 1/--, 1/--, 5/registerP, 1/1 }


\paragraph{Syntax} \texttt{absbx} \emph{registerM} = \emph{registerP}

\paragraph{Description} The absolute values of the \emph{registerP} interpreted as a byte hexadecuple are stored into the \emph{registerM}.


\paragraph{Execution} {Reservation table \textsf{ALU\_LITE} (Section~\ref{sec:reservation-ALU-LITE})}
~
{\begin{lstlisting}
stage RR:
  new argument2 = PGR[registerP];
  for (I = 0; I < 16; I += 1) {
    result1.8[I] = _ABS(_SX_8(argument2.8[I]))
  };
stage E1:
  PGR[registerM] = result1;
\end{lstlisting}}
\newpage
\subsubsection[{\small ABSD: Absolute Value Double Word}]{ABSD\hfill Absolute Value Double Word}
\label{instr:ABSD}\index{absd}
 
\paragraph{Encoding} Format \textsf{ALU\_DBWR} (Section~\ref{sec:format-ALU-DBWR}), \verb|alucode8 = 0001|\\

\bitfields{ 1/P, 2/11, 1/0, 4/1111, 6/registerW, 2/10, 3/101, 1/0, 4/0001, 1/0, 1/0, 6/registerZ }


\paragraph{Syntax} \texttt{absd} \emph{registerW} = \emph{registerZ}

\paragraph{Description} The absolute value of the \emph{registerZ} is stored into the \emph{registerW}.


\paragraph{Execution} {Reservation table \textsf{ALU\_TINY} (Section~\ref{sec:reservation-ALU-TINY})}
~
{\begin{lstlisting}
stage RR:
  new argument2 = GPR[registerZ];
  new result1 = _ABS(_SX_64(argument2));
stage E1:
  GPR[registerW] = result1;
\end{lstlisting}}
\newpage
\subsubsection[{\small ABSDP: Absolute Value Double Word Pair}]{ABSDP\hfill Absolute Value Double Word Pair}
\label{instr:ABSDP}\index{absdp}
 
\paragraph{Encoding} Format \textsf{ALU\_DPBWR} (Section~\ref{sec:format-ALU-DPBWR}), \verb|alucode8 = 0001|\\

\bitfields{ 1/P, 2/11, 1/1, 4/1111, 5/registerM, 1/--, 2/10, 3/101, 1/0, 4/0001, 1/--, 1/--, 5/registerP, 1/0 }


\paragraph{Syntax} \texttt{absdp} \emph{registerM} = \emph{registerP}

\paragraph{Description} The absolute values of the \emph{registerP} interpreted as a double word pair are stored into the \emph{registerM}.


\paragraph{Execution} {Reservation table \textsf{ALU\_LITE} (Section~\ref{sec:reservation-ALU-LITE})}
~
{\begin{lstlisting}
stage RR:
  new argument2 = PGR[registerP];
  for (I = 0; I < 2; I += 1) {
    result1.64[I] = _ABS(_SX_64(argument2.64[I]))
  };
stage E1:
  PGR[registerM] = result1;
\end{lstlisting}}
\newpage
\subsubsection[{\small ABSHO: Absolute Value Half Word Octuple}]{ABSHO\hfill Absolute Value Half Word Octuple}
\label{instr:ABSHO}\index{absho}
 
\paragraph{Encoding} Format \textsf{ALU\_HOBWR} (Section~\ref{sec:format-ALU-HOBWR}), \verb|alucode8 = 0001|\\

\bitfields{ 1/P, 2/11, 1/1, 4/1111, 5/registerM, 1/--, 2/10, 3/101, 1/1, 4/0001, 1/--, 1/--, 5/registerP, 1/0 }


\paragraph{Syntax} \texttt{absho} \emph{registerM} = \emph{registerP}

\paragraph{Description} The absolute values of the \emph{registerP} interpreted as a half word octuple are stored into the \emph{registerM}.


\paragraph{Execution} {Reservation table \textsf{ALU\_LITE} (Section~\ref{sec:reservation-ALU-LITE})}
~
{\begin{lstlisting}
stage RR:
  new argument2 = PGR[registerP];
  for (I = 0; I < 8; I += 1) {
    result1.16[I] = _ABS(_SX_16(argument2.16[I]))
  };
stage E1:
  PGR[registerM] = result1;
\end{lstlisting}}
\newpage
\subsubsection[{\small ABSSBX: Absolute Value Saturated Byte Hexadecuple}]{ABSSBX\hfill Absolute Value Saturated Byte Hexadecuple}
\label{instr:ABSSBX}\index{abssbx}
 
\paragraph{Encoding} Format \textsf{ALU\_BXBWR} (Section~\ref{sec:format-ALU-BXBWR}), \verb|alucode8 = 0011|\\

\bitfields{ 1/P, 2/11, 1/1, 4/1111, 5/registerM, 1/--, 2/10, 3/101, 1/1, 4/0011, 1/--, 1/--, 5/registerP, 1/1 }


\paragraph{Syntax} \texttt{abssbx} \emph{registerM} = \emph{registerP}

\paragraph{Description} The saturated absolute values of the \emph{registerP} interpreted as a byte hexadecuple are stored into the \emph{registerM}.


\paragraph{Execution} {Reservation table \textsf{ALU\_LITE} (Section~\ref{sec:reservation-ALU-LITE})}
~
{\begin{lstlisting}
stage RR:
  new argument2 = PGR[registerP];
  for (I = 0; I < 16; I += 1) {
    result1.8[I] = _SAT_8(_ABS(_SX_8(argument2.8[I])))
  };
stage E1:
  PGR[registerM] = result1;
\end{lstlisting}}
\newpage
\subsubsection[{\small ABSSD: Absolute Value Saturated Double Word}]{ABSSD\hfill Absolute Value Saturated Double Word}
\label{instr:ABSSD}\index{abssd}
 
\paragraph{Encoding} Format \textsf{ALU\_DBWR} (Section~\ref{sec:format-ALU-DBWR}), \verb|alucode8 = 0011|\\

\bitfields{ 1/P, 2/11, 1/0, 4/1111, 6/registerW, 2/10, 3/101, 1/0, 4/0011, 1/0, 1/0, 6/registerZ }


\paragraph{Syntax} \texttt{abssd} \emph{registerW} = \emph{registerZ}

\paragraph{Description} The absolute value of the \emph{registerZ} is 64-bit saturated and stored into the \emph{registerW}.


\paragraph{Execution} {Reservation table \textsf{ALU\_TINY} (Section~\ref{sec:reservation-ALU-TINY})}
~
{\begin{lstlisting}
stage RR:
  new argument2 = GPR[registerZ];
  new result1 = _SAT_64(_ABS(_SX_64(argument2)));
stage E1:
  GPR[registerW] = result1;
\end{lstlisting}}
\newpage
\subsubsection[{\small ABSSDP: Absolute Value Saturated Double Word Pair}]{ABSSDP\hfill Absolute Value Saturated Double Word Pair}
\label{instr:ABSSDP}\index{abssdp}
 
\paragraph{Encoding} Format \textsf{ALU\_DPBWR} (Section~\ref{sec:format-ALU-DPBWR}), \verb|alucode8 = 0011|\\

\bitfields{ 1/P, 2/11, 1/1, 4/1111, 5/registerM, 1/--, 2/10, 3/101, 1/0, 4/0011, 1/--, 1/--, 5/registerP, 1/0 }


\paragraph{Syntax} \texttt{abssdp} \emph{registerM} = \emph{registerP}

\paragraph{Description} The saturated absolute values of the \emph{registerP} interpreted as a double word pair are stored into the \emph{registerM}.


\paragraph{Execution} {Reservation table \textsf{ALU\_LITE} (Section~\ref{sec:reservation-ALU-LITE})}
~
{\begin{lstlisting}
stage RR:
  new argument2 = PGR[registerP];
  for (I = 0; I < 2; I += 1) {
    result1.64[I] = _SAT_64(_ABS(_SX_64(argument2.64[I])))
  };
stage E1:
  PGR[registerM] = result1;
\end{lstlisting}}
\newpage
\subsubsection[{\small ABSSHO: Absolute Value Saturated Half Word Octuple}]{ABSSHO\hfill Absolute Value Saturated Half Word Octuple}
\label{instr:ABSSHO}\index{abssho}
 
\paragraph{Encoding} Format \textsf{ALU\_HOBWR} (Section~\ref{sec:format-ALU-HOBWR}), \verb|alucode8 = 0011|\\

\bitfields{ 1/P, 2/11, 1/1, 4/1111, 5/registerM, 1/--, 2/10, 3/101, 1/1, 4/0011, 1/--, 1/--, 5/registerP, 1/0 }


\paragraph{Syntax} \texttt{abssho} \emph{registerM} = \emph{registerP}

\paragraph{Description} The saturated absolute values of the \emph{registerP} interpreted as a half word octuple are stored into the \emph{registerM}.


\paragraph{Execution} {Reservation table \textsf{ALU\_LITE} (Section~\ref{sec:reservation-ALU-LITE})}
~
{\begin{lstlisting}
stage RR:
  new argument2 = PGR[registerP];
  for (I = 0; I < 8; I += 1) {
    result1.16[I] = _SAT_16(_ABS(_SX_16(argument2.16[I])))
  };
stage E1:
  PGR[registerM] = result1;
\end{lstlisting}}
\newpage
\subsubsection[{\small ABSSW: Absolute Value Saturated Word}]{ABSSW\hfill Absolute Value Saturated Word}
\label{instr:ABSSW}\index{abssw}
 
\paragraph{Encoding} Format \textsf{ALU\_BWRW} (Section~\ref{sec:format-ALU-BWRW}), \verb|alucode8 = 0011|\\

\bitfields{ 1/P, 2/11, 1/0, 4/1111, 6/registerW, 2/11, 3/101, 1/signextw, 4/0011, 1/0, 1/0, 6/registerZ }


\paragraph{Syntax} \texttt{abssw} \emph{registerW} = \emph{registerZ}

\paragraph{Description} The absolute value of the \emph{registerZ} is 32-bit saturated and stored into the \emph{registerW}.


\paragraph{Execution} {Reservation table \textsf{ALU\_TINY} (Section~\ref{sec:reservation-ALU-TINY})}
~
{\begin{lstlisting}
stage RR:
  new argument2 = GPR[registerZ];
  new result1 = _SAT_32(_ABS(_SX_32(argument2)));
stage E1:
  GPR[registerW] = _ZX_32(result1);
\end{lstlisting}}
\newpage
\subsubsection[{\small ABSSWQ: Absolute Value Saturated Word Quadruple}]{ABSSWQ\hfill Absolute Value Saturated Word Quadruple}
\label{instr:ABSSWQ}\index{absswq}
 
\paragraph{Encoding} Format \textsf{ALU\_WQBWR} (Section~\ref{sec:format-ALU-WQBWR}), \verb|alucode8 = 0011|\\

\bitfields{ 1/P, 2/11, 1/1, 4/1111, 5/registerM, 1/--, 2/10, 3/101, 1/0, 4/0011, 1/--, 1/--, 5/registerP, 1/1 }


\paragraph{Syntax} \texttt{absswq} \emph{registerM} = \emph{registerP}

\paragraph{Description} The saturated absolute values of the \emph{registerP} interpreted as a word quadruple are stored into the \emph{registerM}.


\paragraph{Execution} {Reservation table \textsf{ALU\_LITE} (Section~\ref{sec:reservation-ALU-LITE})}
~
{\begin{lstlisting}
stage RR:
  new argument2 = PGR[registerP];
  for (I = 0; I < 4; I += 1) {
    result1.32[I] = _SAT_32(_ABS(_SX_32(argument2.32[I])))
  };
stage E1:
  PGR[registerM] = result1;
\end{lstlisting}}
\newpage
\subsubsection[{\small ABSW: Absolute Value Word}]{ABSW\hfill Absolute Value Word}
\label{instr:ABSW}\index{absw}
 
\paragraph{Encoding} Format \textsf{ALU\_BWRW} (Section~\ref{sec:format-ALU-BWRW}), \verb|alucode8 = 0001|\\

\bitfields{ 1/P, 2/11, 1/0, 4/1111, 6/registerW, 2/11, 3/101, 1/signextw, 4/0001, 1/0, 1/0, 6/registerZ }


\paragraph{Syntax} \texttt{absw} \emph{registerW} = \emph{registerZ}

\paragraph{Description} The absolute value of the \emph{registerZ} is stored into the \emph{registerW}.


\paragraph{Execution} {Reservation table \textsf{ALU\_TINY} (Section~\ref{sec:reservation-ALU-TINY})}
~
{\begin{lstlisting}
stage RR:
  new argument2 = GPR[registerZ];
  new result1 = _ABS(_SX_32(argument2));
stage E1:
  GPR[registerW] = _ZX_32(result1);
\end{lstlisting}}
\newpage
\subsubsection[{\small ABSWQ: Absolute Value Word Quadruple}]{ABSWQ\hfill Absolute Value Word Quadruple}
\label{instr:ABSWQ}\index{abswq}
 
\paragraph{Encoding} Format \textsf{ALU\_WQBWR} (Section~\ref{sec:format-ALU-WQBWR}), \verb|alucode8 = 0001|\\

\bitfields{ 1/P, 2/11, 1/1, 4/1111, 5/registerM, 1/--, 2/10, 3/101, 1/0, 4/0001, 1/--, 1/--, 5/registerP, 1/1 }


\paragraph{Syntax} \texttt{abswq} \emph{registerM} = \emph{registerP}

\paragraph{Description} The absolute values of the \emph{registerP} interpreted as a word quadruple are stored into the \emph{registerM}.


\paragraph{Execution} {Reservation table \textsf{ALU\_LITE} (Section~\ref{sec:reservation-ALU-LITE})}
~
{\begin{lstlisting}
stage RR:
  new argument2 = PGR[registerP];
  for (I = 0; I < 4; I += 1) {
    result1.32[I] = _ABS(_SX_32(argument2.32[I]))
  };
stage E1:
  PGR[registerM] = result1;
\end{lstlisting}}
\newpage
\subsubsection[{\small ADDBX: Add Byte Hexadecuple}]{ADDBX\hfill Add Byte Hexadecuple}
\label{instr:ADDBX}\index{addbx}
 
\paragraph{Encoding} Format \textsf{ALU\_BXWRR0} (Section~\ref{sec:format-ALU-BXWRR0}), \verb|exucode2 = 0010|\\

\bitfields{ 1/P, 2/11, 1/1, 4/0010, 5/registerM, 1/--, 2/10, 3/000, 1/1, 5/registerO, 1/--, 5/registerP, 1/1 }


\paragraph{Syntax} \texttt{addbx} \emph{registerM} = \emph{registerP}, \emph{registerO}

\paragraph{Description} The \emph{registerO} and the \emph{registerP} are interpreted as sixteen bytes packed into 128 bits. The \emph{registerO} bytes and the \emph{registerP} bytes are added. The sixteen resulting bytes are packed and stored into the \emph{registerM}.


\paragraph{Execution} {Reservation table \textsf{ALU\_LITE} (Section~\ref{sec:reservation-ALU-LITE})}
~
{\begin{lstlisting}
stage RR:
  new argument3 = PGR[registerO];
  new argument2 = PGR[registerP];
  for (I = 0; I < 16; I += 1) {
    result1.8[I] = argument3.8[I] + argument2.8[I]
  };
stage E1:
  PGR[registerM] = result1;
\end{lstlisting}}
\newpage

\paragraph{Encoding} Format \textsf{ALU\_BXWRR0.M} (Section~\ref{sec:format-ALU-BXWRR0.M}), \verb|exucode2 = 0010|\\

\bitfields{ 1/P, 2/00, 2/-- --, 27/upper27, 1/1, 2/11, 1/1, 4/0010, 5/registerM, 1/--, 2/10, 3/000, 1/1, 1/splat32, 5/lower5, 5/registerP, 1/1 }


\paragraph{Syntax} \texttt{addbx} \emph{registerM} = \emph{registerP}, \emph{upper27\_\discretionary{}{}{}lower5}\emph{splat32}

\paragraph{Description} The \emph{upper27\_\discretionary{}{}{}lower5} and the \emph{registerP} are interpreted as sixteen bytes packed into 128 bits. The \emph{upper27\_\discretionary{}{}{}lower5} bytes and the \emph{registerP} bytes are added. The sixteen resulting bytes are packed and stored into the \emph{registerM}.


\paragraph{Modifier(s)}
\noindent\begin{itemize}
\item splat32 (cf. splat32 in section \ref{par:modifier-splat32} on page \pageref{par:modifier-splat32})
\end{itemize}

\paragraph{Execution} {Reservation table \textsf{ALU\_LITE.X} (Section~\ref{sec:reservation-ALU-LITE.X})}
~
{\begin{lstlisting}
stage ID:
  new splat32 = _ZX_1(splat32);
stage RR:
  new argument3 = splat32 ? (_WX_32(upper27_lower5) << 32) | _ZX_32(_WX_32(upper27_lower5)) : _SX_32(_WX_32(upper27_lower5));
  new argument2 = PGR[registerP];
  for (I = 0; I < 16; I += 1) {
    result1.8[I] = argument3.8[I] + argument2.8[I]
  };
stage E1:
  PGR[registerM] = result1;
\end{lstlisting}}
\newpage
\subsubsection[{\small ADDD: Add Double Word to Double Word}]{ADDD\hfill Add Double Word to Double Word}
\label{instr:ADDD}\index{addd}
 
\paragraph{Encoding} Format \textsf{ALU\_DWRR0} (Section~\ref{sec:format-ALU-DWRR0}), \verb|exucode2 = 0010|\\

\bitfields{ 1/P, 2/11, 1/0, 4/0010, 6/registerW, 2/10, 3/000, 1/0, 6/registerY, 6/registerZ }


\paragraph{Syntax} \texttt{addd} \emph{registerW} = \emph{registerZ}, \emph{registerY}

\paragraph{Description} The \emph{registerY} and the \emph{registerZ} are added. The result is stored into the \emph{registerW}.


\paragraph{Execution} {Reservation table \textsf{ALU\_TINY} (Section~\ref{sec:reservation-ALU-TINY})}
~
{\begin{lstlisting}
stage RR:
  new argument3 = GPR[registerY];
  new argument2 = GPR[registerZ];
  new result1 = argument3 + argument2;
stage E1:
  GPR[registerW] = result1;
\end{lstlisting}}
\newpage

\paragraph{Encoding} Format \textsf{ALU\_DWRR0.M} (Section~\ref{sec:format-ALU-DWRR0.M}), \verb|exucode2 = 0010|\\

\bitfields{ 1/P, 2/00, 2/-- --, 27/upper27, 1/1, 2/11, 1/0, 4/0010, 6/registerW, 2/10, 3/000, 1/0, 1/splat32, 5/lower5, 6/registerZ }


\paragraph{Syntax} \texttt{addd} \emph{registerW} = \emph{registerZ}, \emph{upper27\_\discretionary{}{}{}lower5}\emph{splat32}

\paragraph{Description} The \emph{upper27\_\discretionary{}{}{}lower5} and the \emph{registerZ} are added. The result is stored into the \emph{registerW}.


\paragraph{Modifier(s)}
\noindent\begin{itemize}
\item splat32 (cf. splat32 in section \ref{par:modifier-splat32} on page \pageref{par:modifier-splat32})
\end{itemize}

\paragraph{Execution} {Reservation table \textsf{ALU\_TINY.X} (Section~\ref{sec:reservation-ALU-TINY.X})}
~
{\begin{lstlisting}
stage ID:
  new splat32 = _ZX_1(splat32);
stage RR:
  new argument3 = splat32 ? (_WX_32(upper27_lower5) << 32) | _ZX_32(_WX_32(upper27_lower5)) : _SX_32(_WX_32(upper27_lower5));
  new argument2 = GPR[registerZ];
  new result1 = argument3 + argument2;
stage E1:
  GPR[registerW] = result1;
\end{lstlisting}}
\newpage

\paragraph{Encoding} Format \textsf{ALU\_DWRI} (Section~\ref{sec:format-ALU-DWRI}), \verb|exucode2 = 0010|\\

\bitfields{ 1/P, 2/11, 1/0, 4/0010, 6/registerW, 2/00, 10/signed10, 6/registerZ }


\paragraph{Syntax} \texttt{addd} \emph{registerW} = \emph{registerZ}, \emph{signed10}

\paragraph{Description} The \emph{signed10} and the \emph{registerZ} are added. The result is stored into the \emph{registerW}.


\paragraph{Execution} {Reservation table \textsf{ALU\_TINY} (Section~\ref{sec:reservation-ALU-TINY})}
~
{\begin{lstlisting}
stage RR:
  new argument3 = _SX_10(signed10);
  new argument2 = GPR[registerZ];
  new result1 = argument3 + argument2;
stage E1:
  GPR[registerW] = result1;
\end{lstlisting}}
\newpage

\paragraph{Encoding} Format \textsf{ALU\_DWRI.X} (Section~\ref{sec:format-ALU-DWRI.X}), \verb|exucode2 = 0010|\\

\bitfields{ 1/P, 2/00, 2/-- --, 27/upper27, 1/1, 2/11, 1/0, 4/0010, 6/registerW, 2/00, 10/lower10, 6/registerZ }


\paragraph{Syntax} \texttt{addd} \emph{registerW} = \emph{registerZ}, \emph{upper27\_\discretionary{}{}{}lower10}

\paragraph{Description} The \emph{upper27\_\discretionary{}{}{}lower10} and the \emph{registerZ} are added. The result is stored into the \emph{registerW}.


\paragraph{Execution} {Reservation table \textsf{ALU\_TINY.X} (Section~\ref{sec:reservation-ALU-TINY.X})}
~
{\begin{lstlisting}
stage RR:
  new argument3 = _SX_37(upper27_lower10);
  new argument2 = GPR[registerZ];
  new result1 = argument3 + argument2;
stage E1:
  GPR[registerW] = result1;
\end{lstlisting}}
\newpage

\paragraph{Encoding} Format \textsf{ALU\_DWRI.Y} (Section~\ref{sec:format-ALU-DWRI.Y}), \verb|exucode2 = 0010|\\

\bitfields{ 1/P, 2/00, 2/-- --, 27/extend27, 1/1, 2/00, 2/-- --, 27/upper27, 1/1, 2/11, 1/0, 4/0010, 6/registerW, 2/00, 10/lower10, 6/registerZ }


\paragraph{Syntax} \texttt{addd} \emph{registerW} = \emph{registerZ}, \emph{extend27\_\discretionary{}{}{}upper27\_\discretionary{}{}{}lower10}

\paragraph{Description} The \emph{extend27\_\discretionary{}{}{}upper27\_\discretionary{}{}{}lower10} and the \emph{registerZ} are added. The result is stored into the \emph{registerW}.


\paragraph{Execution} {Reservation table \textsf{ALU\_TINY.Y} (Section~\ref{sec:reservation-ALU-TINY.Y})}
~
{\begin{lstlisting}
stage RR:
  new argument3 = _WX_64(extend27_upper27_lower10);
  new argument2 = GPR[registerZ];
  new result1 = argument3 + argument2;
stage E1:
  GPR[registerW] = result1;
\end{lstlisting}}
\newpage
\subsubsection[{\small ADDDP: Add Double Word Pair}]{ADDDP\hfill Add Double Word Pair}
\label{instr:ADDDP}\index{adddp}
 
\paragraph{Encoding} Format \textsf{ALU\_DPWRR0} (Section~\ref{sec:format-ALU-DPWRR0}), \verb|exucode2 = 0010|\\

\bitfields{ 1/P, 2/11, 1/1, 4/0010, 5/registerM, 1/--, 2/10, 3/000, 1/0, 5/registerO, 1/--, 5/registerP, 1/0 }


\paragraph{Syntax} \texttt{adddp} \emph{registerM} = \emph{registerP}, \emph{registerO}

\paragraph{Description} The \emph{registerO} and the \emph{registerP} are interpreted as two double words packed into 128 bits. The \emph{registerO} double words and the \emph{registerP} double words are added. The two resulting double words are packed and stored into the \emph{registerM}.


\paragraph{Execution} {Reservation table \textsf{ALU\_LITE} (Section~\ref{sec:reservation-ALU-LITE})}
~
{\begin{lstlisting}
stage RR:
  new argument3 = PGR[registerO];
  new argument2 = PGR[registerP];
  for (I = 0; I < 2; I += 1) {
    result1.32[I] = argument3.32[I] + argument2.32[I]
  };
stage E1:
  PGR[registerM] = result1;
\end{lstlisting}}
\newpage

\paragraph{Encoding} Format \textsf{ALU\_DPWRR0.M} (Section~\ref{sec:format-ALU-DPWRR0.M}), \verb|exucode2 = 0010|\\

\bitfields{ 1/P, 2/00, 2/-- --, 27/upper27, 1/1, 2/11, 1/1, 4/0010, 5/registerM, 1/--, 2/10, 3/000, 1/0, 1/splat32, 5/lower5, 5/registerP, 1/0 }


\paragraph{Syntax} \texttt{adddp} \emph{registerM} = \emph{registerP}, \emph{upper27\_\discretionary{}{}{}lower5}\emph{splat32}

\paragraph{Description} The \emph{upper27\_\discretionary{}{}{}lower5} and the \emph{registerP} are interpreted as two double words packed into 128 bits. The \emph{upper27\_\discretionary{}{}{}lower5} double words and the \emph{registerP} double words are added. The two resulting double words are packed and stored into the \emph{registerM}.


\paragraph{Modifier(s)}
\noindent\begin{itemize}
\item splat32 (cf. splat32 in section \ref{par:modifier-splat32} on page \pageref{par:modifier-splat32})
\end{itemize}

\paragraph{Execution} {Reservation table \textsf{ALU\_LITE.X} (Section~\ref{sec:reservation-ALU-LITE.X})}
~
{\begin{lstlisting}
stage ID:
  new splat32 = _ZX_1(splat32);
stage RR:
  new argument3 = splat32 ? (_WX_32(upper27_lower5) << 32) | _ZX_32(_WX_32(upper27_lower5)) : _SX_32(_WX_32(upper27_lower5));
  new argument2 = PGR[registerP];
  for (I = 0; I < 2; I += 1) {
    result1.32[I] = argument3.32[I] + argument2.32[I]
  };
stage E1:
  PGR[registerM] = result1;
\end{lstlisting}}
\newpage
\subsubsection[{\small ADDHO: Add Half Word Octuple}]{ADDHO\hfill Add Half Word Octuple}
\label{instr:ADDHO}\index{addho}
 
\paragraph{Encoding} Format \textsf{ALU\_HOWRR0} (Section~\ref{sec:format-ALU-HOWRR0}), \verb|exucode2 = 0010|\\

\bitfields{ 1/P, 2/11, 1/1, 4/0010, 5/registerM, 1/--, 2/10, 3/000, 1/1, 5/registerO, 1/--, 5/registerP, 1/0 }


\paragraph{Syntax} \texttt{addho} \emph{registerM} = \emph{registerP}, \emph{registerO}

\paragraph{Description} The \emph{registerO} and the \emph{registerP} are interpreted as eight half words packed into 128 bits. The \emph{registerO} half words and the \emph{registerP} half words are added. The eight resulting half words are packed and stored into the \emph{registerM}.


\paragraph{Execution} {Reservation table \textsf{ALU\_LITE} (Section~\ref{sec:reservation-ALU-LITE})}
~
{\begin{lstlisting}
stage RR:
  new argument3 = PGR[registerO];
  new argument2 = PGR[registerP];
  for (I = 0; I < 8; I += 1) {
    result1.16[I] = argument3.16[I] + argument2.16[I]
  };
stage E1:
  PGR[registerM] = result1;
\end{lstlisting}}
\newpage

\paragraph{Encoding} Format \textsf{ALU\_HOWRR0.M} (Section~\ref{sec:format-ALU-HOWRR0.M}), \verb|exucode2 = 0010|\\

\bitfields{ 1/P, 2/00, 2/-- --, 27/upper27, 1/1, 2/11, 1/1, 4/0010, 5/registerM, 1/--, 2/10, 3/000, 1/1, 1/splat32, 5/lower5, 5/registerP, 1/0 }


\paragraph{Syntax} \texttt{addho} \emph{registerM} = \emph{registerP}, \emph{upper27\_\discretionary{}{}{}lower5}\emph{splat32}

\paragraph{Description} The \emph{upper27\_\discretionary{}{}{}lower5} and the \emph{registerP} are interpreted as eight half words packed into 128 bits. The \emph{upper27\_\discretionary{}{}{}lower5} half words and the \emph{registerP} half words are added. The eight resulting half words are packed and stored into the \emph{registerM}.


\paragraph{Modifier(s)}
\noindent\begin{itemize}
\item splat32 (cf. splat32 in section \ref{par:modifier-splat32} on page \pageref{par:modifier-splat32})
\end{itemize}

\paragraph{Execution} {Reservation table \textsf{ALU\_LITE.X} (Section~\ref{sec:reservation-ALU-LITE.X})}
~
{\begin{lstlisting}
stage ID:
  new splat32 = _ZX_1(splat32);
stage RR:
  new argument3 = splat32 ? (_WX_32(upper27_lower5) << 32) | _ZX_32(_WX_32(upper27_lower5)) : _SX_32(_WX_32(upper27_lower5));
  new argument2 = PGR[registerP];
  for (I = 0; I < 8; I += 1) {
    result1.16[I] = argument3.16[I] + argument2.16[I]
  };
stage E1:
  PGR[registerM] = result1;
\end{lstlisting}}
\newpage
\subsubsection[{\small ADDQ: Add Quadruple Word to Quadruple Word}]{ADDQ\hfill Add Quadruple Word to Quadruple Word}
\label{instr:ADDQ}\index{addq}
 
\paragraph{Encoding} Format \textsf{ALU\_QWRR0} (Section~\ref{sec:format-ALU-QWRR0}), \verb|exucode2 = 0010|\\

\bitfields{ 1/P, 2/11, 1/0, 4/0010, 5/registerM, 1/--, 2/10, 3/000, 1/1, 5/registerO, 1/--, 5/registerP, 1/-- }


\paragraph{Syntax} \texttt{addq} \emph{registerM} = \emph{registerP}, \emph{registerO}

\paragraph{Description} The \emph{registerO} and the \emph{registerP} are added. The result is stored into the \emph{registerM}.


\paragraph{Execution} {Reservation table \textsf{ALU\_LITE} (Section~\ref{sec:reservation-ALU-LITE})}
~
{\begin{lstlisting}
stage RR:
  new argument3 = PGR[registerO];
  new argument2 = PGR[registerP];
  new result1 = argument3 + argument2;
stage E1:
  PGR[registerM] = result1;
\end{lstlisting}}
\newpage

\paragraph{Encoding} Format \textsf{ALU\_QWRR0.M} (Section~\ref{sec:format-ALU-QWRR0.M}), \verb|exucode2 = 0010|\\

\bitfields{ 1/P, 2/00, 2/-- --, 27/upper27, 1/1, 2/11, 1/0, 4/0010, 5/registerM, 1/--, 2/10, 3/000, 1/1, 1/splat32, 5/lower5, 5/registerP, 1/1 }


\paragraph{Syntax} \texttt{addq} \emph{registerM} = \emph{registerP}, \emph{upper27\_\discretionary{}{}{}lower5}\emph{splat32}

\paragraph{Description} The \emph{upper27\_\discretionary{}{}{}lower5} and the \emph{registerP} are added. The result is stored into the \emph{registerM}.


\paragraph{Modifier(s)}
\noindent\begin{itemize}
\item splat32 (cf. splat32 in section \ref{par:modifier-splat32} on page \pageref{par:modifier-splat32})
\end{itemize}

\paragraph{Execution} {Reservation table \textsf{ALU\_LITE.X} (Section~\ref{sec:reservation-ALU-LITE.X})}
~
{\begin{lstlisting}
stage ID:
  new splat32 = _ZX_1(splat32);
stage RR:
  new argument3 = splat32 ? (_WX_32(upper27_lower5) << 32) | _ZX_32(_WX_32(upper27_lower5)) : _SX_32(_WX_32(upper27_lower5));
  new argument2 = PGR[registerP];
  new result1 = argument3 + argument2;
stage E1:
  PGR[registerM] = result1;
\end{lstlisting}}
\newpage
\subsubsection[{\small ADDSBX: Add Saturated Byte Hexadecuple}]{ADDSBX\hfill Add Saturated Byte Hexadecuple}
\label{instr:ADDSBX}\index{addsbx}
 
\paragraph{Encoding} Format \textsf{ALU\_BXWRR1} (Section~\ref{sec:format-ALU-BXWRR1}), \verb|exucode2 = 1000|\\

\bitfields{ 1/P, 2/11, 1/1, 4/1000, 5/registerM, 1/--, 2/10, 3/101, 1/1, 5/registerO, 1/--, 5/registerP, 1/1 }


\paragraph{Syntax} \texttt{addsbx} \emph{registerM} = \emph{registerP}, \emph{registerO}

\paragraph{Description} The \emph{registerO} and the \emph{registerP} are interpreted as sixteen bytes packed into 128 bits. The \emph{registerO} bytes and the \emph{registerP} bytes are added with saturation to 8 bits. The sixteen resulting bytes are packed and stored into the \emph{registerM}.


\paragraph{Execution} {Reservation table \textsf{ALU\_LITE} (Section~\ref{sec:reservation-ALU-LITE})}
~
{\begin{lstlisting}
stage RR:
  new argument3 = PGR[registerO];
  new argument2 = PGR[registerP];
  for (I = 0; I < 16; I += 1) {
    result1.8[I] = _SAT_8(_SX_8(argument3.8[I]) + _SX_8(argument2.8[I]))
  };
stage E1:
  PGR[registerM] = result1;
\end{lstlisting}}
\newpage

\paragraph{Encoding} Format \textsf{ALU\_BXWRR1.M} (Section~\ref{sec:format-ALU-BXWRR1.M}), \verb|exucode2 = 1000|\\

\bitfields{ 1/P, 2/00, 2/-- --, 27/upper27, 1/1, 2/11, 1/1, 4/1000, 5/registerM, 1/--, 2/10, 3/101, 1/1, 1/splat32, 5/lower5, 5/registerP, 1/1 }


\paragraph{Syntax} \texttt{addsbx} \emph{registerM} = \emph{registerP}, \emph{upper27\_\discretionary{}{}{}lower5}\emph{splat32}

\paragraph{Description} The \emph{upper27\_\discretionary{}{}{}lower5} and the \emph{registerP} are interpreted as sixteen bytes packed into 128 bits. The \emph{upper27\_\discretionary{}{}{}lower5} bytes and the \emph{registerP} bytes are added with saturation to 8 bits. The sixteen resulting bytes are packed and stored into the \emph{registerM}.


\paragraph{Modifier(s)}
\noindent\begin{itemize}
\item splat32 (cf. splat32 in section \ref{par:modifier-splat32} on page \pageref{par:modifier-splat32})
\end{itemize}

\paragraph{Execution} {Reservation table \textsf{ALU\_LITE.X} (Section~\ref{sec:reservation-ALU-LITE.X})}
~
{\begin{lstlisting}
stage ID:
  new splat32 = _ZX_1(splat32);
stage RR:
  new argument3 = splat32 ? (_WX_32(upper27_lower5) << 32) | _ZX_32(_WX_32(upper27_lower5)) : _SX_32(_WX_32(upper27_lower5));
  new argument2 = PGR[registerP];
  for (I = 0; I < 16; I += 1) {
    result1.8[I] = _SAT_8(_SX_8(argument3.8[I]) + _SX_8(argument2.8[I]))
  };
stage E1:
  PGR[registerM] = result1;
\end{lstlisting}}
\newpage
\subsubsection[{\small ADDSD: Add Saturated Double Word}]{ADDSD\hfill Add Saturated Double Word}
\label{instr:ADDSD}\index{addsd}
 
\paragraph{Encoding} Format \textsf{ALU\_DWRR1} (Section~\ref{sec:format-ALU-DWRR1}), \verb|exucode2 = 1000|\\

\bitfields{ 1/P, 2/11, 1/0, 4/1000, 6/registerW, 2/10, 3/101, 1/0, 6/registerY, 6/registerZ }


\paragraph{Syntax} \texttt{addsd} \emph{registerW} = \emph{registerZ}, \emph{registerY}

\paragraph{Description} The \emph{registerY} and the \emph{registerZ} are added with 64-bit saturation.


\paragraph{Execution} {Reservation table \textsf{ALU\_TINY} (Section~\ref{sec:reservation-ALU-TINY})}
~
{\begin{lstlisting}
stage RR:
  new argument3 = GPR[registerY];
  new argument2 = GPR[registerZ];
  new result1 = _SAT_64(_SX_64(argument3) + _SX_64(argument2));
stage E1:
  GPR[registerW] = result1;
\end{lstlisting}}
\newpage

\paragraph{Encoding} Format \textsf{ALU\_DWRR1.M} (Section~\ref{sec:format-ALU-DWRR1.M}), \verb|exucode2 = 1000|\\

\bitfields{ 1/P, 2/00, 2/-- --, 27/upper27, 1/1, 2/11, 1/0, 4/1000, 6/registerW, 2/10, 3/101, 1/0, 1/splat32, 5/lower5, 6/registerZ }


\paragraph{Syntax} \texttt{addsd} \emph{registerW} = \emph{registerZ}, \emph{upper27\_\discretionary{}{}{}lower5}\emph{splat32}

\paragraph{Description} The \emph{upper27\_\discretionary{}{}{}lower5} and the \emph{registerZ} are added with 64-bit saturation.


\paragraph{Modifier(s)}
\noindent\begin{itemize}
\item splat32 (cf. splat32 in section \ref{par:modifier-splat32} on page \pageref{par:modifier-splat32})
\end{itemize}

\paragraph{Execution} {Reservation table \textsf{ALU\_TINY.X} (Section~\ref{sec:reservation-ALU-TINY.X})}
~
{\begin{lstlisting}
stage ID:
  new splat32 = _ZX_1(splat32);
stage RR:
  new argument3 = splat32 ? (_WX_32(upper27_lower5) << 32) | _ZX_32(_WX_32(upper27_lower5)) : _SX_32(_WX_32(upper27_lower5));
  new argument2 = GPR[registerZ];
  new result1 = _SAT_64(_SX_64(argument3) + _SX_64(argument2));
stage E1:
  GPR[registerW] = result1;
\end{lstlisting}}
\newpage
\subsubsection[{\small ADDSDP: Add Saturated Double Word Pair}]{ADDSDP\hfill Add Saturated Double Word Pair}
\label{instr:ADDSDP}\index{addsdp}
 
\paragraph{Encoding} Format \textsf{ALU\_DPWRR1} (Section~\ref{sec:format-ALU-DPWRR1}), \verb|exucode2 = 1000|\\

\bitfields{ 1/P, 2/11, 1/1, 4/1000, 5/registerM, 1/--, 2/10, 3/101, 1/0, 5/registerO, 1/--, 5/registerP, 1/0 }


\paragraph{Syntax} \texttt{addsdp} \emph{registerM} = \emph{registerP}, \emph{registerO}

\paragraph{Description} The \emph{registerO} and the \emph{registerP} are interpreted as two double words packed into 128 bits. The \emph{registerO} double words and the \emph{registerP} double words are added with saturation to 64 bits. The two resulting double words are packed and stored into the \emph{registerM}.


\paragraph{Execution} {Reservation table \textsf{ALU\_LITE} (Section~\ref{sec:reservation-ALU-LITE})}
~
{\begin{lstlisting}
stage RR:
  new argument3 = PGR[registerO];
  new argument2 = PGR[registerP];
  for (I = 0; I < 2; I += 1) {
    result1.64[I] = _SAT_64(_SX_64(argument3.64[I]) + _SX_64(argument2.64[I]))
  };
stage E1:
  PGR[registerM] = result1;
\end{lstlisting}}
\newpage

\paragraph{Encoding} Format \textsf{ALU\_DPWRR1.M} (Section~\ref{sec:format-ALU-DPWRR1.M}), \verb|exucode2 = 1000|\\

\bitfields{ 1/P, 2/00, 2/-- --, 27/upper27, 1/1, 2/11, 1/1, 4/1000, 5/registerM, 1/--, 2/10, 3/101, 1/0, 1/splat32, 5/lower5, 5/registerP, 1/0 }


\paragraph{Syntax} \texttt{addsdp} \emph{registerM} = \emph{registerP}, \emph{upper27\_\discretionary{}{}{}lower5}\emph{splat32}

\paragraph{Description} The \emph{upper27\_\discretionary{}{}{}lower5} and the \emph{registerP} are interpreted as two double words packed into 128 bits. The \emph{upper27\_\discretionary{}{}{}lower5} double words and the \emph{registerP} double words are added with saturation to 64 bits. The two resulting double words are packed and stored into the \emph{registerM}.


\paragraph{Modifier(s)}
\noindent\begin{itemize}
\item splat32 (cf. splat32 in section \ref{par:modifier-splat32} on page \pageref{par:modifier-splat32})
\end{itemize}

\paragraph{Execution} {Reservation table \textsf{ALU\_LITE.X} (Section~\ref{sec:reservation-ALU-LITE.X})}
~
{\begin{lstlisting}
stage ID:
  new splat32 = _ZX_1(splat32);
stage RR:
  new argument3 = splat32 ? (_WX_32(upper27_lower5) << 32) | _ZX_32(_WX_32(upper27_lower5)) : _SX_32(_WX_32(upper27_lower5));
  new argument2 = PGR[registerP];
  for (I = 0; I < 2; I += 1) {
    result1.64[I] = _SAT_64(_SX_64(argument3.64[I]) + _SX_64(argument2.64[I]))
  };
stage E1:
  PGR[registerM] = result1;
\end{lstlisting}}
\newpage
\subsubsection[{\small ADDSHO: Add Saturated Half Word Octuple}]{ADDSHO\hfill Add Saturated Half Word Octuple}
\label{instr:ADDSHO}\index{addsho}
 
\paragraph{Encoding} Format \textsf{ALU\_HOWRR1} (Section~\ref{sec:format-ALU-HOWRR1}), \verb|exucode2 = 1000|\\

\bitfields{ 1/P, 2/11, 1/1, 4/1000, 5/registerM, 1/--, 2/10, 3/101, 1/1, 5/registerO, 1/--, 5/registerP, 1/0 }


\paragraph{Syntax} \texttt{addsho} \emph{registerM} = \emph{registerP}, \emph{registerO}

\paragraph{Description} The \emph{registerO} and the \emph{registerP} are interpreted as eight half words packed into 128 bits. The \emph{registerO} half words and the \emph{registerP} half words are added with saturation to 16 bits. The eight resulting half words are packed and stored into the \emph{registerM}.


\paragraph{Execution} {Reservation table \textsf{ALU\_LITE} (Section~\ref{sec:reservation-ALU-LITE})}
~
{\begin{lstlisting}
stage RR:
  new argument3 = PGR[registerO];
  new argument2 = PGR[registerP];
  for (I = 0; I < 8; I += 1) {
    result1.16[I] = _SAT_16(_SX_16(argument3.16[I]) + _SX_16(argument2.16[I]))
  };
stage E1:
  PGR[registerM] = result1;
\end{lstlisting}}
\newpage

\paragraph{Encoding} Format \textsf{ALU\_HOWRR1.M} (Section~\ref{sec:format-ALU-HOWRR1.M}), \verb|exucode2 = 1000|\\

\bitfields{ 1/P, 2/00, 2/-- --, 27/upper27, 1/1, 2/11, 1/1, 4/1000, 5/registerM, 1/--, 2/10, 3/101, 1/1, 1/splat32, 5/lower5, 5/registerP, 1/0 }


\paragraph{Syntax} \texttt{addsho} \emph{registerM} = \emph{registerP}, \emph{upper27\_\discretionary{}{}{}lower5}\emph{splat32}

\paragraph{Description} The \emph{upper27\_\discretionary{}{}{}lower5} and the \emph{registerP} are interpreted as eight half words packed into 128 bits. The \emph{upper27\_\discretionary{}{}{}lower5} half words and the \emph{registerP} half words are added with saturation to 16 bits. The eight resulting half words are packed and stored into the \emph{registerM}.


\paragraph{Modifier(s)}
\noindent\begin{itemize}
\item splat32 (cf. splat32 in section \ref{par:modifier-splat32} on page \pageref{par:modifier-splat32})
\end{itemize}

\paragraph{Execution} {Reservation table \textsf{ALU\_LITE.X} (Section~\ref{sec:reservation-ALU-LITE.X})}
~
{\begin{lstlisting}
stage ID:
  new splat32 = _ZX_1(splat32);
stage RR:
  new argument3 = splat32 ? (_WX_32(upper27_lower5) << 32) | _ZX_32(_WX_32(upper27_lower5)) : _SX_32(_WX_32(upper27_lower5));
  new argument2 = PGR[registerP];
  for (I = 0; I < 8; I += 1) {
    result1.16[I] = _SAT_16(_SX_16(argument3.16[I]) + _SX_16(argument2.16[I]))
  };
stage E1:
  PGR[registerM] = result1;
\end{lstlisting}}
\newpage
\subsubsection[{\small ADDSW: Add Saturated Word}]{ADDSW\hfill Add Saturated Word}
\label{instr:ADDSW}\index{addsw}
 
\paragraph{Encoding} Format \textsf{ALU\_WWRR1} (Section~\ref{sec:format-ALU-WWRR1}), \verb|exucode2 = 1000|\\

\bitfields{ 1/P, 2/11, 1/0, 4/1000, 6/registerW, 2/11, 3/101, 1/signextw, 6/registerY, 6/registerZ }


\paragraph{Syntax} \texttt{addsw}\emph{signextw} \emph{registerW} = \emph{registerZ}, \emph{registerY}

\paragraph{Description} The \emph{registerY} and the \emph{registerZ} are added with 32-bit saturation. The result is stored into the \emph{registerW}.


\paragraph{Modifier(s)}
\noindent\begin{itemize}
\item signextw (cf. signextw in section \ref{par:modifier-signextw} on page \pageref{par:modifier-signextw})
\end{itemize}

\paragraph{Execution} {Reservation table \textsf{ALU\_TINY} (Section~\ref{sec:reservation-ALU-TINY})}
~
{\begin{lstlisting}
stage ID:
  new signextw = _ZX_1(signextw);
stage RR:
  new argument3 = GPR[registerY];
  new argument2 = GPR[registerZ];
  new result1 = _SAT_32(_SX_32(argument3) + _SX_32(argument2));
stage E1:
  GPR[registerW] = signextw ? _SX_32(result1) : _ZX_32(result1);
\end{lstlisting}}
\newpage

\paragraph{Encoding} Format \textsf{ALU\_WWRR1.W} (Section~\ref{sec:format-ALU-WWRR1.W}), \verb|exucode2 = 1000|\\

\bitfields{ 1/P, 2/00, 2/-- --, 27/upper27, 1/1, 2/11, 1/0, 4/1000, 6/registerW, 2/11, 3/101, 1/signextw, 1/0, 5/lower5, 6/registerZ }


\paragraph{Syntax} \texttt{addsw}\emph{signextw} \emph{registerW} = \emph{registerZ}, \emph{upper27\_\discretionary{}{}{}lower5}

\paragraph{Description} The \emph{upper27\_\discretionary{}{}{}lower5} and the \emph{registerZ} are added with 32-bit saturation. The result is stored into the \emph{registerW}.


\paragraph{Modifier(s)}
\noindent\begin{itemize}
\item signextw (cf. signextw in section \ref{par:modifier-signextw} on page \pageref{par:modifier-signextw})
\end{itemize}

\paragraph{Execution} {Reservation table \textsf{ALU\_TINY.X} (Section~\ref{sec:reservation-ALU-TINY.X})}
~
{\begin{lstlisting}
stage ID:
  new signextw = _ZX_1(signextw);
stage RR:
  new argument3 = _WX_32(upper27_lower5);
  new argument2 = GPR[registerZ];
  new result1 = _SAT_32(_SX_32(argument3) + _SX_32(argument2));
stage E1:
  GPR[registerW] = signextw ? _SX_32(result1) : _ZX_32(result1);
\end{lstlisting}}
\newpage
\subsubsection[{\small ADDSWQ: Add Saturated Word Quadruple}]{ADDSWQ\hfill Add Saturated Word Quadruple}
\label{instr:ADDSWQ}\index{addswq}
 
\paragraph{Encoding} Format \textsf{ALU\_WQWRR1} (Section~\ref{sec:format-ALU-WQWRR1}), \verb|exucode2 = 1000|\\

\bitfields{ 1/P, 2/11, 1/1, 4/1000, 5/registerM, 1/--, 2/10, 3/101, 1/0, 5/registerO, 1/--, 5/registerP, 1/1 }


\paragraph{Syntax} \texttt{addswq} \emph{registerM} = \emph{registerP}, \emph{registerO}

\paragraph{Description} The \emph{registerO} and the \emph{registerP} are interpreted as four words packed into 128 bits. The \emph{registerO} words and the \emph{registerP} words are added with saturation to 32 bits. The four resulting words are packed and stored into the \emph{registerM}.


\paragraph{Execution} {Reservation table \textsf{ALU\_LITE} (Section~\ref{sec:reservation-ALU-LITE})}
~
{\begin{lstlisting}
stage RR:
  new argument3 = PGR[registerO];
  new argument2 = PGR[registerP];
  for (I = 0; I < 4; I += 1) {
    result1.32[I] = _SAT_32(_SX_32(argument3.32[I]) + _SX_32(argument2.32[I]))
  };
stage E1:
  PGR[registerM] = result1;
\end{lstlisting}}
\newpage

\paragraph{Encoding} Format \textsf{ALU\_WQWRR1.M} (Section~\ref{sec:format-ALU-WQWRR1.M}), \verb|exucode2 = 1000|\\

\bitfields{ 1/P, 2/00, 2/-- --, 27/upper27, 1/1, 2/11, 1/1, 4/1000, 5/registerM, 1/--, 2/10, 3/101, 1/0, 1/splat32, 5/lower5, 5/registerP, 1/1 }


\paragraph{Syntax} \texttt{addswq} \emph{registerM} = \emph{registerP}, \emph{upper27\_\discretionary{}{}{}lower5}\emph{splat32}

\paragraph{Description} The \emph{upper27\_\discretionary{}{}{}lower5} and the \emph{registerP} are interpreted as four words packed into 128 bits. The \emph{upper27\_\discretionary{}{}{}lower5} words and the \emph{registerP} words are added with saturation to 32 bits. The four resulting words are packed and stored into the \emph{registerM}.


\paragraph{Modifier(s)}
\noindent\begin{itemize}
\item splat32 (cf. splat32 in section \ref{par:modifier-splat32} on page \pageref{par:modifier-splat32})
\end{itemize}

\paragraph{Execution} {Reservation table \textsf{ALU\_LITE.X} (Section~\ref{sec:reservation-ALU-LITE.X})}
~
{\begin{lstlisting}
stage ID:
  new splat32 = _ZX_1(splat32);
stage RR:
  new argument3 = splat32 ? (_WX_32(upper27_lower5) << 32) | _ZX_32(_WX_32(upper27_lower5)) : _SX_32(_WX_32(upper27_lower5));
  new argument2 = PGR[registerP];
  for (I = 0; I < 4; I += 1) {
    result1.32[I] = _SAT_32(_SX_32(argument3.32[I]) + _SX_32(argument2.32[I]))
  };
stage E1:
  PGR[registerM] = result1;
\end{lstlisting}}
\newpage
\subsubsection[{\small ADDUSBX: Add Unsigned Saturated Byte Hexadecuple}]{ADDUSBX\hfill Add Unsigned Saturated Byte Hexadecuple}
\label{instr:ADDUSBX}\index{addusbx}
 
\paragraph{Encoding} Format \textsf{ALU\_BXWRR1} (Section~\ref{sec:format-ALU-BXWRR1}), \verb|exucode2 = 1010|\\

\bitfields{ 1/P, 2/11, 1/1, 4/1010, 5/registerM, 1/--, 2/10, 3/101, 1/1, 5/registerO, 1/--, 5/registerP, 1/1 }


\paragraph{Syntax} \texttt{addusbx} \emph{registerM} = \emph{registerP}, \emph{registerO}

\paragraph{Description} The \emph{registerO} and the \emph{registerP} are interpreted as sixteen bytes packed into 128 bits. The \emph{registerO} bytes and the \emph{registerP} bytes are added with unsigned saturation to 8 bits. The sixteen resulting bytes are packed and stored into the \emph{registerM}.


\paragraph{Execution} {Reservation table \textsf{ALU\_LITE} (Section~\ref{sec:reservation-ALU-LITE})}
~
{\begin{lstlisting}
stage RR:
  new argument3 = PGR[registerO];
  new argument2 = PGR[registerP];
  for (I = 0; I < 16; I += 1) {
    result1.8[I] = _SATU_8(_ZX_8(argument3.8[I]) + _ZX_8(argument2.8[I]))
  };
stage E1:
  PGR[registerM] = result1;
\end{lstlisting}}
\newpage

\paragraph{Encoding} Format \textsf{ALU\_BXWRR1.M} (Section~\ref{sec:format-ALU-BXWRR1.M}), \verb|exucode2 = 1010|\\

\bitfields{ 1/P, 2/00, 2/-- --, 27/upper27, 1/1, 2/11, 1/1, 4/1010, 5/registerM, 1/--, 2/10, 3/101, 1/1, 1/splat32, 5/lower5, 5/registerP, 1/1 }


\paragraph{Syntax} \texttt{addusbx} \emph{registerM} = \emph{registerP}, \emph{upper27\_\discretionary{}{}{}lower5}\emph{splat32}

\paragraph{Description} The \emph{upper27\_\discretionary{}{}{}lower5} and the \emph{registerP} are interpreted as sixteen bytes packed into 128 bits. The \emph{upper27\_\discretionary{}{}{}lower5} bytes and the \emph{registerP} bytes are added with unsigned saturation to 8 bits. The sixteen resulting bytes are packed and stored into the \emph{registerM}.


\paragraph{Modifier(s)}
\noindent\begin{itemize}
\item splat32 (cf. splat32 in section \ref{par:modifier-splat32} on page \pageref{par:modifier-splat32})
\end{itemize}

\paragraph{Execution} {Reservation table \textsf{ALU\_LITE.X} (Section~\ref{sec:reservation-ALU-LITE.X})}
~
{\begin{lstlisting}
stage ID:
  new splat32 = _ZX_1(splat32);
stage RR:
  new argument3 = splat32 ? (_WX_32(upper27_lower5) << 32) | _ZX_32(_WX_32(upper27_lower5)) : _SX_32(_WX_32(upper27_lower5));
  new argument2 = PGR[registerP];
  for (I = 0; I < 16; I += 1) {
    result1.8[I] = _SATU_8(_ZX_8(argument3.8[I]) + _ZX_8(argument2.8[I]))
  };
stage E1:
  PGR[registerM] = result1;
\end{lstlisting}}
\newpage
\subsubsection[{\small ADDUSD: Add Unsigned Saturated Double Word}]{ADDUSD\hfill Add Unsigned Saturated Double Word}
\label{instr:ADDUSD}\index{addusd}
 
\paragraph{Encoding} Format \textsf{ALU\_DWRR1} (Section~\ref{sec:format-ALU-DWRR1}), \verb|exucode2 = 1010|\\

\bitfields{ 1/P, 2/11, 1/0, 4/1010, 6/registerW, 2/10, 3/101, 1/0, 6/registerY, 6/registerZ }


\paragraph{Syntax} \texttt{addusd} \emph{registerW} = \emph{registerZ}, \emph{registerY}

\paragraph{Description} The \emph{registerY} and the \emph{registerZ} are added with 64-bit unsigned saturation.


\paragraph{Execution} {Reservation table \textsf{ALU\_TINY} (Section~\ref{sec:reservation-ALU-TINY})}
~
{\begin{lstlisting}
stage RR:
  new argument3 = GPR[registerY];
  new argument2 = GPR[registerZ];
  new result1 = _SATU_64(_ZX_64(argument3) + _ZX_64(argument2));
stage E1:
  GPR[registerW] = result1;
\end{lstlisting}}
\newpage

\paragraph{Encoding} Format \textsf{ALU\_DWRR1.M} (Section~\ref{sec:format-ALU-DWRR1.M}), \verb|exucode2 = 1010|\\

\bitfields{ 1/P, 2/00, 2/-- --, 27/upper27, 1/1, 2/11, 1/0, 4/1010, 6/registerW, 2/10, 3/101, 1/0, 1/splat32, 5/lower5, 6/registerZ }


\paragraph{Syntax} \texttt{addusd} \emph{registerW} = \emph{registerZ}, \emph{upper27\_\discretionary{}{}{}lower5}\emph{splat32}

\paragraph{Description} The \emph{upper27\_\discretionary{}{}{}lower5} and the \emph{registerZ} are added with 64-bit unsigned saturation.


\paragraph{Modifier(s)}
\noindent\begin{itemize}
\item splat32 (cf. splat32 in section \ref{par:modifier-splat32} on page \pageref{par:modifier-splat32})
\end{itemize}

\paragraph{Execution} {Reservation table \textsf{ALU\_TINY.X} (Section~\ref{sec:reservation-ALU-TINY.X})}
~
{\begin{lstlisting}
stage ID:
  new splat32 = _ZX_1(splat32);
stage RR:
  new argument3 = splat32 ? (_WX_32(upper27_lower5) << 32) | _ZX_32(_WX_32(upper27_lower5)) : _SX_32(_WX_32(upper27_lower5));
  new argument2 = GPR[registerZ];
  new result1 = _SATU_64(_ZX_64(argument3) + _ZX_64(argument2));
stage E1:
  GPR[registerW] = result1;
\end{lstlisting}}
\newpage
\subsubsection[{\small ADDUSDP: Add Unsigned Saturated Double Word Pair}]{ADDUSDP\hfill Add Unsigned Saturated Double Word Pair}
\label{instr:ADDUSDP}\index{addusdp}
 
\paragraph{Encoding} Format \textsf{ALU\_DPWRR1} (Section~\ref{sec:format-ALU-DPWRR1}), \verb|exucode2 = 1010|\\

\bitfields{ 1/P, 2/11, 1/1, 4/1010, 5/registerM, 1/--, 2/10, 3/101, 1/0, 5/registerO, 1/--, 5/registerP, 1/0 }


\paragraph{Syntax} \texttt{addusdp} \emph{registerM} = \emph{registerP}, \emph{registerO}

\paragraph{Description} The \emph{registerO} and the \emph{registerP} are interpreted as two double words packed into 128 bits. The \emph{registerO} double words and the \emph{registerP} double words are added with unsigned saturation to 64 bits. The two resulting double words are packed and stored into the \emph{registerM}.


\paragraph{Execution} {Reservation table \textsf{ALU\_LITE} (Section~\ref{sec:reservation-ALU-LITE})}
~
{\begin{lstlisting}
stage RR:
  new argument3 = PGR[registerO];
  new argument2 = PGR[registerP];
  for (I = 0; I < 2; I += 1) {
    result1.64[I] = _SATU_64(_ZX_64(argument3.64[I]) + _ZX_64(argument2.64[I]))
  };
stage E1:
  PGR[registerM] = result1;
\end{lstlisting}}
\newpage

\paragraph{Encoding} Format \textsf{ALU\_DPWRR1.M} (Section~\ref{sec:format-ALU-DPWRR1.M}), \verb|exucode2 = 1010|\\

\bitfields{ 1/P, 2/00, 2/-- --, 27/upper27, 1/1, 2/11, 1/1, 4/1010, 5/registerM, 1/--, 2/10, 3/101, 1/0, 1/splat32, 5/lower5, 5/registerP, 1/0 }


\paragraph{Syntax} \texttt{addusdp} \emph{registerM} = \emph{registerP}, \emph{upper27\_\discretionary{}{}{}lower5}\emph{splat32}

\paragraph{Description} The \emph{upper27\_\discretionary{}{}{}lower5} and the \emph{registerP} are interpreted as two double words packed into 128 bits. The \emph{upper27\_\discretionary{}{}{}lower5} double words and the \emph{registerP} double words are added with unsigned saturation to 64 bits. The two resulting double words are packed and stored into the \emph{registerM}.


\paragraph{Modifier(s)}
\noindent\begin{itemize}
\item splat32 (cf. splat32 in section \ref{par:modifier-splat32} on page \pageref{par:modifier-splat32})
\end{itemize}

\paragraph{Execution} {Reservation table \textsf{ALU\_LITE.X} (Section~\ref{sec:reservation-ALU-LITE.X})}
~
{\begin{lstlisting}
stage ID:
  new splat32 = _ZX_1(splat32);
stage RR:
  new argument3 = splat32 ? (_WX_32(upper27_lower5) << 32) | _ZX_32(_WX_32(upper27_lower5)) : _SX_32(_WX_32(upper27_lower5));
  new argument2 = PGR[registerP];
  for (I = 0; I < 2; I += 1) {
    result1.64[I] = _SATU_64(_ZX_64(argument3.64[I]) + _ZX_64(argument2.64[I]))
  };
stage E1:
  PGR[registerM] = result1;
\end{lstlisting}}
\newpage
\subsubsection[{\small ADDUSHO: Add Unsigned Saturated Half Word Octuple}]{ADDUSHO\hfill Add Unsigned Saturated Half Word Octuple}
\label{instr:ADDUSHO}\index{addusho}
 
\paragraph{Encoding} Format \textsf{ALU\_HOWRR1} (Section~\ref{sec:format-ALU-HOWRR1}), \verb|exucode2 = 1010|\\

\bitfields{ 1/P, 2/11, 1/1, 4/1010, 5/registerM, 1/--, 2/10, 3/101, 1/1, 5/registerO, 1/--, 5/registerP, 1/0 }


\paragraph{Syntax} \texttt{addusho} \emph{registerM} = \emph{registerP}, \emph{registerO}

\paragraph{Description} The \emph{registerO} and the \emph{registerP} are interpreted as eight half words packed into 128 bits. The \emph{registerO} half words and the \emph{registerP} half words are added with unsigned saturation to 16 bits. The eight resulting half words are packed and stored into the \emph{registerM}.


\paragraph{Execution} {Reservation table \textsf{ALU\_LITE} (Section~\ref{sec:reservation-ALU-LITE})}
~
{\begin{lstlisting}
stage RR:
  new argument3 = PGR[registerO];
  new argument2 = PGR[registerP];
  for (I = 0; I < 8; I += 1) {
    result1.16[I] = _SATU_16(_ZX_16(argument3.16[I]) + _ZX_16(argument2.16[I]))
  };
stage E1:
  PGR[registerM] = result1;
\end{lstlisting}}
\newpage

\paragraph{Encoding} Format \textsf{ALU\_HOWRR1.M} (Section~\ref{sec:format-ALU-HOWRR1.M}), \verb|exucode2 = 1010|\\

\bitfields{ 1/P, 2/00, 2/-- --, 27/upper27, 1/1, 2/11, 1/1, 4/1010, 5/registerM, 1/--, 2/10, 3/101, 1/1, 1/splat32, 5/lower5, 5/registerP, 1/0 }


\paragraph{Syntax} \texttt{addusho} \emph{registerM} = \emph{registerP}, \emph{upper27\_\discretionary{}{}{}lower5}\emph{splat32}

\paragraph{Description} The \emph{upper27\_\discretionary{}{}{}lower5} and the \emph{registerP} are interpreted as eight half words packed into 128 bits. The \emph{upper27\_\discretionary{}{}{}lower5} half words and the \emph{registerP} half words are added with unsigned saturation to 16 bits. The eight resulting half words are packed and stored into the \emph{registerM}.


\paragraph{Modifier(s)}
\noindent\begin{itemize}
\item splat32 (cf. splat32 in section \ref{par:modifier-splat32} on page \pageref{par:modifier-splat32})
\end{itemize}

\paragraph{Execution} {Reservation table \textsf{ALU\_LITE.X} (Section~\ref{sec:reservation-ALU-LITE.X})}
~
{\begin{lstlisting}
stage ID:
  new splat32 = _ZX_1(splat32);
stage RR:
  new argument3 = splat32 ? (_WX_32(upper27_lower5) << 32) | _ZX_32(_WX_32(upper27_lower5)) : _SX_32(_WX_32(upper27_lower5));
  new argument2 = PGR[registerP];
  for (I = 0; I < 8; I += 1) {
    result1.16[I] = _SATU_16(_ZX_16(argument3.16[I]) + _ZX_16(argument2.16[I]))
  };
stage E1:
  PGR[registerM] = result1;
\end{lstlisting}}
\newpage
\subsubsection[{\small ADDUSW: Add Unsigned Saturated Word}]{ADDUSW\hfill Add Unsigned Saturated Word}
\label{instr:ADDUSW}\index{addusw}
 
\paragraph{Encoding} Format \textsf{ALU\_WWRR1} (Section~\ref{sec:format-ALU-WWRR1}), \verb|exucode2 = 1010|\\

\bitfields{ 1/P, 2/11, 1/0, 4/1010, 6/registerW, 2/11, 3/101, 1/signextw, 6/registerY, 6/registerZ }


\paragraph{Syntax} \texttt{addusw}\emph{signextw} \emph{registerW} = \emph{registerZ}, \emph{registerY}

\paragraph{Description} The \emph{registerY} and the \emph{registerZ} are added with 32-bit unsigned saturation. The result with the 32 upper bits cleared is stored into the \emph{registerW}.


\paragraph{Modifier(s)}
\noindent\begin{itemize}
\item signextw (cf. signextw in section \ref{par:modifier-signextw} on page \pageref{par:modifier-signextw})
\end{itemize}

\paragraph{Execution} {Reservation table \textsf{ALU\_TINY} (Section~\ref{sec:reservation-ALU-TINY})}
~
{\begin{lstlisting}
stage ID:
  new signextw = _ZX_1(signextw);
stage RR:
  new argument3 = GPR[registerY];
  new argument2 = GPR[registerZ];
  new result1 = _SATU_32(_ZX_32(argument3) + _ZX_32(argument2));
stage E1:
  GPR[registerW] = signextw ? _SX_32(result1) : _ZX_32(result1);
\end{lstlisting}}
\newpage

\paragraph{Encoding} Format \textsf{ALU\_WWRR1.W} (Section~\ref{sec:format-ALU-WWRR1.W}), \verb|exucode2 = 1010|\\

\bitfields{ 1/P, 2/00, 2/-- --, 27/upper27, 1/1, 2/11, 1/0, 4/1010, 6/registerW, 2/11, 3/101, 1/signextw, 1/0, 5/lower5, 6/registerZ }


\paragraph{Syntax} \texttt{addusw}\emph{signextw} \emph{registerW} = \emph{registerZ}, \emph{upper27\_\discretionary{}{}{}lower5}

\paragraph{Description} The \emph{upper27\_\discretionary{}{}{}lower5} and the \emph{registerZ} are added with 32-bit unsigned saturation. The result with the 32 upper bits cleared is stored into the \emph{registerW}.


\paragraph{Modifier(s)}
\noindent\begin{itemize}
\item signextw (cf. signextw in section \ref{par:modifier-signextw} on page \pageref{par:modifier-signextw})
\end{itemize}

\paragraph{Execution} {Reservation table \textsf{ALU\_TINY.X} (Section~\ref{sec:reservation-ALU-TINY.X})}
~
{\begin{lstlisting}
stage ID:
  new signextw = _ZX_1(signextw);
stage RR:
  new argument3 = _WX_32(upper27_lower5);
  new argument2 = GPR[registerZ];
  new result1 = _SATU_32(_ZX_32(argument3) + _ZX_32(argument2));
stage E1:
  GPR[registerW] = signextw ? _SX_32(result1) : _ZX_32(result1);
\end{lstlisting}}
\newpage
\subsubsection[{\small ADDUSWQ: Add Unsigned Saturated Word Quadruple}]{ADDUSWQ\hfill Add Unsigned Saturated Word Quadruple}
\label{instr:ADDUSWQ}\index{adduswq}
 
\paragraph{Encoding} Format \textsf{ALU\_WQWRR1} (Section~\ref{sec:format-ALU-WQWRR1}), \verb|exucode2 = 1010|\\

\bitfields{ 1/P, 2/11, 1/1, 4/1010, 5/registerM, 1/--, 2/10, 3/101, 1/0, 5/registerO, 1/--, 5/registerP, 1/1 }


\paragraph{Syntax} \texttt{adduswq} \emph{registerM} = \emph{registerP}, \emph{registerO}

\paragraph{Description} The \emph{registerO} and the \emph{registerP} are interpreted as four words packed into 128 bits. The \emph{registerO} words and the \emph{registerP} words are added with unsigned saturation to 32 bits. The four resulting words are packed and stored into the \emph{registerM}.


\paragraph{Execution} {Reservation table \textsf{ALU\_LITE} (Section~\ref{sec:reservation-ALU-LITE})}
~
{\begin{lstlisting}
stage RR:
  new argument3 = PGR[registerO];
  new argument2 = PGR[registerP];
  for (I = 0; I < 4; I += 1) {
    result1.32[I] = _SATU_32(_ZX_32(argument3.32[I]) + _ZX_32(argument2.32[I]))
  };
stage E1:
  PGR[registerM] = result1;
\end{lstlisting}}
\newpage

\paragraph{Encoding} Format \textsf{ALU\_WQWRR1.M} (Section~\ref{sec:format-ALU-WQWRR1.M}), \verb|exucode2 = 1010|\\

\bitfields{ 1/P, 2/00, 2/-- --, 27/upper27, 1/1, 2/11, 1/1, 4/1010, 5/registerM, 1/--, 2/10, 3/101, 1/0, 1/splat32, 5/lower5, 5/registerP, 1/1 }


\paragraph{Syntax} \texttt{adduswq} \emph{registerM} = \emph{registerP}, \emph{upper27\_\discretionary{}{}{}lower5}\emph{splat32}

\paragraph{Description} The \emph{upper27\_\discretionary{}{}{}lower5} and the \emph{registerP} are interpreted as four words packed into 128 bits. The \emph{upper27\_\discretionary{}{}{}lower5} words and the \emph{registerP} words are added with unsigned saturation to 32 bits. The four resulting words are packed and stored into the \emph{registerM}.


\paragraph{Modifier(s)}
\noindent\begin{itemize}
\item splat32 (cf. splat32 in section \ref{par:modifier-splat32} on page \pageref{par:modifier-splat32})
\end{itemize}

\paragraph{Execution} {Reservation table \textsf{ALU\_LITE.X} (Section~\ref{sec:reservation-ALU-LITE.X})}
~
{\begin{lstlisting}
stage ID:
  new splat32 = _ZX_1(splat32);
stage RR:
  new argument3 = splat32 ? (_WX_32(upper27_lower5) << 32) | _ZX_32(_WX_32(upper27_lower5)) : _SX_32(_WX_32(upper27_lower5));
  new argument2 = PGR[registerP];
  for (I = 0; I < 4; I += 1) {
    result1.32[I] = _SATU_32(_ZX_32(argument3.32[I]) + _ZX_32(argument2.32[I]))
  };
stage E1:
  PGR[registerM] = result1;
\end{lstlisting}}
\newpage
\subsubsection[{\small ADDW: Add Word to Word}]{ADDW\hfill Add Word to Word}
\label{instr:ADDW}\index{addw}
 
\paragraph{Encoding} Format \textsf{ALU\_WWRR0} (Section~\ref{sec:format-ALU-WWRR0}), \verb|exucode2 = 0010|\\

\bitfields{ 1/P, 2/11, 1/0, 4/0010, 6/registerW, 2/11, 3/000, 1/signextw, 6/registerY, 6/registerZ }


\paragraph{Syntax} \texttt{addw}\emph{signextw} \emph{registerW} = \emph{registerZ}, \emph{registerY}

\paragraph{Description} The \emph{registerY} and the \emph{registerZ} are added. The result with the 32 upper bits cleared is stored into the \emph{registerW}.


\paragraph{Modifier(s)}
\noindent\begin{itemize}
\item signextw (cf. signextw in section \ref{par:modifier-signextw} on page \pageref{par:modifier-signextw})
\end{itemize}

\paragraph{Execution} {Reservation table \textsf{ALU\_TINY} (Section~\ref{sec:reservation-ALU-TINY})}
~
{\begin{lstlisting}
stage ID:
  new signextw = _ZX_1(signextw);
stage RR:
  new argument3 = GPR[registerY];
  new argument2 = GPR[registerZ];
  new result1 = argument3 + argument2;
stage E1:
  GPR[registerW] = signextw ? _SX_32(result1) : _ZX_32(result1);
\end{lstlisting}}
\newpage

\paragraph{Encoding} Format \textsf{ALU\_WWRR0.W} (Section~\ref{sec:format-ALU-WWRR0.W}), \verb|exucode2 = 0010|\\

\bitfields{ 1/P, 2/00, 2/-- --, 27/upper27, 1/1, 2/11, 1/0, 4/0010, 6/registerW, 2/11, 3/000, 1/signextw, 1/0, 5/lower5, 6/registerZ }


\paragraph{Syntax} \texttt{addw}\emph{signextw} \emph{registerW} = \emph{registerZ}, \emph{upper27\_\discretionary{}{}{}lower5}

\paragraph{Description} The \emph{upper27\_\discretionary{}{}{}lower5} and the \emph{registerZ} are added. The result with the 32 upper bits cleared is stored into the \emph{registerW}.


\paragraph{Modifier(s)}
\noindent\begin{itemize}
\item signextw (cf. signextw in section \ref{par:modifier-signextw} on page \pageref{par:modifier-signextw})
\end{itemize}

\paragraph{Execution} {Reservation table \textsf{ALU\_TINY.X} (Section~\ref{sec:reservation-ALU-TINY.X})}
~
{\begin{lstlisting}
stage ID:
  new signextw = _ZX_1(signextw);
stage RR:
  new argument3 = _WX_32(upper27_lower5);
  new argument2 = GPR[registerZ];
  new result1 = argument3 + argument2;
stage E1:
  GPR[registerW] = signextw ? _SX_32(result1) : _ZX_32(result1);
\end{lstlisting}}
\newpage
\subsubsection[{\small ADDWQ: Add Word Quadruple}]{ADDWQ\hfill Add Word Quadruple}
\label{instr:ADDWQ}\index{addwq}
 
\paragraph{Encoding} Format \textsf{ALU\_WQWRR0} (Section~\ref{sec:format-ALU-WQWRR0}), \verb|exucode2 = 0010|\\

\bitfields{ 1/P, 2/11, 1/1, 4/0010, 5/registerM, 1/--, 2/10, 3/000, 1/0, 5/registerO, 1/--, 5/registerP, 1/1 }


\paragraph{Syntax} \texttt{addwq} \emph{registerM} = \emph{registerP}, \emph{registerO}

\paragraph{Description} The \emph{registerO} and the \emph{registerP} are interpreted as four words packed into 128 bits. The \emph{registerO} words and the \emph{registerP} words are added. The four resulting words are packed and stored into the \emph{registerM}.


\paragraph{Execution} {Reservation table \textsf{ALU\_LITE} (Section~\ref{sec:reservation-ALU-LITE})}
~
{\begin{lstlisting}
stage RR:
  new argument3 = PGR[registerO];
  new argument2 = PGR[registerP];
  for (I = 0; I < 4; I += 1) {
    result1.32[I] = argument3.32[I] + argument2.32[I]
  };
stage E1:
  PGR[registerM] = result1;
\end{lstlisting}}
\newpage

\paragraph{Encoding} Format \textsf{ALU\_WQWRR0.M} (Section~\ref{sec:format-ALU-WQWRR0.M}), \verb|exucode2 = 0010|\\

\bitfields{ 1/P, 2/00, 2/-- --, 27/upper27, 1/1, 2/11, 1/1, 4/0010, 5/registerM, 1/--, 2/10, 3/000, 1/0, 1/splat32, 5/lower5, 5/registerP, 1/1 }


\paragraph{Syntax} \texttt{addwq} \emph{registerM} = \emph{registerP}, \emph{upper27\_\discretionary{}{}{}lower5}\emph{splat32}

\paragraph{Description} The \emph{upper27\_\discretionary{}{}{}lower5} and the \emph{registerP} are interpreted as four words packed into 128 bits. The \emph{upper27\_\discretionary{}{}{}lower5} words and the \emph{registerP} words are added. The four resulting words are packed and stored into the \emph{registerM}.


\paragraph{Modifier(s)}
\noindent\begin{itemize}
\item splat32 (cf. splat32 in section \ref{par:modifier-splat32} on page \pageref{par:modifier-splat32})
\end{itemize}

\paragraph{Execution} {Reservation table \textsf{ALU\_LITE.X} (Section~\ref{sec:reservation-ALU-LITE.X})}
~
{\begin{lstlisting}
stage ID:
  new splat32 = _ZX_1(splat32);
stage RR:
  new argument3 = splat32 ? (_WX_32(upper27_lower5) << 32) | _ZX_32(_WX_32(upper27_lower5)) : _SX_32(_WX_32(upper27_lower5));
  new argument2 = PGR[registerP];
  for (I = 0; I < 4; I += 1) {
    result1.32[I] = argument3.32[I] + argument2.32[I]
  };
stage E1:
  PGR[registerM] = result1;
\end{lstlisting}}
\newpage
\subsubsection[{\small ADDX16BX: Add to Times 16 Byte Hexadecuple}]{ADDX16BX\hfill Add to Times 16 Byte Hexadecuple}
\label{instr:ADDX16BX}\index{addx16bx}
 
\paragraph{Encoding} Format \textsf{ALU\_BXWRR1} (Section~\ref{sec:format-ALU-BXWRR1}), \verb|exucode2 = 0011|\\

\bitfields{ 1/P, 2/11, 1/1, 4/0011, 5/registerM, 1/--, 2/10, 3/101, 1/1, 5/registerO, 1/--, 5/registerP, 1/1 }


\paragraph{Syntax} \texttt{addx16bx} \emph{registerM} = \emph{registerP}, \emph{registerO}

\paragraph{Description} The \emph{registerO} and the \emph{registerP} are interpreted as sixteen bytes packed into 128 bits. The \emph{registerP} bytes are shifted left by four bits and added to the \emph{registerO} bytes. The sixteen resulting bytes are packed and stored into the \emph{registerM}.


\paragraph{Execution} {Reservation table \textsf{ALU\_LITE} (Section~\ref{sec:reservation-ALU-LITE})}
~
{\begin{lstlisting}
stage RR:
  new argument3 = PGR[registerO];
  new argument2 = PGR[registerP];
  for (I = 0; I < 16; I += 1) {
    result1.8[I] = _SX_8(argument3.8[I]) + (_SX_8(argument2.8[I]) << 4)
  };
stage E1:
  PGR[registerM] = result1;
\end{lstlisting}}
\newpage

\paragraph{Encoding} Format \textsf{ALU\_BXWRR1.M} (Section~\ref{sec:format-ALU-BXWRR1.M}), \verb|exucode2 = 0011|\\

\bitfields{ 1/P, 2/00, 2/-- --, 27/upper27, 1/1, 2/11, 1/1, 4/0011, 5/registerM, 1/--, 2/10, 3/101, 1/1, 1/splat32, 5/lower5, 5/registerP, 1/1 }


\paragraph{Syntax} \texttt{addx16bx} \emph{registerM} = \emph{registerP}, \emph{upper27\_\discretionary{}{}{}lower5}\emph{splat32}

\paragraph{Description} The \emph{upper27\_\discretionary{}{}{}lower5} and the \emph{registerP} are interpreted as sixteen bytes packed into 128 bits. The \emph{registerP} bytes are shifted left by four bits and added to the \emph{upper27\_\discretionary{}{}{}lower5} bytes. The sixteen resulting bytes are packed and stored into the \emph{registerM}.


\paragraph{Modifier(s)}
\noindent\begin{itemize}
\item splat32 (cf. splat32 in section \ref{par:modifier-splat32} on page \pageref{par:modifier-splat32})
\end{itemize}

\paragraph{Execution} {Reservation table \textsf{ALU\_LITE.X} (Section~\ref{sec:reservation-ALU-LITE.X})}
~
{\begin{lstlisting}
stage ID:
  new splat32 = _ZX_1(splat32);
stage RR:
  new argument3 = splat32 ? (_WX_32(upper27_lower5) << 32) | _ZX_32(_WX_32(upper27_lower5)) : _SX_32(_WX_32(upper27_lower5));
  new argument2 = PGR[registerP];
  for (I = 0; I < 16; I += 1) {
    result1.8[I] = _SX_8(argument3.8[I]) + (_SX_8(argument2.8[I]) << 4)
  };
stage E1:
  PGR[registerM] = result1;
\end{lstlisting}}
\newpage
\subsubsection[{\small ADDX16D: Add Double Word to Double Word Times 16}]{ADDX16D\hfill Add Double Word to Double Word Times 16}
\label{instr:ADDX16D}\index{addx16d}
 
\paragraph{Encoding} Format \textsf{ALU\_DWRR1} (Section~\ref{sec:format-ALU-DWRR1}), \verb|exucode2 = 0011|\\

\bitfields{ 1/P, 2/11, 1/0, 4/0011, 6/registerW, 2/10, 3/101, 1/0, 6/registerY, 6/registerZ }


\paragraph{Syntax} \texttt{addx16d} \emph{registerW} = \emph{registerZ}, \emph{registerY}

\paragraph{Description} The \emph{registerZ} is shifted left by four bits and added to the \emph{registerY}. The result is stored into the \emph{registerW}.


\paragraph{Execution} {Reservation table \textsf{ALU\_TINY} (Section~\ref{sec:reservation-ALU-TINY})}
~
{\begin{lstlisting}
stage RR:
  new argument3 = GPR[registerY];
  new argument2 = GPR[registerZ];
  new result1 = argument3 + (argument2 << 4);
stage E1:
  GPR[registerW] = result1;
\end{lstlisting}}
\newpage

\paragraph{Encoding} Format \textsf{ALU\_DWRR1.M} (Section~\ref{sec:format-ALU-DWRR1.M}), \verb|exucode2 = 0011|\\

\bitfields{ 1/P, 2/00, 2/-- --, 27/upper27, 1/1, 2/11, 1/0, 4/0011, 6/registerW, 2/10, 3/101, 1/0, 1/splat32, 5/lower5, 6/registerZ }


\paragraph{Syntax} \texttt{addx16d} \emph{registerW} = \emph{registerZ}, \emph{upper27\_\discretionary{}{}{}lower5}\emph{splat32}

\paragraph{Description} The \emph{registerZ} is shifted left by four bits and added to the \emph{upper27\_\discretionary{}{}{}lower5}. The result is stored into the \emph{registerW}.


\paragraph{Modifier(s)}
\noindent\begin{itemize}
\item splat32 (cf. splat32 in section \ref{par:modifier-splat32} on page \pageref{par:modifier-splat32})
\end{itemize}

\paragraph{Execution} {Reservation table \textsf{ALU\_TINY.X} (Section~\ref{sec:reservation-ALU-TINY.X})}
~
{\begin{lstlisting}
stage ID:
  new splat32 = _ZX_1(splat32);
stage RR:
  new argument3 = splat32 ? (_WX_32(upper27_lower5) << 32) | _ZX_32(_WX_32(upper27_lower5)) : _SX_32(_WX_32(upper27_lower5));
  new argument2 = GPR[registerZ];
  new result1 = argument3 + (argument2 << 4);
stage E1:
  GPR[registerW] = result1;
\end{lstlisting}}
\newpage
\subsubsection[{\small ADDX16DP: Add to Times 16 Double Word Pair}]{ADDX16DP\hfill Add to Times 16 Double Word Pair}
\label{instr:ADDX16DP}\index{addx16dp}
 
\paragraph{Encoding} Format \textsf{ALU\_DPWRR1} (Section~\ref{sec:format-ALU-DPWRR1}), \verb|exucode2 = 0011|\\

\bitfields{ 1/P, 2/11, 1/1, 4/0011, 5/registerM, 1/--, 2/10, 3/101, 1/0, 5/registerO, 1/--, 5/registerP, 1/0 }


\paragraph{Syntax} \texttt{addx16dp} \emph{registerM} = \emph{registerP}, \emph{registerO}

\paragraph{Description} The \emph{registerO} and the \emph{registerP} are interpreted as two double words packed into 128 bits. The \emph{registerP} double words are shifted left by four bits and added to the \emph{registerO} double words. The two resulting double words are packed and stored into the \emph{registerM}.


\paragraph{Execution} {Reservation table \textsf{ALU\_LITE} (Section~\ref{sec:reservation-ALU-LITE})}
~
{\begin{lstlisting}
stage RR:
  new argument3 = PGR[registerO];
  new argument2 = PGR[registerP];
  for (I = 0; I < 2; I += 1) {
    result1.64[I] = _SX_64(argument3.64[I]) + (_SX_64(argument2.64[I]) << 4)
  };
stage E1:
  PGR[registerM] = result1;
\end{lstlisting}}
\newpage

\paragraph{Encoding} Format \textsf{ALU\_DPWRR1.M} (Section~\ref{sec:format-ALU-DPWRR1.M}), \verb|exucode2 = 0011|\\

\bitfields{ 1/P, 2/00, 2/-- --, 27/upper27, 1/1, 2/11, 1/1, 4/0011, 5/registerM, 1/--, 2/10, 3/101, 1/0, 1/splat32, 5/lower5, 5/registerP, 1/0 }


\paragraph{Syntax} \texttt{addx16dp} \emph{registerM} = \emph{registerP}, \emph{upper27\_\discretionary{}{}{}lower5}\emph{splat32}

\paragraph{Description} The \emph{upper27\_\discretionary{}{}{}lower5} and the \emph{registerP} are interpreted as two double words packed into 128 bits. The \emph{registerP} double words are shifted left by four bits and added to the \emph{upper27\_\discretionary{}{}{}lower5} double words. The two resulting double words are packed and stored into the \emph{registerM}.


\paragraph{Modifier(s)}
\noindent\begin{itemize}
\item splat32 (cf. splat32 in section \ref{par:modifier-splat32} on page \pageref{par:modifier-splat32})
\end{itemize}

\paragraph{Execution} {Reservation table \textsf{ALU\_LITE.X} (Section~\ref{sec:reservation-ALU-LITE.X})}
~
{\begin{lstlisting}
stage ID:
  new splat32 = _ZX_1(splat32);
stage RR:
  new argument3 = splat32 ? (_WX_32(upper27_lower5) << 32) | _ZX_32(_WX_32(upper27_lower5)) : _SX_32(_WX_32(upper27_lower5));
  new argument2 = PGR[registerP];
  for (I = 0; I < 2; I += 1) {
    result1.64[I] = _SX_64(argument3.64[I]) + (_SX_64(argument2.64[I]) << 4)
  };
stage E1:
  PGR[registerM] = result1;
\end{lstlisting}}
\newpage
\subsubsection[{\small ADDX16HO: Add to Times 16 Half Word Octuple}]{ADDX16HO\hfill Add to Times 16 Half Word Octuple}
\label{instr:ADDX16HO}\index{addx16ho}
 
\paragraph{Encoding} Format \textsf{ALU\_HOWRR1} (Section~\ref{sec:format-ALU-HOWRR1}), \verb|exucode2 = 0011|\\

\bitfields{ 1/P, 2/11, 1/1, 4/0011, 5/registerM, 1/--, 2/10, 3/101, 1/1, 5/registerO, 1/--, 5/registerP, 1/0 }


\paragraph{Syntax} \texttt{addx16ho} \emph{registerM} = \emph{registerP}, \emph{registerO}

\paragraph{Description} The \emph{registerO} and the \emph{registerP} are interpreted as eight half words packed into 128 bits. The \emph{registerP} half words are shifted left by eight bits and added to the \emph{registerO} half words. The eight resulting half words are packed and stored into the \emph{registerM}.


\paragraph{Execution} {Reservation table \textsf{ALU\_LITE} (Section~\ref{sec:reservation-ALU-LITE})}
~
{\begin{lstlisting}
stage RR:
  new argument3 = PGR[registerO];
  new argument2 = PGR[registerP];
  for (I = 0; I < 8; I += 1) {
    result1.16[I] = _SX_16(argument3.16[I]) + (_SX_16(argument2.16[I]) << 4)
  };
stage E1:
  PGR[registerM] = result1;
\end{lstlisting}}
\newpage

\paragraph{Encoding} Format \textsf{ALU\_HOWRR1.M} (Section~\ref{sec:format-ALU-HOWRR1.M}), \verb|exucode2 = 0011|\\

\bitfields{ 1/P, 2/00, 2/-- --, 27/upper27, 1/1, 2/11, 1/1, 4/0011, 5/registerM, 1/--, 2/10, 3/101, 1/1, 1/splat32, 5/lower5, 5/registerP, 1/0 }


\paragraph{Syntax} \texttt{addx16ho} \emph{registerM} = \emph{registerP}, \emph{upper27\_\discretionary{}{}{}lower5}\emph{splat32}

\paragraph{Description} The \emph{upper27\_\discretionary{}{}{}lower5} and the \emph{registerP} are interpreted as eight half words packed into 128 bits. The \emph{registerP} half words are shifted left by eight bits and added to the \emph{upper27\_\discretionary{}{}{}lower5} half words. The eight resulting half words are packed and stored into the \emph{registerM}.


\paragraph{Modifier(s)}
\noindent\begin{itemize}
\item splat32 (cf. splat32 in section \ref{par:modifier-splat32} on page \pageref{par:modifier-splat32})
\end{itemize}

\paragraph{Execution} {Reservation table \textsf{ALU\_LITE.X} (Section~\ref{sec:reservation-ALU-LITE.X})}
~
{\begin{lstlisting}
stage ID:
  new splat32 = _ZX_1(splat32);
stage RR:
  new argument3 = splat32 ? (_WX_32(upper27_lower5) << 32) | _ZX_32(_WX_32(upper27_lower5)) : _SX_32(_WX_32(upper27_lower5));
  new argument2 = PGR[registerP];
  for (I = 0; I < 8; I += 1) {
    result1.16[I] = _SX_16(argument3.16[I]) + (_SX_16(argument2.16[I]) << 4)
  };
stage E1:
  PGR[registerM] = result1;
\end{lstlisting}}
\newpage
\subsubsection[{\small ADDX16W: Add Word to Word Times 16}]{ADDX16W\hfill Add Word to Word Times 16}
\label{instr:ADDX16W}\index{addx16w}
 
\paragraph{Encoding} Format \textsf{ALU\_WWRR1} (Section~\ref{sec:format-ALU-WWRR1}), \verb|exucode2 = 0011|\\

\bitfields{ 1/P, 2/11, 1/0, 4/0011, 6/registerW, 2/11, 3/101, 1/signextw, 6/registerY, 6/registerZ }


\paragraph{Syntax} \texttt{addx16w}\emph{signextw} \emph{registerW} = \emph{registerZ}, \emph{registerY}

\paragraph{Description} The \emph{registerZ} is shifted left by four bits and added to the \emph{registerY}. The result with the 32 upper bits cleared is stored into the \emph{registerW}.


\paragraph{Modifier(s)}
\noindent\begin{itemize}
\item signextw (cf. signextw in section \ref{par:modifier-signextw} on page \pageref{par:modifier-signextw})
\end{itemize}

\paragraph{Execution} {Reservation table \textsf{ALU\_TINY} (Section~\ref{sec:reservation-ALU-TINY})}
~
{\begin{lstlisting}
stage ID:
  new signextw = _ZX_1(signextw);
stage RR:
  new argument3 = GPR[registerY];
  new argument2 = GPR[registerZ];
  new result1 = argument3 + (argument2 << 4);
stage E1:
  GPR[registerW] = signextw ? _SX_32(result1) : _ZX_32(result1);
\end{lstlisting}}
\newpage

\paragraph{Encoding} Format \textsf{ALU\_WWRR1.W} (Section~\ref{sec:format-ALU-WWRR1.W}), \verb|exucode2 = 0011|\\

\bitfields{ 1/P, 2/00, 2/-- --, 27/upper27, 1/1, 2/11, 1/0, 4/0011, 6/registerW, 2/11, 3/101, 1/signextw, 1/0, 5/lower5, 6/registerZ }


\paragraph{Syntax} \texttt{addx16w}\emph{signextw} \emph{registerW} = \emph{registerZ}, \emph{upper27\_\discretionary{}{}{}lower5}

\paragraph{Description} The \emph{registerZ} is shifted left by four bits and added to the \emph{upper27\_\discretionary{}{}{}lower5}. The result with the 32 upper bits cleared is stored into the \emph{registerW}.


\paragraph{Modifier(s)}
\noindent\begin{itemize}
\item signextw (cf. signextw in section \ref{par:modifier-signextw} on page \pageref{par:modifier-signextw})
\end{itemize}

\paragraph{Execution} {Reservation table \textsf{ALU\_TINY.X} (Section~\ref{sec:reservation-ALU-TINY.X})}
~
{\begin{lstlisting}
stage ID:
  new signextw = _ZX_1(signextw);
stage RR:
  new argument3 = _WX_32(upper27_lower5);
  new argument2 = GPR[registerZ];
  new result1 = argument3 + (argument2 << 4);
stage E1:
  GPR[registerW] = signextw ? _SX_32(result1) : _ZX_32(result1);
\end{lstlisting}}
\newpage
\subsubsection[{\small ADDX16WQ: Add to Times 16 Word Quadruple}]{ADDX16WQ\hfill Add to Times 16 Word Quadruple}
\label{instr:ADDX16WQ}\index{addx16wq}
 
\paragraph{Encoding} Format \textsf{ALU\_WQWRR1} (Section~\ref{sec:format-ALU-WQWRR1}), \verb|exucode2 = 0011|\\

\bitfields{ 1/P, 2/11, 1/1, 4/0011, 5/registerM, 1/--, 2/10, 3/101, 1/0, 5/registerO, 1/--, 5/registerP, 1/1 }


\paragraph{Syntax} \texttt{addx16wq} \emph{registerM} = \emph{registerP}, \emph{registerO}

\paragraph{Description} The \emph{registerO} and the \emph{registerP} are interpreted as four words packed into 128 bits. The \emph{registerP} words are shifted left by four bits and added to the \emph{registerO} words. The four resulting words are packed and stored into the \emph{registerM}.


\paragraph{Execution} {Reservation table \textsf{ALU\_LITE} (Section~\ref{sec:reservation-ALU-LITE})}
~
{\begin{lstlisting}
stage RR:
  new argument3 = PGR[registerO];
  new argument2 = PGR[registerP];
  for (I = 0; I < 4; I += 1) {
    result1.32[I] = _SX_32(argument3.32[I]) + (_SX_32(argument2.32[I]) << 4)
  };
stage E1:
  PGR[registerM] = result1;
\end{lstlisting}}
\newpage

\paragraph{Encoding} Format \textsf{ALU\_WQWRR1.M} (Section~\ref{sec:format-ALU-WQWRR1.M}), \verb|exucode2 = 0011|\\

\bitfields{ 1/P, 2/00, 2/-- --, 27/upper27, 1/1, 2/11, 1/1, 4/0011, 5/registerM, 1/--, 2/10, 3/101, 1/0, 1/splat32, 5/lower5, 5/registerP, 1/1 }


\paragraph{Syntax} \texttt{addx16wq} \emph{registerM} = \emph{registerP}, \emph{upper27\_\discretionary{}{}{}lower5}\emph{splat32}

\paragraph{Description} The \emph{upper27\_\discretionary{}{}{}lower5} and the \emph{registerP} are interpreted as four words packed into 128 bits. The \emph{registerP} words are shifted left by four bits and added to the \emph{upper27\_\discretionary{}{}{}lower5} words. The four resulting words are packed and stored into the \emph{registerM}.


\paragraph{Modifier(s)}
\noindent\begin{itemize}
\item splat32 (cf. splat32 in section \ref{par:modifier-splat32} on page \pageref{par:modifier-splat32})
\end{itemize}

\paragraph{Execution} {Reservation table \textsf{ALU\_LITE.X} (Section~\ref{sec:reservation-ALU-LITE.X})}
~
{\begin{lstlisting}
stage ID:
  new splat32 = _ZX_1(splat32);
stage RR:
  new argument3 = splat32 ? (_WX_32(upper27_lower5) << 32) | _ZX_32(_WX_32(upper27_lower5)) : _SX_32(_WX_32(upper27_lower5));
  new argument2 = PGR[registerP];
  for (I = 0; I < 4; I += 1) {
    result1.32[I] = _SX_32(argument3.32[I]) + (_SX_32(argument2.32[I]) << 4)
  };
stage E1:
  PGR[registerM] = result1;
\end{lstlisting}}
\newpage
\subsubsection[{\small ADDX2BX: Add to Times 2 Byte Hexadecuple}]{ADDX2BX\hfill Add to Times 2 Byte Hexadecuple}
\label{instr:ADDX2BX}\index{addx2bx}
 
\paragraph{Encoding} Format \textsf{ALU\_BXWRR1} (Section~\ref{sec:format-ALU-BXWRR1}), \verb|exucode2 = 0000|\\

\bitfields{ 1/P, 2/11, 1/1, 4/0000, 5/registerM, 1/--, 2/10, 3/101, 1/1, 5/registerO, 1/--, 5/registerP, 1/1 }


\paragraph{Syntax} \texttt{addx2bx} \emph{registerM} = \emph{registerP}, \emph{registerO}

\paragraph{Description} The \emph{registerO} and the \emph{registerP} are interpreted as sixteen bytes packed into 128 bits. The \emph{registerP} bytes are shifted left by one bit and added to the \emph{registerO} bytes. The sixteen resulting bytes are packed and stored into the \emph{registerM}.


\paragraph{Execution} {Reservation table \textsf{ALU\_LITE} (Section~\ref{sec:reservation-ALU-LITE})}
~
{\begin{lstlisting}
stage RR:
  new argument3 = PGR[registerO];
  new argument2 = PGR[registerP];
  for (I = 0; I < 16; I += 1) {
    result1.8[I] = _SX_8(argument3.8[I]) + (_SX_8(argument2.8[I]) << 1)
  };
stage E1:
  PGR[registerM] = result1;
\end{lstlisting}}
\newpage

\paragraph{Encoding} Format \textsf{ALU\_BXWRR1.M} (Section~\ref{sec:format-ALU-BXWRR1.M}), \verb|exucode2 = 0000|\\

\bitfields{ 1/P, 2/00, 2/-- --, 27/upper27, 1/1, 2/11, 1/1, 4/0000, 5/registerM, 1/--, 2/10, 3/101, 1/1, 1/splat32, 5/lower5, 5/registerP, 1/1 }


\paragraph{Syntax} \texttt{addx2bx} \emph{registerM} = \emph{registerP}, \emph{upper27\_\discretionary{}{}{}lower5}\emph{splat32}

\paragraph{Description} The \emph{upper27\_\discretionary{}{}{}lower5} and the \emph{registerP} are interpreted as sixteen bytes packed into 128 bits. The \emph{registerP} bytes are shifted left by one bit and added to the \emph{upper27\_\discretionary{}{}{}lower5} bytes. The sixteen resulting bytes are packed and stored into the \emph{registerM}.


\paragraph{Modifier(s)}
\noindent\begin{itemize}
\item splat32 (cf. splat32 in section \ref{par:modifier-splat32} on page \pageref{par:modifier-splat32})
\end{itemize}

\paragraph{Execution} {Reservation table \textsf{ALU\_LITE.X} (Section~\ref{sec:reservation-ALU-LITE.X})}
~
{\begin{lstlisting}
stage ID:
  new splat32 = _ZX_1(splat32);
stage RR:
  new argument3 = splat32 ? (_WX_32(upper27_lower5) << 32) | _ZX_32(_WX_32(upper27_lower5)) : _SX_32(_WX_32(upper27_lower5));
  new argument2 = PGR[registerP];
  for (I = 0; I < 16; I += 1) {
    result1.8[I] = _SX_8(argument3.8[I]) + (_SX_8(argument2.8[I]) << 1)
  };
stage E1:
  PGR[registerM] = result1;
\end{lstlisting}}
\newpage
\subsubsection[{\small ADDX2D: Add to Times 2 Double Word}]{ADDX2D\hfill Add to Times 2 Double Word}
\label{instr:ADDX2D}\index{addx2d}
 
\paragraph{Encoding} Format \textsf{ALU\_DWRR1} (Section~\ref{sec:format-ALU-DWRR1}), \verb|exucode2 = 0000|\\

\bitfields{ 1/P, 2/11, 1/0, 4/0000, 6/registerW, 2/10, 3/101, 1/0, 6/registerY, 6/registerZ }


\paragraph{Syntax} \texttt{addx2d} \emph{registerW} = \emph{registerZ}, \emph{registerY}

\paragraph{Description} The \emph{registerZ} is shifted left by one bit and added to the \emph{registerY}. The result is stored into the \emph{registerW}.


\paragraph{Execution} {Reservation table \textsf{ALU\_TINY} (Section~\ref{sec:reservation-ALU-TINY})}
~
{\begin{lstlisting}
stage RR:
  new argument3 = GPR[registerY];
  new argument2 = GPR[registerZ];
  new result1 = argument3 + (argument2 << 1);
stage E1:
  GPR[registerW] = result1;
\end{lstlisting}}
\newpage

\paragraph{Encoding} Format \textsf{ALU\_DWRR1.M} (Section~\ref{sec:format-ALU-DWRR1.M}), \verb|exucode2 = 0000|\\

\bitfields{ 1/P, 2/00, 2/-- --, 27/upper27, 1/1, 2/11, 1/0, 4/0000, 6/registerW, 2/10, 3/101, 1/0, 1/splat32, 5/lower5, 6/registerZ }


\paragraph{Syntax} \texttt{addx2d} \emph{registerW} = \emph{registerZ}, \emph{upper27\_\discretionary{}{}{}lower5}\emph{splat32}

\paragraph{Description} The \emph{registerZ} is shifted left by one bit and added to the \emph{upper27\_\discretionary{}{}{}lower5}. The result is stored into the \emph{registerW}.


\paragraph{Modifier(s)}
\noindent\begin{itemize}
\item splat32 (cf. splat32 in section \ref{par:modifier-splat32} on page \pageref{par:modifier-splat32})
\end{itemize}

\paragraph{Execution} {Reservation table \textsf{ALU\_TINY.X} (Section~\ref{sec:reservation-ALU-TINY.X})}
~
{\begin{lstlisting}
stage ID:
  new splat32 = _ZX_1(splat32);
stage RR:
  new argument3 = splat32 ? (_WX_32(upper27_lower5) << 32) | _ZX_32(_WX_32(upper27_lower5)) : _SX_32(_WX_32(upper27_lower5));
  new argument2 = GPR[registerZ];
  new result1 = argument3 + (argument2 << 1);
stage E1:
  GPR[registerW] = result1;
\end{lstlisting}}
\newpage
\subsubsection[{\small ADDX2DP: Add to Times 2 Double Word Pair}]{ADDX2DP\hfill Add to Times 2 Double Word Pair}
\label{instr:ADDX2DP}\index{addx2dp}
 
\paragraph{Encoding} Format \textsf{ALU\_DPWRR1} (Section~\ref{sec:format-ALU-DPWRR1}), \verb|exucode2 = 0000|\\

\bitfields{ 1/P, 2/11, 1/1, 4/0000, 5/registerM, 1/--, 2/10, 3/101, 1/0, 5/registerO, 1/--, 5/registerP, 1/0 }


\paragraph{Syntax} \texttt{addx2dp} \emph{registerM} = \emph{registerP}, \emph{registerO}

\paragraph{Description} The \emph{registerO} and the \emph{registerP} are interpreted as two double words packed into 128 bits. The \emph{registerP} double words are shifted left by one bit and added to the \emph{registerO} double words. The two resulting double words are packed and stored into the \emph{registerM}.


\paragraph{Execution} {Reservation table \textsf{ALU\_LITE} (Section~\ref{sec:reservation-ALU-LITE})}
~
{\begin{lstlisting}
stage RR:
  new argument3 = PGR[registerO];
  new argument2 = PGR[registerP];
  for (I = 0; I < 2; I += 1) {
    result1.64[I] = _SX_64(argument3.64[I]) + (_SX_64(argument2.64[I]) << 1)
  };
stage E1:
  PGR[registerM] = result1;
\end{lstlisting}}
\newpage

\paragraph{Encoding} Format \textsf{ALU\_DPWRR1.M} (Section~\ref{sec:format-ALU-DPWRR1.M}), \verb|exucode2 = 0000|\\

\bitfields{ 1/P, 2/00, 2/-- --, 27/upper27, 1/1, 2/11, 1/1, 4/0000, 5/registerM, 1/--, 2/10, 3/101, 1/0, 1/splat32, 5/lower5, 5/registerP, 1/0 }


\paragraph{Syntax} \texttt{addx2dp} \emph{registerM} = \emph{registerP}, \emph{upper27\_\discretionary{}{}{}lower5}\emph{splat32}

\paragraph{Description} The \emph{upper27\_\discretionary{}{}{}lower5} and the \emph{registerP} are interpreted as two double words packed into 128 bits. The \emph{registerP} double words are shifted left by one bit and added to the \emph{upper27\_\discretionary{}{}{}lower5} double words. The two resulting double words are packed and stored into the \emph{registerM}.


\paragraph{Modifier(s)}
\noindent\begin{itemize}
\item splat32 (cf. splat32 in section \ref{par:modifier-splat32} on page \pageref{par:modifier-splat32})
\end{itemize}

\paragraph{Execution} {Reservation table \textsf{ALU\_LITE.X} (Section~\ref{sec:reservation-ALU-LITE.X})}
~
{\begin{lstlisting}
stage ID:
  new splat32 = _ZX_1(splat32);
stage RR:
  new argument3 = splat32 ? (_WX_32(upper27_lower5) << 32) | _ZX_32(_WX_32(upper27_lower5)) : _SX_32(_WX_32(upper27_lower5));
  new argument2 = PGR[registerP];
  for (I = 0; I < 2; I += 1) {
    result1.64[I] = _SX_64(argument3.64[I]) + (_SX_64(argument2.64[I]) << 1)
  };
stage E1:
  PGR[registerM] = result1;
\end{lstlisting}}
\newpage
\subsubsection[{\small ADDX2HO: Add to Times 2 Half Word Octuple}]{ADDX2HO\hfill Add to Times 2 Half Word Octuple}
\label{instr:ADDX2HO}\index{addx2ho}
 
\paragraph{Encoding} Format \textsf{ALU\_HOWRR1} (Section~\ref{sec:format-ALU-HOWRR1}), \verb|exucode2 = 0000|\\

\bitfields{ 1/P, 2/11, 1/1, 4/0000, 5/registerM, 1/--, 2/10, 3/101, 1/1, 5/registerO, 1/--, 5/registerP, 1/0 }


\paragraph{Syntax} \texttt{addx2ho} \emph{registerM} = \emph{registerP}, \emph{registerO}

\paragraph{Description} The \emph{registerO} and the \emph{registerP} are interpreted as eight half words packed into 128 bits. The \emph{registerP} half words are shifted left by one bit and added to the \emph{registerO} half words. The eight resulting half words are packed and stored into the \emph{registerM}.


\paragraph{Execution} {Reservation table \textsf{ALU\_LITE} (Section~\ref{sec:reservation-ALU-LITE})}
~
{\begin{lstlisting}
stage RR:
  new argument3 = PGR[registerO];
  new argument2 = PGR[registerP];
  for (I = 0; I < 8; I += 1) {
    result1.16[I] = _SX_16(argument3.16[I]) + (_SX_16(argument2.16[I]) << 1)
  };
stage E1:
  PGR[registerM] = result1;
\end{lstlisting}}
\newpage

\paragraph{Encoding} Format \textsf{ALU\_HOWRR1.M} (Section~\ref{sec:format-ALU-HOWRR1.M}), \verb|exucode2 = 0000|\\

\bitfields{ 1/P, 2/00, 2/-- --, 27/upper27, 1/1, 2/11, 1/1, 4/0000, 5/registerM, 1/--, 2/10, 3/101, 1/1, 1/splat32, 5/lower5, 5/registerP, 1/0 }


\paragraph{Syntax} \texttt{addx2ho} \emph{registerM} = \emph{registerP}, \emph{upper27\_\discretionary{}{}{}lower5}\emph{splat32}

\paragraph{Description} The \emph{upper27\_\discretionary{}{}{}lower5} and the \emph{registerP} are interpreted as eight half words packed into 128 bits. The \emph{registerP} half words are shifted left by one bit and added to the \emph{upper27\_\discretionary{}{}{}lower5} half words. The eight resulting half words are packed and stored into the \emph{registerM}.


\paragraph{Modifier(s)}
\noindent\begin{itemize}
\item splat32 (cf. splat32 in section \ref{par:modifier-splat32} on page \pageref{par:modifier-splat32})
\end{itemize}

\paragraph{Execution} {Reservation table \textsf{ALU\_LITE.X} (Section~\ref{sec:reservation-ALU-LITE.X})}
~
{\begin{lstlisting}
stage ID:
  new splat32 = _ZX_1(splat32);
stage RR:
  new argument3 = splat32 ? (_WX_32(upper27_lower5) << 32) | _ZX_32(_WX_32(upper27_lower5)) : _SX_32(_WX_32(upper27_lower5));
  new argument2 = PGR[registerP];
  for (I = 0; I < 8; I += 1) {
    result1.16[I] = _SX_16(argument3.16[I]) + (_SX_16(argument2.16[I]) << 1)
  };
stage E1:
  PGR[registerM] = result1;
\end{lstlisting}}
\newpage
\subsubsection[{\small ADDX2W: Add to Times 2 Word}]{ADDX2W\hfill Add to Times 2 Word}
\label{instr:ADDX2W}\index{addx2w}
 
\paragraph{Encoding} Format \textsf{ALU\_WWRR1} (Section~\ref{sec:format-ALU-WWRR1}), \verb|exucode2 = 0000|\\

\bitfields{ 1/P, 2/11, 1/0, 4/0000, 6/registerW, 2/11, 3/101, 1/signextw, 6/registerY, 6/registerZ }


\paragraph{Syntax} \texttt{addx2w}\emph{signextw} \emph{registerW} = \emph{registerZ}, \emph{registerY}

\paragraph{Description} The \emph{registerZ} is shifted left by one bit and added to the \emph{registerY}. The result with the 32 upper bits cleared is stored into the \emph{registerW}.


\paragraph{Modifier(s)}
\noindent\begin{itemize}
\item signextw (cf. signextw in section \ref{par:modifier-signextw} on page \pageref{par:modifier-signextw})
\end{itemize}

\paragraph{Execution} {Reservation table \textsf{ALU\_TINY} (Section~\ref{sec:reservation-ALU-TINY})}
~
{\begin{lstlisting}
stage ID:
  new signextw = _ZX_1(signextw);
stage RR:
  new argument3 = GPR[registerY];
  new argument2 = GPR[registerZ];
  new result1 = argument3 + (argument2 << 1);
stage E1:
  GPR[registerW] = signextw ? _SX_32(result1) : _ZX_32(result1);
\end{lstlisting}}
\newpage

\paragraph{Encoding} Format \textsf{ALU\_WWRR1.W} (Section~\ref{sec:format-ALU-WWRR1.W}), \verb|exucode2 = 0000|\\

\bitfields{ 1/P, 2/00, 2/-- --, 27/upper27, 1/1, 2/11, 1/0, 4/0000, 6/registerW, 2/11, 3/101, 1/signextw, 1/0, 5/lower5, 6/registerZ }


\paragraph{Syntax} \texttt{addx2w}\emph{signextw} \emph{registerW} = \emph{registerZ}, \emph{upper27\_\discretionary{}{}{}lower5}

\paragraph{Description} The \emph{registerZ} is shifted left by one bit and added to the \emph{upper27\_\discretionary{}{}{}lower5}. The result with the 32 upper bits cleared is stored into the \emph{registerW}.


\paragraph{Modifier(s)}
\noindent\begin{itemize}
\item signextw (cf. signextw in section \ref{par:modifier-signextw} on page \pageref{par:modifier-signextw})
\end{itemize}

\paragraph{Execution} {Reservation table \textsf{ALU\_TINY.X} (Section~\ref{sec:reservation-ALU-TINY.X})}
~
{\begin{lstlisting}
stage ID:
  new signextw = _ZX_1(signextw);
stage RR:
  new argument3 = _WX_32(upper27_lower5);
  new argument2 = GPR[registerZ];
  new result1 = argument3 + (argument2 << 1);
stage E1:
  GPR[registerW] = signextw ? _SX_32(result1) : _ZX_32(result1);
\end{lstlisting}}
\newpage
\subsubsection[{\small ADDX2WQ: Add to Times 2 Word Quadruple}]{ADDX2WQ\hfill Add to Times 2 Word Quadruple}
\label{instr:ADDX2WQ}\index{addx2wq}
 
\paragraph{Encoding} Format \textsf{ALU\_WQWRR1} (Section~\ref{sec:format-ALU-WQWRR1}), \verb|exucode2 = 0000|\\

\bitfields{ 1/P, 2/11, 1/1, 4/0000, 5/registerM, 1/--, 2/10, 3/101, 1/0, 5/registerO, 1/--, 5/registerP, 1/1 }


\paragraph{Syntax} \texttt{addx2wq} \emph{registerM} = \emph{registerP}, \emph{registerO}

\paragraph{Description} The \emph{registerO} and the \emph{registerP} are interpreted as four words packed into 128 bits. The \emph{registerP} words are shifted left by one bit and added to the \emph{registerO} words. The four resulting words are packed and stored into the \emph{registerM}.


\paragraph{Execution} {Reservation table \textsf{ALU\_LITE} (Section~\ref{sec:reservation-ALU-LITE})}
~
{\begin{lstlisting}
stage RR:
  new argument3 = PGR[registerO];
  new argument2 = PGR[registerP];
  for (I = 0; I < 4; I += 1) {
    result1.32[I] = _SX_32(argument3.32[I]) + (_SX_32(argument2.32[I]) << 1)
  };
stage E1:
  PGR[registerM] = result1;
\end{lstlisting}}
\newpage

\paragraph{Encoding} Format \textsf{ALU\_WQWRR1.M} (Section~\ref{sec:format-ALU-WQWRR1.M}), \verb|exucode2 = 0000|\\

\bitfields{ 1/P, 2/00, 2/-- --, 27/upper27, 1/1, 2/11, 1/1, 4/0000, 5/registerM, 1/--, 2/10, 3/101, 1/0, 1/splat32, 5/lower5, 5/registerP, 1/1 }


\paragraph{Syntax} \texttt{addx2wq} \emph{registerM} = \emph{registerP}, \emph{upper27\_\discretionary{}{}{}lower5}\emph{splat32}

\paragraph{Description} The \emph{upper27\_\discretionary{}{}{}lower5} and the \emph{registerP} are interpreted as four words packed into 128 bits. The \emph{registerP} words are shifted left by one bit and added to the \emph{upper27\_\discretionary{}{}{}lower5} words. The four resulting words are packed and stored into the \emph{registerM}.


\paragraph{Modifier(s)}
\noindent\begin{itemize}
\item splat32 (cf. splat32 in section \ref{par:modifier-splat32} on page \pageref{par:modifier-splat32})
\end{itemize}

\paragraph{Execution} {Reservation table \textsf{ALU\_LITE.X} (Section~\ref{sec:reservation-ALU-LITE.X})}
~
{\begin{lstlisting}
stage ID:
  new splat32 = _ZX_1(splat32);
stage RR:
  new argument3 = splat32 ? (_WX_32(upper27_lower5) << 32) | _ZX_32(_WX_32(upper27_lower5)) : _SX_32(_WX_32(upper27_lower5));
  new argument2 = PGR[registerP];
  for (I = 0; I < 4; I += 1) {
    result1.32[I] = _SX_32(argument3.32[I]) + (_SX_32(argument2.32[I]) << 1)
  };
stage E1:
  PGR[registerM] = result1;
\end{lstlisting}}
\newpage
\subsubsection[{\small ADDX32D: Add to Times 32 Double Word}]{ADDX32D\hfill Add to Times 32 Double Word}
\label{instr:ADDX32D}\index{addx32d}
 
\paragraph{Encoding} Format \textsf{ALU\_DWRR1} (Section~\ref{sec:format-ALU-DWRR1}), \verb|exucode2 = 0100|\\

\bitfields{ 1/P, 2/11, 1/0, 4/0100, 6/registerW, 2/10, 3/101, 1/0, 6/registerY, 6/registerZ }


\paragraph{Syntax} \texttt{addx32d} \emph{registerW} = \emph{registerZ}, \emph{registerY}

\paragraph{Description} The \emph{registerZ} is shifted left by five bits and added to the \emph{registerY}. The result is stored into the \emph{registerW}.


\paragraph{Execution} {Reservation table \textsf{ALU\_TINY} (Section~\ref{sec:reservation-ALU-TINY})}
~
{\begin{lstlisting}
stage RR:
  new argument3 = GPR[registerY];
  new argument2 = GPR[registerZ];
  new result1 = argument3 + (argument2 << 5);
stage E1:
  GPR[registerW] = result1;
\end{lstlisting}}
\newpage

\paragraph{Encoding} Format \textsf{ALU\_DWRR1.M} (Section~\ref{sec:format-ALU-DWRR1.M}), \verb|exucode2 = 0100|\\

\bitfields{ 1/P, 2/00, 2/-- --, 27/upper27, 1/1, 2/11, 1/0, 4/0100, 6/registerW, 2/10, 3/101, 1/0, 1/splat32, 5/lower5, 6/registerZ }


\paragraph{Syntax} \texttt{addx32d} \emph{registerW} = \emph{registerZ}, \emph{upper27\_\discretionary{}{}{}lower5}\emph{splat32}

\paragraph{Description} The \emph{registerZ} is shifted left by five bits and added to the \emph{upper27\_\discretionary{}{}{}lower5}. The result is stored into the \emph{registerW}.


\paragraph{Modifier(s)}
\noindent\begin{itemize}
\item splat32 (cf. splat32 in section \ref{par:modifier-splat32} on page \pageref{par:modifier-splat32})
\end{itemize}

\paragraph{Execution} {Reservation table \textsf{ALU\_TINY.X} (Section~\ref{sec:reservation-ALU-TINY.X})}
~
{\begin{lstlisting}
stage ID:
  new splat32 = _ZX_1(splat32);
stage RR:
  new argument3 = splat32 ? (_WX_32(upper27_lower5) << 32) | _ZX_32(_WX_32(upper27_lower5)) : _SX_32(_WX_32(upper27_lower5));
  new argument2 = GPR[registerZ];
  new result1 = argument3 + (argument2 << 5);
stage E1:
  GPR[registerW] = result1;
\end{lstlisting}}
\newpage
\subsubsection[{\small ADDX32W: Add to Times 32 Word}]{ADDX32W\hfill Add to Times 32 Word}
\label{instr:ADDX32W}\index{addx32w}
 
\paragraph{Encoding} Format \textsf{ALU\_WWRR1} (Section~\ref{sec:format-ALU-WWRR1}), \verb|exucode2 = 0100|\\

\bitfields{ 1/P, 2/11, 1/0, 4/0100, 6/registerW, 2/11, 3/101, 1/signextw, 6/registerY, 6/registerZ }


\paragraph{Syntax} \texttt{addx32w}\emph{signextw} \emph{registerW} = \emph{registerZ}, \emph{registerY}

\paragraph{Description} The \emph{registerZ} is shifted left by five bits and added to the \emph{registerY}. The result with the 32 upper bits cleared is stored into the \emph{registerW}.


\paragraph{Modifier(s)}
\noindent\begin{itemize}
\item signextw (cf. signextw in section \ref{par:modifier-signextw} on page \pageref{par:modifier-signextw})
\end{itemize}

\paragraph{Execution} {Reservation table \textsf{ALU\_TINY} (Section~\ref{sec:reservation-ALU-TINY})}
~
{\begin{lstlisting}
stage ID:
  new signextw = _ZX_1(signextw);
stage RR:
  new argument3 = GPR[registerY];
  new argument2 = GPR[registerZ];
  new result1 = argument3 + (argument2 << 5);
stage E1:
  GPR[registerW] = signextw ? _SX_32(result1) : _ZX_32(result1);
\end{lstlisting}}
\newpage

\paragraph{Encoding} Format \textsf{ALU\_WWRR1.W} (Section~\ref{sec:format-ALU-WWRR1.W}), \verb|exucode2 = 0100|\\

\bitfields{ 1/P, 2/00, 2/-- --, 27/upper27, 1/1, 2/11, 1/0, 4/0100, 6/registerW, 2/11, 3/101, 1/signextw, 1/0, 5/lower5, 6/registerZ }


\paragraph{Syntax} \texttt{addx32w}\emph{signextw} \emph{registerW} = \emph{registerZ}, \emph{upper27\_\discretionary{}{}{}lower5}

\paragraph{Description} The \emph{registerZ} is shifted left by five bits and added to the \emph{upper27\_\discretionary{}{}{}lower5}. The result with the 32 upper bits cleared is stored into the \emph{registerW}.


\paragraph{Modifier(s)}
\noindent\begin{itemize}
\item signextw (cf. signextw in section \ref{par:modifier-signextw} on page \pageref{par:modifier-signextw})
\end{itemize}

\paragraph{Execution} {Reservation table \textsf{ALU\_TINY.X} (Section~\ref{sec:reservation-ALU-TINY.X})}
~
{\begin{lstlisting}
stage ID:
  new signextw = _ZX_1(signextw);
stage RR:
  new argument3 = _WX_32(upper27_lower5);
  new argument2 = GPR[registerZ];
  new result1 = argument3 + (argument2 << 5);
stage E1:
  GPR[registerW] = signextw ? _SX_32(result1) : _ZX_32(result1);
\end{lstlisting}}
\newpage
\subsubsection[{\small ADDX4BX: Add to Times 4 Byte Hexadecuple}]{ADDX4BX\hfill Add to Times 4 Byte Hexadecuple}
\label{instr:ADDX4BX}\index{addx4bx}
 
\paragraph{Encoding} Format \textsf{ALU\_BXWRR1} (Section~\ref{sec:format-ALU-BXWRR1}), \verb|exucode2 = 0001|\\

\bitfields{ 1/P, 2/11, 1/1, 4/0001, 5/registerM, 1/--, 2/10, 3/101, 1/1, 5/registerO, 1/--, 5/registerP, 1/1 }


\paragraph{Syntax} \texttt{addx4bx} \emph{registerM} = \emph{registerP}, \emph{registerO}

\paragraph{Description} The \emph{registerO} and the \emph{registerP} are interpreted as sixteen bytes packed into 128 bits. The \emph{registerP} bytes are shifted left by two bits and added to the \emph{registerO} bytes. The sixteen resulting bytes are packed and stored into the \emph{registerM}.


\paragraph{Execution} {Reservation table \textsf{ALU\_LITE} (Section~\ref{sec:reservation-ALU-LITE})}
~
{\begin{lstlisting}
stage RR:
  new argument3 = PGR[registerO];
  new argument2 = PGR[registerP];
  for (I = 0; I < 16; I += 1) {
    result1.8[I] = _SX_8(argument3.8[I]) + (_SX_8(argument2.8[I]) << 2)
  };
stage E1:
  PGR[registerM] = result1;
\end{lstlisting}}
\newpage

\paragraph{Encoding} Format \textsf{ALU\_BXWRR1.M} (Section~\ref{sec:format-ALU-BXWRR1.M}), \verb|exucode2 = 0001|\\

\bitfields{ 1/P, 2/00, 2/-- --, 27/upper27, 1/1, 2/11, 1/1, 4/0001, 5/registerM, 1/--, 2/10, 3/101, 1/1, 1/splat32, 5/lower5, 5/registerP, 1/1 }


\paragraph{Syntax} \texttt{addx4bx} \emph{registerM} = \emph{registerP}, \emph{upper27\_\discretionary{}{}{}lower5}\emph{splat32}

\paragraph{Description} The \emph{upper27\_\discretionary{}{}{}lower5} and the \emph{registerP} are interpreted as sixteen bytes packed into 128 bits. The \emph{registerP} bytes are shifted left by two bits and added to the \emph{upper27\_\discretionary{}{}{}lower5} bytes. The sixteen resulting bytes are packed and stored into the \emph{registerM}.


\paragraph{Modifier(s)}
\noindent\begin{itemize}
\item splat32 (cf. splat32 in section \ref{par:modifier-splat32} on page \pageref{par:modifier-splat32})
\end{itemize}

\paragraph{Execution} {Reservation table \textsf{ALU\_LITE.X} (Section~\ref{sec:reservation-ALU-LITE.X})}
~
{\begin{lstlisting}
stage ID:
  new splat32 = _ZX_1(splat32);
stage RR:
  new argument3 = splat32 ? (_WX_32(upper27_lower5) << 32) | _ZX_32(_WX_32(upper27_lower5)) : _SX_32(_WX_32(upper27_lower5));
  new argument2 = PGR[registerP];
  for (I = 0; I < 16; I += 1) {
    result1.8[I] = _SX_8(argument3.8[I]) + (_SX_8(argument2.8[I]) << 2)
  };
stage E1:
  PGR[registerM] = result1;
\end{lstlisting}}
\newpage
\subsubsection[{\small ADDX4D: Add to Times 4 Double Word}]{ADDX4D\hfill Add to Times 4 Double Word}
\label{instr:ADDX4D}\index{addx4d}
 
\paragraph{Encoding} Format \textsf{ALU\_DWRR1} (Section~\ref{sec:format-ALU-DWRR1}), \verb|exucode2 = 0001|\\

\bitfields{ 1/P, 2/11, 1/0, 4/0001, 6/registerW, 2/10, 3/101, 1/0, 6/registerY, 6/registerZ }


\paragraph{Syntax} \texttt{addx4d} \emph{registerW} = \emph{registerZ}, \emph{registerY}

\paragraph{Description} The \emph{registerZ} is shifted left by two bits and added to the \emph{registerY}. The result is stored into the \emph{registerW}.


\paragraph{Execution} {Reservation table \textsf{ALU\_TINY} (Section~\ref{sec:reservation-ALU-TINY})}
~
{\begin{lstlisting}
stage RR:
  new argument3 = GPR[registerY];
  new argument2 = GPR[registerZ];
  new result1 = argument3 + (argument2 << 2);
stage E1:
  GPR[registerW] = result1;
\end{lstlisting}}
\newpage

\paragraph{Encoding} Format \textsf{ALU\_DWRR1.M} (Section~\ref{sec:format-ALU-DWRR1.M}), \verb|exucode2 = 0001|\\

\bitfields{ 1/P, 2/00, 2/-- --, 27/upper27, 1/1, 2/11, 1/0, 4/0001, 6/registerW, 2/10, 3/101, 1/0, 1/splat32, 5/lower5, 6/registerZ }


\paragraph{Syntax} \texttt{addx4d} \emph{registerW} = \emph{registerZ}, \emph{upper27\_\discretionary{}{}{}lower5}\emph{splat32}

\paragraph{Description} The \emph{registerZ} is shifted left by two bits and added to the \emph{upper27\_\discretionary{}{}{}lower5}. The result is stored into the \emph{registerW}.


\paragraph{Modifier(s)}
\noindent\begin{itemize}
\item splat32 (cf. splat32 in section \ref{par:modifier-splat32} on page \pageref{par:modifier-splat32})
\end{itemize}

\paragraph{Execution} {Reservation table \textsf{ALU\_TINY.X} (Section~\ref{sec:reservation-ALU-TINY.X})}
~
{\begin{lstlisting}
stage ID:
  new splat32 = _ZX_1(splat32);
stage RR:
  new argument3 = splat32 ? (_WX_32(upper27_lower5) << 32) | _ZX_32(_WX_32(upper27_lower5)) : _SX_32(_WX_32(upper27_lower5));
  new argument2 = GPR[registerZ];
  new result1 = argument3 + (argument2 << 2);
stage E1:
  GPR[registerW] = result1;
\end{lstlisting}}
\newpage
\subsubsection[{\small ADDX4DP: Add to Times 4 Double Word Pair}]{ADDX4DP\hfill Add to Times 4 Double Word Pair}
\label{instr:ADDX4DP}\index{addx4dp}
 
\paragraph{Encoding} Format \textsf{ALU\_DPWRR1} (Section~\ref{sec:format-ALU-DPWRR1}), \verb|exucode2 = 0001|\\

\bitfields{ 1/P, 2/11, 1/1, 4/0001, 5/registerM, 1/--, 2/10, 3/101, 1/0, 5/registerO, 1/--, 5/registerP, 1/0 }


\paragraph{Syntax} \texttt{addx4dp} \emph{registerM} = \emph{registerP}, \emph{registerO}

\paragraph{Description} The \emph{registerO} and the \emph{registerP} are interpreted as two double words packed into 128 bits. The \emph{registerP} double words are shifted left by two bits and added to the \emph{registerO} double words. The two resulting double words are packed and stored into the \emph{registerM}.


\paragraph{Execution} {Reservation table \textsf{ALU\_LITE} (Section~\ref{sec:reservation-ALU-LITE})}
~
{\begin{lstlisting}
stage RR:
  new argument3 = PGR[registerO];
  new argument2 = PGR[registerP];
  for (I = 0; I < 2; I += 1) {
    result1.64[I] = _SX_64(argument3.64[I]) + (_SX_64(argument2.64[I]) << 2)
  };
stage E1:
  PGR[registerM] = result1;
\end{lstlisting}}
\newpage

\paragraph{Encoding} Format \textsf{ALU\_DPWRR1.M} (Section~\ref{sec:format-ALU-DPWRR1.M}), \verb|exucode2 = 0001|\\

\bitfields{ 1/P, 2/00, 2/-- --, 27/upper27, 1/1, 2/11, 1/1, 4/0001, 5/registerM, 1/--, 2/10, 3/101, 1/0, 1/splat32, 5/lower5, 5/registerP, 1/0 }


\paragraph{Syntax} \texttt{addx4dp} \emph{registerM} = \emph{registerP}, \emph{upper27\_\discretionary{}{}{}lower5}\emph{splat32}

\paragraph{Description} The \emph{upper27\_\discretionary{}{}{}lower5} and the \emph{registerP} are interpreted as two double words packed into 128 bits. The \emph{registerP} double words are shifted left by two bits and added to the \emph{upper27\_\discretionary{}{}{}lower5} double words. The two resulting double words are packed and stored into the \emph{registerM}.


\paragraph{Modifier(s)}
\noindent\begin{itemize}
\item splat32 (cf. splat32 in section \ref{par:modifier-splat32} on page \pageref{par:modifier-splat32})
\end{itemize}

\paragraph{Execution} {Reservation table \textsf{ALU\_LITE.X} (Section~\ref{sec:reservation-ALU-LITE.X})}
~
{\begin{lstlisting}
stage ID:
  new splat32 = _ZX_1(splat32);
stage RR:
  new argument3 = splat32 ? (_WX_32(upper27_lower5) << 32) | _ZX_32(_WX_32(upper27_lower5)) : _SX_32(_WX_32(upper27_lower5));
  new argument2 = PGR[registerP];
  for (I = 0; I < 2; I += 1) {
    result1.64[I] = _SX_64(argument3.64[I]) + (_SX_64(argument2.64[I]) << 2)
  };
stage E1:
  PGR[registerM] = result1;
\end{lstlisting}}
\newpage
\subsubsection[{\small ADDX4HO: Add to Times 4 Half Word Octuple}]{ADDX4HO\hfill Add to Times 4 Half Word Octuple}
\label{instr:ADDX4HO}\index{addx4ho}
 
\paragraph{Encoding} Format \textsf{ALU\_HOWRR1} (Section~\ref{sec:format-ALU-HOWRR1}), \verb|exucode2 = 0001|\\

\bitfields{ 1/P, 2/11, 1/1, 4/0001, 5/registerM, 1/--, 2/10, 3/101, 1/1, 5/registerO, 1/--, 5/registerP, 1/0 }


\paragraph{Syntax} \texttt{addx4ho} \emph{registerM} = \emph{registerP}, \emph{registerO}

\paragraph{Description} The \emph{registerO} and the \emph{registerP} are interpreted as eight half words packed into 128 bits. The \emph{registerP} half words are shifted left by two bits and added to the \emph{registerO} half words. The eight resulting half words are packed and stored into the \emph{registerM}.


\paragraph{Execution} {Reservation table \textsf{ALU\_LITE} (Section~\ref{sec:reservation-ALU-LITE})}
~
{\begin{lstlisting}
stage RR:
  new argument3 = PGR[registerO];
  new argument2 = PGR[registerP];
  for (I = 0; I < 8; I += 1) {
    result1.16[I] = _SX_16(argument3.16[I]) + (_SX_16(argument2.16[I]) << 2)
  };
stage E1:
  PGR[registerM] = result1;
\end{lstlisting}}
\newpage

\paragraph{Encoding} Format \textsf{ALU\_HOWRR1.M} (Section~\ref{sec:format-ALU-HOWRR1.M}), \verb|exucode2 = 0001|\\

\bitfields{ 1/P, 2/00, 2/-- --, 27/upper27, 1/1, 2/11, 1/1, 4/0001, 5/registerM, 1/--, 2/10, 3/101, 1/1, 1/splat32, 5/lower5, 5/registerP, 1/0 }


\paragraph{Syntax} \texttt{addx4ho} \emph{registerM} = \emph{registerP}, \emph{upper27\_\discretionary{}{}{}lower5}\emph{splat32}

\paragraph{Description} The \emph{upper27\_\discretionary{}{}{}lower5} and the \emph{registerP} are interpreted as eight half words packed into 128 bits. The \emph{registerP} half words are shifted left by two bits and added to the \emph{upper27\_\discretionary{}{}{}lower5} half words. The eight resulting half words are packed and stored into the \emph{registerM}.


\paragraph{Modifier(s)}
\noindent\begin{itemize}
\item splat32 (cf. splat32 in section \ref{par:modifier-splat32} on page \pageref{par:modifier-splat32})
\end{itemize}

\paragraph{Execution} {Reservation table \textsf{ALU\_LITE.X} (Section~\ref{sec:reservation-ALU-LITE.X})}
~
{\begin{lstlisting}
stage ID:
  new splat32 = _ZX_1(splat32);
stage RR:
  new argument3 = splat32 ? (_WX_32(upper27_lower5) << 32) | _ZX_32(_WX_32(upper27_lower5)) : _SX_32(_WX_32(upper27_lower5));
  new argument2 = PGR[registerP];
  for (I = 0; I < 8; I += 1) {
    result1.16[I] = _SX_16(argument3.16[I]) + (_SX_16(argument2.16[I]) << 2)
  };
stage E1:
  PGR[registerM] = result1;
\end{lstlisting}}
\newpage
\subsubsection[{\small ADDX4W: Add to Times 4 Word}]{ADDX4W\hfill Add to Times 4 Word}
\label{instr:ADDX4W}\index{addx4w}
 
\paragraph{Encoding} Format \textsf{ALU\_WWRR1} (Section~\ref{sec:format-ALU-WWRR1}), \verb|exucode2 = 0001|\\

\bitfields{ 1/P, 2/11, 1/0, 4/0001, 6/registerW, 2/11, 3/101, 1/signextw, 6/registerY, 6/registerZ }


\paragraph{Syntax} \texttt{addx4w}\emph{signextw} \emph{registerW} = \emph{registerZ}, \emph{registerY}

\paragraph{Description} The \emph{registerZ} is shifted left by two bits and added to the \emph{registerY}. The result with the 32 upper bits cleared is stored into the \emph{registerW}.


\paragraph{Modifier(s)}
\noindent\begin{itemize}
\item signextw (cf. signextw in section \ref{par:modifier-signextw} on page \pageref{par:modifier-signextw})
\end{itemize}

\paragraph{Execution} {Reservation table \textsf{ALU\_TINY} (Section~\ref{sec:reservation-ALU-TINY})}
~
{\begin{lstlisting}
stage ID:
  new signextw = _ZX_1(signextw);
stage RR:
  new argument3 = GPR[registerY];
  new argument2 = GPR[registerZ];
  new result1 = argument3 + (argument2 << 2);
stage E1:
  GPR[registerW] = signextw ? _SX_32(result1) : _ZX_32(result1);
\end{lstlisting}}
\newpage

\paragraph{Encoding} Format \textsf{ALU\_WWRR1.W} (Section~\ref{sec:format-ALU-WWRR1.W}), \verb|exucode2 = 0001|\\

\bitfields{ 1/P, 2/00, 2/-- --, 27/upper27, 1/1, 2/11, 1/0, 4/0001, 6/registerW, 2/11, 3/101, 1/signextw, 1/0, 5/lower5, 6/registerZ }


\paragraph{Syntax} \texttt{addx4w}\emph{signextw} \emph{registerW} = \emph{registerZ}, \emph{upper27\_\discretionary{}{}{}lower5}

\paragraph{Description} The \emph{registerZ} is shifted left by two bits and added to the \emph{upper27\_\discretionary{}{}{}lower5}. The result with the 32 upper bits cleared is stored into the \emph{registerW}.


\paragraph{Modifier(s)}
\noindent\begin{itemize}
\item signextw (cf. signextw in section \ref{par:modifier-signextw} on page \pageref{par:modifier-signextw})
\end{itemize}

\paragraph{Execution} {Reservation table \textsf{ALU\_TINY.X} (Section~\ref{sec:reservation-ALU-TINY.X})}
~
{\begin{lstlisting}
stage ID:
  new signextw = _ZX_1(signextw);
stage RR:
  new argument3 = _WX_32(upper27_lower5);
  new argument2 = GPR[registerZ];
  new result1 = argument3 + (argument2 << 2);
stage E1:
  GPR[registerW] = signextw ? _SX_32(result1) : _ZX_32(result1);
\end{lstlisting}}
\newpage
\subsubsection[{\small ADDX4WQ: Add to Times 4 Word Quadruple}]{ADDX4WQ\hfill Add to Times 4 Word Quadruple}
\label{instr:ADDX4WQ}\index{addx4wq}
 
\paragraph{Encoding} Format \textsf{ALU\_WQWRR1} (Section~\ref{sec:format-ALU-WQWRR1}), \verb|exucode2 = 0001|\\

\bitfields{ 1/P, 2/11, 1/1, 4/0001, 5/registerM, 1/--, 2/10, 3/101, 1/0, 5/registerO, 1/--, 5/registerP, 1/1 }


\paragraph{Syntax} \texttt{addx4wq} \emph{registerM} = \emph{registerP}, \emph{registerO}

\paragraph{Description} The \emph{registerO} and the \emph{registerP} are interpreted as four words packed into 128 bits. The \emph{registerP} words are shifted left by two bits and added to the \emph{registerO} words. The four resulting words are packed and stored into the \emph{registerM}.


\paragraph{Execution} {Reservation table \textsf{ALU\_LITE} (Section~\ref{sec:reservation-ALU-LITE})}
~
{\begin{lstlisting}
stage RR:
  new argument3 = PGR[registerO];
  new argument2 = PGR[registerP];
  for (I = 0; I < 4; I += 1) {
    result1.32[I] = _SX_32(argument3.32[I]) + (_SX_32(argument2.32[I]) << 2)
  };
stage E1:
  PGR[registerM] = result1;
\end{lstlisting}}
\newpage

\paragraph{Encoding} Format \textsf{ALU\_WQWRR1.M} (Section~\ref{sec:format-ALU-WQWRR1.M}), \verb|exucode2 = 0001|\\

\bitfields{ 1/P, 2/00, 2/-- --, 27/upper27, 1/1, 2/11, 1/1, 4/0001, 5/registerM, 1/--, 2/10, 3/101, 1/0, 1/splat32, 5/lower5, 5/registerP, 1/1 }


\paragraph{Syntax} \texttt{addx4wq} \emph{registerM} = \emph{registerP}, \emph{upper27\_\discretionary{}{}{}lower5}\emph{splat32}

\paragraph{Description} The \emph{upper27\_\discretionary{}{}{}lower5} and the \emph{registerP} are interpreted as four words packed into 128 bits. The \emph{registerP} words are shifted left by two bits and added to the \emph{upper27\_\discretionary{}{}{}lower5} words. The four resulting words are packed and stored into the \emph{registerM}.


\paragraph{Modifier(s)}
\noindent\begin{itemize}
\item splat32 (cf. splat32 in section \ref{par:modifier-splat32} on page \pageref{par:modifier-splat32})
\end{itemize}

\paragraph{Execution} {Reservation table \textsf{ALU\_LITE.X} (Section~\ref{sec:reservation-ALU-LITE.X})}
~
{\begin{lstlisting}
stage ID:
  new splat32 = _ZX_1(splat32);
stage RR:
  new argument3 = splat32 ? (_WX_32(upper27_lower5) << 32) | _ZX_32(_WX_32(upper27_lower5)) : _SX_32(_WX_32(upper27_lower5));
  new argument2 = PGR[registerP];
  for (I = 0; I < 4; I += 1) {
    result1.32[I] = _SX_32(argument3.32[I]) + (_SX_32(argument2.32[I]) << 2)
  };
stage E1:
  PGR[registerM] = result1;
\end{lstlisting}}
\newpage
\subsubsection[{\small ADDX64D: Add to Times 64 Double Word}]{ADDX64D\hfill Add to Times 64 Double Word}
\label{instr:ADDX64D}\index{addx64d}
 
\paragraph{Encoding} Format \textsf{ALU\_DWRR1} (Section~\ref{sec:format-ALU-DWRR1}), \verb|exucode2 = 0101|\\

\bitfields{ 1/P, 2/11, 1/0, 4/0101, 6/registerW, 2/10, 3/101, 1/0, 6/registerY, 6/registerZ }


\paragraph{Syntax} \texttt{addx64d} \emph{registerW} = \emph{registerZ}, \emph{registerY}

\paragraph{Description} The \emph{registerZ} is shifted left by six bits and added to the \emph{registerY}. The result is stored into the \emph{registerW}.


\paragraph{Execution} {Reservation table \textsf{ALU\_TINY} (Section~\ref{sec:reservation-ALU-TINY})}
~
{\begin{lstlisting}
stage RR:
  new argument3 = GPR[registerY];
  new argument2 = GPR[registerZ];
  new result1 = argument3 + (argument2 << 6);
stage E1:
  GPR[registerW] = result1;
\end{lstlisting}}
\newpage

\paragraph{Encoding} Format \textsf{ALU\_DWRR1.M} (Section~\ref{sec:format-ALU-DWRR1.M}), \verb|exucode2 = 0101|\\

\bitfields{ 1/P, 2/00, 2/-- --, 27/upper27, 1/1, 2/11, 1/0, 4/0101, 6/registerW, 2/10, 3/101, 1/0, 1/splat32, 5/lower5, 6/registerZ }


\paragraph{Syntax} \texttt{addx64d} \emph{registerW} = \emph{registerZ}, \emph{upper27\_\discretionary{}{}{}lower5}\emph{splat32}

\paragraph{Description} The \emph{registerZ} is shifted left by six bits and added to the \emph{upper27\_\discretionary{}{}{}lower5}. The result is stored into the \emph{registerW}.


\paragraph{Modifier(s)}
\noindent\begin{itemize}
\item splat32 (cf. splat32 in section \ref{par:modifier-splat32} on page \pageref{par:modifier-splat32})
\end{itemize}

\paragraph{Execution} {Reservation table \textsf{ALU\_TINY.X} (Section~\ref{sec:reservation-ALU-TINY.X})}
~
{\begin{lstlisting}
stage ID:
  new splat32 = _ZX_1(splat32);
stage RR:
  new argument3 = splat32 ? (_WX_32(upper27_lower5) << 32) | _ZX_32(_WX_32(upper27_lower5)) : _SX_32(_WX_32(upper27_lower5));
  new argument2 = GPR[registerZ];
  new result1 = argument3 + (argument2 << 6);
stage E1:
  GPR[registerW] = result1;
\end{lstlisting}}
\newpage
\subsubsection[{\small ADDX64W: Add to Times 64 Word}]{ADDX64W\hfill Add to Times 64 Word}
\label{instr:ADDX64W}\index{addx64w}
 
\paragraph{Encoding} Format \textsf{ALU\_WWRR1} (Section~\ref{sec:format-ALU-WWRR1}), \verb|exucode2 = 0101|\\

\bitfields{ 1/P, 2/11, 1/0, 4/0101, 6/registerW, 2/11, 3/101, 1/signextw, 6/registerY, 6/registerZ }


\paragraph{Syntax} \texttt{addx64w}\emph{signextw} \emph{registerW} = \emph{registerZ}, \emph{registerY}

\paragraph{Description} The \emph{registerZ} is shifted left by six bits and added to the \emph{registerY}. The result with the 32 upper bits cleared is stored into the \emph{registerW}.


\paragraph{Modifier(s)}
\noindent\begin{itemize}
\item signextw (cf. signextw in section \ref{par:modifier-signextw} on page \pageref{par:modifier-signextw})
\end{itemize}

\paragraph{Execution} {Reservation table \textsf{ALU\_TINY} (Section~\ref{sec:reservation-ALU-TINY})}
~
{\begin{lstlisting}
stage ID:
  new signextw = _ZX_1(signextw);
stage RR:
  new argument3 = GPR[registerY];
  new argument2 = GPR[registerZ];
  new result1 = argument3 + (argument2 << 6);
stage E1:
  GPR[registerW] = signextw ? _SX_32(result1) : _ZX_32(result1);
\end{lstlisting}}
\newpage

\paragraph{Encoding} Format \textsf{ALU\_WWRR1.W} (Section~\ref{sec:format-ALU-WWRR1.W}), \verb|exucode2 = 0101|\\

\bitfields{ 1/P, 2/00, 2/-- --, 27/upper27, 1/1, 2/11, 1/0, 4/0101, 6/registerW, 2/11, 3/101, 1/signextw, 1/0, 5/lower5, 6/registerZ }


\paragraph{Syntax} \texttt{addx64w}\emph{signextw} \emph{registerW} = \emph{registerZ}, \emph{upper27\_\discretionary{}{}{}lower5}

\paragraph{Description} The \emph{registerZ} is shifted left by six bits and added to the \emph{upper27\_\discretionary{}{}{}lower5}. The result with the 32 upper bits cleared is stored into the \emph{registerW}.


\paragraph{Modifier(s)}
\noindent\begin{itemize}
\item signextw (cf. signextw in section \ref{par:modifier-signextw} on page \pageref{par:modifier-signextw})
\end{itemize}

\paragraph{Execution} {Reservation table \textsf{ALU\_TINY.X} (Section~\ref{sec:reservation-ALU-TINY.X})}
~
{\begin{lstlisting}
stage ID:
  new signextw = _ZX_1(signextw);
stage RR:
  new argument3 = _WX_32(upper27_lower5);
  new argument2 = GPR[registerZ];
  new result1 = argument3 + (argument2 << 6);
stage E1:
  GPR[registerW] = signextw ? _SX_32(result1) : _ZX_32(result1);
\end{lstlisting}}
\newpage
\subsubsection[{\small ADDX8BX: Add to Times 8 Byte Hexadecuple}]{ADDX8BX\hfill Add to Times 8 Byte Hexadecuple}
\label{instr:ADDX8BX}\index{addx8bx}
 
\paragraph{Encoding} Format \textsf{ALU\_BXWRR1} (Section~\ref{sec:format-ALU-BXWRR1}), \verb|exucode2 = 0010|\\

\bitfields{ 1/P, 2/11, 1/1, 4/0010, 5/registerM, 1/--, 2/10, 3/101, 1/1, 5/registerO, 1/--, 5/registerP, 1/1 }


\paragraph{Syntax} \texttt{addx8bx} \emph{registerM} = \emph{registerP}, \emph{registerO}

\paragraph{Description} The \emph{registerO} and the \emph{registerP} are interpreted as sixteen bytes packed into 128 bits. The \emph{registerP} bytes are shifted left by three bits and added to the \emph{registerO} bytes. The sixteen resulting bytes are packed and stored into the \emph{registerM}.


\paragraph{Execution} {Reservation table \textsf{ALU\_LITE} (Section~\ref{sec:reservation-ALU-LITE})}
~
{\begin{lstlisting}
stage RR:
  new argument3 = PGR[registerO];
  new argument2 = PGR[registerP];
  for (I = 0; I < 16; I += 1) {
    result1.8[I] = _SX_8(argument3.8[I]) + (_SX_8(argument2.8[I]) << 3)
  };
stage E1:
  PGR[registerM] = result1;
\end{lstlisting}}
\newpage

\paragraph{Encoding} Format \textsf{ALU\_BXWRR1.M} (Section~\ref{sec:format-ALU-BXWRR1.M}), \verb|exucode2 = 0010|\\

\bitfields{ 1/P, 2/00, 2/-- --, 27/upper27, 1/1, 2/11, 1/1, 4/0010, 5/registerM, 1/--, 2/10, 3/101, 1/1, 1/splat32, 5/lower5, 5/registerP, 1/1 }


\paragraph{Syntax} \texttt{addx8bx} \emph{registerM} = \emph{registerP}, \emph{upper27\_\discretionary{}{}{}lower5}\emph{splat32}

\paragraph{Description} The \emph{upper27\_\discretionary{}{}{}lower5} and the \emph{registerP} are interpreted as sixteen bytes packed into 128 bits. The \emph{registerP} bytes are shifted left by three bits and added to the \emph{upper27\_\discretionary{}{}{}lower5} bytes. The sixteen resulting bytes are packed and stored into the \emph{registerM}.


\paragraph{Modifier(s)}
\noindent\begin{itemize}
\item splat32 (cf. splat32 in section \ref{par:modifier-splat32} on page \pageref{par:modifier-splat32})
\end{itemize}

\paragraph{Execution} {Reservation table \textsf{ALU\_LITE.X} (Section~\ref{sec:reservation-ALU-LITE.X})}
~
{\begin{lstlisting}
stage ID:
  new splat32 = _ZX_1(splat32);
stage RR:
  new argument3 = splat32 ? (_WX_32(upper27_lower5) << 32) | _ZX_32(_WX_32(upper27_lower5)) : _SX_32(_WX_32(upper27_lower5));
  new argument2 = PGR[registerP];
  for (I = 0; I < 16; I += 1) {
    result1.8[I] = _SX_8(argument3.8[I]) + (_SX_8(argument2.8[I]) << 3)
  };
stage E1:
  PGR[registerM] = result1;
\end{lstlisting}}
\newpage
\subsubsection[{\small ADDX8D: Add to Times 8 Double Word}]{ADDX8D\hfill Add to Times 8 Double Word}
\label{instr:ADDX8D}\index{addx8d}
 
\paragraph{Encoding} Format \textsf{ALU\_DWRR1} (Section~\ref{sec:format-ALU-DWRR1}), \verb|exucode2 = 0010|\\

\bitfields{ 1/P, 2/11, 1/0, 4/0010, 6/registerW, 2/10, 3/101, 1/0, 6/registerY, 6/registerZ }


\paragraph{Syntax} \texttt{addx8d} \emph{registerW} = \emph{registerZ}, \emph{registerY}

\paragraph{Description} The \emph{registerZ} is shifted left by three bits and added to the \emph{registerY}. The result is stored into the \emph{registerW}.


\paragraph{Execution} {Reservation table \textsf{ALU\_TINY} (Section~\ref{sec:reservation-ALU-TINY})}
~
{\begin{lstlisting}
stage RR:
  new argument3 = GPR[registerY];
  new argument2 = GPR[registerZ];
  new result1 = argument3 + (argument2 << 3);
stage E1:
  GPR[registerW] = result1;
\end{lstlisting}}
\newpage

\paragraph{Encoding} Format \textsf{ALU\_DWRR1.M} (Section~\ref{sec:format-ALU-DWRR1.M}), \verb|exucode2 = 0010|\\

\bitfields{ 1/P, 2/00, 2/-- --, 27/upper27, 1/1, 2/11, 1/0, 4/0010, 6/registerW, 2/10, 3/101, 1/0, 1/splat32, 5/lower5, 6/registerZ }


\paragraph{Syntax} \texttt{addx8d} \emph{registerW} = \emph{registerZ}, \emph{upper27\_\discretionary{}{}{}lower5}\emph{splat32}

\paragraph{Description} The \emph{registerZ} is shifted left by three bits and added to the \emph{upper27\_\discretionary{}{}{}lower5}. The result is stored into the \emph{registerW}.


\paragraph{Modifier(s)}
\noindent\begin{itemize}
\item splat32 (cf. splat32 in section \ref{par:modifier-splat32} on page \pageref{par:modifier-splat32})
\end{itemize}

\paragraph{Execution} {Reservation table \textsf{ALU\_TINY.X} (Section~\ref{sec:reservation-ALU-TINY.X})}
~
{\begin{lstlisting}
stage ID:
  new splat32 = _ZX_1(splat32);
stage RR:
  new argument3 = splat32 ? (_WX_32(upper27_lower5) << 32) | _ZX_32(_WX_32(upper27_lower5)) : _SX_32(_WX_32(upper27_lower5));
  new argument2 = GPR[registerZ];
  new result1 = argument3 + (argument2 << 3);
stage E1:
  GPR[registerW] = result1;
\end{lstlisting}}
\newpage
\subsubsection[{\small ADDX8DP: Add to Times 8 Double Word Pair}]{ADDX8DP\hfill Add to Times 8 Double Word Pair}
\label{instr:ADDX8DP}\index{addx8dp}
 
\paragraph{Encoding} Format \textsf{ALU\_DPWRR1} (Section~\ref{sec:format-ALU-DPWRR1}), \verb|exucode2 = 0010|\\

\bitfields{ 1/P, 2/11, 1/1, 4/0010, 5/registerM, 1/--, 2/10, 3/101, 1/0, 5/registerO, 1/--, 5/registerP, 1/0 }


\paragraph{Syntax} \texttt{addx8dp} \emph{registerM} = \emph{registerP}, \emph{registerO}

\paragraph{Description} The \emph{registerO} and the \emph{registerP} are interpreted as two double words packed into 128 bits. The \emph{registerP} double words are shifted left by three bits and added to the \emph{registerO} double words. The two resulting double words are packed and stored into the \emph{registerM}.


\paragraph{Execution} {Reservation table \textsf{ALU\_LITE} (Section~\ref{sec:reservation-ALU-LITE})}
~
{\begin{lstlisting}
stage RR:
  new argument3 = PGR[registerO];
  new argument2 = PGR[registerP];
  for (I = 0; I < 2; I += 1) {
    result1.64[I] = _SX_64(argument3.64[I]) + (_SX_64(argument2.64[I]) << 3)
  };
stage E1:
  PGR[registerM] = result1;
\end{lstlisting}}
\newpage

\paragraph{Encoding} Format \textsf{ALU\_DPWRR1.M} (Section~\ref{sec:format-ALU-DPWRR1.M}), \verb|exucode2 = 0010|\\

\bitfields{ 1/P, 2/00, 2/-- --, 27/upper27, 1/1, 2/11, 1/1, 4/0010, 5/registerM, 1/--, 2/10, 3/101, 1/0, 1/splat32, 5/lower5, 5/registerP, 1/0 }


\paragraph{Syntax} \texttt{addx8dp} \emph{registerM} = \emph{registerP}, \emph{upper27\_\discretionary{}{}{}lower5}\emph{splat32}

\paragraph{Description} The \emph{upper27\_\discretionary{}{}{}lower5} and the \emph{registerP} are interpreted as two double words packed into 128 bits. The \emph{registerP} double words are shifted left by three bits and added to the \emph{upper27\_\discretionary{}{}{}lower5} double words. The two resulting double words are packed and stored into the \emph{registerM}.


\paragraph{Modifier(s)}
\noindent\begin{itemize}
\item splat32 (cf. splat32 in section \ref{par:modifier-splat32} on page \pageref{par:modifier-splat32})
\end{itemize}

\paragraph{Execution} {Reservation table \textsf{ALU\_LITE.X} (Section~\ref{sec:reservation-ALU-LITE.X})}
~
{\begin{lstlisting}
stage ID:
  new splat32 = _ZX_1(splat32);
stage RR:
  new argument3 = splat32 ? (_WX_32(upper27_lower5) << 32) | _ZX_32(_WX_32(upper27_lower5)) : _SX_32(_WX_32(upper27_lower5));
  new argument2 = PGR[registerP];
  for (I = 0; I < 2; I += 1) {
    result1.64[I] = _SX_64(argument3.64[I]) + (_SX_64(argument2.64[I]) << 3)
  };
stage E1:
  PGR[registerM] = result1;
\end{lstlisting}}
\newpage
\subsubsection[{\small ADDX8HO: Add to Times 8 Half Word Octuple}]{ADDX8HO\hfill Add to Times 8 Half Word Octuple}
\label{instr:ADDX8HO}\index{addx8ho}
 
\paragraph{Encoding} Format \textsf{ALU\_HOWRR1} (Section~\ref{sec:format-ALU-HOWRR1}), \verb|exucode2 = 0010|\\

\bitfields{ 1/P, 2/11, 1/1, 4/0010, 5/registerM, 1/--, 2/10, 3/101, 1/1, 5/registerO, 1/--, 5/registerP, 1/0 }


\paragraph{Syntax} \texttt{addx8ho} \emph{registerM} = \emph{registerP}, \emph{registerO}

\paragraph{Description} The \emph{registerO} and the \emph{registerP} are interpreted as eight half words packed into 128 bits. The \emph{registerP} half words are shifted left by three bits and added to the \emph{registerO} half words. The eight resulting half words are packed and stored into the \emph{registerM}.


\paragraph{Execution} {Reservation table \textsf{ALU\_LITE} (Section~\ref{sec:reservation-ALU-LITE})}
~
{\begin{lstlisting}
stage RR:
  new argument3 = PGR[registerO];
  new argument2 = PGR[registerP];
  for (I = 0; I < 8; I += 1) {
    result1.16[I] = _SX_16(argument3.16[I]) + (_SX_16(argument2.16[I]) << 3)
  };
stage E1:
  PGR[registerM] = result1;
\end{lstlisting}}
\newpage

\paragraph{Encoding} Format \textsf{ALU\_HOWRR1.M} (Section~\ref{sec:format-ALU-HOWRR1.M}), \verb|exucode2 = 0010|\\

\bitfields{ 1/P, 2/00, 2/-- --, 27/upper27, 1/1, 2/11, 1/1, 4/0010, 5/registerM, 1/--, 2/10, 3/101, 1/1, 1/splat32, 5/lower5, 5/registerP, 1/0 }


\paragraph{Syntax} \texttt{addx8ho} \emph{registerM} = \emph{registerP}, \emph{upper27\_\discretionary{}{}{}lower5}\emph{splat32}

\paragraph{Description} The \emph{upper27\_\discretionary{}{}{}lower5} and the \emph{registerP} are interpreted as eight half words packed into 128 bits. The \emph{registerP} half words are shifted left by three bits and added to the \emph{upper27\_\discretionary{}{}{}lower5} half words. The eight resulting half words are packed and stored into the \emph{registerM}.


\paragraph{Modifier(s)}
\noindent\begin{itemize}
\item splat32 (cf. splat32 in section \ref{par:modifier-splat32} on page \pageref{par:modifier-splat32})
\end{itemize}

\paragraph{Execution} {Reservation table \textsf{ALU\_LITE.X} (Section~\ref{sec:reservation-ALU-LITE.X})}
~
{\begin{lstlisting}
stage ID:
  new splat32 = _ZX_1(splat32);
stage RR:
  new argument3 = splat32 ? (_WX_32(upper27_lower5) << 32) | _ZX_32(_WX_32(upper27_lower5)) : _SX_32(_WX_32(upper27_lower5));
  new argument2 = PGR[registerP];
  for (I = 0; I < 8; I += 1) {
    result1.16[I] = _SX_16(argument3.16[I]) + (_SX_16(argument2.16[I]) << 3)
  };
stage E1:
  PGR[registerM] = result1;
\end{lstlisting}}
\newpage
\subsubsection[{\small ADDX8W: Add to Times 8 Word}]{ADDX8W\hfill Add to Times 8 Word}
\label{instr:ADDX8W}\index{addx8w}
 
\paragraph{Encoding} Format \textsf{ALU\_WWRR1} (Section~\ref{sec:format-ALU-WWRR1}), \verb|exucode2 = 0010|\\

\bitfields{ 1/P, 2/11, 1/0, 4/0010, 6/registerW, 2/11, 3/101, 1/signextw, 6/registerY, 6/registerZ }


\paragraph{Syntax} \texttt{addx8w}\emph{signextw} \emph{registerW} = \emph{registerZ}, \emph{registerY}

\paragraph{Description} The \emph{registerZ} is shifted left by three bits and added to the \emph{registerY}. The result with the 32 upper bits cleared is stored into the \emph{registerW}.


\paragraph{Modifier(s)}
\noindent\begin{itemize}
\item signextw (cf. signextw in section \ref{par:modifier-signextw} on page \pageref{par:modifier-signextw})
\end{itemize}

\paragraph{Execution} {Reservation table \textsf{ALU\_TINY} (Section~\ref{sec:reservation-ALU-TINY})}
~
{\begin{lstlisting}
stage ID:
  new signextw = _ZX_1(signextw);
stage RR:
  new argument3 = GPR[registerY];
  new argument2 = GPR[registerZ];
  new result1 = argument3 + (argument2 << 3);
stage E1:
  GPR[registerW] = signextw ? _SX_32(result1) : _ZX_32(result1);
\end{lstlisting}}
\newpage

\paragraph{Encoding} Format \textsf{ALU\_WWRR1.W} (Section~\ref{sec:format-ALU-WWRR1.W}), \verb|exucode2 = 0010|\\

\bitfields{ 1/P, 2/00, 2/-- --, 27/upper27, 1/1, 2/11, 1/0, 4/0010, 6/registerW, 2/11, 3/101, 1/signextw, 1/0, 5/lower5, 6/registerZ }


\paragraph{Syntax} \texttt{addx8w}\emph{signextw} \emph{registerW} = \emph{registerZ}, \emph{upper27\_\discretionary{}{}{}lower5}

\paragraph{Description} The \emph{registerZ} is shifted left by three bits and added to the \emph{upper27\_\discretionary{}{}{}lower5}. The result with the 32 upper bits cleared is stored into the \emph{registerW}.


\paragraph{Modifier(s)}
\noindent\begin{itemize}
\item signextw (cf. signextw in section \ref{par:modifier-signextw} on page \pageref{par:modifier-signextw})
\end{itemize}

\paragraph{Execution} {Reservation table \textsf{ALU\_TINY.X} (Section~\ref{sec:reservation-ALU-TINY.X})}
~
{\begin{lstlisting}
stage ID:
  new signextw = _ZX_1(signextw);
stage RR:
  new argument3 = _WX_32(upper27_lower5);
  new argument2 = GPR[registerZ];
  new result1 = argument3 + (argument2 << 3);
stage E1:
  GPR[registerW] = signextw ? _SX_32(result1) : _ZX_32(result1);
\end{lstlisting}}
\newpage
\subsubsection[{\small ADDX8WQ: Add to Times 8 Word Quadruple}]{ADDX8WQ\hfill Add to Times 8 Word Quadruple}
\label{instr:ADDX8WQ}\index{addx8wq}
 
\paragraph{Encoding} Format \textsf{ALU\_WQWRR1} (Section~\ref{sec:format-ALU-WQWRR1}), \verb|exucode2 = 0010|\\

\bitfields{ 1/P, 2/11, 1/1, 4/0010, 5/registerM, 1/--, 2/10, 3/101, 1/0, 5/registerO, 1/--, 5/registerP, 1/1 }


\paragraph{Syntax} \texttt{addx8wq} \emph{registerM} = \emph{registerP}, \emph{registerO}

\paragraph{Description} The \emph{registerO} and the \emph{registerP} are interpreted as four words packed into 128 bits. The \emph{registerP} words are shifted left by three bits and added to the \emph{registerO} words. The four resulting words are packed and stored into the \emph{registerM}.


\paragraph{Execution} {Reservation table \textsf{ALU\_LITE} (Section~\ref{sec:reservation-ALU-LITE})}
~
{\begin{lstlisting}
stage RR:
  new argument3 = PGR[registerO];
  new argument2 = PGR[registerP];
  for (I = 0; I < 4; I += 1) {
    result1.32[I] = _SX_32(argument3.32[I]) + (_SX_32(argument2.32[I]) << 3)
  };
stage E1:
  PGR[registerM] = result1;
\end{lstlisting}}
\newpage

\paragraph{Encoding} Format \textsf{ALU\_WQWRR1.M} (Section~\ref{sec:format-ALU-WQWRR1.M}), \verb|exucode2 = 0010|\\

\bitfields{ 1/P, 2/00, 2/-- --, 27/upper27, 1/1, 2/11, 1/1, 4/0010, 5/registerM, 1/--, 2/10, 3/101, 1/0, 1/splat32, 5/lower5, 5/registerP, 1/1 }


\paragraph{Syntax} \texttt{addx8wq} \emph{registerM} = \emph{registerP}, \emph{upper27\_\discretionary{}{}{}lower5}\emph{splat32}

\paragraph{Description} The \emph{upper27\_\discretionary{}{}{}lower5} and the \emph{registerP} are interpreted as four words packed into 128 bits. The \emph{registerP} words are shifted left by three bits and added to the \emph{upper27\_\discretionary{}{}{}lower5} words. The four resulting words are packed and stored into the \emph{registerM}.


\paragraph{Modifier(s)}
\noindent\begin{itemize}
\item splat32 (cf. splat32 in section \ref{par:modifier-splat32} on page \pageref{par:modifier-splat32})
\end{itemize}

\paragraph{Execution} {Reservation table \textsf{ALU\_LITE.X} (Section~\ref{sec:reservation-ALU-LITE.X})}
~
{\begin{lstlisting}
stage ID:
  new splat32 = _ZX_1(splat32);
stage RR:
  new argument3 = splat32 ? (_WX_32(upper27_lower5) << 32) | _ZX_32(_WX_32(upper27_lower5)) : _SX_32(_WX_32(upper27_lower5));
  new argument2 = PGR[registerP];
  for (I = 0; I < 4; I += 1) {
    result1.32[I] = _SX_32(argument3.32[I]) + (_SX_32(argument2.32[I]) << 3)
  };
stage E1:
  PGR[registerM] = result1;
\end{lstlisting}}
\newpage
\subsubsection[{\small ANDD: Bitwise And Double Word}]{ANDD\hfill Bitwise And Double Word}
\label{instr:ANDD}\index{andd}
 
\paragraph{Encoding} Format \textsf{ALU\_DWRR0} (Section~\ref{sec:format-ALU-DWRR0}), \verb|exucode2 = 1000|\\

\bitfields{ 1/P, 2/11, 1/0, 4/1000, 6/registerW, 2/10, 3/000, 1/0, 6/registerY, 6/registerZ }


\paragraph{Syntax} \texttt{andd} \emph{registerW} = \emph{registerZ}, \emph{registerY}

\paragraph{Description} Bitwise and of the \emph{registerZ} with the \emph{registerY}. The result is stored into the \emph{registerW}.


\paragraph{Execution} {Reservation table \textsf{ALU\_TINY} (Section~\ref{sec:reservation-ALU-TINY})}
~
{\begin{lstlisting}
stage RR:
  new argument3 = GPR[registerY];
  new argument2 = GPR[registerZ];
  new result1 = argument2 & argument3;
stage E1:
  GPR[registerW] = result1;
\end{lstlisting}}
\newpage

\paragraph{Encoding} Format \textsf{ALU\_DWRR0.M} (Section~\ref{sec:format-ALU-DWRR0.M}), \verb|exucode2 = 1000|\\

\bitfields{ 1/P, 2/00, 2/-- --, 27/upper27, 1/1, 2/11, 1/0, 4/1000, 6/registerW, 2/10, 3/000, 1/0, 1/splat32, 5/lower5, 6/registerZ }


\paragraph{Syntax} \texttt{andd} \emph{registerW} = \emph{registerZ}, \emph{upper27\_\discretionary{}{}{}lower5}\emph{splat32}

\paragraph{Description} Bitwise and of the \emph{registerZ} with the \emph{upper27\_\discretionary{}{}{}lower5}. The result is stored into the \emph{registerW}.


\paragraph{Modifier(s)}
\noindent\begin{itemize}
\item splat32 (cf. splat32 in section \ref{par:modifier-splat32} on page \pageref{par:modifier-splat32})
\end{itemize}

\paragraph{Execution} {Reservation table \textsf{ALU\_TINY.X} (Section~\ref{sec:reservation-ALU-TINY.X})}
~
{\begin{lstlisting}
stage ID:
  new splat32 = _ZX_1(splat32);
stage RR:
  new argument3 = splat32 ? (_WX_32(upper27_lower5) << 32) | _ZX_32(_WX_32(upper27_lower5)) : _SX_32(_WX_32(upper27_lower5));
  new argument2 = GPR[registerZ];
  new result1 = argument2 & argument3;
stage E1:
  GPR[registerW] = result1;
\end{lstlisting}}
\newpage

\paragraph{Encoding} Format \textsf{ALU\_DWRI} (Section~\ref{sec:format-ALU-DWRI}), \verb|exucode2 = 1000|\\

\bitfields{ 1/P, 2/11, 1/0, 4/1000, 6/registerW, 2/00, 10/signed10, 6/registerZ }


\paragraph{Syntax} \texttt{andd} \emph{registerW} = \emph{registerZ}, \emph{signed10}

\paragraph{Description} Bitwise and of the \emph{registerZ} with the \emph{signed10}. The result is stored into the \emph{registerW}.


\paragraph{Execution} {Reservation table \textsf{ALU\_TINY} (Section~\ref{sec:reservation-ALU-TINY})}
~
{\begin{lstlisting}
stage RR:
  new argument3 = _SX_10(signed10);
  new argument2 = GPR[registerZ];
  new result1 = argument2 & argument3;
stage E1:
  GPR[registerW] = result1;
\end{lstlisting}}
\newpage

\paragraph{Encoding} Format \textsf{ALU\_DWRI.X} (Section~\ref{sec:format-ALU-DWRI.X}), \verb|exucode2 = 1000|\\

\bitfields{ 1/P, 2/00, 2/-- --, 27/upper27, 1/1, 2/11, 1/0, 4/1000, 6/registerW, 2/00, 10/lower10, 6/registerZ }


\paragraph{Syntax} \texttt{andd} \emph{registerW} = \emph{registerZ}, \emph{upper27\_\discretionary{}{}{}lower10}

\paragraph{Description} Bitwise and of the \emph{registerZ} with the \emph{upper27\_\discretionary{}{}{}lower10}. The result is stored into the \emph{registerW}.


\paragraph{Execution} {Reservation table \textsf{ALU\_TINY.X} (Section~\ref{sec:reservation-ALU-TINY.X})}
~
{\begin{lstlisting}
stage RR:
  new argument3 = _SX_37(upper27_lower10);
  new argument2 = GPR[registerZ];
  new result1 = argument2 & argument3;
stage E1:
  GPR[registerW] = result1;
\end{lstlisting}}
\newpage

\paragraph{Encoding} Format \textsf{ALU\_DWRI.Y} (Section~\ref{sec:format-ALU-DWRI.Y}), \verb|exucode2 = 1000|\\

\bitfields{ 1/P, 2/00, 2/-- --, 27/extend27, 1/1, 2/00, 2/-- --, 27/upper27, 1/1, 2/11, 1/0, 4/1000, 6/registerW, 2/00, 10/lower10, 6/registerZ }


\paragraph{Syntax} \texttt{andd} \emph{registerW} = \emph{registerZ}, \emph{extend27\_\discretionary{}{}{}upper27\_\discretionary{}{}{}lower10}

\paragraph{Description} Bitwise and of the \emph{registerZ} with the \emph{extend27\_\discretionary{}{}{}upper27\_\discretionary{}{}{}lower10}. The result is stored into the \emph{registerW}.


\paragraph{Execution} {Reservation table \textsf{ALU\_TINY.Y} (Section~\ref{sec:reservation-ALU-TINY.Y})}
~
{\begin{lstlisting}
stage RR:
  new argument3 = _WX_64(extend27_upper27_lower10);
  new argument2 = GPR[registerZ];
  new result1 = argument2 & argument3;
stage E1:
  GPR[registerW] = result1;
\end{lstlisting}}
\newpage
\subsubsection[{\small ANDND: Bitwise And with Not Operand Double Word}]{ANDND\hfill Bitwise And with Not Operand Double Word}
\label{instr:ANDND}\index{andnd}
 
\paragraph{Encoding} Format \textsf{ALU\_DWRR0} (Section~\ref{sec:format-ALU-DWRR0}), \verb|exucode2 = 1110|\\

\bitfields{ 1/P, 2/11, 1/0, 4/1110, 6/registerW, 2/10, 3/000, 1/0, 6/registerY, 6/registerZ }


\paragraph{Syntax} \texttt{andnd} \emph{registerW} = \emph{registerZ}, \emph{registerY}

\paragraph{Description} Bitwise and of not the \emph{registerZ} with the \emph{registerY}. The result is stored into the \emph{registerW}.


\paragraph{Execution} {Reservation table \textsf{ALU\_TINY} (Section~\ref{sec:reservation-ALU-TINY})}
~
{\begin{lstlisting}
stage RR:
  new argument3 = GPR[registerY];
  new argument2 = GPR[registerZ];
  new result1 = ~argument2 & argument3;
stage E1:
  GPR[registerW] = result1;
\end{lstlisting}}
\newpage

\paragraph{Encoding} Format \textsf{ALU\_DWRR0.M} (Section~\ref{sec:format-ALU-DWRR0.M}), \verb|exucode2 = 1110|\\

\bitfields{ 1/P, 2/00, 2/-- --, 27/upper27, 1/1, 2/11, 1/0, 4/1110, 6/registerW, 2/10, 3/000, 1/0, 1/splat32, 5/lower5, 6/registerZ }


\paragraph{Syntax} \texttt{andnd} \emph{registerW} = \emph{registerZ}, \emph{upper27\_\discretionary{}{}{}lower5}\emph{splat32}

\paragraph{Description} Bitwise and of not the \emph{registerZ} with the \emph{upper27\_\discretionary{}{}{}lower5}. The result is stored into the \emph{registerW}.


\paragraph{Modifier(s)}
\noindent\begin{itemize}
\item splat32 (cf. splat32 in section \ref{par:modifier-splat32} on page \pageref{par:modifier-splat32})
\end{itemize}

\paragraph{Execution} {Reservation table \textsf{ALU\_TINY.X} (Section~\ref{sec:reservation-ALU-TINY.X})}
~
{\begin{lstlisting}
stage ID:
  new splat32 = _ZX_1(splat32);
stage RR:
  new argument3 = splat32 ? (_WX_32(upper27_lower5) << 32) | _ZX_32(_WX_32(upper27_lower5)) : _SX_32(_WX_32(upper27_lower5));
  new argument2 = GPR[registerZ];
  new result1 = ~argument2 & argument3;
stage E1:
  GPR[registerW] = result1;
\end{lstlisting}}
\newpage

\paragraph{Encoding} Format \textsf{ALU\_DWRI} (Section~\ref{sec:format-ALU-DWRI}), \verb|exucode2 = 1110|\\

\bitfields{ 1/P, 2/11, 1/0, 4/1110, 6/registerW, 2/00, 10/signed10, 6/registerZ }


\paragraph{Syntax} \texttt{andnd} \emph{registerW} = \emph{registerZ}, \emph{signed10}

\paragraph{Description} Bitwise and of not the \emph{registerZ} with the \emph{signed10}. The result is stored into the \emph{registerW}.


\paragraph{Execution} {Reservation table \textsf{ALU\_TINY} (Section~\ref{sec:reservation-ALU-TINY})}
~
{\begin{lstlisting}
stage RR:
  new argument3 = _SX_10(signed10);
  new argument2 = GPR[registerZ];
  new result1 = ~argument2 & argument3;
stage E1:
  GPR[registerW] = result1;
\end{lstlisting}}
\newpage

\paragraph{Encoding} Format \textsf{ALU\_DWRI.X} (Section~\ref{sec:format-ALU-DWRI.X}), \verb|exucode2 = 1110|\\

\bitfields{ 1/P, 2/00, 2/-- --, 27/upper27, 1/1, 2/11, 1/0, 4/1110, 6/registerW, 2/00, 10/lower10, 6/registerZ }


\paragraph{Syntax} \texttt{andnd} \emph{registerW} = \emph{registerZ}, \emph{upper27\_\discretionary{}{}{}lower10}

\paragraph{Description} Bitwise and of not the \emph{registerZ} with the \emph{upper27\_\discretionary{}{}{}lower10}. The result is stored into the \emph{registerW}.


\paragraph{Execution} {Reservation table \textsf{ALU\_TINY.X} (Section~\ref{sec:reservation-ALU-TINY.X})}
~
{\begin{lstlisting}
stage RR:
  new argument3 = _SX_37(upper27_lower10);
  new argument2 = GPR[registerZ];
  new result1 = ~argument2 & argument3;
stage E1:
  GPR[registerW] = result1;
\end{lstlisting}}
\newpage

\paragraph{Encoding} Format \textsf{ALU\_DWRI.Y} (Section~\ref{sec:format-ALU-DWRI.Y}), \verb|exucode2 = 1110|\\

\bitfields{ 1/P, 2/00, 2/-- --, 27/extend27, 1/1, 2/00, 2/-- --, 27/upper27, 1/1, 2/11, 1/0, 4/1110, 6/registerW, 2/00, 10/lower10, 6/registerZ }


\paragraph{Syntax} \texttt{andnd} \emph{registerW} = \emph{registerZ}, \emph{extend27\_\discretionary{}{}{}upper27\_\discretionary{}{}{}lower10}

\paragraph{Description} Bitwise and of not the \emph{registerZ} with the \emph{extend27\_\discretionary{}{}{}upper27\_\discretionary{}{}{}lower10}. The result is stored into the \emph{registerW}.


\paragraph{Execution} {Reservation table \textsf{ALU\_TINY.Y} (Section~\ref{sec:reservation-ALU-TINY.Y})}
~
{\begin{lstlisting}
stage RR:
  new argument3 = _WX_64(extend27_upper27_lower10);
  new argument2 = GPR[registerZ];
  new result1 = ~argument2 & argument3;
stage E1:
  GPR[registerW] = result1;
\end{lstlisting}}
\newpage
\subsubsection[{\small ANDNQ: Bitwise And with Not Operand Quadruple Word}]{ANDNQ\hfill Bitwise And with Not Operand Quadruple Word}
\label{instr:ANDNQ}\index{andnq}
 
\paragraph{Encoding} Format \textsf{ALU\_QWRR0} (Section~\ref{sec:format-ALU-QWRR0}), \verb|exucode2 = 1110|\\

\bitfields{ 1/P, 2/11, 1/0, 4/1110, 5/registerM, 1/--, 2/10, 3/000, 1/1, 5/registerO, 1/--, 5/registerP, 1/-- }


\paragraph{Syntax} \texttt{andnq} \emph{registerM} = \emph{registerP}, \emph{registerO}

\paragraph{Description} Bitwise and of not the \emph{registerP} with the \emph{registerO}. The result is stored into the \emph{registerM}.


\paragraph{Execution} {Reservation table \textsf{ALU\_LITE} (Section~\ref{sec:reservation-ALU-LITE})}
~
{\begin{lstlisting}
stage RR:
  new argument3 = PGR[registerO];
  new argument2 = PGR[registerP];
  new result1 = ~argument2 & argument3;
stage E1:
  PGR[registerM] = result1;
\end{lstlisting}}
\newpage

\paragraph{Encoding} Format \textsf{ALU\_QWRR0.M} (Section~\ref{sec:format-ALU-QWRR0.M}), \verb|exucode2 = 1110|\\

\bitfields{ 1/P, 2/00, 2/-- --, 27/upper27, 1/1, 2/11, 1/0, 4/1110, 5/registerM, 1/--, 2/10, 3/000, 1/1, 1/splat32, 5/lower5, 5/registerP, 1/1 }


\paragraph{Syntax} \texttt{andnq} \emph{registerM} = \emph{registerP}, \emph{upper27\_\discretionary{}{}{}lower5}\emph{splat32}

\paragraph{Description} Bitwise and of not the \emph{registerP} with the \emph{upper27\_\discretionary{}{}{}lower5}. The result is stored into the \emph{registerM}.


\paragraph{Modifier(s)}
\noindent\begin{itemize}
\item splat32 (cf. splat32 in section \ref{par:modifier-splat32} on page \pageref{par:modifier-splat32})
\end{itemize}

\paragraph{Execution} {Reservation table \textsf{ALU\_LITE.X} (Section~\ref{sec:reservation-ALU-LITE.X})}
~
{\begin{lstlisting}
stage ID:
  new splat32 = _ZX_1(splat32);
stage RR:
  new argument3 = splat32 ? (_WX_32(upper27_lower5) << 32) | _ZX_32(_WX_32(upper27_lower5)) : _SX_32(_WX_32(upper27_lower5));
  new argument2 = PGR[registerP];
  new result1 = ~argument2 & argument3;
stage E1:
  PGR[registerM] = result1;
\end{lstlisting}}
\newpage
\subsubsection[{\small ANDNW: Bitwise And with Not Operand Word}]{ANDNW\hfill Bitwise And with Not Operand Word}
\label{instr:ANDNW}\index{andnw}
 
\paragraph{Encoding} Format \textsf{ALU\_WWRR0} (Section~\ref{sec:format-ALU-WWRR0}), \verb|exucode2 = 1110|\\

\bitfields{ 1/P, 2/11, 1/0, 4/1110, 6/registerW, 2/11, 3/000, 1/signextw, 6/registerY, 6/registerZ }


\paragraph{Syntax} \texttt{andnw}\emph{signextw} \emph{registerW} = \emph{registerZ}, \emph{registerY}

\paragraph{Description} Bitwise and of not the \emph{registerZ} with the \emph{registerY}. The result with the 32 upper bits cleared is stored into the \emph{registerW}.


\paragraph{Modifier(s)}
\noindent\begin{itemize}
\item signextw (cf. signextw in section \ref{par:modifier-signextw} on page \pageref{par:modifier-signextw})
\end{itemize}

\paragraph{Execution} {Reservation table \textsf{ALU\_TINY} (Section~\ref{sec:reservation-ALU-TINY})}
~
{\begin{lstlisting}
stage ID:
  new signextw = _ZX_1(signextw);
stage RR:
  new argument3 = GPR[registerY];
  new argument2 = GPR[registerZ];
  new result1 = ~argument2 & argument3;
stage E1:
  GPR[registerW] = signextw ? _SX_32(result1) : _ZX_32(result1);
\end{lstlisting}}
\newpage

\paragraph{Encoding} Format \textsf{ALU\_WWRR0.W} (Section~\ref{sec:format-ALU-WWRR0.W}), \verb|exucode2 = 1110|\\

\bitfields{ 1/P, 2/00, 2/-- --, 27/upper27, 1/1, 2/11, 1/0, 4/1110, 6/registerW, 2/11, 3/000, 1/signextw, 1/0, 5/lower5, 6/registerZ }


\paragraph{Syntax} \texttt{andnw}\emph{signextw} \emph{registerW} = \emph{registerZ}, \emph{upper27\_\discretionary{}{}{}lower5}

\paragraph{Description} Bitwise and of not the \emph{registerZ} with the \emph{upper27\_\discretionary{}{}{}lower5}. The result with the 32 upper bits cleared is stored into the \emph{registerW}.


\paragraph{Modifier(s)}
\noindent\begin{itemize}
\item signextw (cf. signextw in section \ref{par:modifier-signextw} on page \pageref{par:modifier-signextw})
\end{itemize}

\paragraph{Execution} {Reservation table \textsf{ALU\_TINY.X} (Section~\ref{sec:reservation-ALU-TINY.X})}
~
{\begin{lstlisting}
stage ID:
  new signextw = _ZX_1(signextw);
stage RR:
  new argument3 = _WX_32(upper27_lower5);
  new argument2 = GPR[registerZ];
  new result1 = ~argument2 & argument3;
stage E1:
  GPR[registerW] = signextw ? _SX_32(result1) : _ZX_32(result1);
\end{lstlisting}}
\newpage
\subsubsection[{\small ANDQ: Bitwise And Quadruple Word}]{ANDQ\hfill Bitwise And Quadruple Word}
\label{instr:ANDQ}\index{andq}
 
\paragraph{Encoding} Format \textsf{ALU\_QWRR0} (Section~\ref{sec:format-ALU-QWRR0}), \verb|exucode2 = 1000|\\

\bitfields{ 1/P, 2/11, 1/0, 4/1000, 5/registerM, 1/--, 2/10, 3/000, 1/1, 5/registerO, 1/--, 5/registerP, 1/-- }


\paragraph{Syntax} \texttt{andq} \emph{registerM} = \emph{registerP}, \emph{registerO}

\paragraph{Description} Bitwise and of the \emph{registerP} with the \emph{registerO}. The result is stored into the \emph{registerM}.


\paragraph{Execution} {Reservation table \textsf{ALU\_LITE} (Section~\ref{sec:reservation-ALU-LITE})}
~
{\begin{lstlisting}
stage RR:
  new argument3 = PGR[registerO];
  new argument2 = PGR[registerP];
  new result1 = argument2 & argument3;
stage E1:
  PGR[registerM] = result1;
\end{lstlisting}}
\newpage

\paragraph{Encoding} Format \textsf{ALU\_QWRR0.M} (Section~\ref{sec:format-ALU-QWRR0.M}), \verb|exucode2 = 1000|\\

\bitfields{ 1/P, 2/00, 2/-- --, 27/upper27, 1/1, 2/11, 1/0, 4/1000, 5/registerM, 1/--, 2/10, 3/000, 1/1, 1/splat32, 5/lower5, 5/registerP, 1/1 }


\paragraph{Syntax} \texttt{andq} \emph{registerM} = \emph{registerP}, \emph{upper27\_\discretionary{}{}{}lower5}\emph{splat32}

\paragraph{Description} Bitwise and of the \emph{registerP} with the \emph{upper27\_\discretionary{}{}{}lower5}. The result is stored into the \emph{registerM}.


\paragraph{Modifier(s)}
\noindent\begin{itemize}
\item splat32 (cf. splat32 in section \ref{par:modifier-splat32} on page \pageref{par:modifier-splat32})
\end{itemize}

\paragraph{Execution} {Reservation table \textsf{ALU\_LITE.X} (Section~\ref{sec:reservation-ALU-LITE.X})}
~
{\begin{lstlisting}
stage ID:
  new splat32 = _ZX_1(splat32);
stage RR:
  new argument3 = splat32 ? (_WX_32(upper27_lower5) << 32) | _ZX_32(_WX_32(upper27_lower5)) : _SX_32(_WX_32(upper27_lower5));
  new argument2 = PGR[registerP];
  new result1 = argument2 & argument3;
stage E1:
  PGR[registerM] = result1;
\end{lstlisting}}
\newpage
\subsubsection[{\small ANDW: Bitwise And Word}]{ANDW\hfill Bitwise And Word}
\label{instr:ANDW}\index{andw}
 
\paragraph{Encoding} Format \textsf{ALU\_WWRR0} (Section~\ref{sec:format-ALU-WWRR0}), \verb|exucode2 = 1000|\\

\bitfields{ 1/P, 2/11, 1/0, 4/1000, 6/registerW, 2/11, 3/000, 1/signextw, 6/registerY, 6/registerZ }


\paragraph{Syntax} \texttt{andw}\emph{signextw} \emph{registerW} = \emph{registerZ}, \emph{registerY}

\paragraph{Description} Bitwise and of the \emph{registerZ} with the \emph{registerY}. The result with the 32 upper bits cleared is stored into the \emph{registerW}.


\paragraph{Modifier(s)}
\noindent\begin{itemize}
\item signextw (cf. signextw in section \ref{par:modifier-signextw} on page \pageref{par:modifier-signextw})
\end{itemize}

\paragraph{Execution} {Reservation table \textsf{ALU\_TINY} (Section~\ref{sec:reservation-ALU-TINY})}
~
{\begin{lstlisting}
stage ID:
  new signextw = _ZX_1(signextw);
stage RR:
  new argument3 = GPR[registerY];
  new argument2 = GPR[registerZ];
  new result1 = argument2 & argument3;
stage E1:
  GPR[registerW] = signextw ? _SX_32(result1) : _ZX_32(result1);
\end{lstlisting}}
\newpage

\paragraph{Encoding} Format \textsf{ALU\_WWRR0.W} (Section~\ref{sec:format-ALU-WWRR0.W}), \verb|exucode2 = 1000|\\

\bitfields{ 1/P, 2/00, 2/-- --, 27/upper27, 1/1, 2/11, 1/0, 4/1000, 6/registerW, 2/11, 3/000, 1/signextw, 1/0, 5/lower5, 6/registerZ }


\paragraph{Syntax} \texttt{andw}\emph{signextw} \emph{registerW} = \emph{registerZ}, \emph{upper27\_\discretionary{}{}{}lower5}

\paragraph{Description} Bitwise and of the \emph{registerZ} with the \emph{upper27\_\discretionary{}{}{}lower5}. The result with the 32 upper bits cleared is stored into the \emph{registerW}.


\paragraph{Modifier(s)}
\noindent\begin{itemize}
\item signextw (cf. signextw in section \ref{par:modifier-signextw} on page \pageref{par:modifier-signextw})
\end{itemize}

\paragraph{Execution} {Reservation table \textsf{ALU\_TINY.X} (Section~\ref{sec:reservation-ALU-TINY.X})}
~
{\begin{lstlisting}
stage ID:
  new signextw = _ZX_1(signextw);
stage RR:
  new argument3 = _WX_32(upper27_lower5);
  new argument2 = GPR[registerZ];
  new result1 = argument2 & argument3;
stage E1:
  GPR[registerW] = signextw ? _SX_32(result1) : _ZX_32(result1);
\end{lstlisting}}
\newpage
\subsubsection[{\small AVGBX: Average Byte Hexadecuple}]{AVGBX\hfill Average Byte Hexadecuple}
\label{instr:AVGBX}\index{avgbx}
 
\paragraph{Encoding} Format \textsf{ALU\_BXAVG} (Section~\ref{sec:format-ALU-BXAVG}), \verb|exucode2 = 1100|\\

\bitfields{ 1/P, 2/11, 1/1, 4/1100, 5/registerM, 1/--, 2/10, 3/001, 1/1, 5/registerO, 1/--, 5/registerP, 1/1 }


\paragraph{Syntax} \texttt{avgbx} \emph{registerM} = \emph{registerP}, \emph{registerO}

\paragraph{Description} The \emph{registerO} and the \emph{registerP} are interpreted as sixteen bytes packed into 128 bits. The \emph{registerO} bytes and the \emph{registerP} bytes are added. Each sum is shifted right arithmetic one bit. The sixteen resulting bytes are packed and stored into the \emph{registerM}.


\paragraph{Execution} {Reservation table \textsf{ALU\_LITE} (Section~\ref{sec:reservation-ALU-LITE})}
~
{\begin{lstlisting}
stage RR:
  new argument3 = PGR[registerO];
  new argument2 = PGR[registerP];
  for (I = 0; I < 16; I += 1) {
    result1.8[I] = (_SX_8(argument3.8[I]) + _SX_8(argument2.8[I])) >> 1
  };
stage E1:
  PGR[registerM] = result1;
\end{lstlisting}}
\newpage

\paragraph{Encoding} Format \textsf{ALU\_BXAVG.M} (Section~\ref{sec:format-ALU-BXAVG.M}), \verb|exucode2 = 1100|\\

\bitfields{ 1/P, 2/00, 2/-- --, 27/upper27, 1/1, 2/11, 1/1, 4/1100, 5/registerM, 1/--, 2/10, 3/001, 1/1, 1/splat32, 5/lower5, 5/registerP, 1/1 }


\paragraph{Syntax} \texttt{avgbx} \emph{registerM} = \emph{registerP}, \emph{upper27\_\discretionary{}{}{}lower5}\emph{splat32}

\paragraph{Description} The \emph{upper27\_\discretionary{}{}{}lower5} and the \emph{registerP} are interpreted as sixteen bytes packed into 128 bits. The \emph{upper27\_\discretionary{}{}{}lower5} bytes and the \emph{registerP} bytes are added. Each sum is shifted right arithmetic one bit. The sixteen resulting bytes are packed and stored into the \emph{registerM}.


\paragraph{Modifier(s)}
\noindent\begin{itemize}
\item splat32 (cf. splat32 in section \ref{par:modifier-splat32} on page \pageref{par:modifier-splat32})
\end{itemize}

\paragraph{Execution} {Reservation table \textsf{ALU\_LITE.X} (Section~\ref{sec:reservation-ALU-LITE.X})}
~
{\begin{lstlisting}
stage ID:
  new splat32 = _ZX_1(splat32);
stage RR:
  new argument3 = splat32 ? (_WX_32(upper27_lower5) << 32) | _ZX_32(_WX_32(upper27_lower5)) : _SX_32(_WX_32(upper27_lower5));
  new argument2 = PGR[registerP];
  for (I = 0; I < 16; I += 1) {
    result1.8[I] = (_SX_8(argument3.8[I]) + _SX_8(argument2.8[I])) >> 1
  };
stage E1:
  PGR[registerM] = result1;
\end{lstlisting}}
\newpage
\subsubsection[{\small AVGHO: Average Half Word Octuple}]{AVGHO\hfill Average Half Word Octuple}
\label{instr:AVGHO}\index{avgho}
 
\paragraph{Encoding} Format \textsf{ALU\_HOAVG} (Section~\ref{sec:format-ALU-HOAVG}), \verb|exucode2 = 1100|\\

\bitfields{ 1/P, 2/11, 1/1, 4/1100, 5/registerM, 1/--, 2/10, 3/001, 1/1, 5/registerO, 1/--, 5/registerP, 1/0 }


\paragraph{Syntax} \texttt{avgho} \emph{registerM} = \emph{registerP}, \emph{registerO}

\paragraph{Description} The \emph{registerO} and the \emph{registerP} are interpreted as eight half words packed into 128 bits. The \emph{registerO} half words and the \emph{registerP} half words are added. Each sum is shifted right arithmetic one bit. The eight resulting half words are packed and stored into the \emph{registerM}.


\paragraph{Execution} {Reservation table \textsf{ALU\_LITE} (Section~\ref{sec:reservation-ALU-LITE})}
~
{\begin{lstlisting}
stage RR:
  new argument3 = PGR[registerO];
  new argument2 = PGR[registerP];
  for (I = 0; I < 8; I += 1) {
    result1.16[I] = (_SX_16(argument3.16[I]) + _SX_16(argument2.16[I])) >> 1
  };
stage E1:
  PGR[registerM] = result1;
\end{lstlisting}}
\newpage

\paragraph{Encoding} Format \textsf{ALU\_HOAVG.M} (Section~\ref{sec:format-ALU-HOAVG.M}), \verb|exucode2 = 1100|\\

\bitfields{ 1/P, 2/00, 2/-- --, 27/upper27, 1/1, 2/11, 1/1, 4/1100, 5/registerM, 1/--, 2/10, 3/001, 1/1, 1/splat32, 5/lower5, 5/registerP, 1/0 }


\paragraph{Syntax} \texttt{avgho} \emph{registerM} = \emph{registerP}, \emph{upper27\_\discretionary{}{}{}lower5}\emph{splat32}

\paragraph{Description} The \emph{upper27\_\discretionary{}{}{}lower5} and the \emph{registerP} are interpreted as eight half words packed into 128 bits. The \emph{upper27\_\discretionary{}{}{}lower5} half words and the \emph{registerP} half words are added. Each sum is shifted right arithmetic one bit. The eight resulting half words are packed and stored into the \emph{registerM}.


\paragraph{Modifier(s)}
\noindent\begin{itemize}
\item splat32 (cf. splat32 in section \ref{par:modifier-splat32} on page \pageref{par:modifier-splat32})
\end{itemize}

\paragraph{Execution} {Reservation table \textsf{ALU\_LITE.X} (Section~\ref{sec:reservation-ALU-LITE.X})}
~
{\begin{lstlisting}
stage ID:
  new splat32 = _ZX_1(splat32);
stage RR:
  new argument3 = splat32 ? (_WX_32(upper27_lower5) << 32) | _ZX_32(_WX_32(upper27_lower5)) : _SX_32(_WX_32(upper27_lower5));
  new argument2 = PGR[registerP];
  for (I = 0; I < 8; I += 1) {
    result1.16[I] = (_SX_16(argument3.16[I]) + _SX_16(argument2.16[I])) >> 1
  };
stage E1:
  PGR[registerM] = result1;
\end{lstlisting}}
\newpage
\subsubsection[{\small AVGRBX: Average Rounded Byte Hexadecuple}]{AVGRBX\hfill Average Rounded Byte Hexadecuple}
\label{instr:AVGRBX}\index{avgrbx}
 
\paragraph{Encoding} Format \textsf{ALU\_BXAVG} (Section~\ref{sec:format-ALU-BXAVG}), \verb|exucode2 = 1110|\\

\bitfields{ 1/P, 2/11, 1/1, 4/1110, 5/registerM, 1/--, 2/10, 3/001, 1/1, 5/registerO, 1/--, 5/registerP, 1/1 }


\paragraph{Syntax} \texttt{avgrbx} \emph{registerM} = \emph{registerP}, \emph{registerO}

\paragraph{Description} The \emph{registerO} and the \emph{registerP} are interpreted as sixteen bytes packed into 128 bits. The \emph{registerO} bytes and the \emph{registerP} bytes are added. Each sum is incremented and shifted right arithmetic one bit. The sixteen resulting bytes are packed and stored into the \emph{registerM}.


\paragraph{Execution} {Reservation table \textsf{ALU\_LITE} (Section~\ref{sec:reservation-ALU-LITE})}
~
{\begin{lstlisting}
stage RR:
  new argument3 = PGR[registerO];
  new argument2 = PGR[registerP];
  for (I = 0; I < 16; I += 1) {
    result1.8[I] = ((_SX_8(argument3.8[I]) + _SX_8(argument2.8[I])) + 1) >> 1
  };
stage E1:
  PGR[registerM] = result1;
\end{lstlisting}}
\newpage

\paragraph{Encoding} Format \textsf{ALU\_BXAVG.M} (Section~\ref{sec:format-ALU-BXAVG.M}), \verb|exucode2 = 1110|\\

\bitfields{ 1/P, 2/00, 2/-- --, 27/upper27, 1/1, 2/11, 1/1, 4/1110, 5/registerM, 1/--, 2/10, 3/001, 1/1, 1/splat32, 5/lower5, 5/registerP, 1/1 }


\paragraph{Syntax} \texttt{avgrbx} \emph{registerM} = \emph{registerP}, \emph{upper27\_\discretionary{}{}{}lower5}\emph{splat32}

\paragraph{Description} The \emph{upper27\_\discretionary{}{}{}lower5} and the \emph{registerP} are interpreted as sixteen bytes packed into 128 bits. The \emph{upper27\_\discretionary{}{}{}lower5} bytes and the \emph{registerP} bytes are added. Each sum is incremented and shifted right arithmetic one bit. The sixteen resulting bytes are packed and stored into the \emph{registerM}.


\paragraph{Modifier(s)}
\noindent\begin{itemize}
\item splat32 (cf. splat32 in section \ref{par:modifier-splat32} on page \pageref{par:modifier-splat32})
\end{itemize}

\paragraph{Execution} {Reservation table \textsf{ALU\_LITE.X} (Section~\ref{sec:reservation-ALU-LITE.X})}
~
{\begin{lstlisting}
stage ID:
  new splat32 = _ZX_1(splat32);
stage RR:
  new argument3 = splat32 ? (_WX_32(upper27_lower5) << 32) | _ZX_32(_WX_32(upper27_lower5)) : _SX_32(_WX_32(upper27_lower5));
  new argument2 = PGR[registerP];
  for (I = 0; I < 16; I += 1) {
    result1.8[I] = ((_SX_8(argument3.8[I]) + _SX_8(argument2.8[I])) + 1) >> 1
  };
stage E1:
  PGR[registerM] = result1;
\end{lstlisting}}
\newpage
\subsubsection[{\small AVGRHO: Average Rounded Half Word Octuple}]{AVGRHO\hfill Average Rounded Half Word Octuple}
\label{instr:AVGRHO}\index{avgrho}
 
\paragraph{Encoding} Format \textsf{ALU\_HOAVG} (Section~\ref{sec:format-ALU-HOAVG}), \verb|exucode2 = 1110|\\

\bitfields{ 1/P, 2/11, 1/1, 4/1110, 5/registerM, 1/--, 2/10, 3/001, 1/1, 5/registerO, 1/--, 5/registerP, 1/0 }


\paragraph{Syntax} \texttt{avgrho} \emph{registerM} = \emph{registerP}, \emph{registerO}

\paragraph{Description} The \emph{registerO} and the \emph{registerP} are interpreted as eight half words packed into 128 bits. The \emph{registerO} half words and the \emph{registerP} half words are added. Each sum is incremented and shifted right arithmetic one bit. The eight resulting half words are packed and stored into the \emph{registerM}.


\paragraph{Execution} {Reservation table \textsf{ALU\_LITE} (Section~\ref{sec:reservation-ALU-LITE})}
~
{\begin{lstlisting}
stage RR:
  new argument3 = PGR[registerO];
  new argument2 = PGR[registerP];
  for (I = 0; I < 8; I += 1) {
    result1.16[I] = ((_SX_16(argument3.16[I]) + _SX_16(argument2.16[I])) + 1) >> 1
  };
stage E1:
  PGR[registerM] = result1;
\end{lstlisting}}
\newpage

\paragraph{Encoding} Format \textsf{ALU\_HOAVG.M} (Section~\ref{sec:format-ALU-HOAVG.M}), \verb|exucode2 = 1110|\\

\bitfields{ 1/P, 2/00, 2/-- --, 27/upper27, 1/1, 2/11, 1/1, 4/1110, 5/registerM, 1/--, 2/10, 3/001, 1/1, 1/splat32, 5/lower5, 5/registerP, 1/0 }


\paragraph{Syntax} \texttt{avgrho} \emph{registerM} = \emph{registerP}, \emph{upper27\_\discretionary{}{}{}lower5}\emph{splat32}

\paragraph{Description} The \emph{upper27\_\discretionary{}{}{}lower5} and the \emph{registerP} are interpreted as eight half words packed into 128 bits. The \emph{upper27\_\discretionary{}{}{}lower5} half words and the \emph{registerP} half words are added. Each sum is incremented and shifted right arithmetic one bit. The eight resulting half words are packed and stored into the \emph{registerM}.


\paragraph{Modifier(s)}
\noindent\begin{itemize}
\item splat32 (cf. splat32 in section \ref{par:modifier-splat32} on page \pageref{par:modifier-splat32})
\end{itemize}

\paragraph{Execution} {Reservation table \textsf{ALU\_LITE.X} (Section~\ref{sec:reservation-ALU-LITE.X})}
~
{\begin{lstlisting}
stage ID:
  new splat32 = _ZX_1(splat32);
stage RR:
  new argument3 = splat32 ? (_WX_32(upper27_lower5) << 32) | _ZX_32(_WX_32(upper27_lower5)) : _SX_32(_WX_32(upper27_lower5));
  new argument2 = PGR[registerP];
  for (I = 0; I < 8; I += 1) {
    result1.16[I] = ((_SX_16(argument3.16[I]) + _SX_16(argument2.16[I])) + 1) >> 1
  };
stage E1:
  PGR[registerM] = result1;
\end{lstlisting}}
\newpage
\subsubsection[{\small AVGRUBX: Average Rounded Unsigned Byte Hexadecuple}]{AVGRUBX\hfill Average Rounded Unsigned Byte Hexadecuple}
\label{instr:AVGRUBX}\index{avgrubx}
 
\paragraph{Encoding} Format \textsf{ALU\_BXAVG} (Section~\ref{sec:format-ALU-BXAVG}), \verb|exucode2 = 1111|\\

\bitfields{ 1/P, 2/11, 1/1, 4/1111, 5/registerM, 1/--, 2/10, 3/001, 1/1, 5/registerO, 1/--, 5/registerP, 1/1 }


\paragraph{Syntax} \texttt{avgrubx} \emph{registerM} = \emph{registerP}, \emph{registerO}

\paragraph{Description} The \emph{registerO} and the \emph{registerP} are interpreted as sixteen bytes packed into 128 bits. The \emph{registerO} bytes and the \emph{registerP} bytes are added. Each sum is incremented and shifted right logical one bit. The sixteen resulting bytes are packed and stored into the \emph{registerM}.


\paragraph{Execution} {Reservation table \textsf{ALU\_LITE} (Section~\ref{sec:reservation-ALU-LITE})}
~
{\begin{lstlisting}
stage RR:
  new argument3 = PGR[registerO];
  new argument2 = PGR[registerP];
  for (I = 0; I < 16; I += 1) {
    result1.8[I] = ((_ZX_8(argument3.8[I]) + _ZX_8(argument2.8[I])) + 1) >> 1
  };
stage E1:
  PGR[registerM] = result1;
\end{lstlisting}}
\newpage

\paragraph{Encoding} Format \textsf{ALU\_BXAVG.M} (Section~\ref{sec:format-ALU-BXAVG.M}), \verb|exucode2 = 1111|\\

\bitfields{ 1/P, 2/00, 2/-- --, 27/upper27, 1/1, 2/11, 1/1, 4/1111, 5/registerM, 1/--, 2/10, 3/001, 1/1, 1/splat32, 5/lower5, 5/registerP, 1/1 }


\paragraph{Syntax} \texttt{avgrubx} \emph{registerM} = \emph{registerP}, \emph{upper27\_\discretionary{}{}{}lower5}\emph{splat32}

\paragraph{Description} The \emph{upper27\_\discretionary{}{}{}lower5} and the \emph{registerP} are interpreted as sixteen bytes packed into 128 bits. The \emph{upper27\_\discretionary{}{}{}lower5} bytes and the \emph{registerP} bytes are added. Each sum is incremented and shifted right logical one bit. The sixteen resulting bytes are packed and stored into the \emph{registerM}.


\paragraph{Modifier(s)}
\noindent\begin{itemize}
\item splat32 (cf. splat32 in section \ref{par:modifier-splat32} on page \pageref{par:modifier-splat32})
\end{itemize}

\paragraph{Execution} {Reservation table \textsf{ALU\_LITE.X} (Section~\ref{sec:reservation-ALU-LITE.X})}
~
{\begin{lstlisting}
stage ID:
  new splat32 = _ZX_1(splat32);
stage RR:
  new argument3 = splat32 ? (_WX_32(upper27_lower5) << 32) | _ZX_32(_WX_32(upper27_lower5)) : _SX_32(_WX_32(upper27_lower5));
  new argument2 = PGR[registerP];
  for (I = 0; I < 16; I += 1) {
    result1.8[I] = ((_ZX_8(argument3.8[I]) + _ZX_8(argument2.8[I])) + 1) >> 1
  };
stage E1:
  PGR[registerM] = result1;
\end{lstlisting}}
\newpage
\subsubsection[{\small AVGRUHO: Average Rounded Unsigned Half Word Octuple}]{AVGRUHO\hfill Average Rounded Unsigned Half Word Octuple}
\label{instr:AVGRUHO}\index{avgruho}
 
\paragraph{Encoding} Format \textsf{ALU\_HOAVG} (Section~\ref{sec:format-ALU-HOAVG}), \verb|exucode2 = 1111|\\

\bitfields{ 1/P, 2/11, 1/1, 4/1111, 5/registerM, 1/--, 2/10, 3/001, 1/1, 5/registerO, 1/--, 5/registerP, 1/0 }


\paragraph{Syntax} \texttt{avgruho} \emph{registerM} = \emph{registerP}, \emph{registerO}

\paragraph{Description} The \emph{registerO} and the \emph{registerP} are interpreted as eight half words packed into 128 bits. The \emph{registerO} half words and the \emph{registerP} half words are added. Each sum is incremented and shifted right logical one bit. The eight resulting half words are packed and stored into the \emph{registerM}.


\paragraph{Execution} {Reservation table \textsf{ALU\_LITE} (Section~\ref{sec:reservation-ALU-LITE})}
~
{\begin{lstlisting}
stage RR:
  new argument3 = PGR[registerO];
  new argument2 = PGR[registerP];
  for (I = 0; I < 8; I += 1) {
    result1.16[I] = ((_ZX_16(argument3.16[I]) + _ZX_16(argument2.16[I])) + 1) >> 1
  };
stage E1:
  PGR[registerM] = result1;
\end{lstlisting}}
\newpage

\paragraph{Encoding} Format \textsf{ALU\_HOAVG.M} (Section~\ref{sec:format-ALU-HOAVG.M}), \verb|exucode2 = 1111|\\

\bitfields{ 1/P, 2/00, 2/-- --, 27/upper27, 1/1, 2/11, 1/1, 4/1111, 5/registerM, 1/--, 2/10, 3/001, 1/1, 1/splat32, 5/lower5, 5/registerP, 1/0 }


\paragraph{Syntax} \texttt{avgruho} \emph{registerM} = \emph{registerP}, \emph{upper27\_\discretionary{}{}{}lower5}\emph{splat32}

\paragraph{Description} The \emph{upper27\_\discretionary{}{}{}lower5} and the \emph{registerP} are interpreted as eight half words packed into 128 bits. The \emph{upper27\_\discretionary{}{}{}lower5} half words and the \emph{registerP} half words are added. Each sum is incremented and shifted right logical one bit. The eight resulting half words are packed and stored into the \emph{registerM}.


\paragraph{Modifier(s)}
\noindent\begin{itemize}
\item splat32 (cf. splat32 in section \ref{par:modifier-splat32} on page \pageref{par:modifier-splat32})
\end{itemize}

\paragraph{Execution} {Reservation table \textsf{ALU\_LITE.X} (Section~\ref{sec:reservation-ALU-LITE.X})}
~
{\begin{lstlisting}
stage ID:
  new splat32 = _ZX_1(splat32);
stage RR:
  new argument3 = splat32 ? (_WX_32(upper27_lower5) << 32) | _ZX_32(_WX_32(upper27_lower5)) : _SX_32(_WX_32(upper27_lower5));
  new argument2 = PGR[registerP];
  for (I = 0; I < 8; I += 1) {
    result1.16[I] = ((_ZX_16(argument3.16[I]) + _ZX_16(argument2.16[I])) + 1) >> 1
  };
stage E1:
  PGR[registerM] = result1;
\end{lstlisting}}
\newpage
\subsubsection[{\small AVGRUW: Average Rounded Unsigned Word}]{AVGRUW\hfill Average Rounded Unsigned Word}
\label{instr:AVGRUW}\index{avgruw}
 
\paragraph{Encoding} Format \textsf{ALU\_WAVG} (Section~\ref{sec:format-ALU-WAVG}), \verb|exucode2 = 1111|\\

\bitfields{ 1/P, 2/11, 1/0, 4/1111, 6/registerW, 2/11, 3/001, 1/signextw, 6/registerY, 6/registerZ }


\paragraph{Syntax} \texttt{avgruw}\emph{signextw} \emph{registerW} = \emph{registerZ}, \emph{registerY}

\paragraph{Description} The \emph{registerY} and the \emph{registerZ} are added. The result is incremented, shifted right logical one bit and stored into the \emph{registerW}.


\paragraph{Modifier(s)}
\noindent\begin{itemize}
\item signextw (cf. signextw in section \ref{par:modifier-signextw} on page \pageref{par:modifier-signextw})
\end{itemize}

\paragraph{Execution} {Reservation table \textsf{ALU\_TINY} (Section~\ref{sec:reservation-ALU-TINY})}
~
{\begin{lstlisting}
stage ID:
  new signextw = _ZX_1(signextw);
stage RR:
  new argument3 = GPR[registerY];
  new argument2 = GPR[registerZ];
  new result1 = ((argument3.32[0] + argument2.32[0]) + 1) >> 1;
stage E1:
  GPR[registerW] = signextw ? _SX_32(result1) : _ZX_32(result1);
\end{lstlisting}}
\newpage

\paragraph{Encoding} Format \textsf{ALU\_WAVG.W} (Section~\ref{sec:format-ALU-WAVG.W}), \verb|exucode2 = 1111|\\

\bitfields{ 1/P, 2/00, 2/-- --, 27/upper27, 1/1, 2/11, 1/0, 4/1111, 6/registerW, 2/11, 3/001, 1/signextw, 1/0, 5/lower5, 6/registerZ }


\paragraph{Syntax} \texttt{avgruw}\emph{signextw} \emph{registerW} = \emph{registerZ}, \emph{upper27\_\discretionary{}{}{}lower5}

\paragraph{Description} The \emph{upper27\_\discretionary{}{}{}lower5} and the \emph{registerZ} are added. The result is incremented, shifted right logical one bit and stored into the \emph{registerW}.


\paragraph{Modifier(s)}
\noindent\begin{itemize}
\item signextw (cf. signextw in section \ref{par:modifier-signextw} on page \pageref{par:modifier-signextw})
\end{itemize}

\paragraph{Execution} {Reservation table \textsf{ALU\_TINY.X} (Section~\ref{sec:reservation-ALU-TINY.X})}
~
{\begin{lstlisting}
stage ID:
  new signextw = _ZX_1(signextw);
stage RR:
  new argument3 = _WX_32(upper27_lower5);
  new argument2 = GPR[registerZ];
  new result1 = ((argument3.32[0] + argument2.32[0]) + 1) >> 1;
stage E1:
  GPR[registerW] = signextw ? _SX_32(result1) : _ZX_32(result1);
\end{lstlisting}}
\newpage
\subsubsection[{\small AVGRUWQ: Average Rounded Unsigned Word Quadruple}]{AVGRUWQ\hfill Average Rounded Unsigned Word Quadruple}
\label{instr:AVGRUWQ}\index{avgruwq}
 
\paragraph{Encoding} Format \textsf{ALU\_WQAVG} (Section~\ref{sec:format-ALU-WQAVG}), \verb|exucode2 = 1111|\\

\bitfields{ 1/P, 2/11, 1/1, 4/1111, 5/registerM, 1/--, 2/10, 3/001, 1/0, 5/registerO, 1/--, 5/registerP, 1/1 }


\paragraph{Syntax} \texttt{avgruwq} \emph{registerM} = \emph{registerP}, \emph{registerO}

\paragraph{Description} The \emph{registerO} and the \emph{registerP} are interpreted as four words packed into 128 bits. The \emph{registerO} words and the \emph{registerP} words are added. Each sum is incremented and shifted right logical one bit. The four resulting words are packed and stored into the \emph{registerM}.


\paragraph{Execution} {Reservation table \textsf{ALU\_LITE} (Section~\ref{sec:reservation-ALU-LITE})}
~
{\begin{lstlisting}
stage RR:
  new argument3 = PGR[registerO];
  new argument2 = PGR[registerP];
  for (I = 0; I < 4; I += 1) {
    result1.32[I] = ((_ZX_32(argument3.32[I]) + _ZX_32(argument2.32[I])) + 1) >> 1
  };
stage E1:
  PGR[registerM] = result1;
\end{lstlisting}}
\newpage

\paragraph{Encoding} Format \textsf{ALU\_WQAVG.M} (Section~\ref{sec:format-ALU-WQAVG.M}), \verb|exucode2 = 1111|\\

\bitfields{ 1/P, 2/00, 2/-- --, 27/upper27, 1/1, 2/11, 1/1, 4/1111, 5/registerM, 1/--, 2/10, 3/001, 1/0, 1/splat32, 5/lower5, 5/registerP, 1/1 }


\paragraph{Syntax} \texttt{avgruwq} \emph{registerM} = \emph{registerP}, \emph{upper27\_\discretionary{}{}{}lower5}\emph{splat32}

\paragraph{Description} The \emph{upper27\_\discretionary{}{}{}lower5} and the \emph{registerP} are interpreted as four words packed into 128 bits. The \emph{upper27\_\discretionary{}{}{}lower5} words and the \emph{registerP} words are added. Each sum is incremented and shifted right logical one bit. The four resulting words are packed and stored into the \emph{registerM}.


\paragraph{Modifier(s)}
\noindent\begin{itemize}
\item splat32 (cf. splat32 in section \ref{par:modifier-splat32} on page \pageref{par:modifier-splat32})
\end{itemize}

\paragraph{Execution} {Reservation table \textsf{ALU\_LITE.X} (Section~\ref{sec:reservation-ALU-LITE.X})}
~
{\begin{lstlisting}
stage ID:
  new splat32 = _ZX_1(splat32);
stage RR:
  new argument3 = splat32 ? (_WX_32(upper27_lower5) << 32) | _ZX_32(_WX_32(upper27_lower5)) : _SX_32(_WX_32(upper27_lower5));
  new argument2 = PGR[registerP];
  for (I = 0; I < 4; I += 1) {
    result1.32[I] = ((_ZX_32(argument3.32[I]) + _ZX_32(argument2.32[I])) + 1) >> 1
  };
stage E1:
  PGR[registerM] = result1;
\end{lstlisting}}
\newpage
\subsubsection[{\small AVGRW: Average Rounded Word}]{AVGRW\hfill Average Rounded Word}
\label{instr:AVGRW}\index{avgrw}
 
\paragraph{Encoding} Format \textsf{ALU\_WAVG} (Section~\ref{sec:format-ALU-WAVG}), \verb|exucode2 = 1110|\\

\bitfields{ 1/P, 2/11, 1/0, 4/1110, 6/registerW, 2/11, 3/001, 1/signextw, 6/registerY, 6/registerZ }


\paragraph{Syntax} \texttt{avgrw}\emph{signextw} \emph{registerW} = \emph{registerZ}, \emph{registerY}

\paragraph{Description} The \emph{registerY} and the \emph{registerZ} are added. The result is incremented, shifted right arithmetic one bit and stored into the \emph{registerW}.


\paragraph{Modifier(s)}
\noindent\begin{itemize}
\item signextw (cf. signextw in section \ref{par:modifier-signextw} on page \pageref{par:modifier-signextw})
\end{itemize}

\paragraph{Execution} {Reservation table \textsf{ALU\_TINY} (Section~\ref{sec:reservation-ALU-TINY})}
~
{\begin{lstlisting}
stage ID:
  new signextw = _ZX_1(signextw);
stage RR:
  new argument3 = GPR[registerY];
  new argument2 = GPR[registerZ];
  new result1 = ((_SX_32(argument3) + _SX_32(argument2)) + 1) >> 1;
stage E1:
  GPR[registerW] = signextw ? _SX_32(result1) : _ZX_32(result1);
\end{lstlisting}}
\newpage

\paragraph{Encoding} Format \textsf{ALU\_WAVG.W} (Section~\ref{sec:format-ALU-WAVG.W}), \verb|exucode2 = 1110|\\

\bitfields{ 1/P, 2/00, 2/-- --, 27/upper27, 1/1, 2/11, 1/0, 4/1110, 6/registerW, 2/11, 3/001, 1/signextw, 1/0, 5/lower5, 6/registerZ }


\paragraph{Syntax} \texttt{avgrw}\emph{signextw} \emph{registerW} = \emph{registerZ}, \emph{upper27\_\discretionary{}{}{}lower5}

\paragraph{Description} The \emph{upper27\_\discretionary{}{}{}lower5} and the \emph{registerZ} are added. The result is incremented, shifted right arithmetic one bit and stored into the \emph{registerW}.


\paragraph{Modifier(s)}
\noindent\begin{itemize}
\item signextw (cf. signextw in section \ref{par:modifier-signextw} on page \pageref{par:modifier-signextw})
\end{itemize}

\paragraph{Execution} {Reservation table \textsf{ALU\_TINY.X} (Section~\ref{sec:reservation-ALU-TINY.X})}
~
{\begin{lstlisting}
stage ID:
  new signextw = _ZX_1(signextw);
stage RR:
  new argument3 = _WX_32(upper27_lower5);
  new argument2 = GPR[registerZ];
  new result1 = ((_SX_32(argument3) + _SX_32(argument2)) + 1) >> 1;
stage E1:
  GPR[registerW] = signextw ? _SX_32(result1) : _ZX_32(result1);
\end{lstlisting}}
\newpage
\subsubsection[{\small AVGRWQ: Average Rounded Word Quadruple}]{AVGRWQ\hfill Average Rounded Word Quadruple}
\label{instr:AVGRWQ}\index{avgrwq}
 
\paragraph{Encoding} Format \textsf{ALU\_WQAVG} (Section~\ref{sec:format-ALU-WQAVG}), \verb|exucode2 = 1110|\\

\bitfields{ 1/P, 2/11, 1/1, 4/1110, 5/registerM, 1/--, 2/10, 3/001, 1/0, 5/registerO, 1/--, 5/registerP, 1/1 }


\paragraph{Syntax} \texttt{avgrwq} \emph{registerM} = \emph{registerP}, \emph{registerO}

\paragraph{Description} The \emph{registerO} and the \emph{registerP} are interpreted as four words packed into 128 bits. The \emph{registerO} words and the \emph{registerP} words are added. Each sum is incremented and shifted right arithmetic one bit. The four resulting words are packed and stored into the \emph{registerM}.


\paragraph{Execution} {Reservation table \textsf{ALU\_LITE} (Section~\ref{sec:reservation-ALU-LITE})}
~
{\begin{lstlisting}
stage RR:
  new argument3 = PGR[registerO];
  new argument2 = PGR[registerP];
  for (I = 0; I < 4; I += 1) {
    result1.32[I] = ((_SX_32(argument3.32[I]) + _SX_32(argument2.32[I])) + 1) >> 1
  };
stage E1:
  PGR[registerM] = result1;
\end{lstlisting}}
\newpage

\paragraph{Encoding} Format \textsf{ALU\_WQAVG.M} (Section~\ref{sec:format-ALU-WQAVG.M}), \verb|exucode2 = 1110|\\

\bitfields{ 1/P, 2/00, 2/-- --, 27/upper27, 1/1, 2/11, 1/1, 4/1110, 5/registerM, 1/--, 2/10, 3/001, 1/0, 1/splat32, 5/lower5, 5/registerP, 1/1 }


\paragraph{Syntax} \texttt{avgrwq} \emph{registerM} = \emph{registerP}, \emph{upper27\_\discretionary{}{}{}lower5}\emph{splat32}

\paragraph{Description} The \emph{upper27\_\discretionary{}{}{}lower5} and the \emph{registerP} are interpreted as four words packed into 128 bits. The \emph{upper27\_\discretionary{}{}{}lower5} words and the \emph{registerP} words are added. Each sum is incremented and shifted right arithmetic one bit. The four resulting words are packed and stored into the \emph{registerM}.


\paragraph{Modifier(s)}
\noindent\begin{itemize}
\item splat32 (cf. splat32 in section \ref{par:modifier-splat32} on page \pageref{par:modifier-splat32})
\end{itemize}

\paragraph{Execution} {Reservation table \textsf{ALU\_LITE.X} (Section~\ref{sec:reservation-ALU-LITE.X})}
~
{\begin{lstlisting}
stage ID:
  new splat32 = _ZX_1(splat32);
stage RR:
  new argument3 = splat32 ? (_WX_32(upper27_lower5) << 32) | _ZX_32(_WX_32(upper27_lower5)) : _SX_32(_WX_32(upper27_lower5));
  new argument2 = PGR[registerP];
  for (I = 0; I < 4; I += 1) {
    result1.32[I] = ((_SX_32(argument3.32[I]) + _SX_32(argument2.32[I])) + 1) >> 1
  };
stage E1:
  PGR[registerM] = result1;
\end{lstlisting}}
\newpage
\subsubsection[{\small AVGUBX: Average Unsigned Byte Hexadecuple}]{AVGUBX\hfill Average Unsigned Byte Hexadecuple}
\label{instr:AVGUBX}\index{avgubx}
 
\paragraph{Encoding} Format \textsf{ALU\_BXAVG} (Section~\ref{sec:format-ALU-BXAVG}), \verb|exucode2 = 1101|\\

\bitfields{ 1/P, 2/11, 1/1, 4/1101, 5/registerM, 1/--, 2/10, 3/001, 1/1, 5/registerO, 1/--, 5/registerP, 1/1 }


\paragraph{Syntax} \texttt{avgubx} \emph{registerM} = \emph{registerP}, \emph{registerO}

\paragraph{Description} The \emph{registerO} and the \emph{registerP} are interpreted as sixteen bytes packed into 128 bits. The \emph{registerO} bytes and the \emph{registerP} bytes are added. Each sum is shifted right logical one bit. The sixteen resulting bytes are packed and stored into the \emph{registerM}.


\paragraph{Execution} {Reservation table \textsf{ALU\_LITE} (Section~\ref{sec:reservation-ALU-LITE})}
~
{\begin{lstlisting}
stage RR:
  new argument3 = PGR[registerO];
  new argument2 = PGR[registerP];
  for (I = 0; I < 16; I += 1) {
    result1.8[I] = (_ZX_8(argument3.8[I]) + _ZX_8(argument2.8[I])) >> 1
  };
stage E1:
  PGR[registerM] = result1;
\end{lstlisting}}
\newpage

\paragraph{Encoding} Format \textsf{ALU\_BXAVG.M} (Section~\ref{sec:format-ALU-BXAVG.M}), \verb|exucode2 = 1101|\\

\bitfields{ 1/P, 2/00, 2/-- --, 27/upper27, 1/1, 2/11, 1/1, 4/1101, 5/registerM, 1/--, 2/10, 3/001, 1/1, 1/splat32, 5/lower5, 5/registerP, 1/1 }


\paragraph{Syntax} \texttt{avgubx} \emph{registerM} = \emph{registerP}, \emph{upper27\_\discretionary{}{}{}lower5}\emph{splat32}

\paragraph{Description} The \emph{upper27\_\discretionary{}{}{}lower5} and the \emph{registerP} are interpreted as sixteen bytes packed into 128 bits. The \emph{upper27\_\discretionary{}{}{}lower5} bytes and the \emph{registerP} bytes are added. Each sum is shifted right logical one bit. The sixteen resulting bytes are packed and stored into the \emph{registerM}.


\paragraph{Modifier(s)}
\noindent\begin{itemize}
\item splat32 (cf. splat32 in section \ref{par:modifier-splat32} on page \pageref{par:modifier-splat32})
\end{itemize}

\paragraph{Execution} {Reservation table \textsf{ALU\_LITE.X} (Section~\ref{sec:reservation-ALU-LITE.X})}
~
{\begin{lstlisting}
stage ID:
  new splat32 = _ZX_1(splat32);
stage RR:
  new argument3 = splat32 ? (_WX_32(upper27_lower5) << 32) | _ZX_32(_WX_32(upper27_lower5)) : _SX_32(_WX_32(upper27_lower5));
  new argument2 = PGR[registerP];
  for (I = 0; I < 16; I += 1) {
    result1.8[I] = (_ZX_8(argument3.8[I]) + _ZX_8(argument2.8[I])) >> 1
  };
stage E1:
  PGR[registerM] = result1;
\end{lstlisting}}
\newpage
\subsubsection[{\small AVGUHO: Average Unsigned Half Word Octuple}]{AVGUHO\hfill Average Unsigned Half Word Octuple}
\label{instr:AVGUHO}\index{avguho}
 
\paragraph{Encoding} Format \textsf{ALU\_HOAVG} (Section~\ref{sec:format-ALU-HOAVG}), \verb|exucode2 = 1101|\\

\bitfields{ 1/P, 2/11, 1/1, 4/1101, 5/registerM, 1/--, 2/10, 3/001, 1/1, 5/registerO, 1/--, 5/registerP, 1/0 }


\paragraph{Syntax} \texttt{avguho} \emph{registerM} = \emph{registerP}, \emph{registerO}

\paragraph{Description} The \emph{registerO} and the \emph{registerP} are interpreted as eight half words packed into 128 bits. The \emph{registerO} half words and the \emph{registerP} half words are added. Each sum is shifted right logical one bit. The eight resulting half words are packed and stored into the \emph{registerM}.


\paragraph{Execution} {Reservation table \textsf{ALU\_LITE} (Section~\ref{sec:reservation-ALU-LITE})}
~
{\begin{lstlisting}
stage RR:
  new argument3 = PGR[registerO];
  new argument2 = PGR[registerP];
  for (I = 0; I < 8; I += 1) {
    result1.16[I] = (_ZX_16(argument3.16[I]) + _ZX_16(argument2.16[I])) >> 1
  };
stage E1:
  PGR[registerM] = result1;
\end{lstlisting}}
\newpage

\paragraph{Encoding} Format \textsf{ALU\_HOAVG.M} (Section~\ref{sec:format-ALU-HOAVG.M}), \verb|exucode2 = 1101|\\

\bitfields{ 1/P, 2/00, 2/-- --, 27/upper27, 1/1, 2/11, 1/1, 4/1101, 5/registerM, 1/--, 2/10, 3/001, 1/1, 1/splat32, 5/lower5, 5/registerP, 1/0 }


\paragraph{Syntax} \texttt{avguho} \emph{registerM} = \emph{registerP}, \emph{upper27\_\discretionary{}{}{}lower5}\emph{splat32}

\paragraph{Description} The \emph{upper27\_\discretionary{}{}{}lower5} and the \emph{registerP} are interpreted as eight half words packed into 128 bits. The \emph{upper27\_\discretionary{}{}{}lower5} half words and the \emph{registerP} half words are added. Each sum is shifted right logical one bit. The eight resulting half words are packed and stored into the \emph{registerM}.


\paragraph{Modifier(s)}
\noindent\begin{itemize}
\item splat32 (cf. splat32 in section \ref{par:modifier-splat32} on page \pageref{par:modifier-splat32})
\end{itemize}

\paragraph{Execution} {Reservation table \textsf{ALU\_LITE.X} (Section~\ref{sec:reservation-ALU-LITE.X})}
~
{\begin{lstlisting}
stage ID:
  new splat32 = _ZX_1(splat32);
stage RR:
  new argument3 = splat32 ? (_WX_32(upper27_lower5) << 32) | _ZX_32(_WX_32(upper27_lower5)) : _SX_32(_WX_32(upper27_lower5));
  new argument2 = PGR[registerP];
  for (I = 0; I < 8; I += 1) {
    result1.16[I] = (_ZX_16(argument3.16[I]) + _ZX_16(argument2.16[I])) >> 1
  };
stage E1:
  PGR[registerM] = result1;
\end{lstlisting}}
\newpage
\subsubsection[{\small AVGUW: Average Unsigned Word}]{AVGUW\hfill Average Unsigned Word}
\label{instr:AVGUW}\index{avguw}
 
\paragraph{Encoding} Format \textsf{ALU\_WAVG} (Section~\ref{sec:format-ALU-WAVG}), \verb|exucode2 = 1101|\\

\bitfields{ 1/P, 2/11, 1/0, 4/1101, 6/registerW, 2/11, 3/001, 1/signextw, 6/registerY, 6/registerZ }


\paragraph{Syntax} \texttt{avguw}\emph{signextw} \emph{registerW} = \emph{registerZ}, \emph{registerY}

\paragraph{Description} The \emph{registerY} and the \emph{registerZ} are added. The result is shifted right logical one bit and stored into the \emph{registerW}.


\paragraph{Modifier(s)}
\noindent\begin{itemize}
\item signextw (cf. signextw in section \ref{par:modifier-signextw} on page \pageref{par:modifier-signextw})
\end{itemize}

\paragraph{Execution} {Reservation table \textsf{ALU\_TINY} (Section~\ref{sec:reservation-ALU-TINY})}
~
{\begin{lstlisting}
stage ID:
  new signextw = _ZX_1(signextw);
stage RR:
  new argument3 = GPR[registerY];
  new argument2 = GPR[registerZ];
  new result1 = (argument3.32[0] + argument2.32[0]) >> 1;
stage E1:
  GPR[registerW] = signextw ? _SX_32(result1) : _ZX_32(result1);
\end{lstlisting}}
\newpage

\paragraph{Encoding} Format \textsf{ALU\_WAVG.W} (Section~\ref{sec:format-ALU-WAVG.W}), \verb|exucode2 = 1101|\\

\bitfields{ 1/P, 2/00, 2/-- --, 27/upper27, 1/1, 2/11, 1/0, 4/1101, 6/registerW, 2/11, 3/001, 1/signextw, 1/0, 5/lower5, 6/registerZ }


\paragraph{Syntax} \texttt{avguw}\emph{signextw} \emph{registerW} = \emph{registerZ}, \emph{upper27\_\discretionary{}{}{}lower5}

\paragraph{Description} The \emph{upper27\_\discretionary{}{}{}lower5} and the \emph{registerZ} are added. The result is shifted right logical one bit and stored into the \emph{registerW}.


\paragraph{Modifier(s)}
\noindent\begin{itemize}
\item signextw (cf. signextw in section \ref{par:modifier-signextw} on page \pageref{par:modifier-signextw})
\end{itemize}

\paragraph{Execution} {Reservation table \textsf{ALU\_TINY.X} (Section~\ref{sec:reservation-ALU-TINY.X})}
~
{\begin{lstlisting}
stage ID:
  new signextw = _ZX_1(signextw);
stage RR:
  new argument3 = _WX_32(upper27_lower5);
  new argument2 = GPR[registerZ];
  new result1 = (argument3.32[0] + argument2.32[0]) >> 1;
stage E1:
  GPR[registerW] = signextw ? _SX_32(result1) : _ZX_32(result1);
\end{lstlisting}}
\newpage
\subsubsection[{\small AVGUWQ: Average Unsigned Word Quadruple}]{AVGUWQ\hfill Average Unsigned Word Quadruple}
\label{instr:AVGUWQ}\index{avguwq}
 
\paragraph{Encoding} Format \textsf{ALU\_WQAVG} (Section~\ref{sec:format-ALU-WQAVG}), \verb|exucode2 = 1101|\\

\bitfields{ 1/P, 2/11, 1/1, 4/1101, 5/registerM, 1/--, 2/10, 3/001, 1/0, 5/registerO, 1/--, 5/registerP, 1/1 }


\paragraph{Syntax} \texttt{avguwq} \emph{registerM} = \emph{registerP}, \emph{registerO}

\paragraph{Description} The \emph{registerO} and the \emph{registerP} are interpreted as four words packed into 128 bits. The \emph{registerO} words and the \emph{registerP} words are added. Each sum is shifted right logical one bit. The four resulting words are packed and stored into the \emph{registerM}.


\paragraph{Execution} {Reservation table \textsf{ALU\_LITE} (Section~\ref{sec:reservation-ALU-LITE})}
~
{\begin{lstlisting}
stage RR:
  new argument3 = PGR[registerO];
  new argument2 = PGR[registerP];
  for (I = 0; I < 4; I += 1) {
    result1.32[I] = (_ZX_32(argument3.32[I]) + _ZX_32(argument2.32[I])) >> 1
  };
stage E1:
  PGR[registerM] = result1;
\end{lstlisting}}
\newpage

\paragraph{Encoding} Format \textsf{ALU\_WQAVG.M} (Section~\ref{sec:format-ALU-WQAVG.M}), \verb|exucode2 = 1101|\\

\bitfields{ 1/P, 2/00, 2/-- --, 27/upper27, 1/1, 2/11, 1/1, 4/1101, 5/registerM, 1/--, 2/10, 3/001, 1/0, 1/splat32, 5/lower5, 5/registerP, 1/1 }


\paragraph{Syntax} \texttt{avguwq} \emph{registerM} = \emph{registerP}, \emph{upper27\_\discretionary{}{}{}lower5}\emph{splat32}

\paragraph{Description} The \emph{upper27\_\discretionary{}{}{}lower5} and the \emph{registerP} are interpreted as four words packed into 128 bits. The \emph{upper27\_\discretionary{}{}{}lower5} words and the \emph{registerP} words are added. Each sum is shifted right logical one bit. The four resulting words are packed and stored into the \emph{registerM}.


\paragraph{Modifier(s)}
\noindent\begin{itemize}
\item splat32 (cf. splat32 in section \ref{par:modifier-splat32} on page \pageref{par:modifier-splat32})
\end{itemize}

\paragraph{Execution} {Reservation table \textsf{ALU\_LITE.X} (Section~\ref{sec:reservation-ALU-LITE.X})}
~
{\begin{lstlisting}
stage ID:
  new splat32 = _ZX_1(splat32);
stage RR:
  new argument3 = splat32 ? (_WX_32(upper27_lower5) << 32) | _ZX_32(_WX_32(upper27_lower5)) : _SX_32(_WX_32(upper27_lower5));
  new argument2 = PGR[registerP];
  for (I = 0; I < 4; I += 1) {
    result1.32[I] = (_ZX_32(argument3.32[I]) + _ZX_32(argument2.32[I])) >> 1
  };
stage E1:
  PGR[registerM] = result1;
\end{lstlisting}}
\newpage
\subsubsection[{\small AVGW: Average Word}]{AVGW\hfill Average Word}
\label{instr:AVGW}\index{avgw}
 
\paragraph{Encoding} Format \textsf{ALU\_WAVG} (Section~\ref{sec:format-ALU-WAVG}), \verb|exucode2 = 1100|\\

\bitfields{ 1/P, 2/11, 1/0, 4/1100, 6/registerW, 2/11, 3/001, 1/signextw, 6/registerY, 6/registerZ }


\paragraph{Syntax} \texttt{avgw}\emph{signextw} \emph{registerW} = \emph{registerZ}, \emph{registerY}

\paragraph{Description} The \emph{registerY} and the \emph{registerZ} are added. The result is shifted right arithmetic one bit and stored into the \emph{registerW}.


\paragraph{Modifier(s)}
\noindent\begin{itemize}
\item signextw (cf. signextw in section \ref{par:modifier-signextw} on page \pageref{par:modifier-signextw})
\end{itemize}

\paragraph{Execution} {Reservation table \textsf{ALU\_TINY} (Section~\ref{sec:reservation-ALU-TINY})}
~
{\begin{lstlisting}
stage ID:
  new signextw = _ZX_1(signextw);
stage RR:
  new argument3 = GPR[registerY];
  new argument2 = GPR[registerZ];
  new result1 = (_SX_32(argument3) + _SX_32(argument2)) >> 1;
stage E1:
  GPR[registerW] = signextw ? _SX_32(result1) : _ZX_32(result1);
\end{lstlisting}}
\newpage

\paragraph{Encoding} Format \textsf{ALU\_WAVG.W} (Section~\ref{sec:format-ALU-WAVG.W}), \verb|exucode2 = 1100|\\

\bitfields{ 1/P, 2/00, 2/-- --, 27/upper27, 1/1, 2/11, 1/0, 4/1100, 6/registerW, 2/11, 3/001, 1/signextw, 1/0, 5/lower5, 6/registerZ }


\paragraph{Syntax} \texttt{avgw}\emph{signextw} \emph{registerW} = \emph{registerZ}, \emph{upper27\_\discretionary{}{}{}lower5}

\paragraph{Description} The \emph{upper27\_\discretionary{}{}{}lower5} and the \emph{registerZ} are added. The result is shifted right arithmetic one bit and stored into the \emph{registerW}.


\paragraph{Modifier(s)}
\noindent\begin{itemize}
\item signextw (cf. signextw in section \ref{par:modifier-signextw} on page \pageref{par:modifier-signextw})
\end{itemize}

\paragraph{Execution} {Reservation table \textsf{ALU\_TINY.X} (Section~\ref{sec:reservation-ALU-TINY.X})}
~
{\begin{lstlisting}
stage ID:
  new signextw = _ZX_1(signextw);
stage RR:
  new argument3 = _WX_32(upper27_lower5);
  new argument2 = GPR[registerZ];
  new result1 = (_SX_32(argument3) + _SX_32(argument2)) >> 1;
stage E1:
  GPR[registerW] = signextw ? _SX_32(result1) : _ZX_32(result1);
\end{lstlisting}}
\newpage
\subsubsection[{\small AVGWQ: Average Word Quadruple}]{AVGWQ\hfill Average Word Quadruple}
\label{instr:AVGWQ}\index{avgwq}
 
\paragraph{Encoding} Format \textsf{ALU\_WQAVG} (Section~\ref{sec:format-ALU-WQAVG}), \verb|exucode2 = 1100|\\

\bitfields{ 1/P, 2/11, 1/1, 4/1100, 5/registerM, 1/--, 2/10, 3/001, 1/0, 5/registerO, 1/--, 5/registerP, 1/1 }


\paragraph{Syntax} \texttt{avgwq} \emph{registerM} = \emph{registerP}, \emph{registerO}

\paragraph{Description} The \emph{registerO} and the \emph{registerP} are interpreted as four words packed into 128 bits. The \emph{registerO} words and the \emph{registerP} words are added. Each sum is shifted right arithmetic one bit. The four resulting words are packed and stored into the \emph{registerM}.


\paragraph{Execution} {Reservation table \textsf{ALU\_LITE} (Section~\ref{sec:reservation-ALU-LITE})}
~
{\begin{lstlisting}
stage RR:
  new argument3 = PGR[registerO];
  new argument2 = PGR[registerP];
  for (I = 0; I < 4; I += 1) {
    result1.32[I] = (_SX_32(argument3.32[I]) + _SX_32(argument2.32[I])) >> 1
  };
stage E1:
  PGR[registerM] = result1;
\end{lstlisting}}
\newpage

\paragraph{Encoding} Format \textsf{ALU\_WQAVG.M} (Section~\ref{sec:format-ALU-WQAVG.M}), \verb|exucode2 = 1100|\\

\bitfields{ 1/P, 2/00, 2/-- --, 27/upper27, 1/1, 2/11, 1/1, 4/1100, 5/registerM, 1/--, 2/10, 3/001, 1/0, 1/splat32, 5/lower5, 5/registerP, 1/1 }


\paragraph{Syntax} \texttt{avgwq} \emph{registerM} = \emph{registerP}, \emph{upper27\_\discretionary{}{}{}lower5}\emph{splat32}

\paragraph{Description} The \emph{upper27\_\discretionary{}{}{}lower5} and the \emph{registerP} are interpreted as four words packed into 128 bits. The \emph{upper27\_\discretionary{}{}{}lower5} words and the \emph{registerP} words are added. Each sum is shifted right arithmetic one bit. The four resulting words are packed and stored into the \emph{registerM}.


\paragraph{Modifier(s)}
\noindent\begin{itemize}
\item splat32 (cf. splat32 in section \ref{par:modifier-splat32} on page \pageref{par:modifier-splat32})
\end{itemize}

\paragraph{Execution} {Reservation table \textsf{ALU\_LITE.X} (Section~\ref{sec:reservation-ALU-LITE.X})}
~
{\begin{lstlisting}
stage ID:
  new splat32 = _ZX_1(splat32);
stage RR:
  new argument3 = splat32 ? (_WX_32(upper27_lower5) << 32) | _ZX_32(_WX_32(upper27_lower5)) : _SX_32(_WX_32(upper27_lower5));
  new argument2 = PGR[registerP];
  for (I = 0; I < 4; I += 1) {
    result1.32[I] = (_SX_32(argument3.32[I]) + _SX_32(argument2.32[I])) >> 1
  };
stage E1:
  PGR[registerM] = result1;
\end{lstlisting}}
\newpage
\subsubsection[{\small CATDQ: Catenate Double to Quadruple Word}]{CATDQ\hfill Catenate Double to Quadruple Word}
\label{instr:CATDQ}\index{catdq}
 
\paragraph{Encoding} Format \textsf{ALU\_CATDQ} (Section~\ref{sec:format-ALU-CATDQ}), \verb|exucode2 = 1111|\\

\bitfields{ 1/P, 2/11, 1/0, 4/1111, 5/registerM, 1/--, 2/10, 3/010, 1/1, 6/registerY, 6/registerZ }


\paragraph{Syntax} \texttt{catdq} \emph{registerM} = \emph{registerZ}, \emph{registerY}

\paragraph{Description} The \emph{registerM} operand receives the concatenation of the \emph{registerZ} and \emph{registerY} values.


\paragraph{Execution} {Reservation table \textsf{ALU\_LITE} (Section~\ref{sec:reservation-ALU-LITE})}
~
{\begin{lstlisting}
stage RR:
  new argument3 = GPR[registerY];
  new argument2 = GPR[registerZ];
  new result1 = _ZX_64(argument2) | (argument3 << 64);
stage E1:
  PGR[registerM] = result1;
\end{lstlisting}}
\newpage
\subsubsection[{\small CBSD: Count Bit Set Double Word}]{CBSD\hfill Count Bit Set Double Word}
\label{instr:CBSD}\index{cbsd}
 
\paragraph{Encoding} Format \textsf{ALU\_DBWR} (Section~\ref{sec:format-ALU-DBWR}), \verb|alucode8 = 0110|\\

\bitfields{ 1/P, 2/11, 1/0, 4/1111, 6/registerW, 2/10, 3/101, 1/0, 4/0110, 1/0, 1/0, 6/registerZ }


\paragraph{Syntax} \texttt{cbsd} \emph{registerW} = \emph{registerZ}

\paragraph{Description} Bit set population count value of the \emph{registerZ} is stored into the \emph{registerW}.


\paragraph{Execution} {Reservation table \textsf{ALU\_TINY} (Section~\ref{sec:reservation-ALU-TINY})}
~
{\begin{lstlisting}
stage RR:
  new argument2 = GPR[registerZ];
  new result1 = _ZX_64(_CBS_64(argument2));
stage E1:
  GPR[registerW] = result1;
\end{lstlisting}}
\newpage
\subsubsection[{\small CBSDP: Count Bit Sets Double Word Pair}]{CBSDP\hfill Count Bit Sets Double Word Pair}
\label{instr:CBSDP}\index{cbsdp}
 
\paragraph{Encoding} Format \textsf{ALU\_DPBWR} (Section~\ref{sec:format-ALU-DPBWR}), \verb|alucode8 = 0110|\\

\bitfields{ 1/P, 2/11, 1/1, 4/1111, 5/registerM, 1/--, 2/10, 3/101, 1/0, 4/0110, 1/--, 1/--, 5/registerP, 1/0 }


\paragraph{Syntax} \texttt{cbsdp} \emph{registerM} = \emph{registerP}

\paragraph{Description} Bit set population count values of the \emph{registerP} interpreted as double word pair are stored into the \emph{registerM}.


\paragraph{Execution} {Reservation table \textsf{ALU\_LITE} (Section~\ref{sec:reservation-ALU-LITE})}
~
{\begin{lstlisting}
stage RR:
  new argument2 = PGR[registerP];
  for (I = 0; I < 2; I += 1) {
    result1.64[I] = _ZX_64(_CBS_64(argument2.64[I]))
  };
stage E1:
  PGR[registerM] = result1;
\end{lstlisting}}
\newpage
\subsubsection[{\small CBSHO: Count Bit Sets Half Word Octuple}]{CBSHO\hfill Count Bit Sets Half Word Octuple}
\label{instr:CBSHO}\index{cbsho}
 
\paragraph{Encoding} Format \textsf{ALU\_HOBWR} (Section~\ref{sec:format-ALU-HOBWR}), \verb|alucode8 = 0110|\\

\bitfields{ 1/P, 2/11, 1/1, 4/1111, 5/registerM, 1/--, 2/10, 3/101, 1/1, 4/0110, 1/--, 1/--, 5/registerP, 1/0 }


\paragraph{Syntax} \texttt{cbsho} \emph{registerM} = \emph{registerP}

\paragraph{Description} Bit set population count values of the \emph{registerP} interpreted as half word octuple are stored into the \emph{registerM}.


\paragraph{Execution} {Reservation table \textsf{ALU\_LITE} (Section~\ref{sec:reservation-ALU-LITE})}
~
{\begin{lstlisting}
stage RR:
  new argument2 = PGR[registerP];
  for (I = 0; I < 8; I += 1) {
    result1.16[I] = _ZX_16(_CBS_16(argument2.16[I]))
  };
stage E1:
  PGR[registerM] = result1;
\end{lstlisting}}
\newpage
\subsubsection[{\small CBSW: Count Bit Set Word}]{CBSW\hfill Count Bit Set Word}
\label{instr:CBSW}\index{cbsw}
 
\paragraph{Encoding} Format \textsf{ALU\_BWRW} (Section~\ref{sec:format-ALU-BWRW}), \verb|alucode8 = 0110|\\

\bitfields{ 1/P, 2/11, 1/0, 4/1111, 6/registerW, 2/11, 3/101, 1/signextw, 4/0110, 1/0, 1/0, 6/registerZ }


\paragraph{Syntax} \texttt{cbsw} \emph{registerW} = \emph{registerZ}

\paragraph{Description} Bit set population count value of the \emph{registerZ} is stored into the \emph{registerW}.


\paragraph{Execution} {Reservation table \textsf{ALU\_TINY} (Section~\ref{sec:reservation-ALU-TINY})}
~
{\begin{lstlisting}
stage RR:
  new argument2 = GPR[registerZ];
  new result1 = _ZX_32(_CBS_32(argument2));
stage E1:
  GPR[registerW] = _ZX_32(result1);
\end{lstlisting}}
\newpage
\subsubsection[{\small CBSWQ: Count Bit Sets Word Quadruple}]{CBSWQ\hfill Count Bit Sets Word Quadruple}
\label{instr:CBSWQ}\index{cbswq}
 
\paragraph{Encoding} Format \textsf{ALU\_WQBWR} (Section~\ref{sec:format-ALU-WQBWR}), \verb|alucode8 = 0110|\\

\bitfields{ 1/P, 2/11, 1/1, 4/1111, 5/registerM, 1/--, 2/10, 3/101, 1/0, 4/0110, 1/--, 1/--, 5/registerP, 1/1 }


\paragraph{Syntax} \texttt{cbswq} \emph{registerM} = \emph{registerP}

\paragraph{Description} Bit set population count values of the \emph{registerP} interpreted as word quadruple are stored into the \emph{registerM}.


\paragraph{Execution} {Reservation table \textsf{ALU\_LITE} (Section~\ref{sec:reservation-ALU-LITE})}
~
{\begin{lstlisting}
stage RR:
  new argument2 = PGR[registerP];
  for (I = 0; I < 4; I += 1) {
    result1.32[I] = _ZX_32(_CBS_32(argument2.32[I]))
  };
stage E1:
  PGR[registerM] = result1;
\end{lstlisting}}
\newpage
\subsubsection[{\small CLSD: Count Leading Sign Double Word}]{CLSD\hfill Count Leading Sign Double Word}
\label{instr:CLSD}\index{clsd}
 
\paragraph{Encoding} Format \textsf{ALU\_DBWR} (Section~\ref{sec:format-ALU-DBWR}), \verb|alucode8 = 0101|\\

\bitfields{ 1/P, 2/11, 1/0, 4/1111, 6/registerW, 2/10, 3/101, 1/0, 4/0101, 1/0, 1/0, 6/registerZ }


\paragraph{Syntax} \texttt{clsd} \emph{registerW} = \emph{registerZ}

\paragraph{Description} The leading sign count value of the \emph{registerZ} is stored into the \emph{registerW}. When applied to zero or -1, the result is 63.


\paragraph{Execution} {Reservation table \textsf{ALU\_TINY} (Section~\ref{sec:reservation-ALU-TINY})}
~
{\begin{lstlisting}
stage RR:
  new argument2 = GPR[registerZ];
  new result1 = _ZX_64(_CLS_64(argument2));
stage E1:
  GPR[registerW] = result1;
\end{lstlisting}}
\newpage
\subsubsection[{\small CLSDP: Count Leading Signs Double Word Pair}]{CLSDP\hfill Count Leading Signs Double Word Pair}
\label{instr:CLSDP}\index{clsdp}
 
\paragraph{Encoding} Format \textsf{ALU\_DPBWR} (Section~\ref{sec:format-ALU-DPBWR}), \verb|alucode8 = 0101|\\

\bitfields{ 1/P, 2/11, 1/1, 4/1111, 5/registerM, 1/--, 2/10, 3/101, 1/0, 4/0101, 1/--, 1/--, 5/registerP, 1/0 }


\paragraph{Syntax} \texttt{clsdp} \emph{registerM} = \emph{registerP}

\paragraph{Description} The leading sign count values of the \emph{registerP} interpreted as double word pair are stored into the \emph{registerM}. When applied to zero or -1 at the word level, the result is 63.


\paragraph{Execution} {Reservation table \textsf{ALU\_LITE} (Section~\ref{sec:reservation-ALU-LITE})}
~
{\begin{lstlisting}
stage RR:
  new argument2 = PGR[registerP];
  for (I = 0; I < 2; I += 1) {
    result1.64[I] = _ZX_64(_CLS_64(argument2.64[I]))
  };
stage E1:
  PGR[registerM] = result1;
\end{lstlisting}}
\newpage
\subsubsection[{\small CLSHO: Count Leading Signs Half Word Octuple}]{CLSHO\hfill Count Leading Signs Half Word Octuple}
\label{instr:CLSHO}\index{clsho}
 
\paragraph{Encoding} Format \textsf{ALU\_HOBWR} (Section~\ref{sec:format-ALU-HOBWR}), \verb|alucode8 = 0101|\\

\bitfields{ 1/P, 2/11, 1/1, 4/1111, 5/registerM, 1/--, 2/10, 3/101, 1/1, 4/0101, 1/--, 1/--, 5/registerP, 1/0 }


\paragraph{Syntax} \texttt{clsho} \emph{registerM} = \emph{registerP}

\paragraph{Description} The leading sign count values of the \emph{registerP} interpreted as half word octuple are stored into the \emph{registerM}. When applied to zero or -1 at the half word level, the result is 31.


\paragraph{Execution} {Reservation table \textsf{ALU\_LITE} (Section~\ref{sec:reservation-ALU-LITE})}
~
{\begin{lstlisting}
stage RR:
  new argument2 = PGR[registerP];
  for (I = 0; I < 8; I += 1) {
    result1.16[I] = _ZX_16(_CLS_16(argument2.16[I]))
  };
stage E1:
  PGR[registerM] = result1;
\end{lstlisting}}
\newpage
\subsubsection[{\small CLSW: Count Leading Sign Word}]{CLSW\hfill Count Leading Sign Word}
\label{instr:CLSW}\index{clsw}
 
\paragraph{Encoding} Format \textsf{ALU\_BWRW} (Section~\ref{sec:format-ALU-BWRW}), \verb|alucode8 = 0101|\\

\bitfields{ 1/P, 2/11, 1/0, 4/1111, 6/registerW, 2/11, 3/101, 1/signextw, 4/0101, 1/0, 1/0, 6/registerZ }


\paragraph{Syntax} \texttt{clsw} \emph{registerW} = \emph{registerZ}

\paragraph{Description} The leading sign count value of the \emph{registerZ} is stored into the \emph{registerW}. When applied to zero or -1, the result is 31.


\paragraph{Execution} {Reservation table \textsf{ALU\_TINY} (Section~\ref{sec:reservation-ALU-TINY})}
~
{\begin{lstlisting}
stage RR:
  new argument2 = GPR[registerZ];
  new result1 = _ZX_32(_CLS_32(argument2));
stage E1:
  GPR[registerW] = _ZX_32(result1);
\end{lstlisting}}
\newpage
\subsubsection[{\small CLSWQ: Count Leading Signs Word Quadruple}]{CLSWQ\hfill Count Leading Signs Word Quadruple}
\label{instr:CLSWQ}\index{clswq}
 
\paragraph{Encoding} Format \textsf{ALU\_WQBWR} (Section~\ref{sec:format-ALU-WQBWR}), \verb|alucode8 = 0101|\\

\bitfields{ 1/P, 2/11, 1/1, 4/1111, 5/registerM, 1/--, 2/10, 3/101, 1/0, 4/0101, 1/--, 1/--, 5/registerP, 1/1 }


\paragraph{Syntax} \texttt{clswq} \emph{registerM} = \emph{registerP}

\paragraph{Description} The leading sign count values of the \emph{registerP} interpreted as word quadruple are stored into the \emph{registerM}. When applied to zero or -1 at the word level, the result is 31.


\paragraph{Execution} {Reservation table \textsf{ALU\_LITE} (Section~\ref{sec:reservation-ALU-LITE})}
~
{\begin{lstlisting}
stage RR:
  new argument2 = PGR[registerP];
  for (I = 0; I < 4; I += 1) {
    result1.32[I] = _ZX_32(_CLS_32(argument2.32[I]))
  };
stage E1:
  PGR[registerM] = result1;
\end{lstlisting}}
\newpage
\subsubsection[{\small CLZD: Count Leading Zero Double Word}]{CLZD\hfill Count Leading Zero Double Word}
\label{instr:CLZD}\index{clzd}
 
\paragraph{Encoding} Format \textsf{ALU\_DBWR} (Section~\ref{sec:format-ALU-DBWR}), \verb|alucode8 = 0100|\\

\bitfields{ 1/P, 2/11, 1/0, 4/1111, 6/registerW, 2/10, 3/101, 1/0, 4/0100, 1/0, 1/0, 6/registerZ }


\paragraph{Syntax} \texttt{clzd} \emph{registerW} = \emph{registerZ}

\paragraph{Description} The leading zero count value of the \emph{registerZ} is stored into the \emph{registerW}. When applied to zero, the result is 64.


\paragraph{Execution} {Reservation table \textsf{ALU\_TINY} (Section~\ref{sec:reservation-ALU-TINY})}
~
{\begin{lstlisting}
stage RR:
  new argument2 = GPR[registerZ];
  new result1 = _ZX_64(_CLZ_64(argument2));
stage E1:
  GPR[registerW] = result1;
\end{lstlisting}}
\newpage
\subsubsection[{\small CLZDP: Count Leading Zeroes Double Word Pair}]{CLZDP\hfill Count Leading Zeroes Double Word Pair}
\label{instr:CLZDP}\index{clzdp}
 
\paragraph{Encoding} Format \textsf{ALU\_DPBWR} (Section~\ref{sec:format-ALU-DPBWR}), \verb|alucode8 = 0100|\\

\bitfields{ 1/P, 2/11, 1/1, 4/1111, 5/registerM, 1/--, 2/10, 3/101, 1/0, 4/0100, 1/--, 1/--, 5/registerP, 1/0 }


\paragraph{Syntax} \texttt{clzdp} \emph{registerM} = \emph{registerP}

\paragraph{Description} The leading zero count values of the \emph{registerP} interpreted as a double word pair are stored into the \emph{registerM}. When applied to zero at the word level, the result is 64.


\paragraph{Execution} {Reservation table \textsf{ALU\_LITE} (Section~\ref{sec:reservation-ALU-LITE})}
~
{\begin{lstlisting}
stage RR:
  new argument2 = PGR[registerP];
  for (I = 0; I < 2; I += 1) {
    result1.64[I] = _ZX_64(_CLZ_64(argument2.64[I]))
  };
stage E1:
  PGR[registerM] = result1;
\end{lstlisting}}
\newpage
\subsubsection[{\small CLZHO: Count Leading Zeroes Half Word Octuple}]{CLZHO\hfill Count Leading Zeroes Half Word Octuple}
\label{instr:CLZHO}\index{clzho}
 
\paragraph{Encoding} Format \textsf{ALU\_HOBWR} (Section~\ref{sec:format-ALU-HOBWR}), \verb|alucode8 = 0100|\\

\bitfields{ 1/P, 2/11, 1/1, 4/1111, 5/registerM, 1/--, 2/10, 3/101, 1/1, 4/0100, 1/--, 1/--, 5/registerP, 1/0 }


\paragraph{Syntax} \texttt{clzho} \emph{registerM} = \emph{registerP}

\paragraph{Description} The leading zero count values of the \emph{registerP} interpreted as a half word octuple are stored into the \emph{registerM}. When applied to zero at the half word level, the result is 16.


\paragraph{Execution} {Reservation table \textsf{ALU\_LITE} (Section~\ref{sec:reservation-ALU-LITE})}
~
{\begin{lstlisting}
stage RR:
  new argument2 = PGR[registerP];
  for (I = 0; I < 8; I += 1) {
    result1.16[I] = _ZX_16(_CLZ_16(argument2.16[I]))
  };
stage E1:
  PGR[registerM] = result1;
\end{lstlisting}}
\newpage
\subsubsection[{\small CLZW: Count Leading Zero Word}]{CLZW\hfill Count Leading Zero Word}
\label{instr:CLZW}\index{clzw}
 
\paragraph{Encoding} Format \textsf{ALU\_BWRW} (Section~\ref{sec:format-ALU-BWRW}), \verb|alucode8 = 0100|\\

\bitfields{ 1/P, 2/11, 1/0, 4/1111, 6/registerW, 2/11, 3/101, 1/signextw, 4/0100, 1/0, 1/0, 6/registerZ }


\paragraph{Syntax} \texttt{clzw} \emph{registerW} = \emph{registerZ}

\paragraph{Description} The leading zero count value of the \emph{registerZ} is stored into the \emph{registerW}. When applied to zero, the result is 32.


\paragraph{Execution} {Reservation table \textsf{ALU\_TINY} (Section~\ref{sec:reservation-ALU-TINY})}
~
{\begin{lstlisting}
stage RR:
  new argument2 = GPR[registerZ];
  new result1 = _ZX_32(_CLZ_32(argument2));
stage E1:
  GPR[registerW] = _ZX_32(result1);
\end{lstlisting}}
\newpage
\subsubsection[{\small CLZWQ: Count Leading Zeroes Word Quadruple}]{CLZWQ\hfill Count Leading Zeroes Word Quadruple}
\label{instr:CLZWQ}\index{clzwq}
 
\paragraph{Encoding} Format \textsf{ALU\_WQBWR} (Section~\ref{sec:format-ALU-WQBWR}), \verb|alucode8 = 0100|\\

\bitfields{ 1/P, 2/11, 1/1, 4/1111, 5/registerM, 1/--, 2/10, 3/101, 1/0, 4/0100, 1/--, 1/--, 5/registerP, 1/1 }


\paragraph{Syntax} \texttt{clzwq} \emph{registerM} = \emph{registerP}

\paragraph{Description} The leading zero count values of the \emph{registerP} interpreted as a word quadruple are stored into the \emph{registerM}. When applied to zero at the word level, the result is 32.


\paragraph{Execution} {Reservation table \textsf{ALU\_LITE} (Section~\ref{sec:reservation-ALU-LITE})}
~
{\begin{lstlisting}
stage RR:
  new argument2 = PGR[registerP];
  for (I = 0; I < 4; I += 1) {
    result1.32[I] = _ZX_32(_CLZ_32(argument2.32[I]))
  };
stage E1:
  PGR[registerM] = result1;
\end{lstlisting}}
\newpage
\subsubsection[{\small CMOVEBX: Conditional Move Byte Hexadecuple}]{CMOVEBX\hfill Conditional Move Byte Hexadecuple}
\label{instr:CMOVEBX}\index{cmovebx}
 
\paragraph{Encoding} Format \textsf{ALU\_BXCMWRR} (Section~\ref{sec:format-ALU-BXCMWRR}), \verb|exubit24 = 0|\\

\bitfields{ 1/P, 2/11, 1/1, 3/lanecond, 1/0, 5/registerM, 1/--, 2/10, 3/011, 1/1, 5/registerO, 1/--, 5/registerP, 1/1 }


\paragraph{Syntax} \texttt{cmovebx}\emph{lanecond} \emph{registerP}? \emph{registerM} = \emph{registerO}

\paragraph{Description} The \emph{registerM}, \emph{registerP} and the \emph{registerO} are interpreted as sixteen bytes packed into 128 bits. Each \emph{registerP} word is tested using the condition \emph{lanecond}. If true, the corresponding byte of the \emph{registerO} is stored into the corresponding word of the \emph{registerM}.


\paragraph{Modifier(s)}
\noindent\begin{itemize}
\item lanecond (cf. lanecond in section \ref{par:modifier-lanecond} on page \pageref{par:modifier-lanecond})
\end{itemize}

\paragraph{Execution} {Reservation table \textsf{ALU\_LITE} (Section~\ref{sec:reservation-ALU-LITE})}
~
{\begin{lstlisting}
stage ID:
  new argument4 = _ZX_3(lanecond);
stage RR:
  new argument3 = PGR[registerO];
  new argument2 = PGR[registerP];
stage E1:
  if (lanecond_8(argument4, argument2.8[0])) {
    PGR[registerM] = argument3;
  }
  if (lanecond_8(argument4, argument2.8[1])) {
    PGR[registerM] = argument3;
  }
  if (lanecond_8(argument4, argument2.8[2])) {
    PGR[registerM] = argument3;
  }
  if (lanecond_8(argument4, argument2.8[3])) {
    PGR[registerM] = argument3;
  }
  if (lanecond_8(argument4, argument2.8[4])) {
    PGR[registerM] = argument3;
  }
  if (lanecond_8(argument4, argument2.8[5])) {
    PGR[registerM] = argument3;
  }
  if (lanecond_8(argument4, argument2.8[6])) {
    PGR[registerM] = argument3;
  }
  if (lanecond_8(argument4, argument2.8[7])) {
    PGR[registerM] = argument3;
  }
  if (lanecond_8(argument4, argument2.8[8])) {
    PGR[registerM] = argument3;
  }
  if (lanecond_8(argument4, argument2.8[9])) {
    PGR[registerM] = argument3;
  }
  if (lanecond_8(argument4, argument2.8[10])) {
    PGR[registerM] = argument3;
  }
  if (lanecond_8(argument4, argument2.8[11])) {
    PGR[registerM] = argument3;
  }
  if (lanecond_8(argument4, argument2.8[12])) {
    PGR[registerM] = argument3;
  }
  if (lanecond_8(argument4, argument2.8[13])) {
    PGR[registerM] = argument3;
  }
  if (lanecond_8(argument4, argument2.8[14])) {
    PGR[registerM] = argument3;
  }
  if (lanecond_8(argument4, argument2.8[15])) {
    PGR[registerM] = argument3;
  }
\end{lstlisting}}
\newpage

\paragraph{Encoding} Format \textsf{ALU\_BXCMWRR.M} (Section~\ref{sec:format-ALU-BXCMWRR.M}), \verb|exubit24 = 0|\\

\bitfields{ 1/P, 2/00, 2/-- --, 27/upper27, 1/1, 2/11, 1/1, 3/lanecond, 1/0, 5/registerM, 1/--, 2/10, 3/011, 1/1, 1/splat32, 5/lower5, 5/registerP, 1/1 }


\paragraph{Syntax} \texttt{cmovebx}\emph{lanecond} \emph{registerP}? \emph{registerM} = \emph{upper27\_\discretionary{}{}{}lower5}\emph{splat32}

\paragraph{Description} The \emph{registerM}, \emph{registerP} and the \emph{upper27\_\discretionary{}{}{}lower5} are interpreted as sixteen bytes packed into 128 bits. Each \emph{registerP} word is tested using the condition \emph{lanecond}. If true, the corresponding byte of the \emph{upper27\_\discretionary{}{}{}lower5} is stored into the corresponding word of the \emph{registerM}.


\paragraph{Modifier(s)}
\noindent\begin{itemize}
\item lanecond (cf. lanecond in section \ref{par:modifier-lanecond} on page \pageref{par:modifier-lanecond})
\item splat32 (cf. splat32 in section \ref{par:modifier-splat32} on page \pageref{par:modifier-splat32})
\end{itemize}

\paragraph{Execution} {Reservation table \textsf{ALU\_LITE.X} (Section~\ref{sec:reservation-ALU-LITE.X})}
~
{\begin{lstlisting}
stage ID:
  new splat32 = _ZX_1(splat32);
  new argument4 = _ZX_3(lanecond);
stage RR:
  new argument3 = splat32 ? (_WX_32(upper27_lower5) << 32) | _ZX_32(_WX_32(upper27_lower5)) : _SX_32(_WX_32(upper27_lower5));
  new argument2 = PGR[registerP];
stage E1:
  if (lanecond_8(argument4, argument2.8[0])) {
    PGR[registerM] = argument3;
  }
  if (lanecond_8(argument4, argument2.8[1])) {
    PGR[registerM] = argument3;
  }
  if (lanecond_8(argument4, argument2.8[2])) {
    PGR[registerM] = argument3;
  }
  if (lanecond_8(argument4, argument2.8[3])) {
    PGR[registerM] = argument3;
  }
  if (lanecond_8(argument4, argument2.8[4])) {
    PGR[registerM] = argument3;
  }
  if (lanecond_8(argument4, argument2.8[5])) {
    PGR[registerM] = argument3;
  }
  if (lanecond_8(argument4, argument2.8[6])) {
    PGR[registerM] = argument3;
  }
  if (lanecond_8(argument4, argument2.8[7])) {
    PGR[registerM] = argument3;
  }
  if (lanecond_8(argument4, argument2.8[8])) {
    PGR[registerM] = argument3;
  }
  if (lanecond_8(argument4, argument2.8[9])) {
    PGR[registerM] = argument3;
  }
  if (lanecond_8(argument4, argument2.8[10])) {
    PGR[registerM] = argument3;
  }
  if (lanecond_8(argument4, argument2.8[11])) {
    PGR[registerM] = argument3;
  }
  if (lanecond_8(argument4, argument2.8[12])) {
    PGR[registerM] = argument3;
  }
  if (lanecond_8(argument4, argument2.8[13])) {
    PGR[registerM] = argument3;
  }
  if (lanecond_8(argument4, argument2.8[14])) {
    PGR[registerM] = argument3;
  }
  if (lanecond_8(argument4, argument2.8[15])) {
    PGR[registerM] = argument3;
  }
\end{lstlisting}}
\newpage
\subsubsection[{\small CMOVED: Conditional Move Double Word}]{CMOVED\hfill Conditional Move Double Word}
\label{instr:CMOVED}\index{cmoved}
 
\paragraph{Encoding} Format \textsf{ALU\_DCMWRR} (Section~\ref{sec:format-ALU-DCMWRR}), \verb|steering = 11|\\

\bitfields{ 1/P, 2/11, 1/0, 4/cmovecond, 6/registerW, 2/10, 3/011, 1/0, 6/registerY, 6/registerZ }


\paragraph{Syntax} \texttt{cmoved}\emph{cmovecond} \emph{registerZ}? \emph{registerW} = \emph{registerY}

\paragraph{Description} The \emph{registerZ} double word is tested using the condition \emph{cmovecond}. If true, the \emph{registerY} double word is stored into the \emph{registerW} double word. Else the \emph{registerW} double word is unchanged.


\paragraph{Modifier(s)}
\noindent\begin{itemize}
\item cmovecond (cf. bcucond in section \ref{par:modifier-bcucond} on page \pageref{par:modifier-bcucond})
\end{itemize}

\paragraph{Execution} {Reservation table \textsf{ALU\_LITE} (Section~\ref{sec:reservation-ALU-LITE})}
~
{\begin{lstlisting}
stage ID:
  new argument4 = _ZX_4(cmovecond);
stage RR:
  new argument3 = GPR[registerY];
  new argument2 = GPR[registerZ];
stage E1:
  if (bcucond(argument4, argument2)) {
    GPR[registerW] = argument3;
  }
\end{lstlisting}}
\newpage

\paragraph{Encoding} Format \textsf{ALU\_DCMWRR.M} (Section~\ref{sec:format-ALU-DCMWRR.M}), \verb|steering = 11|\\

\bitfields{ 1/P, 2/00, 2/-- --, 27/upper27, 1/1, 2/11, 1/0, 4/cmovecond, 6/registerW, 2/10, 3/011, 1/0, 1/splat32, 5/lower5, 6/registerZ }


\paragraph{Syntax} \texttt{cmoved}\emph{cmovecond} \emph{registerZ}? \emph{registerW} = \emph{upper27\_\discretionary{}{}{}lower5}\emph{splat32}

\paragraph{Description} The \emph{registerZ} double word is tested using the condition \emph{cmovecond}. If true, the \emph{upper27\_\discretionary{}{}{}lower5} double word is stored into the \emph{registerW} double word. Else the \emph{registerW} double word is unchanged.


\paragraph{Modifier(s)}
\noindent\begin{itemize}
\item cmovecond (cf. bcucond in section \ref{par:modifier-bcucond} on page \pageref{par:modifier-bcucond})
\item splat32 (cf. splat32 in section \ref{par:modifier-splat32} on page \pageref{par:modifier-splat32})
\end{itemize}

\paragraph{Execution} {Reservation table \textsf{ALU\_LITE.X} (Section~\ref{sec:reservation-ALU-LITE.X})}
~
{\begin{lstlisting}
stage ID:
  new splat32 = _ZX_1(splat32);
  new argument4 = _ZX_4(cmovecond);
stage RR:
  new argument3 = splat32 ? (_WX_32(upper27_lower5) << 32) | _ZX_32(_WX_32(upper27_lower5)) : _SX_32(_WX_32(upper27_lower5));
  new argument2 = GPR[registerZ];
stage E1:
  if (bcucond(argument4, argument2)) {
    GPR[registerW] = argument3;
  }
\end{lstlisting}}
\newpage
\subsubsection[{\small CMOVEDP: Conditional Move Double Word Pair}]{CMOVEDP\hfill Conditional Move Double Word Pair}
\label{instr:CMOVEDP}\index{cmovedp}
 
\paragraph{Encoding} Format \textsf{ALU\_DPCMWRR} (Section~\ref{sec:format-ALU-DPCMWRR}), \verb|exubit24 = 0|\\

\bitfields{ 1/P, 2/11, 1/1, 3/lanecond, 1/0, 5/registerM, 1/--, 2/10, 3/011, 1/0, 5/registerO, 1/--, 5/registerP, 1/0 }


\paragraph{Syntax} \texttt{cmovedp}\emph{lanecond} \emph{registerP}? \emph{registerM} = \emph{registerO}

\paragraph{Description} The \emph{registerM}, \emph{registerP} and the \emph{registerO} are interpreted as two double words packed into 128 bits. Each \emph{registerP} double word is tested using the condition \emph{lanecond}. If true, the corresponding \emph{registerO} double word is stored into the corresponding \emph{registerM} double word. Else, the corresponding \emph{registerM} double word is unchanged.


\paragraph{Modifier(s)}
\noindent\begin{itemize}
\item lanecond (cf. lanecond in section \ref{par:modifier-lanecond} on page \pageref{par:modifier-lanecond})
\end{itemize}

\paragraph{Execution} {Reservation table \textsf{ALU\_LITE} (Section~\ref{sec:reservation-ALU-LITE})}
~
{\begin{lstlisting}
stage ID:
  new argument4 = _ZX_3(lanecond);
stage RR:
  new argument3 = PGR[registerO];
  new argument2 = PGR[registerP];
stage E1:
  if (bcucond(argument4, argument2.64[0])) {
    PGR[registerM] = argument3;
  }
  if (bcucond(argument4, argument2.64[1])) {
    PGR[registerM] = argument3;
  }
\end{lstlisting}}
\newpage

\paragraph{Encoding} Format \textsf{ALU\_DPCMWRR.M} (Section~\ref{sec:format-ALU-DPCMWRR.M}), \verb|exubit24 = 0|\\

\bitfields{ 1/P, 2/00, 2/-- --, 27/upper27, 1/1, 2/11, 1/1, 3/lanecond, 1/0, 5/registerM, 1/--, 2/10, 3/011, 1/0, 1/splat32, 5/lower5, 5/registerP, 1/0 }


\paragraph{Syntax} \texttt{cmovedp}\emph{lanecond} \emph{registerP}? \emph{registerM} = \emph{upper27\_\discretionary{}{}{}lower5}\emph{splat32}

\paragraph{Description} The \emph{registerM}, \emph{registerP} and the \emph{upper27\_\discretionary{}{}{}lower5} are interpreted as two double words packed into 128 bits. Each \emph{registerP} double word is tested using the condition \emph{lanecond}. If true, the corresponding \emph{upper27\_\discretionary{}{}{}lower5} double word is stored into the corresponding \emph{registerM} double word. Else, the corresponding \emph{registerM} double word is unchanged.


\paragraph{Modifier(s)}
\noindent\begin{itemize}
\item lanecond (cf. lanecond in section \ref{par:modifier-lanecond} on page \pageref{par:modifier-lanecond})
\item splat32 (cf. splat32 in section \ref{par:modifier-splat32} on page \pageref{par:modifier-splat32})
\end{itemize}

\paragraph{Execution} {Reservation table \textsf{ALU\_LITE.X} (Section~\ref{sec:reservation-ALU-LITE.X})}
~
{\begin{lstlisting}
stage ID:
  new splat32 = _ZX_1(splat32);
  new argument4 = _ZX_3(lanecond);
stage RR:
  new argument3 = splat32 ? (_WX_32(upper27_lower5) << 32) | _ZX_32(_WX_32(upper27_lower5)) : _SX_32(_WX_32(upper27_lower5));
  new argument2 = PGR[registerP];
stage E1:
  if (bcucond(argument4, argument2.64[0])) {
    PGR[registerM] = argument3;
  }
  if (bcucond(argument4, argument2.64[1])) {
    PGR[registerM] = argument3;
  }
\end{lstlisting}}
\newpage
\subsubsection[{\small CMOVEHO: Conditional Move Half Word Octuple}]{CMOVEHO\hfill Conditional Move Half Word Octuple}
\label{instr:CMOVEHO}\index{cmoveho}
 
\paragraph{Encoding} Format \textsf{ALU\_HOCMWRR} (Section~\ref{sec:format-ALU-HOCMWRR}), \verb|exubit24 = 0|\\

\bitfields{ 1/P, 2/11, 1/1, 3/lanecond, 1/0, 5/registerM, 1/--, 2/10, 3/011, 1/1, 5/registerO, 1/--, 5/registerP, 1/0 }


\paragraph{Syntax} \texttt{cmoveho}\emph{lanecond} \emph{registerP}? \emph{registerM} = \emph{registerO}

\paragraph{Description} The \emph{registerM}, \emph{registerP} and the \emph{registerO} are interpreted as eight half words packed into 128 bits. Each \emph{registerP} half word is tested using the condition \emph{lanecond}. If true, the corresponding half word of the \emph{registerO} is stored into the corresponding half word of the \emph{registerM}.


\paragraph{Modifier(s)}
\noindent\begin{itemize}
\item lanecond (cf. lanecond in section \ref{par:modifier-lanecond} on page \pageref{par:modifier-lanecond})
\end{itemize}

\paragraph{Execution} {Reservation table \textsf{ALU\_LITE} (Section~\ref{sec:reservation-ALU-LITE})}
~
{\begin{lstlisting}
stage ID:
  new argument4 = _ZX_3(lanecond);
stage RR:
  new argument3 = PGR[registerO];
  new argument2 = PGR[registerP];
stage E1:
  if (lanecond_16(argument4, argument2.16[0])) {
    PGR[registerM] = argument3;
  }
  if (lanecond_16(argument4, argument2.16[1])) {
    PGR[registerM] = argument3;
  }
  if (lanecond_16(argument4, argument2.16[2])) {
    PGR[registerM] = argument3;
  }
  if (lanecond_16(argument4, argument2.16[3])) {
    PGR[registerM] = argument3;
  }
  if (lanecond_16(argument4, argument2.16[4])) {
    PGR[registerM] = argument3;
  }
  if (lanecond_16(argument4, argument2.16[5])) {
    PGR[registerM] = argument3;
  }
  if (lanecond_16(argument4, argument2.16[6])) {
    PGR[registerM] = argument3;
  }
  if (lanecond_16(argument4, argument2.16[7])) {
    PGR[registerM] = argument3;
  }
\end{lstlisting}}
\newpage

\paragraph{Encoding} Format \textsf{ALU\_HOCMWRR.M} (Section~\ref{sec:format-ALU-HOCMWRR.M}), \verb|exubit24 = 0|\\

\bitfields{ 1/P, 2/00, 2/-- --, 27/upper27, 1/1, 2/11, 1/1, 3/lanecond, 1/0, 5/registerM, 1/--, 2/10, 3/011, 1/1, 1/splat32, 5/lower5, 5/registerP, 1/0 }


\paragraph{Syntax} \texttt{cmoveho}\emph{lanecond} \emph{registerP}? \emph{registerM} = \emph{upper27\_\discretionary{}{}{}lower5}\emph{splat32}

\paragraph{Description} The \emph{registerM}, \emph{registerP} and the \emph{upper27\_\discretionary{}{}{}lower5} are interpreted as eight half words packed into 128 bits. Each \emph{registerP} half word is tested using the condition \emph{lanecond}. If true, the corresponding half word of the \emph{upper27\_\discretionary{}{}{}lower5} is stored into the corresponding half word of the \emph{registerM}.


\paragraph{Modifier(s)}
\noindent\begin{itemize}
\item lanecond (cf. lanecond in section \ref{par:modifier-lanecond} on page \pageref{par:modifier-lanecond})
\item splat32 (cf. splat32 in section \ref{par:modifier-splat32} on page \pageref{par:modifier-splat32})
\end{itemize}

\paragraph{Execution} {Reservation table \textsf{ALU\_LITE.X} (Section~\ref{sec:reservation-ALU-LITE.X})}
~
{\begin{lstlisting}
stage ID:
  new splat32 = _ZX_1(splat32);
  new argument4 = _ZX_3(lanecond);
stage RR:
  new argument3 = splat32 ? (_WX_32(upper27_lower5) << 32) | _ZX_32(_WX_32(upper27_lower5)) : _SX_32(_WX_32(upper27_lower5));
  new argument2 = PGR[registerP];
stage E1:
  if (lanecond_16(argument4, argument2.16[0])) {
    PGR[registerM] = argument3;
  }
  if (lanecond_16(argument4, argument2.16[1])) {
    PGR[registerM] = argument3;
  }
  if (lanecond_16(argument4, argument2.16[2])) {
    PGR[registerM] = argument3;
  }
  if (lanecond_16(argument4, argument2.16[3])) {
    PGR[registerM] = argument3;
  }
  if (lanecond_16(argument4, argument2.16[4])) {
    PGR[registerM] = argument3;
  }
  if (lanecond_16(argument4, argument2.16[5])) {
    PGR[registerM] = argument3;
  }
  if (lanecond_16(argument4, argument2.16[6])) {
    PGR[registerM] = argument3;
  }
  if (lanecond_16(argument4, argument2.16[7])) {
    PGR[registerM] = argument3;
  }
\end{lstlisting}}
\newpage
\subsubsection[{\small CMOVEQ: Conditional Move Quadruple Word}]{CMOVEQ\hfill Conditional Move Quadruple Word}
\label{instr:CMOVEQ}\index{cmoveq}
 
\paragraph{Encoding} Format \textsf{ALU\_QCMWRR} (Section~\ref{sec:format-ALU-QCMWRR}), \verb|steering = 11|\\

\bitfields{ 1/P, 2/11, 1/0, 4/cmovecond, 5/registerM, 1/--, 2/10, 3/011, 1/1, 5/registerO, 1/--, 6/registerZ }


\paragraph{Syntax} \texttt{cmoveq}\emph{cmovecond} \emph{registerZ}? \emph{registerM} = \emph{registerO}

\paragraph{Description} The \emph{registerZ} double word is tested using the condition \emph{cmovecond}. If true, the \emph{registerO} quadruple word is stored into the \emph{registerM} quadruple word. Else the \emph{registerM} quadruple word is unchanged.


\paragraph{Modifier(s)}
\noindent\begin{itemize}
\item cmovecond (cf. bcucond in section \ref{par:modifier-bcucond} on page \pageref{par:modifier-bcucond})
\end{itemize}

\paragraph{Execution} {Reservation table \textsf{ALU\_LITE} (Section~\ref{sec:reservation-ALU-LITE})}
~
{\begin{lstlisting}
stage ID:
  new argument4 = _ZX_4(cmovecond);
stage RR:
  new argument3 = PGR[registerO];
  new argument2 = GPR[registerZ];
stage E1:
  if (bcucond(argument4, argument2)) {
    PGR[registerM] = argument3;
  }
\end{lstlisting}}
\newpage
\subsubsection[{\small CMOVEWQ: Conditional Move Word Quadruple}]{CMOVEWQ\hfill Conditional Move Word Quadruple}
\label{instr:CMOVEWQ}\index{cmovewq}
 
\paragraph{Encoding} Format \textsf{ALU\_WQCMWRR} (Section~\ref{sec:format-ALU-WQCMWRR}), \verb|exubit24 = 0|\\

\bitfields{ 1/P, 2/11, 1/1, 3/lanecond, 1/0, 5/registerM, 1/--, 2/10, 3/011, 1/0, 5/registerO, 1/--, 5/registerP, 1/1 }


\paragraph{Syntax} \texttt{cmovewq}\emph{lanecond} \emph{registerP}? \emph{registerM} = \emph{registerO}

\paragraph{Description} The \emph{registerM}, \emph{registerP} and the \emph{registerO} are interpreted as four words packed into 128 bits. Each \emph{registerP} word is tested using the condition \emph{lanecond}. If true, the corresponding word of the \emph{registerO} is stored into the corresponding word of the \emph{registerM}.


\paragraph{Modifier(s)}
\noindent\begin{itemize}
\item lanecond (cf. lanecond in section \ref{par:modifier-lanecond} on page \pageref{par:modifier-lanecond})
\end{itemize}

\paragraph{Execution} {Reservation table \textsf{ALU\_LITE} (Section~\ref{sec:reservation-ALU-LITE})}
~
{\begin{lstlisting}
stage ID:
  new argument4 = _ZX_3(lanecond);
stage RR:
  new argument3 = PGR[registerO];
  new argument2 = PGR[registerP];
stage E1:
  if (lanecond_32(argument4, argument2.32[0])) {
    PGR[registerM] = argument3;
  }
  if (lanecond_32(argument4, argument2.32[1])) {
    PGR[registerM] = argument3;
  }
  if (lanecond_32(argument4, argument2.32[2])) {
    PGR[registerM] = argument3;
  }
  if (lanecond_32(argument4, argument2.32[3])) {
    PGR[registerM] = argument3;
  }
\end{lstlisting}}
\newpage

\paragraph{Encoding} Format \textsf{ALU\_WQCMWRR.M} (Section~\ref{sec:format-ALU-WQCMWRR.M}), \verb|exubit24 = 0|\\

\bitfields{ 1/P, 2/00, 2/-- --, 27/upper27, 1/1, 2/11, 1/1, 3/lanecond, 1/0, 5/registerM, 1/--, 2/10, 3/011, 1/0, 1/splat32, 5/lower5, 5/registerP, 1/1 }


\paragraph{Syntax} \texttt{cmovewq}\emph{lanecond} \emph{registerP}? \emph{registerM} = \emph{upper27\_\discretionary{}{}{}lower5}\emph{splat32}

\paragraph{Description} The \emph{registerM}, \emph{registerP} and the \emph{upper27\_\discretionary{}{}{}lower5} are interpreted as four words packed into 128 bits. Each \emph{registerP} word is tested using the condition \emph{lanecond}. If true, the corresponding word of the \emph{upper27\_\discretionary{}{}{}lower5} is stored into the corresponding word of the \emph{registerM}.


\paragraph{Modifier(s)}
\noindent\begin{itemize}
\item lanecond (cf. lanecond in section \ref{par:modifier-lanecond} on page \pageref{par:modifier-lanecond})
\item splat32 (cf. splat32 in section \ref{par:modifier-splat32} on page \pageref{par:modifier-splat32})
\end{itemize}

\paragraph{Execution} {Reservation table \textsf{ALU\_LITE.X} (Section~\ref{sec:reservation-ALU-LITE.X})}
~
{\begin{lstlisting}
stage ID:
  new splat32 = _ZX_1(splat32);
  new argument4 = _ZX_3(lanecond);
stage RR:
  new argument3 = splat32 ? (_WX_32(upper27_lower5) << 32) | _ZX_32(_WX_32(upper27_lower5)) : _SX_32(_WX_32(upper27_lower5));
  new argument2 = PGR[registerP];
stage E1:
  if (lanecond_32(argument4, argument2.32[0])) {
    PGR[registerM] = argument3;
  }
  if (lanecond_32(argument4, argument2.32[1])) {
    PGR[registerM] = argument3;
  }
  if (lanecond_32(argument4, argument2.32[2])) {
    PGR[registerM] = argument3;
  }
  if (lanecond_32(argument4, argument2.32[3])) {
    PGR[registerM] = argument3;
  }
\end{lstlisting}}
\newpage
\subsubsection[{\small COMPBX: Compare Byte Hexadecuple}]{COMPBX\hfill Compare Byte Hexadecuple}
\label{instr:COMPBX}\index{compbx}
 
\paragraph{Encoding} Format \textsf{ALU\_BXCWRR} (Section~\ref{sec:format-ALU-BXCWRR}), \verb|alucode3 = 001|\\

\bitfields{ 1/P, 2/11, 1/1, 4/intcomp, 6/registerW, 2/10, 3/001, 1/1, 5/registerO, 1/--, 5/registerP, 1/1 }


\paragraph{Syntax} \texttt{compbx}\emph{intcomp} \emph{registerW} = \emph{registerP}, \emph{registerO}

\paragraph{Description} The byte hexadecuple \emph{registerP} is compared to the byte hexadecuple \emph{registerO} using condition \emph{intcomp}. The boolean result bits are packed and stored into the \emph{registerW}.


\paragraph{Modifier(s)}
\noindent\begin{itemize}
\item intcomp (cf. intcomp in section \ref{par:modifier-intcomp} on page \pageref{par:modifier-intcomp})
\end{itemize}

\paragraph{Execution} {Reservation table \textsf{ALU\_LITE} (Section~\ref{sec:reservation-ALU-LITE})}
~
{\begin{lstlisting}
stage ID:
  new argument4 = _ZX_4(intcomp);
stage RR:
  new argument3 = PGR[registerO];
  new argument2 = PGR[registerP];
  for (I = 0; I < 16; I += 1) {
    result1.1[I] = intcomp_8(argument4, argument2.8[I], argument3.8[I])
  };
stage E1:
  GPR[registerW] = result1;
\end{lstlisting}}
\newpage

\paragraph{Encoding} Format \textsf{ALU\_BXCWRR.M} (Section~\ref{sec:format-ALU-BXCWRR.M}), \verb|alucode3 = 001|\\

\bitfields{ 1/P, 2/00, 2/-- --, 27/upper27, 1/1, 2/11, 1/1, 4/intcomp, 6/registerW, 2/10, 3/001, 1/1, 1/splat32, 5/lower5, 5/registerP, 1/1 }


\paragraph{Syntax} \texttt{compbx}\emph{intcomp} \emph{registerW} = \emph{registerP}, \emph{upper27\_\discretionary{}{}{}lower5}\emph{splat32}

\paragraph{Description} The byte hexadecuple \emph{registerP} is compared to the byte hexadecuple \emph{upper27\_\discretionary{}{}{}lower5} using condition \emph{intcomp}. The boolean result bits are packed and stored into the \emph{registerW}.


\paragraph{Modifier(s)}
\noindent\begin{itemize}
\item intcomp (cf. intcomp in section \ref{par:modifier-intcomp} on page \pageref{par:modifier-intcomp})
\item splat32 (cf. splat32 in section \ref{par:modifier-splat32} on page \pageref{par:modifier-splat32})
\end{itemize}

\paragraph{Execution} {Reservation table \textsf{ALU\_LITE.X} (Section~\ref{sec:reservation-ALU-LITE.X})}
~
{\begin{lstlisting}
stage ID:
  new splat32 = _ZX_1(splat32);
  new argument4 = _ZX_4(intcomp);
stage RR:
  new argument3 = splat32 ? (_WX_32(upper27_lower5) << 32) | _ZX_32(_WX_32(upper27_lower5)) : _SX_32(_WX_32(upper27_lower5));
  new argument2 = PGR[registerP];
  for (I = 0; I < 16; I += 1) {
    result1.1[I] = intcomp_8(argument4, argument2.8[I], argument3.8[I])
  };
stage E1:
  GPR[registerW] = result1;
\end{lstlisting}}
\newpage
\subsubsection[{\small COMPD: Compare Double Word}]{COMPD\hfill Compare Double Word}
\label{instr:COMPD}\index{compd}
 
\paragraph{Encoding} Format \textsf{ALU\_DCWRR} (Section~\ref{sec:format-ALU-DCWRR}), \verb|steering = 11|\\

\bitfields{ 1/P, 2/11, 1/0, 4/intcomp, 6/registerW, 2/10, 3/001, 1/0, 6/registerY, 6/registerZ }


\paragraph{Syntax} \texttt{compd}\emph{intcomp} \emph{registerW} = \emph{registerZ}, \emph{registerY}

\paragraph{Description} The double word \emph{registerZ} is compared to the double word \emph{registerY} using condition \emph{intcomp}. The boolean result extended to double word is stored into the \emph{registerW}.


\paragraph{Modifier(s)}
\noindent\begin{itemize}
\item intcomp (cf. intcomp in section \ref{par:modifier-intcomp} on page \pageref{par:modifier-intcomp})
\end{itemize}

\paragraph{Execution} {Reservation table \textsf{ALU\_TINY} (Section~\ref{sec:reservation-ALU-TINY})}
~
{\begin{lstlisting}
stage ID:
  new argument4 = _ZX_4(intcomp);
stage RR:
  new argument3 = GPR[registerY];
  new argument2 = GPR[registerZ];
  new result1 = intcomp_64(argument4, argument2, argument3);
stage E1:
  GPR[registerW] = result1;
\end{lstlisting}}
\newpage

\paragraph{Encoding} Format \textsf{ALU\_DCWRR.W} (Section~\ref{sec:format-ALU-DCWRR.W}), \verb|steering = 11|\\

\bitfields{ 1/P, 2/00, 2/-- --, 27/upper27, 1/1, 2/11, 1/0, 4/intcomp, 6/registerW, 2/10, 3/001, 1/0, 1/0, 5/lower5, 6/registerZ }


\paragraph{Syntax} \texttt{compd}\emph{intcomp} \emph{registerW} = \emph{registerZ}, \emph{upper27\_\discretionary{}{}{}lower5}

\paragraph{Description} The double word \emph{registerZ} is compared to the double word \emph{upper27\_\discretionary{}{}{}lower5} using condition \emph{intcomp}. The boolean result extended to double word is stored into the \emph{registerW}.


\paragraph{Modifier(s)}
\noindent\begin{itemize}
\item intcomp (cf. intcomp in section \ref{par:modifier-intcomp} on page \pageref{par:modifier-intcomp})
\end{itemize}

\paragraph{Execution} {Reservation table \textsf{ALU\_TINY.X} (Section~\ref{sec:reservation-ALU-TINY.X})}
~
{\begin{lstlisting}
stage ID:
  new argument4 = _ZX_4(intcomp);
stage RR:
  new argument3 = _WX_32(upper27_lower5);
  new argument2 = GPR[registerZ];
  new result1 = intcomp_64(argument4, argument2, argument3);
stage E1:
  GPR[registerW] = result1;
\end{lstlisting}}
\newpage
\subsubsection[{\small COMPDP: Compare Double Word Pair}]{COMPDP\hfill Compare Double Word Pair}
\label{instr:COMPDP}\index{compdp}
 
\paragraph{Encoding} Format \textsf{ALU\_DPCWRR} (Section~\ref{sec:format-ALU-DPCWRR}), \verb|alucode3 = 001|\\

\bitfields{ 1/P, 2/11, 1/1, 4/intcomp, 6/registerW, 2/10, 3/001, 1/0, 5/registerO, 1/--, 5/registerP, 1/0 }


\paragraph{Syntax} \texttt{compdp}\emph{intcomp} \emph{registerW} = \emph{registerP}, \emph{registerO}

\paragraph{Description} The double word pair \emph{registerP} is compared to the double word pair \emph{registerO} using condition \emph{intcomp}. The two boolean result bits are packed and stored into the \emph{registerW}.


\paragraph{Modifier(s)}
\noindent\begin{itemize}
\item intcomp (cf. intcomp in section \ref{par:modifier-intcomp} on page \pageref{par:modifier-intcomp})
\end{itemize}

\paragraph{Execution} {Reservation table \textsf{ALU\_LITE} (Section~\ref{sec:reservation-ALU-LITE})}
~
{\begin{lstlisting}
stage ID:
  new argument4 = _ZX_4(intcomp);
stage RR:
  new argument3 = PGR[registerO];
  new argument2 = PGR[registerP];
  for (I = 0; I < 2; I += 1) {
    result1.1[I] = intcomp_64(argument4, argument2.64[I], argument3.64[I])
  };
stage E1:
  GPR[registerW] = result1;
\end{lstlisting}}
\newpage

\paragraph{Encoding} Format \textsf{ALU\_DPCWRR.M} (Section~\ref{sec:format-ALU-DPCWRR.M}), \verb|alucode3 = 001|\\

\bitfields{ 1/P, 2/00, 2/-- --, 27/upper27, 1/1, 2/11, 1/1, 4/intcomp, 6/registerW, 2/10, 3/001, 1/0, 1/splat32, 5/lower5, 5/registerP, 1/0 }


\paragraph{Syntax} \texttt{compdp}\emph{intcomp} \emph{registerW} = \emph{registerP}, \emph{upper27\_\discretionary{}{}{}lower5}\emph{splat32}

\paragraph{Description} The double word pair \emph{registerP} is compared to the double word pair \emph{upper27\_\discretionary{}{}{}lower5} using condition \emph{intcomp}. The two boolean result bits are packed and stored into the \emph{registerW}.


\paragraph{Modifier(s)}
\noindent\begin{itemize}
\item intcomp (cf. intcomp in section \ref{par:modifier-intcomp} on page \pageref{par:modifier-intcomp})
\item splat32 (cf. splat32 in section \ref{par:modifier-splat32} on page \pageref{par:modifier-splat32})
\end{itemize}

\paragraph{Execution} {Reservation table \textsf{ALU\_LITE.X} (Section~\ref{sec:reservation-ALU-LITE.X})}
~
{\begin{lstlisting}
stage ID:
  new splat32 = _ZX_1(splat32);
  new argument4 = _ZX_4(intcomp);
stage RR:
  new argument3 = splat32 ? (_WX_32(upper27_lower5) << 32) | _ZX_32(_WX_32(upper27_lower5)) : _SX_32(_WX_32(upper27_lower5));
  new argument2 = PGR[registerP];
  for (I = 0; I < 2; I += 1) {
    result1.1[I] = intcomp_64(argument4, argument2.64[I], argument3.64[I])
  };
stage E1:
  GPR[registerW] = result1;
\end{lstlisting}}
\newpage
\subsubsection[{\small COMPHO: Compare Half Word Octuple}]{COMPHO\hfill Compare Half Word Octuple}
\label{instr:COMPHO}\index{compho}
 
\paragraph{Encoding} Format \textsf{ALU\_HOCWRR} (Section~\ref{sec:format-ALU-HOCWRR}), \verb|alucode3 = 001|\\

\bitfields{ 1/P, 2/11, 1/1, 4/intcomp, 6/registerW, 2/10, 3/001, 1/1, 5/registerO, 1/--, 5/registerP, 1/0 }


\paragraph{Syntax} \texttt{compho}\emph{intcomp} \emph{registerW} = \emph{registerP}, \emph{registerO}

\paragraph{Description} The half word octuple \emph{registerP} is compared to the half word octuple \emph{registerO} using condition \emph{intcomp}. The eight boolean result bits are packed and stored into the \emph{registerW}.


\paragraph{Modifier(s)}
\noindent\begin{itemize}
\item intcomp (cf. intcomp in section \ref{par:modifier-intcomp} on page \pageref{par:modifier-intcomp})
\end{itemize}

\paragraph{Execution} {Reservation table \textsf{ALU\_LITE} (Section~\ref{sec:reservation-ALU-LITE})}
~
{\begin{lstlisting}
stage ID:
  new argument4 = _ZX_4(intcomp);
stage RR:
  new argument3 = PGR[registerO];
  new argument2 = PGR[registerP];
  for (I = 0; I < 8; I += 1) {
    result1.1[I] = intcomp_16(argument4, argument2.16[I], argument3.16[I])
  };
stage E1:
  GPR[registerW] = result1;
\end{lstlisting}}
\newpage

\paragraph{Encoding} Format \textsf{ALU\_HOCWRR.M} (Section~\ref{sec:format-ALU-HOCWRR.M}), \verb|alucode3 = 001|\\

\bitfields{ 1/P, 2/00, 2/-- --, 27/upper27, 1/1, 2/11, 1/1, 4/intcomp, 6/registerW, 2/10, 3/001, 1/1, 1/splat32, 5/lower5, 5/registerP, 1/0 }


\paragraph{Syntax} \texttt{compho}\emph{intcomp} \emph{registerW} = \emph{registerP}, \emph{upper27\_\discretionary{}{}{}lower5}\emph{splat32}

\paragraph{Description} The half word octuple \emph{registerP} is compared to the half word octuple \emph{upper27\_\discretionary{}{}{}lower5} using condition \emph{intcomp}. The eight boolean result bits are packed and stored into the \emph{registerW}.


\paragraph{Modifier(s)}
\noindent\begin{itemize}
\item intcomp (cf. intcomp in section \ref{par:modifier-intcomp} on page \pageref{par:modifier-intcomp})
\item splat32 (cf. splat32 in section \ref{par:modifier-splat32} on page \pageref{par:modifier-splat32})
\end{itemize}

\paragraph{Execution} {Reservation table \textsf{ALU\_LITE.X} (Section~\ref{sec:reservation-ALU-LITE.X})}
~
{\begin{lstlisting}
stage ID:
  new splat32 = _ZX_1(splat32);
  new argument4 = _ZX_4(intcomp);
stage RR:
  new argument3 = splat32 ? (_WX_32(upper27_lower5) << 32) | _ZX_32(_WX_32(upper27_lower5)) : _SX_32(_WX_32(upper27_lower5));
  new argument2 = PGR[registerP];
  for (I = 0; I < 8; I += 1) {
    result1.1[I] = intcomp_16(argument4, argument2.16[I], argument3.16[I])
  };
stage E1:
  GPR[registerW] = result1;
\end{lstlisting}}
\newpage
\subsubsection[{\small COMPNBX: Compare Negate Byte Hexadecuple}]{COMPNBX\hfill Compare Negate Byte Hexadecuple}
\label{instr:COMPNBX}\index{compnbx}
 
\paragraph{Encoding} Format \textsf{ALU\_BXCNWRR} (Section~\ref{sec:format-ALU-BXCNWRR}), \verb|alucode3 = 010|\\

\bitfields{ 1/P, 2/11, 1/1, 4/intcomp, 5/registerM, 1/--, 2/10, 3/010, 1/1, 5/registerO, 1/--, 5/registerP, 1/1 }


\paragraph{Syntax} \texttt{compnbx}\emph{intcomp} \emph{registerM} = \emph{registerP}, \emph{registerO}

\paragraph{Description} The \emph{registerP} and the \emph{registerO} are interpreted as sixteen bytes packed into 128 bits. The \emph{registerP} bytes are compared to the \emph{registerO} bytes using condition \emph{intcomp}. The sixteen resulting boolean bytes are negated, packed and stored into the \emph{registerM}.


\paragraph{Modifier(s)}
\noindent\begin{itemize}
\item intcomp (cf. intcomp in section \ref{par:modifier-intcomp} on page \pageref{par:modifier-intcomp})
\end{itemize}

\paragraph{Execution} {Reservation table \textsf{ALU\_LITE} (Section~\ref{sec:reservation-ALU-LITE})}
~
{\begin{lstlisting}
stage ID:
  new argument4 = _ZX_4(intcomp);
stage RR:
  new argument3 = PGR[registerO];
  new argument2 = PGR[registerP];
  for (I = 0; I < 16; I += 1) {
    result1.8[I] = -intcomp_8(argument4, argument2.8[I], argument3.8[I])
  };
stage E1:
  PGR[registerM] = result1;
\end{lstlisting}}
\newpage

\paragraph{Encoding} Format \textsf{ALU\_BXCNWRR.M} (Section~\ref{sec:format-ALU-BXCNWRR.M}), \verb|alucode3 = 010|\\

\bitfields{ 1/P, 2/00, 2/-- --, 27/upper27, 1/1, 2/11, 1/1, 4/intcomp, 5/registerM, 1/--, 2/10, 3/010, 1/1, 1/splat32, 5/lower5, 5/registerP, 1/1 }


\paragraph{Syntax} \texttt{compnbx}\emph{intcomp} \emph{registerM} = \emph{registerP}, \emph{upper27\_\discretionary{}{}{}lower5}\emph{splat32}

\paragraph{Description} The \emph{registerP} and the \emph{upper27\_\discretionary{}{}{}lower5} are interpreted as sixteen bytes packed into 128 bits. The \emph{registerP} bytes are compared to the \emph{upper27\_\discretionary{}{}{}lower5} bytes using condition \emph{intcomp}. The sixteen resulting boolean bytes are negated, packed and stored into the \emph{registerM}.


\paragraph{Modifier(s)}
\noindent\begin{itemize}
\item intcomp (cf. intcomp in section \ref{par:modifier-intcomp} on page \pageref{par:modifier-intcomp})
\item splat32 (cf. splat32 in section \ref{par:modifier-splat32} on page \pageref{par:modifier-splat32})
\end{itemize}

\paragraph{Execution} {Reservation table \textsf{ALU\_LITE.X} (Section~\ref{sec:reservation-ALU-LITE.X})}
~
{\begin{lstlisting}
stage ID:
  new splat32 = _ZX_1(splat32);
  new argument4 = _ZX_4(intcomp);
stage RR:
  new argument3 = splat32 ? (_WX_32(upper27_lower5) << 32) | _ZX_32(_WX_32(upper27_lower5)) : _SX_32(_WX_32(upper27_lower5));
  new argument2 = PGR[registerP];
  for (I = 0; I < 16; I += 1) {
    result1.8[I] = -intcomp_8(argument4, argument2.8[I], argument3.8[I])
  };
stage E1:
  PGR[registerM] = result1;
\end{lstlisting}}
\newpage
\subsubsection[{\small COMPNDP: Compare Negate Double Word Pair}]{COMPNDP\hfill Compare Negate Double Word Pair}
\label{instr:COMPNDP}\index{compndp}
 
\paragraph{Encoding} Format \textsf{ALU\_DPCNWRR} (Section~\ref{sec:format-ALU-DPCNWRR}), \verb|alucode3 = 010|\\

\bitfields{ 1/P, 2/11, 1/1, 4/intcomp, 5/registerM, 1/--, 2/10, 3/010, 1/0, 5/registerO, 1/--, 5/registerP, 1/0 }


\paragraph{Syntax} \texttt{compndp}\emph{intcomp} \emph{registerM} = \emph{registerP}, \emph{registerO}

\paragraph{Description} The double word pair \emph{registerP} is compared to the double word pair \emph{registerO} using condition \emph{intcomp}. The two resulting boolean double words are negated, packed and stored into the \emph{registerM}.


\paragraph{Modifier(s)}
\noindent\begin{itemize}
\item intcomp (cf. intcomp in section \ref{par:modifier-intcomp} on page \pageref{par:modifier-intcomp})
\end{itemize}

\paragraph{Execution} {Reservation table \textsf{ALU\_LITE} (Section~\ref{sec:reservation-ALU-LITE})}
~
{\begin{lstlisting}
stage ID:
  new argument4 = _ZX_4(intcomp);
stage RR:
  new argument3 = PGR[registerO];
  new argument2 = PGR[registerP];
  for (I = 0; I < 2; I += 1) {
    result1.64[I] = -intcomp_64(argument4, argument2.64[I], argument3.64[I])
  };
stage E1:
  PGR[registerM] = result1;
\end{lstlisting}}
\newpage

\paragraph{Encoding} Format \textsf{ALU\_DPCNWRR.M} (Section~\ref{sec:format-ALU-DPCNWRR.M}), \verb|alucode3 = 010|\\

\bitfields{ 1/P, 2/00, 2/-- --, 27/upper27, 1/1, 2/11, 1/1, 4/intcomp, 5/registerM, 1/--, 2/10, 3/010, 1/0, 1/splat32, 5/lower5, 5/registerP, 1/0 }


\paragraph{Syntax} \texttt{compndp}\emph{intcomp} \emph{registerM} = \emph{registerP}, \emph{upper27\_\discretionary{}{}{}lower5}\emph{splat32}

\paragraph{Description} The double word pair \emph{registerP} is compared to the double word pair \emph{upper27\_\discretionary{}{}{}lower5} using condition \emph{intcomp}. The two resulting boolean double words are negated, packed and stored into the \emph{registerM}.


\paragraph{Modifier(s)}
\noindent\begin{itemize}
\item intcomp (cf. intcomp in section \ref{par:modifier-intcomp} on page \pageref{par:modifier-intcomp})
\item splat32 (cf. splat32 in section \ref{par:modifier-splat32} on page \pageref{par:modifier-splat32})
\end{itemize}

\paragraph{Execution} {Reservation table \textsf{ALU\_LITE.X} (Section~\ref{sec:reservation-ALU-LITE.X})}
~
{\begin{lstlisting}
stage ID:
  new splat32 = _ZX_1(splat32);
  new argument4 = _ZX_4(intcomp);
stage RR:
  new argument3 = splat32 ? (_WX_32(upper27_lower5) << 32) | _ZX_32(_WX_32(upper27_lower5)) : _SX_32(_WX_32(upper27_lower5));
  new argument2 = PGR[registerP];
  for (I = 0; I < 2; I += 1) {
    result1.64[I] = -intcomp_64(argument4, argument2.64[I], argument3.64[I])
  };
stage E1:
  PGR[registerM] = result1;
\end{lstlisting}}
\newpage
\subsubsection[{\small COMPNHO: Compare Negate Half Word Octuple}]{COMPNHO\hfill Compare Negate Half Word Octuple}
\label{instr:COMPNHO}\index{compnho}
 
\paragraph{Encoding} Format \textsf{ALU\_HOCNWRR} (Section~\ref{sec:format-ALU-HOCNWRR}), \verb|alucode3 = 010|\\

\bitfields{ 1/P, 2/11, 1/1, 4/intcomp, 5/registerM, 1/--, 2/10, 3/010, 1/1, 5/registerO, 1/--, 5/registerP, 1/0 }


\paragraph{Syntax} \texttt{compnho}\emph{intcomp} \emph{registerM} = \emph{registerP}, \emph{registerO}

\paragraph{Description} The \emph{registerP} and the \emph{registerO} are interpreted as eight half words packed into 128 bits. The \emph{registerP} half words are compared to the \emph{registerO} half words using condition \emph{intcomp}. The eight resulting boolean half word are negated, packed and stored into the \emph{registerM}.


\paragraph{Modifier(s)}
\noindent\begin{itemize}
\item intcomp (cf. intcomp in section \ref{par:modifier-intcomp} on page \pageref{par:modifier-intcomp})
\end{itemize}

\paragraph{Execution} {Reservation table \textsf{ALU\_LITE} (Section~\ref{sec:reservation-ALU-LITE})}
~
{\begin{lstlisting}
stage ID:
  new argument4 = _ZX_4(intcomp);
stage RR:
  new argument3 = PGR[registerO];
  new argument2 = PGR[registerP];
  for (I = 0; I < 8; I += 1) {
    result1.16[I] = -intcomp_16(argument4, argument2.16[I], argument3.16[I])
  };
stage E1:
  PGR[registerM] = result1;
\end{lstlisting}}
\newpage

\paragraph{Encoding} Format \textsf{ALU\_HOCNWRR.M} (Section~\ref{sec:format-ALU-HOCNWRR.M}), \verb|alucode3 = 010|\\

\bitfields{ 1/P, 2/00, 2/-- --, 27/upper27, 1/1, 2/11, 1/1, 4/intcomp, 5/registerM, 1/--, 2/10, 3/010, 1/1, 1/splat32, 5/lower5, 5/registerP, 1/0 }


\paragraph{Syntax} \texttt{compnho}\emph{intcomp} \emph{registerM} = \emph{registerP}, \emph{upper27\_\discretionary{}{}{}lower5}\emph{splat32}

\paragraph{Description} The \emph{registerP} and the \emph{upper27\_\discretionary{}{}{}lower5} are interpreted as eight half words packed into 128 bits. The \emph{registerP} half words are compared to the \emph{upper27\_\discretionary{}{}{}lower5} half words using condition \emph{intcomp}. The eight resulting boolean half word are negated, packed and stored into the \emph{registerM}.


\paragraph{Modifier(s)}
\noindent\begin{itemize}
\item intcomp (cf. intcomp in section \ref{par:modifier-intcomp} on page \pageref{par:modifier-intcomp})
\item splat32 (cf. splat32 in section \ref{par:modifier-splat32} on page \pageref{par:modifier-splat32})
\end{itemize}

\paragraph{Execution} {Reservation table \textsf{ALU\_LITE.X} (Section~\ref{sec:reservation-ALU-LITE.X})}
~
{\begin{lstlisting}
stage ID:
  new splat32 = _ZX_1(splat32);
  new argument4 = _ZX_4(intcomp);
stage RR:
  new argument3 = splat32 ? (_WX_32(upper27_lower5) << 32) | _ZX_32(_WX_32(upper27_lower5)) : _SX_32(_WX_32(upper27_lower5));
  new argument2 = PGR[registerP];
  for (I = 0; I < 8; I += 1) {
    result1.16[I] = -intcomp_16(argument4, argument2.16[I], argument3.16[I])
  };
stage E1:
  PGR[registerM] = result1;
\end{lstlisting}}
\newpage
\subsubsection[{\small COMPNWQ: Compare Negate Word Quadruple}]{COMPNWQ\hfill Compare Negate Word Quadruple}
\label{instr:COMPNWQ}\index{compnwq}
 
\paragraph{Encoding} Format \textsf{ALU\_WQCNWRR} (Section~\ref{sec:format-ALU-WQCNWRR}), \verb|alucode3 = 010|\\

\bitfields{ 1/P, 2/11, 1/1, 4/intcomp, 5/registerM, 1/--, 2/10, 3/010, 1/0, 5/registerO, 1/--, 5/registerP, 1/1 }


\paragraph{Syntax} \texttt{compnwq}\emph{intcomp} \emph{registerM} = \emph{registerP}, \emph{registerO}

\paragraph{Description} The \emph{registerP} and the \emph{registerO} are interpreted as four words packed into 128 bits. The \emph{registerP} words are compared to the \emph{registerO} words using condition \emph{intcomp}. The four resulting boolean words are negated, packed and stored into the \emph{registerM}.


\paragraph{Modifier(s)}
\noindent\begin{itemize}
\item intcomp (cf. intcomp in section \ref{par:modifier-intcomp} on page \pageref{par:modifier-intcomp})
\end{itemize}

\paragraph{Execution} {Reservation table \textsf{ALU\_LITE} (Section~\ref{sec:reservation-ALU-LITE})}
~
{\begin{lstlisting}
stage ID:
  new argument4 = _ZX_4(intcomp);
stage RR:
  new argument3 = PGR[registerO];
  new argument2 = PGR[registerP];
  for (I = 0; I < 4; I += 1) {
    result1.32[I] = -intcomp_32(argument4, argument2.32[I], argument3.32[I])
  };
stage E1:
  PGR[registerM] = result1;
\end{lstlisting}}
\newpage

\paragraph{Encoding} Format \textsf{ALU\_WQCNWRR.M} (Section~\ref{sec:format-ALU-WQCNWRR.M}), \verb|alucode3 = 010|\\

\bitfields{ 1/P, 2/00, 2/-- --, 27/upper27, 1/1, 2/11, 1/1, 4/intcomp, 5/registerM, 1/--, 2/10, 3/010, 1/0, 1/splat32, 5/lower5, 5/registerP, 1/1 }


\paragraph{Syntax} \texttt{compnwq}\emph{intcomp} \emph{registerM} = \emph{registerP}, \emph{upper27\_\discretionary{}{}{}lower5}\emph{splat32}

\paragraph{Description} The \emph{registerP} and the \emph{upper27\_\discretionary{}{}{}lower5} are interpreted as four words packed into 128 bits. The \emph{registerP} words are compared to the \emph{upper27\_\discretionary{}{}{}lower5} words using condition \emph{intcomp}. The four resulting boolean words are negated, packed and stored into the \emph{registerM}.


\paragraph{Modifier(s)}
\noindent\begin{itemize}
\item intcomp (cf. intcomp in section \ref{par:modifier-intcomp} on page \pageref{par:modifier-intcomp})
\item splat32 (cf. splat32 in section \ref{par:modifier-splat32} on page \pageref{par:modifier-splat32})
\end{itemize}

\paragraph{Execution} {Reservation table \textsf{ALU\_LITE.X} (Section~\ref{sec:reservation-ALU-LITE.X})}
~
{\begin{lstlisting}
stage ID:
  new splat32 = _ZX_1(splat32);
  new argument4 = _ZX_4(intcomp);
stage RR:
  new argument3 = splat32 ? (_WX_32(upper27_lower5) << 32) | _ZX_32(_WX_32(upper27_lower5)) : _SX_32(_WX_32(upper27_lower5));
  new argument2 = PGR[registerP];
  for (I = 0; I < 4; I += 1) {
    result1.32[I] = -intcomp_32(argument4, argument2.32[I], argument3.32[I])
  };
stage E1:
  PGR[registerM] = result1;
\end{lstlisting}}
\newpage
\subsubsection[{\small COMPQ: Compare Quadruple Word}]{COMPQ\hfill Compare Quadruple Word}
\label{instr:COMPQ}\index{compq}
 
\paragraph{Encoding} Format \textsf{ALU\_QCWRR} (Section~\ref{sec:format-ALU-QCWRR}), \verb|steering = 11|\\

\bitfields{ 1/P, 2/11, 1/0, 4/intcomp, 5/registerM, 1/--, 2/10, 3/001, 1/1, 5/registerO, 1/--, 5/registerP, 1/-- }


\paragraph{Syntax} \texttt{compq}\emph{intcomp} \emph{registerM} = \emph{registerP}, \emph{registerO}

\paragraph{Description} The quadruple word \emph{registerP} is compared to the quadruple word \emph{registerO} using condition \emph{intcomp}. The boolean result extended to double word is stored into the \emph{registerM}.


\paragraph{Modifier(s)}
\noindent\begin{itemize}
\item intcomp (cf. intcomp in section \ref{par:modifier-intcomp} on page \pageref{par:modifier-intcomp})
\end{itemize}

\paragraph{Execution} {Reservation table \textsf{ALU\_LITE} (Section~\ref{sec:reservation-ALU-LITE})}
~
{\begin{lstlisting}
stage ID:
  new argument4 = _ZX_4(intcomp);
stage RR:
  new argument3 = PGR[registerO];
  new argument2 = PGR[registerP];
  new result1 = intcomp_128(argument4, argument2, argument3);
stage E1:
  PGR[registerM] = result1;
\end{lstlisting}}
\newpage
\subsubsection[{\small COMPW: Compare Word}]{COMPW\hfill Compare Word}
\label{instr:COMPW}\index{compw}
 
\paragraph{Encoding} Format \textsf{ALU\_WCWRR} (Section~\ref{sec:format-ALU-WCWRR}), \verb|steering = 11|\\

\bitfields{ 1/P, 2/11, 1/0, 4/intcomp, 6/registerW, 2/11, 3/001, 1/signextw, 6/registerY, 6/registerZ }


\paragraph{Syntax} \texttt{compw}\emph{intcomp}\emph{signextw} \emph{registerW} = \emph{registerZ}, \emph{registerY}

\paragraph{Description} The word \emph{registerZ} is compared to the word \emph{registerY} using condition \emph{intcomp}. The boolean result extended to word is stored into the \emph{registerW}.


\paragraph{Modifier(s)}
\noindent\begin{itemize}
\item intcomp (cf. intcomp in section \ref{par:modifier-intcomp} on page \pageref{par:modifier-intcomp})
\item signextw (cf. signextw in section \ref{par:modifier-signextw} on page \pageref{par:modifier-signextw})
\end{itemize}

\paragraph{Execution} {Reservation table \textsf{ALU\_TINY} (Section~\ref{sec:reservation-ALU-TINY})}
~
{\begin{lstlisting}
stage ID:
  new signextw = _ZX_1(signextw);
  new argument4 = _ZX_4(intcomp);
stage RR:
  new argument3 = GPR[registerY];
  new argument2 = GPR[registerZ];
  new result1 = intcomp_32(argument4, argument2, argument3);
stage E1:
  GPR[registerW] = signextw ? _SX_32(result1) : _ZX_32(result1);
\end{lstlisting}}
\newpage

\paragraph{Encoding} Format \textsf{ALU\_WCWRR.W} (Section~\ref{sec:format-ALU-WCWRR.W}), \verb|steering = 11|\\

\bitfields{ 1/P, 2/00, 2/-- --, 27/upper27, 1/1, 2/11, 1/0, 4/intcomp, 6/registerW, 2/11, 3/001, 1/signextw, 1/0, 5/lower5, 6/registerZ }


\paragraph{Syntax} \texttt{compw}\emph{intcomp}\emph{signextw} \emph{registerW} = \emph{registerZ}, \emph{upper27\_\discretionary{}{}{}lower5}

\paragraph{Description} The word \emph{registerZ} is compared to the word \emph{upper27\_\discretionary{}{}{}lower5} using condition \emph{intcomp}. The boolean result extended to word is stored into the \emph{registerW}.


\paragraph{Modifier(s)}
\noindent\begin{itemize}
\item intcomp (cf. intcomp in section \ref{par:modifier-intcomp} on page \pageref{par:modifier-intcomp})
\item signextw (cf. signextw in section \ref{par:modifier-signextw} on page \pageref{par:modifier-signextw})
\end{itemize}

\paragraph{Execution} {Reservation table \textsf{ALU\_TINY.X} (Section~\ref{sec:reservation-ALU-TINY.X})}
~
{\begin{lstlisting}
stage ID:
  new signextw = _ZX_1(signextw);
  new argument4 = _ZX_4(intcomp);
stage RR:
  new argument3 = _WX_32(upper27_lower5);
  new argument2 = GPR[registerZ];
  new result1 = intcomp_32(argument4, argument2, argument3);
stage E1:
  GPR[registerW] = signextw ? _SX_32(result1) : _ZX_32(result1);
\end{lstlisting}}
\newpage
\subsubsection[{\small COMPWQ: Compare Word Quadruple}]{COMPWQ\hfill Compare Word Quadruple}
\label{instr:COMPWQ}\index{compwq}
 
\paragraph{Encoding} Format \textsf{ALU\_WQCWRR} (Section~\ref{sec:format-ALU-WQCWRR}), \verb|alucode3 = 001|\\

\bitfields{ 1/P, 2/11, 1/1, 4/intcomp, 6/registerW, 2/10, 3/001, 1/0, 5/registerO, 1/--, 5/registerP, 1/1 }


\paragraph{Syntax} \texttt{compwq}\emph{intcomp} \emph{registerW} = \emph{registerP}, \emph{registerO}

\paragraph{Description} The word quadruple \emph{registerP} is compared to the word quadruple \emph{registerO} using condition \emph{intcomp}. The four boolean result bits are packed and stored into the \emph{registerW}.


\paragraph{Modifier(s)}
\noindent\begin{itemize}
\item intcomp (cf. intcomp in section \ref{par:modifier-intcomp} on page \pageref{par:modifier-intcomp})
\end{itemize}

\paragraph{Execution} {Reservation table \textsf{ALU\_TINY} (Section~\ref{sec:reservation-ALU-TINY})}
~
{\begin{lstlisting}
stage ID:
  new argument4 = _ZX_4(intcomp);
stage RR:
  new argument3 = PGR[registerO];
  new argument2 = PGR[registerP];
  for (I = 0; I < 4; I += 1) {
    result1.1[I] = intcomp_32(argument4, argument2.32[I], argument3.32[I])
  };
stage E1:
  GPR[registerW] = result1;
\end{lstlisting}}
\newpage

\paragraph{Encoding} Format \textsf{ALU\_WQCWRR.M} (Section~\ref{sec:format-ALU-WQCWRR.M}), \verb|alucode3 = 001|\\

\bitfields{ 1/P, 2/00, 2/-- --, 27/upper27, 1/1, 2/11, 1/1, 4/intcomp, 6/registerW, 2/10, 3/001, 1/0, 1/splat32, 5/lower5, 5/registerP, 1/1 }


\paragraph{Syntax} \texttt{compwq}\emph{intcomp} \emph{registerW} = \emph{registerP}, \emph{upper27\_\discretionary{}{}{}lower5}\emph{splat32}

\paragraph{Description} The word quadruple \emph{registerP} is compared to the word quadruple \emph{upper27\_\discretionary{}{}{}lower5} using condition \emph{intcomp}. The four boolean result bits are packed and stored into the \emph{registerW}.


\paragraph{Modifier(s)}
\noindent\begin{itemize}
\item intcomp (cf. intcomp in section \ref{par:modifier-intcomp} on page \pageref{par:modifier-intcomp})
\item splat32 (cf. splat32 in section \ref{par:modifier-splat32} on page \pageref{par:modifier-splat32})
\end{itemize}

\paragraph{Execution} {Reservation table \textsf{ALU\_TINY.X} (Section~\ref{sec:reservation-ALU-TINY.X})}
~
{\begin{lstlisting}
stage ID:
  new splat32 = _ZX_1(splat32);
  new argument4 = _ZX_4(intcomp);
stage RR:
  new argument3 = splat32 ? (_WX_32(upper27_lower5) << 32) | _ZX_32(_WX_32(upper27_lower5)) : _SX_32(_WX_32(upper27_lower5));
  new argument2 = PGR[registerP];
  for (I = 0; I < 4; I += 1) {
    result1.1[I] = intcomp_32(argument4, argument2.32[I], argument3.32[I])
  };
stage E1:
  GPR[registerW] = result1;
\end{lstlisting}}
\newpage
\subsubsection[{\small CRCBELLW: Cyclic Redundancy Check Big-Endian Least Significant with Least Significant Word}]{CRCBELLW\hfill Cyclic Redundancy Check Big-Endian Least Significant with Least Significant Word}
\label{instr:CRCBELLW}\index{crcbellw}
 
\paragraph{Encoding} Format \textsf{ALU\_WCWRRR} (Section~\ref{sec:format-ALU-WCWRRR}), \verb|exucode2 = 0101|\\

\bitfields{ 1/P, 2/11, 1/0, 4/0101, 6/registerW, 2/11, 3/010, 1/1, 6/registerY, 6/registerZ }


\paragraph{Syntax} \texttt{crcbellw} \emph{registerW} = \emph{registerZ}, \emph{registerY}

\paragraph{Description} The 4 least significant bytes of the \emph{registerW} are XOR-ed with the 4 least significant bytes of the \emph{registerZ} considered in reverse order. The result is reduced modulo the 4-byte polynomial stored in the 4 least significant bytes of the \emph{registerY}. The result is stored back into the \emph{registerW}.


\paragraph{Execution} {Reservation table \textsf{ALU\_LITE} (Section~\ref{sec:reservation-ALU-LITE})}
~
{\begin{lstlisting}
stage RR:
  new argument3 = GPR[registerY];
  new argument2 = GPR[registerZ];
stage E1:
  new argument1 = GPR[registerW];
  new argument1_data = argument1.32[0] ^ _SWAP_16(_SWAP_8(argument2));
  new result1 = crc32_be_u32(argument1_data.32[0], argument3.32[0]);
stage E2:
  GPR[registerW] = result1;
\end{lstlisting}}
\newpage

\paragraph{Encoding} Format \textsf{ALU\_WCWRRR.W} (Section~\ref{sec:format-ALU-WCWRRR.W}), \verb|exucode2 = 0101|\\

\bitfields{ 1/P, 2/00, 2/-- --, 27/upper27, 1/1, 2/11, 1/0, 4/0101, 6/registerW, 2/11, 3/010, 1/1, 1/0, 5/lower5, 6/registerZ }


\paragraph{Syntax} \texttt{crcbellw} \emph{registerW} = \emph{registerZ}, \emph{upper27\_\discretionary{}{}{}lower5}

\paragraph{Description} The 4 least significant bytes of the \emph{registerW} are XOR-ed with the 4 least significant bytes of the \emph{registerZ} considered in reverse order. The result is reduced modulo the 4-byte polynomial stored in the 4 least significant bytes of the \emph{upper27\_\discretionary{}{}{}lower5}. The result is stored back into the \emph{registerW}.


\paragraph{Execution} {Reservation table \textsf{ALU\_LITE.X} (Section~\ref{sec:reservation-ALU-LITE.X})}
~
{\begin{lstlisting}
stage RR:
  new argument3 = _WX_32(upper27_lower5);
  new argument2 = GPR[registerZ];
stage E1:
  new argument1 = GPR[registerW];
  new argument1_data = argument1.32[0] ^ _SWAP_16(_SWAP_8(argument2));
  new result1 = crc32_be_u32(argument1_data.32[0], argument3.32[0]);
stage E2:
  GPR[registerW] = result1;
\end{lstlisting}}
\newpage
\subsubsection[{\small CRCBELMW: Cyclic Redundancy Check Big-Endian Least Significant with Most Significant Word}]{CRCBELMW\hfill Cyclic Redundancy Check Big-Endian Least Significant with Most Significant Word}
\label{instr:CRCBELMW}\index{crcbelmw}
 
\paragraph{Encoding} Format \textsf{ALU\_WCWRRR} (Section~\ref{sec:format-ALU-WCWRRR}), \verb|exucode2 = 0100|\\

\bitfields{ 1/P, 2/11, 1/0, 4/0100, 6/registerW, 2/11, 3/010, 1/1, 6/registerY, 6/registerZ }


\paragraph{Syntax} \texttt{crcbelmw} \emph{registerW} = \emph{registerZ}, \emph{registerY}

\paragraph{Description} The 4 least significant bytes of the \emph{registerW} are XOR-ed with the 4 most significant bytes of the \emph{registerZ} considered in reverse order. The result is reduced modulo the 4-byte polynomial stored in the 4 least significant bytes of the \emph{registerY}. The result is stored back into the \emph{registerW}.


\paragraph{Execution} {Reservation table \textsf{ALU\_LITE} (Section~\ref{sec:reservation-ALU-LITE})}
~
{\begin{lstlisting}
stage RR:
  new argument3 = GPR[registerY];
  new argument2 = GPR[registerZ];
stage E1:
  new argument1 = GPR[registerW];
  new argument1_data = argument1.32[0] ^ _SWAP_16(_SWAP_8(argument2 >> 32));
  new result1 = crc32_be_u32(argument1_data.32[0], argument3.32[0]);
stage E2:
  GPR[registerW] = result1;
\end{lstlisting}}
\newpage

\paragraph{Encoding} Format \textsf{ALU\_WCWRRR.W} (Section~\ref{sec:format-ALU-WCWRRR.W}), \verb|exucode2 = 0100|\\

\bitfields{ 1/P, 2/00, 2/-- --, 27/upper27, 1/1, 2/11, 1/0, 4/0100, 6/registerW, 2/11, 3/010, 1/1, 1/0, 5/lower5, 6/registerZ }


\paragraph{Syntax} \texttt{crcbelmw} \emph{registerW} = \emph{registerZ}, \emph{upper27\_\discretionary{}{}{}lower5}

\paragraph{Description} The 4 least significant bytes of the \emph{registerW} are XOR-ed with the 4 most significant bytes of the \emph{registerZ} considered in reverse order. The result is reduced modulo the 4-byte polynomial stored in the 4 least significant bytes of the \emph{upper27\_\discretionary{}{}{}lower5}. The result is stored back into the \emph{registerW}.


\paragraph{Execution} {Reservation table \textsf{ALU\_LITE.X} (Section~\ref{sec:reservation-ALU-LITE.X})}
~
{\begin{lstlisting}
stage RR:
  new argument3 = _WX_32(upper27_lower5);
  new argument2 = GPR[registerZ];
stage E1:
  new argument1 = GPR[registerW];
  new argument1_data = argument1.32[0] ^ _SWAP_16(_SWAP_8(argument2 >> 32));
  new result1 = crc32_be_u32(argument1_data.32[0], argument3.32[0]);
stage E2:
  GPR[registerW] = result1;
\end{lstlisting}}
\newpage
\subsubsection[{\small CRCLELLW: Cyclic Redundancy Check Little-Endian Least Significant with Least Significant Word}]{CRCLELLW\hfill Cyclic Redundancy Check Little-Endian Least Significant with Least Significant Word}
\label{instr:CRCLELLW}\index{crclellw}
 
\paragraph{Encoding} Format \textsf{ALU\_WCWRRR} (Section~\ref{sec:format-ALU-WCWRRR}), \verb|exucode2 = 0111|\\

\bitfields{ 1/P, 2/11, 1/0, 4/0111, 6/registerW, 2/11, 3/010, 1/1, 6/registerY, 6/registerZ }


\paragraph{Syntax} \texttt{crclellw} \emph{registerW} = \emph{registerZ}, \emph{registerY}

\paragraph{Description} The 4 least significant bytes of the \emph{registerW} are XOR-ed with the 4 least significant bytes of the \emph{registerZ}. The result is bit reversed and then reduced modulo the 4-byte polynomial stored in the 4 least significant bytes of the \emph{registerY} bit reversed beforehand. The result is bit reversed and stored back into the \emph{registerW}.


\paragraph{Execution} {Reservation table \textsf{ALU\_LITE} (Section~\ref{sec:reservation-ALU-LITE})}
~
{\begin{lstlisting}
stage RR:
  new argument3 = GPR[registerY];
  new argument2 = GPR[registerZ];
stage E1:
  new argument1 = GPR[registerW];
  new argument1_data = argument1.32[0] ^ argument2.32[0];
  new result1 = reflect_32(crc32_be_u32(reflect_32(argument1_data.32[0]), reflect_32(argument3.32[0])));
stage E2:
  GPR[registerW] = result1;
\end{lstlisting}}
\newpage

\paragraph{Encoding} Format \textsf{ALU\_WCWRRR.W} (Section~\ref{sec:format-ALU-WCWRRR.W}), \verb|exucode2 = 0111|\\

\bitfields{ 1/P, 2/00, 2/-- --, 27/upper27, 1/1, 2/11, 1/0, 4/0111, 6/registerW, 2/11, 3/010, 1/1, 1/0, 5/lower5, 6/registerZ }


\paragraph{Syntax} \texttt{crclellw} \emph{registerW} = \emph{registerZ}, \emph{upper27\_\discretionary{}{}{}lower5}

\paragraph{Description} The 4 least significant bytes of the \emph{registerW} are XOR-ed with the 4 least significant bytes of the \emph{registerZ}. The result is bit reversed and then reduced modulo the 4-byte polynomial stored in the 4 least significant bytes of the \emph{upper27\_\discretionary{}{}{}lower5} bit reversed beforehand. The result is bit reversed and stored back into the \emph{registerW}.


\paragraph{Execution} {Reservation table \textsf{ALU\_LITE.X} (Section~\ref{sec:reservation-ALU-LITE.X})}
~
{\begin{lstlisting}
stage RR:
  new argument3 = _WX_32(upper27_lower5);
  new argument2 = GPR[registerZ];
stage E1:
  new argument1 = GPR[registerW];
  new argument1_data = argument1.32[0] ^ argument2.32[0];
  new result1 = reflect_32(crc32_be_u32(reflect_32(argument1_data.32[0]), reflect_32(argument3.32[0])));
stage E2:
  GPR[registerW] = result1;
\end{lstlisting}}
\newpage
\subsubsection[{\small CRCLELMW: Cyclic Redundancy Check Little-Endian Least Significant with Most Significant Word}]{CRCLELMW\hfill Cyclic Redundancy Check Little-Endian Least Significant with Most Significant Word}
\label{instr:CRCLELMW}\index{crclelmw}
 
\paragraph{Encoding} Format \textsf{ALU\_WCWRRR} (Section~\ref{sec:format-ALU-WCWRRR}), \verb|exucode2 = 0110|\\

\bitfields{ 1/P, 2/11, 1/0, 4/0110, 6/registerW, 2/11, 3/010, 1/1, 6/registerY, 6/registerZ }


\paragraph{Syntax} \texttt{crclelmw} \emph{registerW} = \emph{registerZ}, \emph{registerY}

\paragraph{Description} The 4 least significant bytes of the \emph{registerW} are XOR-ed with the 4 most significant bytes of the \emph{registerZ}. The result is bit reversed and then reduced modulo the 4-byte polynomial stored in the 4 least significant bytes of the \emph{registerY} bit reversed beforehand. The result is bit reversed and stored back into the \emph{registerW}.


\paragraph{Execution} {Reservation table \textsf{ALU\_LITE} (Section~\ref{sec:reservation-ALU-LITE})}
~
{\begin{lstlisting}
stage RR:
  new argument3 = GPR[registerY];
  new argument2 = GPR[registerZ];
stage E1:
  new argument1 = GPR[registerW];
  new argument1_data = argument1.32[0] ^ argument2.32[1];
  new result1 = reflect_32(crc32_be_u32(reflect_32(argument1_data.32[0]), reflect_32(argument3.32[0])));
stage E2:
  GPR[registerW] = result1;
\end{lstlisting}}
\newpage

\paragraph{Encoding} Format \textsf{ALU\_WCWRRR.W} (Section~\ref{sec:format-ALU-WCWRRR.W}), \verb|exucode2 = 0110|\\

\bitfields{ 1/P, 2/00, 2/-- --, 27/upper27, 1/1, 2/11, 1/0, 4/0110, 6/registerW, 2/11, 3/010, 1/1, 1/0, 5/lower5, 6/registerZ }


\paragraph{Syntax} \texttt{crclelmw} \emph{registerW} = \emph{registerZ}, \emph{upper27\_\discretionary{}{}{}lower5}

\paragraph{Description} The 4 least significant bytes of the \emph{registerW} are XOR-ed with the 4 most significant bytes of the \emph{registerZ}. The result is bit reversed and then reduced modulo the 4-byte polynomial stored in the 4 least significant bytes of the \emph{upper27\_\discretionary{}{}{}lower5} bit reversed beforehand. The result is bit reversed and stored back into the \emph{registerW}.


\paragraph{Execution} {Reservation table \textsf{ALU\_LITE.X} (Section~\ref{sec:reservation-ALU-LITE.X})}
~
{\begin{lstlisting}
stage RR:
  new argument3 = _WX_32(upper27_lower5);
  new argument2 = GPR[registerZ];
stage E1:
  new argument1 = GPR[registerW];
  new argument1_data = argument1.32[0] ^ argument2.32[1];
  new result1 = reflect_32(crc32_be_u32(reflect_32(argument1_data.32[0]), reflect_32(argument3.32[0])));
stage E2:
  GPR[registerW] = result1;
\end{lstlisting}}
\newpage
\subsubsection[{\small CTZD: Count Trailing Zero Double Word}]{CTZD\hfill Count Trailing Zero Double Word}
\label{instr:CTZD}\index{ctzd}
 
\paragraph{Encoding} Format \textsf{ALU\_DBWR} (Section~\ref{sec:format-ALU-DBWR}), \verb|alucode8 = 0111|\\

\bitfields{ 1/P, 2/11, 1/0, 4/1111, 6/registerW, 2/10, 3/101, 1/0, 4/0111, 1/0, 1/0, 6/registerZ }


\paragraph{Syntax} \texttt{ctzd} \emph{registerW} = \emph{registerZ}

\paragraph{Description} The trailing zero count value of the \emph{registerZ} is stored into the \emph{registerW}. When applied to zero, the result is 64.


\paragraph{Execution} {Reservation table \textsf{ALU\_TINY} (Section~\ref{sec:reservation-ALU-TINY})}
~
{\begin{lstlisting}
stage RR:
  new argument2 = GPR[registerZ];
  new result1 = _ZX_64(_CTZ_64(argument2));
stage E1:
  GPR[registerW] = result1;
\end{lstlisting}}
\newpage
\subsubsection[{\small CTZDP: Count Trailing Zeroes Double Word Pair}]{CTZDP\hfill Count Trailing Zeroes Double Word Pair}
\label{instr:CTZDP}\index{ctzdp}
 
\paragraph{Encoding} Format \textsf{ALU\_DPBWR} (Section~\ref{sec:format-ALU-DPBWR}), \verb|alucode8 = 0111|\\

\bitfields{ 1/P, 2/11, 1/1, 4/1111, 5/registerM, 1/--, 2/10, 3/101, 1/0, 4/0111, 1/--, 1/--, 5/registerP, 1/0 }


\paragraph{Syntax} \texttt{ctzdp} \emph{registerM} = \emph{registerP}

\paragraph{Description} The trailing zero count values of the \emph{registerP} interpreted as double word pair are stored into the \emph{registerM}. When applied to zero at the word level, the result is 64.


\paragraph{Execution} {Reservation table \textsf{ALU\_LITE} (Section~\ref{sec:reservation-ALU-LITE})}
~
{\begin{lstlisting}
stage RR:
  new argument2 = PGR[registerP];
  for (I = 0; I < 2; I += 1) {
    result1.64[I] = _ZX_64(_CTZ_64(argument2.64[I]))
  };
stage E1:
  PGR[registerM] = result1;
\end{lstlisting}}
\newpage
\subsubsection[{\small CTZHO: Count Trailing Zeroes Half Word Octuple}]{CTZHO\hfill Count Trailing Zeroes Half Word Octuple}
\label{instr:CTZHO}\index{ctzho}
 
\paragraph{Encoding} Format \textsf{ALU\_HOBWR} (Section~\ref{sec:format-ALU-HOBWR}), \verb|alucode8 = 0111|\\

\bitfields{ 1/P, 2/11, 1/1, 4/1111, 5/registerM, 1/--, 2/10, 3/101, 1/1, 4/0111, 1/--, 1/--, 5/registerP, 1/0 }


\paragraph{Syntax} \texttt{ctzho} \emph{registerM} = \emph{registerP}

\paragraph{Description} The trailing zero count values of the \emph{registerP} interpreted as half word octuple are stored into the \emph{registerM}. When applied to zero at the half word level, the result is 16.


\paragraph{Execution} {Reservation table \textsf{ALU\_LITE} (Section~\ref{sec:reservation-ALU-LITE})}
~
{\begin{lstlisting}
stage RR:
  new argument2 = PGR[registerP];
  for (I = 0; I < 8; I += 1) {
    result1.16[I] = _ZX_16(_CTZ_16(argument2.16[I]))
  };
stage E1:
  PGR[registerM] = result1;
\end{lstlisting}}
\newpage
\subsubsection[{\small CTZW: Count Trailing Zero Word}]{CTZW\hfill Count Trailing Zero Word}
\label{instr:CTZW}\index{ctzw}
 
\paragraph{Encoding} Format \textsf{ALU\_BWRW} (Section~\ref{sec:format-ALU-BWRW}), \verb|alucode8 = 0111|\\

\bitfields{ 1/P, 2/11, 1/0, 4/1111, 6/registerW, 2/11, 3/101, 1/signextw, 4/0111, 1/0, 1/0, 6/registerZ }


\paragraph{Syntax} \texttt{ctzw} \emph{registerW} = \emph{registerZ}

\paragraph{Description} The trailing zero count value of the \emph{registerZ} is stored into the \emph{registerW}. When applied to zero, the result is 32.


\paragraph{Execution} {Reservation table \textsf{ALU\_TINY} (Section~\ref{sec:reservation-ALU-TINY})}
~
{\begin{lstlisting}
stage RR:
  new argument2 = GPR[registerZ];
  new result1 = _ZX_32(_CTZ_32(argument2));
stage E1:
  GPR[registerW] = _ZX_32(result1);
\end{lstlisting}}
\newpage
\subsubsection[{\small CTZWQ: Count Trailing Zeroes Word Quadruple}]{CTZWQ\hfill Count Trailing Zeroes Word Quadruple}
\label{instr:CTZWQ}\index{ctzwq}
 
\paragraph{Encoding} Format \textsf{ALU\_WQBWR} (Section~\ref{sec:format-ALU-WQBWR}), \verb|alucode8 = 0111|\\

\bitfields{ 1/P, 2/11, 1/1, 4/1111, 5/registerM, 1/--, 2/10, 3/101, 1/0, 4/0111, 1/--, 1/--, 5/registerP, 1/1 }


\paragraph{Syntax} \texttt{ctzwq} \emph{registerM} = \emph{registerP}

\paragraph{Description} The trailing zero count values of the \emph{registerP} interpreted as word quadruple are stored into the \emph{registerM}. When applied to zero at the word level, the result is 32.


\paragraph{Execution} {Reservation table \textsf{ALU\_LITE} (Section~\ref{sec:reservation-ALU-LITE})}
~
{\begin{lstlisting}
stage RR:
  new argument2 = PGR[registerP];
  for (I = 0; I < 4; I += 1) {
    result1.32[I] = _ZX_32(_CTZ_32(argument2.32[I]))
  };
stage E1:
  PGR[registerM] = result1;
\end{lstlisting}}
\newpage
\subsubsection[{\small DIVMODD: Integer Division Modulus Double Word}]{DIVMODD\hfill Integer Division Modulus Double Word}
\label{instr:DIVMODD}\index{divmodd}
 
\paragraph{Encoding} Format \textsf{ALU\_DDMWRR} (Section~\ref{sec:format-ALU-DDMWRR}), \verb|exucode2 = 1100|\\

\bitfields{ 1/P, 2/11, 1/0, 4/1100, 5/registerM, 1/--, 2/10, 3/101, 1/0, 6/registerY, 6/registerZ }


\paragraph{Syntax} \texttt{divmodd} \emph{registerM} = \emph{registerZ}, \emph{registerY}

\paragraph{Description} The dividend and the divisor assuming 64-bit signed integers in the \emph{registerZ} and the \emph{registerY}. The quotient and the remainder are computed and the corresponding 64 bits integers and packed into the 128 bits of the \emph{registerM}. Division by zero returns zero. Modulus by zero returns zero. Division of LONG\_MIN by -1 returns LONG\_MIN. Modulus of LONG\_MIN by -1 returns zero.


\paragraph{Execution} {Reservation table \textsf{ALU\_FULL} (Section~\ref{sec:reservation-ALU-FULL})}
~
{\begin{lstlisting}
stage RR:
  new argument3 = GPR[registerY];
  new argument2 = GPR[registerZ];
  if (argument3.64[0] == 0) {
    new result1 = 0;
  } else if ((argument2.64[0] == 0x8000000000000000) && (argument3.64[0] == 0xFFFFFFFFFFFFFFFF)) {
    result1.64[0] = argument2.64[0];
    result1.64[1] = 0;
  } else {
    new dividend = _SX_64(argument2.64[0]);
    new divisor = _SX_64(argument3.64[0]);
    result1.64[0] = dividend / divisor;
    result1.64[1] = dividend % divisor;
  }
stage SF:
  PGR[registerM] = result1;
\end{lstlisting}}
\newpage
\subsubsection[{\small DIVMODUD: Integer Division Modulus Unsigned Double Word}]{DIVMODUD\hfill Integer Division Modulus Unsigned Double Word}
\label{instr:DIVMODUD}\index{divmodud}
 
\paragraph{Encoding} Format \textsf{ALU\_DDMWRR} (Section~\ref{sec:format-ALU-DDMWRR}), \verb|exucode2 = 1101|\\

\bitfields{ 1/P, 2/11, 1/0, 4/1101, 5/registerM, 1/--, 2/10, 3/101, 1/0, 6/registerY, 6/registerZ }


\paragraph{Syntax} \texttt{divmodud} \emph{registerM} = \emph{registerZ}, \emph{registerY}

\paragraph{Description} The dividend and the divisor assuming 64-bit unsigned integers in the \emph{registerZ} and the \emph{registerY}. The quotient and the remainder are computed and the corresponding 64-bit integers and packed into the 128 bits of the \emph{registerM}. Quotient of division by zero returns zero. Modulus by zero returns zero.


\paragraph{Execution} {Reservation table \textsf{ALU\_FULL} (Section~\ref{sec:reservation-ALU-FULL})}
~
{\begin{lstlisting}
stage RR:
  new argument3 = GPR[registerY];
  new argument2 = GPR[registerZ];
  if (argument3.64[0] == 0) {
    new result1 = 0;
  } else {
    new dividend = _ZX_64(argument2.64[0]);
    new divisor = _ZX_64(argument3.64[0]);
    result1.64[0] = dividend / divisor;
    result1.64[1] = dividend % divisor;
  }
stage SF:
  PGR[registerM] = result1;
\end{lstlisting}}
\newpage
\subsubsection[{\small DIVMODUW: Integer Division Modulus Unsigned Word}]{DIVMODUW\hfill Integer Division Modulus Unsigned Word}
\label{instr:DIVMODUW}\index{divmoduw}
 
\paragraph{Encoding} Format \textsf{ALU\_WDMWRR} (Section~\ref{sec:format-ALU-WDMWRR}), \verb|exucode2 = 1101|\\

\bitfields{ 1/P, 2/11, 1/0, 4/1101, 5/registerM, 1/--, 2/11, 3/101, 1/signextw, 6/registerY, 6/registerZ }


\paragraph{Syntax} \texttt{divmoduw}\emph{signextw} \emph{registerM} = \emph{registerZ}, \emph{registerY}

\paragraph{Description} The dividend and the divisor assuming 32-bit unsigned integers in the \emph{registerZ} and the \emph{registerY}. The quotient and the remainder are computed and the corresponding 32-bit integers are zero-extended and packed into the 128 bits of the \emph{registerM}. Quotient of division by zero returns zero. Modulus by zero returns zero.


\paragraph{Modifier(s)}
\noindent\begin{itemize}
\item signextw (cf. signextw in section \ref{par:modifier-signextw} on page \pageref{par:modifier-signextw})
\end{itemize}

\paragraph{Execution} {Reservation table \textsf{ALU\_FULL} (Section~\ref{sec:reservation-ALU-FULL})}
~
{\begin{lstlisting}
stage ID:
  new signextw = _ZX_1(signextw);
stage RR:
  new argument3 = GPR[registerY];
  new argument2 = GPR[registerZ];
  if (argument3.32[0] == 0) {
    new result1 = 0;
  } else {
    new dividend = _ZX_32(argument2.32[0]);
    new divisor = _ZX_32(argument3.32[0]);
    result1.64[0] = _ZX_32(dividend / divisor);
    result1.64[1] = _ZX_32(dividend % divisor);
  }
  if (signextw) {
    result1.64[0] = _SX_32(result1.64[0]);
    result1.64[1] = _SX_32(result1.64[1]);
  }
stage SF:
  PGR[registerM] = result1;
\end{lstlisting}}
\newpage
\subsubsection[{\small DIVMODW: Integer Division Modulus Word}]{DIVMODW\hfill Integer Division Modulus Word}
\label{instr:DIVMODW}\index{divmodw}
 
\paragraph{Encoding} Format \textsf{ALU\_WDMWRR} (Section~\ref{sec:format-ALU-WDMWRR}), \verb|exucode2 = 1100|\\

\bitfields{ 1/P, 2/11, 1/0, 4/1100, 5/registerM, 1/--, 2/11, 3/101, 1/signextw, 6/registerY, 6/registerZ }


\paragraph{Syntax} \texttt{divmodw}\emph{signextw} \emph{registerM} = \emph{registerZ}, \emph{registerY}

\paragraph{Description} The dividend and the divisor assuming 32-bit signed integers in the \emph{registerZ} and the \emph{registerY}. The quotient and the remainder are computed and the corresponding 32-bit integers are zero-extended and packed into the 128 bits of the \emph{registerM}. Division by zero returns zero. Modulus by zero returns zero. Division of INT\_MIN by -1 returns INT\_MIN. Modulus of INT\_MIN by -1 returns zero.


\paragraph{Modifier(s)}
\noindent\begin{itemize}
\item signextw (cf. signextw in section \ref{par:modifier-signextw} on page \pageref{par:modifier-signextw})
\end{itemize}

\paragraph{Execution} {Reservation table \textsf{ALU\_FULL} (Section~\ref{sec:reservation-ALU-FULL})}
~
{\begin{lstlisting}
stage ID:
  new signextw = _ZX_1(signextw);
stage RR:
  new argument3 = GPR[registerY];
  new argument2 = GPR[registerZ];
  if (argument3.32[0] == 0) {
    new result1 = 0;
  } else if ((argument2.32[0] == 0x80000000) && (argument3.32[0] == 0xFFFFFFFF)) {
    result1.64[0] = _ZX_32(argument2.32[0]);
    result1.64[1] = 0;
  } else {
    new dividend = _SX_32(argument2.32[0]);
    new divisor = _SX_32(argument3.32[0]);
    result1.64[0] = _ZX_32(dividend / divisor);
    result1.64[1] = _ZX_32(dividend % divisor);
  }
  if (signextw) {
    result1.64[0] = _SX_32(result1.64[0]);
    result1.64[1] = _SX_32(result1.64[1]);
  }
stage SF:
  PGR[registerM] = result1;
\end{lstlisting}}
\newpage
\subsubsection[{\small EORD: Bitwise Exclusive Or Double Word}]{EORD\hfill Bitwise Exclusive Or Double Word}
\label{instr:EORD}\index{eord}
 
\paragraph{Encoding} Format \textsf{ALU\_DWRR0} (Section~\ref{sec:format-ALU-DWRR0}), \verb|exucode2 = 1100|\\

\bitfields{ 1/P, 2/11, 1/0, 4/1100, 6/registerW, 2/10, 3/000, 1/0, 6/registerY, 6/registerZ }


\paragraph{Syntax} \texttt{eord} \emph{registerW} = \emph{registerZ}, \emph{registerY}

\paragraph{Description} Bitwise exclusive or of the \emph{registerZ} with the \emph{registerY}. The result is stored into the \emph{registerW}.


\paragraph{Execution} {Reservation table \textsf{ALU\_TINY} (Section~\ref{sec:reservation-ALU-TINY})}
~
{\begin{lstlisting}
stage RR:
  new argument3 = GPR[registerY];
  new argument2 = GPR[registerZ];
  new result1 = argument2 ^ argument3;
stage E1:
  GPR[registerW] = result1;
\end{lstlisting}}
\newpage

\paragraph{Encoding} Format \textsf{ALU\_DWRR0.M} (Section~\ref{sec:format-ALU-DWRR0.M}), \verb|exucode2 = 1100|\\

\bitfields{ 1/P, 2/00, 2/-- --, 27/upper27, 1/1, 2/11, 1/0, 4/1100, 6/registerW, 2/10, 3/000, 1/0, 1/splat32, 5/lower5, 6/registerZ }


\paragraph{Syntax} \texttt{eord} \emph{registerW} = \emph{registerZ}, \emph{upper27\_\discretionary{}{}{}lower5}\emph{splat32}

\paragraph{Description} Bitwise exclusive or of the \emph{registerZ} with the \emph{upper27\_\discretionary{}{}{}lower5}. The result is stored into the \emph{registerW}.


\paragraph{Modifier(s)}
\noindent\begin{itemize}
\item splat32 (cf. splat32 in section \ref{par:modifier-splat32} on page \pageref{par:modifier-splat32})
\end{itemize}

\paragraph{Execution} {Reservation table \textsf{ALU\_TINY.X} (Section~\ref{sec:reservation-ALU-TINY.X})}
~
{\begin{lstlisting}
stage ID:
  new splat32 = _ZX_1(splat32);
stage RR:
  new argument3 = splat32 ? (_WX_32(upper27_lower5) << 32) | _ZX_32(_WX_32(upper27_lower5)) : _SX_32(_WX_32(upper27_lower5));
  new argument2 = GPR[registerZ];
  new result1 = argument2 ^ argument3;
stage E1:
  GPR[registerW] = result1;
\end{lstlisting}}
\newpage

\paragraph{Encoding} Format \textsf{ALU\_DWRI} (Section~\ref{sec:format-ALU-DWRI}), \verb|exucode2 = 1100|\\

\bitfields{ 1/P, 2/11, 1/0, 4/1100, 6/registerW, 2/00, 10/signed10, 6/registerZ }


\paragraph{Syntax} \texttt{eord} \emph{registerW} = \emph{registerZ}, \emph{signed10}

\paragraph{Description} Bitwise exclusive or of the \emph{registerZ} with the \emph{signed10}. The result is stored into the \emph{registerW}.


\paragraph{Execution} {Reservation table \textsf{ALU\_TINY} (Section~\ref{sec:reservation-ALU-TINY})}
~
{\begin{lstlisting}
stage RR:
  new argument3 = _SX_10(signed10);
  new argument2 = GPR[registerZ];
  new result1 = argument2 ^ argument3;
stage E1:
  GPR[registerW] = result1;
\end{lstlisting}}
\newpage

\paragraph{Encoding} Format \textsf{ALU\_DWRI.X} (Section~\ref{sec:format-ALU-DWRI.X}), \verb|exucode2 = 1100|\\

\bitfields{ 1/P, 2/00, 2/-- --, 27/upper27, 1/1, 2/11, 1/0, 4/1100, 6/registerW, 2/00, 10/lower10, 6/registerZ }


\paragraph{Syntax} \texttt{eord} \emph{registerW} = \emph{registerZ}, \emph{upper27\_\discretionary{}{}{}lower10}

\paragraph{Description} Bitwise exclusive or of the \emph{registerZ} with the \emph{upper27\_\discretionary{}{}{}lower10}. The result is stored into the \emph{registerW}.


\paragraph{Execution} {Reservation table \textsf{ALU\_TINY.X} (Section~\ref{sec:reservation-ALU-TINY.X})}
~
{\begin{lstlisting}
stage RR:
  new argument3 = _SX_37(upper27_lower10);
  new argument2 = GPR[registerZ];
  new result1 = argument2 ^ argument3;
stage E1:
  GPR[registerW] = result1;
\end{lstlisting}}
\newpage

\paragraph{Encoding} Format \textsf{ALU\_DWRI.Y} (Section~\ref{sec:format-ALU-DWRI.Y}), \verb|exucode2 = 1100|\\

\bitfields{ 1/P, 2/00, 2/-- --, 27/extend27, 1/1, 2/00, 2/-- --, 27/upper27, 1/1, 2/11, 1/0, 4/1100, 6/registerW, 2/00, 10/lower10, 6/registerZ }


\paragraph{Syntax} \texttt{eord} \emph{registerW} = \emph{registerZ}, \emph{extend27\_\discretionary{}{}{}upper27\_\discretionary{}{}{}lower10}

\paragraph{Description} Bitwise exclusive or of the \emph{registerZ} with the \emph{extend27\_\discretionary{}{}{}upper27\_\discretionary{}{}{}lower10}. The result is stored into the \emph{registerW}.


\paragraph{Execution} {Reservation table \textsf{ALU\_TINY.Y} (Section~\ref{sec:reservation-ALU-TINY.Y})}
~
{\begin{lstlisting}
stage RR:
  new argument3 = _WX_64(extend27_upper27_lower10);
  new argument2 = GPR[registerZ];
  new result1 = argument2 ^ argument3;
stage E1:
  GPR[registerW] = result1;
\end{lstlisting}}
\newpage
\subsubsection[{\small EORQ: Bitwise Exclusive Or Quadruple Word}]{EORQ\hfill Bitwise Exclusive Or Quadruple Word}
\label{instr:EORQ}\index{eorq}
 
\paragraph{Encoding} Format \textsf{ALU\_QWRR0} (Section~\ref{sec:format-ALU-QWRR0}), \verb|exucode2 = 1100|\\

\bitfields{ 1/P, 2/11, 1/0, 4/1100, 5/registerM, 1/--, 2/10, 3/000, 1/1, 5/registerO, 1/--, 5/registerP, 1/-- }


\paragraph{Syntax} \texttt{eorq} \emph{registerM} = \emph{registerP}, \emph{registerO}

\paragraph{Description} Bitwise exclusive or of the \emph{registerP} with the \emph{registerO}. The result is stored into the \emph{registerM}.


\paragraph{Execution} {Reservation table \textsf{ALU\_LITE} (Section~\ref{sec:reservation-ALU-LITE})}
~
{\begin{lstlisting}
stage RR:
  new argument3 = PGR[registerO];
  new argument2 = PGR[registerP];
  new result1 = argument2 ^ argument3;
stage E1:
  PGR[registerM] = result1;
\end{lstlisting}}
\newpage

\paragraph{Encoding} Format \textsf{ALU\_QWRR0.M} (Section~\ref{sec:format-ALU-QWRR0.M}), \verb|exucode2 = 1100|\\

\bitfields{ 1/P, 2/00, 2/-- --, 27/upper27, 1/1, 2/11, 1/0, 4/1100, 5/registerM, 1/--, 2/10, 3/000, 1/1, 1/splat32, 5/lower5, 5/registerP, 1/1 }


\paragraph{Syntax} \texttt{eorq} \emph{registerM} = \emph{registerP}, \emph{upper27\_\discretionary{}{}{}lower5}\emph{splat32}

\paragraph{Description} Bitwise exclusive or of the \emph{registerP} with the \emph{upper27\_\discretionary{}{}{}lower5}. The result is stored into the \emph{registerM}.


\paragraph{Modifier(s)}
\noindent\begin{itemize}
\item splat32 (cf. splat32 in section \ref{par:modifier-splat32} on page \pageref{par:modifier-splat32})
\end{itemize}

\paragraph{Execution} {Reservation table \textsf{ALU\_LITE.X} (Section~\ref{sec:reservation-ALU-LITE.X})}
~
{\begin{lstlisting}
stage ID:
  new splat32 = _ZX_1(splat32);
stage RR:
  new argument3 = splat32 ? (_WX_32(upper27_lower5) << 32) | _ZX_32(_WX_32(upper27_lower5)) : _SX_32(_WX_32(upper27_lower5));
  new argument2 = PGR[registerP];
  new result1 = argument2 ^ argument3;
stage E1:
  PGR[registerM] = result1;
\end{lstlisting}}
\newpage
\subsubsection[{\small EORW: Bitwise Exclusive Or Word}]{EORW\hfill Bitwise Exclusive Or Word}
\label{instr:EORW}\index{eorw}
 
\paragraph{Encoding} Format \textsf{ALU\_WWRR0} (Section~\ref{sec:format-ALU-WWRR0}), \verb|exucode2 = 1100|\\

\bitfields{ 1/P, 2/11, 1/0, 4/1100, 6/registerW, 2/11, 3/000, 1/signextw, 6/registerY, 6/registerZ }


\paragraph{Syntax} \texttt{eorw}\emph{signextw} \emph{registerW} = \emph{registerZ}, \emph{registerY}

\paragraph{Description} Bitwise exclusive or of the \emph{registerZ} with the \emph{registerY}. The result with the 32 upper bits cleared is stored into the \emph{registerW}.


\paragraph{Modifier(s)}
\noindent\begin{itemize}
\item signextw (cf. signextw in section \ref{par:modifier-signextw} on page \pageref{par:modifier-signextw})
\end{itemize}

\paragraph{Execution} {Reservation table \textsf{ALU\_TINY} (Section~\ref{sec:reservation-ALU-TINY})}
~
{\begin{lstlisting}
stage ID:
  new signextw = _ZX_1(signextw);
stage RR:
  new argument3 = GPR[registerY];
  new argument2 = GPR[registerZ];
  new result1 = argument2 ^ argument3;
stage E1:
  GPR[registerW] = signextw ? _SX_32(result1) : _ZX_32(result1);
\end{lstlisting}}
\newpage

\paragraph{Encoding} Format \textsf{ALU\_WWRR0.W} (Section~\ref{sec:format-ALU-WWRR0.W}), \verb|exucode2 = 1100|\\

\bitfields{ 1/P, 2/00, 2/-- --, 27/upper27, 1/1, 2/11, 1/0, 4/1100, 6/registerW, 2/11, 3/000, 1/signextw, 1/0, 5/lower5, 6/registerZ }


\paragraph{Syntax} \texttt{eorw}\emph{signextw} \emph{registerW} = \emph{registerZ}, \emph{upper27\_\discretionary{}{}{}lower5}

\paragraph{Description} Bitwise exclusive or of the \emph{registerZ} with the \emph{upper27\_\discretionary{}{}{}lower5}. The result with the 32 upper bits cleared is stored into the \emph{registerW}.


\paragraph{Modifier(s)}
\noindent\begin{itemize}
\item signextw (cf. signextw in section \ref{par:modifier-signextw} on page \pageref{par:modifier-signextw})
\end{itemize}

\paragraph{Execution} {Reservation table \textsf{ALU\_TINY.X} (Section~\ref{sec:reservation-ALU-TINY.X})}
~
{\begin{lstlisting}
stage ID:
  new signextw = _ZX_1(signextw);
stage RR:
  new argument3 = _WX_32(upper27_lower5);
  new argument2 = GPR[registerZ];
  new result1 = argument2 ^ argument3;
stage E1:
  GPR[registerW] = signextw ? _SX_32(result1) : _ZX_32(result1);
\end{lstlisting}}
\newpage
\subsubsection[{\small EXTFS: Extract Bit Field and Sign Extend to Double Word}]{EXTFS\hfill Extract Bit Field and Sign Extend to Double Word}
\label{instr:EXTFS}\index{extfs}
 
\paragraph{Encoding} Format \textsf{ALU\_BFWR} (Section~\ref{sec:format-ALU-BFWR}), \verb|wricode1 = 10|\\

\bitfields{ 1/P, 2/11, 1/0, 2/10, 2/stopbit2, 6/registerW, 2/01, 4/stopbit4, 6/startbit, 6/registerZ }


\paragraph{Syntax} \texttt{extfs} \emph{registerW} = \emph{registerZ}, \emph{stopbit2\_\discretionary{}{}{}stopbit4}, \emph{startbit}

\paragraph{Description} The start bit position and the stop bit position are extracted from the \emph{startbit} and the \emph{stopbit2\_\discretionary{}{}{}stopbit4} respectively. The bit-field of the \emph{registerZ} extending from the start bit position to the stop bit position is extracted, sign-extended from the stop bit position, shifted right by the start bit position and stored into the \emph{registerW}. Behavior is undefined if the start bit position is greater than the stop bit position.


\paragraph{Execution} {Reservation table \textsf{ALU\_TINY} (Section~\ref{sec:reservation-ALU-TINY})}
~
{\begin{lstlisting}
stage ID:
  new startbit = _ZX_6(startbit);
  new stopbit = _ZX_6(stopbit2_stopbit4);
stage RR:
  new bias = startbit <= stopbit;
  new mask = _ZX_64(((2 << stopbit) - ((2 - bias) << startbit)) + (bias - 1));
  new argument2 = GPR[registerZ];
  new masked1 = _ZX_64(argument2 & mask);
  new masked2 = _SX_64(argument2 | ~mask);
  new negative = argument2 & (mask & ~(mask >> 1));
  new masked = negative == 0 ? masked1 : masked2;
  new result1 = masked >> startbit;
stage E1:
  GPR[registerW] = result1;
\end{lstlisting}}
\newpage
\subsubsection[{\small EXTFZ: Extract Bit Field and Zero Extend to Double Word}]{EXTFZ\hfill Extract Bit Field and Zero Extend to Double Word}
\label{instr:EXTFZ}\index{extfz}
 
\paragraph{Encoding} Format \textsf{ALU\_BFWR} (Section~\ref{sec:format-ALU-BFWR}), \verb|wricode1 = 01|\\

\bitfields{ 1/P, 2/11, 1/0, 2/01, 2/stopbit2, 6/registerW, 2/01, 4/stopbit4, 6/startbit, 6/registerZ }


\paragraph{Syntax} \texttt{extfz} \emph{registerW} = \emph{registerZ}, \emph{stopbit2\_\discretionary{}{}{}stopbit4}, \emph{startbit}

\paragraph{Description} The start bit position and the stop bit position are extracted from the \emph{startbit} and the \emph{stopbit2\_\discretionary{}{}{}stopbit4} respectively. The bit-field of the \emph{registerZ} extending from the start bit position to the stop bit position is extracted, zero-extended after the stop bit position, shifted right by the start bit position and stored into the \emph{registerW}. Behavior is undefined if the start bit position is greater than the stop bit position.


\paragraph{Execution} {Reservation table \textsf{ALU\_TINY} (Section~\ref{sec:reservation-ALU-TINY})}
~
{\begin{lstlisting}
stage ID:
  new startbit = _ZX_6(startbit);
  new stopbit = _ZX_6(stopbit2_stopbit4);
stage RR:
  new bias = startbit <= stopbit;
  new mask = _ZX_64(((2 << stopbit) - ((2 - bias) << startbit)) + (bias - 1));
  new argument2 = GPR[registerZ];
  new masked = argument2 & mask;
  new result1 = masked >> startbit;
stage E1:
  GPR[registerW] = result1;
\end{lstlisting}}
\newpage
\subsubsection[{\small EXTLQBHO: Extend Lane Q7 Byte to Half Word Octuple}]{EXTLQBHO\hfill Extend Lane Q7 Byte to Half Word Octuple}
\label{instr:EXTLQBHO}\index{extlqbho}
 
\paragraph{Encoding} Format \textsf{ALU\_HOEXTLWR} (Section~\ref{sec:format-ALU-HOEXTLWR}), \verb|alucode8 = 1110|\\

\bitfields{ 1/P, 2/11, 1/1, 4/1111, 5/registerM, 1/oddlanes, 2/10, 3/101, 1/1, 4/1110, 1/--, 1/--, 5/registerP, 1/0 }


\paragraph{Syntax} \texttt{extlqbho}\emph{oddlanes} \emph{registerM} = \emph{registerP}

\paragraph{Description} The \emph{registerP} is interpreted as sixteen bytes packed into 128 bits. The even or odd bytes of the \emph{registerP} are selected, depending on \emph{oddlanes}. These eight bytes interpreted as Q7 numbers are mapped to Q15 numbers, packed and stored into the \emph{registerM}.


\paragraph{Modifier(s)}
\noindent\begin{itemize}
\item oddlanes (cf. oddlanes in section \ref{par:modifier-oddlanes} on page \pageref{par:modifier-oddlanes})
\end{itemize}

\paragraph{Execution} {Reservation table \textsf{ALU\_LITE} (Section~\ref{sec:reservation-ALU-LITE})}
~
{\begin{lstlisting}
stage ID:
  new oddlanes = _ZX_1(oddlanes);
stage RR:
  new argument2 = PGR[registerP];
  if (oddlanes) {
    argument2 = argument2 >> 8;
  }
  for (I = 0; I < 8; I += 1) {
    result1.16[I] = argument2.16[I] << 8
  };
stage E1:
  PGR[registerM] = result1;
\end{lstlisting}}
\newpage
\subsubsection[{\small EXTLQHWQ: Extend Lane Q15 Half Word to Word Quadruple}]{EXTLQHWQ\hfill Extend Lane Q15 Half Word to Word Quadruple}
\label{instr:EXTLQHWQ}\index{extlqhwq}
 
\paragraph{Encoding} Format \textsf{ALU\_WQEXTLWR} (Section~\ref{sec:format-ALU-WQEXTLWR}), \verb|alucode8 = 1110|\\

\bitfields{ 1/P, 2/11, 1/1, 4/1111, 5/registerM, 1/oddlanes, 2/10, 3/101, 1/0, 4/1110, 1/--, 1/--, 5/registerP, 1/1 }


\paragraph{Syntax} \texttt{extlqhwq}\emph{oddlanes} \emph{registerM} = \emph{registerP}

\paragraph{Description} The \emph{registerP} is interpreted as eight half words packed into 128 bits. The even or odd half words of the \emph{registerP} are selected, depending on \emph{oddlanes}. These four half words interpreted as Q15 numbers are mapped to Q31 numbers, packed and stored into the \emph{registerM}.


\paragraph{Modifier(s)}
\noindent\begin{itemize}
\item oddlanes (cf. oddlanes in section \ref{par:modifier-oddlanes} on page \pageref{par:modifier-oddlanes})
\end{itemize}

\paragraph{Execution} {Reservation table \textsf{ALU\_LITE} (Section~\ref{sec:reservation-ALU-LITE})}
~
{\begin{lstlisting}
stage ID:
  new oddlanes = _ZX_1(oddlanes);
stage RR:
  new argument2 = PGR[registerP];
  if (oddlanes) {
    argument2 = argument2 >> 16;
  }
  for (I = 0; I < 4; I += 1) {
    result1.32[I] = argument2.32[I] << 16
  };
stage E1:
  PGR[registerM] = result1;
\end{lstlisting}}
\newpage
\subsubsection[{\small EXTLQNBX: Extend Lane Q3 Nibble to Byte Hexadecuple}]{EXTLQNBX\hfill Extend Lane Q3 Nibble to Byte Hexadecuple}
\label{instr:EXTLQNBX}\index{extlqnbx}
 
\paragraph{Encoding} Format \textsf{ALU\_BXEXTLWR} (Section~\ref{sec:format-ALU-BXEXTLWR}), \verb|alucode8 = 1110|\\

\bitfields{ 1/P, 2/11, 1/1, 4/1111, 5/registerM, 1/oddlanes, 2/10, 3/101, 1/1, 4/1110, 1/--, 1/--, 5/registerP, 1/1 }


\paragraph{Syntax} \texttt{extlqnbx}\emph{oddlanes} \emph{registerM} = \emph{registerP}

\paragraph{Description} The \emph{registerP} is interpreted as thirty-two nibbles packed into 128 bits. The even or odd nibbles of the \emph{registerP} are selected, depending on \emph{oddlanes}. These sixteen nibbles interpreted as Q4 numbers are mapped to Q7 numbers, packed and stored into the \emph{registerM}.


\paragraph{Modifier(s)}
\noindent\begin{itemize}
\item oddlanes (cf. oddlanes in section \ref{par:modifier-oddlanes} on page \pageref{par:modifier-oddlanes})
\end{itemize}

\paragraph{Execution} {Reservation table \textsf{ALU\_LITE} (Section~\ref{sec:reservation-ALU-LITE})}
~
{\begin{lstlisting}
stage ID:
  new oddlanes = _ZX_1(oddlanes);
stage RR:
  new argument2 = PGR[registerP];
  if (oddlanes) {
    argument2 = argument2 >> 4;
  }
  for (I = 0; I < 16; I += 1) {
    result1.8[I] = argument2.8[I] << 4
  };
stage E1:
  PGR[registerM] = result1;
\end{lstlisting}}
\newpage
\subsubsection[{\small EXTLQWDP: Extend Lane Q31 Word to Double Word Pair}]{EXTLQWDP\hfill Extend Lane Q31 Word to Double Word Pair}
\label{instr:EXTLQWDP}\index{extlqwdp}
 
\paragraph{Encoding} Format \textsf{ALU\_DPEXTLWR} (Section~\ref{sec:format-ALU-DPEXTLWR}), \verb|alucode8 = 1110|\\

\bitfields{ 1/P, 2/11, 1/1, 4/1111, 5/registerM, 1/oddlanes, 2/10, 3/101, 1/0, 4/1110, 1/--, 1/--, 5/registerP, 1/0 }


\paragraph{Syntax} \texttt{extlqwdp}\emph{oddlanes} \emph{registerM} = \emph{registerP}

\paragraph{Description} The \emph{registerP} is interpreted as four words packed into 128 bits. The even or odd words of the \emph{registerP} are selected, depending on \emph{oddlanes}. These two words interpreted as Q31 numbers are mapped to Q63 numbers, packed and stored into the \emph{registerM}.


\paragraph{Modifier(s)}
\noindent\begin{itemize}
\item oddlanes (cf. oddlanes in section \ref{par:modifier-oddlanes} on page \pageref{par:modifier-oddlanes})
\end{itemize}

\paragraph{Execution} {Reservation table \textsf{ALU\_LITE} (Section~\ref{sec:reservation-ALU-LITE})}
~
{\begin{lstlisting}
stage ID:
  new oddlanes = _ZX_1(oddlanes);
stage RR:
  new argument2 = PGR[registerP];
  if (oddlanes) {
    argument2 = argument2 >> 32;
  }
  for (I = 0; I < 2; I += 1) {
    result1.64[I] = argument2.64[I] << 32
  };
stage E1:
  PGR[registerM] = result1;
\end{lstlisting}}
\newpage
\subsubsection[{\small EXTLSBHO: Extend Lane with Sign Byte to Half Word Octuple}]{EXTLSBHO\hfill Extend Lane with Sign Byte to Half Word Octuple}
\label{instr:EXTLSBHO}\index{extlsbho}
 
\paragraph{Encoding} Format \textsf{ALU\_HOEXTLWR} (Section~\ref{sec:format-ALU-HOEXTLWR}), \verb|alucode8 = 1101|\\

\bitfields{ 1/P, 2/11, 1/1, 4/1111, 5/registerM, 1/oddlanes, 2/10, 3/101, 1/1, 4/1101, 1/--, 1/--, 5/registerP, 1/0 }


\paragraph{Syntax} \texttt{extlsbho}\emph{oddlanes} \emph{registerM} = \emph{registerP}

\paragraph{Description} The \emph{registerP} is interpreted as sixteen bytes packed into 128 bits. The even or odd bytes of the \emph{registerP} are selected, depending on \emph{oddlanes}. These eight bytes are sign-extended to 16 bits, packed and stored into the \emph{registerM}.


\paragraph{Modifier(s)}
\noindent\begin{itemize}
\item oddlanes (cf. oddlanes in section \ref{par:modifier-oddlanes} on page \pageref{par:modifier-oddlanes})
\end{itemize}

\paragraph{Execution} {Reservation table \textsf{ALU\_LITE} (Section~\ref{sec:reservation-ALU-LITE})}
~
{\begin{lstlisting}
stage ID:
  new oddlanes = _ZX_1(oddlanes);
stage RR:
  new argument2 = PGR[registerP];
  if (oddlanes) {
    argument2 = argument2 >> 8;
  }
  for (I = 0; I < 8; I += 1) {
    result1.16[I] = _SX_8(argument2.16[I])
  };
stage E1:
  PGR[registerM] = result1;
\end{lstlisting}}
\newpage
\subsubsection[{\small EXTLSHWQ: Extend Lane with Sign Half Word to Word Quadruple}]{EXTLSHWQ\hfill Extend Lane with Sign Half Word to Word Quadruple}
\label{instr:EXTLSHWQ}\index{extlshwq}
 
\paragraph{Encoding} Format \textsf{ALU\_WQEXTLWR} (Section~\ref{sec:format-ALU-WQEXTLWR}), \verb|alucode8 = 1101|\\

\bitfields{ 1/P, 2/11, 1/1, 4/1111, 5/registerM, 1/oddlanes, 2/10, 3/101, 1/0, 4/1101, 1/--, 1/--, 5/registerP, 1/1 }


\paragraph{Syntax} \texttt{extlshwq}\emph{oddlanes} \emph{registerM} = \emph{registerP}

\paragraph{Description} The \emph{registerP} is interpreted as eight half words packed into 128 bits. The even or odd half words of the \emph{registerP} are selected, depending on \emph{oddlanes}. These four half words are sign-extended to 32 bits, packed and stored into the \emph{registerM}.


\paragraph{Modifier(s)}
\noindent\begin{itemize}
\item oddlanes (cf. oddlanes in section \ref{par:modifier-oddlanes} on page \pageref{par:modifier-oddlanes})
\end{itemize}

\paragraph{Execution} {Reservation table \textsf{ALU\_LITE} (Section~\ref{sec:reservation-ALU-LITE})}
~
{\begin{lstlisting}
stage ID:
  new oddlanes = _ZX_1(oddlanes);
stage RR:
  new argument2 = PGR[registerP];
  if (oddlanes) {
    argument2 = argument2 >> 16;
  }
  for (I = 0; I < 4; I += 1) {
    result1.32[I] = _SX_16(argument2.32[I])
  };
stage E1:
  PGR[registerM] = result1;
\end{lstlisting}}
\newpage
\subsubsection[{\small EXTLSNBX: Extend Lane with Sign Nibble to Byte Hexadecuple}]{EXTLSNBX\hfill Extend Lane with Sign Nibble to Byte Hexadecuple}
\label{instr:EXTLSNBX}\index{extlsnbx}
 
\paragraph{Encoding} Format \textsf{ALU\_BXEXTLWR} (Section~\ref{sec:format-ALU-BXEXTLWR}), \verb|alucode8 = 1101|\\

\bitfields{ 1/P, 2/11, 1/1, 4/1111, 5/registerM, 1/oddlanes, 2/10, 3/101, 1/1, 4/1101, 1/--, 1/--, 5/registerP, 1/1 }


\paragraph{Syntax} \texttt{extlsnbx}\emph{oddlanes} \emph{registerM} = \emph{registerP}

\paragraph{Description} The \emph{registerP} is interpreted as thirty-two nibbles packed into 128 bits. The even or odd nibbles of the \emph{registerP} are selected, depending on \emph{oddlanes}. These sixteen nibbles are sign-extended to 8 bits, packed and stored into the \emph{registerM}.


\paragraph{Modifier(s)}
\noindent\begin{itemize}
\item oddlanes (cf. oddlanes in section \ref{par:modifier-oddlanes} on page \pageref{par:modifier-oddlanes})
\end{itemize}

\paragraph{Execution} {Reservation table \textsf{ALU\_LITE} (Section~\ref{sec:reservation-ALU-LITE})}
~
{\begin{lstlisting}
stage ID:
  new oddlanes = _ZX_1(oddlanes);
stage RR:
  new argument2 = PGR[registerP];
  if (oddlanes) {
    argument2 = argument2 >> 4;
  }
  for (I = 0; I < 16; I += 1) {
    result1.8[I] = _SX_4(argument2.8[I])
  };
stage E1:
  PGR[registerM] = result1;
\end{lstlisting}}
\newpage
\subsubsection[{\small EXTLSWDP: Extend Lane with Sign Word to Double Word Pair}]{EXTLSWDP\hfill Extend Lane with Sign Word to Double Word Pair}
\label{instr:EXTLSWDP}\index{extlswdp}
 
\paragraph{Encoding} Format \textsf{ALU\_DPEXTLWR} (Section~\ref{sec:format-ALU-DPEXTLWR}), \verb|alucode8 = 1101|\\

\bitfields{ 1/P, 2/11, 1/1, 4/1111, 5/registerM, 1/oddlanes, 2/10, 3/101, 1/0, 4/1101, 1/--, 1/--, 5/registerP, 1/0 }


\paragraph{Syntax} \texttt{extlswdp}\emph{oddlanes} \emph{registerM} = \emph{registerP}

\paragraph{Description} The \emph{registerP} is interpreted as four words packed into 128 bits. The even or odd words of the \emph{registerP} are selected, depending on \emph{oddlanes}. These two words are sign-extended to 64 bits, packed and stored into the \emph{registerM}.


\paragraph{Modifier(s)}
\noindent\begin{itemize}
\item oddlanes (cf. oddlanes in section \ref{par:modifier-oddlanes} on page \pageref{par:modifier-oddlanes})
\end{itemize}

\paragraph{Execution} {Reservation table \textsf{ALU\_LITE} (Section~\ref{sec:reservation-ALU-LITE})}
~
{\begin{lstlisting}
stage ID:
  new oddlanes = _ZX_1(oddlanes);
stage RR:
  new argument2 = PGR[registerP];
  if (oddlanes) {
    argument2 = argument2 >> 32;
  }
  for (I = 0; I < 2; I += 1) {
    result1.64[I] = _SX_32(argument2.64[I])
  };
stage E1:
  PGR[registerM] = result1;
\end{lstlisting}}
\newpage
\subsubsection[{\small EXTLZBHO: Extend Lane with Zero Byte to Half Word Octuple}]{EXTLZBHO\hfill Extend Lane with Zero Byte to Half Word Octuple}
\label{instr:EXTLZBHO}\index{extlzbho}
 
\paragraph{Encoding} Format \textsf{ALU\_HOEXTLWR} (Section~\ref{sec:format-ALU-HOEXTLWR}), \verb|alucode8 = 1100|\\

\bitfields{ 1/P, 2/11, 1/1, 4/1111, 5/registerM, 1/oddlanes, 2/10, 3/101, 1/1, 4/1100, 1/--, 1/--, 5/registerP, 1/0 }


\paragraph{Syntax} \texttt{extlzbho}\emph{oddlanes} \emph{registerM} = \emph{registerP}

\paragraph{Description} The \emph{registerP} is interpreted as sixteen bytes packed into 128 bits. The even or odd bytes of the \emph{registerP} are selected, depending on \emph{oddlanes}. These eight bytes are zero-extended to 16 bits, packed and stored into the \emph{registerM}.


\paragraph{Modifier(s)}
\noindent\begin{itemize}
\item oddlanes (cf. oddlanes in section \ref{par:modifier-oddlanes} on page \pageref{par:modifier-oddlanes})
\end{itemize}

\paragraph{Execution} {Reservation table \textsf{ALU\_LITE} (Section~\ref{sec:reservation-ALU-LITE})}
~
{\begin{lstlisting}
stage ID:
  new oddlanes = _ZX_1(oddlanes);
stage RR:
  new argument2 = PGR[registerP];
  if (oddlanes) {
    argument2 = argument2 >> 8;
  }
  for (I = 0; I < 8; I += 1) {
    result1.16[I] = _ZX_8(argument2.16[I])
  };
stage E1:
  PGR[registerM] = result1;
\end{lstlisting}}
\newpage
\subsubsection[{\small EXTLZHWQ: Extend Lane with Zero Half Word to Word Quadruple}]{EXTLZHWQ\hfill Extend Lane with Zero Half Word to Word Quadruple}
\label{instr:EXTLZHWQ}\index{extlzhwq}
 
\paragraph{Encoding} Format \textsf{ALU\_WQEXTLWR} (Section~\ref{sec:format-ALU-WQEXTLWR}), \verb|alucode8 = 1100|\\

\bitfields{ 1/P, 2/11, 1/1, 4/1111, 5/registerM, 1/oddlanes, 2/10, 3/101, 1/0, 4/1100, 1/--, 1/--, 5/registerP, 1/1 }


\paragraph{Syntax} \texttt{extlzhwq}\emph{oddlanes} \emph{registerM} = \emph{registerP}

\paragraph{Description} The \emph{registerP} is interpreted as eight half words packed into 128 bits. The even or odd half words of the \emph{registerP} are selected, depending on \emph{oddlanes}. These four half words are zero-extended to 32 bits, packed and stored into the \emph{registerM}.


\paragraph{Modifier(s)}
\noindent\begin{itemize}
\item oddlanes (cf. oddlanes in section \ref{par:modifier-oddlanes} on page \pageref{par:modifier-oddlanes})
\end{itemize}

\paragraph{Execution} {Reservation table \textsf{ALU\_LITE} (Section~\ref{sec:reservation-ALU-LITE})}
~
{\begin{lstlisting}
stage ID:
  new oddlanes = _ZX_1(oddlanes);
stage RR:
  new argument2 = PGR[registerP];
  if (oddlanes) {
    argument2 = argument2 >> 16;
  }
  for (I = 0; I < 4; I += 1) {
    result1.32[I] = _ZX_16(argument2.32[I])
  };
stage E1:
  PGR[registerM] = result1;
\end{lstlisting}}
\newpage
\subsubsection[{\small EXTLZNBX: Extend Lane with Zero Nibble to Byte Hexadecuple}]{EXTLZNBX\hfill Extend Lane with Zero Nibble to Byte Hexadecuple}
\label{instr:EXTLZNBX}\index{extlznbx}
 
\paragraph{Encoding} Format \textsf{ALU\_BXEXTLWR} (Section~\ref{sec:format-ALU-BXEXTLWR}), \verb|alucode8 = 1100|\\

\bitfields{ 1/P, 2/11, 1/1, 4/1111, 5/registerM, 1/oddlanes, 2/10, 3/101, 1/1, 4/1100, 1/--, 1/--, 5/registerP, 1/1 }


\paragraph{Syntax} \texttt{extlznbx}\emph{oddlanes} \emph{registerM} = \emph{registerP}

\paragraph{Description} The \emph{registerP} is interpreted as thirty-two nibbles packed into 128 bits. The even or odd nibbles of the \emph{registerP} are selected, depending on \emph{oddlanes}. These sixteen nibbles are zero-extended to 8 bits, packed and stored into the \emph{registerM}.


\paragraph{Modifier(s)}
\noindent\begin{itemize}
\item oddlanes (cf. oddlanes in section \ref{par:modifier-oddlanes} on page \pageref{par:modifier-oddlanes})
\end{itemize}

\paragraph{Execution} {Reservation table \textsf{ALU\_LITE} (Section~\ref{sec:reservation-ALU-LITE})}
~
{\begin{lstlisting}
stage ID:
  new oddlanes = _ZX_1(oddlanes);
stage RR:
  new argument2 = PGR[registerP];
  if (oddlanes) {
    argument2 = argument2 >> 4;
  }
  for (I = 0; I < 16; I += 1) {
    result1.8[I] = _ZX_4(argument2.8[I])
  };
stage E1:
  PGR[registerM] = result1;
\end{lstlisting}}
\newpage
\subsubsection[{\small EXTLZWDP: Extend Lane with Zero Word to Double Word Pair}]{EXTLZWDP\hfill Extend Lane with Zero Word to Double Word Pair}
\label{instr:EXTLZWDP}\index{extlzwdp}
 
\paragraph{Encoding} Format \textsf{ALU\_DPEXTLWR} (Section~\ref{sec:format-ALU-DPEXTLWR}), \verb|alucode8 = 1100|\\

\bitfields{ 1/P, 2/11, 1/1, 4/1111, 5/registerM, 1/oddlanes, 2/10, 3/101, 1/0, 4/1100, 1/--, 1/--, 5/registerP, 1/0 }


\paragraph{Syntax} \texttt{extlzwdp}\emph{oddlanes} \emph{registerM} = \emph{registerP}

\paragraph{Description} The \emph{registerP} is interpreted as four words packed into 128 bits. The even or odd words of the \emph{registerP} are selected, depending on \emph{oddlanes}. These two words are zero-extended to 64 bits, packed and stored into the \emph{registerM}.


\paragraph{Modifier(s)}
\noindent\begin{itemize}
\item oddlanes (cf. oddlanes in section \ref{par:modifier-oddlanes} on page \pageref{par:modifier-oddlanes})
\end{itemize}

\paragraph{Execution} {Reservation table \textsf{ALU\_LITE} (Section~\ref{sec:reservation-ALU-LITE})}
~
{\begin{lstlisting}
stage ID:
  new oddlanes = _ZX_1(oddlanes);
stage RR:
  new argument2 = PGR[registerP];
  if (oddlanes) {
    argument2 = argument2 >> 32;
  }
  for (I = 0; I < 2; I += 1) {
    result1.64[I] = _ZX_32(argument2.64[I])
  };
stage E1:
  PGR[registerM] = result1;
\end{lstlisting}}
\newpage
\subsubsection[{\small FCOMPHO: Floating-Point Compare Half Word Octuple}]{FCOMPHO\hfill Floating-Point Compare Half Word Octuple}
\label{instr:FCOMPHO}\index{fcompho}
 
\paragraph{Encoding} Format \textsf{ALU\_FHOCWRR} (Section~\ref{sec:format-ALU-FHOCWRR}), \verb|alucode3 = 010|\\

\bitfields{ 1/P, 2/11, 1/1, 1/0, 3/floatcomp, 6/registerW, 2/11, 3/010, 1/1, 5/registerO, 1/--, 5/registerP, 1/1 }


\paragraph{Syntax} \texttt{fcompho}\emph{floatcomp} \emph{registerW} = \emph{registerP}, \emph{registerO}

\paragraph{Description} The \emph{registerP} and the \emph{registerO} are interpreted as eight binary 16 floating-point numbers packed into 128 bits. The \emph{registerP} half words are compared to the \emph{registerO} half words using condition \emph{floatcomp}. The eight boolean result bits are packed and stored into the \emph{registerW}.


\paragraph{Modifier(s)}
\noindent\begin{itemize}
\item floatcomp (cf. floatcomp in section \ref{par:modifier-floatcomp} on page \pageref{par:modifier-floatcomp})
\end{itemize}

\paragraph{Execution} {Reservation table \textsf{ALU\_LITE} (Section~\ref{sec:reservation-ALU-LITE})}
~
{\begin{lstlisting}
stage ID:
  new argument4 = _ZX_3(floatcomp);
stage RR:
  new argument3 = PGR[registerO];
  new argument2 = PGR[registerP];
  for (I = 0; I < 8; I += 1) {
    result1.1[I] = floatcomp_16(argument4, argument2.16[I], argument3.16[I])
  };
stage E1:
  GPR[registerW] = result1;
\end{lstlisting}}
\newpage

\paragraph{Encoding} Format \textsf{ALU\_FHOCWRR.M} (Section~\ref{sec:format-ALU-FHOCWRR.M}), \verb|alucode3 = 010|\\

\bitfields{ 1/P, 2/00, 2/-- --, 27/upper27, 1/1, 2/11, 1/1, 1/0, 3/floatcomp, 6/registerW, 2/11, 3/010, 1/1, 1/splat32, 5/lower5, 5/registerP, 1/1 }


\paragraph{Syntax} \texttt{fcompho}\emph{floatcomp} \emph{registerW} = \emph{registerP}, \emph{upper27\_\discretionary{}{}{}lower5}\emph{splat32}

\paragraph{Description} The \emph{registerP} and the \emph{upper27\_\discretionary{}{}{}lower5} are interpreted as eight binary 16 floating-point numbers packed into 128 bits. The \emph{registerP} half words are compared to the \emph{upper27\_\discretionary{}{}{}lower5} half words using condition \emph{floatcomp}. The eight boolean result bits are packed and stored into the \emph{registerW}.


\paragraph{Modifier(s)}
\noindent\begin{itemize}
\item floatcomp (cf. floatcomp in section \ref{par:modifier-floatcomp} on page \pageref{par:modifier-floatcomp})
\item splat32 (cf. splat32 in section \ref{par:modifier-splat32} on page \pageref{par:modifier-splat32})
\end{itemize}

\paragraph{Execution} {Reservation table \textsf{ALU\_LITE.X} (Section~\ref{sec:reservation-ALU-LITE.X})}
~
{\begin{lstlisting}
stage ID:
  new splat32 = _ZX_1(splat32);
  new argument4 = _ZX_3(floatcomp);
stage RR:
  new argument3 = splat32 ? (_WX_32(upper27_lower5) << 32) | _ZX_32(_WX_32(upper27_lower5)) : _SX_32(_WX_32(upper27_lower5));
  new argument2 = PGR[registerP];
  for (I = 0; I < 8; I += 1) {
    result1.1[I] = floatcomp_16(argument4, argument2.16[I], argument3.16[I])
  };
stage E1:
  GPR[registerW] = result1;
\end{lstlisting}}
\newpage
\subsubsection[{\small FCOMPWQ: Floating-Point Compare Word Quadruple}]{FCOMPWQ\hfill Floating-Point Compare Word Quadruple}
\label{instr:FCOMPWQ}\index{fcompwq}
 
\paragraph{Encoding} Format \textsf{ALU\_FWQCWRR} (Section~\ref{sec:format-ALU-FWQCWRR}), \verb|alucode3 = 010|\\

\bitfields{ 1/P, 2/11, 1/1, 1/0, 3/floatcomp, 6/registerW, 2/11, 3/010, 1/1, 5/registerO, 1/--, 5/registerP, 1/0 }


\paragraph{Syntax} \texttt{fcompwq}\emph{floatcomp} \emph{registerW} = \emph{registerP}, \emph{registerO}

\paragraph{Description} The \emph{registerP} and the \emph{registerO} are interpreted as four binary 32 floating-point numbers packed into 128 bits. The \emph{registerP} words are compared to the \emph{registerO} words using condition \emph{floatcomp}. The four boolean result bits are packed and stored into the \emph{registerW}.


\paragraph{Modifier(s)}
\noindent\begin{itemize}
\item floatcomp (cf. floatcomp in section \ref{par:modifier-floatcomp} on page \pageref{par:modifier-floatcomp})
\end{itemize}

\paragraph{Execution} {Reservation table \textsf{ALU\_LITE} (Section~\ref{sec:reservation-ALU-LITE})}
~
{\begin{lstlisting}
stage ID:
  new argument4 = _ZX_3(floatcomp);
stage RR:
  new argument3 = PGR[registerO];
  new argument2 = PGR[registerP];
  for (I = 0; I < 4; I += 1) {
    result1.1[I] = floatcomp_32(argument4, argument2.32[I], argument3.32[I])
  };
stage E1:
  GPR[registerW] = result1;
\end{lstlisting}}
\newpage

\paragraph{Encoding} Format \textsf{ALU\_FWQCWRR.M} (Section~\ref{sec:format-ALU-FWQCWRR.M}), \verb|alucode3 = 010|\\

\bitfields{ 1/P, 2/00, 2/-- --, 27/upper27, 1/1, 2/11, 1/1, 1/0, 3/floatcomp, 6/registerW, 2/11, 3/010, 1/1, 1/splat32, 5/lower5, 5/registerP, 1/0 }


\paragraph{Syntax} \texttt{fcompwq}\emph{floatcomp} \emph{registerW} = \emph{registerP}, \emph{upper27\_\discretionary{}{}{}lower5}\emph{splat32}

\paragraph{Description} The \emph{registerP} and the \emph{upper27\_\discretionary{}{}{}lower5} are interpreted as four binary 32 floating-point numbers packed into 128 bits. The \emph{registerP} words are compared to the \emph{upper27\_\discretionary{}{}{}lower5} words using condition \emph{floatcomp}. The four boolean result bits are packed and stored into the \emph{registerW}.


\paragraph{Modifier(s)}
\noindent\begin{itemize}
\item floatcomp (cf. floatcomp in section \ref{par:modifier-floatcomp} on page \pageref{par:modifier-floatcomp})
\item splat32 (cf. splat32 in section \ref{par:modifier-splat32} on page \pageref{par:modifier-splat32})
\end{itemize}

\paragraph{Execution} {Reservation table \textsf{ALU\_LITE.X} (Section~\ref{sec:reservation-ALU-LITE.X})}
~
{\begin{lstlisting}
stage ID:
  new splat32 = _ZX_1(splat32);
  new argument4 = _ZX_3(floatcomp);
stage RR:
  new argument3 = splat32 ? (_WX_32(upper27_lower5) << 32) | _ZX_32(_WX_32(upper27_lower5)) : _SX_32(_WX_32(upper27_lower5));
  new argument2 = PGR[registerP];
  for (I = 0; I < 4; I += 1) {
    result1.1[I] = floatcomp_32(argument4, argument2.32[I], argument3.32[I])
  };
stage E1:
  GPR[registerW] = result1;
\end{lstlisting}}
\newpage
\subsubsection[{\small FNARROWDWQ: Floating Point Narrow Double Word to Word Quadruple}]{FNARROWDWQ\hfill Floating Point Narrow Double Word to Word Quadruple}
\label{instr:FNARROWDWQ}\index{fnarrowdwq}
 
\paragraph{Encoding} Format \textsf{ALU\_FWQWR} (Section~\ref{sec:format-ALU-FWQWR}), \verb|fpucode6 = 1100|\\

\bitfields{ 1/P, 2/11, 1/1, 1/1, 3/floatmode, 6/registerW, 2/11, 3/001, 1/1, 4/1100, 1/ziplanes, 1/0, 5/registerP, 1/0 }


\paragraph{Syntax} \texttt{fnarrowdwq}\emph{floatmode} \emph{registerW} = \emph{registerP}

\paragraph{Description} The \emph{floatmode} and the \emph{registerP} are interpreted as two binary 64 floating-point numbers packed into 128 bits. The \emph{floatmode} and the \emph{registerP} double words are converted to binary 32 floating-point numbers according to the rounding mode. If \emph{ziplanes}, the resulting words are interleaved with the \emph{registerP} words in even lanes and the \emph{floatmode} words in odd lanes. The four resulting words are packed and stored back into the \emph{registerW}. This instruction may raise exception bits in the CS register.


\paragraph{Modifier(s)}
\noindent\begin{itemize}
\item floatmode (cf. floatmode in section \ref{par:modifier-floatmode} on page \pageref{par:modifier-floatmode})
\end{itemize}

\paragraph{Execution} {Reservation table \textsf{ALU\_LITE} (Section~\ref{sec:reservation-ALU-LITE})}
~
{\begin{lstlisting}
stage ID:
  new floatmode = _ZX_3(floatmode);
stage RR:
  new argument2 = PGR[registerP];
  new RM = decode_riscv_float_rounding_mode(floatmode);
  if (ziplanes) {
    result1.32[0] = f64_to_f32(RM, argument2.64[0]);
    result1.32[1] = f64_to_f32(RM, argument3.64[0]);
    result1.32[2] = f64_to_f32(RM, argument2.64[1]);
  } else {
    result1.32[0] = f64_to_f32(RM, argument2.64[0]);
    result1.32[1] = f64_to_f32(RM, argument2.64[1]);
    result1.32[2] = f64_to_f32(RM, argument3.64[0]);
  }
  result1.32[3] = f64_to_f32(RM, argument3.64[1]);
stage E1:
  GPR[registerW] = result1;
  commit_float_exception_flags();
\end{lstlisting}}
\newpage
\subsubsection[{\small FNARROWWHQ: Floating Point Narrow Word to Half Word Quadruple}]{FNARROWWHQ\hfill Floating Point Narrow Word to Half Word Quadruple}
\label{instr:FNARROWWHQ}\index{fnarrowwhq}
 
\paragraph{Encoding} Format \textsf{ALU\_FHQWR} (Section~\ref{sec:format-ALU-FHQWR}), \verb|fpucode6 = 1100|\\

\bitfields{ 1/P, 2/11, 1/1, 1/1, 3/floatmode, 6/registerW, 2/11, 3/001, 1/1, 4/1100, 1/ziplanes, 1/0, 5/registerP, 1/1 }


\paragraph{Syntax} \texttt{fnarrowwhq}\emph{ziplanes}\emph{floatmode} \emph{registerW} = \emph{registerP}

\paragraph{Description} The \emph{floatmode} and the \emph{registerP} are interpreted as two binary 32 floating-point numbers packed into 64 bits. The \emph{floatmode} and the \emph{registerP} words are converted to binary 16 floating-point numbers according to the rounding mode. If \emph{ziplanes}, the resulting half words are interleaved with the \emph{registerP} half words in even lanes and the \emph{floatmode} half words in odd lanes. The four resulting half words are packed and stored back into the \emph{registerW}. This instruction may raise exception bits in the CS register.


\paragraph{Modifier(s)}
\noindent\begin{itemize}
\item floatmode (cf. floatmode in section \ref{par:modifier-floatmode} on page \pageref{par:modifier-floatmode})
\item ziplanes (cf. ziplanes in section \ref{par:modifier-ziplanes} on page \pageref{par:modifier-ziplanes})
\end{itemize}

\paragraph{Execution} {Reservation table \textsf{ALU\_LITE} (Section~\ref{sec:reservation-ALU-LITE})}
~
{\begin{lstlisting}
stage ID:
  new ziplanes = _ZX_1(ziplanes);
  new floatmode = _ZX_3(floatmode);
stage RR:
  new argument2 = PGR[registerP];
  new RM = decode_riscv_float_rounding_mode(floatmode);
  if (ziplanes) {
    result1.16[0] = f32_to_f16(RM, argument2.32[0]);
    result1.16[1] = f32_to_f16(RM, argument3.32[0]);
    result1.16[2] = f32_to_f16(RM, argument2.32[1]);
  } else {
    result1.16[0] = f32_to_f16(RM, argument2.32[0]);
    result1.16[1] = f32_to_f16(RM, argument2.32[1]);
    result1.16[2] = f32_to_f16(RM, argument3.32[0]);
  }
  result1.16[3] = f32_to_f16(RM, argument3.32[1]);
stage E1:
  GPR[registerW] = result1;
  commit_float_exception_flags();
\end{lstlisting}}
\newpage
\subsubsection[{\small FRACTDWQ: Fraction Double Word to Word Quadruple}]{FRACTDWQ\hfill Fraction Double Word to Word Quadruple}
\label{instr:FRACTDWQ}\index{fractdwq}
 
\paragraph{Encoding} Format \textsf{ALU\_WQNWRR} (Section~\ref{sec:format-ALU-WQNWRR}), \verb|alucode8 = 1001|\\

\bitfields{ 1/P, 2/11, 1/1, 4/1111, 5/registerM, 1/--, 2/10, 3/101, 1/0, 4/1001, 1/ziplanes, 1/0, 4/registerR, 1/--, 1/1 }


\paragraph{Syntax} \texttt{fractdwq}\emph{ziplanes} \emph{registerM} = \emph{registerR}

\paragraph{Description} The \emph{ziplanes} and the \emph{registerR} are interpreted as two double words packed into 128 bits. The \emph{ziplanes} and the \emph{registerR} double words are converted to 32-bit words by removing the low bits. If \emph{ziplanes}, the resulting words are interleaved with the \emph{registerR} words in even lanes and the \emph{ziplanes} words in odd lanes. The four resulting words are packed and stored back into the \emph{registerM}.


\paragraph{Modifier(s)}
\noindent\begin{itemize}
\item ziplanes (cf. ziplanes in section \ref{par:modifier-ziplanes} on page \pageref{par:modifier-ziplanes})
\end{itemize}

\paragraph{Execution} {Reservation table \textsf{ALU\_LITE} (Section~\ref{sec:reservation-ALU-LITE})}
~
{\begin{lstlisting}
stage ID:
  new ziplanes = _ZX_1(ziplanes);
stage RR:
  new argument2 = QGR[registerR];
  if (ziplanes) {
    result1.32[0] = argument2.32[1];
    result1.32[1] = argument3.32[1];
    result1.32[2] = argument2.32[3];
  } else {
    result1.32[0] = argument2.32[1];
    result1.32[1] = argument2.32[3];
    result1.32[2] = argument3.32[1];
  }
  result1.32[3] = argument3.32[3];
stage E1:
  PGR[registerM] = result1;
\end{lstlisting}}
\newpage
\subsubsection[{\small FRACTHBX: Fraction Half Word to Byte Hexadecuple}]{FRACTHBX\hfill Fraction Half Word to Byte Hexadecuple}
\label{instr:FRACTHBX}\index{fracthbx}
 
\paragraph{Encoding} Format \textsf{ALU\_BXNWRR} (Section~\ref{sec:format-ALU-BXNWRR}), \verb|alucode8 = 1001|\\

\bitfields{ 1/P, 2/11, 1/1, 4/1111, 5/registerM, 1/--, 2/10, 3/101, 1/1, 4/1001, 1/ziplanes, 1/0, 5/registerP, 1/1 }


\paragraph{Syntax} \texttt{fracthbx}\emph{ziplanes} \emph{registerM} = \emph{registerP}

\paragraph{Description} The \emph{ziplanes} and the \emph{registerP} are interpreted as eight half words packed into 128 bits. The \emph{ziplanes} and the \emph{registerP} half words are converted to bytes by removing the low bits. If \emph{ziplanes}, the resulting bytes are interleaved with the \emph{registerP} bytes in even lanes and the \emph{ziplanes} bytes in odd lanes. The sixteen resulting bytes are packed and stored back into the \emph{registerM}.


\paragraph{Modifier(s)}
\noindent\begin{itemize}
\item ziplanes (cf. ziplanes in section \ref{par:modifier-ziplanes} on page \pageref{par:modifier-ziplanes})
\end{itemize}

\paragraph{Execution} {Reservation table \textsf{ALU\_LITE} (Section~\ref{sec:reservation-ALU-LITE})}
~
{\begin{lstlisting}
stage ID:
  new ziplanes = _ZX_1(ziplanes);
stage RR:
  new argument2 = PGR[registerP];
  if (ziplanes) {
    result1.8[0] = argument2.8[1];
    result1.8[1] = argument3.8[1];
    result1.8[2] = argument2.8[3];
    result1.8[3] = argument3.8[3];
    result1.8[4] = argument2.8[5];
    result1.8[5] = argument3.8[5];
    result1.8[6] = argument2.8[7];
    result1.8[7] = argument3.8[7];
    result1.8[8] = argument2.8[9];
    result1.8[9] = argument3.8[9];
    result1.8[10] = argument2.8[11];
    result1.8[11] = argument3.8[11];
    result1.8[12] = argument2.8[13];
    result1.8[13] = argument3.8[13];
    result1.8[14] = argument2.8[15];
  } else {
    result1.8[0] = argument2.8[1];
    result1.8[1] = argument2.8[3];
    result1.8[2] = argument2.8[5];
    result1.8[3] = argument2.8[7];
    result1.8[4] = argument2.8[9];
    result1.8[5] = argument2.8[11];
    result1.8[6] = argument2.8[13];
    result1.8[7] = argument2.8[15];
    result1.8[8] = argument3.8[1];
    result1.8[9] = argument3.8[3];
    result1.8[10] = argument3.8[5];
    result1.8[11] = argument3.8[7];
    result1.8[12] = argument3.8[9];
    result1.8[13] = argument3.8[11];
    result1.8[14] = argument3.8[13];
  }
  result1.8[15] = argument3.8[15];
stage E1:
  PGR[registerM] = result1;
\end{lstlisting}}
\newpage
\subsubsection[{\small FRACTWHO: Fraction Word to Half Word Octuple}]{FRACTWHO\hfill Fraction Word to Half Word Octuple}
\label{instr:FRACTWHO}\index{fractwho}
 
\paragraph{Encoding} Format \textsf{ALU\_HONWRR} (Section~\ref{sec:format-ALU-HONWRR}), \verb|alucode8 = 1001|\\

\bitfields{ 1/P, 2/11, 1/1, 4/1111, 5/registerM, 1/--, 2/10, 3/101, 1/1, 4/1001, 1/ziplanes, 1/0, 5/registerP, 1/0 }


\paragraph{Syntax} \texttt{fractwho}\emph{ziplanes} \emph{registerM} = \emph{registerP}

\paragraph{Description} The \emph{ziplanes} and the \emph{registerP} are interpreted as four words packed into 128 bits. The \emph{ziplanes} and the \emph{registerP} words are converted to 16-bit half words by removing the low bits. If \emph{ziplanes}, the resulting half words are interleaved with the \emph{registerP} half words in even lanes and the \emph{ziplanes} half words in odd lanes. The eight resulting half words are packed and stored back into the \emph{registerM}.


\paragraph{Modifier(s)}
\noindent\begin{itemize}
\item ziplanes (cf. ziplanes in section \ref{par:modifier-ziplanes} on page \pageref{par:modifier-ziplanes})
\end{itemize}

\paragraph{Execution} {Reservation table \textsf{ALU\_LITE} (Section~\ref{sec:reservation-ALU-LITE})}
~
{\begin{lstlisting}
stage ID:
  new ziplanes = _ZX_1(ziplanes);
stage RR:
  new argument2 = PGR[registerP];
  if (ziplanes) {
    result1.16[0] = argument2.16[1];
    result1.16[1] = argument3.16[1];
    result1.16[2] = argument2.16[3];
    result1.16[3] = argument3.16[3];
    result1.16[4] = argument2.16[5];
    result1.16[5] = argument3.16[5];
    result1.16[6] = argument2.16[7];
  } else {
    result1.16[0] = argument2.16[1];
    result1.16[1] = argument2.16[3];
    result1.16[2] = argument2.16[5];
    result1.16[3] = argument2.16[7];
    result1.16[4] = argument3.16[1];
    result1.16[5] = argument3.16[3];
    result1.16[6] = argument3.16[5];
  }
  result1.16[7] = argument3.16[7];
stage E1:
  PGR[registerM] = result1;
\end{lstlisting}}
\newpage
\subsubsection[{\small FWIDENWDP: Floating Point Widen Word to Double Word Pair}]{FWIDENWDP\hfill Floating Point Widen Word to Double Word Pair}
\label{instr:FWIDENWDP}\index{fwidenwdp}
 
\paragraph{Encoding} Format \textsf{ALU\_FDPWIDEWR} (Section~\ref{sec:format-ALU-FDPWIDEWR}), \verb|fpucode6 = 1000|\\

\bitfields{ 1/P, 2/11, 1/1, 1/1, 3/111, 5/registerM, 1/--, 2/11, 3/000, 1/0, 4/1000, 1/mostsig, 1/0, 5/registerP, 1/1 }


\paragraph{Syntax} \texttt{fwidenwdp}\emph{mostsig} \emph{registerM} = \emph{registerP}

\paragraph{Description} The \emph{registerP} is interpreted as four binary 32 floating-point numbers packed into 128 bits. The least or most significant words of the \emph{registerP} are selected, depending on \emph{mostsig}. The two numbers are converted to binary 64 floating-point numbers packed and stored into the \emph{registerM}.


\paragraph{Modifier(s)}
\noindent\begin{itemize}
\item mostsig (cf. mostsig in section \ref{par:modifier-mostsig} on page \pageref{par:modifier-mostsig})
\end{itemize}

\paragraph{Execution} {Reservation table \textsf{ALU\_LITE} (Section~\ref{sec:reservation-ALU-LITE})}
~
{\begin{lstlisting}
stage ID:
  new mostsig = _ZX_1(mostsig);
stage RR:
  new argument2 = PGR[registerP];
  if (mostsig) {
    argument2 = argument2 >> 64;
  }
  for (I = 0; I < 2; I += 1) {
    result1.64[I] = f32_to_f64(argument2.32[I])
  };
stage E1:
  PGR[registerM] = result1;
\end{lstlisting}}
\newpage
\subsubsection[{\small INSF: Insert Bit Field into Word}]{INSF\hfill Insert Bit Field into Word}
\label{instr:INSF}\index{insf}
 
\paragraph{Encoding} Format \textsf{ALU\_BFWR} (Section~\ref{sec:format-ALU-BFWR}), \verb|wricode1 = 00|\\

\bitfields{ 1/P, 2/11, 1/0, 2/00, 2/stopbit2, 6/registerW, 2/01, 4/stopbit4, 6/startbit, 6/registerZ }


\paragraph{Syntax} \texttt{insf} \emph{registerW} = \emph{registerZ}, \emph{stopbit2\_\discretionary{}{}{}stopbit4}, \emph{startbit}

\paragraph{Description} The start bit position and the stop bit position are extracted from the \emph{startbit} and the \emph{stopbit2\_\discretionary{}{}{}stopbit4} respectively. A bit-mask is computed with ones set between the start bit position and the stop bit position inclusive. The \emph{registerZ} is shifted left by the start bit position and inserted into the \emph{registerW} according to the bit-mask. Behavior is undefined if the start bit position is greater than the stop bit position.


\paragraph{Execution} {Reservation table \textsf{ALU\_TINY} (Section~\ref{sec:reservation-ALU-TINY})}
~
{\begin{lstlisting}
stage ID:
  new startbit = _ZX_6(startbit);
  new stopbit = _ZX_6(stopbit2_stopbit4);
stage RR:
  new bias = startbit <= stopbit;
  new mask = _ZX_64(((2 << stopbit) - ((2 - bias) << startbit)) + (bias - 1));
  new argument2 = GPR[registerZ];
  new argument1 = GPR[registerW];
  new result1 = ((argument2 << startbit) & mask) | (argument1 & ~mask);
stage E1:
  GPR[registerW] = result1;
\end{lstlisting}}
\newpage
\subsubsection[{\small IORD: Bitwise Inclusive Or Double Word}]{IORD\hfill Bitwise Inclusive Or Double Word}
\label{instr:IORD}\index{iord}
 
\paragraph{Encoding} Format \textsf{ALU\_DWRR0} (Section~\ref{sec:format-ALU-DWRR0}), \verb|exucode2 = 1010|\\

\bitfields{ 1/P, 2/11, 1/0, 4/1010, 6/registerW, 2/10, 3/000, 1/0, 6/registerY, 6/registerZ }


\paragraph{Syntax} \texttt{iord} \emph{registerW} = \emph{registerZ}, \emph{registerY}

\paragraph{Description} Bitwise inclusive or of the \emph{registerZ} with the \emph{registerY}. The result is stored into the \emph{registerW}.


\paragraph{Execution} {Reservation table \textsf{ALU\_TINY} (Section~\ref{sec:reservation-ALU-TINY})}
~
{\begin{lstlisting}
stage RR:
  new argument3 = GPR[registerY];
  new argument2 = GPR[registerZ];
  new result1 = argument2 | argument3;
stage E1:
  GPR[registerW] = result1;
\end{lstlisting}}
\newpage

\paragraph{Encoding} Format \textsf{ALU\_DWRR0.M} (Section~\ref{sec:format-ALU-DWRR0.M}), \verb|exucode2 = 1010|\\

\bitfields{ 1/P, 2/00, 2/-- --, 27/upper27, 1/1, 2/11, 1/0, 4/1010, 6/registerW, 2/10, 3/000, 1/0, 1/splat32, 5/lower5, 6/registerZ }


\paragraph{Syntax} \texttt{iord} \emph{registerW} = \emph{registerZ}, \emph{upper27\_\discretionary{}{}{}lower5}\emph{splat32}

\paragraph{Description} Bitwise inclusive or of the \emph{registerZ} with the \emph{upper27\_\discretionary{}{}{}lower5}. The result is stored into the \emph{registerW}.


\paragraph{Modifier(s)}
\noindent\begin{itemize}
\item splat32 (cf. splat32 in section \ref{par:modifier-splat32} on page \pageref{par:modifier-splat32})
\end{itemize}

\paragraph{Execution} {Reservation table \textsf{ALU\_TINY.X} (Section~\ref{sec:reservation-ALU-TINY.X})}
~
{\begin{lstlisting}
stage ID:
  new splat32 = _ZX_1(splat32);
stage RR:
  new argument3 = splat32 ? (_WX_32(upper27_lower5) << 32) | _ZX_32(_WX_32(upper27_lower5)) : _SX_32(_WX_32(upper27_lower5));
  new argument2 = GPR[registerZ];
  new result1 = argument2 | argument3;
stage E1:
  GPR[registerW] = result1;
\end{lstlisting}}
\newpage

\paragraph{Encoding} Format \textsf{ALU\_DWRI} (Section~\ref{sec:format-ALU-DWRI}), \verb|exucode2 = 1010|\\

\bitfields{ 1/P, 2/11, 1/0, 4/1010, 6/registerW, 2/00, 10/signed10, 6/registerZ }


\paragraph{Syntax} \texttt{iord} \emph{registerW} = \emph{registerZ}, \emph{signed10}

\paragraph{Description} Bitwise inclusive or of the \emph{registerZ} with the \emph{signed10}. The result is stored into the \emph{registerW}.


\paragraph{Execution} {Reservation table \textsf{ALU\_TINY} (Section~\ref{sec:reservation-ALU-TINY})}
~
{\begin{lstlisting}
stage RR:
  new argument3 = _SX_10(signed10);
  new argument2 = GPR[registerZ];
  new result1 = argument2 | argument3;
stage E1:
  GPR[registerW] = result1;
\end{lstlisting}}
\newpage

\paragraph{Encoding} Format \textsf{ALU\_DWRI.X} (Section~\ref{sec:format-ALU-DWRI.X}), \verb|exucode2 = 1010|\\

\bitfields{ 1/P, 2/00, 2/-- --, 27/upper27, 1/1, 2/11, 1/0, 4/1010, 6/registerW, 2/00, 10/lower10, 6/registerZ }


\paragraph{Syntax} \texttt{iord} \emph{registerW} = \emph{registerZ}, \emph{upper27\_\discretionary{}{}{}lower10}

\paragraph{Description} Bitwise inclusive or of the \emph{registerZ} with the \emph{upper27\_\discretionary{}{}{}lower10}. The result is stored into the \emph{registerW}.


\paragraph{Execution} {Reservation table \textsf{ALU\_TINY.X} (Section~\ref{sec:reservation-ALU-TINY.X})}
~
{\begin{lstlisting}
stage RR:
  new argument3 = _SX_37(upper27_lower10);
  new argument2 = GPR[registerZ];
  new result1 = argument2 | argument3;
stage E1:
  GPR[registerW] = result1;
\end{lstlisting}}
\newpage

\paragraph{Encoding} Format \textsf{ALU\_DWRI.Y} (Section~\ref{sec:format-ALU-DWRI.Y}), \verb|exucode2 = 1010|\\

\bitfields{ 1/P, 2/00, 2/-- --, 27/extend27, 1/1, 2/00, 2/-- --, 27/upper27, 1/1, 2/11, 1/0, 4/1010, 6/registerW, 2/00, 10/lower10, 6/registerZ }


\paragraph{Syntax} \texttt{iord} \emph{registerW} = \emph{registerZ}, \emph{extend27\_\discretionary{}{}{}upper27\_\discretionary{}{}{}lower10}

\paragraph{Description} Bitwise inclusive or of the \emph{registerZ} with the \emph{extend27\_\discretionary{}{}{}upper27\_\discretionary{}{}{}lower10}. The result is stored into the \emph{registerW}.


\paragraph{Execution} {Reservation table \textsf{ALU\_TINY.Y} (Section~\ref{sec:reservation-ALU-TINY.Y})}
~
{\begin{lstlisting}
stage RR:
  new argument3 = _WX_64(extend27_upper27_lower10);
  new argument2 = GPR[registerZ];
  new result1 = argument2 | argument3;
stage E1:
  GPR[registerW] = result1;
\end{lstlisting}}
\newpage
\subsubsection[{\small IORND: Bitwise Inclusive Or with Not Operand Double Word}]{IORND\hfill Bitwise Inclusive Or with Not Operand Double Word}
\label{instr:IORND}\index{iornd}
 
\paragraph{Encoding} Format \textsf{ALU\_DWRR0} (Section~\ref{sec:format-ALU-DWRR0}), \verb|exucode2 = 1111|\\

\bitfields{ 1/P, 2/11, 1/0, 4/1111, 6/registerW, 2/10, 3/000, 1/0, 6/registerY, 6/registerZ }


\paragraph{Syntax} \texttt{iornd} \emph{registerW} = \emph{registerZ}, \emph{registerY}

\paragraph{Description} Bitwise inclusive or of not the \emph{registerZ} with the \emph{registerY}. The result is stored into the \emph{registerW}.


\paragraph{Execution} {Reservation table \textsf{ALU\_TINY} (Section~\ref{sec:reservation-ALU-TINY})}
~
{\begin{lstlisting}
stage RR:
  new argument3 = GPR[registerY];
  new argument2 = GPR[registerZ];
  new result1 = ~argument2 | argument3;
stage E1:
  GPR[registerW] = result1;
\end{lstlisting}}
\newpage

\paragraph{Encoding} Format \textsf{ALU\_DWRR0.M} (Section~\ref{sec:format-ALU-DWRR0.M}), \verb|exucode2 = 1111|\\

\bitfields{ 1/P, 2/00, 2/-- --, 27/upper27, 1/1, 2/11, 1/0, 4/1111, 6/registerW, 2/10, 3/000, 1/0, 1/splat32, 5/lower5, 6/registerZ }


\paragraph{Syntax} \texttt{iornd} \emph{registerW} = \emph{registerZ}, \emph{upper27\_\discretionary{}{}{}lower5}\emph{splat32}

\paragraph{Description} Bitwise inclusive or of not the \emph{registerZ} with the \emph{upper27\_\discretionary{}{}{}lower5}. The result is stored into the \emph{registerW}.


\paragraph{Modifier(s)}
\noindent\begin{itemize}
\item splat32 (cf. splat32 in section \ref{par:modifier-splat32} on page \pageref{par:modifier-splat32})
\end{itemize}

\paragraph{Execution} {Reservation table \textsf{ALU\_TINY.X} (Section~\ref{sec:reservation-ALU-TINY.X})}
~
{\begin{lstlisting}
stage ID:
  new splat32 = _ZX_1(splat32);
stage RR:
  new argument3 = splat32 ? (_WX_32(upper27_lower5) << 32) | _ZX_32(_WX_32(upper27_lower5)) : _SX_32(_WX_32(upper27_lower5));
  new argument2 = GPR[registerZ];
  new result1 = ~argument2 | argument3;
stage E1:
  GPR[registerW] = result1;
\end{lstlisting}}
\newpage

\paragraph{Encoding} Format \textsf{ALU\_DWRI} (Section~\ref{sec:format-ALU-DWRI}), \verb|exucode2 = 1111|\\

\bitfields{ 1/P, 2/11, 1/0, 4/1111, 6/registerW, 2/00, 10/signed10, 6/registerZ }


\paragraph{Syntax} \texttt{iornd} \emph{registerW} = \emph{registerZ}, \emph{signed10}

\paragraph{Description} Bitwise inclusive or of not the \emph{registerZ} with the \emph{signed10}. The result is stored into the \emph{registerW}.


\paragraph{Execution} {Reservation table \textsf{ALU\_TINY} (Section~\ref{sec:reservation-ALU-TINY})}
~
{\begin{lstlisting}
stage RR:
  new argument3 = _SX_10(signed10);
  new argument2 = GPR[registerZ];
  new result1 = ~argument2 | argument3;
stage E1:
  GPR[registerW] = result1;
\end{lstlisting}}
\newpage

\paragraph{Encoding} Format \textsf{ALU\_DWRI.X} (Section~\ref{sec:format-ALU-DWRI.X}), \verb|exucode2 = 1111|\\

\bitfields{ 1/P, 2/00, 2/-- --, 27/upper27, 1/1, 2/11, 1/0, 4/1111, 6/registerW, 2/00, 10/lower10, 6/registerZ }


\paragraph{Syntax} \texttt{iornd} \emph{registerW} = \emph{registerZ}, \emph{upper27\_\discretionary{}{}{}lower10}

\paragraph{Description} Bitwise inclusive or of not the \emph{registerZ} with the \emph{upper27\_\discretionary{}{}{}lower10}. The result is stored into the \emph{registerW}.


\paragraph{Execution} {Reservation table \textsf{ALU\_TINY.X} (Section~\ref{sec:reservation-ALU-TINY.X})}
~
{\begin{lstlisting}
stage RR:
  new argument3 = _SX_37(upper27_lower10);
  new argument2 = GPR[registerZ];
  new result1 = ~argument2 | argument3;
stage E1:
  GPR[registerW] = result1;
\end{lstlisting}}
\newpage

\paragraph{Encoding} Format \textsf{ALU\_DWRI.Y} (Section~\ref{sec:format-ALU-DWRI.Y}), \verb|exucode2 = 1111|\\

\bitfields{ 1/P, 2/00, 2/-- --, 27/extend27, 1/1, 2/00, 2/-- --, 27/upper27, 1/1, 2/11, 1/0, 4/1111, 6/registerW, 2/00, 10/lower10, 6/registerZ }


\paragraph{Syntax} \texttt{iornd} \emph{registerW} = \emph{registerZ}, \emph{extend27\_\discretionary{}{}{}upper27\_\discretionary{}{}{}lower10}

\paragraph{Description} Bitwise inclusive or of not the \emph{registerZ} with the \emph{extend27\_\discretionary{}{}{}upper27\_\discretionary{}{}{}lower10}. The result is stored into the \emph{registerW}.


\paragraph{Execution} {Reservation table \textsf{ALU\_TINY.Y} (Section~\ref{sec:reservation-ALU-TINY.Y})}
~
{\begin{lstlisting}
stage RR:
  new argument3 = _WX_64(extend27_upper27_lower10);
  new argument2 = GPR[registerZ];
  new result1 = ~argument2 | argument3;
stage E1:
  GPR[registerW] = result1;
\end{lstlisting}}
\newpage
\subsubsection[{\small IORNQ: Bitwise Inclusive Or with Not Operand Quadruple Word}]{IORNQ\hfill Bitwise Inclusive Or with Not Operand Quadruple Word}
\label{instr:IORNQ}\index{iornq}
 
\paragraph{Encoding} Format \textsf{ALU\_QWRR0} (Section~\ref{sec:format-ALU-QWRR0}), \verb|exucode2 = 1111|\\

\bitfields{ 1/P, 2/11, 1/0, 4/1111, 5/registerM, 1/--, 2/10, 3/000, 1/1, 5/registerO, 1/--, 5/registerP, 1/-- }


\paragraph{Syntax} \texttt{iornq} \emph{registerM} = \emph{registerP}, \emph{registerO}

\paragraph{Description} Bitwise inclusive or of not the \emph{registerP} with the \emph{registerO}. The result is stored into the \emph{registerM}.


\paragraph{Execution} {Reservation table \textsf{ALU\_LITE} (Section~\ref{sec:reservation-ALU-LITE})}
~
{\begin{lstlisting}
stage RR:
  new argument3 = PGR[registerO];
  new argument2 = PGR[registerP];
  new result1 = ~argument2 | argument3;
stage E1:
  PGR[registerM] = result1;
\end{lstlisting}}
\newpage

\paragraph{Encoding} Format \textsf{ALU\_QWRR0.M} (Section~\ref{sec:format-ALU-QWRR0.M}), \verb|exucode2 = 1111|\\

\bitfields{ 1/P, 2/00, 2/-- --, 27/upper27, 1/1, 2/11, 1/0, 4/1111, 5/registerM, 1/--, 2/10, 3/000, 1/1, 1/splat32, 5/lower5, 5/registerP, 1/1 }


\paragraph{Syntax} \texttt{iornq} \emph{registerM} = \emph{registerP}, \emph{upper27\_\discretionary{}{}{}lower5}\emph{splat32}

\paragraph{Description} Bitwise inclusive or of not the \emph{registerP} with the \emph{upper27\_\discretionary{}{}{}lower5}. The result is stored into the \emph{registerM}.


\paragraph{Modifier(s)}
\noindent\begin{itemize}
\item splat32 (cf. splat32 in section \ref{par:modifier-splat32} on page \pageref{par:modifier-splat32})
\end{itemize}

\paragraph{Execution} {Reservation table \textsf{ALU\_LITE.X} (Section~\ref{sec:reservation-ALU-LITE.X})}
~
{\begin{lstlisting}
stage ID:
  new splat32 = _ZX_1(splat32);
stage RR:
  new argument3 = splat32 ? (_WX_32(upper27_lower5) << 32) | _ZX_32(_WX_32(upper27_lower5)) : _SX_32(_WX_32(upper27_lower5));
  new argument2 = PGR[registerP];
  new result1 = ~argument2 | argument3;
stage E1:
  PGR[registerM] = result1;
\end{lstlisting}}
\newpage
\subsubsection[{\small IORNW: Bitwise Inclusive Or with Not Operand Double Word}]{IORNW\hfill Bitwise Inclusive Or with Not Operand Double Word}
\label{instr:IORNW}\index{iornw}
 
\paragraph{Encoding} Format \textsf{ALU\_WWRR0} (Section~\ref{sec:format-ALU-WWRR0}), \verb|exucode2 = 1111|\\

\bitfields{ 1/P, 2/11, 1/0, 4/1111, 6/registerW, 2/11, 3/000, 1/signextw, 6/registerY, 6/registerZ }


\paragraph{Syntax} \texttt{iornw}\emph{signextw} \emph{registerW} = \emph{registerZ}, \emph{registerY}

\paragraph{Description} Bitwise inclusive or of not the \emph{registerZ} with the \emph{registerY}. The result with the 32 upper bits cleared is stored into the \emph{registerW}.


\paragraph{Modifier(s)}
\noindent\begin{itemize}
\item signextw (cf. signextw in section \ref{par:modifier-signextw} on page \pageref{par:modifier-signextw})
\end{itemize}

\paragraph{Execution} {Reservation table \textsf{ALU\_TINY} (Section~\ref{sec:reservation-ALU-TINY})}
~
{\begin{lstlisting}
stage ID:
  new signextw = _ZX_1(signextw);
stage RR:
  new argument3 = GPR[registerY];
  new argument2 = GPR[registerZ];
  new result1 = ~argument2 | argument3;
stage E1:
  GPR[registerW] = signextw ? _SX_32(result1) : _ZX_32(result1);
\end{lstlisting}}
\newpage

\paragraph{Encoding} Format \textsf{ALU\_WWRR0.W} (Section~\ref{sec:format-ALU-WWRR0.W}), \verb|exucode2 = 1111|\\

\bitfields{ 1/P, 2/00, 2/-- --, 27/upper27, 1/1, 2/11, 1/0, 4/1111, 6/registerW, 2/11, 3/000, 1/signextw, 1/0, 5/lower5, 6/registerZ }


\paragraph{Syntax} \texttt{iornw}\emph{signextw} \emph{registerW} = \emph{registerZ}, \emph{upper27\_\discretionary{}{}{}lower5}

\paragraph{Description} Bitwise inclusive or of not the \emph{registerZ} with the \emph{upper27\_\discretionary{}{}{}lower5}. The result with the 32 upper bits cleared is stored into the \emph{registerW}.


\paragraph{Modifier(s)}
\noindent\begin{itemize}
\item signextw (cf. signextw in section \ref{par:modifier-signextw} on page \pageref{par:modifier-signextw})
\end{itemize}

\paragraph{Execution} {Reservation table \textsf{ALU\_TINY.X} (Section~\ref{sec:reservation-ALU-TINY.X})}
~
{\begin{lstlisting}
stage ID:
  new signextw = _ZX_1(signextw);
stage RR:
  new argument3 = _WX_32(upper27_lower5);
  new argument2 = GPR[registerZ];
  new result1 = ~argument2 | argument3;
stage E1:
  GPR[registerW] = signextw ? _SX_32(result1) : _ZX_32(result1);
\end{lstlisting}}
\newpage
\subsubsection[{\small IORQ: Bitwise Inclusive Or Quadruple Word}]{IORQ\hfill Bitwise Inclusive Or Quadruple Word}
\label{instr:IORQ}\index{iorq}
 
\paragraph{Encoding} Format \textsf{ALU\_QWRR0} (Section~\ref{sec:format-ALU-QWRR0}), \verb|exucode2 = 1010|\\

\bitfields{ 1/P, 2/11, 1/0, 4/1010, 5/registerM, 1/--, 2/10, 3/000, 1/1, 5/registerO, 1/--, 5/registerP, 1/-- }


\paragraph{Syntax} \texttt{iorq} \emph{registerM} = \emph{registerP}, \emph{registerO}

\paragraph{Description} Bitwise inclusive or of the \emph{registerP} with the \emph{registerO}. The result is stored into the \emph{registerM}.


\paragraph{Execution} {Reservation table \textsf{ALU\_LITE} (Section~\ref{sec:reservation-ALU-LITE})}
~
{\begin{lstlisting}
stage RR:
  new argument3 = PGR[registerO];
  new argument2 = PGR[registerP];
  new result1 = argument2 | argument3;
stage E1:
  PGR[registerM] = result1;
\end{lstlisting}}
\newpage

\paragraph{Encoding} Format \textsf{ALU\_QWRR0.M} (Section~\ref{sec:format-ALU-QWRR0.M}), \verb|exucode2 = 1010|\\

\bitfields{ 1/P, 2/00, 2/-- --, 27/upper27, 1/1, 2/11, 1/0, 4/1010, 5/registerM, 1/--, 2/10, 3/000, 1/1, 1/splat32, 5/lower5, 5/registerP, 1/1 }


\paragraph{Syntax} \texttt{iorq} \emph{registerM} = \emph{registerP}, \emph{upper27\_\discretionary{}{}{}lower5}\emph{splat32}

\paragraph{Description} Bitwise inclusive or of the \emph{registerP} with the \emph{upper27\_\discretionary{}{}{}lower5}. The result is stored into the \emph{registerM}.


\paragraph{Modifier(s)}
\noindent\begin{itemize}
\item splat32 (cf. splat32 in section \ref{par:modifier-splat32} on page \pageref{par:modifier-splat32})
\end{itemize}

\paragraph{Execution} {Reservation table \textsf{ALU\_LITE.X} (Section~\ref{sec:reservation-ALU-LITE.X})}
~
{\begin{lstlisting}
stage ID:
  new splat32 = _ZX_1(splat32);
stage RR:
  new argument3 = splat32 ? (_WX_32(upper27_lower5) << 32) | _ZX_32(_WX_32(upper27_lower5)) : _SX_32(_WX_32(upper27_lower5));
  new argument2 = PGR[registerP];
  new result1 = argument2 | argument3;
stage E1:
  PGR[registerM] = result1;
\end{lstlisting}}
\newpage
\subsubsection[{\small IORW: Bitwise Inclusive Or Word}]{IORW\hfill Bitwise Inclusive Or Word}
\label{instr:IORW}\index{iorw}
 
\paragraph{Encoding} Format \textsf{ALU\_WWRR0} (Section~\ref{sec:format-ALU-WWRR0}), \verb|exucode2 = 1010|\\

\bitfields{ 1/P, 2/11, 1/0, 4/1010, 6/registerW, 2/11, 3/000, 1/signextw, 6/registerY, 6/registerZ }


\paragraph{Syntax} \texttt{iorw}\emph{signextw} \emph{registerW} = \emph{registerZ}, \emph{registerY}

\paragraph{Description} Bitwise inclusive or of the \emph{registerZ} with the \emph{registerY}. The result with the 32 upper bits cleared is stored into the \emph{registerW}.


\paragraph{Modifier(s)}
\noindent\begin{itemize}
\item signextw (cf. signextw in section \ref{par:modifier-signextw} on page \pageref{par:modifier-signextw})
\end{itemize}

\paragraph{Execution} {Reservation table \textsf{ALU\_TINY} (Section~\ref{sec:reservation-ALU-TINY})}
~
{\begin{lstlisting}
stage ID:
  new signextw = _ZX_1(signextw);
stage RR:
  new argument3 = GPR[registerY];
  new argument2 = GPR[registerZ];
  new result1 = argument2 | argument3;
stage E1:
  GPR[registerW] = signextw ? _SX_32(result1) : _ZX_32(result1);
\end{lstlisting}}
\newpage

\paragraph{Encoding} Format \textsf{ALU\_WWRR0.W} (Section~\ref{sec:format-ALU-WWRR0.W}), \verb|exucode2 = 1010|\\

\bitfields{ 1/P, 2/00, 2/-- --, 27/upper27, 1/1, 2/11, 1/0, 4/1010, 6/registerW, 2/11, 3/000, 1/signextw, 1/0, 5/lower5, 6/registerZ }


\paragraph{Syntax} \texttt{iorw}\emph{signextw} \emph{registerW} = \emph{registerZ}, \emph{upper27\_\discretionary{}{}{}lower5}

\paragraph{Description} Bitwise inclusive or of the \emph{registerZ} with the \emph{upper27\_\discretionary{}{}{}lower5}. The result with the 32 upper bits cleared is stored into the \emph{registerW}.


\paragraph{Modifier(s)}
\noindent\begin{itemize}
\item signextw (cf. signextw in section \ref{par:modifier-signextw} on page \pageref{par:modifier-signextw})
\end{itemize}

\paragraph{Execution} {Reservation table \textsf{ALU\_TINY.X} (Section~\ref{sec:reservation-ALU-TINY.X})}
~
{\begin{lstlisting}
stage ID:
  new signextw = _ZX_1(signextw);
stage RR:
  new argument3 = _WX_32(upper27_lower5);
  new argument2 = GPR[registerZ];
  new result1 = argument2 | argument3;
stage E1:
  GPR[registerW] = signextw ? _SX_32(result1) : _ZX_32(result1);
\end{lstlisting}}
\newpage
\subsubsection[{\small LANDD: Logical And Double Word}]{LANDD\hfill Logical And Double Word}
\label{instr:LANDD}\index{landd}
 
\paragraph{Encoding} Format \textsf{ALU\_DWRR2} (Section~\ref{sec:format-ALU-DWRR2}), \verb|exucode2 = 0000|\\

\bitfields{ 1/P, 2/11, 1/0, 4/0000, 6/registerW, 2/10, 3/010, 1/0, 6/registerY, 6/registerZ }


\paragraph{Syntax} \texttt{landd} \emph{registerW} = \emph{registerZ}, \emph{registerY}

\paragraph{Description} The \emph{registerY} and the \emph{registerZ} are converted to booleans, which are AND-ed. The result is stored into the \emph{registerW}.


\paragraph{Execution} {Reservation table \textsf{ALU\_TINY} (Section~\ref{sec:reservation-ALU-TINY})}
~
{\begin{lstlisting}
stage RR:
  new argument3 = GPR[registerY];
  new argument2 = GPR[registerZ];
  new result1 = _ZX_64(argument3) && _ZX_64(argument2);
stage E1:
  GPR[registerW] = result1;
\end{lstlisting}}
\newpage
\subsubsection[{\small LANDW: Logical And Word}]{LANDW\hfill Logical And Word}
\label{instr:LANDW}\index{landw}
 
\paragraph{Encoding} Format \textsf{ALU\_WWRR2} (Section~\ref{sec:format-ALU-WWRR2}), \verb|exucode2 = 0000|\\

\bitfields{ 1/P, 2/11, 1/0, 4/0000, 6/registerW, 2/11, 3/010, 1/0, 6/registerY, 6/registerZ }


\paragraph{Syntax} \texttt{landw} \emph{registerW} = \emph{registerZ}, \emph{registerY}

\paragraph{Description} The \emph{registerY} and the \emph{registerZ} are converted to booleans, which are AND-ed. The result is stored into the \emph{registerW}.


\paragraph{Execution} {Reservation table \textsf{ALU\_TINY} (Section~\ref{sec:reservation-ALU-TINY})}
~
{\begin{lstlisting}
stage RR:
  new argument3 = GPR[registerY];
  new argument2 = GPR[registerZ];
  new result1 = argument3.32[0] && argument2.32[0];
stage E1:
  GPR[registerW] = result1;
\end{lstlisting}}
\newpage
\subsubsection[{\small LIORD: Logical Or Double Word}]{LIORD\hfill Logical Or Double Word}
\label{instr:LIORD}\index{liord}
 
\paragraph{Encoding} Format \textsf{ALU\_DWRR2} (Section~\ref{sec:format-ALU-DWRR2}), \verb|exucode2 = 0010|\\

\bitfields{ 1/P, 2/11, 1/0, 4/0010, 6/registerW, 2/10, 3/010, 1/0, 6/registerY, 6/registerZ }


\paragraph{Syntax} \texttt{liord} \emph{registerW} = \emph{registerZ}, \emph{registerY}

\paragraph{Description} The \emph{registerY} and the \emph{registerZ} are converted to booleans, which are OR-ed. The result is stored into the \emph{registerW}.


\paragraph{Execution} {Reservation table \textsf{ALU\_TINY} (Section~\ref{sec:reservation-ALU-TINY})}
~
{\begin{lstlisting}
stage RR:
  new argument3 = GPR[registerY];
  new argument2 = GPR[registerZ];
  new result1 = _ZX_64(argument3) || _ZX_64(argument2);
stage E1:
  GPR[registerW] = result1;
\end{lstlisting}}
\newpage
\subsubsection[{\small LIORW: Logical Or Word}]{LIORW\hfill Logical Or Word}
\label{instr:LIORW}\index{liorw}
 
\paragraph{Encoding} Format \textsf{ALU\_WWRR2} (Section~\ref{sec:format-ALU-WWRR2}), \verb|exucode2 = 0010|\\

\bitfields{ 1/P, 2/11, 1/0, 4/0010, 6/registerW, 2/11, 3/010, 1/0, 6/registerY, 6/registerZ }


\paragraph{Syntax} \texttt{liorw} \emph{registerW} = \emph{registerZ}, \emph{registerY}

\paragraph{Description} The \emph{registerY} and the \emph{registerZ} are converted to booleans, which are OR-ed. The result is stored into the \emph{registerW}.


\paragraph{Execution} {Reservation table \textsf{ALU\_TINY} (Section~\ref{sec:reservation-ALU-TINY})}
~
{\begin{lstlisting}
stage RR:
  new argument3 = GPR[registerY];
  new argument2 = GPR[registerZ];
  new result1 = argument3.32[0] || argument2.32[0];
stage E1:
  GPR[registerW] = result1;
\end{lstlisting}}
\newpage
\subsubsection[{\small LNANDD: Logical Not And Double Word}]{LNANDD\hfill Logical Not And Double Word}
\label{instr:LNANDD}\index{lnandd}
 
\paragraph{Encoding} Format \textsf{ALU\_DWRR2} (Section~\ref{sec:format-ALU-DWRR2}), \verb|exucode2 = 0001|\\

\bitfields{ 1/P, 2/11, 1/0, 4/0001, 6/registerW, 2/10, 3/010, 1/0, 6/registerY, 6/registerZ }


\paragraph{Syntax} \texttt{lnandd} \emph{registerW} = \emph{registerZ}, \emph{registerY}

\paragraph{Description} The \emph{registerY} and the \emph{registerZ} are converted to booleans, which are NOT AND-ed. The result is stored into the \emph{registerW}.


\paragraph{Execution} {Reservation table \textsf{ALU\_TINY} (Section~\ref{sec:reservation-ALU-TINY})}
~
{\begin{lstlisting}
stage RR:
  new argument3 = GPR[registerY];
  new argument2 = GPR[registerZ];
  new result1 = !(_ZX_64(argument3) && _ZX_64(argument2));
stage E1:
  GPR[registerW] = result1;
\end{lstlisting}}
\newpage
\subsubsection[{\small LNANDW: Logical Not And Word}]{LNANDW\hfill Logical Not And Word}
\label{instr:LNANDW}\index{lnandw}
 
\paragraph{Encoding} Format \textsf{ALU\_WWRR2} (Section~\ref{sec:format-ALU-WWRR2}), \verb|exucode2 = 0001|\\

\bitfields{ 1/P, 2/11, 1/0, 4/0001, 6/registerW, 2/11, 3/010, 1/0, 6/registerY, 6/registerZ }


\paragraph{Syntax} \texttt{lnandw} \emph{registerW} = \emph{registerZ}, \emph{registerY}

\paragraph{Description} The \emph{registerY} and the \emph{registerZ} are converted to booleans, which are NOT AND-ed. The result is stored into the \emph{registerW}.


\paragraph{Execution} {Reservation table \textsf{ALU\_TINY} (Section~\ref{sec:reservation-ALU-TINY})}
~
{\begin{lstlisting}
stage RR:
  new argument3 = GPR[registerY];
  new argument2 = GPR[registerZ];
  new result1 = !(argument3.32[0] && argument2.32[0]);
stage E1:
  GPR[registerW] = result1;
\end{lstlisting}}
\newpage
\subsubsection[{\small LNIORD: Logical Not Or Double Word}]{LNIORD\hfill Logical Not Or Double Word}
\label{instr:LNIORD}\index{lniord}
 
\paragraph{Encoding} Format \textsf{ALU\_DWRR2} (Section~\ref{sec:format-ALU-DWRR2}), \verb|exucode2 = 0011|\\

\bitfields{ 1/P, 2/11, 1/0, 4/0011, 6/registerW, 2/10, 3/010, 1/0, 6/registerY, 6/registerZ }


\paragraph{Syntax} \texttt{lniord} \emph{registerW} = \emph{registerZ}, \emph{registerY}

\paragraph{Description} The \emph{registerY} and the \emph{registerZ} are converted to booleans, which are not OR-ed. The result is stored into the \emph{registerW}.


\paragraph{Execution} {Reservation table \textsf{ALU\_TINY} (Section~\ref{sec:reservation-ALU-TINY})}
~
{\begin{lstlisting}
stage RR:
  new argument3 = GPR[registerY];
  new argument2 = GPR[registerZ];
  new result1 = !(_ZX_64(argument3) || _ZX_64(argument2));
stage E1:
  GPR[registerW] = result1;
\end{lstlisting}}
\newpage
\subsubsection[{\small LNIORW: Logical Not Or Word}]{LNIORW\hfill Logical Not Or Word}
\label{instr:LNIORW}\index{lniorw}
 
\paragraph{Encoding} Format \textsf{ALU\_WWRR2} (Section~\ref{sec:format-ALU-WWRR2}), \verb|exucode2 = 0011|\\

\bitfields{ 1/P, 2/11, 1/0, 4/0011, 6/registerW, 2/11, 3/010, 1/0, 6/registerY, 6/registerZ }


\paragraph{Syntax} \texttt{lniorw} \emph{registerW} = \emph{registerZ}, \emph{registerY}

\paragraph{Description} The \emph{registerY} and the \emph{registerZ} are converted to booleans, which are not OR-ed. The result is stored into the \emph{registerW}.


\paragraph{Execution} {Reservation table \textsf{ALU\_TINY} (Section~\ref{sec:reservation-ALU-TINY})}
~
{\begin{lstlisting}
stage RR:
  new argument3 = GPR[registerY];
  new argument2 = GPR[registerZ];
  new result1 = !(argument3.32[0] || argument2.32[0]);
stage E1:
  GPR[registerW] = result1;
\end{lstlisting}}
\newpage
\subsubsection[{\small MADDBHO: Multiply Add Extend Byte to Half Word Octuple}]{MADDBHO\hfill Multiply Add Extend Byte to Half Word Octuple}
\label{instr:MADDBHO}\index{maddbho}
 
\paragraph{Encoding} Format \textsf{ALU\_HOMEWRRRE} (Section~\ref{sec:format-ALU-HOMEWRRRE}), \verb|alucode2 = 10|\\

\bitfields{ 1/P, 2/11, 1/1, 2/10, 2/widemult, 5/registerM, 1/0, 2/10, 3/110, 1/1, 5/registerYe, 1/1, 5/registerZe, 1/0 }


\paragraph{Syntax} \texttt{maddbho}\emph{widemult} \emph{registerM} = \emph{registerZe}, \emph{registerYe}

\paragraph{Description} The \emph{registerYe} and the \emph{registerZe} are interpreted as eight bytes packed into 64 bits. These \emph{registerYe} and \emph{registerZe} bytes are sign-extended or zero-extended from 8 bits, depending on the \emph{widemult}. The \emph{registerYe} bytes are multiplied by the \emph{registerZe} bytes, producing eight 16-bit results. The eight resulting half words are packed, added to and stored into the \emph{registerM}.


\paragraph{Modifier(s)}
\noindent\begin{itemize}
\item widemult (cf. widemult in section \ref{par:modifier-widemult} on page \pageref{par:modifier-widemult})
\end{itemize}

\paragraph{Execution} {Reservation table \textsf{ALU\_LITE} (Section~\ref{sec:reservation-ALU-LITE})}
~
{\begin{lstlisting}
stage ID:
  new widemult = _ZX_2(widemult);
stage RR:
  new argument3 = GPR[registerYe<<1];
  new argument2 = GPR[registerZe<<1];
stage E1:
  new argument1 = PGR[registerM];
  for (I = 0; I < 8; I += 1) {
    if (widemult == 0) {
      new product = _SX_8(argument2.8[I]) * _SX_8(argument3.8[I]);
    }
    if (widemult == 1) {
      product = _ZX_8(argument2.8[I]) * _ZX_8(argument3.8[I]);
    }
    if (widemult == 2) {
      product = _SX_8(argument2.8[I]) * _ZX_8(argument3.8[I]);
    }
    result1.16[I] = argument1.16[I] + product
  };
stage E2:
  PGR[registerM] = result1;
\end{lstlisting}}
\newpage

\paragraph{Encoding} Format \textsf{ALU\_HOMEWRRRO} (Section~\ref{sec:format-ALU-HOMEWRRRO}), \verb|alucode2 = 10|\\

\bitfields{ 1/P, 2/11, 1/1, 2/10, 2/widemult, 5/registerM, 1/1, 2/10, 3/110, 1/1, 5/registerYo, 1/1, 5/registerZo, 1/0 }


\paragraph{Syntax} \texttt{maddbho}\emph{widemult} \emph{registerM} = \emph{registerZo}, \emph{registerYo}

\paragraph{Description} The \emph{registerYo} and the \emph{registerZo} are interpreted as eight bytes packed into 64 bits. These \emph{registerYo} and \emph{registerZo} bytes are sign-extended or zero-extended from 8 bits, depending on the \emph{widemult}. The \emph{registerYo} bytes are multiplied by the \emph{registerZo} bytes, producing eight 16-bit results. The eight resulting half words are packed, added to and stored into the \emph{registerM}.


\paragraph{Modifier(s)}
\noindent\begin{itemize}
\item widemult (cf. widemult in section \ref{par:modifier-widemult} on page \pageref{par:modifier-widemult})
\end{itemize}

\paragraph{Execution} {Reservation table \textsf{ALU\_LITE} (Section~\ref{sec:reservation-ALU-LITE})}
~
{\begin{lstlisting}
stage ID:
  new widemult = _ZX_2(widemult);
stage RR:
  new argument3 = GPR[(registerYo<<1)+1];
  new argument2 = GPR[(registerZo<<1)+1];
stage E1:
  new argument1 = PGR[registerM];
  for (I = 0; I < 8; I += 1) {
    if (widemult == 0) {
      new product = _SX_8(argument2.8[I]) * _SX_8(argument3.8[I]);
    }
    if (widemult == 1) {
      product = _ZX_8(argument2.8[I]) * _ZX_8(argument3.8[I]);
    }
    if (widemult == 2) {
      product = _SX_8(argument2.8[I]) * _ZX_8(argument3.8[I]);
    }
    result1.16[I] = argument1.16[I] + product
  };
stage E2:
  PGR[registerM] = result1;
\end{lstlisting}}
\newpage
\subsubsection[{\small MADDD: Multiply Add Double Word}]{MADDD\hfill Multiply Add Double Word}
\label{instr:MADDD}\index{maddd}
 
\paragraph{Encoding} Format \textsf{ALU\_DMWRRR} (Section~\ref{sec:format-ALU-DMWRRR}), \verb|alucode2 = 10|\\

\bitfields{ 1/P, 2/11, 1/0, 2/10, 2/highmult, 6/registerW, 2/10, 3/110, 1/0, 6/registerY, 6/registerZ }


\paragraph{Syntax} \texttt{maddd}\emph{highmult} \emph{registerW} = \emph{registerZ}, \emph{registerY}

\paragraph{Description} The \emph{registerY} and the \emph{registerZ} double words are sign-extended or zero-extended from 64 bits, depending on the \emph{highmult}. The \emph{registerY} double word is multiplied by the \emph{registerZ} double word, and the product is optionally shifted right by 64 bits, depending on the \emph{highmult}. The resulting double word is added to and stored into the \emph{registerW}.


\paragraph{Modifier(s)}
\noindent\begin{itemize}
\item highmult (cf. highmult in section \ref{par:modifier-highmult} on page \pageref{par:modifier-highmult})
\end{itemize}

\paragraph{Execution} {Reservation table \textsf{ALU\_LITE} (Section~\ref{sec:reservation-ALU-LITE})}
~
{\begin{lstlisting}
stage ID:
  new highmult = _ZX_2(highmult);
stage RR:
  new argument3 = GPR[registerY];
  new argument2 = GPR[registerZ];
stage E1:
  new argument1 = GPR[registerW];
  if ((highmult == 0) || (highmult == 3)) {
    new product = _SX_64(argument2) * _SX_64(argument3);
  }
  if (highmult == 1) {
    product = _ZX_64(argument2) * _ZX_64(argument3);
  }
  if (highmult == 2) {
    product = _SX_64(argument2) * _ZX_64(argument3);
  }
  if (highmult) {
    new result1 = argument1 + product.64[1];
  } else {
    result1 = argument1 + product.64[0];
  }
stage E2:
  GPR[registerW] = result1;
\end{lstlisting}}
\newpage

\paragraph{Encoding} Format \textsf{ALU\_DMWRRR.M} (Section~\ref{sec:format-ALU-DMWRRR.M}), \verb|alucode2 = 10|\\

\bitfields{ 1/P, 2/00, 2/-- --, 27/upper27, 1/1, 2/11, 1/0, 2/10, 2/highmult, 6/registerW, 2/10, 3/110, 1/0, 1/splat32, 5/lower5, 6/registerZ }


\paragraph{Syntax} \texttt{maddd}\emph{highmult} \emph{registerW} = \emph{registerZ}, \emph{upper27\_\discretionary{}{}{}lower5}\emph{splat32}

\paragraph{Description} The \emph{upper27\_\discretionary{}{}{}lower5} and the \emph{registerZ} double words are sign-extended or zero-extended from 64 bits, depending on the \emph{highmult}. The \emph{upper27\_\discretionary{}{}{}lower5} double word is multiplied by the \emph{registerZ} double word, and the product is optionally shifted right by 64 bits, depending on the \emph{highmult}. The resulting double word is added to and stored into the \emph{registerW}.


\paragraph{Modifier(s)}
\noindent\begin{itemize}
\item highmult (cf. highmult in section \ref{par:modifier-highmult} on page \pageref{par:modifier-highmult})
\item splat32 (cf. splat32 in section \ref{par:modifier-splat32} on page \pageref{par:modifier-splat32})
\end{itemize}

\paragraph{Execution} {Reservation table \textsf{ALU\_LITE.X} (Section~\ref{sec:reservation-ALU-LITE.X})}
~
{\begin{lstlisting}
stage ID:
  new splat32 = _ZX_1(splat32);
  new highmult = _ZX_2(highmult);
stage RR:
  new argument3 = splat32 ? (_WX_32(upper27_lower5) << 32) | _ZX_32(_WX_32(upper27_lower5)) : _SX_32(_WX_32(upper27_lower5));
  new argument2 = GPR[registerZ];
stage E1:
  new argument1 = GPR[registerW];
  if ((highmult == 0) || (highmult == 3)) {
    new product = _SX_64(argument2) * _SX_64(argument3);
  }
  if (highmult == 1) {
    product = _ZX_64(argument2) * _ZX_64(argument3);
  }
  if (highmult == 2) {
    product = _SX_64(argument2) * _ZX_64(argument3);
  }
  if (highmult) {
    new result1 = argument1 + product.64[1];
  } else {
    result1 = argument1 + product.64[0];
  }
stage E2:
  GPR[registerW] = result1;
\end{lstlisting}}
\newpage
\subsubsection[{\small MADDDP: Multiply Add Double Word Pair}]{MADDDP\hfill Multiply Add Double Word Pair}
\label{instr:MADDDP}\index{madddp}
 
\paragraph{Encoding} Format \textsf{ALU\_DPMWRRR} (Section~\ref{sec:format-ALU-DPMWRRR}), \verb|alucode2 = 10|\\

\bitfields{ 1/P, 2/11, 1/1, 2/10, 2/highmult, 5/registerM, 1/--, 2/10, 3/110, 1/0, 5/registerO, 1/0, 5/registerP, 1/0 }


\paragraph{Syntax} \texttt{madddp}\emph{highmult} \emph{registerM} = \emph{registerP}, \emph{registerO}

\paragraph{Description} The \emph{registerO} and the \emph{registerP} are interpreted as two double words packed into 128 bits. The \emph{registerO} and the \emph{registerP} double words are sign-extended or zero-extended from 64 bits, depending on the \emph{highmult}. The \emph{registerO} double words are multiplied by the \emph{registerP} double words, and the product is optionally shifted right by 64 bits, depending on the \emph{highmult}. The two resulting double words are added to the \emph{registerM} double words, packed and stored back into the \emph{registerM}.


\paragraph{Modifier(s)}
\noindent\begin{itemize}
\item highmult (cf. highmult in section \ref{par:modifier-highmult} on page \pageref{par:modifier-highmult})
\end{itemize}

\paragraph{Execution} {Reservation table \textsf{ALU\_LITE} (Section~\ref{sec:reservation-ALU-LITE})}
~
{\begin{lstlisting}
stage ID:
  new highmult = _ZX_2(highmult);
stage RR:
  new argument3 = PGR[registerO];
  new argument2 = PGR[registerP];
stage E1:
  new argument1 = PGR[registerM];
  for (I = 0; I < 2; I += 1) {
    if ((highmult == 0) || (highmult == 3)) {
      new product = _SX_64(argument2.64[I]) * _SX_64(argument3.64[I]);
    }
    if (highmult == 1) {
      product = _ZX_64(argument2.64[I]) * _ZX_64(argument3.64[I]);
    }
    if (highmult == 2) {
      product = _SX_64(argument2.64[I]) * _ZX_64(argument3.64[I]);
    }
    if (highmult) {
      result1.64[I] = argument1.64[I] + product.64[1];
    } else {
      result1.64[I] = argument1.64[I] + product.64[0];
    }
  };
stage E2:
  PGR[registerM] = result1;
\end{lstlisting}}
\newpage
\subsubsection[{\small MADDDQ: Multiply Add Extend Double Word to Quadruple Word}]{MADDDQ\hfill Multiply Add Extend Double Word to Quadruple Word}
\label{instr:MADDDQ}\index{madddq}
 
\paragraph{Encoding} Format \textsf{ALU\_QMEWRRR} (Section~\ref{sec:format-ALU-QMEWRRR}), \verb|alucode2 = 10|\\

\bitfields{ 1/P, 2/11, 1/0, 2/10, 2/widemult, 5/registerM, 1/--, 2/10, 3/111, 1/1, 6/registerY, 6/registerZ }


\paragraph{Syntax} \texttt{madddq}\emph{widemult} \emph{registerM} = \emph{registerZ}, \emph{registerY}

\paragraph{Description} The \emph{registerY} and the \emph{registerZ} double words are sign-extended or zero-extended from 64 bits, depending on the \emph{widemult}. The \emph{registerY} is multiplied by the \emph{registerZ}, producing a 128-bit result. The resulting quadruple word is added to and stored into the \emph{registerM}.


\paragraph{Modifier(s)}
\noindent\begin{itemize}
\item widemult (cf. widemult in section \ref{par:modifier-widemult} on page \pageref{par:modifier-widemult})
\end{itemize}

\paragraph{Execution} {Reservation table \textsf{ALU\_LITE} (Section~\ref{sec:reservation-ALU-LITE})}
~
{\begin{lstlisting}
stage ID:
  new widemult = _ZX_2(widemult);
stage RR:
  new argument3 = GPR[registerY];
  new argument2 = GPR[registerZ];
stage E1:
  new argument1 = PGR[registerM];
  if (widemult == 0) {
    new product = _SX_64(argument2) * _SX_64(argument3);
  }
  if (widemult == 1) {
    product = _ZX_64(argument2) * _ZX_64(argument3);
  }
  if (widemult == 2) {
    product = _SX_64(argument2) * _ZX_64(argument3);
  }
  new result1 = argument1 + product;
stage E2:
  PGR[registerM] = result1;
\end{lstlisting}}
\newpage
\subsubsection[{\small MADDHWQ: Multiply Add Extend Half Word to Word Quadruple}]{MADDHWQ\hfill Multiply Add Extend Half Word to Word Quadruple}
\label{instr:MADDHWQ}\index{maddhwq}
 
\paragraph{Encoding} Format \textsf{ALU\_WQMEWRRRE} (Section~\ref{sec:format-ALU-WQMEWRRRE}), \verb|alucode2 = 10|\\

\bitfields{ 1/P, 2/11, 1/1, 2/10, 2/widemult, 5/registerM, 1/0, 2/10, 3/110, 1/0, 5/registerYe, 1/1, 5/registerZe, 1/1 }


\paragraph{Syntax} \texttt{maddhwq}\emph{widemult} \emph{registerM} = \emph{registerZe}, \emph{registerYe}

\paragraph{Description} The \emph{registerYe} and the \emph{registerZe} are interpreted as four half words packed into 64 bits. These \emph{registerYe} and \emph{registerZe} half words are sign-extended or zero-extended from 16 bits, depending on the \emph{widemult}. The \emph{registerYe} half words are multiplied by the \emph{registerZe} half words, producing four 32-bit results. The four resulting words are packed, added to and stored into the \emph{registerM}.


\paragraph{Modifier(s)}
\noindent\begin{itemize}
\item widemult (cf. widemult in section \ref{par:modifier-widemult} on page \pageref{par:modifier-widemult})
\end{itemize}

\paragraph{Execution} {Reservation table \textsf{ALU\_LITE} (Section~\ref{sec:reservation-ALU-LITE})}
~
{\begin{lstlisting}
stage ID:
  new widemult = _ZX_2(widemult);
stage RR:
  new argument3 = GPR[registerYe<<1];
  new argument2 = GPR[registerZe<<1];
stage E1:
  new argument1 = PGR[registerM];
  for (I = 0; I < 4; I += 1) {
    if (widemult == 0) {
      new product = _SX_16(argument2.16[I]) * _SX_16(argument3.16[I]);
    }
    if (widemult == 1) {
      product = _ZX_16(argument2.16[I]) * _ZX_16(argument3.16[I]);
    }
    if (widemult == 2) {
      product = _SX_16(argument2.16[I]) * _ZX_16(argument3.16[I]);
    }
    result1.32[I] = argument1.32[I] + product
  };
stage E2:
  PGR[registerM] = result1;
\end{lstlisting}}
\newpage

\paragraph{Encoding} Format \textsf{ALU\_WQMEWRRRO} (Section~\ref{sec:format-ALU-WQMEWRRRO}), \verb|alucode2 = 10|\\

\bitfields{ 1/P, 2/11, 1/1, 2/10, 2/widemult, 5/registerM, 1/1, 2/10, 3/110, 1/0, 5/registerYo, 1/1, 5/registerZo, 1/1 }


\paragraph{Syntax} \texttt{maddhwq}\emph{widemult} \emph{registerM} = \emph{registerZo}, \emph{registerYo}

\paragraph{Description} The \emph{registerYo} and the \emph{registerZo} are interpreted as four half words packed into 64 bits. These \emph{registerYo} and \emph{registerZo} half words are sign-extended or zero-extended from 16 bits, depending on the \emph{widemult}. The \emph{registerYo} half words are multiplied by the \emph{registerZo} half words, producing four 32-bit results. The four resulting words are packed, added to and stored into the \emph{registerM}.


\paragraph{Modifier(s)}
\noindent\begin{itemize}
\item widemult (cf. widemult in section \ref{par:modifier-widemult} on page \pageref{par:modifier-widemult})
\end{itemize}

\paragraph{Execution} {Reservation table \textsf{ALU\_LITE} (Section~\ref{sec:reservation-ALU-LITE})}
~
{\begin{lstlisting}
stage ID:
  new widemult = _ZX_2(widemult);
stage RR:
  new argument3 = GPR[(registerYo<<1)+1];
  new argument2 = GPR[(registerZo<<1)+1];
stage E1:
  new argument1 = PGR[registerM];
  for (I = 0; I < 4; I += 1) {
    if (widemult == 0) {
      new product = _SX_16(argument2.16[I]) * _SX_16(argument3.16[I]);
    }
    if (widemult == 1) {
      product = _ZX_16(argument2.16[I]) * _ZX_16(argument3.16[I]);
    }
    if (widemult == 2) {
      product = _SX_16(argument2.16[I]) * _ZX_16(argument3.16[I]);
    }
    result1.32[I] = argument1.32[I] + product
  };
stage E2:
  PGR[registerM] = result1;
\end{lstlisting}}
\newpage
\subsubsection[{\small MADDW: Multiply Add Word}]{MADDW\hfill Multiply Add Word}
\label{instr:MADDW}\index{maddw}
 
\paragraph{Encoding} Format \textsf{ALU\_WMWRRR} (Section~\ref{sec:format-ALU-WMWRRR}), \verb|alucode2 = 10|\\

\bitfields{ 1/P, 2/11, 1/0, 2/10, 2/highmult, 6/registerW, 2/11, 3/110, 1/signextw, 6/registerY, 6/registerZ }


\paragraph{Syntax} \texttt{maddw}\emph{highmult}\emph{signextw} \emph{registerW} = \emph{registerZ}, \emph{registerY}

\paragraph{Description} The \emph{registerY} and the \emph{registerZ} words are sign-extended or zero-extended from 32 bits, depending on the \emph{highmult}. The \emph{registerY} word is multiplied by the \emph{registerZ} word, and the product is optionally shifted right by 32 bits, depending on the \emph{highmult}. The resulting word is added to and stored into the \emph{registerW}.


\paragraph{Modifier(s)}
\noindent\begin{itemize}
\item highmult (cf. highmult in section \ref{par:modifier-highmult} on page \pageref{par:modifier-highmult})
\item signextw (cf. signextw in section \ref{par:modifier-signextw} on page \pageref{par:modifier-signextw})
\end{itemize}

\paragraph{Execution} {Reservation table \textsf{ALU\_LITE} (Section~\ref{sec:reservation-ALU-LITE})}
~
{\begin{lstlisting}
stage ID:
  new signextw = _ZX_1(signextw);
  new highmult = _ZX_2(highmult);
stage RR:
  new argument3 = GPR[registerY];
  new argument2 = GPR[registerZ];
stage E1:
  new argument1 = GPR[registerW];
  if ((highmult == 0) || (highmult == 3)) {
    new product = _SX_32(argument2) * _SX_32(argument3);
  }
  if (highmult == 1) {
    product = _ZX_32(argument2) * _ZX_32(argument3);
  }
  if (highmult == 2) {
    product = _SX_32(argument2) * _ZX_32(argument3);
  }
  if (highmult) {
    new result1 = argument1 + product.32[1];
  } else {
    result1 = argument1 + product.32[0];
  }
stage E2:
  GPR[registerW] = signextw ? _SX_32(result1) : _ZX_32(result1);
\end{lstlisting}}
\newpage

\paragraph{Encoding} Format \textsf{ALU\_WMWRRR.W} (Section~\ref{sec:format-ALU-WMWRRR.W}), \verb|alucode2 = 10|\\

\bitfields{ 1/P, 2/00, 2/-- --, 27/upper27, 1/1, 2/11, 1/0, 2/10, 2/highmult, 6/registerW, 2/11, 3/110, 1/signextw, 1/0, 5/lower5, 6/registerZ }


\paragraph{Syntax} \texttt{maddw}\emph{highmult}\emph{signextw} \emph{registerW} = \emph{registerZ}, \emph{upper27\_\discretionary{}{}{}lower5}

\paragraph{Description} The \emph{upper27\_\discretionary{}{}{}lower5} and the \emph{registerZ} words are sign-extended or zero-extended from 32 bits, depending on the \emph{highmult}. The \emph{upper27\_\discretionary{}{}{}lower5} word is multiplied by the \emph{registerZ} word, and the product is optionally shifted right by 32 bits, depending on the \emph{highmult}. The resulting word is added to and stored into the \emph{registerW}.


\paragraph{Modifier(s)}
\noindent\begin{itemize}
\item highmult (cf. highmult in section \ref{par:modifier-highmult} on page \pageref{par:modifier-highmult})
\item signextw (cf. signextw in section \ref{par:modifier-signextw} on page \pageref{par:modifier-signextw})
\end{itemize}

\paragraph{Execution} {Reservation table \textsf{ALU\_LITE.X} (Section~\ref{sec:reservation-ALU-LITE.X})}
~
{\begin{lstlisting}
stage ID:
  new signextw = _ZX_1(signextw);
  new highmult = _ZX_2(highmult);
stage RR:
  new argument3 = _WX_32(upper27_lower5);
  new argument2 = GPR[registerZ];
stage E1:
  new argument1 = GPR[registerW];
  if ((highmult == 0) || (highmult == 3)) {
    new product = _SX_32(argument2) * _SX_32(argument3);
  }
  if (highmult == 1) {
    product = _ZX_32(argument2) * _ZX_32(argument3);
  }
  if (highmult == 2) {
    product = _SX_32(argument2) * _ZX_32(argument3);
  }
  if (highmult) {
    new result1 = argument1 + product.32[1];
  } else {
    result1 = argument1 + product.32[0];
  }
stage E2:
  GPR[registerW] = signextw ? _SX_32(result1) : _ZX_32(result1);
\end{lstlisting}}
\newpage
\subsubsection[{\small MADDWD: Multiply Add Extend Word to Double Word}]{MADDWD\hfill Multiply Add Extend Word to Double Word}
\label{instr:MADDWD}\index{maddwd}
 
\paragraph{Encoding} Format \textsf{ALU\_DMEWRRR} (Section~\ref{sec:format-ALU-DMEWRRR}), \verb|alucode2 = 10|\\

\bitfields{ 1/P, 2/11, 1/0, 2/10, 2/widemult, 6/registerW, 2/10, 3/111, 1/0, 6/registerY, 6/registerZ }


\paragraph{Syntax} \texttt{maddwd}\emph{widemult} \emph{registerW} = \emph{registerZ}, \emph{registerY}

\paragraph{Description} The \emph{registerY} and the \emph{registerZ} words are sign-extended or zero-extended from 32 bits, depending on the \emph{widemult}. The \emph{registerY} is multiplied by the \emph{registerZ}, producing a 64-bit result. The resulting double word is added to and stored into the \emph{registerW}.


\paragraph{Modifier(s)}
\noindent\begin{itemize}
\item widemult (cf. widemult in section \ref{par:modifier-widemult} on page \pageref{par:modifier-widemult})
\end{itemize}

\paragraph{Execution} {Reservation table \textsf{ALU\_LITE} (Section~\ref{sec:reservation-ALU-LITE})}
~
{\begin{lstlisting}
stage ID:
  new widemult = _ZX_2(widemult);
stage RR:
  new argument3 = GPR[registerY];
  new argument2 = GPR[registerZ];
stage E1:
  new argument1 = GPR[registerW];
  if (widemult == 0) {
    new product = _SX_32(argument2) * _SX_32(argument3);
  }
  if (widemult == 1) {
    product = _ZX_32(argument2) * _ZX_32(argument3);
  }
  if (widemult == 2) {
    product = _SX_32(argument2) * _ZX_32(argument3);
  }
  new result1 = argument1 + product;
stage E2:
  GPR[registerW] = result1;
\end{lstlisting}}
\newpage

\paragraph{Encoding} Format \textsf{ALU\_DMEWRRR.M} (Section~\ref{sec:format-ALU-DMEWRRR.M}), \verb|alucode2 = 10|\\

\bitfields{ 1/P, 2/00, 2/-- --, 27/upper27, 1/1, 2/11, 1/0, 2/10, 2/widemult, 6/registerW, 2/10, 3/111, 1/0, 1/splat32, 5/lower5, 6/registerZ }


\paragraph{Syntax} \texttt{maddwd}\emph{widemult} \emph{registerW} = \emph{registerZ}, \emph{upper27\_\discretionary{}{}{}lower5}\emph{splat32}

\paragraph{Description} The \emph{upper27\_\discretionary{}{}{}lower5} and the \emph{registerZ} words are sign-extended or zero-extended from 32 bits, depending on the \emph{widemult}. The \emph{upper27\_\discretionary{}{}{}lower5} is multiplied by the \emph{registerZ}, producing a 64-bit result. The resulting double word is added to and stored into the \emph{registerW}.


\paragraph{Modifier(s)}
\noindent\begin{itemize}
\item widemult (cf. widemult in section \ref{par:modifier-widemult} on page \pageref{par:modifier-widemult})
\item splat32 (cf. splat32 in section \ref{par:modifier-splat32} on page \pageref{par:modifier-splat32})
\end{itemize}

\paragraph{Execution} {Reservation table \textsf{ALU\_LITE.X} (Section~\ref{sec:reservation-ALU-LITE.X})}
~
{\begin{lstlisting}
stage ID:
  new splat32 = _ZX_1(splat32);
  new widemult = _ZX_2(widemult);
stage RR:
  new argument3 = splat32 ? (_WX_32(upper27_lower5) << 32) | _ZX_32(_WX_32(upper27_lower5)) : _SX_32(_WX_32(upper27_lower5));
  new argument2 = GPR[registerZ];
stage E1:
  new argument1 = GPR[registerW];
  if (widemult == 0) {
    new product = _SX_32(argument2) * _SX_32(argument3);
  }
  if (widemult == 1) {
    product = _ZX_32(argument2) * _ZX_32(argument3);
  }
  if (widemult == 2) {
    product = _SX_32(argument2) * _ZX_32(argument3);
  }
  new result1 = argument1 + product;
stage E2:
  GPR[registerW] = result1;
\end{lstlisting}}
\newpage
\subsubsection[{\small MADDWDP: Multiply Add Extend Word to Double Word Pair}]{MADDWDP\hfill Multiply Add Extend Word to Double Word Pair}
\label{instr:MADDWDP}\index{maddwdp}
 
\paragraph{Encoding} Format \textsf{ALU\_DPMEWRRRE} (Section~\ref{sec:format-ALU-DPMEWRRRE}), \verb|alucode2 = 10|\\

\bitfields{ 1/P, 2/11, 1/1, 2/10, 2/widemult, 5/registerM, 1/0, 2/10, 3/110, 1/0, 5/registerYe, 1/1, 5/registerZe, 1/0 }


\paragraph{Syntax} \texttt{maddwdp}\emph{widemult} \emph{registerM} = \emph{registerZe}, \emph{registerYe}

\paragraph{Description} The \emph{registerYe} and the \emph{registerZe} are interpreted as two words packed into 64 bits. These \emph{registerYe} and \emph{registerZe} words are sign-extended or zero-extended from 32 bits, depending on the \emph{widemult}. The \emph{registerYe} words are multiplied by the \emph{registerZe} words, producing two 64-bit results. The two resulting double words are packed, added to and stored into the \emph{registerM}.


\paragraph{Modifier(s)}
\noindent\begin{itemize}
\item widemult (cf. widemult in section \ref{par:modifier-widemult} on page \pageref{par:modifier-widemult})
\end{itemize}

\paragraph{Execution} {Reservation table \textsf{ALU\_LITE} (Section~\ref{sec:reservation-ALU-LITE})}
~
{\begin{lstlisting}
stage ID:
  new widemult = _ZX_2(widemult);
stage RR:
  new argument3 = GPR[registerYe<<1];
  new argument2 = GPR[registerZe<<1];
stage E1:
  new argument1 = PGR[registerM];
  for (I = 0; I < 2; I += 1) {
    if (widemult == 0) {
      new product = _SX_32(argument2.32[I]) * _SX_32(argument3.32[I]);
    }
    if (widemult == 1) {
      product = _ZX_32(argument2.32[I]) * _ZX_32(argument3.32[I]);
    }
    if (widemult == 2) {
      product = _SX_32(argument2.32[I]) * _ZX_32(argument3.32[I]);
    }
    result1.64[I] = argument1.64[I] + product
  };
stage E2:
  PGR[registerM] = result1;
\end{lstlisting}}
\newpage

\paragraph{Encoding} Format \textsf{ALU\_DPMEWRRRO} (Section~\ref{sec:format-ALU-DPMEWRRRO}), \verb|alucode2 = 10|\\

\bitfields{ 1/P, 2/11, 1/1, 2/10, 2/widemult, 5/registerM, 1/1, 2/10, 3/110, 1/0, 5/registerYo, 1/1, 5/registerZo, 1/0 }


\paragraph{Syntax} \texttt{maddwdp}\emph{widemult} \emph{registerM} = \emph{registerZo}, \emph{registerYo}

\paragraph{Description} The \emph{registerYo} and the \emph{registerZo} are interpreted as two words packed into 64 bits. These \emph{registerYo} and \emph{registerZo} words are sign-extended or zero-extended from 32 bits, depending on the \emph{widemult}. The \emph{registerYo} words are multiplied by the \emph{registerZo} words, producing two 64-bit results. The two resulting double words are packed, added to and stored into the \emph{registerM}.


\paragraph{Modifier(s)}
\noindent\begin{itemize}
\item widemult (cf. widemult in section \ref{par:modifier-widemult} on page \pageref{par:modifier-widemult})
\end{itemize}

\paragraph{Execution} {Reservation table \textsf{ALU\_LITE} (Section~\ref{sec:reservation-ALU-LITE})}
~
{\begin{lstlisting}
stage ID:
  new widemult = _ZX_2(widemult);
stage RR:
  new argument3 = GPR[(registerYo<<1)+1];
  new argument2 = GPR[(registerZo<<1)+1];
stage E1:
  new argument1 = PGR[registerM];
  for (I = 0; I < 2; I += 1) {
    if (widemult == 0) {
      new product = _SX_32(argument2.32[I]) * _SX_32(argument3.32[I]);
    }
    if (widemult == 1) {
      product = _ZX_32(argument2.32[I]) * _ZX_32(argument3.32[I]);
    }
    if (widemult == 2) {
      product = _SX_32(argument2.32[I]) * _ZX_32(argument3.32[I]);
    }
    result1.64[I] = argument1.64[I] + product
  };
stage E2:
  PGR[registerM] = result1;
\end{lstlisting}}
\newpage
\subsubsection[{\small MADDXBHO: Multiply Add Extend Byte to Half Word Octuple}]{MADDXBHO\hfill Multiply Add Extend Byte to Half Word Octuple}
\label{instr:MADDXBHO}\index{maddxbho}
 
\paragraph{Encoding} Format \textsf{ALU\_HOMXWRRR} (Section~\ref{sec:format-ALU-HOMXWRRR}), \verb|alucode2 = 10|\\

\bitfields{ 1/P, 2/11, 1/1, 2/10, 2/widemult, 5/registerM, 1/oddlanes, 2/10, 3/111, 1/1, 5/registerO, 1/--, 5/registerP, 1/0 }


\paragraph{Syntax} \texttt{maddxbho}\emph{oddlanes}\emph{widemult} \emph{registerM} = \emph{registerP}, \emph{registerO}

\paragraph{Description} The \emph{registerO} and the \emph{registerP} are interpreted as sixteen bytes packed into 128 bits. The even or odd bytes of the \emph{registerO} and the \emph{registerP} are selected, depending on \emph{oddlanes}. These \emph{registerO} and \emph{registerP} bytes are sign-extended or zero-extended from 8 bits, depending on the \emph{widemult}. The \emph{registerO} bytes are multiplied by the \emph{registerP} bytes, producing eight 16-bit results. The eight resulting half words are packed, added to and stored into the \emph{registerM}.


\paragraph{Modifier(s)}
\noindent\begin{itemize}
\item widemult (cf. widemult in section \ref{par:modifier-widemult} on page \pageref{par:modifier-widemult})
\item oddlanes (cf. oddlanes in section \ref{par:modifier-oddlanes} on page \pageref{par:modifier-oddlanes})
\end{itemize}

\paragraph{Execution} {Reservation table \textsf{ALU\_LITE} (Section~\ref{sec:reservation-ALU-LITE})}
~
{\begin{lstlisting}
stage ID:
  new oddlanes = _ZX_1(oddlanes);
  new widemult = _ZX_2(widemult);
stage RR:
  new argument3 = PGR[registerO];
  new argument2 = PGR[registerP];
stage E1:
  new argument1 = PGR[registerM];
  if (oddlanes) {
    new argument2 = argument2 >> 8;
    new argument3 = argument3 >> 8;
  }
  for (I = 0; I < 8; I += 1) {
    if (widemult == 0) {
      new product = _SX_8(argument2.16[I]) * _SX_8(argument3.16[I]);
    }
    if (widemult == 1) {
      product = _ZX_8(argument2.16[I]) * _ZX_8(argument3.16[I]);
    }
    if (widemult == 2) {
      product = _SX_8(argument2.16[I]) * _ZX_8(argument3.16[I]);
    }
    result1.16[I] = argument1.16[I] + product
  };
stage E2:
  PGR[registerM] = result1;
\end{lstlisting}}
\newpage
\subsubsection[{\small MADDXHWQ: Multiply Add Extend Half Word to Word Quadruple}]{MADDXHWQ\hfill Multiply Add Extend Half Word to Word Quadruple}
\label{instr:MADDXHWQ}\index{maddxhwq}
 
\paragraph{Encoding} Format \textsf{ALU\_WQMXWRRR} (Section~\ref{sec:format-ALU-WQMXWRRR}), \verb|alucode2 = 10|\\

\bitfields{ 1/P, 2/11, 1/1, 2/10, 2/widemult, 5/registerM, 1/oddlanes, 2/10, 3/111, 1/0, 5/registerO, 1/--, 5/registerP, 1/1 }


\paragraph{Syntax} \texttt{maddxhwq}\emph{oddlanes}\emph{widemult} \emph{registerM} = \emph{registerP}, \emph{registerO}

\paragraph{Description} The \emph{registerO} and the \emph{registerP} are interpreted as eight half words packed into 128 bits. The even or odd half words of the \emph{registerO} and the \emph{registerP} are selected, depending on \emph{oddlanes}. These \emph{registerO} and \emph{registerP} half words are sign-extended or zero-extended from 16 bits, depending on the \emph{widemult}. The \emph{registerO} half words are multiplied by the \emph{registerP} half words, producing four 32-bit results. The four resulting words are packed, added to and stored into the \emph{registerM}.


\paragraph{Modifier(s)}
\noindent\begin{itemize}
\item widemult (cf. widemult in section \ref{par:modifier-widemult} on page \pageref{par:modifier-widemult})
\item oddlanes (cf. oddlanes in section \ref{par:modifier-oddlanes} on page \pageref{par:modifier-oddlanes})
\end{itemize}

\paragraph{Execution} {Reservation table \textsf{ALU\_LITE} (Section~\ref{sec:reservation-ALU-LITE})}
~
{\begin{lstlisting}
stage ID:
  new oddlanes = _ZX_1(oddlanes);
  new widemult = _ZX_2(widemult);
stage RR:
  new argument3 = PGR[registerO];
  new argument2 = PGR[registerP];
stage E1:
  new argument1 = PGR[registerM];
  if (oddlanes) {
    new argument2 = argument2 >> 16;
    new argument3 = argument3 >> 16;
  }
  for (I = 0; I < 4; I += 1) {
    if (widemult == 0) {
      new product = _SX_16(argument2.32[I]) * _SX_16(argument3.32[I]);
    }
    if (widemult == 1) {
      product = _ZX_16(argument2.32[I]) * _ZX_16(argument3.32[I]);
    }
    if (widemult == 2) {
      product = _SX_16(argument2.32[I]) * _ZX_16(argument3.32[I]);
    }
    result1.32[I] = argument1.32[I] + product
  };
stage E2:
  PGR[registerM] = result1;
\end{lstlisting}}
\newpage
\subsubsection[{\small MADDXWDP: Multiply Add Extend Word to Double Word Pair}]{MADDXWDP\hfill Multiply Add Extend Word to Double Word Pair}
\label{instr:MADDXWDP}\index{maddxwdp}
 
\paragraph{Encoding} Format \textsf{ALU\_DPMXWRRR} (Section~\ref{sec:format-ALU-DPMXWRRR}), \verb|alucode2 = 10|\\

\bitfields{ 1/P, 2/11, 1/1, 2/10, 2/widemult, 5/registerM, 1/oddlanes, 2/10, 3/111, 1/0, 5/registerO, 1/--, 5/registerP, 1/0 }


\paragraph{Syntax} \texttt{maddxwdp}\emph{oddlanes}\emph{widemult} \emph{registerM} = \emph{registerP}, \emph{registerO}

\paragraph{Description} The \emph{registerO} and the \emph{registerP} are interpreted as four words packed into 128 bits. The even or odd words of the \emph{registerO} and the \emph{registerP} are selected, depending on \emph{oddlanes}. These \emph{registerO} and \emph{registerP} words are sign-extended or zero-extended from 32 bits, depending on the \emph{widemult}. The \emph{registerO} words are multiplied by the \emph{registerP} words, producing two 64-bit results. The two resulting double words are packed, added to and stored into the \emph{registerM}.


\paragraph{Modifier(s)}
\noindent\begin{itemize}
\item widemult (cf. widemult in section \ref{par:modifier-widemult} on page \pageref{par:modifier-widemult})
\item oddlanes (cf. oddlanes in section \ref{par:modifier-oddlanes} on page \pageref{par:modifier-oddlanes})
\end{itemize}

\paragraph{Execution} {Reservation table \textsf{ALU\_LITE} (Section~\ref{sec:reservation-ALU-LITE})}
~
{\begin{lstlisting}
stage ID:
  new oddlanes = _ZX_1(oddlanes);
  new widemult = _ZX_2(widemult);
stage RR:
  new argument3 = PGR[registerO];
  new argument2 = PGR[registerP];
stage E1:
  new argument1 = PGR[registerM];
  if (oddlanes) {
    new argument2 = argument2 >> 32;
    new argument3 = argument3 >> 32;
  }
  for (I = 0; I < 2; I += 1) {
    if (widemult == 0) {
      new product = _SX_32(argument2.64[I]) * _SX_32(argument3.64[I]);
    }
    if (widemult == 1) {
      product = _ZX_32(argument2.64[I]) * _ZX_32(argument3.64[I]);
    }
    if (widemult == 2) {
      product = _SX_32(argument2.64[I]) * _ZX_32(argument3.64[I]);
    }
    result1.64[I] = argument1.64[I] + product
  };
stage E2:
  PGR[registerM] = result1;
\end{lstlisting}}
\newpage
\subsubsection[{\small MAKE: Make Word from Immediate}]{MAKE\hfill Make Word from Immediate}
\label{instr:MAKE}\index{make}
 
\paragraph{Encoding} Format \textsf{ALU\_DWI} (Section~\ref{sec:format-ALU-DWI}), \verb|exucode2 = 0000|\\

\bitfields{ 1/P, 2/11, 1/0, 4/0000, 6/registerW, 2/00, 16/signed16 }


\paragraph{Syntax} \texttt{make} \emph{registerW} = \emph{signed16}

\paragraph{Description} The \emph{signed16} is stored into the \emph{registerW}.


\paragraph{Execution} {Reservation table \textsf{ALU\_TINY} (Section~\ref{sec:reservation-ALU-TINY})}
~
{\begin{lstlisting}
stage RR:
  new argument2 = _SX_16(signed16);
  new result1 = _SX_64(argument2);
stage E1:
  GPR[registerW] = result1;
\end{lstlisting}}
\newpage

\paragraph{Encoding} Format \textsf{ALU\_DWI.X} (Section~\ref{sec:format-ALU-DWI.X}), \verb|exucode2 = 0000|\\

\bitfields{ 1/P, 2/00, 2/-- --, 27/upper27, 1/1, 2/11, 1/0, 4/0000, 6/registerW, 2/00, 10/lower10, 6/extend6 }


\paragraph{Syntax} \texttt{make} \emph{registerW} = \emph{extend6\_\discretionary{}{}{}upper27\_\discretionary{}{}{}lower10}

\paragraph{Description} The \emph{extend6\_\discretionary{}{}{}upper27\_\discretionary{}{}{}lower10} is stored into the \emph{registerW}.


\paragraph{Execution} {Reservation table \textsf{ALU\_TINY.X} (Section~\ref{sec:reservation-ALU-TINY.X})}
~
{\begin{lstlisting}
stage RR:
  new argument2 = _SX_43(extend6_upper27_lower10);
  new result1 = _SX_64(argument2);
stage E1:
  GPR[registerW] = result1;
\end{lstlisting}}
\newpage

\paragraph{Encoding} Format \textsf{ALU\_DWI.Y} (Section~\ref{sec:format-ALU-DWI.Y}), \verb|exucode2 = 0000|\\

\bitfields{ 1/P, 2/00, 2/-- --, 27/extend27, 1/1, 2/00, 2/-- --, 27/upper27, 1/1, 2/11, 1/0, 4/0000, 6/registerW, 2/00, 10/lower10, 6/-- -- -- -- -- -- }


\paragraph{Syntax} \texttt{make} \emph{registerW} = \emph{extend27\_\discretionary{}{}{}upper27\_\discretionary{}{}{}lower10}

\paragraph{Description} The \emph{extend27\_\discretionary{}{}{}upper27\_\discretionary{}{}{}lower10} is stored into the \emph{registerW}.


\paragraph{Execution} {Reservation table \textsf{ALU\_TINY.Y} (Section~\ref{sec:reservation-ALU-TINY.Y})}
~
{\begin{lstlisting}
stage RR:
  new argument2 = _WX_64(extend27_upper27_lower10);
  new result1 = _SX_64(argument2);
stage E1:
  GPR[registerW] = result1;
\end{lstlisting}}
\newpage
\subsubsection[{\small MAXBX: Maximum Byte Hexadecuple}]{MAXBX\hfill Maximum Byte Hexadecuple}
\label{instr:MAXBX}\index{maxbx}
 
\paragraph{Encoding} Format \textsf{ALU\_BXWRR0} (Section~\ref{sec:format-ALU-BXWRR0}), \verb|exucode2 = 0101|\\

\bitfields{ 1/P, 2/11, 1/1, 4/0101, 5/registerM, 1/--, 2/10, 3/000, 1/1, 5/registerO, 1/--, 5/registerP, 1/1 }


\paragraph{Syntax} \texttt{maxbx} \emph{registerM} = \emph{registerP}, \emph{registerO}

\paragraph{Description} The \emph{registerO} and the \emph{registerP} are interpreted as sixteen bytes packed into 128 bits. The maximum of the \emph{registerO} bytes and the \emph{registerP} bytes are computed. The sixteen resulting bytes are packed and stored into the \emph{registerM}.


\paragraph{Execution} {Reservation table \textsf{ALU\_LITE} (Section~\ref{sec:reservation-ALU-LITE})}
~
{\begin{lstlisting}
stage RR:
  new argument3 = PGR[registerO];
  new argument2 = PGR[registerP];
  for (I = 0; I < 16; I += 1) {
    result1.8[I] = _MAX(_SX_8(argument3.8[I]), _SX_8(argument2.8[I]))
  };
stage E1:
  PGR[registerM] = result1;
\end{lstlisting}}
\newpage

\paragraph{Encoding} Format \textsf{ALU\_BXWRR0.M} (Section~\ref{sec:format-ALU-BXWRR0.M}), \verb|exucode2 = 0101|\\

\bitfields{ 1/P, 2/00, 2/-- --, 27/upper27, 1/1, 2/11, 1/1, 4/0101, 5/registerM, 1/--, 2/10, 3/000, 1/1, 1/splat32, 5/lower5, 5/registerP, 1/1 }


\paragraph{Syntax} \texttt{maxbx} \emph{registerM} = \emph{registerP}, \emph{upper27\_\discretionary{}{}{}lower5}\emph{splat32}

\paragraph{Description} The \emph{upper27\_\discretionary{}{}{}lower5} and the \emph{registerP} are interpreted as sixteen bytes packed into 128 bits. The maximum of the \emph{upper27\_\discretionary{}{}{}lower5} bytes and the \emph{registerP} bytes are computed. The sixteen resulting bytes are packed and stored into the \emph{registerM}.


\paragraph{Modifier(s)}
\noindent\begin{itemize}
\item splat32 (cf. splat32 in section \ref{par:modifier-splat32} on page \pageref{par:modifier-splat32})
\end{itemize}

\paragraph{Execution} {Reservation table \textsf{ALU\_LITE.X} (Section~\ref{sec:reservation-ALU-LITE.X})}
~
{\begin{lstlisting}
stage ID:
  new splat32 = _ZX_1(splat32);
stage RR:
  new argument3 = splat32 ? (_WX_32(upper27_lower5) << 32) | _ZX_32(_WX_32(upper27_lower5)) : _SX_32(_WX_32(upper27_lower5));
  new argument2 = PGR[registerP];
  for (I = 0; I < 16; I += 1) {
    result1.8[I] = _MAX(_SX_8(argument3.8[I]), _SX_8(argument2.8[I]))
  };
stage E1:
  PGR[registerM] = result1;
\end{lstlisting}}
\newpage
\subsubsection[{\small MAXD: Maximum Double Word}]{MAXD\hfill Maximum Double Word}
\label{instr:MAXD}\index{maxd}
 
\paragraph{Encoding} Format \textsf{ALU\_DWRR0} (Section~\ref{sec:format-ALU-DWRR0}), \verb|exucode2 = 0101|\\

\bitfields{ 1/P, 2/11, 1/0, 4/0101, 6/registerW, 2/10, 3/000, 1/0, 6/registerY, 6/registerZ }


\paragraph{Syntax} \texttt{maxd} \emph{registerW} = \emph{registerZ}, \emph{registerY}

\paragraph{Description} The maximum of the \emph{registerY} and the \emph{registerZ} is computed. The result is stored into the \emph{registerW}.


\paragraph{Execution} {Reservation table \textsf{ALU\_TINY} (Section~\ref{sec:reservation-ALU-TINY})}
~
{\begin{lstlisting}
stage RR:
  new argument3 = GPR[registerY];
  new argument2 = GPR[registerZ];
  new result1 = _MAX(_SX_64(argument3), _SX_64(argument2));
stage E1:
  GPR[registerW] = result1;
\end{lstlisting}}
\newpage

\paragraph{Encoding} Format \textsf{ALU\_DWRR0.M} (Section~\ref{sec:format-ALU-DWRR0.M}), \verb|exucode2 = 0101|\\

\bitfields{ 1/P, 2/00, 2/-- --, 27/upper27, 1/1, 2/11, 1/0, 4/0101, 6/registerW, 2/10, 3/000, 1/0, 1/splat32, 5/lower5, 6/registerZ }


\paragraph{Syntax} \texttt{maxd} \emph{registerW} = \emph{registerZ}, \emph{upper27\_\discretionary{}{}{}lower5}\emph{splat32}

\paragraph{Description} The maximum of the \emph{upper27\_\discretionary{}{}{}lower5} and the \emph{registerZ} is computed. The result is stored into the \emph{registerW}.


\paragraph{Modifier(s)}
\noindent\begin{itemize}
\item splat32 (cf. splat32 in section \ref{par:modifier-splat32} on page \pageref{par:modifier-splat32})
\end{itemize}

\paragraph{Execution} {Reservation table \textsf{ALU\_TINY.X} (Section~\ref{sec:reservation-ALU-TINY.X})}
~
{\begin{lstlisting}
stage ID:
  new splat32 = _ZX_1(splat32);
stage RR:
  new argument3 = splat32 ? (_WX_32(upper27_lower5) << 32) | _ZX_32(_WX_32(upper27_lower5)) : _SX_32(_WX_32(upper27_lower5));
  new argument2 = GPR[registerZ];
  new result1 = _MAX(_SX_64(argument3), _SX_64(argument2));
stage E1:
  GPR[registerW] = result1;
\end{lstlisting}}
\newpage

\paragraph{Encoding} Format \textsf{ALU\_DWRI} (Section~\ref{sec:format-ALU-DWRI}), \verb|exucode2 = 0101|\\

\bitfields{ 1/P, 2/11, 1/0, 4/0101, 6/registerW, 2/00, 10/signed10, 6/registerZ }


\paragraph{Syntax} \texttt{maxd} \emph{registerW} = \emph{registerZ}, \emph{signed10}

\paragraph{Description} The maximum of the \emph{signed10} and the \emph{registerZ} is computed. The result is stored into the \emph{registerW}.


\paragraph{Execution} {Reservation table \textsf{ALU\_TINY} (Section~\ref{sec:reservation-ALU-TINY})}
~
{\begin{lstlisting}
stage RR:
  new argument3 = _SX_10(signed10);
  new argument2 = GPR[registerZ];
  new result1 = _MAX(_SX_64(argument3), _SX_64(argument2));
stage E1:
  GPR[registerW] = result1;
\end{lstlisting}}
\newpage

\paragraph{Encoding} Format \textsf{ALU\_DWRI.X} (Section~\ref{sec:format-ALU-DWRI.X}), \verb|exucode2 = 0101|\\

\bitfields{ 1/P, 2/00, 2/-- --, 27/upper27, 1/1, 2/11, 1/0, 4/0101, 6/registerW, 2/00, 10/lower10, 6/registerZ }


\paragraph{Syntax} \texttt{maxd} \emph{registerW} = \emph{registerZ}, \emph{upper27\_\discretionary{}{}{}lower10}

\paragraph{Description} The maximum of the \emph{upper27\_\discretionary{}{}{}lower10} and the \emph{registerZ} is computed. The result is stored into the \emph{registerW}.


\paragraph{Execution} {Reservation table \textsf{ALU\_TINY.X} (Section~\ref{sec:reservation-ALU-TINY.X})}
~
{\begin{lstlisting}
stage RR:
  new argument3 = _SX_37(upper27_lower10);
  new argument2 = GPR[registerZ];
  new result1 = _MAX(_SX_64(argument3), _SX_64(argument2));
stage E1:
  GPR[registerW] = result1;
\end{lstlisting}}
\newpage

\paragraph{Encoding} Format \textsf{ALU\_DWRI.Y} (Section~\ref{sec:format-ALU-DWRI.Y}), \verb|exucode2 = 0101|\\

\bitfields{ 1/P, 2/00, 2/-- --, 27/extend27, 1/1, 2/00, 2/-- --, 27/upper27, 1/1, 2/11, 1/0, 4/0101, 6/registerW, 2/00, 10/lower10, 6/registerZ }


\paragraph{Syntax} \texttt{maxd} \emph{registerW} = \emph{registerZ}, \emph{extend27\_\discretionary{}{}{}upper27\_\discretionary{}{}{}lower10}

\paragraph{Description} The maximum of the \emph{extend27\_\discretionary{}{}{}upper27\_\discretionary{}{}{}lower10} and the \emph{registerZ} is computed. The result is stored into the \emph{registerW}.


\paragraph{Execution} {Reservation table \textsf{ALU\_TINY.Y} (Section~\ref{sec:reservation-ALU-TINY.Y})}
~
{\begin{lstlisting}
stage RR:
  new argument3 = _WX_64(extend27_upper27_lower10);
  new argument2 = GPR[registerZ];
  new result1 = _MAX(_SX_64(argument3), _SX_64(argument2));
stage E1:
  GPR[registerW] = result1;
\end{lstlisting}}
\newpage
\subsubsection[{\small MAXDP: Maximum Double Word Pair}]{MAXDP\hfill Maximum Double Word Pair}
\label{instr:MAXDP}\index{maxdp}
 
\paragraph{Encoding} Format \textsf{ALU\_DPWRR0} (Section~\ref{sec:format-ALU-DPWRR0}), \verb|exucode2 = 0101|\\

\bitfields{ 1/P, 2/11, 1/1, 4/0101, 5/registerM, 1/--, 2/10, 3/000, 1/0, 5/registerO, 1/--, 5/registerP, 1/0 }


\paragraph{Syntax} \texttt{maxdp} \emph{registerM} = \emph{registerP}, \emph{registerO}

\paragraph{Description} The \emph{registerO} and the \emph{registerP} are interpreted as two double words packed into 128 bits. The maximum of the \emph{registerO} double words and the \emph{registerP} double words are computed. The two resulting double words are packed and stored into the \emph{registerM}.


\paragraph{Execution} {Reservation table \textsf{ALU\_LITE} (Section~\ref{sec:reservation-ALU-LITE})}
~
{\begin{lstlisting}
stage RR:
  new argument3 = PGR[registerO];
  new argument2 = PGR[registerP];
  for (I = 0; I < 2; I += 1) {
    result1.64[I] = _MAX(_SX_64(argument3.64[I]), _SX_64(argument2.64[I]))
  };
stage E1:
  PGR[registerM] = result1;
\end{lstlisting}}
\newpage

\paragraph{Encoding} Format \textsf{ALU\_DPWRR0.M} (Section~\ref{sec:format-ALU-DPWRR0.M}), \verb|exucode2 = 0101|\\

\bitfields{ 1/P, 2/00, 2/-- --, 27/upper27, 1/1, 2/11, 1/1, 4/0101, 5/registerM, 1/--, 2/10, 3/000, 1/0, 1/splat32, 5/lower5, 5/registerP, 1/0 }


\paragraph{Syntax} \texttt{maxdp} \emph{registerM} = \emph{registerP}, \emph{upper27\_\discretionary{}{}{}lower5}\emph{splat32}

\paragraph{Description} The \emph{upper27\_\discretionary{}{}{}lower5} and the \emph{registerP} are interpreted as two double words packed into 128 bits. The maximum of the \emph{upper27\_\discretionary{}{}{}lower5} double words and the \emph{registerP} double words are computed. The two resulting double words are packed and stored into the \emph{registerM}.


\paragraph{Modifier(s)}
\noindent\begin{itemize}
\item splat32 (cf. splat32 in section \ref{par:modifier-splat32} on page \pageref{par:modifier-splat32})
\end{itemize}

\paragraph{Execution} {Reservation table \textsf{ALU\_LITE.X} (Section~\ref{sec:reservation-ALU-LITE.X})}
~
{\begin{lstlisting}
stage ID:
  new splat32 = _ZX_1(splat32);
stage RR:
  new argument3 = splat32 ? (_WX_32(upper27_lower5) << 32) | _ZX_32(_WX_32(upper27_lower5)) : _SX_32(_WX_32(upper27_lower5));
  new argument2 = PGR[registerP];
  for (I = 0; I < 2; I += 1) {
    result1.64[I] = _MAX(_SX_64(argument3.64[I]), _SX_64(argument2.64[I]))
  };
stage E1:
  PGR[registerM] = result1;
\end{lstlisting}}
\newpage
\subsubsection[{\small MAXHO: Maximum Half Word Octuple}]{MAXHO\hfill Maximum Half Word Octuple}
\label{instr:MAXHO}\index{maxho}
 
\paragraph{Encoding} Format \textsf{ALU\_HOWRR0} (Section~\ref{sec:format-ALU-HOWRR0}), \verb|exucode2 = 0101|\\

\bitfields{ 1/P, 2/11, 1/1, 4/0101, 5/registerM, 1/--, 2/10, 3/000, 1/1, 5/registerO, 1/--, 5/registerP, 1/0 }


\paragraph{Syntax} \texttt{maxho} \emph{registerM} = \emph{registerP}, \emph{registerO}

\paragraph{Description} The \emph{registerO} and the \emph{registerP} are interpreted as eight half words packed into 128 bits. The maximum of the \emph{registerO} half words and the \emph{registerP} half words are computed. The eight resulting half words are packed and stored into the \emph{registerM}.


\paragraph{Execution} {Reservation table \textsf{ALU\_LITE} (Section~\ref{sec:reservation-ALU-LITE})}
~
{\begin{lstlisting}
stage RR:
  new argument3 = PGR[registerO];
  new argument2 = PGR[registerP];
  for (I = 0; I < 8; I += 1) {
    result1.16[I] = _MAX(_SX_16(argument3.16[I]), _SX_16(argument2.16[I]))
  };
stage E1:
  PGR[registerM] = result1;
\end{lstlisting}}
\newpage

\paragraph{Encoding} Format \textsf{ALU\_HOWRR0.M} (Section~\ref{sec:format-ALU-HOWRR0.M}), \verb|exucode2 = 0101|\\

\bitfields{ 1/P, 2/00, 2/-- --, 27/upper27, 1/1, 2/11, 1/1, 4/0101, 5/registerM, 1/--, 2/10, 3/000, 1/1, 1/splat32, 5/lower5, 5/registerP, 1/0 }


\paragraph{Syntax} \texttt{maxho} \emph{registerM} = \emph{registerP}, \emph{upper27\_\discretionary{}{}{}lower5}\emph{splat32}

\paragraph{Description} The \emph{upper27\_\discretionary{}{}{}lower5} and the \emph{registerP} are interpreted as eight half words packed into 128 bits. The maximum of the \emph{upper27\_\discretionary{}{}{}lower5} half words and the \emph{registerP} half words are computed. The eight resulting half words are packed and stored into the \emph{registerM}.


\paragraph{Modifier(s)}
\noindent\begin{itemize}
\item splat32 (cf. splat32 in section \ref{par:modifier-splat32} on page \pageref{par:modifier-splat32})
\end{itemize}

\paragraph{Execution} {Reservation table \textsf{ALU\_LITE.X} (Section~\ref{sec:reservation-ALU-LITE.X})}
~
{\begin{lstlisting}
stage ID:
  new splat32 = _ZX_1(splat32);
stage RR:
  new argument3 = splat32 ? (_WX_32(upper27_lower5) << 32) | _ZX_32(_WX_32(upper27_lower5)) : _SX_32(_WX_32(upper27_lower5));
  new argument2 = PGR[registerP];
  for (I = 0; I < 8; I += 1) {
    result1.16[I] = _MAX(_SX_16(argument3.16[I]), _SX_16(argument2.16[I]))
  };
stage E1:
  PGR[registerM] = result1;
\end{lstlisting}}
\newpage
\subsubsection[{\small MAXUBX: Maximum Unsigned Byte Hexadecuple}]{MAXUBX\hfill Maximum Unsigned Byte Hexadecuple}
\label{instr:MAXUBX}\index{maxubx}
 
\paragraph{Encoding} Format \textsf{ALU\_BXWRR0} (Section~\ref{sec:format-ALU-BXWRR0}), \verb|exucode2 = 0111|\\

\bitfields{ 1/P, 2/11, 1/1, 4/0111, 5/registerM, 1/--, 2/10, 3/000, 1/1, 5/registerO, 1/--, 5/registerP, 1/1 }


\paragraph{Syntax} \texttt{maxubx} \emph{registerM} = \emph{registerP}, \emph{registerO}

\paragraph{Description} The \emph{registerO} and the \emph{registerP} are interpreted as sixteen bytes packed into 128 bits. The unsigned maximum of the \emph{registerO} bytes and the \emph{registerP} bytes are computed. The sixteen resulting bytes are packed and stored into the \emph{registerM}.


\paragraph{Execution} {Reservation table \textsf{ALU\_LITE} (Section~\ref{sec:reservation-ALU-LITE})}
~
{\begin{lstlisting}
stage RR:
  new argument3 = PGR[registerO];
  new argument2 = PGR[registerP];
  for (I = 0; I < 16; I += 1) {
    result1.8[I] = _MAX(_ZX_8(argument3.8[I]), _ZX_8(argument2.8[I]))
  };
stage E1:
  PGR[registerM] = result1;
\end{lstlisting}}
\newpage

\paragraph{Encoding} Format \textsf{ALU\_BXWRR0.M} (Section~\ref{sec:format-ALU-BXWRR0.M}), \verb|exucode2 = 0111|\\

\bitfields{ 1/P, 2/00, 2/-- --, 27/upper27, 1/1, 2/11, 1/1, 4/0111, 5/registerM, 1/--, 2/10, 3/000, 1/1, 1/splat32, 5/lower5, 5/registerP, 1/1 }


\paragraph{Syntax} \texttt{maxubx} \emph{registerM} = \emph{registerP}, \emph{upper27\_\discretionary{}{}{}lower5}\emph{splat32}

\paragraph{Description} The \emph{upper27\_\discretionary{}{}{}lower5} and the \emph{registerP} are interpreted as sixteen bytes packed into 128 bits. The unsigned maximum of the \emph{upper27\_\discretionary{}{}{}lower5} bytes and the \emph{registerP} bytes are computed. The sixteen resulting bytes are packed and stored into the \emph{registerM}.


\paragraph{Modifier(s)}
\noindent\begin{itemize}
\item splat32 (cf. splat32 in section \ref{par:modifier-splat32} on page \pageref{par:modifier-splat32})
\end{itemize}

\paragraph{Execution} {Reservation table \textsf{ALU\_LITE.X} (Section~\ref{sec:reservation-ALU-LITE.X})}
~
{\begin{lstlisting}
stage ID:
  new splat32 = _ZX_1(splat32);
stage RR:
  new argument3 = splat32 ? (_WX_32(upper27_lower5) << 32) | _ZX_32(_WX_32(upper27_lower5)) : _SX_32(_WX_32(upper27_lower5));
  new argument2 = PGR[registerP];
  for (I = 0; I < 16; I += 1) {
    result1.8[I] = _MAX(_ZX_8(argument3.8[I]), _ZX_8(argument2.8[I]))
  };
stage E1:
  PGR[registerM] = result1;
\end{lstlisting}}
\newpage
\subsubsection[{\small MAXUD: Maximum Unsigned Double Word}]{MAXUD\hfill Maximum Unsigned Double Word}
\label{instr:MAXUD}\index{maxud}
 
\paragraph{Encoding} Format \textsf{ALU\_DWRR0} (Section~\ref{sec:format-ALU-DWRR0}), \verb|exucode2 = 0111|\\

\bitfields{ 1/P, 2/11, 1/0, 4/0111, 6/registerW, 2/10, 3/000, 1/0, 6/registerY, 6/registerZ }


\paragraph{Syntax} \texttt{maxud} \emph{registerW} = \emph{registerZ}, \emph{registerY}

\paragraph{Description} The unsigned maximum of the \emph{registerY} and the \emph{registerZ} is computed. The result is stored into the \emph{registerW}.


\paragraph{Execution} {Reservation table \textsf{ALU\_TINY} (Section~\ref{sec:reservation-ALU-TINY})}
~
{\begin{lstlisting}
stage RR:
  new argument3 = GPR[registerY];
  new argument2 = GPR[registerZ];
  new result1 = _MAX(_ZX_64(argument3), _ZX_64(argument2));
stage E1:
  GPR[registerW] = result1;
\end{lstlisting}}
\newpage

\paragraph{Encoding} Format \textsf{ALU\_DWRR0.M} (Section~\ref{sec:format-ALU-DWRR0.M}), \verb|exucode2 = 0111|\\

\bitfields{ 1/P, 2/00, 2/-- --, 27/upper27, 1/1, 2/11, 1/0, 4/0111, 6/registerW, 2/10, 3/000, 1/0, 1/splat32, 5/lower5, 6/registerZ }


\paragraph{Syntax} \texttt{maxud} \emph{registerW} = \emph{registerZ}, \emph{upper27\_\discretionary{}{}{}lower5}\emph{splat32}

\paragraph{Description} The unsigned maximum of the \emph{upper27\_\discretionary{}{}{}lower5} and the \emph{registerZ} is computed. The result is stored into the \emph{registerW}.


\paragraph{Modifier(s)}
\noindent\begin{itemize}
\item splat32 (cf. splat32 in section \ref{par:modifier-splat32} on page \pageref{par:modifier-splat32})
\end{itemize}

\paragraph{Execution} {Reservation table \textsf{ALU\_TINY.X} (Section~\ref{sec:reservation-ALU-TINY.X})}
~
{\begin{lstlisting}
stage ID:
  new splat32 = _ZX_1(splat32);
stage RR:
  new argument3 = splat32 ? (_WX_32(upper27_lower5) << 32) | _ZX_32(_WX_32(upper27_lower5)) : _SX_32(_WX_32(upper27_lower5));
  new argument2 = GPR[registerZ];
  new result1 = _MAX(_ZX_64(argument3), _ZX_64(argument2));
stage E1:
  GPR[registerW] = result1;
\end{lstlisting}}
\newpage

\paragraph{Encoding} Format \textsf{ALU\_DWRI} (Section~\ref{sec:format-ALU-DWRI}), \verb|exucode2 = 0111|\\

\bitfields{ 1/P, 2/11, 1/0, 4/0111, 6/registerW, 2/00, 10/signed10, 6/registerZ }


\paragraph{Syntax} \texttt{maxud} \emph{registerW} = \emph{registerZ}, \emph{signed10}

\paragraph{Description} The unsigned maximum of the \emph{signed10} and the \emph{registerZ} is computed. The result is stored into the \emph{registerW}.


\paragraph{Execution} {Reservation table \textsf{ALU\_TINY} (Section~\ref{sec:reservation-ALU-TINY})}
~
{\begin{lstlisting}
stage RR:
  new argument3 = _SX_10(signed10);
  new argument2 = GPR[registerZ];
  new result1 = _MAX(_ZX_64(argument3), _ZX_64(argument2));
stage E1:
  GPR[registerW] = result1;
\end{lstlisting}}
\newpage

\paragraph{Encoding} Format \textsf{ALU\_DWRI.X} (Section~\ref{sec:format-ALU-DWRI.X}), \verb|exucode2 = 0111|\\

\bitfields{ 1/P, 2/00, 2/-- --, 27/upper27, 1/1, 2/11, 1/0, 4/0111, 6/registerW, 2/00, 10/lower10, 6/registerZ }


\paragraph{Syntax} \texttt{maxud} \emph{registerW} = \emph{registerZ}, \emph{upper27\_\discretionary{}{}{}lower10}

\paragraph{Description} The unsigned maximum of the \emph{upper27\_\discretionary{}{}{}lower10} and the \emph{registerZ} is computed. The result is stored into the \emph{registerW}.


\paragraph{Execution} {Reservation table \textsf{ALU\_TINY.X} (Section~\ref{sec:reservation-ALU-TINY.X})}
~
{\begin{lstlisting}
stage RR:
  new argument3 = _SX_37(upper27_lower10);
  new argument2 = GPR[registerZ];
  new result1 = _MAX(_ZX_64(argument3), _ZX_64(argument2));
stage E1:
  GPR[registerW] = result1;
\end{lstlisting}}
\newpage

\paragraph{Encoding} Format \textsf{ALU\_DWRI.Y} (Section~\ref{sec:format-ALU-DWRI.Y}), \verb|exucode2 = 0111|\\

\bitfields{ 1/P, 2/00, 2/-- --, 27/extend27, 1/1, 2/00, 2/-- --, 27/upper27, 1/1, 2/11, 1/0, 4/0111, 6/registerW, 2/00, 10/lower10, 6/registerZ }


\paragraph{Syntax} \texttt{maxud} \emph{registerW} = \emph{registerZ}, \emph{extend27\_\discretionary{}{}{}upper27\_\discretionary{}{}{}lower10}

\paragraph{Description} The unsigned maximum of the \emph{extend27\_\discretionary{}{}{}upper27\_\discretionary{}{}{}lower10} and the \emph{registerZ} is computed. The result is stored into the \emph{registerW}.


\paragraph{Execution} {Reservation table \textsf{ALU\_TINY.Y} (Section~\ref{sec:reservation-ALU-TINY.Y})}
~
{\begin{lstlisting}
stage RR:
  new argument3 = _WX_64(extend27_upper27_lower10);
  new argument2 = GPR[registerZ];
  new result1 = _MAX(_ZX_64(argument3), _ZX_64(argument2));
stage E1:
  GPR[registerW] = result1;
\end{lstlisting}}
\newpage
\subsubsection[{\small MAXUDP: Maximum Unsigned Double Word Pair}]{MAXUDP\hfill Maximum Unsigned Double Word Pair}
\label{instr:MAXUDP}\index{maxudp}
 
\paragraph{Encoding} Format \textsf{ALU\_DPWRR0} (Section~\ref{sec:format-ALU-DPWRR0}), \verb|exucode2 = 0111|\\

\bitfields{ 1/P, 2/11, 1/1, 4/0111, 5/registerM, 1/--, 2/10, 3/000, 1/0, 5/registerO, 1/--, 5/registerP, 1/0 }


\paragraph{Syntax} \texttt{maxudp} \emph{registerM} = \emph{registerP}, \emph{registerO}

\paragraph{Description} The \emph{registerO} and the \emph{registerP} are interpreted as two double words packed into 128 bits. The unsigned maximum of the \emph{registerO} words and the \emph{registerP} words are computed. The two resulting double words are packed and stored into the \emph{registerM}.


\paragraph{Execution} {Reservation table \textsf{ALU\_LITE} (Section~\ref{sec:reservation-ALU-LITE})}
~
{\begin{lstlisting}
stage RR:
  new argument3 = PGR[registerO];
  new argument2 = PGR[registerP];
  for (I = 0; I < 2; I += 1) {
    result1.64[I] = _MAX(_ZX_64(argument3.64[I]), _ZX_64(argument2.64[I]))
  };
stage E1:
  PGR[registerM] = result1;
\end{lstlisting}}
\newpage

\paragraph{Encoding} Format \textsf{ALU\_DPWRR0.M} (Section~\ref{sec:format-ALU-DPWRR0.M}), \verb|exucode2 = 0111|\\

\bitfields{ 1/P, 2/00, 2/-- --, 27/upper27, 1/1, 2/11, 1/1, 4/0111, 5/registerM, 1/--, 2/10, 3/000, 1/0, 1/splat32, 5/lower5, 5/registerP, 1/0 }


\paragraph{Syntax} \texttt{maxudp} \emph{registerM} = \emph{registerP}, \emph{upper27\_\discretionary{}{}{}lower5}\emph{splat32}

\paragraph{Description} The \emph{upper27\_\discretionary{}{}{}lower5} and the \emph{registerP} are interpreted as two double words packed into 128 bits. The unsigned maximum of the \emph{upper27\_\discretionary{}{}{}lower5} words and the \emph{registerP} words are computed. The two resulting double words are packed and stored into the \emph{registerM}.


\paragraph{Modifier(s)}
\noindent\begin{itemize}
\item splat32 (cf. splat32 in section \ref{par:modifier-splat32} on page \pageref{par:modifier-splat32})
\end{itemize}

\paragraph{Execution} {Reservation table \textsf{ALU\_LITE.X} (Section~\ref{sec:reservation-ALU-LITE.X})}
~
{\begin{lstlisting}
stage ID:
  new splat32 = _ZX_1(splat32);
stage RR:
  new argument3 = splat32 ? (_WX_32(upper27_lower5) << 32) | _ZX_32(_WX_32(upper27_lower5)) : _SX_32(_WX_32(upper27_lower5));
  new argument2 = PGR[registerP];
  for (I = 0; I < 2; I += 1) {
    result1.64[I] = _MAX(_ZX_64(argument3.64[I]), _ZX_64(argument2.64[I]))
  };
stage E1:
  PGR[registerM] = result1;
\end{lstlisting}}
\newpage
\subsubsection[{\small MAXUHO: Maximum Unsigned Half Word Octuple}]{MAXUHO\hfill Maximum Unsigned Half Word Octuple}
\label{instr:MAXUHO}\index{maxuho}
 
\paragraph{Encoding} Format \textsf{ALU\_HOWRR0} (Section~\ref{sec:format-ALU-HOWRR0}), \verb|exucode2 = 0111|\\

\bitfields{ 1/P, 2/11, 1/1, 4/0111, 5/registerM, 1/--, 2/10, 3/000, 1/1, 5/registerO, 1/--, 5/registerP, 1/0 }


\paragraph{Syntax} \texttt{maxuho} \emph{registerM} = \emph{registerP}, \emph{registerO}

\paragraph{Description} The \emph{registerO} and the \emph{registerP} are interpreted as eight half words packed into 128 bits. The unsigned maximum of the \emph{registerO} half words and the \emph{registerP} half words are computed. The eight resulting half words are packed and stored into the \emph{registerM}.


\paragraph{Execution} {Reservation table \textsf{ALU\_LITE} (Section~\ref{sec:reservation-ALU-LITE})}
~
{\begin{lstlisting}
stage RR:
  new argument3 = PGR[registerO];
  new argument2 = PGR[registerP];
  for (I = 0; I < 8; I += 1) {
    result1.16[I] = _MAX(_ZX_16(argument3.16[I]), _ZX_16(argument2.16[I]))
  };
stage E1:
  PGR[registerM] = result1;
\end{lstlisting}}
\newpage

\paragraph{Encoding} Format \textsf{ALU\_HOWRR0.M} (Section~\ref{sec:format-ALU-HOWRR0.M}), \verb|exucode2 = 0111|\\

\bitfields{ 1/P, 2/00, 2/-- --, 27/upper27, 1/1, 2/11, 1/1, 4/0111, 5/registerM, 1/--, 2/10, 3/000, 1/1, 1/splat32, 5/lower5, 5/registerP, 1/0 }


\paragraph{Syntax} \texttt{maxuho} \emph{registerM} = \emph{registerP}, \emph{upper27\_\discretionary{}{}{}lower5}\emph{splat32}

\paragraph{Description} The \emph{upper27\_\discretionary{}{}{}lower5} and the \emph{registerP} are interpreted as eight half words packed into 128 bits. The unsigned maximum of the \emph{upper27\_\discretionary{}{}{}lower5} half words and the \emph{registerP} half words are computed. The eight resulting half words are packed and stored into the \emph{registerM}.


\paragraph{Modifier(s)}
\noindent\begin{itemize}
\item splat32 (cf. splat32 in section \ref{par:modifier-splat32} on page \pageref{par:modifier-splat32})
\end{itemize}

\paragraph{Execution} {Reservation table \textsf{ALU\_LITE.X} (Section~\ref{sec:reservation-ALU-LITE.X})}
~
{\begin{lstlisting}
stage ID:
  new splat32 = _ZX_1(splat32);
stage RR:
  new argument3 = splat32 ? (_WX_32(upper27_lower5) << 32) | _ZX_32(_WX_32(upper27_lower5)) : _SX_32(_WX_32(upper27_lower5));
  new argument2 = PGR[registerP];
  for (I = 0; I < 8; I += 1) {
    result1.16[I] = _MAX(_ZX_16(argument3.16[I]), _ZX_16(argument2.16[I]))
  };
stage E1:
  PGR[registerM] = result1;
\end{lstlisting}}
\newpage
\subsubsection[{\small MAXUW: Maximum Unsigned Word}]{MAXUW\hfill Maximum Unsigned Word}
\label{instr:MAXUW}\index{maxuw}
 
\paragraph{Encoding} Format \textsf{ALU\_WWRR0} (Section~\ref{sec:format-ALU-WWRR0}), \verb|exucode2 = 0111|\\

\bitfields{ 1/P, 2/11, 1/0, 4/0111, 6/registerW, 2/11, 3/000, 1/signextw, 6/registerY, 6/registerZ }


\paragraph{Syntax} \texttt{maxuw}\emph{signextw} \emph{registerW} = \emph{registerZ}, \emph{registerY}

\paragraph{Description} The unsigned maximum of the \emph{registerY} and the \emph{registerZ} is computed. The result with the 32 upper bits cleared is stored into the \emph{registerW}.


\paragraph{Modifier(s)}
\noindent\begin{itemize}
\item signextw (cf. signextw in section \ref{par:modifier-signextw} on page \pageref{par:modifier-signextw})
\end{itemize}

\paragraph{Execution} {Reservation table \textsf{ALU\_TINY} (Section~\ref{sec:reservation-ALU-TINY})}
~
{\begin{lstlisting}
stage ID:
  new signextw = _ZX_1(signextw);
stage RR:
  new argument3 = GPR[registerY];
  new argument2 = GPR[registerZ];
  new result1 = _MAX(argument3.32[0], argument2.32[0]);
stage E1:
  GPR[registerW] = signextw ? _SX_32(result1) : _ZX_32(result1);
\end{lstlisting}}
\newpage

\paragraph{Encoding} Format \textsf{ALU\_WWRR0.W} (Section~\ref{sec:format-ALU-WWRR0.W}), \verb|exucode2 = 0111|\\

\bitfields{ 1/P, 2/00, 2/-- --, 27/upper27, 1/1, 2/11, 1/0, 4/0111, 6/registerW, 2/11, 3/000, 1/signextw, 1/0, 5/lower5, 6/registerZ }


\paragraph{Syntax} \texttt{maxuw}\emph{signextw} \emph{registerW} = \emph{registerZ}, \emph{upper27\_\discretionary{}{}{}lower5}

\paragraph{Description} The unsigned maximum of the \emph{upper27\_\discretionary{}{}{}lower5} and the \emph{registerZ} is computed. The result with the 32 upper bits cleared is stored into the \emph{registerW}.


\paragraph{Modifier(s)}
\noindent\begin{itemize}
\item signextw (cf. signextw in section \ref{par:modifier-signextw} on page \pageref{par:modifier-signextw})
\end{itemize}

\paragraph{Execution} {Reservation table \textsf{ALU\_TINY.X} (Section~\ref{sec:reservation-ALU-TINY.X})}
~
{\begin{lstlisting}
stage ID:
  new signextw = _ZX_1(signextw);
stage RR:
  new argument3 = _WX_32(upper27_lower5);
  new argument2 = GPR[registerZ];
  new result1 = _MAX(argument3.32[0], argument2.32[0]);
stage E1:
  GPR[registerW] = signextw ? _SX_32(result1) : _ZX_32(result1);
\end{lstlisting}}
\newpage
\subsubsection[{\small MAXUWQ: Maximum Unsigned Word Quadruple}]{MAXUWQ\hfill Maximum Unsigned Word Quadruple}
\label{instr:MAXUWQ}\index{maxuwq}
 
\paragraph{Encoding} Format \textsf{ALU\_WQWRR0} (Section~\ref{sec:format-ALU-WQWRR0}), \verb|exucode2 = 0111|\\

\bitfields{ 1/P, 2/11, 1/1, 4/0111, 5/registerM, 1/--, 2/10, 3/000, 1/0, 5/registerO, 1/--, 5/registerP, 1/1 }


\paragraph{Syntax} \texttt{maxuwq} \emph{registerM} = \emph{registerP}, \emph{registerO}

\paragraph{Description} The \emph{registerO} and the \emph{registerP} are interpreted as four words packed into 128 bits. The unsigned maximum of the \emph{registerO} words and the \emph{registerP} words are computed. The four resulting words are packed and stored into the \emph{registerM}.


\paragraph{Execution} {Reservation table \textsf{ALU\_LITE} (Section~\ref{sec:reservation-ALU-LITE})}
~
{\begin{lstlisting}
stage RR:
  new argument3 = PGR[registerO];
  new argument2 = PGR[registerP];
  for (I = 0; I < 4; I += 1) {
    result1.32[I] = _MAX(_ZX_32(argument3.32[I]), _ZX_32(argument2.32[I]))
  };
stage E1:
  PGR[registerM] = result1;
\end{lstlisting}}
\newpage

\paragraph{Encoding} Format \textsf{ALU\_WQWRR0.M} (Section~\ref{sec:format-ALU-WQWRR0.M}), \verb|exucode2 = 0111|\\

\bitfields{ 1/P, 2/00, 2/-- --, 27/upper27, 1/1, 2/11, 1/1, 4/0111, 5/registerM, 1/--, 2/10, 3/000, 1/0, 1/splat32, 5/lower5, 5/registerP, 1/1 }


\paragraph{Syntax} \texttt{maxuwq} \emph{registerM} = \emph{registerP}, \emph{upper27\_\discretionary{}{}{}lower5}\emph{splat32}

\paragraph{Description} The \emph{upper27\_\discretionary{}{}{}lower5} and the \emph{registerP} are interpreted as four words packed into 128 bits. The unsigned maximum of the \emph{upper27\_\discretionary{}{}{}lower5} words and the \emph{registerP} words are computed. The four resulting words are packed and stored into the \emph{registerM}.


\paragraph{Modifier(s)}
\noindent\begin{itemize}
\item splat32 (cf. splat32 in section \ref{par:modifier-splat32} on page \pageref{par:modifier-splat32})
\end{itemize}

\paragraph{Execution} {Reservation table \textsf{ALU\_LITE.X} (Section~\ref{sec:reservation-ALU-LITE.X})}
~
{\begin{lstlisting}
stage ID:
  new splat32 = _ZX_1(splat32);
stage RR:
  new argument3 = splat32 ? (_WX_32(upper27_lower5) << 32) | _ZX_32(_WX_32(upper27_lower5)) : _SX_32(_WX_32(upper27_lower5));
  new argument2 = PGR[registerP];
  for (I = 0; I < 4; I += 1) {
    result1.32[I] = _MAX(_ZX_32(argument3.32[I]), _ZX_32(argument2.32[I]))
  };
stage E1:
  PGR[registerM] = result1;
\end{lstlisting}}
\newpage
\subsubsection[{\small MAXW: Maximum Word}]{MAXW\hfill Maximum Word}
\label{instr:MAXW}\index{maxw}
 
\paragraph{Encoding} Format \textsf{ALU\_WWRR0} (Section~\ref{sec:format-ALU-WWRR0}), \verb|exucode2 = 0101|\\

\bitfields{ 1/P, 2/11, 1/0, 4/0101, 6/registerW, 2/11, 3/000, 1/signextw, 6/registerY, 6/registerZ }


\paragraph{Syntax} \texttt{maxw}\emph{signextw} \emph{registerW} = \emph{registerZ}, \emph{registerY}

\paragraph{Description} The maximum of the \emph{registerY} and the \emph{registerZ} is computed. The result with the 32 upper bits cleared is stored into the \emph{registerW}.


\paragraph{Modifier(s)}
\noindent\begin{itemize}
\item signextw (cf. signextw in section \ref{par:modifier-signextw} on page \pageref{par:modifier-signextw})
\end{itemize}

\paragraph{Execution} {Reservation table \textsf{ALU\_TINY} (Section~\ref{sec:reservation-ALU-TINY})}
~
{\begin{lstlisting}
stage ID:
  new signextw = _ZX_1(signextw);
stage RR:
  new argument3 = GPR[registerY];
  new argument2 = GPR[registerZ];
  new result1 = _MAX(_SX_32(argument3), _SX_32(argument2));
stage E1:
  GPR[registerW] = signextw ? _SX_32(result1) : _ZX_32(result1);
\end{lstlisting}}
\newpage

\paragraph{Encoding} Format \textsf{ALU\_WWRR0.W} (Section~\ref{sec:format-ALU-WWRR0.W}), \verb|exucode2 = 0101|\\

\bitfields{ 1/P, 2/00, 2/-- --, 27/upper27, 1/1, 2/11, 1/0, 4/0101, 6/registerW, 2/11, 3/000, 1/signextw, 1/0, 5/lower5, 6/registerZ }


\paragraph{Syntax} \texttt{maxw}\emph{signextw} \emph{registerW} = \emph{registerZ}, \emph{upper27\_\discretionary{}{}{}lower5}

\paragraph{Description} The maximum of the \emph{upper27\_\discretionary{}{}{}lower5} and the \emph{registerZ} is computed. The result with the 32 upper bits cleared is stored into the \emph{registerW}.


\paragraph{Modifier(s)}
\noindent\begin{itemize}
\item signextw (cf. signextw in section \ref{par:modifier-signextw} on page \pageref{par:modifier-signextw})
\end{itemize}

\paragraph{Execution} {Reservation table \textsf{ALU\_TINY.X} (Section~\ref{sec:reservation-ALU-TINY.X})}
~
{\begin{lstlisting}
stage ID:
  new signextw = _ZX_1(signextw);
stage RR:
  new argument3 = _WX_32(upper27_lower5);
  new argument2 = GPR[registerZ];
  new result1 = _MAX(_SX_32(argument3), _SX_32(argument2));
stage E1:
  GPR[registerW] = signextw ? _SX_32(result1) : _ZX_32(result1);
\end{lstlisting}}
\newpage
\subsubsection[{\small MAXWQ: Maximum Word Quadruple}]{MAXWQ\hfill Maximum Word Quadruple}
\label{instr:MAXWQ}\index{maxwq}
 
\paragraph{Encoding} Format \textsf{ALU\_WQWRR0} (Section~\ref{sec:format-ALU-WQWRR0}), \verb|exucode2 = 0101|\\

\bitfields{ 1/P, 2/11, 1/1, 4/0101, 5/registerM, 1/--, 2/10, 3/000, 1/0, 5/registerO, 1/--, 5/registerP, 1/1 }


\paragraph{Syntax} \texttt{maxwq} \emph{registerM} = \emph{registerP}, \emph{registerO}

\paragraph{Description} The \emph{registerO} and the \emph{registerP} are interpreted as four words packed into 128 bits. The maximum of the \emph{registerO} words and the \emph{registerP} words are computed. The four resulting words are packed and stored into the \emph{registerM}.


\paragraph{Execution} {Reservation table \textsf{ALU\_LITE} (Section~\ref{sec:reservation-ALU-LITE})}
~
{\begin{lstlisting}
stage RR:
  new argument3 = PGR[registerO];
  new argument2 = PGR[registerP];
  for (I = 0; I < 4; I += 1) {
    result1.32[I] = _MAX(_SX_32(argument3.32[I]), _SX_32(argument2.32[I]))
  };
stage E1:
  PGR[registerM] = result1;
\end{lstlisting}}
\newpage

\paragraph{Encoding} Format \textsf{ALU\_WQWRR0.M} (Section~\ref{sec:format-ALU-WQWRR0.M}), \verb|exucode2 = 0101|\\

\bitfields{ 1/P, 2/00, 2/-- --, 27/upper27, 1/1, 2/11, 1/1, 4/0101, 5/registerM, 1/--, 2/10, 3/000, 1/0, 1/splat32, 5/lower5, 5/registerP, 1/1 }


\paragraph{Syntax} \texttt{maxwq} \emph{registerM} = \emph{registerP}, \emph{upper27\_\discretionary{}{}{}lower5}\emph{splat32}

\paragraph{Description} The \emph{upper27\_\discretionary{}{}{}lower5} and the \emph{registerP} are interpreted as four words packed into 128 bits. The maximum of the \emph{upper27\_\discretionary{}{}{}lower5} words and the \emph{registerP} words are computed. The four resulting words are packed and stored into the \emph{registerM}.


\paragraph{Modifier(s)}
\noindent\begin{itemize}
\item splat32 (cf. splat32 in section \ref{par:modifier-splat32} on page \pageref{par:modifier-splat32})
\end{itemize}

\paragraph{Execution} {Reservation table \textsf{ALU\_LITE.X} (Section~\ref{sec:reservation-ALU-LITE.X})}
~
{\begin{lstlisting}
stage ID:
  new splat32 = _ZX_1(splat32);
stage RR:
  new argument3 = splat32 ? (_WX_32(upper27_lower5) << 32) | _ZX_32(_WX_32(upper27_lower5)) : _SX_32(_WX_32(upper27_lower5));
  new argument2 = PGR[registerP];
  for (I = 0; I < 4; I += 1) {
    result1.32[I] = _MAX(_SX_32(argument3.32[I]), _SX_32(argument2.32[I]))
  };
stage E1:
  PGR[registerM] = result1;
\end{lstlisting}}
\newpage
\subsubsection[{\small MINBX: Minimum Byte Hexadecuple}]{MINBX\hfill Minimum Byte Hexadecuple}
\label{instr:MINBX}\index{minbx}
 
\paragraph{Encoding} Format \textsf{ALU\_BXWRR0} (Section~\ref{sec:format-ALU-BXWRR0}), \verb|exucode2 = 0100|\\

\bitfields{ 1/P, 2/11, 1/1, 4/0100, 5/registerM, 1/--, 2/10, 3/000, 1/1, 5/registerO, 1/--, 5/registerP, 1/1 }


\paragraph{Syntax} \texttt{minbx} \emph{registerM} = \emph{registerP}, \emph{registerO}

\paragraph{Description} The \emph{registerO} and the \emph{registerP} are interpreted as sixteen bytes packed into 128 bits. The minimum of the \emph{registerO} bytes and the \emph{registerP} bytes are computed. The sixteen resulting bytes are packed and stored into the \emph{registerM}.


\paragraph{Execution} {Reservation table \textsf{ALU\_LITE} (Section~\ref{sec:reservation-ALU-LITE})}
~
{\begin{lstlisting}
stage RR:
  new argument3 = PGR[registerO];
  new argument2 = PGR[registerP];
  for (I = 0; I < 16; I += 1) {
    result1.8[I] = _MIN(_SX_8(argument3.8[I]), _SX_8(argument2.8[I]))
  };
stage E1:
  PGR[registerM] = result1;
\end{lstlisting}}
\newpage

\paragraph{Encoding} Format \textsf{ALU\_BXWRR0.M} (Section~\ref{sec:format-ALU-BXWRR0.M}), \verb|exucode2 = 0100|\\

\bitfields{ 1/P, 2/00, 2/-- --, 27/upper27, 1/1, 2/11, 1/1, 4/0100, 5/registerM, 1/--, 2/10, 3/000, 1/1, 1/splat32, 5/lower5, 5/registerP, 1/1 }


\paragraph{Syntax} \texttt{minbx} \emph{registerM} = \emph{registerP}, \emph{upper27\_\discretionary{}{}{}lower5}\emph{splat32}

\paragraph{Description} The \emph{upper27\_\discretionary{}{}{}lower5} and the \emph{registerP} are interpreted as sixteen bytes packed into 128 bits. The minimum of the \emph{upper27\_\discretionary{}{}{}lower5} bytes and the \emph{registerP} bytes are computed. The sixteen resulting bytes are packed and stored into the \emph{registerM}.


\paragraph{Modifier(s)}
\noindent\begin{itemize}
\item splat32 (cf. splat32 in section \ref{par:modifier-splat32} on page \pageref{par:modifier-splat32})
\end{itemize}

\paragraph{Execution} {Reservation table \textsf{ALU\_LITE.X} (Section~\ref{sec:reservation-ALU-LITE.X})}
~
{\begin{lstlisting}
stage ID:
  new splat32 = _ZX_1(splat32);
stage RR:
  new argument3 = splat32 ? (_WX_32(upper27_lower5) << 32) | _ZX_32(_WX_32(upper27_lower5)) : _SX_32(_WX_32(upper27_lower5));
  new argument2 = PGR[registerP];
  for (I = 0; I < 16; I += 1) {
    result1.8[I] = _MIN(_SX_8(argument3.8[I]), _SX_8(argument2.8[I]))
  };
stage E1:
  PGR[registerM] = result1;
\end{lstlisting}}
\newpage
\subsubsection[{\small MIND: Minimum Double Word}]{MIND\hfill Minimum Double Word}
\label{instr:MIND}\index{mind}
 
\paragraph{Encoding} Format \textsf{ALU\_DWRR0} (Section~\ref{sec:format-ALU-DWRR0}), \verb|exucode2 = 0100|\\

\bitfields{ 1/P, 2/11, 1/0, 4/0100, 6/registerW, 2/10, 3/000, 1/0, 6/registerY, 6/registerZ }


\paragraph{Syntax} \texttt{mind} \emph{registerW} = \emph{registerZ}, \emph{registerY}

\paragraph{Description} The minimum of the \emph{registerY} and the \emph{registerZ} is computed. The result is stored into the \emph{registerW}.


\paragraph{Execution} {Reservation table \textsf{ALU\_TINY} (Section~\ref{sec:reservation-ALU-TINY})}
~
{\begin{lstlisting}
stage RR:
  new argument3 = GPR[registerY];
  new argument2 = GPR[registerZ];
  new result1 = _MIN(_SX_64(argument3), _SX_64(argument2));
stage E1:
  GPR[registerW] = result1;
\end{lstlisting}}
\newpage

\paragraph{Encoding} Format \textsf{ALU\_DWRR0.M} (Section~\ref{sec:format-ALU-DWRR0.M}), \verb|exucode2 = 0100|\\

\bitfields{ 1/P, 2/00, 2/-- --, 27/upper27, 1/1, 2/11, 1/0, 4/0100, 6/registerW, 2/10, 3/000, 1/0, 1/splat32, 5/lower5, 6/registerZ }


\paragraph{Syntax} \texttt{mind} \emph{registerW} = \emph{registerZ}, \emph{upper27\_\discretionary{}{}{}lower5}\emph{splat32}

\paragraph{Description} The minimum of the \emph{upper27\_\discretionary{}{}{}lower5} and the \emph{registerZ} is computed. The result is stored into the \emph{registerW}.


\paragraph{Modifier(s)}
\noindent\begin{itemize}
\item splat32 (cf. splat32 in section \ref{par:modifier-splat32} on page \pageref{par:modifier-splat32})
\end{itemize}

\paragraph{Execution} {Reservation table \textsf{ALU\_TINY.X} (Section~\ref{sec:reservation-ALU-TINY.X})}
~
{\begin{lstlisting}
stage ID:
  new splat32 = _ZX_1(splat32);
stage RR:
  new argument3 = splat32 ? (_WX_32(upper27_lower5) << 32) | _ZX_32(_WX_32(upper27_lower5)) : _SX_32(_WX_32(upper27_lower5));
  new argument2 = GPR[registerZ];
  new result1 = _MIN(_SX_64(argument3), _SX_64(argument2));
stage E1:
  GPR[registerW] = result1;
\end{lstlisting}}
\newpage

\paragraph{Encoding} Format \textsf{ALU\_DWRI} (Section~\ref{sec:format-ALU-DWRI}), \verb|exucode2 = 0100|\\

\bitfields{ 1/P, 2/11, 1/0, 4/0100, 6/registerW, 2/00, 10/signed10, 6/registerZ }


\paragraph{Syntax} \texttt{mind} \emph{registerW} = \emph{registerZ}, \emph{signed10}

\paragraph{Description} The minimum of the \emph{signed10} and the \emph{registerZ} is computed. The result is stored into the \emph{registerW}.


\paragraph{Execution} {Reservation table \textsf{ALU\_TINY} (Section~\ref{sec:reservation-ALU-TINY})}
~
{\begin{lstlisting}
stage RR:
  new argument3 = _SX_10(signed10);
  new argument2 = GPR[registerZ];
  new result1 = _MIN(_SX_64(argument3), _SX_64(argument2));
stage E1:
  GPR[registerW] = result1;
\end{lstlisting}}
\newpage

\paragraph{Encoding} Format \textsf{ALU\_DWRI.X} (Section~\ref{sec:format-ALU-DWRI.X}), \verb|exucode2 = 0100|\\

\bitfields{ 1/P, 2/00, 2/-- --, 27/upper27, 1/1, 2/11, 1/0, 4/0100, 6/registerW, 2/00, 10/lower10, 6/registerZ }


\paragraph{Syntax} \texttt{mind} \emph{registerW} = \emph{registerZ}, \emph{upper27\_\discretionary{}{}{}lower10}

\paragraph{Description} The minimum of the \emph{upper27\_\discretionary{}{}{}lower10} and the \emph{registerZ} is computed. The result is stored into the \emph{registerW}.


\paragraph{Execution} {Reservation table \textsf{ALU\_TINY.X} (Section~\ref{sec:reservation-ALU-TINY.X})}
~
{\begin{lstlisting}
stage RR:
  new argument3 = _SX_37(upper27_lower10);
  new argument2 = GPR[registerZ];
  new result1 = _MIN(_SX_64(argument3), _SX_64(argument2));
stage E1:
  GPR[registerW] = result1;
\end{lstlisting}}
\newpage

\paragraph{Encoding} Format \textsf{ALU\_DWRI.Y} (Section~\ref{sec:format-ALU-DWRI.Y}), \verb|exucode2 = 0100|\\

\bitfields{ 1/P, 2/00, 2/-- --, 27/extend27, 1/1, 2/00, 2/-- --, 27/upper27, 1/1, 2/11, 1/0, 4/0100, 6/registerW, 2/00, 10/lower10, 6/registerZ }


\paragraph{Syntax} \texttt{mind} \emph{registerW} = \emph{registerZ}, \emph{extend27\_\discretionary{}{}{}upper27\_\discretionary{}{}{}lower10}

\paragraph{Description} The minimum of the \emph{extend27\_\discretionary{}{}{}upper27\_\discretionary{}{}{}lower10} and the \emph{registerZ} is computed. The result is stored into the \emph{registerW}.


\paragraph{Execution} {Reservation table \textsf{ALU\_TINY.Y} (Section~\ref{sec:reservation-ALU-TINY.Y})}
~
{\begin{lstlisting}
stage RR:
  new argument3 = _WX_64(extend27_upper27_lower10);
  new argument2 = GPR[registerZ];
  new result1 = _MIN(_SX_64(argument3), _SX_64(argument2));
stage E1:
  GPR[registerW] = result1;
\end{lstlisting}}
\newpage
\subsubsection[{\small MINDP: Minimum Double Word Pair}]{MINDP\hfill Minimum Double Word Pair}
\label{instr:MINDP}\index{mindp}
 
\paragraph{Encoding} Format \textsf{ALU\_DPWRR0} (Section~\ref{sec:format-ALU-DPWRR0}), \verb|exucode2 = 0100|\\

\bitfields{ 1/P, 2/11, 1/1, 4/0100, 5/registerM, 1/--, 2/10, 3/000, 1/0, 5/registerO, 1/--, 5/registerP, 1/0 }


\paragraph{Syntax} \texttt{mindp} \emph{registerM} = \emph{registerP}, \emph{registerO}

\paragraph{Description} The \emph{registerO} and the \emph{registerP} are interpreted as two double words packed into 128 bits. The minimum of the \emph{registerO} double words and the \emph{registerP} double words are computed. The two resulting double words are packed and stored into the \emph{registerM}.


\paragraph{Execution} {Reservation table \textsf{ALU\_LITE} (Section~\ref{sec:reservation-ALU-LITE})}
~
{\begin{lstlisting}
stage RR:
  new argument3 = PGR[registerO];
  new argument2 = PGR[registerP];
  for (I = 0; I < 2; I += 1) {
    result1.64[I] = _MIN(_SX_64(argument3.64[I]), _SX_64(argument2.64[I]))
  };
stage E1:
  PGR[registerM] = result1;
\end{lstlisting}}
\newpage

\paragraph{Encoding} Format \textsf{ALU\_DPWRR0.M} (Section~\ref{sec:format-ALU-DPWRR0.M}), \verb|exucode2 = 0100|\\

\bitfields{ 1/P, 2/00, 2/-- --, 27/upper27, 1/1, 2/11, 1/1, 4/0100, 5/registerM, 1/--, 2/10, 3/000, 1/0, 1/splat32, 5/lower5, 5/registerP, 1/0 }


\paragraph{Syntax} \texttt{mindp} \emph{registerM} = \emph{registerP}, \emph{upper27\_\discretionary{}{}{}lower5}\emph{splat32}

\paragraph{Description} The \emph{upper27\_\discretionary{}{}{}lower5} and the \emph{registerP} are interpreted as two double words packed into 128 bits. The minimum of the \emph{upper27\_\discretionary{}{}{}lower5} double words and the \emph{registerP} double words are computed. The two resulting double words are packed and stored into the \emph{registerM}.


\paragraph{Modifier(s)}
\noindent\begin{itemize}
\item splat32 (cf. splat32 in section \ref{par:modifier-splat32} on page \pageref{par:modifier-splat32})
\end{itemize}

\paragraph{Execution} {Reservation table \textsf{ALU\_LITE.X} (Section~\ref{sec:reservation-ALU-LITE.X})}
~
{\begin{lstlisting}
stage ID:
  new splat32 = _ZX_1(splat32);
stage RR:
  new argument3 = splat32 ? (_WX_32(upper27_lower5) << 32) | _ZX_32(_WX_32(upper27_lower5)) : _SX_32(_WX_32(upper27_lower5));
  new argument2 = PGR[registerP];
  for (I = 0; I < 2; I += 1) {
    result1.64[I] = _MIN(_SX_64(argument3.64[I]), _SX_64(argument2.64[I]))
  };
stage E1:
  PGR[registerM] = result1;
\end{lstlisting}}
\newpage
\subsubsection[{\small MINHO: Minimum Half Word Octuple}]{MINHO\hfill Minimum Half Word Octuple}
\label{instr:MINHO}\index{minho}
 
\paragraph{Encoding} Format \textsf{ALU\_HOWRR0} (Section~\ref{sec:format-ALU-HOWRR0}), \verb|exucode2 = 0100|\\

\bitfields{ 1/P, 2/11, 1/1, 4/0100, 5/registerM, 1/--, 2/10, 3/000, 1/1, 5/registerO, 1/--, 5/registerP, 1/0 }


\paragraph{Syntax} \texttt{minho} \emph{registerM} = \emph{registerP}, \emph{registerO}

\paragraph{Description} The \emph{registerO} and the \emph{registerP} are interpreted as eight half words packed into 128 bits. The minimum of the \emph{registerO} half words and the \emph{registerP} half words are computed. The eight resulting half words are packed and stored into the \emph{registerM}.


\paragraph{Execution} {Reservation table \textsf{ALU\_LITE} (Section~\ref{sec:reservation-ALU-LITE})}
~
{\begin{lstlisting}
stage RR:
  new argument3 = PGR[registerO];
  new argument2 = PGR[registerP];
  for (I = 0; I < 8; I += 1) {
    result1.16[I] = _MIN(_SX_16(argument3.16[I]), _SX_16(argument2.16[I]))
  };
stage E1:
  PGR[registerM] = result1;
\end{lstlisting}}
\newpage

\paragraph{Encoding} Format \textsf{ALU\_HOWRR0.M} (Section~\ref{sec:format-ALU-HOWRR0.M}), \verb|exucode2 = 0100|\\

\bitfields{ 1/P, 2/00, 2/-- --, 27/upper27, 1/1, 2/11, 1/1, 4/0100, 5/registerM, 1/--, 2/10, 3/000, 1/1, 1/splat32, 5/lower5, 5/registerP, 1/0 }


\paragraph{Syntax} \texttt{minho} \emph{registerM} = \emph{registerP}, \emph{upper27\_\discretionary{}{}{}lower5}\emph{splat32}

\paragraph{Description} The \emph{upper27\_\discretionary{}{}{}lower5} and the \emph{registerP} are interpreted as eight half words packed into 128 bits. The minimum of the \emph{upper27\_\discretionary{}{}{}lower5} half words and the \emph{registerP} half words are computed. The eight resulting half words are packed and stored into the \emph{registerM}.


\paragraph{Modifier(s)}
\noindent\begin{itemize}
\item splat32 (cf. splat32 in section \ref{par:modifier-splat32} on page \pageref{par:modifier-splat32})
\end{itemize}

\paragraph{Execution} {Reservation table \textsf{ALU\_LITE.X} (Section~\ref{sec:reservation-ALU-LITE.X})}
~
{\begin{lstlisting}
stage ID:
  new splat32 = _ZX_1(splat32);
stage RR:
  new argument3 = splat32 ? (_WX_32(upper27_lower5) << 32) | _ZX_32(_WX_32(upper27_lower5)) : _SX_32(_WX_32(upper27_lower5));
  new argument2 = PGR[registerP];
  for (I = 0; I < 8; I += 1) {
    result1.16[I] = _MIN(_SX_16(argument3.16[I]), _SX_16(argument2.16[I]))
  };
stage E1:
  PGR[registerM] = result1;
\end{lstlisting}}
\newpage
\subsubsection[{\small MINUBX: Minimum Unsigned Byte Hexadecuple}]{MINUBX\hfill Minimum Unsigned Byte Hexadecuple}
\label{instr:MINUBX}\index{minubx}
 
\paragraph{Encoding} Format \textsf{ALU\_BXWRR0} (Section~\ref{sec:format-ALU-BXWRR0}), \verb|exucode2 = 0110|\\

\bitfields{ 1/P, 2/11, 1/1, 4/0110, 5/registerM, 1/--, 2/10, 3/000, 1/1, 5/registerO, 1/--, 5/registerP, 1/1 }


\paragraph{Syntax} \texttt{minubx} \emph{registerM} = \emph{registerP}, \emph{registerO}

\paragraph{Description} The \emph{registerO} and the \emph{registerP} are interpreted as sixteen bytes packed into 128 bits. The unsigned minimum of the \emph{registerO} bytes and the \emph{registerP} bytes are computed. The sixteen resulting bytes are packed and stored into the \emph{registerM}.


\paragraph{Execution} {Reservation table \textsf{ALU\_LITE} (Section~\ref{sec:reservation-ALU-LITE})}
~
{\begin{lstlisting}
stage RR:
  new argument3 = PGR[registerO];
  new argument2 = PGR[registerP];
  for (I = 0; I < 16; I += 1) {
    result1.8[I] = _MIN(_ZX_8(argument3.8[I]), _ZX_8(argument2.8[I]))
  };
stage E1:
  PGR[registerM] = result1;
\end{lstlisting}}
\newpage

\paragraph{Encoding} Format \textsf{ALU\_BXWRR0.M} (Section~\ref{sec:format-ALU-BXWRR0.M}), \verb|exucode2 = 0110|\\

\bitfields{ 1/P, 2/00, 2/-- --, 27/upper27, 1/1, 2/11, 1/1, 4/0110, 5/registerM, 1/--, 2/10, 3/000, 1/1, 1/splat32, 5/lower5, 5/registerP, 1/1 }


\paragraph{Syntax} \texttt{minubx} \emph{registerM} = \emph{registerP}, \emph{upper27\_\discretionary{}{}{}lower5}\emph{splat32}

\paragraph{Description} The \emph{upper27\_\discretionary{}{}{}lower5} and the \emph{registerP} are interpreted as sixteen bytes packed into 128 bits. The unsigned minimum of the \emph{upper27\_\discretionary{}{}{}lower5} bytes and the \emph{registerP} bytes are computed. The sixteen resulting bytes are packed and stored into the \emph{registerM}.


\paragraph{Modifier(s)}
\noindent\begin{itemize}
\item splat32 (cf. splat32 in section \ref{par:modifier-splat32} on page \pageref{par:modifier-splat32})
\end{itemize}

\paragraph{Execution} {Reservation table \textsf{ALU\_LITE.X} (Section~\ref{sec:reservation-ALU-LITE.X})}
~
{\begin{lstlisting}
stage ID:
  new splat32 = _ZX_1(splat32);
stage RR:
  new argument3 = splat32 ? (_WX_32(upper27_lower5) << 32) | _ZX_32(_WX_32(upper27_lower5)) : _SX_32(_WX_32(upper27_lower5));
  new argument2 = PGR[registerP];
  for (I = 0; I < 16; I += 1) {
    result1.8[I] = _MIN(_ZX_8(argument3.8[I]), _ZX_8(argument2.8[I]))
  };
stage E1:
  PGR[registerM] = result1;
\end{lstlisting}}
\newpage
\subsubsection[{\small MINUD: Minimum Unsigned Double Word}]{MINUD\hfill Minimum Unsigned Double Word}
\label{instr:MINUD}\index{minud}
 
\paragraph{Encoding} Format \textsf{ALU\_DWRR0} (Section~\ref{sec:format-ALU-DWRR0}), \verb|exucode2 = 0110|\\

\bitfields{ 1/P, 2/11, 1/0, 4/0110, 6/registerW, 2/10, 3/000, 1/0, 6/registerY, 6/registerZ }


\paragraph{Syntax} \texttt{minud} \emph{registerW} = \emph{registerZ}, \emph{registerY}

\paragraph{Description} The unsigned minimum of the \emph{registerY} and the \emph{registerZ} is computed. The result is stored into the \emph{registerW}.


\paragraph{Execution} {Reservation table \textsf{ALU\_TINY} (Section~\ref{sec:reservation-ALU-TINY})}
~
{\begin{lstlisting}
stage RR:
  new argument3 = GPR[registerY];
  new argument2 = GPR[registerZ];
  new result1 = _MIN(_ZX_64(argument3), _ZX_64(argument2));
stage E1:
  GPR[registerW] = result1;
\end{lstlisting}}
\newpage

\paragraph{Encoding} Format \textsf{ALU\_DWRR0.M} (Section~\ref{sec:format-ALU-DWRR0.M}), \verb|exucode2 = 0110|\\

\bitfields{ 1/P, 2/00, 2/-- --, 27/upper27, 1/1, 2/11, 1/0, 4/0110, 6/registerW, 2/10, 3/000, 1/0, 1/splat32, 5/lower5, 6/registerZ }


\paragraph{Syntax} \texttt{minud} \emph{registerW} = \emph{registerZ}, \emph{upper27\_\discretionary{}{}{}lower5}\emph{splat32}

\paragraph{Description} The unsigned minimum of the \emph{upper27\_\discretionary{}{}{}lower5} and the \emph{registerZ} is computed. The result is stored into the \emph{registerW}.


\paragraph{Modifier(s)}
\noindent\begin{itemize}
\item splat32 (cf. splat32 in section \ref{par:modifier-splat32} on page \pageref{par:modifier-splat32})
\end{itemize}

\paragraph{Execution} {Reservation table \textsf{ALU\_TINY.X} (Section~\ref{sec:reservation-ALU-TINY.X})}
~
{\begin{lstlisting}
stage ID:
  new splat32 = _ZX_1(splat32);
stage RR:
  new argument3 = splat32 ? (_WX_32(upper27_lower5) << 32) | _ZX_32(_WX_32(upper27_lower5)) : _SX_32(_WX_32(upper27_lower5));
  new argument2 = GPR[registerZ];
  new result1 = _MIN(_ZX_64(argument3), _ZX_64(argument2));
stage E1:
  GPR[registerW] = result1;
\end{lstlisting}}
\newpage

\paragraph{Encoding} Format \textsf{ALU\_DWRI} (Section~\ref{sec:format-ALU-DWRI}), \verb|exucode2 = 0110|\\

\bitfields{ 1/P, 2/11, 1/0, 4/0110, 6/registerW, 2/00, 10/signed10, 6/registerZ }


\paragraph{Syntax} \texttt{minud} \emph{registerW} = \emph{registerZ}, \emph{signed10}

\paragraph{Description} The unsigned minimum of the \emph{signed10} and the \emph{registerZ} is computed. The result is stored into the \emph{registerW}.


\paragraph{Execution} {Reservation table \textsf{ALU\_TINY} (Section~\ref{sec:reservation-ALU-TINY})}
~
{\begin{lstlisting}
stage RR:
  new argument3 = _SX_10(signed10);
  new argument2 = GPR[registerZ];
  new result1 = _MIN(_ZX_64(argument3), _ZX_64(argument2));
stage E1:
  GPR[registerW] = result1;
\end{lstlisting}}
\newpage

\paragraph{Encoding} Format \textsf{ALU\_DWRI.X} (Section~\ref{sec:format-ALU-DWRI.X}), \verb|exucode2 = 0110|\\

\bitfields{ 1/P, 2/00, 2/-- --, 27/upper27, 1/1, 2/11, 1/0, 4/0110, 6/registerW, 2/00, 10/lower10, 6/registerZ }


\paragraph{Syntax} \texttt{minud} \emph{registerW} = \emph{registerZ}, \emph{upper27\_\discretionary{}{}{}lower10}

\paragraph{Description} The unsigned minimum of the \emph{upper27\_\discretionary{}{}{}lower10} and the \emph{registerZ} is computed. The result is stored into the \emph{registerW}.


\paragraph{Execution} {Reservation table \textsf{ALU\_TINY.X} (Section~\ref{sec:reservation-ALU-TINY.X})}
~
{\begin{lstlisting}
stage RR:
  new argument3 = _SX_37(upper27_lower10);
  new argument2 = GPR[registerZ];
  new result1 = _MIN(_ZX_64(argument3), _ZX_64(argument2));
stage E1:
  GPR[registerW] = result1;
\end{lstlisting}}
\newpage

\paragraph{Encoding} Format \textsf{ALU\_DWRI.Y} (Section~\ref{sec:format-ALU-DWRI.Y}), \verb|exucode2 = 0110|\\

\bitfields{ 1/P, 2/00, 2/-- --, 27/extend27, 1/1, 2/00, 2/-- --, 27/upper27, 1/1, 2/11, 1/0, 4/0110, 6/registerW, 2/00, 10/lower10, 6/registerZ }


\paragraph{Syntax} \texttt{minud} \emph{registerW} = \emph{registerZ}, \emph{extend27\_\discretionary{}{}{}upper27\_\discretionary{}{}{}lower10}

\paragraph{Description} The unsigned minimum of the \emph{extend27\_\discretionary{}{}{}upper27\_\discretionary{}{}{}lower10} and the \emph{registerZ} is computed. The result is stored into the \emph{registerW}.


\paragraph{Execution} {Reservation table \textsf{ALU\_TINY.Y} (Section~\ref{sec:reservation-ALU-TINY.Y})}
~
{\begin{lstlisting}
stage RR:
  new argument3 = _WX_64(extend27_upper27_lower10);
  new argument2 = GPR[registerZ];
  new result1 = _MIN(_ZX_64(argument3), _ZX_64(argument2));
stage E1:
  GPR[registerW] = result1;
\end{lstlisting}}
\newpage
\subsubsection[{\small MINUDP: Minimum Unsigned Double Word Pair}]{MINUDP\hfill Minimum Unsigned Double Word Pair}
\label{instr:MINUDP}\index{minudp}
 
\paragraph{Encoding} Format \textsf{ALU\_DPWRR0} (Section~\ref{sec:format-ALU-DPWRR0}), \verb|exucode2 = 0110|\\

\bitfields{ 1/P, 2/11, 1/1, 4/0110, 5/registerM, 1/--, 2/10, 3/000, 1/0, 5/registerO, 1/--, 5/registerP, 1/0 }


\paragraph{Syntax} \texttt{minudp} \emph{registerM} = \emph{registerP}, \emph{registerO}

\paragraph{Description} The \emph{registerO} and the \emph{registerP} are interpreted as two double words packed into 128 bits. The unsigned minimum of the \emph{registerO} words and the \emph{registerP} words are computed. The two resulting double words are packed and stored into the \emph{registerM}.


\paragraph{Execution} {Reservation table \textsf{ALU\_LITE} (Section~\ref{sec:reservation-ALU-LITE})}
~
{\begin{lstlisting}
stage RR:
  new argument3 = PGR[registerO];
  new argument2 = PGR[registerP];
  for (I = 0; I < 2; I += 1) {
    result1.64[I] = _MIN(_ZX_64(argument3.64[I]), _ZX_64(argument2.64[I]))
  };
stage E1:
  PGR[registerM] = result1;
\end{lstlisting}}
\newpage

\paragraph{Encoding} Format \textsf{ALU\_DPWRR0.M} (Section~\ref{sec:format-ALU-DPWRR0.M}), \verb|exucode2 = 0110|\\

\bitfields{ 1/P, 2/00, 2/-- --, 27/upper27, 1/1, 2/11, 1/1, 4/0110, 5/registerM, 1/--, 2/10, 3/000, 1/0, 1/splat32, 5/lower5, 5/registerP, 1/0 }


\paragraph{Syntax} \texttt{minudp} \emph{registerM} = \emph{registerP}, \emph{upper27\_\discretionary{}{}{}lower5}\emph{splat32}

\paragraph{Description} The \emph{upper27\_\discretionary{}{}{}lower5} and the \emph{registerP} are interpreted as two double words packed into 128 bits. The unsigned minimum of the \emph{upper27\_\discretionary{}{}{}lower5} words and the \emph{registerP} words are computed. The two resulting double words are packed and stored into the \emph{registerM}.


\paragraph{Modifier(s)}
\noindent\begin{itemize}
\item splat32 (cf. splat32 in section \ref{par:modifier-splat32} on page \pageref{par:modifier-splat32})
\end{itemize}

\paragraph{Execution} {Reservation table \textsf{ALU\_LITE.X} (Section~\ref{sec:reservation-ALU-LITE.X})}
~
{\begin{lstlisting}
stage ID:
  new splat32 = _ZX_1(splat32);
stage RR:
  new argument3 = splat32 ? (_WX_32(upper27_lower5) << 32) | _ZX_32(_WX_32(upper27_lower5)) : _SX_32(_WX_32(upper27_lower5));
  new argument2 = PGR[registerP];
  for (I = 0; I < 2; I += 1) {
    result1.64[I] = _MIN(_ZX_64(argument3.64[I]), _ZX_64(argument2.64[I]))
  };
stage E1:
  PGR[registerM] = result1;
\end{lstlisting}}
\newpage
\subsubsection[{\small MINUHO: Minimum Unsigned Half Word Octuple}]{MINUHO\hfill Minimum Unsigned Half Word Octuple}
\label{instr:MINUHO}\index{minuho}
 
\paragraph{Encoding} Format \textsf{ALU\_HOWRR0} (Section~\ref{sec:format-ALU-HOWRR0}), \verb|exucode2 = 0110|\\

\bitfields{ 1/P, 2/11, 1/1, 4/0110, 5/registerM, 1/--, 2/10, 3/000, 1/1, 5/registerO, 1/--, 5/registerP, 1/0 }


\paragraph{Syntax} \texttt{minuho} \emph{registerM} = \emph{registerP}, \emph{registerO}

\paragraph{Description} The \emph{registerO} and the \emph{registerP} are interpreted as eight half words packed into 128 bits. The unsigned minimum of the \emph{registerO} half words and the \emph{registerP} half words are computed. The eight resulting half words are packed and stored into the \emph{registerM}.


\paragraph{Execution} {Reservation table \textsf{ALU\_LITE} (Section~\ref{sec:reservation-ALU-LITE})}
~
{\begin{lstlisting}
stage RR:
  new argument3 = PGR[registerO];
  new argument2 = PGR[registerP];
  for (I = 0; I < 8; I += 1) {
    result1.16[I] = _MIN(_ZX_16(argument3.16[I]), _ZX_16(argument2.16[I]))
  };
stage E1:
  PGR[registerM] = result1;
\end{lstlisting}}
\newpage

\paragraph{Encoding} Format \textsf{ALU\_HOWRR0.M} (Section~\ref{sec:format-ALU-HOWRR0.M}), \verb|exucode2 = 0110|\\

\bitfields{ 1/P, 2/00, 2/-- --, 27/upper27, 1/1, 2/11, 1/1, 4/0110, 5/registerM, 1/--, 2/10, 3/000, 1/1, 1/splat32, 5/lower5, 5/registerP, 1/0 }


\paragraph{Syntax} \texttt{minuho} \emph{registerM} = \emph{registerP}, \emph{upper27\_\discretionary{}{}{}lower5}\emph{splat32}

\paragraph{Description} The \emph{upper27\_\discretionary{}{}{}lower5} and the \emph{registerP} are interpreted as eight half words packed into 128 bits. The unsigned minimum of the \emph{upper27\_\discretionary{}{}{}lower5} half words and the \emph{registerP} half words are computed. The eight resulting half words are packed and stored into the \emph{registerM}.


\paragraph{Modifier(s)}
\noindent\begin{itemize}
\item splat32 (cf. splat32 in section \ref{par:modifier-splat32} on page \pageref{par:modifier-splat32})
\end{itemize}

\paragraph{Execution} {Reservation table \textsf{ALU\_LITE.X} (Section~\ref{sec:reservation-ALU-LITE.X})}
~
{\begin{lstlisting}
stage ID:
  new splat32 = _ZX_1(splat32);
stage RR:
  new argument3 = splat32 ? (_WX_32(upper27_lower5) << 32) | _ZX_32(_WX_32(upper27_lower5)) : _SX_32(_WX_32(upper27_lower5));
  new argument2 = PGR[registerP];
  for (I = 0; I < 8; I += 1) {
    result1.16[I] = _MIN(_ZX_16(argument3.16[I]), _ZX_16(argument2.16[I]))
  };
stage E1:
  PGR[registerM] = result1;
\end{lstlisting}}
\newpage
\subsubsection[{\small MINUW: Minimum Unsigned Word}]{MINUW\hfill Minimum Unsigned Word}
\label{instr:MINUW}\index{minuw}
 
\paragraph{Encoding} Format \textsf{ALU\_WWRR0} (Section~\ref{sec:format-ALU-WWRR0}), \verb|exucode2 = 0110|\\

\bitfields{ 1/P, 2/11, 1/0, 4/0110, 6/registerW, 2/11, 3/000, 1/signextw, 6/registerY, 6/registerZ }


\paragraph{Syntax} \texttt{minuw}\emph{signextw} \emph{registerW} = \emph{registerZ}, \emph{registerY}

\paragraph{Description} The unsigned minimum of the \emph{registerY} and the \emph{registerZ} is computed. The result with the 32 upper bits cleared is stored into the \emph{registerW}.


\paragraph{Modifier(s)}
\noindent\begin{itemize}
\item signextw (cf. signextw in section \ref{par:modifier-signextw} on page \pageref{par:modifier-signextw})
\end{itemize}

\paragraph{Execution} {Reservation table \textsf{ALU\_TINY} (Section~\ref{sec:reservation-ALU-TINY})}
~
{\begin{lstlisting}
stage ID:
  new signextw = _ZX_1(signextw);
stage RR:
  new argument3 = GPR[registerY];
  new argument2 = GPR[registerZ];
  new result1 = _MIN(argument3.32[0], argument2.32[0]);
stage E1:
  GPR[registerW] = signextw ? _SX_32(result1) : _ZX_32(result1);
\end{lstlisting}}
\newpage

\paragraph{Encoding} Format \textsf{ALU\_WWRR0.W} (Section~\ref{sec:format-ALU-WWRR0.W}), \verb|exucode2 = 0110|\\

\bitfields{ 1/P, 2/00, 2/-- --, 27/upper27, 1/1, 2/11, 1/0, 4/0110, 6/registerW, 2/11, 3/000, 1/signextw, 1/0, 5/lower5, 6/registerZ }


\paragraph{Syntax} \texttt{minuw}\emph{signextw} \emph{registerW} = \emph{registerZ}, \emph{upper27\_\discretionary{}{}{}lower5}

\paragraph{Description} The unsigned minimum of the \emph{upper27\_\discretionary{}{}{}lower5} and the \emph{registerZ} is computed. The result with the 32 upper bits cleared is stored into the \emph{registerW}.


\paragraph{Modifier(s)}
\noindent\begin{itemize}
\item signextw (cf. signextw in section \ref{par:modifier-signextw} on page \pageref{par:modifier-signextw})
\end{itemize}

\paragraph{Execution} {Reservation table \textsf{ALU\_TINY.X} (Section~\ref{sec:reservation-ALU-TINY.X})}
~
{\begin{lstlisting}
stage ID:
  new signextw = _ZX_1(signextw);
stage RR:
  new argument3 = _WX_32(upper27_lower5);
  new argument2 = GPR[registerZ];
  new result1 = _MIN(argument3.32[0], argument2.32[0]);
stage E1:
  GPR[registerW] = signextw ? _SX_32(result1) : _ZX_32(result1);
\end{lstlisting}}
\newpage
\subsubsection[{\small MINUWQ: Minimum Unsigned Word Quadruple}]{MINUWQ\hfill Minimum Unsigned Word Quadruple}
\label{instr:MINUWQ}\index{minuwq}
 
\paragraph{Encoding} Format \textsf{ALU\_WQWRR0} (Section~\ref{sec:format-ALU-WQWRR0}), \verb|exucode2 = 0110|\\

\bitfields{ 1/P, 2/11, 1/1, 4/0110, 5/registerM, 1/--, 2/10, 3/000, 1/0, 5/registerO, 1/--, 5/registerP, 1/1 }


\paragraph{Syntax} \texttt{minuwq} \emph{registerM} = \emph{registerP}, \emph{registerO}

\paragraph{Description} The \emph{registerO} and the \emph{registerP} are interpreted as four words packed into 128 bits. The unsigned minimum of the \emph{registerO} words and the \emph{registerP} words are computed. The four resulting words are packed and stored into the \emph{registerM}.


\paragraph{Execution} {Reservation table \textsf{ALU\_LITE} (Section~\ref{sec:reservation-ALU-LITE})}
~
{\begin{lstlisting}
stage RR:
  new argument3 = PGR[registerO];
  new argument2 = PGR[registerP];
  for (I = 0; I < 4; I += 1) {
    result1.32[I] = _MIN(_ZX_32(argument3.32[I]), _ZX_32(argument2.32[I]))
  };
stage E1:
  PGR[registerM] = result1;
\end{lstlisting}}
\newpage

\paragraph{Encoding} Format \textsf{ALU\_WQWRR0.M} (Section~\ref{sec:format-ALU-WQWRR0.M}), \verb|exucode2 = 0110|\\

\bitfields{ 1/P, 2/00, 2/-- --, 27/upper27, 1/1, 2/11, 1/1, 4/0110, 5/registerM, 1/--, 2/10, 3/000, 1/0, 1/splat32, 5/lower5, 5/registerP, 1/1 }


\paragraph{Syntax} \texttt{minuwq} \emph{registerM} = \emph{registerP}, \emph{upper27\_\discretionary{}{}{}lower5}\emph{splat32}

\paragraph{Description} The \emph{upper27\_\discretionary{}{}{}lower5} and the \emph{registerP} are interpreted as four words packed into 128 bits. The unsigned minimum of the \emph{upper27\_\discretionary{}{}{}lower5} words and the \emph{registerP} words are computed. The four resulting words are packed and stored into the \emph{registerM}.


\paragraph{Modifier(s)}
\noindent\begin{itemize}
\item splat32 (cf. splat32 in section \ref{par:modifier-splat32} on page \pageref{par:modifier-splat32})
\end{itemize}

\paragraph{Execution} {Reservation table \textsf{ALU\_LITE.X} (Section~\ref{sec:reservation-ALU-LITE.X})}
~
{\begin{lstlisting}
stage ID:
  new splat32 = _ZX_1(splat32);
stage RR:
  new argument3 = splat32 ? (_WX_32(upper27_lower5) << 32) | _ZX_32(_WX_32(upper27_lower5)) : _SX_32(_WX_32(upper27_lower5));
  new argument2 = PGR[registerP];
  for (I = 0; I < 4; I += 1) {
    result1.32[I] = _MIN(_ZX_32(argument3.32[I]), _ZX_32(argument2.32[I]))
  };
stage E1:
  PGR[registerM] = result1;
\end{lstlisting}}
\newpage
\subsubsection[{\small MINW: Minimum Word}]{MINW\hfill Minimum Word}
\label{instr:MINW}\index{minw}
 
\paragraph{Encoding} Format \textsf{ALU\_WWRR0} (Section~\ref{sec:format-ALU-WWRR0}), \verb|exucode2 = 0100|\\

\bitfields{ 1/P, 2/11, 1/0, 4/0100, 6/registerW, 2/11, 3/000, 1/signextw, 6/registerY, 6/registerZ }


\paragraph{Syntax} \texttt{minw}\emph{signextw} \emph{registerW} = \emph{registerZ}, \emph{registerY}

\paragraph{Description} The minimum of the \emph{registerY} and the \emph{registerZ} is computed. The result with the 32 upper bits cleared is stored into the \emph{registerW}.


\paragraph{Modifier(s)}
\noindent\begin{itemize}
\item signextw (cf. signextw in section \ref{par:modifier-signextw} on page \pageref{par:modifier-signextw})
\end{itemize}

\paragraph{Execution} {Reservation table \textsf{ALU\_TINY} (Section~\ref{sec:reservation-ALU-TINY})}
~
{\begin{lstlisting}
stage ID:
  new signextw = _ZX_1(signextw);
stage RR:
  new argument3 = GPR[registerY];
  new argument2 = GPR[registerZ];
  new result1 = _MIN(_SX_32(argument3), _SX_32(argument2));
stage E1:
  GPR[registerW] = signextw ? _SX_32(result1) : _ZX_32(result1);
\end{lstlisting}}
\newpage

\paragraph{Encoding} Format \textsf{ALU\_WWRR0.W} (Section~\ref{sec:format-ALU-WWRR0.W}), \verb|exucode2 = 0100|\\

\bitfields{ 1/P, 2/00, 2/-- --, 27/upper27, 1/1, 2/11, 1/0, 4/0100, 6/registerW, 2/11, 3/000, 1/signextw, 1/0, 5/lower5, 6/registerZ }


\paragraph{Syntax} \texttt{minw}\emph{signextw} \emph{registerW} = \emph{registerZ}, \emph{upper27\_\discretionary{}{}{}lower5}

\paragraph{Description} The minimum of the \emph{upper27\_\discretionary{}{}{}lower5} and the \emph{registerZ} is computed. The result with the 32 upper bits cleared is stored into the \emph{registerW}.


\paragraph{Modifier(s)}
\noindent\begin{itemize}
\item signextw (cf. signextw in section \ref{par:modifier-signextw} on page \pageref{par:modifier-signextw})
\end{itemize}

\paragraph{Execution} {Reservation table \textsf{ALU\_TINY.X} (Section~\ref{sec:reservation-ALU-TINY.X})}
~
{\begin{lstlisting}
stage ID:
  new signextw = _ZX_1(signextw);
stage RR:
  new argument3 = _WX_32(upper27_lower5);
  new argument2 = GPR[registerZ];
  new result1 = _MIN(_SX_32(argument3), _SX_32(argument2));
stage E1:
  GPR[registerW] = signextw ? _SX_32(result1) : _ZX_32(result1);
\end{lstlisting}}
\newpage
\subsubsection[{\small MINWQ: Minimum Word Quadruple}]{MINWQ\hfill Minimum Word Quadruple}
\label{instr:MINWQ}\index{minwq}
 
\paragraph{Encoding} Format \textsf{ALU\_WQWRR0} (Section~\ref{sec:format-ALU-WQWRR0}), \verb|exucode2 = 0100|\\

\bitfields{ 1/P, 2/11, 1/1, 4/0100, 5/registerM, 1/--, 2/10, 3/000, 1/0, 5/registerO, 1/--, 5/registerP, 1/1 }


\paragraph{Syntax} \texttt{minwq} \emph{registerM} = \emph{registerP}, \emph{registerO}

\paragraph{Description} The \emph{registerO} and the \emph{registerP} are interpreted as four words packed into 128 bits. The minimum of the \emph{registerO} words and the \emph{registerP} words are computed. The four resulting words are packed and stored into the \emph{registerM}.


\paragraph{Execution} {Reservation table \textsf{ALU\_LITE} (Section~\ref{sec:reservation-ALU-LITE})}
~
{\begin{lstlisting}
stage RR:
  new argument3 = PGR[registerO];
  new argument2 = PGR[registerP];
  for (I = 0; I < 4; I += 1) {
    result1.32[I] = _MIN(_SX_32(argument3.32[I]), _SX_32(argument2.32[I]))
  };
stage E1:
  PGR[registerM] = result1;
\end{lstlisting}}
\newpage

\paragraph{Encoding} Format \textsf{ALU\_WQWRR0.M} (Section~\ref{sec:format-ALU-WQWRR0.M}), \verb|exucode2 = 0100|\\

\bitfields{ 1/P, 2/00, 2/-- --, 27/upper27, 1/1, 2/11, 1/1, 4/0100, 5/registerM, 1/--, 2/10, 3/000, 1/0, 1/splat32, 5/lower5, 5/registerP, 1/1 }


\paragraph{Syntax} \texttt{minwq} \emph{registerM} = \emph{registerP}, \emph{upper27\_\discretionary{}{}{}lower5}\emph{splat32}

\paragraph{Description} The \emph{upper27\_\discretionary{}{}{}lower5} and the \emph{registerP} are interpreted as four words packed into 128 bits. The minimum of the \emph{upper27\_\discretionary{}{}{}lower5} words and the \emph{registerP} words are computed. The four resulting words are packed and stored into the \emph{registerM}.


\paragraph{Modifier(s)}
\noindent\begin{itemize}
\item splat32 (cf. splat32 in section \ref{par:modifier-splat32} on page \pageref{par:modifier-splat32})
\end{itemize}

\paragraph{Execution} {Reservation table \textsf{ALU\_LITE.X} (Section~\ref{sec:reservation-ALU-LITE.X})}
~
{\begin{lstlisting}
stage ID:
  new splat32 = _ZX_1(splat32);
stage RR:
  new argument3 = splat32 ? (_WX_32(upper27_lower5) << 32) | _ZX_32(_WX_32(upper27_lower5)) : _SX_32(_WX_32(upper27_lower5));
  new argument2 = PGR[registerP];
  for (I = 0; I < 4; I += 1) {
    result1.32[I] = _MIN(_SX_32(argument3.32[I]), _SX_32(argument2.32[I]))
  };
stage E1:
  PGR[registerM] = result1;
\end{lstlisting}}
\newpage
\subsubsection[{\small MSBFBHO: Multiply Subtract From Extend Byte to Half Word Octuple}]{MSBFBHO\hfill Multiply Subtract From Extend Byte to Half Word Octuple}
\label{instr:MSBFBHO}\index{msbfbho}
 
\paragraph{Encoding} Format \textsf{ALU\_HOMEWRRRE} (Section~\ref{sec:format-ALU-HOMEWRRRE}), \verb|alucode2 = 11|\\

\bitfields{ 1/P, 2/11, 1/1, 2/11, 2/widemult, 5/registerM, 1/0, 2/10, 3/110, 1/1, 5/registerYe, 1/1, 5/registerZe, 1/0 }


\paragraph{Syntax} \texttt{msbfbho}\emph{widemult} \emph{registerM} = \emph{registerZe}, \emph{registerYe}

\paragraph{Description} The \emph{registerYe} and the \emph{registerZe} are interpreted as eight bytes packed into 64 bits. These \emph{registerYe} and \emph{registerZe} bytes are sign-extended or zero-extended from 8 bits, depending on the \emph{widemult}. The \emph{registerYe} bytes are multiplied by the \emph{registerZe} bytes, producing eight 16-bit results. The eight resulting half words are packed, subtracted from and stored into the \emph{registerM}.


\paragraph{Modifier(s)}
\noindent\begin{itemize}
\item widemult (cf. widemult in section \ref{par:modifier-widemult} on page \pageref{par:modifier-widemult})
\end{itemize}

\paragraph{Execution} {Reservation table \textsf{ALU\_LITE} (Section~\ref{sec:reservation-ALU-LITE})}
~
{\begin{lstlisting}
stage ID:
  new widemult = _ZX_2(widemult);
stage RR:
  new argument3 = GPR[registerYe<<1];
  new argument2 = GPR[registerZe<<1];
stage E1:
  new argument1 = PGR[registerM];
  for (I = 0; I < 8; I += 1) {
    if (widemult == 0) {
      new product = _SX_8(argument2.8[I]) * _SX_8(argument3.8[I]);
    }
    if (widemult == 1) {
      product = _ZX_8(argument2.8[I]) * _ZX_8(argument3.8[I]);
    }
    if (widemult == 2) {
      product = _SX_8(argument2.8[I]) * _ZX_8(argument3.8[I]);
    }
    result1.16[I] = argument1.16[I] - product
  };
stage E2:
  PGR[registerM] = result1;
\end{lstlisting}}
\newpage

\paragraph{Encoding} Format \textsf{ALU\_HOMEWRRRO} (Section~\ref{sec:format-ALU-HOMEWRRRO}), \verb|alucode2 = 11|\\

\bitfields{ 1/P, 2/11, 1/1, 2/11, 2/widemult, 5/registerM, 1/1, 2/10, 3/110, 1/1, 5/registerYo, 1/1, 5/registerZo, 1/0 }


\paragraph{Syntax} \texttt{msbfbho}\emph{widemult} \emph{registerM} = \emph{registerZo}, \emph{registerYo}

\paragraph{Description} The \emph{registerYo} and the \emph{registerZo} are interpreted as eight bytes packed into 64 bits. These \emph{registerYo} and \emph{registerZo} bytes are sign-extended or zero-extended from 8 bits, depending on the \emph{widemult}. The \emph{registerYo} bytes are multiplied by the \emph{registerZo} bytes, producing eight 16-bit results. The eight resulting half words are packed, subtracted from and stored into the \emph{registerM}.


\paragraph{Modifier(s)}
\noindent\begin{itemize}
\item widemult (cf. widemult in section \ref{par:modifier-widemult} on page \pageref{par:modifier-widemult})
\end{itemize}

\paragraph{Execution} {Reservation table \textsf{ALU\_LITE} (Section~\ref{sec:reservation-ALU-LITE})}
~
{\begin{lstlisting}
stage ID:
  new widemult = _ZX_2(widemult);
stage RR:
  new argument3 = GPR[(registerYo<<1)+1];
  new argument2 = GPR[(registerZo<<1)+1];
stage E1:
  new argument1 = PGR[registerM];
  for (I = 0; I < 8; I += 1) {
    if (widemult == 0) {
      new product = _SX_8(argument2.8[I]) * _SX_8(argument3.8[I]);
    }
    if (widemult == 1) {
      product = _ZX_8(argument2.8[I]) * _ZX_8(argument3.8[I]);
    }
    if (widemult == 2) {
      product = _SX_8(argument2.8[I]) * _ZX_8(argument3.8[I]);
    }
    result1.16[I] = argument1.16[I] - product
  };
stage E2:
  PGR[registerM] = result1;
\end{lstlisting}}
\newpage
\subsubsection[{\small MSBFD: Multiply Subtract From Double Word}]{MSBFD\hfill Multiply Subtract From Double Word}
\label{instr:MSBFD}\index{msbfd}
 
\paragraph{Encoding} Format \textsf{ALU\_DMWRRR} (Section~\ref{sec:format-ALU-DMWRRR}), \verb|alucode2 = 11|\\

\bitfields{ 1/P, 2/11, 1/0, 2/11, 2/highmult, 6/registerW, 2/10, 3/110, 1/0, 6/registerY, 6/registerZ }


\paragraph{Syntax} \texttt{msbfd}\emph{highmult} \emph{registerW} = \emph{registerZ}, \emph{registerY}

\paragraph{Description} The \emph{registerY} and the \emph{registerZ} double words are sign-extended or zero-extended from 64 bits, depending on the \emph{highmult}. The \emph{registerY} double word is multiplied by the \emph{registerZ} double word, and the product is optionally shifted right by 64 bits, depending on the \emph{highmult}. The resulting double word is subtracted from and stored into the \emph{registerW}.


\paragraph{Modifier(s)}
\noindent\begin{itemize}
\item highmult (cf. highmult in section \ref{par:modifier-highmult} on page \pageref{par:modifier-highmult})
\end{itemize}

\paragraph{Execution} {Reservation table \textsf{ALU\_LITE} (Section~\ref{sec:reservation-ALU-LITE})}
~
{\begin{lstlisting}
stage ID:
  new highmult = _ZX_2(highmult);
stage RR:
  new argument3 = GPR[registerY];
  new argument2 = GPR[registerZ];
stage E1:
  new argument1 = GPR[registerW];
  if ((highmult == 0) || (highmult == 3)) {
    new product = _SX_64(argument2) * _SX_64(argument3);
  }
  if (highmult == 1) {
    product = _ZX_64(argument2) * _ZX_64(argument3);
  }
  if (highmult == 2) {
    product = _SX_64(argument2) * _ZX_64(argument3);
  }
  if (highmult) {
    new result1 = argument1 - product.64[1];
  } else {
    result1 = argument1 - product.64[0];
  }
stage E2:
  GPR[registerW] = result1;
\end{lstlisting}}
\newpage

\paragraph{Encoding} Format \textsf{ALU\_DMWRRR.M} (Section~\ref{sec:format-ALU-DMWRRR.M}), \verb|alucode2 = 11|\\

\bitfields{ 1/P, 2/00, 2/-- --, 27/upper27, 1/1, 2/11, 1/0, 2/11, 2/highmult, 6/registerW, 2/10, 3/110, 1/0, 1/splat32, 5/lower5, 6/registerZ }


\paragraph{Syntax} \texttt{msbfd}\emph{highmult} \emph{registerW} = \emph{registerZ}, \emph{upper27\_\discretionary{}{}{}lower5}\emph{splat32}

\paragraph{Description} The \emph{upper27\_\discretionary{}{}{}lower5} and the \emph{registerZ} double words are sign-extended or zero-extended from 64 bits, depending on the \emph{highmult}. The \emph{upper27\_\discretionary{}{}{}lower5} double word is multiplied by the \emph{registerZ} double word, and the product is optionally shifted right by 64 bits, depending on the \emph{highmult}. The resulting double word is subtracted from and stored into the \emph{registerW}.


\paragraph{Modifier(s)}
\noindent\begin{itemize}
\item highmult (cf. highmult in section \ref{par:modifier-highmult} on page \pageref{par:modifier-highmult})
\item splat32 (cf. splat32 in section \ref{par:modifier-splat32} on page \pageref{par:modifier-splat32})
\end{itemize}

\paragraph{Execution} {Reservation table \textsf{ALU\_LITE.X} (Section~\ref{sec:reservation-ALU-LITE.X})}
~
{\begin{lstlisting}
stage ID:
  new splat32 = _ZX_1(splat32);
  new highmult = _ZX_2(highmult);
stage RR:
  new argument3 = splat32 ? (_WX_32(upper27_lower5) << 32) | _ZX_32(_WX_32(upper27_lower5)) : _SX_32(_WX_32(upper27_lower5));
  new argument2 = GPR[registerZ];
stage E1:
  new argument1 = GPR[registerW];
  if ((highmult == 0) || (highmult == 3)) {
    new product = _SX_64(argument2) * _SX_64(argument3);
  }
  if (highmult == 1) {
    product = _ZX_64(argument2) * _ZX_64(argument3);
  }
  if (highmult == 2) {
    product = _SX_64(argument2) * _ZX_64(argument3);
  }
  if (highmult) {
    new result1 = argument1 - product.64[1];
  } else {
    result1 = argument1 - product.64[0];
  }
stage E2:
  GPR[registerW] = result1;
\end{lstlisting}}
\newpage
\subsubsection[{\small MSBFDP: Multiply Subtract From Double Word Pair}]{MSBFDP\hfill Multiply Subtract From Double Word Pair}
\label{instr:MSBFDP}\index{msbfdp}
 
\paragraph{Encoding} Format \textsf{ALU\_DPMWRRR} (Section~\ref{sec:format-ALU-DPMWRRR}), \verb|alucode2 = 11|\\

\bitfields{ 1/P, 2/11, 1/1, 2/11, 2/highmult, 5/registerM, 1/--, 2/10, 3/110, 1/0, 5/registerO, 1/0, 5/registerP, 1/0 }


\paragraph{Syntax} \texttt{msbfdp}\emph{highmult} \emph{registerM} = \emph{registerP}, \emph{registerO}

\paragraph{Description} The \emph{registerO} and the \emph{registerP} are interpreted as two double words packed into 128 bits. The \emph{registerO} and the \emph{registerP} double words are sign-extended or zero-extended from 64 bits, depending on the \emph{highmult}. The \emph{registerO} double words are multiplied by the \emph{registerP} double words, and the product is optionally shifted right by 64 bits, depending on the \emph{highmult}. The two resulting double words are subtracted from the \emph{registerM} double words, packed and stored back into the \emph{registerM}.


\paragraph{Modifier(s)}
\noindent\begin{itemize}
\item highmult (cf. highmult in section \ref{par:modifier-highmult} on page \pageref{par:modifier-highmult})
\end{itemize}

\paragraph{Execution} {Reservation table \textsf{ALU\_LITE} (Section~\ref{sec:reservation-ALU-LITE})}
~
{\begin{lstlisting}
stage ID:
  new highmult = _ZX_2(highmult);
stage RR:
  new argument3 = PGR[registerO];
  new argument2 = PGR[registerP];
stage E1:
  new argument1 = PGR[registerM];
  for (I = 0; I < 2; I += 1) {
    if ((highmult == 0) || (highmult == 3)) {
      new product = _SX_64(argument2.64[I]) * _SX_64(argument3.64[I]);
    }
    if (highmult == 1) {
      product = _ZX_64(argument2.64[I]) * _ZX_64(argument3.64[I]);
    }
    if (highmult == 2) {
      product = _SX_64(argument2.64[I]) * _ZX_64(argument3.64[I]);
    }
    if (highmult) {
      result1.64[I] = argument1.64[I] - product.64[1];
    } else {
      result1.64[I] = argument1.64[I] - product.64[0];
    }
  };
stage E2:
  PGR[registerM] = result1;
\end{lstlisting}}
\newpage
\subsubsection[{\small MSBFDQ: Multiply Subtract From Extend Double Word to Quadruple Word}]{MSBFDQ\hfill Multiply Subtract From Extend Double Word to Quadruple Word}
\label{instr:MSBFDQ}\index{msbfdq}
 
\paragraph{Encoding} Format \textsf{ALU\_QMEWRRR} (Section~\ref{sec:format-ALU-QMEWRRR}), \verb|alucode2 = 11|\\

\bitfields{ 1/P, 2/11, 1/0, 2/11, 2/widemult, 5/registerM, 1/--, 2/10, 3/111, 1/1, 6/registerY, 6/registerZ }


\paragraph{Syntax} \texttt{msbfdq}\emph{widemult} \emph{registerM} = \emph{registerZ}, \emph{registerY}

\paragraph{Description} The \emph{registerY} and the \emph{registerZ} double words are sign-extended or zero-extended from 64 bits, depending on the \emph{widemult}. The \emph{registerY} is multiplied by the \emph{registerZ}, producing a 128-bit result. The resulting quadruple word subtracted from and stored into the \emph{registerM}.


\paragraph{Modifier(s)}
\noindent\begin{itemize}
\item widemult (cf. widemult in section \ref{par:modifier-widemult} on page \pageref{par:modifier-widemult})
\end{itemize}

\paragraph{Execution} {Reservation table \textsf{ALU\_LITE} (Section~\ref{sec:reservation-ALU-LITE})}
~
{\begin{lstlisting}
stage ID:
  new widemult = _ZX_2(widemult);
stage RR:
  new argument3 = GPR[registerY];
  new argument2 = GPR[registerZ];
stage E1:
  new argument1 = PGR[registerM];
  if (widemult == 0) {
    new product = _SX_64(argument2) * _SX_64(argument3);
  }
  if (widemult == 1) {
    product = _ZX_64(argument2) * _ZX_64(argument3);
  }
  if (widemult == 2) {
    product = _SX_64(argument2) * _ZX_64(argument3);
  }
  new result1 = argument1 - product;
stage E2:
  PGR[registerM] = result1;
\end{lstlisting}}
\newpage
\subsubsection[{\small MSBFHWQ: Multiply Subtract From Extend Half Word to Word Quadruple}]{MSBFHWQ\hfill Multiply Subtract From Extend Half Word to Word Quadruple}
\label{instr:MSBFHWQ}\index{msbfhwq}
 
\paragraph{Encoding} Format \textsf{ALU\_WQMEWRRRE} (Section~\ref{sec:format-ALU-WQMEWRRRE}), \verb|alucode2 = 11|\\

\bitfields{ 1/P, 2/11, 1/1, 2/11, 2/widemult, 5/registerM, 1/0, 2/10, 3/110, 1/0, 5/registerYe, 1/1, 5/registerZe, 1/1 }


\paragraph{Syntax} \texttt{msbfhwq}\emph{widemult} \emph{registerM} = \emph{registerZe}, \emph{registerYe}

\paragraph{Description} The \emph{registerYe} and the \emph{registerZe} are interpreted as eight half words packed into 128 bits. These \emph{registerYe} and \emph{registerZe} half words are sign-extended or zero-extended from 16 bits, depending on the \emph{widemult}. The \emph{registerYe} half words are multiplied by the \emph{registerZe} half words, producing four 32-bit results. The four resulting words are packed, subtracted from and stored into the \emph{registerM}.


\paragraph{Modifier(s)}
\noindent\begin{itemize}
\item widemult (cf. widemult in section \ref{par:modifier-widemult} on page \pageref{par:modifier-widemult})
\end{itemize}

\paragraph{Execution} {Reservation table \textsf{ALU\_LITE} (Section~\ref{sec:reservation-ALU-LITE})}
~
{\begin{lstlisting}
stage ID:
  new widemult = _ZX_2(widemult);
stage RR:
  new argument3 = GPR[registerYe<<1];
  new argument2 = GPR[registerZe<<1];
stage E1:
  new argument1 = PGR[registerM];
  for (I = 0; I < 4; I += 1) {
    if (widemult == 0) {
      new product = _SX_16(argument2.16[I]) * _SX_16(argument3.16[I]);
    }
    if (widemult == 1) {
      product = _ZX_16(argument2.16[I]) * _ZX_16(argument3.16[I]);
    }
    if (widemult == 2) {
      product = _SX_16(argument2.16[I]) * _ZX_16(argument3.16[I]);
    }
    result1.32[I] = argument1.32[I] - product
  };
stage E2:
  PGR[registerM] = result1;
\end{lstlisting}}
\newpage

\paragraph{Encoding} Format \textsf{ALU\_WQMEWRRRO} (Section~\ref{sec:format-ALU-WQMEWRRRO}), \verb|alucode2 = 11|\\

\bitfields{ 1/P, 2/11, 1/1, 2/11, 2/widemult, 5/registerM, 1/1, 2/10, 3/110, 1/0, 5/registerYo, 1/1, 5/registerZo, 1/1 }


\paragraph{Syntax} \texttt{msbfhwq}\emph{widemult} \emph{registerM} = \emph{registerZo}, \emph{registerYo}

\paragraph{Description} The \emph{registerYo} and the \emph{registerZo} are interpreted as eight half words packed into 128 bits. These \emph{registerYo} and \emph{registerZo} half words are sign-extended or zero-extended from 16 bits, depending on the \emph{widemult}. The \emph{registerYo} half words are multiplied by the \emph{registerZo} half words, producing four 32-bit results. The four resulting words are packed, subtracted from and stored into the \emph{registerM}.


\paragraph{Modifier(s)}
\noindent\begin{itemize}
\item widemult (cf. widemult in section \ref{par:modifier-widemult} on page \pageref{par:modifier-widemult})
\end{itemize}

\paragraph{Execution} {Reservation table \textsf{ALU\_LITE} (Section~\ref{sec:reservation-ALU-LITE})}
~
{\begin{lstlisting}
stage ID:
  new widemult = _ZX_2(widemult);
stage RR:
  new argument3 = GPR[(registerYo<<1)+1];
  new argument2 = GPR[(registerZo<<1)+1];
stage E1:
  new argument1 = PGR[registerM];
  for (I = 0; I < 4; I += 1) {
    if (widemult == 0) {
      new product = _SX_16(argument2.16[I]) * _SX_16(argument3.16[I]);
    }
    if (widemult == 1) {
      product = _ZX_16(argument2.16[I]) * _ZX_16(argument3.16[I]);
    }
    if (widemult == 2) {
      product = _SX_16(argument2.16[I]) * _ZX_16(argument3.16[I]);
    }
    result1.32[I] = argument1.32[I] - product
  };
stage E2:
  PGR[registerM] = result1;
\end{lstlisting}}
\newpage
\subsubsection[{\small MSBFW: Multiply Subtract From Word}]{MSBFW\hfill Multiply Subtract From Word}
\label{instr:MSBFW}\index{msbfw}
 
\paragraph{Encoding} Format \textsf{ALU\_WMWRRR} (Section~\ref{sec:format-ALU-WMWRRR}), \verb|alucode2 = 11|\\

\bitfields{ 1/P, 2/11, 1/0, 2/11, 2/highmult, 6/registerW, 2/11, 3/110, 1/signextw, 6/registerY, 6/registerZ }


\paragraph{Syntax} \texttt{msbfw}\emph{highmult}\emph{signextw} \emph{registerW} = \emph{registerZ}, \emph{registerY}

\paragraph{Description} The \emph{registerY} and the \emph{registerZ} words are sign-extended or zero-extended from 32 bits, depending on the \emph{highmult}. The \emph{registerY} word is multiplied by the \emph{registerZ} word, and the product is optionally shifted right by 32 bits, depending on the \emph{highmult}. The resulting word is subtracted from and stored into the \emph{registerW}.


\paragraph{Modifier(s)}
\noindent\begin{itemize}
\item highmult (cf. highmult in section \ref{par:modifier-highmult} on page \pageref{par:modifier-highmult})
\item signextw (cf. signextw in section \ref{par:modifier-signextw} on page \pageref{par:modifier-signextw})
\end{itemize}

\paragraph{Execution} {Reservation table \textsf{ALU\_LITE} (Section~\ref{sec:reservation-ALU-LITE})}
~
{\begin{lstlisting}
stage ID:
  new signextw = _ZX_1(signextw);
  new highmult = _ZX_2(highmult);
stage RR:
  new argument3 = GPR[registerY];
  new argument2 = GPR[registerZ];
stage E1:
  new argument1 = GPR[registerW];
  if ((highmult == 0) || (highmult == 3)) {
    new product = _SX_32(argument2) * _SX_32(argument3);
  }
  if (highmult == 1) {
    product = _ZX_32(argument2) * _ZX_32(argument3);
  }
  if (highmult == 2) {
    product = _SX_32(argument2) * _ZX_32(argument3);
  }
  if (highmult) {
    new result1 = argument1 - product.32[1];
  } else {
    result1 = argument1 - product.32[0];
  }
stage E2:
  GPR[registerW] = signextw ? _SX_32(result1) : _ZX_32(result1);
\end{lstlisting}}
\newpage

\paragraph{Encoding} Format \textsf{ALU\_WMWRRR.W} (Section~\ref{sec:format-ALU-WMWRRR.W}), \verb|alucode2 = 11|\\

\bitfields{ 1/P, 2/00, 2/-- --, 27/upper27, 1/1, 2/11, 1/0, 2/11, 2/highmult, 6/registerW, 2/11, 3/110, 1/signextw, 1/0, 5/lower5, 6/registerZ }


\paragraph{Syntax} \texttt{msbfw}\emph{highmult}\emph{signextw} \emph{registerW} = \emph{registerZ}, \emph{upper27\_\discretionary{}{}{}lower5}

\paragraph{Description} The \emph{upper27\_\discretionary{}{}{}lower5} and the \emph{registerZ} words are sign-extended or zero-extended from 32 bits, depending on the \emph{highmult}. The \emph{upper27\_\discretionary{}{}{}lower5} word is multiplied by the \emph{registerZ} word, and the product is optionally shifted right by 32 bits, depending on the \emph{highmult}. The resulting word is subtracted from and stored into the \emph{registerW}.


\paragraph{Modifier(s)}
\noindent\begin{itemize}
\item highmult (cf. highmult in section \ref{par:modifier-highmult} on page \pageref{par:modifier-highmult})
\item signextw (cf. signextw in section \ref{par:modifier-signextw} on page \pageref{par:modifier-signextw})
\end{itemize}

\paragraph{Execution} {Reservation table \textsf{ALU\_LITE.X} (Section~\ref{sec:reservation-ALU-LITE.X})}
~
{\begin{lstlisting}
stage ID:
  new signextw = _ZX_1(signextw);
  new highmult = _ZX_2(highmult);
stage RR:
  new argument3 = _WX_32(upper27_lower5);
  new argument2 = GPR[registerZ];
stage E1:
  new argument1 = GPR[registerW];
  if ((highmult == 0) || (highmult == 3)) {
    new product = _SX_32(argument2) * _SX_32(argument3);
  }
  if (highmult == 1) {
    product = _ZX_32(argument2) * _ZX_32(argument3);
  }
  if (highmult == 2) {
    product = _SX_32(argument2) * _ZX_32(argument3);
  }
  if (highmult) {
    new result1 = argument1 - product.32[1];
  } else {
    result1 = argument1 - product.32[0];
  }
stage E2:
  GPR[registerW] = signextw ? _SX_32(result1) : _ZX_32(result1);
\end{lstlisting}}
\newpage
\subsubsection[{\small MSBFWD: Multiply Subtract From Extend Word to Double Word}]{MSBFWD\hfill Multiply Subtract From Extend Word to Double Word}
\label{instr:MSBFWD}\index{msbfwd}
 
\paragraph{Encoding} Format \textsf{ALU\_DMEWRRR} (Section~\ref{sec:format-ALU-DMEWRRR}), \verb|alucode2 = 11|\\

\bitfields{ 1/P, 2/11, 1/0, 2/11, 2/widemult, 6/registerW, 2/10, 3/111, 1/0, 6/registerY, 6/registerZ }


\paragraph{Syntax} \texttt{msbfwd}\emph{widemult} \emph{registerW} = \emph{registerZ}, \emph{registerY}

\paragraph{Description} The \emph{registerY} and the \emph{registerZ} words are sign-extended or zero-extended from 32 bits, depending on the \emph{widemult}. The \emph{registerY} is multiplied by the \emph{registerZ}, producing a 64-bit result. The resulting double word is subtracted from and stored into the \emph{registerW}.


\paragraph{Modifier(s)}
\noindent\begin{itemize}
\item widemult (cf. widemult in section \ref{par:modifier-widemult} on page \pageref{par:modifier-widemult})
\end{itemize}

\paragraph{Execution} {Reservation table \textsf{ALU\_LITE} (Section~\ref{sec:reservation-ALU-LITE})}
~
{\begin{lstlisting}
stage ID:
  new widemult = _ZX_2(widemult);
stage RR:
  new argument3 = GPR[registerY];
  new argument2 = GPR[registerZ];
stage E1:
  new argument1 = GPR[registerW];
  if (widemult == 0) {
    new product = _SX_32(argument2) * _SX_32(argument3);
  }
  if (widemult == 1) {
    product = _ZX_32(argument2) * _ZX_32(argument3);
  }
  if (widemult == 2) {
    product = _SX_32(argument2) * _ZX_32(argument3);
  }
  new result1 = argument1 - product;
stage E2:
  GPR[registerW] = result1;
\end{lstlisting}}
\newpage

\paragraph{Encoding} Format \textsf{ALU\_DMEWRRR.M} (Section~\ref{sec:format-ALU-DMEWRRR.M}), \verb|alucode2 = 11|\\

\bitfields{ 1/P, 2/00, 2/-- --, 27/upper27, 1/1, 2/11, 1/0, 2/11, 2/widemult, 6/registerW, 2/10, 3/111, 1/0, 1/splat32, 5/lower5, 6/registerZ }


\paragraph{Syntax} \texttt{msbfwd}\emph{widemult} \emph{registerW} = \emph{registerZ}, \emph{upper27\_\discretionary{}{}{}lower5}\emph{splat32}

\paragraph{Description} The \emph{upper27\_\discretionary{}{}{}lower5} and the \emph{registerZ} words are sign-extended or zero-extended from 32 bits, depending on the \emph{widemult}. The \emph{upper27\_\discretionary{}{}{}lower5} is multiplied by the \emph{registerZ}, producing a 64-bit result. The resulting double word is subtracted from and stored into the \emph{registerW}.


\paragraph{Modifier(s)}
\noindent\begin{itemize}
\item widemult (cf. widemult in section \ref{par:modifier-widemult} on page \pageref{par:modifier-widemult})
\item splat32 (cf. splat32 in section \ref{par:modifier-splat32} on page \pageref{par:modifier-splat32})
\end{itemize}

\paragraph{Execution} {Reservation table \textsf{ALU\_LITE.X} (Section~\ref{sec:reservation-ALU-LITE.X})}
~
{\begin{lstlisting}
stage ID:
  new splat32 = _ZX_1(splat32);
  new widemult = _ZX_2(widemult);
stage RR:
  new argument3 = splat32 ? (_WX_32(upper27_lower5) << 32) | _ZX_32(_WX_32(upper27_lower5)) : _SX_32(_WX_32(upper27_lower5));
  new argument2 = GPR[registerZ];
stage E1:
  new argument1 = GPR[registerW];
  if (widemult == 0) {
    new product = _SX_32(argument2) * _SX_32(argument3);
  }
  if (widemult == 1) {
    product = _ZX_32(argument2) * _ZX_32(argument3);
  }
  if (widemult == 2) {
    product = _SX_32(argument2) * _ZX_32(argument3);
  }
  new result1 = argument1 - product;
stage E2:
  GPR[registerW] = result1;
\end{lstlisting}}
\newpage
\subsubsection[{\small MSBFWDP: Multiply Subtract From Extend Word to Double Word Pair}]{MSBFWDP\hfill Multiply Subtract From Extend Word to Double Word Pair}
\label{instr:MSBFWDP}\index{msbfwdp}
 
\paragraph{Encoding} Format \textsf{ALU\_DPMEWRRRE} (Section~\ref{sec:format-ALU-DPMEWRRRE}), \verb|alucode2 = 11|\\

\bitfields{ 1/P, 2/11, 1/1, 2/11, 2/widemult, 5/registerM, 1/0, 2/10, 3/110, 1/0, 5/registerYe, 1/1, 5/registerZe, 1/0 }


\paragraph{Syntax} \texttt{msbfwdp}\emph{widemult} \emph{registerM} = \emph{registerZe}, \emph{registerYe}

\paragraph{Description} The \emph{registerYe} and the \emph{registerZe} are interpreted as two words packed into 65 bits. These \emph{registerYe} and \emph{registerZe} words are sign-extended or zero-extended from 32 bits, depending on the \emph{widemult}. The \emph{registerYe} words are multiplied by the \emph{registerZe} words, producing two 64-bit results. The two resulting double words are packed, subtracted from and stored into the \emph{registerM}.


\paragraph{Modifier(s)}
\noindent\begin{itemize}
\item widemult (cf. widemult in section \ref{par:modifier-widemult} on page \pageref{par:modifier-widemult})
\end{itemize}

\paragraph{Execution} {Reservation table \textsf{ALU\_LITE} (Section~\ref{sec:reservation-ALU-LITE})}
~
{\begin{lstlisting}
stage ID:
  new widemult = _ZX_2(widemult);
stage RR:
  new argument3 = GPR[registerYe<<1];
  new argument2 = GPR[registerZe<<1];
stage E1:
  new argument1 = PGR[registerM];
  for (I = 0; I < 2; I += 1) {
    if (widemult == 0) {
      new product = _SX_32(argument2.64[I]) * _SX_32(argument3.64[I]);
    }
    if (widemult == 1) {
      product = _ZX_32(argument2.64[I]) * _ZX_32(argument3.64[I]);
    }
    if (widemult == 2) {
      product = _SX_32(argument2.64[I]) * _ZX_32(argument3.64[I]);
    }
    result1.64[I] = argument1.64[I] - product
  };
stage E2:
  PGR[registerM] = result1;
\end{lstlisting}}
\newpage

\paragraph{Encoding} Format \textsf{ALU\_DPMEWRRRO} (Section~\ref{sec:format-ALU-DPMEWRRRO}), \verb|alucode2 = 11|\\

\bitfields{ 1/P, 2/11, 1/1, 2/11, 2/widemult, 5/registerM, 1/1, 2/10, 3/110, 1/0, 5/registerYo, 1/1, 5/registerZo, 1/0 }


\paragraph{Syntax} \texttt{msbfwdp}\emph{widemult} \emph{registerM} = \emph{registerZo}, \emph{registerYo}

\paragraph{Description} The \emph{registerYo} and the \emph{registerZo} are interpreted as two words packed into 65 bits. These \emph{registerYo} and \emph{registerZo} words are sign-extended or zero-extended from 32 bits, depending on the \emph{widemult}. The \emph{registerYo} words are multiplied by the \emph{registerZo} words, producing two 64-bit results. The two resulting double words are packed, subtracted from and stored into the \emph{registerM}.


\paragraph{Modifier(s)}
\noindent\begin{itemize}
\item widemult (cf. widemult in section \ref{par:modifier-widemult} on page \pageref{par:modifier-widemult})
\end{itemize}

\paragraph{Execution} {Reservation table \textsf{ALU\_LITE} (Section~\ref{sec:reservation-ALU-LITE})}
~
{\begin{lstlisting}
stage ID:
  new widemult = _ZX_2(widemult);
stage RR:
  new argument3 = GPR[(registerYo<<1)+1];
  new argument2 = GPR[(registerZo<<1)+1];
stage E1:
  new argument1 = PGR[registerM];
  for (I = 0; I < 2; I += 1) {
    if (widemult == 0) {
      new product = _SX_32(argument2.64[I]) * _SX_32(argument3.64[I]);
    }
    if (widemult == 1) {
      product = _ZX_32(argument2.64[I]) * _ZX_32(argument3.64[I]);
    }
    if (widemult == 2) {
      product = _SX_32(argument2.64[I]) * _ZX_32(argument3.64[I]);
    }
    result1.64[I] = argument1.64[I] - product
  };
stage E2:
  PGR[registerM] = result1;
\end{lstlisting}}
\newpage
\subsubsection[{\small MSBFXBHO: Multiply Subtract From Extend Byte to Half Word Octuple}]{MSBFXBHO\hfill Multiply Subtract From Extend Byte to Half Word Octuple}
\label{instr:MSBFXBHO}\index{msbfxbho}
 
\paragraph{Encoding} Format \textsf{ALU\_HOMXWRRR} (Section~\ref{sec:format-ALU-HOMXWRRR}), \verb|alucode2 = 11|\\

\bitfields{ 1/P, 2/11, 1/1, 2/11, 2/widemult, 5/registerM, 1/oddlanes, 2/10, 3/111, 1/1, 5/registerO, 1/--, 5/registerP, 1/0 }


\paragraph{Syntax} \texttt{msbfxbho}\emph{oddlanes}\emph{widemult} \emph{registerM} = \emph{registerP}, \emph{registerO}

\paragraph{Description} The \emph{registerO} and the \emph{registerP} are interpreted as sixteen bytes packed into 128 bits. The even or odd bytes of the \emph{registerO} and the \emph{registerP} are selected, depending on \emph{oddlanes}. These \emph{registerO} and \emph{registerP} bytes are sign-extended or zero-extended from 8 bits, depending on the \emph{widemult}. The \emph{registerO} bytes are multiplied by the \emph{registerP} bytes, producing eight 16-bit results. The eight resulting half words are packed, subtracted from and stored into the \emph{registerM}.


\paragraph{Modifier(s)}
\noindent\begin{itemize}
\item widemult (cf. widemult in section \ref{par:modifier-widemult} on page \pageref{par:modifier-widemult})
\item oddlanes (cf. oddlanes in section \ref{par:modifier-oddlanes} on page \pageref{par:modifier-oddlanes})
\end{itemize}

\paragraph{Execution} {Reservation table \textsf{ALU\_LITE} (Section~\ref{sec:reservation-ALU-LITE})}
~
{\begin{lstlisting}
stage ID:
  new oddlanes = _ZX_1(oddlanes);
  new widemult = _ZX_2(widemult);
stage RR:
  new argument3 = PGR[registerO];
  new argument2 = PGR[registerP];
stage E1:
  new argument1 = PGR[registerM];
  if (oddlanes) {
    new argument2 = argument2 >> 8;
    new argument3 = argument3 >> 8;
  }
  for (I = 0; I < 8; I += 1) {
    if (widemult == 0) {
      new product = _SX_8(argument2.16[I]) * _SX_8(argument3.16[I]);
    }
    if (widemult == 1) {
      product = _ZX_8(argument2.16[I]) * _ZX_8(argument3.16[I]);
    }
    if (widemult == 2) {
      product = _SX_8(argument2.16[I]) * _ZX_8(argument3.16[I]);
    }
    result1.16[I] = argument1.16[I] - product
  };
stage E2:
  PGR[registerM] = result1;
\end{lstlisting}}
\newpage
\subsubsection[{\small MSBFXHWQ: Multiply Subtract From Extend Half Word to Word Quadruple}]{MSBFXHWQ\hfill Multiply Subtract From Extend Half Word to Word Quadruple}
\label{instr:MSBFXHWQ}\index{msbfxhwq}
 
\paragraph{Encoding} Format \textsf{ALU\_WQMXWRRR} (Section~\ref{sec:format-ALU-WQMXWRRR}), \verb|alucode2 = 11|\\

\bitfields{ 1/P, 2/11, 1/1, 2/11, 2/widemult, 5/registerM, 1/oddlanes, 2/10, 3/111, 1/0, 5/registerO, 1/--, 5/registerP, 1/1 }


\paragraph{Syntax} \texttt{msbfxhwq}\emph{oddlanes}\emph{widemult} \emph{registerM} = \emph{registerP}, \emph{registerO}

\paragraph{Description} The \emph{registerO} and the \emph{registerP} are interpreted as eight half words packed into 128 bits. The even or odd half words of the \emph{registerO} and the \emph{registerP} are selected, depending on \emph{oddlanes}. These \emph{registerO} and \emph{registerP} half words are sign-extended or zero-extended from 16 bits, depending on the \emph{widemult}. The \emph{registerO} half words are multiplied by the \emph{registerP} half words, producing four 32-bit results. The four resulting words are packed, subtracted from and stored into the \emph{registerM}.


\paragraph{Modifier(s)}
\noindent\begin{itemize}
\item widemult (cf. widemult in section \ref{par:modifier-widemult} on page \pageref{par:modifier-widemult})
\item oddlanes (cf. oddlanes in section \ref{par:modifier-oddlanes} on page \pageref{par:modifier-oddlanes})
\end{itemize}

\paragraph{Execution} {Reservation table \textsf{ALU\_LITE} (Section~\ref{sec:reservation-ALU-LITE})}
~
{\begin{lstlisting}
stage ID:
  new oddlanes = _ZX_1(oddlanes);
  new widemult = _ZX_2(widemult);
stage RR:
  new argument3 = PGR[registerO];
  new argument2 = PGR[registerP];
stage E1:
  new argument1 = PGR[registerM];
  if (oddlanes) {
    new argument2 = argument2 >> 16;
    new argument3 = argument3 >> 16;
  }
  for (I = 0; I < 4; I += 1) {
    if (widemult == 0) {
      new product = _SX_16(argument2.32[I]) * _SX_16(argument3.32[I]);
    }
    if (widemult == 1) {
      product = _ZX_16(argument2.32[I]) * _ZX_16(argument3.32[I]);
    }
    if (widemult == 2) {
      product = _SX_16(argument2.32[I]) * _ZX_16(argument3.32[I]);
    }
    result1.32[I] = argument1.32[I] - product
  };
stage E2:
  PGR[registerM] = result1;
\end{lstlisting}}
\newpage
\subsubsection[{\small MSBFXWDP: Multiply Subtract From Extend Word to Double Word Pair}]{MSBFXWDP\hfill Multiply Subtract From Extend Word to Double Word Pair}
\label{instr:MSBFXWDP}\index{msbfxwdp}
 
\paragraph{Encoding} Format \textsf{ALU\_DPMXWRRR} (Section~\ref{sec:format-ALU-DPMXWRRR}), \verb|alucode2 = 11|\\

\bitfields{ 1/P, 2/11, 1/1, 2/11, 2/widemult, 5/registerM, 1/oddlanes, 2/10, 3/111, 1/0, 5/registerO, 1/--, 5/registerP, 1/0 }


\paragraph{Syntax} \texttt{msbfxwdp}\emph{oddlanes}\emph{widemult} \emph{registerM} = \emph{registerP}, \emph{registerO}

\paragraph{Description} The \emph{registerO} and the \emph{registerP} are interpreted as four words packed into 128 bits. The even or odd words of the \emph{registerO} and the \emph{registerP} are selected, depending on \emph{oddlanes}. These \emph{registerO} and \emph{registerP} words are sign-extended or zero-extended from 32 bits, depending on the \emph{widemult}. The \emph{registerO} words are multiplied by the \emph{registerP} words, producing two 64-bit results. The two resulting double words are packed, subtracted from and stored into the \emph{registerM}.


\paragraph{Modifier(s)}
\noindent\begin{itemize}
\item widemult (cf. widemult in section \ref{par:modifier-widemult} on page \pageref{par:modifier-widemult})
\item oddlanes (cf. oddlanes in section \ref{par:modifier-oddlanes} on page \pageref{par:modifier-oddlanes})
\end{itemize}

\paragraph{Execution} {Reservation table \textsf{ALU\_LITE} (Section~\ref{sec:reservation-ALU-LITE})}
~
{\begin{lstlisting}
stage ID:
  new oddlanes = _ZX_1(oddlanes);
  new widemult = _ZX_2(widemult);
stage RR:
  new argument3 = PGR[registerO];
  new argument2 = PGR[registerP];
stage E1:
  new argument1 = PGR[registerM];
  if (oddlanes) {
    new argument2 = argument2 >> 32;
    new argument3 = argument3 >> 32;
  }
  for (I = 0; I < 2; I += 1) {
    if (widemult == 0) {
      new product = _SX_32(argument2.64[I]) * _SX_32(argument3.64[I]);
    }
    if (widemult == 1) {
      product = _ZX_32(argument2.64[I]) * _ZX_32(argument3.64[I]);
    }
    if (widemult == 2) {
      product = _SX_32(argument2.64[I]) * _ZX_32(argument3.64[I]);
    }
    result1.64[I] = argument1.64[I] - product
  };
stage E2:
  PGR[registerM] = result1;
\end{lstlisting}}
\newpage
\subsubsection[{\small MULBHO: Multiply Extend Byte to Half Word Octuple}]{MULBHO\hfill Multiply Extend Byte to Half Word Octuple}
\label{instr:MULBHO}\index{mulbho}
 
\paragraph{Encoding} Format \textsf{ALU\_HOMEWRRE} (Section~\ref{sec:format-ALU-HOMEWRRE}), \verb|alucode2 = 00|\\

\bitfields{ 1/P, 2/11, 1/1, 2/00, 2/widemult, 5/registerM, 1/0, 2/10, 3/110, 1/1, 5/registerYe, 1/1, 5/registerZe, 1/0 }


\paragraph{Syntax} \texttt{mulbho}\emph{widemult} \emph{registerM} = \emph{registerZe}, \emph{registerYe}

\paragraph{Description} The \emph{registerYe} and the \emph{registerZe} are interpreted as eight bytes packed into 64 bits. These \emph{registerYe} and \emph{registerZe} bytes are sign-extended or zero-extended from 8 bits, depending on the \emph{widemult}. The \emph{registerYe} bytes are multiplied by the \emph{registerZe} bytes, producing eight 16-bit results. The eight resulting half words are packed and stored into the \emph{registerM}.


\paragraph{Modifier(s)}
\noindent\begin{itemize}
\item widemult (cf. widemult in section \ref{par:modifier-widemult} on page \pageref{par:modifier-widemult})
\end{itemize}

\paragraph{Execution} {Reservation table \textsf{ALU\_LITE} (Section~\ref{sec:reservation-ALU-LITE})}
~
{\begin{lstlisting}
stage ID:
  new widemult = _ZX_2(widemult);
stage RR:
  new argument3 = GPR[registerYe<<1];
  new argument2 = GPR[registerZe<<1];
  for (I = 0; I < 8; I += 1) {
    if (widemult == 0) {
      new product = _SX_8(argument2.8[I]) * _SX_8(argument3.8[I]);
    }
    if (widemult == 1) {
      product = _ZX_8(argument2.8[I]) * _ZX_8(argument3.8[I]);
    }
    if (widemult == 2) {
      product = _SX_8(argument2.8[I]) * _ZX_8(argument3.8[I]);
    }
    result1.16[I] = product
  };
stage E2:
  PGR[registerM] = result1;
\end{lstlisting}}
\newpage

\paragraph{Encoding} Format \textsf{ALU\_HOMEWRRO} (Section~\ref{sec:format-ALU-HOMEWRRO}), \verb|alucode2 = 00|\\

\bitfields{ 1/P, 2/11, 1/1, 2/00, 2/widemult, 5/registerM, 1/1, 2/10, 3/110, 1/1, 5/registerYo, 1/1, 5/registerZo, 1/0 }


\paragraph{Syntax} \texttt{mulbho}\emph{widemult} \emph{registerM} = \emph{registerZo}, \emph{registerYo}

\paragraph{Description} The \emph{registerYo} and the \emph{registerZo} are interpreted as eight bytes packed into 64 bits. These \emph{registerYo} and \emph{registerZo} bytes are sign-extended or zero-extended from 8 bits, depending on the \emph{widemult}. The \emph{registerYo} bytes are multiplied by the \emph{registerZo} bytes, producing eight 16-bit results. The eight resulting half words are packed and stored into the \emph{registerM}.


\paragraph{Modifier(s)}
\noindent\begin{itemize}
\item widemult (cf. widemult in section \ref{par:modifier-widemult} on page \pageref{par:modifier-widemult})
\end{itemize}

\paragraph{Execution} {Reservation table \textsf{ALU\_LITE} (Section~\ref{sec:reservation-ALU-LITE})}
~
{\begin{lstlisting}
stage ID:
  new widemult = _ZX_2(widemult);
stage RR:
  new argument3 = GPR[(registerYo<<1)+1];
  new argument2 = GPR[(registerZo<<1)+1];
  for (I = 0; I < 8; I += 1) {
    if (widemult == 0) {
      new product = _SX_8(argument2.8[I]) * _SX_8(argument3.8[I]);
    }
    if (widemult == 1) {
      product = _ZX_8(argument2.8[I]) * _ZX_8(argument3.8[I]);
    }
    if (widemult == 2) {
      product = _SX_8(argument2.8[I]) * _ZX_8(argument3.8[I]);
    }
    result1.16[I] = product
  };
stage E2:
  PGR[registerM] = result1;
\end{lstlisting}}
\newpage
\subsubsection[{\small MULD: Multiply Double Word}]{MULD\hfill Multiply Double Word}
\label{instr:MULD}\index{muld}
 
\paragraph{Encoding} Format \textsf{ALU\_DMWRR} (Section~\ref{sec:format-ALU-DMWRR}), \verb|alucode2 = 00|\\

\bitfields{ 1/P, 2/11, 1/0, 2/00, 2/highmult, 6/registerW, 2/10, 3/110, 1/0, 6/registerY, 6/registerZ }


\paragraph{Syntax} \texttt{muld}\emph{highmult} \emph{registerW} = \emph{registerZ}, \emph{registerY}

\paragraph{Description} The \emph{registerY} and the \emph{registerZ} double words are sign-extended or zero-extended from 64 bits, depending on the \emph{highmult}. The \emph{registerY} double word is multiplied by the \emph{registerZ} double word, and the product is optionally shifted right by 64 bits, depending on the \emph{highmult}. The resulting double word is stored into the \emph{registerW}.


\paragraph{Modifier(s)}
\noindent\begin{itemize}
\item highmult (cf. highmult in section \ref{par:modifier-highmult} on page \pageref{par:modifier-highmult})
\end{itemize}

\paragraph{Execution} {Reservation table \textsf{ALU\_LITE} (Section~\ref{sec:reservation-ALU-LITE})}
~
{\begin{lstlisting}
stage ID:
  new highmult = _ZX_2(highmult);
stage RR:
  new argument3 = GPR[registerY];
  new argument2 = GPR[registerZ];
  if ((highmult == 0) || (highmult == 3)) {
    new product = _SX_64(argument2) * _SX_64(argument3);
  }
  if (highmult == 1) {
    product = _ZX_64(argument2) * _ZX_64(argument3);
  }
  if (highmult == 2) {
    product = _SX_64(argument2) * _ZX_64(argument3);
  }
  if (highmult) {
    new result1 = product.64[1];
  } else {
    result1 = product.64[0];
  }
stage E2:
  GPR[registerW] = result1;
\end{lstlisting}}
\newpage

\paragraph{Encoding} Format \textsf{ALU\_DMWRR.M} (Section~\ref{sec:format-ALU-DMWRR.M}), \verb|alucode2 = 00|\\

\bitfields{ 1/P, 2/00, 2/-- --, 27/upper27, 1/1, 2/11, 1/0, 2/00, 2/highmult, 6/registerW, 2/10, 3/110, 1/0, 1/splat32, 5/lower5, 6/registerZ }


\paragraph{Syntax} \texttt{muld}\emph{highmult} \emph{registerW} = \emph{registerZ}, \emph{upper27\_\discretionary{}{}{}lower5}\emph{splat32}

\paragraph{Description} The \emph{upper27\_\discretionary{}{}{}lower5} and the \emph{registerZ} double words are sign-extended or zero-extended from 64 bits, depending on the \emph{highmult}. The \emph{upper27\_\discretionary{}{}{}lower5} double word is multiplied by the \emph{registerZ} double word, and the product is optionally shifted right by 64 bits, depending on the \emph{highmult}. The resulting double word is stored into the \emph{registerW}.


\paragraph{Modifier(s)}
\noindent\begin{itemize}
\item highmult (cf. highmult in section \ref{par:modifier-highmult} on page \pageref{par:modifier-highmult})
\item splat32 (cf. splat32 in section \ref{par:modifier-splat32} on page \pageref{par:modifier-splat32})
\end{itemize}

\paragraph{Execution} {Reservation table \textsf{ALU\_LITE.X} (Section~\ref{sec:reservation-ALU-LITE.X})}
~
{\begin{lstlisting}
stage ID:
  new splat32 = _ZX_1(splat32);
  new highmult = _ZX_2(highmult);
stage RR:
  new argument3 = splat32 ? (_WX_32(upper27_lower5) << 32) | _ZX_32(_WX_32(upper27_lower5)) : _SX_32(_WX_32(upper27_lower5));
  new argument2 = GPR[registerZ];
  if ((highmult == 0) || (highmult == 3)) {
    new product = _SX_64(argument2) * _SX_64(argument3);
  }
  if (highmult == 1) {
    product = _ZX_64(argument2) * _ZX_64(argument3);
  }
  if (highmult == 2) {
    product = _SX_64(argument2) * _ZX_64(argument3);
  }
  if (highmult) {
    new result1 = product.64[1];
  } else {
    result1 = product.64[0];
  }
stage E2:
  GPR[registerW] = result1;
\end{lstlisting}}
\newpage
\subsubsection[{\small MULDP: Multiply Double Word Pair}]{MULDP\hfill Multiply Double Word Pair}
\label{instr:MULDP}\index{muldp}
 
\paragraph{Encoding} Format \textsf{ALU\_DPMWRR} (Section~\ref{sec:format-ALU-DPMWRR}), \verb|alucode2 = 00|\\

\bitfields{ 1/P, 2/11, 1/1, 2/00, 2/highmult, 5/registerM, 1/--, 2/10, 3/110, 1/0, 5/registerO, 1/0, 5/registerP, 1/0 }


\paragraph{Syntax} \texttt{muldp}\emph{highmult} \emph{registerM} = \emph{registerP}, \emph{registerO}

\paragraph{Description} The \emph{registerO} and the \emph{registerP} are interpreted as two double words packed into 128 bits. The \emph{registerO} and the \emph{registerP} double words are sign-extended or zero-extended from 64 bits, depending on the \emph{highmult}. The \emph{registerO} double words are multiplied by the \emph{registerP} double words, and the product is optionally shifted right by 64 bits, depending on the \emph{highmult}. The two resulting double words are packed and stored back into the \emph{registerM}.


\paragraph{Modifier(s)}
\noindent\begin{itemize}
\item highmult (cf. highmult in section \ref{par:modifier-highmult} on page \pageref{par:modifier-highmult})
\end{itemize}

\paragraph{Execution} {Reservation table \textsf{ALU\_LITE} (Section~\ref{sec:reservation-ALU-LITE})}
~
{\begin{lstlisting}
stage ID:
  new highmult = _ZX_2(highmult);
stage RR:
  new argument3 = PGR[registerO];
  new argument2 = PGR[registerP];
  for (I = 0; I < 2; I += 1) {
    if ((highmult == 0) || (highmult == 3)) {
      new product = _SX_64(argument2.64[I]) * _SX_64(argument3.64[I]);
    }
    if (highmult == 1) {
      product = _ZX_64(argument2.64[I]) * _ZX_64(argument3.64[I]);
    }
    if (highmult == 2) {
      product = _SX_64(argument2.64[I]) * _ZX_64(argument3.64[I]);
    }
    if (highmult) {
      result1.64[I] = product.64[1];
    } else {
      result1.64[I] = product.64[0];
    }
  };
stage E2:
  PGR[registerM] = result1;
\end{lstlisting}}
\newpage
\subsubsection[{\small MULDQ: Multiply Extend Double Word to Quadruple Word}]{MULDQ\hfill Multiply Extend Double Word to Quadruple Word}
\label{instr:MULDQ}\index{muldq}
 
\paragraph{Encoding} Format \textsf{ALU\_QMEWRR} (Section~\ref{sec:format-ALU-QMEWRR}), \verb|alucode2 = 00|\\

\bitfields{ 1/P, 2/11, 1/0, 2/00, 2/widemult, 5/registerM, 1/--, 2/10, 3/111, 1/1, 6/registerY, 6/registerZ }


\paragraph{Syntax} \texttt{muldq}\emph{widemult} \emph{registerM} = \emph{registerZ}, \emph{registerY}

\paragraph{Description} The \emph{registerY} and the \emph{registerZ} double words are sign-extended or zero-extended from 64 bits, depending on the \emph{widemult}. The \emph{registerY} is multiplied by the \emph{registerZ}, producing a 128-bit result. The resulting quadruple word is stored into the \emph{registerM}.


\paragraph{Modifier(s)}
\noindent\begin{itemize}
\item widemult (cf. widemult in section \ref{par:modifier-widemult} on page \pageref{par:modifier-widemult})
\end{itemize}

\paragraph{Execution} {Reservation table \textsf{ALU\_LITE} (Section~\ref{sec:reservation-ALU-LITE})}
~
{\begin{lstlisting}
stage ID:
  new widemult = _ZX_2(widemult);
stage RR:
  new argument3 = GPR[registerY];
  new argument2 = GPR[registerZ];
  if (widemult == 0) {
    new product = _SX_64(argument2) * _SX_64(argument3);
  }
  if (widemult == 1) {
    product = _ZX_64(argument2) * _ZX_64(argument3);
  }
  if (widemult == 2) {
    product = _SX_64(argument2) * _ZX_64(argument3);
  }
  new result1 = product;
stage E2:
  PGR[registerM] = result1;
\end{lstlisting}}
\newpage
\subsubsection[{\small MULHO: Multiply Half Word Octuple}]{MULHO\hfill Multiply Half Word Octuple}
\label{instr:MULHO}\index{mulho}
 
\paragraph{Encoding} Format \textsf{ALU\_HOMWRR} (Section~\ref{sec:format-ALU-HOMWRR}), \verb|alucode2 = 00|\\

\bitfields{ 1/P, 2/11, 1/1, 2/00, 2/highmult, 5/registerM, 1/--, 2/10, 3/110, 1/1, 5/registerO, 1/0, 5/registerP, 1/0 }


\paragraph{Syntax} \texttt{mulho}\emph{highmult} \emph{registerM} = \emph{registerP}, \emph{registerO}

\paragraph{Description} The \emph{registerO} and the \emph{registerP} are interpreted as eight half words packed into 128 bits. The \emph{registerO} and the \emph{registerP} half words are sign-extended or zero-extended from 16 bits, depending on the \emph{highmult}. The \emph{registerO} half words are multiplied by the \emph{registerP} half words, and the product is optionally shifted right by 16 bits, depending on the \emph{highmult}. The eight resulting half words are packed and stored back into the \emph{registerM}.


\paragraph{Modifier(s)}
\noindent\begin{itemize}
\item highmult (cf. highmult in section \ref{par:modifier-highmult} on page \pageref{par:modifier-highmult})
\end{itemize}

\paragraph{Execution} {Reservation table \textsf{ALU\_LITE} (Section~\ref{sec:reservation-ALU-LITE})}
~
{\begin{lstlisting}
stage ID:
  new highmult = _ZX_2(highmult);
stage RR:
  new argument3 = PGR[registerO];
  new argument2 = PGR[registerP];
  for (I = 0; I < 8; I += 1) {
    if ((highmult == 0) || (highmult == 3)) {
      new product = _SX_16(argument2.16[I]) * _SX_16(argument3.16[I]);
    }
    if (highmult == 1) {
      product = _ZX_16(argument2.16[I]) * _ZX_16(argument3.16[I]);
    }
    if (highmult == 2) {
      product = _SX_16(argument2.16[I]) * _ZX_16(argument3.16[I]);
    }
    if (highmult) {
      result1.16[I] = product.16[1];
    } else {
      result1.16[I] = product.16[0];
    }
  };
stage E2:
  PGR[registerM] = result1;
\end{lstlisting}}
\newpage
\subsubsection[{\small MULHWQ: Multiply Extend Half Word to Word Quadruple}]{MULHWQ\hfill Multiply Extend Half Word to Word Quadruple}
\label{instr:MULHWQ}\index{mulhwq}
 
\paragraph{Encoding} Format \textsf{ALU\_WQMEWRRE} (Section~\ref{sec:format-ALU-WQMEWRRE}), \verb|alucode2 = 00|\\

\bitfields{ 1/P, 2/11, 1/1, 2/00, 2/widemult, 5/registerM, 1/0, 2/10, 3/110, 1/0, 5/registerYe, 1/1, 5/registerZe, 1/1 }


\paragraph{Syntax} \texttt{mulhwq}\emph{widemult} \emph{registerM} = \emph{registerZe}, \emph{registerYe}

\paragraph{Description} The \emph{registerYe} and the \emph{registerZe} are interpreted as four half words packed into 64 bits. These \emph{registerYe} and \emph{registerZe} half words are sign-extended or zero-extended from 16 bits, depending on the \emph{widemult}. The \emph{registerYe} half words are multiplied by the \emph{registerZe} half words, producing four 32-bit results. The four resulting words are packed and stored into the \emph{registerM}.


\paragraph{Modifier(s)}
\noindent\begin{itemize}
\item widemult (cf. widemult in section \ref{par:modifier-widemult} on page \pageref{par:modifier-widemult})
\end{itemize}

\paragraph{Execution} {Reservation table \textsf{ALU\_LITE} (Section~\ref{sec:reservation-ALU-LITE})}
~
{\begin{lstlisting}
stage ID:
  new widemult = _ZX_2(widemult);
stage RR:
  new argument3 = GPR[registerYe<<1];
  new argument2 = GPR[registerZe<<1];
  for (I = 0; I < 4; I += 1) {
    if (widemult == 0) {
      new product = _SX_16(argument2.16[I]) * _SX_16(argument3.16[I]);
    }
    if (widemult == 1) {
      product = _ZX_16(argument2.16[I]) * _ZX_16(argument3.16[I]);
    }
    if (widemult == 2) {
      product = _SX_16(argument2.16[I]) * _ZX_16(argument3.16[I]);
    }
    result1.32[I] = product
  };
stage E2:
  PGR[registerM] = result1;
\end{lstlisting}}
\newpage

\paragraph{Encoding} Format \textsf{ALU\_WQMEWRRO} (Section~\ref{sec:format-ALU-WQMEWRRO}), \verb|alucode2 = 00|\\

\bitfields{ 1/P, 2/11, 1/1, 2/00, 2/widemult, 5/registerM, 1/1, 2/10, 3/110, 1/0, 5/registerYo, 1/1, 5/registerZo, 1/1 }


\paragraph{Syntax} \texttt{mulhwq}\emph{widemult} \emph{registerM} = \emph{registerZo}, \emph{registerYo}

\paragraph{Description} The \emph{registerYo} and the \emph{registerZo} are interpreted as four half words packed into 64 bits. These \emph{registerYo} and \emph{registerZo} half words are sign-extended or zero-extended from 16 bits, depending on the \emph{widemult}. The \emph{registerYo} half words are multiplied by the \emph{registerZo} half words, producing four 32-bit results. The four resulting words are packed and stored into the \emph{registerM}.


\paragraph{Modifier(s)}
\noindent\begin{itemize}
\item widemult (cf. widemult in section \ref{par:modifier-widemult} on page \pageref{par:modifier-widemult})
\end{itemize}

\paragraph{Execution} {Reservation table \textsf{ALU\_LITE} (Section~\ref{sec:reservation-ALU-LITE})}
~
{\begin{lstlisting}
stage ID:
  new widemult = _ZX_2(widemult);
stage RR:
  new argument3 = GPR[(registerYo<<1)+1];
  new argument2 = GPR[(registerZo<<1)+1];
  for (I = 0; I < 4; I += 1) {
    if (widemult == 0) {
      new product = _SX_16(argument2.16[I]) * _SX_16(argument3.16[I]);
    }
    if (widemult == 1) {
      product = _ZX_16(argument2.16[I]) * _ZX_16(argument3.16[I]);
    }
    if (widemult == 2) {
      product = _SX_16(argument2.16[I]) * _ZX_16(argument3.16[I]);
    }
    result1.32[I] = product
  };
stage E2:
  PGR[registerM] = result1;
\end{lstlisting}}
\newpage
\subsubsection[{\small MULNBHO: Multiply Negate Extend Byte to Half Word Octuple}]{MULNBHO\hfill Multiply Negate Extend Byte to Half Word Octuple}
\label{instr:MULNBHO}\index{mulnbho}
 
\paragraph{Encoding} Format \textsf{ALU\_HOMEWRRE} (Section~\ref{sec:format-ALU-HOMEWRRE}), \verb|alucode2 = 01|\\

\bitfields{ 1/P, 2/11, 1/1, 2/01, 2/widemult, 5/registerM, 1/0, 2/10, 3/110, 1/1, 5/registerYe, 1/1, 5/registerZe, 1/0 }


\paragraph{Syntax} \texttt{mulnbho}\emph{widemult} \emph{registerM} = \emph{registerZe}, \emph{registerYe}

\paragraph{Description} The \emph{registerYe} and the \emph{registerZe} are interpreted as eight bytes packed into 64 bits. These \emph{registerYe} and \emph{registerZe} bytes are sign-extended or zero-extended from 8 bits, depending on the \emph{widemult}. The \emph{registerYe} bytes are multiplied by the \emph{registerZe} bytes, producing eight 16-bit results. The eight resulting half words are negated, packed and stored into the \emph{registerM}.


\paragraph{Modifier(s)}
\noindent\begin{itemize}
\item widemult (cf. widemult in section \ref{par:modifier-widemult} on page \pageref{par:modifier-widemult})
\end{itemize}

\paragraph{Execution} {Reservation table \textsf{ALU\_LITE} (Section~\ref{sec:reservation-ALU-LITE})}
~
{\begin{lstlisting}
stage ID:
  new widemult = _ZX_2(widemult);
stage RR:
  new argument3 = GPR[registerYe<<1];
  new argument2 = GPR[registerZe<<1];
  for (I = 0; I < 8; I += 1) {
    if (widemult == 0) {
      new product = _SX_8(argument2.8[I]) * _SX_8(argument3.8[I]);
    }
    if (widemult == 1) {
      product = _ZX_8(argument2.8[I]) * _ZX_8(argument3.8[I]);
    }
    if (widemult == 2) {
      product = _SX_8(argument2.8[I]) * _ZX_8(argument3.8[I]);
    }
    result1.16[I] = -product
  };
stage E2:
  PGR[registerM] = result1;
\end{lstlisting}}
\newpage

\paragraph{Encoding} Format \textsf{ALU\_HOMEWRRO} (Section~\ref{sec:format-ALU-HOMEWRRO}), \verb|alucode2 = 01|\\

\bitfields{ 1/P, 2/11, 1/1, 2/01, 2/widemult, 5/registerM, 1/1, 2/10, 3/110, 1/1, 5/registerYo, 1/1, 5/registerZo, 1/0 }


\paragraph{Syntax} \texttt{mulnbho}\emph{widemult} \emph{registerM} = \emph{registerZo}, \emph{registerYo}

\paragraph{Description} The \emph{registerYo} and the \emph{registerZo} are interpreted as eight bytes packed into 64 bits. These \emph{registerYo} and \emph{registerZo} bytes are sign-extended or zero-extended from 8 bits, depending on the \emph{widemult}. The \emph{registerYo} bytes are multiplied by the \emph{registerZo} bytes, producing eight 16-bit results. The eight resulting half words are negated, packed and stored into the \emph{registerM}.


\paragraph{Modifier(s)}
\noindent\begin{itemize}
\item widemult (cf. widemult in section \ref{par:modifier-widemult} on page \pageref{par:modifier-widemult})
\end{itemize}

\paragraph{Execution} {Reservation table \textsf{ALU\_LITE} (Section~\ref{sec:reservation-ALU-LITE})}
~
{\begin{lstlisting}
stage ID:
  new widemult = _ZX_2(widemult);
stage RR:
  new argument3 = GPR[(registerYo<<1)+1];
  new argument2 = GPR[(registerZo<<1)+1];
  for (I = 0; I < 8; I += 1) {
    if (widemult == 0) {
      new product = _SX_8(argument2.8[I]) * _SX_8(argument3.8[I]);
    }
    if (widemult == 1) {
      product = _ZX_8(argument2.8[I]) * _ZX_8(argument3.8[I]);
    }
    if (widemult == 2) {
      product = _SX_8(argument2.8[I]) * _ZX_8(argument3.8[I]);
    }
    result1.16[I] = -product
  };
stage E2:
  PGR[registerM] = result1;
\end{lstlisting}}
\newpage
\subsubsection[{\small MULND: Multiply Negate Double Word}]{MULND\hfill Multiply Negate Double Word}
\label{instr:MULND}\index{mulnd}
 
\paragraph{Encoding} Format \textsf{ALU\_DMWRR} (Section~\ref{sec:format-ALU-DMWRR}), \verb|alucode2 = 01|\\

\bitfields{ 1/P, 2/11, 1/0, 2/01, 2/highmult, 6/registerW, 2/10, 3/110, 1/0, 6/registerY, 6/registerZ }


\paragraph{Syntax} \texttt{mulnd}\emph{highmult} \emph{registerW} = \emph{registerZ}, \emph{registerY}

\paragraph{Description} The \emph{registerY} and the \emph{registerZ} double words are sign-extended or zero-extended from 64 bits, depending on the \emph{highmult}. The \emph{registerY} double word is multiplied by the \emph{registerZ} double word, and the product is optionally shifted right by 64 bits, depending on the \emph{highmult}. The resulting double word is negated and stored into the \emph{registerW}.


\paragraph{Modifier(s)}
\noindent\begin{itemize}
\item highmult (cf. highmult in section \ref{par:modifier-highmult} on page \pageref{par:modifier-highmult})
\end{itemize}

\paragraph{Execution} {Reservation table \textsf{ALU\_LITE} (Section~\ref{sec:reservation-ALU-LITE})}
~
{\begin{lstlisting}
stage ID:
  new highmult = _ZX_2(highmult);
stage RR:
  new argument3 = GPR[registerY];
  new argument2 = GPR[registerZ];
  if ((highmult == 0) || (highmult == 3)) {
    new product = _SX_64(argument2) * _SX_64(argument3);
  }
  if (highmult == 1) {
    product = _ZX_64(argument2) * _ZX_64(argument3);
  }
  if (highmult == 2) {
    product = _SX_64(argument2) * _ZX_64(argument3);
  }
  if (highmult) {
    new result1 = -product.64[1];
  } else {
    result1 = -product.64[0];
  }
stage E2:
  GPR[registerW] = result1;
\end{lstlisting}}
\newpage

\paragraph{Encoding} Format \textsf{ALU\_DMWRR.M} (Section~\ref{sec:format-ALU-DMWRR.M}), \verb|alucode2 = 01|\\

\bitfields{ 1/P, 2/00, 2/-- --, 27/upper27, 1/1, 2/11, 1/0, 2/01, 2/highmult, 6/registerW, 2/10, 3/110, 1/0, 1/splat32, 5/lower5, 6/registerZ }


\paragraph{Syntax} \texttt{mulnd}\emph{highmult} \emph{registerW} = \emph{registerZ}, \emph{upper27\_\discretionary{}{}{}lower5}\emph{splat32}

\paragraph{Description} The \emph{upper27\_\discretionary{}{}{}lower5} and the \emph{registerZ} double words are sign-extended or zero-extended from 64 bits, depending on the \emph{highmult}. The \emph{upper27\_\discretionary{}{}{}lower5} double word is multiplied by the \emph{registerZ} double word, and the product is optionally shifted right by 64 bits, depending on the \emph{highmult}. The resulting double word is negated and stored into the \emph{registerW}.


\paragraph{Modifier(s)}
\noindent\begin{itemize}
\item highmult (cf. highmult in section \ref{par:modifier-highmult} on page \pageref{par:modifier-highmult})
\item splat32 (cf. splat32 in section \ref{par:modifier-splat32} on page \pageref{par:modifier-splat32})
\end{itemize}

\paragraph{Execution} {Reservation table \textsf{ALU\_LITE.X} (Section~\ref{sec:reservation-ALU-LITE.X})}
~
{\begin{lstlisting}
stage ID:
  new splat32 = _ZX_1(splat32);
  new highmult = _ZX_2(highmult);
stage RR:
  new argument3 = splat32 ? (_WX_32(upper27_lower5) << 32) | _ZX_32(_WX_32(upper27_lower5)) : _SX_32(_WX_32(upper27_lower5));
  new argument2 = GPR[registerZ];
  if ((highmult == 0) || (highmult == 3)) {
    new product = _SX_64(argument2) * _SX_64(argument3);
  }
  if (highmult == 1) {
    product = _ZX_64(argument2) * _ZX_64(argument3);
  }
  if (highmult == 2) {
    product = _SX_64(argument2) * _ZX_64(argument3);
  }
  if (highmult) {
    new result1 = -product.64[1];
  } else {
    result1 = -product.64[0];
  }
stage E2:
  GPR[registerW] = result1;
\end{lstlisting}}
\newpage
\subsubsection[{\small MULNDP: Multiply Negate Double Word Pair}]{MULNDP\hfill Multiply Negate Double Word Pair}
\label{instr:MULNDP}\index{mulndp}
 
\paragraph{Encoding} Format \textsf{ALU\_DPMWRR} (Section~\ref{sec:format-ALU-DPMWRR}), \verb|alucode2 = 01|\\

\bitfields{ 1/P, 2/11, 1/1, 2/01, 2/highmult, 5/registerM, 1/--, 2/10, 3/110, 1/0, 5/registerO, 1/0, 5/registerP, 1/0 }


\paragraph{Syntax} \texttt{mulndp}\emph{highmult} \emph{registerM} = \emph{registerP}, \emph{registerO}

\paragraph{Description} The \emph{registerO} and the \emph{registerP} are interpreted as two double words packed into 128 bits. The \emph{registerO} and the \emph{registerP} double words are sign-extended or zero-extended from 64 bits, depending on the \emph{highmult}. The \emph{registerO} double words are multiplied by the \emph{registerP} double words, and the product is optionally shifted right by 64 bits, depending on the \emph{highmult}. The two resulting double words are negated, packed and stored into the \emph{registerM}.


\paragraph{Modifier(s)}
\noindent\begin{itemize}
\item highmult (cf. highmult in section \ref{par:modifier-highmult} on page \pageref{par:modifier-highmult})
\end{itemize}

\paragraph{Execution} {Reservation table \textsf{ALU\_LITE} (Section~\ref{sec:reservation-ALU-LITE})}
~
{\begin{lstlisting}
stage ID:
  new highmult = _ZX_2(highmult);
stage RR:
  new argument3 = PGR[registerO];
  new argument2 = PGR[registerP];
  for (I = 0; I < 2; I += 1) {
    if ((highmult == 0) || (highmult == 3)) {
      new product = _SX_64(argument2.64[I]) * _SX_64(argument3.64[I]);
    }
    if (highmult == 1) {
      product = _ZX_64(argument2.64[I]) * _ZX_64(argument3.64[I]);
    }
    if (highmult == 2) {
      product = _SX_64(argument2.64[I]) * _ZX_64(argument3.64[I]);
    }
    if (highmult) {
      result1.64[I] = -product.64[1];
    } else {
      result1.64[I] = -product.64[0];
    }
  };
stage E2:
  PGR[registerM] = result1;
\end{lstlisting}}
\newpage
\subsubsection[{\small MULNDQ: Multiply Negate Extend Double Word to Quadruple Word}]{MULNDQ\hfill Multiply Negate Extend Double Word to Quadruple Word}
\label{instr:MULNDQ}\index{mulndq}
 
\paragraph{Encoding} Format \textsf{ALU\_QMEWRR} (Section~\ref{sec:format-ALU-QMEWRR}), \verb|alucode2 = 01|\\

\bitfields{ 1/P, 2/11, 1/0, 2/01, 2/widemult, 5/registerM, 1/--, 2/10, 3/111, 1/1, 6/registerY, 6/registerZ }


\paragraph{Syntax} \texttt{mulndq}\emph{widemult} \emph{registerM} = \emph{registerZ}, \emph{registerY}

\paragraph{Description} The \emph{registerY} and the \emph{registerZ} double words are sign-extended or zero-extended from 64 bits, depending on the \emph{widemult}. The \emph{registerY} is multiplied by the \emph{registerZ}, producing a 128-bit result. The resulting quadruple word is negated and stored into the \emph{registerM}.


\paragraph{Modifier(s)}
\noindent\begin{itemize}
\item widemult (cf. widemult in section \ref{par:modifier-widemult} on page \pageref{par:modifier-widemult})
\end{itemize}

\paragraph{Execution} {Reservation table \textsf{ALU\_LITE} (Section~\ref{sec:reservation-ALU-LITE})}
~
{\begin{lstlisting}
stage ID:
  new widemult = _ZX_2(widemult);
stage RR:
  new argument3 = GPR[registerY];
  new argument2 = GPR[registerZ];
  if (widemult == 0) {
    new product = _SX_64(argument2) * _SX_64(argument3);
  }
  if (widemult == 1) {
    product = _ZX_64(argument2) * _ZX_64(argument3);
  }
  if (widemult == 2) {
    product = _SX_64(argument2) * _ZX_64(argument3);
  }
  new result1 = -product;
stage E2:
  PGR[registerM] = result1;
\end{lstlisting}}
\newpage
\subsubsection[{\small MULNHO: Multiply Negate Half Word Octuple}]{MULNHO\hfill Multiply Negate Half Word Octuple}
\label{instr:MULNHO}\index{mulnho}
 
\paragraph{Encoding} Format \textsf{ALU\_HOMWRR} (Section~\ref{sec:format-ALU-HOMWRR}), \verb|alucode2 = 01|\\

\bitfields{ 1/P, 2/11, 1/1, 2/01, 2/highmult, 5/registerM, 1/--, 2/10, 3/110, 1/1, 5/registerO, 1/0, 5/registerP, 1/0 }


\paragraph{Syntax} \texttt{mulnho}\emph{highmult} \emph{registerM} = \emph{registerP}, \emph{registerO}

\paragraph{Description} The \emph{registerO} and the \emph{registerP} are interpreted as eight half words packed into 128 bits. The \emph{registerO} and the \emph{registerP} half words are sign-extended or zero-extended from 16 bits, depending on the \emph{highmult}. The \emph{registerO} half words are multiplied by the \emph{registerP} half words, and the product is optionally shifted right by 16 bits, depending on the \emph{highmult}. The eight resulting half words are negated, packed and stored back into the \emph{registerM}.


\paragraph{Modifier(s)}
\noindent\begin{itemize}
\item highmult (cf. highmult in section \ref{par:modifier-highmult} on page \pageref{par:modifier-highmult})
\end{itemize}

\paragraph{Execution} {Reservation table \textsf{ALU\_LITE} (Section~\ref{sec:reservation-ALU-LITE})}
~
{\begin{lstlisting}
stage ID:
  new highmult = _ZX_2(highmult);
stage RR:
  new argument3 = PGR[registerO];
  new argument2 = PGR[registerP];
  for (I = 0; I < 8; I += 1) {
    if ((highmult == 0) || (highmult == 3)) {
      new product = _SX_16(argument2.16[I]) * _SX_16(argument3.16[I]);
    }
    if (highmult == 1) {
      product = _ZX_16(argument2.16[I]) * _ZX_16(argument3.16[I]);
    }
    if (highmult == 2) {
      product = _SX_16(argument2.16[I]) * _ZX_16(argument3.16[I]);
    }
    if (highmult) {
      result1.16[I] = -product.16[1];
    } else {
      result1.16[I] = -product.16[0];
    }
  };
stage E2:
  PGR[registerM] = result1;
\end{lstlisting}}
\newpage
\subsubsection[{\small MULNHWQ: Multiply Negate Extend Half Word to Word Quadruple}]{MULNHWQ\hfill Multiply Negate Extend Half Word to Word Quadruple}
\label{instr:MULNHWQ}\index{mulnhwq}
 
\paragraph{Encoding} Format \textsf{ALU\_WQMEWRRE} (Section~\ref{sec:format-ALU-WQMEWRRE}), \verb|alucode2 = 01|\\

\bitfields{ 1/P, 2/11, 1/1, 2/01, 2/widemult, 5/registerM, 1/0, 2/10, 3/110, 1/0, 5/registerYe, 1/1, 5/registerZe, 1/1 }


\paragraph{Syntax} \texttt{mulnhwq}\emph{widemult} \emph{registerM} = \emph{registerZe}, \emph{registerYe}

\paragraph{Description} The \emph{registerYe} and the \emph{registerZe} are interpreted as four half words packed into 64 bits. These \emph{registerYe} and \emph{registerZe} half words are sign-extended or zero-extended from 16 bits, depending on the \emph{widemult}. The \emph{registerYe} half words are multiplied by the \emph{registerZe} half words, producing four 32-bit results. The four resulting words are negated, packed and stored into the \emph{registerM}.


\paragraph{Modifier(s)}
\noindent\begin{itemize}
\item widemult (cf. widemult in section \ref{par:modifier-widemult} on page \pageref{par:modifier-widemult})
\end{itemize}

\paragraph{Execution} {Reservation table \textsf{ALU\_LITE} (Section~\ref{sec:reservation-ALU-LITE})}
~
{\begin{lstlisting}
stage ID:
  new widemult = _ZX_2(widemult);
stage RR:
  new argument3 = GPR[registerYe<<1];
  new argument2 = GPR[registerZe<<1];
  for (I = 0; I < 4; I += 1) {
    if (widemult == 0) {
      new product = _SX_16(argument2.16[I]) * _SX_16(argument3.16[I]);
    }
    if (widemult == 1) {
      product = _ZX_16(argument2.16[I]) * _ZX_16(argument3.16[I]);
    }
    if (widemult == 2) {
      product = _SX_16(argument2.16[I]) * _ZX_16(argument3.16[I]);
    }
    result1.32[I] = -product
  };
stage E2:
  PGR[registerM] = result1;
\end{lstlisting}}
\newpage

\paragraph{Encoding} Format \textsf{ALU\_WQMEWRRO} (Section~\ref{sec:format-ALU-WQMEWRRO}), \verb|alucode2 = 01|\\

\bitfields{ 1/P, 2/11, 1/1, 2/01, 2/widemult, 5/registerM, 1/1, 2/10, 3/110, 1/0, 5/registerYo, 1/1, 5/registerZo, 1/1 }


\paragraph{Syntax} \texttt{mulnhwq}\emph{widemult} \emph{registerM} = \emph{registerZo}, \emph{registerYo}

\paragraph{Description} The \emph{registerYo} and the \emph{registerZo} are interpreted as four half words packed into 64 bits. These \emph{registerYo} and \emph{registerZo} half words are sign-extended or zero-extended from 16 bits, depending on the \emph{widemult}. The \emph{registerYo} half words are multiplied by the \emph{registerZo} half words, producing four 32-bit results. The four resulting words are negated, packed and stored into the \emph{registerM}.


\paragraph{Modifier(s)}
\noindent\begin{itemize}
\item widemult (cf. widemult in section \ref{par:modifier-widemult} on page \pageref{par:modifier-widemult})
\end{itemize}

\paragraph{Execution} {Reservation table \textsf{ALU\_LITE} (Section~\ref{sec:reservation-ALU-LITE})}
~
{\begin{lstlisting}
stage ID:
  new widemult = _ZX_2(widemult);
stage RR:
  new argument3 = GPR[(registerYo<<1)+1];
  new argument2 = GPR[(registerZo<<1)+1];
  for (I = 0; I < 4; I += 1) {
    if (widemult == 0) {
      new product = _SX_16(argument2.16[I]) * _SX_16(argument3.16[I]);
    }
    if (widemult == 1) {
      product = _ZX_16(argument2.16[I]) * _ZX_16(argument3.16[I]);
    }
    if (widemult == 2) {
      product = _SX_16(argument2.16[I]) * _ZX_16(argument3.16[I]);
    }
    result1.32[I] = -product
  };
stage E2:
  PGR[registerM] = result1;
\end{lstlisting}}
\newpage
\subsubsection[{\small MULNW: Multiply Negate Word}]{MULNW\hfill Multiply Negate Word}
\label{instr:MULNW}\index{mulnw}
 
\paragraph{Encoding} Format \textsf{ALU\_WMWRR} (Section~\ref{sec:format-ALU-WMWRR}), \verb|alucode2 = 01|\\

\bitfields{ 1/P, 2/11, 1/0, 2/01, 2/highmult, 6/registerW, 2/11, 3/110, 1/signextw, 6/registerY, 6/registerZ }


\paragraph{Syntax} \texttt{mulnw}\emph{highmult}\emph{signextw} \emph{registerW} = \emph{registerZ}, \emph{registerY}

\paragraph{Description} The \emph{registerY} and the \emph{registerZ} words are sign-extended or zero-extended from 32 bits, depending on the \emph{highmult}. The \emph{registerY} word is multiplied by the \emph{registerZ} word, and the product is optionally shifted right by 32 bits, depending on the \emph{highmult}. The resulting word is negated and stored into the \emph{registerW}.


\paragraph{Modifier(s)}
\noindent\begin{itemize}
\item highmult (cf. highmult in section \ref{par:modifier-highmult} on page \pageref{par:modifier-highmult})
\item signextw (cf. signextw in section \ref{par:modifier-signextw} on page \pageref{par:modifier-signextw})
\end{itemize}

\paragraph{Execution} {Reservation table \textsf{ALU\_LITE} (Section~\ref{sec:reservation-ALU-LITE})}
~
{\begin{lstlisting}
stage ID:
  new signextw = _ZX_1(signextw);
  new highmult = _ZX_2(highmult);
stage RR:
  new argument3 = GPR[registerY];
  new argument2 = GPR[registerZ];
  if ((highmult == 0) || (highmult == 3)) {
    new product = _SX_32(argument2) * _SX_32(argument3);
  }
  if (highmult == 1) {
    product = _ZX_32(argument2) * _ZX_32(argument3);
  }
  if (highmult == 2) {
    product = _SX_32(argument2) * _ZX_32(argument3);
  }
  if (highmult) {
    new result1 = -product.32[1];
  } else {
    result1 = -product.32[0];
  }
stage E2:
  GPR[registerW] = signextw ? _SX_32(result1) : _ZX_32(result1);
\end{lstlisting}}
\newpage

\paragraph{Encoding} Format \textsf{ALU\_WMWRR.W} (Section~\ref{sec:format-ALU-WMWRR.W}), \verb|alucode2 = 01|\\

\bitfields{ 1/P, 2/00, 2/-- --, 27/upper27, 1/1, 2/11, 1/0, 2/01, 2/highmult, 6/registerW, 2/11, 3/110, 1/signextw, 1/0, 5/lower5, 6/registerZ }


\paragraph{Syntax} \texttt{mulnw}\emph{highmult}\emph{signextw} \emph{registerW} = \emph{registerZ}, \emph{upper27\_\discretionary{}{}{}lower5}

\paragraph{Description} The \emph{upper27\_\discretionary{}{}{}lower5} and the \emph{registerZ} words are sign-extended or zero-extended from 32 bits, depending on the \emph{highmult}. The \emph{upper27\_\discretionary{}{}{}lower5} word is multiplied by the \emph{registerZ} word, and the product is optionally shifted right by 32 bits, depending on the \emph{highmult}. The resulting word is negated and stored into the \emph{registerW}.


\paragraph{Modifier(s)}
\noindent\begin{itemize}
\item highmult (cf. highmult in section \ref{par:modifier-highmult} on page \pageref{par:modifier-highmult})
\item signextw (cf. signextw in section \ref{par:modifier-signextw} on page \pageref{par:modifier-signextw})
\end{itemize}

\paragraph{Execution} {Reservation table \textsf{ALU\_LITE.X} (Section~\ref{sec:reservation-ALU-LITE.X})}
~
{\begin{lstlisting}
stage ID:
  new signextw = _ZX_1(signextw);
  new highmult = _ZX_2(highmult);
stage RR:
  new argument3 = _WX_32(upper27_lower5);
  new argument2 = GPR[registerZ];
  if ((highmult == 0) || (highmult == 3)) {
    new product = _SX_32(argument2) * _SX_32(argument3);
  }
  if (highmult == 1) {
    product = _ZX_32(argument2) * _ZX_32(argument3);
  }
  if (highmult == 2) {
    product = _SX_32(argument2) * _ZX_32(argument3);
  }
  if (highmult) {
    new result1 = -product.32[1];
  } else {
    result1 = -product.32[0];
  }
stage E2:
  GPR[registerW] = signextw ? _SX_32(result1) : _ZX_32(result1);
\end{lstlisting}}
\newpage
\subsubsection[{\small MULNWD: Multiply Negate Extend Word to Double Word}]{MULNWD\hfill Multiply Negate Extend Word to Double Word}
\label{instr:MULNWD}\index{mulnwd}
 
\paragraph{Encoding} Format \textsf{ALU\_DMEWRR} (Section~\ref{sec:format-ALU-DMEWRR}), \verb|alucode2 = 01|\\

\bitfields{ 1/P, 2/11, 1/0, 2/01, 2/widemult, 6/registerW, 2/10, 3/111, 1/0, 6/registerY, 6/registerZ }


\paragraph{Syntax} \texttt{mulnwd}\emph{widemult} \emph{registerW} = \emph{registerZ}, \emph{registerY}

\paragraph{Description} The \emph{registerY} and the \emph{registerZ} words are sign-extended or zero-extended from 32 bits, depending on the \emph{widemult}. The \emph{registerY} is multiplied by the \emph{registerZ}, producing a 64-bit result. The resulting double word is negated and stored into the \emph{registerW}.


\paragraph{Modifier(s)}
\noindent\begin{itemize}
\item widemult (cf. widemult in section \ref{par:modifier-widemult} on page \pageref{par:modifier-widemult})
\end{itemize}

\paragraph{Execution} {Reservation table \textsf{ALU\_LITE} (Section~\ref{sec:reservation-ALU-LITE})}
~
{\begin{lstlisting}
stage ID:
  new widemult = _ZX_2(widemult);
stage RR:
  new argument3 = GPR[registerY];
  new argument2 = GPR[registerZ];
  if (widemult == 0) {
    new product = _SX_32(argument2) * _SX_32(argument3);
  }
  if (widemult == 1) {
    product = _ZX_32(argument2) * _ZX_32(argument3);
  }
  if (widemult == 2) {
    product = _SX_32(argument2) * _ZX_32(argument3);
  }
  new result1 = -product;
stage E2:
  GPR[registerW] = result1;
\end{lstlisting}}
\newpage

\paragraph{Encoding} Format \textsf{ALU\_DMEWRR.M} (Section~\ref{sec:format-ALU-DMEWRR.M}), \verb|alucode2 = 01|\\

\bitfields{ 1/P, 2/00, 2/-- --, 27/upper27, 1/1, 2/11, 1/0, 2/01, 2/widemult, 6/registerW, 2/10, 3/111, 1/0, 1/splat32, 5/lower5, 6/registerZ }


\paragraph{Syntax} \texttt{mulnwd}\emph{widemult} \emph{registerW} = \emph{registerZ}, \emph{upper27\_\discretionary{}{}{}lower5}\emph{splat32}

\paragraph{Description} The \emph{upper27\_\discretionary{}{}{}lower5} and the \emph{registerZ} words are sign-extended or zero-extended from 32 bits, depending on the \emph{widemult}. The \emph{upper27\_\discretionary{}{}{}lower5} is multiplied by the \emph{registerZ}, producing a 64-bit result. The resulting double word is negated and stored into the \emph{registerW}.


\paragraph{Modifier(s)}
\noindent\begin{itemize}
\item widemult (cf. widemult in section \ref{par:modifier-widemult} on page \pageref{par:modifier-widemult})
\item splat32 (cf. splat32 in section \ref{par:modifier-splat32} on page \pageref{par:modifier-splat32})
\end{itemize}

\paragraph{Execution} {Reservation table \textsf{ALU\_LITE.X} (Section~\ref{sec:reservation-ALU-LITE.X})}
~
{\begin{lstlisting}
stage ID:
  new splat32 = _ZX_1(splat32);
  new widemult = _ZX_2(widemult);
stage RR:
  new argument3 = splat32 ? (_WX_32(upper27_lower5) << 32) | _ZX_32(_WX_32(upper27_lower5)) : _SX_32(_WX_32(upper27_lower5));
  new argument2 = GPR[registerZ];
  if (widemult == 0) {
    new product = _SX_32(argument2) * _SX_32(argument3);
  }
  if (widemult == 1) {
    product = _ZX_32(argument2) * _ZX_32(argument3);
  }
  if (widemult == 2) {
    product = _SX_32(argument2) * _ZX_32(argument3);
  }
  new result1 = -product;
stage E2:
  GPR[registerW] = result1;
\end{lstlisting}}
\newpage
\subsubsection[{\small MULNWDP: Multiply Negate Extend Word to Double Word Pair}]{MULNWDP\hfill Multiply Negate Extend Word to Double Word Pair}
\label{instr:MULNWDP}\index{mulnwdp}
 
\paragraph{Encoding} Format \textsf{ALU\_DPMEWRRE} (Section~\ref{sec:format-ALU-DPMEWRRE}), \verb|alucode2 = 01|\\

\bitfields{ 1/P, 2/11, 1/1, 2/01, 2/widemult, 5/registerM, 1/0, 2/10, 3/110, 1/0, 5/registerYe, 1/1, 5/registerZe, 1/0 }


\paragraph{Syntax} \texttt{mulnwdp}\emph{widemult} \emph{registerM} = \emph{registerZe}, \emph{registerYe}

\paragraph{Description} The \emph{registerYe} and the \emph{registerZe} are interpreted as two words packed into 64 bits. These \emph{registerYe} and \emph{registerZe} words are sign-extended or zero-extended from 32 bits, depending on the \emph{widemult}. The \emph{registerYe} words are multiplied by the \emph{registerZe} words, producing two 64-bit results. The two resulting double words are negated, packed and stored into the \emph{registerM}.


\paragraph{Modifier(s)}
\noindent\begin{itemize}
\item widemult (cf. widemult in section \ref{par:modifier-widemult} on page \pageref{par:modifier-widemult})
\end{itemize}

\paragraph{Execution} {Reservation table \textsf{ALU\_LITE} (Section~\ref{sec:reservation-ALU-LITE})}
~
{\begin{lstlisting}
stage ID:
  new widemult = _ZX_2(widemult);
stage RR:
  new argument3 = GPR[registerYe<<1];
  new argument2 = GPR[registerZe<<1];
  for (I = 0; I < 2; I += 1) {
    if (widemult == 0) {
      new product = _SX_32(argument2.32[I]) * _SX_32(argument3.32[I]);
    }
    if (widemult == 1) {
      product = _ZX_32(argument2.32[I]) * _ZX_32(argument3.32[I]);
    }
    if (widemult == 2) {
      product = _SX_32(argument2.32[I]) * _ZX_32(argument3.32[I]);
    }
    result1.64[I] = -product
  };
stage E2:
  PGR[registerM] = result1;
\end{lstlisting}}
\newpage

\paragraph{Encoding} Format \textsf{ALU\_DPMEWRRO} (Section~\ref{sec:format-ALU-DPMEWRRO}), \verb|alucode2 = 01|\\

\bitfields{ 1/P, 2/11, 1/1, 2/01, 2/widemult, 5/registerM, 1/1, 2/10, 3/110, 1/0, 5/registerYo, 1/1, 5/registerZo, 1/0 }


\paragraph{Syntax} \texttt{mulnwdp}\emph{widemult} \emph{registerM} = \emph{registerZo}, \emph{registerYo}

\paragraph{Description} The \emph{registerYo} and the \emph{registerZo} are interpreted as two words packed into 64 bits. These \emph{registerYo} and \emph{registerZo} words are sign-extended or zero-extended from 32 bits, depending on the \emph{widemult}. The \emph{registerYo} words are multiplied by the \emph{registerZo} words, producing two 64-bit results. The two resulting double words are negated, packed and stored into the \emph{registerM}.


\paragraph{Modifier(s)}
\noindent\begin{itemize}
\item widemult (cf. widemult in section \ref{par:modifier-widemult} on page \pageref{par:modifier-widemult})
\end{itemize}

\paragraph{Execution} {Reservation table \textsf{ALU\_LITE} (Section~\ref{sec:reservation-ALU-LITE})}
~
{\begin{lstlisting}
stage ID:
  new widemult = _ZX_2(widemult);
stage RR:
  new argument3 = GPR[(registerYo<<1)+1];
  new argument2 = GPR[(registerZo<<1)+1];
  for (I = 0; I < 2; I += 1) {
    if (widemult == 0) {
      new product = _SX_32(argument2.32[I]) * _SX_32(argument3.32[I]);
    }
    if (widemult == 1) {
      product = _ZX_32(argument2.32[I]) * _ZX_32(argument3.32[I]);
    }
    if (widemult == 2) {
      product = _SX_32(argument2.32[I]) * _ZX_32(argument3.32[I]);
    }
    result1.64[I] = -product
  };
stage E2:
  PGR[registerM] = result1;
\end{lstlisting}}
\newpage
\subsubsection[{\small MULNWQ: Multiply Negate Word Quadruple}]{MULNWQ\hfill Multiply Negate Word Quadruple}
\label{instr:MULNWQ}\index{mulnwq}
 
\paragraph{Encoding} Format \textsf{ALU\_WQMWRR} (Section~\ref{sec:format-ALU-WQMWRR}), \verb|alucode2 = 01|\\

\bitfields{ 1/P, 2/11, 1/1, 2/01, 2/highmult, 5/registerM, 1/--, 2/10, 3/110, 1/0, 5/registerO, 1/0, 5/registerP, 1/1 }


\paragraph{Syntax} \texttt{mulnwq}\emph{highmult} \emph{registerM} = \emph{registerP}, \emph{registerO}

\paragraph{Description} The \emph{registerO} and the \emph{registerP} are interpreted as four words packed into 128 bits. The \emph{registerO} and the \emph{registerP} words are sign-extended or zero-extended from 32 bits, depending on the \emph{highmult}. The \emph{registerO} words are multiplied by the \emph{registerP} words, and the product is optionally shifted right by 32 bits, depending on the \emph{highmult}. The four resulting words are negated, packed and stored back into the \emph{registerM}.


\paragraph{Modifier(s)}
\noindent\begin{itemize}
\item highmult (cf. highmult in section \ref{par:modifier-highmult} on page \pageref{par:modifier-highmult})
\end{itemize}

\paragraph{Execution} {Reservation table \textsf{ALU\_LITE} (Section~\ref{sec:reservation-ALU-LITE})}
~
{\begin{lstlisting}
stage ID:
  new highmult = _ZX_2(highmult);
stage RR:
  new argument3 = PGR[registerO];
  new argument2 = PGR[registerP];
  for (I = 0; I < 4; I += 1) {
    if ((highmult == 0) || (highmult == 3)) {
      new product = _SX_32(argument2.32[I]) * _SX_32(argument3.32[I]);
    }
    if (highmult == 1) {
      product = _ZX_32(argument2.32[I]) * _ZX_32(argument3.32[I]);
    }
    if (highmult == 2) {
      product = _SX_32(argument2.32[I]) * _ZX_32(argument3.32[I]);
    }
    if (highmult) {
      result1.32[I] = -product.32[1];
    } else {
      result1.32[I] = -product.32[0];
    }
  };
stage E2:
  PGR[registerM] = result1;
\end{lstlisting}}
\newpage
\subsubsection[{\small MULNXBHO: Multiply Negate Extend Byte to Half Word Octuple}]{MULNXBHO\hfill Multiply Negate Extend Byte to Half Word Octuple}
\label{instr:MULNXBHO}\index{mulnxbho}
 
\paragraph{Encoding} Format \textsf{ALU\_HOMXWRR} (Section~\ref{sec:format-ALU-HOMXWRR}), \verb|alucode2 = 01|\\

\bitfields{ 1/P, 2/11, 1/1, 2/01, 2/widemult, 5/registerM, 1/oddlanes, 2/10, 3/111, 1/1, 5/registerO, 1/--, 5/registerP, 1/0 }


\paragraph{Syntax} \texttt{mulnxbho}\emph{oddlanes}\emph{widemult} \emph{registerM} = \emph{registerP}, \emph{registerO}

\paragraph{Description} The \emph{registerO} and the \emph{registerP} are interpreted as sixteen bytes packed into 128 bits. The even or odd bytes of the \emph{registerO} and the \emph{registerP} are selected, depending on \emph{oddlanes}. These \emph{registerO} and \emph{registerP} bytes are sign-extended or zero-extended from 8 bits, depending on the \emph{widemult}. The \emph{registerO} bytes are multiplied by the \emph{registerP} bytes, producing eight 16-bit results. The eight resulting half words are negated, packed and stored into the \emph{registerM}.


\paragraph{Modifier(s)}
\noindent\begin{itemize}
\item widemult (cf. widemult in section \ref{par:modifier-widemult} on page \pageref{par:modifier-widemult})
\item oddlanes (cf. oddlanes in section \ref{par:modifier-oddlanes} on page \pageref{par:modifier-oddlanes})
\end{itemize}

\paragraph{Execution} {Reservation table \textsf{ALU\_LITE} (Section~\ref{sec:reservation-ALU-LITE})}
~
{\begin{lstlisting}
stage ID:
  new oddlanes = _ZX_1(oddlanes);
  new widemult = _ZX_2(widemult);
stage RR:
  new argument3 = PGR[registerO];
  new argument2 = PGR[registerP];
  if (oddlanes) {
    argument2 = argument2 >> 8;
    argument3 = argument3 >> 8;
  }
  for (I = 0; I < 8; I += 1) {
    if (widemult == 0) {
      new product = _SX_8(argument2.16[I]) * _SX_8(argument3.16[I]);
    }
    if (widemult == 1) {
      product = _ZX_8(argument2.16[I]) * _ZX_8(argument3.16[I]);
    }
    if (widemult == 2) {
      product = _SX_8(argument2.16[I]) * _ZX_8(argument3.16[I]);
    }
    result1.16[I] = -product
  };
stage E2:
  PGR[registerM] = result1;
\end{lstlisting}}
\newpage
\subsubsection[{\small MULNXHWQ: Multiply Negate Extend Half Word to Word Quadruple}]{MULNXHWQ\hfill Multiply Negate Extend Half Word to Word Quadruple}
\label{instr:MULNXHWQ}\index{mulnxhwq}
 
\paragraph{Encoding} Format \textsf{ALU\_WQMXWRR} (Section~\ref{sec:format-ALU-WQMXWRR}), \verb|alucode2 = 01|\\

\bitfields{ 1/P, 2/11, 1/1, 2/01, 2/widemult, 5/registerM, 1/oddlanes, 2/10, 3/111, 1/0, 5/registerO, 1/--, 5/registerP, 1/1 }


\paragraph{Syntax} \texttt{mulnxhwq}\emph{oddlanes}\emph{widemult} \emph{registerM} = \emph{registerP}, \emph{registerO}

\paragraph{Description} The \emph{registerO} and the \emph{registerP} are interpreted as eight half words packed into 128 bits. The even or odd half words of the \emph{registerO} and the \emph{registerP} are selected, depending on \emph{oddlanes}. These \emph{registerO} and \emph{registerP} half words are sign-extended or zero-extended from 16 bits, depending on the \emph{widemult}. The \emph{registerO} half words are multiplied by the \emph{registerP} half words, producing four 32-bit results. The four resulting words are negated, packed and stored into the \emph{registerM}.


\paragraph{Modifier(s)}
\noindent\begin{itemize}
\item widemult (cf. widemult in section \ref{par:modifier-widemult} on page \pageref{par:modifier-widemult})
\item oddlanes (cf. oddlanes in section \ref{par:modifier-oddlanes} on page \pageref{par:modifier-oddlanes})
\end{itemize}

\paragraph{Execution} {Reservation table \textsf{ALU\_LITE} (Section~\ref{sec:reservation-ALU-LITE})}
~
{\begin{lstlisting}
stage ID:
  new oddlanes = _ZX_1(oddlanes);
  new widemult = _ZX_2(widemult);
stage RR:
  new argument3 = PGR[registerO];
  new argument2 = PGR[registerP];
  if (oddlanes) {
    argument2 = argument2 >> 16;
    argument3 = argument3 >> 16;
  }
  for (I = 0; I < 4; I += 1) {
    if (widemult == 0) {
      new product = _SX_16(argument2.32[I]) * _SX_16(argument3.32[I]);
    }
    if (widemult == 1) {
      product = _ZX_16(argument2.32[I]) * _ZX_16(argument3.32[I]);
    }
    if (widemult == 2) {
      product = _SX_16(argument2.32[I]) * _ZX_16(argument3.32[I]);
    }
    result1.32[I] = -product
  };
stage E2:
  PGR[registerM] = result1;
\end{lstlisting}}
\newpage
\subsubsection[{\small MULNXWDP: Multiply Negate Extend Word to Double Word Pair}]{MULNXWDP\hfill Multiply Negate Extend Word to Double Word Pair}
\label{instr:MULNXWDP}\index{mulnxwdp}
 
\paragraph{Encoding} Format \textsf{ALU\_DPMXWRR} (Section~\ref{sec:format-ALU-DPMXWRR}), \verb|alucode2 = 01|\\

\bitfields{ 1/P, 2/11, 1/1, 2/01, 2/widemult, 5/registerM, 1/oddlanes, 2/10, 3/111, 1/0, 5/registerO, 1/--, 5/registerP, 1/0 }


\paragraph{Syntax} \texttt{mulnxwdp}\emph{oddlanes}\emph{widemult} \emph{registerM} = \emph{registerP}, \emph{registerO}

\paragraph{Description} The \emph{registerO} and the \emph{registerP} are interpreted as four words packed into 128 bits. The even or odd words of the \emph{registerO} and the \emph{registerP} are selected, depending on \emph{oddlanes}. These \emph{registerO} and \emph{registerP} words are sign-extended or zero-extended from 32 bits, depending on the \emph{widemult}. The \emph{registerO} words are multiplied by the \emph{registerP} words, producing two 64-bit results. The two resulting double words are negated, packed and stored into the \emph{registerM}.


\paragraph{Modifier(s)}
\noindent\begin{itemize}
\item widemult (cf. widemult in section \ref{par:modifier-widemult} on page \pageref{par:modifier-widemult})
\item oddlanes (cf. oddlanes in section \ref{par:modifier-oddlanes} on page \pageref{par:modifier-oddlanes})
\end{itemize}

\paragraph{Execution} {Reservation table \textsf{ALU\_LITE} (Section~\ref{sec:reservation-ALU-LITE})}
~
{\begin{lstlisting}
stage ID:
  new oddlanes = _ZX_1(oddlanes);
  new widemult = _ZX_2(widemult);
stage RR:
  new argument3 = PGR[registerO];
  new argument2 = PGR[registerP];
  if (oddlanes) {
    argument2 = argument2 >> 32;
    argument3 = argument3 >> 32;
  }
  for (I = 0; I < 2; I += 1) {
    if (widemult == 0) {
      new product = _SX_32(argument2.64[I]) * _SX_32(argument3.64[I]);
    }
    if (widemult == 1) {
      product = _ZX_32(argument2.64[I]) * _ZX_32(argument3.64[I]);
    }
    if (widemult == 2) {
      product = _SX_32(argument2.64[I]) * _ZX_32(argument3.64[I]);
    }
    result1.64[I] = -product
  };
stage E2:
  PGR[registerM] = result1;
\end{lstlisting}}
\newpage
\subsubsection[{\small MULW: Multiply Word}]{MULW\hfill Multiply Word}
\label{instr:MULW}\index{mulw}
 
\paragraph{Encoding} Format \textsf{ALU\_WMWRR} (Section~\ref{sec:format-ALU-WMWRR}), \verb|alucode2 = 00|\\

\bitfields{ 1/P, 2/11, 1/0, 2/00, 2/highmult, 6/registerW, 2/11, 3/110, 1/signextw, 6/registerY, 6/registerZ }


\paragraph{Syntax} \texttt{mulw}\emph{highmult}\emph{signextw} \emph{registerW} = \emph{registerZ}, \emph{registerY}

\paragraph{Description} The \emph{registerY} and the \emph{registerZ} words are sign-extended or zero-extended from 32 bits, depending on the \emph{highmult}. The \emph{registerY} word is multiplied by the \emph{registerZ} word, and the product is optionally shifted right by 32 bits, depending on the \emph{highmult}. The resulting word is stored into the \emph{registerW}.


\paragraph{Modifier(s)}
\noindent\begin{itemize}
\item highmult (cf. highmult in section \ref{par:modifier-highmult} on page \pageref{par:modifier-highmult})
\item signextw (cf. signextw in section \ref{par:modifier-signextw} on page \pageref{par:modifier-signextw})
\end{itemize}

\paragraph{Execution} {Reservation table \textsf{ALU\_LITE} (Section~\ref{sec:reservation-ALU-LITE})}
~
{\begin{lstlisting}
stage ID:
  new signextw = _ZX_1(signextw);
  new highmult = _ZX_2(highmult);
stage RR:
  new argument3 = GPR[registerY];
  new argument2 = GPR[registerZ];
  if ((highmult == 0) || (highmult == 3)) {
    new product = _SX_32(argument2) * _SX_32(argument3);
  }
  if (highmult == 1) {
    product = _ZX_32(argument2) * _ZX_32(argument3);
  }
  if (highmult == 2) {
    product = _SX_32(argument2) * _ZX_32(argument3);
  }
  if (highmult) {
    new result1 = product.32[1];
  } else {
    result1 = product.32[0];
  }
stage E2:
  GPR[registerW] = signextw ? _SX_32(result1) : _ZX_32(result1);
\end{lstlisting}}
\newpage

\paragraph{Encoding} Format \textsf{ALU\_WMWRR.W} (Section~\ref{sec:format-ALU-WMWRR.W}), \verb|alucode2 = 00|\\

\bitfields{ 1/P, 2/00, 2/-- --, 27/upper27, 1/1, 2/11, 1/0, 2/00, 2/highmult, 6/registerW, 2/11, 3/110, 1/signextw, 1/0, 5/lower5, 6/registerZ }


\paragraph{Syntax} \texttt{mulw}\emph{highmult}\emph{signextw} \emph{registerW} = \emph{registerZ}, \emph{upper27\_\discretionary{}{}{}lower5}

\paragraph{Description} The \emph{upper27\_\discretionary{}{}{}lower5} and the \emph{registerZ} words are sign-extended or zero-extended from 32 bits, depending on the \emph{highmult}. The \emph{upper27\_\discretionary{}{}{}lower5} word is multiplied by the \emph{registerZ} word, and the product is optionally shifted right by 32 bits, depending on the \emph{highmult}. The resulting word is stored into the \emph{registerW}.


\paragraph{Modifier(s)}
\noindent\begin{itemize}
\item highmult (cf. highmult in section \ref{par:modifier-highmult} on page \pageref{par:modifier-highmult})
\item signextw (cf. signextw in section \ref{par:modifier-signextw} on page \pageref{par:modifier-signextw})
\end{itemize}

\paragraph{Execution} {Reservation table \textsf{ALU\_LITE.X} (Section~\ref{sec:reservation-ALU-LITE.X})}
~
{\begin{lstlisting}
stage ID:
  new signextw = _ZX_1(signextw);
  new highmult = _ZX_2(highmult);
stage RR:
  new argument3 = _WX_32(upper27_lower5);
  new argument2 = GPR[registerZ];
  if ((highmult == 0) || (highmult == 3)) {
    new product = _SX_32(argument2) * _SX_32(argument3);
  }
  if (highmult == 1) {
    product = _ZX_32(argument2) * _ZX_32(argument3);
  }
  if (highmult == 2) {
    product = _SX_32(argument2) * _ZX_32(argument3);
  }
  if (highmult) {
    new result1 = product.32[1];
  } else {
    result1 = product.32[0];
  }
stage E2:
  GPR[registerW] = signextw ? _SX_32(result1) : _ZX_32(result1);
\end{lstlisting}}
\newpage
\subsubsection[{\small MULWD: Multiply Extend Word to Double Word}]{MULWD\hfill Multiply Extend Word to Double Word}
\label{instr:MULWD}\index{mulwd}
 
\paragraph{Encoding} Format \textsf{ALU\_DMEWRR} (Section~\ref{sec:format-ALU-DMEWRR}), \verb|alucode2 = 00|\\

\bitfields{ 1/P, 2/11, 1/0, 2/00, 2/widemult, 6/registerW, 2/10, 3/111, 1/0, 6/registerY, 6/registerZ }


\paragraph{Syntax} \texttt{mulwd}\emph{widemult} \emph{registerW} = \emph{registerZ}, \emph{registerY}

\paragraph{Description} The \emph{registerY} and the \emph{registerZ} words are sign-extended or zero-extended from 32 bits, depending on the \emph{widemult}. The \emph{registerY} is multiplied by the \emph{registerZ}, producing a 64-bit result. The resulting double word is stored into the \emph{registerW}.


\paragraph{Modifier(s)}
\noindent\begin{itemize}
\item widemult (cf. widemult in section \ref{par:modifier-widemult} on page \pageref{par:modifier-widemult})
\end{itemize}

\paragraph{Execution} {Reservation table \textsf{ALU\_LITE} (Section~\ref{sec:reservation-ALU-LITE})}
~
{\begin{lstlisting}
stage ID:
  new widemult = _ZX_2(widemult);
stage RR:
  new argument3 = GPR[registerY];
  new argument2 = GPR[registerZ];
  if (widemult == 0) {
    new product = _SX_32(argument2) * _SX_32(argument3);
  }
  if (widemult == 1) {
    product = _ZX_32(argument2) * _ZX_32(argument3);
  }
  if (widemult == 2) {
    product = _SX_32(argument2) * _ZX_32(argument3);
  }
  new result1 = product;
stage E2:
  GPR[registerW] = result1;
\end{lstlisting}}
\newpage

\paragraph{Encoding} Format \textsf{ALU\_DMEWRR.M} (Section~\ref{sec:format-ALU-DMEWRR.M}), \verb|alucode2 = 00|\\

\bitfields{ 1/P, 2/00, 2/-- --, 27/upper27, 1/1, 2/11, 1/0, 2/00, 2/widemult, 6/registerW, 2/10, 3/111, 1/0, 1/splat32, 5/lower5, 6/registerZ }


\paragraph{Syntax} \texttt{mulwd}\emph{widemult} \emph{registerW} = \emph{registerZ}, \emph{upper27\_\discretionary{}{}{}lower5}\emph{splat32}

\paragraph{Description} The \emph{upper27\_\discretionary{}{}{}lower5} and the \emph{registerZ} words are sign-extended or zero-extended from 32 bits, depending on the \emph{widemult}. The \emph{upper27\_\discretionary{}{}{}lower5} is multiplied by the \emph{registerZ}, producing a 64-bit result. The resulting double word is stored into the \emph{registerW}.


\paragraph{Modifier(s)}
\noindent\begin{itemize}
\item widemult (cf. widemult in section \ref{par:modifier-widemult} on page \pageref{par:modifier-widemult})
\item splat32 (cf. splat32 in section \ref{par:modifier-splat32} on page \pageref{par:modifier-splat32})
\end{itemize}

\paragraph{Execution} {Reservation table \textsf{ALU\_LITE.X} (Section~\ref{sec:reservation-ALU-LITE.X})}
~
{\begin{lstlisting}
stage ID:
  new splat32 = _ZX_1(splat32);
  new widemult = _ZX_2(widemult);
stage RR:
  new argument3 = splat32 ? (_WX_32(upper27_lower5) << 32) | _ZX_32(_WX_32(upper27_lower5)) : _SX_32(_WX_32(upper27_lower5));
  new argument2 = GPR[registerZ];
  if (widemult == 0) {
    new product = _SX_32(argument2) * _SX_32(argument3);
  }
  if (widemult == 1) {
    product = _ZX_32(argument2) * _ZX_32(argument3);
  }
  if (widemult == 2) {
    product = _SX_32(argument2) * _ZX_32(argument3);
  }
  new result1 = product;
stage E2:
  GPR[registerW] = result1;
\end{lstlisting}}
\newpage
\subsubsection[{\small MULWDP: Multiply Extend Word to Double Word Pair}]{MULWDP\hfill Multiply Extend Word to Double Word Pair}
\label{instr:MULWDP}\index{mulwdp}
 
\paragraph{Encoding} Format \textsf{ALU\_DPMEWRRE} (Section~\ref{sec:format-ALU-DPMEWRRE}), \verb|alucode2 = 00|\\

\bitfields{ 1/P, 2/11, 1/1, 2/00, 2/widemult, 5/registerM, 1/0, 2/10, 3/110, 1/0, 5/registerYe, 1/1, 5/registerZe, 1/0 }


\paragraph{Syntax} \texttt{mulwdp}\emph{widemult} \emph{registerM} = \emph{registerZe}, \emph{registerYe}

\paragraph{Description} The \emph{registerYe} and the \emph{registerZe} are interpreted as two words packed into 64 bits. These \emph{registerYe} and \emph{registerZe} words are sign-extended or zero-extended from 32 bits, depending on the \emph{widemult}. The \emph{registerYe} words are multiplied by the \emph{registerZe} words, producing two 64-bit results. The two resulting double words are packed and stored into the \emph{registerM}.


\paragraph{Modifier(s)}
\noindent\begin{itemize}
\item widemult (cf. widemult in section \ref{par:modifier-widemult} on page \pageref{par:modifier-widemult})
\end{itemize}

\paragraph{Execution} {Reservation table \textsf{ALU\_LITE} (Section~\ref{sec:reservation-ALU-LITE})}
~
{\begin{lstlisting}
stage ID:
  new widemult = _ZX_2(widemult);
stage RR:
  new argument3 = GPR[registerYe<<1];
  new argument2 = GPR[registerZe<<1];
  for (I = 0; I < 2; I += 1) {
    if (widemult == 0) {
      new product = _SX_32(argument2.32[I]) * _SX_32(argument3.32[I]);
    }
    if (widemult == 1) {
      product = _ZX_32(argument2.32[I]) * _ZX_32(argument3.32[I]);
    }
    if (widemult == 2) {
      product = _SX_32(argument2.32[I]) * _ZX_32(argument3.32[I]);
    }
    result1.64[I] = product
  };
stage E2:
  PGR[registerM] = result1;
\end{lstlisting}}
\newpage

\paragraph{Encoding} Format \textsf{ALU\_DPMEWRRO} (Section~\ref{sec:format-ALU-DPMEWRRO}), \verb|alucode2 = 00|\\

\bitfields{ 1/P, 2/11, 1/1, 2/00, 2/widemult, 5/registerM, 1/1, 2/10, 3/110, 1/0, 5/registerYo, 1/1, 5/registerZo, 1/0 }


\paragraph{Syntax} \texttt{mulwdp}\emph{widemult} \emph{registerM} = \emph{registerZo}, \emph{registerYo}

\paragraph{Description} The \emph{registerYo} and the \emph{registerZo} are interpreted as two words packed into 64 bits. These \emph{registerYo} and \emph{registerZo} words are sign-extended or zero-extended from 32 bits, depending on the \emph{widemult}. The \emph{registerYo} words are multiplied by the \emph{registerZo} words, producing two 64-bit results. The two resulting double words are packed and stored into the \emph{registerM}.


\paragraph{Modifier(s)}
\noindent\begin{itemize}
\item widemult (cf. widemult in section \ref{par:modifier-widemult} on page \pageref{par:modifier-widemult})
\end{itemize}

\paragraph{Execution} {Reservation table \textsf{ALU\_LITE} (Section~\ref{sec:reservation-ALU-LITE})}
~
{\begin{lstlisting}
stage ID:
  new widemult = _ZX_2(widemult);
stage RR:
  new argument3 = GPR[(registerYo<<1)+1];
  new argument2 = GPR[(registerZo<<1)+1];
  for (I = 0; I < 2; I += 1) {
    if (widemult == 0) {
      new product = _SX_32(argument2.32[I]) * _SX_32(argument3.32[I]);
    }
    if (widemult == 1) {
      product = _ZX_32(argument2.32[I]) * _ZX_32(argument3.32[I]);
    }
    if (widemult == 2) {
      product = _SX_32(argument2.32[I]) * _ZX_32(argument3.32[I]);
    }
    result1.64[I] = product
  };
stage E2:
  PGR[registerM] = result1;
\end{lstlisting}}
\newpage
\subsubsection[{\small MULWQ: Multiply Word Quadruple}]{MULWQ\hfill Multiply Word Quadruple}
\label{instr:MULWQ}\index{mulwq}
 
\paragraph{Encoding} Format \textsf{ALU\_WQMWRR} (Section~\ref{sec:format-ALU-WQMWRR}), \verb|alucode2 = 00|\\

\bitfields{ 1/P, 2/11, 1/1, 2/00, 2/highmult, 5/registerM, 1/--, 2/10, 3/110, 1/0, 5/registerO, 1/0, 5/registerP, 1/1 }


\paragraph{Syntax} \texttt{mulwq}\emph{highmult} \emph{registerM} = \emph{registerP}, \emph{registerO}

\paragraph{Description} The \emph{registerO} and the \emph{registerP} are interpreted as four words packed into 128 bits. The \emph{registerO} and the \emph{registerP} words are sign-extended or zero-extended from 32 bits, depending on the \emph{highmult}. The \emph{registerO} words are multiplied by the \emph{registerP} words, and the product is optionally shifted right by 32 bits, depending on the \emph{highmult}. The four resulting words are packed and stored back into the \emph{registerM}.


\paragraph{Modifier(s)}
\noindent\begin{itemize}
\item highmult (cf. highmult in section \ref{par:modifier-highmult} on page \pageref{par:modifier-highmult})
\end{itemize}

\paragraph{Execution} {Reservation table \textsf{ALU\_LITE} (Section~\ref{sec:reservation-ALU-LITE})}
~
{\begin{lstlisting}
stage ID:
  new highmult = _ZX_2(highmult);
stage RR:
  new argument3 = PGR[registerO];
  new argument2 = PGR[registerP];
  for (I = 0; I < 4; I += 1) {
    if ((highmult == 0) || (highmult == 3)) {
      new product = _SX_32(argument2.32[I]) * _SX_32(argument3.32[I]);
    }
    if (highmult == 1) {
      product = _ZX_32(argument2.32[I]) * _ZX_32(argument3.32[I]);
    }
    if (highmult == 2) {
      product = _SX_32(argument2.32[I]) * _ZX_32(argument3.32[I]);
    }
    if (highmult) {
      result1.32[I] = product.32[1];
    } else {
      result1.32[I] = product.32[0];
    }
  };
stage E2:
  PGR[registerM] = result1;
\end{lstlisting}}
\newpage
\subsubsection[{\small MULXBHO: Multiply Extend Byte to Half Word Octuple}]{MULXBHO\hfill Multiply Extend Byte to Half Word Octuple}
\label{instr:MULXBHO}\index{mulxbho}
 
\paragraph{Encoding} Format \textsf{ALU\_HOMXWRR} (Section~\ref{sec:format-ALU-HOMXWRR}), \verb|alucode2 = 00|\\

\bitfields{ 1/P, 2/11, 1/1, 2/00, 2/widemult, 5/registerM, 1/oddlanes, 2/10, 3/111, 1/1, 5/registerO, 1/--, 5/registerP, 1/0 }


\paragraph{Syntax} \texttt{mulxbho}\emph{oddlanes}\emph{widemult} \emph{registerM} = \emph{registerP}, \emph{registerO}

\paragraph{Description} The \emph{registerO} and the \emph{registerP} are interpreted as sixteen bytes packed into 128 bits. The even or odd bytes of the \emph{registerO} and the \emph{registerP} are selected, depending on \emph{oddlanes}. These \emph{registerO} and \emph{registerP} bytes are sign-extended or zero-extended from 8 bits, depending on the \emph{widemult}. The \emph{registerO} bytes are multiplied by the \emph{registerP} bytes, producing eight 16-bit results. The eight resulting half words are packed and stored into the \emph{registerM}.


\paragraph{Modifier(s)}
\noindent\begin{itemize}
\item widemult (cf. widemult in section \ref{par:modifier-widemult} on page \pageref{par:modifier-widemult})
\item oddlanes (cf. oddlanes in section \ref{par:modifier-oddlanes} on page \pageref{par:modifier-oddlanes})
\end{itemize}

\paragraph{Execution} {Reservation table \textsf{ALU\_LITE} (Section~\ref{sec:reservation-ALU-LITE})}
~
{\begin{lstlisting}
stage ID:
  new oddlanes = _ZX_1(oddlanes);
  new widemult = _ZX_2(widemult);
stage RR:
  new argument3 = PGR[registerO];
  new argument2 = PGR[registerP];
  if (oddlanes) {
    argument2 = argument2 >> 8;
    argument3 = argument3 >> 8;
  }
  for (I = 0; I < 8; I += 1) {
    if (widemult == 0) {
      new product = _SX_8(argument2.16[I]) * _SX_8(argument3.16[I]);
    }
    if (widemult == 1) {
      product = _ZX_8(argument2.16[I]) * _ZX_8(argument3.16[I]);
    }
    if (widemult == 2) {
      product = _SX_8(argument2.16[I]) * _ZX_8(argument3.16[I]);
    }
    result1.16[I] = product
  };
stage E2:
  PGR[registerM] = result1;
\end{lstlisting}}
\newpage
\subsubsection[{\small MULXHWQ: Multiply Extend Half Word to Word Quadruple}]{MULXHWQ\hfill Multiply Extend Half Word to Word Quadruple}
\label{instr:MULXHWQ}\index{mulxhwq}
 
\paragraph{Encoding} Format \textsf{ALU\_WQMXWRR} (Section~\ref{sec:format-ALU-WQMXWRR}), \verb|alucode2 = 00|\\

\bitfields{ 1/P, 2/11, 1/1, 2/00, 2/widemult, 5/registerM, 1/oddlanes, 2/10, 3/111, 1/0, 5/registerO, 1/--, 5/registerP, 1/1 }


\paragraph{Syntax} \texttt{mulxhwq}\emph{oddlanes}\emph{widemult} \emph{registerM} = \emph{registerP}, \emph{registerO}

\paragraph{Description} The \emph{registerO} and the \emph{registerP} are interpreted as eight half words packed into 128 bits. The even or odd half words of the \emph{registerO} and the \emph{registerP} are selected, depending on \emph{oddlanes}. These \emph{registerO} and \emph{registerP} half words are sign-extended or zero-extended from 16 bits, depending on the \emph{widemult}. The \emph{registerO} half words are multiplied by the \emph{registerP} half words, producing four 32-bit results. The four resulting words are packed and stored into the \emph{registerM}.


\paragraph{Modifier(s)}
\noindent\begin{itemize}
\item widemult (cf. widemult in section \ref{par:modifier-widemult} on page \pageref{par:modifier-widemult})
\item oddlanes (cf. oddlanes in section \ref{par:modifier-oddlanes} on page \pageref{par:modifier-oddlanes})
\end{itemize}

\paragraph{Execution} {Reservation table \textsf{ALU\_LITE} (Section~\ref{sec:reservation-ALU-LITE})}
~
{\begin{lstlisting}
stage ID:
  new oddlanes = _ZX_1(oddlanes);
  new widemult = _ZX_2(widemult);
stage RR:
  new argument3 = PGR[registerO];
  new argument2 = PGR[registerP];
  if (oddlanes) {
    argument2 = argument2 >> 16;
    argument3 = argument3 >> 16;
  }
  for (I = 0; I < 4; I += 1) {
    if (widemult == 0) {
      new product = _SX_16(argument2.32[I]) * _SX_16(argument3.32[I]);
    }
    if (widemult == 1) {
      product = _ZX_16(argument2.32[I]) * _ZX_16(argument3.32[I]);
    }
    if (widemult == 2) {
      product = _SX_16(argument2.32[I]) * _ZX_16(argument3.32[I]);
    }
    result1.32[I] = product
  };
stage E2:
  PGR[registerM] = result1;
\end{lstlisting}}
\newpage
\subsubsection[{\small MULXWDP: Multiply Extend Word to Double Word Pair}]{MULXWDP\hfill Multiply Extend Word to Double Word Pair}
\label{instr:MULXWDP}\index{mulxwdp}
 
\paragraph{Encoding} Format \textsf{ALU\_DPMXWRR} (Section~\ref{sec:format-ALU-DPMXWRR}), \verb|alucode2 = 00|\\

\bitfields{ 1/P, 2/11, 1/1, 2/00, 2/widemult, 5/registerM, 1/oddlanes, 2/10, 3/111, 1/0, 5/registerO, 1/--, 5/registerP, 1/0 }


\paragraph{Syntax} \texttt{mulxwdp}\emph{oddlanes}\emph{widemult} \emph{registerM} = \emph{registerP}, \emph{registerO}

\paragraph{Description} The \emph{registerO} and the \emph{registerP} are interpreted as four words packed into 128 bits. The even or odd words of the \emph{registerO} and the \emph{registerP} are selected, depending on \emph{oddlanes}. These \emph{registerO} and \emph{registerP} words are sign-extended or zero-extended from 32 bits, depending on the \emph{widemult}. The \emph{registerO} words are multiplied by the \emph{registerP} words, producing two 64-bit results. The two resulting double words are packed and stored into the \emph{registerM}.


\paragraph{Modifier(s)}
\noindent\begin{itemize}
\item widemult (cf. widemult in section \ref{par:modifier-widemult} on page \pageref{par:modifier-widemult})
\item oddlanes (cf. oddlanes in section \ref{par:modifier-oddlanes} on page \pageref{par:modifier-oddlanes})
\end{itemize}

\paragraph{Execution} {Reservation table \textsf{ALU\_LITE} (Section~\ref{sec:reservation-ALU-LITE})}
~
{\begin{lstlisting}
stage ID:
  new oddlanes = _ZX_1(oddlanes);
  new widemult = _ZX_2(widemult);
stage RR:
  new argument3 = PGR[registerO];
  new argument2 = PGR[registerP];
  if (oddlanes) {
    argument2 = argument2 >> 32;
    argument3 = argument3 >> 32;
  }
  for (I = 0; I < 2; I += 1) {
    if (widemult == 0) {
      new product = _SX_32(argument2.64[I]) * _SX_32(argument3.64[I]);
    }
    if (widemult == 1) {
      product = _ZX_32(argument2.64[I]) * _ZX_32(argument3.64[I]);
    }
    if (widemult == 2) {
      product = _SX_32(argument2.64[I]) * _ZX_32(argument3.64[I]);
    }
    result1.64[I] = product
  };
stage E2:
  PGR[registerM] = result1;
\end{lstlisting}}
\newpage
\subsubsection[{\small NANDD: Bitwise Not And Double Word}]{NANDD\hfill Bitwise Not And Double Word}
\label{instr:NANDD}\index{nandd}
 
\paragraph{Encoding} Format \textsf{ALU\_DWRR0} (Section~\ref{sec:format-ALU-DWRR0}), \verb|exucode2 = 1001|\\

\bitfields{ 1/P, 2/11, 1/0, 4/1001, 6/registerW, 2/10, 3/000, 1/0, 6/registerY, 6/registerZ }


\paragraph{Syntax} \texttt{nandd} \emph{registerW} = \emph{registerZ}, \emph{registerY}

\paragraph{Description} Bitwise not of bitwise and of the \emph{registerZ} with the \emph{registerY}. The result is stored into the \emph{registerW}.


\paragraph{Execution} {Reservation table \textsf{ALU\_TINY} (Section~\ref{sec:reservation-ALU-TINY})}
~
{\begin{lstlisting}
stage RR:
  new argument3 = GPR[registerY];
  new argument2 = GPR[registerZ];
  new result1 = ~(argument2 & argument3);
stage E1:
  GPR[registerW] = result1;
\end{lstlisting}}
\newpage

\paragraph{Encoding} Format \textsf{ALU\_DWRR0.M} (Section~\ref{sec:format-ALU-DWRR0.M}), \verb|exucode2 = 1001|\\

\bitfields{ 1/P, 2/00, 2/-- --, 27/upper27, 1/1, 2/11, 1/0, 4/1001, 6/registerW, 2/10, 3/000, 1/0, 1/splat32, 5/lower5, 6/registerZ }


\paragraph{Syntax} \texttt{nandd} \emph{registerW} = \emph{registerZ}, \emph{upper27\_\discretionary{}{}{}lower5}\emph{splat32}

\paragraph{Description} Bitwise not of bitwise and of the \emph{registerZ} with the \emph{upper27\_\discretionary{}{}{}lower5}. The result is stored into the \emph{registerW}.


\paragraph{Modifier(s)}
\noindent\begin{itemize}
\item splat32 (cf. splat32 in section \ref{par:modifier-splat32} on page \pageref{par:modifier-splat32})
\end{itemize}

\paragraph{Execution} {Reservation table \textsf{ALU\_TINY.X} (Section~\ref{sec:reservation-ALU-TINY.X})}
~
{\begin{lstlisting}
stage ID:
  new splat32 = _ZX_1(splat32);
stage RR:
  new argument3 = splat32 ? (_WX_32(upper27_lower5) << 32) | _ZX_32(_WX_32(upper27_lower5)) : _SX_32(_WX_32(upper27_lower5));
  new argument2 = GPR[registerZ];
  new result1 = ~(argument2 & argument3);
stage E1:
  GPR[registerW] = result1;
\end{lstlisting}}
\newpage

\paragraph{Encoding} Format \textsf{ALU\_DWRI} (Section~\ref{sec:format-ALU-DWRI}), \verb|exucode2 = 1001|\\

\bitfields{ 1/P, 2/11, 1/0, 4/1001, 6/registerW, 2/00, 10/signed10, 6/registerZ }


\paragraph{Syntax} \texttt{nandd} \emph{registerW} = \emph{registerZ}, \emph{signed10}

\paragraph{Description} Bitwise not of bitwise and of the \emph{registerZ} with the \emph{signed10}. The result is stored into the \emph{registerW}.


\paragraph{Execution} {Reservation table \textsf{ALU\_TINY} (Section~\ref{sec:reservation-ALU-TINY})}
~
{\begin{lstlisting}
stage RR:
  new argument3 = _SX_10(signed10);
  new argument2 = GPR[registerZ];
  new result1 = ~(argument2 & argument3);
stage E1:
  GPR[registerW] = result1;
\end{lstlisting}}
\newpage

\paragraph{Encoding} Format \textsf{ALU\_DWRI.X} (Section~\ref{sec:format-ALU-DWRI.X}), \verb|exucode2 = 1001|\\

\bitfields{ 1/P, 2/00, 2/-- --, 27/upper27, 1/1, 2/11, 1/0, 4/1001, 6/registerW, 2/00, 10/lower10, 6/registerZ }


\paragraph{Syntax} \texttt{nandd} \emph{registerW} = \emph{registerZ}, \emph{upper27\_\discretionary{}{}{}lower10}

\paragraph{Description} Bitwise not of bitwise and of the \emph{registerZ} with the \emph{upper27\_\discretionary{}{}{}lower10}. The result is stored into the \emph{registerW}.


\paragraph{Execution} {Reservation table \textsf{ALU\_TINY.X} (Section~\ref{sec:reservation-ALU-TINY.X})}
~
{\begin{lstlisting}
stage RR:
  new argument3 = _SX_37(upper27_lower10);
  new argument2 = GPR[registerZ];
  new result1 = ~(argument2 & argument3);
stage E1:
  GPR[registerW] = result1;
\end{lstlisting}}
\newpage

\paragraph{Encoding} Format \textsf{ALU\_DWRI.Y} (Section~\ref{sec:format-ALU-DWRI.Y}), \verb|exucode2 = 1001|\\

\bitfields{ 1/P, 2/00, 2/-- --, 27/extend27, 1/1, 2/00, 2/-- --, 27/upper27, 1/1, 2/11, 1/0, 4/1001, 6/registerW, 2/00, 10/lower10, 6/registerZ }


\paragraph{Syntax} \texttt{nandd} \emph{registerW} = \emph{registerZ}, \emph{extend27\_\discretionary{}{}{}upper27\_\discretionary{}{}{}lower10}

\paragraph{Description} Bitwise not of bitwise and of the \emph{registerZ} with the \emph{extend27\_\discretionary{}{}{}upper27\_\discretionary{}{}{}lower10}. The result is stored into the \emph{registerW}.


\paragraph{Execution} {Reservation table \textsf{ALU\_TINY.Y} (Section~\ref{sec:reservation-ALU-TINY.Y})}
~
{\begin{lstlisting}
stage RR:
  new argument3 = _WX_64(extend27_upper27_lower10);
  new argument2 = GPR[registerZ];
  new result1 = ~(argument2 & argument3);
stage E1:
  GPR[registerW] = result1;
\end{lstlisting}}
\newpage
\subsubsection[{\small NANDQ: Bitwise Not And Quadruple Word}]{NANDQ\hfill Bitwise Not And Quadruple Word}
\label{instr:NANDQ}\index{nandq}
 
\paragraph{Encoding} Format \textsf{ALU\_QWRR0} (Section~\ref{sec:format-ALU-QWRR0}), \verb|exucode2 = 1001|\\

\bitfields{ 1/P, 2/11, 1/0, 4/1001, 5/registerM, 1/--, 2/10, 3/000, 1/1, 5/registerO, 1/--, 5/registerP, 1/-- }


\paragraph{Syntax} \texttt{nandq} \emph{registerM} = \emph{registerP}, \emph{registerO}

\paragraph{Description} Bitwise not of bitwise and of the \emph{registerP} with the \emph{registerO}. The result is stored into the \emph{registerM}.


\paragraph{Execution} {Reservation table \textsf{ALU\_LITE} (Section~\ref{sec:reservation-ALU-LITE})}
~
{\begin{lstlisting}
stage RR:
  new argument3 = PGR[registerO];
  new argument2 = PGR[registerP];
  new result1 = ~(argument2 & argument3);
stage E1:
  PGR[registerM] = result1;
\end{lstlisting}}
\newpage

\paragraph{Encoding} Format \textsf{ALU\_QWRR0.M} (Section~\ref{sec:format-ALU-QWRR0.M}), \verb|exucode2 = 1001|\\

\bitfields{ 1/P, 2/00, 2/-- --, 27/upper27, 1/1, 2/11, 1/0, 4/1001, 5/registerM, 1/--, 2/10, 3/000, 1/1, 1/splat32, 5/lower5, 5/registerP, 1/1 }


\paragraph{Syntax} \texttt{nandq} \emph{registerM} = \emph{registerP}, \emph{upper27\_\discretionary{}{}{}lower5}\emph{splat32}

\paragraph{Description} Bitwise not of bitwise and of the \emph{registerP} with the \emph{upper27\_\discretionary{}{}{}lower5}. The result is stored into the \emph{registerM}.


\paragraph{Modifier(s)}
\noindent\begin{itemize}
\item splat32 (cf. splat32 in section \ref{par:modifier-splat32} on page \pageref{par:modifier-splat32})
\end{itemize}

\paragraph{Execution} {Reservation table \textsf{ALU\_LITE.X} (Section~\ref{sec:reservation-ALU-LITE.X})}
~
{\begin{lstlisting}
stage ID:
  new splat32 = _ZX_1(splat32);
stage RR:
  new argument3 = splat32 ? (_WX_32(upper27_lower5) << 32) | _ZX_32(_WX_32(upper27_lower5)) : _SX_32(_WX_32(upper27_lower5));
  new argument2 = PGR[registerP];
  new result1 = ~(argument2 & argument3);
stage E1:
  PGR[registerM] = result1;
\end{lstlisting}}
\newpage
\subsubsection[{\small NANDW: Bitwise Not And Word}]{NANDW\hfill Bitwise Not And Word}
\label{instr:NANDW}\index{nandw}
 
\paragraph{Encoding} Format \textsf{ALU\_WWRR0} (Section~\ref{sec:format-ALU-WWRR0}), \verb|exucode2 = 1001|\\

\bitfields{ 1/P, 2/11, 1/0, 4/1001, 6/registerW, 2/11, 3/000, 1/signextw, 6/registerY, 6/registerZ }


\paragraph{Syntax} \texttt{nandw}\emph{signextw} \emph{registerW} = \emph{registerZ}, \emph{registerY}

\paragraph{Description} Bitwise not of bitwise and of the \emph{registerZ} with the \emph{registerY}. The result with the 32 upper bits cleared is stored into the \emph{registerW}.


\paragraph{Modifier(s)}
\noindent\begin{itemize}
\item signextw (cf. signextw in section \ref{par:modifier-signextw} on page \pageref{par:modifier-signextw})
\end{itemize}

\paragraph{Execution} {Reservation table \textsf{ALU\_TINY} (Section~\ref{sec:reservation-ALU-TINY})}
~
{\begin{lstlisting}
stage ID:
  new signextw = _ZX_1(signextw);
stage RR:
  new argument3 = GPR[registerY];
  new argument2 = GPR[registerZ];
  new result1 = ~(argument2 & argument3);
stage E1:
  GPR[registerW] = signextw ? _SX_32(result1) : _ZX_32(result1);
\end{lstlisting}}
\newpage

\paragraph{Encoding} Format \textsf{ALU\_WWRR0.W} (Section~\ref{sec:format-ALU-WWRR0.W}), \verb|exucode2 = 1001|\\

\bitfields{ 1/P, 2/00, 2/-- --, 27/upper27, 1/1, 2/11, 1/0, 4/1001, 6/registerW, 2/11, 3/000, 1/signextw, 1/0, 5/lower5, 6/registerZ }


\paragraph{Syntax} \texttt{nandw}\emph{signextw} \emph{registerW} = \emph{registerZ}, \emph{upper27\_\discretionary{}{}{}lower5}

\paragraph{Description} Bitwise not of bitwise and of the \emph{registerZ} with the \emph{upper27\_\discretionary{}{}{}lower5}. The result with the 32 upper bits cleared is stored into the \emph{registerW}.


\paragraph{Modifier(s)}
\noindent\begin{itemize}
\item signextw (cf. signextw in section \ref{par:modifier-signextw} on page \pageref{par:modifier-signextw})
\end{itemize}

\paragraph{Execution} {Reservation table \textsf{ALU\_TINY.X} (Section~\ref{sec:reservation-ALU-TINY.X})}
~
{\begin{lstlisting}
stage ID:
  new signextw = _ZX_1(signextw);
stage RR:
  new argument3 = _WX_32(upper27_lower5);
  new argument2 = GPR[registerZ];
  new result1 = ~(argument2 & argument3);
stage E1:
  GPR[registerW] = signextw ? _SX_32(result1) : _ZX_32(result1);
\end{lstlisting}}
\newpage
\subsubsection[{\small NEGBX: Negate Byte Hexadecuple}]{NEGBX\hfill Negate Byte Hexadecuple}
\label{instr:NEGBX}\index{negbx}
 
\paragraph{Encoding} Format \textsf{ALU\_BXBWR} (Section~\ref{sec:format-ALU-BXBWR}), \verb|alucode8 = 0000|\\

\bitfields{ 1/P, 2/11, 1/1, 4/1111, 5/registerM, 1/--, 2/10, 3/101, 1/1, 4/0000, 1/--, 1/--, 5/registerP, 1/1 }


\paragraph{Syntax} \texttt{negbx} \emph{registerM} = \emph{registerP}

\paragraph{Description} The negated values of the \emph{registerP} interpreted as a byte hexadecuple are stored into the \emph{registerM}.


\paragraph{Execution} {Reservation table \textsf{ALU\_LITE} (Section~\ref{sec:reservation-ALU-LITE})}
~
{\begin{lstlisting}
stage RR:
  new argument2 = PGR[registerP];
  for (I = 0; I < 16; I += 1) {
    result1.8[I] = -_SX_8(argument2.8[I])
  };
stage E1:
  PGR[registerM] = result1;
\end{lstlisting}}
\newpage
\subsubsection[{\small NEGD: Negate Double Word}]{NEGD\hfill Negate Double Word}
\label{instr:NEGD}\index{negd}
 
\paragraph{Encoding} Format \textsf{ALU\_DBWR} (Section~\ref{sec:format-ALU-DBWR}), \verb|alucode8 = 0000|\\

\bitfields{ 1/P, 2/11, 1/0, 4/1111, 6/registerW, 2/10, 3/101, 1/0, 4/0000, 1/0, 1/0, 6/registerZ }


\paragraph{Syntax} \texttt{negd} \emph{registerW} = \emph{registerZ}

\paragraph{Description} The 64-bit negated value of the \emph{registerZ} is stored into the \emph{registerW}.


\paragraph{Execution} {Reservation table \textsf{ALU\_TINY} (Section~\ref{sec:reservation-ALU-TINY})}
~
{\begin{lstlisting}
stage RR:
  new argument2 = GPR[registerZ];
  new result1 = -_SX_64(argument2);
stage E1:
  GPR[registerW] = result1;
\end{lstlisting}}
\newpage
\subsubsection[{\small NEGDP: Negate Double Word Pair}]{NEGDP\hfill Negate Double Word Pair}
\label{instr:NEGDP}\index{negdp}
 
\paragraph{Encoding} Format \textsf{ALU\_DPBWR} (Section~\ref{sec:format-ALU-DPBWR}), \verb|alucode8 = 0000|\\

\bitfields{ 1/P, 2/11, 1/1, 4/1111, 5/registerM, 1/--, 2/10, 3/101, 1/0, 4/0000, 1/--, 1/--, 5/registerP, 1/0 }


\paragraph{Syntax} \texttt{negdp} \emph{registerM} = \emph{registerP}

\paragraph{Description} The negated values of the \emph{registerP} interpreted as a double word pair are stored into the \emph{registerM}.


\paragraph{Execution} {Reservation table \textsf{ALU\_LITE} (Section~\ref{sec:reservation-ALU-LITE})}
~
{\begin{lstlisting}
stage RR:
  new argument2 = PGR[registerP];
  for (I = 0; I < 2; I += 1) {
    result1.64[I] = -_SX_64(argument2.64[I])
  };
stage E1:
  PGR[registerM] = result1;
\end{lstlisting}}
\newpage
\subsubsection[{\small NEGHO: Negate Half Word Octuple}]{NEGHO\hfill Negate Half Word Octuple}
\label{instr:NEGHO}\index{negho}
 
\paragraph{Encoding} Format \textsf{ALU\_HOBWR} (Section~\ref{sec:format-ALU-HOBWR}), \verb|alucode8 = 0000|\\

\bitfields{ 1/P, 2/11, 1/1, 4/1111, 5/registerM, 1/--, 2/10, 3/101, 1/1, 4/0000, 1/--, 1/--, 5/registerP, 1/0 }


\paragraph{Syntax} \texttt{negho} \emph{registerM} = \emph{registerP}

\paragraph{Description} The negated values of the \emph{registerP} interpreted as a half word octuple are stored into the \emph{registerM}.


\paragraph{Execution} {Reservation table \textsf{ALU\_LITE} (Section~\ref{sec:reservation-ALU-LITE})}
~
{\begin{lstlisting}
stage RR:
  new argument2 = PGR[registerP];
  for (I = 0; I < 8; I += 1) {
    result1.16[I] = -_SX_16(argument2.16[I])
  };
stage E1:
  PGR[registerM] = result1;
\end{lstlisting}}
\newpage
\subsubsection[{\small NEGQ: Negate Quadruple Word}]{NEGQ\hfill Negate Quadruple Word}
\label{instr:NEGQ}\index{negq}
 
\paragraph{Encoding} Format \textsf{ALU\_QBWR} (Section~\ref{sec:format-ALU-QBWR}), \verb|alucode8 = 0000|\\

\bitfields{ 1/P, 2/11, 1/0, 4/1111, 5/registerM, 1/--, 2/10, 3/101, 1/1, 4/0000, 1/0, 1/0, 5/registerP, 1/-- }


\paragraph{Syntax} \texttt{negq} \emph{registerM} = \emph{registerP}

\paragraph{Description} The 128-bit negated value of the \emph{registerP} is stored into the \emph{registerM}.


\paragraph{Execution} {Reservation table \textsf{ALU\_LITE} (Section~\ref{sec:reservation-ALU-LITE})}
~
{\begin{lstlisting}
stage RR:
  new argument2 = PGR[registerP];
  new result1 = -_SX_128(argument2);
stage E1:
  PGR[registerM] = result1;
\end{lstlisting}}
\newpage
\subsubsection[{\small NEGSBX: Negate Saturated Byte Hexadecuple}]{NEGSBX\hfill Negate Saturated Byte Hexadecuple}
\label{instr:NEGSBX}\index{negsbx}
 
\paragraph{Encoding} Format \textsf{ALU\_BXBWR} (Section~\ref{sec:format-ALU-BXBWR}), \verb|alucode8 = 0010|\\

\bitfields{ 1/P, 2/11, 1/1, 4/1111, 5/registerM, 1/--, 2/10, 3/101, 1/1, 4/0010, 1/--, 1/--, 5/registerP, 1/1 }


\paragraph{Syntax} \texttt{negsbx} \emph{registerM} = \emph{registerP}

\paragraph{Description} The negated saturated values of the \emph{registerP} interpreted as a byte hexadecuple are stored into the \emph{registerM}.


\paragraph{Execution} {Reservation table \textsf{ALU\_LITE} (Section~\ref{sec:reservation-ALU-LITE})}
~
{\begin{lstlisting}
stage RR:
  new argument2 = PGR[registerP];
  for (I = 0; I < 16; I += 1) {
    result1.8[I] = _SAT_8(-_SX_8(argument2.8[I]))
  };
stage E1:
  PGR[registerM] = result1;
\end{lstlisting}}
\newpage
\subsubsection[{\small NEGSD: Negate Saturated Double Word}]{NEGSD\hfill Negate Saturated Double Word}
\label{instr:NEGSD}\index{negsd}
 
\paragraph{Encoding} Format \textsf{ALU\_DBWR} (Section~\ref{sec:format-ALU-DBWR}), \verb|alucode8 = 0010|\\

\bitfields{ 1/P, 2/11, 1/0, 4/1111, 6/registerW, 2/10, 3/101, 1/0, 4/0010, 1/0, 1/0, 6/registerZ }


\paragraph{Syntax} \texttt{negsd} \emph{registerW} = \emph{registerZ}

\paragraph{Description} The 64-bit negated value of the \emph{registerZ} is saturated and stored into the \emph{registerW}.


\paragraph{Execution} {Reservation table \textsf{ALU\_TINY} (Section~\ref{sec:reservation-ALU-TINY})}
~
{\begin{lstlisting}
stage RR:
  new argument2 = GPR[registerZ];
  new result1 = _SAT_64(-_SX_64(argument2));
stage E1:
  GPR[registerW] = result1;
\end{lstlisting}}
\newpage
\subsubsection[{\small NEGSDP: Negate Saturated Double Word Pair}]{NEGSDP\hfill Negate Saturated Double Word Pair}
\label{instr:NEGSDP}\index{negsdp}
 
\paragraph{Encoding} Format \textsf{ALU\_DPBWR} (Section~\ref{sec:format-ALU-DPBWR}), \verb|alucode8 = 0010|\\

\bitfields{ 1/P, 2/11, 1/1, 4/1111, 5/registerM, 1/--, 2/10, 3/101, 1/0, 4/0010, 1/--, 1/--, 5/registerP, 1/0 }


\paragraph{Syntax} \texttt{negsdp} \emph{registerM} = \emph{registerP}

\paragraph{Description} The negated saturated values of the \emph{registerP} interpreted as a double word pair are stored into the \emph{registerM}.


\paragraph{Execution} {Reservation table \textsf{ALU\_LITE} (Section~\ref{sec:reservation-ALU-LITE})}
~
{\begin{lstlisting}
stage RR:
  new argument2 = PGR[registerP];
  for (I = 0; I < 2; I += 1) {
    result1.64[I] = _SAT_64(-_SX_64(argument2.64[I]))
  };
stage E1:
  PGR[registerM] = result1;
\end{lstlisting}}
\newpage
\subsubsection[{\small NEGSHO: Negate Saturated Half Word Octuple}]{NEGSHO\hfill Negate Saturated Half Word Octuple}
\label{instr:NEGSHO}\index{negsho}
 
\paragraph{Encoding} Format \textsf{ALU\_HOBWR} (Section~\ref{sec:format-ALU-HOBWR}), \verb|alucode8 = 0010|\\

\bitfields{ 1/P, 2/11, 1/1, 4/1111, 5/registerM, 1/--, 2/10, 3/101, 1/1, 4/0010, 1/--, 1/--, 5/registerP, 1/0 }


\paragraph{Syntax} \texttt{negsho} \emph{registerM} = \emph{registerP}

\paragraph{Description} The negated saturated values of the \emph{registerP} interpreted as a half word octuple are stored into the \emph{registerM}.


\paragraph{Execution} {Reservation table \textsf{ALU\_LITE} (Section~\ref{sec:reservation-ALU-LITE})}
~
{\begin{lstlisting}
stage RR:
  new argument2 = PGR[registerP];
  for (I = 0; I < 8; I += 1) {
    result1.16[I] = _SAT_16(-_SX_16(argument2.16[I]))
  };
stage E1:
  PGR[registerM] = result1;
\end{lstlisting}}
\newpage
\subsubsection[{\small NEGSW: Negate Saturated Word}]{NEGSW\hfill Negate Saturated Word}
\label{instr:NEGSW}\index{negsw}
 
\paragraph{Encoding} Format \textsf{ALU\_BWRW} (Section~\ref{sec:format-ALU-BWRW}), \verb|alucode8 = 0010|\\

\bitfields{ 1/P, 2/11, 1/0, 4/1111, 6/registerW, 2/11, 3/101, 1/signextw, 4/0010, 1/0, 1/0, 6/registerZ }


\paragraph{Syntax} \texttt{negsw} \emph{registerW} = \emph{registerZ}

\paragraph{Description} The 32-bit negated value of the \emph{registerZ} saturated and is stored into the \emph{registerW}.


\paragraph{Execution} {Reservation table \textsf{ALU\_TINY} (Section~\ref{sec:reservation-ALU-TINY})}
~
{\begin{lstlisting}
stage RR:
  new argument2 = GPR[registerZ];
  new result1 = _SAT_32(-_SX_32(argument2));
stage E1:
  GPR[registerW] = _ZX_32(result1);
\end{lstlisting}}
\newpage
\subsubsection[{\small NEGSWQ: Negate Saturated Word Quadruple}]{NEGSWQ\hfill Negate Saturated Word Quadruple}
\label{instr:NEGSWQ}\index{negswq}
 
\paragraph{Encoding} Format \textsf{ALU\_WQBWR} (Section~\ref{sec:format-ALU-WQBWR}), \verb|alucode8 = 0010|\\

\bitfields{ 1/P, 2/11, 1/1, 4/1111, 5/registerM, 1/--, 2/10, 3/101, 1/0, 4/0010, 1/--, 1/--, 5/registerP, 1/1 }


\paragraph{Syntax} \texttt{negswq} \emph{registerM} = \emph{registerP}

\paragraph{Description} The negated saturated values of the \emph{registerP} interpreted as a word quadruple are stored into the \emph{registerM}.


\paragraph{Execution} {Reservation table \textsf{ALU\_LITE} (Section~\ref{sec:reservation-ALU-LITE})}
~
{\begin{lstlisting}
stage RR:
  new argument2 = PGR[registerP];
  for (I = 0; I < 4; I += 1) {
    result1.32[I] = _SAT_32(-_SX_32(argument2.32[I]))
  };
stage E1:
  PGR[registerM] = result1;
\end{lstlisting}}
\newpage
\subsubsection[{\small NEGW: Negate Word}]{NEGW\hfill Negate Word}
\label{instr:NEGW}\index{negw}
 
\paragraph{Encoding} Format \textsf{ALU\_BWRW} (Section~\ref{sec:format-ALU-BWRW}), \verb|alucode8 = 0000|\\

\bitfields{ 1/P, 2/11, 1/0, 4/1111, 6/registerW, 2/11, 3/101, 1/signextw, 4/0000, 1/0, 1/0, 6/registerZ }


\paragraph{Syntax} \texttt{negw} \emph{registerW} = \emph{registerZ}

\paragraph{Description} The 32-bit negated value of the \emph{registerZ} is stored into the \emph{registerW}.


\paragraph{Execution} {Reservation table \textsf{ALU\_TINY} (Section~\ref{sec:reservation-ALU-TINY})}
~
{\begin{lstlisting}
stage RR:
  new argument2 = GPR[registerZ];
  new result1 = -_SX_32(argument2);
stage E1:
  GPR[registerW] = _ZX_32(result1);
\end{lstlisting}}
\newpage
\subsubsection[{\small NEGWQ: Negate Word Quadruple}]{NEGWQ\hfill Negate Word Quadruple}
\label{instr:NEGWQ}\index{negwq}
 
\paragraph{Encoding} Format \textsf{ALU\_WQBWR} (Section~\ref{sec:format-ALU-WQBWR}), \verb|alucode8 = 0000|\\

\bitfields{ 1/P, 2/11, 1/1, 4/1111, 5/registerM, 1/--, 2/10, 3/101, 1/0, 4/0000, 1/--, 1/--, 5/registerP, 1/1 }


\paragraph{Syntax} \texttt{negwq} \emph{registerM} = \emph{registerP}

\paragraph{Description} The negated values of the \emph{registerP} interpreted as a word quadruple are stored into the \emph{registerM}.


\paragraph{Execution} {Reservation table \textsf{ALU\_LITE} (Section~\ref{sec:reservation-ALU-LITE})}
~
{\begin{lstlisting}
stage RR:
  new argument2 = PGR[registerP];
  for (I = 0; I < 4; I += 1) {
    result1.32[I] = -_SX_32(argument2.32[I])
  };
stage E1:
  PGR[registerM] = result1;
\end{lstlisting}}
\newpage
\subsubsection[{\small NEORD: Bitwise Not Exclusive Or Double Word}]{NEORD\hfill Bitwise Not Exclusive Or Double Word}
\label{instr:NEORD}\index{neord}
 
\paragraph{Encoding} Format \textsf{ALU\_DWRR0} (Section~\ref{sec:format-ALU-DWRR0}), \verb|exucode2 = 1101|\\

\bitfields{ 1/P, 2/11, 1/0, 4/1101, 6/registerW, 2/10, 3/000, 1/0, 6/registerY, 6/registerZ }


\paragraph{Syntax} \texttt{neord} \emph{registerW} = \emph{registerZ}, \emph{registerY}

\paragraph{Description} Bitwise not of bitwise exclusive or of the \emph{registerZ} with the \emph{registerY}. The result is stored into the \emph{registerW}.


\paragraph{Execution} {Reservation table \textsf{ALU\_TINY} (Section~\ref{sec:reservation-ALU-TINY})}
~
{\begin{lstlisting}
stage RR:
  new argument3 = GPR[registerY];
  new argument2 = GPR[registerZ];
  new result1 = ~(argument2 ^ argument3);
stage E1:
  GPR[registerW] = result1;
\end{lstlisting}}
\newpage

\paragraph{Encoding} Format \textsf{ALU\_DWRR0.M} (Section~\ref{sec:format-ALU-DWRR0.M}), \verb|exucode2 = 1101|\\

\bitfields{ 1/P, 2/00, 2/-- --, 27/upper27, 1/1, 2/11, 1/0, 4/1101, 6/registerW, 2/10, 3/000, 1/0, 1/splat32, 5/lower5, 6/registerZ }


\paragraph{Syntax} \texttt{neord} \emph{registerW} = \emph{registerZ}, \emph{upper27\_\discretionary{}{}{}lower5}\emph{splat32}

\paragraph{Description} Bitwise not of bitwise exclusive or of the \emph{registerZ} with the \emph{upper27\_\discretionary{}{}{}lower5}. The result is stored into the \emph{registerW}.


\paragraph{Modifier(s)}
\noindent\begin{itemize}
\item splat32 (cf. splat32 in section \ref{par:modifier-splat32} on page \pageref{par:modifier-splat32})
\end{itemize}

\paragraph{Execution} {Reservation table \textsf{ALU\_TINY.X} (Section~\ref{sec:reservation-ALU-TINY.X})}
~
{\begin{lstlisting}
stage ID:
  new splat32 = _ZX_1(splat32);
stage RR:
  new argument3 = splat32 ? (_WX_32(upper27_lower5) << 32) | _ZX_32(_WX_32(upper27_lower5)) : _SX_32(_WX_32(upper27_lower5));
  new argument2 = GPR[registerZ];
  new result1 = ~(argument2 ^ argument3);
stage E1:
  GPR[registerW] = result1;
\end{lstlisting}}
\newpage

\paragraph{Encoding} Format \textsf{ALU\_DWRI} (Section~\ref{sec:format-ALU-DWRI}), \verb|exucode2 = 1101|\\

\bitfields{ 1/P, 2/11, 1/0, 4/1101, 6/registerW, 2/00, 10/signed10, 6/registerZ }


\paragraph{Syntax} \texttt{neord} \emph{registerW} = \emph{registerZ}, \emph{signed10}

\paragraph{Description} Bitwise not of bitwise exclusive or of the \emph{registerZ} with the \emph{signed10}. The result is stored into the \emph{registerW}.


\paragraph{Execution} {Reservation table \textsf{ALU\_TINY} (Section~\ref{sec:reservation-ALU-TINY})}
~
{\begin{lstlisting}
stage RR:
  new argument3 = _SX_10(signed10);
  new argument2 = GPR[registerZ];
  new result1 = ~(argument2 ^ argument3);
stage E1:
  GPR[registerW] = result1;
\end{lstlisting}}
\newpage

\paragraph{Encoding} Format \textsf{ALU\_DWRI.X} (Section~\ref{sec:format-ALU-DWRI.X}), \verb|exucode2 = 1101|\\

\bitfields{ 1/P, 2/00, 2/-- --, 27/upper27, 1/1, 2/11, 1/0, 4/1101, 6/registerW, 2/00, 10/lower10, 6/registerZ }


\paragraph{Syntax} \texttt{neord} \emph{registerW} = \emph{registerZ}, \emph{upper27\_\discretionary{}{}{}lower10}

\paragraph{Description} Bitwise not of bitwise exclusive or of the \emph{registerZ} with the \emph{upper27\_\discretionary{}{}{}lower10}. The result is stored into the \emph{registerW}.


\paragraph{Execution} {Reservation table \textsf{ALU\_TINY.X} (Section~\ref{sec:reservation-ALU-TINY.X})}
~
{\begin{lstlisting}
stage RR:
  new argument3 = _SX_37(upper27_lower10);
  new argument2 = GPR[registerZ];
  new result1 = ~(argument2 ^ argument3);
stage E1:
  GPR[registerW] = result1;
\end{lstlisting}}
\newpage

\paragraph{Encoding} Format \textsf{ALU\_DWRI.Y} (Section~\ref{sec:format-ALU-DWRI.Y}), \verb|exucode2 = 1101|\\

\bitfields{ 1/P, 2/00, 2/-- --, 27/extend27, 1/1, 2/00, 2/-- --, 27/upper27, 1/1, 2/11, 1/0, 4/1101, 6/registerW, 2/00, 10/lower10, 6/registerZ }


\paragraph{Syntax} \texttt{neord} \emph{registerW} = \emph{registerZ}, \emph{extend27\_\discretionary{}{}{}upper27\_\discretionary{}{}{}lower10}

\paragraph{Description} Bitwise not of bitwise exclusive or of the \emph{registerZ} with the \emph{extend27\_\discretionary{}{}{}upper27\_\discretionary{}{}{}lower10}. The result is stored into the \emph{registerW}.


\paragraph{Execution} {Reservation table \textsf{ALU\_TINY.Y} (Section~\ref{sec:reservation-ALU-TINY.Y})}
~
{\begin{lstlisting}
stage RR:
  new argument3 = _WX_64(extend27_upper27_lower10);
  new argument2 = GPR[registerZ];
  new result1 = ~(argument2 ^ argument3);
stage E1:
  GPR[registerW] = result1;
\end{lstlisting}}
\newpage
\subsubsection[{\small NEORQ: Bitwise Not Exclusive Or Quadruple Word}]{NEORQ\hfill Bitwise Not Exclusive Or Quadruple Word}
\label{instr:NEORQ}\index{neorq}
 
\paragraph{Encoding} Format \textsf{ALU\_QWRR0} (Section~\ref{sec:format-ALU-QWRR0}), \verb|exucode2 = 1101|\\

\bitfields{ 1/P, 2/11, 1/0, 4/1101, 5/registerM, 1/--, 2/10, 3/000, 1/1, 5/registerO, 1/--, 5/registerP, 1/-- }


\paragraph{Syntax} \texttt{neorq} \emph{registerM} = \emph{registerP}, \emph{registerO}

\paragraph{Description} Bitwise not of bitwise exclusive or of the \emph{registerP} with the \emph{registerO}. The result is stored into the \emph{registerM}.


\paragraph{Execution} {Reservation table \textsf{ALU\_LITE} (Section~\ref{sec:reservation-ALU-LITE})}
~
{\begin{lstlisting}
stage RR:
  new argument3 = PGR[registerO];
  new argument2 = PGR[registerP];
  new result1 = ~(argument2 ^ argument3);
stage E1:
  PGR[registerM] = result1;
\end{lstlisting}}
\newpage

\paragraph{Encoding} Format \textsf{ALU\_QWRR0.M} (Section~\ref{sec:format-ALU-QWRR0.M}), \verb|exucode2 = 1101|\\

\bitfields{ 1/P, 2/00, 2/-- --, 27/upper27, 1/1, 2/11, 1/0, 4/1101, 5/registerM, 1/--, 2/10, 3/000, 1/1, 1/splat32, 5/lower5, 5/registerP, 1/1 }


\paragraph{Syntax} \texttt{neorq} \emph{registerM} = \emph{registerP}, \emph{upper27\_\discretionary{}{}{}lower5}\emph{splat32}

\paragraph{Description} Bitwise not of bitwise exclusive or of the \emph{registerP} with the \emph{upper27\_\discretionary{}{}{}lower5}. The result is stored into the \emph{registerM}.


\paragraph{Modifier(s)}
\noindent\begin{itemize}
\item splat32 (cf. splat32 in section \ref{par:modifier-splat32} on page \pageref{par:modifier-splat32})
\end{itemize}

\paragraph{Execution} {Reservation table \textsf{ALU\_LITE.X} (Section~\ref{sec:reservation-ALU-LITE.X})}
~
{\begin{lstlisting}
stage ID:
  new splat32 = _ZX_1(splat32);
stage RR:
  new argument3 = splat32 ? (_WX_32(upper27_lower5) << 32) | _ZX_32(_WX_32(upper27_lower5)) : _SX_32(_WX_32(upper27_lower5));
  new argument2 = PGR[registerP];
  new result1 = ~(argument2 ^ argument3);
stage E1:
  PGR[registerM] = result1;
\end{lstlisting}}
\newpage
\subsubsection[{\small NEORW: Bitwise Not Exclusive Or Word}]{NEORW\hfill Bitwise Not Exclusive Or Word}
\label{instr:NEORW}\index{neorw}
 
\paragraph{Encoding} Format \textsf{ALU\_WWRR0} (Section~\ref{sec:format-ALU-WWRR0}), \verb|exucode2 = 1101|\\

\bitfields{ 1/P, 2/11, 1/0, 4/1101, 6/registerW, 2/11, 3/000, 1/signextw, 6/registerY, 6/registerZ }


\paragraph{Syntax} \texttt{neorw}\emph{signextw} \emph{registerW} = \emph{registerZ}, \emph{registerY}

\paragraph{Description} Bitwise not of bitwise exclusive or of the \emph{registerZ} with the \emph{registerY}. The result with the 32 upper bits cleared is stored into the \emph{registerW}.


\paragraph{Modifier(s)}
\noindent\begin{itemize}
\item signextw (cf. signextw in section \ref{par:modifier-signextw} on page \pageref{par:modifier-signextw})
\end{itemize}

\paragraph{Execution} {Reservation table \textsf{ALU\_TINY} (Section~\ref{sec:reservation-ALU-TINY})}
~
{\begin{lstlisting}
stage ID:
  new signextw = _ZX_1(signextw);
stage RR:
  new argument3 = GPR[registerY];
  new argument2 = GPR[registerZ];
  new result1 = ~(argument2 ^ argument3);
stage E1:
  GPR[registerW] = signextw ? _SX_32(result1) : _ZX_32(result1);
\end{lstlisting}}
\newpage

\paragraph{Encoding} Format \textsf{ALU\_WWRR0.W} (Section~\ref{sec:format-ALU-WWRR0.W}), \verb|exucode2 = 1101|\\

\bitfields{ 1/P, 2/00, 2/-- --, 27/upper27, 1/1, 2/11, 1/0, 4/1101, 6/registerW, 2/11, 3/000, 1/signextw, 1/0, 5/lower5, 6/registerZ }


\paragraph{Syntax} \texttt{neorw}\emph{signextw} \emph{registerW} = \emph{registerZ}, \emph{upper27\_\discretionary{}{}{}lower5}

\paragraph{Description} Bitwise not of bitwise exclusive or of the \emph{registerZ} with the \emph{upper27\_\discretionary{}{}{}lower5}. The result with the 32 upper bits cleared is stored into the \emph{registerW}.


\paragraph{Modifier(s)}
\noindent\begin{itemize}
\item signextw (cf. signextw in section \ref{par:modifier-signextw} on page \pageref{par:modifier-signextw})
\end{itemize}

\paragraph{Execution} {Reservation table \textsf{ALU\_TINY.X} (Section~\ref{sec:reservation-ALU-TINY.X})}
~
{\begin{lstlisting}
stage ID:
  new signextw = _ZX_1(signextw);
stage RR:
  new argument3 = _WX_32(upper27_lower5);
  new argument2 = GPR[registerZ];
  new result1 = ~(argument2 ^ argument3);
stage E1:
  GPR[registerW] = signextw ? _SX_32(result1) : _ZX_32(result1);
\end{lstlisting}}
\newpage
\subsubsection[{\small NIORD: Bitwise Not Inclusive Or Double Word}]{NIORD\hfill Bitwise Not Inclusive Or Double Word}
\label{instr:NIORD}\index{niord}
 
\paragraph{Encoding} Format \textsf{ALU\_DWRR0} (Section~\ref{sec:format-ALU-DWRR0}), \verb|exucode2 = 1011|\\

\bitfields{ 1/P, 2/11, 1/0, 4/1011, 6/registerW, 2/10, 3/000, 1/0, 6/registerY, 6/registerZ }


\paragraph{Syntax} \texttt{niord} \emph{registerW} = \emph{registerZ}, \emph{registerY}

\paragraph{Description} Bitwise not of bitwise inclusive or of the \emph{registerZ} with the \emph{registerY}. The result is stored into the \emph{registerW}.


\paragraph{Execution} {Reservation table \textsf{ALU\_TINY} (Section~\ref{sec:reservation-ALU-TINY})}
~
{\begin{lstlisting}
stage RR:
  new argument3 = GPR[registerY];
  new argument2 = GPR[registerZ];
  new result1 = ~(argument2 | argument3);
stage E1:
  GPR[registerW] = result1;
\end{lstlisting}}
\newpage

\paragraph{Encoding} Format \textsf{ALU\_DWRR0.M} (Section~\ref{sec:format-ALU-DWRR0.M}), \verb|exucode2 = 1011|\\

\bitfields{ 1/P, 2/00, 2/-- --, 27/upper27, 1/1, 2/11, 1/0, 4/1011, 6/registerW, 2/10, 3/000, 1/0, 1/splat32, 5/lower5, 6/registerZ }


\paragraph{Syntax} \texttt{niord} \emph{registerW} = \emph{registerZ}, \emph{upper27\_\discretionary{}{}{}lower5}\emph{splat32}

\paragraph{Description} Bitwise not of bitwise inclusive or of the \emph{registerZ} with the \emph{upper27\_\discretionary{}{}{}lower5}. The result is stored into the \emph{registerW}.


\paragraph{Modifier(s)}
\noindent\begin{itemize}
\item splat32 (cf. splat32 in section \ref{par:modifier-splat32} on page \pageref{par:modifier-splat32})
\end{itemize}

\paragraph{Execution} {Reservation table \textsf{ALU\_TINY.X} (Section~\ref{sec:reservation-ALU-TINY.X})}
~
{\begin{lstlisting}
stage ID:
  new splat32 = _ZX_1(splat32);
stage RR:
  new argument3 = splat32 ? (_WX_32(upper27_lower5) << 32) | _ZX_32(_WX_32(upper27_lower5)) : _SX_32(_WX_32(upper27_lower5));
  new argument2 = GPR[registerZ];
  new result1 = ~(argument2 | argument3);
stage E1:
  GPR[registerW] = result1;
\end{lstlisting}}
\newpage

\paragraph{Encoding} Format \textsf{ALU\_DWRI} (Section~\ref{sec:format-ALU-DWRI}), \verb|exucode2 = 1011|\\

\bitfields{ 1/P, 2/11, 1/0, 4/1011, 6/registerW, 2/00, 10/signed10, 6/registerZ }


\paragraph{Syntax} \texttt{niord} \emph{registerW} = \emph{registerZ}, \emph{signed10}

\paragraph{Description} Bitwise not of bitwise inclusive or of the \emph{registerZ} with the \emph{signed10}. The result is stored into the \emph{registerW}.


\paragraph{Execution} {Reservation table \textsf{ALU\_TINY} (Section~\ref{sec:reservation-ALU-TINY})}
~
{\begin{lstlisting}
stage RR:
  new argument3 = _SX_10(signed10);
  new argument2 = GPR[registerZ];
  new result1 = ~(argument2 | argument3);
stage E1:
  GPR[registerW] = result1;
\end{lstlisting}}
\newpage

\paragraph{Encoding} Format \textsf{ALU\_DWRI.X} (Section~\ref{sec:format-ALU-DWRI.X}), \verb|exucode2 = 1011|\\

\bitfields{ 1/P, 2/00, 2/-- --, 27/upper27, 1/1, 2/11, 1/0, 4/1011, 6/registerW, 2/00, 10/lower10, 6/registerZ }


\paragraph{Syntax} \texttt{niord} \emph{registerW} = \emph{registerZ}, \emph{upper27\_\discretionary{}{}{}lower10}

\paragraph{Description} Bitwise not of bitwise inclusive or of the \emph{registerZ} with the \emph{upper27\_\discretionary{}{}{}lower10}. The result is stored into the \emph{registerW}.


\paragraph{Execution} {Reservation table \textsf{ALU\_TINY.X} (Section~\ref{sec:reservation-ALU-TINY.X})}
~
{\begin{lstlisting}
stage RR:
  new argument3 = _SX_37(upper27_lower10);
  new argument2 = GPR[registerZ];
  new result1 = ~(argument2 | argument3);
stage E1:
  GPR[registerW] = result1;
\end{lstlisting}}
\newpage

\paragraph{Encoding} Format \textsf{ALU\_DWRI.Y} (Section~\ref{sec:format-ALU-DWRI.Y}), \verb|exucode2 = 1011|\\

\bitfields{ 1/P, 2/00, 2/-- --, 27/extend27, 1/1, 2/00, 2/-- --, 27/upper27, 1/1, 2/11, 1/0, 4/1011, 6/registerW, 2/00, 10/lower10, 6/registerZ }


\paragraph{Syntax} \texttt{niord} \emph{registerW} = \emph{registerZ}, \emph{extend27\_\discretionary{}{}{}upper27\_\discretionary{}{}{}lower10}

\paragraph{Description} Bitwise not of bitwise inclusive or of the \emph{registerZ} with the \emph{extend27\_\discretionary{}{}{}upper27\_\discretionary{}{}{}lower10}. The result is stored into the \emph{registerW}.


\paragraph{Execution} {Reservation table \textsf{ALU\_TINY.Y} (Section~\ref{sec:reservation-ALU-TINY.Y})}
~
{\begin{lstlisting}
stage RR:
  new argument3 = _WX_64(extend27_upper27_lower10);
  new argument2 = GPR[registerZ];
  new result1 = ~(argument2 | argument3);
stage E1:
  GPR[registerW] = result1;
\end{lstlisting}}
\newpage
\subsubsection[{\small NIORQ: Bitwise Not Inclusive Or Quadruple Word}]{NIORQ\hfill Bitwise Not Inclusive Or Quadruple Word}
\label{instr:NIORQ}\index{niorq}
 
\paragraph{Encoding} Format \textsf{ALU\_QWRR0} (Section~\ref{sec:format-ALU-QWRR0}), \verb|exucode2 = 1011|\\

\bitfields{ 1/P, 2/11, 1/0, 4/1011, 5/registerM, 1/--, 2/10, 3/000, 1/1, 5/registerO, 1/--, 5/registerP, 1/-- }


\paragraph{Syntax} \texttt{niorq} \emph{registerM} = \emph{registerP}, \emph{registerO}

\paragraph{Description} Bitwise not of bitwise inclusive or of the \emph{registerP} with the \emph{registerO}. The result is stored into the \emph{registerM}.


\paragraph{Execution} {Reservation table \textsf{ALU\_LITE} (Section~\ref{sec:reservation-ALU-LITE})}
~
{\begin{lstlisting}
stage RR:
  new argument3 = PGR[registerO];
  new argument2 = PGR[registerP];
  new result1 = ~(argument2 | argument3);
stage E1:
  PGR[registerM] = result1;
\end{lstlisting}}
\newpage

\paragraph{Encoding} Format \textsf{ALU\_QWRR0.M} (Section~\ref{sec:format-ALU-QWRR0.M}), \verb|exucode2 = 1011|\\

\bitfields{ 1/P, 2/00, 2/-- --, 27/upper27, 1/1, 2/11, 1/0, 4/1011, 5/registerM, 1/--, 2/10, 3/000, 1/1, 1/splat32, 5/lower5, 5/registerP, 1/1 }


\paragraph{Syntax} \texttt{niorq} \emph{registerM} = \emph{registerP}, \emph{upper27\_\discretionary{}{}{}lower5}\emph{splat32}

\paragraph{Description} Bitwise not of bitwise inclusive or of the \emph{registerP} with the \emph{upper27\_\discretionary{}{}{}lower5}. The result is stored into the \emph{registerM}.


\paragraph{Modifier(s)}
\noindent\begin{itemize}
\item splat32 (cf. splat32 in section \ref{par:modifier-splat32} on page \pageref{par:modifier-splat32})
\end{itemize}

\paragraph{Execution} {Reservation table \textsf{ALU\_LITE.X} (Section~\ref{sec:reservation-ALU-LITE.X})}
~
{\begin{lstlisting}
stage ID:
  new splat32 = _ZX_1(splat32);
stage RR:
  new argument3 = splat32 ? (_WX_32(upper27_lower5) << 32) | _ZX_32(_WX_32(upper27_lower5)) : _SX_32(_WX_32(upper27_lower5));
  new argument2 = PGR[registerP];
  new result1 = ~(argument2 | argument3);
stage E1:
  PGR[registerM] = result1;
\end{lstlisting}}
\newpage
\subsubsection[{\small NIORW: Bitwise Not Inclusive Or Word}]{NIORW\hfill Bitwise Not Inclusive Or Word}
\label{instr:NIORW}\index{niorw}
 
\paragraph{Encoding} Format \textsf{ALU\_WWRR0} (Section~\ref{sec:format-ALU-WWRR0}), \verb|exucode2 = 1011|\\

\bitfields{ 1/P, 2/11, 1/0, 4/1011, 6/registerW, 2/11, 3/000, 1/signextw, 6/registerY, 6/registerZ }


\paragraph{Syntax} \texttt{niorw}\emph{signextw} \emph{registerW} = \emph{registerZ}, \emph{registerY}

\paragraph{Description} Bitwise not of bitwise inclusive or of the \emph{registerZ} with the \emph{registerY}. The result with the 32 upper bits cleared is stored into the \emph{registerW}.


\paragraph{Modifier(s)}
\noindent\begin{itemize}
\item signextw (cf. signextw in section \ref{par:modifier-signextw} on page \pageref{par:modifier-signextw})
\end{itemize}

\paragraph{Execution} {Reservation table \textsf{ALU\_TINY} (Section~\ref{sec:reservation-ALU-TINY})}
~
{\begin{lstlisting}
stage ID:
  new signextw = _ZX_1(signextw);
stage RR:
  new argument3 = GPR[registerY];
  new argument2 = GPR[registerZ];
  new result1 = ~(argument2 | argument3);
stage E1:
  GPR[registerW] = signextw ? _SX_32(result1) : _ZX_32(result1);
\end{lstlisting}}
\newpage

\paragraph{Encoding} Format \textsf{ALU\_WWRR0.W} (Section~\ref{sec:format-ALU-WWRR0.W}), \verb|exucode2 = 1011|\\

\bitfields{ 1/P, 2/00, 2/-- --, 27/upper27, 1/1, 2/11, 1/0, 4/1011, 6/registerW, 2/11, 3/000, 1/signextw, 1/0, 5/lower5, 6/registerZ }


\paragraph{Syntax} \texttt{niorw}\emph{signextw} \emph{registerW} = \emph{registerZ}, \emph{upper27\_\discretionary{}{}{}lower5}

\paragraph{Description} Bitwise not of bitwise inclusive or of the \emph{registerZ} with the \emph{upper27\_\discretionary{}{}{}lower5}. The result with the 32 upper bits cleared is stored into the \emph{registerW}.


\paragraph{Modifier(s)}
\noindent\begin{itemize}
\item signextw (cf. signextw in section \ref{par:modifier-signextw} on page \pageref{par:modifier-signextw})
\end{itemize}

\paragraph{Execution} {Reservation table \textsf{ALU\_TINY.X} (Section~\ref{sec:reservation-ALU-TINY.X})}
~
{\begin{lstlisting}
stage ID:
  new signextw = _ZX_1(signextw);
stage RR:
  new argument3 = _WX_32(upper27_lower5);
  new argument2 = GPR[registerZ];
  new result1 = ~(argument2 | argument3);
stage E1:
  GPR[registerW] = signextw ? _SX_32(result1) : _ZX_32(result1);
\end{lstlisting}}
\newpage
\subsubsection[{\small NOP: No Operation}]{NOP\hfill No Operation}
\label{instr:NOP}\index{nop}
 
\paragraph{Encoding} Format \textsf{ALU\_NOP} (Section~\ref{sec:format-ALU-NOP}), \verb|exucode2 = 1111|\\

\bitfields{ 1/P, 2/11, 1/1, 4/1111, 5/-- -- -- -- --, 1/1, 2/11, 4/1111, 5/11111, 1/1, 5/-- -- -- -- --, 1/1 }


\paragraph{Syntax} \texttt{nop}

\paragraph{Description} No effects except extending the current bundle by one syllable.


\paragraph{Execution} {Reservation table \textsf{ALU\_TINY} (Section~\ref{sec:reservation-ALU-TINY})}
~
{\begin{lstlisting}
stage RR:
  ;
\end{lstlisting}}
\newpage
\subsubsection[{\small NOTD: Negate Double Word}]{NOTD\hfill Negate Double Word}
\label{instr:NOTD}\index{notd}
 
\paragraph{Encoding} Format \textsf{ALU\_DBWR} (Section~\ref{sec:format-ALU-DBWR}), \verb|alucode8 = 1000|\\

\bitfields{ 1/P, 2/11, 1/0, 4/1111, 6/registerW, 2/10, 3/101, 1/0, 4/1000, 1/0, 1/0, 6/registerZ }


\paragraph{Syntax} \texttt{notd} \emph{registerW} = \emph{registerZ}

\paragraph{Description} The bitwise not value of the \emph{registerZ} is stored into the \emph{registerW}.


\paragraph{Execution} {Reservation table \textsf{ALU\_TINY} (Section~\ref{sec:reservation-ALU-TINY})}
~
{\begin{lstlisting}
stage RR:
  new argument2 = GPR[registerZ];
  new result1 = ~argument2;
stage E1:
  GPR[registerW] = result1;
\end{lstlisting}}
\newpage
\subsubsection[{\small NOTQ: Negate Quadruple Word}]{NOTQ\hfill Negate Quadruple Word}
\label{instr:NOTQ}\index{notq}
 
\paragraph{Encoding} Format \textsf{ALU\_QBWR} (Section~\ref{sec:format-ALU-QBWR}), \verb|alucode8 = 1000|\\

\bitfields{ 1/P, 2/11, 1/0, 4/1111, 5/registerM, 1/--, 2/10, 3/101, 1/1, 4/1000, 1/0, 1/0, 5/registerP, 1/-- }


\paragraph{Syntax} \texttt{notq} \emph{registerM} = \emph{registerP}

\paragraph{Description} The bitwise not value of the \emph{registerP} is stored into the \emph{registerM}.


\paragraph{Execution} {Reservation table \textsf{ALU\_LITE} (Section~\ref{sec:reservation-ALU-LITE})}
~
{\begin{lstlisting}
stage RR:
  new argument2 = PGR[registerP];
  new result1 = ~argument2;
stage E1:
  PGR[registerM] = result1;
\end{lstlisting}}
\newpage
\subsubsection[{\small NOTW: Negate Word}]{NOTW\hfill Negate Word}
\label{instr:NOTW}\index{notw}
 
\paragraph{Encoding} Format \textsf{ALU\_BWRW} (Section~\ref{sec:format-ALU-BWRW}), \verb|alucode8 = 1000|\\

\bitfields{ 1/P, 2/11, 1/0, 4/1111, 6/registerW, 2/11, 3/101, 1/signextw, 4/1000, 1/0, 1/0, 6/registerZ }


\paragraph{Syntax} \texttt{notw} \emph{registerW} = \emph{registerZ}

\paragraph{Description} The bitwise not value of the \emph{registerZ} is stored into the \emph{registerW}.


\paragraph{Execution} {Reservation table \textsf{ALU\_TINY} (Section~\ref{sec:reservation-ALU-TINY})}
~
{\begin{lstlisting}
stage RR:
  new argument2 = GPR[registerZ];
  new result1 = ~argument2;
stage E1:
  GPR[registerW] = _ZX_32(result1);
\end{lstlisting}}
\newpage
\subsubsection[{\small PCREL: PC-Relative}]{PCREL\hfill PC-Relative}
\label{instr:PCREL}\index{pcrel}
 
\paragraph{Encoding} Format \textsf{ALU\_DWI} (Section~\ref{sec:format-ALU-DWI}), \verb|exucode2 = 0001|\\

\bitfields{ 1/P, 2/11, 1/0, 4/0001, 6/registerW, 2/00, 16/signed16 }


\paragraph{Syntax} \texttt{pcrel} \emph{registerW} = \emph{signed16}

\paragraph{Description} The \emph{signed16} is is added to the bundle \textbf{PC} and the result is stored into the \emph{registerW}.


\paragraph{Execution} {Reservation table \textsf{ALU\_TINY} (Section~\ref{sec:reservation-ALU-TINY})}
~
{\begin{lstlisting}
stage RR:
  new argument2 = _SX_16(signed16);
  new result1 = PC + _SX_64(argument2);
stage E1:
  GPR[registerW] = result1;
\end{lstlisting}}
\newpage

\paragraph{Encoding} Format \textsf{ALU\_DWI.X} (Section~\ref{sec:format-ALU-DWI.X}), \verb|exucode2 = 0001|\\

\bitfields{ 1/P, 2/00, 2/-- --, 27/upper27, 1/1, 2/11, 1/0, 4/0001, 6/registerW, 2/00, 10/lower10, 6/extend6 }


\paragraph{Syntax} \texttt{pcrel} \emph{registerW} = \emph{extend6\_\discretionary{}{}{}upper27\_\discretionary{}{}{}lower10}

\paragraph{Description} The \emph{extend6\_\discretionary{}{}{}upper27\_\discretionary{}{}{}lower10} is is added to the bundle \textbf{PC} and the result is stored into the \emph{registerW}.


\paragraph{Execution} {Reservation table \textsf{ALU\_TINY.X} (Section~\ref{sec:reservation-ALU-TINY.X})}
~
{\begin{lstlisting}
stage RR:
  new argument2 = _SX_43(extend6_upper27_lower10);
  new result1 = PC + _SX_64(argument2);
stage E1:
  GPR[registerW] = result1;
\end{lstlisting}}
\newpage

\paragraph{Encoding} Format \textsf{ALU\_DWI.Y} (Section~\ref{sec:format-ALU-DWI.Y}), \verb|exucode2 = 0001|\\

\bitfields{ 1/P, 2/00, 2/-- --, 27/extend27, 1/1, 2/00, 2/-- --, 27/upper27, 1/1, 2/11, 1/0, 4/0001, 6/registerW, 2/00, 10/lower10, 6/-- -- -- -- -- -- }


\paragraph{Syntax} \texttt{pcrel} \emph{registerW} = \emph{extend27\_\discretionary{}{}{}upper27\_\discretionary{}{}{}lower10}

\paragraph{Description} The \emph{extend27\_\discretionary{}{}{}upper27\_\discretionary{}{}{}lower10} is is added to the bundle \textbf{PC} and the result is stored into the \emph{registerW}.


\paragraph{Execution} {Reservation table \textsf{ALU\_TINY.Y} (Section~\ref{sec:reservation-ALU-TINY.Y})}
~
{\begin{lstlisting}
stage RR:
  new argument2 = _WX_64(extend27_upper27_lower10);
  new result1 = PC + _SX_64(argument2);
stage E1:
  GPR[registerW] = result1;
\end{lstlisting}}
\newpage
\subsubsection[{\small ROLD: Rotate Left Double Word}]{ROLD\hfill Rotate Left Double Word}
\label{instr:ROLD}\index{rold}
 
\paragraph{Encoding} Format \textsf{ALU\_DSWRR} (Section~\ref{sec:format-ALU-DSWRR}), \verb|exucode1 = 110|\\

\bitfields{ 1/P, 2/11, 1/0, 1/0, 3/110, 6/registerW, 2/10, 3/100, 1/0, 6/registerY, 6/registerZ }


\paragraph{Syntax} \texttt{rold} \emph{registerW} = \emph{registerZ}, \emph{registerY}

\paragraph{Description} The \emph{registerZ} is rotated left by the \emph{registerY} modulo 64. The result is stored into the \emph{registerW}.


\paragraph{Execution} {Reservation table \textsf{ALU\_TINY} (Section~\ref{sec:reservation-ALU-TINY})}
~
{\begin{lstlisting}
stage RR:
  new argument3 = GPR[registerY];
  new argument2 = GPR[registerZ];
  new result1 = _ROL_64(argument2.64[0], _ZX_6(argument3));
stage E1:
  GPR[registerW] = result1;
\end{lstlisting}}
\newpage

\paragraph{Encoding} Format \textsf{ALU\_DSWRI} (Section~\ref{sec:format-ALU-DSWRI}), \verb|exucode1 = 110|\\

\bitfields{ 1/P, 2/11, 1/0, 1/1, 3/110, 6/registerW, 2/10, 3/100, 1/0, 6/unsigned6, 6/registerZ }


\paragraph{Syntax} \texttt{rold} \emph{registerW} = \emph{registerZ}, \emph{unsigned6}

\paragraph{Description} The \emph{registerZ} is rotated left by the \emph{unsigned6} modulo 64. The result is stored into the \emph{registerW}.


\paragraph{Execution} {Reservation table \textsf{ALU\_TINY} (Section~\ref{sec:reservation-ALU-TINY})}
~
{\begin{lstlisting}
stage RR:
  new argument3 = _ZX_6(unsigned6);
  new argument2 = GPR[registerZ];
  new result1 = _ROL_64(argument2.64[0], _ZX_6(argument3));
stage E1:
  GPR[registerW] = result1;
\end{lstlisting}}
\newpage
\subsubsection[{\small ROLW: Rotate Left Word}]{ROLW\hfill Rotate Left Word}
\label{instr:ROLW}\index{rolw}
 
\paragraph{Encoding} Format \textsf{ALU\_WSWRR} (Section~\ref{sec:format-ALU-WSWRR}), \verb|exucode1 = 110|\\

\bitfields{ 1/P, 2/11, 1/0, 1/0, 3/110, 6/registerW, 2/11, 3/100, 1/signextw, 6/registerY, 6/registerZ }


\paragraph{Syntax} \texttt{rolw}\emph{signextw} \emph{registerW} = \emph{registerZ}, \emph{registerY}

\paragraph{Description} The \emph{registerZ} is rotated left by the \emph{registerY} modulo 32. The result is stored into the \emph{registerW}.


\paragraph{Modifier(s)}
\noindent\begin{itemize}
\item signextw (cf. signextw in section \ref{par:modifier-signextw} on page \pageref{par:modifier-signextw})
\end{itemize}

\paragraph{Execution} {Reservation table \textsf{ALU\_TINY} (Section~\ref{sec:reservation-ALU-TINY})}
~
{\begin{lstlisting}
stage ID:
  new signextw = _ZX_1(signextw);
stage RR:
  new argument3 = GPR[registerY];
  new argument2 = GPR[registerZ];
  new result1 = _ROL_32(argument2.32[0], _ZX_5(argument3));
stage E1:
  GPR[registerW] = signextw ? _SX_32(result1) : _ZX_32(result1);
\end{lstlisting}}
\newpage

\paragraph{Encoding} Format \textsf{ALU\_WSWRI} (Section~\ref{sec:format-ALU-WSWRI}), \verb|exucode1 = 110|\\

\bitfields{ 1/P, 2/11, 1/0, 1/1, 3/110, 6/registerW, 2/11, 3/100, 1/signextw, 6/unsigned6, 6/registerZ }


\paragraph{Syntax} \texttt{rolw}\emph{signextw} \emph{registerW} = \emph{registerZ}, \emph{unsigned6}

\paragraph{Description} The \emph{registerZ} is rotated left by the \emph{unsigned6} modulo 32. The result is stored into the \emph{registerW}.


\paragraph{Modifier(s)}
\noindent\begin{itemize}
\item signextw (cf. signextw in section \ref{par:modifier-signextw} on page \pageref{par:modifier-signextw})
\end{itemize}

\paragraph{Execution} {Reservation table \textsf{ALU\_TINY} (Section~\ref{sec:reservation-ALU-TINY})}
~
{\begin{lstlisting}
stage ID:
  new signextw = _ZX_1(signextw);
stage RR:
  new argument3 = _ZX_6(unsigned6);
  new argument2 = GPR[registerZ];
  new result1 = _ROL_32(argument2.32[0], _ZX_5(argument3));
stage E1:
  GPR[registerW] = signextw ? _SX_32(result1) : _ZX_32(result1);
\end{lstlisting}}
\newpage
\subsubsection[{\small ROLWQ: Rotate Left Word Quadruple}]{ROLWQ\hfill Rotate Left Word Quadruple}
\label{instr:ROLWQ}\index{rolwq}
 
\paragraph{Encoding} Format \textsf{ALU\_WQSWRI} (Section~\ref{sec:format-ALU-WQSWRI}), \verb|exucode1 = 110|\\

\bitfields{ 1/P, 2/11, 1/1, 1/1, 3/110, 5/registerM, 1/--, 2/10, 3/100, 1/0, 6/unsigned6, 5/registerP, 1/1 }


\paragraph{Syntax} \texttt{rolwq} \emph{registerM} = \emph{registerP}, \emph{unsigned6}

\paragraph{Description} The \emph{registerP} is interpreted as four words packed into 128 bits. The \emph{registerP} words are rotated left by the \emph{unsigned6} modulo 32. The four resulting words are packed and stored into the \emph{registerM}.


\paragraph{Execution} {Reservation table \textsf{ALU\_LITE} (Section~\ref{sec:reservation-ALU-LITE})}
~
{\begin{lstlisting}
stage RR:
  new argument3 = _ZX_6(unsigned6);
  new argument2 = PGR[registerP];
  new shift = _ZX_5(argument3);
  for (I = 0; I < 4; I += 1) {
    result1.32[I] = _ROL_32(argument2.32[I], shift)
  };
stage E1:
  PGR[registerM] = result1;
\end{lstlisting}}
\newpage

\paragraph{Encoding} Format \textsf{ALU\_WQSWRR} (Section~\ref{sec:format-ALU-WQSWRR}), \verb|exucode1 = 110|\\

\bitfields{ 1/P, 2/11, 1/1, 1/0, 3/110, 5/registerM, 1/--, 2/10, 3/100, 1/0, 6/registerY, 5/registerP, 1/1 }


\paragraph{Syntax} \texttt{rolwq} \emph{registerM} = \emph{registerP}, \emph{registerY}

\paragraph{Description} The \emph{registerP} is interpreted as four words packed into 128 bits. The \emph{registerP} words are rotated left by the \emph{registerY} modulo 32. The four resulting words are packed and stored into the \emph{registerM}.


\paragraph{Execution} {Reservation table \textsf{ALU\_LITE} (Section~\ref{sec:reservation-ALU-LITE})}
~
{\begin{lstlisting}
stage RR:
  new argument3 = GPR[registerY];
  new argument2 = PGR[registerP];
  new shift = _ZX_5(argument3);
  for (I = 0; I < 4; I += 1) {
    result1.32[I] = _ROL_32(argument2.32[I], shift)
  };
stage E1:
  PGR[registerM] = result1;
\end{lstlisting}}
\newpage
\subsubsection[{\small RORD: Rotate Right Double Word}]{RORD\hfill Rotate Right Double Word}
\label{instr:RORD}\index{rord}
 
\paragraph{Encoding} Format \textsf{ALU\_DSWRR} (Section~\ref{sec:format-ALU-DSWRR}), \verb|exucode1 = 111|\\

\bitfields{ 1/P, 2/11, 1/0, 1/0, 3/111, 6/registerW, 2/10, 3/100, 1/0, 6/registerY, 6/registerZ }


\paragraph{Syntax} \texttt{rord} \emph{registerW} = \emph{registerZ}, \emph{registerY}

\paragraph{Description} The \emph{registerZ} is rotated right by the \emph{registerY} modulo 64. The result is stored into the \emph{registerW}.


\paragraph{Execution} {Reservation table \textsf{ALU\_TINY} (Section~\ref{sec:reservation-ALU-TINY})}
~
{\begin{lstlisting}
stage RR:
  new argument3 = GPR[registerY];
  new argument2 = GPR[registerZ];
  new result1 = _ROR_64(argument2.64[0], _ZX_6(argument3));
stage E1:
  GPR[registerW] = result1;
\end{lstlisting}}
\newpage

\paragraph{Encoding} Format \textsf{ALU\_DSWRI} (Section~\ref{sec:format-ALU-DSWRI}), \verb|exucode1 = 111|\\

\bitfields{ 1/P, 2/11, 1/0, 1/1, 3/111, 6/registerW, 2/10, 3/100, 1/0, 6/unsigned6, 6/registerZ }


\paragraph{Syntax} \texttt{rord} \emph{registerW} = \emph{registerZ}, \emph{unsigned6}

\paragraph{Description} The \emph{registerZ} is rotated right by the \emph{unsigned6} modulo 64. The result is stored into the \emph{registerW}.


\paragraph{Execution} {Reservation table \textsf{ALU\_TINY} (Section~\ref{sec:reservation-ALU-TINY})}
~
{\begin{lstlisting}
stage RR:
  new argument3 = _ZX_6(unsigned6);
  new argument2 = GPR[registerZ];
  new result1 = _ROR_64(argument2.64[0], _ZX_6(argument3));
stage E1:
  GPR[registerW] = result1;
\end{lstlisting}}
\newpage
\subsubsection[{\small RORW: Rotate Right Word}]{RORW\hfill Rotate Right Word}
\label{instr:RORW}\index{rorw}
 
\paragraph{Encoding} Format \textsf{ALU\_WSWRR} (Section~\ref{sec:format-ALU-WSWRR}), \verb|exucode1 = 111|\\

\bitfields{ 1/P, 2/11, 1/0, 1/0, 3/111, 6/registerW, 2/11, 3/100, 1/signextw, 6/registerY, 6/registerZ }


\paragraph{Syntax} \texttt{rorw}\emph{signextw} \emph{registerW} = \emph{registerZ}, \emph{registerY}

\paragraph{Description} The \emph{registerZ} is rotated right by the \emph{registerY} modulo 32. The result is stored into the \emph{registerW}.


\paragraph{Modifier(s)}
\noindent\begin{itemize}
\item signextw (cf. signextw in section \ref{par:modifier-signextw} on page \pageref{par:modifier-signextw})
\end{itemize}

\paragraph{Execution} {Reservation table \textsf{ALU\_TINY} (Section~\ref{sec:reservation-ALU-TINY})}
~
{\begin{lstlisting}
stage ID:
  new signextw = _ZX_1(signextw);
stage RR:
  new argument3 = GPR[registerY];
  new argument2 = GPR[registerZ];
  new result1 = _ROR_32(argument2.32[0], _ZX_5(argument3));
stage E1:
  GPR[registerW] = signextw ? _SX_32(result1) : _ZX_32(result1);
\end{lstlisting}}
\newpage

\paragraph{Encoding} Format \textsf{ALU\_WSWRI} (Section~\ref{sec:format-ALU-WSWRI}), \verb|exucode1 = 111|\\

\bitfields{ 1/P, 2/11, 1/0, 1/1, 3/111, 6/registerW, 2/11, 3/100, 1/signextw, 6/unsigned6, 6/registerZ }


\paragraph{Syntax} \texttt{rorw}\emph{signextw} \emph{registerW} = \emph{registerZ}, \emph{unsigned6}

\paragraph{Description} The \emph{registerZ} is rotated right by the \emph{unsigned6} modulo 32. The result is stored into the \emph{registerW}.


\paragraph{Modifier(s)}
\noindent\begin{itemize}
\item signextw (cf. signextw in section \ref{par:modifier-signextw} on page \pageref{par:modifier-signextw})
\end{itemize}

\paragraph{Execution} {Reservation table \textsf{ALU\_TINY} (Section~\ref{sec:reservation-ALU-TINY})}
~
{\begin{lstlisting}
stage ID:
  new signextw = _ZX_1(signextw);
stage RR:
  new argument3 = _ZX_6(unsigned6);
  new argument2 = GPR[registerZ];
  new result1 = _ROR_32(argument2.32[0], _ZX_5(argument3));
stage E1:
  GPR[registerW] = signextw ? _SX_32(result1) : _ZX_32(result1);
\end{lstlisting}}
\newpage
\subsubsection[{\small RORWQ: Rotate Right Word Quadruple}]{RORWQ\hfill Rotate Right Word Quadruple}
\label{instr:RORWQ}\index{rorwq}
 
\paragraph{Encoding} Format \textsf{ALU\_WQSWRI} (Section~\ref{sec:format-ALU-WQSWRI}), \verb|exucode1 = 111|\\

\bitfields{ 1/P, 2/11, 1/1, 1/1, 3/111, 5/registerM, 1/--, 2/10, 3/100, 1/0, 6/unsigned6, 5/registerP, 1/1 }


\paragraph{Syntax} \texttt{rorwq} \emph{registerM} = \emph{registerP}, \emph{unsigned6}

\paragraph{Description} The \emph{registerP} is interpreted as four words packed into 128 bits. The \emph{registerP} words are rotated right by the \emph{unsigned6} modulo 32. The four resulting words are packed and stored into the \emph{registerM}.


\paragraph{Execution} {Reservation table \textsf{ALU\_LITE} (Section~\ref{sec:reservation-ALU-LITE})}
~
{\begin{lstlisting}
stage RR:
  new argument3 = _ZX_6(unsigned6);
  new argument2 = PGR[registerP];
  new shift = _ZX_5(argument3);
  for (I = 0; I < 4; I += 1) {
    result1.32[I] = _ROR_32(argument2.32[I], shift)
  };
stage E1:
  PGR[registerM] = result1;
\end{lstlisting}}
\newpage

\paragraph{Encoding} Format \textsf{ALU\_WQSWRR} (Section~\ref{sec:format-ALU-WQSWRR}), \verb|exucode1 = 111|\\

\bitfields{ 1/P, 2/11, 1/1, 1/0, 3/111, 5/registerM, 1/--, 2/10, 3/100, 1/0, 6/registerY, 5/registerP, 1/1 }


\paragraph{Syntax} \texttt{rorwq} \emph{registerM} = \emph{registerP}, \emph{registerY}

\paragraph{Description} The \emph{registerP} is interpreted as four words packed into 128 bits. The \emph{registerP} words are rotated right by the \emph{registerY} modulo 32. The four resulting words are packed and stored into the \emph{registerM}.


\paragraph{Execution} {Reservation table \textsf{ALU\_LITE} (Section~\ref{sec:reservation-ALU-LITE})}
~
{\begin{lstlisting}
stage RR:
  new argument3 = GPR[registerY];
  new argument2 = PGR[registerP];
  new shift = _ZX_5(argument3);
  for (I = 0; I < 4; I += 1) {
    result1.32[I] = _ROR_32(argument2.32[I], shift)
  };
stage E1:
  PGR[registerM] = result1;
\end{lstlisting}}
\newpage
\subsubsection[{\small SBFBX: Subtract From Byte Hexadecuple}]{SBFBX\hfill Subtract From Byte Hexadecuple}
\label{instr:SBFBX}\index{sbfbx}
 
\paragraph{Encoding} Format \textsf{ALU\_BXWRR0} (Section~\ref{sec:format-ALU-BXWRR0}), \verb|exucode2 = 0011|\\

\bitfields{ 1/P, 2/11, 1/1, 4/0011, 5/registerM, 1/--, 2/10, 3/000, 1/1, 5/registerO, 1/--, 5/registerP, 1/1 }


\paragraph{Syntax} \texttt{sbfbx} \emph{registerM} = \emph{registerP}, \emph{registerO}

\paragraph{Description} The \emph{registerO} and the \emph{registerP} are interpreted as sixteen bytes packed into 128 bits. The \emph{registerP} bytes are subtracted from the \emph{registerO} bytes. The sixteen resulting bytes are packed and stored into the \emph{registerM}.


\paragraph{Execution} {Reservation table \textsf{ALU\_LITE} (Section~\ref{sec:reservation-ALU-LITE})}
~
{\begin{lstlisting}
stage RR:
  new argument3 = PGR[registerO];
  new argument2 = PGR[registerP];
  for (I = 0; I < 16; I += 1) {
    result1.8[I] = argument3.8[I] - argument2.8[I]
  };
stage E1:
  PGR[registerM] = result1;
\end{lstlisting}}
\newpage

\paragraph{Encoding} Format \textsf{ALU\_BXWRR0.M} (Section~\ref{sec:format-ALU-BXWRR0.M}), \verb|exucode2 = 0011|\\

\bitfields{ 1/P, 2/00, 2/-- --, 27/upper27, 1/1, 2/11, 1/1, 4/0011, 5/registerM, 1/--, 2/10, 3/000, 1/1, 1/splat32, 5/lower5, 5/registerP, 1/1 }


\paragraph{Syntax} \texttt{sbfbx} \emph{registerM} = \emph{registerP}, \emph{upper27\_\discretionary{}{}{}lower5}\emph{splat32}

\paragraph{Description} The \emph{upper27\_\discretionary{}{}{}lower5} and the \emph{registerP} are interpreted as sixteen bytes packed into 128 bits. The \emph{registerP} bytes are subtracted from the \emph{upper27\_\discretionary{}{}{}lower5} bytes. The sixteen resulting bytes are packed and stored into the \emph{registerM}.


\paragraph{Modifier(s)}
\noindent\begin{itemize}
\item splat32 (cf. splat32 in section \ref{par:modifier-splat32} on page \pageref{par:modifier-splat32})
\end{itemize}

\paragraph{Execution} {Reservation table \textsf{ALU\_LITE.X} (Section~\ref{sec:reservation-ALU-LITE.X})}
~
{\begin{lstlisting}
stage ID:
  new splat32 = _ZX_1(splat32);
stage RR:
  new argument3 = splat32 ? (_WX_32(upper27_lower5) << 32) | _ZX_32(_WX_32(upper27_lower5)) : _SX_32(_WX_32(upper27_lower5));
  new argument2 = PGR[registerP];
  for (I = 0; I < 16; I += 1) {
    result1.8[I] = argument3.8[I] - argument2.8[I]
  };
stage E1:
  PGR[registerM] = result1;
\end{lstlisting}}
\newpage
\subsubsection[{\small SBFD: Subtract From Double Word}]{SBFD\hfill Subtract From Double Word}
\label{instr:SBFD}\index{sbfd}
 
\paragraph{Encoding} Format \textsf{ALU\_DWRR0} (Section~\ref{sec:format-ALU-DWRR0}), \verb|exucode2 = 0011|\\

\bitfields{ 1/P, 2/11, 1/0, 4/0011, 6/registerW, 2/10, 3/000, 1/0, 6/registerY, 6/registerZ }


\paragraph{Syntax} \texttt{sbfd} \emph{registerW} = \emph{registerZ}, \emph{registerY}

\paragraph{Description} The \emph{registerZ} is subtracted from the \emph{registerY}. The result is stored into the \emph{registerW}.


\paragraph{Execution} {Reservation table \textsf{ALU\_TINY} (Section~\ref{sec:reservation-ALU-TINY})}
~
{\begin{lstlisting}
stage RR:
  new argument3 = GPR[registerY];
  new argument2 = GPR[registerZ];
  new result1 = argument3 - argument2;
stage E1:
  GPR[registerW] = result1;
\end{lstlisting}}
\newpage

\paragraph{Encoding} Format \textsf{ALU\_DWRR0.M} (Section~\ref{sec:format-ALU-DWRR0.M}), \verb|exucode2 = 0011|\\

\bitfields{ 1/P, 2/00, 2/-- --, 27/upper27, 1/1, 2/11, 1/0, 4/0011, 6/registerW, 2/10, 3/000, 1/0, 1/splat32, 5/lower5, 6/registerZ }


\paragraph{Syntax} \texttt{sbfd} \emph{registerW} = \emph{registerZ}, \emph{upper27\_\discretionary{}{}{}lower5}\emph{splat32}

\paragraph{Description} The \emph{registerZ} is subtracted from the \emph{upper27\_\discretionary{}{}{}lower5}. The result is stored into the \emph{registerW}.


\paragraph{Modifier(s)}
\noindent\begin{itemize}
\item splat32 (cf. splat32 in section \ref{par:modifier-splat32} on page \pageref{par:modifier-splat32})
\end{itemize}

\paragraph{Execution} {Reservation table \textsf{ALU\_TINY.X} (Section~\ref{sec:reservation-ALU-TINY.X})}
~
{\begin{lstlisting}
stage ID:
  new splat32 = _ZX_1(splat32);
stage RR:
  new argument3 = splat32 ? (_WX_32(upper27_lower5) << 32) | _ZX_32(_WX_32(upper27_lower5)) : _SX_32(_WX_32(upper27_lower5));
  new argument2 = GPR[registerZ];
  new result1 = argument3 - argument2;
stage E1:
  GPR[registerW] = result1;
\end{lstlisting}}
\newpage

\paragraph{Encoding} Format \textsf{ALU\_DWRI} (Section~\ref{sec:format-ALU-DWRI}), \verb|exucode2 = 0011|\\

\bitfields{ 1/P, 2/11, 1/0, 4/0011, 6/registerW, 2/00, 10/signed10, 6/registerZ }


\paragraph{Syntax} \texttt{sbfd} \emph{registerW} = \emph{registerZ}, \emph{signed10}

\paragraph{Description} The \emph{registerZ} is subtracted from the \emph{signed10}. The result is stored into the \emph{registerW}.


\paragraph{Execution} {Reservation table \textsf{ALU\_TINY} (Section~\ref{sec:reservation-ALU-TINY})}
~
{\begin{lstlisting}
stage RR:
  new argument3 = _SX_10(signed10);
  new argument2 = GPR[registerZ];
  new result1 = argument3 - argument2;
stage E1:
  GPR[registerW] = result1;
\end{lstlisting}}
\newpage

\paragraph{Encoding} Format \textsf{ALU\_DWRI.X} (Section~\ref{sec:format-ALU-DWRI.X}), \verb|exucode2 = 0011|\\

\bitfields{ 1/P, 2/00, 2/-- --, 27/upper27, 1/1, 2/11, 1/0, 4/0011, 6/registerW, 2/00, 10/lower10, 6/registerZ }


\paragraph{Syntax} \texttt{sbfd} \emph{registerW} = \emph{registerZ}, \emph{upper27\_\discretionary{}{}{}lower10}

\paragraph{Description} The \emph{registerZ} is subtracted from the \emph{upper27\_\discretionary{}{}{}lower10}. The result is stored into the \emph{registerW}.


\paragraph{Execution} {Reservation table \textsf{ALU\_TINY.X} (Section~\ref{sec:reservation-ALU-TINY.X})}
~
{\begin{lstlisting}
stage RR:
  new argument3 = _SX_37(upper27_lower10);
  new argument2 = GPR[registerZ];
  new result1 = argument3 - argument2;
stage E1:
  GPR[registerW] = result1;
\end{lstlisting}}
\newpage

\paragraph{Encoding} Format \textsf{ALU\_DWRI.Y} (Section~\ref{sec:format-ALU-DWRI.Y}), \verb|exucode2 = 0011|\\

\bitfields{ 1/P, 2/00, 2/-- --, 27/extend27, 1/1, 2/00, 2/-- --, 27/upper27, 1/1, 2/11, 1/0, 4/0011, 6/registerW, 2/00, 10/lower10, 6/registerZ }


\paragraph{Syntax} \texttt{sbfd} \emph{registerW} = \emph{registerZ}, \emph{extend27\_\discretionary{}{}{}upper27\_\discretionary{}{}{}lower10}

\paragraph{Description} The \emph{registerZ} is subtracted from the \emph{extend27\_\discretionary{}{}{}upper27\_\discretionary{}{}{}lower10}. The result is stored into the \emph{registerW}.


\paragraph{Execution} {Reservation table \textsf{ALU\_TINY.Y} (Section~\ref{sec:reservation-ALU-TINY.Y})}
~
{\begin{lstlisting}
stage RR:
  new argument3 = _WX_64(extend27_upper27_lower10);
  new argument2 = GPR[registerZ];
  new result1 = argument3 - argument2;
stage E1:
  GPR[registerW] = result1;
\end{lstlisting}}
\newpage
\subsubsection[{\small SBFDP: Subtract From Double Word Pair}]{SBFDP\hfill Subtract From Double Word Pair}
\label{instr:SBFDP}\index{sbfdp}
 
\paragraph{Encoding} Format \textsf{ALU\_DPWRR0} (Section~\ref{sec:format-ALU-DPWRR0}), \verb|exucode2 = 0011|\\

\bitfields{ 1/P, 2/11, 1/1, 4/0011, 5/registerM, 1/--, 2/10, 3/000, 1/0, 5/registerO, 1/--, 5/registerP, 1/0 }


\paragraph{Syntax} \texttt{sbfdp} \emph{registerM} = \emph{registerP}, \emph{registerO}

\paragraph{Description} The \emph{registerO} and the \emph{registerP} are interpreted as two double words packed into 128 bits. The \emph{registerP} words are subtracted from the \emph{registerO} words. The two resulting double words are packed and stored into the \emph{registerM}.


\paragraph{Execution} {Reservation table \textsf{ALU\_LITE} (Section~\ref{sec:reservation-ALU-LITE})}
~
{\begin{lstlisting}
stage RR:
  new argument3 = PGR[registerO];
  new argument2 = PGR[registerP];
  for (I = 0; I < 2; I += 1) {
    result1.64[I] = argument3.64[I] - argument2.64[I]
  };
stage E1:
  PGR[registerM] = result1;
\end{lstlisting}}
\newpage

\paragraph{Encoding} Format \textsf{ALU\_DPWRR0.M} (Section~\ref{sec:format-ALU-DPWRR0.M}), \verb|exucode2 = 0011|\\

\bitfields{ 1/P, 2/00, 2/-- --, 27/upper27, 1/1, 2/11, 1/1, 4/0011, 5/registerM, 1/--, 2/10, 3/000, 1/0, 1/splat32, 5/lower5, 5/registerP, 1/0 }


\paragraph{Syntax} \texttt{sbfdp} \emph{registerM} = \emph{registerP}, \emph{upper27\_\discretionary{}{}{}lower5}\emph{splat32}

\paragraph{Description} The \emph{upper27\_\discretionary{}{}{}lower5} and the \emph{registerP} are interpreted as two double words packed into 128 bits. The \emph{registerP} words are subtracted from the \emph{upper27\_\discretionary{}{}{}lower5} words. The two resulting double words are packed and stored into the \emph{registerM}.


\paragraph{Modifier(s)}
\noindent\begin{itemize}
\item splat32 (cf. splat32 in section \ref{par:modifier-splat32} on page \pageref{par:modifier-splat32})
\end{itemize}

\paragraph{Execution} {Reservation table \textsf{ALU\_LITE.X} (Section~\ref{sec:reservation-ALU-LITE.X})}
~
{\begin{lstlisting}
stage ID:
  new splat32 = _ZX_1(splat32);
stage RR:
  new argument3 = splat32 ? (_WX_32(upper27_lower5) << 32) | _ZX_32(_WX_32(upper27_lower5)) : _SX_32(_WX_32(upper27_lower5));
  new argument2 = PGR[registerP];
  for (I = 0; I < 2; I += 1) {
    result1.64[I] = argument3.64[I] - argument2.64[I]
  };
stage E1:
  PGR[registerM] = result1;
\end{lstlisting}}
\newpage
\subsubsection[{\small SBFHO: Subtract From Half Word Octuple}]{SBFHO\hfill Subtract From Half Word Octuple}
\label{instr:SBFHO}\index{sbfho}
 
\paragraph{Encoding} Format \textsf{ALU\_HOWRR0} (Section~\ref{sec:format-ALU-HOWRR0}), \verb|exucode2 = 0011|\\

\bitfields{ 1/P, 2/11, 1/1, 4/0011, 5/registerM, 1/--, 2/10, 3/000, 1/1, 5/registerO, 1/--, 5/registerP, 1/0 }


\paragraph{Syntax} \texttt{sbfho} \emph{registerM} = \emph{registerP}, \emph{registerO}

\paragraph{Description} The \emph{registerO} and the \emph{registerP} are interpreted as eight half words packed into 128 bits. The \emph{registerP} half words are subtracted from the \emph{registerO} half words. The eight resulting half words are packed and stored into the \emph{registerM}.


\paragraph{Execution} {Reservation table \textsf{ALU\_LITE} (Section~\ref{sec:reservation-ALU-LITE})}
~
{\begin{lstlisting}
stage RR:
  new argument3 = PGR[registerO];
  new argument2 = PGR[registerP];
  for (I = 0; I < 8; I += 1) {
    result1.16[I] = argument3.16[I] - argument2.16[I]
  };
stage E1:
  PGR[registerM] = result1;
\end{lstlisting}}
\newpage

\paragraph{Encoding} Format \textsf{ALU\_HOWRR0.M} (Section~\ref{sec:format-ALU-HOWRR0.M}), \verb|exucode2 = 0011|\\

\bitfields{ 1/P, 2/00, 2/-- --, 27/upper27, 1/1, 2/11, 1/1, 4/0011, 5/registerM, 1/--, 2/10, 3/000, 1/1, 1/splat32, 5/lower5, 5/registerP, 1/0 }


\paragraph{Syntax} \texttt{sbfho} \emph{registerM} = \emph{registerP}, \emph{upper27\_\discretionary{}{}{}lower5}\emph{splat32}

\paragraph{Description} The \emph{upper27\_\discretionary{}{}{}lower5} and the \emph{registerP} are interpreted as eight half words packed into 128 bits. The \emph{registerP} half words are subtracted from the \emph{upper27\_\discretionary{}{}{}lower5} half words. The eight resulting half words are packed and stored into the \emph{registerM}.


\paragraph{Modifier(s)}
\noindent\begin{itemize}
\item splat32 (cf. splat32 in section \ref{par:modifier-splat32} on page \pageref{par:modifier-splat32})
\end{itemize}

\paragraph{Execution} {Reservation table \textsf{ALU\_LITE.X} (Section~\ref{sec:reservation-ALU-LITE.X})}
~
{\begin{lstlisting}
stage ID:
  new splat32 = _ZX_1(splat32);
stage RR:
  new argument3 = splat32 ? (_WX_32(upper27_lower5) << 32) | _ZX_32(_WX_32(upper27_lower5)) : _SX_32(_WX_32(upper27_lower5));
  new argument2 = PGR[registerP];
  for (I = 0; I < 8; I += 1) {
    result1.16[I] = argument3.16[I] - argument2.16[I]
  };
stage E1:
  PGR[registerM] = result1;
\end{lstlisting}}
\newpage
\subsubsection[{\small SBFQ: Subtract From Quadruple Word}]{SBFQ\hfill Subtract From Quadruple Word}
\label{instr:SBFQ}\index{sbfq}
 
\paragraph{Encoding} Format \textsf{ALU\_QWRR0} (Section~\ref{sec:format-ALU-QWRR0}), \verb|exucode2 = 0011|\\

\bitfields{ 1/P, 2/11, 1/0, 4/0011, 5/registerM, 1/--, 2/10, 3/000, 1/1, 5/registerO, 1/--, 5/registerP, 1/-- }


\paragraph{Syntax} \texttt{sbfq} \emph{registerM} = \emph{registerP}, \emph{registerO}

\paragraph{Description} The \emph{registerP} is subtracted from the \emph{registerO}. The result is stored into the \emph{registerM}.


\paragraph{Execution} {Reservation table \textsf{ALU\_LITE} (Section~\ref{sec:reservation-ALU-LITE})}
~
{\begin{lstlisting}
stage RR:
  new argument3 = PGR[registerO];
  new argument2 = PGR[registerP];
  new result1 = argument3 - argument2;
stage E1:
  PGR[registerM] = result1;
\end{lstlisting}}
\newpage

\paragraph{Encoding} Format \textsf{ALU\_QWRR0.M} (Section~\ref{sec:format-ALU-QWRR0.M}), \verb|exucode2 = 0011|\\

\bitfields{ 1/P, 2/00, 2/-- --, 27/upper27, 1/1, 2/11, 1/0, 4/0011, 5/registerM, 1/--, 2/10, 3/000, 1/1, 1/splat32, 5/lower5, 5/registerP, 1/1 }


\paragraph{Syntax} \texttt{sbfq} \emph{registerM} = \emph{registerP}, \emph{upper27\_\discretionary{}{}{}lower5}\emph{splat32}

\paragraph{Description} The \emph{registerP} is subtracted from the \emph{upper27\_\discretionary{}{}{}lower5}. The result is stored into the \emph{registerM}.


\paragraph{Modifier(s)}
\noindent\begin{itemize}
\item splat32 (cf. splat32 in section \ref{par:modifier-splat32} on page \pageref{par:modifier-splat32})
\end{itemize}

\paragraph{Execution} {Reservation table \textsf{ALU\_LITE.X} (Section~\ref{sec:reservation-ALU-LITE.X})}
~
{\begin{lstlisting}
stage ID:
  new splat32 = _ZX_1(splat32);
stage RR:
  new argument3 = splat32 ? (_WX_32(upper27_lower5) << 32) | _ZX_32(_WX_32(upper27_lower5)) : _SX_32(_WX_32(upper27_lower5));
  new argument2 = PGR[registerP];
  new result1 = argument3 - argument2;
stage E1:
  PGR[registerM] = result1;
\end{lstlisting}}
\newpage
\subsubsection[{\small SBFSBX: Subtract Saturated Byte Hexadecuple}]{SBFSBX\hfill Subtract Saturated Byte Hexadecuple}
\label{instr:SBFSBX}\index{sbfsbx}
 
\paragraph{Encoding} Format \textsf{ALU\_BXWRR1} (Section~\ref{sec:format-ALU-BXWRR1}), \verb|exucode2 = 1001|\\

\bitfields{ 1/P, 2/11, 1/1, 4/1001, 5/registerM, 1/--, 2/10, 3/101, 1/1, 5/registerO, 1/--, 5/registerP, 1/1 }


\paragraph{Syntax} \texttt{sbfsbx} \emph{registerM} = \emph{registerP}, \emph{registerO}

\paragraph{Description} The \emph{registerO} and the \emph{registerP} are interpreted as sixteen bytes packed into 128 bits. The \emph{registerP} bytes are subtracted with saturation to 8 bits from the \emph{registerO} bytes. The sixteen resulting bytes are packed and stored into the \emph{registerM}.


\paragraph{Execution} {Reservation table \textsf{ALU\_LITE} (Section~\ref{sec:reservation-ALU-LITE})}
~
{\begin{lstlisting}
stage RR:
  new argument3 = PGR[registerO];
  new argument2 = PGR[registerP];
  for (I = 0; I < 16; I += 1) {
    result1.8[I] = _SAT_8(_SX_8(argument3.8[I]) - _SX_8(argument2.8[I]))
  };
stage E1:
  PGR[registerM] = result1;
\end{lstlisting}}
\newpage

\paragraph{Encoding} Format \textsf{ALU\_BXWRR1.M} (Section~\ref{sec:format-ALU-BXWRR1.M}), \verb|exucode2 = 1001|\\

\bitfields{ 1/P, 2/00, 2/-- --, 27/upper27, 1/1, 2/11, 1/1, 4/1001, 5/registerM, 1/--, 2/10, 3/101, 1/1, 1/splat32, 5/lower5, 5/registerP, 1/1 }


\paragraph{Syntax} \texttt{sbfsbx} \emph{registerM} = \emph{registerP}, \emph{upper27\_\discretionary{}{}{}lower5}\emph{splat32}

\paragraph{Description} The \emph{upper27\_\discretionary{}{}{}lower5} and the \emph{registerP} are interpreted as sixteen bytes packed into 128 bits. The \emph{registerP} bytes are subtracted with saturation to 8 bits from the \emph{upper27\_\discretionary{}{}{}lower5} bytes. The sixteen resulting bytes are packed and stored into the \emph{registerM}.


\paragraph{Modifier(s)}
\noindent\begin{itemize}
\item splat32 (cf. splat32 in section \ref{par:modifier-splat32} on page \pageref{par:modifier-splat32})
\end{itemize}

\paragraph{Execution} {Reservation table \textsf{ALU\_LITE.X} (Section~\ref{sec:reservation-ALU-LITE.X})}
~
{\begin{lstlisting}
stage ID:
  new splat32 = _ZX_1(splat32);
stage RR:
  new argument3 = splat32 ? (_WX_32(upper27_lower5) << 32) | _ZX_32(_WX_32(upper27_lower5)) : _SX_32(_WX_32(upper27_lower5));
  new argument2 = PGR[registerP];
  for (I = 0; I < 16; I += 1) {
    result1.8[I] = _SAT_8(_SX_8(argument3.8[I]) - _SX_8(argument2.8[I]))
  };
stage E1:
  PGR[registerM] = result1;
\end{lstlisting}}
\newpage
\subsubsection[{\small SBFSD: Subtract Saturated Double Word}]{SBFSD\hfill Subtract Saturated Double Word}
\label{instr:SBFSD}\index{sbfsd}
 
\paragraph{Encoding} Format \textsf{ALU\_DWRR1} (Section~\ref{sec:format-ALU-DWRR1}), \verb|exucode2 = 1001|\\

\bitfields{ 1/P, 2/11, 1/0, 4/1001, 6/registerW, 2/10, 3/101, 1/0, 6/registerY, 6/registerZ }


\paragraph{Syntax} \texttt{sbfsd} \emph{registerW} = \emph{registerZ}, \emph{registerY}

\paragraph{Description} The \emph{registerZ} is subtracted from the \emph{registerY} with 64-bit saturation.


\paragraph{Execution} {Reservation table \textsf{ALU\_TINY} (Section~\ref{sec:reservation-ALU-TINY})}
~
{\begin{lstlisting}
stage RR:
  new argument3 = GPR[registerY];
  new argument2 = GPR[registerZ];
  new result1 = _SAT_64(_SX_64(argument3) - _SX_64(argument2));
stage E1:
  GPR[registerW] = result1;
\end{lstlisting}}
\newpage

\paragraph{Encoding} Format \textsf{ALU\_DWRR1.M} (Section~\ref{sec:format-ALU-DWRR1.M}), \verb|exucode2 = 1001|\\

\bitfields{ 1/P, 2/00, 2/-- --, 27/upper27, 1/1, 2/11, 1/0, 4/1001, 6/registerW, 2/10, 3/101, 1/0, 1/splat32, 5/lower5, 6/registerZ }


\paragraph{Syntax} \texttt{sbfsd} \emph{registerW} = \emph{registerZ}, \emph{upper27\_\discretionary{}{}{}lower5}\emph{splat32}

\paragraph{Description} The \emph{registerZ} is subtracted from the \emph{upper27\_\discretionary{}{}{}lower5} with 64-bit saturation.


\paragraph{Modifier(s)}
\noindent\begin{itemize}
\item splat32 (cf. splat32 in section \ref{par:modifier-splat32} on page \pageref{par:modifier-splat32})
\end{itemize}

\paragraph{Execution} {Reservation table \textsf{ALU\_TINY.X} (Section~\ref{sec:reservation-ALU-TINY.X})}
~
{\begin{lstlisting}
stage ID:
  new splat32 = _ZX_1(splat32);
stage RR:
  new argument3 = splat32 ? (_WX_32(upper27_lower5) << 32) | _ZX_32(_WX_32(upper27_lower5)) : _SX_32(_WX_32(upper27_lower5));
  new argument2 = GPR[registerZ];
  new result1 = _SAT_64(_SX_64(argument3) - _SX_64(argument2));
stage E1:
  GPR[registerW] = result1;
\end{lstlisting}}
\newpage
\subsubsection[{\small SBFSDP: Subtract Saturated Double Word Pair}]{SBFSDP\hfill Subtract Saturated Double Word Pair}
\label{instr:SBFSDP}\index{sbfsdp}
 
\paragraph{Encoding} Format \textsf{ALU\_DPWRR1} (Section~\ref{sec:format-ALU-DPWRR1}), \verb|exucode2 = 1001|\\

\bitfields{ 1/P, 2/11, 1/1, 4/1001, 5/registerM, 1/--, 2/10, 3/101, 1/0, 5/registerO, 1/--, 5/registerP, 1/0 }


\paragraph{Syntax} \texttt{sbfsdp} \emph{registerM} = \emph{registerP}, \emph{registerO}

\paragraph{Description} The \emph{registerO} and the \emph{registerP} are interpreted as two double words packed into 128 bits. The \emph{registerP} double words are subtracted with saturation to 64 bits from the \emph{registerO} double words. The two resulting double words are packed and stored into the \emph{registerM}.


\paragraph{Execution} {Reservation table \textsf{ALU\_LITE} (Section~\ref{sec:reservation-ALU-LITE})}
~
{\begin{lstlisting}
stage RR:
  new argument3 = PGR[registerO];
  new argument2 = PGR[registerP];
  for (I = 0; I < 2; I += 1) {
    result1.64[I] = _SAT_64(_SX_64(argument3.64[I]) - _SX_64(argument2.64[I]))
  };
stage E1:
  PGR[registerM] = result1;
\end{lstlisting}}
\newpage

\paragraph{Encoding} Format \textsf{ALU\_DPWRR1.M} (Section~\ref{sec:format-ALU-DPWRR1.M}), \verb|exucode2 = 1001|\\

\bitfields{ 1/P, 2/00, 2/-- --, 27/upper27, 1/1, 2/11, 1/1, 4/1001, 5/registerM, 1/--, 2/10, 3/101, 1/0, 1/splat32, 5/lower5, 5/registerP, 1/0 }


\paragraph{Syntax} \texttt{sbfsdp} \emph{registerM} = \emph{registerP}, \emph{upper27\_\discretionary{}{}{}lower5}\emph{splat32}

\paragraph{Description} The \emph{upper27\_\discretionary{}{}{}lower5} and the \emph{registerP} are interpreted as two double words packed into 128 bits. The \emph{registerP} double words are subtracted with saturation to 64 bits from the \emph{upper27\_\discretionary{}{}{}lower5} double words. The two resulting double words are packed and stored into the \emph{registerM}.


\paragraph{Modifier(s)}
\noindent\begin{itemize}
\item splat32 (cf. splat32 in section \ref{par:modifier-splat32} on page \pageref{par:modifier-splat32})
\end{itemize}

\paragraph{Execution} {Reservation table \textsf{ALU\_LITE.X} (Section~\ref{sec:reservation-ALU-LITE.X})}
~
{\begin{lstlisting}
stage ID:
  new splat32 = _ZX_1(splat32);
stage RR:
  new argument3 = splat32 ? (_WX_32(upper27_lower5) << 32) | _ZX_32(_WX_32(upper27_lower5)) : _SX_32(_WX_32(upper27_lower5));
  new argument2 = PGR[registerP];
  for (I = 0; I < 2; I += 1) {
    result1.64[I] = _SAT_64(_SX_64(argument3.64[I]) - _SX_64(argument2.64[I]))
  };
stage E1:
  PGR[registerM] = result1;
\end{lstlisting}}
\newpage
\subsubsection[{\small SBFSHO: Subtract Saturated Half Word Octuple}]{SBFSHO\hfill Subtract Saturated Half Word Octuple}
\label{instr:SBFSHO}\index{sbfsho}
 
\paragraph{Encoding} Format \textsf{ALU\_HOWRR1} (Section~\ref{sec:format-ALU-HOWRR1}), \verb|exucode2 = 1001|\\

\bitfields{ 1/P, 2/11, 1/1, 4/1001, 5/registerM, 1/--, 2/10, 3/101, 1/1, 5/registerO, 1/--, 5/registerP, 1/0 }


\paragraph{Syntax} \texttt{sbfsho} \emph{registerM} = \emph{registerP}, \emph{registerO}

\paragraph{Description} The \emph{registerO} and the \emph{registerP} are interpreted as eight half words packed into 128 bits. The \emph{registerP} half words are subtracted with saturation to 16 bits from the \emph{registerO} half words. The eight resulting half words are packed and stored into the \emph{registerM}.


\paragraph{Execution} {Reservation table \textsf{ALU\_LITE} (Section~\ref{sec:reservation-ALU-LITE})}
~
{\begin{lstlisting}
stage RR:
  new argument3 = PGR[registerO];
  new argument2 = PGR[registerP];
  for (I = 0; I < 8; I += 1) {
    result1.16[I] = _SAT_16(_SX_16(argument3.16[I]) - _SX_16(argument2.16[I]))
  };
stage E1:
  PGR[registerM] = result1;
\end{lstlisting}}
\newpage

\paragraph{Encoding} Format \textsf{ALU\_HOWRR1.M} (Section~\ref{sec:format-ALU-HOWRR1.M}), \verb|exucode2 = 1001|\\

\bitfields{ 1/P, 2/00, 2/-- --, 27/upper27, 1/1, 2/11, 1/1, 4/1001, 5/registerM, 1/--, 2/10, 3/101, 1/1, 1/splat32, 5/lower5, 5/registerP, 1/0 }


\paragraph{Syntax} \texttt{sbfsho} \emph{registerM} = \emph{registerP}, \emph{upper27\_\discretionary{}{}{}lower5}\emph{splat32}

\paragraph{Description} The \emph{upper27\_\discretionary{}{}{}lower5} and the \emph{registerP} are interpreted as eight half words packed into 128 bits. The \emph{registerP} half words are subtracted with saturation to 16 bits from the \emph{upper27\_\discretionary{}{}{}lower5} half words. The eight resulting half words are packed and stored into the \emph{registerM}.


\paragraph{Modifier(s)}
\noindent\begin{itemize}
\item splat32 (cf. splat32 in section \ref{par:modifier-splat32} on page \pageref{par:modifier-splat32})
\end{itemize}

\paragraph{Execution} {Reservation table \textsf{ALU\_LITE.X} (Section~\ref{sec:reservation-ALU-LITE.X})}
~
{\begin{lstlisting}
stage ID:
  new splat32 = _ZX_1(splat32);
stage RR:
  new argument3 = splat32 ? (_WX_32(upper27_lower5) << 32) | _ZX_32(_WX_32(upper27_lower5)) : _SX_32(_WX_32(upper27_lower5));
  new argument2 = PGR[registerP];
  for (I = 0; I < 8; I += 1) {
    result1.16[I] = _SAT_16(_SX_16(argument3.16[I]) - _SX_16(argument2.16[I]))
  };
stage E1:
  PGR[registerM] = result1;
\end{lstlisting}}
\newpage
\subsubsection[{\small SBFSW: Subtract Saturated Word}]{SBFSW\hfill Subtract Saturated Word}
\label{instr:SBFSW}\index{sbfsw}
 
\paragraph{Encoding} Format \textsf{ALU\_WWRR1} (Section~\ref{sec:format-ALU-WWRR1}), \verb|exucode2 = 1001|\\

\bitfields{ 1/P, 2/11, 1/0, 4/1001, 6/registerW, 2/11, 3/101, 1/signextw, 6/registerY, 6/registerZ }


\paragraph{Syntax} \texttt{sbfsw}\emph{signextw} \emph{registerW} = \emph{registerZ}, \emph{registerY}

\paragraph{Description} The \emph{registerZ} is subtracted from the \emph{registerY} with 32-bit saturation. The result with the 32 upper bits cleared is stored into the \emph{registerW}.


\paragraph{Modifier(s)}
\noindent\begin{itemize}
\item signextw (cf. signextw in section \ref{par:modifier-signextw} on page \pageref{par:modifier-signextw})
\end{itemize}

\paragraph{Execution} {Reservation table \textsf{ALU\_TINY} (Section~\ref{sec:reservation-ALU-TINY})}
~
{\begin{lstlisting}
stage ID:
  new signextw = _ZX_1(signextw);
stage RR:
  new argument3 = GPR[registerY];
  new argument2 = GPR[registerZ];
  new result1 = _SAT_32(_SX_32(argument3) - _SX_32(argument2));
stage E1:
  GPR[registerW] = signextw ? _SX_32(result1) : _ZX_32(result1);
\end{lstlisting}}
\newpage

\paragraph{Encoding} Format \textsf{ALU\_WWRR1.W} (Section~\ref{sec:format-ALU-WWRR1.W}), \verb|exucode2 = 1001|\\

\bitfields{ 1/P, 2/00, 2/-- --, 27/upper27, 1/1, 2/11, 1/0, 4/1001, 6/registerW, 2/11, 3/101, 1/signextw, 1/0, 5/lower5, 6/registerZ }


\paragraph{Syntax} \texttt{sbfsw}\emph{signextw} \emph{registerW} = \emph{registerZ}, \emph{upper27\_\discretionary{}{}{}lower5}

\paragraph{Description} The \emph{registerZ} is subtracted from the \emph{upper27\_\discretionary{}{}{}lower5} with 32-bit saturation. The result with the 32 upper bits cleared is stored into the \emph{registerW}.


\paragraph{Modifier(s)}
\noindent\begin{itemize}
\item signextw (cf. signextw in section \ref{par:modifier-signextw} on page \pageref{par:modifier-signextw})
\end{itemize}

\paragraph{Execution} {Reservation table \textsf{ALU\_TINY.X} (Section~\ref{sec:reservation-ALU-TINY.X})}
~
{\begin{lstlisting}
stage ID:
  new signextw = _ZX_1(signextw);
stage RR:
  new argument3 = _WX_32(upper27_lower5);
  new argument2 = GPR[registerZ];
  new result1 = _SAT_32(_SX_32(argument3) - _SX_32(argument2));
stage E1:
  GPR[registerW] = signextw ? _SX_32(result1) : _ZX_32(result1);
\end{lstlisting}}
\newpage
\subsubsection[{\small SBFSWQ: Subtract Saturated Word Quadruple}]{SBFSWQ\hfill Subtract Saturated Word Quadruple}
\label{instr:SBFSWQ}\index{sbfswq}
 
\paragraph{Encoding} Format \textsf{ALU\_WQWRR1} (Section~\ref{sec:format-ALU-WQWRR1}), \verb|exucode2 = 1001|\\

\bitfields{ 1/P, 2/11, 1/1, 4/1001, 5/registerM, 1/--, 2/10, 3/101, 1/0, 5/registerO, 1/--, 5/registerP, 1/1 }


\paragraph{Syntax} \texttt{sbfswq} \emph{registerM} = \emph{registerP}, \emph{registerO}

\paragraph{Description} The \emph{registerO} and the \emph{registerP} are interpreted as four words packed into 128 bits. The \emph{registerP} words are subtracted with saturation to 32 bits from the \emph{registerO} words. The four resulting words are packed and stored into the \emph{registerM}.


\paragraph{Execution} {Reservation table \textsf{ALU\_LITE} (Section~\ref{sec:reservation-ALU-LITE})}
~
{\begin{lstlisting}
stage RR:
  new argument3 = PGR[registerO];
  new argument2 = PGR[registerP];
  for (I = 0; I < 4; I += 1) {
    result1.32[I] = _SAT_32(_SX_32(argument3.32[I]) - _SX_32(argument2.32[I]))
  };
stage E1:
  PGR[registerM] = result1;
\end{lstlisting}}
\newpage

\paragraph{Encoding} Format \textsf{ALU\_WQWRR1.M} (Section~\ref{sec:format-ALU-WQWRR1.M}), \verb|exucode2 = 1001|\\

\bitfields{ 1/P, 2/00, 2/-- --, 27/upper27, 1/1, 2/11, 1/1, 4/1001, 5/registerM, 1/--, 2/10, 3/101, 1/0, 1/splat32, 5/lower5, 5/registerP, 1/1 }


\paragraph{Syntax} \texttt{sbfswq} \emph{registerM} = \emph{registerP}, \emph{upper27\_\discretionary{}{}{}lower5}\emph{splat32}

\paragraph{Description} The \emph{upper27\_\discretionary{}{}{}lower5} and the \emph{registerP} are interpreted as four words packed into 128 bits. The \emph{registerP} words are subtracted with saturation to 32 bits from the \emph{upper27\_\discretionary{}{}{}lower5} words. The four resulting words are packed and stored into the \emph{registerM}.


\paragraph{Modifier(s)}
\noindent\begin{itemize}
\item splat32 (cf. splat32 in section \ref{par:modifier-splat32} on page \pageref{par:modifier-splat32})
\end{itemize}

\paragraph{Execution} {Reservation table \textsf{ALU\_LITE.X} (Section~\ref{sec:reservation-ALU-LITE.X})}
~
{\begin{lstlisting}
stage ID:
  new splat32 = _ZX_1(splat32);
stage RR:
  new argument3 = splat32 ? (_WX_32(upper27_lower5) << 32) | _ZX_32(_WX_32(upper27_lower5)) : _SX_32(_WX_32(upper27_lower5));
  new argument2 = PGR[registerP];
  for (I = 0; I < 4; I += 1) {
    result1.32[I] = _SAT_32(_SX_32(argument3.32[I]) - _SX_32(argument2.32[I]))
  };
stage E1:
  PGR[registerM] = result1;
\end{lstlisting}}
\newpage
\subsubsection[{\small SBFUSBX: Subtract Unsigned Saturated Byte Hexadecuple}]{SBFUSBX\hfill Subtract Unsigned Saturated Byte Hexadecuple}
\label{instr:SBFUSBX}\index{sbfusbx}
 
\paragraph{Encoding} Format \textsf{ALU\_BXWRR1} (Section~\ref{sec:format-ALU-BXWRR1}), \verb|exucode2 = 1011|\\

\bitfields{ 1/P, 2/11, 1/1, 4/1011, 5/registerM, 1/--, 2/10, 3/101, 1/1, 5/registerO, 1/--, 5/registerP, 1/1 }


\paragraph{Syntax} \texttt{sbfusbx} \emph{registerM} = \emph{registerP}, \emph{registerO}

\paragraph{Description} The \emph{registerO} and the \emph{registerP} are interpreted as sixteen bytes packed into 128 bits. The \emph{registerP} bytes are subtracted with unsigned saturation to 8 bits from the \emph{registerO} bytes. The sixteen resulting bytes are packed and stored into the \emph{registerM}.


\paragraph{Execution} {Reservation table \textsf{ALU\_LITE} (Section~\ref{sec:reservation-ALU-LITE})}
~
{\begin{lstlisting}
stage RR:
  new argument3 = PGR[registerO];
  new argument2 = PGR[registerP];
  for (I = 0; I < 16; I += 1) {
    result1.8[I] = _SATU_8(_ZX_8(argument3.8[I]) - _ZX_8(argument2.8[I]))
  };
stage E1:
  PGR[registerM] = result1;
\end{lstlisting}}
\newpage

\paragraph{Encoding} Format \textsf{ALU\_BXWRR1.M} (Section~\ref{sec:format-ALU-BXWRR1.M}), \verb|exucode2 = 1011|\\

\bitfields{ 1/P, 2/00, 2/-- --, 27/upper27, 1/1, 2/11, 1/1, 4/1011, 5/registerM, 1/--, 2/10, 3/101, 1/1, 1/splat32, 5/lower5, 5/registerP, 1/1 }


\paragraph{Syntax} \texttt{sbfusbx} \emph{registerM} = \emph{registerP}, \emph{upper27\_\discretionary{}{}{}lower5}\emph{splat32}

\paragraph{Description} The \emph{upper27\_\discretionary{}{}{}lower5} and the \emph{registerP} are interpreted as sixteen bytes packed into 128 bits. The \emph{registerP} bytes are subtracted with unsigned saturation to 8 bits from the \emph{upper27\_\discretionary{}{}{}lower5} bytes. The sixteen resulting bytes are packed and stored into the \emph{registerM}.


\paragraph{Modifier(s)}
\noindent\begin{itemize}
\item splat32 (cf. splat32 in section \ref{par:modifier-splat32} on page \pageref{par:modifier-splat32})
\end{itemize}

\paragraph{Execution} {Reservation table \textsf{ALU\_LITE.X} (Section~\ref{sec:reservation-ALU-LITE.X})}
~
{\begin{lstlisting}
stage ID:
  new splat32 = _ZX_1(splat32);
stage RR:
  new argument3 = splat32 ? (_WX_32(upper27_lower5) << 32) | _ZX_32(_WX_32(upper27_lower5)) : _SX_32(_WX_32(upper27_lower5));
  new argument2 = PGR[registerP];
  for (I = 0; I < 16; I += 1) {
    result1.8[I] = _SATU_8(_ZX_8(argument3.8[I]) - _ZX_8(argument2.8[I]))
  };
stage E1:
  PGR[registerM] = result1;
\end{lstlisting}}
\newpage
\subsubsection[{\small SBFUSD: Subtract Unsigned Saturated Double Word}]{SBFUSD\hfill Subtract Unsigned Saturated Double Word}
\label{instr:SBFUSD}\index{sbfusd}
 
\paragraph{Encoding} Format \textsf{ALU\_DWRR1} (Section~\ref{sec:format-ALU-DWRR1}), \verb|exucode2 = 1011|\\

\bitfields{ 1/P, 2/11, 1/0, 4/1011, 6/registerW, 2/10, 3/101, 1/0, 6/registerY, 6/registerZ }


\paragraph{Syntax} \texttt{sbfusd} \emph{registerW} = \emph{registerZ}, \emph{registerY}

\paragraph{Description} The \emph{registerZ} is subtracted from the \emph{registerY} with 64-bit unsigned saturation.


\paragraph{Execution} {Reservation table \textsf{ALU\_TINY} (Section~\ref{sec:reservation-ALU-TINY})}
~
{\begin{lstlisting}
stage RR:
  new argument3 = GPR[registerY];
  new argument2 = GPR[registerZ];
  new result1 = _SATU_64(_ZX_64(argument3) - _ZX_64(argument2));
stage E1:
  GPR[registerW] = result1;
\end{lstlisting}}
\newpage

\paragraph{Encoding} Format \textsf{ALU\_DWRR1.M} (Section~\ref{sec:format-ALU-DWRR1.M}), \verb|exucode2 = 1011|\\

\bitfields{ 1/P, 2/00, 2/-- --, 27/upper27, 1/1, 2/11, 1/0, 4/1011, 6/registerW, 2/10, 3/101, 1/0, 1/splat32, 5/lower5, 6/registerZ }


\paragraph{Syntax} \texttt{sbfusd} \emph{registerW} = \emph{registerZ}, \emph{upper27\_\discretionary{}{}{}lower5}\emph{splat32}

\paragraph{Description} The \emph{registerZ} is subtracted from the \emph{upper27\_\discretionary{}{}{}lower5} with 64-bit unsigned saturation.


\paragraph{Modifier(s)}
\noindent\begin{itemize}
\item splat32 (cf. splat32 in section \ref{par:modifier-splat32} on page \pageref{par:modifier-splat32})
\end{itemize}

\paragraph{Execution} {Reservation table \textsf{ALU\_TINY.X} (Section~\ref{sec:reservation-ALU-TINY.X})}
~
{\begin{lstlisting}
stage ID:
  new splat32 = _ZX_1(splat32);
stage RR:
  new argument3 = splat32 ? (_WX_32(upper27_lower5) << 32) | _ZX_32(_WX_32(upper27_lower5)) : _SX_32(_WX_32(upper27_lower5));
  new argument2 = GPR[registerZ];
  new result1 = _SATU_64(_ZX_64(argument3) - _ZX_64(argument2));
stage E1:
  GPR[registerW] = result1;
\end{lstlisting}}
\newpage
\subsubsection[{\small SBFUSDP: Subtract Unsigned Saturated Double Word Pair}]{SBFUSDP\hfill Subtract Unsigned Saturated Double Word Pair}
\label{instr:SBFUSDP}\index{sbfusdp}
 
\paragraph{Encoding} Format \textsf{ALU\_DPWRR1} (Section~\ref{sec:format-ALU-DPWRR1}), \verb|exucode2 = 1011|\\

\bitfields{ 1/P, 2/11, 1/1, 4/1011, 5/registerM, 1/--, 2/10, 3/101, 1/0, 5/registerO, 1/--, 5/registerP, 1/0 }


\paragraph{Syntax} \texttt{sbfusdp} \emph{registerM} = \emph{registerP}, \emph{registerO}

\paragraph{Description} The \emph{registerO} and the \emph{registerP} are interpreted as two double words packed into 128 bits. The \emph{registerP} double words are subtracted with unsigned saturation to 64 bits from the \emph{registerO} double words. The two resulting double words are packed and stored into the \emph{registerM}.


\paragraph{Execution} {Reservation table \textsf{ALU\_LITE} (Section~\ref{sec:reservation-ALU-LITE})}
~
{\begin{lstlisting}
stage RR:
  new argument3 = PGR[registerO];
  new argument2 = PGR[registerP];
  for (I = 0; I < 2; I += 1) {
    result1.64[I] = _SATU_64(_ZX_64(argument3.64[I]) - _ZX_64(argument2.64[I]))
  };
stage E1:
  PGR[registerM] = result1;
\end{lstlisting}}
\newpage

\paragraph{Encoding} Format \textsf{ALU\_DPWRR1.M} (Section~\ref{sec:format-ALU-DPWRR1.M}), \verb|exucode2 = 1011|\\

\bitfields{ 1/P, 2/00, 2/-- --, 27/upper27, 1/1, 2/11, 1/1, 4/1011, 5/registerM, 1/--, 2/10, 3/101, 1/0, 1/splat32, 5/lower5, 5/registerP, 1/0 }


\paragraph{Syntax} \texttt{sbfusdp} \emph{registerM} = \emph{registerP}, \emph{upper27\_\discretionary{}{}{}lower5}\emph{splat32}

\paragraph{Description} The \emph{upper27\_\discretionary{}{}{}lower5} and the \emph{registerP} are interpreted as two double words packed into 128 bits. The \emph{registerP} double words are subtracted with unsigned saturation to 64 bits from the \emph{upper27\_\discretionary{}{}{}lower5} double words. The two resulting double words are packed and stored into the \emph{registerM}.


\paragraph{Modifier(s)}
\noindent\begin{itemize}
\item splat32 (cf. splat32 in section \ref{par:modifier-splat32} on page \pageref{par:modifier-splat32})
\end{itemize}

\paragraph{Execution} {Reservation table \textsf{ALU\_LITE.X} (Section~\ref{sec:reservation-ALU-LITE.X})}
~
{\begin{lstlisting}
stage ID:
  new splat32 = _ZX_1(splat32);
stage RR:
  new argument3 = splat32 ? (_WX_32(upper27_lower5) << 32) | _ZX_32(_WX_32(upper27_lower5)) : _SX_32(_WX_32(upper27_lower5));
  new argument2 = PGR[registerP];
  for (I = 0; I < 2; I += 1) {
    result1.64[I] = _SATU_64(_ZX_64(argument3.64[I]) - _ZX_64(argument2.64[I]))
  };
stage E1:
  PGR[registerM] = result1;
\end{lstlisting}}
\newpage
\subsubsection[{\small SBFUSHO: Subtract Unsigned Saturated Half Word Octuple}]{SBFUSHO\hfill Subtract Unsigned Saturated Half Word Octuple}
\label{instr:SBFUSHO}\index{sbfusho}
 
\paragraph{Encoding} Format \textsf{ALU\_HOWRR1} (Section~\ref{sec:format-ALU-HOWRR1}), \verb|exucode2 = 1011|\\

\bitfields{ 1/P, 2/11, 1/1, 4/1011, 5/registerM, 1/--, 2/10, 3/101, 1/1, 5/registerO, 1/--, 5/registerP, 1/0 }


\paragraph{Syntax} \texttt{sbfusho} \emph{registerM} = \emph{registerP}, \emph{registerO}

\paragraph{Description} The \emph{registerO} and the \emph{registerP} are interpreted as eight half words packed into 128 bits. The \emph{registerP} half words are subtracted with unsigned saturation to 16 bits from the \emph{registerO} half words. The eight resulting half words are packed and stored into the \emph{registerM}.


\paragraph{Execution} {Reservation table \textsf{ALU\_LITE} (Section~\ref{sec:reservation-ALU-LITE})}
~
{\begin{lstlisting}
stage RR:
  new argument3 = PGR[registerO];
  new argument2 = PGR[registerP];
  for (I = 0; I < 8; I += 1) {
    result1.16[I] = _SATU_16(_ZX_16(argument3.16[I]) - _ZX_16(argument2.16[I]))
  };
stage E1:
  PGR[registerM] = result1;
\end{lstlisting}}
\newpage

\paragraph{Encoding} Format \textsf{ALU\_HOWRR1.M} (Section~\ref{sec:format-ALU-HOWRR1.M}), \verb|exucode2 = 1011|\\

\bitfields{ 1/P, 2/00, 2/-- --, 27/upper27, 1/1, 2/11, 1/1, 4/1011, 5/registerM, 1/--, 2/10, 3/101, 1/1, 1/splat32, 5/lower5, 5/registerP, 1/0 }


\paragraph{Syntax} \texttt{sbfusho} \emph{registerM} = \emph{registerP}, \emph{upper27\_\discretionary{}{}{}lower5}\emph{splat32}

\paragraph{Description} The \emph{upper27\_\discretionary{}{}{}lower5} and the \emph{registerP} are interpreted as eight half words packed into 128 bits. The \emph{registerP} half words are subtracted with unsigned saturation to 16 bits from the \emph{upper27\_\discretionary{}{}{}lower5} half words. The eight resulting half words are packed and stored into the \emph{registerM}.


\paragraph{Modifier(s)}
\noindent\begin{itemize}
\item splat32 (cf. splat32 in section \ref{par:modifier-splat32} on page \pageref{par:modifier-splat32})
\end{itemize}

\paragraph{Execution} {Reservation table \textsf{ALU\_LITE.X} (Section~\ref{sec:reservation-ALU-LITE.X})}
~
{\begin{lstlisting}
stage ID:
  new splat32 = _ZX_1(splat32);
stage RR:
  new argument3 = splat32 ? (_WX_32(upper27_lower5) << 32) | _ZX_32(_WX_32(upper27_lower5)) : _SX_32(_WX_32(upper27_lower5));
  new argument2 = PGR[registerP];
  for (I = 0; I < 8; I += 1) {
    result1.16[I] = _SATU_16(_ZX_16(argument3.16[I]) - _ZX_16(argument2.16[I]))
  };
stage E1:
  PGR[registerM] = result1;
\end{lstlisting}}
\newpage
\subsubsection[{\small SBFUSW: Subtract Unsigned Saturated Word}]{SBFUSW\hfill Subtract Unsigned Saturated Word}
\label{instr:SBFUSW}\index{sbfusw}
 
\paragraph{Encoding} Format \textsf{ALU\_WWRR1} (Section~\ref{sec:format-ALU-WWRR1}), \verb|exucode2 = 1011|\\

\bitfields{ 1/P, 2/11, 1/0, 4/1011, 6/registerW, 2/11, 3/101, 1/signextw, 6/registerY, 6/registerZ }


\paragraph{Syntax} \texttt{sbfusw}\emph{signextw} \emph{registerW} = \emph{registerZ}, \emph{registerY}

\paragraph{Description} The \emph{registerZ} is subtracted from the \emph{registerY} with 32-bit unsigned saturation. The result with the 32 upper bits cleared is stored into the \emph{registerW}.


\paragraph{Modifier(s)}
\noindent\begin{itemize}
\item signextw (cf. signextw in section \ref{par:modifier-signextw} on page \pageref{par:modifier-signextw})
\end{itemize}

\paragraph{Execution} {Reservation table \textsf{ALU\_TINY} (Section~\ref{sec:reservation-ALU-TINY})}
~
{\begin{lstlisting}
stage ID:
  new signextw = _ZX_1(signextw);
stage RR:
  new argument3 = GPR[registerY];
  new argument2 = GPR[registerZ];
  new result1 = _SATU_32(_ZX_32(argument3) - _ZX_32(argument2));
stage E1:
  GPR[registerW] = signextw ? _SX_32(result1) : _ZX_32(result1);
\end{lstlisting}}
\newpage

\paragraph{Encoding} Format \textsf{ALU\_WWRR1.W} (Section~\ref{sec:format-ALU-WWRR1.W}), \verb|exucode2 = 1011|\\

\bitfields{ 1/P, 2/00, 2/-- --, 27/upper27, 1/1, 2/11, 1/0, 4/1011, 6/registerW, 2/11, 3/101, 1/signextw, 1/0, 5/lower5, 6/registerZ }


\paragraph{Syntax} \texttt{sbfusw}\emph{signextw} \emph{registerW} = \emph{registerZ}, \emph{upper27\_\discretionary{}{}{}lower5}

\paragraph{Description} The \emph{registerZ} is subtracted from the \emph{upper27\_\discretionary{}{}{}lower5} with 32-bit unsigned saturation. The result with the 32 upper bits cleared is stored into the \emph{registerW}.


\paragraph{Modifier(s)}
\noindent\begin{itemize}
\item signextw (cf. signextw in section \ref{par:modifier-signextw} on page \pageref{par:modifier-signextw})
\end{itemize}

\paragraph{Execution} {Reservation table \textsf{ALU\_TINY.X} (Section~\ref{sec:reservation-ALU-TINY.X})}
~
{\begin{lstlisting}
stage ID:
  new signextw = _ZX_1(signextw);
stage RR:
  new argument3 = _WX_32(upper27_lower5);
  new argument2 = GPR[registerZ];
  new result1 = _SATU_32(_ZX_32(argument3) - _ZX_32(argument2));
stage E1:
  GPR[registerW] = signextw ? _SX_32(result1) : _ZX_32(result1);
\end{lstlisting}}
\newpage
\subsubsection[{\small SBFUSWQ: Subtract Unsigned Saturated Word Quadruple}]{SBFUSWQ\hfill Subtract Unsigned Saturated Word Quadruple}
\label{instr:SBFUSWQ}\index{sbfuswq}
 
\paragraph{Encoding} Format \textsf{ALU\_WQWRR1} (Section~\ref{sec:format-ALU-WQWRR1}), \verb|exucode2 = 1011|\\

\bitfields{ 1/P, 2/11, 1/1, 4/1011, 5/registerM, 1/--, 2/10, 3/101, 1/0, 5/registerO, 1/--, 5/registerP, 1/1 }


\paragraph{Syntax} \texttt{sbfuswq} \emph{registerM} = \emph{registerP}, \emph{registerO}

\paragraph{Description} The \emph{registerO} and the \emph{registerP} are interpreted as four words packed into 128 bits. The \emph{registerP} words are subtracted with unsigned saturation to 32 bits from the \emph{registerO} words. The four resulting words are packed and stored into the \emph{registerM}.


\paragraph{Execution} {Reservation table \textsf{ALU\_LITE} (Section~\ref{sec:reservation-ALU-LITE})}
~
{\begin{lstlisting}
stage RR:
  new argument3 = PGR[registerO];
  new argument2 = PGR[registerP];
  for (I = 0; I < 4; I += 1) {
    result1.32[I] = _SATU_32(_ZX_32(argument3.32[I]) - _ZX_32(argument2.32[I]))
  };
stage E1:
  PGR[registerM] = result1;
\end{lstlisting}}
\newpage

\paragraph{Encoding} Format \textsf{ALU\_WQWRR1.M} (Section~\ref{sec:format-ALU-WQWRR1.M}), \verb|exucode2 = 1011|\\

\bitfields{ 1/P, 2/00, 2/-- --, 27/upper27, 1/1, 2/11, 1/1, 4/1011, 5/registerM, 1/--, 2/10, 3/101, 1/0, 1/splat32, 5/lower5, 5/registerP, 1/1 }


\paragraph{Syntax} \texttt{sbfuswq} \emph{registerM} = \emph{registerP}, \emph{upper27\_\discretionary{}{}{}lower5}\emph{splat32}

\paragraph{Description} The \emph{upper27\_\discretionary{}{}{}lower5} and the \emph{registerP} are interpreted as four words packed into 128 bits. The \emph{registerP} words are subtracted with unsigned saturation to 32 bits from the \emph{upper27\_\discretionary{}{}{}lower5} words. The four resulting words are packed and stored into the \emph{registerM}.


\paragraph{Modifier(s)}
\noindent\begin{itemize}
\item splat32 (cf. splat32 in section \ref{par:modifier-splat32} on page \pageref{par:modifier-splat32})
\end{itemize}

\paragraph{Execution} {Reservation table \textsf{ALU\_LITE.X} (Section~\ref{sec:reservation-ALU-LITE.X})}
~
{\begin{lstlisting}
stage ID:
  new splat32 = _ZX_1(splat32);
stage RR:
  new argument3 = splat32 ? (_WX_32(upper27_lower5) << 32) | _ZX_32(_WX_32(upper27_lower5)) : _SX_32(_WX_32(upper27_lower5));
  new argument2 = PGR[registerP];
  for (I = 0; I < 4; I += 1) {
    result1.32[I] = _SATU_32(_ZX_32(argument3.32[I]) - _ZX_32(argument2.32[I]))
  };
stage E1:
  PGR[registerM] = result1;
\end{lstlisting}}
\newpage
\subsubsection[{\small SBFW: Subtract From Word}]{SBFW\hfill Subtract From Word}
\label{instr:SBFW}\index{sbfw}
 
\paragraph{Encoding} Format \textsf{ALU\_WWRR0} (Section~\ref{sec:format-ALU-WWRR0}), \verb|exucode2 = 0011|\\

\bitfields{ 1/P, 2/11, 1/0, 4/0011, 6/registerW, 2/11, 3/000, 1/signextw, 6/registerY, 6/registerZ }


\paragraph{Syntax} \texttt{sbfw}\emph{signextw} \emph{registerW} = \emph{registerZ}, \emph{registerY}

\paragraph{Description} The \emph{registerZ} is subtracted from the \emph{registerY}. The result with the 32 upper bits cleared is stored into the \emph{registerW}.


\paragraph{Modifier(s)}
\noindent\begin{itemize}
\item signextw (cf. signextw in section \ref{par:modifier-signextw} on page \pageref{par:modifier-signextw})
\end{itemize}

\paragraph{Execution} {Reservation table \textsf{ALU\_TINY} (Section~\ref{sec:reservation-ALU-TINY})}
~
{\begin{lstlisting}
stage ID:
  new signextw = _ZX_1(signextw);
stage RR:
  new argument3 = GPR[registerY];
  new argument2 = GPR[registerZ];
  new result1 = argument3 - argument2;
stage E1:
  GPR[registerW] = signextw ? _SX_32(result1) : _ZX_32(result1);
\end{lstlisting}}
\newpage

\paragraph{Encoding} Format \textsf{ALU\_WWRR0.W} (Section~\ref{sec:format-ALU-WWRR0.W}), \verb|exucode2 = 0011|\\

\bitfields{ 1/P, 2/00, 2/-- --, 27/upper27, 1/1, 2/11, 1/0, 4/0011, 6/registerW, 2/11, 3/000, 1/signextw, 1/0, 5/lower5, 6/registerZ }


\paragraph{Syntax} \texttt{sbfw}\emph{signextw} \emph{registerW} = \emph{registerZ}, \emph{upper27\_\discretionary{}{}{}lower5}

\paragraph{Description} The \emph{registerZ} is subtracted from the \emph{upper27\_\discretionary{}{}{}lower5}. The result with the 32 upper bits cleared is stored into the \emph{registerW}.


\paragraph{Modifier(s)}
\noindent\begin{itemize}
\item signextw (cf. signextw in section \ref{par:modifier-signextw} on page \pageref{par:modifier-signextw})
\end{itemize}

\paragraph{Execution} {Reservation table \textsf{ALU\_TINY.X} (Section~\ref{sec:reservation-ALU-TINY.X})}
~
{\begin{lstlisting}
stage ID:
  new signextw = _ZX_1(signextw);
stage RR:
  new argument3 = _WX_32(upper27_lower5);
  new argument2 = GPR[registerZ];
  new result1 = argument3 - argument2;
stage E1:
  GPR[registerW] = signextw ? _SX_32(result1) : _ZX_32(result1);
\end{lstlisting}}
\newpage
\subsubsection[{\small SBFWQ: Subtract From Word Quadruple}]{SBFWQ\hfill Subtract From Word Quadruple}
\label{instr:SBFWQ}\index{sbfwq}
 
\paragraph{Encoding} Format \textsf{ALU\_WQWRR0} (Section~\ref{sec:format-ALU-WQWRR0}), \verb|exucode2 = 0011|\\

\bitfields{ 1/P, 2/11, 1/1, 4/0011, 5/registerM, 1/--, 2/10, 3/000, 1/0, 5/registerO, 1/--, 5/registerP, 1/1 }


\paragraph{Syntax} \texttt{sbfwq} \emph{registerM} = \emph{registerP}, \emph{registerO}

\paragraph{Description} The \emph{registerO} and the \emph{registerP} are interpreted as four words packed into 128 bits. The \emph{registerP} words are subtracted from the \emph{registerO} words. The four resulting words are packed and stored into the \emph{registerM}.


\paragraph{Execution} {Reservation table \textsf{ALU\_LITE} (Section~\ref{sec:reservation-ALU-LITE})}
~
{\begin{lstlisting}
stage RR:
  new argument3 = PGR[registerO];
  new argument2 = PGR[registerP];
  for (I = 0; I < 4; I += 1) {
    result1.32[I] = argument3.32[I] - argument2.32[I]
  };
stage E1:
  PGR[registerM] = result1;
\end{lstlisting}}
\newpage

\paragraph{Encoding} Format \textsf{ALU\_WQWRR0.M} (Section~\ref{sec:format-ALU-WQWRR0.M}), \verb|exucode2 = 0011|\\

\bitfields{ 1/P, 2/00, 2/-- --, 27/upper27, 1/1, 2/11, 1/1, 4/0011, 5/registerM, 1/--, 2/10, 3/000, 1/0, 1/splat32, 5/lower5, 5/registerP, 1/1 }


\paragraph{Syntax} \texttt{sbfwq} \emph{registerM} = \emph{registerP}, \emph{upper27\_\discretionary{}{}{}lower5}\emph{splat32}

\paragraph{Description} The \emph{upper27\_\discretionary{}{}{}lower5} and the \emph{registerP} are interpreted as four words packed into 128 bits. The \emph{registerP} words are subtracted from the \emph{upper27\_\discretionary{}{}{}lower5} words. The four resulting words are packed and stored into the \emph{registerM}.


\paragraph{Modifier(s)}
\noindent\begin{itemize}
\item splat32 (cf. splat32 in section \ref{par:modifier-splat32} on page \pageref{par:modifier-splat32})
\end{itemize}

\paragraph{Execution} {Reservation table \textsf{ALU\_LITE.X} (Section~\ref{sec:reservation-ALU-LITE.X})}
~
{\begin{lstlisting}
stage ID:
  new splat32 = _ZX_1(splat32);
stage RR:
  new argument3 = splat32 ? (_WX_32(upper27_lower5) << 32) | _ZX_32(_WX_32(upper27_lower5)) : _SX_32(_WX_32(upper27_lower5));
  new argument2 = PGR[registerP];
  for (I = 0; I < 4; I += 1) {
    result1.32[I] = argument3.32[I] - argument2.32[I]
  };
stage E1:
  PGR[registerM] = result1;
\end{lstlisting}}
\newpage
\subsubsection[{\small SBMM8D: Swapped Bit Matrix Multiplication 8\texorpdfstring{$\times$}{x}8 Double Word}]{SBMM8D\hfill Swapped Bit Matrix Multiplication 8$\times$8 Double Word}
\label{instr:SBMM8D}\index{sbmm8d}
 
\paragraph{Encoding} Format \textsf{ALU\_DBMWRR} (Section~\ref{sec:format-ALU-DBMWRR}), \verb|exucode2 = 1110|\\

\bitfields{ 1/P, 2/11, 1/0, 4/1110, 6/registerW, 2/10, 3/001, 1/0, 6/registerY, 6/registerZ }


\paragraph{Syntax} \texttt{sbmm8d} \emph{registerW} = \emph{registerZ}, \emph{registerY}

\paragraph{Description} The \emph{registerY} interpreted as a 8$\times$8 bit matrix is multiplied by the \emph{registerZ} interpreted as a 8$\times$8 bit matrix. The resulting 8$\times$8 bit matrix is stored into the \emph{registerW}.


\paragraph{Execution} {Reservation table \textsf{ALU\_TINY} (Section~\ref{sec:reservation-ALU-TINY})}
~
{\begin{lstlisting}
stage RR:
  new argument3 = GPR[registerY];
  new argument2 = GPR[registerZ];
  new result1 = _BMM_8(_ZX_64(argument3), _ZX_64(argument2));
stage E1:
  GPR[registerW] = result1;
\end{lstlisting}}
\newpage

\paragraph{Encoding} Format \textsf{ALU\_DBMWRR.M} (Section~\ref{sec:format-ALU-DBMWRR.M}), \verb|exucode2 = 1110|\\

\bitfields{ 1/P, 2/00, 2/-- --, 27/upper27, 1/1, 2/11, 1/0, 4/1110, 6/registerW, 2/10, 3/001, 1/0, 1/splat32, 5/lower5, 6/registerZ }


\paragraph{Syntax} \texttt{sbmm8d} \emph{registerW} = \emph{registerZ}, \emph{upper27\_\discretionary{}{}{}lower5}\emph{splat32}

\paragraph{Description} The \emph{upper27\_\discretionary{}{}{}lower5} interpreted as a 8$\times$8 bit matrix is multiplied by the \emph{registerZ} interpreted as a 8$\times$8 bit matrix. The resulting 8$\times$8 bit matrix is stored into the \emph{registerW}.


\paragraph{Modifier(s)}
\noindent\begin{itemize}
\item splat32 (cf. splat32 in section \ref{par:modifier-splat32} on page \pageref{par:modifier-splat32})
\end{itemize}

\paragraph{Execution} {Reservation table \textsf{ALU\_TINY.X} (Section~\ref{sec:reservation-ALU-TINY.X})}
~
{\begin{lstlisting}
stage ID:
  new splat32 = _ZX_1(splat32);
stage RR:
  new argument3 = splat32 ? (_WX_32(upper27_lower5) << 32) | _ZX_32(_WX_32(upper27_lower5)) : _SX_32(_WX_32(upper27_lower5));
  new argument2 = GPR[registerZ];
  new result1 = _BMM_8(_ZX_64(argument3), _ZX_64(argument2));
stage E1:
  GPR[registerW] = result1;
\end{lstlisting}}
\newpage

\paragraph{Encoding} Format \textsf{ALU\_DBMWRI} (Section~\ref{sec:format-ALU-DBMWRI}), \verb|exucode2 = 1110|\\

\bitfields{ 1/P, 2/11, 1/0, 4/1110, 6/registerW, 2/01, 10/signed10, 6/registerZ }


\paragraph{Syntax} \texttt{sbmm8d} \emph{registerW} = \emph{registerZ}, \emph{signed10}

\paragraph{Description} The \emph{signed10} interpreted as a 8$\times$8 bit matrix is multiplied by the \emph{registerZ} interpreted as a 8$\times$8 bit matrix. The resulting 8$\times$8 bit matrix is stored into the \emph{registerW}.


\paragraph{Execution} {Reservation table \textsf{ALU\_TINY} (Section~\ref{sec:reservation-ALU-TINY})}
~
{\begin{lstlisting}
stage RR:
  new argument3 = _SX_10(signed10);
  new argument2 = GPR[registerZ];
  new result1 = _BMM_8(_ZX_64(argument3), _ZX_64(argument2));
stage E1:
  GPR[registerW] = result1;
\end{lstlisting}}
\newpage

\paragraph{Encoding} Format \textsf{ALU\_DBMWRI.X} (Section~\ref{sec:format-ALU-DBMWRI.X}), \verb|exucode2 = 1110|\\

\bitfields{ 1/P, 2/00, 2/-- --, 27/upper27, 1/1, 2/11, 1/0, 4/1110, 6/registerW, 2/01, 10/lower10, 6/registerZ }


\paragraph{Syntax} \texttt{sbmm8d} \emph{registerW} = \emph{registerZ}, \emph{upper27\_\discretionary{}{}{}lower10}

\paragraph{Description} The \emph{upper27\_\discretionary{}{}{}lower10} interpreted as a 8$\times$8 bit matrix is multiplied by the \emph{registerZ} interpreted as a 8$\times$8 bit matrix. The resulting 8$\times$8 bit matrix is stored into the \emph{registerW}.


\paragraph{Execution} {Reservation table \textsf{ALU\_TINY.X} (Section~\ref{sec:reservation-ALU-TINY.X})}
~
{\begin{lstlisting}
stage RR:
  new argument3 = _SX_37(upper27_lower10);
  new argument2 = GPR[registerZ];
  new result1 = _BMM_8(_ZX_64(argument3), _ZX_64(argument2));
stage E1:
  GPR[registerW] = result1;
\end{lstlisting}}
\newpage

\paragraph{Encoding} Format \textsf{ALU\_DBMWRI.Y} (Section~\ref{sec:format-ALU-DBMWRI.Y}), \verb|exucode2 = 1110|\\

\bitfields{ 1/P, 2/00, 2/-- --, 27/extend27, 1/1, 2/00, 2/-- --, 27/upper27, 1/1, 2/11, 1/0, 4/1110, 6/registerW, 2/01, 10/lower10, 6/registerZ }


\paragraph{Syntax} \texttt{sbmm8d} \emph{registerW} = \emph{registerZ}, \emph{extend27\_\discretionary{}{}{}upper27\_\discretionary{}{}{}lower10}

\paragraph{Description} The \emph{extend27\_\discretionary{}{}{}upper27\_\discretionary{}{}{}lower10} interpreted as a 8$\times$8 bit matrix is multiplied by the \emph{registerZ} interpreted as a 8$\times$8 bit matrix. The resulting 8$\times$8 bit matrix is stored into the \emph{registerW}.


\paragraph{Execution} {Reservation table \textsf{ALU\_TINY.Y} (Section~\ref{sec:reservation-ALU-TINY.Y})}
~
{\begin{lstlisting}
stage RR:
  new argument3 = _WX_64(extend27_upper27_lower10);
  new argument2 = GPR[registerZ];
  new result1 = _BMM_8(_ZX_64(argument3), _ZX_64(argument2));
stage E1:
  GPR[registerW] = result1;
\end{lstlisting}}
\newpage
\subsubsection[{\small SBMM8DP: Swapped Bit Matrix Multiplication 8\texorpdfstring{$\times$}{x}8 Double Word Pair}]{SBMM8DP\hfill Swapped Bit Matrix Multiplication 8$\times$8 Double Word Pair}
\label{instr:SBMM8DP}\index{sbmm8dp}
 
\paragraph{Encoding} Format \textsf{ALU\_DPSBMM} (Section~\ref{sec:format-ALU-DPSBMM}), \verb|exucode2 = 1110|\\

\bitfields{ 1/P, 2/11, 1/1, 4/1110, 5/registerM, 1/--, 2/10, 3/001, 1/0, 5/registerO, 1/--, 5/registerP, 1/0 }


\paragraph{Syntax} \texttt{sbmm8dp} \emph{registerM} = \emph{registerP}, \emph{registerO}

\paragraph{Description} The \emph{registerO} interpreted as two 8$\times$8 bit matrices are multiplied by the \emph{registerP} interpreted as two 8$\times$8 bit matrices. The resulting 8$\times$8 bit matrix pair is stored into the \emph{registerM}.


\paragraph{Execution} {Reservation table \textsf{ALU\_LITE} (Section~\ref{sec:reservation-ALU-LITE})}
~
{\begin{lstlisting}
stage RR:
  new argument3 = PGR[registerO];
  new argument2 = PGR[registerP];
  for (I = 0; I < 2; I += 1) {
    result1.64[I] = _BMM_8(_ZX_64(argument3.64[I]), _ZX_64(argument2.64[I]))
  };
stage E1:
  PGR[registerM] = result1;
\end{lstlisting}}
\newpage

\paragraph{Encoding} Format \textsf{ALU\_DPSBMM.M} (Section~\ref{sec:format-ALU-DPSBMM.M}), \verb|exucode2 = 1110|\\

\bitfields{ 1/P, 2/00, 2/-- --, 27/upper27, 1/1, 2/11, 1/1, 4/1110, 5/registerM, 1/--, 2/10, 3/001, 1/0, 1/splat32, 5/lower5, 5/registerP, 1/0 }


\paragraph{Syntax} \texttt{sbmm8dp} \emph{registerM} = \emph{registerP}, \emph{upper27\_\discretionary{}{}{}lower5}\emph{splat32}

\paragraph{Description} The \emph{upper27\_\discretionary{}{}{}lower5} interpreted as two 8$\times$8 bit matrices are multiplied by the \emph{registerP} interpreted as two 8$\times$8 bit matrices. The resulting 8$\times$8 bit matrix pair is stored into the \emph{registerM}.


\paragraph{Modifier(s)}
\noindent\begin{itemize}
\item splat32 (cf. splat32 in section \ref{par:modifier-splat32} on page \pageref{par:modifier-splat32})
\end{itemize}

\paragraph{Execution} {Reservation table \textsf{ALU\_LITE.X} (Section~\ref{sec:reservation-ALU-LITE.X})}
~
{\begin{lstlisting}
stage ID:
  new splat32 = _ZX_1(splat32);
stage RR:
  new argument3 = splat32 ? (_WX_32(upper27_lower5) << 32) | _ZX_32(_WX_32(upper27_lower5)) : _SX_32(_WX_32(upper27_lower5));
  new argument2 = PGR[registerP];
  for (I = 0; I < 2; I += 1) {
    result1.64[I] = _BMM_8(_ZX_64(argument3.64[I]), _ZX_64(argument2.64[I]))
  };
stage E1:
  PGR[registerM] = result1;
\end{lstlisting}}
\newpage
\subsubsection[{\small SBMM8EORD: Swapped Bit Matrix Multiplication 8\texorpdfstring{$\times$}{x}8 Exclusive OR Double Word}]{SBMM8EORD\hfill Swapped Bit Matrix Multiplication 8$\times$8 Exclusive OR Double Word}
\label{instr:SBMM8EORD}\index{sbmm8eord}
 
\paragraph{Encoding} Format \textsf{ALU\_DBMWRRR} (Section~\ref{sec:format-ALU-DBMWRRR}), \verb|exucode2 = 1100|\\

\bitfields{ 1/P, 2/11, 1/0, 4/1100, 6/registerW, 2/10, 3/001, 1/0, 6/registerY, 6/registerZ }


\paragraph{Syntax} \texttt{sbmm8eord} \emph{registerW} = \emph{registerZ}, \emph{registerY}

\paragraph{Description} The \emph{registerY} interpreted as a 8$\times$8 bit matrix is multiplied by the \emph{registerZ} interpreted as a 8$\times$8 bit matrix. The resulting $\times$8 bit matrix is XORed with the value of the \emph{registerW} and stored into the \emph{registerW}.


\paragraph{Execution} {Reservation table \textsf{ALU\_TINY} (Section~\ref{sec:reservation-ALU-TINY})}
~
{\begin{lstlisting}
stage RR:
  new argument3 = GPR[registerY];
  new argument2 = GPR[registerZ];
  new argument1 = GPR[registerW];
  new result1 = _BMM_8(_ZX_64(argument3), _ZX_64(argument2)) ^ _ZX_64(argument1);
stage E1:
  GPR[registerW] = result1;
\end{lstlisting}}
\newpage

\paragraph{Encoding} Format \textsf{ALU\_DBMWRRR.M} (Section~\ref{sec:format-ALU-DBMWRRR.M}), \verb|exucode2 = 1100|\\

\bitfields{ 1/P, 2/00, 2/-- --, 27/upper27, 1/1, 2/11, 1/0, 4/1100, 6/registerW, 2/10, 3/001, 1/0, 1/splat32, 5/lower5, 6/registerZ }


\paragraph{Syntax} \texttt{sbmm8eord} \emph{registerW} = \emph{registerZ}, \emph{upper27\_\discretionary{}{}{}lower5}\emph{splat32}

\paragraph{Description} The \emph{upper27\_\discretionary{}{}{}lower5} interpreted as a 8$\times$8 bit matrix is multiplied by the \emph{registerZ} interpreted as a 8$\times$8 bit matrix. The resulting $\times$8 bit matrix is XORed with the value of the \emph{registerW} and stored into the \emph{registerW}.


\paragraph{Modifier(s)}
\noindent\begin{itemize}
\item splat32 (cf. splat32 in section \ref{par:modifier-splat32} on page \pageref{par:modifier-splat32})
\end{itemize}

\paragraph{Execution} {Reservation table \textsf{ALU\_TINY.X} (Section~\ref{sec:reservation-ALU-TINY.X})}
~
{\begin{lstlisting}
stage ID:
  new splat32 = _ZX_1(splat32);
stage RR:
  new argument3 = splat32 ? (_WX_32(upper27_lower5) << 32) | _ZX_32(_WX_32(upper27_lower5)) : _SX_32(_WX_32(upper27_lower5));
  new argument2 = GPR[registerZ];
  new argument1 = GPR[registerW];
  new result1 = _BMM_8(_ZX_64(argument3), _ZX_64(argument2)) ^ _ZX_64(argument1);
stage E1:
  GPR[registerW] = result1;
\end{lstlisting}}
\newpage
\subsubsection[{\small SBMM8EORDP: Swapped Bit Matrix Multiplication 8\texorpdfstring{$\times$}{x}8 Exclusive OR Double Word Pair}]{SBMM8EORDP\hfill Swapped Bit Matrix Multiplication 8$\times$8 Exclusive OR Double Word Pair}
\label{instr:SBMM8EORDP}\index{sbmm8eordp}
 
\paragraph{Encoding} Format \textsf{ALU\_DPSBMME} (Section~\ref{sec:format-ALU-DPSBMME}), \verb|exucode2 = 1100|\\

\bitfields{ 1/P, 2/11, 1/1, 4/1100, 5/registerM, 1/--, 2/10, 3/001, 1/0, 5/registerO, 1/--, 5/registerP, 1/0 }


\paragraph{Syntax} \texttt{sbmm8eordp} \emph{registerM} = \emph{registerP}, \emph{registerO}

\paragraph{Description} The \emph{registerO} interpreted as two 8$\times$8 bit matrices are multiplied by the \emph{registerP} interpreted as two 8$\times$8 bit matrices. The resulting $\times$8 bit matrix pair is XORed with the two double words of the \emph{registerM} and stored into the \emph{registerM}.


\paragraph{Execution} {Reservation table \textsf{ALU\_LITE} (Section~\ref{sec:reservation-ALU-LITE})}
~
{\begin{lstlisting}
stage RR:
  new argument3 = PGR[registerO];
  new argument2 = PGR[registerP];
  new argument1 = PGR[registerM];
  for (I = 0; I < 2; I += 1) {
    result1.64[I] = _BMM_8(_ZX_64(argument3.64[I]), _ZX_64(argument2.64[I])) ^ _ZX_64(argument1.64[I])
  };
stage E1:
  PGR[registerM] = result1;
\end{lstlisting}}
\newpage

\paragraph{Encoding} Format \textsf{ALU\_DPSBMME.M} (Section~\ref{sec:format-ALU-DPSBMME.M}), \verb|exucode2 = 1100|\\

\bitfields{ 1/P, 2/00, 2/-- --, 27/upper27, 1/1, 2/11, 1/1, 4/1100, 5/registerM, 1/--, 2/10, 3/001, 1/0, 1/splat32, 5/lower5, 5/registerP, 1/0 }


\paragraph{Syntax} \texttt{sbmm8eordp} \emph{registerM} = \emph{registerP}, \emph{upper27\_\discretionary{}{}{}lower5}\emph{splat32}

\paragraph{Description} The \emph{upper27\_\discretionary{}{}{}lower5} interpreted as two 8$\times$8 bit matrices are multiplied by the \emph{registerP} interpreted as two 8$\times$8 bit matrices. The resulting $\times$8 bit matrix pair is XORed with the two double words of the \emph{registerM} and stored into the \emph{registerM}.


\paragraph{Modifier(s)}
\noindent\begin{itemize}
\item splat32 (cf. splat32 in section \ref{par:modifier-splat32} on page \pageref{par:modifier-splat32})
\end{itemize}

\paragraph{Execution} {Reservation table \textsf{ALU\_LITE.X} (Section~\ref{sec:reservation-ALU-LITE.X})}
~
{\begin{lstlisting}
stage ID:
  new splat32 = _ZX_1(splat32);
stage RR:
  new argument3 = splat32 ? (_WX_32(upper27_lower5) << 32) | _ZX_32(_WX_32(upper27_lower5)) : _SX_32(_WX_32(upper27_lower5));
  new argument2 = PGR[registerP];
  new argument1 = PGR[registerM];
  for (I = 0; I < 2; I += 1) {
    result1.64[I] = _BMM_8(_ZX_64(argument3.64[I]), _ZX_64(argument2.64[I])) ^ _ZX_64(argument1.64[I])
  };
stage E1:
  PGR[registerM] = result1;
\end{lstlisting}}
\newpage
\subsubsection[{\small SBMMT8D: Swapped Bit Matrix Multiplication Transposed 8\texorpdfstring{$\times$}{x}8 Double Word}]{SBMMT8D\hfill Swapped Bit Matrix Multiplication Transposed 8$\times$8 Double Word}
\label{instr:SBMMT8D}\index{sbmmt8d}
 
\paragraph{Encoding} Format \textsf{ALU\_DBMWRR} (Section~\ref{sec:format-ALU-DBMWRR}), \verb|exucode2 = 1111|\\

\bitfields{ 1/P, 2/11, 1/0, 4/1111, 6/registerW, 2/10, 3/001, 1/0, 6/registerY, 6/registerZ }


\paragraph{Syntax} \texttt{sbmmt8d} \emph{registerW} = \emph{registerZ}, \emph{registerY}

\paragraph{Description} The \emph{registerY} interpreted as a 8$\times$8 bit matrix is multiplied by the \emph{registerZ} interpreted as a 8$\times$8 bit matrix. The resulting 8$\times$8 bit matrix is transposed, and stored into the \emph{registerW}.


\paragraph{Execution} {Reservation table \textsf{ALU\_TINY} (Section~\ref{sec:reservation-ALU-TINY})}
~
{\begin{lstlisting}
stage RR:
  new argument3 = GPR[registerY];
  new argument2 = GPR[registerZ];
  new result1 = _BMT_8(_BMM_8(_ZX_64(argument3), _ZX_64(argument2)));
stage E1:
  GPR[registerW] = result1;
\end{lstlisting}}
\newpage

\paragraph{Encoding} Format \textsf{ALU\_DBMWRR.M} (Section~\ref{sec:format-ALU-DBMWRR.M}), \verb|exucode2 = 1111|\\

\bitfields{ 1/P, 2/00, 2/-- --, 27/upper27, 1/1, 2/11, 1/0, 4/1111, 6/registerW, 2/10, 3/001, 1/0, 1/splat32, 5/lower5, 6/registerZ }


\paragraph{Syntax} \texttt{sbmmt8d} \emph{registerW} = \emph{registerZ}, \emph{upper27\_\discretionary{}{}{}lower5}\emph{splat32}

\paragraph{Description} The \emph{upper27\_\discretionary{}{}{}lower5} interpreted as a 8$\times$8 bit matrix is multiplied by the \emph{registerZ} interpreted as a 8$\times$8 bit matrix. The resulting 8$\times$8 bit matrix is transposed, and stored into the \emph{registerW}.


\paragraph{Modifier(s)}
\noindent\begin{itemize}
\item splat32 (cf. splat32 in section \ref{par:modifier-splat32} on page \pageref{par:modifier-splat32})
\end{itemize}

\paragraph{Execution} {Reservation table \textsf{ALU\_TINY.X} (Section~\ref{sec:reservation-ALU-TINY.X})}
~
{\begin{lstlisting}
stage ID:
  new splat32 = _ZX_1(splat32);
stage RR:
  new argument3 = splat32 ? (_WX_32(upper27_lower5) << 32) | _ZX_32(_WX_32(upper27_lower5)) : _SX_32(_WX_32(upper27_lower5));
  new argument2 = GPR[registerZ];
  new result1 = _BMT_8(_BMM_8(_ZX_64(argument3), _ZX_64(argument2)));
stage E1:
  GPR[registerW] = result1;
\end{lstlisting}}
\newpage

\paragraph{Encoding} Format \textsf{ALU\_DBMWRI} (Section~\ref{sec:format-ALU-DBMWRI}), \verb|exucode2 = 1111|\\

\bitfields{ 1/P, 2/11, 1/0, 4/1111, 6/registerW, 2/01, 10/signed10, 6/registerZ }


\paragraph{Syntax} \texttt{sbmmt8d} \emph{registerW} = \emph{registerZ}, \emph{signed10}

\paragraph{Description} The \emph{signed10} interpreted as a 8$\times$8 bit matrix is multiplied by the \emph{registerZ} interpreted as a 8$\times$8 bit matrix. The resulting 8$\times$8 bit matrix is transposed, and stored into the \emph{registerW}.


\paragraph{Execution} {Reservation table \textsf{ALU\_TINY} (Section~\ref{sec:reservation-ALU-TINY})}
~
{\begin{lstlisting}
stage RR:
  new argument3 = _SX_10(signed10);
  new argument2 = GPR[registerZ];
  new result1 = _BMT_8(_BMM_8(_ZX_64(argument3), _ZX_64(argument2)));
stage E1:
  GPR[registerW] = result1;
\end{lstlisting}}
\newpage

\paragraph{Encoding} Format \textsf{ALU\_DBMWRI.X} (Section~\ref{sec:format-ALU-DBMWRI.X}), \verb|exucode2 = 1111|\\

\bitfields{ 1/P, 2/00, 2/-- --, 27/upper27, 1/1, 2/11, 1/0, 4/1111, 6/registerW, 2/01, 10/lower10, 6/registerZ }


\paragraph{Syntax} \texttt{sbmmt8d} \emph{registerW} = \emph{registerZ}, \emph{upper27\_\discretionary{}{}{}lower10}

\paragraph{Description} The \emph{upper27\_\discretionary{}{}{}lower10} interpreted as a 8$\times$8 bit matrix is multiplied by the \emph{registerZ} interpreted as a 8$\times$8 bit matrix. The resulting 8$\times$8 bit matrix is transposed, and stored into the \emph{registerW}.


\paragraph{Execution} {Reservation table \textsf{ALU\_TINY.X} (Section~\ref{sec:reservation-ALU-TINY.X})}
~
{\begin{lstlisting}
stage RR:
  new argument3 = _SX_37(upper27_lower10);
  new argument2 = GPR[registerZ];
  new result1 = _BMT_8(_BMM_8(_ZX_64(argument3), _ZX_64(argument2)));
stage E1:
  GPR[registerW] = result1;
\end{lstlisting}}
\newpage

\paragraph{Encoding} Format \textsf{ALU\_DBMWRI.Y} (Section~\ref{sec:format-ALU-DBMWRI.Y}), \verb|exucode2 = 1111|\\

\bitfields{ 1/P, 2/00, 2/-- --, 27/extend27, 1/1, 2/00, 2/-- --, 27/upper27, 1/1, 2/11, 1/0, 4/1111, 6/registerW, 2/01, 10/lower10, 6/registerZ }


\paragraph{Syntax} \texttt{sbmmt8d} \emph{registerW} = \emph{registerZ}, \emph{extend27\_\discretionary{}{}{}upper27\_\discretionary{}{}{}lower10}

\paragraph{Description} The \emph{extend27\_\discretionary{}{}{}upper27\_\discretionary{}{}{}lower10} interpreted as a 8$\times$8 bit matrix is multiplied by the \emph{registerZ} interpreted as a 8$\times$8 bit matrix. The resulting 8$\times$8 bit matrix is transposed, and stored into the \emph{registerW}.


\paragraph{Execution} {Reservation table \textsf{ALU\_TINY.Y} (Section~\ref{sec:reservation-ALU-TINY.Y})}
~
{\begin{lstlisting}
stage RR:
  new argument3 = _WX_64(extend27_upper27_lower10);
  new argument2 = GPR[registerZ];
  new result1 = _BMT_8(_BMM_8(_ZX_64(argument3), _ZX_64(argument2)));
stage E1:
  GPR[registerW] = result1;
\end{lstlisting}}
\newpage
\subsubsection[{\small SBMMT8DP: Swapped Bit Matrix Multiplication Transposed 8\texorpdfstring{$\times$}{x}8 Double Word Pair}]{SBMMT8DP\hfill Swapped Bit Matrix Multiplication Transposed 8$\times$8 Double Word Pair}
\label{instr:SBMMT8DP}\index{sbmmt8dp}
 
\paragraph{Encoding} Format \textsf{ALU\_DPSBMM} (Section~\ref{sec:format-ALU-DPSBMM}), \verb|exucode2 = 1111|\\

\bitfields{ 1/P, 2/11, 1/1, 4/1111, 5/registerM, 1/--, 2/10, 3/001, 1/0, 5/registerO, 1/--, 5/registerP, 1/0 }


\paragraph{Syntax} \texttt{sbmmt8dp} \emph{registerM} = \emph{registerP}, \emph{registerO}

\paragraph{Description} The \emph{registerO} interpreted as two 8$\times$8 bit matrices are multiplied by the \emph{registerP} interpreted as two 8$\times$8 bit matrices. The resulting 8$\times$8 bit matrix pair is stored into the \emph{registerM}.


\paragraph{Execution} {Reservation table \textsf{ALU\_LITE} (Section~\ref{sec:reservation-ALU-LITE})}
~
{\begin{lstlisting}
stage RR:
  new argument3 = PGR[registerO];
  new argument2 = PGR[registerP];
  for (I = 0; I < 2; I += 1) {
    result1.64[I] = _BMT_8(_BMM_8(_ZX_64(argument3.64[I]), _ZX_64(argument2.64[I])))
  };
stage E1:
  PGR[registerM] = result1;
\end{lstlisting}}
\newpage

\paragraph{Encoding} Format \textsf{ALU\_DPSBMM.M} (Section~\ref{sec:format-ALU-DPSBMM.M}), \verb|exucode2 = 1111|\\

\bitfields{ 1/P, 2/00, 2/-- --, 27/upper27, 1/1, 2/11, 1/1, 4/1111, 5/registerM, 1/--, 2/10, 3/001, 1/0, 1/splat32, 5/lower5, 5/registerP, 1/0 }


\paragraph{Syntax} \texttt{sbmmt8dp} \emph{registerM} = \emph{registerP}, \emph{upper27\_\discretionary{}{}{}lower5}\emph{splat32}

\paragraph{Description} The \emph{upper27\_\discretionary{}{}{}lower5} interpreted as two 8$\times$8 bit matrices are multiplied by the \emph{registerP} interpreted as two 8$\times$8 bit matrices. The resulting 8$\times$8 bit matrix pair is stored into the \emph{registerM}.


\paragraph{Modifier(s)}
\noindent\begin{itemize}
\item splat32 (cf. splat32 in section \ref{par:modifier-splat32} on page \pageref{par:modifier-splat32})
\end{itemize}

\paragraph{Execution} {Reservation table \textsf{ALU\_LITE.X} (Section~\ref{sec:reservation-ALU-LITE.X})}
~
{\begin{lstlisting}
stage ID:
  new splat32 = _ZX_1(splat32);
stage RR:
  new argument3 = splat32 ? (_WX_32(upper27_lower5) << 32) | _ZX_32(_WX_32(upper27_lower5)) : _SX_32(_WX_32(upper27_lower5));
  new argument2 = PGR[registerP];
  for (I = 0; I < 2; I += 1) {
    result1.64[I] = _BMT_8(_BMM_8(_ZX_64(argument3.64[I]), _ZX_64(argument2.64[I])))
  };
stage E1:
  PGR[registerM] = result1;
\end{lstlisting}}
\newpage
\subsubsection[{\small SBMMT8EORD: Swapped Bit Matrix Multiplication Transposed 8\texorpdfstring{$\times$}{x}8 Exclusive OR Double Word}]{SBMMT8EORD\hfill Swapped Bit Matrix Multiplication Transposed 8$\times$8 Exclusive OR Double Word}
\label{instr:SBMMT8EORD}\index{sbmmt8eord}
 
\paragraph{Encoding} Format \textsf{ALU\_DBMWRRR} (Section~\ref{sec:format-ALU-DBMWRRR}), \verb|exucode2 = 1101|\\

\bitfields{ 1/P, 2/11, 1/0, 4/1101, 6/registerW, 2/10, 3/001, 1/0, 6/registerY, 6/registerZ }


\paragraph{Syntax} \texttt{sbmmt8eord} \emph{registerW} = \emph{registerZ}, \emph{registerY}

\paragraph{Description} The \emph{registerY} interpreted as a 8$\times$8 bit matrix is multiplied by the \emph{registerZ} interpreted as a 8$\times$8 bit matrix. The resulting 8$\times$8 bit matrix is transposed, XORed with the value of the \emph{registerW} and stored into the \emph{registerW}.


\paragraph{Execution} {Reservation table \textsf{ALU\_TINY} (Section~\ref{sec:reservation-ALU-TINY})}
~
{\begin{lstlisting}
stage RR:
  new argument3 = GPR[registerY];
  new argument2 = GPR[registerZ];
  new argument1 = GPR[registerW];
  new result1 = _BMT_8(_BMM_8(_ZX_64(argument3), _ZX_64(argument2))) ^ _ZX_64(argument1);
stage E1:
  GPR[registerW] = result1;
\end{lstlisting}}
\newpage

\paragraph{Encoding} Format \textsf{ALU\_DBMWRRR.M} (Section~\ref{sec:format-ALU-DBMWRRR.M}), \verb|exucode2 = 1101|\\

\bitfields{ 1/P, 2/00, 2/-- --, 27/upper27, 1/1, 2/11, 1/0, 4/1101, 6/registerW, 2/10, 3/001, 1/0, 1/splat32, 5/lower5, 6/registerZ }


\paragraph{Syntax} \texttt{sbmmt8eord} \emph{registerW} = \emph{registerZ}, \emph{upper27\_\discretionary{}{}{}lower5}\emph{splat32}

\paragraph{Description} The \emph{upper27\_\discretionary{}{}{}lower5} interpreted as a 8$\times$8 bit matrix is multiplied by the \emph{registerZ} interpreted as a 8$\times$8 bit matrix. The resulting 8$\times$8 bit matrix is transposed, XORed with the value of the \emph{registerW} and stored into the \emph{registerW}.


\paragraph{Modifier(s)}
\noindent\begin{itemize}
\item splat32 (cf. splat32 in section \ref{par:modifier-splat32} on page \pageref{par:modifier-splat32})
\end{itemize}

\paragraph{Execution} {Reservation table \textsf{ALU\_TINY.X} (Section~\ref{sec:reservation-ALU-TINY.X})}
~
{\begin{lstlisting}
stage ID:
  new splat32 = _ZX_1(splat32);
stage RR:
  new argument3 = splat32 ? (_WX_32(upper27_lower5) << 32) | _ZX_32(_WX_32(upper27_lower5)) : _SX_32(_WX_32(upper27_lower5));
  new argument2 = GPR[registerZ];
  new argument1 = GPR[registerW];
  new result1 = _BMT_8(_BMM_8(_ZX_64(argument3), _ZX_64(argument2))) ^ _ZX_64(argument1);
stage E1:
  GPR[registerW] = result1;
\end{lstlisting}}
\newpage
\subsubsection[{\small SBMMT8EORDP: Swapped Bit Matrix Multiplication Transposed 8\texorpdfstring{$\times$}{x}8 Exclusive OR Double Word Pair}]{SBMMT8EORDP\hfill Swapped Bit Matrix Multiplication Transposed 8$\times$8 Exclusive OR Double Word Pair}
\label{instr:SBMMT8EORDP}\index{sbmmt8eordp}
 
\paragraph{Encoding} Format \textsf{ALU\_DPSBMME} (Section~\ref{sec:format-ALU-DPSBMME}), \verb|exucode2 = 1101|\\

\bitfields{ 1/P, 2/11, 1/1, 4/1101, 5/registerM, 1/--, 2/10, 3/001, 1/0, 5/registerO, 1/--, 5/registerP, 1/0 }


\paragraph{Syntax} \texttt{sbmmt8eordp} \emph{registerM} = \emph{registerP}, \emph{registerO}

\paragraph{Description} The \emph{registerO} interpreted as two 8$\times$8 bit matrices are multiplied by the \emph{registerP} interpreted as two 8$\times$8 bit matrices. The resulting $\times$8 bit matrix pair is transposed, XORed with the two double words of the \emph{registerM} and stored into the \emph{registerM}.


\paragraph{Execution} {Reservation table \textsf{ALU\_LITE} (Section~\ref{sec:reservation-ALU-LITE})}
~
{\begin{lstlisting}
stage RR:
  new argument3 = PGR[registerO];
  new argument2 = PGR[registerP];
  new argument1 = PGR[registerM];
  for (I = 0; I < 2; I += 1) {
    result1.64[I] = _BMT_8(_BMM_8(_ZX_64(argument3.64[I]), _ZX_64(argument2.64[I]))) ^ _ZX_64(argument1.64[I])
  };
stage E1:
  PGR[registerM] = result1;
\end{lstlisting}}
\newpage

\paragraph{Encoding} Format \textsf{ALU\_DPSBMME.M} (Section~\ref{sec:format-ALU-DPSBMME.M}), \verb|exucode2 = 1101|\\

\bitfields{ 1/P, 2/00, 2/-- --, 27/upper27, 1/1, 2/11, 1/1, 4/1101, 5/registerM, 1/--, 2/10, 3/001, 1/0, 1/splat32, 5/lower5, 5/registerP, 1/0 }


\paragraph{Syntax} \texttt{sbmmt8eordp} \emph{registerM} = \emph{registerP}, \emph{upper27\_\discretionary{}{}{}lower5}\emph{splat32}

\paragraph{Description} The \emph{upper27\_\discretionary{}{}{}lower5} interpreted as two 8$\times$8 bit matrices are multiplied by the \emph{registerP} interpreted as two 8$\times$8 bit matrices. The resulting $\times$8 bit matrix pair is transposed, XORed with the two double words of the \emph{registerM} and stored into the \emph{registerM}.


\paragraph{Modifier(s)}
\noindent\begin{itemize}
\item splat32 (cf. splat32 in section \ref{par:modifier-splat32} on page \pageref{par:modifier-splat32})
\end{itemize}

\paragraph{Execution} {Reservation table \textsf{ALU\_LITE.X} (Section~\ref{sec:reservation-ALU-LITE.X})}
~
{\begin{lstlisting}
stage ID:
  new splat32 = _ZX_1(splat32);
stage RR:
  new argument3 = splat32 ? (_WX_32(upper27_lower5) << 32) | _ZX_32(_WX_32(upper27_lower5)) : _SX_32(_WX_32(upper27_lower5));
  new argument2 = PGR[registerP];
  new argument1 = PGR[registerM];
  for (I = 0; I < 2; I += 1) {
    result1.64[I] = _BMT_8(_BMM_8(_ZX_64(argument3.64[I]), _ZX_64(argument2.64[I]))) ^ _ZX_64(argument1.64[I])
  };
stage E1:
  PGR[registerM] = result1;
\end{lstlisting}}
\newpage
\subsubsection[{\small SIGNBX: Apply Sign Byte Hexadecuple}]{SIGNBX\hfill Apply Sign Byte Hexadecuple}
\label{instr:SIGNBX}\index{signbx}
 
\paragraph{Encoding} Format \textsf{ALU\_BXWRR0} (Section~\ref{sec:format-ALU-BXWRR0}), \verb|exucode2 = 0000|\\

\bitfields{ 1/P, 2/11, 1/1, 4/0000, 5/registerM, 1/--, 2/10, 3/000, 1/1, 5/registerO, 1/--, 5/registerP, 1/1 }


\paragraph{Syntax} \texttt{signbx} \emph{registerM} = \emph{registerP}, \emph{registerO}

\paragraph{Description} The \emph{registerO} and the \emph{registerP} are interpreted as sixteen signed bytes packed into 128 bits. The \emph{registerP} bytes are multiplied by the signs of the \emph{registerO} bytes. The sixteen resulting bytes are packed and stored into the \emph{registerM}.


\paragraph{Execution} {Reservation table \textsf{ALU\_LITE} (Section~\ref{sec:reservation-ALU-LITE})}
~
{\begin{lstlisting}
stage RR:
  new argument3 = PGR[registerO];
  new argument2 = PGR[registerP];
  for (I = 0; I < 16; I += 1) {
    result1.8[I] = argument2.8[I]
    if (_SX_8(argument3.8[I]) < 0) {
      result1.8[I] = -argument2.8[I];
    }
    if (_ZX_8(argument3.8[I]) == 0) {
      result1.8[I] = 0;
    }
  };
stage E1:
  PGR[registerM] = result1;
\end{lstlisting}}
\newpage

\paragraph{Encoding} Format \textsf{ALU\_BXWRR0.M} (Section~\ref{sec:format-ALU-BXWRR0.M}), \verb|exucode2 = 0000|\\

\bitfields{ 1/P, 2/00, 2/-- --, 27/upper27, 1/1, 2/11, 1/1, 4/0000, 5/registerM, 1/--, 2/10, 3/000, 1/1, 1/splat32, 5/lower5, 5/registerP, 1/1 }


\paragraph{Syntax} \texttt{signbx} \emph{registerM} = \emph{registerP}, \emph{upper27\_\discretionary{}{}{}lower5}\emph{splat32}

\paragraph{Description} The \emph{upper27\_\discretionary{}{}{}lower5} and the \emph{registerP} are interpreted as sixteen signed bytes packed into 128 bits. The \emph{registerP} bytes are multiplied by the signs of the \emph{upper27\_\discretionary{}{}{}lower5} bytes. The sixteen resulting bytes are packed and stored into the \emph{registerM}.


\paragraph{Modifier(s)}
\noindent\begin{itemize}
\item splat32 (cf. splat32 in section \ref{par:modifier-splat32} on page \pageref{par:modifier-splat32})
\end{itemize}

\paragraph{Execution} {Reservation table \textsf{ALU\_LITE.X} (Section~\ref{sec:reservation-ALU-LITE.X})}
~
{\begin{lstlisting}
stage ID:
  new splat32 = _ZX_1(splat32);
stage RR:
  new argument3 = splat32 ? (_WX_32(upper27_lower5) << 32) | _ZX_32(_WX_32(upper27_lower5)) : _SX_32(_WX_32(upper27_lower5));
  new argument2 = PGR[registerP];
  for (I = 0; I < 16; I += 1) {
    result1.8[I] = argument2.8[I]
    if (_SX_8(argument3.8[I]) < 0) {
      result1.8[I] = -argument2.8[I];
    }
    if (_ZX_8(argument3.8[I]) == 0) {
      result1.8[I] = 0;
    }
  };
stage E1:
  PGR[registerM] = result1;
\end{lstlisting}}
\newpage
\subsubsection[{\small SIGND: Apply Sign Double Word}]{SIGND\hfill Apply Sign Double Word}
\label{instr:SIGND}\index{signd}
 
\paragraph{Encoding} Format \textsf{ALU\_DWRR0} (Section~\ref{sec:format-ALU-DWRR0}), \verb|exucode2 = 0000|\\

\bitfields{ 1/P, 2/11, 1/0, 4/0000, 6/registerW, 2/10, 3/000, 1/0, 6/registerY, 6/registerZ }


\paragraph{Syntax} \texttt{signd} \emph{registerW} = \emph{registerZ}, \emph{registerY}

\paragraph{Description} The \emph{registerY} and the \emph{registerZ} are interpreted as signed integer double words. The \emph{registerZ} double word is multiplied by the sign of the \emph{registerY} double word. The resulting double word is stored into the \emph{registerW}.


\paragraph{Execution} {Reservation table \textsf{ALU\_TINY} (Section~\ref{sec:reservation-ALU-TINY})}
~
{\begin{lstlisting}
stage RR:
  new argument3 = GPR[registerY];
  new argument2 = GPR[registerZ];
  new result1 = argument2;
  if (_SX_64(argument3) < 0) {
    result1 = -argument2;
  }
  if (_ZX_64(argument3) == 0) {
    result1 = 0;
  }
stage E1:
  GPR[registerW] = result1;
\end{lstlisting}}
\newpage

\paragraph{Encoding} Format \textsf{ALU\_DWRR0.M} (Section~\ref{sec:format-ALU-DWRR0.M}), \verb|exucode2 = 0000|\\

\bitfields{ 1/P, 2/00, 2/-- --, 27/upper27, 1/1, 2/11, 1/0, 4/0000, 6/registerW, 2/10, 3/000, 1/0, 1/splat32, 5/lower5, 6/registerZ }


\paragraph{Syntax} \texttt{signd} \emph{registerW} = \emph{registerZ}, \emph{upper27\_\discretionary{}{}{}lower5}\emph{splat32}

\paragraph{Description} The \emph{upper27\_\discretionary{}{}{}lower5} and the \emph{registerZ} are interpreted as signed integer double words. The \emph{registerZ} double word is multiplied by the sign of the \emph{upper27\_\discretionary{}{}{}lower5} double word. The resulting double word is stored into the \emph{registerW}.


\paragraph{Modifier(s)}
\noindent\begin{itemize}
\item splat32 (cf. splat32 in section \ref{par:modifier-splat32} on page \pageref{par:modifier-splat32})
\end{itemize}

\paragraph{Execution} {Reservation table \textsf{ALU\_TINY.X} (Section~\ref{sec:reservation-ALU-TINY.X})}
~
{\begin{lstlisting}
stage ID:
  new splat32 = _ZX_1(splat32);
stage RR:
  new argument3 = splat32 ? (_WX_32(upper27_lower5) << 32) | _ZX_32(_WX_32(upper27_lower5)) : _SX_32(_WX_32(upper27_lower5));
  new argument2 = GPR[registerZ];
  new result1 = argument2;
  if (_SX_64(argument3) < 0) {
    result1 = -argument2;
  }
  if (_ZX_64(argument3) == 0) {
    result1 = 0;
  }
stage E1:
  GPR[registerW] = result1;
\end{lstlisting}}
\newpage
\subsubsection[{\small SIGNDP: Apply Sign Double Word Pair}]{SIGNDP\hfill Apply Sign Double Word Pair}
\label{instr:SIGNDP}\index{signdp}
 
\paragraph{Encoding} Format \textsf{ALU\_DPWRR0} (Section~\ref{sec:format-ALU-DPWRR0}), \verb|exucode2 = 0000|\\

\bitfields{ 1/P, 2/11, 1/1, 4/0000, 5/registerM, 1/--, 2/10, 3/000, 1/0, 5/registerO, 1/--, 5/registerP, 1/0 }


\paragraph{Syntax} \texttt{signdp} \emph{registerM} = \emph{registerP}, \emph{registerO}

\paragraph{Description} The \emph{registerO} and the \emph{registerP} are interpreted as two signed integer double words packed into 128 bits. The \emph{registerP} double words are multiplied by the signs of the \emph{registerO} double words. The two resulting double words are packed and stored into the \emph{registerM}.


\paragraph{Execution} {Reservation table \textsf{ALU\_LITE} (Section~\ref{sec:reservation-ALU-LITE})}
~
{\begin{lstlisting}
stage RR:
  new argument3 = PGR[registerO];
  new argument2 = PGR[registerP];
  for (I = 0; I < 2; I += 1) {
    result1.64[I] = argument2.64[I]
    if (_SX_64(argument3.64[I]) < 0) {
      result1.64[I] = -argument2.64[I];
    }
    if (_ZX_64(argument3.64[I]) == 0) {
      result1.64[I] = 0;
    }
  };
stage E1:
  PGR[registerM] = result1;
\end{lstlisting}}
\newpage

\paragraph{Encoding} Format \textsf{ALU\_DPWRR0.M} (Section~\ref{sec:format-ALU-DPWRR0.M}), \verb|exucode2 = 0000|\\

\bitfields{ 1/P, 2/00, 2/-- --, 27/upper27, 1/1, 2/11, 1/1, 4/0000, 5/registerM, 1/--, 2/10, 3/000, 1/0, 1/splat32, 5/lower5, 5/registerP, 1/0 }


\paragraph{Syntax} \texttt{signdp} \emph{registerM} = \emph{registerP}, \emph{upper27\_\discretionary{}{}{}lower5}\emph{splat32}

\paragraph{Description} The \emph{upper27\_\discretionary{}{}{}lower5} and the \emph{registerP} are interpreted as two signed integer double words packed into 128 bits. The \emph{registerP} double words are multiplied by the signs of the \emph{upper27\_\discretionary{}{}{}lower5} double words. The two resulting double words are packed and stored into the \emph{registerM}.


\paragraph{Modifier(s)}
\noindent\begin{itemize}
\item splat32 (cf. splat32 in section \ref{par:modifier-splat32} on page \pageref{par:modifier-splat32})
\end{itemize}

\paragraph{Execution} {Reservation table \textsf{ALU\_LITE.X} (Section~\ref{sec:reservation-ALU-LITE.X})}
~
{\begin{lstlisting}
stage ID:
  new splat32 = _ZX_1(splat32);
stage RR:
  new argument3 = splat32 ? (_WX_32(upper27_lower5) << 32) | _ZX_32(_WX_32(upper27_lower5)) : _SX_32(_WX_32(upper27_lower5));
  new argument2 = PGR[registerP];
  for (I = 0; I < 2; I += 1) {
    result1.64[I] = argument2.64[I]
    if (_SX_64(argument3.64[I]) < 0) {
      result1.64[I] = -argument2.64[I];
    }
    if (_ZX_64(argument3.64[I]) == 0) {
      result1.64[I] = 0;
    }
  };
stage E1:
  PGR[registerM] = result1;
\end{lstlisting}}
\newpage
\subsubsection[{\small SIGNHO: Apply Sign Half Word Octuple}]{SIGNHO\hfill Apply Sign Half Word Octuple}
\label{instr:SIGNHO}\index{signho}
 
\paragraph{Encoding} Format \textsf{ALU\_HOWRR0} (Section~\ref{sec:format-ALU-HOWRR0}), \verb|exucode2 = 0000|\\

\bitfields{ 1/P, 2/11, 1/1, 4/0000, 5/registerM, 1/--, 2/10, 3/000, 1/1, 5/registerO, 1/--, 5/registerP, 1/0 }


\paragraph{Syntax} \texttt{signho} \emph{registerM} = \emph{registerP}, \emph{registerO}

\paragraph{Description} The \emph{registerO} and the \emph{registerP} are interpreted as eight signed integer half words packed into 128 bits. The \emph{registerP} half words are multiplied by the signs of the \emph{registerO} half words. The eight resulting half words are packed and stored into the \emph{registerM}.


\paragraph{Execution} {Reservation table \textsf{ALU\_LITE} (Section~\ref{sec:reservation-ALU-LITE})}
~
{\begin{lstlisting}
stage RR:
  new argument3 = PGR[registerO];
  new argument2 = PGR[registerP];
  for (I = 0; I < 8; I += 1) {
    result1.16[I] = argument2.16[I]
    if (_SX_16(argument3.16[I]) < 0) {
      result1.16[I] = -argument2.16[I];
    }
    if (_ZX_16(argument3.16[I]) == 0) {
      result1.16[I] = 0;
    }
  };
stage E1:
  PGR[registerM] = result1;
\end{lstlisting}}
\newpage

\paragraph{Encoding} Format \textsf{ALU\_HOWRR0.M} (Section~\ref{sec:format-ALU-HOWRR0.M}), \verb|exucode2 = 0000|\\

\bitfields{ 1/P, 2/00, 2/-- --, 27/upper27, 1/1, 2/11, 1/1, 4/0000, 5/registerM, 1/--, 2/10, 3/000, 1/1, 1/splat32, 5/lower5, 5/registerP, 1/0 }


\paragraph{Syntax} \texttt{signho} \emph{registerM} = \emph{registerP}, \emph{upper27\_\discretionary{}{}{}lower5}\emph{splat32}

\paragraph{Description} The \emph{upper27\_\discretionary{}{}{}lower5} and the \emph{registerP} are interpreted as eight signed integer half words packed into 128 bits. The \emph{registerP} half words are multiplied by the signs of the \emph{upper27\_\discretionary{}{}{}lower5} half words. The eight resulting half words are packed and stored into the \emph{registerM}.


\paragraph{Modifier(s)}
\noindent\begin{itemize}
\item splat32 (cf. splat32 in section \ref{par:modifier-splat32} on page \pageref{par:modifier-splat32})
\end{itemize}

\paragraph{Execution} {Reservation table \textsf{ALU\_LITE.X} (Section~\ref{sec:reservation-ALU-LITE.X})}
~
{\begin{lstlisting}
stage ID:
  new splat32 = _ZX_1(splat32);
stage RR:
  new argument3 = splat32 ? (_WX_32(upper27_lower5) << 32) | _ZX_32(_WX_32(upper27_lower5)) : _SX_32(_WX_32(upper27_lower5));
  new argument2 = PGR[registerP];
  for (I = 0; I < 8; I += 1) {
    result1.16[I] = argument2.16[I]
    if (_SX_16(argument3.16[I]) < 0) {
      result1.16[I] = -argument2.16[I];
    }
    if (_ZX_16(argument3.16[I]) == 0) {
      result1.16[I] = 0;
    }
  };
stage E1:
  PGR[registerM] = result1;
\end{lstlisting}}
\newpage
\subsubsection[{\small SIGNSBX: Apply Sign Saturated Byte Hexadecuple}]{SIGNSBX\hfill Apply Sign Saturated Byte Hexadecuple}
\label{instr:SIGNSBX}\index{signsbx}
 
\paragraph{Encoding} Format \textsf{ALU\_BXWRR0} (Section~\ref{sec:format-ALU-BXWRR0}), \verb|exucode2 = 0001|\\

\bitfields{ 1/P, 2/11, 1/1, 4/0001, 5/registerM, 1/--, 2/10, 3/000, 1/1, 5/registerO, 1/--, 5/registerP, 1/1 }


\paragraph{Syntax} \texttt{signsbx} \emph{registerM} = \emph{registerP}, \emph{registerO}

\paragraph{Description} The \emph{registerO} and the \emph{registerP} are interpreted as sixteen signed bytes packed into 128 bits. The \emph{registerP} bytes are multiplied by the signs of the \emph{registerO} bytes and 8-bit saturated. The sixteen resulting bytes are packed and stored into the \emph{registerM}.


\paragraph{Execution} {Reservation table \textsf{ALU\_LITE} (Section~\ref{sec:reservation-ALU-LITE})}
~
{\begin{lstlisting}
stage RR:
  new argument3 = PGR[registerO];
  new argument2 = PGR[registerP];
  for (I = 0; I < 16; I += 1) {
    result1.8[I] = argument2.8[I]
    if (_SX_8(argument3.8[I]) < 0) {
      result1.8[I] = _SAT_8(-argument2.8[I]);
    }
    if (_ZX_8(argument3.8[I]) == 0) {
      result1.8[I] = 0;
    }
  };
stage E1:
  PGR[registerM] = result1;
\end{lstlisting}}
\newpage

\paragraph{Encoding} Format \textsf{ALU\_BXWRR0.M} (Section~\ref{sec:format-ALU-BXWRR0.M}), \verb|exucode2 = 0001|\\

\bitfields{ 1/P, 2/00, 2/-- --, 27/upper27, 1/1, 2/11, 1/1, 4/0001, 5/registerM, 1/--, 2/10, 3/000, 1/1, 1/splat32, 5/lower5, 5/registerP, 1/1 }


\paragraph{Syntax} \texttt{signsbx} \emph{registerM} = \emph{registerP}, \emph{upper27\_\discretionary{}{}{}lower5}\emph{splat32}

\paragraph{Description} The \emph{upper27\_\discretionary{}{}{}lower5} and the \emph{registerP} are interpreted as sixteen signed bytes packed into 128 bits. The \emph{registerP} bytes are multiplied by the signs of the \emph{upper27\_\discretionary{}{}{}lower5} bytes and 8-bit saturated. The sixteen resulting bytes are packed and stored into the \emph{registerM}.


\paragraph{Modifier(s)}
\noindent\begin{itemize}
\item splat32 (cf. splat32 in section \ref{par:modifier-splat32} on page \pageref{par:modifier-splat32})
\end{itemize}

\paragraph{Execution} {Reservation table \textsf{ALU\_LITE.X} (Section~\ref{sec:reservation-ALU-LITE.X})}
~
{\begin{lstlisting}
stage ID:
  new splat32 = _ZX_1(splat32);
stage RR:
  new argument3 = splat32 ? (_WX_32(upper27_lower5) << 32) | _ZX_32(_WX_32(upper27_lower5)) : _SX_32(_WX_32(upper27_lower5));
  new argument2 = PGR[registerP];
  for (I = 0; I < 16; I += 1) {
    result1.8[I] = argument2.8[I]
    if (_SX_8(argument3.8[I]) < 0) {
      result1.8[I] = _SAT_8(-argument2.8[I]);
    }
    if (_ZX_8(argument3.8[I]) == 0) {
      result1.8[I] = 0;
    }
  };
stage E1:
  PGR[registerM] = result1;
\end{lstlisting}}
\newpage
\subsubsection[{\small SIGNSD: Apply Sign Saturated Double Word}]{SIGNSD\hfill Apply Sign Saturated Double Word}
\label{instr:SIGNSD}\index{signsd}
 
\paragraph{Encoding} Format \textsf{ALU\_DWRR0} (Section~\ref{sec:format-ALU-DWRR0}), \verb|exucode2 = 0001|\\

\bitfields{ 1/P, 2/11, 1/0, 4/0001, 6/registerW, 2/10, 3/000, 1/0, 6/registerY, 6/registerZ }


\paragraph{Syntax} \texttt{signsd} \emph{registerW} = \emph{registerZ}, \emph{registerY}

\paragraph{Description} The \emph{registerY} and the \emph{registerZ} are interpreted as signed integer double words. The \emph{registerZ} double words is multiplied by the sign of the \emph{registerY} double words and 64-bit saturated. The resulting double word is stored into the \emph{registerW}.


\paragraph{Execution} {Reservation table \textsf{ALU\_TINY} (Section~\ref{sec:reservation-ALU-TINY})}
~
{\begin{lstlisting}
stage RR:
  new argument3 = GPR[registerY];
  new argument2 = GPR[registerZ];
  new result1 = argument2;
  if (_SX_64(argument3) < 0) {
    result1 = _SAT_64(-argument2);
  }
  if (_ZX_64(argument3) == 0) {
    result1 = 0;
  }
stage E1:
  GPR[registerW] = result1;
\end{lstlisting}}
\newpage

\paragraph{Encoding} Format \textsf{ALU\_DWRR0.M} (Section~\ref{sec:format-ALU-DWRR0.M}), \verb|exucode2 = 0001|\\

\bitfields{ 1/P, 2/00, 2/-- --, 27/upper27, 1/1, 2/11, 1/0, 4/0001, 6/registerW, 2/10, 3/000, 1/0, 1/splat32, 5/lower5, 6/registerZ }


\paragraph{Syntax} \texttt{signsd} \emph{registerW} = \emph{registerZ}, \emph{upper27\_\discretionary{}{}{}lower5}\emph{splat32}

\paragraph{Description} The \emph{upper27\_\discretionary{}{}{}lower5} and the \emph{registerZ} are interpreted as signed integer double words. The \emph{registerZ} double words is multiplied by the sign of the \emph{upper27\_\discretionary{}{}{}lower5} double words and 64-bit saturated. The resulting double word is stored into the \emph{registerW}.


\paragraph{Modifier(s)}
\noindent\begin{itemize}
\item splat32 (cf. splat32 in section \ref{par:modifier-splat32} on page \pageref{par:modifier-splat32})
\end{itemize}

\paragraph{Execution} {Reservation table \textsf{ALU\_TINY.X} (Section~\ref{sec:reservation-ALU-TINY.X})}
~
{\begin{lstlisting}
stage ID:
  new splat32 = _ZX_1(splat32);
stage RR:
  new argument3 = splat32 ? (_WX_32(upper27_lower5) << 32) | _ZX_32(_WX_32(upper27_lower5)) : _SX_32(_WX_32(upper27_lower5));
  new argument2 = GPR[registerZ];
  new result1 = argument2;
  if (_SX_64(argument3) < 0) {
    result1 = _SAT_64(-argument2);
  }
  if (_ZX_64(argument3) == 0) {
    result1 = 0;
  }
stage E1:
  GPR[registerW] = result1;
\end{lstlisting}}
\newpage
\subsubsection[{\small SIGNSDP: Apply Sign Saturated Double Word Pair}]{SIGNSDP\hfill Apply Sign Saturated Double Word Pair}
\label{instr:SIGNSDP}\index{signsdp}
 
\paragraph{Encoding} Format \textsf{ALU\_DPWRR0} (Section~\ref{sec:format-ALU-DPWRR0}), \verb|exucode2 = 0001|\\

\bitfields{ 1/P, 2/11, 1/1, 4/0001, 5/registerM, 1/--, 2/10, 3/000, 1/0, 5/registerO, 1/--, 5/registerP, 1/0 }


\paragraph{Syntax} \texttt{signsdp} \emph{registerM} = \emph{registerP}, \emph{registerO}

\paragraph{Description} The \emph{registerO} and the \emph{registerP} are interpreted as two signed integer double words packed into 128 bits. The \emph{registerP} double words are multiplied by the signs of the \emph{registerO} double words and 64-bit saturated. The two resulting double words are packed and stored into the \emph{registerM}.


\paragraph{Execution} {Reservation table \textsf{ALU\_LITE} (Section~\ref{sec:reservation-ALU-LITE})}
~
{\begin{lstlisting}
stage RR:
  new argument3 = PGR[registerO];
  new argument2 = PGR[registerP];
  for (I = 0; I < 2; I += 1) {
    result1.64[I] = argument2.64[I]
    if (_SX_64(argument3.64[I]) < 0) {
      result1.64[I] = _SAT_64(-argument2.64[I]);
    }
    if (_ZX_64(argument3.64[I]) == 0) {
      result1.64[I] = 0;
    }
  };
stage E1:
  PGR[registerM] = result1;
\end{lstlisting}}
\newpage

\paragraph{Encoding} Format \textsf{ALU\_DPWRR0.M} (Section~\ref{sec:format-ALU-DPWRR0.M}), \verb|exucode2 = 0001|\\

\bitfields{ 1/P, 2/00, 2/-- --, 27/upper27, 1/1, 2/11, 1/1, 4/0001, 5/registerM, 1/--, 2/10, 3/000, 1/0, 1/splat32, 5/lower5, 5/registerP, 1/0 }


\paragraph{Syntax} \texttt{signsdp} \emph{registerM} = \emph{registerP}, \emph{upper27\_\discretionary{}{}{}lower5}\emph{splat32}

\paragraph{Description} The \emph{upper27\_\discretionary{}{}{}lower5} and the \emph{registerP} are interpreted as two signed integer double words packed into 128 bits. The \emph{registerP} double words are multiplied by the signs of the \emph{upper27\_\discretionary{}{}{}lower5} double words and 64-bit saturated. The two resulting double words are packed and stored into the \emph{registerM}.


\paragraph{Modifier(s)}
\noindent\begin{itemize}
\item splat32 (cf. splat32 in section \ref{par:modifier-splat32} on page \pageref{par:modifier-splat32})
\end{itemize}

\paragraph{Execution} {Reservation table \textsf{ALU\_LITE.X} (Section~\ref{sec:reservation-ALU-LITE.X})}
~
{\begin{lstlisting}
stage ID:
  new splat32 = _ZX_1(splat32);
stage RR:
  new argument3 = splat32 ? (_WX_32(upper27_lower5) << 32) | _ZX_32(_WX_32(upper27_lower5)) : _SX_32(_WX_32(upper27_lower5));
  new argument2 = PGR[registerP];
  for (I = 0; I < 2; I += 1) {
    result1.64[I] = argument2.64[I]
    if (_SX_64(argument3.64[I]) < 0) {
      result1.64[I] = _SAT_64(-argument2.64[I]);
    }
    if (_ZX_64(argument3.64[I]) == 0) {
      result1.64[I] = 0;
    }
  };
stage E1:
  PGR[registerM] = result1;
\end{lstlisting}}
\newpage
\subsubsection[{\small SIGNSHO: Apply Sign Saturated Half Word Octuple}]{SIGNSHO\hfill Apply Sign Saturated Half Word Octuple}
\label{instr:SIGNSHO}\index{signsho}
 
\paragraph{Encoding} Format \textsf{ALU\_HOWRR0} (Section~\ref{sec:format-ALU-HOWRR0}), \verb|exucode2 = 0001|\\

\bitfields{ 1/P, 2/11, 1/1, 4/0001, 5/registerM, 1/--, 2/10, 3/000, 1/1, 5/registerO, 1/--, 5/registerP, 1/0 }


\paragraph{Syntax} \texttt{signsho} \emph{registerM} = \emph{registerP}, \emph{registerO}

\paragraph{Description} The \emph{registerO} and the \emph{registerP} are interpreted as eight signed integer half words packed into 128 bits. The \emph{registerP} half words are multiplied by the signs of the \emph{registerO} half words and 16-bit saturated. The eight resulting half words are packed and stored into the \emph{registerM}.


\paragraph{Execution} {Reservation table \textsf{ALU\_LITE} (Section~\ref{sec:reservation-ALU-LITE})}
~
{\begin{lstlisting}
stage RR:
  new argument3 = PGR[registerO];
  new argument2 = PGR[registerP];
  for (I = 0; I < 8; I += 1) {
    result1.16[I] = argument2.16[I]
    if (_SX_16(argument3.16[I]) < 0) {
      result1.16[I] = _SAT_16(-argument2.16[I]);
    }
    if (_ZX_16(argument3.16[I]) == 0) {
      result1.16[I] = 0;
    }
  };
stage E1:
  PGR[registerM] = result1;
\end{lstlisting}}
\newpage

\paragraph{Encoding} Format \textsf{ALU\_HOWRR0.M} (Section~\ref{sec:format-ALU-HOWRR0.M}), \verb|exucode2 = 0001|\\

\bitfields{ 1/P, 2/00, 2/-- --, 27/upper27, 1/1, 2/11, 1/1, 4/0001, 5/registerM, 1/--, 2/10, 3/000, 1/1, 1/splat32, 5/lower5, 5/registerP, 1/0 }


\paragraph{Syntax} \texttt{signsho} \emph{registerM} = \emph{registerP}, \emph{upper27\_\discretionary{}{}{}lower5}\emph{splat32}

\paragraph{Description} The \emph{upper27\_\discretionary{}{}{}lower5} and the \emph{registerP} are interpreted as eight signed integer half words packed into 128 bits. The \emph{registerP} half words are multiplied by the signs of the \emph{upper27\_\discretionary{}{}{}lower5} half words and 16-bit saturated. The eight resulting half words are packed and stored into the \emph{registerM}.


\paragraph{Modifier(s)}
\noindent\begin{itemize}
\item splat32 (cf. splat32 in section \ref{par:modifier-splat32} on page \pageref{par:modifier-splat32})
\end{itemize}

\paragraph{Execution} {Reservation table \textsf{ALU\_LITE.X} (Section~\ref{sec:reservation-ALU-LITE.X})}
~
{\begin{lstlisting}
stage ID:
  new splat32 = _ZX_1(splat32);
stage RR:
  new argument3 = splat32 ? (_WX_32(upper27_lower5) << 32) | _ZX_32(_WX_32(upper27_lower5)) : _SX_32(_WX_32(upper27_lower5));
  new argument2 = PGR[registerP];
  for (I = 0; I < 8; I += 1) {
    result1.16[I] = argument2.16[I]
    if (_SX_16(argument3.16[I]) < 0) {
      result1.16[I] = _SAT_16(-argument2.16[I]);
    }
    if (_ZX_16(argument3.16[I]) == 0) {
      result1.16[I] = 0;
    }
  };
stage E1:
  PGR[registerM] = result1;
\end{lstlisting}}
\newpage
\subsubsection[{\small SIGNSW: Apply Sign Saturated Word}]{SIGNSW\hfill Apply Sign Saturated Word}
\label{instr:SIGNSW}\index{signsw}
 
\paragraph{Encoding} Format \textsf{ALU\_WWRR0} (Section~\ref{sec:format-ALU-WWRR0}), \verb|exucode2 = 0001|\\

\bitfields{ 1/P, 2/11, 1/0, 4/0001, 6/registerW, 2/11, 3/000, 1/signextw, 6/registerY, 6/registerZ }


\paragraph{Syntax} \texttt{signsw}\emph{signextw} \emph{registerW} = \emph{registerZ}, \emph{registerY}

\paragraph{Description} The \emph{registerY} and the \emph{registerZ} are interpreted as integer signed words. The \emph{registerZ} word is multiplied by the sign of the \emph{registerY} word and 32-bit saturated. The resulting word is stored into the \emph{registerW}.


\paragraph{Modifier(s)}
\noindent\begin{itemize}
\item signextw (cf. signextw in section \ref{par:modifier-signextw} on page \pageref{par:modifier-signextw})
\end{itemize}

\paragraph{Execution} {Reservation table \textsf{ALU\_TINY} (Section~\ref{sec:reservation-ALU-TINY})}
~
{\begin{lstlisting}
stage ID:
  new signextw = _ZX_1(signextw);
stage RR:
  new argument3 = GPR[registerY];
  new argument2 = GPR[registerZ];
  new result1 = argument2;
  if (_SX_32(argument3) < 0) {
    result1 = _SAT_32(-argument2);
  }
  if (_ZX_32(argument3) == 0) {
    result1 = 0;
  }
stage E1:
  GPR[registerW] = signextw ? _SX_32(result1) : _ZX_32(result1);
\end{lstlisting}}
\newpage

\paragraph{Encoding} Format \textsf{ALU\_WWRR0.W} (Section~\ref{sec:format-ALU-WWRR0.W}), \verb|exucode2 = 0001|\\

\bitfields{ 1/P, 2/00, 2/-- --, 27/upper27, 1/1, 2/11, 1/0, 4/0001, 6/registerW, 2/11, 3/000, 1/signextw, 1/0, 5/lower5, 6/registerZ }


\paragraph{Syntax} \texttt{signsw}\emph{signextw} \emph{registerW} = \emph{registerZ}, \emph{upper27\_\discretionary{}{}{}lower5}

\paragraph{Description} The \emph{upper27\_\discretionary{}{}{}lower5} and the \emph{registerZ} are interpreted as integer signed words. The \emph{registerZ} word is multiplied by the sign of the \emph{upper27\_\discretionary{}{}{}lower5} word and 32-bit saturated. The resulting word is stored into the \emph{registerW}.


\paragraph{Modifier(s)}
\noindent\begin{itemize}
\item signextw (cf. signextw in section \ref{par:modifier-signextw} on page \pageref{par:modifier-signextw})
\end{itemize}

\paragraph{Execution} {Reservation table \textsf{ALU\_TINY.X} (Section~\ref{sec:reservation-ALU-TINY.X})}
~
{\begin{lstlisting}
stage ID:
  new signextw = _ZX_1(signextw);
stage RR:
  new argument3 = _WX_32(upper27_lower5);
  new argument2 = GPR[registerZ];
  new result1 = argument2;
  if (_SX_32(argument3) < 0) {
    result1 = _SAT_32(-argument2);
  }
  if (_ZX_32(argument3) == 0) {
    result1 = 0;
  }
stage E1:
  GPR[registerW] = signextw ? _SX_32(result1) : _ZX_32(result1);
\end{lstlisting}}
\newpage
\subsubsection[{\small SIGNSWQ: Apply Sign Saturated Word Quadruple}]{SIGNSWQ\hfill Apply Sign Saturated Word Quadruple}
\label{instr:SIGNSWQ}\index{signswq}
 
\paragraph{Encoding} Format \textsf{ALU\_WQWRR0} (Section~\ref{sec:format-ALU-WQWRR0}), \verb|exucode2 = 0001|\\

\bitfields{ 1/P, 2/11, 1/1, 4/0001, 5/registerM, 1/--, 2/10, 3/000, 1/0, 5/registerO, 1/--, 5/registerP, 1/1 }


\paragraph{Syntax} \texttt{signswq} \emph{registerM} = \emph{registerP}, \emph{registerO}

\paragraph{Description} The \emph{registerO} and the \emph{registerP} are interpreted as four signed integer words packed into 128 bits. The \emph{registerP} words are multiplied by the signs of the \emph{registerO} words and 32-bit saturated. The four resulting words are packed and stored into the \emph{registerM}.


\paragraph{Execution} {Reservation table \textsf{ALU\_LITE} (Section~\ref{sec:reservation-ALU-LITE})}
~
{\begin{lstlisting}
stage RR:
  new argument3 = PGR[registerO];
  new argument2 = PGR[registerP];
  for (I = 0; I < 4; I += 1) {
    result1.32[I] = argument2.32[I]
    if (_SX_32(argument3.32[I]) < 0) {
      result1.32[I] = _SAT_32(-argument2.32[I]);
    }
    if (_ZX_32(argument3.32[I]) == 0) {
      result1.32[I] = 0;
    }
  };
stage E1:
  PGR[registerM] = result1;
\end{lstlisting}}
\newpage

\paragraph{Encoding} Format \textsf{ALU\_WQWRR0.M} (Section~\ref{sec:format-ALU-WQWRR0.M}), \verb|exucode2 = 0001|\\

\bitfields{ 1/P, 2/00, 2/-- --, 27/upper27, 1/1, 2/11, 1/1, 4/0001, 5/registerM, 1/--, 2/10, 3/000, 1/0, 1/splat32, 5/lower5, 5/registerP, 1/1 }


\paragraph{Syntax} \texttt{signswq} \emph{registerM} = \emph{registerP}, \emph{upper27\_\discretionary{}{}{}lower5}\emph{splat32}

\paragraph{Description} The \emph{upper27\_\discretionary{}{}{}lower5} and the \emph{registerP} are interpreted as four signed integer words packed into 128 bits. The \emph{registerP} words are multiplied by the signs of the \emph{upper27\_\discretionary{}{}{}lower5} words and 32-bit saturated. The four resulting words are packed and stored into the \emph{registerM}.


\paragraph{Modifier(s)}
\noindent\begin{itemize}
\item splat32 (cf. splat32 in section \ref{par:modifier-splat32} on page \pageref{par:modifier-splat32})
\end{itemize}

\paragraph{Execution} {Reservation table \textsf{ALU\_LITE.X} (Section~\ref{sec:reservation-ALU-LITE.X})}
~
{\begin{lstlisting}
stage ID:
  new splat32 = _ZX_1(splat32);
stage RR:
  new argument3 = splat32 ? (_WX_32(upper27_lower5) << 32) | _ZX_32(_WX_32(upper27_lower5)) : _SX_32(_WX_32(upper27_lower5));
  new argument2 = PGR[registerP];
  for (I = 0; I < 4; I += 1) {
    result1.32[I] = argument2.32[I]
    if (_SX_32(argument3.32[I]) < 0) {
      result1.32[I] = _SAT_32(-argument2.32[I]);
    }
    if (_ZX_32(argument3.32[I]) == 0) {
      result1.32[I] = 0;
    }
  };
stage E1:
  PGR[registerM] = result1;
\end{lstlisting}}
\newpage
\subsubsection[{\small SIGNW: Apply Sign Word}]{SIGNW\hfill Apply Sign Word}
\label{instr:SIGNW}\index{signw}
 
\paragraph{Encoding} Format \textsf{ALU\_WWRR0} (Section~\ref{sec:format-ALU-WWRR0}), \verb|exucode2 = 0000|\\

\bitfields{ 1/P, 2/11, 1/0, 4/0000, 6/registerW, 2/11, 3/000, 1/signextw, 6/registerY, 6/registerZ }


\paragraph{Syntax} \texttt{signw}\emph{signextw} \emph{registerW} = \emph{registerZ}, \emph{registerY}

\paragraph{Description} The \emph{registerY} and the \emph{registerZ} are interpreted as integer signed words. The \emph{registerZ} word is multiplied by the sign of the \emph{registerY} word. The resulting word is stored into the \emph{registerW}.


\paragraph{Modifier(s)}
\noindent\begin{itemize}
\item signextw (cf. signextw in section \ref{par:modifier-signextw} on page \pageref{par:modifier-signextw})
\end{itemize}

\paragraph{Execution} {Reservation table \textsf{ALU\_TINY} (Section~\ref{sec:reservation-ALU-TINY})}
~
{\begin{lstlisting}
stage ID:
  new signextw = _ZX_1(signextw);
stage RR:
  new argument3 = GPR[registerY];
  new argument2 = GPR[registerZ];
  new result1 = argument2;
  if (_SX_32(argument3) < 0) {
    result1 = -argument2;
  }
  if (_ZX_32(argument3) == 0) {
    result1 = 0;
  }
stage E1:
  GPR[registerW] = signextw ? _SX_32(result1) : _ZX_32(result1);
\end{lstlisting}}
\newpage

\paragraph{Encoding} Format \textsf{ALU\_WWRR0.W} (Section~\ref{sec:format-ALU-WWRR0.W}), \verb|exucode2 = 0000|\\

\bitfields{ 1/P, 2/00, 2/-- --, 27/upper27, 1/1, 2/11, 1/0, 4/0000, 6/registerW, 2/11, 3/000, 1/signextw, 1/0, 5/lower5, 6/registerZ }


\paragraph{Syntax} \texttt{signw}\emph{signextw} \emph{registerW} = \emph{registerZ}, \emph{upper27\_\discretionary{}{}{}lower5}

\paragraph{Description} The \emph{upper27\_\discretionary{}{}{}lower5} and the \emph{registerZ} are interpreted as integer signed words. The \emph{registerZ} word is multiplied by the sign of the \emph{upper27\_\discretionary{}{}{}lower5} word. The resulting word is stored into the \emph{registerW}.


\paragraph{Modifier(s)}
\noindent\begin{itemize}
\item signextw (cf. signextw in section \ref{par:modifier-signextw} on page \pageref{par:modifier-signextw})
\end{itemize}

\paragraph{Execution} {Reservation table \textsf{ALU\_TINY.X} (Section~\ref{sec:reservation-ALU-TINY.X})}
~
{\begin{lstlisting}
stage ID:
  new signextw = _ZX_1(signextw);
stage RR:
  new argument3 = _WX_32(upper27_lower5);
  new argument2 = GPR[registerZ];
  new result1 = argument2;
  if (_SX_32(argument3) < 0) {
    result1 = -argument2;
  }
  if (_ZX_32(argument3) == 0) {
    result1 = 0;
  }
stage E1:
  GPR[registerW] = signextw ? _SX_32(result1) : _ZX_32(result1);
\end{lstlisting}}
\newpage
\subsubsection[{\small SIGNWQ: Apply Sign Word Quadruple}]{SIGNWQ\hfill Apply Sign Word Quadruple}
\label{instr:SIGNWQ}\index{signwq}
 
\paragraph{Encoding} Format \textsf{ALU\_WQWRR0} (Section~\ref{sec:format-ALU-WQWRR0}), \verb|exucode2 = 0000|\\

\bitfields{ 1/P, 2/11, 1/1, 4/0000, 5/registerM, 1/--, 2/10, 3/000, 1/0, 5/registerO, 1/--, 5/registerP, 1/1 }


\paragraph{Syntax} \texttt{signwq} \emph{registerM} = \emph{registerP}, \emph{registerO}

\paragraph{Description} The \emph{registerO} and the \emph{registerP} are interpreted as four signed integer words packed into 128 bits. The \emph{registerP} words are multiplied by the signs of the \emph{registerO} words. The four resulting words are packed and stored into the \emph{registerM}.


\paragraph{Execution} {Reservation table \textsf{ALU\_LITE} (Section~\ref{sec:reservation-ALU-LITE})}
~
{\begin{lstlisting}
stage RR:
  new argument3 = PGR[registerO];
  new argument2 = PGR[registerP];
  for (I = 0; I < 4; I += 1) {
    result1.32[I] = argument2.32[I]
    if (_SX_32(argument3.32[I]) < 0) {
      result1.32[I] = -argument2.32[I];
    }
    if (_ZX_32(argument3.32[I]) == 0) {
      result1.32[I] = 0;
    }
  };
stage E1:
  PGR[registerM] = result1;
\end{lstlisting}}
\newpage

\paragraph{Encoding} Format \textsf{ALU\_WQWRR0.M} (Section~\ref{sec:format-ALU-WQWRR0.M}), \verb|exucode2 = 0000|\\

\bitfields{ 1/P, 2/00, 2/-- --, 27/upper27, 1/1, 2/11, 1/1, 4/0000, 5/registerM, 1/--, 2/10, 3/000, 1/0, 1/splat32, 5/lower5, 5/registerP, 1/1 }


\paragraph{Syntax} \texttt{signwq} \emph{registerM} = \emph{registerP}, \emph{upper27\_\discretionary{}{}{}lower5}\emph{splat32}

\paragraph{Description} The \emph{upper27\_\discretionary{}{}{}lower5} and the \emph{registerP} are interpreted as four signed integer words packed into 128 bits. The \emph{registerP} words are multiplied by the signs of the \emph{upper27\_\discretionary{}{}{}lower5} words. The four resulting words are packed and stored into the \emph{registerM}.


\paragraph{Modifier(s)}
\noindent\begin{itemize}
\item splat32 (cf. splat32 in section \ref{par:modifier-splat32} on page \pageref{par:modifier-splat32})
\end{itemize}

\paragraph{Execution} {Reservation table \textsf{ALU\_LITE.X} (Section~\ref{sec:reservation-ALU-LITE.X})}
~
{\begin{lstlisting}
stage ID:
  new splat32 = _ZX_1(splat32);
stage RR:
  new argument3 = splat32 ? (_WX_32(upper27_lower5) << 32) | _ZX_32(_WX_32(upper27_lower5)) : _SX_32(_WX_32(upper27_lower5));
  new argument2 = PGR[registerP];
  for (I = 0; I < 4; I += 1) {
    result1.32[I] = argument2.32[I]
    if (_SX_32(argument3.32[I]) < 0) {
      result1.32[I] = -argument2.32[I];
    }
    if (_ZX_32(argument3.32[I]) == 0) {
      result1.32[I] = 0;
    }
  };
stage E1:
  PGR[registerM] = result1;
\end{lstlisting}}
\newpage
\subsubsection[{\small SLLBX: Shift Left Logical Byte Hexadecuple}]{SLLBX\hfill Shift Left Logical Byte Hexadecuple}
\label{instr:SLLBX}\index{sllbx}
 
\paragraph{Encoding} Format \textsf{ALU\_BXSWRI} (Section~\ref{sec:format-ALU-BXSWRI}), \verb|exucode1 = 001|\\

\bitfields{ 1/P, 2/11, 1/1, 1/1, 3/001, 5/registerM, 1/--, 2/10, 3/100, 1/1, 6/unsigned6, 5/registerP, 1/1 }


\paragraph{Syntax} \texttt{sllbx} \emph{registerM} = \emph{registerP}, \emph{unsigned6}

\paragraph{Description} The \emph{registerP} is interpreted as sixteen bytes packed into 128 bits. The \emph{registerP} bytes are shifted left by the \emph{unsigned6} modulo 8. The vacant positions are filled in with zeros. The sixteen resulting bytes are packed and stored into the \emph{registerM}.


\paragraph{Execution} {Reservation table \textsf{ALU\_LITE} (Section~\ref{sec:reservation-ALU-LITE})}
~
{\begin{lstlisting}
stage RR:
  new argument3 = _ZX_6(unsigned6);
  new argument2 = PGR[registerP];
  new shift = _ZX_3(argument3);
  for (I = 0; I < 16; I += 1) {
    result1.8[I] = argument2.8[I] << shift
  };
stage E1:
  PGR[registerM] = result1;
\end{lstlisting}}
\newpage

\paragraph{Encoding} Format \textsf{ALU\_BXSWRR} (Section~\ref{sec:format-ALU-BXSWRR}), \verb|exucode1 = 001|\\

\bitfields{ 1/P, 2/11, 1/1, 1/0, 3/001, 5/registerM, 1/--, 2/10, 3/100, 1/1, 6/registerY, 5/registerP, 1/1 }


\paragraph{Syntax} \texttt{sllbx} \emph{registerM} = \emph{registerP}, \emph{registerY}

\paragraph{Description} The \emph{registerP} is interpreted as sixteen bytes packed into 128 bits. The \emph{registerP} bytes are shifted left by the \emph{registerY} modulo 8. The vacant positions are filled in with zeros. The sixteen resulting bytes are packed and stored into the \emph{registerM}.


\paragraph{Execution} {Reservation table \textsf{ALU\_LITE} (Section~\ref{sec:reservation-ALU-LITE})}
~
{\begin{lstlisting}
stage RR:
  new argument3 = GPR[registerY];
  new argument2 = PGR[registerP];
  new shift = _ZX_3(argument3);
  for (I = 0; I < 16; I += 1) {
    result1.8[I] = argument2.8[I] << shift
  };
stage E1:
  PGR[registerM] = result1;
\end{lstlisting}}
\newpage
\subsubsection[{\small SLLD: Shift Left Logical Double Word}]{SLLD\hfill Shift Left Logical Double Word}
\label{instr:SLLD}\index{slld}
 
\paragraph{Encoding} Format \textsf{ALU\_DSWRI} (Section~\ref{sec:format-ALU-DSWRI}), \verb|exucode1 = 001|\\

\bitfields{ 1/P, 2/11, 1/0, 1/1, 3/001, 6/registerW, 2/10, 3/100, 1/0, 6/unsigned6, 6/registerZ }


\paragraph{Syntax} \texttt{slld} \emph{registerW} = \emph{registerZ}, \emph{unsigned6}

\paragraph{Description} The \emph{registerZ} is shifted left by the \emph{unsigned6} modulo 64. The vacant positions are filled in with zeros. The result is stored into the \emph{registerW}.


\paragraph{Execution} {Reservation table \textsf{ALU\_TINY} (Section~\ref{sec:reservation-ALU-TINY})}
~
{\begin{lstlisting}
stage RR:
  new argument3 = _ZX_6(unsigned6);
  new argument2 = GPR[registerZ];
  new result1 = _ZX_64(argument2) << _ZX_6(argument3);
stage E1:
  GPR[registerW] = result1;
\end{lstlisting}}
\newpage

\paragraph{Encoding} Format \textsf{ALU\_DSWRR} (Section~\ref{sec:format-ALU-DSWRR}), \verb|exucode1 = 001|\\

\bitfields{ 1/P, 2/11, 1/0, 1/0, 3/001, 6/registerW, 2/10, 3/100, 1/0, 6/registerY, 6/registerZ }


\paragraph{Syntax} \texttt{slld} \emph{registerW} = \emph{registerZ}, \emph{registerY}

\paragraph{Description} The \emph{registerZ} is shifted left by the \emph{registerY} modulo 64. The vacant positions are filled in with zeros. The result is stored into the \emph{registerW}.


\paragraph{Execution} {Reservation table \textsf{ALU\_TINY} (Section~\ref{sec:reservation-ALU-TINY})}
~
{\begin{lstlisting}
stage RR:
  new argument3 = GPR[registerY];
  new argument2 = GPR[registerZ];
  new result1 = _ZX_64(argument2) << _ZX_6(argument3);
stage E1:
  GPR[registerW] = result1;
\end{lstlisting}}
\newpage
\subsubsection[{\small SLLDP: Shift Left Logical Word Pair by Scalar}]{SLLDP\hfill Shift Left Logical Word Pair by Scalar}
\label{instr:SLLDP}\index{slldp}
 
\paragraph{Encoding} Format \textsf{ALU\_DPSWRI} (Section~\ref{sec:format-ALU-DPSWRI}), \verb|exucode1 = 001|\\

\bitfields{ 1/P, 2/11, 1/1, 1/1, 3/001, 5/registerM, 1/--, 2/10, 3/100, 1/0, 6/unsigned6, 5/registerP, 1/0 }


\paragraph{Syntax} \texttt{slldp} \emph{registerM} = \emph{registerP}, \emph{unsigned6}

\paragraph{Description} The \emph{registerP} is interpreted as two double words packed into 128 bits. The \emph{registerP} double words are shifted left by the \emph{unsigned6} modulo 64. The vacant positions are filled in with zeros. The two resulting double words are packed and stored into the \emph{registerM}.


\paragraph{Execution} {Reservation table \textsf{ALU\_LITE} (Section~\ref{sec:reservation-ALU-LITE})}
~
{\begin{lstlisting}
stage RR:
  new argument3 = _ZX_6(unsigned6);
  new argument2 = PGR[registerP];
  new shift = _ZX_6(argument3);
  for (I = 0; I < 2; I += 1) {
    result1.64[I] = argument2.64[I] << shift
  };
stage E1:
  PGR[registerM] = result1;
\end{lstlisting}}
\newpage

\paragraph{Encoding} Format \textsf{ALU\_DPSWRR} (Section~\ref{sec:format-ALU-DPSWRR}), \verb|exucode1 = 001|\\

\bitfields{ 1/P, 2/11, 1/1, 1/0, 3/001, 5/registerM, 1/--, 2/10, 3/100, 1/0, 5/registerO, 1/--, 5/registerP, 1/0 }


\paragraph{Syntax} \texttt{slldp} \emph{registerM} = \emph{registerP}, \emph{registerO}

\paragraph{Description} The \emph{registerP} is interpreted as two double words packed into 128 bits. The \emph{registerP} double words are shifted left by the \emph{registerO} modulo 64. The vacant positions are filled in with zeros. The two resulting double words are packed and stored into the \emph{registerM}.


\paragraph{Execution} {Reservation table \textsf{ALU\_LITE} (Section~\ref{sec:reservation-ALU-LITE})}
~
{\begin{lstlisting}
stage RR:
  new argument3 = PGR[registerO];
  new argument2 = PGR[registerP];
  new shift = _ZX_6(argument3);
  for (I = 0; I < 2; I += 1) {
    result1.64[I] = argument2.64[I] << shift
  };
stage E1:
  PGR[registerM] = result1;
\end{lstlisting}}
\newpage
\subsubsection[{\small SLLHO: Shift Left Logical Half Word Octuple}]{SLLHO\hfill Shift Left Logical Half Word Octuple}
\label{instr:SLLHO}\index{sllho}
 
\paragraph{Encoding} Format \textsf{ALU\_HOSWRI} (Section~\ref{sec:format-ALU-HOSWRI}), \verb|exucode1 = 001|\\

\bitfields{ 1/P, 2/11, 1/1, 1/1, 3/001, 5/registerM, 1/--, 2/10, 3/100, 1/1, 6/unsigned6, 5/registerP, 1/0 }


\paragraph{Syntax} \texttt{sllho} \emph{registerM} = \emph{registerP}, \emph{unsigned6}

\paragraph{Description} The \emph{registerP} is interpreted as eight half words packed into 128 bits. The \emph{registerP} half words are shifted left by the \emph{unsigned6} modulo 16. The vacant positions are filled in with zeros. The eight resulting half words are packed and stored into the \emph{registerM}.


\paragraph{Execution} {Reservation table \textsf{ALU\_LITE} (Section~\ref{sec:reservation-ALU-LITE})}
~
{\begin{lstlisting}
stage RR:
  new argument3 = _ZX_6(unsigned6);
  new argument2 = PGR[registerP];
  new shift = _ZX_4(argument3);
  for (I = 0; I < 8; I += 1) {
    result1.16[I] = argument2.16[I] << shift
  };
stage E1:
  PGR[registerM] = result1;
\end{lstlisting}}
\newpage

\paragraph{Encoding} Format \textsf{ALU\_HOSWRR} (Section~\ref{sec:format-ALU-HOSWRR}), \verb|exucode1 = 001|\\

\bitfields{ 1/P, 2/11, 1/1, 1/0, 3/001, 5/registerM, 1/--, 2/10, 3/100, 1/1, 6/registerY, 5/registerP, 1/0 }


\paragraph{Syntax} \texttt{sllho} \emph{registerM} = \emph{registerP}, \emph{registerY}

\paragraph{Description} The \emph{registerP} is interpreted as eight half words packed into 128 bits. The \emph{registerP} half words are shifted left by the \emph{registerY} modulo 16. The vacant positions are filled in with zeros. The eight resulting half words are packed and stored into the \emph{registerM}.


\paragraph{Execution} {Reservation table \textsf{ALU\_LITE} (Section~\ref{sec:reservation-ALU-LITE})}
~
{\begin{lstlisting}
stage RR:
  new argument3 = GPR[registerY];
  new argument2 = PGR[registerP];
  new shift = _ZX_4(argument3);
  for (I = 0; I < 8; I += 1) {
    result1.16[I] = argument2.16[I] << shift
  };
stage E1:
  PGR[registerM] = result1;
\end{lstlisting}}
\newpage
\subsubsection[{\small SLLW: Shift Left Logical Word}]{SLLW\hfill Shift Left Logical Word}
\label{instr:SLLW}\index{sllw}
 
\paragraph{Encoding} Format \textsf{ALU\_WSWRR} (Section~\ref{sec:format-ALU-WSWRR}), \verb|exucode1 = 001|\\

\bitfields{ 1/P, 2/11, 1/0, 1/0, 3/001, 6/registerW, 2/11, 3/100, 1/signextw, 6/registerY, 6/registerZ }


\paragraph{Syntax} \texttt{sllw}\emph{signextw} \emph{registerW} = \emph{registerZ}, \emph{registerY}

\paragraph{Description} The \emph{registerZ} is shifted left by the \emph{registerY} modulo 32. The vacant positions are filled in with zeros. The result with the 32 upper bits cleared is stored into the \emph{registerW}.


\paragraph{Modifier(s)}
\noindent\begin{itemize}
\item signextw (cf. signextw in section \ref{par:modifier-signextw} on page \pageref{par:modifier-signextw})
\end{itemize}

\paragraph{Execution} {Reservation table \textsf{ALU\_TINY} (Section~\ref{sec:reservation-ALU-TINY})}
~
{\begin{lstlisting}
stage ID:
  new signextw = _ZX_1(signextw);
stage RR:
  new argument3 = GPR[registerY];
  new argument2 = GPR[registerZ];
  new result1 = argument2.32[0] << _ZX_5(argument3);
stage E1:
  GPR[registerW] = signextw ? _SX_32(result1) : _ZX_32(result1);
\end{lstlisting}}
\newpage

\paragraph{Encoding} Format \textsf{ALU\_WSWRI} (Section~\ref{sec:format-ALU-WSWRI}), \verb|exucode1 = 001|\\

\bitfields{ 1/P, 2/11, 1/0, 1/1, 3/001, 6/registerW, 2/11, 3/100, 1/signextw, 6/unsigned6, 6/registerZ }


\paragraph{Syntax} \texttt{sllw}\emph{signextw} \emph{registerW} = \emph{registerZ}, \emph{unsigned6}

\paragraph{Description} The \emph{registerZ} is shifted left by the \emph{unsigned6} modulo 32. The vacant positions are filled in with zeros. The result with the 32 upper bits cleared is stored into the \emph{registerW}.


\paragraph{Modifier(s)}
\noindent\begin{itemize}
\item signextw (cf. signextw in section \ref{par:modifier-signextw} on page \pageref{par:modifier-signextw})
\end{itemize}

\paragraph{Execution} {Reservation table \textsf{ALU\_TINY} (Section~\ref{sec:reservation-ALU-TINY})}
~
{\begin{lstlisting}
stage ID:
  new signextw = _ZX_1(signextw);
stage RR:
  new argument3 = _ZX_6(unsigned6);
  new argument2 = GPR[registerZ];
  new result1 = argument2.32[0] << _ZX_5(argument3);
stage E1:
  GPR[registerW] = signextw ? _SX_32(result1) : _ZX_32(result1);
\end{lstlisting}}
\newpage
\subsubsection[{\small SLLWQ: Shift Left Logical Word Quadruple}]{SLLWQ\hfill Shift Left Logical Word Quadruple}
\label{instr:SLLWQ}\index{sllwq}
 
\paragraph{Encoding} Format \textsf{ALU\_WQSWRI} (Section~\ref{sec:format-ALU-WQSWRI}), \verb|exucode1 = 001|\\

\bitfields{ 1/P, 2/11, 1/1, 1/1, 3/001, 5/registerM, 1/--, 2/10, 3/100, 1/0, 6/unsigned6, 5/registerP, 1/1 }


\paragraph{Syntax} \texttt{sllwq} \emph{registerM} = \emph{registerP}, \emph{unsigned6}

\paragraph{Description} The \emph{registerP} is interpreted as four words packed into 128 bits. The \emph{registerP} words are shifted left by the \emph{unsigned6} modulo 32. The vacant positions are filled in with zeros. The four resulting words are packed and stored into the \emph{registerM}.


\paragraph{Execution} {Reservation table \textsf{ALU\_LITE} (Section~\ref{sec:reservation-ALU-LITE})}
~
{\begin{lstlisting}
stage RR:
  new argument3 = _ZX_6(unsigned6);
  new argument2 = PGR[registerP];
  new shift = _ZX_5(argument3);
  for (I = 0; I < 4; I += 1) {
    result1.32[I] = argument2.32[I] << shift
  };
stage E1:
  PGR[registerM] = result1;
\end{lstlisting}}
\newpage

\paragraph{Encoding} Format \textsf{ALU\_WQSWRR} (Section~\ref{sec:format-ALU-WQSWRR}), \verb|exucode1 = 001|\\

\bitfields{ 1/P, 2/11, 1/1, 1/0, 3/001, 5/registerM, 1/--, 2/10, 3/100, 1/0, 6/registerY, 5/registerP, 1/1 }


\paragraph{Syntax} \texttt{sllwq} \emph{registerM} = \emph{registerP}, \emph{registerY}

\paragraph{Description} The \emph{registerP} is interpreted as four words packed into 128 bits. The \emph{registerP} words are shifted left by the \emph{registerY} modulo 32. The vacant positions are filled in with zeros. The four resulting words are packed and stored into the \emph{registerM}.


\paragraph{Execution} {Reservation table \textsf{ALU\_LITE} (Section~\ref{sec:reservation-ALU-LITE})}
~
{\begin{lstlisting}
stage RR:
  new argument3 = GPR[registerY];
  new argument2 = PGR[registerP];
  new shift = _ZX_5(argument3);
  for (I = 0; I < 4; I += 1) {
    result1.32[I] = argument2.32[I] << shift
  };
stage E1:
  PGR[registerM] = result1;
\end{lstlisting}}
\newpage
\subsubsection[{\small SLSBX: Shift Left Saturated Byte Hexadecuple}]{SLSBX\hfill Shift Left Saturated Byte Hexadecuple}
\label{instr:SLSBX}\index{slsbx}
 
\paragraph{Encoding} Format \textsf{ALU\_BXSWRI} (Section~\ref{sec:format-ALU-BXSWRI}), \verb|exucode1 = 100|\\

\bitfields{ 1/P, 2/11, 1/1, 1/1, 3/100, 5/registerM, 1/--, 2/10, 3/100, 1/1, 6/unsigned6, 5/registerP, 1/1 }


\paragraph{Syntax} \texttt{slsbx} \emph{registerM} = \emph{registerP}, \emph{unsigned6}

\paragraph{Description} The \emph{registerP} is interpreted as sixteen bytes packed into 128 bits. The \emph{registerP} bytes are shifted left by the \emph{unsigned6} modulo 8. The sixteen resulting bytes are saturated to 8 bits, packed and stored into the \emph{registerM}.


\paragraph{Execution} {Reservation table \textsf{ALU\_LITE} (Section~\ref{sec:reservation-ALU-LITE})}
~
{\begin{lstlisting}
stage RR:
  new argument3 = _ZX_6(unsigned6);
  new argument2 = PGR[registerP];
  new shift = _ZX_3(argument3);
  for (I = 0; I < 16; I += 1) {
    result1.8[I] = _SAT_8(_SX_8(argument2.8[I]) << shift)
  };
stage E1:
  PGR[registerM] = result1;
\end{lstlisting}}
\newpage

\paragraph{Encoding} Format \textsf{ALU\_BXSWRR} (Section~\ref{sec:format-ALU-BXSWRR}), \verb|exucode1 = 100|\\

\bitfields{ 1/P, 2/11, 1/1, 1/0, 3/100, 5/registerM, 1/--, 2/10, 3/100, 1/1, 6/registerY, 5/registerP, 1/1 }


\paragraph{Syntax} \texttt{slsbx} \emph{registerM} = \emph{registerP}, \emph{registerY}

\paragraph{Description} The \emph{registerP} is interpreted as sixteen bytes packed into 128 bits. The \emph{registerP} bytes are shifted left by the \emph{registerY} modulo 8. The sixteen resulting bytes are saturated to 8 bits, packed and stored into the \emph{registerM}.


\paragraph{Execution} {Reservation table \textsf{ALU\_LITE} (Section~\ref{sec:reservation-ALU-LITE})}
~
{\begin{lstlisting}
stage RR:
  new argument3 = GPR[registerY];
  new argument2 = PGR[registerP];
  new shift = _ZX_3(argument3);
  for (I = 0; I < 16; I += 1) {
    result1.8[I] = _SAT_8(_SX_8(argument2.8[I]) << shift)
  };
stage E1:
  PGR[registerM] = result1;
\end{lstlisting}}
\newpage
\subsubsection[{\small SLSD: Shift Left Saturated Double Word}]{SLSD\hfill Shift Left Saturated Double Word}
\label{instr:SLSD}\index{slsd}
 
\paragraph{Encoding} Format \textsf{ALU\_DSWRI} (Section~\ref{sec:format-ALU-DSWRI}), \verb|exucode1 = 100|\\

\bitfields{ 1/P, 2/11, 1/0, 1/1, 3/100, 6/registerW, 2/10, 3/100, 1/0, 6/unsigned6, 6/registerZ }


\paragraph{Syntax} \texttt{slsd} \emph{registerW} = \emph{registerZ}, \emph{unsigned6}

\paragraph{Description} The \emph{registerZ} is shifted left by the \emph{unsigned6} modulo 64. The result is saturated to 64 bits and stored into the \emph{registerW}.


\paragraph{Execution} {Reservation table \textsf{ALU\_LITE} (Section~\ref{sec:reservation-ALU-LITE})}
~
{\begin{lstlisting}
stage RR:
  new argument3 = _ZX_6(unsigned6);
  new argument2 = GPR[registerZ];
  new shift = _ZX_6(argument3);
  new result1 = _SAT_64(_SX_64(argument2) << shift);
stage E1:
  GPR[registerW] = result1;
\end{lstlisting}}
\newpage

\paragraph{Encoding} Format \textsf{ALU\_DSWRR} (Section~\ref{sec:format-ALU-DSWRR}), \verb|exucode1 = 100|\\

\bitfields{ 1/P, 2/11, 1/0, 1/0, 3/100, 6/registerW, 2/10, 3/100, 1/0, 6/registerY, 6/registerZ }


\paragraph{Syntax} \texttt{slsd} \emph{registerW} = \emph{registerZ}, \emph{registerY}

\paragraph{Description} The \emph{registerZ} is shifted left by the \emph{registerY} modulo 64. The result is saturated to 64 bits and stored into the \emph{registerW}.


\paragraph{Execution} {Reservation table \textsf{ALU\_LITE} (Section~\ref{sec:reservation-ALU-LITE})}
~
{\begin{lstlisting}
stage RR:
  new argument3 = GPR[registerY];
  new argument2 = GPR[registerZ];
  new shift = _ZX_6(argument3);
  new result1 = _SAT_64(_SX_64(argument2) << shift);
stage E1:
  GPR[registerW] = result1;
\end{lstlisting}}
\newpage
\subsubsection[{\small SLSDP: Shift Left Saturated Double Word Pair}]{SLSDP\hfill Shift Left Saturated Double Word Pair}
\label{instr:SLSDP}\index{slsdp}
 
\paragraph{Encoding} Format \textsf{ALU\_DPSWRI} (Section~\ref{sec:format-ALU-DPSWRI}), \verb|exucode1 = 100|\\

\bitfields{ 1/P, 2/11, 1/1, 1/1, 3/100, 5/registerM, 1/--, 2/10, 3/100, 1/0, 6/unsigned6, 5/registerP, 1/0 }


\paragraph{Syntax} \texttt{slsdp} \emph{registerM} = \emph{registerP}, \emph{unsigned6}

\paragraph{Description} The \emph{registerP} is interpreted as two double words packed into 128 bits. The \emph{registerP} words are shifted left by the \emph{unsigned6} modulo 64. The two resulting double words are saturated to 64 bits, packed and stored into the \emph{registerM}.


\paragraph{Execution} {Reservation table \textsf{ALU\_LITE} (Section~\ref{sec:reservation-ALU-LITE})}
~
{\begin{lstlisting}
stage RR:
  new argument3 = _ZX_6(unsigned6);
  new argument2 = PGR[registerP];
  new shift = _ZX_6(argument3);
  for (I = 0; I < 2; I += 1) {
    result1.64[I] = _SAT_64(_SX_64(argument2.64[I]) << shift)
  };
stage E1:
  PGR[registerM] = result1;
\end{lstlisting}}
\newpage

\paragraph{Encoding} Format \textsf{ALU\_DPSWRR} (Section~\ref{sec:format-ALU-DPSWRR}), \verb|exucode1 = 100|\\

\bitfields{ 1/P, 2/11, 1/1, 1/0, 3/100, 5/registerM, 1/--, 2/10, 3/100, 1/0, 5/registerO, 1/--, 5/registerP, 1/0 }


\paragraph{Syntax} \texttt{slsdp} \emph{registerM} = \emph{registerP}, \emph{registerO}

\paragraph{Description} The \emph{registerP} is interpreted as two double words packed into 128 bits. The \emph{registerP} words are shifted left by the \emph{registerO} modulo 64. The two resulting double words are saturated to 64 bits, packed and stored into the \emph{registerM}.


\paragraph{Execution} {Reservation table \textsf{ALU\_LITE} (Section~\ref{sec:reservation-ALU-LITE})}
~
{\begin{lstlisting}
stage RR:
  new argument3 = PGR[registerO];
  new argument2 = PGR[registerP];
  new shift = _ZX_6(argument3);
  for (I = 0; I < 2; I += 1) {
    result1.64[I] = _SAT_64(_SX_64(argument2.64[I]) << shift)
  };
stage E1:
  PGR[registerM] = result1;
\end{lstlisting}}
\newpage
\subsubsection[{\small SLSHO: Shift Left Saturated Half Word Octuple}]{SLSHO\hfill Shift Left Saturated Half Word Octuple}
\label{instr:SLSHO}\index{slsho}
 
\paragraph{Encoding} Format \textsf{ALU\_HOSWRI} (Section~\ref{sec:format-ALU-HOSWRI}), \verb|exucode1 = 100|\\

\bitfields{ 1/P, 2/11, 1/1, 1/1, 3/100, 5/registerM, 1/--, 2/10, 3/100, 1/1, 6/unsigned6, 5/registerP, 1/0 }


\paragraph{Syntax} \texttt{slsho} \emph{registerM} = \emph{registerP}, \emph{unsigned6}

\paragraph{Description} The \emph{registerP} is interpreted as eight half words packed into 128 bits. The \emph{registerP} half words are shifted left by the \emph{unsigned6} modulo 16. The eight resulting half words are saturated to 16 bits, packed and stored into the \emph{registerM}.


\paragraph{Execution} {Reservation table \textsf{ALU\_LITE} (Section~\ref{sec:reservation-ALU-LITE})}
~
{\begin{lstlisting}
stage RR:
  new argument3 = _ZX_6(unsigned6);
  new argument2 = PGR[registerP];
  new shift = _ZX_4(argument3);
  for (I = 0; I < 8; I += 1) {
    result1.16[I] = _SAT_16(_SX_16(argument2.16[I]) << shift)
  };
stage E1:
  PGR[registerM] = result1;
\end{lstlisting}}
\newpage

\paragraph{Encoding} Format \textsf{ALU\_HOSWRR} (Section~\ref{sec:format-ALU-HOSWRR}), \verb|exucode1 = 100|\\

\bitfields{ 1/P, 2/11, 1/1, 1/0, 3/100, 5/registerM, 1/--, 2/10, 3/100, 1/1, 6/registerY, 5/registerP, 1/0 }


\paragraph{Syntax} \texttt{slsho} \emph{registerM} = \emph{registerP}, \emph{registerY}

\paragraph{Description} The \emph{registerP} is interpreted as eight half words packed into 128 bits. The \emph{registerP} half words are shifted left by the \emph{registerY} modulo 16. The eight resulting half words are saturated to 16 bits, packed and stored into the \emph{registerM}.


\paragraph{Execution} {Reservation table \textsf{ALU\_LITE} (Section~\ref{sec:reservation-ALU-LITE})}
~
{\begin{lstlisting}
stage RR:
  new argument3 = GPR[registerY];
  new argument2 = PGR[registerP];
  new shift = _ZX_4(argument3);
  for (I = 0; I < 8; I += 1) {
    result1.16[I] = _SAT_16(_SX_16(argument2.16[I]) << shift)
  };
stage E1:
  PGR[registerM] = result1;
\end{lstlisting}}
\newpage
\subsubsection[{\small SLSW: Shift Left Saturated Word}]{SLSW\hfill Shift Left Saturated Word}
\label{instr:SLSW}\index{slsw}
 
\paragraph{Encoding} Format \textsf{ALU\_WSWRR} (Section~\ref{sec:format-ALU-WSWRR}), \verb|exucode1 = 100|\\

\bitfields{ 1/P, 2/11, 1/0, 1/0, 3/100, 6/registerW, 2/11, 3/100, 1/signextw, 6/registerY, 6/registerZ }


\paragraph{Syntax} \texttt{slsw}\emph{signextw} \emph{registerW} = \emph{registerZ}, \emph{registerY}

\paragraph{Description} The \emph{registerZ} is shifted left by the \emph{registerY} modulo 32 then is saturated to 32 bits. The result with the 32 upper bits cleared is stored into the \emph{registerW}.


\paragraph{Modifier(s)}
\noindent\begin{itemize}
\item signextw (cf. signextw in section \ref{par:modifier-signextw} on page \pageref{par:modifier-signextw})
\end{itemize}

\paragraph{Execution} {Reservation table \textsf{ALU\_LITE} (Section~\ref{sec:reservation-ALU-LITE})}
~
{\begin{lstlisting}
stage ID:
  new signextw = _ZX_1(signextw);
stage RR:
  new argument3 = GPR[registerY];
  new argument2 = GPR[registerZ];
  new shift = _ZX_5(argument3);
  new result1 = _SAT_32(_SX_32(argument2) << shift);
stage E1:
  GPR[registerW] = signextw ? _SX_32(result1) : _ZX_32(result1);
\end{lstlisting}}
\newpage

\paragraph{Encoding} Format \textsf{ALU\_WSWRI} (Section~\ref{sec:format-ALU-WSWRI}), \verb|exucode1 = 100|\\

\bitfields{ 1/P, 2/11, 1/0, 1/1, 3/100, 6/registerW, 2/11, 3/100, 1/signextw, 6/unsigned6, 6/registerZ }


\paragraph{Syntax} \texttt{slsw}\emph{signextw} \emph{registerW} = \emph{registerZ}, \emph{unsigned6}

\paragraph{Description} The \emph{registerZ} is shifted left by the \emph{unsigned6} modulo 32 then is saturated to 32 bits. The result with the 32 upper bits cleared is stored into the \emph{registerW}.


\paragraph{Modifier(s)}
\noindent\begin{itemize}
\item signextw (cf. signextw in section \ref{par:modifier-signextw} on page \pageref{par:modifier-signextw})
\end{itemize}

\paragraph{Execution} {Reservation table \textsf{ALU\_LITE} (Section~\ref{sec:reservation-ALU-LITE})}
~
{\begin{lstlisting}
stage ID:
  new signextw = _ZX_1(signextw);
stage RR:
  new argument3 = _ZX_6(unsigned6);
  new argument2 = GPR[registerZ];
  new shift = _ZX_5(argument3);
  new result1 = _SAT_32(_SX_32(argument2) << shift);
stage E1:
  GPR[registerW] = signextw ? _SX_32(result1) : _ZX_32(result1);
\end{lstlisting}}
\newpage
\subsubsection[{\small SLSWQ: Shift Left Saturated Word Quadruple}]{SLSWQ\hfill Shift Left Saturated Word Quadruple}
\label{instr:SLSWQ}\index{slswq}
 
\paragraph{Encoding} Format \textsf{ALU\_WQSWRI} (Section~\ref{sec:format-ALU-WQSWRI}), \verb|exucode1 = 100|\\

\bitfields{ 1/P, 2/11, 1/1, 1/1, 3/100, 5/registerM, 1/--, 2/10, 3/100, 1/0, 6/unsigned6, 5/registerP, 1/1 }


\paragraph{Syntax} \texttt{slswq} \emph{registerM} = \emph{registerP}, \emph{unsigned6}

\paragraph{Description} The \emph{registerP} is interpreted as four words packed into 128 bits. The \emph{registerP} words are shifted left by the \emph{unsigned6} modulo 32. The four resulting words are saturated to 32 bits, packed and stored into the \emph{registerM}.


\paragraph{Execution} {Reservation table \textsf{ALU\_LITE} (Section~\ref{sec:reservation-ALU-LITE})}
~
{\begin{lstlisting}
stage RR:
  new argument3 = _ZX_6(unsigned6);
  new argument2 = PGR[registerP];
  new shift = _ZX_5(argument3);
  for (I = 0; I < 4; I += 1) {
    result1.32[I] = _SAT_32(_SX_32(argument2.32[I]) << shift)
  };
stage E1:
  PGR[registerM] = result1;
\end{lstlisting}}
\newpage

\paragraph{Encoding} Format \textsf{ALU\_WQSWRR} (Section~\ref{sec:format-ALU-WQSWRR}), \verb|exucode1 = 100|\\

\bitfields{ 1/P, 2/11, 1/1, 1/0, 3/100, 5/registerM, 1/--, 2/10, 3/100, 1/0, 6/registerY, 5/registerP, 1/1 }


\paragraph{Syntax} \texttt{slswq} \emph{registerM} = \emph{registerP}, \emph{registerY}

\paragraph{Description} The \emph{registerP} is interpreted as four words packed into 128 bits. The \emph{registerP} words are shifted left by the \emph{registerY} modulo 32. The four resulting words are saturated to 32 bits, packed and stored into the \emph{registerM}.


\paragraph{Execution} {Reservation table \textsf{ALU\_LITE} (Section~\ref{sec:reservation-ALU-LITE})}
~
{\begin{lstlisting}
stage RR:
  new argument3 = GPR[registerY];
  new argument2 = PGR[registerP];
  new shift = _ZX_5(argument3);
  for (I = 0; I < 4; I += 1) {
    result1.32[I] = _SAT_32(_SX_32(argument2.32[I]) << shift)
  };
stage E1:
  PGR[registerM] = result1;
\end{lstlisting}}
\newpage
\subsubsection[{\small SLUSBX: Shift Left Unsigned Saturated Byte Hexadecuple}]{SLUSBX\hfill Shift Left Unsigned Saturated Byte Hexadecuple}
\label{instr:SLUSBX}\index{slusbx}
 
\paragraph{Encoding} Format \textsf{ALU\_BXSWRI} (Section~\ref{sec:format-ALU-BXSWRI}), \verb|exucode1 = 101|\\

\bitfields{ 1/P, 2/11, 1/1, 1/1, 3/101, 5/registerM, 1/--, 2/10, 3/100, 1/1, 6/unsigned6, 5/registerP, 1/1 }


\paragraph{Syntax} \texttt{slusbx} \emph{registerM} = \emph{registerP}, \emph{unsigned6}

\paragraph{Description} The \emph{registerP} is interpreted as sixteen bytes packed into 128 bits. The \emph{registerP} bytes are shifted left by the \emph{unsigned6} modulo 8. The sixteen resulting bytes are unsigned saturated to 8 bits, packed and stored into the \emph{registerM}.


\paragraph{Execution} {Reservation table \textsf{ALU\_LITE} (Section~\ref{sec:reservation-ALU-LITE})}
~
{\begin{lstlisting}
stage RR:
  new argument3 = _ZX_6(unsigned6);
  new argument2 = PGR[registerP];
  new shift = _ZX_3(argument3);
  for (I = 0; I < 16; I += 1) {
    result1.8[I] = _SATU_8(_ZX_8(argument2.8[I]) << shift)
  };
stage E1:
  PGR[registerM] = result1;
\end{lstlisting}}
\newpage

\paragraph{Encoding} Format \textsf{ALU\_BXSWRR} (Section~\ref{sec:format-ALU-BXSWRR}), \verb|exucode1 = 101|\\

\bitfields{ 1/P, 2/11, 1/1, 1/0, 3/101, 5/registerM, 1/--, 2/10, 3/100, 1/1, 6/registerY, 5/registerP, 1/1 }


\paragraph{Syntax} \texttt{slusbx} \emph{registerM} = \emph{registerP}, \emph{registerY}

\paragraph{Description} The \emph{registerP} is interpreted as sixteen bytes packed into 128 bits. The \emph{registerP} bytes are shifted left by the \emph{registerY} modulo 8. The sixteen resulting bytes are unsigned saturated to 8 bits, packed and stored into the \emph{registerM}.


\paragraph{Execution} {Reservation table \textsf{ALU\_LITE} (Section~\ref{sec:reservation-ALU-LITE})}
~
{\begin{lstlisting}
stage RR:
  new argument3 = GPR[registerY];
  new argument2 = PGR[registerP];
  new shift = _ZX_3(argument3);
  for (I = 0; I < 16; I += 1) {
    result1.8[I] = _SATU_8(_ZX_8(argument2.8[I]) << shift)
  };
stage E1:
  PGR[registerM] = result1;
\end{lstlisting}}
\newpage
\subsubsection[{\small SLUSD: Shift Left Unsigned Saturated Double Word}]{SLUSD\hfill Shift Left Unsigned Saturated Double Word}
\label{instr:SLUSD}\index{slusd}
 
\paragraph{Encoding} Format \textsf{ALU\_DSWRI} (Section~\ref{sec:format-ALU-DSWRI}), \verb|exucode1 = 101|\\

\bitfields{ 1/P, 2/11, 1/0, 1/1, 3/101, 6/registerW, 2/10, 3/100, 1/0, 6/unsigned6, 6/registerZ }


\paragraph{Syntax} \texttt{slusd} \emph{registerW} = \emph{registerZ}, \emph{unsigned6}

\paragraph{Description} The \emph{registerZ} is shifted left by the \emph{unsigned6} modulo 64. The result is unsigned saturated to 64 bits and stored into the \emph{registerW}.


\paragraph{Execution} {Reservation table \textsf{ALU\_LITE} (Section~\ref{sec:reservation-ALU-LITE})}
~
{\begin{lstlisting}
stage RR:
  new argument3 = _ZX_6(unsigned6);
  new argument2 = GPR[registerZ];
  new shift = _ZX_6(argument3);
  new result1 = _SATU_64(_ZX_64(argument2) << shift);
stage E1:
  GPR[registerW] = result1;
\end{lstlisting}}
\newpage

\paragraph{Encoding} Format \textsf{ALU\_DSWRR} (Section~\ref{sec:format-ALU-DSWRR}), \verb|exucode1 = 101|\\

\bitfields{ 1/P, 2/11, 1/0, 1/0, 3/101, 6/registerW, 2/10, 3/100, 1/0, 6/registerY, 6/registerZ }


\paragraph{Syntax} \texttt{slusd} \emph{registerW} = \emph{registerZ}, \emph{registerY}

\paragraph{Description} The \emph{registerZ} is shifted left by the \emph{registerY} modulo 64. The result is unsigned saturated to 64 bits and stored into the \emph{registerW}.


\paragraph{Execution} {Reservation table \textsf{ALU\_LITE} (Section~\ref{sec:reservation-ALU-LITE})}
~
{\begin{lstlisting}
stage RR:
  new argument3 = GPR[registerY];
  new argument2 = GPR[registerZ];
  new shift = _ZX_6(argument3);
  new result1 = _SATU_64(_ZX_64(argument2) << shift);
stage E1:
  GPR[registerW] = result1;
\end{lstlisting}}
\newpage
\subsubsection[{\small SLUSDP: Shift Left Unsigned Saturated Double Word Pair}]{SLUSDP\hfill Shift Left Unsigned Saturated Double Word Pair}
\label{instr:SLUSDP}\index{slusdp}
 
\paragraph{Encoding} Format \textsf{ALU\_DPSWRI} (Section~\ref{sec:format-ALU-DPSWRI}), \verb|exucode1 = 101|\\

\bitfields{ 1/P, 2/11, 1/1, 1/1, 3/101, 5/registerM, 1/--, 2/10, 3/100, 1/0, 6/unsigned6, 5/registerP, 1/0 }


\paragraph{Syntax} \texttt{slusdp} \emph{registerM} = \emph{registerP}, \emph{unsigned6}

\paragraph{Description} The \emph{registerP} is interpreted as two double words packed into 128 bits. The \emph{registerP} words are shifted left by the \emph{unsigned6} modulo 64. The two resulting double words are unsigned saturated to 64 bits, packed and stored into the \emph{registerM}.


\paragraph{Execution} {Reservation table \textsf{ALU\_LITE} (Section~\ref{sec:reservation-ALU-LITE})}
~
{\begin{lstlisting}
stage RR:
  new argument3 = _ZX_6(unsigned6);
  new argument2 = PGR[registerP];
  new shift = _ZX_6(argument3);
  for (I = 0; I < 2; I += 1) {
    result1.64[I] = _SATU_64(_ZX_64(argument2.64[I]) << shift)
  };
stage E1:
  PGR[registerM] = result1;
\end{lstlisting}}
\newpage

\paragraph{Encoding} Format \textsf{ALU\_DPSWRR} (Section~\ref{sec:format-ALU-DPSWRR}), \verb|exucode1 = 101|\\

\bitfields{ 1/P, 2/11, 1/1, 1/0, 3/101, 5/registerM, 1/--, 2/10, 3/100, 1/0, 5/registerO, 1/--, 5/registerP, 1/0 }


\paragraph{Syntax} \texttt{slusdp} \emph{registerM} = \emph{registerP}, \emph{registerO}

\paragraph{Description} The \emph{registerP} is interpreted as two double words packed into 128 bits. The \emph{registerP} words are shifted left by the \emph{registerO} modulo 64. The two resulting double words are unsigned saturated to 64 bits, packed and stored into the \emph{registerM}.


\paragraph{Execution} {Reservation table \textsf{ALU\_LITE} (Section~\ref{sec:reservation-ALU-LITE})}
~
{\begin{lstlisting}
stage RR:
  new argument3 = PGR[registerO];
  new argument2 = PGR[registerP];
  new shift = _ZX_6(argument3);
  for (I = 0; I < 2; I += 1) {
    result1.64[I] = _SATU_64(_ZX_64(argument2.64[I]) << shift)
  };
stage E1:
  PGR[registerM] = result1;
\end{lstlisting}}
\newpage
\subsubsection[{\small SLUSHO: Shift Left Unsigned Saturated Half Word Octuple}]{SLUSHO\hfill Shift Left Unsigned Saturated Half Word Octuple}
\label{instr:SLUSHO}\index{slusho}
 
\paragraph{Encoding} Format \textsf{ALU\_HOSWRI} (Section~\ref{sec:format-ALU-HOSWRI}), \verb|exucode1 = 101|\\

\bitfields{ 1/P, 2/11, 1/1, 1/1, 3/101, 5/registerM, 1/--, 2/10, 3/100, 1/1, 6/unsigned6, 5/registerP, 1/0 }


\paragraph{Syntax} \texttt{slusho} \emph{registerM} = \emph{registerP}, \emph{unsigned6}

\paragraph{Description} The \emph{registerP} is interpreted as eight half words packed into 128 bits. The \emph{registerP} half words are shifted left by the \emph{unsigned6} modulo 16. The eight resulting half words are unsigned saturated to 16 bits, packed and stored into the \emph{registerM}.


\paragraph{Execution} {Reservation table \textsf{ALU\_LITE} (Section~\ref{sec:reservation-ALU-LITE})}
~
{\begin{lstlisting}
stage RR:
  new argument3 = _ZX_6(unsigned6);
  new argument2 = PGR[registerP];
  new shift = _ZX_4(argument3);
  for (I = 0; I < 8; I += 1) {
    result1.16[I] = _SATU_16(_ZX_16(argument2.16[I]) << shift)
  };
stage E1:
  PGR[registerM] = result1;
\end{lstlisting}}
\newpage

\paragraph{Encoding} Format \textsf{ALU\_HOSWRR} (Section~\ref{sec:format-ALU-HOSWRR}), \verb|exucode1 = 101|\\

\bitfields{ 1/P, 2/11, 1/1, 1/0, 3/101, 5/registerM, 1/--, 2/10, 3/100, 1/1, 6/registerY, 5/registerP, 1/0 }


\paragraph{Syntax} \texttt{slusho} \emph{registerM} = \emph{registerP}, \emph{registerY}

\paragraph{Description} The \emph{registerP} is interpreted as eight half words packed into 128 bits. The \emph{registerP} half words are shifted left by the \emph{registerY} modulo 16. The eight resulting half words are unsigned saturated to 16 bits, packed and stored into the \emph{registerM}.


\paragraph{Execution} {Reservation table \textsf{ALU\_LITE} (Section~\ref{sec:reservation-ALU-LITE})}
~
{\begin{lstlisting}
stage RR:
  new argument3 = GPR[registerY];
  new argument2 = PGR[registerP];
  new shift = _ZX_4(argument3);
  for (I = 0; I < 8; I += 1) {
    result1.16[I] = _SATU_16(_ZX_16(argument2.16[I]) << shift)
  };
stage E1:
  PGR[registerM] = result1;
\end{lstlisting}}
\newpage
\subsubsection[{\small SLUSW: Shift Left Unsigned Saturated Word}]{SLUSW\hfill Shift Left Unsigned Saturated Word}
\label{instr:SLUSW}\index{slusw}
 
\paragraph{Encoding} Format \textsf{ALU\_WSWRR} (Section~\ref{sec:format-ALU-WSWRR}), \verb|exucode1 = 101|\\

\bitfields{ 1/P, 2/11, 1/0, 1/0, 3/101, 6/registerW, 2/11, 3/100, 1/signextw, 6/registerY, 6/registerZ }


\paragraph{Syntax} \texttt{slusw}\emph{signextw} \emph{registerW} = \emph{registerZ}, \emph{registerY}

\paragraph{Description} The \emph{registerZ} is shifted left by the \emph{registerY} modulo 32 then is unsigned saturated to 32 bits. The result with the 32 upper bits cleared is stored into the \emph{registerW}.


\paragraph{Modifier(s)}
\noindent\begin{itemize}
\item signextw (cf. signextw in section \ref{par:modifier-signextw} on page \pageref{par:modifier-signextw})
\end{itemize}

\paragraph{Execution} {Reservation table \textsf{ALU\_LITE} (Section~\ref{sec:reservation-ALU-LITE})}
~
{\begin{lstlisting}
stage ID:
  new signextw = _ZX_1(signextw);
stage RR:
  new argument3 = GPR[registerY];
  new argument2 = GPR[registerZ];
  new shift = _ZX_5(argument3);
  new result1 = _SATU_32(_ZX_32(argument2) << shift);
stage E1:
  GPR[registerW] = signextw ? _SX_32(result1) : _ZX_32(result1);
\end{lstlisting}}
\newpage

\paragraph{Encoding} Format \textsf{ALU\_WSWRI} (Section~\ref{sec:format-ALU-WSWRI}), \verb|exucode1 = 101|\\

\bitfields{ 1/P, 2/11, 1/0, 1/1, 3/101, 6/registerW, 2/11, 3/100, 1/signextw, 6/unsigned6, 6/registerZ }


\paragraph{Syntax} \texttt{slusw}\emph{signextw} \emph{registerW} = \emph{registerZ}, \emph{unsigned6}

\paragraph{Description} The \emph{registerZ} is shifted left by the \emph{unsigned6} modulo 32 then is unsigned saturated to 32 bits. The result with the 32 upper bits cleared is stored into the \emph{registerW}.


\paragraph{Modifier(s)}
\noindent\begin{itemize}
\item signextw (cf. signextw in section \ref{par:modifier-signextw} on page \pageref{par:modifier-signextw})
\end{itemize}

\paragraph{Execution} {Reservation table \textsf{ALU\_LITE} (Section~\ref{sec:reservation-ALU-LITE})}
~
{\begin{lstlisting}
stage ID:
  new signextw = _ZX_1(signextw);
stage RR:
  new argument3 = _ZX_6(unsigned6);
  new argument2 = GPR[registerZ];
  new shift = _ZX_5(argument3);
  new result1 = _SATU_32(_ZX_32(argument2) << shift);
stage E1:
  GPR[registerW] = signextw ? _SX_32(result1) : _ZX_32(result1);
\end{lstlisting}}
\newpage
\subsubsection[{\small SLUSWQ: Shift Left Unsigned Saturated Word Quadruple}]{SLUSWQ\hfill Shift Left Unsigned Saturated Word Quadruple}
\label{instr:SLUSWQ}\index{sluswq}
 
\paragraph{Encoding} Format \textsf{ALU\_WQSWRI} (Section~\ref{sec:format-ALU-WQSWRI}), \verb|exucode1 = 101|\\

\bitfields{ 1/P, 2/11, 1/1, 1/1, 3/101, 5/registerM, 1/--, 2/10, 3/100, 1/0, 6/unsigned6, 5/registerP, 1/1 }


\paragraph{Syntax} \texttt{sluswq} \emph{registerM} = \emph{registerP}, \emph{unsigned6}

\paragraph{Description} The \emph{registerP} is interpreted as four words packed into 128 bits. The \emph{registerP} words are shifted left by the \emph{unsigned6} modulo 32. The four resulting words are unsigned saturated to 32 bits, packed and stored into the \emph{registerM}.


\paragraph{Execution} {Reservation table \textsf{ALU\_LITE} (Section~\ref{sec:reservation-ALU-LITE})}
~
{\begin{lstlisting}
stage RR:
  new argument3 = _ZX_6(unsigned6);
  new argument2 = PGR[registerP];
  new shift = _ZX_5(argument3);
  for (I = 0; I < 4; I += 1) {
    result1.32[I] = _SATU_32(_ZX_32(argument2.32[I]) << shift)
  };
stage E1:
  PGR[registerM] = result1;
\end{lstlisting}}
\newpage

\paragraph{Encoding} Format \textsf{ALU\_WQSWRR} (Section~\ref{sec:format-ALU-WQSWRR}), \verb|exucode1 = 101|\\

\bitfields{ 1/P, 2/11, 1/1, 1/0, 3/101, 5/registerM, 1/--, 2/10, 3/100, 1/0, 6/registerY, 5/registerP, 1/1 }


\paragraph{Syntax} \texttt{sluswq} \emph{registerM} = \emph{registerP}, \emph{registerY}

\paragraph{Description} The \emph{registerP} is interpreted as four words packed into 128 bits. The \emph{registerP} words are shifted left by the \emph{registerY} modulo 32. The four resulting words are unsigned saturated to 32 bits, packed and stored into the \emph{registerM}.


\paragraph{Execution} {Reservation table \textsf{ALU\_LITE} (Section~\ref{sec:reservation-ALU-LITE})}
~
{\begin{lstlisting}
stage RR:
  new argument3 = GPR[registerY];
  new argument2 = PGR[registerP];
  new shift = _ZX_5(argument3);
  for (I = 0; I < 4; I += 1) {
    result1.32[I] = _SATU_32(_ZX_32(argument2.32[I]) << shift)
  };
stage E1:
  PGR[registerM] = result1;
\end{lstlisting}}
\newpage
\subsubsection[{\small SRABX: Shift Right Arithmetic Byte Hexadecuple}]{SRABX\hfill Shift Right Arithmetic Byte Hexadecuple}
\label{instr:SRABX}\index{srabx}
 
\paragraph{Encoding} Format \textsf{ALU\_BXSWRI} (Section~\ref{sec:format-ALU-BXSWRI}), \verb|exucode1 = 010|\\

\bitfields{ 1/P, 2/11, 1/1, 1/1, 3/010, 5/registerM, 1/--, 2/10, 3/100, 1/1, 6/unsigned6, 5/registerP, 1/1 }


\paragraph{Syntax} \texttt{srabx} \emph{registerM} = \emph{registerP}, \emph{unsigned6}

\paragraph{Description} The \emph{registerP} is interpreted as sixteen bytes packed into 128 bits. The \emph{registerP} bytes are shifted right by the \emph{unsigned6} modulo 8. The leftmost bit are replicated to fill in the vacant positions, thereby extending the signs of the shifted operands. The sixteen resulting bytes are packed and stored into the \emph{registerM}.


\paragraph{Execution} {Reservation table \textsf{ALU\_LITE} (Section~\ref{sec:reservation-ALU-LITE})}
~
{\begin{lstlisting}
stage RR:
  new argument3 = _ZX_6(unsigned6);
  new argument2 = PGR[registerP];
  new shift = _ZX_3(argument3);
  for (I = 0; I < 16; I += 1) {
    result1.8[I] = _SX_8(argument2.8[I]) >> shift
  };
stage E1:
  PGR[registerM] = result1;
\end{lstlisting}}
\newpage

\paragraph{Encoding} Format \textsf{ALU\_BXSWRR} (Section~\ref{sec:format-ALU-BXSWRR}), \verb|exucode1 = 010|\\

\bitfields{ 1/P, 2/11, 1/1, 1/0, 3/010, 5/registerM, 1/--, 2/10, 3/100, 1/1, 6/registerY, 5/registerP, 1/1 }


\paragraph{Syntax} \texttt{srabx} \emph{registerM} = \emph{registerP}, \emph{registerY}

\paragraph{Description} The \emph{registerP} is interpreted as sixteen bytes packed into 128 bits. The \emph{registerP} bytes are shifted right by the \emph{registerY} modulo 8. The leftmost bit are replicated to fill in the vacant positions, thereby extending the signs of the shifted operands. The sixteen resulting bytes are packed and stored into the \emph{registerM}.


\paragraph{Execution} {Reservation table \textsf{ALU\_LITE} (Section~\ref{sec:reservation-ALU-LITE})}
~
{\begin{lstlisting}
stage RR:
  new argument3 = GPR[registerY];
  new argument2 = PGR[registerP];
  new shift = _ZX_3(argument3);
  for (I = 0; I < 16; I += 1) {
    result1.8[I] = _SX_8(argument2.8[I]) >> shift
  };
stage E1:
  PGR[registerM] = result1;
\end{lstlisting}}
\newpage
\subsubsection[{\small SRAD: Shift Right Arithmetic Double Word}]{SRAD\hfill Shift Right Arithmetic Double Word}
\label{instr:SRAD}\index{srad}
 
\paragraph{Encoding} Format \textsf{ALU\_DSWRI} (Section~\ref{sec:format-ALU-DSWRI}), \verb|exucode1 = 010|\\

\bitfields{ 1/P, 2/11, 1/0, 1/1, 3/010, 6/registerW, 2/10, 3/100, 1/0, 6/unsigned6, 6/registerZ }


\paragraph{Syntax} \texttt{srad} \emph{registerW} = \emph{registerZ}, \emph{unsigned6}

\paragraph{Description} The \emph{registerZ} is shifted right by the \emph{unsigned6} modulo 64. The leftmost bit is replicated to fill in the vacant positions, thereby extending the sign of the shifted operand. The result is stored into the \emph{registerW}.


\paragraph{Execution} {Reservation table \textsf{ALU\_TINY} (Section~\ref{sec:reservation-ALU-TINY})}
~
{\begin{lstlisting}
stage RR:
  new argument3 = _ZX_6(unsigned6);
  new argument2 = GPR[registerZ];
  new result1 = _SX_64(argument2) >> _ZX_6(argument3);
stage E1:
  GPR[registerW] = result1;
\end{lstlisting}}
\newpage

\paragraph{Encoding} Format \textsf{ALU\_DSWRR} (Section~\ref{sec:format-ALU-DSWRR}), \verb|exucode1 = 010|\\

\bitfields{ 1/P, 2/11, 1/0, 1/0, 3/010, 6/registerW, 2/10, 3/100, 1/0, 6/registerY, 6/registerZ }


\paragraph{Syntax} \texttt{srad} \emph{registerW} = \emph{registerZ}, \emph{registerY}

\paragraph{Description} The \emph{registerZ} is shifted right by the \emph{registerY} modulo 64. The leftmost bit is replicated to fill in the vacant positions, thereby extending the sign of the shifted operand. The result is stored into the \emph{registerW}.


\paragraph{Execution} {Reservation table \textsf{ALU\_TINY} (Section~\ref{sec:reservation-ALU-TINY})}
~
{\begin{lstlisting}
stage RR:
  new argument3 = GPR[registerY];
  new argument2 = GPR[registerZ];
  new result1 = _SX_64(argument2) >> _ZX_6(argument3);
stage E1:
  GPR[registerW] = result1;
\end{lstlisting}}
\newpage
\subsubsection[{\small SRADP: Shift Right Arithmetic Double Word Pair}]{SRADP\hfill Shift Right Arithmetic Double Word Pair}
\label{instr:SRADP}\index{sradp}
 
\paragraph{Encoding} Format \textsf{ALU\_DPSWRI} (Section~\ref{sec:format-ALU-DPSWRI}), \verb|exucode1 = 010|\\

\bitfields{ 1/P, 2/11, 1/1, 1/1, 3/010, 5/registerM, 1/--, 2/10, 3/100, 1/0, 6/unsigned6, 5/registerP, 1/0 }


\paragraph{Syntax} \texttt{sradp} \emph{registerM} = \emph{registerP}, \emph{unsigned6}

\paragraph{Description} The \emph{registerP} is interpreted as two double words packed into 128 bits. The \emph{registerP} words are shifted right by the \emph{unsigned6} modulo 64. The leftmost bit are replicated to fill in the vacant positions, thereby extending the signs of the shifted operands. The two resulting double words are packed and stored into the \emph{registerM}.


\paragraph{Execution} {Reservation table \textsf{ALU\_LITE} (Section~\ref{sec:reservation-ALU-LITE})}
~
{\begin{lstlisting}
stage RR:
  new argument3 = _ZX_6(unsigned6);
  new argument2 = PGR[registerP];
  new shift = _ZX_6(argument3);
  for (I = 0; I < 2; I += 1) {
    result1.64[I] = _SX_64(argument2.64[I]) >> shift
  };
stage E1:
  PGR[registerM] = result1;
\end{lstlisting}}
\newpage

\paragraph{Encoding} Format \textsf{ALU\_DPSWRR} (Section~\ref{sec:format-ALU-DPSWRR}), \verb|exucode1 = 010|\\

\bitfields{ 1/P, 2/11, 1/1, 1/0, 3/010, 5/registerM, 1/--, 2/10, 3/100, 1/0, 5/registerO, 1/--, 5/registerP, 1/0 }


\paragraph{Syntax} \texttt{sradp} \emph{registerM} = \emph{registerP}, \emph{registerO}

\paragraph{Description} The \emph{registerP} is interpreted as two double words packed into 128 bits. The \emph{registerP} words are shifted right by the \emph{registerO} modulo 64. The leftmost bit are replicated to fill in the vacant positions, thereby extending the signs of the shifted operands. The two resulting double words are packed and stored into the \emph{registerM}.


\paragraph{Execution} {Reservation table \textsf{ALU\_LITE} (Section~\ref{sec:reservation-ALU-LITE})}
~
{\begin{lstlisting}
stage RR:
  new argument3 = PGR[registerO];
  new argument2 = PGR[registerP];
  new shift = _ZX_6(argument3);
  for (I = 0; I < 2; I += 1) {
    result1.64[I] = _SX_64(argument2.64[I]) >> shift
  };
stage E1:
  PGR[registerM] = result1;
\end{lstlisting}}
\newpage
\subsubsection[{\small SRAHO: Shift Right Arithmetic Half Word Octuple}]{SRAHO\hfill Shift Right Arithmetic Half Word Octuple}
\label{instr:SRAHO}\index{sraho}
 
\paragraph{Encoding} Format \textsf{ALU\_HOSWRI} (Section~\ref{sec:format-ALU-HOSWRI}), \verb|exucode1 = 010|\\

\bitfields{ 1/P, 2/11, 1/1, 1/1, 3/010, 5/registerM, 1/--, 2/10, 3/100, 1/1, 6/unsigned6, 5/registerP, 1/0 }


\paragraph{Syntax} \texttt{sraho} \emph{registerM} = \emph{registerP}, \emph{unsigned6}

\paragraph{Description} The \emph{registerP} is interpreted as eight half words packed into 128 bits. The \emph{registerP} half words are shifted right by the \emph{unsigned6} modulo 16. The leftmost bit are replicated to fill in the vacant positions, thereby extending the signs of the shifted operands. The eight resulting half words are packed and stored into the \emph{registerM}.


\paragraph{Execution} {Reservation table \textsf{ALU\_LITE} (Section~\ref{sec:reservation-ALU-LITE})}
~
{\begin{lstlisting}
stage RR:
  new argument3 = _ZX_6(unsigned6);
  new argument2 = PGR[registerP];
  new shift = _ZX_4(argument3);
  for (I = 0; I < 8; I += 1) {
    result1.16[I] = _SX_16(argument2.16[I]) >> shift
  };
stage E1:
  PGR[registerM] = result1;
\end{lstlisting}}
\newpage

\paragraph{Encoding} Format \textsf{ALU\_HOSWRR} (Section~\ref{sec:format-ALU-HOSWRR}), \verb|exucode1 = 010|\\

\bitfields{ 1/P, 2/11, 1/1, 1/0, 3/010, 5/registerM, 1/--, 2/10, 3/100, 1/1, 6/registerY, 5/registerP, 1/0 }


\paragraph{Syntax} \texttt{sraho} \emph{registerM} = \emph{registerP}, \emph{registerY}

\paragraph{Description} The \emph{registerP} is interpreted as eight half words packed into 128 bits. The \emph{registerP} half words are shifted right by the \emph{registerY} modulo 16. The leftmost bit are replicated to fill in the vacant positions, thereby extending the signs of the shifted operands. The eight resulting half words are packed and stored into the \emph{registerM}.


\paragraph{Execution} {Reservation table \textsf{ALU\_LITE} (Section~\ref{sec:reservation-ALU-LITE})}
~
{\begin{lstlisting}
stage RR:
  new argument3 = GPR[registerY];
  new argument2 = PGR[registerP];
  new shift = _ZX_4(argument3);
  for (I = 0; I < 8; I += 1) {
    result1.16[I] = _SX_16(argument2.16[I]) >> shift
  };
stage E1:
  PGR[registerM] = result1;
\end{lstlisting}}
\newpage
\subsubsection[{\small SRAW: Shift Right Arithmetic Word}]{SRAW\hfill Shift Right Arithmetic Word}
\label{instr:SRAW}\index{sraw}
 
\paragraph{Encoding} Format \textsf{ALU\_WSWRR} (Section~\ref{sec:format-ALU-WSWRR}), \verb|exucode1 = 010|\\

\bitfields{ 1/P, 2/11, 1/0, 1/0, 3/010, 6/registerW, 2/11, 3/100, 1/signextw, 6/registerY, 6/registerZ }


\paragraph{Syntax} \texttt{sraw}\emph{signextw} \emph{registerW} = \emph{registerZ}, \emph{registerY}

\paragraph{Description} The \emph{registerZ} is shifted right by the \emph{registerY} modulo 32. The leftmost bit is replicated to fill in the vacant positions, thereby extending the sign of the shifted operand. The result with the 32 upper bits cleared is stored into the \emph{registerW}.


\paragraph{Modifier(s)}
\noindent\begin{itemize}
\item signextw (cf. signextw in section \ref{par:modifier-signextw} on page \pageref{par:modifier-signextw})
\end{itemize}

\paragraph{Execution} {Reservation table \textsf{ALU\_TINY} (Section~\ref{sec:reservation-ALU-TINY})}
~
{\begin{lstlisting}
stage ID:
  new signextw = _ZX_1(signextw);
stage RR:
  new argument3 = GPR[registerY];
  new argument2 = GPR[registerZ];
  new result1 = _SX_32(argument2) >> _ZX_5(argument3);
stage E1:
  GPR[registerW] = signextw ? _SX_32(result1) : _ZX_32(result1);
\end{lstlisting}}
\newpage

\paragraph{Encoding} Format \textsf{ALU\_WSWRI} (Section~\ref{sec:format-ALU-WSWRI}), \verb|exucode1 = 010|\\

\bitfields{ 1/P, 2/11, 1/0, 1/1, 3/010, 6/registerW, 2/11, 3/100, 1/signextw, 6/unsigned6, 6/registerZ }


\paragraph{Syntax} \texttt{sraw}\emph{signextw} \emph{registerW} = \emph{registerZ}, \emph{unsigned6}

\paragraph{Description} The \emph{registerZ} is shifted right by the \emph{unsigned6} modulo 32. The leftmost bit is replicated to fill in the vacant positions, thereby extending the sign of the shifted operand. The result with the 32 upper bits cleared is stored into the \emph{registerW}.


\paragraph{Modifier(s)}
\noindent\begin{itemize}
\item signextw (cf. signextw in section \ref{par:modifier-signextw} on page \pageref{par:modifier-signextw})
\end{itemize}

\paragraph{Execution} {Reservation table \textsf{ALU\_TINY} (Section~\ref{sec:reservation-ALU-TINY})}
~
{\begin{lstlisting}
stage ID:
  new signextw = _ZX_1(signextw);
stage RR:
  new argument3 = _ZX_6(unsigned6);
  new argument2 = GPR[registerZ];
  new result1 = _SX_32(argument2) >> _ZX_5(argument3);
stage E1:
  GPR[registerW] = signextw ? _SX_32(result1) : _ZX_32(result1);
\end{lstlisting}}
\newpage
\subsubsection[{\small SRAWQ: Shift Right Arithmetic Word Quadruple}]{SRAWQ\hfill Shift Right Arithmetic Word Quadruple}
\label{instr:SRAWQ}\index{srawq}
 
\paragraph{Encoding} Format \textsf{ALU\_WQSWRI} (Section~\ref{sec:format-ALU-WQSWRI}), \verb|exucode1 = 010|\\

\bitfields{ 1/P, 2/11, 1/1, 1/1, 3/010, 5/registerM, 1/--, 2/10, 3/100, 1/0, 6/unsigned6, 5/registerP, 1/1 }


\paragraph{Syntax} \texttt{srawq} \emph{registerM} = \emph{registerP}, \emph{unsigned6}

\paragraph{Description} The \emph{registerP} is interpreted as four words packed into 128 bits. The \emph{registerP} words are shifted right by the \emph{unsigned6} modulo 32. The leftmost bit are replicated to fill in the vacant positions, thereby extending the signs of the shifted operands. The four resulting words are packed and stored into the \emph{registerM}.


\paragraph{Execution} {Reservation table \textsf{ALU\_LITE} (Section~\ref{sec:reservation-ALU-LITE})}
~
{\begin{lstlisting}
stage RR:
  new argument3 = _ZX_6(unsigned6);
  new argument2 = PGR[registerP];
  new shift = _ZX_5(argument3);
  for (I = 0; I < 4; I += 1) {
    result1.32[I] = _SX_32(argument2.32[I]) >> shift
  };
stage E1:
  PGR[registerM] = result1;
\end{lstlisting}}
\newpage

\paragraph{Encoding} Format \textsf{ALU\_WQSWRR} (Section~\ref{sec:format-ALU-WQSWRR}), \verb|exucode1 = 010|\\

\bitfields{ 1/P, 2/11, 1/1, 1/0, 3/010, 5/registerM, 1/--, 2/10, 3/100, 1/0, 6/registerY, 5/registerP, 1/1 }


\paragraph{Syntax} \texttt{srawq} \emph{registerM} = \emph{registerP}, \emph{registerY}

\paragraph{Description} The \emph{registerP} is interpreted as four words packed into 128 bits. The \emph{registerP} words are shifted right by the \emph{registerY} modulo 32. The leftmost bit are replicated to fill in the vacant positions, thereby extending the signs of the shifted operands. The four resulting words are packed and stored into the \emph{registerM}.


\paragraph{Execution} {Reservation table \textsf{ALU\_LITE} (Section~\ref{sec:reservation-ALU-LITE})}
~
{\begin{lstlisting}
stage RR:
  new argument3 = GPR[registerY];
  new argument2 = PGR[registerP];
  new shift = _ZX_5(argument3);
  for (I = 0; I < 4; I += 1) {
    result1.32[I] = _SX_32(argument2.32[I]) >> shift
  };
stage E1:
  PGR[registerM] = result1;
\end{lstlisting}}
\newpage
\subsubsection[{\small SRLBX: Shift Right Logical Byte Hexadecuple}]{SRLBX\hfill Shift Right Logical Byte Hexadecuple}
\label{instr:SRLBX}\index{srlbx}
 
\paragraph{Encoding} Format \textsf{ALU\_BXSWRI} (Section~\ref{sec:format-ALU-BXSWRI}), \verb|exucode1 = 011|\\

\bitfields{ 1/P, 2/11, 1/1, 1/1, 3/011, 5/registerM, 1/--, 2/10, 3/100, 1/1, 6/unsigned6, 5/registerP, 1/1 }


\paragraph{Syntax} \texttt{srlbx} \emph{registerM} = \emph{registerP}, \emph{unsigned6}

\paragraph{Description} The \emph{registerP} is interpreted as sixteen bytes packed into 128 bits. The \emph{registerP} bytes are shifted right by the \emph{unsigned6} modulo 8. The vacant positions are filled in with zeros. The sixteen resulting bytes are packed and stored into the \emph{registerM}.


\paragraph{Execution} {Reservation table \textsf{ALU\_LITE} (Section~\ref{sec:reservation-ALU-LITE})}
~
{\begin{lstlisting}
stage RR:
  new argument3 = _ZX_6(unsigned6);
  new argument2 = PGR[registerP];
  new shift = _ZX_3(argument3);
  for (I = 0; I < 16; I += 1) {
    result1.8[I] = _ZX_8(argument2.8[I]) >> shift
  };
stage E1:
  PGR[registerM] = result1;
\end{lstlisting}}
\newpage

\paragraph{Encoding} Format \textsf{ALU\_BXSWRR} (Section~\ref{sec:format-ALU-BXSWRR}), \verb|exucode1 = 011|\\

\bitfields{ 1/P, 2/11, 1/1, 1/0, 3/011, 5/registerM, 1/--, 2/10, 3/100, 1/1, 6/registerY, 5/registerP, 1/1 }


\paragraph{Syntax} \texttt{srlbx} \emph{registerM} = \emph{registerP}, \emph{registerY}

\paragraph{Description} The \emph{registerP} is interpreted as sixteen bytes packed into 128 bits. The \emph{registerP} bytes are shifted right by the \emph{registerY} modulo 8. The vacant positions are filled in with zeros. The sixteen resulting bytes are packed and stored into the \emph{registerM}.


\paragraph{Execution} {Reservation table \textsf{ALU\_LITE} (Section~\ref{sec:reservation-ALU-LITE})}
~
{\begin{lstlisting}
stage RR:
  new argument3 = GPR[registerY];
  new argument2 = PGR[registerP];
  new shift = _ZX_3(argument3);
  for (I = 0; I < 16; I += 1) {
    result1.8[I] = _ZX_8(argument2.8[I]) >> shift
  };
stage E1:
  PGR[registerM] = result1;
\end{lstlisting}}
\newpage
\subsubsection[{\small SRLD: Shift Right Logical Double Word}]{SRLD\hfill Shift Right Logical Double Word}
\label{instr:SRLD}\index{srld}
 
\paragraph{Encoding} Format \textsf{ALU\_DSWRI} (Section~\ref{sec:format-ALU-DSWRI}), \verb|exucode1 = 011|\\

\bitfields{ 1/P, 2/11, 1/0, 1/1, 3/011, 6/registerW, 2/10, 3/100, 1/0, 6/unsigned6, 6/registerZ }


\paragraph{Syntax} \texttt{srld} \emph{registerW} = \emph{registerZ}, \emph{unsigned6}

\paragraph{Description} The \emph{registerZ} is shifted right by the \emph{unsigned6} modulo 64. The vacant positions are filled in with zeros. The result is stored into the \emph{registerW}.


\paragraph{Execution} {Reservation table \textsf{ALU\_TINY} (Section~\ref{sec:reservation-ALU-TINY})}
~
{\begin{lstlisting}
stage RR:
  new argument3 = _ZX_6(unsigned6);
  new argument2 = GPR[registerZ];
  new result1 = _ZX_64(argument2) >> _ZX_6(argument3);
stage E1:
  GPR[registerW] = result1;
\end{lstlisting}}
\newpage

\paragraph{Encoding} Format \textsf{ALU\_DSWRR} (Section~\ref{sec:format-ALU-DSWRR}), \verb|exucode1 = 011|\\

\bitfields{ 1/P, 2/11, 1/0, 1/0, 3/011, 6/registerW, 2/10, 3/100, 1/0, 6/registerY, 6/registerZ }


\paragraph{Syntax} \texttt{srld} \emph{registerW} = \emph{registerZ}, \emph{registerY}

\paragraph{Description} The \emph{registerZ} is shifted right by the \emph{registerY} modulo 64. The vacant positions are filled in with zeros. The result is stored into the \emph{registerW}.


\paragraph{Execution} {Reservation table \textsf{ALU\_TINY} (Section~\ref{sec:reservation-ALU-TINY})}
~
{\begin{lstlisting}
stage RR:
  new argument3 = GPR[registerY];
  new argument2 = GPR[registerZ];
  new result1 = _ZX_64(argument2) >> _ZX_6(argument3);
stage E1:
  GPR[registerW] = result1;
\end{lstlisting}}
\newpage
\subsubsection[{\small SRLDP: Shift Right Logical Double Word Pair}]{SRLDP\hfill Shift Right Logical Double Word Pair}
\label{instr:SRLDP}\index{srldp}
 
\paragraph{Encoding} Format \textsf{ALU\_DPSWRI} (Section~\ref{sec:format-ALU-DPSWRI}), \verb|exucode1 = 011|\\

\bitfields{ 1/P, 2/11, 1/1, 1/1, 3/011, 5/registerM, 1/--, 2/10, 3/100, 1/0, 6/unsigned6, 5/registerP, 1/0 }


\paragraph{Syntax} \texttt{srldp} \emph{registerM} = \emph{registerP}, \emph{unsigned6}

\paragraph{Description} The \emph{registerP} is interpreted as two double words packed into 128 bits. The \emph{registerP} double words are shifted right by the \emph{unsigned6} modulo 64. The vacant positions are filled in with zeros. The two resulting double words are packed and stored into the \emph{registerM}.


\paragraph{Execution} {Reservation table \textsf{ALU\_LITE} (Section~\ref{sec:reservation-ALU-LITE})}
~
{\begin{lstlisting}
stage RR:
  new argument3 = _ZX_6(unsigned6);
  new argument2 = PGR[registerP];
  new shift = _ZX_6(argument3);
  for (I = 0; I < 2; I += 1) {
    result1.64[I] = _ZX_64(argument2.64[I]) >> shift
  };
stage E1:
  PGR[registerM] = result1;
\end{lstlisting}}
\newpage

\paragraph{Encoding} Format \textsf{ALU\_DPSWRR} (Section~\ref{sec:format-ALU-DPSWRR}), \verb|exucode1 = 011|\\

\bitfields{ 1/P, 2/11, 1/1, 1/0, 3/011, 5/registerM, 1/--, 2/10, 3/100, 1/0, 5/registerO, 1/--, 5/registerP, 1/0 }


\paragraph{Syntax} \texttt{srldp} \emph{registerM} = \emph{registerP}, \emph{registerO}

\paragraph{Description} The \emph{registerP} is interpreted as two double words packed into 128 bits. The \emph{registerP} double words are shifted right by the \emph{registerO} modulo 64. The vacant positions are filled in with zeros. The two resulting double words are packed and stored into the \emph{registerM}.


\paragraph{Execution} {Reservation table \textsf{ALU\_LITE} (Section~\ref{sec:reservation-ALU-LITE})}
~
{\begin{lstlisting}
stage RR:
  new argument3 = PGR[registerO];
  new argument2 = PGR[registerP];
  new shift = _ZX_6(argument3);
  for (I = 0; I < 2; I += 1) {
    result1.64[I] = _ZX_64(argument2.64[I]) >> shift
  };
stage E1:
  PGR[registerM] = result1;
\end{lstlisting}}
\newpage
\subsubsection[{\small SRLHO: Shift Right Logical Half Word Octuple}]{SRLHO\hfill Shift Right Logical Half Word Octuple}
\label{instr:SRLHO}\index{srlho}
 
\paragraph{Encoding} Format \textsf{ALU\_HOSWRI} (Section~\ref{sec:format-ALU-HOSWRI}), \verb|exucode1 = 011|\\

\bitfields{ 1/P, 2/11, 1/1, 1/1, 3/011, 5/registerM, 1/--, 2/10, 3/100, 1/1, 6/unsigned6, 5/registerP, 1/0 }


\paragraph{Syntax} \texttt{srlho} \emph{registerM} = \emph{registerP}, \emph{unsigned6}

\paragraph{Description} The \emph{registerP} is interpreted as eight half words packed into 128 bits. The \emph{registerP} half words are shifted right by the \emph{unsigned6} modulo 16. The vacant positions are filled in with zeros. The eight resulting half words are packed and stored into the \emph{registerM}.


\paragraph{Execution} {Reservation table \textsf{ALU\_LITE} (Section~\ref{sec:reservation-ALU-LITE})}
~
{\begin{lstlisting}
stage RR:
  new argument3 = _ZX_6(unsigned6);
  new argument2 = PGR[registerP];
  new shift = _ZX_4(argument3);
  for (I = 0; I < 8; I += 1) {
    result1.16[I] = _ZX_16(argument2.16[I]) >> shift
  };
stage E1:
  PGR[registerM] = result1;
\end{lstlisting}}
\newpage

\paragraph{Encoding} Format \textsf{ALU\_HOSWRR} (Section~\ref{sec:format-ALU-HOSWRR}), \verb|exucode1 = 011|\\

\bitfields{ 1/P, 2/11, 1/1, 1/0, 3/011, 5/registerM, 1/--, 2/10, 3/100, 1/1, 6/registerY, 5/registerP, 1/0 }


\paragraph{Syntax} \texttt{srlho} \emph{registerM} = \emph{registerP}, \emph{registerY}

\paragraph{Description} The \emph{registerP} is interpreted as eight half words packed into 128 bits. The \emph{registerP} half words are shifted right by the \emph{registerY} modulo 16. The vacant positions are filled in with zeros. The eight resulting half words are packed and stored into the \emph{registerM}.


\paragraph{Execution} {Reservation table \textsf{ALU\_LITE} (Section~\ref{sec:reservation-ALU-LITE})}
~
{\begin{lstlisting}
stage RR:
  new argument3 = GPR[registerY];
  new argument2 = PGR[registerP];
  new shift = _ZX_4(argument3);
  for (I = 0; I < 8; I += 1) {
    result1.16[I] = _ZX_16(argument2.16[I]) >> shift
  };
stage E1:
  PGR[registerM] = result1;
\end{lstlisting}}
\newpage
\subsubsection[{\small SRLW: Shift Right Logical Word}]{SRLW\hfill Shift Right Logical Word}
\label{instr:SRLW}\index{srlw}
 
\paragraph{Encoding} Format \textsf{ALU\_WSWRR} (Section~\ref{sec:format-ALU-WSWRR}), \verb|exucode1 = 011|\\

\bitfields{ 1/P, 2/11, 1/0, 1/0, 3/011, 6/registerW, 2/11, 3/100, 1/signextw, 6/registerY, 6/registerZ }


\paragraph{Syntax} \texttt{srlw}\emph{signextw} \emph{registerW} = \emph{registerZ}, \emph{registerY}

\paragraph{Description} The \emph{registerZ} is shifted right by the \emph{registerY} modulo 32. The vacant positions are filled in with zeros. The result with the 32 upper bits cleared is stored into the \emph{registerW}.


\paragraph{Modifier(s)}
\noindent\begin{itemize}
\item signextw (cf. signextw in section \ref{par:modifier-signextw} on page \pageref{par:modifier-signextw})
\end{itemize}

\paragraph{Execution} {Reservation table \textsf{ALU\_TINY} (Section~\ref{sec:reservation-ALU-TINY})}
~
{\begin{lstlisting}
stage ID:
  new signextw = _ZX_1(signextw);
stage RR:
  new argument3 = GPR[registerY];
  new argument2 = GPR[registerZ];
  new result1 = argument2.32[0] >> _ZX_5(argument3);
stage E1:
  GPR[registerW] = signextw ? _SX_32(result1) : _ZX_32(result1);
\end{lstlisting}}
\newpage

\paragraph{Encoding} Format \textsf{ALU\_WSWRI} (Section~\ref{sec:format-ALU-WSWRI}), \verb|exucode1 = 011|\\

\bitfields{ 1/P, 2/11, 1/0, 1/1, 3/011, 6/registerW, 2/11, 3/100, 1/signextw, 6/unsigned6, 6/registerZ }


\paragraph{Syntax} \texttt{srlw}\emph{signextw} \emph{registerW} = \emph{registerZ}, \emph{unsigned6}

\paragraph{Description} The \emph{registerZ} is shifted right by the \emph{unsigned6} modulo 32. The vacant positions are filled in with zeros. The result with the 32 upper bits cleared is stored into the \emph{registerW}.


\paragraph{Modifier(s)}
\noindent\begin{itemize}
\item signextw (cf. signextw in section \ref{par:modifier-signextw} on page \pageref{par:modifier-signextw})
\end{itemize}

\paragraph{Execution} {Reservation table \textsf{ALU\_TINY} (Section~\ref{sec:reservation-ALU-TINY})}
~
{\begin{lstlisting}
stage ID:
  new signextw = _ZX_1(signextw);
stage RR:
  new argument3 = _ZX_6(unsigned6);
  new argument2 = GPR[registerZ];
  new result1 = argument2.32[0] >> _ZX_5(argument3);
stage E1:
  GPR[registerW] = signextw ? _SX_32(result1) : _ZX_32(result1);
\end{lstlisting}}
\newpage
\subsubsection[{\small SRLWQ: Shift Right Logical Word Quadruple}]{SRLWQ\hfill Shift Right Logical Word Quadruple}
\label{instr:SRLWQ}\index{srlwq}
 
\paragraph{Encoding} Format \textsf{ALU\_WQSWRI} (Section~\ref{sec:format-ALU-WQSWRI}), \verb|exucode1 = 011|\\

\bitfields{ 1/P, 2/11, 1/1, 1/1, 3/011, 5/registerM, 1/--, 2/10, 3/100, 1/0, 6/unsigned6, 5/registerP, 1/1 }


\paragraph{Syntax} \texttt{srlwq} \emph{registerM} = \emph{registerP}, \emph{unsigned6}

\paragraph{Description} The \emph{registerP} is interpreted as four words packed into 128 bits. The \emph{registerP} words are shifted right by the \emph{unsigned6} modulo 32. The vacant positions are filled in with zeros. The four resulting words are packed and stored into the \emph{registerM}.


\paragraph{Execution} {Reservation table \textsf{ALU\_LITE} (Section~\ref{sec:reservation-ALU-LITE})}
~
{\begin{lstlisting}
stage RR:
  new argument3 = _ZX_6(unsigned6);
  new argument2 = PGR[registerP];
  new shift = _ZX_5(argument3);
  for (I = 0; I < 4; I += 1) {
    result1.32[I] = _ZX_32(argument2.32[I]) >> shift
  };
stage E1:
  PGR[registerM] = result1;
\end{lstlisting}}
\newpage

\paragraph{Encoding} Format \textsf{ALU\_WQSWRR} (Section~\ref{sec:format-ALU-WQSWRR}), \verb|exucode1 = 011|\\

\bitfields{ 1/P, 2/11, 1/1, 1/0, 3/011, 5/registerM, 1/--, 2/10, 3/100, 1/0, 6/registerY, 5/registerP, 1/1 }


\paragraph{Syntax} \texttt{srlwq} \emph{registerM} = \emph{registerP}, \emph{registerY}

\paragraph{Description} The \emph{registerP} is interpreted as four words packed into 128 bits. The \emph{registerP} words are shifted right by the \emph{registerY} modulo 32. The vacant positions are filled in with zeros. The four resulting words are packed and stored into the \emph{registerM}.


\paragraph{Execution} {Reservation table \textsf{ALU\_LITE} (Section~\ref{sec:reservation-ALU-LITE})}
~
{\begin{lstlisting}
stage RR:
  new argument3 = GPR[registerY];
  new argument2 = PGR[registerP];
  new shift = _ZX_5(argument3);
  for (I = 0; I < 4; I += 1) {
    result1.32[I] = _ZX_32(argument2.32[I]) >> shift
  };
stage E1:
  PGR[registerM] = result1;
\end{lstlisting}}
\newpage
\subsubsection[{\small SRSBX: Shift Right Symmetric Byte Hexadecuple}]{SRSBX\hfill Shift Right Symmetric Byte Hexadecuple}
\label{instr:SRSBX}\index{srsbx}
 
\paragraph{Encoding} Format \textsf{ALU\_BXSWRI} (Section~\ref{sec:format-ALU-BXSWRI}), \verb|exucode1 = 000|\\

\bitfields{ 1/P, 2/11, 1/1, 1/1, 3/000, 5/registerM, 1/--, 2/10, 3/100, 1/1, 6/unsigned6, 5/registerP, 1/1 }


\paragraph{Syntax} \texttt{srsbx} \emph{registerM} = \emph{registerP}, \emph{unsigned6}

\paragraph{Description} The \emph{registerP} is interpreted as sixteen bytes packed into 128 bits. If non-negative, each \emph{registerP} word is shifted right by the \emph{unsigned6} modulo 8. Else, each \emph{registerP} word is biased by (2**(\emph{unsigned6} modulo 8)) - 1 and shifted right by the \emph{unsigned6} modulo 8. Each result is stored into the corresponding \emph{registerM} byte. This instruction implements a signed division of the \emph{registerP} bytes by the power of two given by the \emph{unsigned6} modulo 8.


\paragraph{Execution} {Reservation table \textsf{ALU\_LITE} (Section~\ref{sec:reservation-ALU-LITE})}
~
{\begin{lstlisting}
stage RR:
  new argument3 = _ZX_6(unsigned6);
  new argument2 = PGR[registerP];
  new shift = _ZX_3(argument3);
  new bias = (1 << shift) - 1;
  for (I = 0; I < 16; I += 1) {
    result1.8[I] = (_SX_8(argument2.8[I]) + _SX_8(argument2.8[I]) < 0 ? bias : 0) >> shift
  };
stage E1:
  PGR[registerM] = result1;
\end{lstlisting}}
\newpage

\paragraph{Encoding} Format \textsf{ALU\_BXSWRR} (Section~\ref{sec:format-ALU-BXSWRR}), \verb|exucode1 = 000|\\

\bitfields{ 1/P, 2/11, 1/1, 1/0, 3/000, 5/registerM, 1/--, 2/10, 3/100, 1/1, 6/registerY, 5/registerP, 1/1 }


\paragraph{Syntax} \texttt{srsbx} \emph{registerM} = \emph{registerP}, \emph{registerY}

\paragraph{Description} The \emph{registerP} is interpreted as sixteen bytes packed into 128 bits. If non-negative, each \emph{registerP} word is shifted right by the \emph{registerY} modulo 8. Else, each \emph{registerP} word is biased by (2**(\emph{registerY} modulo 8)) - 1 and shifted right by the \emph{registerY} modulo 8. Each result is stored into the corresponding \emph{registerM} byte. This instruction implements a signed division of the \emph{registerP} bytes by the power of two given by the \emph{registerY} modulo 8.


\paragraph{Execution} {Reservation table \textsf{ALU\_LITE} (Section~\ref{sec:reservation-ALU-LITE})}
~
{\begin{lstlisting}
stage RR:
  new argument3 = GPR[registerY];
  new argument2 = PGR[registerP];
  new shift = _ZX_3(argument3);
  new bias = (1 << shift) - 1;
  for (I = 0; I < 16; I += 1) {
    result1.8[I] = (_SX_8(argument2.8[I]) + _SX_8(argument2.8[I]) < 0 ? bias : 0) >> shift
  };
stage E1:
  PGR[registerM] = result1;
\end{lstlisting}}
\newpage
\subsubsection[{\small SRSD: Shift Right Symmetric Double Word}]{SRSD\hfill Shift Right Symmetric Double Word}
\label{instr:SRSD}\index{srsd}
 
\paragraph{Encoding} Format \textsf{ALU\_DSWRI} (Section~\ref{sec:format-ALU-DSWRI}), \verb|exucode1 = 000|\\

\bitfields{ 1/P, 2/11, 1/0, 1/1, 3/000, 6/registerW, 2/10, 3/100, 1/0, 6/unsigned6, 6/registerZ }


\paragraph{Syntax} \texttt{srsd} \emph{registerW} = \emph{registerZ}, \emph{unsigned6}

\paragraph{Description} If non-negative, the \emph{registerZ} is shifted right by the \emph{unsigned6} modulo 64. Else, the \emph{registerZ} biased by (2**(\emph{unsigned6} modulo 64)) - 1 and shifted right by the \emph{unsigned6} modulo 64. The result is stored into the \emph{registerW}. This instruction implements a signed division of the \emph{registerZ} double word by a power of two.


\paragraph{Execution} {Reservation table \textsf{ALU\_LITE} (Section~\ref{sec:reservation-ALU-LITE})}
~
{\begin{lstlisting}
stage RR:
  new argument3 = _ZX_6(unsigned6);
  new argument2 = GPR[registerZ];
  new shift = _ZX_6(argument3);
  new bias = _SX_64(argument2) < 0 ? (1 << shift) - 1 : 0;
  new result1 = (_SX_64(argument2) + bias) >> shift;
stage E1:
  GPR[registerW] = result1;
\end{lstlisting}}
\newpage

\paragraph{Encoding} Format \textsf{ALU\_DSWRR} (Section~\ref{sec:format-ALU-DSWRR}), \verb|exucode1 = 000|\\

\bitfields{ 1/P, 2/11, 1/0, 1/0, 3/000, 6/registerW, 2/10, 3/100, 1/0, 6/registerY, 6/registerZ }


\paragraph{Syntax} \texttt{srsd} \emph{registerW} = \emph{registerZ}, \emph{registerY}

\paragraph{Description} If non-negative, the \emph{registerZ} is shifted right by the \emph{registerY} modulo 64. Else, the \emph{registerZ} biased by (2**(\emph{registerY} modulo 64)) - 1 and shifted right by the \emph{registerY} modulo 64. The result is stored into the \emph{registerW}. This instruction implements a signed division of the \emph{registerZ} double word by a power of two.


\paragraph{Execution} {Reservation table \textsf{ALU\_LITE} (Section~\ref{sec:reservation-ALU-LITE})}
~
{\begin{lstlisting}
stage RR:
  new argument3 = GPR[registerY];
  new argument2 = GPR[registerZ];
  new shift = _ZX_6(argument3);
  new bias = _SX_64(argument2) < 0 ? (1 << shift) - 1 : 0;
  new result1 = (_SX_64(argument2) + bias) >> shift;
stage E1:
  GPR[registerW] = result1;
\end{lstlisting}}
\newpage
\subsubsection[{\small SRSDP: Shift Right Symmetric Double Word Pair}]{SRSDP\hfill Shift Right Symmetric Double Word Pair}
\label{instr:SRSDP}\index{srsdp}
 
\paragraph{Encoding} Format \textsf{ALU\_DPSWRI} (Section~\ref{sec:format-ALU-DPSWRI}), \verb|exucode1 = 000|\\

\bitfields{ 1/P, 2/11, 1/1, 1/1, 3/000, 5/registerM, 1/--, 2/10, 3/100, 1/0, 6/unsigned6, 5/registerP, 1/0 }


\paragraph{Syntax} \texttt{srsdp} \emph{registerM} = \emph{registerP}, \emph{unsigned6}

\paragraph{Description} The \emph{registerP} is interpreted as two double words packed into 128 bits. If non-negative, each \emph{registerP} double word is shifted right by the \emph{unsigned6} modulo 64. Else, each \emph{registerP} double word is biased by (2**(\emph{unsigned6} modulo 64)) - 1 and shifted right by the \emph{unsigned6} modulo 64. Each result is stored into the corresponding \emph{registerM} double word. This instruction implements a signed division of the \emph{registerP} words by the power of two given by the \emph{unsigned6} modulo 64.


\paragraph{Execution} {Reservation table \textsf{ALU\_LITE} (Section~\ref{sec:reservation-ALU-LITE})}
~
{\begin{lstlisting}
stage RR:
  new argument3 = _ZX_6(unsigned6);
  new argument2 = PGR[registerP];
  new shift = _ZX_6(argument3);
  new bias = (1 << shift) - 1;
  for (I = 0; I < 2; I += 1) {
    result1.64[I] = (_SX_64(argument2.64[I]) + _SX_64(argument2.64[I]) < 0 ? bias : 0) >> shift
  };
stage E1:
  PGR[registerM] = result1;
\end{lstlisting}}
\newpage

\paragraph{Encoding} Format \textsf{ALU\_DPSWRR} (Section~\ref{sec:format-ALU-DPSWRR}), \verb|exucode1 = 000|\\

\bitfields{ 1/P, 2/11, 1/1, 1/0, 3/000, 5/registerM, 1/--, 2/10, 3/100, 1/0, 5/registerO, 1/--, 5/registerP, 1/0 }


\paragraph{Syntax} \texttt{srsdp} \emph{registerM} = \emph{registerP}, \emph{registerO}

\paragraph{Description} The \emph{registerP} is interpreted as two double words packed into 128 bits. If non-negative, each \emph{registerP} double word is shifted right by the \emph{registerO} modulo 64. Else, each \emph{registerP} double word is biased by (2**(\emph{registerO} modulo 64)) - 1 and shifted right by the \emph{registerO} modulo 64. Each result is stored into the corresponding \emph{registerM} double word. This instruction implements a signed division of the \emph{registerP} words by the power of two given by the \emph{registerO} modulo 64.


\paragraph{Execution} {Reservation table \textsf{ALU\_LITE} (Section~\ref{sec:reservation-ALU-LITE})}
~
{\begin{lstlisting}
stage RR:
  new argument3 = PGR[registerO];
  new argument2 = PGR[registerP];
  new shift = _ZX_6(argument3);
  new bias = (1 << shift) - 1;
  for (I = 0; I < 2; I += 1) {
    result1.64[I] = (_SX_64(argument2.64[I]) + _SX_64(argument2.64[I]) < 0 ? bias : 0) >> shift
  };
stage E1:
  PGR[registerM] = result1;
\end{lstlisting}}
\newpage
\subsubsection[{\small SRSHO: Shift Right Symmetric Half Word Octuple}]{SRSHO\hfill Shift Right Symmetric Half Word Octuple}
\label{instr:SRSHO}\index{srsho}
 
\paragraph{Encoding} Format \textsf{ALU\_HOSWRI} (Section~\ref{sec:format-ALU-HOSWRI}), \verb|exucode1 = 000|\\

\bitfields{ 1/P, 2/11, 1/1, 1/1, 3/000, 5/registerM, 1/--, 2/10, 3/100, 1/1, 6/unsigned6, 5/registerP, 1/0 }


\paragraph{Syntax} \texttt{srsho} \emph{registerM} = \emph{registerP}, \emph{unsigned6}

\paragraph{Description} The \emph{registerP} is interpreted as eight half words packed into 128 bits. If non-negative, each \emph{registerP} half word is shifted right by the \emph{unsigned6} modulo 16. Else, each \emph{registerP} half word is biased by (2**(\emph{unsigned6} modulo 16)) - 1 and shifted right by the \emph{unsigned6} modulo 16. Each result is stored into the corresponding \emph{registerM} half word. This instruction implements a signed division of the \emph{registerP} half words by the power of two given by the \emph{unsigned6} modulo 16.


\paragraph{Execution} {Reservation table \textsf{ALU\_LITE} (Section~\ref{sec:reservation-ALU-LITE})}
~
{\begin{lstlisting}
stage RR:
  new argument3 = _ZX_6(unsigned6);
  new argument2 = PGR[registerP];
  new shift = _ZX_4(argument3);
  new bias = (1 << shift) - 1;
  for (I = 0; I < 8; I += 1) {
    result1.16[I] = (_SX_16(argument2.16[I]) + _SX_16(argument2.16[I]) < 0 ? bias : 0) >> shift
  };
stage E1:
  PGR[registerM] = result1;
\end{lstlisting}}
\newpage

\paragraph{Encoding} Format \textsf{ALU\_HOSWRR} (Section~\ref{sec:format-ALU-HOSWRR}), \verb|exucode1 = 000|\\

\bitfields{ 1/P, 2/11, 1/1, 1/0, 3/000, 5/registerM, 1/--, 2/10, 3/100, 1/1, 6/registerY, 5/registerP, 1/0 }


\paragraph{Syntax} \texttt{srsho} \emph{registerM} = \emph{registerP}, \emph{registerY}

\paragraph{Description} The \emph{registerP} is interpreted as eight half words packed into 128 bits. If non-negative, each \emph{registerP} half word is shifted right by the \emph{registerY} modulo 16. Else, each \emph{registerP} half word is biased by (2**(\emph{registerY} modulo 16)) - 1 and shifted right by the \emph{registerY} modulo 16. Each result is stored into the corresponding \emph{registerM} half word. This instruction implements a signed division of the \emph{registerP} half words by the power of two given by the \emph{registerY} modulo 16.


\paragraph{Execution} {Reservation table \textsf{ALU\_LITE} (Section~\ref{sec:reservation-ALU-LITE})}
~
{\begin{lstlisting}
stage RR:
  new argument3 = GPR[registerY];
  new argument2 = PGR[registerP];
  new shift = _ZX_4(argument3);
  new bias = (1 << shift) - 1;
  for (I = 0; I < 8; I += 1) {
    result1.16[I] = (_SX_16(argument2.16[I]) + _SX_16(argument2.16[I]) < 0 ? bias : 0) >> shift
  };
stage E1:
  PGR[registerM] = result1;
\end{lstlisting}}
\newpage
\subsubsection[{\small SRSW: Shift Right Symmetric Word}]{SRSW\hfill Shift Right Symmetric Word}
\label{instr:SRSW}\index{srsw}
 
\paragraph{Encoding} Format \textsf{ALU\_WSWRR} (Section~\ref{sec:format-ALU-WSWRR}), \verb|exucode1 = 000|\\

\bitfields{ 1/P, 2/11, 1/0, 1/0, 3/000, 6/registerW, 2/11, 3/100, 1/signextw, 6/registerY, 6/registerZ }


\paragraph{Syntax} \texttt{srsw}\emph{signextw} \emph{registerW} = \emph{registerZ}, \emph{registerY}

\paragraph{Description} If non-negative, the \emph{registerZ} is shifted right by the \emph{registerY} modulo 32. Else, the \emph{registerZ} biased by (2**(\emph{registerY} modulo 32)) - 1 and shifted right by the \emph{registerY} modulo 32. The result with the 32 upper bits cleared is stored into the \emph{registerW}. This instruction implements a signed division of the \emph{registerZ} least significant word by a power of two.


\paragraph{Modifier(s)}
\noindent\begin{itemize}
\item signextw (cf. signextw in section \ref{par:modifier-signextw} on page \pageref{par:modifier-signextw})
\end{itemize}

\paragraph{Execution} {Reservation table \textsf{ALU\_LITE} (Section~\ref{sec:reservation-ALU-LITE})}
~
{\begin{lstlisting}
stage ID:
  new signextw = _ZX_1(signextw);
stage RR:
  new argument3 = GPR[registerY];
  new argument2 = GPR[registerZ];
  new shift = _ZX_5(argument3);
  new bias = _SX_32(argument2) < 0 ? (1 << shift) - 1 : 0;
  new result1 = (_SX_32(argument2) + bias) >> shift;
stage E1:
  GPR[registerW] = signextw ? _SX_32(result1) : _ZX_32(result1);
\end{lstlisting}}
\newpage

\paragraph{Encoding} Format \textsf{ALU\_WSWRI} (Section~\ref{sec:format-ALU-WSWRI}), \verb|exucode1 = 000|\\

\bitfields{ 1/P, 2/11, 1/0, 1/1, 3/000, 6/registerW, 2/11, 3/100, 1/signextw, 6/unsigned6, 6/registerZ }


\paragraph{Syntax} \texttt{srsw}\emph{signextw} \emph{registerW} = \emph{registerZ}, \emph{unsigned6}

\paragraph{Description} If non-negative, the \emph{registerZ} is shifted right by the \emph{unsigned6} modulo 32. Else, the \emph{registerZ} biased by (2**(\emph{unsigned6} modulo 32)) - 1 and shifted right by the \emph{unsigned6} modulo 32. The result with the 32 upper bits cleared is stored into the \emph{registerW}. This instruction implements a signed division of the \emph{registerZ} least significant word by a power of two.


\paragraph{Modifier(s)}
\noindent\begin{itemize}
\item signextw (cf. signextw in section \ref{par:modifier-signextw} on page \pageref{par:modifier-signextw})
\end{itemize}

\paragraph{Execution} {Reservation table \textsf{ALU\_LITE} (Section~\ref{sec:reservation-ALU-LITE})}
~
{\begin{lstlisting}
stage ID:
  new signextw = _ZX_1(signextw);
stage RR:
  new argument3 = _ZX_6(unsigned6);
  new argument2 = GPR[registerZ];
  new shift = _ZX_5(argument3);
  new bias = _SX_32(argument2) < 0 ? (1 << shift) - 1 : 0;
  new result1 = (_SX_32(argument2) + bias) >> shift;
stage E1:
  GPR[registerW] = signextw ? _SX_32(result1) : _ZX_32(result1);
\end{lstlisting}}
\newpage
\subsubsection[{\small SRSWQ: Shift Right Symmetric Word Quadruple}]{SRSWQ\hfill Shift Right Symmetric Word Quadruple}
\label{instr:SRSWQ}\index{srswq}
 
\paragraph{Encoding} Format \textsf{ALU\_WQSWRI} (Section~\ref{sec:format-ALU-WQSWRI}), \verb|exucode1 = 000|\\

\bitfields{ 1/P, 2/11, 1/1, 1/1, 3/000, 5/registerM, 1/--, 2/10, 3/100, 1/0, 6/unsigned6, 5/registerP, 1/1 }


\paragraph{Syntax} \texttt{srswq} \emph{registerM} = \emph{registerP}, \emph{unsigned6}

\paragraph{Description} The \emph{registerP} is interpreted as four words packed into 128 bits. If non-negative, each \emph{registerP} word is shifted right by the \emph{unsigned6} modulo 32. Else, each \emph{registerP} word is biased by (2**(\emph{unsigned6} modulo 32)) - 1 and shifted right by the \emph{unsigned6} modulo 32. The four resulting words are packed and stored into the \emph{registerM}. This instruction implements a signed division of the \emph{registerP} words by the power of two given by the \emph{unsigned6} modulo 32.


\paragraph{Execution} {Reservation table \textsf{ALU\_LITE} (Section~\ref{sec:reservation-ALU-LITE})}
~
{\begin{lstlisting}
stage RR:
  new argument3 = _ZX_6(unsigned6);
  new argument2 = PGR[registerP];
  new shift = _ZX_5(argument3);
  new bias = (1 << shift) - 1;
  for (I = 0; I < 4; I += 1) {
    result1.32[I] = (_SX_32(argument2.32[I]) + _SX_32(argument2.32[I]) < 0 ? bias : 0) >> shift
  };
stage E1:
  PGR[registerM] = result1;
\end{lstlisting}}
\newpage

\paragraph{Encoding} Format \textsf{ALU\_WQSWRR} (Section~\ref{sec:format-ALU-WQSWRR}), \verb|exucode1 = 000|\\

\bitfields{ 1/P, 2/11, 1/1, 1/0, 3/000, 5/registerM, 1/--, 2/10, 3/100, 1/0, 6/registerY, 5/registerP, 1/1 }


\paragraph{Syntax} \texttt{srswq} \emph{registerM} = \emph{registerP}, \emph{registerY}

\paragraph{Description} The \emph{registerP} is interpreted as four words packed into 128 bits. If non-negative, each \emph{registerP} word is shifted right by the \emph{registerY} modulo 32. Else, each \emph{registerP} word is biased by (2**(\emph{registerY} modulo 32)) - 1 and shifted right by the \emph{registerY} modulo 32. The four resulting words are packed and stored into the \emph{registerM}. This instruction implements a signed division of the \emph{registerP} words by the power of two given by the \emph{registerY} modulo 32.


\paragraph{Execution} {Reservation table \textsf{ALU\_LITE} (Section~\ref{sec:reservation-ALU-LITE})}
~
{\begin{lstlisting}
stage RR:
  new argument3 = GPR[registerY];
  new argument2 = PGR[registerP];
  new shift = _ZX_5(argument3);
  new bias = (1 << shift) - 1;
  for (I = 0; I < 4; I += 1) {
    result1.32[I] = (_SX_32(argument2.32[I]) + _SX_32(argument2.32[I]) < 0 ? bias : 0) >> shift
  };
stage E1:
  PGR[registerM] = result1;
\end{lstlisting}}
\newpage
\subsubsection[{\small STSUD: Subtract, Test and Shift Unsigned Double Word}]{STSUD\hfill Subtract, Test and Shift Unsigned Double Word}
\label{instr:STSUD}\index{stsud}
 
\paragraph{Encoding} Format \textsf{ALU\_DWRRS} (Section~\ref{sec:format-ALU-DWRRS}), \verb|exucode2 = 1111|\\

\bitfields{ 1/P, 2/11, 1/0, 4/1111, 6/registerW, 2/10, 3/010, 1/0, 6/registerY, 6/registerZ }


\paragraph{Syntax} \texttt{stsud} \emph{registerW} = \emph{registerZ}, \emph{registerY}

\paragraph{Description} The \emph{registerZ} is subtracted from the \emph{registerY} If positive, the result is shifted left one bit and OR-ed with 0x1. Else, the \emph{registerY} is shifted left one bit. The result is stored into the \emph{registerW}. This instruction is provided to support double word unsigned integer division.


\paragraph{Execution} {Reservation table \textsf{ALU\_TINY} (Section~\ref{sec:reservation-ALU-TINY})}
~
{\begin{lstlisting}
stage RR:
  new argument3 = GPR[registerY];
  new argument2 = GPR[registerZ];
  new result1 = _ZX_64(argument3) >= _ZX_64(argument2) ? ((_ZX_64(argument3) - _ZX_64(argument2)) << 1) | 1 : _ZX_64(argument3) << 1;
stage E1:
  GPR[registerW] = result1;
\end{lstlisting}}
\newpage
\subsubsection[{\small STSUDP: Subtract, Test and Shift Unsigned Double Word Pair}]{STSUDP\hfill Subtract, Test and Shift Unsigned Double Word Pair}
\label{instr:STSUDP}\index{stsudp}
 
\paragraph{Encoding} Format \textsf{ALU\_DPWRRS} (Section~\ref{sec:format-ALU-DPWRRS}), \verb|exucode2 = 1111|\\

\bitfields{ 1/P, 2/11, 1/1, 4/1111, 5/registerM, 1/--, 2/10, 3/010, 1/0, 5/registerO, 1/--, 5/registerP, 1/0 }


\paragraph{Syntax} \texttt{stsudp} \emph{registerM} = \emph{registerP}, \emph{registerO}

\paragraph{Description} The \emph{registerP} interpreted as a double word pair is subtracted from the \emph{registerO} interpreted as a double word pair. For each result, if the difference is positive it is shifted left one bit and OR-ed with 0x1. Else, the result is the corresponding word of the \emph{registerO} is shifted left one bit. The two resulting double words are packed and stored into the \emph{registerM}.


\paragraph{Execution} {Reservation table \textsf{ALU\_LITE} (Section~\ref{sec:reservation-ALU-LITE})}
~
{\begin{lstlisting}
stage RR:
  new argument3 = PGR[registerO];
  new argument2 = PGR[registerP];
  for (I = 0; I < 2; I += 1) {
    result1.64[I] = _ZX_64(argument3.64[I]) >= _ZX_64(argument2.64[I]) ? ((_ZX_64(argument3.64[I]) - _ZX_64(argument2.64[I])) << 1) | 1 : _ZX_64(argument3.64[I]) << 1
  };
stage E1:
  PGR[registerM] = result1;
\end{lstlisting}}
\newpage
\subsubsection[{\small STSUHO: Subtract, Test and Shift Unsigned Half Word Octuple}]{STSUHO\hfill Subtract, Test and Shift Unsigned Half Word Octuple}
\label{instr:STSUHO}\index{stsuho}
 
\paragraph{Encoding} Format \textsf{ALU\_HOWRRS} (Section~\ref{sec:format-ALU-HOWRRS}), \verb|exucode2 = 1111|\\

\bitfields{ 1/P, 2/11, 1/1, 4/1111, 5/registerM, 1/--, 2/10, 3/010, 1/1, 5/registerO, 1/--, 5/registerP, 1/0 }


\paragraph{Syntax} \texttt{stsuho} \emph{registerM} = \emph{registerP}, \emph{registerO}

\paragraph{Description} The \emph{registerP} interpreted as a half word octuple is subtracted from the \emph{registerO} interpreted as a half word octuple. For each result, if the difference is positive it is shifted left one bit and OR-ed with 0x1. Else, the result is the corresponding half word of the \emph{registerO} is shifted left one bit. The eight resulting half word are packed and stored into the \emph{registerM}.


\paragraph{Execution} {Reservation table \textsf{ALU\_LITE} (Section~\ref{sec:reservation-ALU-LITE})}
~
{\begin{lstlisting}
stage RR:
  new argument3 = PGR[registerO];
  new argument2 = PGR[registerP];
  for (I = 0; I < 8; I += 1) {
    result1.16[I] = _ZX_16(argument3.16[I]) >= _ZX_16(argument2.16[I]) ? ((_ZX_16(argument3.16[I]) - _ZX_16(argument2.16[I])) << 1) | 1 : _ZX_16(argument3.16[I]) << 1
  };
stage E1:
  PGR[registerM] = result1;
\end{lstlisting}}
\newpage
\subsubsection[{\small STSUW: Subtract, Test and Shift Unsigned Word}]{STSUW\hfill Subtract, Test and Shift Unsigned Word}
\label{instr:STSUW}\index{stsuw}
 
\paragraph{Encoding} Format \textsf{ALU\_WWRRS} (Section~\ref{sec:format-ALU-WWRRS}), \verb|exucode2 = 1111|\\

\bitfields{ 1/P, 2/11, 1/0, 4/1111, 6/registerW, 2/11, 3/010, 1/signextw, 6/registerY, 6/registerZ }


\paragraph{Syntax} \texttt{stsuw}\emph{signextw} \emph{registerW} = \emph{registerZ}, \emph{registerY}

\paragraph{Description} The \emph{registerZ} is subtracted from the \emph{registerY}. If positive, the result is shifted left one bit and OR-ed with 0x1. Else, the \emph{registerY} is shifted left one bit. The result with the 32 upper bits cleared is stored into the \emph{registerW}. This instruction is provided to support unsigned integer division.


\paragraph{Modifier(s)}
\noindent\begin{itemize}
\item signextw (cf. signextw in section \ref{par:modifier-signextw} on page \pageref{par:modifier-signextw})
\end{itemize}

\paragraph{Execution} {Reservation table \textsf{ALU\_TINY} (Section~\ref{sec:reservation-ALU-TINY})}
~
{\begin{lstlisting}
stage ID:
  new signextw = _ZX_1(signextw);
stage RR:
  new argument3 = GPR[registerY];
  new argument2 = GPR[registerZ];
  new result1 = argument3.32[0] >= argument2.32[0] ? ((argument3.32[0] - argument2.32[0]) << 1) | 1 : argument3.32[0] << 1;
stage E1:
  GPR[registerW] = signextw ? _SX_32(result1) : _ZX_32(result1);
\end{lstlisting}}
\newpage
\subsubsection[{\small STSUWQ: Subtract, Test and Shift Unsigned Word Quadruple}]{STSUWQ\hfill Subtract, Test and Shift Unsigned Word Quadruple}
\label{instr:STSUWQ}\index{stsuwq}
 
\paragraph{Encoding} Format \textsf{ALU\_WQWRRS} (Section~\ref{sec:format-ALU-WQWRRS}), \verb|exucode2 = 1111|\\

\bitfields{ 1/P, 2/11, 1/1, 4/1111, 5/registerM, 1/--, 2/10, 3/010, 1/0, 5/registerO, 1/--, 5/registerP, 1/1 }


\paragraph{Syntax} \texttt{stsuwq} \emph{registerM} = \emph{registerP}, \emph{registerO}

\paragraph{Description} The \emph{registerP} interpreted as a word quadruple is subtracted from the \emph{registerO} interpreted as a word quadruple. For each result, if the difference is positive it is shifted left one bit and OR-ed with 0x1. Else, the result is the corresponding word of the \emph{registerO} is shifted left one bit. The four resulting word are packed and stored into the \emph{registerM}.


\paragraph{Execution} {Reservation table \textsf{ALU\_LITE} (Section~\ref{sec:reservation-ALU-LITE})}
~
{\begin{lstlisting}
stage RR:
  new argument3 = PGR[registerO];
  new argument2 = PGR[registerP];
  for (I = 0; I < 4; I += 1) {
    result1.32[I] = _ZX_32(argument3.32[I]) >= _ZX_32(argument2.32[I]) ? ((_ZX_32(argument3.32[I]) - _ZX_32(argument2.32[I])) << 1) | 1 : _ZX_32(argument3.32[I]) << 1
  };
stage E1:
  PGR[registerM] = result1;
\end{lstlisting}}
\newpage
\subsubsection[{\small TRUNCDWQ: Truncate Double Word to Word Quadruple}]{TRUNCDWQ\hfill Truncate Double Word to Word Quadruple}
\label{instr:TRUNCDWQ}\index{truncdwq}
 
\paragraph{Encoding} Format \textsf{ALU\_WQNWRR} (Section~\ref{sec:format-ALU-WQNWRR}), \verb|alucode8 = 1000|\\

\bitfields{ 1/P, 2/11, 1/1, 4/1111, 5/registerM, 1/--, 2/10, 3/101, 1/0, 4/1000, 1/ziplanes, 1/0, 4/registerR, 1/--, 1/1 }


\paragraph{Syntax} \texttt{truncdwq}\emph{ziplanes} \emph{registerM} = \emph{registerR}

\paragraph{Description} The \emph{ziplanes} and the \emph{registerR} are interpreted as two double words packed into 128 bits. The \emph{ziplanes} and the \emph{registerR} double words are converted to 32-bit words by removing the high bits. If \emph{ziplanes}, the resulting words are interleaved with the \emph{registerR} words in even lanes and the \emph{ziplanes} words in odd lanes. The four resulting words are packed and stored back into the \emph{registerM}.


\paragraph{Modifier(s)}
\noindent\begin{itemize}
\item ziplanes (cf. ziplanes in section \ref{par:modifier-ziplanes} on page \pageref{par:modifier-ziplanes})
\end{itemize}

\paragraph{Execution} {Reservation table \textsf{ALU\_LITE} (Section~\ref{sec:reservation-ALU-LITE})}
~
{\begin{lstlisting}
stage ID:
  new ziplanes = _ZX_1(ziplanes);
stage RR:
  new argument2 = QGR[registerR];
  if (ziplanes) {
    result1.32[0] = argument2.64[0];
    result1.32[1] = argument3.64[0];
    result1.32[2] = argument2.64[1];
  } else {
    result1.32[0] = argument2.64[0];
    result1.32[1] = argument2.64[1];
    result1.32[2] = argument3.64[0];
  }
  result1.32[3] = argument3.64[1];
stage E1:
  PGR[registerM] = result1;
\end{lstlisting}}
\newpage
\subsubsection[{\small TRUNCHBX: Truncate Half Word to Byte Hexadecuple}]{TRUNCHBX\hfill Truncate Half Word to Byte Hexadecuple}
\label{instr:TRUNCHBX}\index{trunchbx}
 
\paragraph{Encoding} Format \textsf{ALU\_BXNWRR} (Section~\ref{sec:format-ALU-BXNWRR}), \verb|alucode8 = 1000|\\

\bitfields{ 1/P, 2/11, 1/1, 4/1111, 5/registerM, 1/--, 2/10, 3/101, 1/1, 4/1000, 1/ziplanes, 1/0, 5/registerP, 1/1 }


\paragraph{Syntax} \texttt{trunchbx}\emph{ziplanes} \emph{registerM} = \emph{registerP}

\paragraph{Description} The \emph{ziplanes} and the \emph{registerP} are interpreted as eight half words packed into 128 bits. The \emph{ziplanes} and the \emph{registerP} half words are converted to bytes by removing the high bits. If \emph{ziplanes}, the resulting bytes are interleaved with the \emph{registerP} bytes in even lanes and the \emph{ziplanes} bytes in odd lanes. The sixteen resulting bytes are packed and stored back into the \emph{registerM}.


\paragraph{Modifier(s)}
\noindent\begin{itemize}
\item ziplanes (cf. ziplanes in section \ref{par:modifier-ziplanes} on page \pageref{par:modifier-ziplanes})
\end{itemize}

\paragraph{Execution} {Reservation table \textsf{ALU\_LITE} (Section~\ref{sec:reservation-ALU-LITE})}
~
{\begin{lstlisting}
stage ID:
  new ziplanes = _ZX_1(ziplanes);
stage RR:
  new argument2 = PGR[registerP];
  if (ziplanes) {
    result1.8[0] = argument2.16[0];
    result1.8[1] = argument3.16[0];
    result1.8[2] = argument2.16[1];
    result1.8[3] = argument3.16[1];
    result1.8[4] = argument2.16[2];
    result1.8[5] = argument3.16[2];
    result1.8[6] = argument2.16[3];
    result1.8[7] = argument3.16[3];
    result1.8[8] = argument2.16[4];
    result1.8[9] = argument3.16[4];
    result1.8[10] = argument2.16[5];
    result1.8[11] = argument3.16[5];
    result1.8[12] = argument2.16[6];
    result1.8[13] = argument3.16[6];
    result1.8[14] = argument2.16[7];
  } else {
    result1.8[0] = argument2.16[0];
    result1.8[1] = argument2.16[1];
    result1.8[2] = argument2.16[2];
    result1.8[3] = argument2.16[3];
    result1.8[4] = argument2.16[4];
    result1.8[5] = argument2.16[5];
    result1.8[6] = argument2.16[6];
    result1.8[7] = argument2.16[7];
    result1.8[8] = argument3.16[0];
    result1.8[9] = argument3.16[1];
    result1.8[10] = argument3.16[2];
    result1.8[11] = argument3.16[3];
    result1.8[12] = argument3.16[4];
    result1.8[13] = argument3.16[5];
    result1.8[14] = argument3.16[6];
  }
  result1.8[15] = argument3.16[7];
stage E1:
  PGR[registerM] = result1;
\end{lstlisting}}
\newpage
\subsubsection[{\small TRUNCWHO: Truncate Word to Half Word Octuple}]{TRUNCWHO\hfill Truncate Word to Half Word Octuple}
\label{instr:TRUNCWHO}\index{truncwho}
 
\paragraph{Encoding} Format \textsf{ALU\_HONWRR} (Section~\ref{sec:format-ALU-HONWRR}), \verb|alucode8 = 1000|\\

\bitfields{ 1/P, 2/11, 1/1, 4/1111, 5/registerM, 1/--, 2/10, 3/101, 1/1, 4/1000, 1/ziplanes, 1/0, 5/registerP, 1/0 }


\paragraph{Syntax} \texttt{truncwho}\emph{ziplanes} \emph{registerM} = \emph{registerP}

\paragraph{Description} The \emph{ziplanes} and the \emph{registerP} are interpreted as four words packed into 128 bits. The \emph{ziplanes} and the \emph{registerP} words are converted to 16-bit half words by removing the high bits. If \emph{ziplanes}, the resulting half words are interleaved with the \emph{registerP} half words in even lanes and the \emph{ziplanes} half words in odd lanes. The eight resulting half words are packed and stored back into the \emph{registerM}.


\paragraph{Modifier(s)}
\noindent\begin{itemize}
\item ziplanes (cf. ziplanes in section \ref{par:modifier-ziplanes} on page \pageref{par:modifier-ziplanes})
\end{itemize}

\paragraph{Execution} {Reservation table \textsf{ALU\_LITE} (Section~\ref{sec:reservation-ALU-LITE})}
~
{\begin{lstlisting}
stage ID:
  new ziplanes = _ZX_1(ziplanes);
stage RR:
  new argument2 = PGR[registerP];
  if (ziplanes) {
    result1.16[0] = argument2.32[0];
    result1.16[1] = argument3.32[0];
    result1.16[2] = argument2.32[1];
    result1.16[3] = argument3.32[1];
    result1.16[4] = argument2.32[2];
    result1.16[5] = argument3.32[2];
    result1.16[6] = argument2.32[3];
  } else {
    result1.16[0] = argument2.32[0];
    result1.16[1] = argument2.32[1];
    result1.16[2] = argument2.32[2];
    result1.16[3] = argument2.32[3];
    result1.16[4] = argument3.32[0];
    result1.16[5] = argument3.32[1];
    result1.16[6] = argument3.32[2];
  }
  result1.16[7] = argument3.32[3];
stage E1:
  PGR[registerM] = result1;
\end{lstlisting}}
\newpage
\subsubsection[{\small WIDENQBHO: Widen Q7 Byte to Half Word Octuple}]{WIDENQBHO\hfill Widen Q7 Byte to Half Word Octuple}
\label{instr:WIDENQBHO}\index{widenqbho}
 
\paragraph{Encoding} Format \textsf{ALU\_HOWIDEWR} (Section~\ref{sec:format-ALU-HOWIDEWR}), \verb|alucode8 = 1010|\\

\bitfields{ 1/P, 2/11, 1/1, 4/1111, 5/registerM, 1/--, 2/10, 3/101, 1/1, 4/1010, 1/mostsig, 1/1, 5/registerP, 1/0 }


\paragraph{Syntax} \texttt{widenqbho}\emph{mostsig} \emph{registerM} = \emph{registerP}

\paragraph{Description} The \emph{registerP} is interpreted as sixteen bytes packed into 128 bits. The least or most significant bytes of the \emph{registerP} are selected, depending on \emph{mostsig}. These eight bytes interpreted as Q7 numbers are mapped to Q15 numbers, packed and stored into the \emph{registerM}.


\paragraph{Modifier(s)}
\noindent\begin{itemize}
\item mostsig (cf. mostsig in section \ref{par:modifier-mostsig} on page \pageref{par:modifier-mostsig})
\end{itemize}

\paragraph{Execution} {Reservation table \textsf{ALU\_LITE} (Section~\ref{sec:reservation-ALU-LITE})}
~
{\begin{lstlisting}
stage ID:
  new mostsig = _ZX_1(mostsig);
stage RR:
  new argument2 = PGR[registerP];
  if (mostsig) {
    argument2 = argument2 >> 64;
  }
  for (I = 0; I < 8; I += 1) {
    result1.16[I] = argument2.8[I] << 8
  };
stage E1:
  PGR[registerM] = result1;
\end{lstlisting}}
\newpage
\subsubsection[{\small WIDENQHWQ: Widen Q15 Half Word to Word Quadruple}]{WIDENQHWQ\hfill Widen Q15 Half Word to Word Quadruple}
\label{instr:WIDENQHWQ}\index{widenqhwq}
 
\paragraph{Encoding} Format \textsf{ALU\_WQWIDEWR} (Section~\ref{sec:format-ALU-WQWIDEWR}), \verb|alucode8 = 1010|\\

\bitfields{ 1/P, 2/11, 1/1, 4/1111, 5/registerM, 1/--, 2/10, 3/101, 1/0, 4/1010, 1/mostsig, 1/1, 5/registerP, 1/1 }


\paragraph{Syntax} \texttt{widenqhwq}\emph{mostsig} \emph{registerM} = \emph{registerP}

\paragraph{Description} The \emph{registerP} is interpreted as eight half words packed into 128 bits. The least or most significant half words of the \emph{registerP} are selected, depending on \emph{mostsig}. These four half words interpreted as Q15 numbers are mapped to Q31 numbers, packed and stored into the \emph{registerM}.


\paragraph{Modifier(s)}
\noindent\begin{itemize}
\item mostsig (cf. mostsig in section \ref{par:modifier-mostsig} on page \pageref{par:modifier-mostsig})
\end{itemize}

\paragraph{Execution} {Reservation table \textsf{ALU\_LITE} (Section~\ref{sec:reservation-ALU-LITE})}
~
{\begin{lstlisting}
stage ID:
  new mostsig = _ZX_1(mostsig);
stage RR:
  new argument2 = PGR[registerP];
  if (mostsig) {
    argument2 = argument2 >> 64;
  }
  for (I = 0; I < 4; I += 1) {
    result1.32[I] = argument2.16[I] << 16
  };
stage E1:
  PGR[registerM] = result1;
\end{lstlisting}}
\newpage
\subsubsection[{\small WIDENQWDP: Widen Q31 Word to Double Word Pair}]{WIDENQWDP\hfill Widen Q31 Word to Double Word Pair}
\label{instr:WIDENQWDP}\index{widenqwdp}
 
\paragraph{Encoding} Format \textsf{ALU\_DPWIDEWR} (Section~\ref{sec:format-ALU-DPWIDEWR}), \verb|alucode8 = 1010|\\

\bitfields{ 1/P, 2/11, 1/1, 4/1111, 5/registerM, 1/--, 2/10, 3/101, 1/0, 4/1010, 1/mostsig, 1/1, 5/registerP, 1/0 }


\paragraph{Syntax} \texttt{widenqwdp}\emph{mostsig} \emph{registerM} = \emph{registerP}

\paragraph{Description} The \emph{registerP} is interpreted as four words packed into 128 bits. The least or most significant words of the \emph{registerP} are selected, depending on \emph{mostsig}. These two words interpreted as Q31 numbers are mapped to Q63 numbers, packed and stored into the \emph{registerM}.


\paragraph{Modifier(s)}
\noindent\begin{itemize}
\item mostsig (cf. mostsig in section \ref{par:modifier-mostsig} on page \pageref{par:modifier-mostsig})
\end{itemize}

\paragraph{Execution} {Reservation table \textsf{ALU\_LITE} (Section~\ref{sec:reservation-ALU-LITE})}
~
{\begin{lstlisting}
stage ID:
  new mostsig = _ZX_1(mostsig);
stage RR:
  new argument2 = PGR[registerP];
  if (mostsig) {
    argument2 = argument2 >> 64;
  }
  for (I = 0; I < 2; I += 1) {
    result1.64[I] = argument2.32[I] << 32
  };
stage E1:
  PGR[registerM] = result1;
\end{lstlisting}}
\newpage
\subsubsection[{\small WIDENSBHO: Sign Widen Byte to Half Word Octuple}]{WIDENSBHO\hfill Sign Widen Byte to Half Word Octuple}
\label{instr:WIDENSBHO}\index{widensbho}
 
\paragraph{Encoding} Format \textsf{ALU\_HOWIDEWR} (Section~\ref{sec:format-ALU-HOWIDEWR}), \verb|alucode8 = 1001|\\

\bitfields{ 1/P, 2/11, 1/1, 4/1111, 5/registerM, 1/--, 2/10, 3/101, 1/1, 4/1001, 1/mostsig, 1/1, 5/registerP, 1/0 }


\paragraph{Syntax} \texttt{widensbho}\emph{mostsig} \emph{registerM} = \emph{registerP}

\paragraph{Description} The \emph{registerP} is interpreted as sixteen bytes packed into 128 bits. The least or most significant bytes of the \emph{registerP} are selected, depending on \emph{mostsig}. These eight bytes are sign-extended to 16 bits, packed and stored into the \emph{registerM}.


\paragraph{Modifier(s)}
\noindent\begin{itemize}
\item mostsig (cf. mostsig in section \ref{par:modifier-mostsig} on page \pageref{par:modifier-mostsig})
\end{itemize}

\paragraph{Execution} {Reservation table \textsf{ALU\_LITE} (Section~\ref{sec:reservation-ALU-LITE})}
~
{\begin{lstlisting}
stage ID:
  new mostsig = _ZX_1(mostsig);
stage RR:
  new argument2 = PGR[registerP];
  if (mostsig) {
    argument2 = argument2 >> 64;
  }
  for (I = 0; I < 8; I += 1) {
    result1.16[I] = _SX_8(argument2.8[I])
  };
stage E1:
  PGR[registerM] = result1;
\end{lstlisting}}
\newpage
\subsubsection[{\small WIDENSHWQ: Sign Widen Half Word to Word Quadruple}]{WIDENSHWQ\hfill Sign Widen Half Word to Word Quadruple}
\label{instr:WIDENSHWQ}\index{widenshwq}
 
\paragraph{Encoding} Format \textsf{ALU\_WQWIDEWR} (Section~\ref{sec:format-ALU-WQWIDEWR}), \verb|alucode8 = 1001|\\

\bitfields{ 1/P, 2/11, 1/1, 4/1111, 5/registerM, 1/--, 2/10, 3/101, 1/0, 4/1001, 1/mostsig, 1/1, 5/registerP, 1/1 }


\paragraph{Syntax} \texttt{widenshwq}\emph{mostsig} \emph{registerM} = \emph{registerP}

\paragraph{Description} The \emph{registerP} is interpreted as eight half words packed into 128 bits. The least or most significant half words of the \emph{registerP} are selected, depending on \emph{mostsig}. These four half words are sign-extended to 32 bits, packed and stored into the \emph{registerM}.


\paragraph{Modifier(s)}
\noindent\begin{itemize}
\item mostsig (cf. mostsig in section \ref{par:modifier-mostsig} on page \pageref{par:modifier-mostsig})
\end{itemize}

\paragraph{Execution} {Reservation table \textsf{ALU\_LITE} (Section~\ref{sec:reservation-ALU-LITE})}
~
{\begin{lstlisting}
stage ID:
  new mostsig = _ZX_1(mostsig);
stage RR:
  new argument2 = PGR[registerP];
  if (mostsig) {
    argument2 = argument2 >> 64;
  }
  for (I = 0; I < 4; I += 1) {
    result1.32[I] = _SX_16(argument2.16[I])
  };
stage E1:
  PGR[registerM] = result1;
\end{lstlisting}}
\newpage
\subsubsection[{\small WIDENSWDP: Sign Widen Word to Double Word Pair}]{WIDENSWDP\hfill Sign Widen Word to Double Word Pair}
\label{instr:WIDENSWDP}\index{widenswdp}
 
\paragraph{Encoding} Format \textsf{ALU\_DPWIDEWR} (Section~\ref{sec:format-ALU-DPWIDEWR}), \verb|alucode8 = 1001|\\

\bitfields{ 1/P, 2/11, 1/1, 4/1111, 5/registerM, 1/--, 2/10, 3/101, 1/0, 4/1001, 1/mostsig, 1/1, 5/registerP, 1/0 }


\paragraph{Syntax} \texttt{widenswdp}\emph{mostsig} \emph{registerM} = \emph{registerP}

\paragraph{Description} The \emph{registerP} is interpreted as four words packed into 128 bits. The least or most significant words of the \emph{registerP} are selected, depending on \emph{mostsig}. These two words are sign-extended to 64 bits, packed and stored into the \emph{registerM}.


\paragraph{Modifier(s)}
\noindent\begin{itemize}
\item mostsig (cf. mostsig in section \ref{par:modifier-mostsig} on page \pageref{par:modifier-mostsig})
\end{itemize}

\paragraph{Execution} {Reservation table \textsf{ALU\_LITE} (Section~\ref{sec:reservation-ALU-LITE})}
~
{\begin{lstlisting}
stage ID:
  new mostsig = _ZX_1(mostsig);
stage RR:
  new argument2 = PGR[registerP];
  if (mostsig) {
    argument2 = argument2 >> 64;
  }
  for (I = 0; I < 2; I += 1) {
    result1.64[I] = _SX_32(argument2.32[I])
  };
stage E1:
  PGR[registerM] = result1;
\end{lstlisting}}
\newpage
\subsubsection[{\small WIDENZBHO: Zero Widen Byte to Half Word Octuple}]{WIDENZBHO\hfill Zero Widen Byte to Half Word Octuple}
\label{instr:WIDENZBHO}\index{widenzbho}
 
\paragraph{Encoding} Format \textsf{ALU\_HOWIDEWR} (Section~\ref{sec:format-ALU-HOWIDEWR}), \verb|alucode8 = 1000|\\

\bitfields{ 1/P, 2/11, 1/1, 4/1111, 5/registerM, 1/--, 2/10, 3/101, 1/1, 4/1000, 1/mostsig, 1/1, 5/registerP, 1/0 }


\paragraph{Syntax} \texttt{widenzbho}\emph{mostsig} \emph{registerM} = \emph{registerP}

\paragraph{Description} The \emph{registerP} is interpreted as sixteen bytes packed into 128 bits. The least or most significant bytes of the \emph{registerP} are selected, depending on \emph{mostsig}. These eight bytes are zero-extended to 16 bits, packed and stored into the \emph{registerM}.


\paragraph{Modifier(s)}
\noindent\begin{itemize}
\item mostsig (cf. mostsig in section \ref{par:modifier-mostsig} on page \pageref{par:modifier-mostsig})
\end{itemize}

\paragraph{Execution} {Reservation table \textsf{ALU\_LITE} (Section~\ref{sec:reservation-ALU-LITE})}
~
{\begin{lstlisting}
stage ID:
  new mostsig = _ZX_1(mostsig);
stage RR:
  new argument2 = PGR[registerP];
  if (mostsig) {
    argument2 = argument2 >> 64;
  }
  for (I = 0; I < 8; I += 1) {
    result1.16[I] = _ZX_8(argument2.8[I])
  };
stage E1:
  PGR[registerM] = result1;
\end{lstlisting}}
\newpage
\subsubsection[{\small WIDENZHWQ: Zero Widen Half Word to Word Quadruple}]{WIDENZHWQ\hfill Zero Widen Half Word to Word Quadruple}
\label{instr:WIDENZHWQ}\index{widenzhwq}
 
\paragraph{Encoding} Format \textsf{ALU\_WQWIDEWR} (Section~\ref{sec:format-ALU-WQWIDEWR}), \verb|alucode8 = 1000|\\

\bitfields{ 1/P, 2/11, 1/1, 4/1111, 5/registerM, 1/--, 2/10, 3/101, 1/0, 4/1000, 1/mostsig, 1/1, 5/registerP, 1/1 }


\paragraph{Syntax} \texttt{widenzhwq}\emph{mostsig} \emph{registerM} = \emph{registerP}

\paragraph{Description} The \emph{registerP} is interpreted as eight half words packed into 128 bits. The least or most significant half words of the \emph{registerP} are selected, depending on \emph{mostsig}. These four half words are zero-extended to 32 bits, packed and stored into the \emph{registerM}.


\paragraph{Modifier(s)}
\noindent\begin{itemize}
\item mostsig (cf. mostsig in section \ref{par:modifier-mostsig} on page \pageref{par:modifier-mostsig})
\end{itemize}

\paragraph{Execution} {Reservation table \textsf{ALU\_LITE} (Section~\ref{sec:reservation-ALU-LITE})}
~
{\begin{lstlisting}
stage ID:
  new mostsig = _ZX_1(mostsig);
stage RR:
  new argument2 = PGR[registerP];
  if (mostsig) {
    argument2 = argument2 >> 64;
  }
  for (I = 0; I < 4; I += 1) {
    result1.32[I] = _ZX_16(argument2.16[I])
  };
stage E1:
  PGR[registerM] = result1;
\end{lstlisting}}
\newpage
\subsubsection[{\small WIDENZWDP: Zero Widen Word to Double Word Pair}]{WIDENZWDP\hfill Zero Widen Word to Double Word Pair}
\label{instr:WIDENZWDP}\index{widenzwdp}
 
\paragraph{Encoding} Format \textsf{ALU\_DPWIDEWR} (Section~\ref{sec:format-ALU-DPWIDEWR}), \verb|alucode8 = 1000|\\

\bitfields{ 1/P, 2/11, 1/1, 4/1111, 5/registerM, 1/--, 2/10, 3/101, 1/0, 4/1000, 1/mostsig, 1/1, 5/registerP, 1/0 }


\paragraph{Syntax} \texttt{widenzwdp}\emph{mostsig} \emph{registerM} = \emph{registerP}

\paragraph{Description} The \emph{registerP} is interpreted as four words packed into 128 bits. The least or most significant words of the \emph{registerP} are selected, depending on \emph{mostsig}. These two words are zero-extended to 64 bits, packed and stored into the \emph{registerM}.


\paragraph{Modifier(s)}
\noindent\begin{itemize}
\item mostsig (cf. mostsig in section \ref{par:modifier-mostsig} on page \pageref{par:modifier-mostsig})
\end{itemize}

\paragraph{Execution} {Reservation table \textsf{ALU\_LITE} (Section~\ref{sec:reservation-ALU-LITE})}
~
{\begin{lstlisting}
stage ID:
  new mostsig = _ZX_1(mostsig);
stage RR:
  new argument2 = PGR[registerP];
  if (mostsig) {
    argument2 = argument2 >> 64;
  }
  for (I = 0; I < 2; I += 1) {
    result1.64[I] = _ZX_32(argument2.32[I])
  };
stage E1:
  PGR[registerM] = result1;
\end{lstlisting}}
\newpage
\subsubsection[{\small XMOVETD: Move To Extension Double Word}]{XMOVETD\hfill Move To Extension Double Word}
\label{instr:XMOVETD}\index{xmovetd}
 
\paragraph{Encoding} Format \textsf{ALU\_XMOVETD0} (Section~\ref{sec:format-ALU-XMOVETD0}), \verb|alucode3 = 010|\\

\bitfields{ 1/P, 2/11, 1/0, 2/10, 2/00, 6/registerAx, 2/10, 3/010, 1/1, 6/-- -- -- -- -- --, 6/registerZ }


\paragraph{Syntax} \texttt{xmovetd} \emph{registerAx} = \emph{registerZ}

\paragraph{Description} The \emph{registerAx} operand receives the \emph{registerZ} value.


\paragraph{Execution} {Reservation table \textsf{ALU\_LITE\_MISC} (Section~\ref{sec:reservation-ALU-LITE-MISC})}
~
{\begin{lstlisting}
stage RR:
  new argument2 = GPR[registerZ];
  new result1 = _ZX_64(argument2);
stage E1:
  CS.XMF = 1;
  XCR[registerAx<<2] = result1;
\end{lstlisting}}
\newpage

\paragraph{Encoding} Format \textsf{ALU\_XMOVETD1} (Section~\ref{sec:format-ALU-XMOVETD1}), \verb|alucode3 = 010|\\

\bitfields{ 1/P, 2/11, 1/0, 2/10, 2/01, 6/registerAy, 2/10, 3/010, 1/1, 6/-- -- -- -- -- --, 6/registerZ }


\paragraph{Syntax} \texttt{xmovetd} \emph{registerAy} = \emph{registerZ}

\paragraph{Description} The \emph{registerAy} operand receives the \emph{registerZ} value.


\paragraph{Execution} {Reservation table \textsf{ALU\_LITE\_MISC} (Section~\ref{sec:reservation-ALU-LITE-MISC})}
~
{\begin{lstlisting}
stage RR:
  new argument2 = GPR[registerZ];
  new result1 = _ZX_64(argument2);
stage E1:
  CS.XMF = 1;
  XCR[(registerAy<<2)+1] = result1;
\end{lstlisting}}
\newpage

\paragraph{Encoding} Format \textsf{ALU\_XMOVETD2} (Section~\ref{sec:format-ALU-XMOVETD2}), \verb|alucode3 = 010|\\

\bitfields{ 1/P, 2/11, 1/0, 2/10, 2/10, 6/registerAz, 2/10, 3/010, 1/1, 6/-- -- -- -- -- --, 6/registerZ }


\paragraph{Syntax} \texttt{xmovetd} \emph{registerAz} = \emph{registerZ}

\paragraph{Description} The \emph{registerAz} operand receives the \emph{registerZ} value.


\paragraph{Execution} {Reservation table \textsf{ALU\_LITE\_MISC} (Section~\ref{sec:reservation-ALU-LITE-MISC})}
~
{\begin{lstlisting}
stage RR:
  new argument2 = GPR[registerZ];
  new result1 = _ZX_64(argument2);
stage E1:
  CS.XMF = 1;
  XCR[(registerAz<<2)+2] = result1;
\end{lstlisting}}
\newpage

\paragraph{Encoding} Format \textsf{ALU\_XMOVETD3} (Section~\ref{sec:format-ALU-XMOVETD3}), \verb|alucode3 = 010|\\

\bitfields{ 1/P, 2/11, 1/0, 2/10, 2/11, 6/registerAt, 2/10, 3/010, 1/1, 6/-- -- -- -- -- --, 6/registerZ }


\paragraph{Syntax} \texttt{xmovetd} \emph{registerAt} = \emph{registerZ}

\paragraph{Description} The \emph{registerAt} operand receives the \emph{registerZ} value.


\paragraph{Execution} {Reservation table \textsf{ALU\_LITE\_MISC} (Section~\ref{sec:reservation-ALU-LITE-MISC})}
~
{\begin{lstlisting}
stage RR:
  new argument2 = GPR[registerZ];
  new result1 = _ZX_64(argument2);
stage E1:
  CS.XMF = 1;
  XCR[(registerAt<<2)+3] = result1;
\end{lstlisting}}
\newpage
\subsubsection[{\small XMOVETO: Move To Extension Octuple Word}]{XMOVETO\hfill Move To Extension Octuple Word}
\label{instr:XMOVETO}\index{xmoveto}
 
\paragraph{Encoding} Format \textsf{ALU\_XMOVETO} (Section~\ref{sec:format-ALU-XMOVETO}), \verb|alucode3 = 010|\\

\bitfields{ 1/P, 2/11, 1/0, 4/1110, 6/registerA, 2/10, 3/010, 1/1, 5/registerO, 1/--, 5/registerP, 1/-- }


\paragraph{Syntax} \texttt{xmoveto} \emph{registerA} = \emph{registerP}, \emph{registerO}

\paragraph{Description} The \emph{registerA} operand receives the concatenation of the \emph{registerP} value and the \emph{registerO} value.


\paragraph{Execution} {Reservation table \textsf{ALU\_LITE\_MISC} (Section~\ref{sec:reservation-ALU-LITE-MISC})}
~
{\begin{lstlisting}
stage RR:
  new argument3 = PGR[registerO];
  new argument2 = PGR[registerP];
  new result1 = 0;
  result1.128[0] = argument2;
  result1.128[1] = argument3;
stage E1:
  CS.XMF = 1;
  XVR[registerA] = result1;
\end{lstlisting}}
\newpage
\subsubsection[{\small XMOVETQ: Move To Extension Quadruple Word}]{XMOVETQ\hfill Move To Extension Quadruple Word}
\label{instr:XMOVETQ}\index{xmovetq}
 
\paragraph{Encoding} Format \textsf{ALU\_XMOVETQ0} (Section~\ref{sec:format-ALU-XMOVETQ0}), \verb|alucode3 = 010|\\

\bitfields{ 1/P, 2/11, 1/0, 3/110, 1/0, 6/registerAE, 2/10, 3/010, 1/1, 6/registerY, 6/registerZ }


\paragraph{Syntax} \texttt{xmovetq} \emph{registerAE} = \emph{registerZ}, \emph{registerY}

\paragraph{Description} The \emph{registerAE} operand receives the concatenation of the \emph{registerZ} value and the \emph{registerY} value.


\paragraph{Execution} {Reservation table \textsf{ALU\_LITE\_MISC} (Section~\ref{sec:reservation-ALU-LITE-MISC})}
~
{\begin{lstlisting}
stage RR:
  new argument3 = GPR[registerY];
  new argument2 = GPR[registerZ];
  new result1 = 0;
  result1.64[0] = argument2;
  result1.64[1] = argument3;
stage E1:
  CS.XMF = 1;
  XBR[registerAE<<1] = result1;
\end{lstlisting}}
\newpage

\paragraph{Encoding} Format \textsf{ALU\_XMOVETQ1} (Section~\ref{sec:format-ALU-XMOVETQ1}), \verb|alucode3 = 010|\\

\bitfields{ 1/P, 2/11, 1/0, 3/110, 1/1, 6/registerAO, 2/10, 3/010, 1/1, 6/registerY, 6/registerZ }


\paragraph{Syntax} \texttt{xmovetq} \emph{registerAO} = \emph{registerZ}, \emph{registerY}

\paragraph{Description} The \emph{registerAO} operand receives the concatenation of the \emph{registerZ} value and the \emph{registerY} value.


\paragraph{Execution} {Reservation table \textsf{ALU\_LITE\_MISC} (Section~\ref{sec:reservation-ALU-LITE-MISC})}
~
{\begin{lstlisting}
stage RR:
  new argument3 = GPR[registerY];
  new argument2 = GPR[registerZ];
  new result1 = 0;
  result1.64[0] = argument2;
  result1.64[1] = argument3;
stage E1:
  CS.XMF = 1;
  XBR[(registerAO<<1)+1] = result1;
\end{lstlisting}}
\newpage
\subsection{FPU Instructions}
\label{sec:FPU-instructions}

\subsubsection[{\small FABSD: Floating-Point Absolute Value Double Word}]{FABSD\hfill Floating-Point Absolute Value Double Word}
\label{instr:FABSD}\index{fabsd}
 
\paragraph{Encoding} Format \textsf{ALU\_FDMO1} (Section~\ref{sec:format-ALU-FDMO1}), \verb|fpucode6 = 0011|\\

\bitfields{ 1/P, 2/11, 1/1, 1/1, 3/111, 6/registerW, 2/00, 3/000, 1/0, 4/0011, 1/--, 1/0, 6/registerZ }


\paragraph{Syntax} \texttt{fabsd} \emph{registerW} = \emph{registerZ}

\paragraph{Description} The absolute value the \emph{registerZ}, assuming binary 64 floating-point number, is stored into the \emph{registerW}. Implemented as a bitwise operation, does not raise FP exceptions.


\paragraph{Execution} {Reservation table \textsf{ALU\_LITE} (Section~\ref{sec:reservation-ALU-LITE})}
~
{\begin{lstlisting}
stage RR:
  new argument2 = GPR[registerZ];
  new result1 = argument2 & 0x7FFFFFFFFFFFFFFF;
stage E1:
  GPR[registerW] = result1;
\end{lstlisting}}
\newpage
\subsubsection[{\small FABSDP: Floating-Point Absolute Value Double Word Pair}]{FABSDP\hfill Floating-Point Absolute Value Double Word Pair}
\label{instr:FABSDP}\index{fabsdp}
 
\paragraph{Encoding} Format \textsf{ALU\_FDPMO1} (Section~\ref{sec:format-ALU-FDPMO1}), \verb|fpucode6 = 0011|\\

\bitfields{ 1/P, 2/11, 1/1, 1/1, 3/111, 5/registerM, 1/--, 2/11, 3/000, 1/0, 4/0011, 1/--, 1/0, 5/registerP, 1/1 }


\paragraph{Syntax} \texttt{fabsdp} \emph{registerM} = \emph{registerP}

\paragraph{Description} The \emph{registerP} is interpreted as two binary 64 floating-point numbers packed into 128 bits. The two absolute values the \emph{registerP} double words are packed and stored into the \emph{registerM}.


\paragraph{Execution} {Reservation table \textsf{ALU\_LITE} (Section~\ref{sec:reservation-ALU-LITE})}
~
{\begin{lstlisting}
stage RR:
  new argument2 = PGR[registerP];
  for (I = 0; I < 2; I += 1) {
    result1.64[I] = argument2.64[I] & 0x7FFFFFFFFFFFFFFF
  };
stage E1:
  PGR[registerM] = result1;
\end{lstlisting}}
\newpage
\subsubsection[{\small FABSH: Floating-Point Absolute Value Half Word}]{FABSH\hfill Floating-Point Absolute Value Half Word}
\label{instr:FABSH}\index{fabsh}
 
\paragraph{Encoding} Format \textsf{ALU\_FHMO1} (Section~\ref{sec:format-ALU-FHMO1}), \verb|fpucode6 = 0011|\\

\bitfields{ 1/P, 2/11, 1/1, 1/1, 3/111, 6/registerW, 2/01, 3/000, 1/1, 4/0011, 1/--, 1/0, 6/registerZ }


\paragraph{Syntax} \texttt{fabsh} \emph{registerW} = \emph{registerZ}

\paragraph{Description} The absolute value the \emph{registerZ}, assuming binary 16 floating-point number, is stored into the \emph{registerW}. Implemented as a bitwise operation, does not raise FP exceptions.


\paragraph{Execution} {Reservation table \textsf{ALU\_LITE} (Section~\ref{sec:reservation-ALU-LITE})}
~
{\begin{lstlisting}
stage RR:
  new argument2 = GPR[registerZ];
  new result1 = _ZX_16(argument2 & 0x7FFF);
stage E1:
  GPR[registerW] = result1;
\end{lstlisting}}
\newpage
\subsubsection[{\small FABSHO: Floating-Point Absolute Value Half Word Octuple}]{FABSHO\hfill Floating-Point Absolute Value Half Word Octuple}
\label{instr:FABSHO}\index{fabsho}
 
\paragraph{Encoding} Format \textsf{ALU\_FHOMO1} (Section~\ref{sec:format-ALU-FHOMO1}), \verb|fpucode6 = 0011|\\

\bitfields{ 1/P, 2/11, 1/1, 1/1, 3/111, 5/registerM, 1/--, 2/11, 3/000, 1/1, 4/0011, 1/--, 1/0, 5/registerP, 1/1 }


\paragraph{Syntax} \texttt{fabsho} \emph{registerM} = \emph{registerP}

\paragraph{Description} The \emph{registerP} is interpreted as eight binary 16 floating-point numbers packed into 128 bits. The eight absolute values the \emph{registerP} half words are packed and stored into the \emph{registerM}.


\paragraph{Execution} {Reservation table \textsf{ALU\_LITE} (Section~\ref{sec:reservation-ALU-LITE})}
~
{\begin{lstlisting}
stage RR:
  new argument2 = PGR[registerP];
  for (I = 0; I < 8; I += 1) {
    result1.16[I] = argument2.16[I] & 0x7FFF
  };
stage E1:
  PGR[registerM] = result1;
\end{lstlisting}}
\newpage
\subsubsection[{\small FABSW: Floating-Point Absolute Value Word}]{FABSW\hfill Floating-Point Absolute Value Word}
\label{instr:FABSW}\index{fabsw}
 
\paragraph{Encoding} Format \textsf{ALU\_FWMO1} (Section~\ref{sec:format-ALU-FWMO1}), \verb|fpucode6 = 0011|\\

\bitfields{ 1/P, 2/11, 1/1, 1/1, 3/111, 6/registerW, 2/01, 3/000, 1/0, 4/0011, 1/--, 1/0, 6/registerZ }


\paragraph{Syntax} \texttt{fabsw} \emph{registerW} = \emph{registerZ}

\paragraph{Description} The absolute value the \emph{registerZ}, assuming binary 32 floating-point number, is stored into the \emph{registerW}. Implemented as a bitwise operation, does not raise FP exceptions.


\paragraph{Execution} {Reservation table \textsf{ALU\_LITE} (Section~\ref{sec:reservation-ALU-LITE})}
~
{\begin{lstlisting}
stage RR:
  new argument2 = GPR[registerZ];
  new result1 = _ZX_32(argument2 & 0x7FFFFFFF);
stage E1:
  GPR[registerW] = result1;
\end{lstlisting}}
\newpage
\subsubsection[{\small FABSWP: Floating-Point Absolute Value Word Pair}]{FABSWP\hfill Floating-Point Absolute Value Word Pair}
\label{instr:FABSWP}\index{fabswp}
 
\paragraph{Encoding} Format \textsf{ALU\_FWPMO1} (Section~\ref{sec:format-ALU-FWPMO1}), \verb|fpucode6 = 0011|\\

\bitfields{ 1/P, 2/11, 1/1, 1/1, 3/111, 6/registerW, 2/00, 3/000, 1/0, 4/0011, 1/--, 1/1, 6/registerZ }


\paragraph{Syntax} \texttt{fabswp} \emph{registerW} = \emph{registerZ}

\paragraph{Description} The \emph{registerZ} is interpreted as two binary 32 floating-point numbers packed into 64 bits. The two absolute values the \emph{registerZ} words are packed and stored into the \emph{registerW}.


\paragraph{Execution} {Reservation table \textsf{ALU\_LITE} (Section~\ref{sec:reservation-ALU-LITE})}
~
{\begin{lstlisting}
stage RR:
  new argument2 = GPR[registerZ];
  for (I = 0; I < 2; I += 1) {
    result1.32[I] = argument2.32[I] & 0x7FFFFFFF
  };
stage E1:
  GPR[registerW] = result1;
\end{lstlisting}}
\newpage
\subsubsection[{\small FABSWQ: Floating-Point Absolute Value Word Quadruple}]{FABSWQ\hfill Floating-Point Absolute Value Word Quadruple}
\label{instr:FABSWQ}\index{fabswq}
 
\paragraph{Encoding} Format \textsf{ALU\_FWQMO1} (Section~\ref{sec:format-ALU-FWQMO1}), \verb|fpucode6 = 0011|\\

\bitfields{ 1/P, 2/11, 1/1, 1/1, 3/111, 5/registerM, 1/--, 2/11, 3/000, 1/1, 4/0011, 1/--, 1/0, 5/registerP, 1/0 }


\paragraph{Syntax} \texttt{fabswq} \emph{registerM} = \emph{registerP}

\paragraph{Description} The \emph{registerP} is interpreted as four binary 32 floating-point numbers packed into 128 bits. The four absolute values the \emph{registerP} words are packed and stored into the \emph{registerM}.


\paragraph{Execution} {Reservation table \textsf{ALU\_LITE} (Section~\ref{sec:reservation-ALU-LITE})}
~
{\begin{lstlisting}
stage RR:
  new argument2 = PGR[registerP];
  for (I = 0; I < 4; I += 1) {
    result1.32[I] = argument2.32[I] & 0x7FFFFFFF
  };
stage E1:
  PGR[registerM] = result1;
\end{lstlisting}}
\newpage
\subsubsection[{\small FADDD: Floating-Point Add Double Word}]{FADDD\hfill Floating-Point Add Double Word}
\label{instr:FADDD}\index{faddd}
 
\paragraph{Encoding} Format \textsf{ALU\_FDADD} (Section~\ref{sec:format-ALU-FDADD}), \verb|alucode3 = 000|\\

\bitfields{ 1/P, 2/11, 1/1, 1/0, 3/floatmode, 6/registerW, 2/00, 3/000, 1/0, 6/registerY, 6/registerZ }


\paragraph{Syntax} \texttt{faddd}\emph{floatmode} \emph{registerW} = \emph{registerZ}, \emph{registerY}

\paragraph{Description} The \emph{registerZ} is added to the \emph{registerY}, assuming binary 64 floating-point numbers. The result is rounded to the binary 64 floating-point format according to the rounding mode and stored into the \emph{registerW}. This instruction may raise exception bits in the CS register.


\paragraph{Modifier(s)}
\noindent\begin{itemize}
\item floatmode (cf. floatmode in section \ref{par:modifier-floatmode} on page \pageref{par:modifier-floatmode})
\end{itemize}

\paragraph{Execution} {Reservation table \textsf{ALU\_LITE} (Section~\ref{sec:reservation-ALU-LITE})}
~
{\begin{lstlisting}
stage ID:
  new floatmode = _ZX_3(floatmode);
stage RR:
  new argument3 = GPR[registerY];
  new argument2 = GPR[registerZ];
  new RM = decode_riscv_float_rounding_mode(floatmode);
  new result1 = f64_add(RM, argument2, argument3);
stage E4:
  GPR[registerW] = result1;
  commit_float_exception_flags();
\end{lstlisting}}
\newpage
\subsubsection[{\small FADDDP: Floating-Point Add Double Word Pair}]{FADDDP\hfill Floating-Point Add Double Word Pair}
\label{instr:FADDDP}\index{fadddp}
 
\paragraph{Encoding} Format \textsf{ALU\_FDPADD} (Section~\ref{sec:format-ALU-FDPADD}), \verb|alucode3 = 000|\\

\bitfields{ 1/P, 2/11, 1/1, 1/0, 3/floatmode, 5/registerM, 1/--, 2/11, 3/000, 1/0, 5/registerO, 1/--, 5/registerP, 1/1 }


\paragraph{Syntax} \texttt{fadddp}\emph{floatmode} \emph{registerM} = \emph{registerP}, \emph{registerO}

\paragraph{Description} The \emph{registerP} and the \emph{registerO} are interpreted as two binary 64 floating-point numbers packed into 128 bits. The \emph{registerP} double words are added to the \emph{registerO} double words. The two resulting double words are packed and stored into the \emph{registerM}.


\paragraph{Modifier(s)}
\noindent\begin{itemize}
\item floatmode (cf. floatmode in section \ref{par:modifier-floatmode} on page \pageref{par:modifier-floatmode})
\end{itemize}

\paragraph{Execution} {Reservation table \textsf{ALU\_LITE} (Section~\ref{sec:reservation-ALU-LITE})}
~
{\begin{lstlisting}
stage ID:
  new floatmode = _ZX_3(floatmode);
stage RR:
  new argument3 = PGR[registerO];
  new argument2 = PGR[registerP];
  new RM = decode_riscv_float_rounding_mode(floatmode);
  for (I = 0; I < 2; I += 1) {
    result1.64[I] = f64_add(RM, argument2.64[I], argument3.64[I])
  };
stage E4:
  PGR[registerM] = result1;
  commit_float_exception_flags();
\end{lstlisting}}
\newpage
\subsubsection[{\small FADDH: Floating-Point Add Half Word}]{FADDH\hfill Floating-Point Add Half Word}
\label{instr:FADDH}\index{faddh}
 
\paragraph{Encoding} Format \textsf{ALU\_FHADD} (Section~\ref{sec:format-ALU-FHADD}), \verb|alucode3 = 000|\\

\bitfields{ 1/P, 2/11, 1/1, 1/0, 3/floatmode, 6/registerW, 2/01, 3/000, 1/1, 6/registerY, 6/registerZ }


\paragraph{Syntax} \texttt{faddh}\emph{floatmode} \emph{registerW} = \emph{registerZ}, \emph{registerY}

\paragraph{Description} The \emph{registerZ} is added to the \emph{registerY}, assuming binary 16 floating-point numbers. The result is rounded to the binary 16 floating-point format according to the rounding mode and stored into the \emph{registerW}. This instruction may raise exception bits in the CS register.


\paragraph{Modifier(s)}
\noindent\begin{itemize}
\item floatmode (cf. floatmode in section \ref{par:modifier-floatmode} on page \pageref{par:modifier-floatmode})
\end{itemize}

\paragraph{Execution} {Reservation table \textsf{ALU\_LITE} (Section~\ref{sec:reservation-ALU-LITE})}
~
{\begin{lstlisting}
stage ID:
  new floatmode = _ZX_3(floatmode);
stage RR:
  new argument3 = GPR[registerY];
  new argument2 = GPR[registerZ];
  new RM = decode_riscv_float_rounding_mode(floatmode);
  new result1 = f16_add(RM, argument2, argument3);
stage E3:
  GPR[registerW] = result1;
  commit_float_exception_flags();
\end{lstlisting}}
\newpage
\subsubsection[{\small FADDHO: Floating-Point Add Half Word Octuple}]{FADDHO\hfill Floating-Point Add Half Word Octuple}
\label{instr:FADDHO}\index{faddho}
 
\paragraph{Encoding} Format \textsf{ALU\_FHOADD} (Section~\ref{sec:format-ALU-FHOADD}), \verb|alucode3 = 000|\\

\bitfields{ 1/P, 2/11, 1/1, 1/0, 3/floatmode, 5/registerM, 1/--, 2/11, 3/000, 1/1, 5/registerO, 1/--, 5/registerP, 1/1 }


\paragraph{Syntax} \texttt{faddho}\emph{floatmode} \emph{registerM} = \emph{registerP}, \emph{registerO}

\paragraph{Description} The \emph{registerP} and the \emph{registerO} are interpreted as eight binary 16 floating-point numbers packed into 128 bits. The \emph{registerP} half words are added to the \emph{registerO} half words. The eight resulting half words are packed and stored into the \emph{registerM}.


\paragraph{Modifier(s)}
\noindent\begin{itemize}
\item floatmode (cf. floatmode in section \ref{par:modifier-floatmode} on page \pageref{par:modifier-floatmode})
\end{itemize}

\paragraph{Execution} {Reservation table \textsf{ALU\_LITE} (Section~\ref{sec:reservation-ALU-LITE})}
~
{\begin{lstlisting}
stage ID:
  new floatmode = _ZX_3(floatmode);
stage RR:
  new argument3 = PGR[registerO];
  new argument2 = PGR[registerP];
  new RM = decode_riscv_float_rounding_mode(floatmode);
  for (I = 0; I < 8; I += 1) {
    result1.16[I] = f16_add(RM, argument2.16[I], argument3.16[I])
  };
stage E4:
  PGR[registerM] = result1;
  commit_float_exception_flags();
\end{lstlisting}}
\newpage
\subsubsection[{\small FADDW: Floating-Point Add Word}]{FADDW\hfill Floating-Point Add Word}
\label{instr:FADDW}\index{faddw}
 
\paragraph{Encoding} Format \textsf{ALU\_FWADD} (Section~\ref{sec:format-ALU-FWADD}), \verb|alucode3 = 000|\\

\bitfields{ 1/P, 2/11, 1/1, 1/0, 3/floatmode, 6/registerW, 2/01, 3/000, 1/0, 6/registerY, 6/registerZ }


\paragraph{Syntax} \texttt{faddw}\emph{floatmode} \emph{registerW} = \emph{registerZ}, \emph{registerY}

\paragraph{Description} The \emph{registerZ} is added to the \emph{registerY}, assuming binary 32 floating-point numbers. The result is rounded to the binary 32 floating-point format according to the rounding mode and stored into the \emph{registerW}. This instruction may raise exception bits in the CS register.


\paragraph{Modifier(s)}
\noindent\begin{itemize}
\item floatmode (cf. floatmode in section \ref{par:modifier-floatmode} on page \pageref{par:modifier-floatmode})
\end{itemize}

\paragraph{Execution} {Reservation table \textsf{ALU\_LITE} (Section~\ref{sec:reservation-ALU-LITE})}
~
{\begin{lstlisting}
stage ID:
  new floatmode = _ZX_3(floatmode);
stage RR:
  new argument3 = GPR[registerY];
  new argument2 = GPR[registerZ];
  new RM = decode_riscv_float_rounding_mode(floatmode);
  new result1 = f32_add(RM, argument2, argument3);
stage E4:
  GPR[registerW] = result1;
  commit_float_exception_flags();
\end{lstlisting}}
\newpage
\subsubsection[{\small FADDWC: Floating-Point Add to Word Complex}]{FADDWC\hfill Floating-Point Add to Word Complex}
\label{instr:FADDWC}\index{faddwc}
 
\paragraph{Encoding} Format \textsf{ALU\_FWCOP} (Section~\ref{sec:format-ALU-FWCOP}), \verb|alucode6 = 00|\\

\bitfields{ 1/P, 2/11, 1/1, 1/imultiply, 3/floatmode, 6/registerW, 2/00, 2/00, 1/conjugate, 1/1, 6/registerY, 6/registerZ }


\paragraph{Syntax} \texttt{faddwc}\emph{conjugate}\emph{imultiply}\emph{floatmode} \emph{registerW} = \emph{registerZ}, \emph{registerY}

\paragraph{Description} The \emph{registerY} and the \emph{registerZ} are interpreted as binary 32 floating-point complex numbers. If \emph{conjugate} is set, the complex conjugate of the \emph{registerZ} is used in the operation. The real and imaginary parts are respectively added to and rounded according to the rounding mode. If \emph{imultiply} is set, the result is multiplied by the imaginary unit. The resulting binary 32 floating-point complex number is stored back into the \emph{registerW}. This instruction may raise exception bits in the CS register.


\paragraph{Modifier(s)}
\noindent\begin{itemize}
\item floatmode (cf. floatmode in section \ref{par:modifier-floatmode} on page \pageref{par:modifier-floatmode})
\item imultiply (cf. imultiply in section \ref{par:modifier-imultiply} on page \pageref{par:modifier-imultiply})
\item conjugate (cf. conjugate in section \ref{par:modifier-conjugate} on page \pageref{par:modifier-conjugate})
\end{itemize}

\paragraph{Execution} {Reservation table \textsf{ALU\_LITE} (Section~\ref{sec:reservation-ALU-LITE})}
~
{\begin{lstlisting}
stage ID:
  new conjugate = _ZX_1(conjugate);
  new imultiply = _ZX_1(imultiply);
  new floatmode = _ZX_3(floatmode);
stage RR:
  new argument3 = GPR[registerY];
  new argument2 = GPR[registerZ];
  new argument1 = GPR[registerW];
  new RM = decode_riscv_float_rounding_mode(floatmode);
  new fconj = conjugate ? 0x80000000 : 0;
  result1.32[0] = f32_add(RM, argument2.32[0], argument3.32[0]);
  result1.32[1] = f32_add(RM, argument2.32[1] ^ fconj, argument3.32[1]);
  if (imultiply) {
    new result1 = _SWAP_32(result1);
    result1.32[0] = result1.32[0] ^ 0x80000000;
  }
stage E4:
  GPR[registerW] = result1;
  commit_float_exception_flags();
\end{lstlisting}}
\newpage
\subsubsection[{\small FADDWQ: Floating-Point Add Word Quadruple}]{FADDWQ\hfill Floating-Point Add Word Quadruple}
\label{instr:FADDWQ}\index{faddwq}
 
\paragraph{Encoding} Format \textsf{ALU\_FWQADD} (Section~\ref{sec:format-ALU-FWQADD}), \verb|alucode3 = 000|\\

\bitfields{ 1/P, 2/11, 1/1, 1/0, 3/floatmode, 5/registerM, 1/--, 2/11, 3/000, 1/1, 5/registerO, 1/--, 5/registerP, 1/0 }


\paragraph{Syntax} \texttt{faddwq}\emph{floatmode} \emph{registerM} = \emph{registerP}, \emph{registerO}

\paragraph{Description} The \emph{registerP} and the \emph{registerO} are interpreted as four binary 32 floating-point numbers packed into 128 bits. The \emph{registerP} words are added to the \emph{registerO} words. The four resulting words are packed and stored into the \emph{registerM}.


\paragraph{Modifier(s)}
\noindent\begin{itemize}
\item floatmode (cf. floatmode in section \ref{par:modifier-floatmode} on page \pageref{par:modifier-floatmode})
\end{itemize}

\paragraph{Execution} {Reservation table \textsf{ALU\_LITE} (Section~\ref{sec:reservation-ALU-LITE})}
~
{\begin{lstlisting}
stage ID:
  new floatmode = _ZX_3(floatmode);
stage RR:
  new argument3 = PGR[registerO];
  new argument2 = PGR[registerP];
  new RM = decode_riscv_float_rounding_mode(floatmode);
  for (I = 0; I < 4; I += 1) {
    result1.32[I] = f32_add(RM, argument2.32[I], argument3.32[I])
  };
stage E4:
  PGR[registerM] = result1;
  commit_float_exception_flags();
\end{lstlisting}}
\newpage
\subsubsection[{\small FCLASSD: Floating-Point Classify Double Word}]{FCLASSD\hfill Floating-Point Classify Double Word}
\label{instr:FCLASSD}\index{fclassd}
 
\paragraph{Encoding} Format \textsf{ALU\_FDMO1} (Section~\ref{sec:format-ALU-FDMO1}), \verb|fpucode6 = 0110|\\

\bitfields{ 1/P, 2/11, 1/1, 1/1, 3/111, 6/registerW, 2/00, 3/000, 1/0, 4/0110, 1/--, 1/0, 6/registerZ }


\paragraph{Syntax} \texttt{fclassd} \emph{registerW} = \emph{registerZ}

\paragraph{Description} Classify the \emph{registerZ}, assuming a binary 64 floating-point number. is stored into the \emph{registerW}.


\paragraph{Execution} {Reservation table \textsf{ALU\_LITE} (Section~\ref{sec:reservation-ALU-LITE})}
~
{\begin{lstlisting}
stage RR:
  new argument2 = GPR[registerZ];
  new result1 = f64_classify(argument2);
stage E1:
  GPR[registerW] = result1;
\end{lstlisting}}
\newpage
\subsubsection[{\small FCLASSH: Floating-Point Classify Half Word}]{FCLASSH\hfill Floating-Point Classify Half Word}
\label{instr:FCLASSH}\index{fclassh}
 
\paragraph{Encoding} Format \textsf{ALU\_FHMO1} (Section~\ref{sec:format-ALU-FHMO1}), \verb|fpucode6 = 0110|\\

\bitfields{ 1/P, 2/11, 1/1, 1/1, 3/111, 6/registerW, 2/01, 3/000, 1/1, 4/0110, 1/--, 1/0, 6/registerZ }


\paragraph{Syntax} \texttt{fclassh} \emph{registerW} = \emph{registerZ}

\paragraph{Description} Classify the \emph{registerZ}, assuming a binary 16 floating-point number. is stored into the \emph{registerW}.


\paragraph{Execution} {Reservation table \textsf{ALU\_LITE} (Section~\ref{sec:reservation-ALU-LITE})}
~
{\begin{lstlisting}
stage RR:
  new argument2 = GPR[registerZ];
  new result1 = f16_classify(argument2);
stage E1:
  GPR[registerW] = result1;
\end{lstlisting}}
\newpage
\subsubsection[{\small FCLASSW: Floating-Point Classify Word}]{FCLASSW\hfill Floating-Point Classify Word}
\label{instr:FCLASSW}\index{fclassw}
 
\paragraph{Encoding} Format \textsf{ALU\_FWMO1} (Section~\ref{sec:format-ALU-FWMO1}), \verb|fpucode6 = 0110|\\

\bitfields{ 1/P, 2/11, 1/1, 1/1, 3/111, 6/registerW, 2/01, 3/000, 1/0, 4/0110, 1/--, 1/0, 6/registerZ }


\paragraph{Syntax} \texttt{fclassw} \emph{registerW} = \emph{registerZ}

\paragraph{Description} Classify the \emph{registerZ}, assuming a binary 32 floating-point number. is stored into the \emph{registerW}.


\paragraph{Execution} {Reservation table \textsf{ALU\_LITE} (Section~\ref{sec:reservation-ALU-LITE})}
~
{\begin{lstlisting}
stage RR:
  new argument2 = GPR[registerZ];
  new result1 = f32_classify(argument2);
stage E1:
  GPR[registerW] = result1;
\end{lstlisting}}
\newpage
\subsubsection[{\small FCLASSWP: Floating-Point Classify Word Pair}]{FCLASSWP\hfill Floating-Point Classify Word Pair}
\label{instr:FCLASSWP}\index{fclasswp}
 
\paragraph{Encoding} Format \textsf{ALU\_FWPMO1} (Section~\ref{sec:format-ALU-FWPMO1}), \verb|fpucode6 = 0110|\\

\bitfields{ 1/P, 2/11, 1/1, 1/1, 3/111, 6/registerW, 2/00, 3/000, 1/0, 4/0110, 1/--, 1/1, 6/registerZ }


\paragraph{Syntax} \texttt{fclasswp} \emph{registerW} = \emph{registerZ}

\paragraph{Description} The \emph{registerZ} is interpreted as two binary 32 floating-point numbers packed into 64 bits. The two binary 32 floating-point number are classified. The resulting two words are packed and stored into the \emph{registerW}.


\paragraph{Execution} {Reservation table \textsf{ALU\_LITE} (Section~\ref{sec:reservation-ALU-LITE})}
~
{\begin{lstlisting}
stage RR:
  new argument2 = GPR[registerZ];
  for (I = 0; I < 2; I += 1) {
    result1.32[I] = f32_classify(argument2.32[I])
  };
stage E1:
  GPR[registerW] = result1;
\end{lstlisting}}
\newpage
\subsubsection[{\small FCLASSWQ: Floating-Point Classify Word Quadruple}]{FCLASSWQ\hfill Floating-Point Classify Word Quadruple}
\label{instr:FCLASSWQ}\index{fclasswq}
 
\paragraph{Encoding} Format \textsf{ALU\_FWQMO1} (Section~\ref{sec:format-ALU-FWQMO1}), \verb|fpucode6 = 0110|\\

\bitfields{ 1/P, 2/11, 1/1, 1/1, 3/111, 5/registerM, 1/--, 2/11, 3/000, 1/1, 4/0110, 1/--, 1/0, 5/registerP, 1/0 }


\paragraph{Syntax} \texttt{fclasswq} \emph{registerM} = \emph{registerP}

\paragraph{Description} The \emph{registerP} is interpreted as four binary 32 floating-point numbers packed into 128 bits. The four binary 32 floating-point number are classified. The resulting four words are packed and stored into the \emph{registerM}.


\paragraph{Execution} {Reservation table \textsf{ALU\_LITE} (Section~\ref{sec:reservation-ALU-LITE})}
~
{\begin{lstlisting}
stage RR:
  new argument2 = PGR[registerP];
  for (I = 0; I < 4; I += 1) {
    result1.32[I] = f32_classify(argument2.32[I])
  };
stage E1:
  PGR[registerM] = result1;
\end{lstlisting}}
\newpage
\subsubsection[{\small FCOMPD: Floating-Point Compare Double Word}]{FCOMPD\hfill Floating-Point Compare Double Word}
\label{instr:FCOMPD}\index{fcompd}
 
\paragraph{Encoding} Format \textsf{ALU\_FDCWRR} (Section~\ref{sec:format-ALU-FDCWRR}), \verb|alucode3 = 010|\\

\bitfields{ 1/P, 2/11, 1/1, 1/0, 3/floatcomp, 6/registerW, 2/00, 3/010, 1/0, 6/registerY, 6/registerZ }


\paragraph{Syntax} \texttt{fcompd}\emph{floatcomp} \emph{registerW} = \emph{registerZ}, \emph{registerY}

\paragraph{Description} The \emph{registerZ} is compared to the \emph{registerY} using condition \emph{floatcomp}, assuming binary 64 floating-point numbers. The boolean result is stored into the \emph{registerW}. This instruction does not raise an invalid exception when handling a signaling NaN.


\paragraph{Modifier(s)}
\noindent\begin{itemize}
\item floatcomp (cf. floatcomp in section \ref{par:modifier-floatcomp} on page \pageref{par:modifier-floatcomp})
\end{itemize}

\paragraph{Execution} {Reservation table \textsf{ALU\_LITE} (Section~\ref{sec:reservation-ALU-LITE})}
~
{\begin{lstlisting}
stage ID:
  new argument4 = _ZX_3(floatcomp);
stage RR:
  new argument3 = GPR[registerY];
  new argument2 = GPR[registerZ];
  new result1 = floatcomp_64(argument4, argument2, argument3);
stage E1:
  GPR[registerW] = result1;
\end{lstlisting}}
\newpage

\paragraph{Encoding} Format \textsf{ALU\_FDCWRR.W} (Section~\ref{sec:format-ALU-FDCWRR.W}), \verb|alucode3 = 010|\\

\bitfields{ 1/P, 2/00, 2/-- --, 27/upper27, 1/1, 2/11, 1/1, 1/0, 3/floatcomp, 6/registerW, 2/00, 3/010, 1/0, 1/0, 5/lower5, 6/registerZ }


\paragraph{Syntax} \texttt{fcompd}\emph{floatcomp} \emph{registerW} = \emph{registerZ}, \emph{upper27\_\discretionary{}{}{}lower5}

\paragraph{Description} The \emph{registerZ} is compared to the \emph{upper27\_\discretionary{}{}{}lower5} using condition \emph{floatcomp}, assuming binary 64 floating-point numbers. The boolean result is stored into the \emph{registerW}. This instruction does not raise an invalid exception when handling a signaling NaN.


\paragraph{Modifier(s)}
\noindent\begin{itemize}
\item floatcomp (cf. floatcomp in section \ref{par:modifier-floatcomp} on page \pageref{par:modifier-floatcomp})
\end{itemize}

\paragraph{Execution} {Reservation table \textsf{ALU\_LITE.X} (Section~\ref{sec:reservation-ALU-LITE.X})}
~
{\begin{lstlisting}
stage ID:
  new argument4 = _ZX_3(floatcomp);
stage RR:
  new argument3 = _WX_32(upper27_lower5);
  new argument2 = GPR[registerZ];
  new result1 = floatcomp_64(argument4, argument2, argument3);
stage E1:
  GPR[registerW] = result1;
\end{lstlisting}}
\newpage
\subsubsection[{\small FCOMPDP: Floating-Point Compare Double Word Pair}]{FCOMPDP\hfill Floating-Point Compare Double Word Pair}
\label{instr:FCOMPDP}\index{fcompdp}
 
\paragraph{Encoding} Format \textsf{ALU\_FDPCWRR} (Section~\ref{sec:format-ALU-FDPCWRR}), \verb|alucode3 = 010|\\

\bitfields{ 1/P, 2/11, 1/1, 1/0, 3/floatcomp, 6/registerW, 2/11, 3/010, 1/0, 5/registerO, 1/--, 5/registerP, 1/1 }


\paragraph{Syntax} \texttt{fcompdp}\emph{floatcomp} \emph{registerW} = \emph{registerP}, \emph{registerO}

\paragraph{Description} The \emph{registerP} and the \emph{registerO} are interpreted as two binary 64 floating-point numbers packed into 128 bits. The \emph{registerP} double words are compared to the \emph{registerO} double words using condition \emph{floatcomp}. The two boolean result bits are packed and stored into the \emph{registerW}.


\paragraph{Modifier(s)}
\noindent\begin{itemize}
\item floatcomp (cf. floatcomp in section \ref{par:modifier-floatcomp} on page \pageref{par:modifier-floatcomp})
\end{itemize}

\paragraph{Execution} {Reservation table \textsf{ALU\_LITE} (Section~\ref{sec:reservation-ALU-LITE})}
~
{\begin{lstlisting}
stage ID:
  new argument4 = _ZX_3(floatcomp);
stage RR:
  new argument3 = PGR[registerO];
  new argument2 = PGR[registerP];
  for (I = 0; I < 2; I += 1) {
    result1.1[I] = floatcomp_64(argument4, argument2.64[I], argument3.64[I])
  };
stage E1:
  GPR[registerW] = result1;
\end{lstlisting}}
\newpage

\paragraph{Encoding} Format \textsf{ALU\_FDPCWRR.M} (Section~\ref{sec:format-ALU-FDPCWRR.M}), \verb|alucode3 = 010|\\

\bitfields{ 1/P, 2/00, 2/-- --, 27/upper27, 1/1, 2/11, 1/1, 1/0, 3/floatcomp, 6/registerW, 2/11, 3/010, 1/0, 1/splat32, 5/lower5, 5/registerP, 1/1 }


\paragraph{Syntax} \texttt{fcompdp}\emph{floatcomp} \emph{registerW} = \emph{registerP}, \emph{upper27\_\discretionary{}{}{}lower5}\emph{splat32}

\paragraph{Description} The \emph{registerP} and the \emph{upper27\_\discretionary{}{}{}lower5} are interpreted as two binary 64 floating-point numbers packed into 128 bits. The \emph{registerP} double words are compared to the \emph{upper27\_\discretionary{}{}{}lower5} double words using condition \emph{floatcomp}. The two boolean result bits are packed and stored into the \emph{registerW}.


\paragraph{Modifier(s)}
\noindent\begin{itemize}
\item floatcomp (cf. floatcomp in section \ref{par:modifier-floatcomp} on page \pageref{par:modifier-floatcomp})
\item splat32 (cf. splat32 in section \ref{par:modifier-splat32} on page \pageref{par:modifier-splat32})
\end{itemize}

\paragraph{Execution} {Reservation table \textsf{ALU\_LITE.X} (Section~\ref{sec:reservation-ALU-LITE.X})}
~
{\begin{lstlisting}
stage ID:
  new splat32 = _ZX_1(splat32);
  new argument4 = _ZX_3(floatcomp);
stage RR:
  new argument3 = splat32 ? (_WX_32(upper27_lower5) << 32) | _ZX_32(_WX_32(upper27_lower5)) : _SX_32(_WX_32(upper27_lower5));
  new argument2 = PGR[registerP];
  for (I = 0; I < 2; I += 1) {
    result1.1[I] = floatcomp_64(argument4, argument2.64[I], argument3.64[I])
  };
stage E1:
  GPR[registerW] = result1;
\end{lstlisting}}
\newpage
\subsubsection[{\small FCOMPH: Floating-Point Compare Half Word}]{FCOMPH\hfill Floating-Point Compare Half Word}
\label{instr:FCOMPH}\index{fcomph}
 
\paragraph{Encoding} Format \textsf{ALU\_FHCWRR} (Section~\ref{sec:format-ALU-FHCWRR}), \verb|alucode3 = 010|\\

\bitfields{ 1/P, 2/11, 1/1, 1/0, 3/floatcomp, 6/registerW, 2/01, 3/010, 1/1, 6/registerY, 6/registerZ }


\paragraph{Syntax} \texttt{fcomph}\emph{floatcomp} \emph{registerW} = \emph{registerZ}, \emph{registerY}

\paragraph{Description} The \emph{registerZ} is compared to the \emph{registerY} using condition \emph{floatcomp}, assuming binary 16 floating-point numbers. The boolean result is stored into the \emph{registerW}. This instruction does not raise an invalid exception when handling a signaling NaN.


\paragraph{Modifier(s)}
\noindent\begin{itemize}
\item floatcomp (cf. floatcomp in section \ref{par:modifier-floatcomp} on page \pageref{par:modifier-floatcomp})
\end{itemize}

\paragraph{Execution} {Reservation table \textsf{ALU\_LITE} (Section~\ref{sec:reservation-ALU-LITE})}
~
{\begin{lstlisting}
stage ID:
  new argument4 = _ZX_3(floatcomp);
stage RR:
  new argument3 = GPR[registerY];
  new argument2 = GPR[registerZ];
  new result1 = floatcomp_16(argument4, argument2, argument3);
stage E1:
  GPR[registerW] = _ZX_32(result1);
\end{lstlisting}}
\newpage

\paragraph{Encoding} Format \textsf{ALU\_FHCWRR.W} (Section~\ref{sec:format-ALU-FHCWRR.W}), \verb|alucode3 = 010|\\

\bitfields{ 1/P, 2/00, 2/-- --, 27/upper27, 1/1, 2/11, 1/1, 1/0, 3/floatcomp, 6/registerW, 2/01, 3/010, 1/1, 1/0, 5/lower5, 6/registerZ }


\paragraph{Syntax} \texttt{fcomph}\emph{floatcomp} \emph{registerW} = \emph{registerZ}, \emph{upper27\_\discretionary{}{}{}lower5}

\paragraph{Description} The \emph{registerZ} is compared to the \emph{upper27\_\discretionary{}{}{}lower5} using condition \emph{floatcomp}, assuming binary 16 floating-point numbers. The boolean result is stored into the \emph{registerW}. This instruction does not raise an invalid exception when handling a signaling NaN.


\paragraph{Modifier(s)}
\noindent\begin{itemize}
\item floatcomp (cf. floatcomp in section \ref{par:modifier-floatcomp} on page \pageref{par:modifier-floatcomp})
\end{itemize}

\paragraph{Execution} {Reservation table \textsf{ALU\_LITE.X} (Section~\ref{sec:reservation-ALU-LITE.X})}
~
{\begin{lstlisting}
stage ID:
  new argument4 = _ZX_3(floatcomp);
stage RR:
  new argument3 = _WX_32(upper27_lower5);
  new argument2 = GPR[registerZ];
  new result1 = floatcomp_16(argument4, argument2, argument3);
stage E1:
  GPR[registerW] = _ZX_32(result1);
\end{lstlisting}}
\newpage
\subsubsection[{\small FCOMPNDP: Floating-Point Compare Negate Double Word Pair}]{FCOMPNDP\hfill Floating-Point Compare Negate Double Word Pair}
\label{instr:FCOMPNDP}\index{fcompndp}
 
\paragraph{Encoding} Format \textsf{ALU\_FDPCNWRR} (Section~\ref{sec:format-ALU-FDPCNWRR}), \verb|alucode3 = 010|\\

\bitfields{ 1/P, 2/11, 1/1, 1/1, 3/floatcomp, 5/registerM, 1/--, 2/11, 3/010, 1/0, 5/registerO, 1/0, 5/registerP, 1/1 }


\paragraph{Syntax} \texttt{fcompndp}\emph{floatcomp} \emph{registerM} = \emph{registerP}, \emph{registerO}

\paragraph{Description} The \emph{registerP} and the \emph{registerO} are interpreted as two binary 64 floating-point numbers packed into 128 bits. The \emph{registerP} double words are compared to the \emph{registerO} double words using condition \emph{floatcomp}. The two negated boolean results extended to double words are packed and stored into the \emph{registerM}. This instruction does not raise an invalid exception when handling a signaling NaN.


\paragraph{Modifier(s)}
\noindent\begin{itemize}
\item floatcomp (cf. floatcomp in section \ref{par:modifier-floatcomp} on page \pageref{par:modifier-floatcomp})
\end{itemize}

\paragraph{Execution} {Reservation table \textsf{ALU\_LITE} (Section~\ref{sec:reservation-ALU-LITE})}
~
{\begin{lstlisting}
stage ID:
  new argument4 = _ZX_3(floatcomp);
stage RR:
  new argument3 = PGR[registerO];
  new argument2 = PGR[registerP];
  for (I = 0; I < 2; I += 1) {
    result1.64[I] = -floatcomp_64(argument4, argument2.64[I], argument3.64[I])
  };
stage E1:
  PGR[registerM] = result1;
\end{lstlisting}}
\newpage

\paragraph{Encoding} Format \textsf{ALU\_FDPCNWRR.M} (Section~\ref{sec:format-ALU-FDPCNWRR.M}), \verb|alucode3 = 011|\\

\bitfields{ 1/P, 2/00, 2/-- --, 27/upper27, 1/1, 2/11, 1/1, 1/1, 3/floatcomp, 5/registerM, 1/--, 2/11, 3/011, 1/0, 1/splat32, 5/lower5, 5/registerP, 1/1 }


\paragraph{Syntax} \texttt{fcompndp}\emph{floatcomp} \emph{registerM} = \emph{registerP}, \emph{upper27\_\discretionary{}{}{}lower5}\emph{splat32}

\paragraph{Description} The \emph{registerP} and the \emph{upper27\_\discretionary{}{}{}lower5} are interpreted as two binary 64 floating-point numbers packed into 128 bits. The \emph{registerP} double words are compared to the \emph{upper27\_\discretionary{}{}{}lower5} double words using condition \emph{floatcomp}. The two negated boolean results extended to double words are packed and stored into the \emph{registerM}. This instruction does not raise an invalid exception when handling a signaling NaN.


\paragraph{Modifier(s)}
\noindent\begin{itemize}
\item floatcomp (cf. floatcomp in section \ref{par:modifier-floatcomp} on page \pageref{par:modifier-floatcomp})
\item splat32 (cf. splat32 in section \ref{par:modifier-splat32} on page \pageref{par:modifier-splat32})
\end{itemize}

\paragraph{Execution} {Reservation table \textsf{ALU\_LITE.X} (Section~\ref{sec:reservation-ALU-LITE.X})}
~
{\begin{lstlisting}
stage ID:
  new splat32 = _ZX_1(splat32);
  new argument4 = _ZX_3(floatcomp);
stage RR:
  new argument3 = splat32 ? (_WX_32(upper27_lower5) << 32) | _ZX_32(_WX_32(upper27_lower5)) : _SX_32(_WX_32(upper27_lower5));
  new argument2 = PGR[registerP];
  for (I = 0; I < 2; I += 1) {
    result1.64[I] = -floatcomp_64(argument4, argument2.64[I], argument3.64[I])
  };
stage E1:
  PGR[registerM] = result1;
\end{lstlisting}}
\newpage
\subsubsection[{\small FCOMPNHO: Floating-Point Compare Negate Half Word Octuple}]{FCOMPNHO\hfill Floating-Point Compare Negate Half Word Octuple}
\label{instr:FCOMPNHO}\index{fcompnho}
 
\paragraph{Encoding} Format \textsf{ALU\_FHOCNWRR} (Section~\ref{sec:format-ALU-FHOCNWRR}), \verb|alucode3 = 010|\\

\bitfields{ 1/P, 2/11, 1/1, 1/1, 3/floatcomp, 5/registerM, 1/--, 2/11, 3/010, 1/1, 5/registerO, 1/0, 5/registerP, 1/1 }


\paragraph{Syntax} \texttt{fcompnho}\emph{floatcomp} \emph{registerM} = \emph{registerP}, \emph{registerO}

\paragraph{Description} The \emph{registerP} and the \emph{registerO} are interpreted as eight binary 16 floating-point numbers packed into 128 bits. The \emph{registerP} half words are compared to the \emph{registerO} half words using condition \emph{floatcomp}. The eight negated boolean results extended to half words are packed and stored into the \emph{registerM}. This instruction does not raise an invalid exception when handling a signaling NaN.


\paragraph{Modifier(s)}
\noindent\begin{itemize}
\item floatcomp (cf. floatcomp in section \ref{par:modifier-floatcomp} on page \pageref{par:modifier-floatcomp})
\end{itemize}

\paragraph{Execution} {Reservation table \textsf{ALU\_LITE} (Section~\ref{sec:reservation-ALU-LITE})}
~
{\begin{lstlisting}
stage ID:
  new argument4 = _ZX_3(floatcomp);
stage RR:
  new argument3 = PGR[registerO];
  new argument2 = PGR[registerP];
  for (I = 0; I < 8; I += 1) {
    result1.16[I] = -floatcomp_16(argument4, argument2.16[I], argument3.16[I])
  };
stage E1:
  PGR[registerM] = result1;
\end{lstlisting}}
\newpage

\paragraph{Encoding} Format \textsf{ALU\_FHOCNWRR.M} (Section~\ref{sec:format-ALU-FHOCNWRR.M}), \verb|alucode3 = 010|\\

\bitfields{ 1/P, 2/00, 2/-- --, 27/upper27, 1/1, 2/11, 1/1, 1/1, 3/floatcomp, 5/registerM, 1/--, 2/11, 3/010, 1/1, 1/splat32, 5/lower5, 5/registerP, 1/1 }


\paragraph{Syntax} \texttt{fcompnho}\emph{floatcomp} \emph{registerM} = \emph{registerP}, \emph{upper27\_\discretionary{}{}{}lower5}\emph{splat32}

\paragraph{Description} The \emph{registerP} and the \emph{upper27\_\discretionary{}{}{}lower5} are interpreted as eight binary 16 floating-point numbers packed into 128 bits. The \emph{registerP} half words are compared to the \emph{upper27\_\discretionary{}{}{}lower5} half words using condition \emph{floatcomp}. The eight negated boolean results extended to half words are packed and stored into the \emph{registerM}. This instruction does not raise an invalid exception when handling a signaling NaN.


\paragraph{Modifier(s)}
\noindent\begin{itemize}
\item floatcomp (cf. floatcomp in section \ref{par:modifier-floatcomp} on page \pageref{par:modifier-floatcomp})
\item splat32 (cf. splat32 in section \ref{par:modifier-splat32} on page \pageref{par:modifier-splat32})
\end{itemize}

\paragraph{Execution} {Reservation table \textsf{ALU\_LITE.X} (Section~\ref{sec:reservation-ALU-LITE.X})}
~
{\begin{lstlisting}
stage ID:
  new splat32 = _ZX_1(splat32);
  new argument4 = _ZX_3(floatcomp);
stage RR:
  new argument3 = splat32 ? (_WX_32(upper27_lower5) << 32) | _ZX_32(_WX_32(upper27_lower5)) : _SX_32(_WX_32(upper27_lower5));
  new argument2 = PGR[registerP];
  for (I = 0; I < 8; I += 1) {
    result1.16[I] = -floatcomp_16(argument4, argument2.16[I], argument3.16[I])
  };
stage E1:
  PGR[registerM] = result1;
\end{lstlisting}}
\newpage
\subsubsection[{\small FCOMPNWQ: Floating-Point Compare Negate Word Quadruple}]{FCOMPNWQ\hfill Floating-Point Compare Negate Word Quadruple}
\label{instr:FCOMPNWQ}\index{fcompnwq}
 
\paragraph{Encoding} Format \textsf{ALU\_FWQCNWRR} (Section~\ref{sec:format-ALU-FWQCNWRR}), \verb|alucode3 = 010|\\

\bitfields{ 1/P, 2/11, 1/1, 1/1, 3/floatcomp, 5/registerM, 1/--, 2/11, 3/010, 1/1, 5/registerO, 1/0, 5/registerP, 1/0 }


\paragraph{Syntax} \texttt{fcompnwq}\emph{floatcomp} \emph{registerM} = \emph{registerP}, \emph{registerO}

\paragraph{Description} The \emph{registerP} and the \emph{registerO} are interpreted as four binary 32 floating-point numbers packed into 128 bits. The \emph{registerP} words are compared to the \emph{registerO} words using condition \emph{floatcomp}. The four negated boolean results extended to words are packed and stored into the \emph{registerM}. This instruction does not raise an invalid exception when handling a signaling NaN.


\paragraph{Modifier(s)}
\noindent\begin{itemize}
\item floatcomp (cf. floatcomp in section \ref{par:modifier-floatcomp} on page \pageref{par:modifier-floatcomp})
\end{itemize}

\paragraph{Execution} {Reservation table \textsf{ALU\_LITE} (Section~\ref{sec:reservation-ALU-LITE})}
~
{\begin{lstlisting}
stage ID:
  new argument4 = _ZX_3(floatcomp);
stage RR:
  new argument3 = PGR[registerO];
  new argument2 = PGR[registerP];
  for (I = 0; I < 4; I += 1) {
    result1.32[I] = -floatcomp_32(argument4, argument2.32[I], argument3.32[I])
  };
stage E1:
  PGR[registerM] = result1;
\end{lstlisting}}
\newpage

\paragraph{Encoding} Format \textsf{ALU\_FWQCNWRR.M} (Section~\ref{sec:format-ALU-FWQCNWRR.M}), \verb|alucode3 = 001|\\

\bitfields{ 1/P, 2/00, 2/-- --, 27/upper27, 1/1, 2/11, 1/1, 1/1, 3/floatcomp, 5/registerM, 1/--, 2/11, 3/001, 1/1, 1/splat32, 5/lower5, 5/registerP, 1/0 }


\paragraph{Syntax} \texttt{fcompnwq}\emph{floatcomp} \emph{registerM} = \emph{registerP}, \emph{upper27\_\discretionary{}{}{}lower5}\emph{splat32}

\paragraph{Description} The \emph{registerP} and the \emph{upper27\_\discretionary{}{}{}lower5} are interpreted as four binary 32 floating-point numbers packed into 128 bits. The \emph{registerP} words are compared to the \emph{upper27\_\discretionary{}{}{}lower5} words using condition \emph{floatcomp}. The four negated boolean results extended to words are packed and stored into the \emph{registerM}. This instruction does not raise an invalid exception when handling a signaling NaN.


\paragraph{Modifier(s)}
\noindent\begin{itemize}
\item floatcomp (cf. floatcomp in section \ref{par:modifier-floatcomp} on page \pageref{par:modifier-floatcomp})
\item splat32 (cf. splat32 in section \ref{par:modifier-splat32} on page \pageref{par:modifier-splat32})
\end{itemize}

\paragraph{Execution} {Reservation table \textsf{ALU\_LITE.X} (Section~\ref{sec:reservation-ALU-LITE.X})}
~
{\begin{lstlisting}
stage ID:
  new splat32 = _ZX_1(splat32);
  new argument4 = _ZX_3(floatcomp);
stage RR:
  new argument3 = splat32 ? (_WX_32(upper27_lower5) << 32) | _ZX_32(_WX_32(upper27_lower5)) : _SX_32(_WX_32(upper27_lower5));
  new argument2 = PGR[registerP];
  for (I = 0; I < 4; I += 1) {
    result1.32[I] = -floatcomp_32(argument4, argument2.32[I], argument3.32[I])
  };
stage E1:
  PGR[registerM] = result1;
\end{lstlisting}}
\newpage
\subsubsection[{\small FCOMPW: Floating-Point Compare Word}]{FCOMPW\hfill Floating-Point Compare Word}
\label{instr:FCOMPW}\index{fcompw}
 
\paragraph{Encoding} Format \textsf{ALU\_FWCWRR} (Section~\ref{sec:format-ALU-FWCWRR}), \verb|alucode3 = 010|\\

\bitfields{ 1/P, 2/11, 1/1, 1/0, 3/floatcomp, 6/registerW, 2/01, 3/010, 1/0, 6/registerY, 6/registerZ }


\paragraph{Syntax} \texttt{fcompw}\emph{floatcomp} \emph{registerW} = \emph{registerZ}, \emph{registerY}

\paragraph{Description} The \emph{registerZ} is compared to the \emph{registerY} using condition \emph{floatcomp}, assuming binary 32 floating-point numbers. The boolean result is stored into the \emph{registerW}. This instruction does not raise an invalid exception when handling a signaling NaN.


\paragraph{Modifier(s)}
\noindent\begin{itemize}
\item floatcomp (cf. floatcomp in section \ref{par:modifier-floatcomp} on page \pageref{par:modifier-floatcomp})
\end{itemize}

\paragraph{Execution} {Reservation table \textsf{ALU\_LITE} (Section~\ref{sec:reservation-ALU-LITE})}
~
{\begin{lstlisting}
stage ID:
  new argument4 = _ZX_3(floatcomp);
stage RR:
  new argument3 = GPR[registerY];
  new argument2 = GPR[registerZ];
  new result1 = floatcomp_32(argument4, argument2, argument3);
stage E1:
  GPR[registerW] = _ZX_32(result1);
\end{lstlisting}}
\newpage

\paragraph{Encoding} Format \textsf{ALU\_FWCWRR.W} (Section~\ref{sec:format-ALU-FWCWRR.W}), \verb|alucode3 = 010|\\

\bitfields{ 1/P, 2/00, 2/-- --, 27/upper27, 1/1, 2/11, 1/1, 1/0, 3/floatcomp, 6/registerW, 2/01, 3/010, 1/0, 1/0, 5/lower5, 6/registerZ }


\paragraph{Syntax} \texttt{fcompw}\emph{floatcomp} \emph{registerW} = \emph{registerZ}, \emph{upper27\_\discretionary{}{}{}lower5}

\paragraph{Description} The \emph{registerZ} is compared to the \emph{upper27\_\discretionary{}{}{}lower5} using condition \emph{floatcomp}, assuming binary 32 floating-point numbers. The boolean result is stored into the \emph{registerW}. This instruction does not raise an invalid exception when handling a signaling NaN.


\paragraph{Modifier(s)}
\noindent\begin{itemize}
\item floatcomp (cf. floatcomp in section \ref{par:modifier-floatcomp} on page \pageref{par:modifier-floatcomp})
\end{itemize}

\paragraph{Execution} {Reservation table \textsf{ALU\_LITE.X} (Section~\ref{sec:reservation-ALU-LITE.X})}
~
{\begin{lstlisting}
stage ID:
  new argument4 = _ZX_3(floatcomp);
stage RR:
  new argument3 = _WX_32(upper27_lower5);
  new argument2 = GPR[registerZ];
  new result1 = floatcomp_32(argument4, argument2, argument3);
stage E1:
  GPR[registerW] = _ZX_32(result1);
\end{lstlisting}}
\newpage
\subsubsection[{\small FDIVD: Floating-Point Divide Double Word}]{FDIVD\hfill Floating-Point Divide Double Word}
\label{instr:FDIVD}\index{fdivd}
 
\paragraph{Encoding} Format \textsf{ALU\_FDMUL} (Section~\ref{sec:format-ALU-FDMUL}), \verb|alucode3 = 111|\\

\bitfields{ 1/P, 2/11, 1/1, 1/fnegate, 3/floatmode, 6/registerW, 2/00, 3/111, 1/0, 6/registerY, 6/registerZ }


\paragraph{Syntax} \texttt{fdivd}\emph{fnegate}\emph{floatmode} \emph{registerW} = \emph{registerZ}, \emph{registerY}

\paragraph{Description} The \emph{registerZ} is divided by the \emph{registerY}, assuming binary 64 floating-point numbers. If \emph{fnegate} is set, the \emph{registerZ} double word is negated prior to the operation. The result is rounded to the binary 64 floating-point format according to the rounding mode and stored into the \emph{registerW}. This instruction may raise exception bits in the CS register.


\paragraph{Modifier(s)}
\noindent\begin{itemize}
\item floatmode (cf. floatmode in section \ref{par:modifier-floatmode} on page \pageref{par:modifier-floatmode})
\item fnegate (cf. fnegate in section \ref{par:modifier-fnegate} on page \pageref{par:modifier-fnegate})
\end{itemize}

\paragraph{Execution} {Reservation table \textsf{ALU\_FULL} (Section~\ref{sec:reservation-ALU-FULL})}
~
{\begin{lstlisting}
stage ID:
  new fnegate = _ZX_1(fnegate);
  new floatmode = _ZX_3(floatmode);
stage RR:
  new argument3 = GPR[registerY];
  new argument2 = GPR[registerZ];
  new RM = decode_riscv_float_rounding_mode(floatmode);
  new fsign = fnegate ? 0x8000000000000000 : 0;
  new result1 = f64_div(RM, argument2 ^ fsign, argument3);
stage E4:
  GPR[registerW] = result1;
  commit_float_exception_flags();
\end{lstlisting}}
\newpage
\subsubsection[{\small FDIVH: Floating-Point Divide Half Word}]{FDIVH\hfill Floating-Point Divide Half Word}
\label{instr:FDIVH}\index{fdivh}
 
\paragraph{Encoding} Format \textsf{ALU\_FHMUL} (Section~\ref{sec:format-ALU-FHMUL}), \verb|alucode3 = 111|\\

\bitfields{ 1/P, 2/11, 1/1, 1/fnegate, 3/floatmode, 6/registerW, 2/01, 3/111, 1/1, 6/registerY, 6/registerZ }


\paragraph{Syntax} \texttt{fdivh}\emph{fnegate}\emph{floatmode} \emph{registerW} = \emph{registerZ}, \emph{registerY}

\paragraph{Description} The \emph{registerZ} is divided by the \emph{registerY}, assuming binary 16 floating-point numbers. If \emph{fnegate} is set, the \emph{registerZ} half word is negated prior to the operation. The result is rounded to the binary 16 floating-point format according to the rounding mode and stored into the \emph{registerW}. This instruction may raise exception bits in the CS register.


\paragraph{Modifier(s)}
\noindent\begin{itemize}
\item floatmode (cf. floatmode in section \ref{par:modifier-floatmode} on page \pageref{par:modifier-floatmode})
\item fnegate (cf. fnegate in section \ref{par:modifier-fnegate} on page \pageref{par:modifier-fnegate})
\end{itemize}

\paragraph{Execution} {Reservation table \textsf{ALU\_FULL} (Section~\ref{sec:reservation-ALU-FULL})}
~
{\begin{lstlisting}
stage ID:
  new fnegate = _ZX_1(fnegate);
  new floatmode = _ZX_3(floatmode);
stage RR:
  new argument3 = GPR[registerY];
  new argument2 = GPR[registerZ];
  new RM = decode_riscv_float_rounding_mode(floatmode);
  new fsign = fnegate ? 0x8000 : 0;
  new result1 = f16_div(RM, argument2 ^ fsign, argument3);
stage E3:
  GPR[registerW] = result1;
  commit_float_exception_flags();
\end{lstlisting}}
\newpage
\subsubsection[{\small FDIVW: Floating-Point Divide Word}]{FDIVW\hfill Floating-Point Divide Word}
\label{instr:FDIVW}\index{fdivw}
 
\paragraph{Encoding} Format \textsf{ALU\_FWMUL} (Section~\ref{sec:format-ALU-FWMUL}), \verb|alucode3 = 111|\\

\bitfields{ 1/P, 2/11, 1/1, 1/fnegate, 3/floatmode, 6/registerW, 2/01, 3/111, 1/0, 6/registerY, 6/registerZ }


\paragraph{Syntax} \texttt{fdivw}\emph{fnegate}\emph{floatmode} \emph{registerW} = \emph{registerZ}, \emph{registerY}

\paragraph{Description} The \emph{registerZ} is divided by the \emph{registerY}, assuming binary 32 floating-point numbers. If \emph{fnegate} is set, the \emph{registerZ} word is negated prior to the operation. The result is rounded to the binary 32 floating-point format according to the rounding mode and stored into the \emph{registerW}. This instruction may raise exception bits in the CS register.


\paragraph{Modifier(s)}
\noindent\begin{itemize}
\item floatmode (cf. floatmode in section \ref{par:modifier-floatmode} on page \pageref{par:modifier-floatmode})
\item fnegate (cf. fnegate in section \ref{par:modifier-fnegate} on page \pageref{par:modifier-fnegate})
\end{itemize}

\paragraph{Execution} {Reservation table \textsf{ALU\_FULL} (Section~\ref{sec:reservation-ALU-FULL})}
~
{\begin{lstlisting}
stage ID:
  new fnegate = _ZX_1(fnegate);
  new floatmode = _ZX_3(floatmode);
stage RR:
  new argument3 = GPR[registerY];
  new argument2 = GPR[registerZ];
  new RM = decode_riscv_float_rounding_mode(floatmode);
  new fsign = fnegate ? 0x80000000 : 0;
  new result1 = f32_div(RM, argument2 ^ fsign, argument3);
stage E4:
  GPR[registerW] = result1;
  commit_float_exception_flags();
\end{lstlisting}}
\newpage
\subsubsection[{\small FEXTLHWQ: Floating Point Extend Lane Half Word to Word Quadruple}]{FEXTLHWQ\hfill Floating Point Extend Lane Half Word to Word Quadruple}
\label{instr:FEXTLHWQ}\index{fextlhwq}
 
\paragraph{Encoding} Format \textsf{ALU\_FWQEXTLWR} (Section~\ref{sec:format-ALU-FWQEXTLWR}), \verb|fpucode6 = 1001|\\

\bitfields{ 1/P, 2/11, 1/1, 1/1, 3/111, 5/registerM, 1/oddlanes, 2/11, 3/000, 1/1, 4/1001, 1/--, 1/0, 5/registerP, 1/0 }


\paragraph{Syntax} \texttt{fextlhwq}\emph{oddlanes} \emph{registerM} = \emph{registerP}

\paragraph{Description} The \emph{registerP} interpreted as eight binary 16 floating-point numbers packed into 128 bits. The even or odd half words of the \emph{registerP} are selected, depending on \emph{oddlanes}. The four selected half words are converted to binary 32 floating-point numbers, packed and stored into the \emph{registerM}. This instruction may raise invalid exception bit in the CS register.


\paragraph{Modifier(s)}
\noindent\begin{itemize}
\item oddlanes (cf. oddlanes in section \ref{par:modifier-oddlanes} on page \pageref{par:modifier-oddlanes})
\end{itemize}

\paragraph{Execution} {Reservation table \textsf{ALU\_LITE} (Section~\ref{sec:reservation-ALU-LITE})}
~
{\begin{lstlisting}
stage ID:
  new oddlanes = _ZX_1(oddlanes);
stage RR:
  new argument2 = PGR[registerP];
  if (oddlanes) {
    argument2 = argument2 >> 16;
  }
  for (I = 0; I < 4; I += 1) {
    result1.32[I] = f16_to_f32(_ZX_16(argument2.32[I]))
  };
stage E1:
  PGR[registerM] = result1;
\end{lstlisting}}
\newpage
\subsubsection[{\small FFMAD: Floating-Point Fused Multiply Add Double Word}]{FFMAD\hfill Floating-Point Fused Multiply Add Double Word}
\label{instr:FFMAD}\index{ffmad}
 
\paragraph{Encoding} Format \textsf{ALU\_FDFMA} (Section~\ref{sec:format-ALU-FDFMA}), \verb|alucode3 = 100|\\

\bitfields{ 1/P, 2/11, 1/1, 1/fnegate, 3/floatmode, 6/registerW, 2/00, 3/100, 1/0, 6/registerY, 6/registerZ }


\paragraph{Syntax} \texttt{ffmad}\emph{fnegate}\emph{floatmode} \emph{registerW} = \emph{registerZ}, \emph{registerY}

\paragraph{Description} The \emph{registerY} is multiplied by the \emph{registerZ}, assuming binary 64 floating-point numbers, and the product is added to the \emph{registerW}. If \emph{fnegate} is set, the \emph{registerW} and the \emph{registerZ} double words are negated prior to the operation. The result is rounded to the binary 64 floating-point format according to the rounding mode and stored into the \emph{registerW}. This instruction may raise exception bits in the CS register.


\paragraph{Modifier(s)}
\noindent\begin{itemize}
\item floatmode (cf. floatmode in section \ref{par:modifier-floatmode} on page \pageref{par:modifier-floatmode})
\item fnegate (cf. fnegate in section \ref{par:modifier-fnegate} on page \pageref{par:modifier-fnegate})
\end{itemize}

\paragraph{Execution} {Reservation table \textsf{ALU\_LITE} (Section~\ref{sec:reservation-ALU-LITE})}
~
{\begin{lstlisting}
stage ID:
  new fnegate = _ZX_1(fnegate);
  new floatmode = _ZX_3(floatmode);
stage RR:
  new argument3 = GPR[registerY];
  new argument2 = GPR[registerZ];
  new argument1 = GPR[registerW];
  new RM = decode_riscv_float_rounding_mode(floatmode);
  new fsign = fnegate ? 0x8000000000000000 : 0;
  new result1 = f64_mulAdd(RM, argument3, argument2 ^ fsign, argument1 ^ fsign);
stage E4:
  GPR[registerW] = result1;
  commit_float_exception_flags();
\end{lstlisting}}
\newpage
\subsubsection[{\small FFMADP: Floating-Point Fused Multiply Add Double Word Pair}]{FFMADP\hfill Floating-Point Fused Multiply Add Double Word Pair}
\label{instr:FFMADP}\index{ffmadp}
 
\paragraph{Encoding} Format \textsf{ALU\_FDPFMA} (Section~\ref{sec:format-ALU-FDPFMA}), \verb|alucode3 = 100|\\

\bitfields{ 1/P, 2/11, 1/1, 1/fnegate, 3/floatmode, 5/registerM, 1/--, 2/11, 3/100, 1/0, 5/registerO, 1/0, 5/registerP, 1/1 }


\paragraph{Syntax} \texttt{ffmadp}\emph{fnegate}\emph{floatmode} \emph{registerM} = \emph{registerP}, \emph{registerO}

\paragraph{Description} The \emph{registerM}, the \emph{registerP} and the \emph{registerO} are interpreted as two binary 64 floating-point numbers packed into 128 bits. If \emph{fnegate} is set, the \emph{registerM} and the \emph{registerP} double words are negated prior to the operation. The \emph{registerO} double words are multiplied by the \emph{registerP} double words, assuming binary 64 floating-point numbers, and the product double words are added to the \emph{registerM} double words. The results are rounded to the binary 64 floating-point format according to the rounding mode and stored into the \emph{registerM}. This instruction may raise exception bits in the CS register.


\paragraph{Modifier(s)}
\noindent\begin{itemize}
\item floatmode (cf. floatmode in section \ref{par:modifier-floatmode} on page \pageref{par:modifier-floatmode})
\item fnegate (cf. fnegate in section \ref{par:modifier-fnegate} on page \pageref{par:modifier-fnegate})
\end{itemize}

\paragraph{Execution} {Reservation table \textsf{ALU\_LITE} (Section~\ref{sec:reservation-ALU-LITE})}
~
{\begin{lstlisting}
stage ID:
  new fnegate = _ZX_1(fnegate);
  new floatmode = _ZX_3(floatmode);
stage RR:
  new argument3 = PGR[registerO];
  new argument2 = PGR[registerP];
  new argument1 = PGR[registerM];
  new RM = decode_riscv_float_rounding_mode(floatmode);
  new fsign = fnegate ? 0x8000000000000000 : 0;
  for (I = 0; I < 2; I += 1) {
    result1.64[I] = f64_mulAdd(RM, argument3.64[I], argument2.64[I] ^ fsign, argument1.64[I] ^ fsign)
  };
stage E4:
  PGR[registerM] = result1;
  commit_float_exception_flags();
\end{lstlisting}}
\newpage
\subsubsection[{\small FFMAH: Floating-Point Fused Multiply Add Half Word}]{FFMAH\hfill Floating-Point Fused Multiply Add Half Word}
\label{instr:FFMAH}\index{ffmah}
 
\paragraph{Encoding} Format \textsf{ALU\_FHFMA} (Section~\ref{sec:format-ALU-FHFMA}), \verb|alucode3 = 100|\\

\bitfields{ 1/P, 2/11, 1/1, 1/fnegate, 3/floatmode, 6/registerW, 2/01, 3/100, 1/1, 6/registerY, 6/registerZ }


\paragraph{Syntax} \texttt{ffmah}\emph{fnegate}\emph{floatmode} \emph{registerW} = \emph{registerZ}, \emph{registerY}

\paragraph{Description} The \emph{registerY} is multiplied by the \emph{registerZ}, assuming binary 16 floating-point numbers, and the product is added to the \emph{registerW}. If \emph{fnegate} is set, the \emph{registerW} and the \emph{registerZ} half words are negated prior to the operation. The result is rounded to the binary 16 floating-point format according to the rounding mode and stored into the \emph{registerW}. This instruction may raise exception bits in the CS register.


\paragraph{Modifier(s)}
\noindent\begin{itemize}
\item floatmode (cf. floatmode in section \ref{par:modifier-floatmode} on page \pageref{par:modifier-floatmode})
\item fnegate (cf. fnegate in section \ref{par:modifier-fnegate} on page \pageref{par:modifier-fnegate})
\end{itemize}

\paragraph{Execution} {Reservation table \textsf{ALU\_LITE} (Section~\ref{sec:reservation-ALU-LITE})}
~
{\begin{lstlisting}
stage ID:
  new fnegate = _ZX_1(fnegate);
  new floatmode = _ZX_3(floatmode);
stage RR:
  new argument3 = GPR[registerY];
  new argument2 = GPR[registerZ];
  new argument1 = GPR[registerW];
  new RM = decode_riscv_float_rounding_mode(floatmode);
  new fsign = fnegate ? 0x8000 : 0;
  new result1 = f16_mulAdd(RM, argument3, argument2 ^ fsign, argument1 ^ fsign);
stage E3:
  GPR[registerW] = result1;
  commit_float_exception_flags();
\end{lstlisting}}
\newpage
\subsubsection[{\small FFMAHO: Floating-Point Fused Multiply Add Half Word Octuple}]{FFMAHO\hfill Floating-Point Fused Multiply Add Half Word Octuple}
\label{instr:FFMAHO}\index{ffmaho}
 
\paragraph{Encoding} Format \textsf{ALU\_FHOFMA} (Section~\ref{sec:format-ALU-FHOFMA}), \verb|alucode3 = 100|\\

\bitfields{ 1/P, 2/11, 1/1, 1/fnegate, 3/floatmode, 5/registerM, 1/--, 2/11, 3/100, 1/1, 5/registerO, 1/0, 5/registerP, 1/1 }


\paragraph{Syntax} \texttt{ffmaho}\emph{fnegate}\emph{floatmode} \emph{registerM} = \emph{registerP}, \emph{registerO}

\paragraph{Description} The \emph{registerM}, the \emph{registerP} and the \emph{registerO} are interpreted as eight binary 16 floating-point numbers packed into 128 bits. If \emph{fnegate} is set, the \emph{registerM} and the \emph{registerP} half words are negated prior to the operation. The \emph{registerO} half words are multiplied by the \emph{registerP} half words, assuming binary 16 floating-point numbers, and the product half words are added to the \emph{registerM} half words. The eight results are rounded to the binary 16 floating-point format according to the rounding mode and stored into the \emph{registerM}. This instruction may raise exception bits in the CS register.


\paragraph{Modifier(s)}
\noindent\begin{itemize}
\item floatmode (cf. floatmode in section \ref{par:modifier-floatmode} on page \pageref{par:modifier-floatmode})
\item fnegate (cf. fnegate in section \ref{par:modifier-fnegate} on page \pageref{par:modifier-fnegate})
\end{itemize}

\paragraph{Execution} {Reservation table \textsf{ALU\_LITE} (Section~\ref{sec:reservation-ALU-LITE})}
~
{\begin{lstlisting}
stage ID:
  new fnegate = _ZX_1(fnegate);
  new floatmode = _ZX_3(floatmode);
stage RR:
  new argument3 = PGR[registerO];
  new argument2 = PGR[registerP];
  new argument1 = PGR[registerM];
  new RM = decode_riscv_float_rounding_mode(floatmode);
  new fsign = fnegate ? 0x8000 : 0;
  for (I = 0; I < 8; I += 1) {
    result1.16[I] = f16_mulAdd(RM, argument3.16[I], argument2.16[I] ^ fsign, argument1.16[I] ^ fsign)
  };
stage E4:
  PGR[registerM] = result1;
  commit_float_exception_flags();
\end{lstlisting}}
\newpage
\subsubsection[{\small FFMAW: Floating-Point Fused Multiply Add Word}]{FFMAW\hfill Floating-Point Fused Multiply Add Word}
\label{instr:FFMAW}\index{ffmaw}
 
\paragraph{Encoding} Format \textsf{ALU\_FWFMA} (Section~\ref{sec:format-ALU-FWFMA}), \verb|alucode3 = 100|\\

\bitfields{ 1/P, 2/11, 1/1, 1/fnegate, 3/floatmode, 6/registerW, 2/01, 3/100, 1/0, 6/registerY, 6/registerZ }


\paragraph{Syntax} \texttt{ffmaw}\emph{fnegate}\emph{floatmode} \emph{registerW} = \emph{registerZ}, \emph{registerY}

\paragraph{Description} The \emph{registerY} is multiplied by the \emph{registerZ}, assuming binary 32 floating-point numbers, and the product is added to the \emph{registerW}. If \emph{fnegate} is set, the \emph{registerW} and the \emph{registerZ} words are negated prior to the operation. The result is rounded to the binary 32 floating-point format according to the rounding mode and stored into the \emph{registerW}. This instruction may raise exception bits in the CS register.


\paragraph{Modifier(s)}
\noindent\begin{itemize}
\item floatmode (cf. floatmode in section \ref{par:modifier-floatmode} on page \pageref{par:modifier-floatmode})
\item fnegate (cf. fnegate in section \ref{par:modifier-fnegate} on page \pageref{par:modifier-fnegate})
\end{itemize}

\paragraph{Execution} {Reservation table \textsf{ALU\_LITE} (Section~\ref{sec:reservation-ALU-LITE})}
~
{\begin{lstlisting}
stage ID:
  new fnegate = _ZX_1(fnegate);
  new floatmode = _ZX_3(floatmode);
stage RR:
  new argument3 = GPR[registerY];
  new argument2 = GPR[registerZ];
  new argument1 = GPR[registerW];
  new RM = decode_riscv_float_rounding_mode(floatmode);
  new fsign = fnegate ? 0x80000000 : 0;
  new result1 = f32_mulAdd(RM, argument3, argument2 ^ fsign, argument1 ^ fsign);
stage E4:
  GPR[registerW] = result1;
  commit_float_exception_flags();
\end{lstlisting}}
\newpage
\subsubsection[{\small FFMAWC: Floating-Point Multiply Add Word Complex}]{FFMAWC\hfill Floating-Point Multiply Add Word Complex}
\label{instr:FFMAWC}\index{ffmawc}
 
\paragraph{Encoding} Format \textsf{ALU\_FWCOP} (Section~\ref{sec:format-ALU-FWCOP}), \verb|alucode6 = 11|\\

\bitfields{ 1/P, 2/11, 1/1, 1/imultiply, 3/floatmode, 6/registerW, 2/00, 2/11, 1/conjugate, 1/1, 6/registerY, 6/registerZ }


\paragraph{Syntax} \texttt{ffmawc}\emph{conjugate}\emph{imultiply}\emph{floatmode} \emph{registerW} = \emph{registerZ}, \emph{registerY}

\paragraph{Description} The \emph{registerW}, the \emph{registerY} and the \emph{registerZ} are interpreted as binary 32 floating-point complex numbers. If \emph{conjugate} is set, the complex conjugate of the \emph{registerZ} is used in the multiplication. The complex product is added to the \emph{registerW} and rounded according to the rounding mode. If \emph{imultiply} is set, the result is multiplied by the imaginary unit. The resulting binary 32 floating-point complex number is stored back into the \emph{registerW}. This instruction may raise exception bits in the CS register.


\paragraph{Modifier(s)}
\noindent\begin{itemize}
\item floatmode (cf. floatmode in section \ref{par:modifier-floatmode} on page \pageref{par:modifier-floatmode})
\item imultiply (cf. imultiply in section \ref{par:modifier-imultiply} on page \pageref{par:modifier-imultiply})
\item conjugate (cf. conjugate in section \ref{par:modifier-conjugate} on page \pageref{par:modifier-conjugate})
\end{itemize}

\paragraph{Execution} {Reservation table \textsf{ALU\_LITE} (Section~\ref{sec:reservation-ALU-LITE})}
~
{\begin{lstlisting}
stage ID:
  new conjugate = _ZX_1(conjugate);
  new imultiply = _ZX_1(imultiply);
  new floatmode = _ZX_3(floatmode);
stage RR:
  new argument3 = GPR[registerY];
  new argument2 = GPR[registerZ];
  new argument1 = GPR[registerW];
  new RM = decode_riscv_float_rounding_mode(floatmode);
  if (conjugate) {
    argument2.64[0] = fconj_32_32(argument2.64[0]);
  }
  result1.64[0] = ffmac_32_32(RM, argument2.64[0], argument3.64[0], argument1.64[0]);
  if (imultiply) {
    new result1 = _SWAP_32(result1);
    result1.32[0] = result1.32[0] ^ 0x80000000;
  }
stage E4:
  GPR[registerW] = result1;
  commit_float_exception_flags();
\end{lstlisting}}
\newpage
\subsubsection[{\small FFMAWQ: Floating-Point Fused Multiply Add Word Quadruple}]{FFMAWQ\hfill Floating-Point Fused Multiply Add Word Quadruple}
\label{instr:FFMAWQ}\index{ffmawq}
 
\paragraph{Encoding} Format \textsf{ALU\_FWQFMA} (Section~\ref{sec:format-ALU-FWQFMA}), \verb|alucode3 = 100|\\

\bitfields{ 1/P, 2/11, 1/1, 1/fnegate, 3/floatmode, 5/registerM, 1/--, 2/11, 3/100, 1/1, 5/registerO, 1/0, 5/registerP, 1/0 }


\paragraph{Syntax} \texttt{ffmawq}\emph{fnegate}\emph{floatmode} \emph{registerM} = \emph{registerP}, \emph{registerO}

\paragraph{Description} The \emph{registerM}, the \emph{registerP} and the \emph{registerO} are interpreted as four binary 32 floating-point numbers packed into 128 bits. If \emph{fnegate} is set, the \emph{registerM} and the \emph{registerP} words are negated prior to the operation. The \emph{registerO} words are multiplied by the \emph{registerP} words, assuming binary 32 floating-point numbers, and the product words are added to the \emph{registerM} words. The four results are rounded to the binary 32 floating-point format according to the rounding mode and stored into the \emph{registerM}. This instruction may raise exception bits in the CS register.


\paragraph{Modifier(s)}
\noindent\begin{itemize}
\item floatmode (cf. floatmode in section \ref{par:modifier-floatmode} on page \pageref{par:modifier-floatmode})
\item fnegate (cf. fnegate in section \ref{par:modifier-fnegate} on page \pageref{par:modifier-fnegate})
\end{itemize}

\paragraph{Execution} {Reservation table \textsf{ALU\_LITE} (Section~\ref{sec:reservation-ALU-LITE})}
~
{\begin{lstlisting}
stage ID:
  new fnegate = _ZX_1(fnegate);
  new floatmode = _ZX_3(floatmode);
stage RR:
  new argument3 = PGR[registerO];
  new argument2 = PGR[registerP];
  new argument1 = PGR[registerM];
  new RM = decode_riscv_float_rounding_mode(floatmode);
  new fsign = fnegate ? 0x80000000 : 0;
  for (I = 0; I < 4; I += 1) {
    result1.32[I] = f32_mulAdd(RM, argument3.32[I], argument2.32[I] ^ fsign, argument1.32[I] ^ fsign)
  };
stage E4:
  PGR[registerM] = result1;
  commit_float_exception_flags();
\end{lstlisting}}
\newpage
\subsubsection[{\small FFMSD: Floating-Point Fused Multiply Subtract From Double Word}]{FFMSD\hfill Floating-Point Fused Multiply Subtract From Double Word}
\label{instr:FFMSD}\index{ffmsd}
 
\paragraph{Encoding} Format \textsf{ALU\_FDFMA} (Section~\ref{sec:format-ALU-FDFMA}), \verb|alucode3 = 101|\\

\bitfields{ 1/P, 2/11, 1/1, 1/fnegate, 3/floatmode, 6/registerW, 2/00, 3/101, 1/0, 6/registerY, 6/registerZ }


\paragraph{Syntax} \texttt{ffmsd}\emph{fnegate}\emph{floatmode} \emph{registerW} = \emph{registerZ}, \emph{registerY}

\paragraph{Description} The \emph{registerY} is multiplied by the \emph{registerZ}, assuming binary 64 floating-point numbers, and the product is subtracted from the \emph{registerW}. If \emph{fnegate} is set, the \emph{registerW} and the \emph{registerZ} double words are negated prior to the operation. The result is rounded to the binary 64 floating-point format according to the rounding mode and stored into the \emph{registerW}. This instruction may raise exception bits in the CS register.


\paragraph{Modifier(s)}
\noindent\begin{itemize}
\item floatmode (cf. floatmode in section \ref{par:modifier-floatmode} on page \pageref{par:modifier-floatmode})
\item fnegate (cf. fnegate in section \ref{par:modifier-fnegate} on page \pageref{par:modifier-fnegate})
\end{itemize}

\paragraph{Execution} {Reservation table \textsf{ALU\_LITE} (Section~\ref{sec:reservation-ALU-LITE})}
~
{\begin{lstlisting}
stage ID:
  new fnegate = _ZX_1(fnegate);
  new floatmode = _ZX_3(floatmode);
stage RR:
  new argument3 = GPR[registerY];
  new argument2 = GPR[registerZ];
  new argument1 = GPR[registerW];
  new RM = decode_riscv_float_rounding_mode(floatmode);
  new fsign = fnegate ? 0x8000000000000000 : 0;
  new result1 = f64_mulnAdd(RM, argument3, argument2 ^ fsign, argument1 ^ fsign);
stage E4:
  GPR[registerW] = result1;
  commit_float_exception_flags();
\end{lstlisting}}
\newpage
\subsubsection[{\small FFMSDP: Floating-Point Fused Multiply Subtract From Double Word Pair}]{FFMSDP\hfill Floating-Point Fused Multiply Subtract From Double Word Pair}
\label{instr:FFMSDP}\index{ffmsdp}
 
\paragraph{Encoding} Format \textsf{ALU\_FDPFMA} (Section~\ref{sec:format-ALU-FDPFMA}), \verb|alucode3 = 101|\\

\bitfields{ 1/P, 2/11, 1/1, 1/fnegate, 3/floatmode, 5/registerM, 1/--, 2/11, 3/101, 1/0, 5/registerO, 1/0, 5/registerP, 1/1 }


\paragraph{Syntax} \texttt{ffmsdp}\emph{fnegate}\emph{floatmode} \emph{registerM} = \emph{registerP}, \emph{registerO}

\paragraph{Description} The \emph{registerM}, the \emph{registerP} and the \emph{registerO} are interpreted as two binary 64 floating-point numbers packed into 128 bits. If \emph{fnegate} is set, the \emph{registerM} and the \emph{registerP} double words are negated prior to the operation. The \emph{registerO} double words are multiplied by the \emph{registerP} double words, assuming binary 64 floating-point numbers, and the product double words are subtracted from the \emph{registerM} double words. The results are rounded to the binary 64 floating-point format according to the rounding mode and stored into the \emph{registerM}. This instruction may raise exception bits in the CS register.


\paragraph{Modifier(s)}
\noindent\begin{itemize}
\item floatmode (cf. floatmode in section \ref{par:modifier-floatmode} on page \pageref{par:modifier-floatmode})
\item fnegate (cf. fnegate in section \ref{par:modifier-fnegate} on page \pageref{par:modifier-fnegate})
\end{itemize}

\paragraph{Execution} {Reservation table \textsf{ALU\_LITE} (Section~\ref{sec:reservation-ALU-LITE})}
~
{\begin{lstlisting}
stage ID:
  new fnegate = _ZX_1(fnegate);
  new floatmode = _ZX_3(floatmode);
stage RR:
  new argument3 = PGR[registerO];
  new argument2 = PGR[registerP];
  new argument1 = PGR[registerM];
  new RM = decode_riscv_float_rounding_mode(floatmode);
  new fsign = fnegate ? 0x8000000000000000 : 0;
  for (I = 0; I < 2; I += 1) {
    result1.64[I] = f64_mulnAdd(RM, argument3.64[I], argument2.64[I] ^ fsign, argument1.64[I] ^ fsign)
  };
stage E4:
  PGR[registerM] = result1;
  commit_float_exception_flags();
\end{lstlisting}}
\newpage
\subsubsection[{\small FFMSH: Floating-Point Fused Multiply Subtract From Half Word}]{FFMSH\hfill Floating-Point Fused Multiply Subtract From Half Word}
\label{instr:FFMSH}\index{ffmsh}
 
\paragraph{Encoding} Format \textsf{ALU\_FHFMA} (Section~\ref{sec:format-ALU-FHFMA}), \verb|alucode3 = 101|\\

\bitfields{ 1/P, 2/11, 1/1, 1/fnegate, 3/floatmode, 6/registerW, 2/01, 3/101, 1/1, 6/registerY, 6/registerZ }


\paragraph{Syntax} \texttt{ffmsh}\emph{fnegate}\emph{floatmode} \emph{registerW} = \emph{registerZ}, \emph{registerY}

\paragraph{Description} The \emph{registerY} is multiplied by the \emph{registerZ}, assuming binary 16 floating-point numbers, and the product is subtracted from the \emph{registerW}. If \emph{fnegate} is set, the \emph{registerW} and the \emph{registerZ} half words are negated prior to the operation. The result is rounded to the binary 16 floating-point format according to the rounding mode and stored into the \emph{registerW}. This instruction may raise exception bits in the CS register.


\paragraph{Modifier(s)}
\noindent\begin{itemize}
\item floatmode (cf. floatmode in section \ref{par:modifier-floatmode} on page \pageref{par:modifier-floatmode})
\item fnegate (cf. fnegate in section \ref{par:modifier-fnegate} on page \pageref{par:modifier-fnegate})
\end{itemize}

\paragraph{Execution} {Reservation table \textsf{ALU\_LITE} (Section~\ref{sec:reservation-ALU-LITE})}
~
{\begin{lstlisting}
stage ID:
  new fnegate = _ZX_1(fnegate);
  new floatmode = _ZX_3(floatmode);
stage RR:
  new argument3 = GPR[registerY];
  new argument2 = GPR[registerZ];
  new argument1 = GPR[registerW];
  new RM = decode_riscv_float_rounding_mode(floatmode);
  new fsign = fnegate ? 0x8000 : 0;
  new result1 = f16_mulnAdd(RM, argument3, argument2 ^ fsign, argument1 ^ fsign);
stage E3:
  GPR[registerW] = result1;
  commit_float_exception_flags();
\end{lstlisting}}
\newpage
\subsubsection[{\small FFMSHO: Floating-Point Fused Multiply Subtract From Half Word Octuple}]{FFMSHO\hfill Floating-Point Fused Multiply Subtract From Half Word Octuple}
\label{instr:FFMSHO}\index{ffmsho}
 
\paragraph{Encoding} Format \textsf{ALU\_FHOFMA} (Section~\ref{sec:format-ALU-FHOFMA}), \verb|alucode3 = 101|\\

\bitfields{ 1/P, 2/11, 1/1, 1/fnegate, 3/floatmode, 5/registerM, 1/--, 2/11, 3/101, 1/1, 5/registerO, 1/0, 5/registerP, 1/1 }


\paragraph{Syntax} \texttt{ffmsho}\emph{fnegate}\emph{floatmode} \emph{registerM} = \emph{registerP}, \emph{registerO}

\paragraph{Description} The \emph{registerM}, the \emph{registerP} and the \emph{registerO} are interpreted as eight binary 16 floating-point numbers packed into 128 bits. If \emph{fnegate} is set, the \emph{registerM} and the \emph{registerP} words are negated prior to the operation. The \emph{registerO} half words are multiplied by the \emph{registerP} half words, assuming binary 16 floating-point numbers, and the product half words are subtracted from the \emph{registerM} half words. The eight results are rounded to the binary 16 floating-point format according to the rounding mode and stored into the \emph{registerM}. This instruction may raise exception bits in the CS register.


\paragraph{Modifier(s)}
\noindent\begin{itemize}
\item floatmode (cf. floatmode in section \ref{par:modifier-floatmode} on page \pageref{par:modifier-floatmode})
\item fnegate (cf. fnegate in section \ref{par:modifier-fnegate} on page \pageref{par:modifier-fnegate})
\end{itemize}

\paragraph{Execution} {Reservation table \textsf{ALU\_LITE} (Section~\ref{sec:reservation-ALU-LITE})}
~
{\begin{lstlisting}
stage ID:
  new fnegate = _ZX_1(fnegate);
  new floatmode = _ZX_3(floatmode);
stage RR:
  new argument3 = PGR[registerO];
  new argument2 = PGR[registerP];
  new argument1 = PGR[registerM];
  new RM = decode_riscv_float_rounding_mode(floatmode);
  new fsign = fnegate ? 0x8000 : 0;
  for (I = 0; I < 8; I += 1) {
    result1.16[I] = f16_mulnAdd(RM, argument3.16[I], argument2.16[I] ^ fsign, argument1.16[I] ^ fsign)
  };
stage E4:
  PGR[registerM] = result1;
  commit_float_exception_flags();
\end{lstlisting}}
\newpage
\subsubsection[{\small FFMSW: Floating-Point Fused Multiply Subtract From Word}]{FFMSW\hfill Floating-Point Fused Multiply Subtract From Word}
\label{instr:FFMSW}\index{ffmsw}
 
\paragraph{Encoding} Format \textsf{ALU\_FWFMA} (Section~\ref{sec:format-ALU-FWFMA}), \verb|alucode3 = 101|\\

\bitfields{ 1/P, 2/11, 1/1, 1/fnegate, 3/floatmode, 6/registerW, 2/01, 3/101, 1/0, 6/registerY, 6/registerZ }


\paragraph{Syntax} \texttt{ffmsw}\emph{fnegate}\emph{floatmode} \emph{registerW} = \emph{registerZ}, \emph{registerY}

\paragraph{Description} The \emph{registerY} is multiplied by the \emph{registerZ}, assuming binary 32 floating-point numbers, and the product is subtracted from the \emph{registerW}. If \emph{fnegate} is set, the \emph{registerW} and the \emph{registerZ} words are negated prior to the operation. The result is rounded to the binary 32 floating-point format according to the rounding mode and stored into the \emph{registerW}. This instruction may raise exception bits in the CS register.


\paragraph{Modifier(s)}
\noindent\begin{itemize}
\item floatmode (cf. floatmode in section \ref{par:modifier-floatmode} on page \pageref{par:modifier-floatmode})
\item fnegate (cf. fnegate in section \ref{par:modifier-fnegate} on page \pageref{par:modifier-fnegate})
\end{itemize}

\paragraph{Execution} {Reservation table \textsf{ALU\_LITE} (Section~\ref{sec:reservation-ALU-LITE})}
~
{\begin{lstlisting}
stage ID:
  new fnegate = _ZX_1(fnegate);
  new floatmode = _ZX_3(floatmode);
stage RR:
  new argument3 = GPR[registerY];
  new argument2 = GPR[registerZ];
  new argument1 = GPR[registerW];
  new RM = decode_riscv_float_rounding_mode(floatmode);
  new fsign = fnegate ? 0x80000000 : 0;
  new result1 = f32_mulnAdd(RM, argument3, argument2 ^ fsign, argument1 ^ fsign);
stage E4:
  GPR[registerW] = result1;
  commit_float_exception_flags();
\end{lstlisting}}
\newpage
\subsubsection[{\small FFMSWQ: Floating-Point Fused Multiply Subtract From Word Quadruple}]{FFMSWQ\hfill Floating-Point Fused Multiply Subtract From Word Quadruple}
\label{instr:FFMSWQ}\index{ffmswq}
 
\paragraph{Encoding} Format \textsf{ALU\_FWQFMA} (Section~\ref{sec:format-ALU-FWQFMA}), \verb|alucode3 = 101|\\

\bitfields{ 1/P, 2/11, 1/1, 1/fnegate, 3/floatmode, 5/registerM, 1/--, 2/11, 3/101, 1/1, 5/registerO, 1/0, 5/registerP, 1/0 }


\paragraph{Syntax} \texttt{ffmswq}\emph{fnegate}\emph{floatmode} \emph{registerM} = \emph{registerP}, \emph{registerO}

\paragraph{Description} The \emph{registerM}, the \emph{registerP} and the \emph{registerO} are interpreted as four binary 32 floating-point numbers packed into 128 bits. If \emph{fnegate} is set, the \emph{registerM} and the \emph{registerP} words are negated prior to the operation. The \emph{registerO} words are multiplied by the \emph{registerP} words, assuming binary 32 floating-point numbers, and the product words are subtracted from the \emph{registerM} words. The four results are rounded to the binary 32 floating-point format according to the rounding mode and stored into the \emph{registerM}. This instruction may raise exception bits in the CS register.


\paragraph{Modifier(s)}
\noindent\begin{itemize}
\item floatmode (cf. floatmode in section \ref{par:modifier-floatmode} on page \pageref{par:modifier-floatmode})
\item fnegate (cf. fnegate in section \ref{par:modifier-fnegate} on page \pageref{par:modifier-fnegate})
\end{itemize}

\paragraph{Execution} {Reservation table \textsf{ALU\_LITE} (Section~\ref{sec:reservation-ALU-LITE})}
~
{\begin{lstlisting}
stage ID:
  new fnegate = _ZX_1(fnegate);
  new floatmode = _ZX_3(floatmode);
stage RR:
  new argument3 = PGR[registerO];
  new argument2 = PGR[registerP];
  new argument1 = PGR[registerM];
  new RM = decode_riscv_float_rounding_mode(floatmode);
  new fsign = fnegate ? 0x80000000 : 0;
  for (I = 0; I < 4; I += 1) {
    result1.32[I] = f32_mulnAdd(RM, argument3.32[I], argument2.32[I] ^ fsign, argument1.32[I] ^ fsign)
  };
stage E4:
  PGR[registerM] = result1;
  commit_float_exception_flags();
\end{lstlisting}}
\newpage
\subsubsection[{\small FIXEDD: Floating Point Conversion to Fixed-Point Double Word}]{FIXEDD\hfill Floating Point Conversion to Fixed-Point Double Word}
\label{instr:FIXEDD}\index{fixedd}
 
\paragraph{Encoding} Format \textsf{ALU\_FDFIXED} (Section~\ref{sec:format-ALU-FDFIXED}), \verb|fpucode6 = 0100|\\

\bitfields{ 1/P, 2/11, 1/1, 1/1, 3/floatmode, 6/registerW, 2/00, 3/001, 1/0, 4/0100, 1/--, 1/0, 6/registerZ }


\paragraph{Syntax} \texttt{fixedd}\emph{floatmode} \emph{registerW} = \emph{registerZ}

\paragraph{Description} The \emph{registerZ} interpreted as a 64-bit floating-point number is converted to a 64-bit signed integer according to the rounding mode. This instruction may raise exception bits in the CS register.


\paragraph{Modifier(s)}
\noindent\begin{itemize}
\item floatmode (cf. floatmode in section \ref{par:modifier-floatmode} on page \pageref{par:modifier-floatmode})
\end{itemize}

\paragraph{Execution} {Reservation table \textsf{ALU\_LITE} (Section~\ref{sec:reservation-ALU-LITE})}
~
{\begin{lstlisting}
stage ID:
  new floatmode = _ZX_3(floatmode);
stage RR:
  new argument2 = GPR[registerZ];
  new RM = decode_riscv_float_rounding_mode(floatmode);
  new result1 = f64_to_i64(RM, argument2);
stage E4:
  GPR[registerW] = result1;
  commit_float_exception_flags();
\end{lstlisting}}
\newpage
\subsubsection[{\small FIXEDDP: Floating Point Conversion to Fixed-Point Double Word Pair}]{FIXEDDP\hfill Floating Point Conversion to Fixed-Point Double Word Pair}
\label{instr:FIXEDDP}\index{fixeddp}
 
\paragraph{Encoding} Format \textsf{ALU\_FDPFIXED} (Section~\ref{sec:format-ALU-FDPFIXED}), \verb|fpucode6 = 0100|\\

\bitfields{ 1/P, 2/11, 1/1, 1/1, 3/floatmode, 5/registerM, 1/--, 2/11, 3/001, 1/0, 4/0100, 1/--, 1/0, 5/registerP, 1/1 }


\paragraph{Syntax} \texttt{fixeddp}\emph{floatmode} \emph{registerM} = \emph{registerP}

\paragraph{Description} The \emph{registerP} is interpreted as two binary 64 floating-point numbers packed into 128 bits. The \emph{registerP} double words are converted to 64-bit signed integers according to the rounding mode. The two resulting double words are packed and stored into the \emph{registerM}. This instruction may raise exception bits in the CS register.


\paragraph{Modifier(s)}
\noindent\begin{itemize}
\item floatmode (cf. floatmode in section \ref{par:modifier-floatmode} on page \pageref{par:modifier-floatmode})
\end{itemize}

\paragraph{Execution} {Reservation table \textsf{ALU\_LITE} (Section~\ref{sec:reservation-ALU-LITE})}
~
{\begin{lstlisting}
stage ID:
  new floatmode = _ZX_3(floatmode);
stage RR:
  new argument2 = PGR[registerP];
  new RM = decode_riscv_float_rounding_mode(floatmode);
  for (I = 0; I < 2; I += 1) {
    result1.64[I] = f64_to_i64(RM, argument2.64[I])
  };
stage E4:
  PGR[registerM] = result1;
  commit_float_exception_flags();
\end{lstlisting}}
\newpage
\subsubsection[{\small FIXEDDW: Floating Point Conversion from Double Word to Fixed-Point Word}]{FIXEDDW\hfill Floating Point Conversion from Double Word to Fixed-Point Word}
\label{instr:FIXEDDW}\index{fixeddw}
 
\paragraph{Encoding} Format \textsf{ALU\_FWFIXED} (Section~\ref{sec:format-ALU-FWFIXED}), \verb|fpucode6 = 0110|\\

\bitfields{ 1/P, 2/11, 1/1, 1/1, 3/floatmode, 6/registerW, 2/01, 3/001, 1/0, 4/0110, 1/--, 1/0, 6/registerZ }


\paragraph{Syntax} \texttt{fixeddw}\emph{floatmode} \emph{registerW} = \emph{registerZ}

\paragraph{Description} The \emph{registerZ} interpreted as a 64-bit floating-point number is converted to a 32-bit signed integer according to the rounding mode. This instruction may raise exception bits in the CS register.


\paragraph{Modifier(s)}
\noindent\begin{itemize}
\item floatmode (cf. floatmode in section \ref{par:modifier-floatmode} on page \pageref{par:modifier-floatmode})
\end{itemize}

\paragraph{Execution} {Reservation table \textsf{ALU\_LITE} (Section~\ref{sec:reservation-ALU-LITE})}
~
{\begin{lstlisting}
stage ID:
  new floatmode = _ZX_3(floatmode);
stage RR:
  new argument2 = GPR[registerZ];
  new RM = decode_riscv_float_rounding_mode(floatmode);
  new result1 = f64_to_i32(RM, argument2);
stage E4:
  GPR[registerW] = result1;
  commit_float_exception_flags();
\end{lstlisting}}
\newpage
\subsubsection[{\small FIXEDHO: Floating Point Conversion to Fixed-Point Half Word Octuple}]{FIXEDHO\hfill Floating Point Conversion to Fixed-Point Half Word Octuple}
\label{instr:FIXEDHO}\index{fixedho}
 
\paragraph{Encoding} Format \textsf{ALU\_FHOFIXED} (Section~\ref{sec:format-ALU-FHOFIXED}), \verb|fpucode5 = 00100|\\

\bitfields{ 1/P, 2/11, 1/1, 1/1, 3/floatmode, 5/registerM, 1/--, 2/11, 3/001, 1/1, 5/00100, 1/0, 5/registerP, 1/1 }


\paragraph{Syntax} \texttt{fixedho}\emph{floatmode} \emph{registerM} = \emph{registerP}

\paragraph{Description} The \emph{registerP} is interpreted as eight binary 16 floating-point numbers packed into 128 bits. The \emph{registerP} half words are converted to 16-bit signed integers according to the rounding mode. The eight resulting half words are packed and stored into the \emph{registerM}. This instruction may raise exception bits in the CS register.


\paragraph{Modifier(s)}
\noindent\begin{itemize}
\item floatmode (cf. floatmode in section \ref{par:modifier-floatmode} on page \pageref{par:modifier-floatmode})
\end{itemize}

\paragraph{Execution} {Reservation table \textsf{ALU\_LITE} (Section~\ref{sec:reservation-ALU-LITE})}
~
{\begin{lstlisting}
stage ID:
  new floatmode = _ZX_3(floatmode);
stage RR:
  new argument2 = PGR[registerP];
  new RM = decode_riscv_float_rounding_mode(floatmode);
  for (I = 0; I < 8; I += 1) {
    result1.16[I] = f16_to_i16(RM, argument2.16[I])
  };
stage E4:
  PGR[registerM] = result1;
  commit_float_exception_flags();
\end{lstlisting}}
\newpage
\subsubsection[{\small FIXEDUD: Floating Point Conversion to Unsigned Fixed-Point Double Word}]{FIXEDUD\hfill Floating Point Conversion to Unsigned Fixed-Point Double Word}
\label{instr:FIXEDUD}\index{fixedud}
 
\paragraph{Encoding} Format \textsf{ALU\_FDFIXED} (Section~\ref{sec:format-ALU-FDFIXED}), \verb|fpucode6 = 0101|\\

\bitfields{ 1/P, 2/11, 1/1, 1/1, 3/floatmode, 6/registerW, 2/00, 3/001, 1/0, 4/0101, 1/--, 1/0, 6/registerZ }


\paragraph{Syntax} \texttt{fixedud}\emph{floatmode} \emph{registerW} = \emph{registerZ}

\paragraph{Description} The \emph{registerZ} interpreted as a 64-bit floating-point number is converted to a 64-bit unsigned integer according to the rounding mode. This instruction may raise exception bits in the CS register.


\paragraph{Modifier(s)}
\noindent\begin{itemize}
\item floatmode (cf. floatmode in section \ref{par:modifier-floatmode} on page \pageref{par:modifier-floatmode})
\end{itemize}

\paragraph{Execution} {Reservation table \textsf{ALU\_LITE} (Section~\ref{sec:reservation-ALU-LITE})}
~
{\begin{lstlisting}
stage ID:
  new floatmode = _ZX_3(floatmode);
stage RR:
  new argument2 = GPR[registerZ];
  new RM = decode_riscv_float_rounding_mode(floatmode);
  new result1 = f64_to_ui64(RM, argument2);
stage E4:
  GPR[registerW] = result1;
  commit_float_exception_flags();
\end{lstlisting}}
\newpage
\subsubsection[{\small FIXEDUDP: Floating Point Conversion to Unsigned Fixed-Point Double Word Pair}]{FIXEDUDP\hfill Floating Point Conversion to Unsigned Fixed-Point Double Word Pair}
\label{instr:FIXEDUDP}\index{fixedudp}
 
\paragraph{Encoding} Format \textsf{ALU\_FDPFIXED} (Section~\ref{sec:format-ALU-FDPFIXED}), \verb|fpucode6 = 0101|\\

\bitfields{ 1/P, 2/11, 1/1, 1/1, 3/floatmode, 5/registerM, 1/--, 2/11, 3/001, 1/0, 4/0101, 1/--, 1/0, 5/registerP, 1/1 }


\paragraph{Syntax} \texttt{fixedudp}\emph{floatmode} \emph{registerM} = \emph{registerP}

\paragraph{Description} The \emph{registerP} is interpreted as two binary 64 floating-point numbers packed into 128 bits. The \emph{registerP} double words are converted to 64-bit unsigned integers according to the rounding mode. The two resulting double words are packed and stored into the \emph{registerM}. This instruction may raise exception bits in the CS register.


\paragraph{Modifier(s)}
\noindent\begin{itemize}
\item floatmode (cf. floatmode in section \ref{par:modifier-floatmode} on page \pageref{par:modifier-floatmode})
\end{itemize}

\paragraph{Execution} {Reservation table \textsf{ALU\_LITE} (Section~\ref{sec:reservation-ALU-LITE})}
~
{\begin{lstlisting}
stage ID:
  new floatmode = _ZX_3(floatmode);
stage RR:
  new argument2 = PGR[registerP];
  new RM = decode_riscv_float_rounding_mode(floatmode);
  for (I = 0; I < 2; I += 1) {
    result1.64[I] = f64_to_ui64(RM, argument2.64[I])
  };
stage E4:
  PGR[registerM] = result1;
  commit_float_exception_flags();
\end{lstlisting}}
\newpage
\subsubsection[{\small FIXEDUDW: Floating Point Conversion from Double Word to Unsigned Fixed-Point Word}]{FIXEDUDW\hfill Floating Point Conversion from Double Word to Unsigned Fixed-Point Word}
\label{instr:FIXEDUDW}\index{fixedudw}
 
\paragraph{Encoding} Format \textsf{ALU\_FWFIXED} (Section~\ref{sec:format-ALU-FWFIXED}), \verb|fpucode6 = 0111|\\

\bitfields{ 1/P, 2/11, 1/1, 1/1, 3/floatmode, 6/registerW, 2/01, 3/001, 1/0, 4/0111, 1/--, 1/0, 6/registerZ }


\paragraph{Syntax} \texttt{fixedudw}\emph{floatmode} \emph{registerW} = \emph{registerZ}

\paragraph{Description} The \emph{registerZ} interpreted as a 64-bit floating-point number is converted to a 32-bit unsigned integer according to the rounding mode. This instruction may raise exception bits in the CS register.


\paragraph{Modifier(s)}
\noindent\begin{itemize}
\item floatmode (cf. floatmode in section \ref{par:modifier-floatmode} on page \pageref{par:modifier-floatmode})
\end{itemize}

\paragraph{Execution} {Reservation table \textsf{ALU\_LITE} (Section~\ref{sec:reservation-ALU-LITE})}
~
{\begin{lstlisting}
stage ID:
  new floatmode = _ZX_3(floatmode);
stage RR:
  new argument2 = GPR[registerZ];
  new RM = decode_riscv_float_rounding_mode(floatmode);
  new result1 = f64_to_ui32(RM, argument2);
stage E4:
  GPR[registerW] = result1;
  commit_float_exception_flags();
\end{lstlisting}}
\newpage
\subsubsection[{\small FIXEDUHO: Floating Point Conversion to Unsigned Fixed-Point Half Word Octuple}]{FIXEDUHO\hfill Floating Point Conversion to Unsigned Fixed-Point Half Word Octuple}
\label{instr:FIXEDUHO}\index{fixeduho}
 
\paragraph{Encoding} Format \textsf{ALU\_FHOFIXED} (Section~\ref{sec:format-ALU-FHOFIXED}), \verb|fpucode5 = 00101|\\

\bitfields{ 1/P, 2/11, 1/1, 1/1, 3/floatmode, 5/registerM, 1/--, 2/11, 3/001, 1/1, 5/00101, 1/0, 5/registerP, 1/1 }


\paragraph{Syntax} \texttt{fixeduho}\emph{floatmode} \emph{registerM} = \emph{registerP}

\paragraph{Description} The \emph{registerP} is interpreted as eight binary 16 floating-point numbers packed into 128 bits. The \emph{registerP} half words are converted to 16-bit unsigned integers according to the rounding mode. The eight resulting half words are packed and stored into the \emph{registerM}. This instruction may raise exception bits in the CS register.


\paragraph{Modifier(s)}
\noindent\begin{itemize}
\item floatmode (cf. floatmode in section \ref{par:modifier-floatmode} on page \pageref{par:modifier-floatmode})
\end{itemize}

\paragraph{Execution} {Reservation table \textsf{ALU\_LITE} (Section~\ref{sec:reservation-ALU-LITE})}
~
{\begin{lstlisting}
stage ID:
  new floatmode = _ZX_3(floatmode);
stage RR:
  new argument2 = PGR[registerP];
  new RM = decode_riscv_float_rounding_mode(floatmode);
  for (I = 0; I < 8; I += 1) {
    result1.16[I] = f16_to_ui16(RM, argument2.16[I])
  };
stage E4:
  PGR[registerM] = result1;
  commit_float_exception_flags();
\end{lstlisting}}
\newpage
\subsubsection[{\small FIXEDUW: Floating Point Conversion to Unsigned Fixed-Point}]{FIXEDUW\hfill Floating Point Conversion to Unsigned Fixed-Point}
\label{instr:FIXEDUW}\index{fixeduw}
 
\paragraph{Encoding} Format \textsf{ALU\_FWFIXED} (Section~\ref{sec:format-ALU-FWFIXED}), \verb|fpucode6 = 0101|\\

\bitfields{ 1/P, 2/11, 1/1, 1/1, 3/floatmode, 6/registerW, 2/01, 3/001, 1/0, 4/0101, 1/--, 1/0, 6/registerZ }


\paragraph{Syntax} \texttt{fixeduw}\emph{floatmode} \emph{registerW} = \emph{registerZ}

\paragraph{Description} The \emph{registerZ} interpreted as a 32-bit floating-point number is converted to a 32-bit unsigned integer according to the rounding mode. This instruction may raise exception bits in the CS register.


\paragraph{Modifier(s)}
\noindent\begin{itemize}
\item floatmode (cf. floatmode in section \ref{par:modifier-floatmode} on page \pageref{par:modifier-floatmode})
\end{itemize}

\paragraph{Execution} {Reservation table \textsf{ALU\_LITE} (Section~\ref{sec:reservation-ALU-LITE})}
~
{\begin{lstlisting}
stage ID:
  new floatmode = _ZX_3(floatmode);
stage RR:
  new argument2 = GPR[registerZ];
  new RM = decode_riscv_float_rounding_mode(floatmode);
  new result1 = f32_to_ui32(RM, argument2);
stage E4:
  GPR[registerW] = result1;
  commit_float_exception_flags();
\end{lstlisting}}
\newpage
\subsubsection[{\small FIXEDUWD: Floating Point Conversion from Word to Unsigned Fixed-Point Double Word}]{FIXEDUWD\hfill Floating Point Conversion from Word to Unsigned Fixed-Point Double Word}
\label{instr:FIXEDUWD}\index{fixeduwd}
 
\paragraph{Encoding} Format \textsf{ALU\_FDFIXED} (Section~\ref{sec:format-ALU-FDFIXED}), \verb|fpucode6 = 0111|\\

\bitfields{ 1/P, 2/11, 1/1, 1/1, 3/floatmode, 6/registerW, 2/00, 3/001, 1/0, 4/0111, 1/--, 1/0, 6/registerZ }


\paragraph{Syntax} \texttt{fixeduwd}\emph{floatmode} \emph{registerW} = \emph{registerZ}

\paragraph{Description} The \emph{registerZ} interpreted as a 32-bit floating-point number is converted to a 64-bit unsigned integer according to the rounding mode. This instruction may raise exception bits in the CS register.


\paragraph{Modifier(s)}
\noindent\begin{itemize}
\item floatmode (cf. floatmode in section \ref{par:modifier-floatmode} on page \pageref{par:modifier-floatmode})
\end{itemize}

\paragraph{Execution} {Reservation table \textsf{ALU\_LITE} (Section~\ref{sec:reservation-ALU-LITE})}
~
{\begin{lstlisting}
stage ID:
  new floatmode = _ZX_3(floatmode);
stage RR:
  new argument2 = GPR[registerZ];
  new RM = decode_riscv_float_rounding_mode(floatmode);
  new result1 = f32_to_ui64(RM, argument2);
stage E4:
  GPR[registerW] = result1;
  commit_float_exception_flags();
\end{lstlisting}}
\newpage
\subsubsection[{\small FIXEDUWP: Floating Point Conversion to Unsigned Fixed-Point Word Pair}]{FIXEDUWP\hfill Floating Point Conversion to Unsigned Fixed-Point Word Pair}
\label{instr:FIXEDUWP}\index{fixeduwp}
 
\paragraph{Encoding} Format \textsf{ALU\_FWPFIXED} (Section~\ref{sec:format-ALU-FWPFIXED}), \verb|fpucode6 = 0101|\\

\bitfields{ 1/P, 2/11, 1/1, 1/1, 3/floatmode, 6/registerW, 2/00, 3/001, 1/0, 4/0101, 1/--, 1/1, 6/registerZ }


\paragraph{Syntax} \texttt{fixeduwp}\emph{floatmode} \emph{registerW} = \emph{registerZ}

\paragraph{Description} The \emph{registerZ} is interpreted as two binary 32 floating-point numbers packed into 64 bits. The \emph{registerZ} words are converted to 32-bit unsigned integers according to the rounding mode. The two resulting words are packed and stored into the \emph{registerW}. This instruction may raise exception bits in the CS register.


\paragraph{Modifier(s)}
\noindent\begin{itemize}
\item floatmode (cf. floatmode in section \ref{par:modifier-floatmode} on page \pageref{par:modifier-floatmode})
\end{itemize}

\paragraph{Execution} {Reservation table \textsf{ALU\_LITE} (Section~\ref{sec:reservation-ALU-LITE})}
~
{\begin{lstlisting}
stage ID:
  new floatmode = _ZX_3(floatmode);
stage RR:
  new argument2 = GPR[registerZ];
  new RM = decode_riscv_float_rounding_mode(floatmode);
  for (I = 0; I < 2; I += 1) {
    result1.32[I] = f32_to_ui32(RM, argument2.32[I])
  };
stage E4:
  GPR[registerW] = result1;
  commit_float_exception_flags();
\end{lstlisting}}
\newpage
\subsubsection[{\small FIXEDUWQ: Floating Point Conversion to Unsigned Fixed-Point Word Quadruple}]{FIXEDUWQ\hfill Floating Point Conversion to Unsigned Fixed-Point Word Quadruple}
\label{instr:FIXEDUWQ}\index{fixeduwq}
 
\paragraph{Encoding} Format \textsf{ALU\_FWQFIXED} (Section~\ref{sec:format-ALU-FWQFIXED}), \verb|fpucode6 = 0101|\\

\bitfields{ 1/P, 2/11, 1/1, 1/1, 3/floatmode, 5/registerM, 1/--, 2/11, 3/001, 1/1, 4/0101, 1/--, 1/0, 5/registerP, 1/0 }


\paragraph{Syntax} \texttt{fixeduwq}\emph{floatmode} \emph{registerM} = \emph{registerP}

\paragraph{Description} The \emph{registerP} is interpreted as four binary 32 floating-point numbers packed into 128 bits. The \emph{registerP} words are converted to 32-bit unsigned integers according to the rounding mode. The four resulting words are packed and stored into the \emph{registerM}. This instruction may raise exception bits in the CS register.


\paragraph{Modifier(s)}
\noindent\begin{itemize}
\item floatmode (cf. floatmode in section \ref{par:modifier-floatmode} on page \pageref{par:modifier-floatmode})
\end{itemize}

\paragraph{Execution} {Reservation table \textsf{ALU\_LITE} (Section~\ref{sec:reservation-ALU-LITE})}
~
{\begin{lstlisting}
stage ID:
  new floatmode = _ZX_3(floatmode);
stage RR:
  new argument2 = PGR[registerP];
  new RM = decode_riscv_float_rounding_mode(floatmode);
  for (I = 0; I < 4; I += 1) {
    result1.32[I] = f32_to_ui32(RM, argument2.32[I])
  };
stage E4:
  PGR[registerM] = result1;
  commit_float_exception_flags();
\end{lstlisting}}
\newpage
\subsubsection[{\small FIXEDW: Floating Point Conversion to Fixed-Point}]{FIXEDW\hfill Floating Point Conversion to Fixed-Point}
\label{instr:FIXEDW}\index{fixedw}
 
\paragraph{Encoding} Format \textsf{ALU\_FWFIXED} (Section~\ref{sec:format-ALU-FWFIXED}), \verb|fpucode6 = 0100|\\

\bitfields{ 1/P, 2/11, 1/1, 1/1, 3/floatmode, 6/registerW, 2/01, 3/001, 1/0, 4/0100, 1/--, 1/0, 6/registerZ }


\paragraph{Syntax} \texttt{fixedw}\emph{floatmode} \emph{registerW} = \emph{registerZ}

\paragraph{Description} The \emph{registerZ} interpreted as a 32-bit floating-point number is converted to a 32-bit signed integer according to the rounding mode. This instruction may raise exception bits in the CS register.


\paragraph{Modifier(s)}
\noindent\begin{itemize}
\item floatmode (cf. floatmode in section \ref{par:modifier-floatmode} on page \pageref{par:modifier-floatmode})
\end{itemize}

\paragraph{Execution} {Reservation table \textsf{ALU\_LITE} (Section~\ref{sec:reservation-ALU-LITE})}
~
{\begin{lstlisting}
stage ID:
  new floatmode = _ZX_3(floatmode);
stage RR:
  new argument2 = GPR[registerZ];
  new RM = decode_riscv_float_rounding_mode(floatmode);
  new result1 = f32_to_i32(RM, argument2);
stage E4:
  GPR[registerW] = result1;
  commit_float_exception_flags();
\end{lstlisting}}
\newpage
\subsubsection[{\small FIXEDWD: Floating Point Conversion from Word to Fixed-Point Double Word}]{FIXEDWD\hfill Floating Point Conversion from Word to Fixed-Point Double Word}
\label{instr:FIXEDWD}\index{fixedwd}
 
\paragraph{Encoding} Format \textsf{ALU\_FDFIXED} (Section~\ref{sec:format-ALU-FDFIXED}), \verb|fpucode6 = 0110|\\

\bitfields{ 1/P, 2/11, 1/1, 1/1, 3/floatmode, 6/registerW, 2/00, 3/001, 1/0, 4/0110, 1/--, 1/0, 6/registerZ }


\paragraph{Syntax} \texttt{fixedwd}\emph{floatmode} \emph{registerW} = \emph{registerZ}

\paragraph{Description} The \emph{registerZ} interpreted as a 32-bit floating-point number is converted to a 64-bit signed integer according to the rounding mode. This instruction may raise exception bits in the CS register.


\paragraph{Modifier(s)}
\noindent\begin{itemize}
\item floatmode (cf. floatmode in section \ref{par:modifier-floatmode} on page \pageref{par:modifier-floatmode})
\end{itemize}

\paragraph{Execution} {Reservation table \textsf{ALU\_LITE} (Section~\ref{sec:reservation-ALU-LITE})}
~
{\begin{lstlisting}
stage ID:
  new floatmode = _ZX_3(floatmode);
stage RR:
  new argument2 = GPR[registerZ];
  new RM = decode_riscv_float_rounding_mode(floatmode);
  new result1 = f32_to_i64(RM, argument2);
stage E4:
  GPR[registerW] = result1;
  commit_float_exception_flags();
\end{lstlisting}}
\newpage
\subsubsection[{\small FIXEDWP: Floating Point Conversion to Fixed-Point Word Pair}]{FIXEDWP\hfill Floating Point Conversion to Fixed-Point Word Pair}
\label{instr:FIXEDWP}\index{fixedwp}
 
\paragraph{Encoding} Format \textsf{ALU\_FWPFIXED} (Section~\ref{sec:format-ALU-FWPFIXED}), \verb|fpucode6 = 0100|\\

\bitfields{ 1/P, 2/11, 1/1, 1/1, 3/floatmode, 6/registerW, 2/00, 3/001, 1/0, 4/0100, 1/--, 1/1, 6/registerZ }


\paragraph{Syntax} \texttt{fixedwp}\emph{floatmode} \emph{registerW} = \emph{registerZ}

\paragraph{Description} The \emph{registerZ} is interpreted as two binary 32 floating-point numbers packed into 64 bits. The \emph{registerZ} words are converted to 32-bit signed integers according to the rounding mode. The two resulting words are packed and stored into the \emph{registerW}. This instruction may raise exception bits in the CS register.


\paragraph{Modifier(s)}
\noindent\begin{itemize}
\item floatmode (cf. floatmode in section \ref{par:modifier-floatmode} on page \pageref{par:modifier-floatmode})
\end{itemize}

\paragraph{Execution} {Reservation table \textsf{ALU\_LITE} (Section~\ref{sec:reservation-ALU-LITE})}
~
{\begin{lstlisting}
stage ID:
  new floatmode = _ZX_3(floatmode);
stage RR:
  new argument2 = GPR[registerZ];
  new RM = decode_riscv_float_rounding_mode(floatmode);
  for (I = 0; I < 2; I += 1) {
    result1.32[I] = f32_to_i32(RM, argument2.32[I])
  };
stage E4:
  GPR[registerW] = result1;
  commit_float_exception_flags();
\end{lstlisting}}
\newpage
\subsubsection[{\small FIXEDWQ: Floating Point Conversion to Fixed-Point Word Quadruple}]{FIXEDWQ\hfill Floating Point Conversion to Fixed-Point Word Quadruple}
\label{instr:FIXEDWQ}\index{fixedwq}
 
\paragraph{Encoding} Format \textsf{ALU\_FWQFIXED} (Section~\ref{sec:format-ALU-FWQFIXED}), \verb|fpucode6 = 0100|\\

\bitfields{ 1/P, 2/11, 1/1, 1/1, 3/floatmode, 5/registerM, 1/--, 2/11, 3/001, 1/1, 4/0100, 1/--, 1/0, 5/registerP, 1/0 }


\paragraph{Syntax} \texttt{fixedwq}\emph{floatmode} \emph{registerM} = \emph{registerP}

\paragraph{Description} The \emph{registerP} is interpreted as four binary 32 floating-point numbers packed into 128 bits. The \emph{registerP} words are converted to 32-bit signed integers according to the rounding mode. The four resulting words are packed and stored into the \emph{registerM}. This instruction may raise exception bits in the CS register.


\paragraph{Modifier(s)}
\noindent\begin{itemize}
\item floatmode (cf. floatmode in section \ref{par:modifier-floatmode} on page \pageref{par:modifier-floatmode})
\end{itemize}

\paragraph{Execution} {Reservation table \textsf{ALU\_LITE} (Section~\ref{sec:reservation-ALU-LITE})}
~
{\begin{lstlisting}
stage ID:
  new floatmode = _ZX_3(floatmode);
stage RR:
  new argument2 = PGR[registerP];
  new RM = decode_riscv_float_rounding_mode(floatmode);
  for (I = 0; I < 4; I += 1) {
    result1.32[I] = f32_to_i32(RM, argument2.32[I])
  };
stage E4:
  PGR[registerM] = result1;
  commit_float_exception_flags();
\end{lstlisting}}
\newpage
\subsubsection[{\small FLOATD: Floating Point Conversion from Fixed-Point Double Word}]{FLOATD\hfill Floating Point Conversion from Fixed-Point Double Word}
\label{instr:FLOATD}\index{floatd}
 
\paragraph{Encoding} Format \textsf{ALU\_FDFLOAT} (Section~\ref{sec:format-ALU-FDFLOAT}), \verb|fpucode6 = 0000|\\

\bitfields{ 1/P, 2/11, 1/1, 1/1, 3/floatmode, 6/registerW, 2/00, 3/001, 1/0, 4/0000, 1/--, 1/0, 6/registerZ }


\paragraph{Syntax} \texttt{floatd}\emph{floatmode} \emph{registerW} = \emph{registerZ}

\paragraph{Description} The \emph{registerZ} interpreted as a 64-bit signed integer is converted to a 64-bit floating-point number according to the rounding mode. This instruction may raise exception bits in the CS register.


\paragraph{Modifier(s)}
\noindent\begin{itemize}
\item floatmode (cf. floatmode in section \ref{par:modifier-floatmode} on page \pageref{par:modifier-floatmode})
\end{itemize}

\paragraph{Execution} {Reservation table \textsf{ALU\_LITE} (Section~\ref{sec:reservation-ALU-LITE})}
~
{\begin{lstlisting}
stage ID:
  new floatmode = _ZX_3(floatmode);
stage RR:
  new argument2 = GPR[registerZ];
  new RM = decode_riscv_float_rounding_mode(floatmode);
  new result1 = i64_to_f64(RM, argument2);
stage E4:
  GPR[registerW] = result1;
  commit_float_exception_flags();
\end{lstlisting}}
\newpage
\subsubsection[{\small FLOATDP: Floating Point Conversion from Fixed-Point Double Word Pair}]{FLOATDP\hfill Floating Point Conversion from Fixed-Point Double Word Pair}
\label{instr:FLOATDP}\index{floatdp}
 
\paragraph{Encoding} Format \textsf{ALU\_FDPFLOAT} (Section~\ref{sec:format-ALU-FDPFLOAT}), \verb|fpucode6 = 0000|\\

\bitfields{ 1/P, 2/11, 1/1, 1/1, 3/floatmode, 5/registerM, 1/--, 2/11, 3/001, 1/0, 4/0000, 1/--, 1/0, 5/registerP, 1/1 }


\paragraph{Syntax} \texttt{floatdp}\emph{floatmode} \emph{registerM} = \emph{registerP}

\paragraph{Description} The \emph{registerP} is interpreted as two 64-bit signed integers packed into 128 bits. The \emph{registerP} double words are converted to binary 64 floating-point numbers according to the rounding mode. The two resulting double words are packed and stored into the \emph{registerM}. This instruction may raise exception bits in the CS register.


\paragraph{Modifier(s)}
\noindent\begin{itemize}
\item floatmode (cf. floatmode in section \ref{par:modifier-floatmode} on page \pageref{par:modifier-floatmode})
\end{itemize}

\paragraph{Execution} {Reservation table \textsf{ALU\_LITE} (Section~\ref{sec:reservation-ALU-LITE})}
~
{\begin{lstlisting}
stage ID:
  new floatmode = _ZX_3(floatmode);
stage RR:
  new argument2 = PGR[registerP];
  new RM = decode_riscv_float_rounding_mode(floatmode);
  for (I = 0; I < 2; I += 1) {
    result1.64[I] = i64_to_f64(RM, argument2.64[I])
  };
stage E4:
  PGR[registerM] = result1;
  commit_float_exception_flags();
\end{lstlisting}}
\newpage
\subsubsection[{\small FLOATDW: Floating Point Conversion from Fixed-Point Double Word to Word}]{FLOATDW\hfill Floating Point Conversion from Fixed-Point Double Word to Word}
\label{instr:FLOATDW}\index{floatdw}
 
\paragraph{Encoding} Format \textsf{ALU\_FWFLOAT} (Section~\ref{sec:format-ALU-FWFLOAT}), \verb|fpucode6 = 0010|\\

\bitfields{ 1/P, 2/11, 1/1, 1/1, 3/floatmode, 6/registerW, 2/01, 3/001, 1/0, 4/0010, 1/--, 1/0, 6/registerZ }


\paragraph{Syntax} \texttt{floatdw}\emph{floatmode} \emph{registerW} = \emph{registerZ}

\paragraph{Description} The \emph{registerZ} interpreted as a 64-bit signed integer is converted to a 32-bit floating-point number according to the rounding mode. This instruction may raise exception bits in the CS register.


\paragraph{Modifier(s)}
\noindent\begin{itemize}
\item floatmode (cf. floatmode in section \ref{par:modifier-floatmode} on page \pageref{par:modifier-floatmode})
\end{itemize}

\paragraph{Execution} {Reservation table \textsf{ALU\_LITE} (Section~\ref{sec:reservation-ALU-LITE})}
~
{\begin{lstlisting}
stage ID:
  new floatmode = _ZX_3(floatmode);
stage RR:
  new argument2 = GPR[registerZ];
  new RM = decode_riscv_float_rounding_mode(floatmode);
  new result1 = i64_to_f32(RM, argument2);
stage E4:
  GPR[registerW] = result1;
  commit_float_exception_flags();
\end{lstlisting}}
\newpage
\subsubsection[{\small FLOATHO: Floating Point Conversion from Fixed-Point Half Word Octuple}]{FLOATHO\hfill Floating Point Conversion from Fixed-Point Half Word Octuple}
\label{instr:FLOATHO}\index{floatho}
 
\paragraph{Encoding} Format \textsf{ALU\_FHOFLOAT} (Section~\ref{sec:format-ALU-FHOFLOAT}), \verb|fpucode6 = 0000|\\

\bitfields{ 1/P, 2/11, 1/1, 1/1, 3/floatmode, 5/registerM, 1/--, 2/11, 3/001, 1/1, 4/0000, 1/--, 1/0, 5/registerP, 1/1 }


\paragraph{Syntax} \texttt{floatho}\emph{floatmode} \emph{registerM} = \emph{registerP}

\paragraph{Description} The \emph{registerP} is interpreted as eight 16-bit signed integers packed into 128 bits. The \emph{registerP} half words are converted to binary 16 floating-point numbers according to the rounding mode. The eight resulting half words are packed and stored into the \emph{registerM}. This instruction may raise exception bits in the CS register.


\paragraph{Modifier(s)}
\noindent\begin{itemize}
\item floatmode (cf. floatmode in section \ref{par:modifier-floatmode} on page \pageref{par:modifier-floatmode})
\end{itemize}

\paragraph{Execution} {Reservation table \textsf{ALU\_LITE} (Section~\ref{sec:reservation-ALU-LITE})}
~
{\begin{lstlisting}
stage ID:
  new floatmode = _ZX_3(floatmode);
stage RR:
  new argument2 = PGR[registerP];
  new RM = decode_riscv_float_rounding_mode(floatmode);
  for (I = 0; I < 8; I += 1) {
    result1.16[I] = i16_to_f16(RM, argument2.16[I])
  };
stage E4:
  PGR[registerM] = result1;
  commit_float_exception_flags();
\end{lstlisting}}
\newpage
\subsubsection[{\small FLOATUD: Floating Point Conversion from Unsigned Fixed-Point Double Word}]{FLOATUD\hfill Floating Point Conversion from Unsigned Fixed-Point Double Word}
\label{instr:FLOATUD}\index{floatud}
 
\paragraph{Encoding} Format \textsf{ALU\_FDFLOAT} (Section~\ref{sec:format-ALU-FDFLOAT}), \verb|fpucode6 = 0001|\\

\bitfields{ 1/P, 2/11, 1/1, 1/1, 3/floatmode, 6/registerW, 2/00, 3/001, 1/0, 4/0001, 1/--, 1/0, 6/registerZ }


\paragraph{Syntax} \texttt{floatud}\emph{floatmode} \emph{registerW} = \emph{registerZ}

\paragraph{Description} The \emph{registerZ} interpreted as a 64-bit unsigned integer is converted to a 64-bit floating-point number according to the rounding mode. This instruction may raise exception bits in the CS register.


\paragraph{Modifier(s)}
\noindent\begin{itemize}
\item floatmode (cf. floatmode in section \ref{par:modifier-floatmode} on page \pageref{par:modifier-floatmode})
\end{itemize}

\paragraph{Execution} {Reservation table \textsf{ALU\_LITE} (Section~\ref{sec:reservation-ALU-LITE})}
~
{\begin{lstlisting}
stage ID:
  new floatmode = _ZX_3(floatmode);
stage RR:
  new argument2 = GPR[registerZ];
  new RM = decode_riscv_float_rounding_mode(floatmode);
  new result1 = ui64_to_f64(RM, argument2);
stage E4:
  GPR[registerW] = result1;
  commit_float_exception_flags();
\end{lstlisting}}
\newpage
\subsubsection[{\small FLOATUDP: Floating Point Conversion from Unsigned Fixed-Point Double Word Pair}]{FLOATUDP\hfill Floating Point Conversion from Unsigned Fixed-Point Double Word Pair}
\label{instr:FLOATUDP}\index{floatudp}
 
\paragraph{Encoding} Format \textsf{ALU\_FDPFLOAT} (Section~\ref{sec:format-ALU-FDPFLOAT}), \verb|fpucode6 = 0001|\\

\bitfields{ 1/P, 2/11, 1/1, 1/1, 3/floatmode, 5/registerM, 1/--, 2/11, 3/001, 1/0, 4/0001, 1/--, 1/0, 5/registerP, 1/1 }


\paragraph{Syntax} \texttt{floatudp}\emph{floatmode} \emph{registerM} = \emph{registerP}

\paragraph{Description} The \emph{registerP} is interpreted as two 64-bit unsigned integers packed into 128 bits. The \emph{registerP} double words are converted to binary 64 floating-point numbers according to the rounding mode. The two resulting double words are packed and stored into the \emph{registerM}. This instruction may raise exception bits in the CS register.


\paragraph{Modifier(s)}
\noindent\begin{itemize}
\item floatmode (cf. floatmode in section \ref{par:modifier-floatmode} on page \pageref{par:modifier-floatmode})
\end{itemize}

\paragraph{Execution} {Reservation table \textsf{ALU\_LITE} (Section~\ref{sec:reservation-ALU-LITE})}
~
{\begin{lstlisting}
stage ID:
  new floatmode = _ZX_3(floatmode);
stage RR:
  new argument2 = PGR[registerP];
  new RM = decode_riscv_float_rounding_mode(floatmode);
  for (I = 0; I < 2; I += 1) {
    result1.64[I] = ui64_to_f64(RM, argument2.64[I])
  };
stage E4:
  PGR[registerM] = result1;
  commit_float_exception_flags();
\end{lstlisting}}
\newpage
\subsubsection[{\small FLOATUDW: Floating Point Conversion from Unsigned Fixed-Point Double Word to Word}]{FLOATUDW\hfill Floating Point Conversion from Unsigned Fixed-Point Double Word to Word}
\label{instr:FLOATUDW}\index{floatudw}
 
\paragraph{Encoding} Format \textsf{ALU\_FWFLOAT} (Section~\ref{sec:format-ALU-FWFLOAT}), \verb|fpucode6 = 0011|\\

\bitfields{ 1/P, 2/11, 1/1, 1/1, 3/floatmode, 6/registerW, 2/01, 3/001, 1/0, 4/0011, 1/--, 1/0, 6/registerZ }


\paragraph{Syntax} \texttt{floatudw}\emph{floatmode} \emph{registerW} = \emph{registerZ}

\paragraph{Description} The \emph{registerZ} interpreted as a 64-bit unsigned integer is converted to a 32-bit floating-point number according to the rounding mode. This instruction may raise exception bits in the CS register.


\paragraph{Modifier(s)}
\noindent\begin{itemize}
\item floatmode (cf. floatmode in section \ref{par:modifier-floatmode} on page \pageref{par:modifier-floatmode})
\end{itemize}

\paragraph{Execution} {Reservation table \textsf{ALU\_LITE} (Section~\ref{sec:reservation-ALU-LITE})}
~
{\begin{lstlisting}
stage ID:
  new floatmode = _ZX_3(floatmode);
stage RR:
  new argument2 = GPR[registerZ];
  new RM = decode_riscv_float_rounding_mode(floatmode);
  new result1 = ui64_to_f32(RM, argument2);
stage E4:
  GPR[registerW] = result1;
  commit_float_exception_flags();
\end{lstlisting}}
\newpage
\subsubsection[{\small FLOATUHO: Floating Point Conversion from Unsigned Fixed-Point Half Word Octuple}]{FLOATUHO\hfill Floating Point Conversion from Unsigned Fixed-Point Half Word Octuple}
\label{instr:FLOATUHO}\index{floatuho}
 
\paragraph{Encoding} Format \textsf{ALU\_FHOFLOAT} (Section~\ref{sec:format-ALU-FHOFLOAT}), \verb|fpucode6 = 0001|\\

\bitfields{ 1/P, 2/11, 1/1, 1/1, 3/floatmode, 5/registerM, 1/--, 2/11, 3/001, 1/1, 4/0001, 1/--, 1/0, 5/registerP, 1/1 }


\paragraph{Syntax} \texttt{floatuho}\emph{floatmode} \emph{registerM} = \emph{registerP}

\paragraph{Description} The \emph{registerP} is interpreted as eight 16-bit unsigned integers packed into 128 bits. The \emph{registerP} half words are converted to binary 16 floating-point numbers according to the rounding mode. The eight resulting half words are packed and stored into the \emph{registerM}. This instruction may raise exception bits in the CS register.


\paragraph{Modifier(s)}
\noindent\begin{itemize}
\item floatmode (cf. floatmode in section \ref{par:modifier-floatmode} on page \pageref{par:modifier-floatmode})
\end{itemize}

\paragraph{Execution} {Reservation table \textsf{ALU\_LITE} (Section~\ref{sec:reservation-ALU-LITE})}
~
{\begin{lstlisting}
stage ID:
  new floatmode = _ZX_3(floatmode);
stage RR:
  new argument2 = PGR[registerP];
  new RM = decode_riscv_float_rounding_mode(floatmode);
  for (I = 0; I < 8; I += 1) {
    result1.16[I] = ui16_to_f16(RM, argument2.16[I])
  };
stage E4:
  PGR[registerM] = result1;
  commit_float_exception_flags();
\end{lstlisting}}
\newpage
\subsubsection[{\small FLOATUW: Floating Point Conversion from Unsigned Fixed-Point}]{FLOATUW\hfill Floating Point Conversion from Unsigned Fixed-Point}
\label{instr:FLOATUW}\index{floatuw}
 
\paragraph{Encoding} Format \textsf{ALU\_FWFLOAT} (Section~\ref{sec:format-ALU-FWFLOAT}), \verb|fpucode6 = 0001|\\

\bitfields{ 1/P, 2/11, 1/1, 1/1, 3/floatmode, 6/registerW, 2/01, 3/001, 1/0, 4/0001, 1/--, 1/0, 6/registerZ }


\paragraph{Syntax} \texttt{floatuw}\emph{floatmode} \emph{registerW} = \emph{registerZ}

\paragraph{Description} The \emph{registerZ} interpreted as a 32-bit unsigned integer is converted to a 32-bit floating-point number according to the rounding mode. This instruction may raise exception bits in the CS register.


\paragraph{Modifier(s)}
\noindent\begin{itemize}
\item floatmode (cf. floatmode in section \ref{par:modifier-floatmode} on page \pageref{par:modifier-floatmode})
\end{itemize}

\paragraph{Execution} {Reservation table \textsf{ALU\_LITE} (Section~\ref{sec:reservation-ALU-LITE})}
~
{\begin{lstlisting}
stage ID:
  new floatmode = _ZX_3(floatmode);
stage RR:
  new argument2 = GPR[registerZ];
  new RM = decode_riscv_float_rounding_mode(floatmode);
  new result1 = ui32_to_f32(RM, argument2);
stage E4:
  GPR[registerW] = result1;
  commit_float_exception_flags();
\end{lstlisting}}
\newpage
\subsubsection[{\small FLOATUWD: Floating Point Conversion from Unsigned Fixed-Point Word to Double Word}]{FLOATUWD\hfill Floating Point Conversion from Unsigned Fixed-Point Word to Double Word}
\label{instr:FLOATUWD}\index{floatuwd}
 
\paragraph{Encoding} Format \textsf{ALU\_FDFLOAT} (Section~\ref{sec:format-ALU-FDFLOAT}), \verb|fpucode6 = 0011|\\

\bitfields{ 1/P, 2/11, 1/1, 1/1, 3/floatmode, 6/registerW, 2/00, 3/001, 1/0, 4/0011, 1/--, 1/0, 6/registerZ }


\paragraph{Syntax} \texttt{floatuwd}\emph{floatmode} \emph{registerW} = \emph{registerZ}

\paragraph{Description} The \emph{registerZ} interpreted as a 32-bit unsigned integer is converted to a 64-bit floating-point number according to the rounding mode. This instruction may raise exception bits in the CS register.


\paragraph{Modifier(s)}
\noindent\begin{itemize}
\item floatmode (cf. floatmode in section \ref{par:modifier-floatmode} on page \pageref{par:modifier-floatmode})
\end{itemize}

\paragraph{Execution} {Reservation table \textsf{ALU\_LITE} (Section~\ref{sec:reservation-ALU-LITE})}
~
{\begin{lstlisting}
stage ID:
  new floatmode = _ZX_3(floatmode);
stage RR:
  new argument2 = GPR[registerZ];
  new RM = decode_riscv_float_rounding_mode(floatmode);
  new result1 = ui32_to_f64(RM, argument2);
stage E4:
  GPR[registerW] = result1;
  commit_float_exception_flags();
\end{lstlisting}}
\newpage
\subsubsection[{\small FLOATUWP: Floating Point Conversion from Unsigned Fixed-Point Word Pair}]{FLOATUWP\hfill Floating Point Conversion from Unsigned Fixed-Point Word Pair}
\label{instr:FLOATUWP}\index{floatuwp}
 
\paragraph{Encoding} Format \textsf{ALU\_FWPFLOAT} (Section~\ref{sec:format-ALU-FWPFLOAT}), \verb|fpucode6 = 0001|\\

\bitfields{ 1/P, 2/11, 1/1, 1/1, 3/floatmode, 6/registerW, 2/00, 3/001, 1/0, 4/0001, 1/--, 1/1, 6/registerZ }


\paragraph{Syntax} \texttt{floatuwp}\emph{floatmode} \emph{registerW} = \emph{registerZ}

\paragraph{Description} The \emph{registerZ} is interpreted as two 32-bit unsigned integers packed into 64 bits. The \emph{registerZ} words are converted to binary 32 floating-point numbers according to the rounding mode. The two resulting words are packed and stored into the \emph{registerW}. This instruction may raise exception bits in the CS register.


\paragraph{Modifier(s)}
\noindent\begin{itemize}
\item floatmode (cf. floatmode in section \ref{par:modifier-floatmode} on page \pageref{par:modifier-floatmode})
\end{itemize}

\paragraph{Execution} {Reservation table \textsf{ALU\_LITE} (Section~\ref{sec:reservation-ALU-LITE})}
~
{\begin{lstlisting}
stage ID:
  new floatmode = _ZX_3(floatmode);
stage RR:
  new argument2 = GPR[registerZ];
  new RM = decode_riscv_float_rounding_mode(floatmode);
  for (I = 0; I < 2; I += 1) {
    result1.32[I] = ui32_to_f32(RM, argument2.32[I])
  };
stage E4:
  GPR[registerW] = result1;
  commit_float_exception_flags();
\end{lstlisting}}
\newpage
\subsubsection[{\small FLOATUWQ: Floating Point Conversion from Unsigned Fixed-Point Word Quadruple}]{FLOATUWQ\hfill Floating Point Conversion from Unsigned Fixed-Point Word Quadruple}
\label{instr:FLOATUWQ}\index{floatuwq}
 
\paragraph{Encoding} Format \textsf{ALU\_FWQFLOAT} (Section~\ref{sec:format-ALU-FWQFLOAT}), \verb|fpucode6 = 0001|\\

\bitfields{ 1/P, 2/11, 1/1, 1/1, 3/floatmode, 5/registerM, 1/--, 2/11, 3/001, 1/1, 4/0001, 1/--, 1/0, 5/registerP, 1/0 }


\paragraph{Syntax} \texttt{floatuwq}\emph{floatmode} \emph{registerM} = \emph{registerP}

\paragraph{Description} The \emph{registerP} is interpreted as four 32-bit unsigned integers packed into 128 bits. The \emph{registerP} words are converted to binary 32 floating-point numbers according to the rounding mode. The four resulting words are packed and stored into the \emph{registerM}. This instruction may raise exception bits in the CS register.


\paragraph{Modifier(s)}
\noindent\begin{itemize}
\item floatmode (cf. floatmode in section \ref{par:modifier-floatmode} on page \pageref{par:modifier-floatmode})
\end{itemize}

\paragraph{Execution} {Reservation table \textsf{ALU\_LITE} (Section~\ref{sec:reservation-ALU-LITE})}
~
{\begin{lstlisting}
stage ID:
  new floatmode = _ZX_3(floatmode);
stage RR:
  new argument2 = PGR[registerP];
  new RM = decode_riscv_float_rounding_mode(floatmode);
  for (I = 0; I < 4; I += 1) {
    result1.32[I] = ui32_to_f32(RM, argument2.32[I])
  };
stage E4:
  PGR[registerM] = result1;
  commit_float_exception_flags();
\end{lstlisting}}
\newpage
\subsubsection[{\small FLOATW: Floating Point Conversion from Fixed-Point}]{FLOATW\hfill Floating Point Conversion from Fixed-Point}
\label{instr:FLOATW}\index{floatw}
 
\paragraph{Encoding} Format \textsf{ALU\_FWFLOAT} (Section~\ref{sec:format-ALU-FWFLOAT}), \verb|fpucode6 = 0000|\\

\bitfields{ 1/P, 2/11, 1/1, 1/1, 3/floatmode, 6/registerW, 2/01, 3/001, 1/0, 4/0000, 1/--, 1/0, 6/registerZ }


\paragraph{Syntax} \texttt{floatw}\emph{floatmode} \emph{registerW} = \emph{registerZ}

\paragraph{Description} The \emph{registerZ} interpreted as a 32-bit signed integer is converted to a 32-bit floating-point number according to the rounding mode. This instruction may raise exception bits in the CS register.


\paragraph{Modifier(s)}
\noindent\begin{itemize}
\item floatmode (cf. floatmode in section \ref{par:modifier-floatmode} on page \pageref{par:modifier-floatmode})
\end{itemize}

\paragraph{Execution} {Reservation table \textsf{ALU\_LITE} (Section~\ref{sec:reservation-ALU-LITE})}
~
{\begin{lstlisting}
stage ID:
  new floatmode = _ZX_3(floatmode);
stage RR:
  new argument2 = GPR[registerZ];
  new RM = decode_riscv_float_rounding_mode(floatmode);
  new result1 = i32_to_f32(RM, argument2);
stage E4:
  GPR[registerW] = result1;
  commit_float_exception_flags();
\end{lstlisting}}
\newpage
\subsubsection[{\small FLOATWD: Floating Point Conversion from Fixed-Point Word to Double Word}]{FLOATWD\hfill Floating Point Conversion from Fixed-Point Word to Double Word}
\label{instr:FLOATWD}\index{floatwd}
 
\paragraph{Encoding} Format \textsf{ALU\_FDFLOAT} (Section~\ref{sec:format-ALU-FDFLOAT}), \verb|fpucode6 = 0010|\\

\bitfields{ 1/P, 2/11, 1/1, 1/1, 3/floatmode, 6/registerW, 2/00, 3/001, 1/0, 4/0010, 1/--, 1/0, 6/registerZ }


\paragraph{Syntax} \texttt{floatwd}\emph{floatmode} \emph{registerW} = \emph{registerZ}

\paragraph{Description} The \emph{registerZ} interpreted as a 32-bit signed integer is converted to a 64-bit floating-point number according to the rounding mode. This instruction may raise exception bits in the CS register.


\paragraph{Modifier(s)}
\noindent\begin{itemize}
\item floatmode (cf. floatmode in section \ref{par:modifier-floatmode} on page \pageref{par:modifier-floatmode})
\end{itemize}

\paragraph{Execution} {Reservation table \textsf{ALU\_LITE} (Section~\ref{sec:reservation-ALU-LITE})}
~
{\begin{lstlisting}
stage ID:
  new floatmode = _ZX_3(floatmode);
stage RR:
  new argument2 = GPR[registerZ];
  new RM = decode_riscv_float_rounding_mode(floatmode);
  new result1 = i32_to_f64(RM, argument2);
stage E4:
  GPR[registerW] = result1;
  commit_float_exception_flags();
\end{lstlisting}}
\newpage
\subsubsection[{\small FLOATWP: Floating Point Conversion from Fixed-Point Word Pair}]{FLOATWP\hfill Floating Point Conversion from Fixed-Point Word Pair}
\label{instr:FLOATWP}\index{floatwp}
 
\paragraph{Encoding} Format \textsf{ALU\_FWPFLOAT} (Section~\ref{sec:format-ALU-FWPFLOAT}), \verb|fpucode6 = 0000|\\

\bitfields{ 1/P, 2/11, 1/1, 1/1, 3/floatmode, 6/registerW, 2/00, 3/001, 1/0, 4/0000, 1/--, 1/1, 6/registerZ }


\paragraph{Syntax} \texttt{floatwp}\emph{floatmode} \emph{registerW} = \emph{registerZ}

\paragraph{Description} The \emph{registerZ} is interpreted as two 32-bit signed integers packed into 64 bits. The \emph{registerZ} words are converted to binary 32 floating-point numbers according to the rounding mode. The two resulting words are packed and stored into the \emph{registerW}. This instruction may raise exception bits in the CS register.


\paragraph{Modifier(s)}
\noindent\begin{itemize}
\item floatmode (cf. floatmode in section \ref{par:modifier-floatmode} on page \pageref{par:modifier-floatmode})
\end{itemize}

\paragraph{Execution} {Reservation table \textsf{ALU\_LITE} (Section~\ref{sec:reservation-ALU-LITE})}
~
{\begin{lstlisting}
stage ID:
  new floatmode = _ZX_3(floatmode);
stage RR:
  new argument2 = GPR[registerZ];
  new RM = decode_riscv_float_rounding_mode(floatmode);
  for (I = 0; I < 2; I += 1) {
    result1.32[I] = i32_to_f32(RM, argument2.32[I])
  };
stage E4:
  GPR[registerW] = result1;
  commit_float_exception_flags();
\end{lstlisting}}
\newpage
\subsubsection[{\small FLOATWQ: Floating Point Conversion from Fixed-Point Word Quadruple}]{FLOATWQ\hfill Floating Point Conversion from Fixed-Point Word Quadruple}
\label{instr:FLOATWQ}\index{floatwq}
 
\paragraph{Encoding} Format \textsf{ALU\_FWQFLOAT} (Section~\ref{sec:format-ALU-FWQFLOAT}), \verb|fpucode6 = 0000|\\

\bitfields{ 1/P, 2/11, 1/1, 1/1, 3/floatmode, 5/registerM, 1/--, 2/11, 3/001, 1/1, 4/0000, 1/--, 1/0, 5/registerP, 1/0 }


\paragraph{Syntax} \texttt{floatwq}\emph{floatmode} \emph{registerM} = \emph{registerP}

\paragraph{Description} The \emph{registerP} is interpreted as four 32-bit signed integers packed into 128 bits. The \emph{registerP} words are converted to binary 32 floating-point numbers according to the rounding mode. The four resulting words are packed and stored into the \emph{registerM}. This instruction may raise exception bits in the CS register.


\paragraph{Modifier(s)}
\noindent\begin{itemize}
\item floatmode (cf. floatmode in section \ref{par:modifier-floatmode} on page \pageref{par:modifier-floatmode})
\end{itemize}

\paragraph{Execution} {Reservation table \textsf{ALU\_LITE} (Section~\ref{sec:reservation-ALU-LITE})}
~
{\begin{lstlisting}
stage ID:
  new floatmode = _ZX_3(floatmode);
stage RR:
  new argument2 = PGR[registerP];
  new RM = decode_riscv_float_rounding_mode(floatmode);
  for (I = 0; I < 4; I += 1) {
    result1.32[I] = i32_to_f32(RM, argument2.32[I])
  };
stage E4:
  PGR[registerM] = result1;
  commit_float_exception_flags();
\end{lstlisting}}
\newpage
\subsubsection[{\small FMAXD: Floating-Point Maximum Double Word}]{FMAXD\hfill Floating-Point Maximum Double Word}
\label{instr:FMAXD}\index{fmaxd}
 
\paragraph{Encoding} Format \textsf{ALU\_FDMMWR} (Section~\ref{sec:format-ALU-FDMMWR}), \verb|floatmode = 001|\\

\bitfields{ 1/P, 2/11, 1/1, 1/1, 3/001, 6/registerW, 2/00, 3/000, 1/0, 6/registerY, 6/registerZ }


\paragraph{Syntax} \texttt{fmaxd} \emph{registerW} = \emph{registerZ}, \emph{registerY}

\paragraph{Description} The \emph{registerZ} is compared to the \emph{registerY}, assuming binary 64 floating-point numbers. The maximum value is stored into the \emph{registerW}.


\paragraph{Execution} {Reservation table \textsf{ALU\_LITE} (Section~\ref{sec:reservation-ALU-LITE})}
~
{\begin{lstlisting}
stage RR:
  new argument3 = GPR[registerY];
  new argument2 = GPR[registerZ];
  new result1 = f64_max(argument2, argument3);
stage E1:
  GPR[registerW] = result1;
\end{lstlisting}}
\newpage
\subsubsection[{\small FMAXDP: Floating-Point Maximum Double Word Pair}]{FMAXDP\hfill Floating-Point Maximum Double Word Pair}
\label{instr:FMAXDP}\index{fmaxdp}
 
\paragraph{Encoding} Format \textsf{ALU\_FDPMMWR} (Section~\ref{sec:format-ALU-FDPMMWR}), \verb|floatmode = 001|\\

\bitfields{ 1/P, 2/11, 1/1, 1/1, 3/001, 5/registerM, 1/--, 2/11, 3/000, 1/0, 5/registerO, 1/0, 5/registerP, 1/1 }


\paragraph{Syntax} \texttt{fmaxdp} \emph{registerM} = \emph{registerP}, \emph{registerO}

\paragraph{Description} The \emph{registerP} and the \emph{registerO} are interpreted as two binary 64 floating-point numbers packed into 128 bits. The \emph{registerP} double words are compared to the \emph{registerO} double words. The two maximum double words are packed and stored into the \emph{registerM}.


\paragraph{Execution} {Reservation table \textsf{ALU\_LITE} (Section~\ref{sec:reservation-ALU-LITE})}
~
{\begin{lstlisting}
stage RR:
  new argument3 = PGR[registerO];
  new argument2 = PGR[registerP];
  for (I = 0; I < 2; I += 1) {
    result1.64[I] = f64_max(argument2.64[I], argument3.64[I])
  };
stage E1:
  PGR[registerM] = result1;
\end{lstlisting}}
\newpage
\subsubsection[{\small FMAXH: Floating-Point Maximum Half Word}]{FMAXH\hfill Floating-Point Maximum Half Word}
\label{instr:FMAXH}\index{fmaxh}
 
\paragraph{Encoding} Format \textsf{ALU\_FHMMWR} (Section~\ref{sec:format-ALU-FHMMWR}), \verb|floatmode = 001|\\

\bitfields{ 1/P, 2/11, 1/1, 1/1, 3/001, 6/registerW, 2/01, 3/011, 1/1, 6/registerY, 6/registerZ }


\paragraph{Syntax} \texttt{fmaxh} \emph{registerW} = \emph{registerZ}, \emph{registerY}

\paragraph{Description} The \emph{registerZ} is compared to the \emph{registerY}, assuming binary 16 floating-point numbers. The maximum value is stored into the \emph{registerW}.


\paragraph{Execution} {Reservation table \textsf{ALU\_LITE} (Section~\ref{sec:reservation-ALU-LITE})}
~
{\begin{lstlisting}
stage RR:
  new argument3 = GPR[registerY];
  new argument2 = GPR[registerZ];
  new result1 = f16_max(argument2, argument3);
stage E1:
  GPR[registerW] = result1;
\end{lstlisting}}
\newpage
\subsubsection[{\small FMAXHO: Floating-Point Maximum Half Word Octuple}]{FMAXHO\hfill Floating-Point Maximum Half Word Octuple}
\label{instr:FMAXHO}\index{fmaxho}
 
\paragraph{Encoding} Format \textsf{ALU\_FHOMMWR} (Section~\ref{sec:format-ALU-FHOMMWR}), \verb|floatmode = 001|\\

\bitfields{ 1/P, 2/11, 1/1, 1/1, 3/001, 5/registerM, 1/--, 2/11, 3/000, 1/1, 5/registerO, 1/0, 5/registerP, 1/1 }


\paragraph{Syntax} \texttt{fmaxho} \emph{registerM} = \emph{registerP}, \emph{registerO}

\paragraph{Description} The \emph{registerP} and the \emph{registerO} are interpreted as eight binary 16 floating-point numbers packed into 128 bits. The \emph{registerP} half words are compared to the \emph{registerO} half words. The eight maximum half words are packed and stored into the \emph{registerM}.


\paragraph{Execution} {Reservation table \textsf{ALU\_LITE} (Section~\ref{sec:reservation-ALU-LITE})}
~
{\begin{lstlisting}
stage RR:
  new argument3 = PGR[registerO];
  new argument2 = PGR[registerP];
  for (I = 0; I < 8; I += 1) {
    result1.16[I] = f16_max(argument2.16[I], argument3.16[I])
  };
stage E1:
  PGR[registerM] = result1;
\end{lstlisting}}
\newpage
\subsubsection[{\small FMAXND: Floating-Point Maximum Number Double Word}]{FMAXND\hfill Floating-Point Maximum Number Double Word}
\label{instr:FMAXND}\index{fmaxnd}
 
\paragraph{Encoding} Format \textsf{ALU\_FDMMWR} (Section~\ref{sec:format-ALU-FDMMWR}), \verb|floatmode = 011|\\

\bitfields{ 1/P, 2/11, 1/1, 1/1, 3/011, 6/registerW, 2/00, 3/000, 1/0, 6/registerY, 6/registerZ }


\paragraph{Syntax} \texttt{fmaxnd} \emph{registerW} = \emph{registerZ}, \emph{registerY}

\paragraph{Description} The \emph{registerZ} is compared to the \emph{registerY}, assuming binary 64 floating-point numbers. The maximum number value is stored into the \emph{registerW}.


\paragraph{Execution} {Reservation table \textsf{ALU\_LITE} (Section~\ref{sec:reservation-ALU-LITE})}
~
{\begin{lstlisting}
stage RR:
  new argument3 = GPR[registerY];
  new argument2 = GPR[registerZ];
  new result1 = f64_maxNum(argument2, argument3);
stage E1:
  GPR[registerW] = result1;
\end{lstlisting}}
\newpage
\subsubsection[{\small FMAXNDP: Floating-Point Maximum Number Double Word Pair}]{FMAXNDP\hfill Floating-Point Maximum Number Double Word Pair}
\label{instr:FMAXNDP}\index{fmaxndp}
 
\paragraph{Encoding} Format \textsf{ALU\_FDPMMWR} (Section~\ref{sec:format-ALU-FDPMMWR}), \verb|floatmode = 011|\\

\bitfields{ 1/P, 2/11, 1/1, 1/1, 3/011, 5/registerM, 1/--, 2/11, 3/000, 1/0, 5/registerO, 1/0, 5/registerP, 1/1 }


\paragraph{Syntax} \texttt{fmaxndp} \emph{registerM} = \emph{registerP}, \emph{registerO}

\paragraph{Description} The \emph{registerP} and the \emph{registerO} are interpreted as two binary 64 floating-point numbers packed into 128 bits. The \emph{registerP} double words are compared to the \emph{registerO} double words. The two maximum number double words are packed and stored into the \emph{registerM}.


\paragraph{Execution} {Reservation table \textsf{ALU\_LITE} (Section~\ref{sec:reservation-ALU-LITE})}
~
{\begin{lstlisting}
stage RR:
  new argument3 = PGR[registerO];
  new argument2 = PGR[registerP];
  for (I = 0; I < 2; I += 1) {
    result1.64[I] = f64_maxNum(argument2.64[I], argument3.64[I])
  };
stage E1:
  PGR[registerM] = result1;
\end{lstlisting}}
\newpage
\subsubsection[{\small FMAXNH: Floating-Point Maximum Number Half Word}]{FMAXNH\hfill Floating-Point Maximum Number Half Word}
\label{instr:FMAXNH}\index{fmaxnh}
 
\paragraph{Encoding} Format \textsf{ALU\_FHMMWR} (Section~\ref{sec:format-ALU-FHMMWR}), \verb|floatmode = 011|\\

\bitfields{ 1/P, 2/11, 1/1, 1/1, 3/011, 6/registerW, 2/01, 3/011, 1/1, 6/registerY, 6/registerZ }


\paragraph{Syntax} \texttt{fmaxnh} \emph{registerW} = \emph{registerZ}, \emph{registerY}

\paragraph{Description} The \emph{registerZ} is compared to the \emph{registerY}, assuming binary 16 floating-point numbers. The maximum number value is stored into the \emph{registerW}.


\paragraph{Execution} {Reservation table \textsf{ALU\_LITE} (Section~\ref{sec:reservation-ALU-LITE})}
~
{\begin{lstlisting}
stage RR:
  new argument3 = GPR[registerY];
  new argument2 = GPR[registerZ];
  new result1 = f16_maxNum(argument2, argument3);
stage E1:
  GPR[registerW] = result1;
\end{lstlisting}}
\newpage
\subsubsection[{\small FMAXNHO: Floating-Point Maximum Number Half Word Octuple}]{FMAXNHO\hfill Floating-Point Maximum Number Half Word Octuple}
\label{instr:FMAXNHO}\index{fmaxnho}
 
\paragraph{Encoding} Format \textsf{ALU\_FHOMMWR} (Section~\ref{sec:format-ALU-FHOMMWR}), \verb|floatmode = 011|\\

\bitfields{ 1/P, 2/11, 1/1, 1/1, 3/011, 5/registerM, 1/--, 2/11, 3/000, 1/1, 5/registerO, 1/0, 5/registerP, 1/1 }


\paragraph{Syntax} \texttt{fmaxnho} \emph{registerM} = \emph{registerP}, \emph{registerO}

\paragraph{Description} The \emph{registerP} and the \emph{registerO} are interpreted as eight binary 16 floating-point numbers packed into 128 bits. The \emph{registerP} half words are compared to the \emph{registerO} half words. The eight maximum number half words are packed and stored into the \emph{registerM}.


\paragraph{Execution} {Reservation table \textsf{ALU\_LITE} (Section~\ref{sec:reservation-ALU-LITE})}
~
{\begin{lstlisting}
stage RR:
  new argument3 = PGR[registerO];
  new argument2 = PGR[registerP];
  for (I = 0; I < 8; I += 1) {
    result1.16[I] = f16_maxNum(argument2.16[I], argument3.16[I])
  };
stage E1:
  PGR[registerM] = result1;
\end{lstlisting}}
\newpage
\subsubsection[{\small FMAXNW: Floating-Point Maximum Number Word}]{FMAXNW\hfill Floating-Point Maximum Number Word}
\label{instr:FMAXNW}\index{fmaxnw}
 
\paragraph{Encoding} Format \textsf{ALU\_FWMMWR} (Section~\ref{sec:format-ALU-FWMMWR}), \verb|floatmode = 011|\\

\bitfields{ 1/P, 2/11, 1/1, 1/1, 3/011, 6/registerW, 2/01, 3/000, 1/0, 6/registerY, 6/registerZ }


\paragraph{Syntax} \texttt{fmaxnw} \emph{registerW} = \emph{registerZ}, \emph{registerY}

\paragraph{Description} The \emph{registerZ} is compared to the \emph{registerY}, assuming binary 32 floating-point numbers. The maximum number value is stored into the \emph{registerW}.


\paragraph{Execution} {Reservation table \textsf{ALU\_LITE} (Section~\ref{sec:reservation-ALU-LITE})}
~
{\begin{lstlisting}
stage RR:
  new argument3 = GPR[registerY];
  new argument2 = GPR[registerZ];
  new result1 = f32_maxNum(argument2, argument3);
stage E1:
  GPR[registerW] = result1;
\end{lstlisting}}
\newpage
\subsubsection[{\small FMAXNWQ: Floating-Point Maximum Number Word Quadruple}]{FMAXNWQ\hfill Floating-Point Maximum Number Word Quadruple}
\label{instr:FMAXNWQ}\index{fmaxnwq}
 
\paragraph{Encoding} Format \textsf{ALU\_FWQMMWR} (Section~\ref{sec:format-ALU-FWQMMWR}), \verb|floatmode = 011|\\

\bitfields{ 1/P, 2/11, 1/1, 1/1, 3/011, 5/registerM, 1/--, 2/11, 3/000, 1/1, 5/registerO, 1/0, 5/registerP, 1/0 }


\paragraph{Syntax} \texttt{fmaxnwq} \emph{registerM} = \emph{registerP}, \emph{registerO}

\paragraph{Description} The \emph{registerP} and the \emph{registerO} are interpreted as four binary 32 floating-point numbers packed into 128 bits. The \emph{registerP} words are compared to the \emph{registerO} words. The four maximum number words are packed and stored into the \emph{registerM}.


\paragraph{Execution} {Reservation table \textsf{ALU\_LITE} (Section~\ref{sec:reservation-ALU-LITE})}
~
{\begin{lstlisting}
stage RR:
  new argument3 = PGR[registerO];
  new argument2 = PGR[registerP];
  for (I = 0; I < 4; I += 1) {
    result1.32[I] = f32_maxNum(argument2.32[I], argument3.32[I])
  };
stage E1:
  PGR[registerM] = result1;
\end{lstlisting}}
\newpage
\subsubsection[{\small FMAXW: Floating-Point Maximum Word}]{FMAXW\hfill Floating-Point Maximum Word}
\label{instr:FMAXW}\index{fmaxw}
 
\paragraph{Encoding} Format \textsf{ALU\_FWMMWR} (Section~\ref{sec:format-ALU-FWMMWR}), \verb|floatmode = 001|\\

\bitfields{ 1/P, 2/11, 1/1, 1/1, 3/001, 6/registerW, 2/01, 3/000, 1/0, 6/registerY, 6/registerZ }


\paragraph{Syntax} \texttt{fmaxw} \emph{registerW} = \emph{registerZ}, \emph{registerY}

\paragraph{Description} The \emph{registerZ} is compared to the \emph{registerY}, assuming binary 32 floating-point numbers. The maximum value is stored into the \emph{registerW}.


\paragraph{Execution} {Reservation table \textsf{ALU\_LITE} (Section~\ref{sec:reservation-ALU-LITE})}
~
{\begin{lstlisting}
stage RR:
  new argument3 = GPR[registerY];
  new argument2 = GPR[registerZ];
  new result1 = f32_max(argument2, argument3);
stage E1:
  GPR[registerW] = result1;
\end{lstlisting}}
\newpage
\subsubsection[{\small FMAXWQ: Floating-Point Maximum Word Quadruple}]{FMAXWQ\hfill Floating-Point Maximum Word Quadruple}
\label{instr:FMAXWQ}\index{fmaxwq}
 
\paragraph{Encoding} Format \textsf{ALU\_FWQMMWR} (Section~\ref{sec:format-ALU-FWQMMWR}), \verb|floatmode = 001|\\

\bitfields{ 1/P, 2/11, 1/1, 1/1, 3/001, 5/registerM, 1/--, 2/11, 3/000, 1/1, 5/registerO, 1/0, 5/registerP, 1/0 }


\paragraph{Syntax} \texttt{fmaxwq} \emph{registerM} = \emph{registerP}, \emph{registerO}

\paragraph{Description} The \emph{registerP} and the \emph{registerO} are interpreted as four binary 32 floating-point numbers packed into 128 bits. The \emph{registerP} words are compared to the \emph{registerO} words. The four maximum words are packed and stored into the \emph{registerM}.


\paragraph{Execution} {Reservation table \textsf{ALU\_LITE} (Section~\ref{sec:reservation-ALU-LITE})}
~
{\begin{lstlisting}
stage RR:
  new argument3 = PGR[registerO];
  new argument2 = PGR[registerP];
  for (I = 0; I < 4; I += 1) {
    result1.32[I] = f32_max(argument2.32[I], argument3.32[I])
  };
stage E1:
  PGR[registerM] = result1;
\end{lstlisting}}
\newpage
\subsubsection[{\small FMIND: Floating-Point Minimum Double Word}]{FMIND\hfill Floating-Point Minimum Double Word}
\label{instr:FMIND}\index{fmind}
 
\paragraph{Encoding} Format \textsf{ALU\_FDMMWR} (Section~\ref{sec:format-ALU-FDMMWR}), \verb|floatmode = 000|\\

\bitfields{ 1/P, 2/11, 1/1, 1/1, 3/000, 6/registerW, 2/00, 3/000, 1/0, 6/registerY, 6/registerZ }


\paragraph{Syntax} \texttt{fmind} \emph{registerW} = \emph{registerZ}, \emph{registerY}

\paragraph{Description} The \emph{registerZ} is compared to the \emph{registerY}, assuming binary 64 floating-point numbers. The minimum is stored into the \emph{registerW}.


\paragraph{Execution} {Reservation table \textsf{ALU\_LITE} (Section~\ref{sec:reservation-ALU-LITE})}
~
{\begin{lstlisting}
stage RR:
  new argument3 = GPR[registerY];
  new argument2 = GPR[registerZ];
  new result1 = f64_min(argument2, argument3);
stage E1:
  GPR[registerW] = result1;
\end{lstlisting}}
\newpage
\subsubsection[{\small FMINDP: Floating-Point Minimum Double Word Pair}]{FMINDP\hfill Floating-Point Minimum Double Word Pair}
\label{instr:FMINDP}\index{fmindp}
 
\paragraph{Encoding} Format \textsf{ALU\_FDPMMWR} (Section~\ref{sec:format-ALU-FDPMMWR}), \verb|floatmode = 000|\\

\bitfields{ 1/P, 2/11, 1/1, 1/1, 3/000, 5/registerM, 1/--, 2/11, 3/000, 1/0, 5/registerO, 1/0, 5/registerP, 1/1 }


\paragraph{Syntax} \texttt{fmindp} \emph{registerM} = \emph{registerP}, \emph{registerO}

\paragraph{Description} The \emph{registerP} and the \emph{registerO} are interpreted as two binary 64 floating-point numbers packed into 128 bits. The \emph{registerP} double words are compared to the \emph{registerO} double words. The two minimum double words are packed and stored into the \emph{registerM}.


\paragraph{Execution} {Reservation table \textsf{ALU\_LITE} (Section~\ref{sec:reservation-ALU-LITE})}
~
{\begin{lstlisting}
stage RR:
  new argument3 = PGR[registerO];
  new argument2 = PGR[registerP];
  for (I = 0; I < 2; I += 1) {
    result1.64[I] = f64_min(argument2.64[I], argument3.64[I])
  };
stage E1:
  PGR[registerM] = result1;
\end{lstlisting}}
\newpage
\subsubsection[{\small FMINH: Floating-Point Minimum Half Word}]{FMINH\hfill Floating-Point Minimum Half Word}
\label{instr:FMINH}\index{fminh}
 
\paragraph{Encoding} Format \textsf{ALU\_FHMMWR} (Section~\ref{sec:format-ALU-FHMMWR}), \verb|floatmode = 000|\\

\bitfields{ 1/P, 2/11, 1/1, 1/1, 3/000, 6/registerW, 2/01, 3/011, 1/1, 6/registerY, 6/registerZ }


\paragraph{Syntax} \texttt{fminh} \emph{registerW} = \emph{registerZ}, \emph{registerY}

\paragraph{Description} The \emph{registerZ} is compared to the \emph{registerY}, assuming binary 16 floating-point numbers. The minimum value is stored into the \emph{registerW}.


\paragraph{Execution} {Reservation table \textsf{ALU\_LITE} (Section~\ref{sec:reservation-ALU-LITE})}
~
{\begin{lstlisting}
stage RR:
  new argument3 = GPR[registerY];
  new argument2 = GPR[registerZ];
  new result1 = f16_min(argument2, argument3);
stage E1:
  GPR[registerW] = result1;
\end{lstlisting}}
\newpage
\subsubsection[{\small FMINHO: Floating-Point Minimum Half Word Octuple}]{FMINHO\hfill Floating-Point Minimum Half Word Octuple}
\label{instr:FMINHO}\index{fminho}
 
\paragraph{Encoding} Format \textsf{ALU\_FHOMMWR} (Section~\ref{sec:format-ALU-FHOMMWR}), \verb|floatmode = 000|\\

\bitfields{ 1/P, 2/11, 1/1, 1/1, 3/000, 5/registerM, 1/--, 2/11, 3/000, 1/1, 5/registerO, 1/0, 5/registerP, 1/1 }


\paragraph{Syntax} \texttt{fminho} \emph{registerM} = \emph{registerP}, \emph{registerO}

\paragraph{Description} The \emph{registerP} and the \emph{registerO} are interpreted as eight binary 16 floating-point numbers packed into 128 bits. The \emph{registerP} half words are compared to the \emph{registerO} half words. The eight minimum half words are packed and stored into the \emph{registerM}.


\paragraph{Execution} {Reservation table \textsf{ALU\_LITE} (Section~\ref{sec:reservation-ALU-LITE})}
~
{\begin{lstlisting}
stage RR:
  new argument3 = PGR[registerO];
  new argument2 = PGR[registerP];
  for (I = 0; I < 8; I += 1) {
    result1.16[I] = f16_min(argument2.16[I], argument3.16[I])
  };
stage E1:
  PGR[registerM] = result1;
\end{lstlisting}}
\newpage
\subsubsection[{\small FMINND: Floating-Point Minimum Number Double Word}]{FMINND\hfill Floating-Point Minimum Number Double Word}
\label{instr:FMINND}\index{fminnd}
 
\paragraph{Encoding} Format \textsf{ALU\_FDMMWR} (Section~\ref{sec:format-ALU-FDMMWR}), \verb|floatmode = 010|\\

\bitfields{ 1/P, 2/11, 1/1, 1/1, 3/010, 6/registerW, 2/00, 3/000, 1/0, 6/registerY, 6/registerZ }


\paragraph{Syntax} \texttt{fminnd} \emph{registerW} = \emph{registerZ}, \emph{registerY}

\paragraph{Description} The \emph{registerZ} is compared to the \emph{registerY}, assuming binary 64 floating-point numbers. The minimum value is stored into the \emph{registerW}.


\paragraph{Execution} {Reservation table \textsf{ALU\_LITE} (Section~\ref{sec:reservation-ALU-LITE})}
~
{\begin{lstlisting}
stage RR:
  new argument3 = GPR[registerY];
  new argument2 = GPR[registerZ];
  new result1 = f64_minNum(argument2, argument3);
stage E1:
  GPR[registerW] = result1;
\end{lstlisting}}
\newpage
\subsubsection[{\small FMINNDP: Floating-Point Minimum Number Double Word Pair}]{FMINNDP\hfill Floating-Point Minimum Number Double Word Pair}
\label{instr:FMINNDP}\index{fminndp}
 
\paragraph{Encoding} Format \textsf{ALU\_FDPMMWR} (Section~\ref{sec:format-ALU-FDPMMWR}), \verb|floatmode = 010|\\

\bitfields{ 1/P, 2/11, 1/1, 1/1, 3/010, 5/registerM, 1/--, 2/11, 3/000, 1/0, 5/registerO, 1/0, 5/registerP, 1/1 }


\paragraph{Syntax} \texttt{fminndp} \emph{registerM} = \emph{registerP}, \emph{registerO}

\paragraph{Description} The \emph{registerP} and the \emph{registerO} are interpreted as two binary 64 floating-point numbers packed into 128 bits. The \emph{registerP} double words are compared to the \emph{registerO} double words. The two minimum number double words are packed and stored into the \emph{registerM}.


\paragraph{Execution} {Reservation table \textsf{ALU\_LITE} (Section~\ref{sec:reservation-ALU-LITE})}
~
{\begin{lstlisting}
stage RR:
  new argument3 = PGR[registerO];
  new argument2 = PGR[registerP];
  for (I = 0; I < 2; I += 1) {
    result1.64[I] = f64_minNum(argument2.64[I], argument3.64[I])
  };
stage E1:
  PGR[registerM] = result1;
\end{lstlisting}}
\newpage
\subsubsection[{\small FMINNH: Floating-Point Minimum Number Half Word}]{FMINNH\hfill Floating-Point Minimum Number Half Word}
\label{instr:FMINNH}\index{fminnh}
 
\paragraph{Encoding} Format \textsf{ALU\_FHMMWR} (Section~\ref{sec:format-ALU-FHMMWR}), \verb|floatmode = 010|\\

\bitfields{ 1/P, 2/11, 1/1, 1/1, 3/010, 6/registerW, 2/01, 3/011, 1/1, 6/registerY, 6/registerZ }


\paragraph{Syntax} \texttt{fminnh} \emph{registerW} = \emph{registerZ}, \emph{registerY}

\paragraph{Description} The \emph{registerZ} is compared to the \emph{registerY}, assuming binary 16 floating-point numbers. The minimum number value is stored into the \emph{registerW}.


\paragraph{Execution} {Reservation table \textsf{ALU\_LITE} (Section~\ref{sec:reservation-ALU-LITE})}
~
{\begin{lstlisting}
stage RR:
  new argument3 = GPR[registerY];
  new argument2 = GPR[registerZ];
  new result1 = f16_minNum(argument2, argument3);
stage E1:
  GPR[registerW] = result1;
\end{lstlisting}}
\newpage
\subsubsection[{\small FMINNHO: Floating-Point Minimum Number Half Word Octuple}]{FMINNHO\hfill Floating-Point Minimum Number Half Word Octuple}
\label{instr:FMINNHO}\index{fminnho}
 
\paragraph{Encoding} Format \textsf{ALU\_FHOMMWR} (Section~\ref{sec:format-ALU-FHOMMWR}), \verb|floatmode = 010|\\

\bitfields{ 1/P, 2/11, 1/1, 1/1, 3/010, 5/registerM, 1/--, 2/11, 3/000, 1/1, 5/registerO, 1/0, 5/registerP, 1/1 }


\paragraph{Syntax} \texttt{fminnho} \emph{registerM} = \emph{registerP}, \emph{registerO}

\paragraph{Description} The \emph{registerP} and the \emph{registerO} are interpreted as eight binary 16 floating-point numbers packed into 128 bits. The \emph{registerP} half words are compared to the \emph{registerO} half words. The eight minimum number half words are packed and stored into the \emph{registerM}.


\paragraph{Execution} {Reservation table \textsf{ALU\_LITE} (Section~\ref{sec:reservation-ALU-LITE})}
~
{\begin{lstlisting}
stage RR:
  new argument3 = PGR[registerO];
  new argument2 = PGR[registerP];
  for (I = 0; I < 8; I += 1) {
    result1.16[I] = f16_minNum(argument2.16[I], argument3.16[I])
  };
stage E1:
  PGR[registerM] = result1;
\end{lstlisting}}
\newpage
\subsubsection[{\small FMINNW: Floating-Point Minimum Number Word}]{FMINNW\hfill Floating-Point Minimum Number Word}
\label{instr:FMINNW}\index{fminnw}
 
\paragraph{Encoding} Format \textsf{ALU\_FWMMWR} (Section~\ref{sec:format-ALU-FWMMWR}), \verb|floatmode = 010|\\

\bitfields{ 1/P, 2/11, 1/1, 1/1, 3/010, 6/registerW, 2/01, 3/000, 1/0, 6/registerY, 6/registerZ }


\paragraph{Syntax} \texttt{fminnw} \emph{registerW} = \emph{registerZ}, \emph{registerY}

\paragraph{Description} The \emph{registerZ} is compared to the \emph{registerY}, assuming binary 32 floating-point numbers. The minimum number value is stored into the \emph{registerW}.


\paragraph{Execution} {Reservation table \textsf{ALU\_LITE} (Section~\ref{sec:reservation-ALU-LITE})}
~
{\begin{lstlisting}
stage RR:
  new argument3 = GPR[registerY];
  new argument2 = GPR[registerZ];
  new result1 = f32_minNum(argument2, argument3);
stage E1:
  GPR[registerW] = result1;
\end{lstlisting}}
\newpage
\subsubsection[{\small FMINNWQ: Floating-Point Minimum Number Word Quadruple}]{FMINNWQ\hfill Floating-Point Minimum Number Word Quadruple}
\label{instr:FMINNWQ}\index{fminnwq}
 
\paragraph{Encoding} Format \textsf{ALU\_FWQMMWR} (Section~\ref{sec:format-ALU-FWQMMWR}), \verb|floatmode = 010|\\

\bitfields{ 1/P, 2/11, 1/1, 1/1, 3/010, 5/registerM, 1/--, 2/11, 3/000, 1/1, 5/registerO, 1/0, 5/registerP, 1/0 }


\paragraph{Syntax} \texttt{fminnwq} \emph{registerM} = \emph{registerP}, \emph{registerO}

\paragraph{Description} The \emph{registerP} and the \emph{registerO} are interpreted as four binary 32 floating-point numbers packed into 128 bits. The \emph{registerP} words are compared to the \emph{registerO} words. The four minimum number words are packed and stored into the \emph{registerM}.


\paragraph{Execution} {Reservation table \textsf{ALU\_LITE} (Section~\ref{sec:reservation-ALU-LITE})}
~
{\begin{lstlisting}
stage RR:
  new argument3 = PGR[registerO];
  new argument2 = PGR[registerP];
  for (I = 0; I < 4; I += 1) {
    result1.32[I] = f32_minNum(argument2.32[I], argument3.32[I])
  };
stage E1:
  PGR[registerM] = result1;
\end{lstlisting}}
\newpage
\subsubsection[{\small FMINW: Floating-Point Minimum Word}]{FMINW\hfill Floating-Point Minimum Word}
\label{instr:FMINW}\index{fminw}
 
\paragraph{Encoding} Format \textsf{ALU\_FWMMWR} (Section~\ref{sec:format-ALU-FWMMWR}), \verb|floatmode = 000|\\

\bitfields{ 1/P, 2/11, 1/1, 1/1, 3/000, 6/registerW, 2/01, 3/000, 1/0, 6/registerY, 6/registerZ }


\paragraph{Syntax} \texttt{fminw} \emph{registerW} = \emph{registerZ}, \emph{registerY}

\paragraph{Description} The \emph{registerZ} is compared to the \emph{registerY}, assuming binary 32 floating-point numbers. The minimum value is stored into the \emph{registerW}.


\paragraph{Execution} {Reservation table \textsf{ALU\_LITE} (Section~\ref{sec:reservation-ALU-LITE})}
~
{\begin{lstlisting}
stage RR:
  new argument3 = GPR[registerY];
  new argument2 = GPR[registerZ];
  new result1 = f32_min(argument2, argument3);
stage E1:
  GPR[registerW] = result1;
\end{lstlisting}}
\newpage
\subsubsection[{\small FMINWQ: Floating-Point Minimum Word Quadruple}]{FMINWQ\hfill Floating-Point Minimum Word Quadruple}
\label{instr:FMINWQ}\index{fminwq}
 
\paragraph{Encoding} Format \textsf{ALU\_FWQMMWR} (Section~\ref{sec:format-ALU-FWQMMWR}), \verb|floatmode = 000|\\

\bitfields{ 1/P, 2/11, 1/1, 1/1, 3/000, 5/registerM, 1/--, 2/11, 3/000, 1/1, 5/registerO, 1/0, 5/registerP, 1/0 }


\paragraph{Syntax} \texttt{fminwq} \emph{registerM} = \emph{registerP}, \emph{registerO}

\paragraph{Description} The \emph{registerP} and the \emph{registerO} are interpreted as four binary 32 floating-point numbers packed into 128 bits. The \emph{registerP} words are compared to the \emph{registerO} words. The four minimum words are packed and stored into the \emph{registerM}.


\paragraph{Execution} {Reservation table \textsf{ALU\_LITE} (Section~\ref{sec:reservation-ALU-LITE})}
~
{\begin{lstlisting}
stage RR:
  new argument3 = PGR[registerO];
  new argument2 = PGR[registerP];
  for (I = 0; I < 4; I += 1) {
    result1.32[I] = f32_min(argument2.32[I], argument3.32[I])
  };
stage E1:
  PGR[registerM] = result1;
\end{lstlisting}}
\newpage
\subsubsection[{\small FMULD: Floating-Point Multiply Double Word}]{FMULD\hfill Floating-Point Multiply Double Word}
\label{instr:FMULD}\index{fmuld}
 
\paragraph{Encoding} Format \textsf{ALU\_FDMUL} (Section~\ref{sec:format-ALU-FDMUL}), \verb|alucode3 = 110|\\

\bitfields{ 1/P, 2/11, 1/1, 1/fnegate, 3/floatmode, 6/registerW, 2/00, 3/110, 1/0, 6/registerY, 6/registerZ }


\paragraph{Syntax} \texttt{fmuld}\emph{fnegate}\emph{floatmode} \emph{registerW} = \emph{registerZ}, \emph{registerY}

\paragraph{Description} The \emph{registerZ} is multiplied by the \emph{registerY}, assuming binary 64 floating-point numbers. If \emph{fnegate} is set, the \emph{registerZ} double word is negated prior to the operation. The result is rounded to the binary 64 floating-point format according to the rounding mode and stored into the \emph{registerW}. This instruction may raise exception bits in the CS register.


\paragraph{Modifier(s)}
\noindent\begin{itemize}
\item floatmode (cf. floatmode in section \ref{par:modifier-floatmode} on page \pageref{par:modifier-floatmode})
\item fnegate (cf. fnegate in section \ref{par:modifier-fnegate} on page \pageref{par:modifier-fnegate})
\end{itemize}

\paragraph{Execution} {Reservation table \textsf{ALU\_LITE} (Section~\ref{sec:reservation-ALU-LITE})}
~
{\begin{lstlisting}
stage ID:
  new fnegate = _ZX_1(fnegate);
  new floatmode = _ZX_3(floatmode);
stage RR:
  new argument3 = GPR[registerY];
  new argument2 = GPR[registerZ];
  new RM = decode_riscv_float_rounding_mode(floatmode);
  new fsign = fnegate ? 0x8000000000000000 : 0;
  new result1 = f64_mul(RM, argument2 ^ fsign, argument3);
stage E4:
  GPR[registerW] = result1;
  commit_float_exception_flags();
\end{lstlisting}}
\newpage
\subsubsection[{\small FMULDP: Floating-Point Multiply Double Word Pair}]{FMULDP\hfill Floating-Point Multiply Double Word Pair}
\label{instr:FMULDP}\index{fmuldp}
 
\paragraph{Encoding} Format \textsf{ALU\_FDPMUL} (Section~\ref{sec:format-ALU-FDPMUL}), \verb|alucode3 = 110|\\

\bitfields{ 1/P, 2/11, 1/1, 1/fnegate, 3/floatmode, 5/registerM, 1/--, 2/11, 3/110, 1/0, 5/registerO, 1/0, 5/registerP, 1/1 }


\paragraph{Syntax} \texttt{fmuldp}\emph{fnegate}\emph{floatmode} \emph{registerM} = \emph{registerP}, \emph{registerO}

\paragraph{Description} The \emph{registerP} and the \emph{registerO} are interpreted as two binary 64 floating-point numbers packed into 128 bits. If \emph{fnegate} is set, the \emph{registerP} double words are negated prior to the operation. The \emph{registerP} double words are multiplied by the \emph{registerO} double words. The two resulting double words are packed and stored into the \emph{registerM}.


\paragraph{Modifier(s)}
\noindent\begin{itemize}
\item floatmode (cf. floatmode in section \ref{par:modifier-floatmode} on page \pageref{par:modifier-floatmode})
\item fnegate (cf. fnegate in section \ref{par:modifier-fnegate} on page \pageref{par:modifier-fnegate})
\end{itemize}

\paragraph{Execution} {Reservation table \textsf{ALU\_LITE} (Section~\ref{sec:reservation-ALU-LITE})}
~
{\begin{lstlisting}
stage ID:
  new fnegate = _ZX_1(fnegate);
  new floatmode = _ZX_3(floatmode);
stage RR:
  new argument3 = PGR[registerO];
  new argument2 = PGR[registerP];
  new RM = decode_riscv_float_rounding_mode(floatmode);
  new fsign = fnegate ? 0x8000000000000000 : 0;
  for (I = 0; I < 2; I += 1) {
    result1.64[I] = f64_mul(RM, argument2.64[I] ^ fsign, argument3.64[I])
  };
stage E4:
  PGR[registerM] = result1;
  commit_float_exception_flags();
\end{lstlisting}}
\newpage
\subsubsection[{\small FMULH: Floating-Point Multiply Half Word}]{FMULH\hfill Floating-Point Multiply Half Word}
\label{instr:FMULH}\index{fmulh}
 
\paragraph{Encoding} Format \textsf{ALU\_FHMUL} (Section~\ref{sec:format-ALU-FHMUL}), \verb|alucode3 = 110|\\

\bitfields{ 1/P, 2/11, 1/1, 1/fnegate, 3/floatmode, 6/registerW, 2/01, 3/110, 1/1, 6/registerY, 6/registerZ }


\paragraph{Syntax} \texttt{fmulh}\emph{fnegate}\emph{floatmode} \emph{registerW} = \emph{registerZ}, \emph{registerY}

\paragraph{Description} The \emph{registerZ} is multiplied by the \emph{registerY}, assuming binary 16 floating-point numbers. If \emph{fnegate} is set, the \emph{registerZ} half word is negated prior to the operation. The result is rounded to the binary 16 floating-point format according to the rounding mode and stored into the \emph{registerW}. This instruction may raise exception bits in the CS register.


\paragraph{Modifier(s)}
\noindent\begin{itemize}
\item floatmode (cf. floatmode in section \ref{par:modifier-floatmode} on page \pageref{par:modifier-floatmode})
\item fnegate (cf. fnegate in section \ref{par:modifier-fnegate} on page \pageref{par:modifier-fnegate})
\end{itemize}

\paragraph{Execution} {Reservation table \textsf{ALU\_LITE} (Section~\ref{sec:reservation-ALU-LITE})}
~
{\begin{lstlisting}
stage ID:
  new fnegate = _ZX_1(fnegate);
  new floatmode = _ZX_3(floatmode);
stage RR:
  new argument3 = GPR[registerY];
  new argument2 = GPR[registerZ];
  new RM = decode_riscv_float_rounding_mode(floatmode);
  new fsign = fnegate ? 0x8000 : 0;
  new result1 = f16_mul(RM, argument2 ^ fsign, argument3);
stage E3:
  GPR[registerW] = result1;
  commit_float_exception_flags();
\end{lstlisting}}
\newpage
\subsubsection[{\small FMULHO: Floating-Point Multiply Half Word Octuple}]{FMULHO\hfill Floating-Point Multiply Half Word Octuple}
\label{instr:FMULHO}\index{fmulho}
 
\paragraph{Encoding} Format \textsf{ALU\_FHOMUL} (Section~\ref{sec:format-ALU-FHOMUL}), \verb|alucode3 = 110|\\

\bitfields{ 1/P, 2/11, 1/1, 1/fnegate, 3/floatmode, 5/registerM, 1/--, 2/11, 3/110, 1/1, 5/registerO, 1/0, 5/registerP, 1/1 }


\paragraph{Syntax} \texttt{fmulho}\emph{fnegate}\emph{floatmode} \emph{registerM} = \emph{registerP}, \emph{registerO}

\paragraph{Description} The \emph{registerP} and the \emph{registerO} are interpreted as eight binary 16 floating-point numbers packed into 128 bits. If \emph{fnegate} is set, the \emph{registerP} half words are negated prior to the operation. The \emph{registerP} half words are multiplied by the \emph{registerO} half words. The eight resulting half words are packed and stored into the \emph{registerM}.


\paragraph{Modifier(s)}
\noindent\begin{itemize}
\item floatmode (cf. floatmode in section \ref{par:modifier-floatmode} on page \pageref{par:modifier-floatmode})
\item fnegate (cf. fnegate in section \ref{par:modifier-fnegate} on page \pageref{par:modifier-fnegate})
\end{itemize}

\paragraph{Execution} {Reservation table \textsf{ALU\_LITE} (Section~\ref{sec:reservation-ALU-LITE})}
~
{\begin{lstlisting}
stage ID:
  new fnegate = _ZX_1(fnegate);
  new floatmode = _ZX_3(floatmode);
stage RR:
  new argument3 = PGR[registerO];
  new argument2 = PGR[registerP];
  new RM = decode_riscv_float_rounding_mode(floatmode);
  new fsign = fnegate ? 0x8000 : 0;
  for (I = 0; I < 8; I += 1) {
    result1.16[I] = f16_mul(RM, argument2.16[I] ^ fsign, argument3.16[I])
  };
stage E4:
  PGR[registerM] = result1;
  commit_float_exception_flags();
\end{lstlisting}}
\newpage
\subsubsection[{\small FMULW: Floating-Point Multiply Word}]{FMULW\hfill Floating-Point Multiply Word}
\label{instr:FMULW}\index{fmulw}
 
\paragraph{Encoding} Format \textsf{ALU\_FWMUL} (Section~\ref{sec:format-ALU-FWMUL}), \verb|alucode3 = 110|\\

\bitfields{ 1/P, 2/11, 1/1, 1/fnegate, 3/floatmode, 6/registerW, 2/01, 3/110, 1/0, 6/registerY, 6/registerZ }


\paragraph{Syntax} \texttt{fmulw}\emph{fnegate}\emph{floatmode} \emph{registerW} = \emph{registerZ}, \emph{registerY}

\paragraph{Description} The \emph{registerZ} is multiplied by the \emph{registerY}, assuming binary 32 floating-point numbers. If \emph{fnegate} is set, the \emph{registerZ} word is negated prior to the operation. The result is rounded to the binary 32 floating-point format according to the rounding mode and stored into the \emph{registerW}. This instruction may raise exception bits in the CS register.


\paragraph{Modifier(s)}
\noindent\begin{itemize}
\item floatmode (cf. floatmode in section \ref{par:modifier-floatmode} on page \pageref{par:modifier-floatmode})
\item fnegate (cf. fnegate in section \ref{par:modifier-fnegate} on page \pageref{par:modifier-fnegate})
\end{itemize}

\paragraph{Execution} {Reservation table \textsf{ALU\_LITE} (Section~\ref{sec:reservation-ALU-LITE})}
~
{\begin{lstlisting}
stage ID:
  new fnegate = _ZX_1(fnegate);
  new floatmode = _ZX_3(floatmode);
stage RR:
  new argument3 = GPR[registerY];
  new argument2 = GPR[registerZ];
  new RM = decode_riscv_float_rounding_mode(floatmode);
  new fsign = fnegate ? 0x80000000 : 0;
  new result1 = f32_mul(RM, argument2 ^ fsign, argument3);
stage E4:
  GPR[registerW] = result1;
  commit_float_exception_flags();
\end{lstlisting}}
\newpage
\subsubsection[{\small FMULWC: Floating-Point Multiply Word Complex}]{FMULWC\hfill Floating-Point Multiply Word Complex}
\label{instr:FMULWC}\index{fmulwc}
 
\paragraph{Encoding} Format \textsf{ALU\_FWCOP} (Section~\ref{sec:format-ALU-FWCOP}), \verb|alucode6 = 10|\\

\bitfields{ 1/P, 2/11, 1/1, 1/imultiply, 3/floatmode, 6/registerW, 2/00, 2/10, 1/conjugate, 1/1, 6/registerY, 6/registerZ }


\paragraph{Syntax} \texttt{fmulwc}\emph{conjugate}\emph{imultiply}\emph{floatmode} \emph{registerW} = \emph{registerZ}, \emph{registerY}

\paragraph{Description} The \emph{registerY} and the \emph{registerZ} are interpreted as binary 32 floating-point complex numbers. If \emph{conjugate} is set, the complex conjugate of the \emph{registerZ} is used in the multiplication. The complex product is and rounded according to the rounding mode. If \emph{imultiply} is set, the result is multiplied by the imaginary unit. The resulting binary 32 floating-point complex number is stored back into the \emph{registerW}. This instruction may raise exception bits in the CS register.


\paragraph{Modifier(s)}
\noindent\begin{itemize}
\item floatmode (cf. floatmode in section \ref{par:modifier-floatmode} on page \pageref{par:modifier-floatmode})
\item imultiply (cf. imultiply in section \ref{par:modifier-imultiply} on page \pageref{par:modifier-imultiply})
\item conjugate (cf. conjugate in section \ref{par:modifier-conjugate} on page \pageref{par:modifier-conjugate})
\end{itemize}

\paragraph{Execution} {Reservation table \textsf{ALU\_LITE} (Section~\ref{sec:reservation-ALU-LITE})}
~
{\begin{lstlisting}
stage ID:
  new conjugate = _ZX_1(conjugate);
  new imultiply = _ZX_1(imultiply);
  new floatmode = _ZX_3(floatmode);
stage RR:
  new argument3 = GPR[registerY];
  new argument2 = GPR[registerZ];
  new argument1 = GPR[registerW];
  new RM = decode_riscv_float_rounding_mode(floatmode);
  if (conjugate) {
    argument2.64[0] = fconj_32_32(argument2.64[0]);
  }
  result1.64[0] = fmulc_32_32(RM, argument2.64[0], argument3.64[0]);
  if (imultiply) {
    new result1 = _SWAP_32(result1);
    result1.32[0] = result1.32[0] ^ 0x80000000;
  }
stage E4:
  GPR[registerW] = result1;
  commit_float_exception_flags();
\end{lstlisting}}
\newpage
\subsubsection[{\small FMULWQ: Floating-Point Multiply Word Quadruple}]{FMULWQ\hfill Floating-Point Multiply Word Quadruple}
\label{instr:FMULWQ}\index{fmulwq}
 
\paragraph{Encoding} Format \textsf{ALU\_FWQMUL} (Section~\ref{sec:format-ALU-FWQMUL}), \verb|alucode3 = 110|\\

\bitfields{ 1/P, 2/11, 1/1, 1/fnegate, 3/floatmode, 5/registerM, 1/--, 2/11, 3/110, 1/1, 5/registerO, 1/0, 5/registerP, 1/0 }


\paragraph{Syntax} \texttt{fmulwq}\emph{fnegate}\emph{floatmode} \emph{registerM} = \emph{registerP}, \emph{registerO}

\paragraph{Description} The \emph{registerP} and the \emph{registerO} are interpreted as four binary 32 floating-point numbers packed into 128 bits. If \emph{fnegate} is set, the \emph{registerP} words are negated prior to the operation. The \emph{registerP} words are multiplied by the \emph{registerO} words. The four resulting words are packed and stored into the \emph{registerM}.


\paragraph{Modifier(s)}
\noindent\begin{itemize}
\item floatmode (cf. floatmode in section \ref{par:modifier-floatmode} on page \pageref{par:modifier-floatmode})
\item fnegate (cf. fnegate in section \ref{par:modifier-fnegate} on page \pageref{par:modifier-fnegate})
\end{itemize}

\paragraph{Execution} {Reservation table \textsf{ALU\_LITE} (Section~\ref{sec:reservation-ALU-LITE})}
~
{\begin{lstlisting}
stage ID:
  new fnegate = _ZX_1(fnegate);
  new floatmode = _ZX_3(floatmode);
stage RR:
  new argument3 = PGR[registerO];
  new argument2 = PGR[registerP];
  new RM = decode_riscv_float_rounding_mode(floatmode);
  new fsign = fnegate ? 0x80000000 : 0;
  for (I = 0; I < 4; I += 1) {
    result1.32[I] = f32_mul(RM, argument2.32[I] ^ fsign, argument3.32[I])
  };
stage E4:
  PGR[registerM] = result1;
  commit_float_exception_flags();
\end{lstlisting}}
\newpage
\subsubsection[{\small FNARROWDW: Floating Point Narrow Double Word to Word}]{FNARROWDW\hfill Floating Point Narrow Double Word to Word}
\label{instr:FNARROWDW}\index{fnarrowdw}
 
\paragraph{Encoding} Format \textsf{ALU\_FWWR} (Section~\ref{sec:format-ALU-FWWR}), \verb|fpucode6 = 1100|\\

\bitfields{ 1/P, 2/11, 1/1, 1/1, 3/floatmode, 6/registerW, 2/01, 3/001, 1/0, 4/1100, 1/--, 1/0, 6/registerZ }


\paragraph{Syntax} \texttt{fnarrowdw}\emph{floatmode} \emph{registerW} = \emph{registerZ}

\paragraph{Description} The \emph{registerZ} interpreted as a binary 64 floating-point number is converted to a binary 32 floating-point number according to the rounding mode and stored into the \emph{registerW}. This instruction may raise exception bits in the CS register.


\paragraph{Modifier(s)}
\noindent\begin{itemize}
\item floatmode (cf. floatmode in section \ref{par:modifier-floatmode} on page \pageref{par:modifier-floatmode})
\end{itemize}

\paragraph{Execution} {Reservation table \textsf{ALU\_LITE} (Section~\ref{sec:reservation-ALU-LITE})}
~
{\begin{lstlisting}
stage ID:
  new floatmode = _ZX_3(floatmode);
stage RR:
  new argument2 = GPR[registerZ];
  new RM = decode_riscv_float_rounding_mode(floatmode);
  new result1 = f64_to_f32(RM, argument2.64[0]);
stage E1:
  GPR[registerW] = result1;
  commit_float_exception_flags();
\end{lstlisting}}
\newpage
\subsubsection[{\small FNARROWWH: Floating Point Narrow Word to Half Word}]{FNARROWWH\hfill Floating Point Narrow Word to Half Word}
\label{instr:FNARROWWH}\index{fnarrowwh}
 
\paragraph{Encoding} Format \textsf{ALU\_FHWR} (Section~\ref{sec:format-ALU-FHWR}), \verb|fpucode6 = 1100|\\

\bitfields{ 1/P, 2/11, 1/1, 1/1, 3/floatmode, 6/registerW, 2/01, 3/001, 1/1, 4/1100, 1/--, 1/0, 6/registerZ }


\paragraph{Syntax} \texttt{fnarrowwh}\emph{floatmode} \emph{registerW} = \emph{registerZ}

\paragraph{Description} The \emph{registerZ} interpreted as a binary 32 floating-point number is converted to a binary 16 floating-point number according to the rounding mode and stored into the \emph{registerW}. This instruction may raise exception bits in the CS register.


\paragraph{Modifier(s)}
\noindent\begin{itemize}
\item floatmode (cf. floatmode in section \ref{par:modifier-floatmode} on page \pageref{par:modifier-floatmode})
\end{itemize}

\paragraph{Execution} {Reservation table \textsf{ALU\_LITE} (Section~\ref{sec:reservation-ALU-LITE})}
~
{\begin{lstlisting}
stage ID:
  new floatmode = _ZX_3(floatmode);
stage RR:
  new argument2 = GPR[registerZ];
  new RM = decode_riscv_float_rounding_mode(floatmode);
  new result1 = f32_to_f16(RM, argument2.32[0]);
stage E1:
  GPR[registerW] = result1;
  commit_float_exception_flags();
\end{lstlisting}}
\newpage
\subsubsection[{\small FNEGD: Floating-Point Negate Double Word}]{FNEGD\hfill Floating-Point Negate Double Word}
\label{instr:FNEGD}\index{fnegd}
 
\paragraph{Encoding} Format \textsf{ALU\_FDMO1} (Section~\ref{sec:format-ALU-FDMO1}), \verb|fpucode6 = 0010|\\

\bitfields{ 1/P, 2/11, 1/1, 1/1, 3/111, 6/registerW, 2/00, 3/000, 1/0, 4/0010, 1/--, 1/0, 6/registerZ }


\paragraph{Syntax} \texttt{fnegd} \emph{registerW} = \emph{registerZ}

\paragraph{Description} The negated value the \emph{registerZ}, assuming binary 64 floating-point number, is stored into the \emph{registerW}. Implemented as a bitwise operation, does not raise FP exceptions.


\paragraph{Execution} {Reservation table \textsf{ALU\_LITE} (Section~\ref{sec:reservation-ALU-LITE})}
~
{\begin{lstlisting}
stage RR:
  new argument2 = GPR[registerZ];
  new result1 = argument2 ^ 0x8000000000000000;
stage E1:
  GPR[registerW] = result1;
\end{lstlisting}}
\newpage
\subsubsection[{\small FNEGDP: Floating-Point Negate Double Word Pair}]{FNEGDP\hfill Floating-Point Negate Double Word Pair}
\label{instr:FNEGDP}\index{fnegdp}
 
\paragraph{Encoding} Format \textsf{ALU\_FDPMO1} (Section~\ref{sec:format-ALU-FDPMO1}), \verb|fpucode6 = 0010|\\

\bitfields{ 1/P, 2/11, 1/1, 1/1, 3/111, 5/registerM, 1/--, 2/11, 3/000, 1/0, 4/0010, 1/--, 1/0, 5/registerP, 1/1 }


\paragraph{Syntax} \texttt{fnegdp} \emph{registerM} = \emph{registerP}

\paragraph{Description} The \emph{registerP} is interpreted as two binary 64 floating-point numbers packed into 128 bits. The two negated values the \emph{registerP} double words are packed and stored into the \emph{registerM}.


\paragraph{Execution} {Reservation table \textsf{ALU\_LITE} (Section~\ref{sec:reservation-ALU-LITE})}
~
{\begin{lstlisting}
stage RR:
  new argument2 = PGR[registerP];
  for (I = 0; I < 2; I += 1) {
    result1.64[I] = argument2.64[I] ^ 0x8000000000000000
  };
stage E1:
  PGR[registerM] = result1;
\end{lstlisting}}
\newpage
\subsubsection[{\small FNEGH: Floating-Point Negate Half Word}]{FNEGH\hfill Floating-Point Negate Half Word}
\label{instr:FNEGH}\index{fnegh}
 
\paragraph{Encoding} Format \textsf{ALU\_FHMO1} (Section~\ref{sec:format-ALU-FHMO1}), \verb|fpucode6 = 0010|\\

\bitfields{ 1/P, 2/11, 1/1, 1/1, 3/111, 6/registerW, 2/01, 3/000, 1/1, 4/0010, 1/--, 1/0, 6/registerZ }


\paragraph{Syntax} \texttt{fnegh} \emph{registerW} = \emph{registerZ}

\paragraph{Description} The negated value the \emph{registerZ}, assuming binary 16 floating-point number, is stored into the \emph{registerW}. Implemented as a bitwise operation, does not raise FP exceptions.


\paragraph{Execution} {Reservation table \textsf{ALU\_LITE} (Section~\ref{sec:reservation-ALU-LITE})}
~
{\begin{lstlisting}
stage RR:
  new argument2 = GPR[registerZ];
  new result1 = _ZX_32(argument2 ^ 0x8000);
stage E1:
  GPR[registerW] = result1;
\end{lstlisting}}
\newpage
\subsubsection[{\small FNEGHO: Floating-Point Negate Half Word Octuple}]{FNEGHO\hfill Floating-Point Negate Half Word Octuple}
\label{instr:FNEGHO}\index{fnegho}
 
\paragraph{Encoding} Format \textsf{ALU\_FHOMO1} (Section~\ref{sec:format-ALU-FHOMO1}), \verb|fpucode6 = 0010|\\

\bitfields{ 1/P, 2/11, 1/1, 1/1, 3/111, 5/registerM, 1/--, 2/11, 3/000, 1/1, 4/0010, 1/--, 1/0, 5/registerP, 1/1 }


\paragraph{Syntax} \texttt{fnegho} \emph{registerM} = \emph{registerP}

\paragraph{Description} The \emph{registerP} is interpreted as eight binary 16 floating-point numbers packed into 128 bits. The eight negated values the \emph{registerP} half words are packed and stored into the \emph{registerM}.


\paragraph{Execution} {Reservation table \textsf{ALU\_LITE} (Section~\ref{sec:reservation-ALU-LITE})}
~
{\begin{lstlisting}
stage RR:
  new argument2 = PGR[registerP];
  for (I = 0; I < 8; I += 1) {
    result1.16[I] = argument2.16[I] ^ 0x8000
  };
stage E1:
  PGR[registerM] = result1;
\end{lstlisting}}
\newpage
\subsubsection[{\small FNEGW: Floating-Point Negate Word}]{FNEGW\hfill Floating-Point Negate Word}
\label{instr:FNEGW}\index{fnegw}
 
\paragraph{Encoding} Format \textsf{ALU\_FWMO1} (Section~\ref{sec:format-ALU-FWMO1}), \verb|fpucode6 = 0010|\\

\bitfields{ 1/P, 2/11, 1/1, 1/1, 3/111, 6/registerW, 2/01, 3/000, 1/0, 4/0010, 1/--, 1/0, 6/registerZ }


\paragraph{Syntax} \texttt{fnegw} \emph{registerW} = \emph{registerZ}

\paragraph{Description} The negated value the \emph{registerZ}, assuming binary 32 floating-point number, is stored into the \emph{registerW}. Implemented as a bitwise operation, does not raise FP exceptions.


\paragraph{Execution} {Reservation table \textsf{ALU\_LITE} (Section~\ref{sec:reservation-ALU-LITE})}
~
{\begin{lstlisting}
stage RR:
  new argument2 = GPR[registerZ];
  new result1 = _ZX_32(argument2 ^ 0x80000000);
stage E1:
  GPR[registerW] = result1;
\end{lstlisting}}
\newpage
\subsubsection[{\small FNEGWP: Floating-Point Negate Word Pair}]{FNEGWP\hfill Floating-Point Negate Word Pair}
\label{instr:FNEGWP}\index{fnegwp}
 
\paragraph{Encoding} Format \textsf{ALU\_FWPMO1} (Section~\ref{sec:format-ALU-FWPMO1}), \verb|fpucode6 = 0010|\\

\bitfields{ 1/P, 2/11, 1/1, 1/1, 3/111, 6/registerW, 2/00, 3/000, 1/0, 4/0010, 1/--, 1/1, 6/registerZ }


\paragraph{Syntax} \texttt{fnegwp} \emph{registerW} = \emph{registerZ}

\paragraph{Description} The \emph{registerZ} is interpreted as two binary 32 floating-point numbers packed into 64 bits. The two negated values the \emph{registerZ} words are packed and stored into the \emph{registerW}.


\paragraph{Execution} {Reservation table \textsf{ALU\_LITE} (Section~\ref{sec:reservation-ALU-LITE})}
~
{\begin{lstlisting}
stage RR:
  new argument2 = GPR[registerZ];
  for (I = 0; I < 2; I += 1) {
    result1.32[I] = argument2.32[I] ^ 0x80000000
  };
stage E1:
  GPR[registerW] = result1;
\end{lstlisting}}
\newpage
\subsubsection[{\small FNEGWQ: Floating-Point Negate Word Quadruple}]{FNEGWQ\hfill Floating-Point Negate Word Quadruple}
\label{instr:FNEGWQ}\index{fnegwq}
 
\paragraph{Encoding} Format \textsf{ALU\_FWQMO1} (Section~\ref{sec:format-ALU-FWQMO1}), \verb|fpucode6 = 0010|\\

\bitfields{ 1/P, 2/11, 1/1, 1/1, 3/111, 5/registerM, 1/--, 2/11, 3/000, 1/1, 4/0010, 1/--, 1/0, 5/registerP, 1/0 }


\paragraph{Syntax} \texttt{fnegwq} \emph{registerM} = \emph{registerP}

\paragraph{Description} The \emph{registerP} is interpreted as four binary 32 floating-point numbers packed into 128 bits. The four negated values the \emph{registerP} words are packed and stored into the \emph{registerM}.


\paragraph{Execution} {Reservation table \textsf{ALU\_LITE} (Section~\ref{sec:reservation-ALU-LITE})}
~
{\begin{lstlisting}
stage RR:
  new argument2 = PGR[registerP];
  for (I = 0; I < 4; I += 1) {
    result1.32[I] = argument2.32[I] ^ 0x80000000
  };
stage E1:
  PGR[registerM] = result1;
\end{lstlisting}}
\newpage
\subsubsection[{\small FRINTD: Floating Point Round to Int Double Word}]{FRINTD\hfill Floating Point Round to Int Double Word}
\label{instr:FRINTD}\index{frintd}
 
\paragraph{Encoding} Format \textsf{ALU\_FDMO0} (Section~\ref{sec:format-ALU-FDMO0}), \verb|fpucode6 = 1001|\\

\bitfields{ 1/P, 2/11, 1/1, 1/1, 3/floatmode, 6/registerW, 2/00, 3/001, 1/0, 4/1001, 1/--, 1/0, 6/registerZ }


\paragraph{Syntax} \texttt{frintd}\emph{floatmode} \emph{registerW} = \emph{registerZ}

\paragraph{Description} The the \emph{registerZ} interpreted as a 64-bit floating-point number is rounded to integer in the 64-bit floating-point reprsentation according to the rounding mode and stored into the \emph{registerW}. This instruction may raise exception bits in the CS register.


\paragraph{Modifier(s)}
\noindent\begin{itemize}
\item floatmode (cf. floatmode in section \ref{par:modifier-floatmode} on page \pageref{par:modifier-floatmode})
\end{itemize}

\paragraph{Execution} {Reservation table \textsf{ALU\_LITE} (Section~\ref{sec:reservation-ALU-LITE})}
~
{\begin{lstlisting}
stage ID:
  new floatmode = _ZX_3(floatmode);
stage RR:
  new argument2 = GPR[registerZ];
  new RM = decode_riscv_float_rounding_mode(floatmode);
  new result1 = f64_rint(RM, argument2);
stage E4:
  GPR[registerW] = result1;
  commit_float_exception_flags();
\end{lstlisting}}
\newpage
\subsubsection[{\small FRINTH: Floating Point Round to Int Half Word}]{FRINTH\hfill Floating Point Round to Int Half Word}
\label{instr:FRINTH}\index{frinth}
 
\paragraph{Encoding} Format \textsf{ALU\_FHMO0} (Section~\ref{sec:format-ALU-FHMO0}), \verb|fpucode6 = 1001|\\

\bitfields{ 1/P, 2/11, 1/1, 1/1, 3/floatmode, 6/registerW, 2/01, 3/001, 1/1, 4/1001, 1/--, 1/0, 6/registerZ }


\paragraph{Syntax} \texttt{frinth}\emph{floatmode} \emph{registerW} = \emph{registerZ}

\paragraph{Description} The the \emph{registerZ} interpreted as a 16-bit floating-point number is rounded to integer in the 16-bit floating-point reprsentation according to the rounding mode and stored into the \emph{registerW}. This instruction may raise exception bits in the CS register.


\paragraph{Modifier(s)}
\noindent\begin{itemize}
\item floatmode (cf. floatmode in section \ref{par:modifier-floatmode} on page \pageref{par:modifier-floatmode})
\end{itemize}

\paragraph{Execution} {Reservation table \textsf{ALU\_LITE} (Section~\ref{sec:reservation-ALU-LITE})}
~
{\begin{lstlisting}
stage ID:
  new floatmode = _ZX_3(floatmode);
stage RR:
  new argument2 = GPR[registerZ];
  new RM = decode_riscv_float_rounding_mode(floatmode);
  new result1 = f16_rint(RM, argument2);
stage SF:
  GPR[registerW] = result1;
  commit_float_exception_flags();
\end{lstlisting}}
\newpage
\subsubsection[{\small FRINTW: Floating Point Round to Int Word}]{FRINTW\hfill Floating Point Round to Int Word}
\label{instr:FRINTW}\index{frintw}
 
\paragraph{Encoding} Format \textsf{ALU\_FWMO0} (Section~\ref{sec:format-ALU-FWMO0}), \verb|fpucode6 = 1001|\\

\bitfields{ 1/P, 2/11, 1/1, 1/1, 3/floatmode, 6/registerW, 2/01, 3/001, 1/0, 4/1001, 1/--, 1/0, 6/registerZ }


\paragraph{Syntax} \texttt{frintw}\emph{floatmode} \emph{registerW} = \emph{registerZ}

\paragraph{Description} The the \emph{registerZ} interpreted as a 32-bit floating-point number is rounded to integer in the 32-bit floating-point reprsentation according to the rounding mode and stored into the \emph{registerW}. This instruction may raise exception bits in the CS register.


\paragraph{Modifier(s)}
\noindent\begin{itemize}
\item floatmode (cf. floatmode in section \ref{par:modifier-floatmode} on page \pageref{par:modifier-floatmode})
\end{itemize}

\paragraph{Execution} {Reservation table \textsf{ALU\_LITE} (Section~\ref{sec:reservation-ALU-LITE})}
~
{\begin{lstlisting}
stage ID:
  new floatmode = _ZX_3(floatmode);
stage RR:
  new argument2 = GPR[registerZ];
  new RM = decode_riscv_float_rounding_mode(floatmode);
  new result1 = f32_rint(RM, argument2);
stage SF:
  GPR[registerW] = result1;
  commit_float_exception_flags();
\end{lstlisting}}
\newpage
\subsubsection[{\small FSBFD: Floating-Point Subtract From Double Word}]{FSBFD\hfill Floating-Point Subtract From Double Word}
\label{instr:FSBFD}\index{fsbfd}
 
\paragraph{Encoding} Format \textsf{ALU\_FDADD} (Section~\ref{sec:format-ALU-FDADD}), \verb|alucode3 = 001|\\

\bitfields{ 1/P, 2/11, 1/1, 1/0, 3/floatmode, 6/registerW, 2/00, 3/001, 1/0, 6/registerY, 6/registerZ }


\paragraph{Syntax} \texttt{fsbfd}\emph{floatmode} \emph{registerW} = \emph{registerZ}, \emph{registerY}

\paragraph{Description} The \emph{registerZ} is subtracted from the \emph{registerY}, assuming binary 64 floating-point numbers. The result is rounded to the binary 64 floating-point format according to the rounding mode and stored into the \emph{registerW}. This instruction may raise exception bits in the CS register.


\paragraph{Modifier(s)}
\noindent\begin{itemize}
\item floatmode (cf. floatmode in section \ref{par:modifier-floatmode} on page \pageref{par:modifier-floatmode})
\end{itemize}

\paragraph{Execution} {Reservation table \textsf{ALU\_LITE} (Section~\ref{sec:reservation-ALU-LITE})}
~
{\begin{lstlisting}
stage ID:
  new floatmode = _ZX_3(floatmode);
stage RR:
  new argument3 = GPR[registerY];
  new argument2 = GPR[registerZ];
  new RM = decode_riscv_float_rounding_mode(floatmode);
  new result1 = f64_sub(RM, argument3, argument2);
stage E4:
  GPR[registerW] = result1;
  commit_float_exception_flags();
\end{lstlisting}}
\newpage
\subsubsection[{\small FSBFDP: Floating-Point Subtract From Double Word Pair}]{FSBFDP\hfill Floating-Point Subtract From Double Word Pair}
\label{instr:FSBFDP}\index{fsbfdp}
 
\paragraph{Encoding} Format \textsf{ALU\_FDPADD} (Section~\ref{sec:format-ALU-FDPADD}), \verb|alucode3 = 001|\\

\bitfields{ 1/P, 2/11, 1/1, 1/0, 3/floatmode, 5/registerM, 1/--, 2/11, 3/001, 1/0, 5/registerO, 1/--, 5/registerP, 1/1 }


\paragraph{Syntax} \texttt{fsbfdp}\emph{floatmode} \emph{registerM} = \emph{registerP}, \emph{registerO}

\paragraph{Description} The \emph{registerP} and the \emph{registerO} are interpreted as two binary 64 floating-point numbers packed into 128 bits. The \emph{registerP} double words are subtracted from the \emph{registerO} double words. The two resulting double words are packed and stored into the \emph{registerM}.


\paragraph{Modifier(s)}
\noindent\begin{itemize}
\item floatmode (cf. floatmode in section \ref{par:modifier-floatmode} on page \pageref{par:modifier-floatmode})
\end{itemize}

\paragraph{Execution} {Reservation table \textsf{ALU\_LITE} (Section~\ref{sec:reservation-ALU-LITE})}
~
{\begin{lstlisting}
stage ID:
  new floatmode = _ZX_3(floatmode);
stage RR:
  new argument3 = PGR[registerO];
  new argument2 = PGR[registerP];
  new RM = decode_riscv_float_rounding_mode(floatmode);
  for (I = 0; I < 2; I += 1) {
    result1.64[I] = f64_sub(RM, argument3.64[I], argument2.64[I])
  };
stage E4:
  PGR[registerM] = result1;
  commit_float_exception_flags();
\end{lstlisting}}
\newpage
\subsubsection[{\small FSBFH: Floating-Point Subtract From Half Word}]{FSBFH\hfill Floating-Point Subtract From Half Word}
\label{instr:FSBFH}\index{fsbfh}
 
\paragraph{Encoding} Format \textsf{ALU\_FHADD} (Section~\ref{sec:format-ALU-FHADD}), \verb|alucode3 = 001|\\

\bitfields{ 1/P, 2/11, 1/1, 1/0, 3/floatmode, 6/registerW, 2/01, 3/001, 1/1, 6/registerY, 6/registerZ }


\paragraph{Syntax} \texttt{fsbfh}\emph{floatmode} \emph{registerW} = \emph{registerZ}, \emph{registerY}

\paragraph{Description} The \emph{registerZ} is subtracted from the \emph{registerY}, assuming binary 16 floating-point numbers. The result is rounded to the binary 16 floating-point format according to the rounding mode and stored into the \emph{registerW}. This instruction may raise exception bits in the CS register.


\paragraph{Modifier(s)}
\noindent\begin{itemize}
\item floatmode (cf. floatmode in section \ref{par:modifier-floatmode} on page \pageref{par:modifier-floatmode})
\end{itemize}

\paragraph{Execution} {Reservation table \textsf{ALU\_LITE} (Section~\ref{sec:reservation-ALU-LITE})}
~
{\begin{lstlisting}
stage ID:
  new floatmode = _ZX_3(floatmode);
stage RR:
  new argument3 = GPR[registerY];
  new argument2 = GPR[registerZ];
  new RM = decode_riscv_float_rounding_mode(floatmode);
  new result1 = f16_sub(RM, argument3, argument2);
stage E3:
  GPR[registerW] = result1;
  commit_float_exception_flags();
\end{lstlisting}}
\newpage
\subsubsection[{\small FSBFHO: Floating-Point Subtract From Half Word Octuple}]{FSBFHO\hfill Floating-Point Subtract From Half Word Octuple}
\label{instr:FSBFHO}\index{fsbfho}
 
\paragraph{Encoding} Format \textsf{ALU\_FHOADD} (Section~\ref{sec:format-ALU-FHOADD}), \verb|alucode3 = 001|\\

\bitfields{ 1/P, 2/11, 1/1, 1/0, 3/floatmode, 5/registerM, 1/--, 2/11, 3/001, 1/1, 5/registerO, 1/--, 5/registerP, 1/1 }


\paragraph{Syntax} \texttt{fsbfho}\emph{floatmode} \emph{registerM} = \emph{registerP}, \emph{registerO}

\paragraph{Description} The \emph{registerP} and the \emph{registerO} are interpreted as eight binary 16 floating-point numbers packed into 128 bits. The \emph{registerP} half words are subtracted from the \emph{registerO} half words. The eight resulting half words are packed and stored into the \emph{registerM}.


\paragraph{Modifier(s)}
\noindent\begin{itemize}
\item floatmode (cf. floatmode in section \ref{par:modifier-floatmode} on page \pageref{par:modifier-floatmode})
\end{itemize}

\paragraph{Execution} {Reservation table \textsf{ALU\_LITE} (Section~\ref{sec:reservation-ALU-LITE})}
~
{\begin{lstlisting}
stage ID:
  new floatmode = _ZX_3(floatmode);
stage RR:
  new argument3 = PGR[registerO];
  new argument2 = PGR[registerP];
  new RM = decode_riscv_float_rounding_mode(floatmode);
  for (I = 0; I < 8; I += 1) {
    result1.16[I] = f16_sub(RM, argument3.16[I], argument2.16[I])
  };
stage E4:
  PGR[registerM] = result1;
  commit_float_exception_flags();
\end{lstlisting}}
\newpage
\subsubsection[{\small FSBFW: Floating-Point Subtract From Word}]{FSBFW\hfill Floating-Point Subtract From Word}
\label{instr:FSBFW}\index{fsbfw}
 
\paragraph{Encoding} Format \textsf{ALU\_FWADD} (Section~\ref{sec:format-ALU-FWADD}), \verb|alucode3 = 001|\\

\bitfields{ 1/P, 2/11, 1/1, 1/0, 3/floatmode, 6/registerW, 2/01, 3/001, 1/0, 6/registerY, 6/registerZ }


\paragraph{Syntax} \texttt{fsbfw}\emph{floatmode} \emph{registerW} = \emph{registerZ}, \emph{registerY}

\paragraph{Description} The \emph{registerZ} is subtracted from the \emph{registerY}, assuming binary 32 floating-point numbers. The result is rounded to the binary 32 floating-point format according to the rounding mode and stored into the \emph{registerW}. This instruction may raise exception bits in the CS register.


\paragraph{Modifier(s)}
\noindent\begin{itemize}
\item floatmode (cf. floatmode in section \ref{par:modifier-floatmode} on page \pageref{par:modifier-floatmode})
\end{itemize}

\paragraph{Execution} {Reservation table \textsf{ALU\_LITE} (Section~\ref{sec:reservation-ALU-LITE})}
~
{\begin{lstlisting}
stage ID:
  new floatmode = _ZX_3(floatmode);
stage RR:
  new argument3 = GPR[registerY];
  new argument2 = GPR[registerZ];
  new RM = decode_riscv_float_rounding_mode(floatmode);
  new result1 = f32_sub(RM, argument3, argument2);
stage E4:
  GPR[registerW] = result1;
  commit_float_exception_flags();
\end{lstlisting}}
\newpage
\subsubsection[{\small FSBFWC: Floating-Point Subtract from Word Complex}]{FSBFWC\hfill Floating-Point Subtract from Word Complex}
\label{instr:FSBFWC}\index{fsbfwc}
 
\paragraph{Encoding} Format \textsf{ALU\_FWCOP} (Section~\ref{sec:format-ALU-FWCOP}), \verb|alucode6 = 01|\\

\bitfields{ 1/P, 2/11, 1/1, 1/imultiply, 3/floatmode, 6/registerW, 2/00, 2/01, 1/conjugate, 1/1, 6/registerY, 6/registerZ }


\paragraph{Syntax} \texttt{fsbfwc}\emph{conjugate}\emph{imultiply}\emph{floatmode} \emph{registerW} = \emph{registerZ}, \emph{registerY}

\paragraph{Description} The \emph{registerY} and the \emph{registerZ} are interpreted as binary 32 floating-point complex numbers. If \emph{conjugate} is set, the complex conjugate of the \emph{registerZ} is used in the operation. The real and imaginary parts are respectively subtracted from and rounded according to the rounding mode. If \emph{imultiply} is set, the result is multiplied by the imaginary unit. The resulting binary 32 floating-point complex number is stored back into the \emph{registerW}. This instruction may raise exception bits in the CS register.


\paragraph{Modifier(s)}
\noindent\begin{itemize}
\item floatmode (cf. floatmode in section \ref{par:modifier-floatmode} on page \pageref{par:modifier-floatmode})
\item imultiply (cf. imultiply in section \ref{par:modifier-imultiply} on page \pageref{par:modifier-imultiply})
\item conjugate (cf. conjugate in section \ref{par:modifier-conjugate} on page \pageref{par:modifier-conjugate})
\end{itemize}

\paragraph{Execution} {Reservation table \textsf{ALU\_LITE} (Section~\ref{sec:reservation-ALU-LITE})}
~
{\begin{lstlisting}
stage ID:
  new conjugate = _ZX_1(conjugate);
  new imultiply = _ZX_1(imultiply);
  new floatmode = _ZX_3(floatmode);
stage RR:
  new argument3 = GPR[registerY];
  new argument2 = GPR[registerZ];
  new argument1 = GPR[registerW];
  new RM = decode_riscv_float_rounding_mode(floatmode);
  new fconj = conjugate ? 0x80000000 : 0;
  result1.32[0] = f32_sub(RM, argument3.32[0], argument2.32[0]);
  result1.32[1] = f32_sub(RM, argument3.32[1], argument2.32[1] ^ fconj);
  if (imultiply) {
    new result1 = _SWAP_32(result1);
    result1.32[0] = result1.32[0] ^ 0x80000000;
  }
stage E4:
  GPR[registerW] = result1;
  commit_float_exception_flags();
\end{lstlisting}}
\newpage
\subsubsection[{\small FSBFWQ: Floating-Point Subtract From Word Quadruple}]{FSBFWQ\hfill Floating-Point Subtract From Word Quadruple}
\label{instr:FSBFWQ}\index{fsbfwq}
 
\paragraph{Encoding} Format \textsf{ALU\_FWQADD} (Section~\ref{sec:format-ALU-FWQADD}), \verb|alucode3 = 001|\\

\bitfields{ 1/P, 2/11, 1/1, 1/0, 3/floatmode, 5/registerM, 1/--, 2/11, 3/001, 1/1, 5/registerO, 1/--, 5/registerP, 1/0 }


\paragraph{Syntax} \texttt{fsbfwq}\emph{floatmode} \emph{registerM} = \emph{registerP}, \emph{registerO}

\paragraph{Description} The \emph{registerP} and the \emph{registerO} are interpreted as four binary 32 floating-point numbers packed into 128 bits. The \emph{registerP} words are subtracted from to the \emph{registerO} words. The four resulting words are packed and stored into the \emph{registerM}.


\paragraph{Modifier(s)}
\noindent\begin{itemize}
\item floatmode (cf. floatmode in section \ref{par:modifier-floatmode} on page \pageref{par:modifier-floatmode})
\end{itemize}

\paragraph{Execution} {Reservation table \textsf{ALU\_LITE} (Section~\ref{sec:reservation-ALU-LITE})}
~
{\begin{lstlisting}
stage ID:
  new floatmode = _ZX_3(floatmode);
stage RR:
  new argument3 = PGR[registerO];
  new argument2 = PGR[registerP];
  new RM = decode_riscv_float_rounding_mode(floatmode);
  for (I = 0; I < 4; I += 1) {
    result1.32[I] = f32_sub(RM, argument3.32[I], argument2.32[I])
  };
stage E4:
  PGR[registerM] = result1;
  commit_float_exception_flags();
\end{lstlisting}}
\newpage
\subsubsection[{\small FSIGND: Floating-Point Sign Transfer Double Word}]{FSIGND\hfill Floating-Point Sign Transfer Double Word}
\label{instr:FSIGND}\index{fsignd}
 
\paragraph{Encoding} Format \textsf{ALU\_FDMMWR} (Section~\ref{sec:format-ALU-FDMMWR}), \verb|floatmode = 100|\\

\bitfields{ 1/P, 2/11, 1/1, 1/1, 3/100, 6/registerW, 2/00, 3/000, 1/0, 6/registerY, 6/registerZ }


\paragraph{Syntax} \texttt{fsignd} \emph{registerW} = \emph{registerZ}, \emph{registerY}

\paragraph{Description} The \emph{registerZ} and the \emph{registerY} are interpreted as binary 64 floating-point numbers. The non-sign bits of the \emph{registerZ} are concatenated with the sign bit of the \emph{registerY}. The result is stored into the \emph{registerW}.


\paragraph{Execution} {Reservation table \textsf{ALU\_LITE} (Section~\ref{sec:reservation-ALU-LITE})}
~
{\begin{lstlisting}
stage RR:
  new argument3 = GPR[registerY];
  new argument2 = GPR[registerZ];
  new result1 = (argument2 & 0x3FFFFFFFFFFFFFFF) | (argument3 & 0x8000000000000000);
stage E1:
  GPR[registerW] = result1;
\end{lstlisting}}
\newpage
\subsubsection[{\small FSIGNDP: Floating-Point Sign Transfer Double Word Pair}]{FSIGNDP\hfill Floating-Point Sign Transfer Double Word Pair}
\label{instr:FSIGNDP}\index{fsigndp}
 
\paragraph{Encoding} Format \textsf{ALU\_FDPMMWR} (Section~\ref{sec:format-ALU-FDPMMWR}), \verb|floatmode = 100|\\

\bitfields{ 1/P, 2/11, 1/1, 1/1, 3/100, 5/registerM, 1/--, 2/11, 3/000, 1/0, 5/registerO, 1/0, 5/registerP, 1/1 }


\paragraph{Syntax} \texttt{fsigndp} \emph{registerM} = \emph{registerP}, \emph{registerO}

\paragraph{Description} The \emph{registerP} and the \emph{registerO} are interpreted as two binary 64 floating-point numbers packed into 128 bits. The non-sign bits of the \emph{registerP} double words are concatenated with the sign bit of the \emph{registerO} double words. The two resulting double words are packed and stored into the \emph{registerM}.


\paragraph{Execution} {Reservation table \textsf{ALU\_LITE} (Section~\ref{sec:reservation-ALU-LITE})}
~
{\begin{lstlisting}
stage RR:
  new argument3 = PGR[registerO];
  new argument2 = PGR[registerP];
  for (I = 0; I < 2; I += 1) {
    result1.64[I] = (argument2.64[I] & 0x3FFFFFFFFFFFFFFF) | (argument3.64[I] & 0x8000000000000000)
  };
stage E1:
  PGR[registerM] = result1;
\end{lstlisting}}
\newpage
\subsubsection[{\small FSIGNH: Floating-Point Sign Transfer Half Word}]{FSIGNH\hfill Floating-Point Sign Transfer Half Word}
\label{instr:FSIGNH}\index{fsignh}
 
\paragraph{Encoding} Format \textsf{ALU\_FHMMWR} (Section~\ref{sec:format-ALU-FHMMWR}), \verb|floatmode = 100|\\

\bitfields{ 1/P, 2/11, 1/1, 1/1, 3/100, 6/registerW, 2/01, 3/011, 1/1, 6/registerY, 6/registerZ }


\paragraph{Syntax} \texttt{fsignh} \emph{registerW} = \emph{registerZ}, \emph{registerY}

\paragraph{Description} The \emph{registerZ} and the \emph{registerY} are interpreted as binary 16 floating-point numbers. The non-sign bits of the \emph{registerZ} are concatenated with the sign bit of the \emph{registerY}. The result is stored into the \emph{registerW}.


\paragraph{Execution} {Reservation table \textsf{ALU\_LITE} (Section~\ref{sec:reservation-ALU-LITE})}
~
{\begin{lstlisting}
stage RR:
  new argument3 = GPR[registerY];
  new argument2 = GPR[registerZ];
  new result1 = (argument2 & 0x3FFF) | (argument3 & 0x8000);
stage E1:
  GPR[registerW] = result1;
\end{lstlisting}}
\newpage
\subsubsection[{\small FSIGNHO: Floating-Point Sign Transfer Half Word Octuple}]{FSIGNHO\hfill Floating-Point Sign Transfer Half Word Octuple}
\label{instr:FSIGNHO}\index{fsignho}
 
\paragraph{Encoding} Format \textsf{ALU\_FHOMMWR} (Section~\ref{sec:format-ALU-FHOMMWR}), \verb|floatmode = 100|\\

\bitfields{ 1/P, 2/11, 1/1, 1/1, 3/100, 5/registerM, 1/--, 2/11, 3/000, 1/1, 5/registerO, 1/0, 5/registerP, 1/1 }


\paragraph{Syntax} \texttt{fsignho} \emph{registerM} = \emph{registerP}, \emph{registerO}

\paragraph{Description} The \emph{registerP} and the \emph{registerO} are interpreted as eight binary 16 floating-point numbers packed into 128 bits. The non-sign bits of the \emph{registerP} half words are concatenated with the sign bit of the \emph{registerO} half words. The eight resulting half words are packed and stored into the \emph{registerM}.


\paragraph{Execution} {Reservation table \textsf{ALU\_LITE} (Section~\ref{sec:reservation-ALU-LITE})}
~
{\begin{lstlisting}
stage RR:
  new argument3 = PGR[registerO];
  new argument2 = PGR[registerP];
  for (I = 0; I < 8; I += 1) {
    result1.16[I] = (argument2.16[I] & 0x3FFF) | (argument3.16[I] & 0x8000)
  };
stage E1:
  PGR[registerM] = result1;
\end{lstlisting}}
\newpage
\subsubsection[{\small FSIGNMD: Floating-Point Sign Combine Transfer Double Word}]{FSIGNMD\hfill Floating-Point Sign Combine Transfer Double Word}
\label{instr:FSIGNMD}\index{fsignmd}
 
\paragraph{Encoding} Format \textsf{ALU\_FDMMWR} (Section~\ref{sec:format-ALU-FDMMWR}), \verb|floatmode = 110|\\

\bitfields{ 1/P, 2/11, 1/1, 1/1, 3/110, 6/registerW, 2/00, 3/000, 1/0, 6/registerY, 6/registerZ }


\paragraph{Syntax} \texttt{fsignmd} \emph{registerW} = \emph{registerZ}, \emph{registerY}

\paragraph{Description} The \emph{registerZ} and the \emph{registerY} are interpreted as binary 64 floating-point numbers. The sign bits of the \emph{registerZ} is combined with the sign bit of the \emph{registerY} as if multiplying. The result is stored into the \emph{registerW}.


\paragraph{Execution} {Reservation table \textsf{ALU\_LITE} (Section~\ref{sec:reservation-ALU-LITE})}
~
{\begin{lstlisting}
stage RR:
  new argument3 = GPR[registerY];
  new argument2 = GPR[registerZ];
  new result1 = argument2 ^ (~argument3 & 0x8000000000000000);
stage E1:
  GPR[registerW] = result1;
\end{lstlisting}}
\newpage
\subsubsection[{\small FSIGNMDP: Floating-Point Sign Combine Transfer Double Word Pair}]{FSIGNMDP\hfill Floating-Point Sign Combine Transfer Double Word Pair}
\label{instr:FSIGNMDP}\index{fsignmdp}
 
\paragraph{Encoding} Format \textsf{ALU\_FDPMMWR} (Section~\ref{sec:format-ALU-FDPMMWR}), \verb|floatmode = 110|\\

\bitfields{ 1/P, 2/11, 1/1, 1/1, 3/110, 5/registerM, 1/--, 2/11, 3/000, 1/0, 5/registerO, 1/0, 5/registerP, 1/1 }


\paragraph{Syntax} \texttt{fsignmdp} \emph{registerM} = \emph{registerP}, \emph{registerO}

\paragraph{Description} The \emph{registerP} and the \emph{registerO} are interpreted as two binary 64 floating-point numbers packed into 128 bits. The sign bits of the \emph{registerP} double words are combined with the sign bit of the \emph{registerO} double words as if multiplying. The two resulting double words are packed and stored into the \emph{registerM}.


\paragraph{Execution} {Reservation table \textsf{ALU\_LITE} (Section~\ref{sec:reservation-ALU-LITE})}
~
{\begin{lstlisting}
stage RR:
  new argument3 = PGR[registerO];
  new argument2 = PGR[registerP];
  for (I = 0; I < 2; I += 1) {
    result1.64[I] = argument2.64[I] ^ (~argument3.64[I] & 0x8000000000000000)
  };
stage E1:
  PGR[registerM] = result1;
\end{lstlisting}}
\newpage
\subsubsection[{\small FSIGNMH: Floating-Point Sign Combine Transfer Half Word}]{FSIGNMH\hfill Floating-Point Sign Combine Transfer Half Word}
\label{instr:FSIGNMH}\index{fsignmh}
 
\paragraph{Encoding} Format \textsf{ALU\_FHMMWR} (Section~\ref{sec:format-ALU-FHMMWR}), \verb|floatmode = 110|\\

\bitfields{ 1/P, 2/11, 1/1, 1/1, 3/110, 6/registerW, 2/01, 3/011, 1/1, 6/registerY, 6/registerZ }


\paragraph{Syntax} \texttt{fsignmh} \emph{registerW} = \emph{registerZ}, \emph{registerY}

\paragraph{Description} The \emph{registerZ} and the \emph{registerY} are interpreted as binary 16 floating-point numbers. The sign bits of the \emph{registerZ} is combined with the sign bit of the \emph{registerY} as if multiplying. The result is stored into the \emph{registerW}.


\paragraph{Execution} {Reservation table \textsf{ALU\_LITE} (Section~\ref{sec:reservation-ALU-LITE})}
~
{\begin{lstlisting}
stage RR:
  new argument3 = GPR[registerY];
  new argument2 = GPR[registerZ];
  new result1 = argument2 ^ (~argument3 & 0x8000);
stage E1:
  GPR[registerW] = result1;
\end{lstlisting}}
\newpage
\subsubsection[{\small FSIGNMHO: Floating-Point Sign Combine Transfer Half Word Octuple}]{FSIGNMHO\hfill Floating-Point Sign Combine Transfer Half Word Octuple}
\label{instr:FSIGNMHO}\index{fsignmho}
 
\paragraph{Encoding} Format \textsf{ALU\_FHOMMWR} (Section~\ref{sec:format-ALU-FHOMMWR}), \verb|floatmode = 110|\\

\bitfields{ 1/P, 2/11, 1/1, 1/1, 3/110, 5/registerM, 1/--, 2/11, 3/000, 1/1, 5/registerO, 1/0, 5/registerP, 1/1 }


\paragraph{Syntax} \texttt{fsignmho} \emph{registerM} = \emph{registerP}, \emph{registerO}

\paragraph{Description} The \emph{registerP} and the \emph{registerO} are interpreted as eight binary 16 floating-point numbers packed into 128 bits. The sign bits of the \emph{registerP} half words are combined with the sign bit of the \emph{registerO} half words as if multiplying. The eight resulting half words are packed and stored into the \emph{registerM}.


\paragraph{Execution} {Reservation table \textsf{ALU\_LITE} (Section~\ref{sec:reservation-ALU-LITE})}
~
{\begin{lstlisting}
stage RR:
  new argument3 = PGR[registerO];
  new argument2 = PGR[registerP];
  for (I = 0; I < 8; I += 1) {
    result1.16[I] = argument2.16[I] ^ (~argument3.16[I] & 0x8000000000000000)
  };
stage E1:
  PGR[registerM] = result1;
\end{lstlisting}}
\newpage
\subsubsection[{\small FSIGNMW: Floating-Point Sign Combine Transfer Word}]{FSIGNMW\hfill Floating-Point Sign Combine Transfer Word}
\label{instr:FSIGNMW}\index{fsignmw}
 
\paragraph{Encoding} Format \textsf{ALU\_FWMMWR} (Section~\ref{sec:format-ALU-FWMMWR}), \verb|floatmode = 110|\\

\bitfields{ 1/P, 2/11, 1/1, 1/1, 3/110, 6/registerW, 2/01, 3/000, 1/0, 6/registerY, 6/registerZ }


\paragraph{Syntax} \texttt{fsignmw} \emph{registerW} = \emph{registerZ}, \emph{registerY}

\paragraph{Description} The \emph{registerZ} and the \emph{registerY} are interpreted as binary 32 floating-point numbers. The sign bits of the \emph{registerZ} is combined with the sign bit of the \emph{registerY} as if multiplying. The result is stored into the \emph{registerW}.


\paragraph{Execution} {Reservation table \textsf{ALU\_LITE} (Section~\ref{sec:reservation-ALU-LITE})}
~
{\begin{lstlisting}
stage RR:
  new argument3 = GPR[registerY];
  new argument2 = GPR[registerZ];
  new result1 = argument2 ^ (~argument3 & 0x80000000);
stage E1:
  GPR[registerW] = result1;
\end{lstlisting}}
\newpage
\subsubsection[{\small FSIGNMWQ: Floating-Point Sign Combine Transfer Double Word Quadruple}]{FSIGNMWQ\hfill Floating-Point Sign Combine Transfer Double Word Quadruple}
\label{instr:FSIGNMWQ}\index{fsignmwq}
 
\paragraph{Encoding} Format \textsf{ALU\_FWQMMWR} (Section~\ref{sec:format-ALU-FWQMMWR}), \verb|floatmode = 110|\\

\bitfields{ 1/P, 2/11, 1/1, 1/1, 3/110, 5/registerM, 1/--, 2/11, 3/000, 1/1, 5/registerO, 1/0, 5/registerP, 1/0 }


\paragraph{Syntax} \texttt{fsignmwq} \emph{registerM} = \emph{registerP}, \emph{registerO}

\paragraph{Description} The \emph{registerP} and the \emph{registerO} are interpreted as four binary 32 floating-point numbers packed into 128 bits. The sign bits of the \emph{registerP} words are combined with the sign bit of the \emph{registerO} words as if multiplying. The four resulting words are packed and stored into the \emph{registerM}.


\paragraph{Execution} {Reservation table \textsf{ALU\_LITE} (Section~\ref{sec:reservation-ALU-LITE})}
~
{\begin{lstlisting}
stage RR:
  new argument3 = PGR[registerO];
  new argument2 = PGR[registerP];
  for (I = 0; I < 4; I += 1) {
    result1.32[I] = argument2.32[I] ^ (~argument3.32[I] & 0x80000000)
  };
stage E1:
  PGR[registerM] = result1;
\end{lstlisting}}
\newpage
\subsubsection[{\small FSIGNND: Floating-Point Sign Negate Transfer Double Word}]{FSIGNND\hfill Floating-Point Sign Negate Transfer Double Word}
\label{instr:FSIGNND}\index{fsignnd}
 
\paragraph{Encoding} Format \textsf{ALU\_FDMMWR} (Section~\ref{sec:format-ALU-FDMMWR}), \verb|floatmode = 101|\\

\bitfields{ 1/P, 2/11, 1/1, 1/1, 3/101, 6/registerW, 2/00, 3/000, 1/0, 6/registerY, 6/registerZ }


\paragraph{Syntax} \texttt{fsignnd} \emph{registerW} = \emph{registerZ}, \emph{registerY}

\paragraph{Description} The \emph{registerZ} and the \emph{registerY} are interpreted as binary 64 floating-point numbers. The non-sign bits of the \emph{registerZ} are concatenated with the negated sign bit of the \emph{registerY}. The result is stored into the \emph{registerW}.


\paragraph{Execution} {Reservation table \textsf{ALU\_LITE} (Section~\ref{sec:reservation-ALU-LITE})}
~
{\begin{lstlisting}
stage RR:
  new argument3 = GPR[registerY];
  new argument2 = GPR[registerZ];
  new result1 = (argument2 & 0x3FFFFFFFFFFFFFFF) | (~argument3 & 0x8000000000000000);
stage E1:
  GPR[registerW] = result1;
\end{lstlisting}}
\newpage
\subsubsection[{\small FSIGNNDP: Floating-Point Sign Negate Transfer Double Word Pair}]{FSIGNNDP\hfill Floating-Point Sign Negate Transfer Double Word Pair}
\label{instr:FSIGNNDP}\index{fsignndp}
 
\paragraph{Encoding} Format \textsf{ALU\_FDPMMWR} (Section~\ref{sec:format-ALU-FDPMMWR}), \verb|floatmode = 101|\\

\bitfields{ 1/P, 2/11, 1/1, 1/1, 3/101, 5/registerM, 1/--, 2/11, 3/000, 1/0, 5/registerO, 1/0, 5/registerP, 1/1 }


\paragraph{Syntax} \texttt{fsignndp} \emph{registerM} = \emph{registerP}, \emph{registerO}

\paragraph{Description} The \emph{registerP} and the \emph{registerO} are interpreted as two binary 64 floating-point numbers packed into 128 bits. The non-sign bits of the \emph{registerP} double words are concatenated with the negated sign bit of the \emph{registerO} double words. The two resulting double words are packed and stored into the \emph{registerM}.


\paragraph{Execution} {Reservation table \textsf{ALU\_LITE} (Section~\ref{sec:reservation-ALU-LITE})}
~
{\begin{lstlisting}
stage RR:
  new argument3 = PGR[registerO];
  new argument2 = PGR[registerP];
  for (I = 0; I < 2; I += 1) {
    result1.64[I] = (argument2.64[I] & 0x3FFFFFFFFFFFFFFF) | (~argument3.64[I] & 0x8000000000000000)
  };
stage E1:
  PGR[registerM] = result1;
\end{lstlisting}}
\newpage
\subsubsection[{\small FSIGNNH: Floating-Point Sign Negate Transfer Half Word}]{FSIGNNH\hfill Floating-Point Sign Negate Transfer Half Word}
\label{instr:FSIGNNH}\index{fsignnh}
 
\paragraph{Encoding} Format \textsf{ALU\_FHMMWR} (Section~\ref{sec:format-ALU-FHMMWR}), \verb|floatmode = 101|\\

\bitfields{ 1/P, 2/11, 1/1, 1/1, 3/101, 6/registerW, 2/01, 3/011, 1/1, 6/registerY, 6/registerZ }


\paragraph{Syntax} \texttt{fsignnh} \emph{registerW} = \emph{registerZ}, \emph{registerY}

\paragraph{Description} The \emph{registerZ} and the \emph{registerY} are interpreted as binary 16 floating-point numbers. The non-sign bits of the \emph{registerZ} are concatenated with the negated sign bit of the \emph{registerY}. The result is stored into the \emph{registerW}.


\paragraph{Execution} {Reservation table \textsf{ALU\_LITE} (Section~\ref{sec:reservation-ALU-LITE})}
~
{\begin{lstlisting}
stage RR:
  new argument3 = GPR[registerY];
  new argument2 = GPR[registerZ];
  new result1 = (argument2 & 0x3FFF) | (~argument3 & 0x8000);
stage E1:
  GPR[registerW] = result1;
\end{lstlisting}}
\newpage
\subsubsection[{\small FSIGNNHO: Floating-Point Sign Negate Transfer Half Word Octuple}]{FSIGNNHO\hfill Floating-Point Sign Negate Transfer Half Word Octuple}
\label{instr:FSIGNNHO}\index{fsignnho}
 
\paragraph{Encoding} Format \textsf{ALU\_FHOMMWR} (Section~\ref{sec:format-ALU-FHOMMWR}), \verb|floatmode = 101|\\

\bitfields{ 1/P, 2/11, 1/1, 1/1, 3/101, 5/registerM, 1/--, 2/11, 3/000, 1/1, 5/registerO, 1/0, 5/registerP, 1/1 }


\paragraph{Syntax} \texttt{fsignnho} \emph{registerM} = \emph{registerP}, \emph{registerO}

\paragraph{Description} The \emph{registerP} and the \emph{registerO} are interpreted as eight binary 16 floating-point numbers packed into 128 bits. The non-sign bits of the \emph{registerP} half words are concatenated with the negated sign bit of the \emph{registerO} half words. The eight resulting half words are packed and stored into the \emph{registerM}.


\paragraph{Execution} {Reservation table \textsf{ALU\_LITE} (Section~\ref{sec:reservation-ALU-LITE})}
~
{\begin{lstlisting}
stage RR:
  new argument3 = PGR[registerO];
  new argument2 = PGR[registerP];
  for (I = 0; I < 8; I += 1) {
    result1.16[I] = (argument2.16[I] & 0x3FFF) | (~argument3.16[I] & 0x8000)
  };
stage E1:
  PGR[registerM] = result1;
\end{lstlisting}}
\newpage
\subsubsection[{\small FSIGNNW: Floating-Point Sign Negate Transfer Word}]{FSIGNNW\hfill Floating-Point Sign Negate Transfer Word}
\label{instr:FSIGNNW}\index{fsignnw}
 
\paragraph{Encoding} Format \textsf{ALU\_FWMMWR} (Section~\ref{sec:format-ALU-FWMMWR}), \verb|floatmode = 101|\\

\bitfields{ 1/P, 2/11, 1/1, 1/1, 3/101, 6/registerW, 2/01, 3/000, 1/0, 6/registerY, 6/registerZ }


\paragraph{Syntax} \texttt{fsignnw} \emph{registerW} = \emph{registerZ}, \emph{registerY}

\paragraph{Description} The \emph{registerZ} and the \emph{registerY} are interpreted as binary 32 floating-point numbers. The non-sign bits of the \emph{registerZ} are concatenated with the negated sign bit of the \emph{registerY}. The result is stored into the \emph{registerW}.


\paragraph{Execution} {Reservation table \textsf{ALU\_LITE} (Section~\ref{sec:reservation-ALU-LITE})}
~
{\begin{lstlisting}
stage RR:
  new argument3 = GPR[registerY];
  new argument2 = GPR[registerZ];
  new result1 = (argument2 & 0x3FFFFFFF) | (~argument3 & 0x80000000);
stage E1:
  GPR[registerW] = result1;
\end{lstlisting}}
\newpage
\subsubsection[{\small FSIGNNWQ: Floating-Point Sign Negate Transfer Word Quadruple}]{FSIGNNWQ\hfill Floating-Point Sign Negate Transfer Word Quadruple}
\label{instr:FSIGNNWQ}\index{fsignnwq}
 
\paragraph{Encoding} Format \textsf{ALU\_FWQMMWR} (Section~\ref{sec:format-ALU-FWQMMWR}), \verb|floatmode = 101|\\

\bitfields{ 1/P, 2/11, 1/1, 1/1, 3/101, 5/registerM, 1/--, 2/11, 3/000, 1/1, 5/registerO, 1/0, 5/registerP, 1/0 }


\paragraph{Syntax} \texttt{fsignnwq} \emph{registerM} = \emph{registerP}, \emph{registerO}

\paragraph{Description} The \emph{registerP} and the \emph{registerO} are interpreted as four binary 32 floating-point numbers packed into 128 bits. The non-sign bits of the \emph{registerP} words are concatenated with the negated sign bit of the \emph{registerO} words. The four resulting words are packed and stored into the \emph{registerM}.


\paragraph{Execution} {Reservation table \textsf{ALU\_LITE} (Section~\ref{sec:reservation-ALU-LITE})}
~
{\begin{lstlisting}
stage RR:
  new argument3 = PGR[registerO];
  new argument2 = PGR[registerP];
  for (I = 0; I < 4; I += 1) {
    result1.32[I] = (argument2.32[I] & 0x3FFFFFFF) | (~argument3.32[I] & 0x80000000)
  };
stage E1:
  PGR[registerM] = result1;
\end{lstlisting}}
\newpage
\subsubsection[{\small FSIGNW: Floating-Point Sign Transfer Word}]{FSIGNW\hfill Floating-Point Sign Transfer Word}
\label{instr:FSIGNW}\index{fsignw}
 
\paragraph{Encoding} Format \textsf{ALU\_FWMMWR} (Section~\ref{sec:format-ALU-FWMMWR}), \verb|floatmode = 100|\\

\bitfields{ 1/P, 2/11, 1/1, 1/1, 3/100, 6/registerW, 2/01, 3/000, 1/0, 6/registerY, 6/registerZ }


\paragraph{Syntax} \texttt{fsignw} \emph{registerW} = \emph{registerZ}, \emph{registerY}

\paragraph{Description} The \emph{registerZ} and the \emph{registerY} are interpreted as binary 32 floating-point numbers. The non-sign bits of the \emph{registerZ} are concatenated with the sign bit of the \emph{registerY}. The result is stored into the \emph{registerW}.


\paragraph{Execution} {Reservation table \textsf{ALU\_LITE} (Section~\ref{sec:reservation-ALU-LITE})}
~
{\begin{lstlisting}
stage RR:
  new argument3 = GPR[registerY];
  new argument2 = GPR[registerZ];
  new result1 = (argument2 & 0x3FFFFFFF) | (argument3 & 0x80000000);
stage E1:
  GPR[registerW] = result1;
\end{lstlisting}}
\newpage
\subsubsection[{\small FSIGNWQ: Floating-Point Sign Transfer Word Quadruple}]{FSIGNWQ\hfill Floating-Point Sign Transfer Word Quadruple}
\label{instr:FSIGNWQ}\index{fsignwq}
 
\paragraph{Encoding} Format \textsf{ALU\_FWQMMWR} (Section~\ref{sec:format-ALU-FWQMMWR}), \verb|floatmode = 100|\\

\bitfields{ 1/P, 2/11, 1/1, 1/1, 3/100, 5/registerM, 1/--, 2/11, 3/000, 1/1, 5/registerO, 1/0, 5/registerP, 1/0 }


\paragraph{Syntax} \texttt{fsignwq} \emph{registerM} = \emph{registerP}, \emph{registerO}

\paragraph{Description} The \emph{registerP} and the \emph{registerO} are interpreted as four binary 32 floating-point numbers packed into 128 bits. The non-sign bits of the \emph{registerP} words are concatenated with the sign bit of the \emph{registerO} words. The four resulting words are packed and stored into the \emph{registerM}.


\paragraph{Execution} {Reservation table \textsf{ALU\_LITE} (Section~\ref{sec:reservation-ALU-LITE})}
~
{\begin{lstlisting}
stage RR:
  new argument3 = PGR[registerO];
  new argument2 = PGR[registerP];
  for (I = 0; I < 4; I += 1) {
    result1.32[I] = (argument2.32[I] & 0x3FFFFFFF) | (argument3.32[I] & 0x80000000)
  };
stage E1:
  PGR[registerM] = result1;
\end{lstlisting}}
\newpage
\subsubsection[{\small FSQRTD: Floating Point Square Root Double Word}]{FSQRTD\hfill Floating Point Square Root Double Word}
\label{instr:FSQRTD}\index{fsqrtd}
 
\paragraph{Encoding} Format \textsf{ALU\_FDMO0} (Section~\ref{sec:format-ALU-FDMO0}), \verb|fpucode6 = 1000|\\

\bitfields{ 1/P, 2/11, 1/1, 1/1, 3/floatmode, 6/registerW, 2/00, 3/001, 1/0, 4/1000, 1/--, 1/0, 6/registerZ }


\paragraph{Syntax} \texttt{fsqrtd}\emph{floatmode} \emph{registerW} = \emph{registerZ}

\paragraph{Description} The square root of the \emph{registerZ} interpreted as a 64-bit floating-point number is computed according to the rounding mode and stored into the \emph{registerW}. This instruction may raise exception bits in the CS register.


\paragraph{Modifier(s)}
\noindent\begin{itemize}
\item floatmode (cf. floatmode in section \ref{par:modifier-floatmode} on page \pageref{par:modifier-floatmode})
\end{itemize}

\paragraph{Execution} {Reservation table \textsf{ALU\_FULL} (Section~\ref{sec:reservation-ALU-FULL})}
~
{\begin{lstlisting}
stage ID:
  new floatmode = _ZX_3(floatmode);
stage RR:
  new argument2 = GPR[registerZ];
  new RM = decode_riscv_float_rounding_mode(floatmode);
  new result1 = f64_sqrt(RM, argument2);
stage E4:
  GPR[registerW] = result1;
  commit_float_exception_flags();
\end{lstlisting}}
\newpage
\subsubsection[{\small FSQRTH: Floating Point Square Root Half Word}]{FSQRTH\hfill Floating Point Square Root Half Word}
\label{instr:FSQRTH}\index{fsqrth}
 
\paragraph{Encoding} Format \textsf{ALU\_FHMO0} (Section~\ref{sec:format-ALU-FHMO0}), \verb|fpucode6 = 1000|\\

\bitfields{ 1/P, 2/11, 1/1, 1/1, 3/floatmode, 6/registerW, 2/01, 3/001, 1/1, 4/1000, 1/--, 1/0, 6/registerZ }


\paragraph{Syntax} \texttt{fsqrth}\emph{floatmode} \emph{registerW} = \emph{registerZ}

\paragraph{Description} The square root of the \emph{registerZ} interpreted as a 16-bit floating-point number is computed according to the rounding mode and stored into the \emph{registerW}. This instruction may raise exception bits in the CS register.


\paragraph{Modifier(s)}
\noindent\begin{itemize}
\item floatmode (cf. floatmode in section \ref{par:modifier-floatmode} on page \pageref{par:modifier-floatmode})
\end{itemize}

\paragraph{Execution} {Reservation table \textsf{ALU\_FULL} (Section~\ref{sec:reservation-ALU-FULL})}
~
{\begin{lstlisting}
stage ID:
  new floatmode = _ZX_3(floatmode);
stage RR:
  new argument2 = GPR[registerZ];
  new RM = decode_riscv_float_rounding_mode(floatmode);
  new result1 = f16_sqrt(RM, argument2);
stage SF:
  GPR[registerW] = result1;
  commit_float_exception_flags();
\end{lstlisting}}
\newpage
\subsubsection[{\small FSQRTW: Floating Point Square Root Word}]{FSQRTW\hfill Floating Point Square Root Word}
\label{instr:FSQRTW}\index{fsqrtw}
 
\paragraph{Encoding} Format \textsf{ALU\_FWMO0} (Section~\ref{sec:format-ALU-FWMO0}), \verb|fpucode6 = 1000|\\

\bitfields{ 1/P, 2/11, 1/1, 1/1, 3/floatmode, 6/registerW, 2/01, 3/001, 1/0, 4/1000, 1/--, 1/0, 6/registerZ }


\paragraph{Syntax} \texttt{fsqrtw}\emph{floatmode} \emph{registerW} = \emph{registerZ}

\paragraph{Description} The square root of the \emph{registerZ} interpreted as a 32-bit floating-point number is computed according to the rounding mode and stored into the \emph{registerW}. This instruction may raise exception bits in the CS register.


\paragraph{Modifier(s)}
\noindent\begin{itemize}
\item floatmode (cf. floatmode in section \ref{par:modifier-floatmode} on page \pageref{par:modifier-floatmode})
\end{itemize}

\paragraph{Execution} {Reservation table \textsf{ALU\_FULL} (Section~\ref{sec:reservation-ALU-FULL})}
~
{\begin{lstlisting}
stage ID:
  new floatmode = _ZX_3(floatmode);
stage RR:
  new argument2 = GPR[registerZ];
  new RM = decode_riscv_float_rounding_mode(floatmode);
  new result1 = f32_sqrt(RM, argument2);
stage SF:
  GPR[registerW] = result1;
  commit_float_exception_flags();
\end{lstlisting}}
\newpage
\subsubsection[{\small FSRECD: Floating-Point Seed for Reciprocal Double Word}]{FSRECD\hfill Floating-Point Seed for Reciprocal Double Word}
\label{instr:FSRECD}\index{fsrecd}
 
\paragraph{Encoding} Format \textsf{ALU\_FDMO1} (Section~\ref{sec:format-ALU-FDMO1}), \verb|fpucode6 = 0100|\\

\bitfields{ 1/P, 2/11, 1/1, 1/1, 3/111, 6/registerW, 2/00, 3/000, 1/0, 4/0100, 1/--, 1/0, 6/registerZ }


\paragraph{Syntax} \texttt{fsrecd} \emph{registerW} = \emph{registerZ}

\paragraph{Description} An approximation of the reciprocal of the \emph{registerZ}, assuming binary 64 floating-point numbers, is stored into the \emph{registerW}. This instruction may raise exception bits in the CS register.


\paragraph{Execution} {Reservation table \textsf{ALU\_LITE} (Section~\ref{sec:reservation-ALU-LITE})}
~
{\begin{lstlisting}
stage RR:
  new argument2 = GPR[registerZ];
  new result1 = f64_fast_rec(argument2);
stage E1:
  GPR[registerW] = result1;
\end{lstlisting}}
\newpage
\subsubsection[{\small FSRECW: Floating-Point Seed for Reciprocal Word}]{FSRECW\hfill Floating-Point Seed for Reciprocal Word}
\label{instr:FSRECW}\index{fsrecw}
 
\paragraph{Encoding} Format \textsf{ALU\_FWMO1} (Section~\ref{sec:format-ALU-FWMO1}), \verb|fpucode6 = 0100|\\

\bitfields{ 1/P, 2/11, 1/1, 1/1, 3/111, 6/registerW, 2/01, 3/000, 1/0, 4/0100, 1/--, 1/0, 6/registerZ }


\paragraph{Syntax} \texttt{fsrecw} \emph{registerW} = \emph{registerZ}

\paragraph{Description} An approximation of the reciprocal of the \emph{registerZ}, assuming binary 32 floating-point numbers, is stored into the \emph{registerW}. This instruction may raise exception bits in the CS register.


\paragraph{Execution} {Reservation table \textsf{ALU\_LITE} (Section~\ref{sec:reservation-ALU-LITE})}
~
{\begin{lstlisting}
stage RR:
  new argument2 = GPR[registerZ];
  new result1 = f32_fast_rec(argument2);
stage E1:
  GPR[registerW] = result1;
\end{lstlisting}}
\newpage
\subsubsection[{\small FSRECWP: Floating-Point Seed for Reciprocal Word Pair}]{FSRECWP\hfill Floating-Point Seed for Reciprocal Word Pair}
\label{instr:FSRECWP}\index{fsrecwp}
 
\paragraph{Encoding} Format \textsf{ALU\_FWPMO1} (Section~\ref{sec:format-ALU-FWPMO1}), \verb|fpucode6 = 0100|\\

\bitfields{ 1/P, 2/11, 1/1, 1/1, 3/111, 6/registerW, 2/00, 3/000, 1/0, 4/0100, 1/--, 1/1, 6/registerZ }


\paragraph{Syntax} \texttt{fsrecwp} \emph{registerW} = \emph{registerZ}

\paragraph{Description} An approximation of the reciprocal of the \emph{registerZ}, assuming a pair of binary 32 floating-point numbers, is stored into the \emph{registerW}. This instruction may raise exception bits in the CS register.


\paragraph{Execution} {Reservation table \textsf{ALU\_LITE} (Section~\ref{sec:reservation-ALU-LITE})}
~
{\begin{lstlisting}
stage RR:
  new argument2 = GPR[registerZ];
  for (I = 0; I < 2; I += 1) {
    result1.32[I] = f32_fast_rec(argument2.32[I])
  };
stage E1:
  GPR[registerW] = result1;
\end{lstlisting}}
\newpage
\subsubsection[{\small FSRECWQ: Floating-Point Seed for Reciprocal Word Quadruple}]{FSRECWQ\hfill Floating-Point Seed for Reciprocal Word Quadruple}
\label{instr:FSRECWQ}\index{fsrecwq}
 
\paragraph{Encoding} Format \textsf{ALU\_FWQMO1} (Section~\ref{sec:format-ALU-FWQMO1}), \verb|fpucode6 = 0100|\\

\bitfields{ 1/P, 2/11, 1/1, 1/1, 3/111, 5/registerM, 1/--, 2/11, 3/000, 1/1, 4/0100, 1/--, 1/0, 5/registerP, 1/0 }


\paragraph{Syntax} \texttt{fsrecwq} \emph{registerM} = \emph{registerP}

\paragraph{Description} An approximation of the reciprocal of the \emph{registerP}, assuming a quadruple of binary 32 floating-point numbers, is stored into the \emph{registerM}. This instruction may raise exception bits in the CS register.


\paragraph{Execution} {Reservation table \textsf{ALU\_LITE} (Section~\ref{sec:reservation-ALU-LITE})}
~
{\begin{lstlisting}
stage RR:
  new argument2 = PGR[registerP];
  for (I = 0; I < 4; I += 1) {
    result1.32[I] = f32_fast_rec(argument2.32[I])
  };
stage E1:
  PGR[registerM] = result1;
\end{lstlisting}}
\newpage
\subsubsection[{\small FSRSRD: Floating-Point Seed for Reciprocal Square Root Double Word}]{FSRSRD\hfill Floating-Point Seed for Reciprocal Square Root Double Word}
\label{instr:FSRSRD}\index{fsrsrd}
 
\paragraph{Encoding} Format \textsf{ALU\_FDMO1} (Section~\ref{sec:format-ALU-FDMO1}), \verb|fpucode6 = 0101|\\

\bitfields{ 1/P, 2/11, 1/1, 1/1, 3/111, 6/registerW, 2/00, 3/000, 1/0, 4/0101, 1/--, 1/0, 6/registerZ }


\paragraph{Syntax} \texttt{fsrsrd} \emph{registerW} = \emph{registerZ}

\paragraph{Description} An approximation of the reciprocal of square root of the \emph{registerZ}, assuming binary 64 floating-point numbers, is stored into the \emph{registerW}.


\paragraph{Execution} {Reservation table \textsf{ALU\_LITE} (Section~\ref{sec:reservation-ALU-LITE})}
~
{\begin{lstlisting}
stage RR:
  new argument2 = GPR[registerZ];
  new result1 = f64_fast_rsqrt(argument2);
stage E1:
  GPR[registerW] = result1;
\end{lstlisting}}
\newpage
\subsubsection[{\small FSRSRW: Floating-Point Seed for Reciprocal Square Root Word}]{FSRSRW\hfill Floating-Point Seed for Reciprocal Square Root Word}
\label{instr:FSRSRW}\index{fsrsrw}
 
\paragraph{Encoding} Format \textsf{ALU\_FWMO1} (Section~\ref{sec:format-ALU-FWMO1}), \verb|fpucode6 = 0101|\\

\bitfields{ 1/P, 2/11, 1/1, 1/1, 3/111, 6/registerW, 2/01, 3/000, 1/0, 4/0101, 1/--, 1/0, 6/registerZ }


\paragraph{Syntax} \texttt{fsrsrw} \emph{registerW} = \emph{registerZ}

\paragraph{Description} An approximation of the reciprocal of square root of the \emph{registerZ}, assuming binary 32 floating-point numbers, is stored into the \emph{registerW}.


\paragraph{Execution} {Reservation table \textsf{ALU\_LITE} (Section~\ref{sec:reservation-ALU-LITE})}
~
{\begin{lstlisting}
stage RR:
  new argument2 = GPR[registerZ];
  new result1 = f32_fast_rsqrt(argument2);
stage E1:
  GPR[registerW] = result1;
\end{lstlisting}}
\newpage
\subsubsection[{\small FSRSRWP: Floating-Point Seed for Reciprocal Square Root Word Pair}]{FSRSRWP\hfill Floating-Point Seed for Reciprocal Square Root Word Pair}
\label{instr:FSRSRWP}\index{fsrsrwp}
 
\paragraph{Encoding} Format \textsf{ALU\_FWPMO1} (Section~\ref{sec:format-ALU-FWPMO1}), \verb|fpucode6 = 0101|\\

\bitfields{ 1/P, 2/11, 1/1, 1/1, 3/111, 6/registerW, 2/00, 3/000, 1/0, 4/0101, 1/--, 1/1, 6/registerZ }


\paragraph{Syntax} \texttt{fsrsrwp} \emph{registerW} = \emph{registerZ}

\paragraph{Description} An approximation of the reciprocal of square root of the \emph{registerZ}, assuming a pair of binary 32 floating-point numbers, is stored into the \emph{registerW}.


\paragraph{Execution} {Reservation table \textsf{ALU\_LITE} (Section~\ref{sec:reservation-ALU-LITE})}
~
{\begin{lstlisting}
stage RR:
  new argument2 = GPR[registerZ];
  for (I = 0; I < 2; I += 1) {
    result1.32[I] = f32_fast_rsqrt(argument2.32[I])
  };
stage E1:
  GPR[registerW] = result1;
\end{lstlisting}}
\newpage
\subsubsection[{\small FSRSRWQ: Floating-Point Seed for Reciprocal Square Root Word Quadruple}]{FSRSRWQ\hfill Floating-Point Seed for Reciprocal Square Root Word Quadruple}
\label{instr:FSRSRWQ}\index{fsrsrwq}
 
\paragraph{Encoding} Format \textsf{ALU\_FWQMO1} (Section~\ref{sec:format-ALU-FWQMO1}), \verb|fpucode6 = 0101|\\

\bitfields{ 1/P, 2/11, 1/1, 1/1, 3/111, 5/registerM, 1/--, 2/11, 3/000, 1/1, 4/0101, 1/--, 1/0, 5/registerP, 1/0 }


\paragraph{Syntax} \texttt{fsrsrwq} \emph{registerM} = \emph{registerP}

\paragraph{Description} An approximation of the reciprocal of square root of the \emph{registerP}, assuming a quadruple of binary 32 floating-point numbers, is stored into the \emph{registerM}.


\paragraph{Execution} {Reservation table \textsf{ALU\_LITE} (Section~\ref{sec:reservation-ALU-LITE})}
~
{\begin{lstlisting}
stage RR:
  new argument2 = PGR[registerP];
  for (I = 0; I < 4; I += 1) {
    result1.32[I] = f32_fast_rsqrt(argument2.32[I])
  };
stage E1:
  PGR[registerM] = result1;
\end{lstlisting}}
\newpage
\subsubsection[{\small FWIDENHW: Floating Point Widen Half Word to Word}]{FWIDENHW\hfill Floating Point Widen Half Word to Word}
\label{instr:FWIDENHW}\index{fwidenhw}
 
\paragraph{Encoding} Format \textsf{ALU\_FWWIDEN} (Section~\ref{sec:format-ALU-FWWIDEN}), \verb|fpucode6 = 1000|\\

\bitfields{ 1/P, 2/11, 1/1, 1/1, 3/111, 6/registerW, 2/01, 3/000, 1/0, 4/1000, 1/mostsig, 1/0, 6/registerZ }


\paragraph{Syntax} \texttt{fwidenhw}\emph{mostsig} \emph{registerW} = \emph{registerZ}

\paragraph{Description} The \emph{registerZ} is interpreted as a binary 16 floating-point number. The least or most significant half word of the \emph{registerZ} word is selected, depending on \emph{mostsig}. The number is converted to a binary 32 floating-point number and stored into the \emph{registerW}.


\paragraph{Modifier(s)}
\noindent\begin{itemize}
\item mostsig (cf. mostsig in section \ref{par:modifier-mostsig} on page \pageref{par:modifier-mostsig})
\end{itemize}

\paragraph{Execution} {Reservation table \textsf{ALU\_LITE} (Section~\ref{sec:reservation-ALU-LITE})}
~
{\begin{lstlisting}
stage ID:
  new mostsig = _ZX_1(mostsig);
stage RR:
  new argument2 = GPR[registerZ];
  if (mostsig) {
    argument2 = argument2 >> 16;
  }
  new result1 = f16_to_f32(argument2.16[0]);
stage E1:
  GPR[registerW] = result1;
\end{lstlisting}}
\newpage
\subsubsection[{\small FWIDENHWQ: Floating Point Widen Half Word to Word Quadruple}]{FWIDENHWQ\hfill Floating Point Widen Half Word to Word Quadruple}
\label{instr:FWIDENHWQ}\index{fwidenhwq}
 
\paragraph{Encoding} Format \textsf{ALU\_FWQWIDEWR} (Section~\ref{sec:format-ALU-FWQWIDEWR}), \verb|fpucode6 = 1000|\\

\bitfields{ 1/P, 2/11, 1/1, 1/1, 3/111, 5/registerM, 1/--, 2/11, 3/000, 1/1, 4/1000, 1/mostsig, 1/0, 5/registerP, 1/0 }


\paragraph{Syntax} \texttt{fwidenhwq}\emph{mostsig} \emph{registerM} = \emph{registerP}

\paragraph{Description} The \emph{registerP} interpreted as eight binary 16 floating-point numbers packed into 128 bits. The least or most significant half words of the \emph{registerP} are selected, depending on \emph{mostsig}. The four selected half words are converted to binary 32 floating-point numbers, packed and stored into the \emph{registerM}. This instruction may raise invalid exception bit in the CS register.


\paragraph{Modifier(s)}
\noindent\begin{itemize}
\item mostsig (cf. mostsig in section \ref{par:modifier-mostsig} on page \pageref{par:modifier-mostsig})
\end{itemize}

\paragraph{Execution} {Reservation table \textsf{ALU\_LITE} (Section~\ref{sec:reservation-ALU-LITE})}
~
{\begin{lstlisting}
stage ID:
  new mostsig = _ZX_1(mostsig);
stage RR:
  new argument2 = PGR[registerP];
  if (mostsig) {
    argument2 = argument2 >> 64;
  }
  for (I = 0; I < 4; I += 1) {
    result1.32[I] = f16_to_f32(_ZX_16(argument2.16[I]))
  };
stage E1:
  PGR[registerM] = result1;
\end{lstlisting}}
\newpage
\subsubsection[{\small FWIDENWD: Floating Point Widen Word to Double Word}]{FWIDENWD\hfill Floating Point Widen Word to Double Word}
\label{instr:FWIDENWD}\index{fwidenwd}
 
\paragraph{Encoding} Format \textsf{ALU\_FDWIDEN} (Section~\ref{sec:format-ALU-FDWIDEN}), \verb|fpucode6 = 1000|\\

\bitfields{ 1/P, 2/11, 1/1, 1/1, 3/111, 6/registerW, 2/00, 3/000, 1/0, 4/1000, 1/mostsig, 1/0, 6/registerZ }


\paragraph{Syntax} \texttt{fwidenwd}\emph{mostsig} \emph{registerW} = \emph{registerZ}

\paragraph{Description} The \emph{registerZ} is interpreted as a binary 32 floating-point number. The least or most significant word of the \emph{registerZ} double word is selected, depending on \emph{mostsig}. The number is converted to a binary 64 floating-point number and stored into the \emph{registerW}.


\paragraph{Modifier(s)}
\noindent\begin{itemize}
\item mostsig (cf. mostsig in section \ref{par:modifier-mostsig} on page \pageref{par:modifier-mostsig})
\end{itemize}

\paragraph{Execution} {Reservation table \textsf{ALU\_LITE} (Section~\ref{sec:reservation-ALU-LITE})}
~
{\begin{lstlisting}
stage ID:
  new mostsig = _ZX_1(mostsig);
stage RR:
  new argument2 = GPR[registerZ];
  if (mostsig) {
    argument2 = argument2 >> 32;
  }
  new result1 = f32_to_f64(argument2.32[0]);
stage E1:
  GPR[registerW] = result1;
\end{lstlisting}}
\newpage
\subsection{EXT Instructions}
\label{sec:EXT-instructions}

\subsubsection[{\small XACCESSO: Access from Extension into Octuple Word}]{XACCESSO\hfill Access from Extension into Octuple Word}
\label{instr:XACCESSO}\index{xaccesso}
 
\paragraph{Encoding} Format \textsf{EXT\_XACCESSO2} (Section~\ref{sec:format-EXT-XACCESSO2}), \verb|log2size = 000|\\

\bitfields{ 1/P, 2/10, 1/0, 1/0, 3/000, 4/registerN, 2/11, 2/00, 1/1, 3/000, 5/registerCg, 1/0, 6/registerZ }


\paragraph{Syntax} \texttt{xaccesso} \emph{registerN} = \emph{registerCg}, \emph{registerZ}

\paragraph{Description} The 5 least significant bits of the \emph{registerZ} are interpreted as a byte count. The most significant bits of the \emph{registerZ} are interpreted as a vector register index into the buffer specified by the \emph{registerCg}. The vector register indexed and its successor modulo the buffer size are concatenated, and shifted right bt the byte count. The resulting 32 bytes are written to the \emph{registerN}.


\paragraph{Execution} {Reservation table \textsf{EXT\_MISC\_AUXW} (Section~\ref{sec:reservation-EXT-MISC-AUXW})}
~
{\begin{lstlisting}
stage RR:
  new mask = 1;
  new buffer = registerCg << 1;
  new address = GPR[registerZ];
  new shift = (address & 31) * 8;
  new index_0 = address >> 5;
  new index_1 = index_0 + 1;
  new where_0 = buffer + (index_0 & mask);
  new where_1 = buffer + (index_1 & mask);
  new result1_0 = XVR[where_0];
  new result1_1 = XVR[where_1];
  new result1 = result1_0;
  if (shift) {
    result1 = (_ZX_256(result1_0) >> shift) | _ZX_256(_ZX_256(result1_1) << (256 - shift));
  }
stage E1:
  CS.XMF = 1;
stage E3:
  QGR[registerN] = result1;
\end{lstlisting}}
\newpage

\paragraph{Encoding} Format \textsf{EXT\_XACCESSO4} (Section~\ref{sec:format-EXT-XACCESSO4}), \verb|log2size = 000|\\

\bitfields{ 1/P, 2/10, 1/0, 1/0, 3/000, 4/registerN, 2/11, 2/00, 1/1, 3/000, 4/registerCh, 2/01, 6/registerZ }


\paragraph{Syntax} \texttt{xaccesso} \emph{registerN} = \emph{registerCh}, \emph{registerZ}

\paragraph{Description} The 5 least significant bits of the \emph{registerZ} are interpreted as a byte count. The most significant bits of the \emph{registerZ} are interpreted as a vector register index into the buffer specified by the \emph{registerCh}. The vector register indexed and its successor modulo the buffer size are concatenated, and shifted right bt the byte count. The resulting 32 bytes are written to the \emph{registerN}.


\paragraph{Execution} {Reservation table \textsf{EXT\_MISC\_AUXW} (Section~\ref{sec:reservation-EXT-MISC-AUXW})}
~
{\begin{lstlisting}
stage RR:
  new mask = 3;
  new buffer = registerCh << 2;
  new address = GPR[registerZ];
  new shift = (address & 31) * 8;
  new index_0 = address >> 5;
  new index_1 = index_0 + 1;
  new where_0 = buffer + (index_0 & mask);
  new where_1 = buffer + (index_1 & mask);
  new result1_0 = XVR[where_0];
  new result1_1 = XVR[where_1];
  new result1 = result1_0;
  if (shift) {
    result1 = (_ZX_256(result1_0) >> shift) | _ZX_256(_ZX_256(result1_1) << (256 - shift));
  }
stage E1:
  CS.XMF = 1;
stage E3:
  QGR[registerN] = result1;
\end{lstlisting}}
\newpage

\paragraph{Encoding} Format \textsf{EXT\_XACCESSO8} (Section~\ref{sec:format-EXT-XACCESSO8}), \verb|log2size = 000|\\

\bitfields{ 1/P, 2/10, 1/0, 1/0, 3/000, 4/registerN, 2/11, 2/00, 1/1, 3/000, 3/registerCi, 3/011, 6/registerZ }


\paragraph{Syntax} \texttt{xaccesso} \emph{registerN} = \emph{registerCi}, \emph{registerZ}

\paragraph{Description} The 5 least significant bits of the \emph{registerZ} are interpreted as a byte count. The most significant bits of the \emph{registerZ} are interpreted as a vector register index into the buffer specified by the \emph{registerCi}. The vector register indexed and its successor modulo the buffer size are concatenated, and shifted right bt the byte count. The resulting 32 bytes are written to the \emph{registerN}.


\paragraph{Execution} {Reservation table \textsf{EXT\_MISC\_AUXW} (Section~\ref{sec:reservation-EXT-MISC-AUXW})}
~
{\begin{lstlisting}
stage RR:
  new mask = 7;
  new buffer = registerCi << 3;
  new address = GPR[registerZ];
  new shift = (address & 31) * 8;
  new index_0 = address >> 5;
  new index_1 = index_0 + 1;
  new where_0 = buffer + (index_0 & mask);
  new where_1 = buffer + (index_1 & mask);
  new result1_0 = XVR[where_0];
  new result1_1 = XVR[where_1];
  new result1 = result1_0;
  if (shift) {
    result1 = (_ZX_256(result1_0) >> shift) | _ZX_256(_ZX_256(result1_1) << (256 - shift));
  }
stage E1:
  CS.XMF = 1;
stage E3:
  QGR[registerN] = result1;
\end{lstlisting}}
\newpage

\paragraph{Encoding} Format \textsf{EXT\_XACCESSO16} (Section~\ref{sec:format-EXT-XACCESSO16}), \verb|log2size = 000|\\

\bitfields{ 1/P, 2/10, 1/0, 1/0, 3/000, 4/registerN, 2/11, 2/00, 1/1, 3/000, 2/registerCj, 4/0111, 6/registerZ }


\paragraph{Syntax} \texttt{xaccesso} \emph{registerN} = \emph{registerCj}, \emph{registerZ}

\paragraph{Description} The 5 least significant bits of the \emph{registerZ} are interpreted as a byte count. The most significant bits of the \emph{registerZ} are interpreted as a vector register index into the buffer specified by the \emph{registerCj}. The vector register indexed and its successor modulo the buffer size are concatenated, and shifted right bt the byte count. The resulting 32 bytes are written to the \emph{registerN}.


\paragraph{Execution} {Reservation table \textsf{EXT\_MISC\_AUXW} (Section~\ref{sec:reservation-EXT-MISC-AUXW})}
~
{\begin{lstlisting}
stage RR:
  new mask = 15;
  new buffer = registerCj << 4;
  new address = GPR[registerZ];
  new shift = (address & 31) * 8;
  new index_0 = address >> 5;
  new index_1 = index_0 + 1;
  new where_0 = buffer + (index_0 & mask);
  new where_1 = buffer + (index_1 & mask);
  new result1_0 = XVR[where_0];
  new result1_1 = XVR[where_1];
  new result1 = result1_0;
  if (shift) {
    result1 = (_ZX_256(result1_0) >> shift) | _ZX_256(_ZX_256(result1_1) << (256 - shift));
  }
stage E1:
  CS.XMF = 1;
stage E3:
  QGR[registerN] = result1;
\end{lstlisting}}
\newpage

\paragraph{Encoding} Format \textsf{EXT\_XACCESSO32} (Section~\ref{sec:format-EXT-XACCESSO32}), \verb|log2size = 000|\\

\bitfields{ 1/P, 2/10, 1/0, 1/0, 3/000, 4/registerN, 2/11, 2/00, 1/1, 3/000, 1/registerCk, 5/01111, 6/registerZ }


\paragraph{Syntax} \texttt{xaccesso} \emph{registerN} = \emph{registerCk}, \emph{registerZ}

\paragraph{Description} The 5 least significant bits of the \emph{registerZ} are interpreted as a byte count. The most significant bits of the \emph{registerZ} are interpreted as a vector register index into the buffer specified by the \emph{registerCk}. The vector register indexed and its successor modulo the buffer size are concatenated, and shifted right bt the byte count. The resulting 32 bytes are written to the \emph{registerN}.


\paragraph{Execution} {Reservation table \textsf{EXT\_MISC\_AUXW} (Section~\ref{sec:reservation-EXT-MISC-AUXW})}
~
{\begin{lstlisting}
stage RR:
  new mask = 31;
  new buffer = registerCk << 5;
  new address = GPR[registerZ];
  new shift = (address & 31) * 8;
  new index_0 = address >> 5;
  new index_1 = index_0 + 1;
  new where_0 = buffer + (index_0 & mask);
  new where_1 = buffer + (index_1 & mask);
  new result1_0 = XVR[where_0];
  new result1_1 = XVR[where_1];
  new result1 = result1_0;
  if (shift) {
    result1 = (_ZX_256(result1_0) >> shift) | _ZX_256(_ZX_256(result1_1) << (256 - shift));
  }
stage E1:
  CS.XMF = 1;
stage E3:
  QGR[registerN] = result1;
\end{lstlisting}}
\newpage

\paragraph{Encoding} Format \textsf{EXT\_XACCESSO64} (Section~\ref{sec:format-EXT-XACCESSO64}), \verb|log2size = 000|\\

\bitfields{ 1/P, 2/10, 1/0, 1/0, 3/000, 4/registerN, 2/11, 2/00, 1/1, 3/000, 1/registerCl, 5/11111, 6/registerZ }


\paragraph{Syntax} \texttt{xaccesso} \emph{registerN} = \emph{registerCl}, \emph{registerZ}

\paragraph{Description} The 5 least significant bits of the \emph{registerZ} are interpreted as a byte count. The most significant bits of the \emph{registerZ} are interpreted as a vector register index into the buffer specified by the \emph{registerCl}. The vector register indexed and its successor modulo the buffer size are concatenated, and shifted right bt the byte count. The resulting 32 bytes are written to the \emph{registerN}.


\paragraph{Execution} {Reservation table \textsf{EXT\_MISC\_AUXW} (Section~\ref{sec:reservation-EXT-MISC-AUXW})}
~
{\begin{lstlisting}
stage RR:
  new mask = 63;
  new buffer = registerCl;
  new address = GPR[registerZ];
  new shift = (address & 31) * 8;
  new index_0 = address >> 5;
  new index_1 = index_0 + 1;
  new where_0 = buffer + (index_0 & mask);
  new where_1 = buffer + (index_1 & mask);
  new result1_0 = XVR[where_0];
  new result1_1 = XVR[where_1];
  new result1 = result1_0;
  if (shift) {
    result1 = (_ZX_256(result1_0) >> shift) | _ZX_256(_ZX_256(result1_1) << (256 - shift));
  }
stage E1:
  CS.XMF = 1;
stage E3:
  QGR[registerN] = result1;
\end{lstlisting}}
\newpage
\subsubsection[{\small XALIGNO: Align Extension into Octuple Word}]{XALIGNO\hfill Align Extension into Octuple Word}
\label{instr:XALIGNO}\index{xaligno}
 
\paragraph{Encoding} Format \textsf{EXT\_XALIGNO2} (Section~\ref{sec:format-EXT-XALIGNO2}), \verb|log2size = 000|\\

\bitfields{ 1/P, 2/10, 1/0, 1/0, 3/001, 6/registerA, 2/00, 1/1, 3/000, 5/registerCg, 1/0, 6/registerZ }


\paragraph{Syntax} \texttt{xaligno} \emph{registerA} = \emph{registerCg}, \emph{registerZ}

\paragraph{Description} The 5 least significant bits of the \emph{registerZ} are interpreted as a byte count. The most significant bits of the \emph{registerZ} are interpreted as a vector register index into the buffer specified by the \emph{registerCg}. The vector register indexed and its successor modulo the buffer size are concatenated, and shifted right bt the byte count. The resulting 32 bytes are written to the \emph{registerA}.


\paragraph{Execution} {Reservation table \textsf{EXT\_MISC\_AUXW} (Section~\ref{sec:reservation-EXT-MISC-AUXW})}
~
{\begin{lstlisting}
stage RR:
  new mask = 1;
  new buffer = registerCg << 1;
  new address = GPR[registerZ];
  new shift = (address & 31) * 8;
  new index_0 = address >> 5;
  new index_1 = index_0 + 1;
  new where_0 = buffer + (index_0 & mask);
  new where_1 = buffer + (index_1 & mask);
  new result1_0 = XVR[where_0];
  new result1_1 = XVR[where_1];
  new result1 = result1_0;
  if (shift) {
    result1 = (_ZX_256(result1_0) >> shift) | _ZX_256(_ZX_256(result1_1) << (256 - shift));
  }
stage E1:
  CS.XMF = 1;
stage E3:
  XVR[registerA] = result1;
\end{lstlisting}}
\newpage

\paragraph{Encoding} Format \textsf{EXT\_XALIGNO4} (Section~\ref{sec:format-EXT-XALIGNO4}), \verb|log2size = 000|\\

\bitfields{ 1/P, 2/10, 1/0, 1/0, 3/001, 6/registerA, 2/00, 1/1, 3/000, 4/registerCh, 2/01, 6/registerZ }


\paragraph{Syntax} \texttt{xaligno} \emph{registerA} = \emph{registerCh}, \emph{registerZ}

\paragraph{Description} The 5 least significant bits of the \emph{registerZ} are interpreted as a byte count. The most significant bits of the \emph{registerZ} are interpreted as a vector register index into the buffer specified by the \emph{registerCh}. The vector register indexed and its successor modulo the buffer size are concatenated, and shifted right bt the byte count. The resulting 32 bytes are written to the \emph{registerA}.


\paragraph{Execution} {Reservation table \textsf{EXT\_MISC\_AUXW} (Section~\ref{sec:reservation-EXT-MISC-AUXW})}
~
{\begin{lstlisting}
stage RR:
  new mask = 3;
  new buffer = registerCh << 2;
  new address = GPR[registerZ];
  new shift = (address & 31) * 8;
  new index_0 = address >> 5;
  new index_1 = index_0 + 1;
  new where_0 = buffer + (index_0 & mask);
  new where_1 = buffer + (index_1 & mask);
  new result1_0 = XVR[where_0];
  new result1_1 = XVR[where_1];
  new result1 = result1_0;
  if (shift) {
    result1 = (_ZX_256(result1_0) >> shift) | _ZX_256(_ZX_256(result1_1) << (256 - shift));
  }
stage E1:
  CS.XMF = 1;
stage E3:
  XVR[registerA] = result1;
\end{lstlisting}}
\newpage

\paragraph{Encoding} Format \textsf{EXT\_XALIGNO8} (Section~\ref{sec:format-EXT-XALIGNO8}), \verb|log2size = 000|\\

\bitfields{ 1/P, 2/10, 1/0, 1/0, 3/001, 6/registerA, 2/00, 1/1, 3/000, 3/registerCi, 3/011, 6/registerZ }


\paragraph{Syntax} \texttt{xaligno} \emph{registerA} = \emph{registerCi}, \emph{registerZ}

\paragraph{Description} The 5 least significant bits of the \emph{registerZ} are interpreted as a byte count. The most significant bits of the \emph{registerZ} are interpreted as a vector register index into the buffer specified by the \emph{registerCi}. The vector register indexed and its successor modulo the buffer size are concatenated, and shifted right bt the byte count. The resulting 32 bytes are written to the \emph{registerA}.


\paragraph{Execution} {Reservation table \textsf{EXT\_MISC\_AUXW} (Section~\ref{sec:reservation-EXT-MISC-AUXW})}
~
{\begin{lstlisting}
stage RR:
  new mask = 7;
  new buffer = registerCi << 3;
  new address = GPR[registerZ];
  new shift = (address & 31) * 8;
  new index_0 = address >> 5;
  new index_1 = index_0 + 1;
  new where_0 = buffer + (index_0 & mask);
  new where_1 = buffer + (index_1 & mask);
  new result1_0 = XVR[where_0];
  new result1_1 = XVR[where_1];
  new result1 = result1_0;
  if (shift) {
    result1 = (_ZX_256(result1_0) >> shift) | _ZX_256(_ZX_256(result1_1) << (256 - shift));
  }
stage E1:
  CS.XMF = 1;
stage E3:
  XVR[registerA] = result1;
\end{lstlisting}}
\newpage

\paragraph{Encoding} Format \textsf{EXT\_XALIGNO16} (Section~\ref{sec:format-EXT-XALIGNO16}), \verb|log2size = 000|\\

\bitfields{ 1/P, 2/10, 1/0, 1/0, 3/001, 6/registerA, 2/00, 1/1, 3/000, 2/registerCj, 4/0111, 6/registerZ }


\paragraph{Syntax} \texttt{xaligno} \emph{registerA} = \emph{registerCj}, \emph{registerZ}

\paragraph{Description} The 5 least significant bits of the \emph{registerZ} are interpreted as a byte count. The most significant bits of the \emph{registerZ} are interpreted as a vector register index into the buffer specified by the \emph{registerCj}. The vector register indexed and its successor modulo the buffer size are concatenated, and shifted right bt the byte count. The resulting 32 bytes are written to the \emph{registerA}.


\paragraph{Execution} {Reservation table \textsf{EXT\_MISC\_AUXW} (Section~\ref{sec:reservation-EXT-MISC-AUXW})}
~
{\begin{lstlisting}
stage RR:
  new mask = 15;
  new buffer = registerCj << 4;
  new address = GPR[registerZ];
  new shift = (address & 31) * 8;
  new index_0 = address >> 5;
  new index_1 = index_0 + 1;
  new where_0 = buffer + (index_0 & mask);
  new where_1 = buffer + (index_1 & mask);
  new result1_0 = XVR[where_0];
  new result1_1 = XVR[where_1];
  new result1 = result1_0;
  if (shift) {
    result1 = (_ZX_256(result1_0) >> shift) | _ZX_256(_ZX_256(result1_1) << (256 - shift));
  }
stage E1:
  CS.XMF = 1;
stage E3:
  XVR[registerA] = result1;
\end{lstlisting}}
\newpage

\paragraph{Encoding} Format \textsf{EXT\_XALIGNO32} (Section~\ref{sec:format-EXT-XALIGNO32}), \verb|log2size = 000|\\

\bitfields{ 1/P, 2/10, 1/0, 1/0, 3/001, 6/registerA, 2/00, 1/1, 3/000, 1/registerCk, 5/01111, 6/registerZ }


\paragraph{Syntax} \texttt{xaligno} \emph{registerA} = \emph{registerCk}, \emph{registerZ}

\paragraph{Description} The 5 least significant bits of the \emph{registerZ} are interpreted as a byte count. The most significant bits of the \emph{registerZ} are interpreted as a vector register index into the buffer specified by the \emph{registerCk}. The vector register indexed and its successor modulo the buffer size are concatenated, and shifted right bt the byte count. The resulting 32 bytes are written to the \emph{registerA}.


\paragraph{Execution} {Reservation table \textsf{EXT\_MISC\_AUXW} (Section~\ref{sec:reservation-EXT-MISC-AUXW})}
~
{\begin{lstlisting}
stage RR:
  new mask = 31;
  new buffer = registerCk << 5;
  new address = GPR[registerZ];
  new shift = (address & 31) * 8;
  new index_0 = address >> 5;
  new index_1 = index_0 + 1;
  new where_0 = buffer + (index_0 & mask);
  new where_1 = buffer + (index_1 & mask);
  new result1_0 = XVR[where_0];
  new result1_1 = XVR[where_1];
  new result1 = result1_0;
  if (shift) {
    result1 = (_ZX_256(result1_0) >> shift) | _ZX_256(_ZX_256(result1_1) << (256 - shift));
  }
stage E1:
  CS.XMF = 1;
stage E3:
  XVR[registerA] = result1;
\end{lstlisting}}
\newpage

\paragraph{Encoding} Format \textsf{EXT\_XALIGNO64} (Section~\ref{sec:format-EXT-XALIGNO64}), \verb|log2size = 000|\\

\bitfields{ 1/P, 2/10, 1/0, 1/0, 3/001, 6/registerA, 2/00, 1/1, 3/000, 1/registerCl, 5/11111, 6/registerZ }


\paragraph{Syntax} \texttt{xaligno} \emph{registerA} = \emph{registerCl}, \emph{registerZ}

\paragraph{Description} The 5 least significant bits of the \emph{registerZ} are interpreted as a byte count. The most significant bits of the \emph{registerZ} are interpreted as a vector register index into the buffer specified by the \emph{registerCl}. The vector register indexed and its successor modulo the buffer size are concatenated, and shifted right bt the byte count. The resulting 32 bytes are written to the \emph{registerA}.


\paragraph{Execution} {Reservation table \textsf{EXT\_MISC\_AUXW} (Section~\ref{sec:reservation-EXT-MISC-AUXW})}
~
{\begin{lstlisting}
stage RR:
  new mask = 63;
  new buffer = registerCl;
  new address = GPR[registerZ];
  new shift = (address & 31) * 8;
  new index_0 = address >> 5;
  new index_1 = index_0 + 1;
  new where_0 = buffer + (index_0 & mask);
  new where_1 = buffer + (index_1 & mask);
  new result1_0 = XVR[where_0];
  new result1_1 = XVR[where_1];
  new result1 = result1_0;
  if (shift) {
    result1 = (_ZX_256(result1_0) >> shift) | _ZX_256(_ZX_256(result1_1) << (256 - shift));
  }
stage E1:
  CS.XMF = 1;
stage E3:
  XVR[registerA] = result1;
\end{lstlisting}}
\newpage
\subsubsection[{\small XCOPYO: Copy Extension Octuple Word}]{XCOPYO\hfill Copy Extension Octuple Word}
\label{instr:XCOPYO}\index{xcopyo}
 
\paragraph{Encoding} Format \textsf{EXT\_XCOPYO} (Section~\ref{sec:format-EXT-XCOPYO}), \verb|log2size = 000|\\

\bitfields{ 1/P, 2/10, 1/0, 1/0, 3/001, 6/registerA, 2/00, 1/0, 3/000, 6/registerC, 6/-- -- -- -- -- -- }


\paragraph{Syntax} \texttt{xcopyo} \emph{registerA} = \emph{registerC}

\paragraph{Description} The \emph{registerA} operand receives the value of the \emph{registerC} operand.


\paragraph{Execution} {Reservation table \textsf{EXT\_MISC\_AUXW} (Section~\ref{sec:reservation-EXT-MISC-AUXW})}
~
{\begin{lstlisting}
stage RR:
  new argument2 = XVR[registerC];
  new result1 = argument2;
stage E1:
  CS.XMF = 1;
stage E3:
  XVR[registerA] = result1;
\end{lstlisting}}
\newpage
\subsubsection[{\small XMOVEFD: Move From Extension Double Word}]{XMOVEFD\hfill Move From Extension Double Word}
\label{instr:XMOVEFD}\index{xmovefd}
 
\paragraph{Encoding} Format \textsf{EXT\_XMOVEFD} (Section~\ref{sec:format-EXT-XMOVEFD}), \verb|tcacode1 = 000|\\

\bitfields{ 1/P, 2/10, 1/0, 1/0, 3/000, 6/registerW, 2/00, 1/1, 1/1, 2/qselectC, 6/registerCc, 6/-- -- -- -- -- -- }


\paragraph{Syntax} \texttt{xmovefd} \emph{registerW} = \emph{registerCc\_\discretionary{}{}{}qselectC}

\paragraph{Description} The \emph{registerW} operand receives the contents of the \emph{registerCc\_\discretionary{}{}{}qselectC}.


\paragraph{Execution} {Reservation table \textsf{EXT\_MISC\_AUXW} (Section~\ref{sec:reservation-EXT-MISC-AUXW})}
~
{\begin{lstlisting}
stage RR:
  new argument2 = XCR[registerCc_qselectC];
  new result1 = _ZX_64(argument2);
stage E1:
  CS.XMF = 1;
stage E3:
  GPR[registerW] = result1;
\end{lstlisting}}
\newpage
\subsubsection[{\small XMOVEFO: Move From Extension Octuple Word}]{XMOVEFO\hfill Move From Extension Octuple Word}
\label{instr:XMOVEFO}\index{xmovefo}
 
\paragraph{Encoding} Format \textsf{EXT\_XMOVEFO} (Section~\ref{sec:format-EXT-XMOVEFO}), \verb|tcacode1 = 000|\\

\bitfields{ 1/P, 2/10, 1/0, 1/0, 3/000, 4/registerN, 2/01, 2/00, 1/1, 3/000, 6/registerC, 6/-- -- -- -- -- -- }


\paragraph{Syntax} \texttt{xmovefo} \emph{registerN} = \emph{registerC}

\paragraph{Description} The \emph{registerN} operand receives the contents of the \emph{registerC}.


\paragraph{Execution} {Reservation table \textsf{EXT\_MISC\_AUXW} (Section~\ref{sec:reservation-EXT-MISC-AUXW})}
~
{\begin{lstlisting}
stage RR:
  new argument2 = XVR[registerC];
  new result1 = _ZX_256(argument2);
stage E1:
  CS.XMF = 1;
stage E3:
  QGR[registerN] = result1;
\end{lstlisting}}
\newpage
\subsubsection[{\small XMOVEFQ: Move From Extension Quadruple Word}]{XMOVEFQ\hfill Move From Extension Quadruple Word}
\label{instr:XMOVEFQ}\index{xmovefq}
 
\paragraph{Encoding} Format \textsf{EXT\_XMOVEFQ} (Section~\ref{sec:format-EXT-XMOVEFQ}), \verb|tcacode1 = 000|\\

\bitfields{ 1/P, 2/10, 1/0, 1/0, 3/000, 5/registerM, 1/1, 2/00, 1/0, 1/0, 1/0, 1/hselectC, 6/registerCb, 6/-- -- -- -- -- -- }


\paragraph{Syntax} \texttt{xmovefq} \emph{registerM} = \emph{registerCb\_\discretionary{}{}{}hselectC}

\paragraph{Description} The \emph{registerM} operand receives the contents of the \emph{registerCb\_\discretionary{}{}{}hselectC}.


\paragraph{Execution} {Reservation table \textsf{EXT\_MISC\_AUXW} (Section~\ref{sec:reservation-EXT-MISC-AUXW})}
~
{\begin{lstlisting}
stage RR:
  new argument2 = XBR[registerCb_hselectC];
  new result1 = _ZX_128(argument2);
stage E1:
  CS.XMF = 1;
stage E3:
  PGR[registerM] = result1;
\end{lstlisting}}
\newpage
\subsection{Instruction operands}
\label{sec:Instruction-operands}

\subsubsection{Immediate classes}
\label{opnd:Immediate-classes}
{\scriptsize
\begin{tabular}{|l|l|l|} \hline
\bf{Immediate classes} & \bf{Operands} & \bf{What} \\ \hline
brknumber & brknumber & \verb| Break Number |\\
pcrel11s2 & pcrel11s2 & \verb| CCB Offset |\\
pcrel17s2 & pcrel17s2 & \verb| CB Offset |\\
pcrel27s2 & pcrel27s2 & \verb| GOTO / CALL Offset |\\
pcrel38s2 & upper27\_\discretionary{}{}{}lower11 & \verb| CCB Offset Extended |\\
pcrel44s2 & upper27\_\discretionary{}{}{}lower17 & \verb| CB Offset Extended |\\
pcrel54s2 & upper27\_\discretionary{}{}{}lower27 & \verb| GOTO Offset Extended |\\
signed10 & signed10 & \verb| Sign Extended Immediate 10 bits |\\
signed16 & signed16 & \verb| Sign Extended Immediate 16 bits |\\
signed27 & offset27 & \verb| Sign Extended Immediate 27 bits |\\
signed37 & upper27\_\discretionary{}{}{}lower10 & \verb| Sign Extended Immediate 37 bits |\\
signed43 & extend6\_\discretionary{}{}{}upper27\_\discretionary{}{}{}lower10 & \verb| Sign Extended Immediate 43 bits |\\
signed54 & extend27\_\discretionary{}{}{}offset27 & \verb| Sign Extended Immediate 54 bits |\\
sysnumber & sysnumber & \verb| System Number |\\
unsigned6 & activate, startbit, stopbit2\_\discretionary{}{}{}stopbit4, unsigned6 & \verb| Zero Extended Immediate 6 bits |\\
wrapped32 & upper27\_\discretionary{}{}{}lower5 & \verb| Wrapped Immediate 32 bits |\\
wrapped64 & extend27\_\discretionary{}{}{}upper27\_\discretionary{}{}{}lower10 & \verb| Wrapped Immediate 64 bits |\\
\hline \end{tabular}
}
\newpage
{\scriptsize
\begin{tabular}{|l|l|l|l|l|} \hline
\bf{Immediate classes} & \bf{Extend} & \bf{Shift} & \bf{Width} & \bf{Values} \\ \hline
brknumber & Unsigned & 0 & 2 & 0x0 to 0x3 \\
pcrel11s2 & Signed & 2 & 11 & 0x1000 to 0xffc \\
pcrel17s2 & Signed & 2 & 17 & 0x40000 to 0x3fffc \\
pcrel27s2 & Signed & 2 & 27 & 0x10000000 to 0xffffffc \\
pcrel38s2 & Signed & 2 & 38 & 0x8000000000 to 0x7ffffffffc \\
pcrel44s2 & Signed & 2 & 44 & 0x200000000000 to 0x1ffffffffffc \\
pcrel54s2 & Signed & 2 & 54 & 0x80000000000000 to 0x7ffffffffffffc \\
signed10 & Signed & 0 & 10 & 0x200 to 0x1ff \\
signed16 & Signed & 0 & 16 & 0x8000 to 0x7fff \\
signed27 & Signed & 0 & 27 & 0x4000000 to 0x3ffffff \\
signed37 & Signed & 0 & 37 & 0x1000000000 to 0xfffffffff \\
signed43 & Signed & 0 & 43 & 0x40000000000 to 0x3ffffffffff \\
signed54 & Signed & 0 & 54 & 0x20000000000000 to 0x1fffffffffffff \\
sysnumber & Unsigned & 0 & 12 & 0x0 to 0xfff \\
unsigned6 & Unsigned & 0 & 6 & 0x0 to 0x3f \\
wrapped32 & Wrap & 0 & 32 & 0x80000000 to 0xffffffff \\
wrapped64 & Wrap & 0 & 64 & 0x8000000000000000 to 0x7fffffffffffffff \\
\hline \end{tabular}
}
\newpage
\subsubsection{Operand Modifiers}
\label{opnd:Operand-Modifiers}
{\scriptsize
\begin{tabular}{|l|l|l|} \hline
\bf{Operand Modifiers} & \bf{Instances} \\ \hline
accesses &  \\
bcucond & , ,  \\
boolcas &  \\
cachelev &  \\
ccbcomp &  \\
coherency &  \\
conjugate &  \\
floatcomp &  \\
floatmode &  \\
fnegate &  \\
highmult &  \\
imultiply &  \\
intcomp &  \\
lanecond &  \\
lanesize &  \\
lanetodo &  \\
mostsig &  \\
oddlanes &  \\
qindex &  \\
signextw &  \\
splat32 &  \\
variant &  \\
widemult &  \\
ziplanes &  \\
\hline \end{tabular}
}
\newpage

\paragraph{Modifier} \textsf{accesses}:\\
\phantomsection
\label{par:modifier-accesses}
\noindent \begin{tabular}{|l|l|r|} \hline
\bf{Name} & \bf{Case} & \bf{Value} \\ \hline
. & All Accesses & 0 \\
.W & Write Accesses & 1 \\
.R & Read Accesses & 2 \\
.WA & Write Accesses (asynchronous) & 3 \\
\hline \end{tabular}
\newpage

\paragraph{Modifier} \textsf{bcucond}:\\
\phantomsection
\label{par:modifier-bcucond}
\noindent \begin{tabular}{|l|l|r|} \hline
\bf{Name} & \bf{Case} & \bf{Value} \\ \hline
.DLTZ & Double Less Than Zero & 0 \\
.DGEZ & Double Greater Than or Equal to Zero & 1 \\
.DLEZ & Double Less Than or Equal to Zero & 2 \\
.DGTZ & Double Greater Than Zero & 3 \\
.DEQZ & Double Equal to Zero & 4 \\
.DNEZ & Double Not Equal to Zero & 5 \\
.ODD & Odd (LSB Set) & 6 \\
.EVEN & Even (LSB Clear) & 7 \\
.WLTZ & Word Less Than Zero & 8 \\
.WGEZ & Word Greater Than or Equal to Zero & 9 \\
.WLEZ & Word Less Than or Equal to Zero & 10 \\
.WGTZ & Word Greater Than Zero & 11 \\
.WEQZ & Word Equal to Zero & 12 \\
.WNEZ & Word Not Equal to Zero & 13 \\
\hline \end{tabular}
\newpage

\paragraph{Modifier} \textsf{boolcas}:\\
\phantomsection
\label{par:modifier-boolcas}
\noindent \begin{tabular}{|l|l|r|} \hline
\bf{Name} & \bf{Case} & \bf{Value} \\ \hline
.V & Valued CAS & 0 \\
. & Boolean CAS & 1 \\
\hline \end{tabular}
\newpage

\paragraph{Modifier} \textsf{cachelev}:\\
\phantomsection
\label{par:modifier-cachelev}
\noindent \begin{tabular}{|l|l|r|} \hline
\bf{Name} & \bf{Case} & \bf{Value} \\ \hline
.L1 & L1 Cache & 0 \\
.L2 & L2 Cache & 1 \\
\hline \end{tabular}
\newpage

\paragraph{Modifier} \textsf{ccbcomp}:\\
\phantomsection
\label{par:modifier-ccbcomp}
\noindent \begin{tabular}{|l|l|r|} \hline
\bf{Name} & \bf{Case} & \bf{Value} \\ \hline
.DLT & Double Word Less Than & 0 \\
.DGE & Double Word Greater Than or Equal & 1 \\
.DLTU & Double Word Less Than Unsigned & 2 \\
.DGEU & Double Word Greater Than or Equal Unsigned & 3 \\
.DEQ & Double Word Equal & 4 \\
.DNE & Double Word Not Equal & 5 \\
.DANY & Double Any bit set after AND-ing & 6 \\
.DNONE & Double No bits set after AND-ing & 7 \\
.WLT & Word Less Than & 8 \\
.WGE & Word Greater Than or Equal & 9 \\
.WLTU & Word Less Than Unsigned & 10 \\
.WGEU & Word Greater Than or Equal Unsigned & 11 \\
.WEQ & Word Equal & 12 \\
.WNE & Word Not Equal & 13 \\
.WANY & Word Any bit set after AND-ing & 14 \\
.WNONE & Word No bits set after AND-ing & 15 \\
\hline \end{tabular}
\newpage

\paragraph{Modifier} \textsf{coherency}:\\
\phantomsection
\label{par:modifier-coherency}
\noindent \begin{tabular}{|l|l|r|} \hline
\bf{Name} & \bf{Case} & \bf{Value} \\ \hline
. & Local & 0 \\
.G & Global & 1 \\
.S & Reserved & 2 \\
\hline \end{tabular}
\newpage

\paragraph{Modifier} \textsf{conjugate}:\\
\phantomsection
\label{par:modifier-conjugate}
\noindent \begin{tabular}{|l|l|r|} \hline
\bf{Name} & \bf{Case} & \bf{Value} \\ \hline
. & Default & 0 \\
.C & Conjugate & 1 \\
\hline \end{tabular}
\newpage

\paragraph{Modifier} \textsf{floatcomp}:\\
\phantomsection
\label{par:modifier-floatcomp}
\noindent \begin{tabular}{|l|l|r|} \hline
\bf{Name} & \bf{Case} & \bf{Value} \\ \hline
.ONE & Ordered and Not Equal & 0 \\
.UEQ & Unordered or Equal & 1 \\
.OEQ & Ordered and Equal & 2 \\
.UNE & Unordered or Not Equal & 3 \\
.OLT & Ordered and Less Than & 4 \\
.UGE & Unordered or Greater Than or Equal & 5 \\
.OGE & Ordered and Greater Than or Equal & 6 \\
.ULT & Unordered or Less Than & 7 \\
\hline \end{tabular}
\newpage

\paragraph{Modifier} \textsf{floatmode}:\\
\phantomsection
\label{par:modifier-floatmode}
\noindent \begin{tabular}{|l|l|r|} \hline
\bf{Name} & \bf{Case} & \bf{Value} \\ \hline
.RN & Round to Nearest, ties to Even & 0 \\
.RZ & Round toward Zero & 1 \\
.RD & Round Downward & 2 \\
.RU & Round Upward & 3 \\
.RM & Round to Nearest, ties to Max Magnitude & 4 \\
.R5 & Reserved 5 & 5 \\
.RO & Round to Odd & 6 \\
. & Use CS rounding & 7 \\
\hline \end{tabular}
\newpage

\paragraph{Modifier} \textsf{fnegate}:\\
\phantomsection
\label{par:modifier-fnegate}
\noindent \begin{tabular}{|l|l|r|} \hline
\bf{Name} & \bf{Case} & \bf{Value} \\ \hline
. & Default Result & 0 \\
.N & Negate Result & 1 \\
\hline \end{tabular}
\newpage

\paragraph{Modifier} \textsf{highmult}:\\
\phantomsection
\label{par:modifier-highmult}
\noindent \begin{tabular}{|l|l|r|} \hline
\bf{Name} & \bf{Case} & \bf{Value} \\ \hline
. & Low & 0 \\
.H & High & 1 \\
.HU & High Unsigned & 2 \\
.HSU & High Signed Unsigned & 3 \\
\hline \end{tabular}
\newpage

\paragraph{Modifier} \textsf{imultiply}:\\
\phantomsection
\label{par:modifier-imultiply}
\noindent \begin{tabular}{|l|l|r|} \hline
\bf{Name} & \bf{Case} & \bf{Value} \\ \hline
. & Default & 0 \\
.MI & Multiply by I & 1 \\
\hline \end{tabular}
\newpage

\paragraph{Modifier} \textsf{intcomp}:\\
\phantomsection
\label{par:modifier-intcomp}
\noindent \begin{tabular}{|l|l|r|} \hline
\bf{Name} & \bf{Case} & \bf{Value} \\ \hline
.LT & Less Than & 0 \\
.GE & Greater Than or Equal & 1 \\
.LTU & Less Than Unsigned & 2 \\
.GEU & Greater Than or Equal Unsigned & 3 \\
.EQ & Equal & 4 \\
.NE & Not Equal & 5 \\
.ANY & Any bits set in after AND-ing & 6 \\
.NONE & No bits set after AND-ing & 7 \\
.LE & Less Than or Equal & 8 \\
.GT & Greater Than & 9 \\
.LEU & Less Than or Equal Unsigned & 10 \\
.GTU & Greater Than Unsigned & 11 \\
\hline \end{tabular}
\newpage

\paragraph{Modifier} \textsf{lanecond}:\\
\phantomsection
\label{par:modifier-lanecond}
\noindent \begin{tabular}{|l|l|r|} \hline
\bf{Name} & \bf{Case} & \bf{Value} \\ \hline
.LTZ & Less Than Zero & 0 \\
.GEZ & Greater Than or Equal to Zero & 1 \\
.LEZ & Less Than or Equal to Zero & 2 \\
.GTZ & Greater Than Zero & 3 \\
.EQZ & Equal to Zero & 4 \\
.NEZ & Not Equal to Zero & 5 \\
.ODD & Odd (LSB Set) & 6 \\
.EVEN & Even (LSB Clear) & 7 \\
\hline \end{tabular}
\newpage

\paragraph{Modifier} \textsf{lanesize}:\\
\phantomsection
\label{par:modifier-lanesize}
\noindent \begin{tabular}{|l|l|r|} \hline
\bf{Name} & \bf{Case} & \bf{Value} \\ \hline
. & Byte Size & 0 \\
.H & Half Word Size & 1 \\
.W & Word Size & 2 \\
.D & Double Word Size & 3 \\
\hline \end{tabular}
\newpage

\paragraph{Modifier} \textsf{lanetodo}:\\
\phantomsection
\label{par:modifier-lanetodo}
\noindent \begin{tabular}{|l|l|r|} \hline
\bf{Name} & \bf{Case} & \bf{Value} \\ \hline
.MT & Mask True & 0 \\
.MF & Mask False & 1 \\
.MTC & Mask True or Clear & 2 \\
.MFC & Mask False or Clear & 3 \\
\hline \end{tabular}
\newpage

\paragraph{Modifier} \textsf{mostsig}:\\
\phantomsection
\label{par:modifier-mostsig}
\noindent \begin{tabular}{|l|l|r|} \hline
\bf{Name} & \bf{Case} & \bf{Value} \\ \hline
. & Least Significant Bits & 0 \\
.M & Most Significant Bits & 1 \\
\hline \end{tabular}
\newpage

\paragraph{Modifier} \textsf{oddlanes}:\\
\phantomsection
\label{par:modifier-oddlanes}
\noindent \begin{tabular}{|l|l|r|} \hline
\bf{Name} & \bf{Case} & \bf{Value} \\ \hline
. & Even Lanes & 0 \\
.O & Odd Lanes & 1 \\
\hline \end{tabular}
\newpage

\paragraph{Modifier} \textsf{qindex}:\\
\phantomsection
\label{par:modifier-qindex}
\noindent \begin{tabular}{|l|l|r|} \hline
\bf{Name} & \bf{Case} & \bf{Value} \\ \hline
.Q0 & Quarter 0 & 0 \\
.Q1 & Quarter 1 & 1 \\
.Q2 & Quarter 2 & 2 \\
.Q3 & Quarter 3 & 3 \\
\hline \end{tabular}
\newpage

\paragraph{Modifier} \textsf{signextw}:\\
\phantomsection
\label{par:modifier-signextw}
\noindent \begin{tabular}{|l|l|r|} \hline
\bf{Name} & \bf{Case} & \bf{Value} \\ \hline
. & Zero Extend & 0 \\
.SX & Sign Extend & 1 \\
\hline \end{tabular}
\newpage

\paragraph{Modifier} \textsf{splat32}:\\
\phantomsection
\label{par:modifier-splat32}
\noindent \begin{tabular}{|l|l|r|} \hline
\bf{Name} & \bf{Case} & \bf{Value} \\ \hline
. & No splat & 0 \\
.@ & Splat x2 & 1 \\
\hline \end{tabular}
\newpage

\paragraph{Modifier} \textsf{variant}:\\
\phantomsection
\label{par:modifier-variant}
\noindent \begin{tabular}{|l|l|r|} \hline
\bf{Name} & \bf{Case} & \bf{Value} \\ \hline
. & Cached & 0 \\
.S & Speculative & 1 \\
.U & Uncached & 2 \\
.US & Uncached Speculative & 3 \\
\hline \end{tabular}
\newpage

\paragraph{Modifier} \textsf{widemult}:\\
\phantomsection
\label{par:modifier-widemult}
\noindent \begin{tabular}{|l|l|r|} \hline
\bf{Name} & \bf{Case} & \bf{Value} \\ \hline
. & Signed & 0 \\
.U & Unsigned & 1 \\
.SU & Signed by Unsigned & 2 \\
\hline \end{tabular}
\newpage

\paragraph{Modifier} \textsf{ziplanes}:\\
\phantomsection
\label{par:modifier-ziplanes}
\noindent \begin{tabular}{|l|l|r|} \hline
\bf{Name} & \bf{Case} & \bf{Value} \\ \hline
. & No zip & 0 \\
.Z & Zip lanes & 1 \\
\hline \end{tabular}
\newpage
\subsubsection{Register Classes}
\label{opnd:Register-Classes}
{\scriptsize
\begin{tabular}{|l|l|l|} \hline
\bf{Register Classes} & \bf{Instances} & \bf{Width} \\ \hline
aloneReg & systemAlone & 64 \\
buffer16Reg & registerCj, registerGj & 0 \\
buffer2Reg & registerCg, registerGg & 0 \\
buffer32Reg & registerCk, registerGk & 0 \\
buffer4Reg & registerCh, registerGh & 0 \\
buffer64Reg & registerCl, registerGl & 0 \\
buffer8Reg & registerCi, registerGi & 0 \\
onlyfxReg & systemT2 & 64 \\
onlygetReg & systemS2 & 64 \\
onlyraReg & systemRA & 64 \\
onlysetReg & systemT3 & 64 \\
onlyswapReg & systemS4 & 64 \\
pairedReg & registerM, registerO, registerP, registerU & 128 \\
quadReg & registerN, registerQ, registerR, registerV & 256 \\
singleReg & registerT, registerW, registerY, registerZ & 64 \\
worddRegE & registerYe, registerZe & 64 \\
worddRegO & registerYo, registerZo & 64 \\
xworddReg & registerCc\_\discretionary{}{}{}qselectC & 64 \\
xworddReg0M4 & registerAx & 64 \\
xworddReg1M4 & registerAy & 64 \\
xworddReg2M4 & registerAz & 64 \\
xworddReg3M4 & registerAt & 64 \\
xwordoReg & registerA, registerC, registerE, registerG & 256 \\
xwordqReg & registerCb\_\discretionary{}{}{}hselectC & 128 \\
xwordqRegE & registerAE & 128 \\
xwordqRegO & registerAO & 128 \\
xwordvReg & registerEq, registerGq & 1024 \\
\hline \end{tabular}
}
\newpage
{\scriptsize
\begin{tabular}{|l|l|l|} \hline
\bf{Register Class} & \bf{Register File} & \bf{Registers} \\ \hline
 aloneReg & SFR & \texttt{ } \\
 & & \texttt{ } \\
 & & \texttt{ } \\
 & & \texttt{ } \\
 & & \texttt{ } \\
 & & \texttt{ } \\
 & & \texttt{ } \\
 & & \texttt{ } \\
 & & \texttt{ } \\
 & & \texttt{1 } \\
 & & \texttt{1 } \\
 & & \texttt{1 } \\
 & & \texttt{1 } \\
\hline \end{tabular}
}
\newpage
{\scriptsize
\begin{tabular}{|l|l|l|} \hline
\bf{Register Class} & \bf{Register File} & \bf{Registers} \\ \hline
 buffer16Reg & X16R & \texttt{} \\
 & & \texttt{} \\
 & & \texttt{} \\
 & & \texttt{} \\
\hline \end{tabular}
}
\newpage
{\scriptsize
\begin{tabular}{|l|l|l|} \hline
\bf{Register Class} & \bf{Register File} & \bf{Registers} \\ \hline
 buffer2Reg & X2R & \texttt{} \\
 & & \texttt{} \\
 & & \texttt{} \\
 & & \texttt{} \\
 & & \texttt{} \\
 & & \texttt{} \\
 & & \texttt{} \\
 & & \texttt{} \\
 & & \texttt{} \\
 & & \texttt{} \\
 & & \texttt{} \\
 & & \texttt{} \\
 & & \texttt{} \\
 & & \texttt{} \\
 & & \texttt{} \\
 & & \texttt{} \\
 & & \texttt{} \\
 & & \texttt{} \\
 & & \texttt{} \\
 & & \texttt{} \\
 & & \texttt{} \\
 & & \texttt{} \\
 & & \texttt{} \\
 & & \texttt{} \\
 & & \texttt{} \\
 & & \texttt{} \\
 & & \texttt{} \\
 & & \texttt{} \\
 & & \texttt{} \\
 & & \texttt{} \\
 & & \texttt{} \\
 & & \texttt{} \\
\hline \end{tabular}
}
\newpage
{\scriptsize
\begin{tabular}{|l|l|l|} \hline
\bf{Register Class} & \bf{Register File} & \bf{Registers} \\ \hline
 buffer32Reg & X32R & \texttt{} \\
 & & \texttt{} \\
\hline \end{tabular}
}
\newpage
{\scriptsize
\begin{tabular}{|l|l|l|} \hline
\bf{Register Class} & \bf{Register File} & \bf{Registers} \\ \hline
 buffer4Reg & X4R & \texttt{} \\
 & & \texttt{} \\
 & & \texttt{} \\
 & & \texttt{} \\
 & & \texttt{} \\
 & & \texttt{} \\
 & & \texttt{} \\
 & & \texttt{} \\
 & & \texttt{} \\
 & & \texttt{} \\
 & & \texttt{} \\
 & & \texttt{} \\
 & & \texttt{} \\
 & & \texttt{} \\
 & & \texttt{} \\
 & & \texttt{} \\
\hline \end{tabular}
}
\newpage
{\scriptsize
\begin{tabular}{|l|l|l|} \hline
\bf{Register Class} & \bf{Register File} & \bf{Registers} \\ \hline
 buffer64Reg & X64R & \texttt{} \\
\hline \end{tabular}
}
\newpage
{\scriptsize
\begin{tabular}{|l|l|l|} \hline
\bf{Register Class} & \bf{Register File} & \bf{Registers} \\ \hline
 buffer8Reg & X8R & \texttt{} \\
 & & \texttt{} \\
 & & \texttt{} \\
 & & \texttt{} \\
 & & \texttt{} \\
 & & \texttt{} \\
 & & \texttt{} \\
 & & \texttt{} \\
\hline \end{tabular}
}
\newpage
{\scriptsize
\begin{tabular}{|l|l|l|} \hline
\bf{Register Class} & \bf{Register File} & \bf{Registers} \\ \hline
 onlyfxReg & SFR & \texttt{ } \\
 & & \texttt{ } \\
 & & \texttt{ } \\
 & & \texttt{ } \\
 & & \texttt{ } \\
 & & \texttt{ } \\
 & & \texttt{ } \\
 & & \texttt{ } \\
 & & \texttt{ } \\
 & & \texttt{ } \\
 & & \texttt{ } \\
 & & \texttt{ } \\
 & & \texttt{ } \\
 & & \texttt{ } \\
 & & \texttt{ } \\
 & & \texttt{ } \\
 & & \texttt{ } \\
 & & \texttt{ } \\
 & & \texttt{ } \\
 & & \texttt{ } \\
 & & \texttt{ } \\
 & & \texttt{ } \\
 & & \texttt{1 } \\
 & & \texttt{1 } \\
 & & \texttt{1 } \\
 & & \texttt{1 } \\
 & & \texttt{1 } \\
 & & \texttt{1 } \\
 & & \texttt{1 } \\
 & & \texttt{1 } \\
 & & \texttt{ } \\
 & & \texttt{ } \\
 & & \texttt{1 } \\
 & & \texttt{1 } \\
 & & \texttt{1 } \\
 & & \texttt{1 } \\
 & & \texttt{ } \\
\hline \end{tabular}
}
\newpage
{\scriptsize
\begin{tabular}{|l|l|l|} \hline
\bf{Register Class} & \bf{Register File} & \bf{Registers} \\ \hline
 onlygetReg & SFR & \texttt{ } \\
 & & \texttt{ } \\
 & & \texttt{ } \\
 & & \texttt{ } \\
 & & \texttt{ } \\
 & & \texttt{ } \\
 & & \texttt{ } \\
 & & \texttt{ } \\
 & & \texttt{ } \\
 & & \texttt{ } \\
 & & \texttt{ } \\
 & & \texttt{ } \\
 & & \texttt{ } \\
 & & \texttt{ } \\
 & & \texttt{ } \\
 & & \texttt{ } \\
 & & \texttt{ } \\
 & & \texttt{ } \\
 & & \texttt{ } \\
 & & \texttt{ } \\
 & & \texttt{ } \\
 & & \texttt{ } \\
 & & \texttt{ } \\
 & & \texttt{ } \\
 & & \texttt{ } \\
 & & \texttt{ } \\
 & & \texttt{ } \\
 & & \texttt{ } \\
 & & \texttt{ } \\
 & & \texttt{ } \\
 & & \texttt{ } \\
 & & \texttt{ } \\
 & & \texttt{ } \\
 & & \texttt{ } \\
 & & \texttt{ } \\
 & & \texttt{ } \\
 & & \texttt{ } \\
 & & \texttt{ } \\
 & & \texttt{ } \\
 & & \texttt{ } \\
 & & \texttt{ } \\
 & & \texttt{ } \\
 & & \texttt{ } \\
 & & \texttt{ } \\
 & & \texttt{ } \\
 & & \texttt{ } \\
 & & \texttt{ } \\
 & & \texttt{ } \\
 & & \texttt{ } \\
 & & \texttt{ } \\
 & & \texttt{ } \\
 & & \texttt{ } \\
 & & \texttt{ } \\
 & & \texttt{ } \\
 & & \texttt{ } \\
 & & \texttt{ } \\
 & & \texttt{ } \\
 & & \texttt{ } \\
 & & \texttt{ } \\
 & & \texttt{ } \\
 & & \texttt{ } \\
 & & \texttt{ } \\
 & & \texttt{ } \\
 & & \texttt{ } \\
 & & \texttt{1 } \\
 & & \texttt{1 } \\
 & & \texttt{1 } \\
 & & \texttt{1 } \\
 & & \texttt{1 } \\
 & & \texttt{1 } \\
 & & \texttt{1 } \\
 & & \texttt{1 } \\
 & & \texttt{1 } \\
 & & \texttt{1 } \\
 & & \texttt{1 } \\
 & & \texttt{1 } \\
 & & \texttt{1 } \\
 & & \texttt{1 } \\
 & & \texttt{1 } \\
 & & \texttt{1 } \\
 & & \texttt{1 } \\
 & & \texttt{1 } \\
 & & \texttt{1 } \\
 & & \texttt{1 } \\
 & & \texttt{1 } \\
 & & \texttt{1 } \\
 & & \texttt{1 } \\
 & & \texttt{1 } \\
 & & \texttt{1 } \\
 & & \texttt{1 } \\
 & & \texttt{1 } \\
 & & \texttt{1 } \\
 & & \texttt{1 } \\
 & & \texttt{1 } \\
 & & \texttt{1 } \\
 & & \texttt{1 } \\
 & & \texttt{1 } \\
 & & \texttt{1 } \\
 & & \texttt{1 } \\
 & & \texttt{1 } \\
 & & \texttt{ } \\
 & & \texttt{ } \\
 & & \texttt{ } \\
 & & \texttt{ } \\
 & & \texttt{ } \\
 & & \texttt{ } \\
 & & \texttt{ } \\
 & & \texttt{ } \\
 & & \texttt{ } \\
 & & \texttt{ } \\
 & & \texttt{ } \\
 & & \texttt{ } \\
 & & \texttt{ } \\
 & & \texttt{ } \\
 & & \texttt{ } \\
 & & \texttt{ } \\
 & & \texttt{ } \\
 & & \texttt{ } \\
 & & \texttt{ } \\
 & & \texttt{ } \\
 & & \texttt{ } \\
 & & \texttt{ } \\
 & & \texttt{ } \\
 & & \texttt{ } \\
 & & \texttt{ } \\
 & & \texttt{ } \\
 & & \texttt{ } \\
 & & \texttt{ } \\
 & & \texttt{ } \\
 & & \texttt{ } \\
 & & \texttt{ } \\
 & & \texttt{ } \\
 & & \texttt{ } \\
 & & \texttt{ } \\
 & & \texttt{ } \\
 & & \texttt{ } \\
 & & \texttt{ } \\
 & & \texttt{ } \\
 & & \texttt{ } \\
 & & \texttt{ } \\
 & & \texttt{ } \\
 & & \texttt{ } \\
 & & \texttt{ } \\
 & & \texttt{ } \\
 & & \texttt{ } \\
 & & \texttt{ } \\
 & & \texttt{ } \\
 & & \texttt{ } \\
 & & \texttt{ } \\
 & & \texttt{ } \\
 & & \texttt{ } \\
 & & \texttt{ } \\
 & & \texttt{ } \\
 & & \texttt{ } \\
 & & \texttt{ } \\
 & & \texttt{ } \\
 & & \texttt{ } \\
 & & \texttt{ } \\
 & & \texttt{ } \\
 & & \texttt{ } \\
 & & \texttt{ } \\
 & & \texttt{ } \\
 & & \texttt{ } \\
 & & \texttt{ } \\
 & & \texttt{ } \\
 & & \texttt{ } \\
 & & \texttt{ } \\
 & & \texttt{ } \\
 & & \texttt{ } \\
 & & \texttt{ } \\
 & & \texttt{ } \\
 & & \texttt{ } \\
 & & \texttt{ } \\
 & & \texttt{ } \\
 & & \texttt{ } \\
 & & \texttt{ } \\
 & & \texttt{ } \\
 & & \texttt{ } \\
 & & \texttt{ } \\
 & & \texttt{ } \\
 & & \texttt{ } \\
 & & \texttt{ } \\
 & & \texttt{ } \\
 & & \texttt{ } \\
 & & \texttt{ } \\
 & & \texttt{ } \\
 & & \texttt{ } \\
 & & \texttt{ } \\
 & & \texttt{ } \\
 & & \texttt{ } \\
 & & \texttt{ } \\
 & & \texttt{ } \\
 & & \texttt{ } \\
 & & \texttt{ } \\
 & & \texttt{ } \\
 & & \texttt{ } \\
 & & \texttt{ } \\
 & & \texttt{ } \\
 & & \texttt{ } \\
 & & \texttt{ } \\
 & & \texttt{ } \\
 & & \texttt{ } \\
 & & \texttt{ } \\
 & & \texttt{ } \\
 & & \texttt{ } \\
 & & \texttt{ } \\
 & & \texttt{ } \\
 & & \texttt{ } \\
 & & \texttt{ } \\
 & & \texttt{ } \\
 & & \texttt{ } \\
 & & \texttt{ } \\
 & & \texttt{ } \\
 & & \texttt{ } \\
 & & \texttt{ } \\
 & & \texttt{ } \\
 & & \texttt{ } \\
 & & \texttt{ } \\
 & & \texttt{ } \\
 & & \texttt{ } \\
 & & \texttt{ } \\
 & & \texttt{ } \\
 & & \texttt{ } \\
 & & \texttt{ } \\
 & & \texttt{ } \\
 & & \texttt{ } \\
 & & \texttt{ } \\
 & & \texttt{ } \\
 & & \texttt{ } \\
 & & \texttt{ } \\
 & & \texttt{ } \\
 & & \texttt{ } \\
 & & \texttt{ } \\
 & & \texttt{ } \\
 & & \texttt{ } \\
 & & \texttt{ } \\
 & & \texttt{ } \\
 & & \texttt{ } \\
 & & \texttt{ } \\
 & & \texttt{ } \\
 & & \texttt{ } \\
 & & \texttt{ } \\
 & & \texttt{ } \\
 & & \texttt{ } \\
 & & \texttt{ } \\
 & & \texttt{ } \\
 & & \texttt{ } \\
 & & \texttt{ } \\
 & & \texttt{ } \\
 & & \texttt{ } \\
 & & \texttt{ } \\
 & & \texttt{ } \\
 & & \texttt{ } \\
 & & \texttt{ } \\
 & & \texttt{ } \\
 & & \texttt{ } \\
 & & \texttt{ } \\
 & & \texttt{ } \\
 & & \texttt{ } \\
 & & \texttt{ } \\
 & & \texttt{ } \\
 & & \texttt{ } \\
 & & \texttt{ } \\
 & & \texttt{ } \\
 & & \texttt{ } \\
 & & \texttt{ } \\
 & & \texttt{ } \\
 & & \texttt{ } \\
 & & \texttt{ } \\
 & & \texttt{ } \\
 & & \texttt{ } \\
 & & \texttt{ } \\
 & & \texttt{ } \\
 & & \texttt{ } \\
 & & \texttt{ } \\
 & & \texttt{ } \\
 & & \texttt{ } \\
 & & \texttt{ } \\
 & & \texttt{ } \\
 & & \texttt{ } \\
 & & \texttt{ } \\
 & & \texttt{ } \\
 & & \texttt{ } \\
 & & \texttt{ } \\
 & & \texttt{ } \\
 & & \texttt{ } \\
 & & \texttt{ } \\
 & & \texttt{ } \\
 & & \texttt{ } \\
 & & \texttt{ } \\
 & & \texttt{ } \\
 & & \texttt{ } \\
 & & \texttt{ } \\
 & & \texttt{ } \\
 & & \texttt{ } \\
 & & \texttt{ } \\
 & & \texttt{ } \\
 & & \texttt{ } \\
 & & \texttt{ } \\
 & & \texttt{ } \\
 & & \texttt{ } \\
 & & \texttt{ } \\
 & & \texttt{ } \\
 & & \texttt{ } \\
 & & \texttt{ } \\
 & & \texttt{ } \\
 & & \texttt{ } \\
 & & \texttt{ } \\
 & & \texttt{ } \\
 & & \texttt{ } \\
 & & \texttt{ } \\
 & & \texttt{ } \\
 & & \texttt{ } \\
 & & \texttt{ } \\
 & & \texttt{ } \\
 & & \texttt{ } \\
 & & \texttt{ } \\
 & & \texttt{ } \\
 & & \texttt{ } \\
 & & \texttt{ } \\
 & & \texttt{ } \\
 & & \texttt{ } \\
 & & \texttt{ } \\
 & & \texttt{ } \\
 & & \texttt{ } \\
 & & \texttt{ } \\
 & & \texttt{ } \\
 & & \texttt{ } \\
 & & \texttt{ } \\
 & & \texttt{ } \\
 & & \texttt{ } \\
 & & \texttt{ } \\
 & & \texttt{ } \\
 & & \texttt{ } \\
 & & \texttt{ } \\
 & & \texttt{ } \\
 & & \texttt{ } \\
 & & \texttt{ } \\
 & & \texttt{ } \\
 & & \texttt{ } \\
 & & \texttt{ } \\
 & & \texttt{ } \\
 & & \texttt{ } \\
 & & \texttt{ } \\
 & & \texttt{ } \\
 & & \texttt{ } \\
 & & \texttt{ } \\
 & & \texttt{ } \\
 & & \texttt{ } \\
 & & \texttt{ } \\
 & & \texttt{ } \\
 & & \texttt{ } \\
 & & \texttt{ } \\
 & & \texttt{ } \\
 & & \texttt{ } \\
 & & \texttt{ } \\
 & & \texttt{ } \\
 & & \texttt{ } \\
 & & \texttt{ } \\
 & & \texttt{ } \\
 & & \texttt{ } \\
 & & \texttt{ } \\
 & & \texttt{ } \\
 & & \texttt{ } \\
 & & \texttt{ } \\
\hline \end{tabular}
}
\newpage
{\scriptsize
\begin{tabular}{|l|l|l|} \hline
\bf{Register Class} & \bf{Register File} & \bf{Registers} \\ \hline
 onlyraReg & SFR & \texttt{ } \\
\hline \end{tabular}
}
\newpage
{\scriptsize
\begin{tabular}{|l|l|l|} \hline
\bf{Register Class} & \bf{Register File} & \bf{Registers} \\ \hline
 onlysetReg & SFR & \texttt{ } \\
 & & \texttt{ } \\
 & & \texttt{ } \\
 & & \texttt{ } \\
 & & \texttt{ } \\
 & & \texttt{ } \\
 & & \texttt{ } \\
 & & \texttt{ } \\
 & & \texttt{ } \\
 & & \texttt{ } \\
 & & \texttt{ } \\
 & & \texttt{ } \\
 & & \texttt{ } \\
 & & \texttt{ } \\
 & & \texttt{ } \\
 & & \texttt{ } \\
 & & \texttt{ } \\
 & & \texttt{ } \\
 & & \texttt{ } \\
 & & \texttt{ } \\
 & & \texttt{ } \\
 & & \texttt{ } \\
 & & \texttt{ } \\
 & & \texttt{ } \\
 & & \texttt{ } \\
 & & \texttt{ } \\
 & & \texttt{ } \\
 & & \texttt{ } \\
 & & \texttt{ } \\
 & & \texttt{ } \\
 & & \texttt{ } \\
 & & \texttt{ } \\
 & & \texttt{ } \\
 & & \texttt{ } \\
 & & \texttt{ } \\
 & & \texttt{ } \\
 & & \texttt{ } \\
 & & \texttt{ } \\
 & & \texttt{ } \\
 & & \texttt{ } \\
 & & \texttt{ } \\
 & & \texttt{ } \\
 & & \texttt{ } \\
 & & \texttt{ } \\
 & & \texttt{ } \\
 & & \texttt{ } \\
 & & \texttt{ } \\
 & & \texttt{ } \\
 & & \texttt{ } \\
 & & \texttt{ } \\
 & & \texttt{ } \\
 & & \texttt{ } \\
 & & \texttt{ } \\
 & & \texttt{ } \\
 & & \texttt{ } \\
 & & \texttt{ } \\
 & & \texttt{ } \\
 & & \texttt{1 } \\
 & & \texttt{1 } \\
 & & \texttt{1 } \\
 & & \texttt{1 } \\
 & & \texttt{1 } \\
 & & \texttt{1 } \\
 & & \texttt{1 } \\
 & & \texttt{1 } \\
 & & \texttt{1 } \\
 & & \texttt{1 } \\
 & & \texttt{1 } \\
 & & \texttt{1 } \\
 & & \texttt{1 } \\
 & & \texttt{1 } \\
 & & \texttt{1 } \\
 & & \texttt{1 } \\
 & & \texttt{1 } \\
 & & \texttt{1 } \\
 & & \texttt{1 } \\
 & & \texttt{1 } \\
 & & \texttt{1 } \\
 & & \texttt{1 } \\
 & & \texttt{1 } \\
 & & \texttt{1 } \\
 & & \texttt{1 } \\
 & & \texttt{1 } \\
 & & \texttt{1 } \\
 & & \texttt{1 } \\
 & & \texttt{1 } \\
 & & \texttt{1 } \\
 & & \texttt{1 } \\
 & & \texttt{1 } \\
 & & \texttt{1 } \\
 & & \texttt{1 } \\
 & & \texttt{1 } \\
 & & \texttt{1 } \\
 & & \texttt{ } \\
 & & \texttt{ } \\
 & & \texttt{ } \\
 & & \texttt{ } \\
 & & \texttt{ } \\
 & & \texttt{ } \\
 & & \texttt{ } \\
 & & \texttt{ } \\
 & & \texttt{ } \\
 & & \texttt{ } \\
 & & \texttt{ } \\
 & & \texttt{ } \\
 & & \texttt{ } \\
 & & \texttt{ } \\
 & & \texttt{ } \\
 & & \texttt{ } \\
 & & \texttt{ } \\
 & & \texttt{ } \\
 & & \texttt{ } \\
 & & \texttt{ } \\
 & & \texttt{ } \\
 & & \texttt{ } \\
 & & \texttt{ } \\
 & & \texttt{ } \\
 & & \texttt{ } \\
 & & \texttt{ } \\
 & & \texttt{ } \\
 & & \texttt{ } \\
 & & \texttt{ } \\
 & & \texttt{ } \\
 & & \texttt{ } \\
 & & \texttt{ } \\
 & & \texttt{ } \\
 & & \texttt{ } \\
 & & \texttt{ } \\
 & & \texttt{ } \\
 & & \texttt{ } \\
 & & \texttt{ } \\
 & & \texttt{ } \\
 & & \texttt{ } \\
 & & \texttt{ } \\
 & & \texttt{ } \\
 & & \texttt{ } \\
 & & \texttt{ } \\
 & & \texttt{ } \\
 & & \texttt{ } \\
 & & \texttt{ } \\
 & & \texttt{ } \\
 & & \texttt{ } \\
 & & \texttt{ } \\
 & & \texttt{ } \\
 & & \texttt{ } \\
 & & \texttt{ } \\
 & & \texttt{ } \\
 & & \texttt{ } \\
 & & \texttt{ } \\
 & & \texttt{ } \\
 & & \texttt{ } \\
 & & \texttt{ } \\
 & & \texttt{ } \\
 & & \texttt{ } \\
 & & \texttt{ } \\
 & & \texttt{ } \\
 & & \texttt{ } \\
 & & \texttt{ } \\
 & & \texttt{ } \\
 & & \texttt{ } \\
 & & \texttt{ } \\
 & & \texttt{ } \\
 & & \texttt{ } \\
 & & \texttt{ } \\
 & & \texttt{ } \\
 & & \texttt{ } \\
 & & \texttt{ } \\
 & & \texttt{ } \\
 & & \texttt{ } \\
 & & \texttt{ } \\
 & & \texttt{ } \\
 & & \texttt{ } \\
 & & \texttt{ } \\
 & & \texttt{ } \\
 & & \texttt{ } \\
 & & \texttt{ } \\
 & & \texttt{ } \\
 & & \texttt{ } \\
 & & \texttt{ } \\
 & & \texttt{ } \\
 & & \texttt{ } \\
 & & \texttt{ } \\
 & & \texttt{ } \\
 & & \texttt{ } \\
 & & \texttt{ } \\
 & & \texttt{ } \\
 & & \texttt{ } \\
 & & \texttt{ } \\
 & & \texttt{ } \\
 & & \texttt{ } \\
 & & \texttt{ } \\
 & & \texttt{ } \\
 & & \texttt{ } \\
 & & \texttt{ } \\
 & & \texttt{ } \\
 & & \texttt{ } \\
 & & \texttt{ } \\
 & & \texttt{ } \\
 & & \texttt{ } \\
 & & \texttt{ } \\
 & & \texttt{ } \\
 & & \texttt{ } \\
 & & \texttt{ } \\
 & & \texttt{ } \\
 & & \texttt{ } \\
 & & \texttt{ } \\
 & & \texttt{ } \\
 & & \texttt{ } \\
 & & \texttt{ } \\
 & & \texttt{ } \\
 & & \texttt{ } \\
 & & \texttt{ } \\
 & & \texttt{ } \\
 & & \texttt{ } \\
 & & \texttt{ } \\
 & & \texttt{ } \\
 & & \texttt{ } \\
 & & \texttt{ } \\
 & & \texttt{ } \\
 & & \texttt{ } \\
 & & \texttt{ } \\
 & & \texttt{ } \\
 & & \texttt{ } \\
 & & \texttt{ } \\
 & & \texttt{ } \\
 & & \texttt{ } \\
 & & \texttt{ } \\
 & & \texttt{ } \\
 & & \texttt{ } \\
 & & \texttt{ } \\
 & & \texttt{ } \\
 & & \texttt{ } \\
 & & \texttt{ } \\
 & & \texttt{ } \\
 & & \texttt{ } \\
 & & \texttt{ } \\
 & & \texttt{ } \\
 & & \texttt{ } \\
 & & \texttt{ } \\
 & & \texttt{ } \\
 & & \texttt{ } \\
 & & \texttt{ } \\
 & & \texttt{ } \\
 & & \texttt{ } \\
 & & \texttt{ } \\
 & & \texttt{ } \\
 & & \texttt{ } \\
 & & \texttt{ } \\
 & & \texttt{ } \\
 & & \texttt{ } \\
 & & \texttt{ } \\
 & & \texttt{ } \\
 & & \texttt{ } \\
 & & \texttt{ } \\
 & & \texttt{ } \\
 & & \texttt{ } \\
 & & \texttt{ } \\
 & & \texttt{ } \\
 & & \texttt{ } \\
 & & \texttt{ } \\
 & & \texttt{ } \\
 & & \texttt{ } \\
 & & \texttt{ } \\
 & & \texttt{ } \\
 & & \texttt{ } \\
 & & \texttt{ } \\
 & & \texttt{ } \\
 & & \texttt{ } \\
 & & \texttt{ } \\
 & & \texttt{ } \\
 & & \texttt{ } \\
 & & \texttt{ } \\
 & & \texttt{ } \\
 & & \texttt{ } \\
 & & \texttt{ } \\
 & & \texttt{ } \\
 & & \texttt{ } \\
 & & \texttt{ } \\
 & & \texttt{ } \\
 & & \texttt{ } \\
 & & \texttt{ } \\
 & & \texttt{ } \\
 & & \texttt{ } \\
 & & \texttt{ } \\
 & & \texttt{ } \\
 & & \texttt{ } \\
 & & \texttt{ } \\
 & & \texttt{ } \\
 & & \texttt{ } \\
 & & \texttt{ } \\
 & & \texttt{ } \\
 & & \texttt{ } \\
 & & \texttt{ } \\
 & & \texttt{ } \\
 & & \texttt{ } \\
 & & \texttt{ } \\
 & & \texttt{ } \\
 & & \texttt{ } \\
 & & \texttt{ } \\
 & & \texttt{ } \\
 & & \texttt{ } \\
 & & \texttt{ } \\
 & & \texttt{ } \\
 & & \texttt{ } \\
 & & \texttt{ } \\
 & & \texttt{ } \\
 & & \texttt{ } \\
 & & \texttt{ } \\
 & & \texttt{ } \\
 & & \texttt{ } \\
 & & \texttt{ } \\
 & & \texttt{ } \\
 & & \texttt{ } \\
 & & \texttt{ } \\
 & & \texttt{ } \\
 & & \texttt{ } \\
 & & \texttt{ } \\
 & & \texttt{ } \\
 & & \texttt{ } \\
 & & \texttt{ } \\
 & & \texttt{ } \\
 & & \texttt{ } \\
 & & \texttt{ } \\
 & & \texttt{ } \\
 & & \texttt{ } \\
 & & \texttt{ } \\
 & & \texttt{ } \\
 & & \texttt{ } \\
 & & \texttt{ } \\
 & & \texttt{ } \\
 & & \texttt{ } \\
 & & \texttt{ } \\
 & & \texttt{ } \\
 & & \texttt{ } \\
 & & \texttt{ } \\
 & & \texttt{ } \\
 & & \texttt{ } \\
 & & \texttt{ } \\
 & & \texttt{ } \\
 & & \texttt{ } \\
 & & \texttt{ } \\
 & & \texttt{ } \\
 & & \texttt{ } \\
 & & \texttt{ } \\
 & & \texttt{ } \\
 & & \texttt{ } \\
 & & \texttt{ } \\
 & & \texttt{ } \\
 & & \texttt{ } \\
 & & \texttt{ } \\
 & & \texttt{ } \\
 & & \texttt{ } \\
 & & \texttt{ } \\
 & & \texttt{ } \\
 & & \texttt{ } \\
 & & \texttt{ } \\
 & & \texttt{ } \\
 & & \texttt{ } \\
\hline \end{tabular}
}
\newpage
{\scriptsize
\begin{tabular}{|l|l|l|} \hline
\bf{Register Class} & \bf{Register File} & \bf{Registers} \\ \hline
 onlyswapReg & SFR & \texttt{ } \\
 & & \texttt{ } \\
 & & \texttt{ } \\
 & & \texttt{ } \\
 & & \texttt{ } \\
 & & \texttt{ } \\
 & & \texttt{ } \\
 & & \texttt{ } \\
 & & \texttt{ } \\
 & & \texttt{ } \\
 & & \texttt{ } \\
 & & \texttt{ } \\
 & & \texttt{ } \\
 & & \texttt{ } \\
 & & \texttt{ } \\
 & & \texttt{ } \\
 & & \texttt{ } \\
 & & \texttt{ } \\
 & & \texttt{ } \\
 & & \texttt{ } \\
 & & \texttt{ } \\
 & & \texttt{ } \\
 & & \texttt{ } \\
 & & \texttt{ } \\
 & & \texttt{ } \\
 & & \texttt{ } \\
 & & \texttt{ } \\
 & & \texttt{ } \\
 & & \texttt{ } \\
 & & \texttt{ } \\
 & & \texttt{ } \\
 & & \texttt{ } \\
 & & \texttt{ } \\
 & & \texttt{ } \\
 & & \texttt{ } \\
 & & \texttt{ } \\
 & & \texttt{ } \\
 & & \texttt{ } \\
 & & \texttt{ } \\
 & & \texttt{ } \\
 & & \texttt{ } \\
 & & \texttt{ } \\
 & & \texttt{ } \\
 & & \texttt{ } \\
 & & \texttt{ } \\
 & & \texttt{ } \\
 & & \texttt{ } \\
 & & \texttt{ } \\
 & & \texttt{ } \\
 & & \texttt{ } \\
 & & \texttt{ } \\
 & & \texttt{ } \\
 & & \texttt{ } \\
 & & \texttt{ } \\
 & & \texttt{ } \\
 & & \texttt{ } \\
 & & \texttt{ } \\
 & & \texttt{ } \\
 & & \texttt{ } \\
 & & \texttt{ } \\
 & & \texttt{ } \\
 & & \texttt{ } \\
 & & \texttt{ } \\
 & & \texttt{1 } \\
 & & \texttt{1 } \\
 & & \texttt{1 } \\
 & & \texttt{1 } \\
 & & \texttt{1 } \\
 & & \texttt{1 } \\
 & & \texttt{1 } \\
 & & \texttt{1 } \\
 & & \texttt{1 } \\
 & & \texttt{1 } \\
 & & \texttt{1 } \\
 & & \texttt{1 } \\
 & & \texttt{1 } \\
 & & \texttt{1 } \\
 & & \texttt{1 } \\
 & & \texttt{1 } \\
 & & \texttt{1 } \\
 & & \texttt{1 } \\
 & & \texttt{1 } \\
 & & \texttt{1 } \\
 & & \texttt{1 } \\
 & & \texttt{1 } \\
 & & \texttt{1 } \\
 & & \texttt{1 } \\
 & & \texttt{1 } \\
 & & \texttt{1 } \\
 & & \texttt{1 } \\
 & & \texttt{1 } \\
 & & \texttt{1 } \\
 & & \texttt{1 } \\
 & & \texttt{1 } \\
 & & \texttt{1 } \\
 & & \texttt{ } \\
 & & \texttt{ } \\
 & & \texttt{ } \\
 & & \texttt{ } \\
 & & \texttt{ } \\
 & & \texttt{ } \\
 & & \texttt{ } \\
 & & \texttt{ } \\
 & & \texttt{1 } \\
 & & \texttt{1 } \\
 & & \texttt{1 } \\
 & & \texttt{1 } \\
 & & \texttt{ } \\
 & & \texttt{ } \\
 & & \texttt{ } \\
 & & \texttt{ } \\
 & & \texttt{ } \\
 & & \texttt{ } \\
 & & \texttt{ } \\
 & & \texttt{ } \\
 & & \texttt{ } \\
 & & \texttt{ } \\
 & & \texttt{ } \\
 & & \texttt{ } \\
 & & \texttt{ } \\
 & & \texttt{ } \\
 & & \texttt{ } \\
 & & \texttt{ } \\
 & & \texttt{ } \\
 & & \texttt{ } \\
 & & \texttt{ } \\
 & & \texttt{ } \\
 & & \texttt{ } \\
 & & \texttt{ } \\
 & & \texttt{ } \\
 & & \texttt{ } \\
 & & \texttt{ } \\
 & & \texttt{ } \\
 & & \texttt{ } \\
 & & \texttt{ } \\
 & & \texttt{ } \\
 & & \texttt{ } \\
 & & \texttt{ } \\
 & & \texttt{ } \\
 & & \texttt{ } \\
 & & \texttt{ } \\
 & & \texttt{ } \\
 & & \texttt{ } \\
 & & \texttt{ } \\
 & & \texttt{ } \\
 & & \texttt{ } \\
 & & \texttt{ } \\
 & & \texttt{ } \\
 & & \texttt{ } \\
 & & \texttt{ } \\
 & & \texttt{ } \\
 & & \texttt{ } \\
 & & \texttt{ } \\
 & & \texttt{ } \\
 & & \texttt{ } \\
 & & \texttt{ } \\
 & & \texttt{ } \\
 & & \texttt{ } \\
 & & \texttt{ } \\
 & & \texttt{ } \\
 & & \texttt{ } \\
 & & \texttt{ } \\
 & & \texttt{ } \\
 & & \texttt{ } \\
 & & \texttt{ } \\
 & & \texttt{ } \\
 & & \texttt{ } \\
 & & \texttt{ } \\
 & & \texttt{ } \\
 & & \texttt{ } \\
 & & \texttt{ } \\
 & & \texttt{ } \\
 & & \texttt{ } \\
 & & \texttt{ } \\
 & & \texttt{ } \\
 & & \texttt{ } \\
 & & \texttt{ } \\
 & & \texttt{ } \\
 & & \texttt{ } \\
 & & \texttt{ } \\
 & & \texttt{ } \\
 & & \texttt{ } \\
 & & \texttt{ } \\
 & & \texttt{ } \\
 & & \texttt{ } \\
 & & \texttt{ } \\
 & & \texttt{ } \\
 & & \texttt{ } \\
 & & \texttt{ } \\
 & & \texttt{ } \\
 & & \texttt{ } \\
 & & \texttt{ } \\
 & & \texttt{ } \\
 & & \texttt{ } \\
 & & \texttt{ } \\
 & & \texttt{ } \\
 & & \texttt{ } \\
 & & \texttt{ } \\
 & & \texttt{ } \\
 & & \texttt{ } \\
 & & \texttt{ } \\
 & & \texttt{ } \\
 & & \texttt{ } \\
 & & \texttt{ } \\
 & & \texttt{ } \\
 & & \texttt{ } \\
 & & \texttt{ } \\
 & & \texttt{ } \\
 & & \texttt{ } \\
 & & \texttt{ } \\
 & & \texttt{ } \\
 & & \texttt{ } \\
 & & \texttt{ } \\
 & & \texttt{ } \\
 & & \texttt{ } \\
 & & \texttt{ } \\
 & & \texttt{ } \\
 & & \texttt{ } \\
 & & \texttt{ } \\
 & & \texttt{ } \\
 & & \texttt{ } \\
 & & \texttt{ } \\
 & & \texttt{ } \\
 & & \texttt{ } \\
 & & \texttt{ } \\
 & & \texttt{ } \\
 & & \texttt{ } \\
 & & \texttt{ } \\
 & & \texttt{ } \\
 & & \texttt{ } \\
 & & \texttt{ } \\
 & & \texttt{ } \\
 & & \texttt{ } \\
 & & \texttt{ } \\
 & & \texttt{ } \\
 & & \texttt{ } \\
 & & \texttt{ } \\
 & & \texttt{ } \\
 & & \texttt{ } \\
 & & \texttt{ } \\
 & & \texttt{ } \\
 & & \texttt{ } \\
 & & \texttt{ } \\
 & & \texttt{ } \\
 & & \texttt{ } \\
 & & \texttt{ } \\
 & & \texttt{ } \\
 & & \texttt{ } \\
 & & \texttt{ } \\
 & & \texttt{ } \\
 & & \texttt{ } \\
 & & \texttt{ } \\
 & & \texttt{ } \\
 & & \texttt{ } \\
 & & \texttt{ } \\
 & & \texttt{ } \\
 & & \texttt{ } \\
 & & \texttt{ } \\
 & & \texttt{ } \\
 & & \texttt{ } \\
 & & \texttt{ } \\
 & & \texttt{ } \\
 & & \texttt{ } \\
 & & \texttt{ } \\
 & & \texttt{ } \\
 & & \texttt{ } \\
 & & \texttt{ } \\
 & & \texttt{ } \\
 & & \texttt{ } \\
 & & \texttt{ } \\
 & & \texttt{ } \\
 & & \texttt{ } \\
 & & \texttt{ } \\
 & & \texttt{ } \\
 & & \texttt{ } \\
 & & \texttt{ } \\
 & & \texttt{ } \\
 & & \texttt{ } \\
 & & \texttt{ } \\
 & & \texttt{ } \\
 & & \texttt{ } \\
 & & \texttt{ } \\
 & & \texttt{ } \\
 & & \texttt{ } \\
 & & \texttt{ } \\
 & & \texttt{ } \\
 & & \texttt{ } \\
 & & \texttt{ } \\
 & & \texttt{ } \\
 & & \texttt{ } \\
 & & \texttt{ } \\
 & & \texttt{ } \\
 & & \texttt{ } \\
 & & \texttt{ } \\
 & & \texttt{ } \\
 & & \texttt{ } \\
 & & \texttt{ } \\
 & & \texttt{ } \\
 & & \texttt{ } \\
 & & \texttt{ } \\
 & & \texttt{ } \\
 & & \texttt{ } \\
 & & \texttt{ } \\
 & & \texttt{ } \\
 & & \texttt{ } \\
 & & \texttt{ } \\
 & & \texttt{ } \\
 & & \texttt{ } \\
 & & \texttt{ } \\
 & & \texttt{ } \\
 & & \texttt{ } \\
 & & \texttt{ } \\
 & & \texttt{ } \\
 & & \texttt{ } \\
 & & \texttt{ } \\
 & & \texttt{ } \\
 & & \texttt{ } \\
 & & \texttt{ } \\
 & & \texttt{ } \\
 & & \texttt{ } \\
 & & \texttt{ } \\
 & & \texttt{ } \\
 & & \texttt{ } \\
 & & \texttt{ } \\
 & & \texttt{ } \\
 & & \texttt{ } \\
 & & \texttt{ } \\
 & & \texttt{ } \\
 & & \texttt{ } \\
 & & \texttt{ } \\
 & & \texttt{ } \\
 & & \texttt{ } \\
 & & \texttt{ } \\
 & & \texttt{ } \\
 & & \texttt{ } \\
 & & \texttt{ } \\
 & & \texttt{ } \\
 & & \texttt{ } \\
 & & \texttt{ } \\
 & & \texttt{ } \\
 & & \texttt{ } \\
 & & \texttt{ } \\
 & & \texttt{ } \\
 & & \texttt{ } \\
 & & \texttt{ } \\
 & & \texttt{ } \\
 & & \texttt{ } \\
 & & \texttt{ } \\
 & & \texttt{ } \\
 & & \texttt{ } \\
 & & \texttt{ } \\
 & & \texttt{ } \\
 & & \texttt{ } \\
 & & \texttt{ } \\
 & & \texttt{ } \\
 & & \texttt{ } \\
 & & \texttt{ } \\
 & & \texttt{ } \\
 & & \texttt{ } \\
 & & \texttt{ } \\
 & & \texttt{ } \\
 & & \texttt{ } \\
 & & \texttt{ } \\
 & & \texttt{ } \\
 & & \texttt{ } \\
 & & \texttt{ } \\
 & & \texttt{ } \\
 & & \texttt{ } \\
 & & \texttt{ } \\
 & & \texttt{ } \\
 & & \texttt{ } \\
 & & \texttt{ } \\
 & & \texttt{ } \\
 & & \texttt{ } \\
 & & \texttt{ } \\
 & & \texttt{ } \\
 & & \texttt{ } \\
 & & \texttt{ } \\
 & & \texttt{ } \\
 & & \texttt{ } \\
 & & \texttt{ } \\
 & & \texttt{ } \\
 & & \texttt{ } \\
 & & \texttt{ } \\
 & & \texttt{ } \\
 & & \texttt{ } \\
 & & \texttt{ } \\
 & & \texttt{ } \\
 & & \texttt{ } \\
 & & \texttt{ } \\
 & & \texttt{ } \\
 & & \texttt{ } \\
 & & \texttt{ } \\
 & & \texttt{ } \\
 & & \texttt{ } \\
 & & \texttt{ } \\
 & & \texttt{ } \\
 & & \texttt{ } \\
 & & \texttt{ } \\
 & & \texttt{ } \\
 & & \texttt{ } \\
 & & \texttt{ } \\
 & & \texttt{ } \\
 & & \texttt{ } \\
 & & \texttt{ } \\
 & & \texttt{ } \\
 & & \texttt{ } \\
 & & \texttt{ } \\
 & & \texttt{ } \\
 & & \texttt{ } \\
 & & \texttt{ } \\
 & & \texttt{ } \\
 & & \texttt{ } \\
 & & \texttt{ } \\
 & & \texttt{ } \\
 & & \texttt{ } \\
 & & \texttt{ } \\
 & & \texttt{ } \\
 & & \texttt{ } \\
 & & \texttt{ } \\
 & & \texttt{ } \\
 & & \texttt{ } \\
 & & \texttt{ } \\
 & & \texttt{ } \\
 & & \texttt{ } \\
 & & \texttt{ } \\
 & & \texttt{ } \\
 & & \texttt{ } \\
 & & \texttt{ } \\
 & & \texttt{ } \\
 & & \texttt{ } \\
 & & \texttt{ } \\
 & & \texttt{ } \\
 & & \texttt{ } \\
 & & \texttt{ } \\
 & & \texttt{ } \\
 & & \texttt{ } \\
 & & \texttt{ } \\
 & & \texttt{ } \\
 & & \texttt{ } \\
 & & \texttt{ } \\
 & & \texttt{ } \\
 & & \texttt{ } \\
 & & \texttt{ } \\
 & & \texttt{ } \\
 & & \texttt{ } \\
 & & \texttt{ } \\
 & & \texttt{ } \\
 & & \texttt{ } \\
 & & \texttt{ } \\
 & & \texttt{ } \\
 & & \texttt{ } \\
 & & \texttt{ } \\
 & & \texttt{ } \\
 & & \texttt{ } \\
 & & \texttt{ } \\
 & & \texttt{ } \\
 & & \texttt{ } \\
 & & \texttt{ } \\
 & & \texttt{ } \\
 & & \texttt{ } \\
 & & \texttt{ } \\
 & & \texttt{ } \\
 & & \texttt{ } \\
 & & \texttt{ } \\
 & & \texttt{ } \\
 & & \texttt{ } \\
 & & \texttt{ } \\
 & & \texttt{ } \\
 & & \texttt{ } \\
 & & \texttt{ } \\
 & & \texttt{ } \\
 & & \texttt{ } \\
 & & \texttt{ } \\
 & & \texttt{ } \\
 & & \texttt{ } \\
 & & \texttt{ } \\
 & & \texttt{ } \\
 & & \texttt{ } \\
 & & \texttt{ } \\
 & & \texttt{ } \\
 & & \texttt{ } \\
 & & \texttt{ } \\
 & & \texttt{ } \\
 & & \texttt{ } \\
 & & \texttt{ } \\
 & & \texttt{ } \\
 & & \texttt{ } \\
 & & \texttt{ } \\
 & & \texttt{ } \\
 & & \texttt{ } \\
 & & \texttt{ } \\
 & & \texttt{ } \\
 & & \texttt{ } \\
 & & \texttt{ } \\
 & & \texttt{ } \\
 & & \texttt{ } \\
 & & \texttt{ } \\
 & & \texttt{ } \\
 & & \texttt{ } \\
 & & \texttt{ } \\
 & & \texttt{ } \\
 & & \texttt{ } \\
 & & \texttt{ } \\
 & & \texttt{ } \\
 & & \texttt{ } \\
 & & \texttt{ } \\
 & & \texttt{ } \\
 & & \texttt{ } \\
 & & \texttt{ } \\
 & & \texttt{ } \\
\hline \end{tabular}
}
\newpage
{\scriptsize
\begin{tabular}{|l|l|l|} \hline
\bf{Register Class} & \bf{Register File} & \bf{Registers} \\ \hline
 pairedReg & PGR & \texttt{ } \\
 & & \texttt{ } \\
 & & \texttt{ } \\
 & & \texttt{ } \\
 & & \texttt{ } \\
 & & \texttt{ } \\
 & & \texttt{ } \\
 & & \texttt{ } \\
 & & \texttt{ } \\
 & & \texttt{ } \\
 & & \texttt{ } \\
 & & \texttt{ } \\
 & & \texttt{ } \\
 & & \texttt{ } \\
 & & \texttt{ } \\
 & & \texttt{ } \\
 & & \texttt{ } \\
 & & \texttt{ } \\
 & & \texttt{ } \\
 & & \texttt{ } \\
 & & \texttt{ } \\
 & & \texttt{ } \\
 & & \texttt{ } \\
 & & \texttt{ } \\
 & & \texttt{ } \\
 & & \texttt{ } \\
 & & \texttt{ } \\
 & & \texttt{ } \\
 & & \texttt{ } \\
 & & \texttt{ } \\
 & & \texttt{ } \\
 & & \texttt{ } \\
\hline \end{tabular}
}
\newpage
{\scriptsize
\begin{tabular}{|l|l|l|} \hline
\bf{Register Class} & \bf{Register File} & \bf{Registers} \\ \hline
 quadReg & QGR & \texttt{} \\
 & & \texttt{} \\
 & & \texttt{} \\
 & & \texttt{} \\
 & & \texttt{} \\
 & & \texttt{} \\
 & & \texttt{} \\
 & & \texttt{} \\
 & & \texttt{} \\
 & & \texttt{} \\
 & & \texttt{} \\
 & & \texttt{} \\
 & & \texttt{} \\
 & & \texttt{} \\
 & & \texttt{} \\
 & & \texttt{} \\
\hline \end{tabular}
}
\newpage
{\scriptsize
\begin{tabular}{|l|l|l|} \hline
\bf{Register Class} & \bf{Register File} & \bf{Registers} \\ \hline
 singleReg & GPR & \texttt{  } \\
 & & \texttt{  } \\
 & & \texttt{  } \\
 & & \texttt{  } \\
 & & \texttt{  } \\
 & & \texttt{  } \\
 & & \texttt{  } \\
 & & \texttt{  } \\
 & & \texttt{  } \\
 & & \texttt{  } \\
 & & \texttt{  } \\
 & & \texttt{  } \\
 & & \texttt{ } \\
 & & \texttt{ } \\
 & & \texttt{  } \\
 & & \texttt{  } \\
 & & \texttt{  } \\
 & & \texttt{  } \\
 & & \texttt{  } \\
 & & \texttt{  } \\
 & & \texttt{  } \\
 & & \texttt{  } \\
 & & \texttt{  } \\
 & & \texttt{  } \\
 & & \texttt{  } \\
 & & \texttt{  } \\
 & & \texttt{  } \\
 & & \texttt{  } \\
 & & \texttt{  } \\
 & & \texttt{  } \\
 & & \texttt{  } \\
 & & \texttt{  } \\
 & & \texttt{  } \\
 & & \texttt{  } \\
 & & \texttt{  } \\
 & & \texttt{  } \\
 & & \texttt{  } \\
 & & \texttt{  } \\
 & & \texttt{  } \\
 & & \texttt{  } \\
 & & \texttt{  } \\
 & & \texttt{  } \\
 & & \texttt{  } \\
 & & \texttt{  } \\
 & & \texttt{  } \\
 & & \texttt{  } \\
 & & \texttt{  } \\
 & & \texttt{  } \\
 & & \texttt{  } \\
 & & \texttt{  } \\
 & & \texttt{  } \\
 & & \texttt{  } \\
 & & \texttt{  } \\
 & & \texttt{  } \\
 & & \texttt{  } \\
 & & \texttt{  } \\
 & & \texttt{  } \\
 & & \texttt{  } \\
 & & \texttt{  } \\
 & & \texttt{  } \\
 & & \texttt{  } \\
 & & \texttt{  } \\
 & & \texttt{  } \\
 & & \texttt{  } \\
\hline \end{tabular}
}
\newpage
{\scriptsize
\begin{tabular}{|l|l|l|} \hline
\bf{Register Class} & \bf{Register File} & \bf{Registers} \\ \hline
 worddRegE & GPR & \texttt{  } \\
 & & \texttt{  } \\
 & & \texttt{  } \\
 & & \texttt{  } \\
 & & \texttt{  } \\
 & & \texttt{  } \\
 & & \texttt{ } \\
 & & \texttt{  } \\
 & & \texttt{  } \\
 & & \texttt{  } \\
 & & \texttt{  } \\
 & & \texttt{  } \\
 & & \texttt{  } \\
 & & \texttt{  } \\
 & & \texttt{  } \\
 & & \texttt{  } \\
 & & \texttt{  } \\
 & & \texttt{  } \\
 & & \texttt{  } \\
 & & \texttt{  } \\
 & & \texttt{  } \\
 & & \texttt{  } \\
 & & \texttt{  } \\
 & & \texttt{  } \\
 & & \texttt{  } \\
 & & \texttt{  } \\
 & & \texttt{  } \\
 & & \texttt{  } \\
 & & \texttt{  } \\
 & & \texttt{  } \\
 & & \texttt{  } \\
 & & \texttt{  } \\
\hline \end{tabular}
}
\newpage
{\scriptsize
\begin{tabular}{|l|l|l|} \hline
\bf{Register Class} & \bf{Register File} & \bf{Registers} \\ \hline
 worddRegO & GPR & \texttt{  } \\
 & & \texttt{  } \\
 & & \texttt{  } \\
 & & \texttt{  } \\
 & & \texttt{  } \\
 & & \texttt{  } \\
 & & \texttt{ } \\
 & & \texttt{  } \\
 & & \texttt{  } \\
 & & \texttt{  } \\
 & & \texttt{  } \\
 & & \texttt{  } \\
 & & \texttt{  } \\
 & & \texttt{  } \\
 & & \texttt{  } \\
 & & \texttt{  } \\
 & & \texttt{  } \\
 & & \texttt{  } \\
 & & \texttt{  } \\
 & & \texttt{  } \\
 & & \texttt{  } \\
 & & \texttt{  } \\
 & & \texttt{  } \\
 & & \texttt{  } \\
 & & \texttt{  } \\
 & & \texttt{  } \\
 & & \texttt{  } \\
 & & \texttt{  } \\
 & & \texttt{  } \\
 & & \texttt{  } \\
 & & \texttt{  } \\
 & & \texttt{  } \\
\hline \end{tabular}
}
\newpage
{\scriptsize
\begin{tabular}{|l|l|l|} \hline
\bf{Register Class} & \bf{Register File} & \bf{Registers} \\ \hline
 xworddReg & XCR & \texttt{1 } \\
 & & \texttt{1 } \\
 & & \texttt{1 } \\
 & & \texttt{1 } \\
 & & \texttt{1 } \\
 & & \texttt{1 } \\
 & & \texttt{1 } \\
 & & \texttt{1 } \\
 & & \texttt{1 } \\
 & & \texttt{1 } \\
 & & \texttt{1 } \\
 & & \texttt{1 } \\
 & & \texttt{1 } \\
 & & \texttt{1 } \\
 & & \texttt{1 } \\
 & & \texttt{1 } \\
 & & \texttt{1 } \\
 & & \texttt{1 } \\
 & & \texttt{1 } \\
 & & \texttt{1 } \\
 & & \texttt{1 } \\
 & & \texttt{1 } \\
 & & \texttt{1 } \\
 & & \texttt{1 } \\
 & & \texttt{1 } \\
 & & \texttt{1 } \\
 & & \texttt{1 } \\
 & & \texttt{1 } \\
 & & \texttt{1 } \\
 & & \texttt{1 } \\
 & & \texttt{1 } \\
 & & \texttt{1 } \\
 & & \texttt{1 } \\
 & & \texttt{1 } \\
 & & \texttt{1 } \\
 & & \texttt{1 } \\
 & & \texttt{1 } \\
 & & \texttt{1 } \\
 & & \texttt{1 } \\
 & & \texttt{1 } \\
 & & \texttt{1 } \\
 & & \texttt{1 } \\
 & & \texttt{1 } \\
 & & \texttt{1 } \\
 & & \texttt{1 } \\
 & & \texttt{1 } \\
 & & \texttt{1 } \\
 & & \texttt{1 } \\
 & & \texttt{1 } \\
 & & \texttt{1 } \\
 & & \texttt{1 } \\
 & & \texttt{1 } \\
 & & \texttt{1 } \\
 & & \texttt{1 } \\
 & & \texttt{1 } \\
 & & \texttt{1 } \\
 & & \texttt{1 } \\
 & & \texttt{1 } \\
 & & \texttt{1 } \\
 & & \texttt{1 } \\
 & & \texttt{1 } \\
 & & \texttt{1 } \\
 & & \texttt{1 } \\
 & & \texttt{1 } \\
 & & \texttt{1 } \\
 & & \texttt{1 } \\
 & & \texttt{1 } \\
 & & \texttt{1 } \\
 & & \texttt{1 } \\
 & & \texttt{1 } \\
 & & \texttt{1 } \\
 & & \texttt{1 } \\
 & & \texttt{1 } \\
 & & \texttt{1 } \\
 & & \texttt{1 } \\
 & & \texttt{1 } \\
 & & \texttt{1 } \\
 & & \texttt{1 } \\
 & & \texttt{1 } \\
 & & \texttt{1 } \\
 & & \texttt{1 } \\
 & & \texttt{1 } \\
 & & \texttt{1 } \\
 & & \texttt{1 } \\
 & & \texttt{1 } \\
 & & \texttt{1 } \\
 & & \texttt{1 } \\
 & & \texttt{1 } \\
 & & \texttt{1 } \\
 & & \texttt{1 } \\
 & & \texttt{1 } \\
 & & \texttt{1 } \\
 & & \texttt{1 } \\
 & & \texttt{1 } \\
 & & \texttt{1 } \\
 & & \texttt{1 } \\
 & & \texttt{1 } \\
 & & \texttt{1 } \\
 & & \texttt{1 } \\
 & & \texttt{1 } \\
 & & \texttt{1 } \\
 & & \texttt{1 } \\
 & & \texttt{1 } \\
 & & \texttt{1 } \\
 & & \texttt{1 } \\
 & & \texttt{1 } \\
 & & \texttt{1 } \\
 & & \texttt{1 } \\
 & & \texttt{1 } \\
 & & \texttt{1 } \\
 & & \texttt{1 } \\
 & & \texttt{1 } \\
 & & \texttt{1 } \\
 & & \texttt{1 } \\
 & & \texttt{1 } \\
 & & \texttt{1 } \\
 & & \texttt{1 } \\
 & & \texttt{1 } \\
 & & \texttt{1 } \\
 & & \texttt{1 } \\
 & & \texttt{1 } \\
 & & \texttt{1 } \\
 & & \texttt{1 } \\
 & & \texttt{1 } \\
 & & \texttt{1 } \\
 & & \texttt{1 } \\
 & & \texttt{1 } \\
 & & \texttt{1 } \\
 & & \texttt{1 } \\
 & & \texttt{1 } \\
 & & \texttt{1 } \\
 & & \texttt{1 } \\
 & & \texttt{1 } \\
 & & \texttt{1 } \\
 & & \texttt{1 } \\
 & & \texttt{1 } \\
 & & \texttt{1 } \\
 & & \texttt{1 } \\
 & & \texttt{1 } \\
 & & \texttt{1 } \\
 & & \texttt{1 } \\
 & & \texttt{1 } \\
 & & \texttt{1 } \\
 & & \texttt{1 } \\
 & & \texttt{1 } \\
 & & \texttt{1 } \\
 & & \texttt{1 } \\
 & & \texttt{1 } \\
 & & \texttt{1 } \\
 & & \texttt{1 } \\
 & & \texttt{1 } \\
 & & \texttt{1 } \\
 & & \texttt{1 } \\
 & & \texttt{1 } \\
 & & \texttt{1 } \\
 & & \texttt{1 } \\
 & & \texttt{1 } \\
 & & \texttt{1 } \\
 & & \texttt{1 } \\
 & & \texttt{1 } \\
 & & \texttt{1 } \\
 & & \texttt{1 } \\
 & & \texttt{1 } \\
 & & \texttt{1 } \\
 & & \texttt{1 } \\
 & & \texttt{1 } \\
 & & \texttt{1 } \\
 & & \texttt{1 } \\
 & & \texttt{1 } \\
 & & \texttt{1 } \\
 & & \texttt{1 } \\
 & & \texttt{1 } \\
 & & \texttt{1 } \\
 & & \texttt{1 } \\
 & & \texttt{1 } \\
 & & \texttt{1 } \\
 & & \texttt{1 } \\
 & & \texttt{1 } \\
 & & \texttt{1 } \\
 & & \texttt{1 } \\
 & & \texttt{1 } \\
 & & \texttt{1 } \\
 & & \texttt{1 } \\
 & & \texttt{1 } \\
 & & \texttt{1 } \\
 & & \texttt{1 } \\
 & & \texttt{1 } \\
 & & \texttt{1 } \\
 & & \texttt{1 } \\
 & & \texttt{1 } \\
 & & \texttt{1 } \\
 & & \texttt{1 } \\
 & & \texttt{1 } \\
 & & \texttt{1 } \\
 & & \texttt{1 } \\
 & & \texttt{1 } \\
 & & \texttt{1 } \\
 & & \texttt{1 } \\
 & & \texttt{1 } \\
 & & \texttt{1 } \\
 & & \texttt{1 } \\
 & & \texttt{1 } \\
 & & \texttt{1 } \\
 & & \texttt{1 } \\
 & & \texttt{1 } \\
 & & \texttt{1 } \\
 & & \texttt{1 } \\
 & & \texttt{1 } \\
 & & \texttt{1 } \\
 & & \texttt{1 } \\
 & & \texttt{1 } \\
 & & \texttt{1 } \\
 & & \texttt{1 } \\
 & & \texttt{1 } \\
 & & \texttt{1 } \\
 & & \texttt{1 } \\
 & & \texttt{1 } \\
 & & \texttt{1 } \\
 & & \texttt{1 } \\
 & & \texttt{1 } \\
 & & \texttt{1 } \\
 & & \texttt{1 } \\
 & & \texttt{1 } \\
 & & \texttt{1 } \\
 & & \texttt{1 } \\
 & & \texttt{1 } \\
 & & \texttt{1 } \\
 & & \texttt{1 } \\
 & & \texttt{1 } \\
 & & \texttt{1 } \\
 & & \texttt{1 } \\
 & & \texttt{1 } \\
 & & \texttt{1 } \\
 & & \texttt{1 } \\
 & & \texttt{1 } \\
 & & \texttt{1 } \\
 & & \texttt{1 } \\
 & & \texttt{1 } \\
 & & \texttt{1 } \\
 & & \texttt{1 } \\
 & & \texttt{1 } \\
 & & \texttt{1 } \\
 & & \texttt{1 } \\
 & & \texttt{1 } \\
 & & \texttt{1 } \\
 & & \texttt{1 } \\
 & & \texttt{1 } \\
 & & \texttt{1 } \\
 & & \texttt{1 } \\
 & & \texttt{1 } \\
 & & \texttt{1 } \\
 & & \texttt{1 } \\
 & & \texttt{1 } \\
 & & \texttt{1 } \\
 & & \texttt{1 } \\
 & & \texttt{1 } \\
\hline \end{tabular}
}
\newpage
{\scriptsize
\begin{tabular}{|l|l|l|} \hline
\bf{Register Class} & \bf{Register File} & \bf{Registers} \\ \hline
 xworddReg0M4 & XCR & \texttt{1 } \\
 & & \texttt{1 } \\
 & & \texttt{1 } \\
 & & \texttt{1 } \\
 & & \texttt{1 } \\
 & & \texttt{1 } \\
 & & \texttt{1 } \\
 & & \texttt{1 } \\
 & & \texttt{1 } \\
 & & \texttt{1 } \\
 & & \texttt{1 } \\
 & & \texttt{1 } \\
 & & \texttt{1 } \\
 & & \texttt{1 } \\
 & & \texttt{1 } \\
 & & \texttt{1 } \\
 & & \texttt{1 } \\
 & & \texttt{1 } \\
 & & \texttt{1 } \\
 & & \texttt{1 } \\
 & & \texttt{1 } \\
 & & \texttt{1 } \\
 & & \texttt{1 } \\
 & & \texttt{1 } \\
 & & \texttt{1 } \\
 & & \texttt{1 } \\
 & & \texttt{1 } \\
 & & \texttt{1 } \\
 & & \texttt{1 } \\
 & & \texttt{1 } \\
 & & \texttt{1 } \\
 & & \texttt{1 } \\
 & & \texttt{1 } \\
 & & \texttt{1 } \\
 & & \texttt{1 } \\
 & & \texttt{1 } \\
 & & \texttt{1 } \\
 & & \texttt{1 } \\
 & & \texttt{1 } \\
 & & \texttt{1 } \\
 & & \texttt{1 } \\
 & & \texttt{1 } \\
 & & \texttt{1 } \\
 & & \texttt{1 } \\
 & & \texttt{1 } \\
 & & \texttt{1 } \\
 & & \texttt{1 } \\
 & & \texttt{1 } \\
 & & \texttt{1 } \\
 & & \texttt{1 } \\
 & & \texttt{1 } \\
 & & \texttt{1 } \\
 & & \texttt{1 } \\
 & & \texttt{1 } \\
 & & \texttt{1 } \\
 & & \texttt{1 } \\
 & & \texttt{1 } \\
 & & \texttt{1 } \\
 & & \texttt{1 } \\
 & & \texttt{1 } \\
 & & \texttt{1 } \\
 & & \texttt{1 } \\
 & & \texttt{1 } \\
 & & \texttt{1 } \\
\hline \end{tabular}
}
\newpage
{\scriptsize
\begin{tabular}{|l|l|l|} \hline
\bf{Register Class} & \bf{Register File} & \bf{Registers} \\ \hline
 xworddReg1M4 & XCR & \texttt{1 } \\
 & & \texttt{1 } \\
 & & \texttt{1 } \\
 & & \texttt{1 } \\
 & & \texttt{1 } \\
 & & \texttt{1 } \\
 & & \texttt{1 } \\
 & & \texttt{1 } \\
 & & \texttt{1 } \\
 & & \texttt{1 } \\
 & & \texttt{1 } \\
 & & \texttt{1 } \\
 & & \texttt{1 } \\
 & & \texttt{1 } \\
 & & \texttt{1 } \\
 & & \texttt{1 } \\
 & & \texttt{1 } \\
 & & \texttt{1 } \\
 & & \texttt{1 } \\
 & & \texttt{1 } \\
 & & \texttt{1 } \\
 & & \texttt{1 } \\
 & & \texttt{1 } \\
 & & \texttt{1 } \\
 & & \texttt{1 } \\
 & & \texttt{1 } \\
 & & \texttt{1 } \\
 & & \texttt{1 } \\
 & & \texttt{1 } \\
 & & \texttt{1 } \\
 & & \texttt{1 } \\
 & & \texttt{1 } \\
 & & \texttt{1 } \\
 & & \texttt{1 } \\
 & & \texttt{1 } \\
 & & \texttt{1 } \\
 & & \texttt{1 } \\
 & & \texttt{1 } \\
 & & \texttt{1 } \\
 & & \texttt{1 } \\
 & & \texttt{1 } \\
 & & \texttt{1 } \\
 & & \texttt{1 } \\
 & & \texttt{1 } \\
 & & \texttt{1 } \\
 & & \texttt{1 } \\
 & & \texttt{1 } \\
 & & \texttt{1 } \\
 & & \texttt{1 } \\
 & & \texttt{1 } \\
 & & \texttt{1 } \\
 & & \texttt{1 } \\
 & & \texttt{1 } \\
 & & \texttt{1 } \\
 & & \texttt{1 } \\
 & & \texttt{1 } \\
 & & \texttt{1 } \\
 & & \texttt{1 } \\
 & & \texttt{1 } \\
 & & \texttt{1 } \\
 & & \texttt{1 } \\
 & & \texttt{1 } \\
 & & \texttt{1 } \\
 & & \texttt{1 } \\
\hline \end{tabular}
}
\newpage
{\scriptsize
\begin{tabular}{|l|l|l|} \hline
\bf{Register Class} & \bf{Register File} & \bf{Registers} \\ \hline
 xworddReg2M4 & XCR & \texttt{1 } \\
 & & \texttt{1 } \\
 & & \texttt{1 } \\
 & & \texttt{1 } \\
 & & \texttt{1 } \\
 & & \texttt{1 } \\
 & & \texttt{1 } \\
 & & \texttt{1 } \\
 & & \texttt{1 } \\
 & & \texttt{1 } \\
 & & \texttt{1 } \\
 & & \texttt{1 } \\
 & & \texttt{1 } \\
 & & \texttt{1 } \\
 & & \texttt{1 } \\
 & & \texttt{1 } \\
 & & \texttt{1 } \\
 & & \texttt{1 } \\
 & & \texttt{1 } \\
 & & \texttt{1 } \\
 & & \texttt{1 } \\
 & & \texttt{1 } \\
 & & \texttt{1 } \\
 & & \texttt{1 } \\
 & & \texttt{1 } \\
 & & \texttt{1 } \\
 & & \texttt{1 } \\
 & & \texttt{1 } \\
 & & \texttt{1 } \\
 & & \texttt{1 } \\
 & & \texttt{1 } \\
 & & \texttt{1 } \\
 & & \texttt{1 } \\
 & & \texttt{1 } \\
 & & \texttt{1 } \\
 & & \texttt{1 } \\
 & & \texttt{1 } \\
 & & \texttt{1 } \\
 & & \texttt{1 } \\
 & & \texttt{1 } \\
 & & \texttt{1 } \\
 & & \texttt{1 } \\
 & & \texttt{1 } \\
 & & \texttt{1 } \\
 & & \texttt{1 } \\
 & & \texttt{1 } \\
 & & \texttt{1 } \\
 & & \texttt{1 } \\
 & & \texttt{1 } \\
 & & \texttt{1 } \\
 & & \texttt{1 } \\
 & & \texttt{1 } \\
 & & \texttt{1 } \\
 & & \texttt{1 } \\
 & & \texttt{1 } \\
 & & \texttt{1 } \\
 & & \texttt{1 } \\
 & & \texttt{1 } \\
 & & \texttt{1 } \\
 & & \texttt{1 } \\
 & & \texttt{1 } \\
 & & \texttt{1 } \\
 & & \texttt{1 } \\
 & & \texttt{1 } \\
\hline \end{tabular}
}
\newpage
{\scriptsize
\begin{tabular}{|l|l|l|} \hline
\bf{Register Class} & \bf{Register File} & \bf{Registers} \\ \hline
 xworddReg3M4 & XCR & \texttt{1 } \\
 & & \texttt{1 } \\
 & & \texttt{1 } \\
 & & \texttt{1 } \\
 & & \texttt{1 } \\
 & & \texttt{1 } \\
 & & \texttt{1 } \\
 & & \texttt{1 } \\
 & & \texttt{1 } \\
 & & \texttt{1 } \\
 & & \texttt{1 } \\
 & & \texttt{1 } \\
 & & \texttt{1 } \\
 & & \texttt{1 } \\
 & & \texttt{1 } \\
 & & \texttt{1 } \\
 & & \texttt{1 } \\
 & & \texttt{1 } \\
 & & \texttt{1 } \\
 & & \texttt{1 } \\
 & & \texttt{1 } \\
 & & \texttt{1 } \\
 & & \texttt{1 } \\
 & & \texttt{1 } \\
 & & \texttt{1 } \\
 & & \texttt{1 } \\
 & & \texttt{1 } \\
 & & \texttt{1 } \\
 & & \texttt{1 } \\
 & & \texttt{1 } \\
 & & \texttt{1 } \\
 & & \texttt{1 } \\
 & & \texttt{1 } \\
 & & \texttt{1 } \\
 & & \texttt{1 } \\
 & & \texttt{1 } \\
 & & \texttt{1 } \\
 & & \texttt{1 } \\
 & & \texttt{1 } \\
 & & \texttt{1 } \\
 & & \texttt{1 } \\
 & & \texttt{1 } \\
 & & \texttt{1 } \\
 & & \texttt{1 } \\
 & & \texttt{1 } \\
 & & \texttt{1 } \\
 & & \texttt{1 } \\
 & & \texttt{1 } \\
 & & \texttt{1 } \\
 & & \texttt{1 } \\
 & & \texttt{1 } \\
 & & \texttt{1 } \\
 & & \texttt{1 } \\
 & & \texttt{1 } \\
 & & \texttt{1 } \\
 & & \texttt{1 } \\
 & & \texttt{1 } \\
 & & \texttt{1 } \\
 & & \texttt{1 } \\
 & & \texttt{1 } \\
 & & \texttt{1 } \\
 & & \texttt{1 } \\
 & & \texttt{1 } \\
 & & \texttt{1 } \\
\hline \end{tabular}
}
\newpage
{\scriptsize
\begin{tabular}{|l|l|l|} \hline
\bf{Register Class} & \bf{Register File} & \bf{Registers} \\ \hline
 xwordoReg & XVR & \texttt{  } \\
 & & \texttt{  } \\
 & & \texttt{  } \\
 & & \texttt{  } \\
 & & \texttt{  } \\
 & & \texttt{  } \\
 & & \texttt{  } \\
 & & \texttt{  } \\
 & & \texttt{  } \\
 & & \texttt{  } \\
 & & \texttt{  } \\
 & & \texttt{  } \\
 & & \texttt{  } \\
 & & \texttt{  } \\
 & & \texttt{  } \\
 & & \texttt{  } \\
 & & \texttt{  } \\
 & & \texttt{  } \\
 & & \texttt{  } \\
 & & \texttt{  } \\
 & & \texttt{  } \\
 & & \texttt{  } \\
 & & \texttt{  } \\
 & & \texttt{  } \\
 & & \texttt{  } \\
 & & \texttt{  } \\
 & & \texttt{  } \\
 & & \texttt{  } \\
 & & \texttt{  } \\
 & & \texttt{  } \\
 & & \texttt{  } \\
 & & \texttt{  } \\
 & & \texttt{  } \\
 & & \texttt{  } \\
 & & \texttt{  } \\
 & & \texttt{  } \\
 & & \texttt{  } \\
 & & \texttt{  } \\
 & & \texttt{  } \\
 & & \texttt{  } \\
 & & \texttt{  } \\
 & & \texttt{  } \\
 & & \texttt{  } \\
 & & \texttt{  } \\
 & & \texttt{  } \\
 & & \texttt{  } \\
 & & \texttt{  } \\
 & & \texttt{  } \\
 & & \texttt{  } \\
 & & \texttt{  } \\
 & & \texttt{  } \\
 & & \texttt{  } \\
 & & \texttt{  } \\
 & & \texttt{  } \\
 & & \texttt{  } \\
 & & \texttt{  } \\
 & & \texttt{  } \\
 & & \texttt{  } \\
 & & \texttt{  } \\
 & & \texttt{  } \\
 & & \texttt{  } \\
 & & \texttt{  } \\
 & & \texttt{  } \\
 & & \texttt{  } \\
\hline \end{tabular}
}
\newpage
{\scriptsize
\begin{tabular}{|l|l|l|} \hline
\bf{Register Class} & \bf{Register File} & \bf{Registers} \\ \hline
 xwordqReg & XBR & \texttt{1 } \\
 & & \texttt{1 } \\
 & & \texttt{1 } \\
 & & \texttt{1 } \\
 & & \texttt{1 } \\
 & & \texttt{1 } \\
 & & \texttt{1 } \\
 & & \texttt{1 } \\
 & & \texttt{1 } \\
 & & \texttt{1 } \\
 & & \texttt{1 } \\
 & & \texttt{1 } \\
 & & \texttt{1 } \\
 & & \texttt{1 } \\
 & & \texttt{1 } \\
 & & \texttt{1 } \\
 & & \texttt{1 } \\
 & & \texttt{1 } \\
 & & \texttt{1 } \\
 & & \texttt{1 } \\
 & & \texttt{1 } \\
 & & \texttt{1 } \\
 & & \texttt{1 } \\
 & & \texttt{1 } \\
 & & \texttt{1 } \\
 & & \texttt{1 } \\
 & & \texttt{1 } \\
 & & \texttt{1 } \\
 & & \texttt{1 } \\
 & & \texttt{1 } \\
 & & \texttt{1 } \\
 & & \texttt{1 } \\
 & & \texttt{1 } \\
 & & \texttt{1 } \\
 & & \texttt{1 } \\
 & & \texttt{1 } \\
 & & \texttt{1 } \\
 & & \texttt{1 } \\
 & & \texttt{1 } \\
 & & \texttt{1 } \\
 & & \texttt{1 } \\
 & & \texttt{1 } \\
 & & \texttt{1 } \\
 & & \texttt{1 } \\
 & & \texttt{1 } \\
 & & \texttt{1 } \\
 & & \texttt{1 } \\
 & & \texttt{1 } \\
 & & \texttt{1 } \\
 & & \texttt{1 } \\
 & & \texttt{1 } \\
 & & \texttt{1 } \\
 & & \texttt{1 } \\
 & & \texttt{1 } \\
 & & \texttt{1 } \\
 & & \texttt{1 } \\
 & & \texttt{1 } \\
 & & \texttt{1 } \\
 & & \texttt{1 } \\
 & & \texttt{1 } \\
 & & \texttt{1 } \\
 & & \texttt{1 } \\
 & & \texttt{1 } \\
 & & \texttt{1 } \\
 & & \texttt{1 } \\
 & & \texttt{1 } \\
 & & \texttt{1 } \\
 & & \texttt{1 } \\
 & & \texttt{1 } \\
 & & \texttt{1 } \\
 & & \texttt{1 } \\
 & & \texttt{1 } \\
 & & \texttt{1 } \\
 & & \texttt{1 } \\
 & & \texttt{1 } \\
 & & \texttt{1 } \\
 & & \texttt{1 } \\
 & & \texttt{1 } \\
 & & \texttt{1 } \\
 & & \texttt{1 } \\
 & & \texttt{1 } \\
 & & \texttt{1 } \\
 & & \texttt{1 } \\
 & & \texttt{1 } \\
 & & \texttt{1 } \\
 & & \texttt{1 } \\
 & & \texttt{1 } \\
 & & \texttt{1 } \\
 & & \texttt{1 } \\
 & & \texttt{1 } \\
 & & \texttt{1 } \\
 & & \texttt{1 } \\
 & & \texttt{1 } \\
 & & \texttt{1 } \\
 & & \texttt{1 } \\
 & & \texttt{1 } \\
 & & \texttt{1 } \\
 & & \texttt{1 } \\
 & & \texttt{1 } \\
 & & \texttt{1 } \\
 & & \texttt{1 } \\
 & & \texttt{1 } \\
 & & \texttt{1 } \\
 & & \texttt{1 } \\
 & & \texttt{1 } \\
 & & \texttt{1 } \\
 & & \texttt{1 } \\
 & & \texttt{1 } \\
 & & \texttt{1 } \\
 & & \texttt{1 } \\
 & & \texttt{1 } \\
 & & \texttt{1 } \\
 & & \texttt{1 } \\
 & & \texttt{1 } \\
 & & \texttt{1 } \\
 & & \texttt{1 } \\
 & & \texttt{1 } \\
 & & \texttt{1 } \\
 & & \texttt{1 } \\
 & & \texttt{1 } \\
 & & \texttt{1 } \\
 & & \texttt{1 } \\
 & & \texttt{1 } \\
 & & \texttt{1 } \\
 & & \texttt{1 } \\
 & & \texttt{1 } \\
 & & \texttt{1 } \\
 & & \texttt{1 } \\
\hline \end{tabular}
}
\newpage
{\scriptsize
\begin{tabular}{|l|l|l|} \hline
\bf{Register Class} & \bf{Register File} & \bf{Registers} \\ \hline
 xwordqRegE & XBR & \texttt{1 } \\
 & & \texttt{1 } \\
 & & \texttt{1 } \\
 & & \texttt{1 } \\
 & & \texttt{1 } \\
 & & \texttt{1 } \\
 & & \texttt{1 } \\
 & & \texttt{1 } \\
 & & \texttt{1 } \\
 & & \texttt{1 } \\
 & & \texttt{1 } \\
 & & \texttt{1 } \\
 & & \texttt{1 } \\
 & & \texttt{1 } \\
 & & \texttt{1 } \\
 & & \texttt{1 } \\
 & & \texttt{1 } \\
 & & \texttt{1 } \\
 & & \texttt{1 } \\
 & & \texttt{1 } \\
 & & \texttt{1 } \\
 & & \texttt{1 } \\
 & & \texttt{1 } \\
 & & \texttt{1 } \\
 & & \texttt{1 } \\
 & & \texttt{1 } \\
 & & \texttt{1 } \\
 & & \texttt{1 } \\
 & & \texttt{1 } \\
 & & \texttt{1 } \\
 & & \texttt{1 } \\
 & & \texttt{1 } \\
 & & \texttt{1 } \\
 & & \texttt{1 } \\
 & & \texttt{1 } \\
 & & \texttt{1 } \\
 & & \texttt{1 } \\
 & & \texttt{1 } \\
 & & \texttt{1 } \\
 & & \texttt{1 } \\
 & & \texttt{1 } \\
 & & \texttt{1 } \\
 & & \texttt{1 } \\
 & & \texttt{1 } \\
 & & \texttt{1 } \\
 & & \texttt{1 } \\
 & & \texttt{1 } \\
 & & \texttt{1 } \\
 & & \texttt{1 } \\
 & & \texttt{1 } \\
 & & \texttt{1 } \\
 & & \texttt{1 } \\
 & & \texttt{1 } \\
 & & \texttt{1 } \\
 & & \texttt{1 } \\
 & & \texttt{1 } \\
 & & \texttt{1 } \\
 & & \texttt{1 } \\
 & & \texttt{1 } \\
 & & \texttt{1 } \\
 & & \texttt{1 } \\
 & & \texttt{1 } \\
 & & \texttt{1 } \\
 & & \texttt{1 } \\
\hline \end{tabular}
}
\newpage
{\scriptsize
\begin{tabular}{|l|l|l|} \hline
\bf{Register Class} & \bf{Register File} & \bf{Registers} \\ \hline
 xwordqRegO & XBR & \texttt{1 } \\
 & & \texttt{1 } \\
 & & \texttt{1 } \\
 & & \texttt{1 } \\
 & & \texttt{1 } \\
 & & \texttt{1 } \\
 & & \texttt{1 } \\
 & & \texttt{1 } \\
 & & \texttt{1 } \\
 & & \texttt{1 } \\
 & & \texttt{1 } \\
 & & \texttt{1 } \\
 & & \texttt{1 } \\
 & & \texttt{1 } \\
 & & \texttt{1 } \\
 & & \texttt{1 } \\
 & & \texttt{1 } \\
 & & \texttt{1 } \\
 & & \texttt{1 } \\
 & & \texttt{1 } \\
 & & \texttt{1 } \\
 & & \texttt{1 } \\
 & & \texttt{1 } \\
 & & \texttt{1 } \\
 & & \texttt{1 } \\
 & & \texttt{1 } \\
 & & \texttt{1 } \\
 & & \texttt{1 } \\
 & & \texttt{1 } \\
 & & \texttt{1 } \\
 & & \texttt{1 } \\
 & & \texttt{1 } \\
 & & \texttt{1 } \\
 & & \texttt{1 } \\
 & & \texttt{1 } \\
 & & \texttt{1 } \\
 & & \texttt{1 } \\
 & & \texttt{1 } \\
 & & \texttt{1 } \\
 & & \texttt{1 } \\
 & & \texttt{1 } \\
 & & \texttt{1 } \\
 & & \texttt{1 } \\
 & & \texttt{1 } \\
 & & \texttt{1 } \\
 & & \texttt{1 } \\
 & & \texttt{1 } \\
 & & \texttt{1 } \\
 & & \texttt{1 } \\
 & & \texttt{1 } \\
 & & \texttt{1 } \\
 & & \texttt{1 } \\
 & & \texttt{1 } \\
 & & \texttt{1 } \\
 & & \texttt{1 } \\
 & & \texttt{1 } \\
 & & \texttt{1 } \\
 & & \texttt{1 } \\
 & & \texttt{1 } \\
 & & \texttt{1 } \\
 & & \texttt{1 } \\
 & & \texttt{1 } \\
 & & \texttt{1 } \\
 & & \texttt{1 } \\
\hline \end{tabular}
}
\newpage
{\scriptsize
\begin{tabular}{|l|l|l|} \hline
\bf{Register Class} & \bf{Register File} & \bf{Registers} \\ \hline
 xwordvReg & XMR & \texttt{} \\
 & & \texttt{} \\
 & & \texttt{} \\
 & & \texttt{} \\
 & & \texttt{} \\
 & & \texttt{} \\
 & & \texttt{} \\
 & & \texttt{} \\
 & & \texttt{} \\
 & & \texttt{} \\
 & & \texttt{} \\
 & & \texttt{} \\
 & & \texttt{} \\
 & & \texttt{} \\
 & & \texttt{} \\
 & & \texttt{} \\
\hline \end{tabular}
}
\newpage
